\documentclass[11pt,a4paper,oneside]{book}
\usepackage[T1]{fontenc}
\usepackage[utf8]{inputenc}
\usepackage{lmodern}
\usepackage{microtype}
\usepackage[a4paper,top=2.6cm,bottom=2.8cm,left=3.0cm,right=2.6cm]{geometry}
\usepackage{amsmath}
\usepackage{amssymb}
\usepackage{amsthm}
\usepackage{enumitem}
\usepackage{booktabs}
\usepackage{array}
\usepackage{xcolor}
\usepackage{graphicx}
\usepackage{tikz}
\usepackage{pgfplots}
\pgfplotsset{compat=1.18}
\usetikzlibrary{arrows.meta,positioning,calc,patterns}
\usepackage[font=small,labelfont=bf,justification=centering]{caption}
\usepackage[most]{tcolorbox}
\usepackage{makeidx}
\usepackage[hidelinks,bookmarksnumbered,bookmarksopenlevel=1,plainpages=false]{hyperref}

\hypersetup{
  pdftitle={Kenapa Mesin Bisa Salah? Dari Aljabar Linear sampai Keamanan Machine Learning (Edisi Pertama)},
  pdfauthor={Afiq Andico Pangimpian},
  pdfsubject={Buku ajar bertingkat: matematika, komputasi, pembelajaran mesin, dan keamanan},
  pdfkeywords={aljabar linear, analisis real, peluang, optimasi, jaringan saraf, adversarial ML, Lean 4},
  pdfcreator={pdfLaTeX},
}

\definecolor{ink}{HTML}{1A1A1A}
\definecolor{accent}{HTML}{8C2F39}
\definecolor{layerorange}{HTML}{B45309}
\definecolor{layerteal}{HTML}{0F766E}
\definecolor{layergrey}{HTML}{374151}
\definecolor{layerindigo}{HTML}{4338CA}
\definecolor{layerjade}{HTML}{15803D}
\definecolor{layerbrown}{HTML}{92400E}
\definecolor{layerplum}{HTML}{86198F}
\color{ink}

\theoremstyle{definition}
\newtheorem{definisi}{Definisi}[chapter]
\theoremstyle{plain}
\newtheorem{teorema}{Teorema}[chapter]
\newtheorem{lemma}[teorema]{Lemma}
\newtheorem{akibat}[teorema]{Akibat}
\newtheorem{proposisi}[teorema]{Proposisi}
\theoremstyle{remark}
\newtheorem{contoh}{Contoh}[chapter]
\newtheorem{latihan}{Latihan}[chapter]
\newtheorem{catatan}{Catatan}[chapter]
\newtheorem*{ringkasan}{Ringkasan}

\renewcommand{\chaptername}{Bab}
\renewcommand{\partname}{Bagian}
\renewcommand{\contentsname}{Daftar Isi}
\renewcommand{\bibname}{Daftar Rujukan}
\renewcommand{\indexname}{Indeks}
\renewcommand{\figurename}{Gambar}
\renewcommand{\tablename}{Tabel}
\renewcommand{\listfigurename}{Daftar Gambar}

\setlist[itemize]{topsep=0.35em,itemsep=0.15em,parsep=0pt}
\setlist[enumerate]{topsep=0.35em,itemsep=0.15em,parsep=0pt}

% ---------------- empat lapisan pembaca ----------------
\newtcolorbox{untukanak}{colback=orange!6,colframe=layerorange,
  title={\small\bfseries Untuk pembaca muda (10 tahun ke atas)},
  fonttitle=\bfseries,breakable,left=6pt,right=6pt,top=4pt,bottom=4pt}
\newtcolorbox{untukdosen}{colback=teal!5,colframe=layerteal,
  title={\small\bfseries Catatan untuk dosen},
  fonttitle=\bfseries,breakable,left=6pt,right=6pt,top=4pt,bottom=4pt}
\newtcolorbox{untukpeneliti}{colback=black!4,colframe=layergrey,
  title={\small\bfseries Untuk peneliti: batas dan lanjutan},
  fonttitle=\bfseries,breakable,left=6pt,right=6pt,top=4pt,bottom=4pt}
% ---- kanal gaya belajar: auditori, kinestetik, prosedural, pemecah masalah ----
\newtcolorbox{bacakeras}{colback=violet!5,colframe=layerindigo,
  title={\small\bfseries Bacakan keras-keras (pembelajar auditori)},
  fonttitle=\bfseries,breakable,left=6pt,right=6pt,top=4pt,bottom=4pt}
\newtcolorbox{tangan}{colback=green!5,colframe=layerjade,
  title={\small\bfseries Kerjakan dengan tangan (pembelajar kinestetik)},
  fonttitle=\bfseries,breakable,left=6pt,right=6pt,top=4pt,bottom=4pt}
\newtcolorbox{resep}{colback=orange!4,colframe=layerbrown,
  title={\small\bfseries Resep langkah (pembelajar prosedural)},
  fonttitle=\bfseries,breakable,left=6pt,right=6pt,top=4pt,bottom=4pt}
\newtcolorbox{tantangan}{colback=magenta!4,colframe=layerplum,
  title={\small\bfseries Tantangan (pemecah masalah)},
  fonttitle=\bfseries,breakable,left=6pt,right=6pt,top=4pt,bottom=4pt}

\newcommand{\kamus}[3]{%
  \par\smallskip\noindent
  \colorbox{blue!5}{\parbox{\dimexpr\linewidth-2\fboxsep}{%
    \small\textbf{#1} (\emph{#2}) --- #3}}\par\smallskip}
\newcommand{\tingkat}[1]{%
  \par\noindent{\footnotesize\scshape Tingkat: #1}\par\smallskip}
\newcommand{\kertas}[1]{%
  \par\medskip\noindent
  \fbox{\parbox{\dimexpr\linewidth-2\fboxsep-2\fboxrule}{\small #1}}%
  \par\medskip}
\newcommand{\peta}[1]{%
  \par\medskip\noindent{\small\itshape Peta ke kitab sumber: #1}\par\medskip}

\newenvironment{lab}{\par\medskip\noindent\textbf{Lab komputasi}\par\nopagebreak}{\par\medskip}
% Kode lebih kecil dari teks: tanpa ini baris selebar 80+ kolom meluber keluar
% margin. Ambang aman pada tata letak ini ~90 kolom (dijaga gerbang).
\makeatletter
\renewcommand{\verbatim@font}{\footnotesize\ttfamily}
\makeatother

\makeindex

% Baris yang memuat kata tanpa titik potong (mis. "Cauchy-Schwarz" di daftar isi)
% diberi ruang lebih untuk meregang, bukan dibiarkan menembus margin. Tanpa ini
% ada empat baris meluber 1,87-7,99 pt ke margin kanan.
\setlength{\emergencystretch}{3em}

\begin{document}

\frontmatter

% ---------------------------- halaman judul ----------------------------
\hypersetup{pageanchor=false}
\begin{titlepage}
\centering
\vspace*{1.6cm}
{\color{accent}\rule{\textwidth}{1.2pt}}\\[0.45cm]
{\large\scshape Kitab Ajar Bertingkat}\\[0.15cm]
{\itshape Matematika $\cdot$ Komputasi $\cdot$ Pembelajaran Mesin $\cdot$ Keamanan}\\[0.45cm]
{\color{accent}\rule{\textwidth}{0.6pt}}\\[1.8cm]
{\Huge\bfseries Kenapa Mesin\\[0.25cm] Bisa Salah?}\\[0.7cm]
{\Large\itshape Dari Aljabar Linear\\[0.18cm] sampai Keamanan Machine Learning}\\[0.7cm]
{\large\itshape Edisi Pertama}\\[0.5cm]
{\large Enam bagian, tiga puluh dua bab, empat jalur baca}\\[0.35cm]
{\normalsize\itshape untuk pembaca usia 10 tahun, mahasiswa, dosen, dan peneliti}\\[2.0cm]
{\large Disusun oleh Afiq Andico Pangimpian}\\[0.3cm]
{\large Oktober 2026}
\vfill
{\footnotesize\itshape Buku ajar tentang cara kita tahu bahwa sebuah jawaban benar --- dan
cara mesin gagal tahu. Sumber lengkapnya ada di Daftar Rujukan.}
\end{titlepage}
\hypersetup{pageanchor=true}

% ---------------------------- kata pengantar ----------------------------
\chapter*{Kata Pengantar}
\addcontentsline{toc}{chapter}{Kata Pengantar}
\markboth{Kata Pengantar}{Kata Pengantar}

Edisi pertama buku ini tumbuh dari satu masalah sederhana: koleksi kitab ajar yang
saya kumpulkan berisi lebih dari sepuluh ribu halaman, tetapi halaman-halaman itu
tidak pernah terbaca berurutan. Membuka Rudin pada bab tiga tanpa pernah menyentuh
definisi supremum hanya menghasilkan frustrasi. Yang kurang bukan bahan, melainkan
urutan.

Sebelum naskah ini, ada versi tiga belas bab. Setelah dibaca ulang, dua kelemahan
kelihatan. Pertama, tiga belas bab terlalu cepat sehingga bagian tersulit hanya lewat
sekilas. Kedua, buku itu hanya bisa dibaca satu jenis pembaca: mahasiswa yang sudah
terbiasa dengan notasi formal. Siswa SMP yang penasaran tidak punya pintu masuk, dan
dosen tidak menemukan catatan mengajar di dalamnya.

Versi ini memperbaiki keduanya. Cakupannya diperluas menjadi enam bagian dan tiga puluh
dua bab, dengan materi yang di versi tiga belas bab tidak ada: himpunan dan
kardinalitas, deret tak hingga, sistem persamaan linear dan determinan, integral lipat
serta teorema lapangan vektor, persamaan diferensial dan sistem dinamik, regresi,
generalisasi dan regularisasi model. Setiap bab kini disusun berlapis.

Lapisan pertama adalah kotak \emph{Untuk pembaca muda}. Isinya analogi benda
sehari-hari, tanpa satu pun simbol matematika, ditutup satu pertanyaan yang membuat
penasaran. Lapisan kedua adalah badan bab itu sendiri: definisi, teorema, bukti,
contoh hitung, dan latihan, seperti buku ajar biasa. Lapisan ketiga adalah kotak
\emph{Catatan untuk dosen}: urutan mengajar yang saya sarankan, miskonsepsi yang
hampir selalu muncul, dan satu aktivitas kelas. Lapisan keempat adalah kotak
\emph{Untuk peneliti}: apa yang benar-benar terbukti, apa yang hanya dipopulerkan,
sumber primernya, dan pertanyaan yang masih terbuka.

Keempat lapisan tidak saling menggantikan: pembaca muda dapat membaca satu bab hanya
dari kotak oranye dan ringkasannya, mahasiswa membaca semuanya kecuali kotak peneliti,
dosen memakai kotak hijau sebagai rencana pertemuan, dan peneliti memulai dari kotak
abu-abu. Satu buku, empat jalur.

Lapisan itu mengatur \emph{kedalaman}. Yang mengatur \emph{cara masuk} adalah empat
kotak kedua: \emph{Bacakan keras} untuk pembaca yang paham dengan mendengar,
\emph{Kerjakan dengan tangan} untuk yang paham dengan bergerak dan meraba,
\emph{Resep langkah} untuk yang paham dengan menghafal prosedur lebih dahulu, dan
\emph{Tantangan} untuk yang paham justru ketika soalnya belum jelas. Keduanya hidup
bersama: satu bab yang sama, dengan enam pintu.

Tiga prinsip penulisan tidak berubah. \emph{Tidak ada klaim tanpa alat uji}: setiap
teorema disertai bukti atau pengakuan jujur bahwa bukti dilewati beserta lokasinya;
setiap angka berasal dari eksekusi nyata yang pembaca dapat ulangi. \emph{Notasi
diperkenalkan sebelum dipakai}. \emph{Kejujuran evaluasi adalah bagian dari
matematika}: bab-bab terakhir tidak berhenti pada cara membangun model, melainkan
pada cara mengetahui apakah klaim tentang model itu terbukti.

Satu hal yang sengaja tidak ada di buku ini: kunci jawaban lengkap. Memberi angka
akhir tanpa jejak pemeriksaan akan melanggar prinsip pertama. Sebagai gantinya,
Lampiran~C memberi kriteria penilaian dan alat verifikasi untuk setiap jenis latihan,
sehingga pembaca dapat memeriksa dirinya sendiri dengan cara yang sama seperti
penulis memeriksa naskah ini.

Buku ini bukan pengganti kitab sumbernya, melainkan pintu masuk yang tertata: jika
satu bab membuat pembaca membuka Rudin, Axler, Boyd, atau makalah Athalye lalu membaca
bagian yang ditunjuk, buku ini sudah berhasil. Kesalahan yang tersisa milik penulis,
bukan kitab-kitab itu.

\enlargethispage{3\baselineskip}
\vspace{0.35cm}
\begin{flushright}
Afiq Andico Pangimpian\\
Oktober 2026
\end{flushright}

% ---------------------------- cara membaca ----------------------------
\chapter*{Cara Membaca: Empat Jalur}
\addcontentsline{toc}{chapter}{Cara Membaca: Empat Jalur}
\markboth{Cara Membaca}{Cara Membaca}

Setiap bab punya bentuk yang sama, sehingga pembaca tahu apa yang diharapkan dan
boleh melompat tanpa kehilangan arah.

\begin{center}
\small
\begin{tabular}{@{}p{3.1cm}p{4.6cm}p{4.4cm}@{}}
\toprule
\textbf{Jalur} & \textbf{Yang dibaca} & \textbf{Hasil yang didapat} \\
\midrule
Pembaca muda (10-14) & kotak \emph{Untuk pembaca muda}, Ringkasan, satu contoh
  hitung & gagasan inti tanpa notasi; tahu apa yang sedang dihitung \\
Mahasiswa S1 & seluruh badan bab, Lab komputasi, Latihan; lewati kotak peneliti &
  definisi, bukti, kemampuan menghitung dan memeriksa \\
Dosen & kotak \emph{Catatan untuk dosen}, peta sumber, Latihan & rencana
  pertemuan, miskonsepsi yang harus diantisipasi, bahan papan tulis \\
Peneliti & kotak \emph{Untuk peneliti}, peta sumber, Lampiran~A dan~B & batas
  klaim, sumber primer, pertanyaan terbuka, notasi yang konsisten \\
\bottomrule
\end{tabular}
\end{center}

\medskip
Tanda \textsc{tingkat} di awal setiap bagian menunjukkan kesulitan bacaan:
\emph{dasar} dapat dibaca tanpa bab sebelumnya selain yang ditunjuk;
\emph{menengah} menuntut aljabar dan notasi buku ini;
\emph{lanjut} boleh dilewati pada pembacaan pertama tanpa memutus alur bab.

\medskip
\noindent\textbf{Enam cara belajar, enam pintu masuk.} Orang tidak belajar dengan cara
yang sama. Karena itu setiap bab memuat lebih dari satu rute menuju gagasan yang sama,
dan tidak ada rute yang wajib: pilih yang paling cepat membuatmu paham, lalu kembali ke
badan bab untuk melihat mengapa cara itu bekerja.

\begin{center}
\small
\begin{tabular}{@{}p{3.4cm}p{4.3cm}p{4.6cm}@{}}
\toprule
\textbf{Tipe} & \textbf{Cari ini di setiap bab} & \textbf{Cara memakainya} \\
\midrule
Visual & 2-3 gambar bernomor (di Daftar Gambar), tabel, dan warna kotak & baca
  gambar sebelum teksnya; gambar memuat angka yang sama dengan hitungan di teks \\[2pt]
Auditori & kotak \emph{Bacakan keras-keras} (indigo) & bacakan naskahnya keras-keras
  atau rekam suaramu; ulangi bagian yang berirama \\[2pt]
Kinestetik & kotak \emph{Kerjakan dengan tangan} (hijau) & lakukan dengan benda nyata:
  kertas, kartu, lantai, uang, atau badan sendiri \\[2pt]
Pemaham konsep & definisi, teorema, bukti, kotak \emph{Untuk peneliti} & jangan pakai
  rumus sebelum tahu asalnya; bukti adalah tempat asal itu ditulis \\[2pt]
Penghafal prosedural & kotak \emph{Resep langkah} (cokelat) dan Latihan & hafalkan
  urutan kartu resep, lalu terapkan pada angka berbeda sampai otomatis \\[2pt]
Pemecah masalah & kotak \emph{Tantangan} (ungu) dan butir latihan bertanda &
  kerjakan tantangan lebih dulu; petunjuknya bertingkat, jangan dibuka sekaligus \\
\bottomrule
\end{tabular}
\end{center}

Empat jalur di bawah ini mengatur \emph{berapa banyak} yang dibaca; enam tipe di atas
mengatur \emph{lewat pintu mana} masuk. Keduanya boleh dikombinasikan: seorang pembaca
visual yang mengejar jalur peneliti tetap membaca semua bukti, hanya saja ia mulai dari
gambarnya.

Tiga jalur belajar yang sudah dicoba, dari yang paling ringan:

\begin{enumerate}
  \item \textbf{Jalur rasa ingin tahu (8-12 minggu).} Baca kotak oranye, Ringkasan,
        dan satu contoh hitung dari setiap bab, berurutan dari Bab~1 sampai Bab~32.
        Tanpa latihan. Tujuannya satu: tahu peta medan.
  \item \textbf{Jalur mahasiswa (2 semester).} Semester pertama: Bab~1-10 dan
        Bab~17. Semester kedua: Bab~11-16 lalu 18-32. Itu dua puluh satu bab, lebih
        banyak dari satu semester biasa; cara memangkasnya tanpa memutus prasyarat
        adalah menunda Bab~31-32 (kriptografi dan privasi) ke mata kuliah pilihan,
        sehingga tersisa sembilan belas bab. Kerjakan Lab komputasi setiap bab;
        latihan minimal empat butir per bab. Urutan ini menjaga janji prasyarat:
        Bab~18 menuntut Bab~13 dan~17, jadi ketiganya tidak boleh berjauhan; Bab~21
        menuntut Bab~10 dan~20; Bab~23 menuntut Bab~12.
  \item \textbf{Jalur peneliti keamanan ML (9 minggu).} Sembilan belas bab, bukan
        tujuh: Bab~3, 4, 9, 10, 12, 13, 18, 19, 20, 22, 23, 25, lalu 26-30, ditutup
        Bab~31 dan~32. Bab~13, 18, 22, dan~23 ikut karena masing-masing adalah
        prasyarat Bab~20 (uji hipotesis), Bab~29 (adversarial), dan Bab~25 (jaringan
        saraf); Bab~31 dan~32 melengkapi sisi kriptografi dan privasi yang dituntut
        judul buku. Membaca jalur ini dalam enam minggu berarti melewati prasyarat
        dan kehilangan alasan mengapa klaimnya berlaku.
\end{enumerate}

Simbol dipakai konsisten di seluruh buku: $\mathbb{N}$, $\mathbb{Z}$, $\mathbb{Q}$,
$\mathbb{R}$ untuk himpunan bilangan; $\mathbf{x}$ untuk vektor dan $x_i$ untuk
komponennya; $\|\mathbf{x}\|$ untuk norma Euklides; $\nabla f$ untuk gradien;
$\varepsilon$ dan $\delta$ selalu bilangan positif kecil; $\mathbb{E}[X]$ untuk
nilai harapan; $\mathcal{L}(V,W)$ untuk ruang transformasi linear; $\top$ untuk
transpos. Lampiran~A memuat daftar lengkapnya.

% ---------------------------- peta ketergantungan ----------------------------
\chapter*{Peta Ketergantungan Bab}
\addcontentsline{toc}{chapter}{Peta Ketergantungan Bab}
\markboth{Peta Ketergantungan Bab}{Peta Ketergantungan Bab}

Buku ini boleh dibaca tidak berurutan, tetapi tidak semua bab berdiri sendiri.
Tabel berikut menyebut prasyarat minimum setiap bab; semua prasyarat lain boleh
diabaikan pada pembacaan pertama.

\begin{center}
\small
\begin{tabular}{@{}r p{9.0cm} p{2.8cm}@{}}
\toprule
\textbf{Bab} & \textbf{Judul} & \textbf{Prasyarat} \\
\midrule
1 & Bahasa Matematika dan Seni Membuktikan & --- \\
2 & Himpunan, Relasi, Fungsi, dan Kardinalitas & 1 \\
3 & Aksioma Bilangan Real dan Supremum & 1, 2 \\
4 & Barisan dan Limit & 3 \\
5 & Deret Tak Hingga & 4 \\
6 & Ruang Vektor dan Subruang & 2 \\
7 & Basis, Dimensi, dan Transformasi Linear & 6 \\
8 & Sistem Persamaan Linear, Determinan, dan Invers & 6 \\
9 & Nilai Eigen dan Diagonalisasi & 7, 8 \\
10 & Ortogonalitas, SVD, dan Kuadrat Terkecil & 7, 9 \\
11 & Kekontinuan dan Kekontinuan Seragam & 3, 4 \\
12 & Turunan: Satu dan Banyak Variabel & 11, 7 \\
13 & Integral Riemann dan Teorema Dasar Kalkulus & 11 \\
14 & Deret Fungsi dan Pertukaran Limit & 5, 11, 13 \\
15 & Integral Lipat, Perubahan Variabel, dan Teorema Lapangan Vektor & 10, 12, 13 \\
16 & Persamaan Diferensial dan Sistem Dinamik & 12, 14, 15 \\
17 & Peluang dan Peluang Bersyarat & 2, 5 \\
18 & Variabel Acak dan Distribusi & 17, 13 \\
19 & Hukum Bilangan Besar dan Teorema Limit Pusat & 18, 5 \\
20 & Estimasi, Interval Kepercayaan, dan Uji Hipotesis & 19 \\
21 & Regresi dan Statistika Terapan & 10, 20 \\
22 & Aritmetika Floating-Point dan Analisis Galat & 3, 5 \\
23 & Optimasi Konveks dan Dualitas & 6, 12 \\
24 & Pemrograman Linear dan Integer & 8, 23 \\
25 & Jaringan Saraf dan Turunan Balik & 10, 12, 23 \\
26 & Generalisasi, Regularisasi, dan Pelatihan & 18, 25 \\
27 & Evaluasi Model yang Jujur & 20, 26 \\
28 & Metode Formal dengan Lean 4 & 1, 2 \\
29 & Adversarial Machine Learning: Serangan dan Pertahanan & 25, 22 \\
30 & Batas Serangan: Gradien Dikaburkan dan Evasi NIDS & 27, 29 \\
31 & Kriptografi: Bilangan Bulat Modular sampai Ancaman Kuantum & 2, 3, 28 \\
32 & Privasi Diferensial dan Fondasi Matematis QML & 10, 17, 18 \\
\bottomrule
\end{tabular}
\end{center}

\phantomsection
\pdfbookmark[0]{Daftar Isi}{toc}
\tableofcontents
\cleardoublepage
\phantomsection
\pdfbookmark[0]{Daftar Gambar}{lof}
\listoffigures

\mainmatter

\part{Fondasi: Bahasa, Himpunan, dan Bilangan}

Bagian ini menjawab satu pertanyaan yang biasanya dilewati: apa arti kata
``berlaku untuk semua'' dan bagaimana sebuah kalimat bisa dibuktikan, bukan hanya
diyakini. Prasyaratnya hanya aljabar sekolah. Setelah lima bab, kamu bisa menulis
bukti langsung, bukti dengan kontradiksi, dan induksi, serta tahu kapan sebuah
argumen belum lengkap.
\chapter{Bahasa Matematika dan Seni Membuktikan}

Di suatu kelas, seorang mahasiswa menulis di papan: ``$n^{2}+n+41$ prima untuk setiap
bilangan bulat $n \geq 0$''. Ia lalu memeriksa satu per satu. Untuk $n=0$ diperoleh
$41$; untuk $n=1$ diperoleh $43$; untuk $n=2$ diperoleh $47$. Ia melanjutkan sampai
$n=39$, mendapat $1601$, dan $1601$ memang prima. Empat puluh nilai, empat puluh
kegagalan yang tidak terjadi. Di $n=40$ klaim itu patah:
$40^{2}+40+41 = 1681 = 41^{2}$. Satu baris tambahan membatalkan seluruh halaman
sebelumnya.

Kejadian itu bukan kecelakaan aritmetika, melainkan kesalahan tata bahasa. ``Untuk setiap
$n \geq 0$'' berbicara tentang tak berhingga banyak bilangan, sedangkan daftar hanya
menyentuh berhingga banyak bilangan; tidak ada aturan logika yang memindahkan kebenaran
dari yang berhingga ke yang tak berhingga. Bab ini karena itu menawarkan cara berbicara
sehingga pemeriksaan tidak perlu.

Kebutuhan itu bukan kemewahan. Di Bab~29 kita akan menyimpulkan bahwa sebuah
pengklasifikasi dapat dikelilingi dengan menggeser satu masukan sejauh $\varepsilon$,
dan kesimpulan itu berlaku untuk semua masukan. Di Bab~27 laporan akurasi yang bagus pada
ribuan data uji tidak membuktikan apa pun tentang data esok hari. Bab ini menyiapkan
perangkat untuk keduanya.

\begin{untukanak}
Di tepi sebuah danau hiduplah sekawanan angsa. Seorang anak datang setiap pagi selama
hampir setahun dan menghitung angsa yang dilihatnya: putih, semuanya, selama dua ratus
tujuh puluh pagi. Ia pun menulis di buku catatannya: semua angsa di danau ini putih.

Kakeknya bertanya berapa angsa yang belum pernah dilihat anak itu. Anak itu mengaku
tidak tahu, dan kakeknya berkata bahwa catatan itu bercerita tentang angsa yang sudah
dilihat, bukan tentang semua angsa di danau.

Cerita ini bukan tentang angsa, melainkan tentang beda antara ``sudah kuperiksa
semuanya'' dan ``berlaku untuk semuanya''. Yang pertama hanya sekuat ingatan kita. Yang kedua tidak bergantung pada siapa yang memeriksa. Angsa hitam
cukup satu untuk meruntuhkan catatan setebal apa pun, sedangkan seribu angsa putih
tidak pernah cukup untuk menegakkannya.

Angsa hitam pada bab ini bernama empat puluh. Pada bilangan itu rumus yang tampak rapi
mendadak menghasilkan bilangan yang bisa dipecah. Di bab ini kita akan tahu mengapa satu
contoh meruntuhkan klaim, mengapa banyak contoh tidak menegakkannya, dan bagaimana celah
itu ditutup tanpa memeriksa angsa terakhir.
\end{untukanak}

\section{Pernyataan, penghubung, dan tabel kebenaran}
\tingkat{dasar}

Kalimat baru bisa dinilai setelah kita tahu apa yang dijanjikannya. ``Tujuh adalah
bilangan prima'' menjanjikan sesuatu yang bisa diperiksa, dan janji itu benar. ``Tutup
jendelanya'' tidak menjanjikan apa pun. ``$x+1=3$'' belum bisa dinilai, karena nilainya
bergantung pada $x$. Ketiganya berbentuk sama di mata ejaan, tetapi hanya yang pertama
termasuk bahan matematika.

\begin{definisi}[Pernyataan]
\index{pernyataan}
\emph{Pernyataan} adalah kalimat yang bernilai benar atau bernilai salah, tetapi tidak
keduanya sekaligus.
\end{definisi}

Pernyataan tunggal menjadi berguna setelah digabung. Setiap penghubung punya satu arti
yang ditetapkan, bukan disimpulkan dari kebiasaan berbahasa.

\kamus{penghubung logika}{logical connective}{seperti tombol pada mesin hitung: dua
lampu masuk, satu lampu keluar, dan hubungan antarlampu ditetapkan pabrik, bukan
ditebak pemakainya.}

Dengan gambar tombol itu di kepala, kita tetapkan arti kelima penghubung secara resmi.

\begin{definisi}[Penghubung logika]
Misalkan $P$ dan $Q$ pernyataan.

Lima penghubung berikut kita pakai.
\begin{itemize}
  \item $\neg P$ (\emph{negasi}, \emph{not}): benar tepat ketika $P$ salah.
  \item $P \wedge Q$ (\emph{konjungsi}, \emph{and}): benar tepat ketika keduanya benar.
  \item $P \vee Q$ (\emph{disjungsi}, \emph{or}): benar tepat ketika paling sedikit
        satu benar; kata ``atau'' di sini inklusif.
  \item $P \Rightarrow Q$ (\emph{implikasi}, \emph{implication}): salah tepat ketika
        $P$ benar dan $Q$ salah; benar pada tiga kemungkinan lainnya.
  \item $P \Leftrightarrow Q$ (\emph{bikondisional}, \emph{iff}): benar tepat ketika
        $P$ dan $Q$ bernilai sama.
\end{itemize}
\end{definisi}

Seluruh isi definisi itu dapat dipadatkan menjadi satu tabel: barisnya semua kombinasi
nilai, kolomnya nilai setiap penghubung.

{\centering
\begin{tabular}{cc|ccccc}
\hline
$P$ & $Q$ & $\neg P$ & $P \wedge Q$ & $P \vee Q$ & $P \Rightarrow Q$ & $P \Leftrightarrow Q$ \\
\hline
B & B & S & B & B & B & B \\
B & S & S & S & B & S & S \\
S & B & B & S & B & B & S \\
S & S & B & S & S & B & B \\
\hline
\end{tabular}\par}

\begin{figure}[htbp]
\centering
\begin{tikzpicture}[font=\small]
  \fill[layerjade!16] (0,0) rectangle (3,1.8);
  \fill[accent!16] (3,0) rectangle (6,1.8);
  \fill[layerjade!16] (0,-1.8) rectangle (3,0);
  \fill[layerjade!16] (3,-1.8) rectangle (6,0);
  \draw[layergrey!70,thick] (0,-1.8) grid[xstep=3,ystep=1.8] (6,1.8);
  \draw[accent,very thick] (3,0) rectangle (6,1.8);
  \node at (1.5,0.9) {B};
  \node at (4.5,0.9) {S};
  \node at (1.5,-0.9) {B};
  \node at (4.5,-0.9) {B};
  \node[above,font=\footnotesize] at (1.5,1.8) {$Q$ benar};
  \node[above,font=\footnotesize] at (4.5,1.8) {$Q$ salah};
  \node[left,font=\footnotesize] at (0,0.9) {$P$ benar};
  \node[left,font=\footnotesize] at (0,-0.9) {$P$ salah};
  \node[below,font=\footnotesize,text=accent] at (4.5,-1.8) {satu-satunya sel yang salah};
\end{tikzpicture}
\caption{Kisi kebenaran $P \Rightarrow Q$: baris berisi nilai $P$, kolom berisi nilai $Q$.
Tiga sel bernilai B dan satu sel bernilai S, yaitu ketika $P$ benar sementara $Q$ salah.
Siapa pun yang ingin menggugurkan implikasi cukup menunjuk sel yang dibingkai tebal.}
\label{fig:kisi-implikasi}
\end{figure}

Baris ketiga dan keempat adalah yang paling sering mengejutkan. Pada baris keempat,
$P$ salah dan $Q$ salah, dan implikasinya dinilai benar. Pada baris ketiga, $P$ salah
dan $Q$ benar, dan implikasinya juga benar. Jadi satu-satunya cara menggugurkan
$P \Rightarrow Q$ adalah menunjukkan $P$ benar sementara $Q$ salah; kisi pada
Gambar~\ref{fig:kisi-implikasi} memuat sel itu sebagai satu-satunya sel yang salah.

Ketetapan itu bukan pilihan gaya. Perhatikan ``jika $n$ habis dibagi $4$, maka $n$
genap'', yang harus berlaku untuk setiap bilangan bulat $n$, termasuk $n=3$. Di situ
antesedennya salah; kalau implikasi dinilai salah, pernyataan yang jelas benar tadi akan
gugur karena satu bilangan yang tidak relevan. Implikasi mencatat hubungan nilai, bukan
sebab-akibat.

\kamus{implikasi}{implication}{seperti janji ``kalau hujan, jalan basah'': janji itu
hanya dilanggar bila hujan turun dan jalan tetap kering; bila tidak hujan, si pemilik
janji belum berbuat salah.}

Dua pernyataan majemuk bisa bernilai sama pada semua kemungkinan; kesamaan semacam itu
punya nama sendiri.

\begin{definisi}[Tautologi dan kesetaraan]
\index{tautologi}
Pernyataan majemuk disebut \emph{tautologi} jika bernilai benar pada setiap penetapan
nilai pernyataan-pernyataan penyusunnya. Dua pernyataan majemuk $A$ dan $B$ disebut
\emph{setara secara logis}, ditulis $A \equiv B$, jika $A \Leftrightarrow B$ tautologi.
\end{definisi}

Untuk dua pernyataan penyusun ada $2^{2}=4$ penetapan nilai; untuk tiga penyusun ada
$2^{3}=8$. Kesetaraan yang hanya melibatkan dua atau tiga pernyataan karena itu dapat
dibuktikan dengan memeriksa seluruh baris, dan pemeriksaan itu lengkap. Inilah satu-satunya tempat dalam bab ini di mana
pemeriksaan habis-habisan diterima sebagai bukti, sebab himpunan kasusnya berhingga dan
sudah dihitung.

\begin{proposisi}
Untuk setiap pernyataan $P$ dan $Q$ berlaku $P \Rightarrow Q \equiv \neg P \vee Q$.
\end{proposisi}
\par\smallskip\noindent\textit{Bukti.} Kita bandingkan kolom $P \Rightarrow Q$ dan kolom $\neg P \vee Q$ pada keempat baris
tabel kebenaran. Pada baris pertama keduanya benar. Pada baris kedua $\neg P$ salah dan
$Q$ salah, sehingga keduanya salah. Pada baris ketiga dan keempat $\neg P$ benar,
sehingga keduanya benar. Keempat baris memberi nilai yang sama, jadi kedua pernyataan
setara.
\hfill$\blacksquare${\parfillskip=0pt\par}

Dua kesetaraan turunan dipakai sepanjang buku ini. Yang pertama membolehkan kita membalik
arah pembicaraan tanpa mengubah nilai.

\begin{akibat}[Kontraposisi]
\index{kontraposisi}
Untuk setiap pernyataan $P$ dan $Q$ berlaku
$P \Rightarrow Q \equiv \neg Q \Rightarrow \neg P$.
\end{akibat}
\par\smallskip\noindent\textit{Bukti.} Menurut proposisi sebelumnya, kedua ruas sama dengan $\neg P \vee Q$ dan
$\neg(\neg Q) \vee \neg P$. Karena $\neg(\neg Q) \equiv Q$ dan disjungsi tidak
memperhatikan urutan, kedua bentuk itu identik.
\hfill$\blacksquare${\parfillskip=0pt\par}

Yang \emph{tidak} setara adalah konversnya, yaitu $Q \Rightarrow P$. Pada baris ketiga
tabel, $Q \Rightarrow P$ benar sementara $P \Rightarrow Q$ salah. Menukar anteseden
dengan konsekuen adalah kesalahan tersering, dan sering tidak terlihat karena kalimatnya
tetap enak dibaca.

\begin{proposisi}
Untuk setiap pernyataan $P$ dan $Q$ berlaku
$\neg(P \Rightarrow Q) \equiv P \wedge \neg Q$.
\end{proposisi}
\par\smallskip\noindent\textit{Bukti.} Menurut definisi negasi dan proposisi pertama,
$\neg(P \Rightarrow Q) \equiv \neg(\neg P \vee Q)$. Hukum De Morgan yang dibuktikan di
Bab~2 memberi $\neg(\neg P \vee Q) \equiv \neg(\neg P) \wedge \neg Q$, dan
$\neg(\neg P) \equiv P$. Jadi $\neg(P \Rightarrow Q) \equiv P \wedge \neg Q$.
Pembaca yang belum menerima hukum De Morgan dapat memeriksa keempat baris tabel secara
langsung.
\hfill$\blacksquare${\parfillskip=0pt\par}

\begin{tangan}
Sediakan delapan kartu kecil: empat bertuliskan B, empat bertuliskan S. Susun empat baris,
setiap baris satu kartu untuk $P$ dan satu untuk $Q$: B-B, B-S, S-B, S-S. Di kanan setiap
baris letakkan kartu ketiga yang Anda tulis sendiri: B kalau ``jika $P$ maka $Q$'' Anda
nilai benar, S kalau salah.

Kerjakan tanpa melihat tabel, lalu bandingkan. Hanya satu baris mendapat kartu S, yaitu
baris yang $P$-nya B dan $Q$-nya S. Geser kartu itu ke atas dan ke bawah; ia tetap
satu-satunya. Sekarang tukar kartu P dan Q pada baris B-S sehingga menjadi S-B, lalu nilai
ulang: kartunya berubah dari S menjadi B. Perpindahan satu kartu itulah konvers.
\end{tangan}

\section{Kuantor dan cara menegasikannya}
\tingkat{dasar}

Sebagian besar kalimat matematika berbicara tentang banyak objek, dan banyaknya itu
ditandai kuantor. Ada dua kuantor, dan keduanya membaca urutan sebagai bagian dari arti.

\kamus{kuantor}{quantifier}{seperti dua jenis pengumuman di sekolah: ``semua murid
wajib membawa buku'' dan ``ada murid yang membawa buku''; keduanya terdengar mirip
tetapi menuntut bukti yang sangat berbeda.}

Kita tuliskan artinya secara resmi.

\begin{definisi}[Kuantor universal dan eksistensial]
\index{kuantor}
Misalkan $P(x)$ pernyataan yang bergantung pada $x$ di dalam himpunan $A$.
\begin{itemize}
  \item $(\forall x \in A)\,P(x)$ berarti $P(x)$ benar untuk setiap $x \in A$.
  \item $(\exists x \in A)\,P(x)$ berarti terdapat paling sedikit satu $x \in A$
        dengan $P(x)$ benar.
\end{itemize}
Pernyataan pertama salah jika ada satu $x \in A$ dengan $P(x)$ salah; satu objek
semacam itu disebut \emph{contoh penyangkal} (\emph{counterexample}).
\end{definisi}

Kuantor eksistensial menuntut paling sedikit satu, bukan tepat satu; menambahkan tuntutan
yang tidak tercantum dalam definisi adalah sumber kesalahan tersering.

Urutan dua kuantor yang berbeda jenis tidak boleh ditukar. Di dalam $\mathbb{R}$,
pernyataan
\[
  (\forall x)(\exists y)\; y > x
\]
benar, sedangkan
\[
  (\exists y)(\forall x)\; y > x
\]
salah. Bukti keduanya pendek.

\begin{proposisi}
Di dalam $\mathbb{R}$ berlaku $(\forall x)(\exists y)\, y > x$, tetapi tidak berlaku
$(\exists y)(\forall x)\, y > x$.
\end{proposisi}
\par\smallskip\noindent\textit{Bukti.} Untuk bagian pertama, ambil $x \in \mathbb{R}$
sebarang dan pilih $y = x+1$. Maka $y > x$, sehingga $y$ menjadi saksi bagi $x$ itu.
Karena $x$ sebarang, setiap $x$ punya saksi; saksinya boleh berbeda untuk $x$ yang
berbeda.

Untuk bagian kedua, andaikan ada $y_{0} \in \mathbb{R}$ dengan $y_{0} > x$ untuk
setiap $x \in \mathbb{R}$. Ambil $x = y_{0}+1$, yang juga bilangan real. Maka
$y_{0} > y_{0}+1$, yaitu $0 > 1$, suatu hal yang tidak mungkin. Jadi tidak ada $y_{0}$
semacam itu.
\hfill$\blacksquare${\parfillskip=0pt\par}

\index{negasi}Negasi bekerja dari luar ke dalam: setiap kuantor berganti jenis, dan negasi berjalan masuk
sampai menempel pada pernyataan terdalam. Untuk pernyataan majemuk, hukum De Morgan
memindahkan negasi melewati konjungsi dan disjungsi.

{\centering
\begin{tabular}{ll}
\hline
Pernyataan & Negasinya \\
\hline
$(\forall x \in A)\,P(x)$ & $(\exists x \in A)\,\neg P(x)$ \\
$(\exists x \in A)\,P(x)$ & $(\forall x \in A)\,\neg P(x)$ \\
$P \wedge Q$ & $\neg P \vee \neg Q$ \\
$P \vee Q$ & $\neg P \wedge \neg Q$ \\
$P \Rightarrow Q$ & $P \wedge \neg Q$ \\
\hline
\end{tabular}\par}

Baris terakhir itu yang paling sering dilanggar. Menyangkal ``jika hujan maka jalan
basah'' bukan ``jika hujan maka jalan kering'', melainkan ``hujan dan jalan tidak basah''.
Kesalahan yang sama muncul saat orang menyangkal definisi: kata kerja di dalamnya
dinegasikan, kuantornya dibiarkan utuh.

\begin{proposisi}
Untuk setiap pernyataan $P(x)$ dan $Q(x)$ berlaku
\[
  \neg(\forall x)\bigl(P(x) \Rightarrow Q(x)\bigr)
  \;\equiv\; (\exists x)\bigl(P(x) \wedge \neg Q(x)\bigr).
\]
\end{proposisi}
\par\smallskip\noindent\textit{Bukti.} Kuantor universal dinegasikan menjadi kuantor
eksistensial atas negasi badan:
\[
  \neg(\forall x)R(x) \;\equiv\; (\exists x)\neg R(x).
\]
Badan $R(x)$ di sini adalah $P(x) \Rightarrow Q(x)$, yang negasinya $P(x) \wedge
\neg Q(x)$ menurut proposisi bagian sebelumnya. Jadi negasinya berbentuk
$(\exists x)(P(x) \wedge \neg Q(x))$.
\hfill$\blacksquare${\parfillskip=0pt\par}

Bentuk itu memberi tugas yang jelas ketika hendak menolak pernyataan berkuantor: tunjukkan
satu objek yang memenuhi anteseden dan gagal pada konsekuen.

Dua negasi berikut dipakai berulang kali; yang pertama menjadi bahan Bab~4.

\begin{contoh}[Negasi definisi limit]
Definisi limit berbunyi: barisan $(a_{n})$ menuju $L$ jika
\[
  (\forall \varepsilon > 0)(\exists N \in \mathbb{N})(\forall n \geq N)\;
  |a_{n} - L| < \varepsilon .
\]
Menegasikan dari luar ke dalam memberi
\[
  (\exists \varepsilon > 0)(\forall N \in \mathbb{N})(\exists n \geq N)\;
  |a_{n} - L| \geq \varepsilon .
\]
Bacaan kalimat kedua: ada satu toleransi $\varepsilon$ yang tidak bisa diperbaiki oleh
$N$ mana pun, karena selalu ada suku setelah $N$ yang menyimpang paling sedikit sejauh
$\varepsilon$. Untuk membuktikan bahwa $(a_{n})$ tidak menuju $L$, kita harus menemukan
$\varepsilon$ itu, lalu untuk setiap $N$ menemukan $n \geq N$ yang melanggar; $n$ boleh
bergantung pada $N$.
\end{contoh}

Negasi kedua menyangkut keberadaan bilangan prima. Euklides membuktikan bahwa untuk
setiap bilangan asli $n$ ada bilangan prima yang lebih besar daripada $n$:
\[
  (\forall n \in \mathbb{N})(\exists p)\bigl(p \text{ prima} \wedge p > n\bigr).
\]
Negasinya adalah
\[
  (\exists n \in \mathbb{N})(\forall p)\bigl(p \text{ prima} \Rightarrow p \leq n\bigr),
\]
yaitu: ada satu batas $n$ sehingga semua bilangan prima tidak melebihi $n$. Menyangkal
klaim Euklides menuntut batas semacam itu, bukan sekadar gagal menemukan prima besar.

\begin{untukanak}
Bayangkan sebuah kelas dengan tiga puluh murid. Guru menulis di papan: ``semua murid
membawa pensil.'' Untuk membuktikan kalimat itu benar, guru harus memeriksa ketiga puluh meja; untuk
membuktikannya salah, cukup satu murid yang tidak membawa pensil.

Sekarang guru menulis kalimat lain: ``ada murid yang membawa pensil.'' Kali ini
keadaannya terbalik. Satu murid yang membawa pensil sudah cukup untuk membuktikan
kalimat itu. Tetapi untuk membuktikannya salah, guru harus memeriksa ketiga puluh meja
dan memastikan tidak ada satu pun pensil.

Dua kalimat itu hanya bertukar kata di depan, tetapi pekerjaan memeriksanya bertukar
seluruhnya. Ahli matematika berhati-hati dengan ``semua'' dan ``ada'' karena keduanya
menghitung berapa banyak yang harus diperiksa.

Ada satu jebakan lagi. Kalau ada murid yang membawa dua pensil, apakah ``ada murid yang
membawa pensil'' masih benar? Tentu benar: kalimat itu menuntut paling sedikit satu,
bukan tepat satu. Di bab ini kita akan tahu mengapa tuntutan yang tidak tertulis tidak
boleh ditambahkan sendiri.
\end{untukanak}

\section{Membaca definisi secara harfiah}
\tingkat{menengah}

Definisi matematika bukan penjelasan singkat tentang sebuah kata; ia kontrak yang
berlaku dua arah. Kalau sebuah bilangan disebut ganjil apabila berbentuk $2k+1$ untuk
suatu bilangan bulat $k$, kita terikat pada dua hal: setiap bilangan berbentuk $2k+1$
harus disebut ganjil, dan setiap bilangan yang disebut ganjil harus berbentuk $2k+1$.

Membaca definisi berarti membaca satu kalimat berbentuk ``$x$ disebut $B$ jika dan hanya
jika $C(x)$'', walaupun kata ``hanya jika'' jarang ditulis; frasa ``disebut'' sudah
mengandung kedua arah. Arah pertama memberi izin memakai nama $B$, arah kedua memberi izin
memakai syarat $C(x)$, dan membalik salah satunya merusak seluruh argumen.

\kamus{syarat perlu}{necessary condition}{seperti syarat masuk gedung: tanpa kartu
anggota kita tidak boleh masuk, tetapi punya kartu anggota belum berarti kita sudah
duduk di dalam.}

Bandingkan dua pernyataan tentang bilangan bulat $n$: ``jika $6$ membagi $n$ maka $3$
membagi $n$'' benar, sedangkan ``jika $3$ membagi $n$ maka $6$ membagi $n$'' salah,
dengan contoh penyangkal $n=3$. Jadi ``$3$ membagi $n$'' syarat perlu bagi ``$6$ membagi
$n$'', tetapi bukan syarat cukup. Definisi yang baik menyebut keduanya, sebab hanya
dengan begitu nama dan syaratnya dapat ditukar.

\begin{bacakeras}
Bacakan sebagai dialog dua suara, bergantian, dengan satu tarikan napas di setiap titik.

Suara pertama: ``Kalau hujan, jalan basah.''
Suara kedua: ``Jalan basah. Berarti hujan.''
Suara pertama: ``Belum tentu. Jalan juga basah kalau disiram.''
Suara kedua: ``Tetap basah. Tetap hujan.''
Suara pertama: ``Yang kujanjikan hanya satu arah. Hujan dulu, baru basah. Arah balik tidak pernah kujanjikan.''

Ulangi dialog itu sekali lagi, lebih cepat, sampai telinga menangkap bedanya. Kalimat
pertama menaruh hujan di depan; kalimat kedua memindahkan hujan ke belakang. Bunyinya
berputar, janjinya tidak.

Baris yang harus diulang tiga kali dengan suara keras: ``jika hujan maka basah, bukan jika
basah maka hujan.''
\end{bacakeras}

Sekarang perhatikan definisi yang paling banyak dipakai di buku ini.

\begin{definisi}[Bilangan prima]
\index{bilangan prima}
Bilangan bulat $p$ disebut \emph{prima} jika $p \geq 2$ dan satu-satunya pembagi
positif $p$ adalah $1$ dan $p$ sendiri.
\end{definisi}

Kalimat sehari-harinya: bilangan prima lebih besar daripada satu dan tidak dapat dipecah
menjadi hasil kali dua bilangan yang lebih kecil; kedua klausanya diuji terpisah.

\begin{proposisi}
Bilangan $1$ bukan bilangan prima.
\end{proposisi}
\par\smallskip\noindent\textit{Bukti.} Definisi menuntut $p \geq 2$. Untuk $p = 1$
syarat itu gagal, sehingga $1$ tidak memenuhi definisi, apa pun yang dapat kita
katakan tentang pembaginya. Perhatikan bahwa $1$ sebenarnya lolos pada klausa kedua:
satu-satunya pembagi positif $1$ memang $1$ dan $1$ sendiri. Jadi yang menyingkirkannya
adalah klausa pertama, bukan klausa kedua.
\hfill$\blacksquare${\parfillskip=0pt\par}

Kejadian semacam ini menjelaskan mengapa definisi perlu dibaca sampai klausa terakhir.
Pembaca yang melewatkan klausa pertama akan menyimpulkan bahwa $1$ prima, dan
ketunggalan faktorisasi $12 = 2^{2}\cdot 3$ di Bab~28 akan hilang, karena $12$ juga
boleh ditulis $1 \cdot 2^{2} \cdot 3$.

Definisi juga menentukan cara menyangkal keanggotaan. Untuk menunjukkan bahwa sebuah
bilangan bukan prima, kita harus menemukan pembagi positif yang bukan $1$ dan bukan
bilangan itu sendiri. Berikut satu contoh yang sering lolos dari pemeriksaan cepat.

\begin{proposisi}
Bilangan $91$ bukan bilangan prima. Lebih umum, pernyataan ``setiap bilangan ganjil
yang tidak habis dibagi $3$ dan tidak habis dibagi $5$ adalah prima'' salah.
\end{proposisi}
\par\smallskip\noindent\textit{Bukti.} Faktanya $91 = 7 \cdot 13$, dan kedua faktor itu
berada di antara $1$ dan $91$, sehingga $91$ punya pembagi positif selain $1$ dan
dirinya. Bilangan $91$ ganjil, tidak habis dibagi $3$ karena $9+1 = 10$ bukan kelipatan
$3$, dan tidak habis dibagi $5$ karena angka terakhirnya bukan $0$ atau $5$. Jadi $91$
adalah contoh penyangkal bagi pernyataan umum tadi. Sebagai pembanding, $97$ memang prima; di antara bilangan bulat $2$ dan $99$ terdapat
tepat $25$ bilangan prima, hasil hitungan program pada bagian lab.
\hfill$\blacksquare${\parfillskip=0pt\par}

\begin{definisi}[Kuadrat sempurna]
Bilangan bulat $n$ disebut \emph{kuadrat sempurna} jika terdapat bilangan bulat $k$
dengan $n = k^{2}$.
\end{definisi}

Menurut definisi ini, $0$ adalah kuadrat sempurna, karena $0 = 0^{2}$, dan $1$ juga,
karena $1 = 1^{2}$. Sementara itu $2$ bukan kuadrat sempurna, dan pernyataan itu butuh
bukti, bukan pengakuan bahwa kita tidak menemukan $k$.

\begin{proposisi}
Bilangan $2$ bukan kuadrat sempurna.
\end{proposisi}
\par\smallskip\noindent\textit{Bukti.} Andaikan $2 = k^{2}$ untuk suatu $k \in
\mathbb{Z}$. Karena $k^{2} = 2 > 0$, berlaku $k \neq 0$. Jika $|k| = 1$ maka
$k^{2} = 1 \neq 2$. Jika $|k| \geq 2$ maka $k^{2} \geq 4 > 2$. Karena $|k|$ bilangan
bulat positif, hanya dua kemungkinan itu yang tersisa, dan keduanya gagal. Jadi tidak
ada $k$ semacam itu.
\hfill$\blacksquare${\parfillskip=0pt\par}

\index{contoh penyangkal}Kebiasaan yang layak dipelihara: setiap kali bertemu
definisi, ajukan tiga pertanyaan. Apa saja klausanya? Objek apa yang memenuhi
seluruhnya? Objek apa yang sering dianggap memenuhi padahal gagal pada satu klausa?
Pertanyaan ketiga menghasilkan contoh penyangkal, dan contoh penyangkal tidak dapat
dibantah maupun diperdebatkan.

\section{Empat pola bukti}
\tingkat{menengah}

Bukti adalah pengangkutan kebenaran: dari pernyataan yang sudah kita terima menuju
pernyataan yang ingin kita terima, lewat langkah yang masing-masing dapat diperiksa.
Empat pola di bawah ini menutup sebagian besar kebutuhan, dan
Gambar~\ref{fig:peta-tiga-bukti} memetakan tiga jalur pertamanya. Pola kelima yang lazim,
pembagian kasus, hanyalah pola pertama yang diterapkan pada kasus-kasus yang saling
meniadakan.

\begin{figure}[htbp]
\centering
\begin{tikzpicture}[font=\small,>=Stealth,
  jalur/.style={draw,rounded corners=2pt,align=center,inner sep=5pt,text width=4.6cm},
  ujung/.style={draw,rounded corners=2pt,align=center,inner sep=4pt,text width=2.2cm}]
  \node[jalur,fill=layerorange!12] (klaim) at (0,0) {klaim: $P \Rightarrow Q$};
  \node[jalur,fill=layerteal!12] (lgs) at (6.1,2.4)
    {\textbf{langsung}\\ uraikan $P$, hitung, muncullah $Q$};
  \node[jalur,fill=layerindigo!10] (ktp) at (6.1,0)
    {\textbf{kontraposisi}\\ buktikan $\neg Q \Rightarrow \neg P$; nilainya sama};
  \node[jalur,fill=layerplum!10] (ktr) at (6.1,-2.4)
    {\textbf{kontradiksi}\\ andaikan $P$ benar dan $Q$ salah, turunkan pertentangan};
  \node[ujung,fill=layergrey!10] (selesai) at (11.6,0)
    {tidak ada\\penetapan yang\\menggugurkan};
  \draw[->,thick] (klaim.east) -- (lgs.west);
  \draw[->,thick] (klaim.east) -- (ktp.west);
  \draw[->,thick] (klaim.east) -- (ktr.west);
  \draw[->,thick] (lgs.east) -- (selesai.north west);
  \draw[->,thick] (ktp.east) -- (selesai.west);
  \draw[->,thick] (ktr.east) -- (selesai.south west);
\end{tikzpicture}
\caption{Peta jalan tiga bentuk bukti implikasi. Semua jalur berangkat dari klaim yang sama
dan berakhir pada kenyataan bahwa tidak ada penetapan nilai yang membuat $P$ benar sementara
$Q$ salah. Bukti kontraposisi sah karena $P \Rightarrow Q \equiv \neg Q \Rightarrow \neg P$.
Bentuk keempat, bukti konstruktif dengan saksi, tidak masuk peta ini karena ia membuktikan
pernyataan berbentuk $\exists$, bukan implikasi.}
\label{fig:peta-tiga-bukti}
\end{figure}

\paragraph{Pola pertama: bukti langsung.} Kita berangkat dari hipotesis, menguraikannya
menurut definisi, lalu menghitung sampai kesimpulan muncul. Contoh berikut dipilih
karena kesimpulannya lebih kuat daripada yang biasanya dibuktikan orang.

\begin{teorema}
Untuk setiap bilangan bulat ganjil $n$, bilangan $n^{2}-1$ habis dibagi $8$.
\end{teorema}
\par\smallskip\noindent\textit{Bukti.} Karena $n$ ganjil, $n = 2k+1$ untuk suatu
$k \in \mathbb{Z}$. Maka
\[
  n^{2} - 1 = (2k+1)^{2} - 1 = 4k^{2} + 4k = 4k(k+1).
\]
Di antara dua bilangan bulat berurutan $k$ dan $k+1$, tepat satu di antaranya genap,
sehingga $k(k+1) = 2m$ untuk suatu $m \in \mathbb{Z}$. Substitusi memberi
$n^{2}-1 = 4 \cdot 2m = 8m$, yang berarti $n^{2}-1$ habis dibagi $8$.
\hfill$\blacksquare${\parfillskip=0pt\par}

Akibat langsungnya, kuadrat bilangan ganjil selalu bersisa $1$ bila dibagi $8$. Untuk
$n = 1,3,5,7,9,11$ berturut-turut diperoleh $n^{2}-1 = 0, 8, 24, 48, 80, 120$, dan
keenamnya habis dibagi $8$. Pemeriksaan itu bukan bukti, tetapi ia memastikan bahwa
bukti di atas tidak salah ketik.

\begin{figure}[htbp]
\centering
\begin{tikzpicture}[font=\small,>=Stealth,
  tahap/.style={draw,rounded corners=2pt,align=left,inner sep=5pt,text width=6.6cm},
  nama/.style={font=\footnotesize\itshape,text=layergrey,align=right,text width=2.4cm}]
  \node[tahap,fill=layerteal!10] (asumsi) at (0,0)
    {$n$ ganjil, jadi $n = 2k+1$ untuk suatu $k \in \mathbb{Z}$};
  \node[nama] at (-4.9,0) {asumsi (diberikan)};
  \node[tahap,fill=layerorange!10] (langkah) at (0,-2.1)
    {$n^{2}-1 = 4k^{2}+4k = 4k(k+1) = 8m$};
  \node[nama] at (-4.9,-2.1) {langkah (dihitung)};
  \node[tahap,fill=layerjade!10] (kesimpulan) at (0,-4.2)
    {$8$ membagi $n^{2}-1$};
  \node[nama] at (-4.9,-4.2) {kesimpulan (diminta)};
  \draw[->,thick] (asumsi.south) -- (langkah.north);
  \draw[->,thick] (langkah.south) -- (kesimpulan.north);
\end{tikzpicture}
\caption{Anatomi satu bukti langsung untuk pernyataan ``jika $n$ ganjil maka $8$ membagi
$n^{2}-1$''. Kotak atas satu-satunya yang boleh dipakai tanpa alasan; kotak tengah memuat
alasannya, yaitu definisi bilangan ganjil dan aljabar; kotak bawah adalah yang diminta.
Panahnya hanya boleh mengalir ke bawah.}
\label{fig:anatomi-bukti}
\end{figure}

Gambar~\ref{fig:anatomi-bukti} memberi nama ketiga bagian bukti itu: asumsi, langkah,
kesimpulan.

\paragraph{Pola kedua: bukti kontraposisi.} Kita membuktikan $\neg Q \Rightarrow \neg P$
sebagai ganti $P \Rightarrow Q$. Menurut akibat pada bagian pertama, keduanya setara,
sehingga sah. Pola ini dipilih ketika negasi kesimpulan lebih mudah dipegang daripada
hipotesisnya.

\begin{proposisi}
Untuk setiap bilangan bulat $n$, jika $5n+3$ genap maka $n$ ganjil.
\end{proposisi}
\par\smallskip\noindent\textit{Bukti.} Kontraposisinya berbunyi: jika $n$ genap maka
$5n+3$ ganjil. Misalkan $n = 2k$. Maka $5n+3 = 10k+3 = 2(5k+1)+1$, yang berbentuk dua
kali bilangan bulat ditambah satu, jadi ganjil. Karena kontraposisi terbukti,
implikasi aslinya terbukti.
\hfill$\blacksquare${\parfillskip=0pt\par}

Membuktikan langsung akan menuntut pembagian oleh $5$ dari $5n+3 = 2m$; bukan mustahil,
tetapi lebih panjang. Pola kontraposisi dipilih karena menghemat.

Sebelum pola ketiga, kita butuh satu lemma tentang paritas. Lemma ini akan dipakai dua
kali: sekali di sini, sekali pada bukti ketakrasionalan akar dua.

\begin{lemma}
\index{lemma paritas}
Untuk setiap bilangan bulat $n$, jika $n^{2}$ genap maka $n$ genap.
\end{lemma}
\par\smallskip\noindent\textit{Bukti.} Kita buktikan kontraposisinya: jika $n$ ganjil
maka $n^{2}$ ganjil. Menurut teorema pertama bagian ini, untuk $n$ ganjil berlaku
$n^{2} = 8m+1$ untuk suatu $m$, sehingga $n^{2}$ bersisa $1$ bila dibagi $2$, yaitu
ganjil.
\hfill$\blacksquare${\parfillskip=0pt\par}

\paragraph{Pola ketiga: bukti kontradiksi.} Kita mengandaikan kesimpulan salah, lalu menurunkan dua pernyataan yang tidak dapat benar
bersama. Pola ini dipakai ketika objek yang disangkal tidak punya bentuk yang nyaman
dihitung, misalnya ketika yang disangkal adalah keberadaan sesuatu.

\begin{teorema}
\index{akar dua}
Tidak ada bilangan rasional $r$ dengan $r^{2} = 2$.
\end{teorema}
\par\smallskip\noindent\textit{Bukti.} Andaikan ada $r \in \mathbb{Q}$ dengan
$r^{2} = 2$. Setiap bilangan rasional dapat ditulis sebagai $p/q$ dengan $p,q$ bilangan
bulat, $q \neq 0$, dan $p,q$ tidak memiliki faktor persekutuan selain $1$; kita
membuang faktor persekutuan itu dengan membagi keduanya, sehingga bentuk tanpa faktor
persekutuan selalu dapat dicapai.

Dari $p^{2}/q^{2} = 2$ diperoleh $p^{2} = 2q^{2}$, sehingga $p^{2}$ genap. Menurut
lemma di atas, $p$ genap; tulis $p = 2m$ dengan $m \in \mathbb{Z}$. Substitusi memberi
$4m^{2} = 2q^{2}$, yaitu $q^{2} = 2m^{2}$, sehingga $q^{2}$ genap dan, sekali lagi
menurut lemma, $q$ genap. Kini $p$ dan $q$ keduanya genap, jadi keduanya habis dibagi
$2$, bertentangan dengan pilihan $p/q$ tanpa faktor persekutuan. Pengandaian salah,
maka tidak ada $r$ rasional dengan $r^{2} = 2$.
\hfill$\blacksquare${\parfillskip=0pt\par}

Kontradiksinya bukan antara $r^{2}=2$ dan aritmetika biasa, melainkan antara ``$p,q$
keduanya genap'' dan janji awal ``$p,q$ tanpa faktor persekutuan''. Karena itu langkah
membuang faktor persekutuan bukan hiasan; ia bagian dari hipotesis yang nanti dilanggar.

\paragraph{Pola keempat: bukti konstruktif dengan saksi.} Untuk membuktikan pernyataan
berbentuk $\exists$, kita menunjuk satu objek dan memeriksa bahwa objek itu memenuhi
syarat. Objek yang ditunjuk disebut \emph{saksi} (\emph{witness}).

\begin{teorema}
Untuk setiap bilangan asli $n$ terdapat $n$ bilangan bulat berurutan yang semuanya
bukan prima.
\end{teorema}
\par\smallskip\noindent\textit{Bukti.} Ambil $N = (n+1)!$ dan perhatikan $n$ bilangan
\[
  N+2,\; N+3,\; \dots,\; N+(n+1).
\]
Untuk setiap $k$ dengan $2 \leq k \leq n+1$, bilangan $k$ membagi $(n+1)!$ karena $k$
adalah salah satu faktornya, sehingga $k$ juga membagi $N+k$. Selain itu
$N+k > k > 1$, jadi $N+k$ memiliki pembagi positif $k$ yang bukan $1$ dan bukan
dirinya sendiri. Menurut definisi, $N+k$ bukan prima. Kita telah menunjuk $n$ bilangan
yang diminta, dan tidak perlu mencari yang terkecil.
\hfill$\blacksquare${\parfillskip=0pt\par}

Untuk $n = 5$ bukti itu menunjuk $6! = 720$, lalu $722, 723, 724, 725, 726$. Kelimanya
memang terurai: $722 = 2 \cdot 19^{2}$, $723 = 3 \cdot 241$, $724 = 2^{2} \cdot 181$,
$725 = 5^{2} \cdot 29$, dan $726 = 2 \cdot 3 \cdot 11^{2}$. Bilangan $721 = 7 \cdot 103$ juga terurai, tetapi ia tidak termasuk daftar kita.

Bukti konstruktif tidak menjanjikan saksi terkecil atau tercepat; yang dinilai hanya
apakah objek yang ditunjuk memenuhi syarat. Menyangkal pernyataan berbentuk $\exists$
menuntut ditunjukkannya bahwa tidak satu pun objek memenuhi, dan di situlah kontradiksi
dibutuhkan.

\begin{resep}
Empat langkah baku membuktikan ``jika $P$ maka $Q$''.
\begin{enumerate}
  \item Tulis $P$ dan $Q$ menurut definisinya, masing-masing satu baris.
  \item Buka dengan ``andaikan $P$ benar''. Yang boleh diandaikan hanya $P$; $Q$ belum
        boleh dipakai.
  \item Turunkan $Q$ dengan substitusi dan aljabar, satu langkah per baris, dan sebut
        definisi yang dipakai.
  \item Tutup dengan kalimat yang menyebut $Q$ persis seperti bentuknya di langkah 1.
\end{enumerate}
Peringatan urutan: langkah 2 tidak boleh ditukar dengan langkah 3. Coretan dari kesimpulan
ke belakang boleh, tetapi bukti yang diserahkan harus berjalan dari $P$ ke $Q$; kalau
tidak, yang terbukti adalah konversnya.
\end{resep}

\section{Induksi dan induksi kuat}
\tingkat{menengah}

Induksi adalah pola bukti untuk pernyataan yang berderet, satu untuk setiap bilangan
asli. Kita menegakkan kartu pertama dan memastikan setiap kartu yang jatuh menjatuhkan
kartu berikutnya; setelah keduanya, tidak ada kartu yang boleh berdiri.

\kamus{induksi}{induction}{seperti deretan domino: cukup pastikan kartu pertama jatuh
dan setiap kartu menjatuhkan kartu di depannya, lalu seluruh deret jatuh tanpa perlu
didorong satu per satu.}

Prinsip induksi kita ambil sebagai sifat dasar $\mathbb{N} = \{1,2,3,\dots\}$, setara
dengan pernyataan bahwa setiap himpunan bagian tak kosong dari $\mathbb{N}$ mempunyai
anggota terkecil; tidak ada yang dibuktikan secara melingkar, kita hanya memilih titik
berangkat seperti pada Bab~3.

\begin{catatan}[Prinsip induksi]
Misalkan $P(n)$ pernyataan tentang bilangan asli. Jika $P(1)$ benar dan, untuk setiap
$k \in \mathbb{N}$, kebenaran $P(k)$ memaksa kebenaran $P(k+1)$, maka $P(n)$ benar
untuk setiap $n \in \mathbb{N}$.
\end{catatan}

Dengan prinsip itu, kita buktikan rumus yang sudah ditebak pada pembuka bagian ini.

\begin{teorema}
Untuk setiap bilangan asli $n$ berlaku
\[
  1 + 3 + 5 + \cdots + (2n-1) = n^{2}.
\]
\end{teorema}
\par\smallskip\noindent\textit{Bukti.} Basis: untuk $n = 1$ ruas kiri bernilai $1$ dan
ruas kanan bernilai $1^{2} = 1$. Langkah: andaikan rumus benar untuk $n = k$, yaitu
$1+3+\cdots+(2k-1) = k^{2}$. Suku berikutnya adalah $2(k+1)-1 = 2k+1$, sehingga
\[
  1+3+\cdots+(2k-1)+(2k+1) = k^{2} + 2k + 1 = (k+1)^{2}.
\]
Jadi rumus benar untuk $k+1$. Menurut prinsip induksi, rumus benar untuk setiap
$n \in \mathbb{N}$.
\hfill$\blacksquare${\parfillskip=0pt\par}

Contoh kedua berbentuk pertidaksamaan. Dari $2^{k} \geq k+1$ dan syarat $k \geq 0$
diperoleh $2^{k+1} = 2 \cdot 2^{k} \geq 2(k+1) \geq k+2$, sehingga dengan basis $n = 0$
berlaku $2^{n} \geq n+1$ untuk setiap $n \geq 0$; hipotesis dan syarat samping dipakai
sekaligus.

Induksi biasa menuntut satu langkah mundur. Kadang kesimpulan untuk $n$ hanya turun dari
banyak nilai sebelumnya sekaligus, dan di situ induksi biasa tidak cukup.

\begin{catatan}[Prinsip induksi kuat]
Jika untuk setiap $n \in \mathbb{N}$ kebenaran $P(1), P(2), \dots, P(n-1)$ memaksa
kebenaran $P(n)$, maka $P(n)$ benar untuk setiap $n \in \mathbb{N}$.
\end{catatan}

\kamus{induksi kuat}{strong induction}{seperti menyusun menara dari balok: untuk
meletakkan balok kesepuluh, kita boleh memakai seluruh balok di bawahnya, bukan hanya
balok kesembilan.}

Dengan izin itu kita dapat membuktikan hal yang tidak terjangkau induksi biasa.

\begin{teorema}
Setiap bilangan bulat $n \geq 2$ dapat ditulis sebagai hasil kali bilangan-bilangan
prima.
\end{teorema}
\par\smallskip\noindent\textit{Bukti.} Kita pakai induksi kuat pada $n$. Untuk $n = 2$:
bilangan $2$ sendiri prima, jadi ia hasil kali dengan satu faktor. Sekarang ambil
$n \geq 3$ dan andaikan pernyataan benar untuk semua bilangan bulat $m$ dengan
$2 \leq m < n$. Jika $n$ prima, selesai. Jika $n$ bukan prima, menurut definisi terdapat
pembagi positif $a$ dengan $a \neq 1$ dan $a \neq n$; tulis $n = ab$ dengan
$b = n/a$. Karena $a \neq 1$ dan $a \neq n$, berlaku $2 \leq a \leq n-1$, dan begitu
pula $2 \leq b \leq n-1$. Hipotesis induksi berlaku untuk $a$ maupun $b$, jadi keduanya
hasil kali bilangan prima, dan hasil kali keduanya adalah $n$.
\hfill$\blacksquare${\parfillskip=0pt\par}

Hipotesis induksi kuat dipakai dua kali, pada $a$ dan pada $b$. Di situlah induksi biasa
tidak nyaman: keduanya bisa jauh lebih kecil daripada $n-1$, sehingga pernyataan untuk
$n-1$ saja tidak menutupnya.

Teorema itu menutup janji bagian sebelumnya: daftar prima tidak pernah lengkap.

\begin{akibat}[Tak berhingga banyak bilangan prima]
\index{prima tak berhingga}
Untuk setiap daftar berhingga bilangan prima $p_{1}, p_{2}, \dots, p_{k}$ terdapat
bilangan prima yang tidak ada di dalam daftar itu.
\end{akibat}
\par\smallskip\noindent\textit{Bukti.} Misalkan $N = p_{1}p_{2}\cdots p_{k} + 1$. Karena
$N \geq 2$, teorema sebelumnya memberi bahwa $N$ adalah hasil kali bilangan prima;
ambil satu faktor prima $q$ dari $N$. Andaikan $q$ termasuk daftar, misalnya
$q = p_{i}$. Maka $q$ membagi hasil kali $p_{1}\cdots p_{k}$, dan karena $q$ juga
membagi $N$, bilangan $q$ membagi selisih $N - p_{1}\cdots p_{k} = 1$. Tidak ada
bilangan prima yang membagi $1$, jadi pengandaian itu salah. Maka $q$ tidak ada di
dalam daftar.
\hfill$\blacksquare${\parfillskip=0pt\par}

Bukti itu tidak membuktikan bahwa $N$ prima, hanya bahwa $N$ punya faktor prima di luar
daftar: untuk $2,3,5,7$ diperoleh $N = 211$ yang kebetulan prima, untuk $2,3,5,7,11,13$
diperoleh $N = 30031 = 59 \cdot 509$. Bukti itu konstruktif; ``andaikan hanya ada
berhingga banyak prima'' hanyalah pembungkus retorisnya.

\begin{untukanak}
Bayangkan deretan seribu kotak, dan kita ingin membuktikan setiap kotak berisi
kelereng. Kita tidak perlu membuka seribu kotak. Cukup buka kotak pertama, lalu buktikan
satu hal: kalau sebuah kotak berisi kelereng, kotak di sebelahnya juga berisi kelereng.
Kalau dua hal itu benar, seribu kotak selesai diperiksa tanpa dibuka satu per satu.

Inilah induksi, dan ia terasa seperti sulap. Kuncinya pada janji kedua: kalau hubungan
antarkotak bocor di satu tempat saja, seluruh deretan roboh.

Sekarang bayangkan kotak pertama ternyata kosong. Hubungan antarkotak tetap benar, tetapi
tidak menolong, karena tidak ada yang pertama untuk dijatuhkan. Inilah kesalahan yang
paling sering terjadi.

Di bab ini kita akan tahu mengapa dua langkah itu tidak boleh ditukar, dan mengapa
pemeriksaan pada banyak kotak bukan pengganti keduanya.
\end{untukanak}

\section{Ketika bukti gagal}
\tingkat{lanjut}

Kembali ke papan tulis pada pembuka bab. Empat puluh nilai pertama prima, dan pada
$n = 40$ pola itu patah:
\[
  40^{2} + 40 + 41 = 1600 + 40 + 41 = 1681 = 41^{2}.
\]
Kegagalan itu bukan kebetulan, dan bentuknya dapat dijelaskan tanpa komputer. Untuk
setiap bilangan bulat $p \geq 2$, substitusi $n = p-1$ memberi
\[
  (p-1)^{2} + (p-1) + p = p^{2} - 2p + 1 + p - 1 + p = p^{2},
\]
sedangkan substitusi $n = p$ memberi $p^{2} + p + p = p(p+2)$. Jadi setiap bilangan
$p \geq 2$ menghasilkan dua titik gagal, yaitu $n = p-1$ dan $n = p$. Dengan $p = 41$
keduanya menghasilkan $1681 = 41^{2}$ dan $1763 = 41 \cdot 43$.

Pola yang sama muncul pada polinomial yang lebih kecil. Untuk $n^{2}+n+11$ titik gagal
pertama $n = 10$, yang memberi $121 = 11^{2}$; untuk $n^{2}+n+17$ titik gagal pertama
$n = 16$, yang memberi $289 = 17^{2}$; untuk $n^{2}+n+29$ kegagalan datang pada $n = 2$,
karena $35 = 5 \cdot 7$ bukan prima. Semakin besar bilangan pada posisi konstanta,
semakin panjang rentang awal yang bersih; itulah seluruh rahasia reputasi $41$.

Kegagalan itu memberi tiga pelajaran yang sering tertukar.

Pelajaran pertama bersifat asimetris. Satu contoh penyangkal menumbangkan pernyataan
``untuk semua'', sedangkan empat puluh contoh yang berhasil tidak membuktikan apa pun
tentang nilai keempat puluh satu. Pemeriksaan komputer karena itu alat menemukan dugaan,
bukan alat membuktikan.

Pelajaran kedua menyangkut letak cacat. Cacat pertama yang paling sering lolos muncul
pada bukti induksi. Perhatikan klaim yang salah berikut: untuk setiap $n \geq 1$
berlaku $1+3+\cdots+(2n-1) = n^{2}+1$. Langkah induksinya, bila diperiksa, ternyata
sah:
\[
  \bigl(n^{2}+1\bigr) + (2n+1) = n^{2} + 2n + 2 = (n+1)^{2} + 1 .
\]
Jadi dari kebenaran rumus pada $n$ memang mengalir kebenaran pada $n+1$. Yang gagal
adalah basis: untuk $n = 1$ ruas kiri bernilai $1$ sedangkan ruas kanan bernilai $2$.
Langkah induksi memindahkan kebenaran, ia tidak menciptakannya.

Pelajaran ketiga menyangkut batas pemeriksaan. Uji pembagian lazim: $m > 1$ prima bila
tidak habis dibagi bilangan prima mana pun yang lebih kecil daripada $\sqrt{m}$. Untuk $m = 1681$ diperoleh $\sqrt{1681} = 41$, dan prima yang
lebih kecil daripada $41$ adalah $2, 3, 5, 7, 11, 13, 17, 19, 23, 29, 31, 37$; tidak
satu pun membagi $1681$. Dengan aturan itu $1681$ lolos sebagai prima. Kesalahannya ada
pada kata ``lebih kecil'': aturan yang benar memakai ``lebih kecil atau sama dengan'',
karena $1681 = 41 \cdot 41$ dan pembaginya tepat sebesar akar. Satu tanda pertidaksamaan
yang longgar menghasilkan satu prima palsu, tepat di $n = 40$.

Tiga cacat lain berulang di tulisan yang lebih panjang. Cacat penalaran melingkar:
``karena $n^{2}+n+41$ prima untuk semua $n$, maka prima juga untuk $n = 40$'' memakai
yang harus dibuktikan sebagai alasan. Cacat membagi dengan nol: dari $a = b$ orang
menyimpulkan $a+b = b$ setelah membagi dengan $a-b$, padahal $a-b = 0$. Cacat menukar
arah: dari $P \Rightarrow Q$ dan $Q$ orang menyimpulkan $P$; pada tabel kebenaran,
baris ketiga adalah tempat kesalahan itu bersembunyi.

Keempat cacat itu berbagi satu akar: sebuah langkah dipakai di luar syaratnya.
Kebiasaan yang berguna bukan menghafal pola, melainkan memeriksa syarat setiap langkah:
pada kontradiksi pengandaiannya harus negasi kesimpulan, pada induksi ada basis dan
langkah, pada pembagian pembaginya tidak nol.

\begin{tantangan}
Klaim: untuk setiap bilangan bulat $k \geq 0$, bilangan $2^{2^{k}}+1$ prima. Untuk
$k = 0,1,2,3,4$ nilainya $3$, $5$, $17$, $257$, $65537$, dan kelimanya memang prima.
Tentukan $k$ terkecil yang menggugurkan klaim itu, lalu tunjukkan pembaginya; klaim itu
salah, dan kesalahannya dapat ditunjukkan tanpa komputer.

Petunjuk 1: lima nilai bukan bukti. Periksa $k = 5$; bilangannya $4294967297$, dan pembagi
terkecilnya jauh lebih kecil daripada akar kuadratnya.

Petunjuk 2: pembagi itu berbentuk $5 \cdot 2^{7}+1 = 641$, dan
$4294967297 = 641 \cdot 6700417$.

Kalau mentok, kembali ke bagian ``Ketika bukti gagal''.
\end{tantangan}

\section{Contoh hitung}
\tingkat{menengah}

Tiga perhitungan berikut menopang bab ini dan dapat diulang di kertas. Yang pertama
jumlah bilangan ganjil pertama, suku demi suku dari $n = 1$.

{\centering
\begin{tabular}{cccc}
\hline
$n$ & suku-suku & jumlah & $n^{2}$ \\
\hline
$1$ & $1$ & $1$ & $1$ \\
$2$ & $1+3$ & $4$ & $4$ \\
$3$ & $1+3+5$ & $9$ & $9$ \\
$4$ & $1+3+5+7$ & $16$ & $16$ \\
$5$ & $1+3+5+7+9$ & $25$ & $25$ \\
$6$ & $1+3+5+7+9+11$ & $36$ & $36$ \\
\hline
\end{tabular}\par}

Lanjutannya: $36+13 = 49 = 7^{2}$, lalu $49+15 = 64$, $64+17 = 81$, $81+19 = 100$.
Selisih dua kuadrat berurutan memang $(n+1)^{2}-n^{2} = 2n+1$, dan di situlah induksi
berpijak.

Yang kedua uji primalitas pada $n^{2}+n+41$, hanya pada nilai yang menarik.

{\centering
\begin{tabular}{ccc}
\hline
$n$ & $n^{2}+n+41$ & keterangan \\
\hline
$0$ & $41$ & prima \\
$1$ & $43$ & prima \\
$39$ & $1601$ & prima \\
$40$ & $1681 = 41^{2}$ & bukan prima \\
$41$ & $1763 = 41 \cdot 43$ & bukan prima \\
\hline
\end{tabular}\par}

Pemeriksaan menyeluruh untuk $n = 0$ sampai $n = 40$ memberi $40$ nilai prima dari $41$;
satu-satunya yang gagal adalah $n = 40$, dan $\sqrt{1681} = 41$ tepat, sehingga pembagi
yang menentukan berada di batas pemeriksaan.

Yang ketiga pemeriksaan kontraposisi pada tabel kebenaran, ditulis lengkap.

{\centering
\begin{tabular}{cc|cc|c}
\hline
$P$ & $Q$ & $P \Rightarrow Q$ & $\neg Q \Rightarrow \neg P$ & setara \\
\hline
B & B & B & B & ya \\
B & S & S & S & ya \\
S & B & B & B & ya \\
S & S & B & B & ya \\
\hline
\end{tabular}\par}

Keempat baris memberi kesamaan: pola kontraposisi bukan trik retoris melainkan penggantian
bentuk yang tidak mengubah nilai. Kolom $Q \Rightarrow P$ akan berbeda pada baris ketiga,
dan perbedaan itulah yang membuat konvers tidak boleh dipakai.

\kertas{Hitung $1+3+\cdots+(2n-1)$ untuk $n = 11$ dan $n = 12$ dengan bertolak dari
$100 = 10^{2}$, yaitu tambahkan dua bilangan ganjil berikutnya satu per satu. Lalu
hitung $n^{2}+n+17$ untuk $n = 16$ dan periksa dengan pembagian apakah $289$ sama
dengan $17^{2}$. Terakhir, cari titik gagal pertama $n^{2}+n+11$ dengan menguji
$n = 0, 1, 2, \dots$ dan jelaskan mengapa titik itu dapat diramalkan dari bentuk
polinomialnya.}

\section{Lab komputasi}
\tingkat{lanjut}

Skrip berikut hanya memakai pustaka standar Python dan dapat dijalankan ulang tanpa
pemasangan apa pun; setiap baris keluarannya adalah keluaran asli.

\begin{lab}
\begin{verbatim}
from itertools import product

def prima(n):
    if n < 2:
        return False
    d = 2
    while d * d <= n:
        if n % d == 0:
            return False
        d += 1
    return True

pasang = product([False, True], repeat=2)
print("kontraposisi:", all((not P or Q) == (Q or not P)
                          for P, Q in pasang))

D = range(1, 11)
print("y>x:", sum(y > x for x in D for y in D), "dari 100")
print("x tanpa saksi:", [x for x in D if all(y <= x for y in D)])

nilai = [(n, n*n + n + 41) for n in range(41)]
print("prima:", sum(prima(v) for _, v in nilai), "dari 41")
print("gagal:", [(n, v, int(v**0.5)) for n, v in nilai if not prima(v)])
print("41*41, 41*43:", 41*41, 41*43)

sisa = {(n*n - 1) % 8 for n in range(1, 40, 2)}
print("sisa mod 8:", sisa)

langkah = all(n*n + 1 + 2*n + 1 == (n+1)**2 + 1 for n in range(1, 500))
print("langkah palsu:", langkah)
\end{verbatim}
\end{lab}

Keluaran yang diperoleh:

\begin{verbatim}
kontraposisi: True
y>x: 45 dari 100
x tanpa saksi: [10]
prima: 40 dari 41
gagal: [(40, 1681, 41)]
41*41, 41*43: 1681 1763
sisa mod 8: {0}
langkah palsu: True
\end{verbatim}

Baris pertama memeriksa keempat baris tabel sekaligus. Baris kedua dan ketiga menghitung
yang tidak nyaman dikerjakan di kertas: dari seratus pasangan $(x,y)$ pada
$\{1,\dots,10\}$, $45$ memenuhi $y > x$, dan satu-satunya $x$ tanpa saksi adalah $10$;
itulah kegagalan $(\forall x)(\exists y)$ pada himpunan berhingga.

Baris keempat dan kelima memberi $40$ prima dari $41$ nilai, gagal hanya di
$(40,1681,41)$; angka $41$ itu adalah $\sqrt{1681}$ yang dipotong. Baris keenam
memastikan $1681 = 41^{2}$ dan $1763 = 41 \cdot 43$. Baris ketujuh menyatakan seluruh
$n$ ganjil di bawah $40$ memberi sisa $0$ modulo $8$, sesuai teorema pola pertama. Baris
terakhir peringatan: langkah induksi klaim yang salah itu sah untuk lima ratus nilai
pertama, sehingga pemeriksaan numerik tidak menemukan cacatnya.

\begin{untukdosen}
Urutan mengajar yang berhasil dipakai: mulai dari $n = 40$ dan $1681$, bukan dari
definisi. Setelah mahasiswa kesal karena empat puluh pemeriksaan ternyata sia-sia, baru
masuk tabel kebenaran, dengan penekanan pada baris ketiga dan keempat. Lanjutkan ke
negasi kuantor lewat definisi limit barisan, lalu ke definisi harfiah dengan contoh
$91 = 7 \cdot 13$; tutup dengan induksi, kegagalan basis, dan lab.

Dua miskonsepsi tersering. Pertama, implikasi dibaca sebagai sebab-akibat, sehingga baris
dengan anteseden salah dinilai salah dan bukti berkasus hampa ditolak. Kedua, induksi
dianggap selesai setelah langkah terbukti, dan induksi kuat dianggap induksi biasa yang
diulang; akibatnya basis tidak diperiksa dan hipotesis dipakai tanpa menyebut
rentangnya.

Aktivitas papan tulis: bagi kelas menjadi dua kelompok. Kelompok pertama menulis klaim
berkuantor, kelompok kedua menulis negasinya, lalu masing-masing menyebut satu objek yang
menyangkal klaim lawan. Permainan berhenti ketika sebuah negasi tidak dapat disangkal.

Mahasiswa siap lanjut apabila dapat menegasikan definisi limit tanpa membuka catatan
dan dapat menjelaskan mengapa basis pada induksi tidak boleh dilewati.
\end{untukdosen}

\section{Ringkasan}
\tingkat{dasar}

\begin{ringkasan}
\begin{itemize}
  \item Pernyataan bernilai benar atau salah; $P \Rightarrow Q$ hanya salah ketika $P$
        benar dan $Q$ salah.
  \item $P \Rightarrow Q \equiv \neg Q \Rightarrow \neg P$, tetapi konversnya tidak
        setara; $\neg(P \Rightarrow Q) \equiv P \wedge \neg Q$.
  \item Negasi berjalan dari luar ke dalam: setiap kuantor berganti jenis, sehingga
        menyangkal ``untuk semua'' berarti menunjuk satu contoh penyangkal.
  \item Definisi berlaku dua arah: nama hanya boleh dipakai bila seluruh klausa
        terpenuhi, dan seluruh klausa dipakai untuk memeriksanya.
  \item Empat pola bukti: langsung, kontraposisi, kontradiksi, konstruktif. Pilih
        yang paling pendek, bukan yang paling megah.
  \item Induksi menuntut basis dan langkah; induksi kuat menuntut seluruh nilai
        sebelumnya. Langkah tanpa basis tidak membuktikan apa pun.
\end{itemize}
\end{ringkasan}

\section{Latihan}
\tingkat{menengah}

\begin{latihan}
\begin{enumerate}
  \item Tulis negasi dari pernyataan berikut, lalu terjemahkan hasilnya menjadi
        kalimat Indonesia yang menyatakan apa yang harus ditunjukkan untuk
        menyanggahnya: ``untuk setiap $\varepsilon > 0$ terdapat bilangan asli $N$
        sehingga untuk setiap $n \geq N$ berlaku $|a_{n} - L| < \varepsilon$''.

  \item Tulis negasi dari ``setiap bilangan prima lebih besar daripada $2$ adalah
        ganjil'', lalu sebutkan satu objek yang menyangkal pernyataan aslinya.
        Kerjakan hal yang sama untuk ``terdapat bilangan asli $n$ sehingga
        $n^{2}+n+41$ tidak prima''.

  \item Buktikan bahwa untuk setiap bilangan bulat $n$, bilangan $n^{2}+n$ selalu
        genap. Kerjakan dengan dua cara: pembagian kasus paritas, dan kontraposisi.

  \item Buktikan dengan induksi bahwa untuk setiap bilangan asli $n$ berlaku
        \[
          1^{2} + 2^{2} + \cdots + n^{2} = \frac{n(n+1)(2n+1)}{6}.
        \]
        Periksa rumus itu untuk $n = 1, 2, 3$ sebelum menyusun buktinya.

  \item Buktikan dengan induksi kuat bahwa setiap bilangan asli $n$ dapat ditulis
        dalam bentuk $n = 2^{k}m$ dengan $k \geq 0$ bilangan bulat dan $m$ bilangan
        ganjil. Tunjukkan di bagian mana hipotesis induksi kuat dipakai.

  \item Nilai klaim berikut: benar, salah, atau belum cukup bukti.
        \begin{enumerate}
          \item Jika $P(1)$ benar dan $P(k) \Rightarrow P(k+1)$ untuk setiap
                $k \in \mathbb{N}$, maka $P(n)$ benar untuk setiap $n \in
                \mathbb{N}$.
          \item Jika $P(n)$ benar untuk $n = 1, 2, \dots, 10^{6}$, maka $P(n)$ benar
                untuk setiap $n \in \mathbb{N}$.
          \item Jika $n^{2}+n+41$ prima untuk $n = 0, 1, \dots, 39$, maka klaim itu
                benar untuk setiap $n \geq 0$.
          \item Jika $3$ membagi $n$, maka $6$ membagi $n$.
        \end{enumerate}
        Untuk setiap butir yang salah, berikan satu contoh penyangkal konkret.

  \item Hitung dengan Python banyaknya bilangan prima di antara $n^{2}+n+41$ untuk
        $n = 0$ sampai $n = 100$, lalu tentukan semua $n$ pada rentang itu yang
        membuat nilainya bukan prima. Periksa apakah setiap titik gagal itu berbentuk
        $n = p-1$ atau $n = p$ untuk suatu bilangan prima $p$.

  \item Untuk $n = 1$ sampai $n = 12$, hitung $2^{n}$ dan $n^{2}$, lalu tentukan
        bilangan asli terkecil $n_{0}$ sehingga $2^{n} > n^{2}$ berlaku untuk setiap
        $n \geq n_{0}$. Buktikan klaim itu dengan induksi, dengan basis $n_{0}$.

  \item Buktikan bahwa pernyataan ``terdapat bilangan bulat $n > 1$ dengan
        $n^{2}-1$ prima'' salah. Sebutkan lebih dulu bentuk faktorisasi $n^{2}-1$
        untuk $n > 1$, lalu jelaskan mengapa faktor-faktornya tidak mungkin sama
        dengan $1$.

  \item Tunjukkan bahwa pernyataan ``jika $n^{2}$ habis dibagi $4$ maka $n$ habis
        dibagi $4$'' salah, lalu perbaiki kesimpulannya menjadi klaim yang benar dan
        buktikan.
\end{enumerate}

\par\smallskip\noindent\textit{Petunjuk untuk tiga butir tersulit.} Butir 5: pisahkan
kasus $n$ ganjil, yang langsung selesai dengan $k=0$ dan $m=n$, dari kasus $n$ genap,
yang memakai hipotesis induksi pada $n/2$. Butir 7: hitung juga $\sqrt{n^{2}+n+41}$
untuk setiap titik gagal, dan bandingkan dengan pembagi terkecilnya. Butir 9: $n^{2}-1 = (n-1)(n+1)$; untuk $n > 1$ kedua faktor itu terletak di antara $1$
dan $n^{2}-1$.
\end{latihan}

\begin{untukpeneliti}
Yang benar-benar terbukti di bab ini tiga hal klasik yang bertahan tanpa komputer:
ketakrasionalan $\sqrt{2}$ dengan bukti paritas, tak berhingga banyak prima lewat
konstruksi Euklides, dan ketunggalan bentuk setiap bilangan asli sebagai hasil kali prima.
Yang hanya dipopulerkan adalah cerita bahwa $41$ ``hampir'' rumus prima; yang terjadi
hanyalah rentang awal yang panjang, dan panjangnya dapat diramalkan dari $n^{2}+n+p$.

Sumber primer: Abbott, \emph{Understanding Analysis}, Bab~1 Bagian~1.1 (irasionalitas
$\sqrt{2}$) dan Bagian~1.2 (pendahuluan, termasuk induksi); Avigad dan Massot,
\emph{Mathematics in Lean}, Bab~3 (logika) serta Bab~5 Bagian~5.1 dan 5.2
(akar irasional, induksi dan rekursi); Euklides, \emph{Elements}, Buku~IX
Proposisi~20 untuk tak berhingga prima.

Pertanyaan yang masih terbuka: tidak ada polinomial berderajat positif dengan koefisien
bulat yang bernilai prima untuk semua $n \geq 0$ (bukti satu baris: $f(0)$ membagi
$f(t\,f(0))$), tetapi seberapa panjang rentang prima yang dapat dicapai polinomial
berderajat tinggi masih dipelajari, dan distribusi prima terikat pada hipotesis Riemann
yang belum terbukti.

Batas keberlakuan: semua hasil bab ini elementer dan tidak bergantung pada pilihan aksioma
himpunan. Verifikasi formal di Lean (Bab~28) memeriksa keabsahan setiap langkah deduktif,
tetapi tidak memeriksa apakah pernyataan formal yang dibuktikan memang yang dimaksud
penulisnya.
\end{untukpeneliti}

\peta{Bab 1 pada \emph{Understanding Analysis} (Abbott), khususnya Bagian 1.1 tentang
irasionalitas $\sqrt{2}$ dan Bagian 1.2 tentang pendahuluan; Bab 3 (logika), Bab 5
Bagian 5.1 dan 5.2 (akar irasional serta induksi dan rekursi) pada
\emph{Mathematics in Lean} (Avigad dan Massot); dan Buku IX Proposisi 20
\emph{Elements} (Euklides) untuk tak berhingga banyak bilangan prima.}

\chapter{Himpunan, Relasi, Fungsi, dan Kardinalitas}

Sebuah penguji keamanan menerima spesifikasi model: pengklasifikasi dengan tiga puluh
fitur biner, masing-masing hanya bernilai nol atau satu. Untuk menguji menyeluruh, ia
harus mencoba setiap masukan yang mungkin. Setiap fitur menggandakan jumlah
kemungkinan, sehingga jawabannya $2^{30} = 1.073.741.824$. Pada kecepatan satu juta
masukan per detik, pengujian itu memakan $1073{,}74$ detik, sekitar $17{,}9$ menit.
Tambahkan sepuluh fitur lagi: jumlahnya menjadi $2^{40} = 1.099.511.627.776$, dan
mesin yang seribu kali lebih cepat memerlukan $1099{,}51$ detik, sekitar $18{,}3$
menit. Kita menggandakan masalah seribu kali dan mempercepat mesin seribu kali, lalu
kembali ke titik awal.

Kesulitan yang sama muncul lebih tajam begitu fitur boleh bernilai sebarang bilangan
real. Jumlah masukan berhenti bersifat "banyak" dan mulai bersifat "tak terhingga",
dan pertanyaannya berubah: apakah dua himpunan tak terhingga masih dapat dibandingkan
ukurannya? Naluri mengatakan bahwa bilangan rasional lebih banyak daripada bilangan
asli, sebab di antara dua bilangan asli selalu ada tak berhingga bilangan rasional.
Naluri itu keliru, dan bab ini menunjukkan mengapa.

Bab sebelumnya memberi kita bahasa: kuantor, implikasi, pola bukti. Yang belum ada
adalah objek yang dibicarakan. Bab ini memasang perbendaharaan dasarnya, yaitu
himpunan dan operasinya, relasi, fungsi, dan satu cara membandingkan ukuran himpunan
yang tidak bergantung pada kemampuan menghitung. Perbendaharaan itu dipakai di
seluruh buku, dari subruang di Bab 6 sampai pembuktian formal di Bab 28.

\begin{untukanak}
Bayangkan dua keranjang kelereng, satu merah dan satu biru. Kita ingin tahu keranjang
mana yang lebih banyak, tetapi tidak boleh menghitung. Ambil satu kelereng merah,
pasangkan dengan satu kelereng biru, sisihkan keduanya, dan ulangi sampai salah satu
keranjang kosong: bila keduanya kosong bersamaan isinya sama banyak, bila biru bersisa
biru lebih banyak. Cara ini tidak pernah menyebut angka; ia hanya memasangkan.

Sekarang misalkan keranjangnya tak berujung, sehingga kita tidak akan pernah selesai
memindahkan kelereng. Apakah memasangkan masih bisa dipakai? Bisa, dan itulah kejutan
bab ini: dengan aturan pemasangan yang tepat, seluruh bilangan cacah dapat dipasangkan
habis dengan seluruh bilangan genap, tanpa sisa di kedua pihak, dan pasangan bilangan
dapat dipasangkan dengan satu bilangan tunggal. Kita akan tahu mengapa memasangkan
lebih kuat daripada menghitung, dan mengapa ada himpunan tak berujung yang tetap tidak
dapat dipasangkan dengan bilangan cacah, sekeras apa pun kita mencoba.
\end{untukanak}

\section{Operasi pada himpunan dan hukum De Morgan}
\tingkat{dasar}

Satu himpunan belum berguna sampai kita dapat menyusun himpunan baru darinya. Empat
cara penyusunan didefinisikan dengan pola sama: sebut syarat keanggotaan, lalu himpunan
baru adalah kumpulan semua objek yang memenuhinya. Misalkan $X$ himpunan semesta dan
$A, B \subseteq X$; dari Bab 1 kita ingat bahwa $A \subseteq B$ berarti setiap anggota
$A$ juga anggota $B$, dan kesamaan himpunan dibuktikan dengan dua inklusi.

\begin{definisi}[Gabungan, irisan, selisih, komplemen]
Gabungan $A \cup B$ adalah himpunan semua $x \in X$ yang menjadi anggota $A$ atau
anggota $B$. Irisan $A \cap B$ adalah himpunan semua $x \in X$ yang menjadi anggota
$A$ dan sekaligus anggota $B$. Selisih $A \setminus B$ adalah himpunan semua
$x \in X$ yang menjadi anggota $A$ tetapi bukan anggota $B$. Komplemen
$A^{c} = X \setminus A$ adalah himpunan semua $x \in X$ yang bukan anggota $A$.
\end{definisi}

Kata "atau" pada gabungan bersifat inklusif: $x$ boleh berada di kedua himpunan, dan
itulah syarat agar $A \cap B \subseteq A \cup B$ selalu benar, seperti dicatat Hoffman
dan Kunze di Apendiks. Himpunan kosong $\varnothing$ diperlukan supaya irisan dua
himpunan tak bersekutu tetap berupa himpunan.

\begin{proposisi}[Hukum De Morgan]\index{hukum De Morgan}
Untuk sebarang $A, B \subseteq X$ berlaku
\[
  X \setminus (A \cup B) = (X \setminus A) \cap (X \setminus B),
  \qquad
  X \setminus (A \cap B) = (X \setminus A) \cup (X \setminus B).
\]
\end{proposisi}
\par\smallskip\noindent\textit{Bukti.} Ambil $x \in X$ sebarang. Menurut definisi
komplemen dan gabungan,
\[
  x \in X \setminus (A \cup B)
  \iff x \in X \ \text{dan}\ \text{tidak}\ (x \in A \ \text{atau}\ x \in B).
\]
Menurut hukum De Morgan untuk pernyataan dari Bab 1, ingkaran "salah satu" adalah
"keduanya salah", sehingga
\[
  x \in X \setminus (A \cup B)
  \iff x \in X \ \text{dan}\ x \notin A \ \text{dan}\ x \notin B
  \iff x \in (X \setminus A) \cap (X \setminus B).
\]
Karena $x$ sebarang, kedua himpunan sama. Kesamaan kedua diperoleh dengan mengganti
setiap "dan" menjadi "atau" pada rantai yang sama.
\hfill$\blacksquare${\parfillskip=0pt\par}

Hukum itu berlaku untuk dua himpunan; yang kita perlukan nanti adalah versi untuk
sebanyak apa pun himpunan.

\begin{definisi}[Keluarga terindeks]
Misalkan $I$ himpunan indeks dan untuk setiap $i \in I$ diberikan $A_i \subseteq X$.
\begin{align*}
  \bigcup_{i \in I} A_i &= \{\, x \in X : x \in A_i \ \text{untuk suatu}\ i \in I \,\}, \\
  \bigcap_{i \in I} A_i &= \{\, x \in X : x \in A_i \ \text{untuk setiap}\ i \in I \,\}.
\end{align*}
\end{definisi}

Versi umumnya berbunyi $X \setminus \bigcup_{i} A_i = \bigcap_{i} (X \setminus A_i)$
dan $X \setminus \bigcap_{i} A_i = \bigcup_{i} (X \setminus A_i)$. Buktinya tidak
menambah gagasan baru: hanya "untuk suatu" berubah menjadi "untuk setiap" ketika
ingkaran ditarik ke dalam, persis seperti pada Bab 1.

Angka membuat definisi itu dapat diperiksa. Untuk $X = \{1,\dots,10\}$,
$A = \{1,\dots,5\}$, dan $B = \{4,\dots,8\}$, Gambar~\ref{fig:venn} menampilkan
ketiga operasi atas himpunan yang sama: gabungan beranggota $3 + 2 + 3 = 8$, irisan
beranggota $2$, dan komplemen $B$ beranggota $3 + 2 = 5$.

\begin{figure}[htbp]
\centering
\begin{tikzpicture}[scale=0.82,every node/.style={font=\footnotesize}]
  \begin{scope}
    \begin{scope}
      \clip (-0.5,0) circle (0.85);
      \fill[layerteal!30] (-1.55,-1.25) rectangle (1.55,1.25);
    \end{scope}
    \begin{scope}
      \clip (0.5,0) circle (0.85);
      \fill[layerteal!30] (-1.55,-1.25) rectangle (1.55,1.25);
    \end{scope}
    \draw[ink,thick] (-0.5,0) circle (0.85);
    \draw[ink,thick] (0.5,0) circle (0.85);
    \draw[ink,thick] (-1.55,-1.25) rectangle (1.55,1.25);
    \node at (-1.0,0) {3};
    \node at (0,0) {2};
    \node at (1.0,0) {3};
    \node at (0,-1.0) {2};
    \node[above] at (-0.5,0.86) {$A$};
    \node[above] at (0.5,0.86) {$B$};
    \node[below] at (0,-1.25) {$A \cup B$: 8};
  \end{scope}
  \begin{scope}[xshift=3.9cm]
    \begin{scope}
      \clip (-0.5,0) circle (0.85);
      \clip (0.5,0) circle (0.85);
      \fill[accent!35] (-1.55,-1.25) rectangle (1.55,1.25);
    \end{scope}
    \draw[ink,thick] (-0.5,0) circle (0.85);
    \draw[ink,thick] (0.5,0) circle (0.85);
    \draw[ink,thick] (-1.55,-1.25) rectangle (1.55,1.25);
    \node at (-1.0,0) {3};
    \node at (0,0) {2};
    \node at (1.0,0) {3};
    \node at (0,-1.0) {2};
    \node[above] at (-0.5,0.86) {$A$};
    \node[above] at (0.5,0.86) {$B$};
    \node[below] at (0,-1.25) {$A \cap B$: 2};
  \end{scope}
  \begin{scope}[xshift=7.8cm]
    \fill[layergrey!22] (-1.55,-1.25) rectangle (1.55,1.25);
    \fill[white] (0.5,0) circle (0.85);
    \draw[ink,thick] (-0.5,0) circle (0.85);
    \draw[ink,thick] (0.5,0) circle (0.85);
    \draw[ink,thick] (-1.55,-1.25) rectangle (1.55,1.25);
    \node at (-1.0,0) {3};
    \node[ink!40] at (0,0) {2};
    \node[ink!40] at (1.0,0) {3};
    \node at (0,-1.0) {2};
    \node[above] at (-0.5,0.86) {$A$};
    \node[above] at (0.5,0.86) {$B$};
    \node[below] at (0,-1.25) {$B^{c}$: 5};
  \end{scope}
\end{tikzpicture}
\caption{Tiga diagram Venn atas himpunan yang sama, dengan banyaknya anggota pada
tiap daerah. Kiri: $A \cup B$ beranggota $3 + 2 + 3 = 8$; tengah: $A \cap B$
beranggota $2$; kanan: $B^{c}$ beranggota $3 + 2 = 5$, yaitu daerah berarsir,
sedangkan anggota di dalam $B$ ditulis pudar.}
\label{fig:venn}
\end{figure}

\begin{tangan}
Bahan: dua puluh kartu nama bekas, dua potong tali sepanjang dua meter, lantai lapang.

Bentuk dua lingkaran tali yang bertumpang tindih kira-kira seperempat lebarnya. Tulis
nama lima teman penyuka kopi dan lima teman penyuka teh. Letakkan kartu kopi di
lingkaran pertama, kartu teh di lingkaran kedua, lalu pindahkan kartu nama teman yang
suka keduanya ke daerah tumpang tindih: itulah irisannya. Sebut satu nama yang tidak
suka keduanya dan letakkan kartunya di luar kedua lingkaran.

Hitung tiga angka: kartu di dalam kedua lingkaran, kartu di tumpang tindih, dan kartu
di luar lingkaran kedua. Yang harus terlihat: jumlah kartu di gabungan bukan jumlah
kartu kopi ditambah jumlah kartu teh, sebab kartu di tumpang tindih terhitung dua
kali, dan kelebihan itu tepat sebesar irisannya.
\end{tangan}

Himpunan yang paling sering kita bentuk dari $A$ adalah himpunan semua himpunan
bagiannya, dan ukurannya tumbuh jauh lebih cepat daripada $A$ sendiri.

\begin{definisi}[Himpunan kuasa]\index{himpunan kuasa}
Himpunan kuasa $\mathcal{P}(A)$ adalah himpunan semua himpunan bagian $A$, yaitu
$\mathcal{P}(A) = \{\, B : B \subseteq A \,\}$.
\end{definisi}

\kamus{himpunan kuasa}{power set}{Daftar semua cara memilih sebagian anggota sebuah kantong, termasuk memilih kosong dan memilih semuanya.}

Sebagai contoh, $\mathcal{P}(\varnothing) = \{\varnothing\}$ beranggota satu, dan
$\mathcal{P}(\{1,2\}) = \{\varnothing, \{1\}, \{2\}, \{1,2\}\}$ beranggota empat.

\begin{proposisi}
Jika $A$ berhingga dengan $|A| = n$, maka $|\mathcal{P}(A)| = 2^{n}$.
\end{proposisi}
\par\smallskip\noindent\textit{Bukti.} Untuk $B \subseteq A$ definisikan fungsi
karakteristik $\chi_B : A \to \{0,1\}$ oleh $\chi_B(a) = 1$ bila $a \in B$ dan
$\chi_B(a) = 0$ bila $a \notin B$. Pemetaan $B \mapsto \chi_B$ adalah bijeksi antara
$\mathcal{P}(A)$ dan himpunan semua fungsi $A \to \{0,1\}$: dari $\chi$ kita pulihkan
$B = \{a : \chi(a) = 1\}$, dan dua himpunan bagian berbeda memberi fungsi berbeda.
Jadi $|\mathcal{P}(A)|$ sama dengan banyaknya fungsi $A \to \{0,1\}$.

Fungsi itu kita hitung dengan induksi pada $n$. Untuk $n = 0$ hanya ada fungsi
kosong, dan $2^{0} = 1$. Misalkan $|A| = n+1$, pilih $a_0 \in A$, dan tulis
$A' = A \setminus \{a_0\}$. Setiap $f : A \to \{0,1\}$ ditentukan tunggal oleh batasan
$f|_{A'}$ dan nilai $f(a_0)$; hipotesis induksi memberi $2^{n}$ pilihan batasan dan $2$
pilihan nilai, seluruhnya $2^{n+1}$.
\hfill$\blacksquare${\parfillskip=0pt\par}

Bukti itu tidak menyebut kombinasi: bijeksi memindahkan hitungan, induksi memecahnya
menjadi dua pilihan bebas.

\begin{bacakeras}
Bacakan berpasangan, satu orang suara anggota dan satu orang suara himpunan bagian.
Mulai lambat, lalu percepat.

Anggota: Satu jeruk masuk keranjang. Saya di dalam.
Himpunan bagian: Satu keranjang masuk gudang. Saya di dalam.
Anggota: Saya barang.
Himpunan bagian: Saya wadah.
Anggota: Saya ditunjuk satu per satu.
Himpunan bagian: Saya dibuktikan sekaligus.
Anggota: Kalau saya masuk, saya satu tanda.
Himpunan bagian: Kalau saya masuk, saya satu bukti.
Anggota: Saya kecil dan banyak.
Himpunan bagian: Saya besar dan kadang hanya satu.

Ulangi sampai kalimat terakhir keluar tanpa berpikir. Baris yang diulang tiga kali
bersama: barang masuk kantong, kantong masuk kantong, dan keduanya memakai tanda
berbeda.
\end{bacakeras}

\section{Hasilkali kartesius dan relasi}
\tingkat{dasar}

Urutan penting. Himpunan $\{1,2\}$ sama dengan $\{2,1\}$, tetapi koordinat titik pada
bidang tidak boleh ditukar tanpa memindahkan titiknya. Karena itu kita memerlukan
objek yang mengingat urutan, dan definisinya harus dibuat dari himpunan saja.

\begin{definisi}[Pasangan terurut dan hasilkali kartesius]\index{hasilkali kartesius}
Pasangan terurut $(a,b)$ didefinisikan sebagai $\{\{a\}, \{a,b\}\}$. Hasilkali
kartesius $A \times B$ adalah himpunan semua pasangan terurut $(a,b)$ dengan $a \in A$
dan $b \in B$.
\end{definisi}

Definisi itu terlihat aneh, tetapi yang dibutuhkan hanya satu sifat, yaitu
$(a,b) = (c,d)$ jika dan hanya jika $a = c$ dan $b = d$. Pemeriksaan sifat itu
memerlukan empat kasus dan kita tinggalkan sebagai Latihan 4.

Ukuran hasilkali mudah ditebak, dan buktinya memakai pemotongan menurut koordinat
pertama.

\begin{proposisi}
Jika $A$ dan $B$ berhingga, maka $|A \times B| = |A| \cdot |B|$.
\end{proposisi}
\par\smallskip\noindent\textit{Bukti.} Untuk setiap $a \in A$, himpunan
$\{a\} \times B$ beranggota $|B|$, sebab $b \mapsto (a,b)$ bijeksi antara $B$ dan
$\{a\} \times B$. Keluarga $\{\{a\} \times B\}_{a \in A}$ terdiri atas himpunan
saling lepas, karena $(a,b) = (a',b')$ memaksa $a = a'$. Gabungannya $A \times B$,
jadi $|A \times B|$ adalah jumlah $|A|$ suku bernilai $|B|$.
\hfill$\blacksquare${\parfillskip=0pt\par}

Sekarang objek yang menjadi judul bab ini. Sebuah relasi hanyalah pernyataan yang
mungkin benar atau salah untuk setiap pasangan objek, sehingga dapat disimpan sebagai
himpunan pasangan yang pernyataannya benar.

\begin{definisi}[Relasi]\index{relasi}
Relasi dari $A$ ke $B$ adalah sebarang $R \subseteq A \times B$; bila $A = B$ kita
menyebutnya relasi pada $A$. Kita menulis $a \mathrel{R} b$ untuk $(a,b) \in R$.
Daerah asal $R$ adalah $\{\, a \in A : (a,b) \in R\ \text{untuk suatu}\ b \in B \,\}$
dan daerah hasilnya $\{\, b \in B : (a,b) \in R\ \text{untuk suatu}\ a \in A \,\}$.
\end{definisi}

Karena relasi pada $A$ adalah himpunan bagian $A \times A$, banyaknya relasi pada
himpunan $n$ anggota adalah $2^{n^{2}}$. Untuk $n = 3$ angkanya $2^{9} = 512$ dan
untuk $n = 4$ menjadi $2^{16} = 65.536$; keduanya diperiksa dengan skrip di bagian
lab.

Empat sifat menyaringnya menjadi objek berguna. Relasi $R$ pada $A$ disebut refleksif
bila $a \mathrel{R} a$ untuk setiap $a \in A$; simetris bila $a \mathrel{R} b$
berakibat $b \mathrel{R} a$; transitif bila $a \mathrel{R} b$ dan $b \mathrel{R} c$
berakibat $a \mathrel{R} c$; dan antisimetris bila $a \mathrel{R} b$ bersama
$b \mathrel{R} a$ hanya mungkin ketika $a = b$.

\begin{definisi}[Urutan parsial dan relasi ekuivalensi]\index{relasi ekuivalensi}
Relasi pada $A$ yang refleksif, antisimetris, dan transitif disebut urutan parsial.
Relasi pada $A$ yang refleksif, simetris, dan transitif disebut relasi ekuivalensi.
\end{definisi}

Relasi $\le$ pada $\mathbb{R}$ adalah urutan parsial, begitu pula relasi "membagi"
pada bilangan asli positif, dan $\subseteq$ pada $\mathcal{P}(A)$ memberi contoh yang
kelak menjadi kerangka subruang pada Bab 6. Kata "parsial" dipakai karena dua anggota
tidak harus dapat dibandingkan: pada $\mathcal{P}(\{1,2,3\})$ himpunan $\{1\}$ dan
$\{2\}$ tidak termuat satu sama lain. Urutan yang setiap dua anggotanya dapat
dibandingkan disebut urutan total.

\section{Relasi ekuivalensi dan partisi}
\tingkat{menengah}

Relasi ekuivalensi memotong himpunan menjadi kelompok yang tidak tumpang tindih.
Sebelum sampai ke teoremanya, kita perlu nama untuk kelompoknya.

\begin{definisi}[Kelas ekuivalensi]\index{kelas ekuivalensi}
Misalkan $\sim$ relasi ekuivalensi pada $A$ dan $x \in A$. Kelas ekuivalensi $x$
adalah $[x] = \{\, y \in A : y \sim x \,\}$.
\end{definisi}

Karena $x \sim x$, tidak ada kelas yang kosong. Sifat yang lebih menarik: dua kelas
tidak dapat bertumpang sedikit pun.

\begin{lemma}
Misalkan $\sim$ relasi ekuivalensi pada $A$ dan $x, y \in A$. Maka (i) $x \in [x]$;
(ii) $[x] = [y]$ jika dan hanya jika $x \sim y$; dan (iii) bila
$[x] \cap [y] \neq \varnothing$, maka $[x] = [y]$.
\end{lemma}
\par\smallskip\noindent\textit{Bukti.} Bagian (i) adalah pernyataan $x \sim x$, benar
karena refleksif. Untuk (ii), misalkan $[x] = [y]$. Menurut (i),
$y \in [y] = [x]$, jadi $y \sim x$ dan karena simetris $x \sim y$. Sebaliknya,
misalkan $x \sim y$ dan ambil $z \in [x]$, artinya $z \sim x$. Dari $z \sim x$ dan
$x \sim y$ kita peroleh $z \sim y$, jadi $z \in [y]$; ini memberi
$[x] \subseteq [y]$, dan argumen dengan peran bertukar memberi $[y] \subseteq [x]$.
Untuk (iii), misalkan $z \in [x] \cap [y]$. Maka $z \sim x$ dan $z \sim y$; dari
$z \sim x$ dan simetri diperoleh $x \sim z$, dan bersama $z \sim y$ serta
transitivitas diperoleh $x \sim y$, sehingga $[x] = [y]$ menurut (ii).
\hfill$\blacksquare${\parfillskip=0pt\par}

Akibatnya himpunan semua kelas membagi $A$ menjadi potongan yang saling lepas dan
menutupi seluruh $A$.

\begin{definisi}[Partisi]\index{partisi}
Partisi himpunan $A$ adalah keluarga $\mathcal{B}$ himpunan bagian tak kosong dari
$A$ sedemikian sehingga setiap $x \in A$ berada tepat di satu anggota $\mathcal{B}$.
\end{definisi}

\begin{teorema}
Setiap relasi ekuivalensi pada $A$ menentukan partisi $A$, setiap partisi $A$
menentukan relasi ekuivalensi pada $A$, dan kedua proses itu saling membalik.
\end{teorema}
\par\smallskip\noindent\textit{Bukti.} Misalkan $\sim$ relasi ekuivalensi dan
$\mathcal{B} = \{\, [x] : x \in A \,\}$. Menurut Lemma, setiap $[x]$ tidak kosong dan
dua anggotanya saling lepas atau identik; setiap $x$ berada di $[x]$. Jadi
$\mathcal{B}$ partisi.

Sebaliknya, misalkan $\mathcal{B}$ partisi $A$ dan tetapkan $x \sim y$ bila $x$ dan
$y$ berada di anggota $\mathcal{B}$ yang sama. Refleksif dan simetris jelas. Untuk
transitivitas, misalkan $x \sim y$ dan $y \sim z$, dan misalkan $B$ anggota
$\mathcal{B}$ yang memuat $y$; maka $x \in B$ dan $z \in B$, jadi $x \sim z$.

Kedua proses saling membalik: dari $\sim$ kita memperoleh partisi kelas, dan partisi
itu mengembalikan $\sim$ semula karena $x \sim y$ tepat ketika $x \in [y]$; dari
partisi kita memperoleh $\sim$, dan kelas $[x]$ adalah anggota $\mathcal{B}$ yang
memuat $x$.
\hfill$\blacksquare${\parfillskip=0pt\par}

Banyaknya kelas disebut indeks. Himpunan semua kelas ditulis $A/{\sim}$ dan disebut
himpunan kuosien, sedangkan $q : A \to A/{\sim}$ dengan $q(x) = [x]$ disebut pemetaan
kuosien. Hoffman dan Kunze mencatat bahwa setiap relasi ekuivalensi dapat dipandang
berasal dari sebuah fungsi: dua anggota setara tepat ketika fungsi itu memberi nilai
sama.

\begin{contoh}
Pada bilangan bulat, tetapkan $x \equiv y \pmod{n}$ bila $n$ membagi $x-y$. Relasi
ini refleksif karena $n \mid 0$, simetris karena $n \mid (x-y)$ berakibat
$n \mid (y-x)$, dan transitif karena $n \mid (x-y)$ bersama $n \mid (y-z)$ berakibat
$n \mid (x-z)$. Kelasnya himpunan bilangan bersisa sama, dan ada tepat $n$ kelas,
yaitu $[0], \dots, [n-1]$, sebab setiap bilangan bulat dapat ditulis $x = qn + r$
dengan $0 \le r \le n-1$ secara tunggal.
\end{contoh}

Contoh yang lebih dekat dengan pekerjaan kita: untuk $f : A \to B$, tetapkan
$x \sim y$ bila $f(x) = f(y)$. Kelas $[x]$ adalah himpunan masukan dengan keluaran sama;
dalam pembelajaran mesin itulah "satu berkas prediksi", dan banyaknya kelas sama dengan
banyaknya label berbeda.

\begin{untukanak}
Lemari pakaian di rumah biasanya tidak diurutkan menurut aturan resmi. Kaus putih
menumpuk di satu rak, kaus bergambar di rak lain, kaus olahraga di keranjang ketiga.
Tidak ada yang menulis aturannya, tetapi setiap kaus berada di tepat satu tempat.

Itulah yang dilakukan sebuah relasi ekuivalensi: ia menetapkan kapan dua barang
dianggap sekamar, lalu mengelompokkan seluruh barang. Hasilnya pembagian rapi, setiap
barang masuk tepat satu kelompok, dan kelompok berbeda tidak berbagi anggota.

Sekarang pertanyaan yang membuat penasaran. Bagaimana kalau satu kaus menurut satu
aturan masuk kelompok putih, tetapi menurut aturan lain masuk kelompok olahraga? Apakah
pengelompokan masih rapi? Kerapian itu tidak bergantung pada aturan yang dipilih,
melainkan pada tiga sifat: setiap barang sekamar dengan dirinya sendiri, sekamar
berlaku dua arah, dan sekamar dapat dipindahkan lewat barang ketiga.
\end{untukanak}

\section{Fungsi: injektif, surjektif, bijektif}
\tingkat{menengah}

Kita tetapkan sekarang arti kata "fungsi", karena bab ini bergantung pada perbedaan
halus antara kodomain dan daerah hasil.

\begin{definisi}[Fungsi]\index{fungsi injektif}
Fungsi dari $A$ ke $B$ adalah aturan $f$ yang memasangkan setiap $x \in A$ dengan
tepat satu $f(x) \in B$. Himpunan $A$ disebut daerah asal, $B$ disebut kodomain, dan
$f(A) = \{\, f(x) : x \in A \,\}$ disebut daerah hasil.
\end{definisi}

Kodomain dan daerah hasil berbeda, dan perbedaan itu penting. Fungsi
$f : \mathbb{R} \to \mathbb{R}$ dengan $f(x) = x^{2}$ berkodomain $\mathbb{R}$ tetapi
berdaerah hasil $[0,\infty)$. Hoffman dan Kunze memakai tiga nama bergantian untuk
objek yang sama: fungsi, pemetaan, transformasi.

Untuk $S \subseteq B$ kita definisikan prapeta
$f^{-1}(S) = \{\, x \in A : f(x) \in S \,\}$. Notasi $f^{-1}$ di sini tidak
mengandaikan $f$ punya invers: prapeta selalu terdefinisi, invers belum tentu ada.

\begin{definisi}[Injektif, surjektif, bijektif]\index{fungsi surjektif}\index{bijeksi}
Fungsi $f : A \to B$ disebut injektif bila $f(x) = f(y)$ hanya mungkin ketika
$x = y$; surjektif bila $f(A) = B$; dan bijektif bila kedua-duanya.
\end{definisi}

\kamus{bijeksi}{bijection}{Pemasangan sempurna: setiap barang di kotak pertama punya satu pasangan di kotak kedua, tanpa sisa di kedua pihak.}

Gambar~\ref{fig:bijeksi} membandingkan ketiganya pada dua kolom simpul: panel kiri
tidak menyisakan simpul di pihak mana pun, panel tengah menyisakan $d$ di kanan,
panel kanan mengirim dua simpul kiri ke satu simpul kanan.

\begin{figure}[htbp]
\centering
\begin{tikzpicture}[scale=0.95,every node/.style={font=\footnotesize},>=Stealth,
  kotak/.style={draw=ink,thick,circle,inner sep=1.2pt,minimum size=4.4mm}]
  \begin{scope}
    \foreach \i/\y in {1/1.0,2/0,3/-1.0} {
      \node[kotak] (a\i) at (0,\y) {\i};
    }
    \node[kotak] (ba) at (1.9,1.0) {$a$};
    \node[kotak] (bb) at (1.9,0) {$b$};
    \node[kotak] (bc) at (1.9,-1.0) {$c$};
    \draw[->,thick,layerjade] (a1) -- (ba);
    \draw[->,thick,layerjade] (a2) -- (bb);
    \draw[->,thick,layerjade] (a3) -- (bc);
    \node[below] at (0.95,-1.75) {bijektif};
  \end{scope}
  \begin{scope}[xshift=3.4cm]
    \foreach \i/\y in {1/1.2,2/0.4,3/-0.4} {
      \node[kotak] (c\i) at (0,\y) {\i};
    }
    \node[kotak] (da) at (1.9,1.2) {$a$};
    \node[kotak] (db) at (1.9,0.4) {$b$};
    \node[kotak] (dc) at (1.9,-0.4) {$c$};
    \node[kotak] (dd) at (1.9,-1.2) {$d$};
    \draw[->,thick,layerindigo] (c1) -- (da);
    \draw[->,thick,layerindigo] (c2) -- (db);
    \draw[->,thick,layerindigo] (c3) -- (dc);
    \node[below] at (0.95,-1.75) {injektif saja};
  \end{scope}
  \begin{scope}[xshift=6.9cm]
    \foreach \i/\y in {1/1.2,2/0,3/-1.2} {
      \node[kotak] (e\i) at (0,\y) {\i};
    }
    \node[kotak] (fa) at (1.9,0.8) {$a$};
    \node[kotak] (fb) at (1.9,-0.8) {$b$};
    \draw[->,thick,layerbrown] (e1) -- (fa);
    \draw[->,thick,layerbrown] (e2) -- (fa);
    \draw[->,thick,layerbrown] (e3) -- (fb);
    \node[below] at (0.95,-1.75) {surjektif saja};
  \end{scope}
\end{tikzpicture}
\caption{Tiga pemetaan antara dua kolom simpul. Kiri: setiap simpul terpakai tepat
sekali, jadi bijektif. Tengah: kolom kanan beranggota empat dan simpul $d$ tidak
mendapat panah masuk, jadi injektif tetapi tidak surjektif. Kanan: dua simpul kiri
menuju $a$, jadi surjektif tetapi tidak injektif.}
\label{fig:bijeksi}
\end{figure}

Kata "bijektif" hanya nama pendek untuk gabungan dua sifat sebelumnya. Kekuatannya
muncul dalam proposisi berikut: bijeksi adalah syarat tepat agar fungsi dapat dibalik.

\begin{proposisi}
$f : A \to B$ bijektif jika dan hanya jika terdapat $g : B \to A$ dengan
$g \circ f = \mathrm{id}_A$ dan $f \circ g = \mathrm{id}_B$. Bila $g$ ada, ia tunggal.
\end{proposisi}
\par\smallskip\noindent\textit{Bukti.} Misalkan $g$ memenuhi kedua kesamaan. Bila
$f(x) = f(y)$ maka $x = g(f(x)) = g(f(y)) = y$, jadi $f$ injektif. Untuk sebarang
$b \in B$ berlaku $b = f(g(b))$, jadi $f$ surjektif.

Sebaliknya, misalkan $f$ bijektif. Untuk setiap $b \in B$ ada tepat satu $x \in A$
dengan $f(x) = b$: keberadaannya dari surjektif, ketunggalannya dari injektif.
Definisikan $g(b) = x$ untuk $x$ itu, maka $g(f(x)) = x$ dan $f(g(b)) = b$.

Untuk ketunggalan, misalkan $g$ dan $h$ memenuhi kedua kesamaan. Maka
\[
  g = g \circ \mathrm{id}_B = g \circ (f \circ h) = (g \circ f) \circ h
  = \mathrm{id}_A \circ h = h,
\]
yang hanya memakai assosiativitas komposisi dan kedua kesamaan itu.
\hfill$\blacksquare${\parfillskip=0pt\par}

Komposisi berkelakuan baik terhadap ketiga sifat: bila $f$ dan $g$ injektif maka
$g \circ f$ injektif, dan serupa untuk surjektif serta bijektif. Kebalikannya tidak
berlaku: dengan $f : \{1\} \to \{1,2\}$, $f(1) = 1$, dan $g : \{1,2\} \to \{1\}$,
$g(1) = g(2) = 1$, komposisi $g \circ f$ bijektif, padahal $g$ tidak injektif dan $f$
tidak surjektif.

Pada himpunan berhingga ketiga sifat saling mengunci, dan itulah sebabnya "satu-satu"
dan "pada" baru dapat dibedakan pada himpunan tak berhingga.

\begin{lemma}
Misalkan $A$ dan $B$ berhingga, $f : A \to B$, dan $C \subseteq B$. Maka
$|f(A)| \le |A|$ dengan kesamaan jika dan hanya jika $f$ injektif, dan
$|C| \le |B|$ dengan kesamaan jika dan hanya jika $C = B$.
\end{lemma}
\par\smallskip\noindent\textit{Bukti.} Kita buktikan bagian kedua dengan induksi pada
$n = |B|$. Untuk $n = 0$ berlaku $B = \varnothing$, sehingga $C = B$. Andaikan benar
untuk himpunan $n$ anggota dan misalkan $|B| = n+1$; pilih $b \in B$ dan tulis
$B' = B \setminus \{b\}$.

Bila $b \notin C$, maka $C \subseteq B'$ dan hipotesis induksi memberi
$|C| \le |B'| = n < n+1$, dengan kesamaan mustahil. Bila $b \in C$, tulis
$C' = C \setminus \{b\} \subseteq B'$. Maka $|C| = |C'| + 1 \le |B'| + 1 = |B|$;
kesamaan berlaku tepat ketika $|C'| = |B'|$, yang menurut hipotesis induksi berarti
$C' = B'$, yaitu $C = B$.

Bagian pertama diperoleh dengan induksi yang sama pada $|A|$: tulis
$A = A' \cup \{a\}$ dengan $a \notin A'$ dan perhatikan
$f(A) = f(A') \cup \{f(a)\}$, sehingga $|f(A)| \le |A'| + 1 = |A|$. Kesamaan menuntut
$f(a) \notin f(A')$ dan $|f(A')| = |A'|$; yang terakhir berarti $f|_{A'}$ injektif
menurut hipotesis induksi, dan yang pertama berarti $f$ tidak mengulang nilai pada
$a$.
\hfill$\blacksquare${\parfillskip=0pt\par}

\begin{proposisi}
Misalkan $A$ dan $B$ berhingga dengan $|A| = |B|$, dan $f : A \to B$. Maka $f$
injektif, $f$ surjektif, dan $f$ bijektif adalah tiga pernyataan yang setara.
\end{proposisi}
\par\smallskip\noindent\textit{Bukti.} Bila $f$ injektif, Lemma memberi
$|f(A)| = |A| = |B|$; karena $f(A) \subseteq B$ dan keduanya berukuran sama, Lemma
memberi $f(A) = B$, jadi $f$ surjektif. Bila $f$ surjektif, maka $f(A) = B$ sehingga
$|f(A)| = |B| = |A|$, dan Lemma memaksa $f$ injektif. Karena bijektif berarti kedua
sifat sekaligus, kesetaraan ketiganya selesai.
\hfill$\blacksquare${\parfillskip=0pt\par}

Peringatan yang harus dibawa ke bab berikutnya: proposisi itu memakai kata
"berhingga" secara esensial. Pada $\mathbb{N}$, fungsi $f(x) = 2x$ injektif tetapi
tidak surjektif, dan fungsi yang memetakan $1$ ke $1$ serta $n \ge 2$ ke $n-1$
surjektif tetapi tidak injektif.

\section{Himpunan berhingga, terhitung, dan tak terhitung}
\tingkat{lanjut}

Sekarang kita kembali ke pertanyaan pembuka. Untuk membandingkan dua himpunan tanpa
menghitung kita memasangkan anggotanya, dan yang perlu ditetapkan hanyalah bahwa
pemasangan itu memakai fungsi yang dapat dibalik.

\begin{definisi}[Berkekuatan sama]
Himpunan $A$ dan $B$ berkekuatan sama, ditulis $A \sim B$, bila terdapat bijeksi
$A \to B$.
\end{definisi}

Relasi $\sim$ refleksif (identitas bijeksi), simetris (invers bijeksi adalah
bijeksi), dan transitif (komposisi dua bijeksi bijeksi). Jadi ia berperilaku sebagai
relasi ekuivalensi, tetapi tidak pada "himpunan semua himpunan", sebab kumpulan itu
bukan himpunan; ini dibahas pada kotak untuk peneliti.

\begin{definisi}[Berhingga, terhitung, tak terhitung]\index{himpunan terhitung}
Himpunan $A$ berhingga bila $A = \varnothing$ atau terdapat $n \in \mathbb{N}$ dan
bijeksi $A \to \{1,2,\dots,n\}$; dalam hal kedua kita tulis $|A| = n$. Himpunan $A$
terhitung bila $A \sim \mathbb{N}$. Himpunan tak berhingga yang tidak terhitung
disebut tak terhitung.
\end{definisi}

\kamus{terhitung}{countable}{Semua anggotanya bisa diberi nomor urut tanpa ada yang terlewat, seperti setiap penumpang mendapat nomor tiket yang berbeda.}

Nama "terhitung" tepat: bila $A \sim \mathbb{N}$, bijeksinya adalah daftar. Kita
dapat menyebut anggota pertama, kedua, dan seterusnya, dan setiap anggota muncul tepat
sekali. Yang mengejutkan adalah betapa banyak himpunan yang dapat didaftar begitu.
Pasangan bilangan, misalnya.

\begin{proposisi}[Pasangan Cantor]\index{pasangan Cantor}
$\mathbb{N} \times \mathbb{N}$ terhitung.
\end{proposisi}
\par\smallskip\noindent\textit{Bukti.} Definisikan
$\pi(m,n) = \frac{(m+n)(m+n+1)}{2} + n$ untuk $m, n \in \mathbb{N} =
\{0,1,2,\dots\}$. Tulis $T_s = \frac{s(s+1)}{2}$, yaitu banyaknya pasangan dengan
$m + n < s$. Untuk pasangan dengan $m + n = s$ berlaku $0 \le n \le s$, sehingga
$\pi(m,n) = T_s + n$ mengambil nilai $T_s, \dots, T_s + s$ tepat satu kali
masing-masing. Karena $T_s + s = T_{s+1} - 1$, blok nilai untuk diagonal $s$ adalah
$\{T_s, \dots, T_{s+1}-1\}$. Blok-blok itu bersambung tanpa celah dan tanpa tumpang
tindih mulai dari $T_0 = 0$. Jadi setiap bilangan cacah berada di tepat satu blok,
dan di dalam bloknya berasal dari tepat satu pasangan; karena itu $\pi$ bijektif.
\hfill$\blacksquare${\parfillskip=0pt\par}

Inversnya eksplisit: $s = \lfloor (\sqrt{8q+1}-1)/2 \rfloor$, $n = q - T_s$, dan
$m = s - n$. Akar kuadrat itulah satu-satunya tempat aritmetika mesin dapat menipu
kita, dan contohnya ada di bagian lab.

Gambar~\ref{fig:kisi} menggambar kisi itu: sel pada baris $m$ dan kolom $n$ berlabel
$\pi(m,n)$, dan warna menandai diagonal $m + n = s$, sehingga membaca menurut warna
memberi $0$ sampai $14$ tanpa celah dan tanpa pengulangan. Panel kanan memakai arah
baca yang sama pada daftar empat bilangan real, dengan baris $d$ yang membalik digit
diagonal sehingga $d$ tidak ada di daftar.

\begin{figure}[htbp]
\centering
\begin{tikzpicture}[scale=0.95,every node/.style={font=\footnotesize}]
  \foreach \s/\warna in {0/layerteal,1/layerorange,2/layerindigo,3/layerjade,4/layerplum} {
    \foreach \n in {0,...,\s} {
      \pgfmathtruncatemacro{\m}{\s-\n}
      \pgfmathtruncatemacro{\v}{int(\s*(\s+1)/2+\n)}
      \draw[fill=\warna!30,draw=ink!45] (\n*0.86,-\m*0.72) rectangle ++(0.8,0.66);
      \node at (\n*0.86+0.4,-\m*0.72+0.33) {\v};
    }
  }
  \draw[->,thick] (4.35,0.33) -- (4.95,0.33);
  \node[right] at (4.95,0.33) {$n$};
  \draw[->,thick] (-0.45,-0.05) -- (-0.45,-3.0) node[below] {$m$};
  \begin{scope}[xshift=6.6cm,yshift=0.45cm]
    \foreach \i/\j/\b in {1/1/0,1/2/1,1/3/1,1/4/0, 2/1/1,2/2/0,2/3/0,2/4/1,
                          3/1/0,3/2/0,3/3/1,3/4/0, 4/1/1,4/2/1,4/3/0,4/4/0} {
      \ifnum\i=\j
        \fill[accent!28] (\j*0.52-0.52,-\i*0.42) rectangle ++(0.5,0.38);
      \fi
      \draw[ink!45] (\j*0.52-0.52,-\i*0.42) rectangle ++(0.5,0.38);
    }
    \draw[accent,thick] (0.25,-0.23) -- (1.79,-1.49);
    \foreach \i/\j/\b in {1/1/0,1/2/1,1/3/1,1/4/0, 2/1/1,2/2/0,2/3/0,2/4/1,
                          3/1/0,3/2/0,3/3/1,3/4/0, 4/1/1,4/2/1,4/3/0,4/4/0} {
      \node at (\j*0.52-0.27,-\i*0.42+0.19) {\b};
    }
    \foreach \i in {1,...,4} {
      \node[left] at (-0.62,-\i*0.42+0.19) {$x_{\i}$};
    }
    \foreach \j/\b in {1/1,2/1,3/0,4/1} {
      \fill[layergrey!25] (\j*0.52-0.52,-2.34) rectangle ++(0.5,0.38);
      \draw[ink!45] (\j*0.52-0.52,-2.34) rectangle ++(0.5,0.38);
      \node at (\j*0.52-0.27,-2.34+0.19) {\b};
    }
    \node[left] at (-0.10,-2.15) {$d$};
    \node[below] at (0.52,-2.42) {digit dibalik};
  \end{scope}
\end{tikzpicture}
\caption{Kiri: kisi pasangan $(m,n)$ dengan $m + n \le 4$, setiap sel berlabel
$\pi(m,n)$ dan warna menandai diagonal $m + n = s$; lima diagonal pertama memuat
nilai $0$ sampai $14$. Kanan: tabel digit empat bilangan real, dengan sel diagonal
diwarnai dan baris $d$ membalik digit diagonal $0,0,1,0$ menjadi $1,1,0,1$.}
\label{fig:kisi}
\end{figure}

Karena setiap rasional positif adalah pasangan bilangan asli yang sudah
disederhanakan, himpunan rasional juga terhitung. Yang perlu dijaga hanyalah agar
tidak ada bilangan yang didaftar dua kali.

\begin{proposisi}
Himpunan $\mathbb{Q}_{>0}$ bilangan rasional positif terhitung.
\end{proposisi}
\par\smallskip\noindent\textit{Bukti.} Setiap $q \in \mathbb{Q}_{>0}$ dapat ditulis
tunggal sebagai $q = p/r$ dengan $p, r$ bilangan asli dan $\gcd(p,r) = 1$. Definisikan
tinggi $q$ sebagai $p + r$. Untuk setiap $k \ge 2$ hanya ada berhingga rasional
bertinggi $k$, sebab $p$ dan $r$ berada di antara $1$ dan $k-1$. Maka daftar berikut
mencacah $\mathbb{Q}_{>0}$ tanpa pengulangan: untuk $k = 2, 3, 4, \dots$ tuliskan
semua $p/r$ bertinggi $k$ dengan $\gcd(p,r) = 1$, urut menaik menurut $p$. Setiap
rasional positif muncul tepat sekali, yaitu pada tingkat $k = p + r$ dengan bentuk
sederhana yang tunggal, sehingga daftar itu bijeksi $\mathbb{N} \to \mathbb{Q}_{>0}$.
\hfill$\blacksquare${\parfillskip=0pt\par}

\begin{akibat}
$\mathbb{Q}$ dan $\mathbb{Z}$ terhitung.
\end{akibat}
\par\smallskip\noindent\textit{Bukti.} Daftar $0, 1, -1, 2, -2, \dots$ adalah
bijeksi $\mathbb{N} \to \mathbb{Z}$. Untuk $\mathbb{Q}$, cacah setiap rasional tak
nol sebagai $\pm p/r$ dalam bentuk sederhana menurut $(p+r)$ dengan tanda bergantian,
lalu letakkan $0$ di awal; satu anggota tambahan tidak merusak sifat terhitung.
\hfill$\blacksquare${\parfillskip=0pt\par}

Jadi $\mathbb{N}$, $\mathbb{Z}$, dan $\mathbb{Q}$ berkekuatan sama, meskipun
$\mathbb{Q}$ terasa jauh lebih rapat. Rapat bukan ukuran besar; ukurannya adalah ada
atau tidaknya daftar.

\begin{tantangan}
Hotel di tepi pantai punya kamar bernomor $1, 2, 3, \dots$ tanpa ujung, dan malam
itu setiap kamar terisi tepat satu tamu. Seorang pelancong datang tanpa pesanan.
Resepsionis berkata hotel penuh; manajer berkata masih ada kamar untuknya tanpa
memulangkan seorang tamu pun.

Tunjukkan aturan pemindahan tamu yang memenuhi janji manajer, lalu jelaskan mengapa
aturan itu bijeksi antara tamu lama ditambah satu tamu baru dan seluruh kamar.
Mengapa hotel yang sama tidak dapat menampung tamu baru sebanyak bilangan real?

Petunjuk 1: jangan mencari kamar kosong, sebab tidak ada. Ciptakan kamar kosong dengan
memindahkan penghuninya.

Petunjuk 2: geser penghuni kamar $n$ ke kamar $n+1$ untuk setiap $n$, lalu taruh
pelancong di kamar $1$; periksa bahwa setiap kamar terisi tepat satu orang dan tidak
ada tamu yang hilang.

Kalau mentok, kembali ke bagian Himpunan berhingga, terhitung, tak terhitung,
khususnya Definisi berkekuatan sama.
\end{tantangan}

Cantor menunjukkan bahwa $\mathcal{P}(A)$ selalu lebih besar daripada $A$, dengan trik
yang tidak bergantung pada ukuran $A$.

\begin{teorema}[Cantor]\index{diagonal Cantor}
Untuk setiap himpunan $A$ tidak ada fungsi surjektif $A \to \mathcal{P}(A)$.
\end{teorema}
\par\smallskip\noindent\textit{Bukti.} Misalkan $f : A \to \mathcal{P}(A)$ sebarang
dan definisikan $D = \{\, a \in A : a \notin f(a) \,\}$. Maka $D \subseteq A$, jadi
$D \in \mathcal{P}(A)$. Andaikan $D = f(a_0)$ untuk suatu $a_0 \in A$. Bila
$a_0 \in D$, maka menurut definisi $D$ berlaku $a_0 \notin f(a_0) = D$, kontradiksi.
Bila $a_0 \notin D$, maka $a_0 \notin f(a_0)$, yang justru syarat keanggotaan $D$,
sehingga $a_0 \in D$, kontradiksi lagi. Jadi $D$ tidak berada di daerah hasil $f$.
\hfill$\blacksquare${\parfillskip=0pt\par}

Himpunan $D$ dibentuk dengan menanyakan, untuk setiap $a$, apakah $a$ menjadi anggota
himpunan yang dipasangkan kepadanya, lalu membalik jawabannya; pembalikan itulah yang
membuat $D$ berbeda dari setiap nilai $f$. Pada $A = \mathbb{N}$ trik ini muncul sebagai
pembacaan diagonal tabel dengan setiap entrinya diubah, karena itu disebut argumen
diagonal.

\begin{akibat}
$\mathcal{P}(\mathbb{N})$ tak terhitung, dan $\mathbb{R}$ tak terhitung.
\end{akibat}
\par\smallskip\noindent\textit{Bukti.} Menurut Teorema Cantor dengan $A = \mathbb{N}$,
tidak ada surjeksi $\mathbb{N} \to \mathcal{P}(\mathbb{N})$, jadi
$\mathcal{P}(\mathbb{N})$ tak terhitung.

Untuk $\mathbb{R}$ cukup ditunjukkan bahwa $[0,1]$ tak terhitung. Setiap $x \in [0,1]$
memiliki ekspansi biner $x = \sum_{k \ge 1} b_k 2^{-k}$ dengan $b_k \in \{0,1\}$;
keberadaannya dibuktikan pada Bab 3 dan Bab 5. Pemetaan dari himpunan barisan biner ke
$[0,1]$ karena itu surjektif, dan himpunan barisan biner berkekuatan sama dengan
$\mathcal{P}(\mathbb{N})$ sebab barisan biner adalah fungsi $\mathbb{N} \to \{0,1\}$.
Bila $[0,1]$ terhitung, bayangan prapetanya di himpunan barisan biner akan terhitung
pula, bertentangan dengan paragraf sebelumnya.
\hfill$\blacksquare${\parfillskip=0pt\par}

Satu catatan kejujuran: ekspansi biner tidak tunggal, sebab
$0{,}0111\dots = 0{,}1000\dots$ dalam basis dua. Argumen di atas tidak terganggu karena
hanya memakai sifat surjektif, tetapi $|\mathbb{R}| = |\mathcal{P}(\mathbb{N})|$
menuntut alat berikut.

\begin{proposisi}[Cantor--Bernstein--Schroeder]
Bila terdapat fungsi injektif $A \to B$ dan injektif $B \to A$, maka $A \sim B$.
\end{proposisi}
\par\smallskip\noindent\textit{Sketsa bukti.} Misalkan $f : A \to B$ dan
$g : B \to A$ injektif. Tetapkan $A_0$ sebagai himpunan titik $A$ yang tidak punya
pendahulu di $B$, lalu $A_{k+1} = g(f(A_k))$, dan tulis $A^{+}$ untuk gabungan semua
$A_k$ dengan $k$ genap. Definisikan $h(x) = f(x)$ bila $x \in A^{+}$ dan
$h(x) = g^{-1}(x)$ bila $x \notin A^{+}$. Pada bagian pertama $h$ injektif karena $f$
injektif, pada bagian kedua karena $g^{-1}$ injektif, dan bayangan kedua bagian
terpisah menurut cara $A^{+}$ dibentuk, sehingga $h$ injektif. Untuk kelinearan
bayangan, ambil $y \in B$. Bila $g(y) \in A^{+}$, maka $g(y)$ berada di suatu $A_k$
dengan $k$ genap sehingga $g(y) = g(f(x))$ untuk $x \in A_{k-1}$ dan $h(x) = y$; bila
$g(y) \notin A^{+}$, maka $h(g(y)) = y$. Jadi $h$ surjektif.

Pembukuan kedalaman rantai itulah bagian yang halus, dan rinciannya dilewati di sini;
versi formal lengkapnya ada pada \emph{Mathematics in Lean} bagian 4.3.
\hfill$\blacksquare${\parfillskip=0pt\par}

Dengan alat itu kita boleh menyimpulkan $|\mathbb{R}| = |\mathcal{P}(\mathbb{N})|$:
injeksi dari $\mathcal{P}(\mathbb{N})$ ke $\mathbb{R}$ diberikan ekspansi biner, dan
injeksi sebaliknya oleh fungsi tangen. Kekuatan $\mathbb{R}$ biasa ditulis
$2^{\aleph_0}$, dengan $\aleph_0$ kekuatan $\mathbb{N}$.

\begin{untukanak}
Bayangkan mesin tiket di pintu masuk pameran. Setiap pengunjung mendapat satu tiket
bernomor: 1, 2, 3, dan seterusnya, tidak ada nomor terpakai dua kali, dan setiap nomor
yang mungkin akan suatu saat diberikan. Kumpulan benda yang bisa diberi tiket seperti
itu disebut terhitung.

Kelihatannya semua kumpulan tak berujung bisa diberi tiket. Tetapi tidak. Ada
kumpulan yang anggotanya begitu banyak sehingga sebanyak apa pun tiket yang kita
cetak, selalu ada anggota yang tidak kebagian. Mesin tiketnya tidak rusak; kumpulan
itu memang lebih besar daripada persediaan nomor.

Pertanyaan yang membuat penasaran: bagaimana kita bisa tahu sebuah kumpulan tidak
bisa diberi tiket, padahal belum semua cara dicoba? Jawabannya trik pembalikan Cantor.
Setiap kali seseorang mengaku sudah memberi tiket kepada semua anggota, kita susun satu
anggota baru yang sengaja berbeda dari setiap anggota bertiket; anggota itu pasti
terlewat.
\end{untukanak}

\section{Ketaksamaan, nilai mutlak, dan batas atas}
\tingkat{menengah}

Bagian terakhir memasang dua alat untuk Bab 3. Yang pertama mengukur jarak antara dua
bilangan tanpa memandang arah, yaitu nilai mutlak.

\begin{definisi}[Nilai mutlak]\index{nilai mutlak}
Untuk $x \in \mathbb{R}$, $|x| = x$ bila $x \ge 0$ dan $|x| = -x$ bila $x < 0$.
\end{definisi}

Jarak melalui titik ketiga tidak pernah lebih pendek daripada jarak langsung.

\begin{proposisi}[Ketaksamaan segitiga]
Untuk setiap $x, y \in \mathbb{R}$ berlaku $|x + y| \le |x| + |y|$ dan
$\bigl| |x| - |y| \bigr| \le |x - y|$.
\end{proposisi}
\par\smallskip\noindent\textit{Bukti.} Perhatikan $-|x| \le x \le |x|$ dan
$-|y| \le y \le |y|$. Menjumlahkan keduanya memberi
$-(|x|+|y|) \le x + y \le |x|+|y|$, yang menurut definisi nilai mutlak berarti
$|x+y| \le |x|+|y|$. Untuk yang kedua, tulis $x = (x-y) + y$ dan terapkan hasil
pertama: $|x| \le |x-y| + |y|$, jadi $|x| - |y| \le |x-y|$. Menukar peran $x$ dan $y$
memberi $|y| - |x| \le |x-y|$, sehingga $\bigl| |x| - |y| \bigr| \le |x-y|$.
\hfill$\blacksquare${\parfillskip=0pt\par}

Sekarang batas. Himpunan bilangan dapat terbatas di atas tanpa memiliki anggota
terbesar, dan justru kasus itulah yang paling penting.

\begin{definisi}[Batas atas dan supremum]\index{supremum}
Misalkan $S \subseteq \mathbb{R}$ tak kosong. Bilangan $u$ disebut batas atas $S$
bila $s \le u$ untuk setiap $s \in S$; bila ada, $S$ disebut terbatas di atas.
Bilangan $u$ disebut supremum $S$, ditulis $u = \sup S$, bila $u$ batas atas $S$ dan
$u \le v$ untuk setiap batas atas $v$ dari $S$.
\end{definisi}

\kamus{supremum}{supremum (least upper bound)}{Pagar paling rapat: batas tertinggi yang tidak dapat diturunkan lagi tanpa menyisakan anggota di atasnya.}

Supremum, bila ada, tunggal, sebab dua batas atas terkecil saling mendominasi. Yang
tidak dijamin definisi adalah keberadaannya: $\{x \in \mathbb{Q} : x^{2} < 2\}$
terbatas di atas dalam $\mathbb{Q}$ tetapi tidak punya supremum di sana.

\begin{catatan}[Aksioma kelengkapan]
Setiap himpunan tak kosong $S \subseteq \mathbb{R}$ yang terbatas di atas memiliki
supremum di $\mathbb{R}$. Aksioma ini tidak dibuktikan di sini; pembangunannya
beserta akibat-akibatnya adalah isi Bab 3, termasuk ciri epsilon yang mengubah
pernyataan tentang batas terkecil menjadi pernyataan tentang $\varepsilon$, dan
contoh himpunan yang memiliki supremum tanpa maksimum.
\end{catatan}

Pada himpunan berhingga supremum dan maksimum berimpit, dan menukar keduanya sering
lolos pada contoh kecil.

\section{Contoh hitung}
\tingkat{menengah}

Empat hitungan berikut dikerjakan tangan sampai angka, lalu diperiksa ulang dengan
skrip. Semuanya kecil, sehingga kesalahan aritmetika langsung terlihat.

\begin{resep}
Untuk dua himpunan hingga $A$ dan $B$:
\begin{enumerate}
  \item Hitung $|A|$, banyaknya anggota $A$.
  \item Hitung $|B|$ dengan cara yang sama.
  \item Hitung $|A \cap B|$, anggota yang muncul di kedua daftar.
  \item Jumlahkan lalu kurangi: $|A \cup B| = |A| + |B| - |A \cap B|$.
\end{enumerate}
Peringatan urutan: langkah 3 tidak boleh ditunda sampai sesudah langkah 4. Kalau
penjumlahan dikerjakan lebih dulu, anggota irisan terhitung dua kali dan hasilnya
kelebihan tepat $|A \cap B|$. Untuk $A = \{1,\dots,5\}$ dan $B = \{4,\dots,8\}$:
$5 + 5 - 2 = 8$, sama dengan angka pada diagram Venn sebelumnya.
\end{resep}

Pertama, himpunan kuasa. Untuk $A = \{1,2,3\}$ himpunan bagiannya berukuran
$0,1,2,3$ dengan cacah $1, 3, 3, 1$, seluruhnya $1 + 3 + 3 + 1 = 8 = 2^{3}$. Baris
$n = 4$ memberi $1+4+6+4+1 = 16$, $n = 2$ memberi $1+2+1 = 4$, dan $n = 1$ memberi
$1+1 = 2$: segitiga Pascal yang sama dengan Proposisi tentang $2^{n}$.

Kedua, pasangan Cantor. Dengan $\pi(m,n) = T_{m+n} + n$ dan $T_0 = 0$, $T_1 = 1$,
$T_2 = 3$, $T_3 = 6$: diagonal $s = 0$ memuat $(0,0) \mapsto 0$; $s = 1$ memuat
$(1,0) \mapsto 1$ dan $(0,1) \mapsto 2$; $s = 2$ memuat $(2,0) \mapsto 3$,
$(1,1) \mapsto 4$, $(0,2) \mapsto 5$; $s = 3$ memuat $(3,0) \mapsto 6$,
$(2,1) \mapsto 7$, $(1,2) \mapsto 8$, $(0,3) \mapsto 9$. Sepuluh nilai itu tepat
$0, \dots, 9$ tanpa pengulangan dan tanpa celah.

Arah sebaliknya lebih mudah salah. Ambil $q = 1000$. Kita cari $s$ dengan
$T_s \le 1000 < T_{s+1}$. Karena $T_{44} = 44 \cdot 45/2 = 990$ dan $T_{45} = 1035$,
maka $s = 44$, lalu $n = 1000 - 990 = 10$ dan $m = 44 - 10 = 34$. Pemeriksaan:
$\pi(34,10) = 44 \cdot 45/2 + 10 = 990 + 10 = 1000$, benar.

Ketiga, bilangan rasional. Menurut aturan tinggi, tingkat $k = 2$ memuat $1/1$;
$k = 3$ memuat $1/2, 2/1$; $k = 5$ memuat $1/4, 2/3, 3/2, 4/1$. Banyaknya suku tiap
tingkat adalah nilai fungsi totien Euler, $1, 2, 2, 4, 2, 6$ untuk $k = 2, \dots, 7$.
Kumulatifnya sampai $k = 7$ adalah $1+2+2+4+2+6 = 17$, jadi $6/1$ menempati posisi
ke-17. Karena kumulatif sampai $k = 9$ adalah $27$ dan tingkat $k = 10$ memuat
$1/9, 3/7, 7/3, 9/1$, bilangan $7/3$ menempati posisi ke-30.

Keempat, banyaknya relasi. Pada himpunan $3$ anggota ada $2^{9} = 512$ relasi, dan
hanya $5$ yang ekuivalensi: $1$ dengan satu kelas, $3$ dengan dua kelas, $1$ dengan tiga
kelas. Pada himpunan $4$ anggota ada $2^{16} = 65.536$ relasi dan $15$ yang ekuivalen.
Keempat angka itu muncul lagi sebagai empat baris terakhir keluaran skrip; barisan
$1, 2, 5, 15, 52$ dikenal sebagai bilangan Bell.

\kertas{Empat angka yang sudah diperiksa ulang. (i) $|\mathcal{P}(A)| = 2^{n}$ untuk
$n = 0,1,2,3,4$ memberi $1, 2, 4, 8, 16$.
(ii) $\pi(34,10) = 990 + 10 = 1000$, dan $\pi$ pada diagonal $s = 0,1,2,3$ memberi
$0,1,2,3,4,5,6,7,8,9$. (iii) Banyak rasional positif bertinggi $k$ adalah
nilai fungsi totien di $k$; kumulatif sampai $k = 9$ adalah $27$, jadi $7/3$ di
posisi ke-30. (iv) Relasi pada himpunan $3$ anggota: $512$ seluruhnya, $5$ ekuivalen.}

\section{Lab komputasi}
\tingkat{lanjut}

Skrip berikut menjalankan keempat hitungan itu dan menambahkan satu pemeriksaan yang
tidak dapat dikerjakan tangan: ia membandingkan invers pasangan Cantor yang dihitung
dengan pecahan mesin dan yang dihitung dengan akar bilangan bulat. Keduanya
seharusnya memberi jawaban sama; pada satu titik jawabannya berbeda.

\begin{lab}
\begin{verbatim}
from itertools import combinations
from math import gcd, isqrt, sqrt, floor

# (1) banyak himpunan bagian untuk n kecil
for n in range(5):
    A = range(n)
    subs = [S for k in range(n + 1) for S in combinations(A, k)]
    print("n =", n, "|P(A)| =", len(subs), "2^n =", 2 ** n)

# (2) pasangan Cantor pi(m, n) = (m+n)(m+n+1)/2 + n
def pi(m, n):
    s = m + n
    return s * (s + 1) // 2 + n

pasangan = [(0, 0), (1, 0), (0, 1), (2, 0), (1, 1),
            (0, 2), (3, 0), (2, 1), (1, 2), (0, 3)]
print("pi:", [pi(m, n) for m, n in pasangan])
keys = {pi(m, n) for m in range(61) for n in range(61)}
print("kunci unik pada 61x61 pasangan:", len(keys) == 61 * 61)

# (3) invers: float vs isqrt
s, n = 2 * 10**8, 2 * 10**8 - 1
q = s * (s + 1) // 2 + n
print("q =", q, "s benar =", s)
print("s dari float:", int(floor((sqrt(8 * q + 1) - 1) / 2)))
print("s dari isqrt:", (isqrt(8 * q + 1) - 1) // 2)

# (4) enumerasi Q positif menurut tinggi k = p + q
seq = []
for k in range(2, 10):
    row = [(p, k - p) for p in range(1, k) if gcd(p, k - p) == 1]
    print("k =", k, "banyak =", len(row), row)
    seq += row
print("suku ke-1..8:", seq[:8])
print("total sampai k = 9:", len(seq))

# (5) relasi pada himpunan n anggota: 2^(n*n) seluruhnya, Bell yang ekuivalen
def ekuivalen(R, n):
    refl = all((a, a) in R for a in range(n))
    sim = all((b, a) in R for a, b in R)
    trans = all((a, c) in R for a, b in R for c in range(n)
                if (b, c) in R)
    return refl and sim and trans

for n in range(1, 5):
    semua = 2 ** (n * n)
    hitung = 0
    for bit in range(semua):
        R = {(a, b) for a in range(n) for b in range(n)
             if bit >> (a * n + b) & 1}
        if ekuivalen(R, n):
            hitung += 1
    print("n =", n, "relasi =", semua, "ekuivalen =", hitung)
\end{verbatim}
\end{lab}

Keluaran yang tercetak, apa adanya:

\begin{verbatim}
n = 0 |P(A)| = 1 2^n = 1
n = 1 |P(A)| = 2 2^n = 2
n = 2 |P(A)| = 4 2^n = 4
n = 3 |P(A)| = 8 2^n = 8
n = 4 |P(A)| = 16 2^n = 16
pi: [0, 1, 2, 3, 4, 5, 6, 7, 8, 9]
kunci unik pada 61x61 pasangan: True
q = 20000000299999999 s benar = 200000000
s dari float: 200000001
s dari isqrt: 200000000
k = 2 banyak = 1 [(1, 1)]
k = 3 banyak = 2 [(1, 2), (2, 1)]
k = 4 banyak = 2 [(1, 3), (3, 1)]
k = 5 banyak = 4 [(1, 4), (2, 3), (3, 2), (4, 1)]
k = 6 banyak = 2 [(1, 5), (5, 1)]
k = 7 banyak = 6 [(1, 6), (2, 5), (3, 4), (4, 3), (5, 2), (6, 1)]
k = 8 banyak = 4 [(1, 7), (3, 5), (5, 3), (7, 1)]
k = 9 banyak = 6 [(1, 8), (2, 7), (4, 5), (5, 4), (7, 2), (8, 1)]
suku ke-1..8: [(1, 1), (1, 2), (2, 1), (1, 3), (3, 1), (1, 4), (2, 3), (3, 2)]
total sampai k = 9: 27
n = 1 relasi = 2 ekuivalen = 1
n = 2 relasi = 16 ekuivalen = 2
n = 3 relasi = 512 ekuivalen = 5
n = 4 relasi = 65536 ekuivalen = 15
\end{verbatim}

Yang perlu dibaca dari keluaran itu: lima baris pertama mengonfirmasi $2^{n}$ untuk
$n = 0$ sampai $4$; sepuluh nilai $\pi$ berurutan tanpa celah, dan $61 \times 61 =
3721$ pasangan memberi kunci unik sebanyak $3721$, yang merupakan bukti numerik, bukan
bukti matematis, sebab bukti matematisnya argumen blok $T_s$. Empat baris terakhir
menyapu seluruh $2^{n^{2}}$ relasi untuk $n = 1$ sampai $4$ dan mencacah yang
refleksif, simetris, dan transitif: hasilnya $1, 2, 5, 15$. Baris $n = 5$ tidak ada
karena menyapu $2^{25}$ relasi terlalu lama; nilainya $52$.

Ketiga, invers dengan pecahan salah. Pada $q = 20.000.000.299.999.999$ rumus
$\lfloor (\sqrt{8q+1}-1)/2 \rfloor$ dengan pecahan mesin 64-bit memberi
$s = 200.000.001$, sedangkan nilai benarnya $200.000.000$. Kesalahannya satu, dan itu
cukup untuk mengembalikan pasangan yang sama sekali berbeda: dengan $s = 200.000.001$
nilai $n = q - T_s$ menjadi negatif. Angka $8q+1$ berorde $10^{17}$, sedangkan pecahan
64-bit hanya menyimpan sekitar $16$ angka desimal dengan tepat; \texttt{isqrt} bekerja
pada bilangan bulat dan memberi $200.000.000$. Pelajarannya umum: membalik rumus
penomoran yang eksak menuntut aritmetika bilangan bulat, seperti dibahas pada Bab 22.

\begin{untukdosen}
Empat pertemuan. Pertama: operasi himpunan dan De Morgan, tiap bukti ditulis sebagai
rantai ekuivalensi, bukan dua inklusi terpisah. Kedua: hasilkali kartesius, relasi, lalu
penyaringan menjadi urutan parsial dan relasi ekuivalensi. Ketiga: kelas, partisi, dan
fungsi, dengan Lemma kelas dibuktikan di papan. Keempat: kardinalitas sampai Teorema
Cantor, ditambah sketsa Cantor--Bernstein bila waktu memungkinkan.

Dua miskonsepsi yang paling sering muncul. Pertama, mahasiswa mengira relasi yang
simetris dan transitif otomatis refleksif; penyangkal $R = \{(1,1),(2,2),(1,2),(2,1)\}$
pada $\{1,2,3\}$ biasanya cukup, dan Latihan 2 meminta penilaian klaim itu. Kedua,
mahasiswa mengira $\mathbb{Q}$ lebih besar daripada $\mathbb{N}$ karena kerapatannya;
mintalah sepuluh suku pertama enumerasi tinggi di papan, lalu tanyakan posisi $7/3$.

Aktivitas yang berhasil: bagi papan menjadi "bilangan asli" dan "rasional positif", lalu
minta dua relawan memasangkan anggota selama dua menit. Pemasangan naif akan gagal,
lalu bersama-sama susun pemasangan menurut tinggi. Mahasiswa siap lanjut bila dapat
menjelaskan mengapa daftar itu tidak melewatkan satu pun rasional, dan mengapa daftar
serupa gagal untuk $[0,1]$.
\end{untukdosen}

\section{Ringkasan}
\tingkat{dasar}

\begin{ringkasan}
\begin{itemize}
  \item Gabungan, irisan, selisih, dan komplemen didefinisikan lewat syarat
        keanggotaan; hukum De Morgan adalah hukum logika yang dipindahkan ke
        himpunan, dan versinya berlaku untuk keluarga terindeks apa pun.
  \item $|\mathcal{P}(A)| = 2^{|A|}$ untuk $A$ berhingga, lewat bijeksi ke fungsi
        karakteristik $A \to \{0,1\}$.
  \item Relasi adalah himpunan bagian $A \times B$; pada himpunan $n$ anggota ada
        $2^{n^{2}}$ relasi, dan empat sifat menyaringnya menjadi relasi ekuivalensi
        dan urutan parsial.
  \item Relasi ekuivalensi dan partisi adalah dua wajah objek yang sama; kelasnya
        selalu tidak kosong, saling lepas, dan menutupi seluruh himpunan.
  \item Fungsi bijektif tepat ketika ia punya invers dua sisi yang tunggal, dan pada
        himpunan berhingga berukuran sama ketiga sifat saling mengunci.
  \item $\mathbb{N} \times \mathbb{N}$ dan $\mathbb{Q}$ terhitung lewat penomoran
        eksplisit, sedangkan $\mathcal{P}(\mathbb{N})$ dan $\mathbb{R}$ tak terhitung
        lewat argumen diagonal Cantor.
  \item Supremum adalah batas atas terkecil; keberadaannya dijamin aksioma
        kelengkapan, sedangkan ciri epsilon dan pembangunannya dilanjutkan di Bab 3.
\end{itemize}
\end{ringkasan}

\section{Latihan}
\tingkat{menengah}

\begin{latihan}
\begin{enumerate}
  \item Buktikan $A \cap (B \cup C) = (A \cap B) \cup (A \cap C)$ dan
        $A \setminus (B \setminus C) = (A \setminus B) \cup (A \cap C)$ dengan rantai
        ekuivalensi keanggotaan.
  \item Nilai klaim ini: benar, salah, atau belum cukup bukti. "Setiap relasi pada
        himpunan tak kosong yang simetris dan transitif pasti refleksif." Bila salah,
        berikan contoh penyangkal sekecil mungkin; bila belum cukup bukti, sebutkan
        bukti tambahan apa yang diperlukan.
  \item Buktikan bahwa bila $f : A \to B$ injektif, maka $S \mapsto f(S)$ adalah
        fungsi injektif dari $\mathcal{P}(A)$ ke $\mathcal{P}(B)$. Jelaskan di mana
        sifat injektif dipakai dan mengapa kesimpulannya gagal tanpanya.
  \item Buktikan lemma pasangan terurut Kuratowski:
        $\{\{a\},\{a,b\}\} = \{\{c\},\{c,d\}\}$ jika dan hanya jika $a = c$ dan
        $b = d$. Tangani kasus $a = b$ secara terpisah.
  \item Lengkapi bagian pertama Lemma hitungan: untuk $f : A \to B$ dengan $A$
        berhingga, $|f(A)| \le |A|$ dan kesamaan berlaku hanya bila $f$ injektif.
  \item Hitung banyaknya relasi ekuivalensi pada himpunan $5$ anggota yang memiliki
        tepat dua kelas, lalu periksa jawabannya dengan Python.
  \item Hitung $\pi(7,5)$ dengan rumus pasangan Cantor, lalu tentukan $(m,n)$ dengan
        $\pi(m,n) = 1000$. Periksa kedua jawaban dengan Python.
  \item Tentukan posisi $7/3$ dalam enumerasi rasional positif menurut tinggi, lalu
        tentukan bilangan pada posisi ke-30 dan ke-41. Periksa dengan Python.
  \item Buktikan bahwa bila $A_1, A_2, A_3, \dots$ masing-masing terhitung dan
        penomoran setiap $A_n$ diberikan eksplisit, maka $\bigcup_{n \ge 1} A_n$
        terhitung.
  \item Tunjukkan $(0,1)$ berkekuatan sama dengan $\mathbb{R}$ melalui satu bijeksi
        eksplisit, dan jelaskan mengapa fungsi yang kita pilih harus dibatasi daerah
        asalnya.
\end{enumerate}
\end{latihan}

Petunjuk untuk tiga butir tersulit. Latihan 5: pada langkah induksi tulis
$A = A' \cup \{a\}$ dengan $a \notin A'$, lalu bandingkan $f(A)$ dengan
$f(A') \cup \{f(a)\}$; kesamaan ukuran menuntut $f(a) \notin f(A')$ dan
$|f(A')| = |A'|$ sekaligus. Latihan 6: banyaknya cara memisahkan $n$ anggota menjadi
dua bagian tak kosong adalah $(2^{n}-2)/2$, sebab memilih bagian pertama atau kedua
memberi pemisahan yang sama. Latihan 9: susun penomoran seluruh $A_n$ sebagai tabel dua
kolom, baca diagonal demi diagonal seperti pada bukti bahwa
$\mathbb{N} \times \mathbb{N}$ terhitung, lalu buang suku berulang.

\begin{untukpeneliti}
Yang benar-benar terbukti di sini ada empat. Teorema Cantor tidak memakai aksioma
pilihan dan bersifat konstruktif: dari sebarang $f$ kita membangun $D$ yang tidak
berada di daerah hasilnya. Ketakterhitungan $\mathbb{R}$ bergantung pada ekspansi biner
setiap titik $[0,1]$; fakta itu memerlukan kelengkapan dan di sini hanya dirujuk ke
Bab 3 dan Bab 5.

Teorema Cantor--Bernstein--Schroeder di bab ini hanya sketsa. Pembuktian penuh
memerlukan pembukuan kedalaman rantai yang hati-hati; versi formal lengkapnya ada
pada Avigad--Massot, \emph{Mathematics in Lean}, bagian 4.3, sedangkan perlakuan
klasik untuk himpunan, fungsi, dan relasi ekuivalensi ada pada Hoffman dan Kunze,
\emph{Linear Algebra}, Apendiks A.1--A.3. Rujukan primer untuk diagonal dan
ketakterhitungan adalah Cantor, "Ueber eine Eigenschaft des Inbegriffes aller
reellen algebraischen Zahlen" (1874) dan "Ueber eine elementare Frage der
Mannigfaltigkeitslehre" (1891).

Pertanyaan terbuka yang paling terkenal: adakah kardinal di antara $|\mathbb{N}|$ dan
$|\mathbb{R}|$? Hipotesis kontinum menyatakan tidak ada. Godel (1940) menunjukkan
hipotesis itu konsisten dengan ZFC, Cohen (1963) menunjukkan ingkarannya juga konsisten.
Jadi pertanyaan itu tidak dapat dijawab di dalam ZFC.

Dua batas teknis. Pertama, $\sim$ tidak boleh disebut relasi ekuivalensi pada
"himpunan semua himpunan", sebab kumpulan itu bukan himpunan; perlakuan yang benar
memakai kelas sebagai objek metateori. Kedua, akibat "gabungan terhitung dari himpunan
terhitung adalah terhitung" memerlukan pemilihan terhitung bila penomorannya tidak
diberikan, dan Latihan 9 sengaja mensyaratkan penomoran eksplisit.
\end{untukpeneliti}

\peta{Bab 2 pada Abbott, \emph{Understanding Analysis}, bagian "Some Preliminaries"
(himpunan, fungsi, kardinalitas) dan "Cantor's Theorem"; Hoffman dan Kunze,
\emph{Linear Algebra}, Apendiks A.1 (himpunan), A.2 (fungsi dan invers), A.3 (relasi
ekuivalensi, kelas, kuosien); Avigad--Massot, \emph{Mathematics in Lean}, Bab 4
("Sets and Functions") untuk ekstensionalitas, prapeta, dan
Cantor--Bernstein--Schroeder, serta Bab 6.1 untuk himpunan berhingga.}

\chapter{Aksioma Bilangan Real dan Supremum}

Sebuah program mencari panjang diagonal persegi bersisi satu. Ia memelihara dua
bilangan, $a=1$ dan $b=2$, dengan $a^2-2<0<b^2-2$. Setiap langkah ia menghitung titik
tengah $m=(a+b)/2$, menilai $m^2-2$, lalu mengganti $a$ atau $b$ dengan $m$ supaya
tanda kedua ujung tetap berbeda. Setelah 40 langkah lebar selangnya di bawah $10^{-12}$,
setelah 100 langkah di bawah $10^{-30}$. Tetapi program itu tidak pernah mengeluarkan
jawaban: tidak ada langkah yang menghasilkan $x^2-2=0$ persis, dan kedua ujung selang
terus mengapit sesuatu yang bukan bilangan rasional.

Kegagalan itu bukan kesalahan algoritme. Ia menunjukkan bahwa bilangan rasional punya
lubang, dan bahwa himpunan $\{q\in\mathbb{Q}: q^2<2\}$ tidak mempunyai batas atas
terkecil di dalam $\mathbb{Q}$.

Bab 1 mengajari kita menulis bukti. Bab 2 memberi bahasa himpunan dan fungsi, sekaligus
memakai $|x|$ dan $\sup$ sebagai alat hitung tanpa mempertanyakan asalnya. Yang belum
ada adalah objek yang menjamin alat itu sah. Bab ini membangunnya, bukan dari desimal,
melainkan dari tiga kelompok pernyataan yang diterima tanpa bukti: hukum lapangan,
hukum urutan, dan kelengkapan. Sisanya, termasuk keberadaan $\sqrt{2}$, diturunkan
dari situ.

\begin{untukanak}
Bayangkan sebuah penggaris yang hanya punya tanda pecahan: setengah, seperempat,
seperdelapan, dan terus makin rapat. Tidak ada angka lain di penggaris itu. Sekarang
ukur diagonal sebuah persegi yang sisinya satu jengkal. Ujung diagonal jatuh di antara
dua tanda, dan selalu begitu, tidak peduli seberapa rapat tandanya. Kita bisa mendekat
sejauh yang kita mau, tetapi tidak pernah tepat.

Pertanyaannya: apakah panjang itu benar-benar ada? Penggarisnya jelas tidak punya
jawaban. Namun diagonal itu, kalau dijinjing atau dipasang di lantai, punya panjang
tertentu. Panjang itu tidak hilang hanya karena alat ukur kita tidak memuatnya.

Cerita yang sama terjadi pada bilangan. Kalau kita hanya punya bilangan pecahan, ada
titik di garis yang tidak punya nama. Matematikawan menutup lubang itu dengan satu
janji: setiap kumpulan bilangan yang terbatas ke atas punya atap paling rendah. Di bab
ini kita akan tahu mengapa janji itu perlu, apa akibatnya, dan bagaimana ia memberi
nama pada panjang diagonal tadi, yaitu akar dua.
\end{untukanak}

\section{Aksioma lapangan}
\tingkat{dasar}

Kita tidak mendefinisikan bilangan real sebagai desimal tak berujung: cara itu membuat
penjumlahan rumit, dan membuat definisi batas atas terkecil berputar-putar, sebab
``supremum semua pemotongan desimal'' sudah memakai bilangan yang sedang dibangun.
Lebih bersih menyebut sifatnya lebih dulu, lalu menurunkan sisanya.

Lapisan pertama bersifat aljabar. Kita mengandaikan dua operasi yang menaati aturan
yang kita pakai sejak sekolah; definisi berikut hanya memberi nama aturan itu.

\begin{definisi}[Lapangan]\index{lapangan}
Himpunan $F$ disebut \emph{lapangan} jika pada anggotanya terdefinisi penjumlahan dan
perkalian sehingga: hasil kedua operasi selalu berada di $F$; kedua operasi komutatif
dan asosiatif; ada unsur $0$ dengan $0+x=x$ dan unsur $1\neq 0$ dengan $1\cdot x=x$;
setiap $x$ punya lawan $-x$ dengan $x+(-x)=0$; setiap $x\neq 0$ punya kebalikan $1/x$
dengan $x\cdot(1/x)=1$; dan berlaku $x(y+z)=xy+xz$.
\end{definisi}

Aksioma itu tidak menyebut urutan, tidak menyebut jarak, dan tidak menyebut
kelengkapan. Karena itu himpunan lain juga memenuhinya. Bilangan rasional dengan
penjumlahan dan perkalian biasa adalah lapangan. Pertanyaan yang wajar: apa saja yang
sudah terkandung di dalam aturan itu?

Tiga pernyataan aljabar yang biasa kita hafal ternyata akibat, bukan aksioma; yang
paling sering dipakai adalah hasil kali dengan nol yang memusnahkan.

\begin{teorema}
Dalam sebarang lapangan, untuk semua $x,y$ berlaku: (a) $0\cdot x=0$;
(b) $(-x)y=-(xy)$; (c) $(-x)(-y)=xy$; (d) jika $xy=0$ maka $x=0$ atau $y=0$.
\end{teorema}

\begin{proof}
(a) Dari $0+0=0$ dan hukum distributif, $0\cdot x=(0+0)x=0\cdot x+0\cdot x$.
Tambahkan $-(0\cdot x)$ pada kedua ruas. Ruas kiri menjadi $0$, ruas kanan menjadi
$0\cdot x$. Jadi $0=0\cdot x$.

(b) Hitung $(-x)y+xy=(-x+x)y=0\cdot y=0$, memakai (a). Jadi $(-x)y$ lawan dari $xy$,
dan lawan itu tunggal.

(c) Terapkan (b) dua kali: $(-x)(-y)=-\bigl(x(-y)\bigr)=-\bigl(-(xy)\bigr)=xy$.
Langkah terakhir memakai $--a=a$, yang juga akibat ketunggalan lawan.

(d) Andaikan $x\neq 0$ dan $xy=0$. Kalikan dengan $1/x$: $(1/x)(xy)=0$, sehingga
$y=0$. Jadi salah satu faktornya nol.
\end{proof}

Tidak satu pun dari keempat pernyataan itu diandaikan; semuanya turun dari hukum
distributif dan keberadaan lawan. Pola ini dipakai sepanjang buku: sebut aturan
sesedikit mungkin, lalu panen akibatnya.

\kamus{lapangan}{field}{Himpunan dengan dua operasi hitung yang aturannya tetap, seperti
kotak alat yang penutupnya selalu rapat: hasil hitungan tidak pernah keluar dari kotak.}

Bab 2 memakai $\mathbb{Q}$ sebagai himpunan terhitung. Di sini ia muncul dalam peran
lain: lapangan yang terurut. Dua sifat itu belum cukup, dan itu isi bagian berikutnya.

\section{Aksioma urutan}
\tingkat{dasar}

Lapangan saja tidak bisa membandingkan dua bilangan. Kita memerlukan relasi ``lebih
kecil dari'' yang menghormati penjumlahan dan perkalian, sebab tanpa itu pertidaksamaan
tidak bisa diolah. Aksioma urutan mengatakan tiga hal: setiap dua bilangan bisa
dibandingkan, perbandingan menular, dan menambah atau mengalikan dengan bilangan positif
tidak merusaknya.

\begin{definisi}[Lapangan terurut]\index{lapangan terurut}
Lapangan $F$ disebut \emph{terurut} jika ada relasi $<$ pada $F$ sehingga untuk setiap
$a,b,c$ berlaku: tepat satu di antara $a<b$, $a=b$, $b<a$ benar; jika $a<b$ dan $b<c$
maka $a<c$; jika $a<b$ maka $a+c<b+c$; dan jika $a<b$ serta $c>0$ maka $ac<bc$.
\end{definisi}

Hampir semua aturan tanda yang kita hafal adalah akibat definisi itu, bukan tambahan.
Bukti berikut menurunkan lima di antaranya sekaligus; polanya sama: tambahkan sesuatu
pada kedua ruas, atau kalikan dengan bilangan yang tandanya diketahui.

\begin{proposisi}
Dalam setiap lapangan terurut berlaku: (a) jika $x>0$ maka $-x<0$, dan sebaliknya;
(b) jika $x>0$ dan $y<z$ maka $xy<xz$; (c) jika $x<0$ dan $y<z$ maka $xy>xz$;
(d) jika $x\neq 0$ maka $x^2>0$, khususnya $1>0$; (e) jika $0<x<y$ maka
$0<1/y<1/x$.
\end{proposisi}

\begin{proof}
(a) Jika $x>0$, tambahkan $-x$ pada kedua ruas: $0>-x$. Sebaliknya, jika $x<0$ maka
$0-x>0$, yaitu $-x>0$.

(b) Karena $z>y$, kita punya $z-y>0$. Kalikan dengan $x>0$: $x(z-y)>0$. Jadi
$xz=x(z-y)+xy>xy$.

(c) Dari (a), $-x>0$. Terapkan (b) pada $-x$ dan $y<z$: $(-x)y<(-x)z$. Menurut
Teorema sebelumnya, ruas kiri adalah $-(xy)$ dan ruas kanan $-(xz)$. Jadi
$-(xy)<-(xz)$, dan dengan (a) lagi $xy>xz$.

(d) Jika $x>0$, aksioma urutan pada $y=x>0$ memberi $x^2>0$. Jika $x<0$ maka $-x>0$,
sehingga $(-x)^2>0$; padahal $(-x)^2=x^2$. Untuk $x=0$ jelas $x^2=0$. Karena
$1=1^2$, kita peroleh $1>0$.

(e) Karena $y>0$ dan $1/y\neq 0$, dari (d) $1/y>0$; sama untuk $1/x$. Kalikan
$x<y$ dengan bilangan positif $(1/x)(1/y)$: hasilnya $1/y<1/x$.
\end{proof}

Dua akibat yang sering dipakai: $a<b$ setara dengan $b-a>0$, dan $a<b$ setara dengan
$-b<-a$; yang kedua menjelaskan mengapa mengalikan dengan bilangan negatif membalik
arah, sebab mengalikan dengan $-1$ mencerminkan garis bilangan terhadap nol. Akibat lain:
jika $0<a<b$ maka $a^2<b^2$, cukup dengan mengalikan $a<b$ dengan $a>0$ lalu dengan
$b>0$ dan menjumlahkannya. Jadi membandingkan kuadrat setara dengan membandingkan
bilangan positif aslinya; setiap kali kita menulis $q^2<2$ untuk $q$ rasional, kita
membandingkan $q$ dengan sesuatu yang belum tentu rasional.

\kamus{lapangan terurut}{ordered field}{Lapangan yang anggotanya bisa diurutkan seperti
anak tangga, dan aturan hitungnya tidak merusak urutan itu.}

Perhatikan apa yang belum kita punya. Kita belum tahu apakah ada bilangan positif yang
lebih kecil dari semua $1/n$, dan belum tahu apakah $\{q\in\mathbb{Q}: q^2<2\}$ punya
atap. Itu bukan soal urutan, melainkan soal kelengkapan.

\section{Nilai mutlak dan ketaksamaan segitiga}
\tingkat{menengah}

Di Bab 2 kita menulis $|x-y|$ sebagai jarak, dan memakainya untuk membandingkan galat,
tanpa membuktikan bahwa lambang itu berperilaku seperti jarak. Sekarang sifatnya
diturunkan dari urutan. Nilai mutlak membuang tanda dan menyisakan jarak ke titik nol;
definisi berikut menulis dua cabang, dan cabang keduanya yang biasa dipakai.

\begin{definisi}[Nilai mutlak]\index{nilai mutlak}
Untuk $x\in\mathbb{R}$, $|x|=x$ jika $x\ge 0$ dan $|x|=-x$ jika $x<0$. Secara setara,
$|x|=\max\{x,-x\}$, dan untuk bilangan real $|x|=\sqrt{x^2}$ dengan akar yang
dibahas di bagian berikutnya.
\end{definisi}

Dari definisi itu langsung: $|x|\ge 0$, $|{-x}|=|x|$, $|x|^2=x^2$, dan $|xy|=|x|\,|y|$.
Yang terakhir terbukti dengan memeriksa empat kombinasi tanda, atau dengan mencatat
bahwa kedua ruas tidak negatif dan kuadratnya sama.

Pertidaksamaan yang memuat nilai mutlak paling mudah diolah setelah diubah menjadi dua
pertidaksamaan biasa. Lemma berikut adalah alasan mengapa perubahan itu sah.

\begin{lemma}
Untuk $c\ge 0$, $|x|\le c$ jika dan hanya jika $-c\le x\le c$.
\end{lemma}

\begin{proof}
Jika $|x|\le c$: untuk $x\ge 0$ kita punya $x=|x|\le c$ dan $-c\le 0\le x$; untuk
$x<0$ kita punya $-x=|x|\le c$, sehingga $x\ge -c$, dan $x<0\le c$. Jadi
$-c\le x\le c$.

Sebaliknya, jika $-c\le x\le c$: untuk $x\ge 0$, $|x|=x\le c$; untuk $x<0$,
$|x|=-x\le c$ karena $x\ge -c$. Dalam kedua kasus $|x|\le c$.
\end{proof}

Ketaksamaan segitiga adalah satu-satunya alat yang kita butuhkan untuk mengendalikan
selisih dua bilangan. Namanya berasal dari gambar: panjang satu sisi segitiga tidak
melebihi jumlah panjang dua sisi lainnya.

\begin{teorema}[Ketaksamaan segitiga]\index{ketaksamaan segitiga}
Untuk semua $x,y\in\mathbb{R}$, $|x+y|\le |x|+|y|$. Selain itu
$\bigl||x|-|y|\bigr|\le |x-y|$.
\end{teorema}

\begin{proof}
Dari definisi, $-|x|\le x\le |x|$ dan $-|y|\le y\le |y|$. Jumlahkan kedua pasangan:
$-(|x|+|y|)\le x+y\le |x|+|y|$. Menurut lemma sebelumnya dengan $c=|x|+|y|$, diperoleh
$|x+y|\le |x|+|y|$.

Untuk bagian kedua, tulis $x=(x-y)+y$ dan terapkan bagian pertama:
$|x|\le |x-y|+|y|$, sehingga $|x|-|y|\le |x-y|$. Dengan menukar peran $x$ dan $y$
diperoleh $|y|-|x|\le |x-y|$. Karena $c\le d$ dan $-c\le d$ setara dengan $|c|\le d$,
maka $\bigl||x|-|y|\bigr|\le |x-y|$.
\end{proof}

Sebagai pemeriksaan, ambil $x=3$ dan $y=-5$: $|x+y|=2$ sedangkan $|x|+|y|=8$, jadi
ketaksamaannya longgar. Untuk $x=3$ dan $y=5$ kedua ruas bernilai $8$ dan ketaksamaan
menjadi persamaan; itu terjadi tepat ketika tandanya sama atau salah satunya nol.

Kegunaan praktisnya: batas pada jarak menjadi batas pada fungsi. Jika $|x-\sqrt2|<\delta$,
maka
\[
|x^2-2|=|x-\sqrt2|\,|x+\sqrt2|<\delta\,(2\sqrt2+\delta),
\]
karena $|x+\sqrt2|\le|x-\sqrt2|+2\sqrt2$. Untuk $\delta=10^{-3}$ batas itu bernilai
$0.002829427$, dan nilai sebenarnya di $x=\sqrt2+10^{-3}$ juga $0.002829427$ sebab
$x^2-2=2\sqrt2\,\delta+\delta^2$ tepat: batas itu ketat. Untuk $\delta=10^{-6}$ batas
yang sama mengecil menjadi $2.828428\times10^{-6}$.

\kamus{nilai mutlak}{absolute value}{Jarak sebuah titik ke nol pada garis bilangan,
tanpa memedulikan sebelah mana titik itu berada.}

Ketaksamaan segitiga hanya memakai aksioma urutan, bukan kelengkapan, jadi ia berlaku
juga di $\mathbb{Q}$ dan tidak bisa membedakan keduanya. Alat pembeda muncul di bagian
berikutnya.

\section{Kelengkapan: batas atas terkecil}
\tingkat{menengah}

Urutan memberi perbandingan, tetapi tidak memberi atap. Himpunan
$A=\{1-1/n: n\in\mathbb{N}\}$ memuat bilangan yang makin dekat ke satu tanpa pernah
sama dengan satu: satu bukan anggotanya, tetapi batasnya paling rapat. Himpunan
$B=\{q\in\mathbb{Q}: q^2<2\}$ terbatas, tetapi tidak punya batas rapat di
$\mathbb{Q}$. Keduanya memperlihatkan bahwa batas longgar dan batas paling rapat perlu
nama berbeda.

\begin{untukanak}
Bayangkan kita harus memasang atap tenda untuk sekelompok anak. Atap terlalu rendah:
ada kepala yang menonjol. Atap terlalu tinggi: tenda boros dan goyah. Yang kita cari
adalah atap paling rendah yang masih membuat semua anak muat. Sebut saja atap pas.

Sekarang ada yang aneh. Di kelompok itu tidak ada anak tertinggi; selalu ada yang
sedikit lebih tinggi. Namun atap pas tetap ada, karena tinggi mereka terbatas: tidak
ada yang setinggi pohon. Atap pas ditentukan oleh seluruh kelompok, bukan oleh satu
anak.

Di bab ini atap pas punya nama: batas atas terkecil. Sifat yang menjamin atap pas
selalu ada disebut kelengkapan. Bilangan pecahan tidak lengkap: ada kelompok yang
tingginya terbatas, tetapi atap pasnya tidak ada di antara bilangan pecahan. Pertanyaan
yang bisa dibawa pulang: kalau tidak ada anak tertinggi, bagaimana caranya membuktikan
bahwa atap pas itu memang yang paling rendah?
\end{untukanak}

Batas atas adalah langit-langit yang memuat semua anggota; batas atas terkecil adalah
langit-langit terendah di antaranya. Definisi berikut menamai keduanya beserta cerminnya
di bawah.

\begin{definisi}[Batas atas dan supremum]\index{batas atas}\index{supremum}
Misalkan $A\subseteq\mathbb{R}$ dan $A\neq\varnothing$. Bilangan $u$ disebut
\emph{batas atas} $A$ jika $a\le u$ untuk setiap $a\in A$. Jika $A$ punya
sekurang-kurangnya satu batas atas, $A$ disebut terbatas di atas. Bilangan $s$ disebut
\emph{supremum} $A$, ditulis $s=\sup A$, jika $s$ batas atas $A$ dan $s\le u$ untuk
setiap batas atas $u$. Secara kembar, $t$ disebut \emph{infimum} $A$, ditulis
$t=\inf A$, jika $t$ batas bawah dan $t\ge v$ untuk setiap batas bawah $v$.
\end{definisi}

Dua syarat pada definisi supremum sering tertukar: yang pertama mengatakan $s$ masih
batas atas, yang kedua bahwa tidak ada batas atas yang lebih kecil. Pada $A=\{1-1/n\}$,
bilangan $1$ memenuhi keduanya sedangkan $2$ hanya yang pertama; himpunan itu juga tidak
punya anggota terbesar, jadi ``maksimum'' dan ``supremum'' bukan sinonim.

\begin{bacakeras}
Dua suara. Suara pertama bertanya, suara kedua menjawab. Bacakan bergantian dengan teman.

``Pagar itu apa?'' ``Pagar adalah batas atas: semua anggota kumpulan ada di bawahnya.''
``Berapa banyak pagarnya?'' ``Banyak. Setiap pagar yang lebih tinggi juga pagar.''
``Lalu yang mana yang kita pakai?'' ``Yang paling rendah. Namanya batas atas terkecil.''
``Kalau pagar terendah itu kebetulan ditempati anggota kumpulan?'' ``Maka ia juga maksimum.
Kalau tidak ditempati, ia tetap batas atas terkecil. Dua nama, dua pekerjaan.''

Ulangi tiga kali, makin cepat: pagar tinggi bukan jawaban, pagar terendah baru jawaban.

Baris yang harus diulang pembaca, sekali, dengan suara jelas: batas atas itu pagar; batas
atas terkecil itu pagar terendah yang masih menahan semua anggota.
\end{bacakeras}

Gambar~\ref{fig:sup-garis} menaruh kedua gagasan itu pada satu garis. Anggota
$A=\{1-1/n\}$ mengerumun ke satu arah tanpa pernah mencapai ujungnya; semua bilangan di
kanan satu adalah batas atas, dan yang paling kiri di antara mereka adalah supremum.

\begin{figure}[htbp]
\centering
\begin{tikzpicture}[x=1.5cm,y=1cm,font=\small]
  \draw[-{Latex[length=2mm]},ink,thick] (-0.35,0) -- (5.5,0);
  \foreach \t in {0,1,2,3,4,5}{
    \draw[ink] (\t,0.07) -- (\t,-0.07);
    \node[below=2pt,ink] at (\t,0) {$\t$};
  }
  \foreach \n in {1,...,6}{
    \fill[accent] ({1-1/\n},1.0) circle (1.7pt);
  }
  \draw[accent,thick,fill=white] (1,1.0) circle (1.7pt);
  \node[anchor=west,accent,font=\footnotesize] at (1.18,1.0) {anggota $A=\{1-1/n\}$};
  \draw[layergrey,dashed] (1,0.95) -- (1,0);
  \node[accent,anchor=west,font=\footnotesize] at (1.12,0.45)
    {$\sup A=1$, dan $1$ bukan anggota $A$};
  \fill[layerteal!14] (1,-0.95) rectangle (5.35,-0.6);
  \draw[layerteal,thick] (1,-0.95) rectangle (5.35,-0.6);
  \node[layerteal!80!black,font=\footnotesize] at (3.2,-0.775) {semua batas atas $A$};
  \draw[accent,line width=1.4pt] (1,-1.02) -- (1,-0.52);
\end{tikzpicture}
\caption{Himpunan $A=\{1-1/n:n\in\mathbb{N}\}$ pada garis bilangan. Setiap $u\ge 1$
adalah batas atas; yang terkecil, $u=1$, adalah $\sup A$. Bilangan $1$ bukan anggota
$A$, jadi $A$ tidak mempunyai maksimum.}
\label{fig:sup-garis}
\end{figure}

\kamus{supremum}{supremum (least upper bound)}{Atap paling rendah yang masih menutupi
seluruh anggota himpunan; atap itu tidak harus bagian dari himpunan.}

Supremum, kalau ada, hanya ada satu; syarat keduanya bisa ditulis dengan epsilon, dan
bentuk itu yang paling sering dipakai dalam bukti.

\begin{lemma}
(i) Jika $A$ mempunyai supremum, maka supremumnya tunggal. (ii) Bilangan $s$ adalah
supremum $A$ jika dan hanya jika $s$ batas atas $A$ dan untuk setiap
$\varepsilon>0$ terdapat $a\in A$ dengan $a>s-\varepsilon$.
\end{lemma}

\begin{proof}
(i) Misalkan $s$ dan $s'$ keduanya supremum. Karena $s$ batas atas dan $s'$ supremum,
$s'\le s$. Karena $s'$ batas atas dan $s$ supremum, $s\le s'$. Jadi $s=s'$.

(ii) Andaikan $s$ batas atas dan syarat epsilon berlaku, tetapi ada batas atas $u<s$.
Ambil $\varepsilon=s-u>0$. Menurut syarat epsilon ada $a\in A$ dengan
$a>s-\varepsilon=u$, bertentangan dengan $u$ batas atas. Sebaliknya, jika $s$ supremum
dan ada $\varepsilon>0$ sehingga $a\le s-\varepsilon$ untuk semua $a\in A$, maka
$s-\varepsilon$ batas atas yang lebih kecil dari $s$, bertentangan dengan $s$ batas
atas terkecil.
\end{proof}

\begin{resep}
Untuk membuktikan bahwa bilangan $s$ adalah supremum himpunan $A$:
\begin{enumerate}
\item Tunjukkan $A\neq\varnothing$ dengan menyebut satu anggotanya.
\item Buktikan $s$ batas atas: ambil $a\in A$ sebarang, tunjukkan $a\le s$.
\item Buktikan tidak ada batas atas yang lebih kecil: ambil $\varepsilon>0$, hasilkan
      $a\in A$ dengan $a>s-\varepsilon$, misalnya $a=1-1/n$ dengan $n>1/\varepsilon$.
\item Simpulkan: $s$ batas atas dan setiap bilangan di bawah $s$ bukan batas atas, jadi
      $s=\sup A$.
\end{enumerate}
Peringatan urutan: langkah 2 harus selesai sebelum langkah 3. Memilih $a>s-\varepsilon$
lebih dulu tidak membuktikan apa pun, sebab $s$ belum diketahui batas atas. Pada
$A=\{1-1/n\}$, bilangan $1/2$ lulus langkah 3 tetapi gagal langkah 2.
\end{resep}

Aksioma berikut membedakan $\mathbb{R}$ dari $\mathbb{Q}$. Ia tidak bisa dibuktikan dari
hukum lapangan dan urutan, sebab ia pernyataan tentang kelengkapan garis.

\begin{teorema}[Sifat batas atas terkecil]\index{aksioma kelengkapan}
Setiap himpunan tak kosong $A\subseteq\mathbb{R}$ yang terbatas di atas mempunyai
supremum di $\mathbb{R}$.
\end{teorema}

\begin{proof}[Bukti]
Pernyataan ini diterima sebagai aksioma kelengkapan. Rumusan lain yang setara: irisan
selang bersarang, dan kekonvergenan barisan Cauchy. Kita memilih supremum karena paling
langsung dipakai; pembangunan $\mathbb{R}$ dari $\mathbb{Q}$ yang membuktikannya ada
pada Apendiks Bab 1 Rudin.
\end{proof}

Infimum tidak perlu diaksiomakan: ia supremum yang dilihat di cermin.

\begin{proposisi}
Jika $A$ tak kosong dan terbatas di bawah, maka $\inf A=-\sup(-A)$, dengan
$-A=\{-a: a\in A\}$.
\end{proposisi}

\begin{proof}
Himpunan $-A$ tak kosong dan terbatas di atas, jadi $s=\sup(-A)$ ada. Untuk setiap
$a\in A$ berlaku $-a\le s$, sehingga $a\ge -s$: jadi $-s$ batas bawah $A$. Sekarang
ambil $t>-s$. Maka $-t<s$, sehingga $-t$ bukan batas atas $-A$; ada $a\in A$ dengan
$-a>-t$, yaitu $a<t$. Jadi $t$ bukan batas bawah $A$. Karena itu $-s$ batas bawah
terbesar.
\end{proof}

Akibat pertama kelengkapan: bilangan asli tidak punya langit-langit. Kedengarannya
sepele, tetapi tanpa kelengkapan ia tidak bisa dibuktikan.

\begin{teorema}[Sifat Archimedes]\index{sifat archimedes}
Untuk setiap $x>0$ dan setiap $y\in\mathbb{R}$ terdapat $n\in\mathbb{N}$ dengan
$nx>y$.
\end{teorema}

\begin{proof}
Andaikan sebaliknya: ada $x>0$ dan $y$ sehingga $nx\le y$ untuk semua $n$. Maka himpunan
$S=\{nx: n\in\mathbb{N}\}$ tak kosong dan terbatas di atas oleh $y$, sehingga
$\alpha=\sup S$ ada. Karena $x>0$, bilangan $\alpha-x$ lebih kecil dari $\alpha$ dan
karena itu bukan batas atas $S$. Jadi ada $m$ dengan $mx>\alpha-x$, yang memberi
$(m+1)x>\alpha$. Padahal $(m+1)x\in S$, dan $\alpha$ batas atas $S$. Kontradiksi.
\end{proof}

Versi resiprokalnya yang paling sering dipakai: untuk setiap $\varepsilon>0$ ada
$n$ dengan $1/n<\varepsilon$, cukup dengan menerapkan teorema pada $x=\varepsilon$ dan
$y=1$. Kalimat ``untuk $n$ cukup besar'' adalah singkatan dari akibat itu, dan sekarang
kita tahu ia sah.

Sifat itu juga memusnahkan bilangan positif yang sangat kecil: tidak ada $\delta>0$
yang lebih kecil dari $1/n$ untuk semua $n$, sebab $1/\delta$ akan menjadi batas atas
$\mathbb{N}$. Karena itu analisis real tidak memerlukan ``infinitesimal''.

Akibat kedua: setiap selang terbuka, sekecil apa pun, memuat bilangan rasional.

\begin{teorema}[Kerapatan $\mathbb{Q}$]\index{kerapatan}
Jika $x<y$ bilangan real, maka ada $p\in\mathbb{Q}$ dengan $x<p<y$.
\end{teorema}

\begin{proof}
Karena $y-x>0$, Sifat Archimedes memberi $n\in\mathbb{N}$ dengan $n(y-x)>1$. Terapkan
Archimedes lagi untuk memperoleh bilangan bulat $m_1>nx$ dan $m_2>-nx$; jadi
$-m_2<nx<m_1$. Karena $nx$ terkurung di antara dua bilangan bulat, ada bilangan bulat
$m$ dengan $m-1\le nx<m$. Maka $nx<m$, dan dari $m-1\le nx$ diperoleh
$m\le nx+1<ny$. Karena $n>0$, pembagian memberi $x<m/n<y$.
\end{proof}

Bukti itu konstruktif: ia memberi resep untuk menemukan $p$, yaitu memilih $n$ cukup
besar lalu membulatkan $nx$ ke atas. Ia tidak memberi penyebut yang kecil, jadi
aproksimasi efisien masih menuntut gagasan lain.

\section{Mengapa bilangan rasional belum cukup}
\tingkat{lanjut}

Program di pembuka gagal karena himpunan yang ia apit tidak punya atap di
$\mathbb{Q}$. Mari kita buktikan, bukan hanya menunjukkan angkanya. Ambil
$A=\{q\in\mathbb{Q}: q^2<2\}$. Himpunan ini tidak kosong, sebab $1\in A$, dan terbatas
di atas di dalam $\mathbb{Q}$: jika $q\ge 3/2$ maka $q^2\ge 9/4>2$, jadi $3/2$ batas
atas.

Selang-selang yang dipertahankan program itu digambar sebagai tumpukan batang yang makin
pendek (Gambar~\ref{fig:selang-bersarang}). Setiap batang masih memuat akar dua dan
panjangnya menjadi separuh, sementara semua ujungnya tetap rasional; batang-batang itu
menyempit ke satu titik tanpa pernah menyentuhnya.

\begin{figure}[htbp]
\centering
\begin{tikzpicture}[x=9cm,y=1cm,font=\small]
  \foreach \y/\a/\b/\lab/\pan in {%
    0/1/2/{$[1,\;2]$}/{$1$},
    0.55/1/1.5/{$[1,\;\tfrac32]$}/{$1/2$},
    1.1/1.25/1.5/{$[\tfrac54,\;\tfrac32]$}/{$1/4$},
    1.65/1.375/1.5/{$[\tfrac{11}{8},\;\tfrac32]$}/{$1/8$},
    2.2/1.375/1.4375/{$[\tfrac{11}{8},\;\tfrac{23}{16}]$}/{$1/16$},
    2.75/1.40625/1.4375/{$[\tfrac{45}{32},\;\tfrac{23}{16}]$}/{$1/32$}}{
    \fill[layerindigo!12] (\a,\y) rectangle (\b,\y+0.3);
    \draw[layerindigo,thick] (\a,\y) rectangle (\b,\y+0.3);
    \node[layerindigo,anchor=east,font=\footnotesize] at (\a-0.018,\y+0.15) {\lab};
    \node[layergrey,anchor=west,font=\footnotesize] at (\b+0.015,\y+0.15) {\pan};
  }
  \draw[accent,dashed,thick] (1.41421,-0.25) -- (1.41421,3.15);
  \node[accent,font=\footnotesize,above] at (1.41421,3.15) {$\sqrt2$};
  \draw[-{Latex[length=1.6mm]},layergrey] (1.0,-0.12) -- (2.0,-0.12);
  \node[layergrey,font=\footnotesize,below] at (1.5,-0.14) {panjang $=1$};
\end{tikzpicture}
\caption{Selang yang dipertahankan program pembuka: $[1,2]$, $[1,\tfrac32]$,
$[\tfrac54,\tfrac32]$, $[\tfrac{11}{8},\tfrac32]$, $[\tfrac{11}{8},\tfrac{23}{16}]$,
$[\tfrac{45}{32},\tfrac{23}{16}]$. Panjangnya turun dari $1$ ke $1/32$; diteruskan tanpa
henti, irisannya menyusut ke satu titik, yaitu $\sqrt2$, yang bukan bilangan rasional.}
\label{fig:selang-bersarang}
\end{figure}

Langkah pertama: tidak ada anggota terbesar. Setiap rasional dengan $q^2<2$ masih bisa
digeser sedikit ke atas tanpa melewati dua, dan besar geserannya punya rumus.

\begin{lemma}
Jika $q\in\mathbb{Q}$, $q>0$, dan $q^2<2$, maka terdapat $h\in\mathbb{Q}$, $h>0$,
dengan $(q+h)^2<2$. Khususnya $q+h\in A$ dan $q+h>q$.
\end{lemma}

\begin{proof}
Ambil $h=\dfrac{2-q^2}{2q+2}$. Karena $q^2<2$ dan $2q+2>0$, maka $h>0$ dan $h$
rasional. Juga $h<1$, sebab $2-q^2<2q+2$ setara dengan $q^2+2q>0$, yang benar untuk
$q>0$. Maka $h^2\le h$, sehingga
\[
(q+h)^2=q^2+2qh+h^2\le q^2+h(2q+1)
=q^2+(2-q^2)\frac{2q+1}{2q+2}<q^2+(2-q^2)=2,
\]
karena $\frac{2q+1}{2q+2}<1$.
\end{proof}

Rumus itu bisa diperiksa dengan angka. Untuk $q=1393/985$ berlaku $2-q^2=1/970225$,
sehingga $h=1/4684660$ dan
\[
q+h=\frac{6625109}{4684660}=1.414213411432206\ldots,
\]
yang kuadratnya $1.999999573074719\ldots<2$. Jadi $1393/985$ bukan anggota terbesar,
meskipun galat kuadratnya sudah sekitar satu per sejuta. Prosedur ini bisa diulang tanpa
henti: tidak ada langkah terakhir.

\begin{tangan}
Bahan: tali 1,5 m, enam lakban, meteran.

Tempel lakban di lantai tegak lurus dinding, pada jarak $60/n$ cm untuk $n=1,\ldots,6$:
60, 30, 20, 15, 12, dan 10 cm. Kamu baru menempel enam, tetapi aturannya tidak berhenti.
Bentangkan tali melintang sebagai pagar, lalu geser ke arah dinding. Setiap kali
berhenti, periksa satu hal: apakah semua lakban ada di sisi tali yang sama dengan
dinding? Selama ya, tali masih pagar.

Tali bisa didekatkan ke dinding tanpa batas dan tetap pagar: 5 cm, 1 cm, 1 mm. Yang
menghentikannya bukan lakban, melainkan dinding. Kumpulan lakbannya tidak punya anggota
terkecil, dan tidak satu pun menempel di dinding. Pagar terendah tetap ada---dinding---
dan ia bukan anggota kumpulan. Itulah supremum.
\end{tangan}

Bukan hanya maksimumnya yang tidak ada. Batas atas terkecilnya pun tidak ada di dalam
bilangan rasional, dan buktinya memerlukan satu fakta tentang akar dua.

\begin{proposisi}
Himpunan $A=\{q\in\mathbb{Q}: q^2<2\}$ tidak mempunyai supremum di $\mathbb{Q}$.
\end{proposisi}

\begin{proof}
Andaikan $p\in\mathbb{Q}$ supremum $A$. Karena $1\in A$, maka $p\ge 1>0$. Kita
tunjukkan $p^2=2$, yang mustahil karena tidak ada bilangan rasional yang kuadratnya dua.

Jika $p^2<2$, lemma sebelumnya memberi $q=p+h\in A$ dengan $q>p$, bertentangan dengan
$p$ batas atas. Jika $p^2>2$, ambil $k=(p^2-2)/(2p)>0$, yang rasional dan lebih kecil
dari $p$ sebab $p^2-2<2p^2$. Untuk setiap $q\in A$,
\[
(p-k)^2=p^2-2pk+k^2>p^2-2pk=2>q^2,
\]
jadi $q<p-k$ untuk $q>0$, dan untuk $q\le 0$ hal itu jelas karena $p-k>0$: bilangan
$p-k$ masih batas atas $A$, tetapi lebih kecil dari $p$.
Bertentangan dengan $p$ batas atas terkecil. Jadi $p^2=2$.

Sekarang $p=m/n$ dengan $m,n$ bulat tanpa faktor bersama. Dari $m^2=2n^2$, $m^2$ genap,
sehingga $m$ genap; tulis $m=2r$. Maka $4r^2=2n^2$, yaitu $n^2=2r^2$, sehingga $n$
juga genap. Keduanya genap, bertentangan dengan pilihan $m,n$.
\end{proof}

Di dalam $\mathbb{R}$ hal itu tidak terjadi. Himpunan yang sama, kali ini dipandang
sebagai himpunan bilangan real, terbatas di atas oleh $2$, jadi menurut aksioma
kelengkapan ia punya supremum. Supremum itu adalah bilangan yang kita namai $\sqrt{2}$.
Bukti berikut menutup dua kemungkinan yang tersisa, dan bentuknya meniru bukti lemma
tadi dengan $q$ diganti bilangan real.

\begin{teorema}[Keberadaan akar dua]\index{akar dua}
Terdapat tepat satu $s>0$ dengan $s^2=2$.
\end{teorema}

\begin{proof}
Misalkan $E=\{t\in\mathbb{R}: t>0,\ t^2<2\}$. Himpunan ini tak kosong karena $1\in E$,
dan terbatas di atas oleh $2$: jika $t>2$ maka $t^2>4>2$, jadi $t\notin E$. Menurut
sifat batas atas terkecil, $s=\sup E$ ada, dan $s\ge 1>0$ karena $1\in E$.

Andaikan $s^2<2$. Ambil $h=\frac12\min\{1,(2-s^2)/(2s+1)\}$. Maka $0<h<1$ dan
$h(2s+1)<2-s^2$: jika nilai minimum yang pertama, $h=\frac12$ dan
$\frac12(2s+1)\le\frac12(2-s^2)<2-s^2$; jika yang kedua, $h(2s+1)=\frac12(2-s^2)$.
Karena $h<1$, kita punya $h^2<h$, sehingga
\[
(s+h)^2=s^2+2sh+h^2<s^2+h(2s+1)<s^2+(2-s^2)=2 .
\]
Jadi $s+h\in E$ padahal $s+h>s$, bertentangan dengan $s$ batas atas $E$. Karena itu
$s^2\ge 2$.

Andaikan $s^2>2$. Ambil $k=(s^2-2)/(2s)>0$; maka $k<s$ karena $s^2-2<2s^2$. Untuk
setiap $t\in E$ dengan $t\ge s-k$ berlaku
\[
t^2\ge(s-k)^2=s^2-2sk+k^2>s^2-2sk=2,
\]
sehingga $t\notin E$. Jadi $s-k$ batas atas $E$ yang lebih kecil dari $s$, bertentangan
dengan $s$ batas atas terkecil. Karena $s^2<2$ dan $s^2>2$ keduanya mustahil, $s^2=2$.

Ketunggalan: jika $s,t>0$ dan $s^2=t^2=2$, maka $(s-t)(s+t)=0$. Karena $s+t>0$,
$s=t$.
\end{proof}

Bilangan $s$ irasional, menurut argumen paritas tadi: lubang di $\mathbb{Q}$ nyata, dan
$\mathbb{R}$ menutupnya dengan tepat satu bilangan. Yang menutupnya bukan urutan dan
bukan hukum lapangan, melainkan kelengkapan. Program pembuka tidak gagal karena kurang
teliti, melainkan karena sistem bilangannya tidak memuat objek yang dicarinya.

Satu hal perlu dinyatakan jujur: kita belum membangun $\mathbb{R}$ dari $\mathbb{Q}$,
hanya mengandaikan $\mathbb{R}$ memenuhi tiga kelompok aksioma. Pembangunannya dikenal.
Dedekind memotong $\mathbb{Q}$ menjadi dua: sebuah \emph{potongan} adalah himpunan
bagian $\alpha\subset\mathbb{Q}$ yang tidak kosong, tidak seluruh $\mathbb{Q}$, turun
tertutup, dan tanpa anggota terbesar. Sifat terakhir persis lemma kita tadi. Cantor
memakai kelas ekuivalen barisan Cauchy rasional; keduanya terbit pada 1872. Rinciannya
ada pada Apendiks Bab 1 Rudin.

Pembangunan itu menjelaskan kedudukan desimal. Untuk $x>0$, ambil $n_0$ bilangan bulat
terbesar dengan $n_0\le x$, lalu $n_1$ angka terbesar dengan $n_0+n_1/10\le x$, lalu
$n_2$ dengan $n_0+n_1/10+n_2/100\le x$, dan seterusnya. Pemotongan itu terbatas di atas
oleh $x$, jadi supremumnya ada dan sama dengan $x$. Desimal bukan bahan dasar bilangan
real, melainkan cara menamai supremum; kelengkapan menjamin namanya ada, bukan bahwa
namanya bisa ditulis dengan pecahan.

\kamus{aksioma kelengkapan}{completeness axiom}{Janji bahwa garis bilangan tidak
berlubang: setiap kumpulan yang terbatas ke atas selalu punya atap paling rendah.}

Keberadaan $\sqrt2$ membuat soal komputasi masuk akal: $s$ bisa dihitung tanpa tahu
digitnya lebih dulu. Iterasi Newton, $x_{k+1}=\frac12(x_k+2/x_k)$, tetap rasional bila
$x_k$ rasional, sehingga seluruh proses berjalan di dalam $\mathbb{Q}$. Dari $s^2=2$
diperoleh
\[
x_{k+1}-s=\frac{(x_k-s)^2}{2x_k},\qquad e_{k+1}=e_k^2\cdot\frac{s}{2x_k},
\]
dengan $e_k=|x_k-s|/s$. Ruas kanan bentuk pertama tidak negatif, jadi setiap langkah
tetap di atas $s$; di sana faktor $\frac{s}{2x_k}$ tidak melebihi $1/2$, sehingga
$e_{k+1}\le e_k^2/2$ dan jumlah digit benar kira-kira berlipat dua setiap langkah.

Satu peringatan: kita belum membuktikan $x_k\to s$. Untuk itu diperlukan alat Bab 4.
Yang bisa dikatakan sekarang: karena $e_1<0.061$, dari $e_{k+1}\le e_k^2/2$ dapat
diturunkan $e_k\le 2^{-(2^k-1)}$ untuk $k\ge 1$, dan itu cukup untuk merancang uji
numerik yang jujur.

\section{Interval, titik akumulasi, dan himpunan terbuka}
\tingkat{lanjut}

Dengan supremum di tangan, kita bisa memberi nama pada potongan garis bilangan, seperti
aturan finis lomba lari: garisnya sama, yang berbeda hanya apakah garis itu ikut
dihitung. Nama-nama itu dipakai lagi di Bab 11 dan Bab 13, jadi kita ambil yang perlu
saja tanpa membangun topologi secara penuh.

Interval adalah potongan garis yang tidak berlubang. Definisi berikut menuliskan
sembilan bentuk yang mungkin; semuanya ditentukan oleh dua ujung dan oleh keputusan
apakah ujung itu ikut dihitung.

\begin{definisi}[Interval]
Misalkan $a<b$ bilangan real. Himpunan $(a,b)=\{x: a<x<b\}$ disebut interval terbuka,
$[a,b]=\{x: a\le x\le b\}$ interval tertutup, dan $[a,b)$, $(a,b]$ adalah bentuk
campuran. Himpunan $(a,\infty)=\{x: x>a\}$, $[a,\infty)$, $(-\infty,b)$,
$(-\infty,b]$, dan $\mathbb{R}$ sendiri disebut interval tak terbatas. Kesembilan
bentuk itu disebut interval.
\end{definisi}

Ciri yang menandai semuanya: jika dua titik ada di dalam interval, seluruh ruas di
antaranya juga di dalam, sehingga interval disebut tidak berlubang. Daftar tadi tidak
dibuktikan lengkap; klasifikasinya hanya memeriksa batas atas, batas bawah, dan apakah
batas itu termasuk. Titik ``dekat'' dinyatakan dengan persekitaran, yaitu interval
terbuka kecil di sekitar satu titik.

\begin{definisi}[Persekitaran, titik dalam, terbuka, tertutup]\index{himpunan terbuka}\index{himpunan tertutup}\index{titik akumulasi}
Untuk $\varepsilon>0$, persekitaran titik $x$ adalah
$V_\varepsilon(x)=(x-\varepsilon,x+\varepsilon)$. Titik $x$ disebut titik dalam
$S\subseteq\mathbb{R}$ jika $V_\varepsilon(x)\subseteq S$ untuk suatu $\varepsilon>0$;
$S$ disebut terbuka jika setiap anggotanya titik dalam. Titik $x$ disebut titik
akumulasi $S$ jika setiap $V_\varepsilon(x)$ memuat anggota $S$ yang berbeda dari $x$.
Himpunan $S$ disebut tertutup jika komplemennya terbuka.
\end{definisi}

Beberapa contoh menata intuisi. Interval $(0,1)$ terbuka dan tidak tertutup, sebab
setiap persekitaran $0$ memotongnya. Interval $[0,1]$ tertutup dan tidak terbuka,
sedangkan $[0,1)$ bukan keduanya. Himpunan $\mathbb{Q}$ juga bukan keduanya, sebab
setiap persekitaran memuat rasional dan irasional. Himpunan $\{1/n:n\in\mathbb{N}\}$
tidak tertutup: $0$ ada di komplemennya, tetapi setiap persekitaran $0$ memuat $1/n$.
Himpunan $\mathbb{Z}$ tertutup, sebab komplemennya gabungan interval terbuka $(k,k+1)$.

Jadi ``tertutup'' bukan negasi dari ``terbuka'': $\mathbb{R}$ dan $\varnothing$
keduanya terbuka sekaligus tertutup, dan komplemen himpunan terbuka selalu tertutup.
Kerapatan $\mathbb{Q}$ membuat contoh itu lebih tajam. Setiap persekitaran memuat
rasional dan irasional, jadi tidak ada titik yang menjadi titik dalam $\mathbb{Q}$
maupun komplemennya: himpunan itu padat sekaligus berlubang di mana-mana. Karena itu
kelengkapan tidak bisa diganti dengan kerapatan.

Persekitaran sendiri terbuka, dan itu membuat definisi titik dalam bisa dipakai
berulang.

\begin{proposisi}
(i) Setiap persekitaran $V_\varepsilon(x)$ terbuka. (ii) Gabungan sebarang keluarga
himpunan terbuka adalah terbuka. (iii) Irisan berhingga himpunan terbuka adalah
terbuka.
\end{proposisi}

\begin{proof}
(i) Ambil $y\in V_\varepsilon(x)$, jadi $|y-x|<\varepsilon$. Tetapkan
$\delta=\varepsilon-|y-x|>0$. Untuk setiap $z\in V_\delta(y)$,
$|z-x|\le|z-y|+|y-x|<\delta+|y-x|=\varepsilon$ menurut ketaksamaan segitiga. Jadi
$V_\delta(y)\subseteq V_\varepsilon(x)$.

(ii) Jika $x$ berada di gabungan, ia berada di salah satu anggotanya, dan karena
anggota itu terbuka ada persekitaran $x$ yang termuat di dalamnya, sehingga juga
termuat di dalam gabungan.

(iii) Jika $x$ berada di irisan berhingga $S_1\cap\cdots\cap S_n$, untuk setiap $i$ ada
$\varepsilon_i>0$ dengan $V_{\varepsilon_i}(x)\subseteq S_i$. Ambil
$\varepsilon=\min\{\varepsilon_1,\ldots,\varepsilon_n\}>0$. Maka
$V_\varepsilon(x)$ termuat di setiap $S_i$, jadi termuat di irisannya.
\end{proof}

Butir (iii) hanya untuk irisan berhingga: $\bigcap_n(-1/n,1/n)=\{0\}$ tidak terbuka,
sebab untuk $x\neq0$ Sifat Archimedes memberi $n$ dengan $1/n<|x|$.

\section{Contoh hitung}
\tingkat{menengah}

Tiga perhitungan yang dijanjikan dijalankan di sini: iterasi Newton, aproksimasi
rasional sampai galat di bawah $10^{-6}$, dan pemeriksaan bahwa $\{q:q^2<2\}$ tidak
punya maksimum. Semua angkanya dihitung dengan pecahan eksak.

Mulai dari $x_0=1$, iterasi $x_{k+1}=(x_k+2/x_k)/2$ memberi pecahan berikut; penyebutnya
tumbuh, dan itu harga yang dibayar untuk ketepatan.

\begin{center}\small
\begin{tabular}{rlll}
\hline
$k$ & $x_k$ & desimal & galat relatif \\
\hline
0 & $1$ & $1$ & $2.92893\times10^{-1}$ \\
1 & $3/2$ & $1.5$ & $6.06602\times10^{-2}$ \\
2 & $17/12$ & $1.41666666667$ & $1.73461\times10^{-3}$ \\
3 & $577/408$ & $1.41421568627$ & $1.50183\times10^{-6}$ \\
4 & $665857/470832$ & $1.41421356237$ & $1.12774\times10^{-12}$ \\
5 & $886731088897/627013566048$ & $1.41421356237$ & $6.35896\times10^{-25}$ \\
\hline
\end{tabular}
\end{center}

Galat relatif dihitung terhadap $\sqrt{2}=1.414213562373095\ldots$. Faktor
penurunannya naik sendiri: dari langkah 0 ke 1 sekitar $4.8$, lalu $35$, $1155$, dan
$1.3\times10^{6}$ --- itulah tanda pengkuadratan. Nisbah $e_{k+1}/e_k^2$ bernilai
$0.7071$, $0.4714$, $0.4991$, $0.49999925$, dan $0.5$, sesuai rumus
$e_{k+1}=e_k^2\cdot\frac{s}{2x_k}$ karena $x_k$ mendekat ke $s$ dari atas. Pada langkah
4 dan 5 desimalnya sama sampai sebelas angka; angka berikutnya tidak muncul karena kita
mencetak dua belas angka penting, bukan karena galatnya berhenti turun. Pecahan langkah
5 punya dua belas digit dengan galat $6.36\times10^{-25}$, dan satu angka tambahan hanya
menuntut kira-kira satu langkah lagi.

Cara kedua memakai aproksimasi rasional yang lebih hemat penyebut, yaitu pecahan lanjar
$\sqrt2$:
\[
1,\quad \frac32,\quad \frac75,\quad \frac{17}{12},\quad \frac{41}{29},\quad \ldots
\]
dengan $p_{k+1}=p_k+2q_k$ dan $q_{k+1}=p_k+q_k$ mulai dari $(p_0,q_0)=(1,1)$. Tabel
berikut memuat galat mutlaknya.

\begin{center}\small
\begin{tabular}{rlll}
\hline
$p/q$ & desimal & galat mutlak & $p^2-2q^2$ \\
\hline
$1/1$ & $1$ & $4.14214\times10^{-1}$ & $-1$ \\
$3/2$ & $1.5$ & $8.57864\times10^{-2}$ & $1$ \\
$7/5$ & $1.4$ & $1.42136\times10^{-2}$ & $-1$ \\
$17/12$ & $1.41666666667$ & $2.45310\times10^{-3}$ & $1$ \\
$41/29$ & $1.41379310345$ & $4.20459\times10^{-4}$ & $-1$ \\
$99/70$ & $1.41428571429$ & $7.21519\times10^{-5}$ & $1$ \\
$239/169$ & $1.41420118343$ & $1.23789\times10^{-5}$ & $-1$ \\
$577/408$ & $1.41421568627$ & $2.12390\times10^{-6}$ & $1$ \\
$1393/985$ & $1.41421319797$ & $3.64404\times10^{-7}$ & $-1$ \\
$3363/2378$ & $1.41421362489$ & $6.25218\times10^{-8}$ & $1$ \\
\hline
\end{tabular}
\end{center}

Aproksimasi pertama yang galatnya di bawah $10^{-6}$ adalah $1393/985$ dengan galat
$3.64\times10^{-7}$; angka sebelumnya, $577/408$, masih $2.12\times10^{-6}$. Kolom
terakhir menjelaskan sebabnya: karena $p^2-2q^2=\pm1$,
\[
\left|\frac pq-\sqrt2\right|
=\frac{|p^2-2q^2|}{q^2\left(\frac pq+\sqrt2\right)}
<\frac{1}{2\sqrt2\,q^2},
\]
sehingga galat turun seperti kebalikan kuadrat penyebut. Untuk $q=985$ batas itu
$3.644\times10^{-7}$ dan galat sebenarnya $3.64404\times10^{-7}$: hampir tepat, karena
$\frac pq$ sudah sangat dekat ke $\sqrt2$. Penyebut 985 jauh lebih kecil daripada
$627013566048$ pada langkah 5 Newton, meskipun galatnya hanya sekitar sejuta kali lebih
besar --- itu beda metode yang mengejar galat kuadrat dan yang mengejar galat rasional
terkecil.

Gambar~\ref{fig:tangga-akar} menampilkan kolom galat itu sebagai tangga menurun. Tanda
$p^2-2q^2$ bergantian, sehingga pecahannya melompat dari bawah ke atas akar dua, dan
nisbah dua galat berurutan menuju $5{,}8284=(1+\sqrt2)^2$.

\begin{figure}[htbp]
\centering
\begin{tikzpicture}[x=1.02cm,y=1cm,font=\small]
  \foreach \lv in {1,...,7}{
    \draw[layergrey!20] (0.55,\lv) -- (10.75,\lv);
    \node[layergrey,anchor=east,font=\scriptsize] at (0.45,\lv) {$10^{-\lv}$};
  }
  \node[layergrey,anchor=east,font=\scriptsize] at (0.45,0) {$10^{0}$};
  \draw[-{Latex[length=2mm]},ink,thick] (0.55,0) -- (10.95,0);
  \draw[-{Latex[length=2mm]},ink,thick] (0.55,0) -- (0.55,7.85);
  \foreach \k/\lab in {%
    1/{$1/1$},2/{$3/2$},3/{$7/5$},4/{$17/12$},5/{$41/29$},
    6/{$99/70$},7/{$239/169$},8/{$577/408$},9/{$1393/985$},10/{$3363/2378$}}{
    \draw[ink] (\k,0.07) -- (\k,-0.07);
    \node[rotate=32,anchor=north east,font=\scriptsize,layergrey] at (\k,-0.12) {\lab};
  }
  \foreach \k/\v in {1/0.3828,3/1.8473,5/3.3763,7/4.9073,9/6.4384}{
    \fill[accent] (\k,\v) circle (1.8pt);
  }
  \foreach \k/\v in {2/1.0666,4/2.6103,6/4.1418,8/5.6729,10/7.2040}{
    \fill[layerteal] (\k,\v) circle (1.8pt);
  }
  \draw[layergrey,dashed,thick] (0.55,6) -- (10.75,6);
  \node[layergrey,font=\scriptsize,anchor=south east] at (10.7,6.02) {ambang $10^{-6}$};
  \fill[accent] (0.95,7.55) circle (1.8pt);
  \node[right,font=\scriptsize] at (1.05,7.55) {di bawah $\sqrt2$};
  \fill[layerteal] (4.35,7.55) circle (1.8pt);
  \node[right,font=\scriptsize] at (4.45,7.55) {di atas $\sqrt2$};
  \node[font=\scriptsize,anchor=north] at (5.7,-1.18) {pecahan berurutan};
  \node[font=\scriptsize,rotate=90,anchor=south] at (-0.35,3.9) {$|p/q-\sqrt2|$};
\end{tikzpicture}
\caption{Tangga aproksimasi rasional: galat $|p/q-\sqrt2|$ bagi sepuluh pecahan
berurutan, skala logaritmik. Ambang $10^{-6}$ pertama dilewati $1393/985$, dengan galat
$3{,}644\times10^{-7}$; sesudahnya tiap anak tangga turun dengan faktor $5{,}8284$.}
\label{fig:tangga-akar}
\end{figure}

Pemeriksaan ketiga menyangkut ketiadaan maksimum. Satu langkah lagi dari $q$ sebelumnya
memberi
\[
q'=\frac{149856914148961}{105964844887080}=1.41421349985133266\ldots,
\]
dan kuadratnya masih di bawah dua. Kedua bilangan itu rasional, keduanya di dalam $A$,
dan masing-masing lebih besar dari yang sebelumnya: tidak ada alasan proses ini berhenti.

\begin{tantangan}
Tinjau $S=\{q\in\mathbb{Q}: q\ge 0,\ q^2\le 2\}$. Sekilas himpunan ini punya anggota
terbesar: bilangan yang kuadratnya tepat dua, sebab tanda $\le$ seolah memasukkan batas
itu ke dalam himpunan. Buktikan bahwa $S$ tidak punya maksimum, lalu tentukan $\sup S$
di $\mathbb{R}$ dan jelaskan mengapa jawabannya bukan anggota $S$.

Petunjuk 1: anggota $S$ harus rasional. Periksa dulu apakah calon terbesarmu boleh
menjadi anggota himpunan.

Petunjuk 2: untuk setiap $q\in S$ dengan $q>0$ berlaku $q^2<2$ secara tegas, sebab
$q^2=2$ mustahil bagi rasional. Jadi lemma geseran $h=(2-q^2)/(2q+2)$ bekerja dan
menghasilkan anggota $S$ yang lebih besar; batas yang tak pernah tercapai itu $\sqrt2$.

Kalau mentok, kembali ke bagian ``Mengapa bilangan rasional belum cukup'' dan ke lemma
tentang geseran.
\end{tantangan}

Terakhir, Sifat Archimedes pada angka. Untuk $\varepsilon=10^{-6}$ kita perlu
$n>10^{6}$. Pada $n=10^{6}$ resiprokalnya tepat $10^{-6}$, yang tidak lebih kecil dari
$\varepsilon$; pada $n=10^{6}+1$ resiprokalnya $9.99999\times10^{-7}$, baru memenuhi.
Perbedaan satu pada $n$ menentukan sah tidaknya pilihan, dan itulah sebabnya
pertidaksamaan tegas harus dihormati.

\kertas{Untuk $x=1393/985$, hitung $h=(2-x^2)/(2x+2)$ dan periksa dengan tangan bahwa
$(x+h)^2<2$. Lalu hitung galat relatif $|x-\sqrt2|/\sqrt2$ dan bandingkan dengan
$1/(2\sqrt2\,x^2)$. Jelaskan mengapa keduanya hampir sama, dan mengapa fakta itu tidak
bertentangan dengan ketiadaan maksimum pada $A$.}

\section{Lab komputasi}
\tingkat{lanjut}

Kode berikut dijalankan dengan Python 3 dan \texttt{mpmath} untuk aritmetika presisi
tinggi; iterasi Newton memakai \texttt{Fraction}, jadi tidak ada pembulatan sebelum
pencetakan.

\begin{lab}
\begin{verbatim}
from fractions import Fraction as F
import mpmath as mp
mp.mp.dps = 40
akar2 = mp.sqrt(2)

print("Newton untuk akar dua, x0 = 1")
x = F(1)
for k in range(6):
    v = mp.mpf(x.numerator) / x.denominator
    print(k, f"{x.numerator}/{x.denominator}", mp.nstr(v, 12),
          mp.nstr(abs(v - akar2) / akar2, 6))
    x = (x + F(2) / x) / 2

print()
print("Aproksimasi rasional berurutan, galat < 1e-6 ?")
p, q = 1, 1
for k in range(10):
    v = mp.mpf(p) / q
    print(f"{p}/{q}", mp.nstr(v, 12), mp.nstr(abs(v - akar2), 6),
          abs(v - akar2) < mp.mpf(10) ** -6)
    p, q = p + 2 * q, p + q

print()
print("Uji: {q rasional : q^2 < 2} tidak punya maksimum")
q = F(1393, 985)
for k in range(2):
    h = (2 - q * q) / (2 * q + 2)
    q = q + h
    print(q, mp.nstr(mp.mpf(q.numerator) / q.denominator, 18), q * q < 2)
\end{verbatim}
\end{lab}

Keluaran yang diperoleh, apa adanya:

\begin{verbatim}
Newton untuk akar dua, x0 = 1
0 1/1 1.0 0.292893
1 3/2 1.5 0.0606602
2 17/12 1.41666666667 0.00173461
3 577/408 1.41421568627 1.50183e-6
4 665857/470832 1.41421356237 1.12774e-12
5 886731088897/627013566048 1.41421356237 6.35896e-25

Aproksimasi rasional berurutan, galat < 1e-6 ?
1/1 1.0 0.414214 False
3/2 1.5 0.0857864 False
7/5 1.4 0.0142136 False
17/12 1.41666666667 0.0024531 False
41/29 1.41379310345 0.000420459 False
99/70 1.41428571429 7.21519e-5 False
239/169 1.41420118343 1.23789e-5 False
577/408 1.41421568627 2.1239e-6 False
1393/985 1.41421319797 3.64404e-7 True
3363/2378 1.41421362489 6.25218e-8 True

Uji: {q rasional : q^2 < 2} tidak punya maksimum
6625109/4684660 1.41421341143220639 True
149856914148961/105964844887080 1.41421349985133266 True
\end{verbatim}

Tiga hal terbaca dari keluaran itu. Galat Newton berhenti tercetak pada dua belas angka
di langkah 4, walaupun nilai sebenarnya masih turun sampai $6.4\times10^{-25}$: yang
membatasi lebar cetakan, bukan matematikanya. Ambang $10^{-6}$ dilewati tepat pada
$1393/985$, tempat kolom \texttt{True} mulai muncul. Kedua bilangan pada uji maksimum
juga \texttt{True}: keduanya rasional dan kuadratnya masih di bawah dua.

Kode itu tidak membuktikan apa pun. Ia memeriksa bahwa rumus di bab ini memberi angka
yang disebutkan dan bahwa arah pertidaksamaannya benar; buktinya tetap di lemma dan
teorema. Keluaran program adalah data, bukan argumen.

\begin{untukdosen}
Urutan mengajar yang saya sarankan: program pencarian akar di pembuka, lalu aksioma
lapangan dengan $0\cdot x=0$ dibuktikan di papan, lalu urutan, nilai mutlak, dan
kelengkapan. Jangan menunda kelengkapan; seluruh isi bab bergantung padanya.

Dua miskonsepsi yang sering muncul. Pertama, mahasiswa mengira supremum harus anggota
himpunan, sehingga mereka menolak $\sup\{1-1/n\}=1$. Kedua, mahasiswa mengira
kelengkapan bisa dibuktikan dari hukum lapangan dan urutan, padahal $\mathbb{Q}$
memenuhi keduanya dan gagal pada kelengkapan. Keduanya hilang setelah mahasiswa melihat
dua contoh itu berdampingan.

Aktivitas papan tulis yang berhasil: setiap mahasiswa mengambil $q=1.4$, menghitung
$h=(2-q^2)/(2q+2)$ dan memeriksa $(q+h)^2<2$; hasil yang berbeda-beda justru menunjukkan
rumusnya bekerja untuk setiap $q$. Tanda siap lanjut: mahasiswa dapat menjelaskan mengapa
$A$ tidak punya maksimum tanpa menghitung.
\end{untukdosen}

\section{Ringkasan}
\tingkat{dasar}

\begin{ringkasan}
\begin{itemize}
\item Bilangan real dilandasi tiga kelompok aksioma: lapangan, urutan, kelengkapan.
      Yang pertama memberi aljabar, yang kedua perbandingan, yang ketiga menutup lubang.
\item Dari lapangan diturunkan $0\cdot x=0$, $(-x)y=-(xy)$, dan ketiadaan pembagi nol.
\item Dari urutan diturunkan aturan tanda, termasuk kuadrat tak negatif dan $1>0$.
\item Nilai mutlak mengubah jarak menjadi pertidaksamaan: $|x|\le c$ setara dengan
      $-c\le x\le c$, dan $|x+y|\le|x|+|y|$.
\item Kelengkapan menjamin supremum setiap himpunan tak kosong yang terbatas di atas,
      lalu menurunkan Sifat Archimedes dan kerapatan $\mathbb{Q}$.
\item $\mathbb{Q}$ memenuhi lapangan dan urutan tetapi tidak lengkap: $\{q: q^2<2\}$
      tidak punya maksimum; di $\mathbb{R}$ supremumnya $\sqrt{2}$.
\end{itemize}
\end{ringkasan}

\section{Latihan}
\tingkat{menengah}

\begin{latihan}
\begin{enumerate}
  \item Buktikan bahwa dalam lapangan terurut, $a<b$ setara dengan $-b<-a$, dan
        simpulkan bahwa mengalikan pertidaksamaan dengan $-1$ membalik arahnya.
  \item Buktikan $a<\frac{a+b}{2}<b$ untuk $a<b$, lalu jelaskan mengapa hasil itu tidak
        memerlukan kelengkapan sedangkan $\sup\{q:q^2<2\}$ memerlukannya.
  \item Hitung $\sup$ dan $\inf$ himpunan $S=\{2-1/n : n\in\mathbb{N}\}$, lalu periksa
        kedua syarat definisi. Tentukan juga apakah $S$ punya maksimum.
  \item Nilai klaim ini: ``Karena setiap bilangan positif lebih besar dari $1/n$ untuk
        suatu $n$, bilangan real positif terkecil tidak ada.'' Benar, salah, atau belum
        cukup bukti? Sebutkan bagian yang didukung bukti.
  \item Dengan $h=(2-q^2)/(2q+2)$, hitung $q+h$ untuk $q=1$ dan $q=7/5$, lalu periksa
        bahwa keduanya masih di dalam $A=\{q\in\mathbb{Q}:q^2<2\}$. Bandingkan besar
        kenaikannya.
  \item Tunjukkan bahwa jika $|x|\le 2$ dan $|y|\le 2$ maka
        $|x^3-y^3|\le 12\,|x-y|$, dan sebutkan di mana ketaksamaan segitiga dipakai.
  \item Buktikan bahwa jika $S$ terbuka maka komplemennya tertutup, dan sebaliknya,
        dengan memakai definisi persekitaran.
  \item Hitung lima langkah iterasi Newton untuk $\sqrt{3}$ mulai dari $x_0=2$, dan
        laporkan galat relatif tiap langkah. Periksa apakah nisbah $e_{k+1}/e_k^2$
        mendekat ke $1/2$.
  \item Tunjukkan $\bigcap_{n=1}^{\infty}(-1/n,1/n)=\{0\}$ dan jelaskan mengapa itu tidak
        bertentangan dengan ketertutupan irisan berhingga.
  \item Buktikan bahwa di antara dua bilangan real selalu ada bilangan irasional.
\end{enumerate}

\noindent\textbf{Petunjuk.} Butir 4: klaim itu memuat dua pernyataan berbeda, satu
tentang setiap $\varepsilon$ dan satu tentang keberadaan minimum. Butir 8: gunakan
$e_{k+1}=e_k^2\cdot\frac{\sqrt3}{2x_k}$ dan $x_k\ge\sqrt3$. Butir 10: pakai kerapatan
$\mathbb{Q}$ pada $x/\sqrt2$ dan $y/\sqrt2$, lalu kalikan kembali.
\end{latihan}

\begin{untukpeneliti}
Yang benar-benar terbukti: aksioma kelengkapan tidak bergantung pada lapangan dan urutan,
sebab $\mathbb{Q}$ memenuhi keduanya dan gagal pada yang ketiga. Keberadaan lapangan
terurut lengkap dibuktikan Rudin lewat potongan Dedekind, dan di sana juga dinyatakan
tanpa bukti bahwa setiap dua lapangan terurut lengkap isomorfik; klaim itu tidak dipakai
di sini.

Sumber primer: Rudin, \emph{Principles of Mathematical Analysis}, Bab 1 (Definisi 1.12,
Proposisi 1.18, Teorema 1.19--1.21, Apendiks tentang potongan); Abbott,
\emph{Understanding Analysis}, Bab 1 (Aksioma Kelengkapan dan akibatnya).

Yang masih terbuka: isi efektif iterasi Newton. Tanpa batas awal eksplisit,
$e_{k+1}\le e_k^2/2$ belum memberi jumlah langkah; itu memerlukan Bab 4 dan, untuk
fungsi umum, teori Kantorovich. Batas keberlakuan bab: topologi sebatas interval dan
persekitaran, dan $\mathbb{R}$ diandaikan, bukan dibangun.
\end{untukpeneliti}

\peta{Bab 3 pada Rudin, \emph{Principles of Mathematical Analysis}: Bab 1 (Ordered Sets,
Fields, The Real Field, The Least-Upper-Bound Property, Sifat Archimedes, kerapatan
$\mathbb{Q}$, keberadaan akar ke-$n$, Apendiks potongan Dedekind). Pada Abbott,
\emph{Understanding Analysis}: Bab 1 (The Axiom of Completeness dan Consequences of
Completeness); himpunan terbuka dan titik akumulasi dari Bab 3, tidak dikembangkan di
sini.}

\chapter{Barisan dan Limit}

Sebuah program kecil mencetak satu bilangan per baris. Aturannya hanya satu: suku
berikutnya adalah akar dari dua ditambah suku sebelumnya, dimulai dari satu. Yang
muncul di layar: 1.000000, lalu 1.732051, 1.931852, 1.982890, 1.995718, 1.998929,
dan seterusnya. Pada baris kedua belas layar menunjukkan 1.999999738545. Angka itu
tampak berhenti bergerak. Tetapi tampak dan benar adalah dua hal berbeda: program
tidak tahu apakah baris keseribu masih akan bergeser, dan ia juga tidak tahu ke
bilangan mana deretan ini menuju.

Contoh kedua lebih nakal. Barisan $(1+1/n)^n$ untuk $n=10$ bernilai 2.5937424601,
untuk $n=100$ bernilai 2.7048138294, untuk $n=1000$ bernilai 2.7169239322, dan
untuk $n=10^7$ bernilai 2.7182816941. Enam angka pertama tampak mantap pada 2.71828.
Tetapi nilai sebenarnya pada $n=10^7$ adalah 2.7182816925, sehingga tujuh digit
terakhir pada keluaran mesin itu adalah galat pembulatan, bukan tanda kedekatan.
Tabel dapat menipu ke dua arah sekaligus: membuat barisan yang divergen tampak
tenang, dan membuat barisan yang konvergen tampak melenceng. Bab 22 nanti
menjelaskan mengapa aritmetika mesin berperilaku begitu; bab ini menyiapkan alat
yang membuat kita tidak perlu bergantung pada tabel.

Bab 3 sudah memberi bahan bakunya: aksioma lapangan, aksioma urutan, dan sifat
supremum. Tetapi supremum hanya bekerja setelah himpunan sukunya terkumpul; ia
tidak memberi tahu apakah barisan berhenti bergerak. Yang kita perlukan adalah
definisi yang bisa diperiksa dengan kertas, lalu teorema yang mengubah pemeriksaan
itu menjadi pekerjaan terkira: definisi $\varepsilon$-$N$, sifat aljabar limit,
teorema apit, barisan monoton, subbarisan, kriteria Cauchy, dan limit superior.

\begin{untukanak}
Pernahkah kalian berjalan menuju pintu dengan aturan aneh: setiap langkah hanya
menempuh separuh dari jarak yang tersisa? Langkah pertama membawa kalian separuh
jalan. Langkah kedua memotong separuh dari sisa itu, jadi kalian sudah di tiga
perempat jalan. Langkah ketiga menambah seperdelapan lagi. Setiap langkah pasti
maju, tetapi selalu ada sisa yang belum terjangkau. Kalau aturan itu dijalankan
terus, kalian tidak pernah benar-benar menyentuh pintu.

Meski begitu, ada pertanyaan yang bisa dijawab dengan pasti. Berapa langkah yang
diperlukan supaya jarak ke pintu kurang dari sejengkal? Jawabannya bisa dihitung
sebelum melangkah, dan tidak berubah karena ruangan menjadi dua kali lebih panjang:
kalian hanya perlu beberapa langkah tambahan. Jadi ada dua hal berbeda di sini.
"Tidak pernah sampai" dan "sedekat apa pun bisa didekati" bukan dua pernyataan
yang bertentangan. Justru keduanya benar sekaligus.

Sekarang pertanyaan yang membuat penasaran. Bagaimana kalau aturannya diubah
sedikit, misalnya setiap langkah menempuh sepertiga sisa jarak, atau separuh
ditambah sedikit? Apakah kalian masih bisa menebak di mana perjalanan itu berakhir,
dan apakah kecepatannya masih bisa dihitung? Di bab ini kita akan tahu mengapa
pertanyaan semacam itu punya jawaban yang pasti, dan mengapa jawabannya tidak
pernah kita peroleh hanya dengan menonton angka di layar.
\end{untukanak}

\section{Definisi limit dan cara memakainya}
\tingkat{dasar}

\kamus{barisan}{sequence}{deretan bilangan yang diberi nomor urut, seperti nomor
rumah di satu sisi jalan yang tidak pernah habis}

Limit menjawab satu pertanyaan: seberapa dekat suku-suku akhir berkumpul, dan mulai
kapan. Dalam bahasa sehari-hari, barisan menuju $L$ bila setelah beberapa suku awal
dibuang, semua sisa sukunya sedekat apa pun yang diminta kepada $L$.

\begin{definisi}[Kekonvergenan]
Barisan $(x_n)$ bilangan real \emph{konvergen} ke $L \in \mathbb{R}$ jika untuk
setiap $\varepsilon>0$ terdapat $N \in \mathbb{N}$ sehingga untuk setiap $n>N$
berlaku $|x_n-L|<\varepsilon$. Kita tulis $x_n \to L$, dan $L$ disebut
\emph{limit}\index{limit barisan} barisan itu. Barisan yang tidak konvergen ke
bilangan real mana pun disebut \emph{divergen}\index{barisan divergen}.
\end{definisi}

Urutan kuantor pada definisi itu menentukan segalanya. Nilai $\varepsilon$ dipilih
lebih dahulu, lalu $N$ menyusul dan boleh bergantung pada $\varepsilon$. Membalik
urutan itu menghasilkan pernyataan yang jauh lebih kuat dan biasanya salah: ada
barisan yang memenuhi "untuk setiap $n$ ada $N$" secara hampa, karena $N=n$ selalu
bisa dipakai. Cara membaca yang sehat adalah sebagai permainan dua pihak. Penantang
menyebut satu $\varepsilon$, sekecil apa pun. Penjawab harus membalas dengan satu
$N$ yang berlaku untuk semua suku setelahnya. Barisan konvergen bila penjawab
selalu punya balasan, untuk setiap tantangan.

Gambar~\ref{fig:pita-epsilon} memotret permainan itu untuk $x_n=1+1/n$. Dua pita
berpusat di $L=1$ digambar dengan lebar berbeda, dan masing-masing punya ambangnya
sendiri: $N=2$ untuk $\varepsilon=0.5$, $N=10$ untuk $\varepsilon=0.1$. Tidak ada
yang berubah pada barisannya; yang berubah hanya permintaan penantang.

\begin{figure}[htbp]
\centering
\begin{tikzpicture}[x=0.72cm,y=3.0cm,>=Stealth,font=\footnotesize]
  \fill[layerteal!14] (0,0.5) rectangle (14,1.5);
  \fill[layerorange!20] (0,0.9) rectangle (14,1.1);
  \draw[->,ink] (-0.1,0.35) -- (14.6,0.35) node[right] {$n$};
  \draw[->,ink] (0,0.35) -- (0,2.15) node[above] {$x_n$};
  \draw[dashed,layergrey,thick] (0,1) -- (14.1,1);
  \node[layergrey,anchor=south west] at (0.08,1.02) {$L=1$};
  \foreach \n/\y in {1/2.0000,2/1.5000,3/1.3333,4/1.2500,5/1.2000,6/1.1667,%
      7/1.1429,8/1.1250,9/1.1111,10/1.1000,11/1.0909,12/1.0833,13/1.0769,14/1.0714}
    \fill[accent] (\n,\y) circle (2.2pt);
  \draw[dashed,layerteal,thick] (2,0.35) -- (2,1.52);
  \draw[dashed,layerorange,thick] (10,0.35) -- (10,1.12);
  \node[anchor=north] at (2,0.30) {$N_1=2$};
  \node[anchor=north] at (10,0.30) {$N_2=10$};
  \node[anchor=north west,align=left] at (4.6,2.28)
    {pita luar: $\varepsilon=0.5$\\ pita dalam: $\varepsilon=0.1$};
\end{tikzpicture}
\caption{Barisan $x_n=1+1/n$ masuk ke pita $\varepsilon$ di sekitar $L=1$. Dengan
$\varepsilon=0.5$ semua suku sesudah $n=2$ sudah di dalam pita; dengan
$\varepsilon=0.1$ ambangnya bergeser ke $n=10$. Pita yang lebih sempit menuntut
ambang yang lebih jauh.}
\label{fig:pita-epsilon}
\end{figure}

\begin{bacakeras}
Dua suara bergantian. Suara pertama bertanya, suara kedua menjawab. Bacakan pelan,
beri jeda pada setiap titik, dan jangan menambah kata di antaranya.

``Semakin lama?'' \ ``Semakin dekat.'' \ ``Seberapa dekat?'' \ ``Sedekat yang
diminta.'' \ ``Kalau yang diminta seperseribu?'' \ ``Tunggu seribu langkah, lalu
ukur.'' \ ``Kalau seperseratus juta?'' \ ``Tunggu seratus juta langkah. Sabar tidak
apa-apa.'' \ ``Jadi akan sampai?'' \ ``Tidak. Tidak pernah sampai.'' \ ``Lalu apa
gunanya menunggu?'' \ ``Sebab jaraknya mengecil, dan yang meminta bebas memilih
seberapa kecil.'' \ ``Dan kalau permintaannya berhenti?'' \ ``Maka tidak ada yang
perlu dibuktikan.''

Baris yang diulang tiga kali, makin pelan setiap kali: dekat bukan sampai; dekat
adalah jarak yang boleh kita tentukan sendiri.
\end{bacakeras}

Untuk $x_n=1/n$, balasannya langsung: ambil $N$ bilangan asli dengan $N>1/\varepsilon$,
yang ada menurut sifat Archimedes di Bab 3. Untuk $n>N$ berlaku
$|1/n-0|=1/n<1/N<\varepsilon$. Bila penantang meminta $\varepsilon=10^{-3}$,
jawabannya $N=1000$; pada $n=1000$ galatnya tepat $10^{-3}$ dan belum memenuhi
syarat ketat, sedangkan pada $n=1001$ galatnya $0.000999$. Perbedaan satu suku ini
bukan keisengan: definisi memakai $n>N$, bukan $n \ge N$, sehingga $N$ yang dipilih
harus benar-benar melewati ambang.

\begin{tangan}
Bahan: lakban, meteran, kapur, dan satu lembar catatan. Bentuk garis bilangan di
lantai sepanjang dua meter: nol di ujung kiri, dua di ujung kanan, dan angka satu
tepat di tengah. Tempel dua pita mendatar selebar telapak kaki, berpusat di angka
satu: pita lebar dari setengah sampai satu setengah meter, pita sempit dari sembilan
persepuluh sampai satu koma satu meter.

Sekarang berjalanlah ke posisi satu ditambah satu per n, untuk n sama dengan satu,
dua, tiga, sampai empat belas. Setiap kali, catat apakah tumit masuk pita lebar,
pita sempit, atau di luar keduanya, lalu berjalan mundur supaya tanda tidak
terinjak.

Yang harus terlihat: dua langkah pertama masih di luar pita sempit, langkah
kesepuluh berhenti tepat di tepinya, dan sesudah itu tidak ada lagi tumit yang
keluar. Pita yang lebih sempit menuntut langkah yang lebih banyak, sedangkan angka
satu tidak pernah tersentuh tumit.
\end{tangan}

Kegunaan definisi justru paling terasa ketika kita ingin membuktikan bahwa sesuatu
\emph{tidak} konvergen. Negasi definisi berbunyi: ada $\varepsilon>0$ sehingga untuk
setiap $N$ terdapat $n>N$ dengan $|x_n-L| \ge \varepsilon$. Perhatikan bahwa $L$
masih nama bebas di situ; untuk menyangkal kekonvergenan kita harus melakukannya
untuk setiap $L$.

\begin{contoh}
Barisan $x_n=(-1)^n$ divergen.
\end{contoh}

\begin{proof}[Bukti]
Andaikan $x_n \to L$ untuk suatu $L$. Ambil $\varepsilon=1$; ada $N$ sehingga
$|x_n-L|<1$ untuk semua $n>N$. Pilih $n>N$ yang genap; maka $x_n=1$ dan
$|1-L|<1$. Suku berikutnya $n+1$ juga melebihi $N$ dan ganjil, jadi $x_{n+1}=-1$
dan $|-1-L|<1$. Ketaksamaan segitiga memberi
\[
2 = |1-(-1)| \le |1-L| + |L-(-1)| < 1+1 = 2,
\]
suatu kemustahilan. Jadi tidak ada $L$ yang memenuhi definisi.
\end{proof}

Bukti tadi memperlihatkan pola yang akan sering muncul: untuk menyangkal, kita
memilih $\varepsilon$ yang lebih kecil daripada jarak antar dua kelompok suku.
Kelompok genap dan ganjil berjarak tetap $2$, sehingga tidak ada satu $N$ pun yang
dapat memisahkan keduanya bila $\varepsilon<1$.

Sifat pertama yang patut dicatat adalah ketunggalan. Kalimat sehari-harinya: sebuah
barisan tidak dapat berhenti di dua tempat berbeda sekaligus.

\begin{teorema}[Ketunggalan limit]
Jika $x_n \to L$ dan $x_n \to M$, maka $L=M$.
\end{teorema}

\begin{proof}[Bukti]
Andaikan $L \neq M$ dan misalkan $\varepsilon=|L-M|/2>0$. Ada $N_1$ dengan
$|x_n-L|<\varepsilon$ untuk $n>N_1$, dan $N_2$ dengan $|x_n-M|<\varepsilon$ untuk
$n>N_2$. Untuk $n>\max(N_1,N_2)$ kedua ketaksamaan berlaku sekaligus, sehingga
\[
|L-M| \le |L-x_n| + |x_n-M| < 2\varepsilon = |L-M| .
\]
Ketaksamaan $|L-M|<|L-M|$ mustahil, jadi pengandaian salah.
\end{proof}

Frasa "untuk $n$ yang cukup besar" berarti: ada $N$ sehingga sifat yang disebut
berlaku untuk semua $n>N$. Dengan kosakata itu, definisi limit ringkasnya: $x_n$
sedekat apa pun ke $L$ bila $n$ cukup besar.

Sifat kedua: barisan yang konvergen tidak dapat lari ke tak hingga.

\begin{teorema}
Setiap barisan konvergen terbatas.
\end{teorema}

\begin{proof}[Bukti]
Misalkan $x_n \to L$. Ambil $\varepsilon=1$; ada $N$ dengan $|x_n-L|<1$ untuk
$n>N$, sehingga $|x_n| \le |x_n-L|+|L| < |L|+1$. Hanya suku pertama sampai suku
ke-$N$ yang belum terkendali, dan jumlahnya berhingga. Maka
\[
M=\max\{|x_1|,|x_2|,\dots,|x_N|,|L|+1\}
\]
adalah batas yang berlaku untuk semua $n$.
\end{proof}

Kebalikannya tidak berlaku: $(-1)^n$ terbatas tetapi divergen. Keterbatasan adalah
syarat perlu, bukan syarat cukup; dari tabel yang tampak tenang kita hanya boleh
menyimpulkan keterbatasan.

\section{Sifat aljabar limit}
\tingkat{menengah}

Kalimat sehari-harinya: limit sopan terhadap penjumlahan, perkalian, dan pembagian,
selama pembaginya tidak menuju nol.

Sebelum membuktikan sifat-sifat itu, kita perlu satu kendali kecil: barisan yang
menuju bilangan tak nol akhirnya menjauh dari nol.

\begin{lemma}
Jika $y_n \to y$ dengan $y \neq 0$, maka ada $N$ sehingga $|y_n|>|y|/2$ untuk
semua $n>N$.
\end{lemma}

\begin{proof}[Bukti]
Ambil $\varepsilon=|y|/2>0$. Ada $N$ dengan $|y_n-y|<|y|/2$ untuk $n>N$. Maka
$|y_n| \ge |y|-|y_n-y| > |y|-|y|/2 = |y|/2$.
\end{proof}

Konsekuensinya untuk berhitung: limit dapat dihitung suku demi suku, dan satu-satunya
pengecualian adalah pembagian dengan nol.

\begin{teorema}[Sifat aljabar limit]
Jika $x_n \to x$ dan $y_n \to y$, maka
$x_n+y_n \to x+y$, \ $x_n y_n \to xy$, dan $c x_n \to cx$ untuk setiap
$c \in \mathbb{R}$. Jika $y \neq 0$ dan $y_n \neq 0$ untuk semua $n$, maka
$x_n/y_n \to x/y$.
\end{teorema}

\begin{proof}[Bukti]
Untuk jumlah, ambil $\varepsilon>0$. Pilih $N_1$ dengan $|x_n-x|<\varepsilon/2$
untuk $n>N_1$ dan $N_2$ dengan $|y_n-y|<\varepsilon/2$ untuk $n>N_2$. Untuk
$n>\max(N_1,N_2)$,
$|(x_n+y_n)-(x+y)| \le |x_n-x|+|y_n-y| < \varepsilon$. Kasus kelipatan skalar
serupa dan lebih singkat: untuk $c=0$ barisan konstan; untuk $c \neq 0$ pakai
$\varepsilon/|c|$.

Untuk perkalian, tulis selisihnya sebagai dua bagian,
\[
x_n y_n - xy = x_n(y_n-y) + (x_n-x)y .
\]
Karena $(x_n)$ konvergen, ia terbatas: ada $M>0$ dengan $|x_n| \le M$ untuk semua
$n$. Pilih $N$ sehingga untuk $n>N$ berlaku $|y_n-y|<\varepsilon/(2M)$ dan
$|x_n-x|<\varepsilon/(2(|y|+1))$. Maka
\[
|x_n y_n - xy| \le M|y_n-y| + |x_n-x|\,|y| < \frac{\varepsilon}{2}
+ \frac{\varepsilon}{2} = \varepsilon .
\]
Penyebut $|y|+1$ dipakai supaya tidak pernah nol ketika $y=0$; ia tetap konstanta
positif, jadi tidak mengubah kesimpulan.

Untuk hasil bagi, cukup ditunjukkan $1/y_n \to 1/y$, lalu gunakan hasil perkalian.
Menurut lemma, untuk $n>N_1$ berlaku $|y_n|>|y|/2$, sehingga
\[
\left|\frac{1}{y_n}-\frac{1}{y}\right| = \frac{|y-y_n|}{|y_n y|}
\le \frac{2}{|y|^2}\,|y_n-y| .
\]
Ruas kanan lebih kecil dari $\varepsilon$ asalkan $|y_n-y|<\varepsilon|y|^2/2$,
dan itu dapat dijamin dengan memilih $N_2$.
\end{proof}

Perhatikan teknik yang baru muncul: menambah dan mengurangkan suku yang sama untuk
memecah selisih menjadi dua bagian kecil. Hampir setiap bukti kekonvergenan yang
melibatkan perkalian memakainya, dan ambang yang diambil selalu yang terbesar.

Sifat berikutnya menyangkut urutan, bukan operasi. Kalimat sehari-harinya: limit
tidak dapat melompati pagar yang dipatuhi semua suku.

\begin{teorema}[Sifat urutan]
Jika $x_n \to x$ dan $x_n \ge 0$ untuk semua $n$, maka $x \ge 0$. Lebih umum,
jika $x_n \le y_n$ untuk semua $n$ yang cukup besar, $x_n \to x$, dan $y_n \to y$,
maka $x \le y$.
\end{teorema}

\begin{proof}[Bukti]
Andaikan $x<0$ dan ambil $\varepsilon=-x>0$. Ada $N$ dengan $|x_n-x|<\varepsilon$
untuk $n>N$, sehingga khususnya $x_n < x+\varepsilon = 0$. Itu bertentangan dengan
$x_n \ge 0$. Untuk pernyataan kedua, terapkan bagian pertama pada $y_n-x_n \ge 0$,
yang limitnya $y-x$ menurut teorema sebelumnya.
\end{proof}

Perhatikan bahwa kesimpulan yang sah adalah $x \ge 0$, bukan $x>0$: barisan
$1/n$ seluruhnya positif tetapi limitnya nol. Sifat urutan juga tidak boleh dipakai
mundur. Dari $x_n<y_n$ kita hanya memperoleh $x \le y$; kesenjangan ketat bisa
lenyap di limit.

Dengan sifat aljabar, limit barisan rasional dapat dihitung seperti biasa. Untuk
\[
x_n=\frac{3n^2+2n}{n^2-5},
\]
bagi pembilang dan penyebut dengan $n^2$:
$x_n=\dfrac{3+2/n}{1-5/n^2}$. Karena $1/n \to 0$ dan $1/n^2 \to 0$, pembilangnya
menuju $3$ dan penyebutnya menuju $1 \neq 0$, jadi $x_n \to 3$. Nilai numeriknya
memang merayap ke sana: $x_{10}=3.368421$, $x_{100}=3.021511$. Contoh lain,
\[
y_n=\frac{2n^2-n}{3n^2+7} \longrightarrow \frac{2}{3},
\]
dengan $y_{10}=0.618893$ dan $y_{100}=0.663179$; keduanya masih jauh dari $2/3$
pada $n=100$ karena galatnya sekitar $1/n$.

Sifat aljabar hanya berbicara tentang limit yang sudah ada; menerapkannya pada
barisan yang belum diketahui konvergen adalah kekeliruan yang mahal.

\section{Teorema apit dan barisan pembanding}
\tingkat{menengah}

Kadang limit sebuah barisan tidak dapat dihitung langsung, tetapi dapat dikurung
dari atas dan dari bawah oleh dua barisan yang limitnya sama. Kalimat sehari-harinya:
bila dua penjaga berjalan berdekatan sambil mengapit seorang anak, anak itu harus
berhenti di tempat yang sama dengan mereka.

\begin{teorema}[Teorema apit]
\index{teorema apit}
Misalkan $a_n \le x_n \le b_n$ untuk semua $n$ yang cukup besar. Jika
$a_n \to L$ dan $b_n \to L$, maka $x_n \to L$.
\end{teorema}

\begin{proof}[Bukti]
Ambil $\varepsilon>0$. Pilih $N_1$ sehingga $|a_n-L|<\varepsilon$ untuk $n>N_1$,
yakni $a_n>L-\varepsilon$, dan $N_2$ sehingga $b_n<L+\varepsilon$ untuk $n>N_2$.
Ambil juga $N_3$ sehingga pengapitan berlaku untuk $n>N_3$. Untuk
$n>\max(N_1,N_2,N_3)$,
\[
L-\varepsilon < a_n \le x_n \le b_n < L+\varepsilon,
\]
jadi $|x_n-L|<\varepsilon$.
\end{proof}

Sebelum memakainya, kita perlu satu barisan pembanding yang akan muncul terus:
pangkat yang mengecil menuju nol.

\begin{lemma}
Jika $0<q<1$, maka $q^n \to 0$.
\end{lemma}

\begin{proof}[Bukti]
Tulis $1/q=1+h$ dengan $h>0$. Ketaksamaan Bernoulli, yang dibuktikan dengan induksi
di Bab 1, memberi $(1+h)^n \ge 1+nh$. Karena kedua ruas positif,
\[
0 < q^n = \frac{1}{(1+h)^n} \le \frac{1}{1+nh} .
\]
Diberikan $\varepsilon>0$, sifat Archimedes memberi $n$ dengan $1+nh>1/\varepsilon$,
yakni $1/(1+nh)<\varepsilon$. Menurut teorema apit, $q^n \to 0$.
\end{proof}

Sekarang contoh yang tidak dapat diselesaikan dengan membagi pangkat tertinggi.
Untuk $x_n=(2^n+3^n)^{1/n}$, suku terbesar di dalam kurung adalah $3^n$; keluarkan
pangkat itu:
\[
x_n = 3\left(1+\left(\tfrac{2}{3}\right)^{n}\right)^{1/n}.
\]
Karena $(2/3)^n \to 0$ dan $1 \le 1+(2/3)^n \le 2$, kita peroleh
$1 \le x_n/3 \le 2^{1/n}$. Tinggal ditunjukkan $2^{1/n} \to 1$: tulis
$2^{1/n}=1+a_n$ dengan $a_n \ge 0$, maka
$2=(1+a_n)^n \ge 1+na_n$, sehingga $0 \le a_n \le 1/n$. Teorema apit menutup
kedua ujung, dan $x_n \to 3$. Angka Python menegaskan arahnya:
$x_1=5$, $x_2=3.605551$, $x_5=3.075152$, $x_{10}=3.005162$, dan $x_{100}=3$.

\begin{untukanak}
Dua penjaga sekolah berjalan berdampingan menyusuri koridor, dan seorang murid
berjalan di antara mereka. Penjaga di kiri melangkah ke arah pintu keluar, penjaga
di kanan juga, dan keduanya makin rapat sampai bahu mereka hampir bersentuhan.
Murid yang terjepit di tengah tidak perlu tahu jalan; ia hanya perlu tetap di
antara kedua penjaga. Ke mana pun penjaga itu berhenti, di situ pula murid
berhenti.

Sekarang pertanyaan yang membuat penasaran. Bagaimana kalau penjaga kiri berhenti
lebih awal daripada penjaga kanan: masih bisakah kita menebak di mana murid
berhenti? Lalu bagaimana kalau jarak mereka tidak pernah nol, hanya makin kecil
dari hari ke hari? Apakah murid itu tetap punya tempat berhenti yang pasti, atau
terus berjalan tanpa ujung?

Ternyata jawabannya tidak bergantung pada wajah penjaganya, hanya pada dua hal: ke
mana penjaga kiri menuju, dan ke mana penjaga kanan menuju. Kalau keduanya menuju
tempat yang sama, murid di tengah tidak punya pilihan lain. Cara berpikir "terjepit
di antara dua pihak" itulah yang dipakai matematikawan untuk barisan yang rumusnya
terlalu rumit dihitung langsung.
\end{untukanak}

\section{Barisan monoton dan teorema kekonvergenan monoton}
\tingkat{menengah}

Kalimat sehari-harinya: barisan yang hanya bergerak satu arah, dan tidak pernah
melewati pagar tertentu, pasti berhenti di suatu tempat. Inilah tempat kelengkapan
bekerja paling murah, karena limitnya tidak perlu ditebak lebih dahulu.

\kamus{barisan monoton}{monotone sequence}{tangga yang anak tangganya hanya naik,
tidak pernah turun}

Bahasa sehari-harinya sederhana: satu arah saja, tidak pernah berbalik.

\begin{definisi}[Barisan monoton]
Barisan $(x_n)$ disebut \emph{naik}\index{barisan monoton} jika $x_n \le x_{n+1}$
untuk semua $n$, dan \emph{turun} jika $x_n \ge x_{n+1}$ untuk semua $n$.
Barisan disebut monoton jika ia naik atau turun.
\end{definisi}

Gambaran kasarnya: yang melangkah satu arah tanpa melewati pagar pasti berhenti, dan
berhentinya tepat di pagar terdekat yang tidak pernah dilewati.

\begin{teorema}[Kekonvergenan monoton]
\index{teorema kekonvergenan monoton}
Barisan naik yang terbatas di atas konvergen, dan limitnya adalah
$\sup\{x_n : n \in \mathbb{N}\}$.
\end{teorema}

\begin{proof}[Bukti]
Misalkan $(x_n)$ naik dan terbatas di atas. Himpunan sukunya tidak kosong dan
terbatas di atas, jadi menurut sifat supremum di Bab 3 ia mempunyai supremum;
sebut $s=\sup\{x_n\}$. Ambil $\varepsilon>0$. Karena $s-\varepsilon$ bukan batas
atas, ada $N$ dengan $x_N>s-\varepsilon$. Untuk $n>N$, kemonotonan memberi
$x_n \ge x_N>s-\varepsilon$, sedangkan $x_n \le s$ menurut definisi supremum. Jadi
$|x_n-s|<\varepsilon$.
\end{proof}

Akibat langsungnya: untuk barisan seperti itu limitnya dapat dibaca dari himpunan
sukunya, tanpa menghitung suku demi suku.

\begin{akibat}
Barisan turun yang terbatas di bawah konvergen, dan limitnya adalah infimum
himpunan sukunya.
\end{akibat}

\begin{proof}[Bukti]
Terapkan teorema pada $(-x_n)$, yang naik dan terbatas di atas, lalu gunakan sifat
kelipatan skalar pada limit dan kesamaan
$\sup\{-x_n\}=-\inf\{x_n\}$.
\end{proof}

Perhatikan bahwa teorema ini tidak mengatakan limitnya sama dengan batas atas yang
kita pakai. Batas $3$ pada contoh berikut hanya alat bukti; limitnya justru $e$.

Aplikasi pertama adalah barisan rekursif yang membuka bab ini. Misalkan $x_1=1$ dan
$x_{n+1}=\sqrt{2+x_n}$.

Pertama, semua suku berada di antara $1$ dan $2$. Untuk $n=1$ jelas. Bila
$1 \le x_n \le 2$, maka $3 \le 2+x_n \le 4$, sehingga
$1 \le \sqrt{3} \le x_{n+1} \le 2$. Dengan induksi, $1 \le x_n \le 2$ untuk semua
$n$.

Kedua, barisan itu naik. Karena $x_n \le 2$, berlaku
\[
x_{n+1} \ge x_n \iff 2+x_n \ge x_n^2 \iff (2-x_n)(1+x_n) \ge 0,
\]
dan ruas terakhir benar sebab $x_n \le 2$ dan $x_n \ge 1$.

Menurut teorema kekonvergenan monoton, $(x_n)$ konvergen; sebut limitnya $L$.
Karena $x_n \le 2$ untuk semua $n$, sifat urutan memberi $L \le 2$. Sekarang
kuadratkan relasinya: $x_{n+1}^2=2+x_n$. Ruas kiri menuju $L^2$ menurut sifat
perkalian, ruas kanan menuju $2+L$, jadi $L^2=2+L$, atau
$(L-2)(L+1)=0$. Karena $L \ge 1$, satu-satunya kemungkinan adalah $L=2$.

Perhatikan bahwa kita tidak pernah memakai kesinambungan akar; cukup sifat aljabar
limit dan sifat urutan, sehingga peralatan Bab 11 belum diperlukan di sini.

Aplikasi kedua adalah barisan $y_n=(1+1/n)^n$. Ekspansi binomial memberi
\[
y_n = \sum_{k=0}^{n} \binom{n}{k}\frac{1}{n^k}
= \sum_{k=0}^{n} \frac{1}{k!}\prod_{j=1}^{k-1}\left(1-\frac{j}{n}\right),
\]
dengan hasil kali kosong bernilai $1$ untuk $k=0$ dan $k=1$. Setiap faktor
$1-j/n$ naik ketika $n$ bertambah, sehingga setiap suku dengan indeks $k$ naik;
jumlahnya juga mendapat satu suku baru yang positif. Jadi $(y_n)$ naik. Setiap
faktor paling besar $1$, sehingga
\[
y_n \le \sum_{k=0}^{n}\frac{1}{k!} \le 1+1+\frac12+\frac14+\cdots = 3,
\]
karena $1/k! \le 1/2^{k-1}$ untuk $k \ge 1$, yang diperiksa dengan induksi.
Barisan naik dan terbatas di atas oleh $3$, jadi ia konvergen. Limitnya kita beri
nama $e$.

Gambar~\ref{fig:monoton-e} menampilkan barisan itu untuk $n=1$ sampai $n=30$.
Titik-titiknya naik tanpa pernah berbalik, garis putus-putus di $3$ adalah pagar
yang dipakai dalam bukti, dan garis titik-titik di $2.71828$ adalah tempat ia
benar-benar berhenti. Jarak antara pagar dan tujuan tetap ada sepanjang gambar;
pagar bekerja karena tidak pernah dilewati, bukan karena dicapai.

\begin{figure}[htbp]
\centering
\begin{tikzpicture}[x=0.32cm,y=3.0cm,>=Stealth,font=\footnotesize]
  \draw[->,ink] (-0.6,1.85) -- (31.2,1.85) node[right] {$n$};
  \draw[->,ink] (0,1.85) -- (0,3.12) node[above] {$y_n$};
  \draw[dashed,layergrey,thick] (0,3.0) -- (30.6,3.0);
  \node[layergrey,anchor=south] at (20,3.02) {batas bukti $B=3$};
  \draw[densely dotted,layerindigo,very thick] (0,2.71828) -- (30.4,2.71828);
  \node[layerindigo,anchor=west] at (31.3,2.71828) {limit $e\approx2.71828$};
  \foreach \n/\y in {1/2.0000,2/2.2500,3/2.3704,4/2.4414,5/2.4883,6/2.5216,%
      7/2.5465,8/2.5658,9/2.5812,10/2.5937,11/2.6042,12/2.6130,13/2.6206,%
      14/2.6272,15/2.6329,16/2.6379,17/2.6424,18/2.6464,19/2.6500,20/2.6533,%
      21/2.6563,22/2.6590,23/2.6615,24/2.6637,25/2.6658,26/2.6678,27/2.6696,%
      28/2.6713,29/2.6728,30/2.6743}
    \fill[accent] (\n,\y) circle (1.7pt);
  \draw[->,layerteal,thick] (9.5,2.42) -- (15.5,2.60);
  \node[layerteal,anchor=north] at (12,2.40) {naik, tidak pernah berbalik};
\end{tikzpicture}
\caption{Barisan $y_n=(1+1/n)^n$ untuk $n=1$ sampai $n=30$. Titik-titiknya naik dan
tidak pernah berbalik; pagar $B=3$ tidak pernah disentuh; nilai yang dituju adalah
$e\approx2.71828$, bukan pagarnya.}
\label{fig:monoton-e}
\end{figure}

\section{Subbarisan, Bolzano--Weierstrass, dan kriteria Cauchy}
\tingkat{lanjut}

Kalimat sehari-harinya: dari barisan yang terbatas selalu dapat dipetik satu
deretan anak yang berhenti di suatu tempat, walaupun deretan induknya kacau.

\kamus{subbarisan}{subsequence}{menyortir sebagian surat dari tumpukan tanpa
mengubah urutannya}

Subbarisan dibentuk dengan memilih sebagian suku saja: nomor yang dipakai makin
besar, tetapi suku yang dilewati boleh sebanyak apa pun.

\begin{definisi}[Subbarisan]
Misalkan $n_1<n_2<n_3<\cdots$ bilangan asli naik. Barisan $(x_{n_k})_{k \in
\mathbb{N}}$ disebut \emph{subbarisan}\index{subbarisan} dari $(x_n)$.
\end{definisi}

Karena $n_k \ge k$ untuk setiap $k$, yang diperiksa dengan induksi, subbarisan
tidak dapat "lebih lambat" daripada barisan asalnya dalam hal indeks.

Kalimat sehari-harinya: memilih sebagian suku tidak mengubah tempat berhenti, hanya
memperlambat atau mempercepat kedatangannya.

\begin{proposisi}
Jika $x_n \to L$, maka setiap subbarisan $(x_{n_k})$ juga konvergen ke $L$.
\end{proposisi}

\begin{proof}[Bukti]
Ambil $\varepsilon>0$ dan $N$ dengan $|x_n-L|<\varepsilon$ untuk $n>N$. Untuk
$k>N$ berlaku $n_k \ge k>N$, jadi $|x_{n_k}-L|<\varepsilon$.
\end{proof}

Akibatnya sederhana tetapi paling sering dipakai untuk menyangkal: dua tempat
berhenti yang berbeda sudah cukup.

\begin{akibat}
Jika $(x_n)$ mempunyai dua subbarisan dengan limit berbeda, maka $(x_n)$ divergen.
\end{akibat}

\begin{proof}[Bukti]
Bila $(x_n)$ konvergen, semua subbarisannya konvergen ke limit yang sama, dan
ketunggalan limit melarang dua nilai berbeda.
\end{proof}

Akibat ini menjelaskan contoh $(-1)^n$ tanpa argumen kemustahilan: subbarisan
genap menuju $1$, subbarisan ganjil menuju $-1$. Cara kerjanya juga menjelaskan
mengapa divergensi lebih sering dapat dibuktikan daripada kekonvergenan.

Untuk membuktikan Bolzano--Weierstrass, kita butuh lemma berikut. Isinya: setiap
barisan, betapa pun kacau, memuat subbarisan yang monoton.

\begin{lemma}
Setiap barisan real mempunyai subbarisan monoton.
\end{lemma}

\begin{proof}[Bukti]
Sebut indeks $n$ sebuah \emph{puncak} bila $x_n \ge x_m$ untuk semua $m>n$, yakni
bila tidak ada suku sesudahnya yang lebih besar. Ada dua kemungkinan.

Bila puncak tak berhingga banyak, susun $n_1<n_2<\cdots$ sebagai daftar puncak.
Karena $n_1$ puncak, $x_{n_1} \ge x_{n_2}$; karena $n_2$ puncak,
$x_{n_2} \ge x_{n_3}$; dan seterusnya. Hasilnya subbarisan turun.

Bila puncak hanya berhingga, ambil $n_1$ melewati puncak terakhir. Karena $n_1$
bukan puncak, ada $n_2>n_1$ dengan $x_{n_2}>x_{n_1}$. Karena $n_2$ juga melewati
puncak terakhir, ia bukan puncak, jadi ada $n_3>n_2$ dengan $x_{n_3}>x_{n_2}$.
Proses ini berlanjut tanpa henti dan menghasilkan subbarisan naik.
\end{proof}

Isi teorema ini: dari barisan yang terkurung selalu dapat dipetik sebagian suku yang
benar-benar menuju satu titik.

\begin{teorema}[Bolzano--Weierstrass]
\index{bolzano-weierstrass}
Setiap barisan real yang terbatas mempunyai subbarisan yang konvergen.
\end{teorema}

\begin{proof}[Bukti]
Menurut lemma, $(x_n)$ mempunyai subbarisan monoton. Subbarisan dari barisan
terbatas juga terbatas. Barisan monoton yang terbatas konvergen menurut teorema
kekonvergenan monoton. Jadi subbarisan itu konvergen.
\end{proof}

Bukti alternatif memakai pembagian dua interval dan sifat interval bersarang; kedua
bukti bersandar pada kelengkapan, yang tadi dipinjam lewat supremum.

\kamus{barisan Cauchy}{Cauchy sequence}{penumpang yang makin berdesakan di lorong
kereta, tanpa tahu pintu keluarnya di mana}

Sekarang perhatikan masalah yang berbeda. Definisi limit menuntut kita menyebut
$L$ lebih dahulu, padahal dalam banyak perhitungan justru $L$ yang dicari. Yang
kita punya hanya daftar angka. Karena itu kita perlu ukuran kedekatan antar suku
sendiri, tanpa menyebut tujuan.

\begin{definisi}[Barisan Cauchy]
Barisan $(x_n)$ disebut \emph{Cauchy}\index{barisan cauchy} jika untuk setiap
$\varepsilon>0$ terdapat $N$ sehingga untuk semua $m,n>N$ berlaku
$|x_n-x_m|<\varepsilon$.
\end{definisi}

Perhatikan perbedaan urutan gagasannya. Pada definisi limit semua suku dibandingkan
dengan satu bilangan tetap $L$; pada definisi Cauchy suku dibandingkan dengan suku
lain, tanpa bilangan luar. Karena itu definisi Cauchy masuk akal di ruang mana pun
yang punya jarak, sedangkan definisi limit menuntut calon tujuan.

\begin{teorema}[Kriteria Cauchy]
\index{barisan cauchy}
Barisan real konvergen jika dan hanya jika ia Cauchy.
\end{teorema}

\begin{proof}[Bukti]
Misalkan $x_n \to L$ dan ambil $\varepsilon>0$. Ada $N$ dengan
$|x_n-L|<\varepsilon/2$ untuk $n>N$. Untuk $m,n>N$,
\[
|x_n-x_m| \le |x_n-L| + |L-x_m| < \varepsilon .
\]
Jadi konvergen memberi Cauchy.

Sebaliknya, misalkan $(x_n)$ Cauchy. Ambil $\varepsilon=1$; ada $N_0$ sehingga
$|x_n-x_m|<1$ untuk $m,n>N_0$. Dengan memilih $m=N_0+1$ diperoleh
$|x_n|<|x_{N_0+1}|+1$ untuk $n>N_0$, jadi $(x_n)$ terbatas. Menurut
Bolzano--Weierstrass ada subbarisan $x_{n_k} \to L$. Ambil $\varepsilon>0$; pilih
$N_1$ dengan $|x_n-x_m|<\varepsilon/2$ untuk $m,n>N_1$, dan $K$ dengan
$|x_{n_k}-L|<\varepsilon/2$ untuk $k>K$. Pilih $k>K$ yang juga memenuhi
$n_k>N_1$; untuk $n>N_1$,
\[
|x_n-L| \le |x_n-x_{n_k}| + |x_{n_k}-L| < \varepsilon .
\]
Jadi $x_n \to L$.
\end{proof}

Arah kedua memakai Bolzano--Weierstrass, yang memakai teorema kekonvergenan
monoton, yang memakai supremum. Rantai itu tidak dapat diputus: di dalam
$\mathbb{Q}$ teorema ini gagal. Barisan pemotongan desimal
$1;\ 1.4;\ 1.41;\ 1.414;\ 1.4142;\ \dots$ menuju akar dua adalah Cauchy, sebab
selisih dua sukunya paling besar $10^{-k}$ bila keduanya melewati suku ke-$k$,
tetapi tidak ada bilangan rasional yang menjadi limitnya. Bila barisan itu
konvergen ke $L$, maka $a_k^2 \to L^2$; di pihak lain $a_k^2 \to 2$, sehingga
$L^2=2$. Bab 3 sudah menunjukkan tidak ada rasional yang kuadratnya $2$. Jadi
"Cauchy" dan "konvergen" adalah dua sifat berbeda; yang menyamakan keduanya adalah
kelengkapan, dan itulah isi aksioma supremum.

\begin{tantangan}
Barisan $h_n=1+\frac12+\cdots+\frac1n$ tumbuh sangat lambat: $h_{10}=2.928968$,
$h_{100}=5.187378$, $h_{1000}=7.485471$, dan $h_{10^6}=14.392727$. Selisih dua suku
berurutan, yaitu $1/(n+1)$, menuju nol, sehingga barisan ini lolos pemeriksaan
dangkal yang paling sering menipu: suku-sukunya makin rapat. Pertanyaannya, apakah
$(h_n)$ konvergen? Buktikan jawabanmu, dan jelaskan mengapa $h_{10^6}<15$ bukan
bukti bahwa barisan itu terbatas.

Petunjuk 1: jangan periksa suku satu per satu. Jumlahkan dalam blok yang panjangnya
berlipat dua, yaitu satu suku, dua suku, empat suku, delapan suku, dan seterusnya.

Petunjuk 2: blok ke-$k$ berisi $2^{k-1}$ suku, dan suku terkecil di dalamnya paling
kecil $1/2^k$; jadi setiap blok menyumbang sedikitnya $1/2$. Bandingkan hasilnya
dengan $h_{2^m}$ dan lihat apa yang terjadi bila $m$ membesar.

Kalau mentok, kembali ke Teorema Kekonvergenan Monoton: di sana barisan naik
dibuktikan konvergen justru karena ada pagar, dan di sini justru pagarnya yang
tidak ada.
\end{tantangan}

\section{Limit superior dan limit inferior}
\tingkat{lanjut}

Kalimat sehari-harinya: untuk barisan yang berayun, kita masih dapat menyebut
tempat berayun yang paling tinggi dan yang paling rendah, walaupun barisan itu
sendiri tidak punya limit.

\kamus{limit superior}{upper limit}{puncak tertinggi yang masih dikunjungi
barisan sesering apa pun setelah melewati sebarang titik}

Sepanjang bagian ini $(x_n)$ terbatas. Untuk setiap $n$ bentuklah himpunan ekor
$\{x_k : k \ge n\}$ dan tulis
\[
s_n=\sup\{x_k : k \ge n\}, \qquad t_n=\inf\{x_k : k \ge n\}.
\]
Karena ekornya menyusut, $(s_n)$ turun dan $(t_n)$ naik; keduanya terbatas, jadi
keduanya konvergen menurut akibat teorema kekonvergenan monoton.

Kalimat sehari-harinya: yang kita lacak bukan di mana barisan berada, melainkan
seberapa tinggi dan seberapa rendah ia masih berkunjung setelah titik mana pun.

\begin{definisi}[Limit superior dan inferior]
Untuk barisan terbatas $(x_n)$,
\[
\limsup_{n\to\infty} x_n = \lim_{n\to\infty} s_n
\qquad \text{dan} \qquad
\liminf_{n\to\infty} x_n = \lim_{n\to\infty} t_n .
\]
Bilangan pertama disebut \emph{limit superior}\index{limit superior}, yang kedua
\emph{limit inferior}\index{limit inferior}.
\end{definisi}

Karena $s_n \ge t_n$ untuk setiap $n$ dan limit mempertahankan urutan, selalu
berlaku $\liminf x_n \le \limsup x_n$. Keduanya dapat berbeda jauh, dan justru
selisih itulah ukuran kekacauan barisan.

\begin{contoh}
Untuk $x_n=(-1)^n+1/n$, ekornya memuat suku genap $1+1/k$ dan suku ganjil
$-1+1/k$. Supremum ekor ke-$n$ datang dari suku genap terkecil yang masih
tersedia, jadi $s_1=1.5$, $s_2=1.5$, $s_3=1.25$, $s_4=1.25$, $s_5=1.166667$,
dan seterusnya menurun menuju $1$. Infimum ekornya berbeda sifat: suku ganjil
$-1+1/k$ selalu hadir di setiap ekor dan turun menuju $-1$ tanpa pernah
mencapainya, sehingga $t_n=-1$ untuk semua $n$, walaupun tidak ada satu suku pun
yang bernilai $-1$. Maka $\limsup x_n=1$ dan $\liminf x_n=-1$, sedangkan
barisannya divergen.
\end{contoh}

Gambar~\ref{fig:divergen-ayun} memplot suku ke-$n$ terhadap $n$ untuk barisan yang
sama. Deretan titik genap menempel di dekat $1$, deretan titik ganjil menempel di
dekat $-1$, dan jarak kedua kumpulan itu tetap $2$. Tidak ada satu garis mendatar
yang dapat didekati seluruh titik, sejauh apa pun kita berjalan.

\begin{figure}[htbp]
\centering
\begin{tikzpicture}[x=0.62cm,y=1.9cm,>=Stealth,font=\footnotesize]
  \draw[->,ink] (-0.3,0) -- (13.4,0) node[right] {$n$};
  \draw[->,ink] (0,-1.3) -- (0,2.0) node[above] {$x_n$};
  \draw[densely dotted,layerindigo,very thick] (0,1) -- (13.0,1);
  \draw[densely dotted,layerplum,very thick] (0,-1) -- (13.0,-1);
  \node[layerindigo,anchor=south west] at (0.1,1.03) {$\limsup x_n=1$};
  \node[layerplum,anchor=north west] at (0.1,-1.03) {$\liminf x_n=-1$};
  \foreach \n/\y in {1/0.0000,2/1.5000,3/-0.6667,4/1.2500,5/-0.8000,6/1.1667,%
      7/-0.8571,8/1.1250,9/-0.8889,10/1.1000,11/-0.9091,12/1.0833}
    \fill[accent] (\n,\y) circle (2.2pt);
  \draw[<->,layerbrown,thick] (12.5,-1) -- (12.5,1) node[midway,right] {jarak $2$};
\end{tikzpicture}
\caption{Grafik suku ke-$n$ terhadap $n$ untuk $x_n=(-1)^n+1/n$. Titik genap
menempel di dekat $1$, titik ganjil di dekat $-1$, dan jarak kedua kumpulan itu
tetap $2$. Tidak ada pita mendatar yang memuat semua suku, jadi barisan ini
divergen.}
\label{fig:divergen-ayun}
\end{figure}

Kalimat sehari-harinya: nilai tertinggi yang masih dikunjungi barisan sesering apa
pun memang dikunjungi oleh sebagian sukunya, bukan sekadar didekati dari jauh.

\begin{teorema}
Untuk barisan terbatas $(x_n)$, limit superiornya adalah limit subbarisan
terbesar: $\limsup x_n$ sama dengan supremum himpunan semua limit subbarisan yang
konvergen, dan nilai itu sendiri dicapai oleh suatu subbarisan.
\end{teorema}

\begin{proof}[Bukti]
Tulis $S=\limsup x_n$. Pertama, setiap limit subbarisan paling besar $S$. Bila
$x_{n_k} \to L$, maka untuk setiap $k$ berlaku $x_{n_k} \le s_{n_k}$, dan karena
$n_k \ge k$ serta $(s_n)$ turun, $s_{n_k} \le s_k$. Jadi $x_{n_k} \le s_k$, dan
mengambil limit memberi $L \le S$.

Kedua, kita bangun subbarisan yang menuju $S$, mulai dari $n_0=0$. Untuk setiap $k$,
karena $s_n \to S$ dan $(s_n)$ turun, ada $M \ge k$ dengan $s_M<S+1/k$. Ekor
$\{x_n : n \ge M\}$ mempunyai supremum $s_M$ dan memuat indeks sebesar apa pun, jadi
ada $n_k>\max(k,n_{k-1})$ dengan $n_k \ge M$ dan $x_{n_k}>s_M-1/k$. Untuk pilihan itu
\[
S - \frac1k \le s_M - \frac1k < x_{n_k} \le s_M < S + \frac1k ,
\]
sehingga $|x_{n_k}-S|<1/k$. Jadi $x_{n_k} \to S$.
\end{proof}

Akibatnya: kesamaan limit inferior dan limit superior menjadi cara termurah
memeriksa kekonvergenan tanpa menebak limitnya lebih dahulu.

\begin{akibat}
Barisan terbatas $(x_n)$ konvergen ke $L$ jika dan hanya jika
$\liminf x_n = \limsup x_n = L$.
\end{akibat}

\begin{proof}[Bukti]
Bila $x_n \to L$, setiap subbarisan juga menuju $L$, jadi himpunan limit
subbarisan hanya berisi $L$ dan supremumnya $L$; dengan cara yang sama infimumnya
$L$. Sebaliknya, bila kedua nilai sama dengan $L$, maka $t_n \le x_n \le s_n$ dan
kedua pagar menuju $L$, sehingga teorema apit memberi $x_n \to L$.
\end{proof}

Bagi barisan tak terbatas dipakai kesepakatan $\limsup x_n=+\infty$ bila $(x_n)$
tidak terbatas di atas; hukum aritmetika untuk $\infty$ tidak berlaku otomatis, dan
kita tidak memakainya di sini.

\section{Contoh hitung}
\tingkat{menengah}

Kita kembali ke dua barisan pembuka, kali ini dengan angka yang bisa diperiksa
ulang dan dengan langkah antara yang ditulis.

\begin{resep}
\begin{enumerate}
  \item Tulis selisihnya. Bentuk $|x_n-L|$, lalu sederhanakan sampai tinggal satu
        pecahan dalam $n$; untuk contoh di bawah, $|x_n-3|=5/(n+2)$.
  \item Selesaikan untuk $n$. Pecahkan $5/(n+2)<\varepsilon$ menjadi
        $n>5/\varepsilon-2$, dan namai ambang itu $g(\varepsilon)$.
  \item Pilih $N$. Ambil $N=\lfloor g(\varepsilon)\rfloor$, bilangan terkecil yang
        sah, lalu periksa galat pada $n=N$ dan $n=N+1$ sebelum menulis apa pun.
  \item Tutup dengan kalimat definisi: untuk setiap $\varepsilon>0$ dipilih $N$
        itu, maka setiap $n>N$ memenuhi $|x_n-3|<\varepsilon$.
\end{enumerate}
Peringatan urutan: $\varepsilon$ disebut lebih dahulu, $N$ menjawab. Menyusun $N$
lalu mencari $\varepsilon$ yang cocok membalik permainan; hasilnya biasanya $N$
yang bergantung pada $n$, bukan pada $\varepsilon$, dan di titik itu buktinya batal.
\end{resep}

Untuk $x_n=(3n+1)/(n+2)$,
\[
|x_n-3| = \left|\frac{3n+1-3n-6}{n+2}\right| = \frac{5}{n+2}.
\]
Pertidaksamaan $5/(n+2)<\varepsilon$ setara dengan $n>5/\varepsilon-2$. Karena
syaratnya ketat, bilangan bulat terbesar yang masih di bawah ambang itu sudah sah:
$N=\lfloor 5/\varepsilon-2 \rfloor$ adalah pilihan terkecil yang memenuhi definisi.
Untuk $\varepsilon=10^{-2}$ ambangnya $498$ dan $N=498$. Uji batasnya: pada $n=498$
galatnya $5/500=0.01$, yang \emph{tidak} lebih kecil dari $\varepsilon$; pada $n=499$
galatnya $5/501=0.00998004$, baru memenuhi. Untuk $\varepsilon=10^{-4}$ diperoleh
$N=49998$, dengan galat $0.0001$ pada $n=49998$ dan $9.9998 \times 10^{-5}$ pada
$n=49999$.

Nilai sukunya bergerak sangat lambat: $n=10$ memberi $2.583333$, $n=1000$ memberi
$2.995010$, dan galat $10^{-4}$ baru tercapai pada suku ke-$49999$. Empat angka di
depan koma sudah benar jauh sebelum galatnya kecil.

\kertas{Untuk $x_n=(3n+1)/(n+2)$, hitung $N$ yang sah untuk $\varepsilon=10^{-3}$
dengan rumus $N=\lfloor 5/\varepsilon-2 \rfloor$, lalu periksa galat pada $n=N$ dan
$n=N+1$. Jelaskan mengapa galat pada $n=N$ boleh tepat sama dengan $\varepsilon$
sementara definisi tetap terpenuhi.}

Untuk $y_n=(1+1/n)^n$, aritmetika 40 digit memberi galat $0.124539$ pada $n=10$,
$0.00135790$ pada $n=10^3$, dan $1.35914 \times 10^{-7}$ pada $n=10^7$. Hasil kali
$n$ dengan galatnya hampir tetap: $1.245394$, $1.357896$, lalu merayap ke
$1.3591409 = e/2$. Jadi galatnya turun seperti $1/n$, bukan $1/n^2$; satu angka
tambahan di belakang koma menuntut $n$ sepuluh kali lebih besar. Di sinilah tabel
mesin pada pembuka bab rusak: pada $n=10^7$ float64 melaporkan
$1.34327 \times 10^{-7}$, karena galat relatif kecil pada $1+1/n$
terlipatgandakan oleh pangkat $n$.

Untuk rekursi $x_{n+1}=\sqrt{2+x_n}$ dengan $x_1=1$, galatnya terhadap $2$ adalah
$2.679 \times 10^{-1}$ pada $n=2$, $1.711 \times 10^{-2}$ pada $n=4$, dan
$2.615 \times 10^{-7}$ pada $n=12$. Nisbah dua galat berurutan menetap di
$0.250000$.

Angka $1/4$ itu bukan kebetulan. Dari relasinya,
\[
x_{n+1}-2 = \frac{(2+x_n)-4}{\sqrt{2+x_n}+2}
= \frac{x_n-2}{\sqrt{2+x_n}+2},
\]
sehingga
\[
|x_{n+1}-2| = \frac{|x_n-2|}{\sqrt{2+x_n}+2}
\le \frac{|x_n-2|}{2+\sqrt{2}} = 0.292893\,|x_n-2| ,
\]
karena $x_n \ge 1$ memberi $\sqrt{2+x_n} \ge \sqrt{3}$. Batas $0.292893$ berlaku
untuk semua $n$, dan nisbah sebenarnya menuju $1/4$ karena penyebutnya menuju $4$.
Akibat praktisnya: galat $10^{-6}$ tercapai pada $n=12$, $10^{-9}$ pada $n=17$, dan
ketelitian mesin sekitar $n=27$; jauh lebih murah daripada membagi interval dua kali
lipat.

\section{Lab komputasi: empat pemeriksaan}
\tingkat{lanjut}

\kamus{galat pemotongan}{truncation error}{selisih antara potret dan wajah aslinya}

Empat klaim di bab ini dapat diperiksa dengan satu skrip pendek: ambang $N$ yang
eksplisit, laju galat $(1+1/n)^n$, ketelitian mesin, dan nisbah galat rekursi.

\begin{lab}
\begin{verbatim}
import math, mpmath as mp

def N_for(eps):
    N = int(5/eps - 2)
    while 5/(N+2) >= eps: N += 1
    while N > 0 and 5/(N+2) < eps: N -= 1
    return N

print("== 1. N eksplisit untuk (3n+1)/(n+2) -> 3 ==")
for eps in (1e-2, 1e-4):
    N = N_for(eps)
    print("eps=%.0e  N=%d  err(N)=%.6g  err(N+1)=%.6g"
          % (eps, N, abs((3*N+1)/(N+2)-3), abs((3*N+4)/(N+3)-3)))

print("== 2. (1+1/n)^n dalam float64 ==")
for k in range(1, 8):
    n = 10**k
    print("n=1e%d  %.16f  e-y=%.3e" % (k, (1+1/n)**n, math.e-(1+1/n)**n))

print("== 3. galat sebenarnya (mpmath, 40 digit) ==")
mp.mp.dps = 40
for k in (5, 6, 7):
    n = 10**k
    y = (1+mp.mpf(1)/n)**n
    print("n=1e%d  e-y=%s  (float64: %.3e)"
          % (k, mp.nstr(mp.e-y, 6), math.e-(1+1/n)**n))

print("== 4. rekursi x_{n+1}=sqrt(2+x_n), x_1=1 ==")
x, xs = 1.0, []
for _ in range(30):
    xs.append(x); x = math.sqrt(2+x)
for n in (1, 2, 3, 4, 5, 6, 11, 12, 13, 14):
    e = abs(xs[n-1]-2)
    r = "" if n == 1 else "  rasio=%.6f" % (e/abs(xs[n-2]-2))
    print("n=%2d  x=%.12f  galat=%.3e%s" % (n, xs[n-1], e, r))
e6 = next(n for n in range(1, 31) if abs(xs[n-1]-2) < 1e-6)
e9 = next(n for n in range(1, 31) if abs(xs[n-1]-2) < 1e-9)
print("galat < 1e-6 pertama kali pada n=%d ; galat < 1e-9 pada n=%d" % (e6, e9))
\end{verbatim}
\end{lab}

\newpage
Keluarannya, apa adanya:

\begin{verbatim}
== 1. N eksplisit untuk (3n+1)/(n+2) -> 3 ==
eps=1e-02  N=498  err(N)=0.01  err(N+1)=0.00998004
eps=1e-04  N=49998  err(N)=0.0001  err(N+1)=9.9998e-05
== 2. (1+1/n)^n dalam float64 ==
n=1e1  2.5937424601000023  e-y=1.245e-01
n=1e2  2.7048138294215285  e-y=1.347e-02
n=1e3  2.7169239322355936  e-y=1.358e-03
n=1e4  2.7181459268249255  e-y=1.359e-04
n=1e5  2.7182682371922975  e-y=1.359e-05
n=1e6  2.7182804690957534  e-y=1.359e-06
n=1e7  2.7182816941320818  e-y=1.343e-07
== 3. galat sebenarnya (mpmath, 40 digit) ==
n=1e5  e-y=1.35913e-5  (float64: 1.359e-05)
n=1e6  e-y=1.35914e-6  (float64: 1.359e-06)
n=1e7  e-y=1.35914e-7  (float64: 1.343e-07)
== 4. rekursi x_{n+1}=sqrt(2+x_n), x_1=1 ==
n= 1  x=1.000000000000  galat=1.000e+00
n= 2  x=1.732050807569  galat=2.679e-01  rasio=0.267949
n= 3  x=1.931851652578  galat=6.815e-02  rasio=0.254333
n= 4  x=1.982889722748  galat=1.711e-02  rasio=0.251074
n= 5  x=1.995717846477  galat=4.282e-03  rasio=0.250268
n= 6  x=1.998929174953  galat=1.071e-03  rasio=0.250067
n=11  x=1.999998954179  galat=1.046e-06  rasio=0.250000
n=12  x=1.999999738545  galat=2.615e-07  rasio=0.250000
n=13  x=1.999999934636  galat=6.536e-08  rasio=0.250000
n=14  x=1.999999983659  galat=1.634e-08  rasio=0.250000
galat < 1e-6 pertama kali pada n=12 ; galat < 1e-9 pada n=17
\end{verbatim}

Tiga hal yang perlu dibaca dari keluaran itu. Pertama, galat pada $n=N$ tepat sama
dengan $\varepsilon$ di kedua baris; itulah sebabnya definisi memakai syarat ketat
$n>N$. Kedua, kolom $e-y$ berlipat sepersepuluh setiap kali $n$ berlipat sepuluh
sampai $n=10^6$, lalu berhenti mengikuti pola; bagian 3 menunjukkan bahwa yang
berhenti adalah ketelitian float64, bukan kekonvergenan barisan. Galat pemotongan
terus mengecil seperti $e/(2n)$, tetapi galat pembulatan mengambil alih begitu
keduanya bertemu. Ketiga, nisbah galat pada bagian 4 turun ke $0.250000$ dan tinggal
di sana, tepat seperti perhitungan tangan sebelumnya. Program tidak membuktikan apa
pun: ia hanya memeriksa klaim yang sudah kita buktikan.

\begin{untukdosen}
Urutan mengajar: mulai dari permainan $\varepsilon$-$N$ di kelas, bukan dari
teorema. Tunjuk dua baris mahasiswa; satu menyebut $\varepsilon$, satu menjawab $N$
untuk $x_n=(3n+1)/(n+2)$, dan dosen mencatat sampai rumusnya muncul sendiri.
Sesudah itu ketunggalan, sifat aljabar, teorema apit, barisan monoton, lalu
subbarisan dan Cauchy. Limit superior ditaruh di akhir, sebagai "limit untuk
barisan yang tidak punya limit".

Dua miskonsepsi khas. Pertama, mahasiswa menganggap $N$ boleh bergantung pada $n$,
lalu menulis "untuk setiap $n$ ada $N=n$" dan menyangka definisi terpenuhi;
gejalanya tampak ketika mereka tidak dapat menghasilkan satu $N$ untuk seluruh
ekor. Kedua, mereka menyamakan "suku makin rapat" dengan "konvergen" dan memakai
kriteria Cauchy tanpa memeriksa ruangnya; pemotongan desimal akar dua di dalam
$\mathbb{Q}$ adalah penawar yang efektif.

Aktivitas papan tulis: minta mahasiswa menuliskan negasi definisi limit, lalu
buktikan divergensi $(-1)^n$ dengan negasi itu, bukan dengan kontradiksi. Kelas
siap lanjut bila mereka dapat menghasilkan $N$ eksplisit tanpa diberi rumus dan
menyebut satu barisan terbatas yang divergen.
\end{untukdosen}

\section{Ringkasan}
\tingkat{dasar}

\begin{ringkasan}
\begin{itemize}
  \item $x_n \to L$ berarti untuk setiap $\varepsilon>0$ ada $N$ sehingga
        $|x_n-L|<\varepsilon$ untuk semua $n>N$. $\varepsilon$ dipilih lebih
        dahulu; $N$ menjawab, dan boleh bergantung pada $\varepsilon$ saja.
  \item Limit tunggal, dan barisan konvergen terbatas. Kebalikannya tidak berlaku:
        keterbatasan hanya syarat perlu.
  \item Limit menghormati $+$, $\times$, dan pembagian dengan penyebut tak nol;
        urutan tidak pernah dilompati, hanya dipertahankan lemah.
  \item Teorema apit mengubah pengurungan menjadi limit, dan lemma $q^n \to 0$
        memberi pembanding yang dipakai berulang. Barisan monoton yang terbatas
        selalu konvergen, dan limitnya supremum atau infimum.
  \item Setiap barisan terbatas punya subbarisan konvergen; barisan real konvergen
        jika dan hanya jika Cauchy, dan kesetaraan itu adalah kelengkapan.
  \item $\liminf \le \limsup$, dan keduanya sama dengan $L$ tepat ketika barisan
        konvergen ke $L$.
\end{itemize}
\end{ringkasan}

\section{Latihan}
\tingkat{menengah}

\begin{latihan}
\begin{enumerate}
  \item Buktikan dengan definisi bahwa $x_n=(2n-1)/(3n+4)$ konvergen ke $2/3$, dan
        berikan $N$ eksplisit sebagai fungsi $\varepsilon$. Bandingkan dengan
        ambang $n$ terkecil yang memenuhi.
  \item Buktikan: jika $x_n \to L$ maka $|x_n| \to |L|$. Tunjukkan dengan contoh
        bahwa kebalikannya salah.
  \item Hitung $\limsup$ dan $\liminf$ dari $x_n=1+(-1)^n+1/n$ dan dari
        $y_n=(-1)^n(1+1/n)$. Periksa jawabanmu dengan Python untuk $n$ sampai
        $1000$.
  \item Nilai klaim ini: "Karena $(1+1/n)^n$ naik dan terbatas di atas oleh $3$,
        limitnya adalah $3$." Benar, salah, atau belum cukup bukti? Sebutkan
        bagian mana yang benar dan bagian mana yang keliru.
  \item Buktikan bahwa $(2^n+3^n)^{1/n} \to 3$ dengan menuliskan kedua pagar
        secara lengkap, termasuk bukti $2^{1/n} \to 1$.
  \item Untuk rekursi $x_{n+1}=\sqrt{2+x_n}$ dengan $x_1=1$, hitung dengan Python
        suku terkecil dengan galat di bawah $10^{-6}$ dan $10^{-12}$. Bandingkan
        dengan tebakan dari nisbah $1/4$.
  \item Buktikan bahwa barisan Cauchy terbatas, tanpa memakai Bolzano--Weierstrass.
        Lalu buktikan bahwa barisan monoton yang terbatas adalah Cauchy.
  \item Untuk $x_n=(3n+1)/(n+2)$ dan $\varepsilon=10^{-3}$, tentukan $N$ terkecil
        yang sah, lalu hitung galat pada $n=N$, $N+1$, dan $N+2$ dengan Python.
        Jelaskan mengapa $N-1$ tidak sah.
  \item Buktikan: jika setiap subbarisan dari $(x_n)$ mempunyai subbarisan yang
        konvergen ke $L$, maka $x_n \to L$.
\end{enumerate}

Petunjuk untuk tiga butir tersulit. Butir 5: gunakan
$1 \le (1+(2/3)^n)^{1/n} \le 2^{1/n}$ lalu lemma $2^{1/n} \to 1$ lewat
ketaksamaan Bernoulli. Butir 7: untuk bagian kedua, bandingkan $|x_n-x_m|$ dengan
$x_m-x_n$ ketika $m>n$ dan gunakan kekonvergenan monoton. Butir 9: andaikan
$x_n$ tidak menuju $L$, bentuk subbarisan yang berjarak tetap dari $L$, lalu pakai
Bolzano--Weierstrass.
\end{latihan}

\begin{untukpeneliti}
Yang benar-benar terbukti di sini hanyalah kesetaraan berikut, dan semuanya
bersandar pada satu aksioma: sifat supremum. Dalam lapangan terurut, pernyataan
berikut ekuivalen: setiap himpunan tak kosong yang terbatas di atas punya
supremum; setiap barisan naik yang terbatas konvergen; setiap barisan Cauchy
konvergen; setiap barisan terbatas punya subbarisan konvergen. Rantai buktinya ada
di Rudin, \emph{Principles of Mathematical Analysis}, Bab 3 (bagian Convergent
Sequences, Subsequences, Cauchy Sequences) dan Abbott, \emph{Understanding
Analysis}, Bab 2 (bagian 2.2 sampai 2.6). Yang populer tetapi tidak terbukti:
"barisan Cauchy selalu konvergen". Itu salah tanpa kelengkapan, dan contoh
pemotongan desimal akar dua di dalam $\mathbb{Q}$ adalah penyangkalnya.

Pertanyaan yang masih terbuka: seberapa cepat sebuah barisan konvergen, dan apakah
laju itu selalu dapat dihitung. Jawabannya tidak. Hasil klasik Specker (1949)
memberi barisan rasional yang dapat dihitung, naik, dan terbatas, tetapi limitnya
tidak dapat dihitung; jadi "konvergen" tidak otomatis memberi batas galat yang bisa
diprogram.

Batas keberlakuan: seluruh bab ini bekerja di $\mathbb{R}$ dengan jarak biasa.
Di ruang metrik umum, Bolzano--Weierstrass berubah menjadi pernyataan tentang
kekompakan, dan kriteria Cauchy tetap menjadi syarat perlu tetapi belum cukup.
Teorema apit juga diam bila kedua pagar tidak menuju nilai yang sama; di situ kita
butuh pembanding yang lebih halus, dan itulah bahan Bab 5.
\end{untukpeneliti}

\peta{Bab 4 pada Rudin, \emph{Principles of Mathematical Analysis}: Bab 3 (Numerical
Sequences and Series, dengan bagian Convergent Sequences, Subsequences, Cauchy
Sequences, Upper and Lower Limits, dan Some Special Sequences). Pada Abbott,
\emph{Understanding Analysis}: Bab 2 (The Limit of a Sequence, The Algebraic and
Order Limit Theorems, The Monotone Convergence Theorem and a First Look at Infinite
Series, Subsequences and the Bolzano--Weierstrass Theorem, The Cauchy Criterion).
Aksioma dan akibat kelengkapan yang dipakai di sini dibangun di Bab 3 buku ini,
bukan diulang.}

\chapter{Deret Tak Hingga}

Sebuah tim melatih pengklasifikasi pada satu juta contoh. Fungsi objektifnya jumlah
kerugian tiap contoh, dan pelatihan berhenti otomatis ketika kenaikan total turun di
bawah $10^{-7}$ per epoch. Pada epoch keempat puluh proses berhenti sendiri. Modelnya
tetap buruk: setiap contoh baru hanya menyumbang sekitar $10^{-6}$ pada total, jadi
total hampir tidak bergerak, sementara total itu sendiri belum tentu menuju nilai
berhingga.

Kecurigaan itu bisa diuji dengan deret paling sederhana. Jumlah
$H_n = 1 + \tfrac{1}{2} + \cdots + \tfrac{1}{n}$ pada $n=10^6$ bernilai
$14.3927267229$. Menambah satu juta suku lagi hanya menaikkannya menjadi
$15.0858736534$; tambahannya $0.6931469306$, hampir tepat $\ln 2 = 0.6931471806$.
Tiap kali banyaknya suku berlipat dua, deret harmonik naik sekitar $\ln 2$, tanpa
batas: taksiran $\ln n + \gamma$ memberi $H_{10^9} \approx 21.3004815$. Aturan
berhenti berbasis "total sudah tidak bergerak" karena itu tidak membedakan konvergen
dari divergen; yang membedakan adalah laju pengecilan suku.

Bab 4 sudah memberi limit, kriteria Cauchy, kekonvergenan monoton, dan limit superior.
Deret tidak menambah gagasan baru, sebab jumlah parsial hanyalah barisan. Tetapi jumlah
parsial jarang punya rumus tertutup, dan menghitung $s_{10^6}$ tidak memberi tahu apakah
$s_{10^7}$ masih bergerak. Bab ini mengubah pertanyaan itu menjadi pertanyaan yang bisa
dijawab dengan memeriksa $a_n$ saja.

\begin{untukanak}
Bayangkan dua toples recehan. Di toples pertama kamu menabung dengan aturan aneh:
hari pertama satu koin, hari kedua setengah koin, hari ketiga seperempat koin, dan
seterusnya, selalu separuh dari hari sebelumnya. Kamu menabung setiap hari, tanpa
pernah libur, sampai kapan pun. Berapa isi toples itu setelah seratus tahun? Jawabannya
mengejutkan: tidak pernah lebih dari dua koin. Setiap tambahan memang mengecil, dan
jumlahnya berhenti di suatu angka yang bisa kamu sebut dengan pasti.

Di toples kedua aturannya berbeda. Koin hari pertama besar, lalu makin kecil, dan
makin kecil lagi. Setelah sebulan isinya baru segepok. Setelah setahun agak lebih
tebal. Kamu mungkin menyimpulkan toples ini juga akan berhenti di suatu angka, karena
tambahannya makin tipis. Ternyata tidak. Kalau kamu menabung cukup lama, isinya akan
melewati angka berapa pun yang kamu sebut, walaupun untuk naik satu angka saja kamu
butuh bertahun-tahun. Dua toples, dua aturan yang kelihatan mirip, dua nasib yang
berlawanan.

Pertanyaannya: tanpa menunggu bertahun-tahun, bagaimana tahu toples mana yang berhenti
dan mana yang tidak? Adakah ciri pada koinnya sendiri, bukan pada hasil akhirnya? Bab
ini menjawabnya, dan jawabannya dipakai setiap hari oleh orang yang menulis program.
\end{untukanak}

\section{Deret, Jumlah Parsial, dan Deret yang Bisa Dihitung Langsung}
\tingkat{dasar}

Deret\index{deret} adalah penjumlahan yang tidak pernah selesai. Karena itu
definisinya tidak memakai kata "jumlah" secara langsung, melainkan memakai barisan
jumlah sementara\index{jumlah parsial}.
Kita membaca hasil setelah sekian suku, lalu menanyakan ke mana bacaan itu menuju.

\begin{definisi}
Diberikan barisan $(a_n)$ di $\mathbb{R}$. Jumlah parsial ke-$n$ adalah
\[
  s_n = a_1 + a_2 + \cdots + a_n = \sum_{k=1}^{n} a_k .
\]
Lambang $\sum_{n=1}^{\infty} a_n$ disebut deret, dan deret itu \emph{konvergen} ke
$s$ bila barisan $(s_n)$ konvergen ke $s$; dalam hal ini kita menulis
$\sum_{n=1}^{\infty} a_n = s$ dan menyebut $s$ sebagai jumlah deret. Bila $(s_n)$
divergen, deret disebut \emph{divergen}.
\end{definisi}

\kamus{deret}{series}{penjumlahan yang tidak pernah selesai, yang nilainya dibaca dari jumlah sementara setelah sekian suku.}

Perhatikan $s_n - s_{n-1} = a_n$. Jadi deret dan barisan adalah benda yang sama; yang
berubah hanya pertanyaannya: pada barisan rumus $s_n$ diketahui, pada deret hanya $a_n$
yang kita pegang.

Kriteria Cauchy dari Bab 4 bisa langsung diterjemahkan, dan terjemahannya adalah alat
paling dasar di bab ini.

\begin{teorema}[Kriteria Cauchy untuk deret]
$\sum a_n$ konvergen jika dan hanya jika untuk setiap $\varepsilon > 0$ terdapat
bilangan bulat $N$ sehingga
\[
  m \ge n \ge N \quad \Longrightarrow \quad
  \left| \sum_{k=n}^{m} a_k \right| < \varepsilon .
\]
\end{teorema}

\noindent\textbf{Bukti.} Karena $\sum_{k=n}^{m} a_k = s_m - s_{n-1}$, syarat di atas tidak lain adalah syarat
Cauchy untuk barisan $(s_n)$. Menurut Bab 4, barisan real konvergen jika dan hanya
jika ia barisan Cauchy.

Kriteria ini bicara tentang potongan ekor deret, bukan tentang suku tunggal. Dengan
mengambil $m = n$ kita memperoleh syarat perlu yang sangat murah diperiksa.

\begin{akibat}[Uji suku ke-$n$]
Jika $\sum a_n$ konvergen, maka $a_n \to 0$.
\end{akibat}

\noindent\textbf{Bukti.} Ambil $m=n$ pada kriteria Cauchy: untuk setiap $\varepsilon>0$ terdapat $N$ sehingga
$n \ge N$ berakibat $|a_n| < \varepsilon$. Itu persis definisi $a_n \to 0$.

Uji ini hanya menyaring. Ia mengatakan bahwa suku yang tidak menuju nol pasti
menggagalkan konvergensi, dan tidak mengatakan apa pun tentang suku yang menuju nol.
Dua contoh cepat untuk sisi yang mudah: $\sum \frac{n}{n+1}$ divergen karena
$\frac{n}{n+1} \to 1$, dan $\sum (-1)^n$ divergen karena sukunya tidak menuju nol.
Sisi yang sulit adalah deret harmonik: $\frac{1}{n} \to 0$ tetapi $\sum \frac{1}{n}$
divergen. Kita buktikan itu di bagian kedua.

\begin{bacakeras}
Dua suara. Bacakan bergantian, lalu tukar perannya.

\textbf{Suara A:} Saya mengecil. Setiap hari lebih kecil daripada kemarin.

\textbf{Suara B:} Bagus. Tetapi mengecil belum berarti berhenti.

\textbf{A:} Saya mengecil separuh. Separuh dari separuh. Separuh lagi.

\textbf{B:} Maka kamu punya pagar. Jumlahmu akan berhenti.

\textbf{A:} Saya mengecil, tetapi hanya sedikit. Sepertiga. Seperempat. Seperlima.

\textbf{B:} Awas. Tipis bukan berarti habis.

\textbf{A:} Berapa pun kecilnya tambahan saya, saya tetap menambah.

\textbf{B:} Dan berapa pun lamanya, jumlahmu tetap naik.

Ulangi tiga kali makin cepat: mengecil itu urusan suku, berhenti itu urusan jumlah.
Hafalkan: yang menentukan bukan suku yang mengecil, melainkan sisa yang tidak pernah
habis.
\end{bacakeras}

Deret pertama yang bisa dihitung langsung adalah deret geometri\index{deret geometri}. Setiap sukunya
diperoleh dengan mengalikan suku sebelumnya dengan bilangan tetap.

\begin{teorema}[Deret geometri]
Untuk $|r| < 1$,
\[
  \sum_{k=0}^{\infty} r^k = \frac{1}{1-r},
\]
dan untuk $|r| \ge 1$ deret itu divergen.
\end{teorema}

\noindent\textbf{Bukti.} Untuk $r \neq 1$ berlaku identitas
\[
  (1-r)\left(1 + r + \cdots + r^n\right) = 1 - r^{n+1},
\]
sehingga $s_n = \frac{1-r^{n+1}}{1-r}$. Bila $|r|<1$ maka $r^{n+1} \to 0$ (Bab 4,
karena $|r|^{n} \to 0$ untuk $0 \le |r| < 1$), jadi $s_n \to \frac{1}{1-r}$. Bila
$|r| \ge 1$ maka $|r|^k \ge 1$ untuk semua $k$, sehingga $s_n$ tidak mungkin konvergen:
untuk $r=1$ jelas $s_n = n+1$, dan untuk $|r|\ge 1$ dengan $r\neq 1$ suku $r^k$ tidak
menuju nol.

Dengan $r = \tfrac{1}{2}$ kita memperoleh $1 + \tfrac{1}{2} + \tfrac{1}{4} + \cdots = 2$,
dan dengan $r = -\tfrac{1}{2}$ kita memperoleh $1 - \tfrac{1}{2} + \tfrac{1}{4} - \cdots
= \tfrac{2}{3}$. Yang kedua menarik: suku-sukunya berganti tanda dan jumlahnya lebih
kecil daripada suku pertama dikali dua, karena tanda yang berlawanan saling menghapus.

Kasus $r=\tfrac{1}{2}$ punya gambar yang menjelaskan dirinya sendiri: potong persegi
satuan menjadi dua, potong lagi sisanya, lalu teruskan. Gambar~\ref{fig:potong-geometri}
menunjukkan potongan $\tfrac12$, $\tfrac14$, $\tfrac18$ dan sisa $2^{-5} = 0.03125$
setelah lima potong, dengan luas terpakai $1 - 2^{-5} = 0.96875$. Sisa itu menyusut
tanpa lenyap; itulah rupa geometri dari kalimat "jumlahnya berhenti".

\begin{figure}[htbp]
\centering
\begin{tikzpicture}[scale=6.0]
  \draw[fill=layerteal!45,draw=ink] (0.5,0) rectangle (1,1);
  \draw[fill=layerindigo!40,draw=ink] (0,0.5) rectangle (0.5,1);
  \draw[fill=layerorange!50,draw=ink] (0.25,0) rectangle (0.5,0.5);
  \draw[fill=layerjade!45,draw=ink] (0,0.25) rectangle (0.25,0.5);
  \draw[fill=layerplum!40,draw=ink] (0.125,0) rectangle (0.25,0.25);
  \draw[fill=white,draw=ink,densely dashed] (0,0) rectangle (0.125,0.25);
  \draw[ink,thick] (0,0) rectangle (1,1);
  \node[font=\small] at (0.75,0.5) {$\tfrac{1}{2}$};
  \node[font=\small] at (0.25,0.75) {$\tfrac{1}{4}$};
  \node[font=\small] at (0.375,0.25) {$\tfrac{1}{8}$};
  \node[font=\small] at (0.125,0.375) {$\tfrac{1}{16}$};
  \draw[ink,thin,-{Stealth}] (0.06,0.10) -- (-0.30,-0.05);
  \node[font=\scriptsize,anchor=east] at (-0.30,-0.05) {sisa $\tfrac{1}{32}$};
\end{tikzpicture}
\caption{Deret geometri separuh sebagai potongan persegi satuan. Potongan berukuran
$\tfrac12$, $\tfrac14$, $\tfrac18$, $\tfrac1{16}$ dipasang berurutan, dan sisa yang
belum tertutup setelah empat potong adalah $\tfrac1{16}$; setelah lima potong sisanya
$2^{-5} = 0.03125$ dan luas terpakai $1 - 2^{-5} = 0.96875$. Sisa selalu menyusut,
tidak pernah lenyap.}
\label{fig:potong-geometri}
\end{figure}

\begin{tangan}
Bahan: satu lembar kertas persegi, pensil, penggaris.

\begin{enumerate}
\item Lipat kertas menjadi dua, buka, arsir separuhnya. Itu potongan pertama.
\item Lipat lagi bagian yang belum diarsir menjadi dua, buka, arsir separuhnya.
\item Ulangi tiga kali lagi, selalu hanya pada bagian yang belum diarsir.
\item Periksa bagian putih yang tersisa: setelah lima lipatan tinggal sepertiga puluh
      dua kertas semula.
\end{enumerate}

Yang harus kamu lihat: setiap arsiran baru selalu separuh dari yang sebelumnya, dan
bagian putih tidak pernah lenyap, hanya menguncup ke satu sudut. Kalau kertasnya cukup
tipis untuk dilipat sepuluh kali, bagian putih tinggal $\tfrac{1}{1024} = 0.0009765625$;
menurut hitungan ia menuju nol, sehingga arsiran akhirnya menutup seluruh kertas.
Sekarang ganti aturannya: arsir $\tfrac12$ kertas, lalu $\tfrac13$, lalu $\tfrac14$.
Lipatan tidak lagi bisa mengikuti, dan kertas habis sebelum deretnya berhenti.
\end{tangan}

Keluarga kedua yang bisa dihitung langsung adalah deret teleskopik\index{deret teleskopik}, yaitu deret yang
sukunya dapat ditulis sebagai selisih dua suku berurutan dari barisan lain sehingga
hampir semua pengurangan saling membatalkan.

\begin{contoh}
Untuk $n \ge 1$,
\[
  \frac{1}{n(n+1)} = \frac{1}{n} - \frac{1}{n+1}.
\]
Karena itu
\[
  s_N = \sum_{n=1}^{N} \left(\frac{1}{n} - \frac{1}{n+1}\right)
      = 1 - \frac{1}{N+1},
\]
sebab semua suku antara saling menghapus dan hanya suku pertama serta suku terakhir
yang tersisa. Jadi $\sum_{n=1}^{\infty} \frac{1}{n(n+1)} = 1$. Pada $N=100$ rumus itu
memberi $s_{100} = \tfrac{100}{101} = 0.9900990099$, dan sisa ekornya adalah
$\tfrac{1}{101} = 0.0099009901$; keduanya berjumlah tepat $1$.
\end{contoh}

Perhatikan bahwa pada contoh teleskopik kita tidak perlu uji apa pun. Rumus tertutup
untuk $s_N$ langsung memberi limit. Inilah yang jarang terjadi. Untuk hampir semua
deret, termasuk deret harmonik dan $\sum 1/n^2$, rumus seperti itu tidak ada, dan kita
harus memutuskan konvergensi tanpa pernah mengetahui jumlahnya secara tepat.

Sebelum masuk ke uji-uji, kita catat bahwa deret berperilaku linear, seperti halnya
limit barisan.

\begin{proposisi}
Jika $\sum a_n = A$ dan $\sum b_n = B$ konvergen, maka untuk setiap $c, d \in
\mathbb{R}$ deret $\sum (c a_n + d b_n)$ konvergen dengan jumlah $cA + dB$.
Menambahkan atau membuang berhingga banyak suku tidak mengubah konvergensi, hanya
mengubah jumlahnya.
\end{proposisi}

\noindent\textbf{Bukti.} Jumlah parsial deret gabungan adalah $c s_n^{(a)} + d s_n^{(b)}$, dan menurut sifat
aljabar limit di Bab 4 barisan itu konvergen ke $cA + dB$. Untuk pernyataan kedua,
misalkan suku dibuang mulai indeks $M$. Jumlah parsial deret baru adalah
$s_{N} - s_{M-1}$ untuk suatu $N$, yaitu selisih barisan konvergen dengan konstanta,
sehingga kekonvergenannya sama.

Sebagai pemakaian, $\sum_{n=1}^{\infty} \left(\frac{1}{2^n} +
\frac{1}{n(n+1)}\right) = 1 + 1 = 2$, karena bagian pertama adalah deret geometri
dengan $r=\tfrac{1}{2}$ yang dimulai dari $n=1$.

\section{Deret Suku Tak Negatif: Pembanding dan Deret-$p$}
\tingkat{menengah}

Bila semua suku tidak negatif, jumlah parsial hanya bisa naik atau diam. Pertanyaan
konvergen lalu berubah menjadi pertanyaan yang jauh lebih sederhana: apakah jumlah
parsial berhenti di bawah suatu pagar?

\begin{teorema}
Jika $a_n \ge 0$ untuk semua $n$, maka $\sum a_n$ konvergen jika dan hanya jika
barisan jumlah parsialnya terbatas.
\end{teorema}

\noindent\textbf{Bukti.} Dengan $a_n \ge 0$ berlaku $s_{n+1} = s_n + a_{n+1} \ge s_n$, jadi $(s_n)$ monoton
tak menurun. Menurut teorema kekonvergenan monoton di Bab 4, barisan monoton yang
terbatas di atas konvergen, dan barisan monoton yang tidak terbatas di atas divergen
menuju $+\infty$.

Teorema ini terlihat sederhana, tetapi seluruh bagian ini hanya variasi darinya. Yang
kita butuhkan hanyalah cara membandingkan deret yang belum dikenal dengan deret yang
sudah dikenal, dan persediaan deret yang sudah dikenal masih tipis: deret geometri.

\begin{teorema}[Uji pembanding]
Misalkan $0 \le a_n \le c_n$ untuk $n \ge N_0$.
\begin{enumerate}
\item Jika $\sum c_n$ konvergen, maka $\sum a_n$ konvergen.
\item Jika $a_n \ge d_n \ge 0$ untuk $n \ge N_0$ dan $\sum d_n$ divergen, maka
      $\sum a_n$ divergen.
\end{enumerate}
\end{teorema}

\noindent\textbf{Bukti.} Untuk (1), ambil $\varepsilon>0$. Menurut kriteria Cauchy ada $N \ge N_0$ sehingga
$m \ge n \ge N$ berakibat $\sum_{k=n}^{m} c_k < \varepsilon$. Karena setiap suku tidak
negatif dan $a_k \le c_k$,
\[
  \left| \sum_{k=n}^{m} a_k \right| = \sum_{k=n}^{m} a_k \le \sum_{k=n}^{m} c_k
  < \varepsilon ,
\]
jadi $\sum a_n$ memenuhi kriteria Cauchy. Untuk (2), bila $\sum a_n$ konvergen maka
dengan (1) dan $d_n \le a_n$ deret $\sum d_n$ ikut konvergen, bertentangan dengan
hipotesis.

Uji pembanding\index{uji pembanding} menuntut kita memilih pembanding yang tepat. Untuk suku berbentuk
pecahan rasional, pembanding yang wajar adalah suku terbesar di pembilang dibagi suku
terbesar di penyebut. Misalnya $\frac{1}{n^2 + 3n} \le \frac{1}{n^2}$, sehingga
$\sum \frac{1}{n^2+3n}$ konvergen asalkan $\sum \frac{1}{n^2}$ konvergen. Karena itu
seluruh pertanyaan berpindah ke satu keluarga deret.

Keluarga itu adalah deret-$p$\index{deret-p}, dan alat untuk menanganinya adalah pengelompokan suku
menjadi blok berukuran pangkat dua. Gagasan pengelompokan ini dipakai berulang kali,
termasuk pada uji kondensasi\index{kondensasi Cauchy} berikut.

\begin{lemma}[Kondensasi Cauchy]
Misalkan $a_1 \ge a_2 \ge a_3 \ge \cdots \ge 0$. Maka $\sum_{n=1}^{\infty} a_n$
konvergen jika dan hanya jika $\sum_{k=0}^{\infty} 2^k a_{2^k}$ konvergen.
\end{lemma}

\noindent\textbf{Bukti.} Karena sukunya tak negatif, kedua deret konvergen tepat ketika jumlah parsialnya
terbatas (teorema pertama bagian ini). Kita bandingkan blok-blok. Untuk $k \ge 0$,
blok $a_{2^k+1} + \cdots + a_{2^{k+1}}$ memuat $2^k$ suku, masing-masing tidak lebih
besar dari $a_{2^k}$ dan tidak lebih kecil dari $a_{2^{k+1}}$. Jadi
\[
  \sum_{n=1}^{2^m} a_n \;\le\; a_1 + \sum_{k=0}^{m-1} 2^k a_{2^k},
  \qquad
  \sum_{n=1}^{2^m} a_n \;\ge\; \frac{1}{2}\sum_{k=1}^{m} 2^k a_{2^k}.
\]
Pertidaksamaan kiri menunjukkan bahwa jumlah parsial $\sum a_n$ terbatas bila
$\sum 2^k a_{2^k}$ terbatas. Pertidaksamaan kanan menunjukkan sebaliknya, karena
jumlah parsial $\sum 2^k a_{2^k}$ terbatas bila $\sum a_n$ terbatas. Kedua arah
selesai.

\begin{teorema}[Deret-$p$]
$\sum_{n=1}^{\infty} \frac{1}{n^p}$ konvergen jika $p > 1$ dan divergen jika
$p \le 1$.
\end{teorema}

\noindent\textbf{Bukti.} Untuk $p \le 0$ suku $\frac{1}{n^p} = n^{-p}$ tidak menuju nol, jadi deret divergen
menurut uji suku ke-$n$. Untuk $p>0$ barisan $a_n = n^{-p}$ menurun menuju nol,
sehingga lemma kondensasi berlaku dan
\[
  \sum_{k=0}^{\infty} 2^k a_{2^k}
  = \sum_{k=0}^{\infty} 2^k \cdot 2^{-kp}
  = \sum_{k=0}^{\infty} \left(2^{1-p}\right)^{k}.
\]
Itu deret geometri dengan rasio $r = 2^{1-p}$. Karena $p>1$ tepat ketika $1-p<0$,
yaitu tepat ketika $r<1$, teorema deret geometri memberi: konvergen untuk $p>1$,
divergen untuk $p \le 1$.

Kasus $p=1$ adalah deret harmonik\index{deret harmonik}, dan buktinya sekarang hanya
satu baris: rasio
geometrinya $2^{1-1} = 1$, sehingga deret kondensasinya $\sum 1$ divergen. Kita bisa
membaca kerapuhannya langsung dari blok: setiap blok berukuran $2^k$ menyumbang paling
sedikit $2^k \cdot \frac{1}{2^{k+1}} = \tfrac{1}{2}$. Karena itu
\[
  H_{2^m} \;\ge\; 1 + \frac{m}{2},
\]
sehingga $H_n \to \infty$, tetapi untuk menaikkan $H$ sebesar satu perlu menggandakan
$n$ berkali-kali: laju pertumbuhannya logaritmik. Inilah alasan angka $14.39$ pada
pembukaan bab ini terlihat seperti batas padahal bukan.

Gambar~\ref{fig:batang-suku} menaruh kedua jenis suku berdampingan. Batang harmonik
tetap gemuk: suku kesepuluhnya masih $0.1$, sedangkan suku geometri kesepuluh
$2^{-10} = 0.0009765625$. Jumlah sepuluh suku pertama: $H_{10} = 2.928968254$ melawan
$1 - 2^{-10} = 0.9990234375$. Suku harmonik lebih besar pada setiap indeks, jadi besar
suku menyesatkan; yang menentukan adalah apakah pengecilannya meninggalkan ekor
berhingga.

\begin{figure}[htbp]
\centering
\begin{tikzpicture}
\begin{axis}[
  width=0.93\linewidth, height=6.0cm,
  ybar, bar width=3.4pt,
  xlabel={indeks $n$}, ylabel={besar suku $a_n$},
  xmin=0.3, xmax=10.7, ymin=0, ymax=1.12,
  xtick={1,2,3,4,5,6,7,8,9,10},
  scaled y ticks=false,
  tick label style={font=\scriptsize},
  label style={font=\small},
  ymajorgrids, grid style={layergrey!20},
  legend style={at={(0.98,0.96)},anchor=north east,draw=none,fill=none,
                font=\scriptsize},
]
\addplot[fill=layerteal!85,draw=layerteal] coordinates {
  (1,1) (2,0.5) (3,0.3333333) (4,0.25) (5,0.2)
  (6,0.1666667) (7,0.1428571) (8,0.125) (9,0.1111111) (10,0.1)};
\addlegendentry{harmonik $1/n$}
\addplot[fill=layerorange!85,draw=layerorange] coordinates {
  (1,0.5) (2,0.25) (3,0.125) (4,0.0625) (5,0.03125)
  (6,0.015625) (7,0.0078125) (8,0.0039063) (9,0.0019531) (10,0.0009766)};
\addlegendentry{geometri $2^{-n}$}
\end{axis}
\end{tikzpicture}
\caption{Batang suku harmonik $1/n$ dan geometri $2^{-n}$ untuk $n=1,\ldots,10$.
Batang harmonik lebih tinggi pada setiap indeks dan suku kesepuluhnya masih $0.1$,
sedangkan suku geometri kesepuluh $2^{-10} = 0.0009765625$. Jumlah sepuluh suku
pertama: $H_{10} = 2.928968254$ melawan $1 - 2^{-10} = 0.9990234375$.}
\label{fig:batang-suku}
\end{figure}

\kamus{deret harmonik}{harmonic series}{tumpukan yang terus bertambah tetapi tiap tambahannya makin tipis, sehingga tinggi totalnya naik sangat lambat tanpa batas.}

Deret-$p$ bukan batas akhir. Dengan kondensasi yang sama, Rudin menunjukkan bahwa
$\sum \frac{1}{n(\log n)^p}$ konvergen untuk $p>1$ dan divergen untuk $p<1$, dan bahwa
menambah satu logaritma lagi menggeser batasnya. Deret-deret itu berbeda sangat sedikit
tetapi berada di sisi berlawanan: tidak ada garis batas yang rapi, yang ada hanyalah uji
yang berlaku pada kelas tertentu.

\begin{tantangan}
Ambil $k(n)$ bilangan bulat dengan $2^{k-1} \le n < 2^k$, lalu tetapkan
$a_n = \frac{1}{k(n)}$. Jadi $a_1 = 1$, $a_2 = a_3 = \tfrac12$,
$a_4 = \cdots = a_7 = \tfrac13$, dan seterusnya. Tunjukkan bahwa $a_n \to 0$ tetapi
$\sum a_n$ divergen, lalu jelaskan mengapa uji rasio memberi limit superior $1$ di
sini sehingga tidak berguna.

\textbf{Petunjuk 1.} Jangan menjumlahkan suku satu per satu. Jumlahkan satu blok, dan
hitung dulu ada berapa suku di dalam blok itu.

\textbf{Petunjuk 2.} Blok ke-$k$ memuat $2^{k-1}$ suku yang semuanya bernilai
$\frac{1}{k}$, jadi sumbangannya $\frac{2^{k-1}}{k}$. Sepuluh blok pertama sudah
berjumlah $118.653968$.

Kalau mentok, kembali ke bagian Deret Suku Tak Negatif, tempat blok berukuran pangkat
dua dipakai pada Lemma kondensasi Cauchy.
\end{tantangan}

\begin{untukanak}
Coba bayangkan menumpuk koin di atas meja. Koin pertama tebal. Koin kedua setengah
tebalnya. Koin ketiga sepertiga, koin keempat seperempat, dan seterusnya: koin
ke-sejuta sangat tipis, setipis sepersejuta koin pertama. Kalau kamu menumpuk sejuta
koin, tinggi tumpukan itu kira-kira empat belas setebal koin pertama, bukan sejuta
kali. Kelihatannya tumpukan itu akan berhenti tumbuh, karena koin berikutnya hampir
tidak menambah apa pun.

Sekarang pertanyaan yang membuat banyak orang salah: berapa koin agar tumpukan
mencapai dua puluh kali tebal koin pertama? Jawabannya sekitar dua ratus tujuh puluh
dua juta koin. Untuk dua puluh satu kali, perlu sekitar sembilan miliar koin. Setiap
tingkat menuntut koin berkali-kali lipat, tetapi tidak ada tingkat yang mustahil.
Tumpukan itu tidak pernah berhenti, hanya melambat.

Inilah sebabnya program komputer mudah tertipu. Kalau kamu berhenti menumpuk setelah
sejuta koin, tumpukan tampak tidak bergerak lagi, padahal ia belum selesai. Di bab ini
kita akan tahu mengapa lambat bukan berarti berhenti, dan bagaimana membedakannya
sebelum menghabiskan waktu menumpuk.
\end{untukanak}

\section{Uji Rasio dan Uji Akar}
\tingkat{menengah}

Uji pembanding menuntut kita sudah punya deret pembanding. Untuk suku yang memuat
faktorial atau pangkat $n$, pembanding seperti itu sulit ditebak, tetapi perbandingan
antar suku berurutan mudah dihitung. Gagasannya: bila setiap suku paling banyak
$\beta$ kali suku sebelumnya dengan $\beta<1$, deretnya mengecil paling lambat seperti
deret geometri, jadi konvergen.

\begin{teorema}[Uji rasio]
Misalkan $a_n \neq 0$ untuk semua $n$ yang cukup besar.
\begin{enumerate}
\item Jika $\limsup_{n\to\infty} \left| \frac{a_{n+1}}{a_n} \right| < 1$, maka
      $\sum a_n$ konvergen, bahkan konvergen mutlak.
\item Jika terdapat $n_0$ sehingga $\left| \frac{a_{n+1}}{a_n} \right| \ge 1$ untuk
      semua $n \ge n_0$, maka $\sum a_n$ divergen.
\end{enumerate}
\end{teorema}

\noindent\textbf{Bukti.} Untuk (1), tulis $L = \limsup |a_{n+1}/a_n|$ dan pilih $\beta$ dengan $L < \beta < 1$.
Menurut sifat limit superior di Bab 4, hanya berhingga banyak $n$ yang memenuhi
$|a_{n+1}/a_n| > \beta$; jadi ada $N$ sehingga $n \ge N$ berakibat
$|a_{n+1}| \le \beta |a_n|$. Dengan induksi,
\[
  |a_{N+k}| \le \beta^{k} |a_N| \qquad (k \ge 0).
\]
Deret $\sum \beta^k |a_N|$ konvergen karena $0<\beta<1$, jadi menurut uji pembanding
$\sum |a_n|$ konvergen, dan menurut teorema konvergensi mutlak di bagian kelima bab ini,
kekonvergenan mutlak mengimplikasikan kekonvergenan.

Untuk (2), syarat $|a_{n+1}| \ge |a_n|$ untuk $n \ge n_0$ membuat $(|a_n|)$ tak
menurun dan tak negatif, sehingga tidak mungkin menuju nol. Maka $a_n$ tidak menuju
nol, dan uji suku ke-$n$ menggugurkan kekonvergenan.

\kamus{uji rasio}{ratio test}{mengukur apakah tiap suku menyusut paling sedikit sebesar sebagian tetap dari suku sebelumnya, sehingga deretnya mengecil seperti deret geometri.}
\index{uji rasio}

Batas $1$ pada uji rasio adalah titik buta. Untuk $1/n$ rasionya $\frac{n}{n+1} \to 1$
dan deretnya divergen; untuk $1/n^2$ rasionya $\left(\frac{n}{n+1}\right)^2 \to 1$
dan deretnya konvergen. Jadi rasio menuju $1$ tidak memberi informasi; orde pengecilan
yang berbeda jauh bisa memberi rasio yang sama.

Dua deret berikut memperlihatkan uji rasio bekerja pada arah yang berlawanan.
Untuk $a_n = \frac{n!}{10^n}$ berlaku
$\frac{a_{n+1}}{a_n} = \frac{n+1}{10}$, yang menuju $+\infty$; deretnya divergen.
Untuk $a_n = \frac{10^n}{n!}$ berlaku $\frac{a_{n+1}}{a_n} = \frac{10}{n+1}$, yang
menuju $0$; deretnya konvergen. Jumlah deret kedua tepat $e^{10} = 22026.4657948$,
sedangkan jumlah $20$ suku pertamanya baru $21950.3788494$.

Uji akar\index{uji akar} memakai besaran lain: akar ke-$n$ dari suku ke-$n$. Secara intuitif ia
mengukur "rasio geometri rata-rata" dari seluruh suku, bukan rasio antar dua suku
berurutan.

\begin{teorema}[Uji akar]
\begin{enumerate}
\item Jika $\limsup_{n\to\infty} \sqrt[n]{|a_n|} < 1$, maka $\sum a_n$ konvergen,
      bahkan konvergen mutlak.
\item Jika $\limsup_{n\to\infty} \sqrt[n]{|a_n|} > 1$, maka $\sum a_n$ divergen.
\end{enumerate}
\end{teorema}

\noindent\textbf{Bukti.} Untuk (1), pilih $\beta$ dengan $\limsup \sqrt[n]{|a_n|} < \beta < 1$. Hanya
berhingga banyak $n$ dengan $\sqrt[n]{|a_n|} > \beta$, jadi ada $N$ sehingga
$|a_n| \le \beta^n$ untuk $n \ge N$. Uji pembanding dengan deret geometri
$\sum \beta^n$ menutup kasus ini.

Untuk (2), bila limit superiornya lebih besar dari $1$, maka
$\sqrt[n]{|a_n|} > 1$, yaitu $|a_n| > 1$, untuk tak berhingga banyak $n$. Suku-suku
itu tidak menuju nol, jadi deretnya divergen menurut uji suku ke-$n$.

Uji akar lebih kuat daripada uji rasio. Misalkan $a_n = 2^{-n}$ untuk $n$ genap dan
$a_n = 3^{-n}$ untuk $n$ ganjil. Maka $\sqrt[n]{a_n}$ bernilai $\tfrac{1}{2}$ atau
$\tfrac{1}{3}$, sehingga limit superiornya $\tfrac{1}{2}$ dan uji akar menyatakan
deretnya konvergen. Rasionya berayun: pada $n=100$, $\frac{a_{101}}{a_{100}}$ bernilai
sekitar $3.05 \times 10^{17}$, dengan limit superior $+\infty$ dan limit inferior $0$.
Uji rasio tidak menyimpulkan apa pun, tepat ketika sukunya berayun dalam skala berbeda.

Uji rasio tetap populer karena faktorial dan pangkat $n$ saling membatalkan pada
rasio, sedangkan akar ke-$n$ dari faktorial menuntut aproksimasi Stirling. Urutan
pemakaian: uji suku ke-$n$, pembanding dengan deret-$p$ atau geometri, lalu uji rasio;
uji akar dipakai bila rasio berayun atau memberi $1$.

Sekarang uji-uji itu bertumpuk, dan urutan pemeriksaan menjadi keterampilan tersendiri.

\begin{resep}
\begin{enumerate}
\item Periksa $a_n \to 0$. Bila tidak, deret divergen; selesai.
\item Kenali bentuknya: faktorial memanggil uji rasio, pangkat $n$ memanggil uji akar,
      pecahan rasional memanggil pembanding deret-$p$.
\item Bila uji pilihanmu memberi batas $1$, jangan menyimpulkan apa pun. Pindah ke
      pembanding atau kondensasi.
\item Bila tanda berganti dan $|a_n|$ menurun ke nol, pakai Leibniz dan laporkan batas
      galat $b_{N+1}$.
\end{enumerate}

Peringatan urutan: langkah 1 jangan ditukar dengan langkah 2. Uji rasio pada deret yang
sukunya tidak menuju nol bisa memberi batas $1$ dan menyesatkan, sedangkan uji suku
ke-$n$ menyelesaikannya tanpa hitungan.
\end{resep}

\section{Deret Berganti Tanda dan Uji Leibniz}
\tingkat{menengah}

Sampai di sini semua deret bersuku tak negatif, dan di sana konvergensi berarti
jumlah parsial berhenti naik. Begitu tanda berganti-ganti, jumlah parsial bisa naik
dan turun, dan pembatalan antar tanda dapat menghasilkan konvergensi yang tidak
dijamin oleh besarnya suku.

\begin{teorema}[Uji Leibniz]
Misalkan $(b_n)$ barisan tak negatif yang tak menaik, yaitu $b_1 \ge b_2 \ge \cdots
\ge 0$, dan $b_n \to 0$. Maka $\sum_{n=1}^{\infty} (-1)^{n+1} b_n$ konvergen.
\end{teorema}

\noindent\textbf{Bukti.} Tulis $s_N$ untuk jumlah parsial ke-$N$. Pandang jumlah parsial berindeks genap dan
berindeks ganjil secara terpisah. Karena
\[
  s_{2m+2} - s_{2m} = b_{2m+1} - b_{2m+2} \ge 0,
\]
barisan $(s_{2m})$ tak menurun. Karena
\[
  s_{2m+1} - s_{2m-1} = -b_{2m} + b_{2m+1} \le 0,
\]
barisan $(s_{2m-1})$ tak menaik. Selain itu $s_{2m} \le s_{2m-1}$ dan
$s_{2m-1} - s_{2m} = b_{2m} \to 0$. Jadi $(s_{2m})$ tak menurun dan terbatas di atas
oleh $s_1 = b_1$, sehingga konvergen ke suatu $s$; $(s_{2m-1})$ tak menaik dan
terbatas di bawah oleh $s_2$, sehingga konvergen ke suatu $t$; dan
$s - t = \lim (s_{2m} - s_{2m-1}) = 0$. Kedua subbarisan menuju nilai yang sama, dan
karena subbarisan genap dan ganjil bersama-sama mencakup semua indeks, seluruh barisan
$(s_N)$ konvergen ke $s$ itu.

Uji Leibniz\index{uji Leibniz} memberi lebih dari konvergensi: batas galat yang bisa
dipakai tanpa mengetahui jumlah sebenarnya, dan itulah yang membuat deret berganti tanda
berguna untuk menghitung konstanta.

\begin{teorema}[Batas galat Leibniz]
Dengan hipotesis uji Leibniz dan $s = \sum (-1)^{n+1} b_n$,
\[
  \left| s - s_N \right| \le b_{N+1}.
\]
Lebih tajam: $s$ terletak di antara $s_N$ dan $s_{N+1}$, sehingga tanda galat
$s - s_N$ sama dengan tanda suku pertama yang dibuang, yaitu $(-1)^{N+2}$.
\end{teorema}

\noindent\textbf{Bukti.} Dari bukti sebelumnya, $(s_{2m})$ naik menuju $s$ dan $(s_{2m-1})$ turun menuju $s$.
Jadi untuk setiap $N$, bilangan $s$ berada di antara $s_N$ dan $s_{N+1}$. Karena
$|s_{N+1} - s_N| = b_{N+1}$, jarak $s$ ke $s_N$ tidak melebihi $b_{N+1}$, dan arah
pergeserannya ditentukan oleh tanda $s_{N+1} - s_N = (-1)^{N+2} b_{N+1}$.

Contoh klasiknya adalah deret harmonik berganti tanda,
\[
  1 - \frac{1}{2} + \frac{1}{3} - \frac{1}{4} + \cdots = \ln 2 .
\]
Kesamaan dengan $\ln 2$ dibuktikan lewat deret Taylor untuk $\ln(1+x)$ pada Bab 14;
di sini kita hanya memakai konvergensinya, yang dijamin uji Leibniz karena
$b_n = \frac{1}{n}$ menurun menuju nol. Angka $\ln 2$ muncul lagi, kini sebagai jumlah deret yang konvergen; selisih
$H_{2n} - H_n \to \ln 2$ adalah versi lain perhitungan yang sama.

Gambar~\ref{fig:tangga-leibniz} menggambar jumlah parsial deret itu sebagai dua tangga:
ganjil turun dari $s_1 = 1$, genap naik dari $s_2 = 0.5$, keduanya mengapit
$\ln 2 = 0.6931471806$. Pada $n = 10$ dan $n = 11$ apitannya
$0.6456349206 < \ln 2 < 0.7365440115$, selebar $b_{11} = \tfrac{1}{11} = 0.0909090909$,
tepat seperti yang dijanjikan teorema batas galat.

\begin{figure}[htbp]
\centering
\begin{tikzpicture}
\begin{axis}[
  width=0.93\linewidth, height=6.2cm,
  xlabel={banyak suku $n$}, ylabel={jumlah parsial $s_n$},
  xmin=0.4, xmax=15.0, ymin=0.45, ymax=1.06,
  xtick={1,2,3,4,5,6,7,8,9,10,11,12,13,14},
  tick label style={font=\scriptsize},
  label style={font=\small},
  ymajorgrids, grid style={layergrey!20},
  legend style={at={(0.98,0.04)},anchor=south east,draw=layergrey!40,fill=white,
                fill opacity=0.9,text opacity=1,font=\scriptsize},
]
\addplot[color=layerindigo,mark=*,mark size=1.6pt,thick] coordinates {
  (1,1.0) (3,0.8333333) (5,0.7833333) (7,0.7595238)
  (9,0.7456349) (11,0.7365440) (13,0.7301338)};
\addlegendentry{$n$ ganjil, turun}
\addplot[color=layerteal,mark=square*,mark size=1.6pt,thick] coordinates {
  (2,0.5) (4,0.5833333) (6,0.6166667) (8,0.6345238)
  (10,0.6456349) (12,0.6532107) (14,0.6587052)};
\addlegendentry{$n$ genap, naik}
\draw[dashed,accent,thick] (axis cs:0.4,0.6931472) -- (axis cs:15.0,0.6931472);
\node[font=\scriptsize,text=accent,fill=white,inner sep=1pt]
  at (axis cs:6.4,0.668) {$\ln 2 = 0.6931471806$};
\end{axis}
\end{tikzpicture}
\caption{Jumlah parsial deret harmonik berganti tanda. Subbarisan ganjil turun,
subbarisan genap naik, dan keduanya mengapit garis batas $\ln 2 = 0.6931471806$.
Pada $n = 10$ dan $n = 11$: $s_{10} = 0.6456349206$ dan $s_{11} = 0.7365440115$,
selebar $b_{11} = \tfrac{1}{11} = 0.0909090909$.}
\label{fig:tangga-leibniz}
\end{figure}

Batas galat Leibniz berguna dan sering ditaksir terlalu optimistis. Dengan $1000$
suku, jumlah parsialnya $0.7851481635$ sedangkan $\pi/4 = 0.7853981634$; galat
sebenarnya $2.5000 \times 10^{-4}$, batas yang dijanjikan
$b_{1001} = \tfrac{1}{2001} = 4.9975 \times 10^{-4}$. Batas itu benar tetapi dua kali
galat nyata, karena galat sebenarnya kira-kira $\tfrac{1}{4N}$ sedangkan batasnya
$\tfrac{1}{2N}$. Untuk galat di bawah $10^{-6}$ perlu $N \ge 500000$, dan pada $N$ itu
galatnya $5.0000 \times 10^{-7}$.

\section{Konvergensi Mutlak, Pengelompokan, dan Susunan Ulang}
\tingkat{lanjut}

Deret yang konvergen karena tandanya saling menghapus jauh lebih rapuh daripada deret
yang konvergen karena sukunya cukup cepat mengecil. Kerapuhan itu punya nama, dan
bab ini berutang satu penjelasan tentangnya.

\begin{definisi}
Deret $\sum a_n$ \emph{konvergen mutlak} bila $\sum |a_n|$ konvergen. Bila
$\sum a_n$ konvergen tetapi $\sum |a_n|$ divergen, deret itu disebut
\emph{konvergen bersyarat}\index{konvergensi bersyarat}.
\end{definisi}

\kamus{konvergensi mutlak}{absolute convergence}{deret yang tetap konvergen meskipun semua tandanya dipaksa menjadi positif, jadi tidak bergantung pada saling menghapus.}
\index{konvergensi mutlak}

\begin{teorema}
Jika $\sum a_n$ konvergen mutlak, maka $\sum a_n$ konvergen dan
$\left| \sum a_n \right| \le \sum |a_n|$.
\end{teorema}

\noindent\textbf{Bukti.} Untuk setiap $n \le m$ berlaku
$\left| \sum_{k=n}^{m} a_k \right| \le \sum_{k=n}^{m} |a_k|$. Karena $\sum |a_n|$
konvergen, ekspresi paling kanan bisa dibuat lebih kecil dari $\varepsilon$ menurut
kriteria Cauchy, sehingga $\sum a_n$ juga memenuhi kriteria Cauchy. Ketaksamaan
nilai mutlak jumlah diperoleh dengan mengambil limit pada
$\left| s_n \right| \le \sum_{k=1}^{n} |a_k|$.

Kebalikannya tidak berlaku: deret harmonik berganti tanda konvergen bersyarat,
karena $\sum \frac{1}{n}$ divergen. Pada deret seperti ini, urutan penjumlahan mulai
berperan. Sebelum sampai ke sana, kita lihat dulu operasi yang aman, yaitu
pengelompokan\index{pengelompokan}.

\begin{proposisi}[Pengelompokan]
Misalkan $\sum a_n$ konvergen ke $s$, dan $n_1 < n_2 < n_3 < \cdots$ barisan indeks.
Maka deret yang sukunya
\[
  A_k = a_{n_{k-1}+1} + a_{n_{k-1}+2} + \cdots + a_{n_k}
\]
(dengan $n_0 = 0$) juga konvergen ke $s$. Kebalikannya salah: pengelompokan yang
konvergen tidak menjamin deret asalnya konvergen.
\end{proposisi}

\noindent\textbf{Bukti.} Jumlah parsial deret kelompok adalah $S_K = A_1 + \cdots + A_K = s_{n_K}$, yaitu
subbarisan dari $(s_n)$. Subbarisan dari barisan konvergen menuju limit yang sama
(Bab 4), jadi $S_K \to s$.

Untuk kebalikannya, ambil $1 - 1 + 1 - 1 + \cdots$. Deret ini divergen karena sukunya
tidak menuju nol, tetapi dengan mengelompokkan berpasangan diperoleh
$(1-1) + (1-1) + \cdots = 0 + 0 + \cdots$, yang konvergen ke $0$.

Contoh penyangkal itu praktis penting: tanda kurung boleh dipasang pada deret yang sudah
diketahui konvergen, tetapi tidak boleh dipakai untuk membuktikannya.

Sekarang kita ubah urutannya, bukan hanya mengelompokkannya: susunan ulang menomori
kembali suku yang sama sehingga setiap suku muncul tepat sekali, tetapi pada posisi
berbeda.

\begin{definisi}
Misalkan $\sigma$ suatu bijeksi dari himpunan bilangan bulat positif ke dirinya
sendiri. Deret $\sum a_{\sigma(n)}$ disebut \emph{susunan ulang} dari $\sum a_n$.
\end{definisi}

\kamus{susunan ulang}{rearrangement}{menukar urutan kedatangan barang tanpa mengubah barangnya, sehingga total sementara bisa berbeda meski isi gudang sama.}
\index{susunan ulang}

Untuk deret yang konvergen mutlak, susunan ulang tidak berbahaya.

\begin{teorema}
Jika $\sum a_n$ konvergen mutlak, maka setiap susunan ulangnya konvergen, dan
semuanya menuju jumlah yang sama.
\end{teorema}

\noindent\textbf{Bukti.} Tulis $s = \sum a_n$ dan misalkan $\sum a_{\sigma(n)}$ susunan ulang dengan jumlah
parsial $s'_N$. Ambil $\varepsilon > 0$. Menurut kriteria Cauchy ada $N_0$ sehingga
$m \ge n \ge N_0$ berakibat $\sum_{k=n}^{m} |a_k| < \varepsilon$; khususnya
$\sum_{k=N_0}^{\infty} |a_k| \le \varepsilon$.

Pilih $P$ sehingga bilangan $1, 2, \ldots, N_0-1$ semuanya muncul di antara
$\sigma(1), \ldots, \sigma(P)$. Untuk $N > P$, suku-suku $a_1, \ldots, a_{N_0-1}$
saling menghapus pada selisih $s_N - s'_N$ (keduanya memuat suku itu), sehingga yang
tersisa hanyalah suku-suku berindeks $k \ge N_0$. Maka
\[
  \left| s_N - s'_N \right| \;\le\; \sum_{k \ge N_0} |a_k| \;\le\; \varepsilon
  \qquad (N > P).
\]
Karena $s_N \to s$ dan $\left| s'_N - s \right| \le \left| s'_N - s_N \right| +
\left| s_N - s \right|$, kita peroleh $s'_N \to s$.

Untuk deret yang hanya konvergen bersyarat, kesimpulan itu runtuh sepenuhnya. Teorema
berikut dirumuskan tanpa bukti; buktinya panjang dan teknis, dan dapat dibaca pada
Rudin, bab tentang barisan dan deret bilangan, bagian "Rearrangements".

\begin{teorema}[Susunan ulang Riemann, pernyataan]
Misalkan $\sum a_n$ konvergen tetapi tidak konvergen mutlak, dan misalkan $\alpha \le
\beta$ dua bilangan real (boleh $\pm\infty$). Maka ada susunan ulang $\sum
a_{\sigma(n)}$ yang jumlah parsialnya mempunyai limit inferior $\alpha$ dan limit
superior $\beta$. Khususnya, untuk setiap $c \in \mathbb{R}$ ada susunan ulang yang
konvergen ke $c$, dan ada susunan ulang yang divergen.
\end{teorema}

Gagasan buktinya begini. Bila $\sum a_n$ konvergen bersyarat, deret suku positifnya
dan deret nilai mutlak suku negatifnya sama-sama divergen. Kita menyusun ulang dengan
aturan: tambahkan suku positif sampai melewati $\beta$, lalu suku negatif sampai turun
di bawah $\alpha$, lalu ulangi. Karena kedua deret divergen, langkah itu selalu bisa
dilakukan, dan karena $a_n \to 0$, loncatan melewati batas semakin kecil. Hasilnya
osilasi dengan limit inferior $\alpha$ dan limit superior $\beta$. Konstruksi
lengkapnya ada pada Rudin, Teorema 3.54.

Kasus paling terkenal adalah deret harmonik berganti tanda yang disusun ulang dengan
pola dua suku positif lalu satu suku negatif,
\[
  1 + \frac{1}{3} - \frac{1}{2} + \frac{1}{5} + \frac{1}{7} - \frac{1}{4} + \cdots
\]
Deret ini konvergen, tetapi bukan ke $\ln 2$: jumlahnya $\tfrac{3}{2}\ln 2 =
1.0397207708$, dan perhitungannya ada di bagian berikut. Mengubah urutan mengubah
jawaban meskipun sukunya sama persis. Penjumlahan bilangan real tidak komutatif tanpa
syarat, dan itu peringatan konkret bagi siapa pun yang menjumlahkan pecahan di
komputer.

Ini juga menjelaskan mengapa uji di bagian kedua dan ketiga menuntut suku tak negatif.
Ketika tanda berganti, mengecil tidak cukup: pengecilannya harus cukup cepat mengalahkan
pembatalan.

\begin{untukanak}
Bayangkan dua tumpuk kartu: hijau bertanda "masuk", merah bertanda "keluar",
masing-masing dengan angka. Kamu memindahkan kartu satu per satu, dan saldo di layar
berubah: hijau menambah, merah mengurangi. Kalau urutan pemindahannya kamu ubah,
apakah saldo akhirnya berubah?

Kalau jumlah hijau dan merah sama-sama besar sehingga saling menghapus, urutannya
penting. Memindahkan hijau dulu membuat saldo melonjak sebelum turun; memindahkan
merah dulu membuat saldo sempat negatif. Saldo akhirnya pun bisa berbeda.

Ada satu keadaan aman: bila jumlah seluruh angka tanpa memandang warna tetap terbatas,
urutan tidak berpengaruh. Itulah konvergensi mutlak. Di bab ini kita akan tahu mengapa,
dan mengapa komputer bisa memberi jawaban berbeda untuk penjumlahan yang tampak sama.
\end{untukanak}

\section{Contoh Hitung}
\tingkat{menengah}

Bagian ini mengerjakan tiga hitungan yang menjadi tulang punggung bab. Semua angka
diperiksa dengan Python pada bagian lab; di sini langkah antaranya ditulis agar bisa
diulang dengan kertas dan pensil.

\medskip\noindent\textbf{1. Deret harmonik pada satu juta suku}

Kita ingin $H_{10^6}$ tanpa menjumlahkan satu juta suku. Alatnya adalah bentuk
asimtotik
\[
  H_n = \ln n + \gamma + \frac{1}{2n} - \frac{1}{12n^2} + O\left(\frac{1}{n^4}\right),
\]
dengan $\gamma = 0.5772156649$ adalah konstanta Euler--Mascheroni. Untuk $n = 10^6$:
\[
  \ln 10^6 = 13.8155105580, \qquad \ln 10^6 + \gamma = 14.3927262229 .
\]
Suku koreksi pertama $\frac{1}{2n} = 5 \times 10^{-7}$ memberi
$14.3927267229$, dan koreksi kedua $-\frac{1}{12n^2} = -8.33 \times 10^{-14}$
jauh di bawah presisi yang kita butuhkan. Jadi
\[
  H_{10^6} = 14.3927267228657\ldots
\]
Penjumlahan langsung dengan \texttt{math.fsum} memberi $14.392726722865724$, sepakat
sampai $10^{-15}$.

Yang lebih penting adalah berapa banyak suku yang dibutuhkan untuk mencapai batas
tertentu. Dari $H_n \approx \ln n + \gamma$ kita peroleh $n \approx e^{H - \gamma}$.
Untuk $H = 20$:
\[
  n \approx e^{20 - 0.5772156649} = e^{19.4227843351} = 2.724006 \times 10^{8}.
\]
Pencarian langsung memberi $n = 272400600$ sebagai indeks pertama dengan $H_n > 20$,
dan $H$ di titik itu $20.0000000016$. Untuk $H = 21$ perlu sekitar
$7.4 \times 10^8$ suku; taksiran pada $n = 10^9$ memberi $21.3004815$. Setiap
tambahan satu satuan menuntut penggandaan $n$ sebesar $e \approx 2.718$.

\medskip\noindent\textbf{2. Jumlah $\sum 1/n^2$ dan $\pi^2/6$}

Dengan $N = 10^6$ suku, penjumlahan langsung memberi
\[
  s_{10^6} = 1.6449330668487263, \qquad \frac{\pi^2}{6} = 1.6449340668482264 .
\]
Selisihnya $9.9999950 \times 10^{-7}$, sangat dekat dengan $\frac{1}{N} = 10^{-6}$;
itu masuk akal karena ekor deret $\sum_{n>N} n^{-2}$ kira-kira $\frac{1}{N}$.
Koreksi orde berikutnya adalah $-\frac{1}{2N^2} = -5 \times 10^{-13}$, sehingga
\[
  s_{10^6} + \frac{1}{N} - \frac{1}{2N^2} = 1.6449340668482262,
\]
yang cocok dengan $\pi^2/6$ sampai $2 \times 10^{-16}$, batas presisi pecahan
64-bit. Satu juta suku ditambah satu koreksi sudah memberi ketelitian mesin.

\medskip\noindent\textbf{3. Susunan ulang harmonik berganti tanda}

Ambil pola dua positif satu negatif: pada blok ke-$k$ kita menjumlahkan
$\frac{1}{4k-3} + \frac{1}{4k-1} - \frac{1}{2k}$. Setelah $K$ blok, seluruh suku
positif dengan penyebut $1, 3, \ldots, 4K-1$ sudah terpakai, dan suku negatif dengan
penyebut $2, 4, \ldots, 2K$ juga sudah terpakai. Jadi
\[
  S_K = \sum_{j=1}^{2K} \frac{1}{2j-1} - \sum_{j=1}^{K} \frac{1}{2j}
      = \left( H_{4K} - \frac{1}{2}H_{2K} \right) - \frac{1}{2}H_K .
\]
Dengan $H_m \approx \ln m + \gamma$,
\[
  S_K \approx \left(\ln 4K - \tfrac{1}{2}\ln 2K\right) - \tfrac{1}{2}\ln K
  = \ln \frac{4K}{\sqrt{2K}\sqrt{K}} = \ln\left(2\sqrt{2}\right) = \tfrac{3}{2}\ln 2 .
\]
Semua suku $\gamma$ menghapus, dan hasilnya tidak memuat $K$. Jadi
$S_K \to 1.0397207708$, bukan $\ln 2 = 0.6931471806$. Pada $K = 10^6$ blok nilai
numeriknya $1.0397205208$, galat $2.5 \times 10^{-7}$ terhadap $\tfrac{3}{2}\ln 2$,
sesuai orde $\frac{1}{4K}$ seperti deret Leibniz.

\kertas{Ringkasan angka. $H_{10^6} = 14.3927267229$ dan
$\ln 10^6 + \gamma = 14.3927262229$, selisihnya $\frac{1}{2n} = 5 \times 10^{-7}$.
$H_n > 20$ pertama kali tercapai pada $n = 272400600$, dekat
$e^{20-\gamma} = 2.724006 \times 10^8$. Deret Leibniz $1000$ suku memberi galat
$2.5000 \times 10^{-4}$ terhadap $\pi/4$, batasnya $\frac{1}{2001} =
4.9975 \times 10^{-4}$. Jumlah $\sum 1/n^2$ sampai $10^6$ suku kurang
$9.9999950 \times 10^{-7}$ dari $\pi^2/6$; satu koreksi Euler--Maclaurin menutupnya
sampai presisi mesin. Susunan ulang dua positif satu negatif menuju
$\frac{3}{2}\ln 2 = 1.0397207708$, bukan $\ln 2$.}

\section{Lab Komputasi}
\tingkat{lanjut}

Kode berikut dijalankan apa adanya pada Python 3.14.4 dengan NumPy 2.5.2. Bagian pertama
menguji ketelitian penjumlahan harmonik dan taksiran asimtotiknya.

\begin{lab}
\begin{verbatim}
import math
import numpy as np

n = 10**6
a = 1.0 / np.arange(1, n + 1)
vals = a.tolist()
naive = 0.0
for x in vals:
    naive += x

print("np.sum   ", repr(float(np.sum(a))))
print("for-loop ", repr(naive))
print("fsum     ", repr(math.fsum(vals)))
small = a.astype(np.float32)
print("float32  ", repr(float(np.sum(small, dtype=np.float32))))
g = 0.5772156649015328606
print("ln n + g ", repr(math.log(n) + g))
print("1/(2n)   ", 1 / (2 * n))
\end{verbatim}
\end{lab}

Keluaran yang diperoleh:
\begin{verbatim}
np.sum    14.392726722865726
for-loop  14.392726722864989
fsum      14.392726722865724
float32   14.392724990844727
ln n + g  14.392726222865807
1/(2n)    5e-07
\end{verbatim}

Empat hasil pertama seharusnya sama, dan tiga di antaranya berbeda pada digit
terakhir: galat pembulatan. Menjumlahkan sejuta bilangan berukuran $10^{-6}$ ke
akumulator berukuran $14$ membuang bit rendah. \texttt{float32} kehilangan
$1.7 \times 10^{-6}$, lebih besar daripada koreksi $\frac{1}{12n^2}$ yang ingin diukur;
bila galat yang dicari lebih kecil daripada galat penjumlahannya, hasilnya tidak
bermakna.

Bagian kedua memeriksa tiga angka dari bagian sebelumnya.

\begin{lab}
\begin{verbatim}
import math
import numpy as np

pi4 = math.pi / 4
s = 0.0
for k in range(1000):
    s += (-1)**k / (2*k + 1)
print("Leibniz 1000 suku", repr(s))
print("galat", repr(pi4 - s))
print("batas 1/2001", repr(1/2001))

N = 10**6
m = np.arange(1, N + 1, dtype=np.float64)
S = float(np.sum(1.0 / m**2))
print("jumlah 1/n^2", repr(S))
print("pi^2/6", repr(math.pi**2/6))
print("terkoreksi", repr(S + 1/N - 1/(2*N**2)))

S = 0.0
for k in range(1, 10**6 + 1):
    S += 1.0/(4*k-3) + 1.0/(4*k-1) - 1.0/(2*k)
print("susunan ulang", repr(S))
print("1.5 ln 2", repr(1.5*math.log(2)))
\end{verbatim}
\end{lab}

Keluarannya:
\begin{verbatim}
Leibniz 1000 suku 0.7851481634599485
galat 0.00024999993749974525
batas 1/2001 0.0004997501249375312
jumlah 1/n^2 1.6449330668487263
pi^2/6 1.6449340668482264
terkoreksi 1.6449340668482262
susunan ulang 1.0397205208399414
1.5 ln 2 1.0397207708399179
\end{verbatim}

Tiga hal perlu dibaca dari keluaran ini. Pertama, galat Leibniz $2.5000 \times
10^{-4}$ persis separuh batas $4.9975 \times 10^{-4}$; itu bukan kebetulan, karena
sisa deret berganti tanda yang monoton mendekati setengah suku pertama yang dibuang.
Kedua, koreksi Euler--Maclaurin menutup selisih sampai $2 \times 10^{-16}$, satu unit
pembulatan pada besaran $1.64$. Ketiga, susunan ulang memberi $1.0397205$, bukan
$0.6931472$, dan galatnya terhadap $\tfrac{3}{2}\ln 2$ hanya $2.5 \times 10^{-7}$.

\begin{untukdosen}
Urutan mengajar yang saya sarankan: mulai dari masalah epoch pada pembukaan, lalu
definisi jumlah parsial dan dua deret yang bisa dihitung (geometri, teleskopik).
Uji-uji masuk sebagai urutan pemeriksaan, bukan daftar hafalan: suku ke-$n$, pembanding,
rasio, akar. Deret-$p$ menyusul setelah kondensasi, yang menjelaskan mengapa harmonik
divergen dan mengapa lambat.

Dua miskonsepsi yang paling sering muncul. Pertama, mahasiswa menyimpulkan
konvergensi dari $a_n \to 0$; deret harmonik adalah obatnya, dan sebaiknya
ditunjukkan dengan tabel $H_{10} = 2.9290$, $H_{100} = 5.1874$, dan
$H_{1000} = 7.4855$ supaya lambatnya terlihat.
Kedua, mahasiswa menganggap "deret konvergen" berarti "jumlah parsial berhenti
berubah"; padahal perubahan per suku bisa jauh di bawah presisi mesin sementara
deretnya divergen. Ini tempat yang tepat untuk menyebut galat pembulatan.

Aktivitas kelas yang berhasil: bagi kelas menjadi kelompok, tiap kelompok diberi
empat deret tanpa nama, dan tugasnya hanya menentukan konvergen atau divergen
beserta ujinya dalam sepuluh menit. Deret yang dipilih harus memuat satu jebakan
rasio $=1$. Diskusikan kelompok yang memakai uji berbeda dan sama-sama benar.
Mahasiswa siap lanjut bila bisa menjelaskan mengapa $\sum 1/n^2$ konvergen sementara
$\sum 1/n$ tidak, cukup dengan pengelompokan blok.
\end{untukdosen}

\section{Ringkasan}
\tingkat{dasar}

\begin{ringkasan}
\begin{itemize}
\item Deret adalah barisan jumlah parsial; konvergen berarti barisan itu konvergen.
\item Untuk suku tak negatif, konvergen setara dengan jumlah parsial terbatas. Dari
      sini lahir uji pembanding, kondensasi, dan klasifikasi deret-$p$: konvergen
      tepat ketika $p>1$.
\item Uji rasio dan uji akar mengukur kecepatan pengecilan relatif terhadap deret
      geometri. Keduanya buta pada batas $1$, dan uji akar lebih kuat daripada uji
      rasio.
\item Uji Leibniz memberi konvergensi untuk deret berganti tanda yang menurun menuju
      nol, sekaligus batas galat $b_{N+1}$ yang tandanya diketahui.
\item Konvergensi mutlak kebal terhadap pengelompokan dan susunan ulang. Konvergensi
      bersyarat tidak: susunan ulang Riemann dapat mengubah jumlahnya menjadi
      bilangan apa pun, atau membuatnya divergen.
\item Deret harmonik divergen secara logaritmik: $H_{10^6} \approx 14.39$ dan perlu
      $272$ juta suku untuk mencapai $20$. Aturan berhenti berbasis "total tidak
      berubah lagi" tidak dapat dipercaya.
\end{itemize}
\end{ringkasan}

\section{Latihan}
\tingkat{menengah}

\begin{latihan}
\begin{enumerate}
\item Hitung $s_{20}$ untuk deret geometri $\sum_{k=0}^{\infty} \left(\frac{3}{5}\right)^k$,
      lalu bandingkan dengan jumlah sebenarnya. Berapa galatnya, dan berapa batas
      galat yang bisa kamu janjikan dari rumus $\frac{r^{21}}{1-r}$?

\item Buktikan bahwa $\sum_{n=1}^{\infty} \frac{1}{n(n+2)} = \frac{3}{4}$ dengan
      menguraikan suku menjadi pecahan parsial dan menunjukkan deretnya teleskopik.

\item Buktikan: bila $a_n \ge 0$ dan $\sum a_n$ konvergen, maka $\sum a_n^2$
      konvergen. Tunjukkan dengan contoh bahwa kebalikannya salah.

\item Nilai klaim ini: benar, salah, atau belum cukup bukti. "Karena
      $\frac{1}{n} \to 0$ dan setiap sukunya lebih kecil daripada suku deret
      geometri dengan rasio $\frac{1}{2}$ untuk $n$ cukup besar, deret harmonik
      konvergen." Sebutkan tepat di mana penalaran itu patah.

\item Tentukan konvergen atau divergen, dengan menyebut uji yang dipakai:
      (a) $\sum \frac{n^2}{2^n}$, (b) $\sum \frac{2^n}{n^2}$,
      (c) $\sum \frac{n!}{n^n}$, (d) $\sum \frac{1}{n^{1/3}}$.

\item Hitung jumlah $\sum_{n=1}^{10^6} \frac{1}{n^2}$ dengan NumPy, lalu taksir
      ekornya dengan $\frac{1}{N} - \frac{1}{2N^2}$ dan bandingkan hasilnya dengan
      $\pi^2/6$. Berapa angka di belakang koma yang benar sebelum koreksi, dan
      sesudahnya?

\item Deret $\sum \frac{(-1)^{n+1}}{\sqrt{n}}$ konvergen. Tentukan $N$ terkecil
      sehingga batas Leibniz menjamin galat di bawah $10^{-3}$, lalu hitung galat
      sebenarnya pada $N$ itu dengan Python. Apakah batasnya longgar atau ketat?

\item Buktikan bahwa bila $\sum a_n$ konvergen mutlak dan $(b_n)$ terbatas, maka
      $\sum a_n b_n$ konvergen mutlak.

\item Susun ulang deret harmonik berganti tanda dengan pola tiga positif satu
      negatif. Tebak jumlahnya dengan metode bagian contoh hitung, lalu periksa
      secara numerik. Apakah hasilnya masih berbentuk kelipatan $\ln 2$?

\item Buktikan bahwa bila $\sum a_n$ konvergen bersyarat, maka deret suku
      positifnya divergen dan deret nilai mutlak suku negatifnya juga divergen.
      Jelaskan mengapa fakta ini menjadi bahan bakar teorema susunan ulang Riemann.
\end{enumerate}

\noindent\emph{Petunjuk untuk tiga butir tersulit.} Butir 7: batas Leibniz memakai
suku pertama yang dibuang, jadi selesaikan $\frac{1}{\sqrt{N+1}} < 10^{-3}$ dan
periksa apakah ketaksamaan itu cukup tajam. Butir 9: setelah $K$ blok, suku positif
terpakai sampai penyebut $6K-1$, suku negatifnya sampai $\frac{1}{2K}$; tulis jumlah
parsialnya dalam $H_{6K}$, $H_{3K}$, dan $H_{K}$, lalu ambil limit. Butir 10: tulis
$p_n = \max(a_n, 0)$ dan $q_n = \max(-a_n, 0)$; maka $a_n = p_n - q_n$ dan
$|a_n| = p_n + q_n$, sehingga konvergensi salah satu deret itu menyimpulkan tentang
$\sum |a_n|$.
\end{latihan}

\begin{untukpeneliti}
Yang benar-benar terbukti di bab ini adalah pernyataan dengan bukti penuh: kriteria
Cauchy untuk deret, uji pembanding, kondensasi Cauchy, klasifikasi deret-$p$, uji rasio
dan uji akar dalam bentuk limit superior, uji Leibniz beserta batas galatnya, dan
ketertutupan deret konvergen mutlak terhadap susunan ulang. Pernyataan susunan ulang
Riemann saya kutip tanpa bukti. Sumber primernya Rudin, \emph{Principles of Mathematical
Analysis}, bab "Numerical Sequences and Series": Teorema 3.22--3.23, 3.25, 3.27--3.28,
3.33--3.34, 3.42--3.43, 3.45--3.46, dan 3.54--3.55. Abbott, \emph{Understanding
Analysis}, bab "Sequences and Series", membahas hal yang sama dengan penekanan pada
susunan ulang di bagian pembuka bab itu.

Yang sering dipopulerkan secara berlebihan adalah gagasan ambang batas antara deret
konvergen dan divergen. Rudin menunjukkan tidak ada ambang seperti itu: dua deret yang
berbeda sangat sedikit bisa berada di sisi berlawanan, dan prosedur itu bisa diteruskan
tanpa henti.

Pertanyaan terbuka: seberapa andal keputusan konvergensi bila $a_n$ hanya tersedia
sebagai taksiran berbising, seperti jumlah kerugian pada pembelajaran mesin? Untuk
deret semacam itu belum ada teori yang rapi.

Batas keberlakuan: seluruh bab membahas deret bilangan real. Di ruang bernorma lain
definisinya sama tetapi ujinya berbeda; deret fungsi menuntut kekonvergenan seragam
(Bab 14). Semua angka berasal dari NumPy 2.5.2, aritmetika 64-bit.
\end{untukpeneliti}

\peta{Bab 5 pada Rudin, \emph{Principles of Mathematical Analysis}, bab "Numerical
Sequences and Series": deret dan jumlah parsial (3.21--3.24), deret suku tak negatif
dan pembanding (3.25--3.29), uji akar dan rasio (3.33--3.37), konvergensi mutlak dan
susunan ulang (3.45--3.55). Pada Abbott, \emph{Understanding Analysis}, bab "Sequences
and Series": susunan ulang di bagian pembuka. Bab ini tidak memakai ulang prosa Bab 4;
dari sana hanya dirujuk definisi limit, kriteria Cauchy, dan kekonvergenan monoton.}


\part{Aljabar Linear}

Bagian ini memasang kerangka yang dipakai hampir setiap bab sesudahnya: ruang
vektor, basis, determinan, nilai eigen, dan dekomposisi nilai singular. Prasyaratnya
Bab~1--2. Setelah lima bab, kamu bisa menerjemahkan soal menjadi matriks,
menyelesaikannya, dan menjelaskan mengapa jawabannya tunggal atau tidak.
\chapter{Ruang Vektor dan Subruang}

Tabel berikut memuat tiga rekaman aliran jaringan; setiap baris satu rekaman, setiap kolom satu
fitur (paket masuk, paket keluar, durasi, koneksi gagal). Kolom-kolomnya, dibaca dari atas ke
bawah, adalah
\[
  \mathbf{c}_{1}=(1,2,3),\quad \mathbf{c}_{2}=(2,4,6),\quad
  \mathbf{c}_{3}=(0,1,2),\quad \mathbf{c}_{4}=(1,3,5),
\]
sedangkan baris-barisnya $\mathbf{r}_{1}=(1,2,0,1)$, $\mathbf{r}_{2}=(2,4,1,3)$, dan
$\mathbf{r}_{3}=(3,6,2,5)$.

Berapa kolom yang benar-benar membawa informasi baru? Menghitung kolom tak nol memberi empat,
dan hitungan itu keliru: kolom kedua dua kali kolom pertama, kolom keempat jumlah kolom
pertama dan ketiga. Tabel itu hanya memuat dua arah data mandiri. Kalimat ``dapat dihitung
ulang'' baru punya arti setelah kita dapat mengukur berapa vektor minimal yang membangunnya.

Bab 1 sampai Bab 5 menyiapkan bahasa himpunan, aksioma bilangan real, lalu limit barisan dan
deret; semuanya bekerja dengan kedekatan. Persoalan di atas tidak memuat kedekatan sama sekali:
entrinya eksak, dan pertanyaannya murni aljabar. Himpunan mana yang tertutup terhadap kedua
operasi, dan berapa vektor paling sedikit yang membangunnya. Satu kebiasaan dari Bab 5 tetap
terpakai: di sini setiap kombinasi berhingga banyaknya.

Ada pertanyaan kedua yang menuntut alat yang sama. Penilai linear memberi skor
$s(\mathbf{x})=\mathbf{w}\cdot\mathbf{x}$ dengan $\mathbf{w}=(1,1,1)$. Arah
$\boldsymbol{\delta}$ yang tidak mengubah skor memenuhi
$\mathbf{w}\cdot\boldsymbol{\delta}=0$, yaitu sebuah bidang; setiap perturbasi terpecah
menjadi bagian yang mengubah skor dan bagian yang tidak. Apakah pecahan itu tunggal? Ya, dan
sifat itu bernama jumlah langsung.

\begin{untukanak}
Bayangkan sekotak cat poster berisi empat kaleng: merah, kuning, biru, dan jingga. Kita
boleh mencampur isi kaleng dengan takaran apa pun, dan hasil campurannya boleh dipakai
untuk mencampur lagi. Jingga bisa dibuat dari merah dan kuning. Jadi kaleng jingga tidak
menambah satu kemampuan pun yang belum ada di kotak itu; ia hanya memakan tempat.

Pertanyaan yang membuat penasaran: berapa kaleng paling sedikit yang harus tetap ada agar
semua warna yang bisa kita buat tetap bisa kita buat? Jawabannya tidak sama dengan banyaknya
kaleng yang tersisa setelah kita membuang satu, dan urutan membuangnya tidak boleh mengubah
angka itu. Itu bukan kebetulan, dan membuktikannya tidak sesederhana kelihatannya.

Ada pertanyaan yang lebih aneh. Misalkan dua orang mencampur dengan resep berbeda, lalu
hasilnya sama persis. Apakah resepnya sama? Kadang ya, kadang tidak. Di bab ini kita akan tahu
mengapa angka itu punya nama sendiri, mengapa ia tidak bergantung pada urutan membuang, dan
mengapa dua campuran berbeda kadang menghasilkan warna yang sama.
\end{untukanak}

\section{Aksioma ruang vektor dan contohnya}
\tingkat{dasar}

Ruang vektor adalah himpunan yang anggotanya boleh dijumlahkan dan boleh diperbesar atau
diperkecil dengan bilangan, dengan hasil yang tetap berada di dalam himpunan itu, dan dengan
aturan hitung yang sama seperti pada pasangan bilangan biasa.

\begin{definisi}[Ruang vektor]
\index{ruang vektor}
Himpunan tak kosong $V$ bersama operasi penjumlahan $V\times V\to V$ dan perkalian skalar
$\mathbb{R}\times V\to V$ disebut \emph{ruang vektor atas $\mathbb{R}$} jika untuk semua
$\mathbf{u},\mathbf{v},\mathbf{w}\in V$ dan semua $a,b\in\mathbb{R}$ berlaku:
\begin{enumerate}
  \item $\mathbf{u}+\mathbf{v}=\mathbf{v}+\mathbf{u}$;
  \item $(\mathbf{u}+\mathbf{v})+\mathbf{w}=\mathbf{u}+(\mathbf{v}+\mathbf{w})$;
  \item ada $\mathbf{0}\in V$ dengan $\mathbf{v}+\mathbf{0}=\mathbf{v}$ untuk setiap
        $\mathbf{v}$;
  \item ada $-\mathbf{v}\in V$ dengan $\mathbf{v}+(-\mathbf{v})=\mathbf{0}$;
  \item $a(\mathbf{u}+\mathbf{v})=a\mathbf{u}+a\mathbf{v}$;
  \item $(a+b)\mathbf{v}=a\mathbf{v}+b\mathbf{v}$;
  \item $a(b\mathbf{v})=(ab)\mathbf{v}$;
  \item $1\mathbf{v}=\mathbf{v}$.
\end{enumerate}
\end{definisi}

Kedua operasi itu membawa syarat ketertutupan di dalam tipenya: hasilnya selalu di $V$, jadi
definisi ini tidak memuat aksioma penutupan tersendiri. Delapan aksioma itu juga tidak saling
bebas, dan kita tidak mengejar daftar terpendek; yang dibutuhkan akibat-akibat kecilnya,
semuanya diturunkan tanpa melihat bentuk anggota $V$.

\begin{lemma}
Dalam sembarang ruang vektor, vektor nol tunggal, invers penjumlahan setiap vektor tunggal,
$0\mathbf{v}=\mathbf{0}$, $(-1)\mathbf{v}=-\mathbf{v}$, dan $c\mathbf{v}=\mathbf{0}$
mengakibatkan $c=0$ atau $\mathbf{v}=\mathbf{0}$.
\end{lemma}

\par\medskip\noindent\textbf{Bukti.}\ 
Untuk ketunggalan nol, misalkan $\mathbf{0}$ dan $\mathbf{0}'$ keduanya netral. Maka
$\mathbf{0}=\mathbf{0}+\mathbf{0}'=\mathbf{0}'$. Untuk ketunggalan invers, misalkan
$\mathbf{v}+\mathbf{a}=\mathbf{0}$ dan $\mathbf{v}+\mathbf{b}=\mathbf{0}$. Maka
\[
  \mathbf{a}=\mathbf{a}+\mathbf{0}=\mathbf{a}+(\mathbf{v}+\mathbf{b})
  =(\mathbf{a}+\mathbf{v})+\mathbf{b}=\mathbf{0}+\mathbf{b}=\mathbf{b}.
\]
Selanjutnya $0\mathbf{v}=(0+0)\mathbf{v}=0\mathbf{v}+0\mathbf{v}$. Menambahkan
$-(0\mathbf{v})$ pada kedua ruas memberi $\mathbf{0}=0\mathbf{v}$. Dari situ
$\mathbf{v}+(-1)\mathbf{v}=1\mathbf{v}+(-1)\mathbf{v}=(1-1)\mathbf{v}=0\mathbf{v}=\mathbf{0}$,
jadi $(-1)\mathbf{v}$ adalah invers $\mathbf{v}$, dan karena invers tunggal,
$(-1)\mathbf{v}=-\mathbf{v}$. Terakhir, jika $c\mathbf{v}=\mathbf{0}$ dan $c\neq 0$, maka
\[
  \mathbf{v}=1\mathbf{v}=(c^{-1}c)\mathbf{v}=c^{-1}(c\mathbf{v})=c^{-1}\mathbf{0}=\mathbf{0}.
\]
\hfill$\square$\par\medskip

Empat contoh berikut cukup untuk seluruh buku ini.

\emph{(i) Ruang koordinat.} $\mathbb{R}^{n}$ dengan penjumlahan dan perkalian skalar
koordinat demi koordinat. Di sini $\mathbf{0}=(0,\dots,0)$, dan kedelapan aksioma tinggal
sifat lapangan bilangan real yang sudah dibahas di Bab 3.

\emph{(ii) Ruang matriks.} $M_{2\times3}(\mathbb{R})$, himpunan matriks $2\times3$ berentri
real, dengan operasi pada setiap entri dan matriks nol sebagai $\mathbf{0}$. Tidak ada yang
berubah bila $2\times3$ diganti $m\times n$.

\emph{(iii) Ruang polinomial.} $P_{3}$ adalah himpunan polinomial berkoefisien real
berderajat paling tinggi $3$ beserta polinomial nol, dengan operasi suku demi suku. Jebakan
yang layak dicatat: polinomial berderajat \emph{tepat} $3$ tidak membentuk ruang vektor, sebab
$x^{3}$ dan $-x^{3}+x$ berderajat tiga tetapi jumlahnya $x$ berderajat satu. Kata ``paling
tinggi'' bukan hiasan.

\emph{(iv) Ruang fungsi.} Himpunan semua fungsi $\mathbb{R}\to\mathbb{R}$, ditulis
$\mathbb{R}^{\mathbb{R}}$, dengan $(f+g)(x)=f(x)+g(x)$ dan $(af)(x)=af(x)$. Fungsi nol adalah
$\mathbf{0}$, dan aksioma ruang vektor diwarisi dari $\mathbb{R}$ di setiap titik $x$.

Dari contoh (iv) muncul contoh yang lebih halus. Himpunan $C[0,1]$ fungsi kontinu pada
$[0,1]$ adalah ruang vektor, karena jumlah dan kelipatan skalar fungsi kontinu tetap kontinu.
Begitu pula barisan real yang konvergen, memakai sifat limit jumlah dari Bab 4. Keduanya
bagian dari ruang yang lebih besar, dan bagian semacam itu kita definisikan sekarang.

Kembali ke $\mathbb{R}^{2}$, tempat kedua operasi itu terlihat langsung.
Gambar~\ref{fig:ch06-dua-operasi} memakai $\mathbf{u}=(2,1)$ dan $\mathbf{v}=(1,2)$: jumlah
$\mathbf{u}+\mathbf{v}=(3,3)$ tercapai lewat dua urutan, dan $2\mathbf{u}$ serta
$-\mathbf{u}$ tetap segaris dengan $\mathbf{u}$.

\begin{figure}[htbp]
\centering
\begin{tikzpicture}[scale=0.85,>=Stealth,font=\small]
  \draw[->,layergrey] (-2.8,0) -- (5.0,0) node[below] {$x$};
  \draw[->,layergrey] (0,-1.8) -- (0,4.2) node[left] {$y$};
  \draw[dashed,layergrey] (2,1) -- (3,3);
  \draw[dashed,layergrey] (1,2) -- (3,3);
  \draw[->,thick,accent] (0,0) -- (2,1) node[below right] {$\mathbf{u}$};
  \draw[->,thick,layerindigo] (0,0) -- (1,2) node[above left] {$\mathbf{v}$};
  \draw[->,thick,layerteal] (0,0) -- (3,3) node[right] {$\mathbf{u}+\mathbf{v}$};
  \draw[->,layerorange] (0,0) -- (4,2) node[right] {$2\mathbf{u}$};
  \draw[->,layerbrown] (0,0) -- (-2,-1) node[left] {$-\mathbf{u}$};
  \fill[ink] (0,0) circle (1.2pt);
  \node[below left] at (0,0) {$\mathbf{0}$};
\end{tikzpicture}
\caption{Dua operasi di $\mathbb{R}^{2}$. Panah tebal menunjukkan $\mathbf{u}=(2,1)$ dan
$\mathbf{v}=(1,2)$; garis putus-putus melengkapi jajar genjang, sehingga
$\mathbf{u}+\mathbf{v}=(3,3)$ tercapai lewat dua urutan yang sama. Panah tipis menunjukkan
$2\mathbf{u}=(4,2)$ dan $-\mathbf{u}=(-2,-1)$, keduanya tetap pada satu garis dengan
$\mathbf{u}$.}
\label{fig:ch06-dua-operasi}
\end{figure}

\kamus{ruang vektor}{vector space}{kotak pencampur: isinya boleh dijumlahkan dan diperbesar, dan hasilnya tidak pernah keluar dari kotak.}
\begin{bacakeras}
Delapan aksioma, dibaca seperti daftar periksa. Dua suara, bergantian; suara kedua selalu
menjawab dengan tiga kata.

Suara pertama: Tukar tempat. Suara kedua: Hasilnya sama.
Suara pertama: Pindah kurung. Suara kedua: Hasilnya sama.
Suara pertama: Ada vektor nol. Suara kedua: Tidak menambah apa pun.
Suara pertama: Ada vektor lawan. Suara kedua: Kembali ke nol.
Suara pertama: Skalar menyebar ke dua suku. Suara kedua: Kiri dan kanan sama.
Suara pertama: Dua skalar menjumlah. Suara kedua: Satu vektor menanggung.
Suara pertama: Dua skalar mengali. Suara kedua: Satu kali saja.
Suara pertama: Skalar satu. Suara kedua: Diam, tidak berubah.

Dua baris terakhir berbunyi hampir sama padahal berbeda: yang satu menjumlah takaran, yang
lain mengalikan takaran. Di situlah kekeliruan hitung paling sering lahir.

Ucapkan sekali lagi, pelan, sebagai satu tarikan napas: tukar tempat, pindah kurung, ada nol,
ada lawan, menyebar, menjumlah, mengali, diam.
\end{bacakeras}


\section{Subruang dan uji subruang}
\tingkat{dasar}

Subruang adalah bagian dari ruang vektor yang tetap berperilaku sebagai ruang vektor walaupun
kita hanya memakai anggota bagian itu.

\begin{definisi}[Subruang]
\index{subruang}
Himpunan $W\subseteq V$ disebut \emph{subruang} dari $V$ jika $W$ merupakan ruang vektor
dengan operasi yang sama seperti pada $V$.
\end{definisi}

Definisi itu menuntut delapan aksioma, dan memeriksa delapan hal setiap kali terlalu mahal.
Enam di antaranya gratis: komutatif, asosiatif, distributif berlaku di $W$ karena berlaku di
$V$. Yang tersisa tiga hal, dan itulah isi teorema berikut.

Jika $W$ tidak kosong dan tertutup terhadap kedua operasi, maka $W$ otomatis subruang.

\begin{teorema}[Uji subruang]
\index{uji subruang}
Misalkan $V$ ruang vektor dan $W\subseteq V$ dengan $W\neq\varnothing$. Maka $W$ subruang
dari $V$ jika dan hanya jika $W$ tertutup terhadap penjumlahan dan perkalian skalar.
\end{teorema}

\par\medskip\noindent\textbf{Bukti.}\ 
Jika $W$ subruang, hasil kedua operasi pada anggota $W$ harus menjadi anggota $W$, jadi
kedua penutupan berlaku. Sebaliknya, andaikan $W$ tidak kosong dan tertutup terhadap kedua
operasi. Ambil $\mathbf{w}\in W$. Menurut penutupan skalar, $\mathbf{0}=0\mathbf{w}\in W$,
jadi $W$ memuat netral. Juga $-\mathbf{w}=(-1)\mathbf{w}\in W$, jadi setiap anggota $W$
punya invers di dalam $W$. Aksioma yang tersisa berlaku karena berlaku di $V$ dan operasi
pada $W$ hanyalah pembatasan operasi pada $V$.
\hfill$\square$\par\medskip

Dalam praktik kita memeriksa tiga hal: $\mathbf{0}\in W$ dan kedua ketertutupan. Ketiganya
perlu; tiga contoh berikut menunjukkan bahwa kegagalan satu syarat tidak dapat ditutupi
syarat lain.
\begin{resep}
\begin{enumerate}
  \item Tulis syarat keanggotaan $W$ sebagai persamaan, bukan sebagai gambar atau nama.
  \item Masukkan $\mathbf{v}=\mathbf{0}$ ke syarat itu. Bila gagal, berhenti: bukan subruang.
  \item Ambil $\mathbf{u},\mathbf{v}\in W$ sembarang, periksa $\mathbf{u}+\mathbf{v}$ memakai
        syarat yang sama, bukan memakai contoh angka.
  \item Ambil $c\in\mathbb{R}$ sembarang dan $\mathbf{u}\in W$, periksa $c\mathbf{u}$; sertakan
        $c<0$ dan $0<c<1$.
\end{enumerate}
Peringatan urutan: jangan menukar langkah 2 dengan langkah 3 dan 4. Himpunan yang tidak memuat
$\mathbf{0}$ tidak perlu diuji ketertutupannya, tetapi berhenti setelah langkah 2 juga keliru:
memuat $\mathbf{0}$ belum cukup, seperti ditunjukkan contoh-contoh berikut.
\end{resep}


\emph{Garis tak melalui titik asal.} $W=\{(x,y)\in\mathbb{R}^{2}:x+y=1\}$ tidak memuat
$\mathbf{0}$, jadi bukan subruang; ia juga tidak tertutup terhadap penjumlahan maupun
perkalian skalar, sebab $(1,0)+(0,1)$ dan $2(1,0)$ keduanya keluar dari $W$.

\emph{Gabungan dua sumbu.} $W=\{(x,0)\}\cup\{(0,y)\}$ memuat $\mathbf{0}$ dan tertutup
terhadap perkalian skalar, tetapi $(1,0)+(0,1)=(1,1)$ keluar dari $W$.

\emph{Kuadran pertama.} $W=\{(x,y):x\ge 0,\ y\ge 0\}$ tertutup terhadap penjumlahan,
tetapi $(-1)(1,1)=(-1,-1)$ keluar dari $W$.

Ketiga contoh gagal pada syarat berbeda; Gambar~\ref{fig:ch06-garis} memperlihatkan pasangan
paling sederhana, yaitu garis melalui titik asal dan garis yang melewatinya.

\begin{figure}[htbp]
\centering
\begin{tikzpicture}[scale=0.72,>=Stealth,font=\small]
  \draw[->,layergrey] (-1.6,0) -- (3.2,0) node[below] {$x$};
  \draw[->,layergrey] (0,-1.6) -- (0,3.2) node[left] {$y$};
  \draw[very thick,layerteal] (-1.4,-1.4) -- (2.8,2.8);
  \fill[ink] (0,0) circle (1.3pt) node[below left] {$\mathbf{0}$};
  \fill[accent] (2,2) circle (1.5pt) node[above right] {$(2,2)$};
  \node at (0.8,-2.4) {\emph{garis melalui titik asal}};
  \begin{scope}[xshift=7cm]
    \draw[->,layergrey] (-1.6,0) -- (3.2,0) node[below] {$x$};
    \draw[->,layergrey] (0,-1.6) -- (0,3.2) node[left] {$y$};
    \draw[very thick,accent] (-0.8,1.8) -- (2.4,-1.4);
    \fill[accent] (1,0) circle (1.5pt) node[below right] {$(1,0)$};
    \fill[accent] (0,1) circle (1.5pt) node[above left] {$(0,1)$};
    \draw[ink,thick] (0,0) circle (3pt);
    \node[below left] at (0,0) {$\mathbf{0}$};
    \node at (0.8,-2.4) {\emph{garis tanpa titik asal}};
  \end{scope}
\end{tikzpicture}
\caption{Dua himpunan di $\mathbb{R}^{2}$. Kiri: $\{(x,y):y=x\}$ memuat $\mathbf{0}$ dan
$(2,2)$, jadi tertutup terhadap kedua operasi. Kanan: $\{(x,y):x+y=1\}$ memuat $(1,0)$ dan
$(0,1)$, tetapi titik asal ditandai lingkaran kosong karena tidak termasuk; satu titik yang
hilang sudah cukup membatalkan subruang.}
\label{fig:ch06-garis}
\end{figure}



Sebaliknya, dua kelas contoh berikut selalu lolos.

\emph{Ruang nol.} Untuk matriks $A$ berukuran $m\times n$, himpunan
$\{\mathbf{x}:A\mathbf{x}=\mathbf{0}\}$ adalah subruang, sebab $A\mathbf{0}=\mathbf{0}$,
$A(\mathbf{x}+\mathbf{y})=A\mathbf{x}+A\mathbf{y}=\mathbf{0}$, dan
$A(c\mathbf{x})=cA\mathbf{x}=\mathbf{0}$. Himpunan itu disebut ruang nol atau kernel $A$,
ditulis $\ker A$, dan dibahas lanjut di Bab 7.

\emph{Polinomial yang lenyap di satu titik.} Himpunan
$\{p\in P_{3}:p(1)=0\}$ adalah subruang $P_{3}$, karena evaluasi di titik $1$ linear:
$(p+q)(1)=p(1)+q(1)$ dan $(cp)(1)=c\,p(1)$.

Kedua contoh itu sama-sama mempertahankan himpunan solusi persamaan homogen. Itu bukan
kebetulan.
\begin{tantangan}
Contoh kuadran pertama gagal karena satu skalar negatif. Carilah himpunan
$W\subseteq\mathbb{R}^{2}$ yang memuat $\mathbf{0}$, tertutup terhadap penjumlahan, tidak
tertutup terhadap perkalian skalar, dan bukan gabungan dua sumbu, bukan kuadran, bukan pula
himpunan berkoordinat bulat.

Petunjuk 1: jangan mengatur arah panahnya; atur bilangan apa yang boleh muncul pada jumlah
kedua koordinat.

Petunjuk 2: minta $x+y$ bulat. Maka $(1,0)+(0,1)=(1,1)$ tetap di dalam, sedangkan
$\tfrac{1}{2}(1,0)$ keluar.

Kalau mentok, kembali ke bagian ``Subruang dan uji subruang'' dan bandingkan dengan contoh
garis $x+y=1$; bedanya hanya pada ruas kanan persamaan.
\end{tantangan}


\kamus{subruang}{subspace}{lorong di dalam gedung: kita boleh berjalan bolak-balik di lorong itu, tetapi tidak pernah keluar darinya.}

\begin{untukanak}
Bayangkan dua lorong lurus di sebuah gedung yang bertemu di satu titik, seperti huruf X.
Setiap lorong, kalau kita berjalan di dalamnya, tetap merupakan lorong: kita bisa maju,
mundur, dan berhenti di mana saja tanpa meninggalkannya. Begitu juga kalau kita berjalan
dua kali lebih jauh di lorong yang sama; kita masih di lorong itu.

Sekarang perhatikan gabungan kedua lorong. Ambil satu langkah di lorong pertama, lalu satu
langkah di lorong kedua. Tempat kita berhenti tidak berada di lorong mana pun. Jadi gabungan
dua lorong bukan lorong, walaupun masing-masing lorong adalah lorong. Inilah hal yang paling
sering dilupakan tentang kata ``tertutup'': yang diperiksa bukan anggotanya satu per satu,
melainkan hasil mencampurnya.

Pertanyaan yang membuat penasaran: kalau ada tiga lorong bertemu di satu titik, apakah
gabungannya lorong? Kadang ya, kalau ketiganya lorong yang sama. Di bab ini kita akan tahu
mengapa ``tertutup terhadap pencampuran'' jauh lebih kuat daripada ``memuat titik temu'', dan
mengapa satu contoh langkah yang keluar sudah cukup membatalkan semuanya.
\end{untukanak}

\section{Kombinasi linear dan rentang}
\tingkat{menengah}

Kombinasi linear adalah hasil mencampur beberapa vektor dengan takaran tertentu. Adapun
rentang adalah himpunan semua hasil campuran itu, untuk segala takaran yang mungkin.

\begin{definisi}[Kombinasi linear dan rentang]
\index{kombinasi linear}
\index{rentang}
Vektor $\mathbf{v}$ adalah \emph{kombinasi linear} dari
$\mathbf{v}_{1},\dots,\mathbf{v}_{k}$ jika ada skalar $c_{1},\dots,c_{k}$ dengan
$\mathbf{v}=c_{1}\mathbf{v}_{1}+\dots+c_{k}\mathbf{v}_{k}$. \emph{Rentang} himpunan
$S\subseteq V$, ditulis $\operatorname{span}(S)$, adalah himpunan semua kombinasi linear
berhingga anggota $S$. Untuk $S=\varnothing$ kita tetapkan
$\operatorname{span}(S)=\{\mathbf{0}\}$.
\end{definisi}

Rentang selalu subruang, dan ia subruang terkecil yang memuat $S$. Pernyataan pertama
mengatakan rentang layak disebut ruang; yang kedua, tidak ada ruang lain yang lebih hemat.

\begin{lemma}
Untuk setiap $S\subseteq V$, $\operatorname{span}(S)$ adalah subruang $V$, dan setiap
subruang yang memuat $S$ juga memuat $\operatorname{span}(S)$.
\end{lemma}

\par\medskip\noindent\textbf{Bukti.}\ 
Untuk $S=\varnothing$ pernyataan benar karena $\{\mathbf{0}\}$ subruang. Untuk $S$ tak
kosong, ambil $\mathbf{u}=\sum a_{i}\mathbf{s}_{i}$ dan $\mathbf{v}=\sum b_{j}\mathbf{t}_{j}$
dua kombinasi linear berhingga anggota $S$. Maka
$\mathbf{u}+\mathbf{v}=\sum a_{i}\mathbf{s}_{i}+\sum b_{j}\mathbf{t}_{j}$ dan
$c\mathbf{u}=\sum (ca_{i})\mathbf{s}_{i}$ juga kombinasi linear berhingga anggota $S$, jadi
rentang tertutup terhadap kedua operasi; ia juga memuat $\mathbf{0}=0\mathbf{s}$ untuk
sembarang $\mathbf{s}\in S$. Jika $W$ subruang dan $S\subseteq W$, maka setiap
$\mathbf{s}\in S$ ada di $W$, dan ketertutupan $W$ memaksa setiap kombinasi linear berhingga
anggota $S$ juga ada di $W$.
\hfill$\square$\par\medskip

Karena itu $\operatorname{span}(S)$ adalah irisan semua subruang yang memuat $S$, dan
menambahkan vektor yang sudah di dalam rentang tidak mengubah apa pun:
$\operatorname{span}(\operatorname{span}(S))=\operatorname{span}(S)$.

Jika $\mathbf{b}$ berada di rentang kolom-kolom $A$, maka sistem $A\mathbf{x}=\mathbf{b}$
punya solusi.

\begin{proposisi}
Misalkan $A$ matriks $m\times n$ dengan kolom $\mathbf{a}_{1},\dots,\mathbf{a}_{n}$ dan
$\mathbf{b}\in\mathbb{R}^{m}$. Sistem $A\mathbf{x}=\mathbf{b}$ konsisten jika dan hanya jika
$\mathbf{b}\in\operatorname{span}(\mathbf{a}_{1},\dots,\mathbf{a}_{n})$.
\end{proposisi}

\par\medskip\noindent\textbf{Bukti.}\ 
Menulis $A\mathbf{x}=\mathbf{b}$ koordinat demi koordinat sama dengan menulis
$x_{1}\mathbf{a}_{1}+\dots+x_{n}\mathbf{a}_{n}=\mathbf{b}$. Jadi solusi $\mathbf{x}$ ada
tepat ketika $\mathbf{b}$ merupakan kombinasi linear kolom-kolom $A$ dengan koefisien
$x_{1},\dots,x_{n}$.
\hfill$\square$\par\medskip

Himpunan $\operatorname{span}(\mathbf{a}_{1},\dots,\mathbf{a}_{n})$ disebut
\emph{ruang kolom} $A$, dan ia subruang $\mathbb{R}^{m}$ menurut lemma di atas. Untuk tabel
pembuka, ruang kolom itu dibangun empat kolom, dan dua di antaranya berlebih.

Sebagai contoh kecil, di $\mathbb{R}^{3}$ tinjau $\mathbf{s}_{1}=(1,0,1)$ dan
$\mathbf{s}_{2}=(0,1,1)$. Setiap kombinasi linear berbentuk $(a,b,a+b)$, jadi rentangnya
bidang $z=x+y$. Vektor $(1,2,4)$ di luar bidang itu karena $4\neq 1+2$, sedangkan $(1,2,3)$
di dalamnya dengan $a=1$, $b=2$.

\begin{tangan}
Bahan: dua pensil sama panjang dan satu meja datar. Tandai satu titik di meja dengan ujung
jari; sebut itu titik asal.

Letakkan pensil pertama dengan ekor di titik asal. Letakkan pensil kedua dengan ekor di ujung
pensil pertama, menyiku dengan yang pertama. Sekarang geser pensil kedua tanpa memutarnya
sampai ekornya kembali ke titik asal. Ujungnya berpindah ke tempat lain, dan tempat itu adalah
jumlah kedua pensil: dua urutan penjumlahan berhenti di titik yang sama.

Sekarang skalarnya. Balik satu pensil, arahnya menjadi negatif. Sambung dua pensil yang sama
dari titik asal, panjangnya menjadi dua kali. Yang harus pembaca lihat: semua ujung itu tetap
pada satu garis lurus, sejauh apa pun pensil disambung. Satu pensil hanya menutupi satu garis
di meja; dua pensil yang tidak sejajar menutupi seluruh permukaan meja. Itu bedanya garis dan
bidang, dan tidak ada pensil ketiga yang menambah permukaan baru.
\end{tangan}



\section{Kebebasan linear dan ketergantungan}
\tingkat{menengah}

Himpunan vektor disebut bebas linear jika tidak ada satu pun anggotanya yang dapat dicampur
dari anggota lain. Ia bergantung linear jika ada anggota yang berlebih.

\begin{definisi}[Bebas dan bergantung linear]
\index{kebebasan linear}
Himpunan $\{\mathbf{v}_{1},\dots,\mathbf{v}_{k}\}\subseteq V$ \emph{bebas linear} jika
satu-satunya pilihan skalar yang memenuhi
$c_{1}\mathbf{v}_{1}+\dots+c_{k}\mathbf{v}_{k}=\mathbf{0}$ adalah
$c_{1}=\dots=c_{k}=0$. Jika ada pilihan lain, himpunan itu \emph{bergantung linear}.
\end{definisi}

Dua kesepakatan kecil: himpunan kosong bebas linear karena tidak ada koefisien yang perlu
diperiksa, dan $\{\mathbf{0}\}$ bergantung linear karena $1\cdot\mathbf{0}=\mathbf{0}$
adalah relasi nontrivial.

Dalam himpunan yang bergantung linear selalu ada satu vektor yang terletak di rentang
vektor-vektor sebelumnya.

\begin{lemma}[Ketergantungan]
Misalkan $\mathbf{v}_{1}\neq\mathbf{0}$ dan $S=\{\mathbf{v}_{1},\dots,\mathbf{v}_{k}\}$.
Himpunan $S$ bergantung linear jika dan hanya jika ada indeks $j\ge 2$ dengan
$\mathbf{v}_{j}\in\operatorname{span}(\mathbf{v}_{1},\dots,\mathbf{v}_{j-1})$.
\end{lemma}

\par\medskip\noindent\textbf{Bukti.}\ 
Jika $\mathbf{v}_{j}=\sum_{i<j}c_{i}\mathbf{v}_{i}$, maka
$\sum_{i<j}c_{i}\mathbf{v}_{i}-\mathbf{v}_{j}=\mathbf{0}$ adalah relasi nontrivial, karena
koefisien $\mathbf{v}_{j}$ adalah $-1$; jadi $S$ bergantung. Sebaliknya, jika $S$ bergantung
maka ada relasi nontrivial $\sum_{i=1}^{k}c_{i}\mathbf{v}_{i}=\mathbf{0}$. Pilih $j$ terbesar
dengan $c_{j}\neq 0$. Karena $\mathbf{v}_{1}\neq\mathbf{0}$, indeks $j$ bukan $1$, jadi
$j\ge 2$ dan
\[
  \mathbf{v}_{j}=-\frac{1}{c_{j}}\sum_{i<j}c_{i}\mathbf{v}_{i}
  \in\operatorname{span}(\mathbf{v}_{1},\dots,\mathbf{v}_{j-1}).
\]
\hfill$\square$\par\medskip

Tiga akibatnya kita pakai terus: himpunan yang memuat $\mathbf{0}$ selalu bergantung; dua
vektor bergantung tepat ketika salah satunya kelipatan yang lain; dan setiap himpunan bagian
dari himpunan bebas linear juga bebas linear.

Sebagai contoh, tinjau $\mathbf{v}_{1}=(1,2,-1)$, $\mathbf{v}_{2}=(0,1,3)$, dan
$\mathbf{v}_{3}=(2,5,1)$ di $\mathbb{R}^{3}$. Karena
$2\mathbf{v}_{1}+\mathbf{v}_{2}=(2,5,1)=\mathbf{v}_{3}$, relasi
$2\mathbf{v}_{1}+\mathbf{v}_{2}-\mathbf{v}_{3}=\mathbf{0}$ nontrivial dan ketiganya
bergantung. Tidak satu pun vektor itu nol, dan tidak ada dua yang berkelipatan; ketergantungan
baru muncul pada tiga vektor.

Gambar~\ref{fig:ch06-bidang} menaruh keduanya berdampingan: rentang dua vektor sebagai
bidang, dan tiga vektor yang salah satunya sudah di rentang dua lainnya.

\begin{figure}[htbp]
\centering
\begin{tikzpicture}[scale=0.8,>=Stealth,font=\small]
  \begin{scope}[x={(1cm,0cm)}, y={(0.40cm,0.28cm)}, z={(0cm,1cm)}]
    \fill[layerteal!10] (0,0,0) -- (2.4,0,2.4) -- (2.4,1.8,4.2) -- (0,1.8,2.4) -- cycle;
    \draw[layerteal] (0,0,0) -- (2.4,0,2.4) -- (2.4,1.8,4.2) -- (0,1.8,2.4) -- cycle;
    \draw[->,layergrey] (0,0,0) -- (3.0,0,0) node[below] {$x$};
    \draw[->,layergrey] (0,0,0) -- (0,2.2,0) node[right] {$y$};
    \draw[->,layergrey] (0,0,0) -- (0,0,2.6) node[left] {$z$};
    \draw[->,thick,accent] (0,0,0) -- (1,0,1)
      node[above left,font=\scriptsize,inner sep=1pt] {$\mathbf{s}_{1}$};
    \draw[->,thick,layerindigo] (0,0,0) -- (0,1,1)
      node[right,font=\scriptsize,inner sep=1pt] {$\mathbf{s}_{2}$};
    \fill[layerjade] (1,2,3) circle (1.6pt);
    \node[right,font=\scriptsize] at (1,2,3) {$(1,2,3)$};
    \draw[ink] (1,2,4) circle (2pt);
    \node[right,font=\scriptsize] at (1,2,4) {$(1,2,4)$};
  \end{scope}
  \begin{scope}[xshift=6.4cm]
    \draw[->,layergrey] (-0.4,0) -- (3.2,0) node[below] {$x$};
    \draw[->,layergrey] (0,-0.4) -- (0,2.6) node[left] {$y$};
    \draw[dashed,layergrey] (0,0) -- (2,1);
    \draw[dashed,layergrey] (2,1) -- (2.4,2);
    \draw[->,thick,accent] (0,0) -- (1,0.5) node[below right] {$\mathbf{v}_{1}$};
    \draw[->,thick,layerindigo] (0,0) -- (0.4,1) node[above left] {$\mathbf{v}_{2}$};
    \draw[->,very thick,layerorange] (0,0) -- (2.4,2) node[right] {$\mathbf{v}_{3}$};
    \fill[ink] (0,0) circle (1.2pt);
  \end{scope}
\end{tikzpicture}
\caption{Kiri: $\operatorname{span}\{(1,0,1),(0,1,1)\}$ adalah bidang $z=x+y$; titik
$(1,2,3)$ berada di dalamnya, titik $(1,2,4)$ di luarnya. Kanan: sketsa dua dimensi dari
$\mathbf{v}_{1}=(1,2,-1)$ dan $\mathbf{v}_{2}=(0,1,3)$ dengan
$\mathbf{v}_{3}=2\mathbf{v}_{1}+\mathbf{v}_{2}=(2,5,1)$; garis putus-putus menunjukkan
$2\mathbf{v}_{1}$ ditambah $\mathbf{v}_{2}$ mendarat tepat di ujung $\mathbf{v}_{3}$, sehingga
ketiganya bergantung.}
\label{fig:ch06-bidang}
\end{figure}



Di $\mathbb{R}^{2}$ setiap tiga vektor selalu bergantung, tanpa perlu menghitung apa pun.

\begin{proposisi}
Setiap himpunan tiga vektor di $\mathbb{R}^{2}$ bergantung linear.
\end{proposisi}

\par\medskip\noindent\textbf{Bukti.}\ 
Misalkan $S=\{\mathbf{u},\mathbf{v},\mathbf{w}\}\subseteq\mathbb{R}^{2}$. Jika dua di
antaranya berkelipatan, $S$ bergantung. Jika tidak, $\mathbf{u}$ dan $\mathbf{v}$ bebas
linear. Ambil $\mathbf{w}=(w_{1},w_{2})$ dan cari $a,b$ dengan
$a\mathbf{u}+b\mathbf{v}=\mathbf{w}$. Menuliskan komponennya memberi dua persamaan linear
dalam $a$ dan $b$, dan karena $\mathbf{u}$, $\mathbf{v}$ tidak berkelipatan, sistem itu punya
solusi tunggal. Jadi $\mathbf{w}\in\operatorname{span}(\mathbf{u},\mathbf{v})$.
\hfill$\square$\par\medskip

Di $\mathbb{R}^{3}$ pernyataan serupa berlaku untuk empat vektor, tetapi cara di atas akan
menjadi pengecekan kasus yang panjang; kita butuh argumen seragam.

Untuk memeriksa kebebasan secara hitung, susun vektor sebagai kolom matriks lalu selesaikan
$A\mathbf{c}=\mathbf{0}$; himpunan bebas linear tepat ketika sistem itu hanya punya solusi
trivial. Determinan memberi jalan pintas di $\mathbb{R}^{3}$, tetapi alat itu baru dibangun di
Bab 8.

\section{Basis dan dimensi}
\tingkat{lanjut}

Basis adalah daftar vektor yang cukup untuk membangun seluruh ruang dan tidak memuat satu pun
vektor berlebih. Kata ``cukup'' dan ``tidak berlebih'' persis dua syarat yang sudah kita punya.

\begin{definisi}[Basis dan dimensi]
\index{basis}
Himpunan $B\subseteq V$ disebut \emph{basis} dari $V$ jika $B$ bebas linear dan
$\operatorname{span}(B)=V$. Jika $V$ mempunyai basis berhingga, $V$ disebut
\emph{berdimensi hingga}; jika tidak, $V$ berdimensi tak hingga.
\end{definisi}

Definisi dimensi belum bisa ditulis, karena ia bergantung pada satu kenyataan yang belum
dibuktikan: dua basis dari ruang yang sama selalu berukuran sama.

Langkah pertama: penulisan dalam basis selalu tunggal.

\begin{lemma}[Ketunggalan koordinat]
Jika $B=\{\mathbf{b}_{1},\dots,\mathbf{b}_{n}\}$ basis $V$, maka setiap $\mathbf{v}\in V$
dapat ditulis secara tunggal sebagai $\mathbf{v}=a_{1}\mathbf{b}_{1}+\dots+a_{n}\mathbf{b}_{n}$.
\end{lemma}

\par\medskip\noindent\textbf{Bukti.}\ 
Keberadaan penulisan datang dari $\operatorname{span}(B)=V$. Andaikan ada dua penulisan
dengan koefisien $a_{i}$ dan $a_{i}'$. Pengurangan keduanya memberi
$\sum_{i}(a_{i}-a_{i}')\mathbf{b}_{i}=\mathbf{0}$. Karena $B$ bebas linear, semua koefisien
selisih itu nol, jadi $a_{i}=a_{i}'$ untuk setiap $i$.
\hfill$\square$\par\medskip

Angka $a_{1},\dots,a_{n}$ disebut koordinat $\mathbf{v}$ terhadap $B$; pertukaran basis
adalah pokok bahasan Bab 7, dan di sini kita hanya memerlukan ketunggalannya.

Langkah kedua: satu vektor dapat ditukar dengan vektor lain tanpa merusak sifat merentang.
Menukar satu vektor dalam himpunan merentang dengan kombinasi yang memuatnya tidak mengurangi
apa yang dapat dibangun.

\begin{lemma}[Penukaran]
\index{lemma penukaran}
Misalkan $\mathbf{u}=\sum_{\mathbf{s}\in S}a_{\mathbf{s}}\mathbf{s}$ jumlah berhingga dengan
$a_{\mathbf{s}_{0}}\neq 0$ untuk suatu $\mathbf{s}_{0}\in S$. Jika $S\cup\{\mathbf{u}\}$
merentang $V$, maka $(S\setminus\{\mathbf{s}_{0}\})\cup\{\mathbf{u}\}$ juga merentang $V$.
\end{lemma}

\par\medskip\noindent\textbf{Bukti.}\ 
Dari relasi itu,
$\mathbf{s}_{0}=a_{\mathbf{s}_{0}}^{-1}\bigl(\mathbf{u}-\sum_{\mathbf{s}\neq\mathbf{s}_{0}}a_{\mathbf{s}}\mathbf{s}\bigr)$,
jadi $\mathbf{s}_{0}$ adalah kombinasi linear anggota himpunan kecil
$T=(S\setminus\{\mathbf{s}_{0}\})\cup\{\mathbf{u}\}$. Setiap anggota $S\cup\{\mathbf{u}\}$
karena itu berada di $\operatorname{span}(T)$, sehingga
$\operatorname{span}(S\cup\{\mathbf{u}\})\subseteq\operatorname{span}(T)$. Arah sebaliknya
jelas karena $T\subseteq S\cup\{\mathbf{u}\}$. Jadi kedua rentang sama dengan $V$.
\hfill$\square$\par\medskip

Langkah ketiga adalah inti bab ini.

Jika ruang dapat dibangun oleh $n$ vektor, maka mustahil menemukan $n+1$ vektor yang bebas
linear di dalamnya.

\begin{teorema}[Pertukaran]
\index{teorema pertukaran}
Jika $V$ direntang oleh $n$ vektor dan $\mathbf{w}_{1},\dots,\mathbf{w}_{m}$ bebas linear di
$V$, maka $m\le n$.
\end{teorema}

\par\medskip\noindent\textbf{Bukti.}\ 
Misalkan $S_{0}=\{\mathbf{v}_{1},\dots,\mathbf{v}_{n}\}$ merentang $V$. Untuk $j=1$:
$\mathbf{w}_{1}\neq\mathbf{0}$ adalah kombinasi linear anggota $S_{0}$ dengan sekurangnya
satu koefisien tak nol, sehingga lemma penukaran memberi $T_{1}\subseteq S_{0}$ berukuran
$n-1$ dengan $\{\mathbf{w}_{1}\}\cup T_{1}$ merentang.

Andaikan $\{\mathbf{w}_{1},\dots,\mathbf{w}_{j}\}\cup T_{j}$ merentang dengan
$|T_{j}|=n-j$. Vektor $\mathbf{w}_{j+1}$ adalah kombinasi linear himpunan itu. Jika semua
koefisien pada $T_{j}$ nol, maka $\mathbf{w}_{j+1}$ menjadi kombinasi linear
$\mathbf{w}_{1},\dots,\mathbf{w}_{j}$, yang mustahil karena himpunan itu bebas linear. Jadi
ada anggota $T_{j}$ berkoefisien tak nol, dan lemma penukaran memberi $T_{j+1}$ berukuran
$n-j-1$ yang tetap merentang bersama $\mathbf{w}_{1},\dots,\mathbf{w}_{j+1}$.

Setelah $m$ langkah kita memperoleh himpunan merentang berukuran $n-m$; ukuran himpunan tidak
boleh negatif, jadi $m\le n$.
\hfill$\square$\par\medskip

Sekarang ukuran sebuah ruang dapat ditulis: banyaknya vektor dalam daftar paling hematnya,
dan angka itu tidak bergantung pada daftar yang dipilih.

Dua basis dari ruang berdimensi hingga selalu sama banyak.

\begin{akibat}
Jika $V$ berdimensi hingga dengan basis $B_{1}$ berukuran $n$ dan basis $B_{2}$ berukuran
$m$, maka $n=m$.
\end{akibat}

\par\medskip\noindent\textbf{Bukti.}\ 
$B_{1}$ merentang dan $B_{2}$ bebas linear, jadi teorema pertukaran memberi $m\le n$.
Sebaliknya $B_{2}$ merentang dan $B_{1}$ bebas linear, sehingga $n\le m$.
\hfill$\square$\par\medskip

Banyaknya vektor dalam daftar paling hemat itu cukup penting untuk diberi nama sendiri.

\begin{definisi}[Dimensi]
\index{dimensi}
Untuk $V$ berdimensi hingga, $\dim V$ adalah banyaknya anggota sembarang basis $V$. Ruang
nol $\{\mathbf{0}\}$ ditetapkan berdimensi $0$.
\end{definisi}

Ukuran subruang tidak pernah melebihi ukuran ruangnya; daftar yang panjangnya pas dengan
dimensi langsung menjadi basis; dan daftar yang bebas linear selalu bisa dilengkapi sampai
penuh.

\begin{akibat}
Misalkan $\dim V=n$. (a) Setiap himpunan $n$ vektor yang bebas linear adalah basis.
(b) Setiap himpunan $n$ vektor yang merentang adalah basis. (c) Setiap himpunan bebas linear
dapat diperluas menjadi basis. (d) Jika $U$ subruang $V$, maka $\dim U\le\dim V$, dan
kesamaan berlaku hanya bila $U=V$.
\end{akibat}

\par\medskip\noindent\textbf{Bukti.}\ 
Untuk (a), himpunan bebas linear berukuran $n$ yang tidak merentang dapat ditambah satu
vektor di luar rentangnya, dan hasilnya tetap bebas linear; jadi ada $n+1$ vektor bebas
linear, dan teorema pertukaran melawannya. Untuk (b), himpunan merentang berukuran $n$ yang
bergantung dapat kehilangan satu vektor lewat lemma ketergantungan dan tetap merentang dengan
$n-1$ vektor; lagi-lagi bertentangan dengan pertukaran. Untuk (c), tambahkan vektor di luar
rentang selama masih ada; proses berhenti setelah paling banyak $n$ langkah, dan saat berhenti
rentangnya sudah seluruh $V$. Untuk (d), basis $U$ adalah himpunan bebas linear di $V$, jadi
ukurannya paling banyak $n$; bila ukurannya $n$, bagian (a) membuatnya basis $V$, sehingga
$U=V$.
\hfill$\square$\par\medskip

Dengan demikian $\dim\mathbb{R}^{n}=n$, $\dim P_{n}=n+1$ dengan basis $1,x,\dots,x^{n}$,
$\dim M_{m\times n}(\mathbb{R})=mn$ dengan basis matriks yang tepat satu entrinya $1$, dan
$\dim\{\mathbf{0}\}=0$. Untuk ruang fungsi seperti $\mathbb{R}^{\mathbb{R}}$ tidak ada basis
berhingga; buktinya memerlukan aksioma pilihan dan di luar cakupan buku ini.

\kamus{basis}{basis}{daftar kaleng cat paling hemat: cukup untuk membuat semua warna, dan tidak ada satu kaleng pun yang bisa dibuang.}

\kamus{dimensi}{dimension}{banyaknya kaleng dalam daftar paling hemat itu; angkanya milik ruang, bukan milik daftar yang kebetulan kita pilih.}

\section{Jumlah langsung dan subruang komplemen}
\tingkat{lanjut}

Jumlah dua subruang adalah himpunan semua hasil menambah satu anggota dari masing-masing
subruang.

\begin{definisi}[Jumlah subruang]
\index{jumlah subruang}
Untuk subruang $U,W\subseteq V$, \emph{jumlah} $U+W$ adalah
$\{\mathbf{u}+\mathbf{w}:\mathbf{u}\in U,\ \mathbf{w}\in W\}$.
\end{definisi}

Ia berbeda dari gabungan, yang hampir selalu bukan subruang seperti contoh dua sumbu.

Jumlahnya sendiri sebuah subruang, dan tidak ada subruang lain yang lebih kecil yang
memuat keduanya.

\begin{proposisi}
Untuk subruang $U,W\subseteq V$, himpunan $U+W$ adalah subruang $V$, dan ia termuat di
setiap subruang yang memuat $U\cup W$.
\end{proposisi}

\par\medskip\noindent\textbf{Bukti.}\ 
Himpunan $U+W$ memuat $\mathbf{0}=\mathbf{0}+\mathbf{0}$. Jika
$\mathbf{u}_{1}+\mathbf{w}_{1}$ dan $\mathbf{u}_{2}+\mathbf{w}_{2}$ dua anggotanya, jumlahnya
$(\mathbf{u}_{1}+\mathbf{u}_{2})+(\mathbf{w}_{1}+\mathbf{w}_{2})$ dengan suku pertama di $U$
dan kedua di $W$, dan $c(\mathbf{u}_{1}+\mathbf{w}_{1})=c\mathbf{u}_{1}+c\mathbf{w}_{1}$ juga
demikian; jadi $U+W$ subruang. Jika $Z$ subruang memuat $U$ dan $W$, maka $Z$ memuat setiap
$\mathbf{u}+\mathbf{w}$ karena tertutup terhadap penjumlahan, jadi $U+W\subseteq Z$.
\hfill$\square$\par\medskip

Ukuran jumlah bukan jumlah ukuran, karena bagian yang beririsan terhitung dua kali.

$\dim(U+W)$ sama dengan $\dim U$ ditambah $\dim W$ dikurangi dimensi irisan.

\begin{teorema}[Dimensi jumlah]
\index{dimensi jumlah}
Jika $U$ dan $W$ subruang ruang berdimensi hingga $V$, maka
\[
  \dim(U+W)=\dim U+\dim W-\dim(U\cap W).
\]
\end{teorema}

\par\medskip\noindent\textbf{Bukti.}\ 
Irisan $U\cap W$ subruang, jadi berdimensi hingga. Ambil basis
$\{\mathbf{z}_{1},\dots,\mathbf{z}_{m}\}$ dari $U\cap W$. Karena bebas linear di $U$,
himpunan itu diperluas menjadi basis $U$ dengan menambahkan $\mathbf{u}_{1},\dots,
\mathbf{u}_{j}$, sehingga $\dim U=m+j$; dengan cara sama, menambahkan
$\mathbf{w}_{1},\dots,\mathbf{w}_{k}$ menghasilkan basis $W$, sehingga $\dim W=m+k$.

Klaim: $\{\mathbf{z}_{1},\dots,\mathbf{z}_{m},\mathbf{u}_{1},\dots,\mathbf{u}_{j},
\mathbf{w}_{1},\dots,\mathbf{w}_{k}\}$ basis $U+W$. Himpunan itu memuat basis $U$ dan basis
$W$, jadi rentangnya memuat $U+W$; arah sebaliknya jelas. Untuk kebebasan, andaikan
\[
  \sum_{i}a_{i}\mathbf{z}_{i}+\sum_{r}b_{r}\mathbf{u}_{r}
  +\sum_{s}c_{s}\mathbf{w}_{s}=\mathbf{0}.
\]
Maka $\sum_{s}c_{s}\mathbf{w}_{s}=-\sum_{i}a_{i}\mathbf{z}_{i}-\sum_{r}b_{r}\mathbf{u}_{r}$
berada di $U$ sekaligus di $W$; jadi ia anggota $U\cap W$ dan dapat ditulis
$\sum_{i}d_{i}\mathbf{z}_{i}$. Maka $\sum_{s}c_{s}\mathbf{w}_{s}-\sum_{i}d_{i}\mathbf{z}_{i}
=\mathbf{0}$ adalah relasi linear dalam basis $W$, sehingga semua $c_{s}$ dan $d_{i}$ nol;
persamaan awal menyusut menjadi relasi linear dalam basis $U$, sehingga $a_{i}$ dan $b_{r}$
juga nol.

Jadi $\dim(U+W)=m+j+k=(m+j)+(m+k)-m=\dim U+\dim W-\dim(U\cap W)$.
\hfill$\square$\par\medskip

Sekarang pertanyaan kedua pembuka bab: kapan penguraian menjadi dua suku itu tunggal?

Jumlah $U+W$ disebut \emph{jumlah langsung}, ditulis $U\oplus W$, jika setiap anggota
$U+W$ dapat ditulis sebagai $\mathbf{u}+\mathbf{w}$ dengan cara tepat satu.

\begin{definisi}[Jumlah langsung]
\index{jumlah langsung}
Jika setiap $\mathbf{x}\in U+W$ memiliki tepat satu penulisan
$\mathbf{x}=\mathbf{u}+\mathbf{w}$ dengan $\mathbf{u}\in U$ dan $\mathbf{w}\in W$, maka
$U+W$ disebut jumlah langsung dan ditulis $U\oplus W$.
\end{definisi}

Ketunggalan seluruh penulisan dapat diperiksa hanya dengan melihat penulisan nol.

$U+W$ jumlah langsung tepat ketika $U\cap W=\{\mathbf{0}\}$.

\begin{teorema}[Kriteria jumlah langsung]
Misalkan $U$ dan $W$ subruang $V$. Maka $U+W=U\oplus W$ jika dan hanya jika
$U\cap W=\{\mathbf{0}\}$.
\end{teorema}

\par\medskip\noindent\textbf{Bukti.}\ 
Andaikan $U+W$ jumlah langsung dan ambil $\mathbf{x}\in U\cap W$. Maka
$\mathbf{0}=\mathbf{x}+(-\mathbf{x})$ dengan $\mathbf{x}\in U$ dan $-\mathbf{x}\in W$, dan
juga $\mathbf{0}=\mathbf{0}+\mathbf{0}$. Ketunggalan memberi $\mathbf{x}=\mathbf{0}$.

Sebaliknya andaikan $U\cap W=\{\mathbf{0}\}$. Ambil $\mathbf{x}\in U+W$ dengan dua penulisan
$\mathbf{x}=\mathbf{u}_{1}+\mathbf{w}_{1}=\mathbf{u}_{2}+\mathbf{w}_{2}$. Maka
$\mathbf{u}_{1}-\mathbf{u}_{2}=\mathbf{w}_{2}-\mathbf{w}_{1}$ berada di $U$ dan di $W$,
jadi sama dengan $\mathbf{0}$. Akibatnya $\mathbf{u}_{1}=\mathbf{u}_{2}$ dan
$\mathbf{w}_{1}=\mathbf{w}_{2}$.
\hfill$\square$\par\medskip

Bagi dua subruang kriteria itu cukup. Bagi tiga subruang atau lebih, memeriksa irisan
berpasangan tidak cukup, dan contoh berikut menunjukkan mengapa.

Di $\mathbb{R}^{3}$ ambil $V_{1}=\operatorname{span}\{(1,0,0)\}$,
$V_{2}=\operatorname{span}\{(0,1,0)\}$, dan $V_{3}=\operatorname{span}\{(1,1,0)\}$. Setiap
irisan berpasangan hanya memuat $\mathbf{0}$, tetapi $V_{1}+V_{2}+V_{3}$ bukan jumlah
langsung: $(1,1,0)$ dapat ditulis $(1,0,0)+(0,1,0)+\mathbf{0}$ dan juga
$\mathbf{0}+\mathbf{0}+(1,1,0)$. Nol pun punya dua penulisan, yaitu
$\mathbf{0}+\mathbf{0}+\mathbf{0}$ dan $(1,0,0)+(0,1,0)+(-1,-1,0)$.

Untuk banyak suku, kriterianya kembali ke definisi: $\mathbf{0}=\mathbf{v}_{1}+\dots+
\mathbf{v}_{k}$ dengan $\mathbf{v}_{i}\in V_{i}$ hanya mungkin bila semua $\mathbf{v}_{i}$
nol.

Terakhir, setiap subruang dapat dipasangkan dengan subruang lain sehingga keduanya
menguraikan seluruh ruang.

\begin{teorema}[Komplemen]
\index{komplemen}
Jika $V$ berdimensi hingga dan $U$ subruang $V$, maka ada subruang $W\subseteq V$ dengan
$V=U\oplus W$. Dalam hal itu $\dim V=\dim U+\dim W$.
\end{teorema}

\par\medskip\noindent\textbf{Bukti.}\ 
Jika $U=\{\mathbf{0}\}$, ambil $W=V$. Jika tidak, ambil basis
$\{\mathbf{u}_{1},\dots,\mathbf{u}_{m}\}$ dari $U$ dan perluas menjadi basis
$\{\mathbf{u}_{1},\dots,\mathbf{u}_{m},\mathbf{w}_{1},\dots,\mathbf{w}_{k}\}$ dari $V$;
perluasan ini dimungkinkan oleh akibat (c) di bagian sebelumnya. Tetapkan
$W=\operatorname{span}\{\mathbf{w}_{1},\dots,\mathbf{w}_{k}\}$.

Setiap $\mathbf{x}\in V$ adalah kombinasi linear seluruh basis, dan memisahkan suku
$\mathbf{u}$ dari suku $\mathbf{w}$ menunjukkannya sebagai anggota $U+W$. Untuk irisan, ambil
$\mathbf{x}\in U\cap W$. Maka
$\mathbf{x}=\sum a_{i}\mathbf{u}_{i}=\sum b_{s}\mathbf{w}_{s}$, sehingga
$\sum a_{i}\mathbf{u}_{i}-\sum b_{s}\mathbf{w}_{s}=\mathbf{0}$. Ini relasi linear dalam basis
$V$, jadi semua koefisiennya nol, khususnya $\mathbf{x}=\mathbf{0}$. Menurut kriteria di
atas, $V=U\oplus W$, dan dari teorema dimensi jumlah,
$\dim V=\dim U+\dim W$.

Komplemen itu tidak tunggal. Di $\mathbb{R}^{2}$ dengan $U$ sumbu-$x$, setiap garis melalui
titik asal selain sumbu-$x$ adalah komplemen $U$; yang tunggal bukan pilihannya, melainkan
dimensinya. Di Bab 10, komplemen ortogonal membuat penguraian itu sekaligus terpendek.

Sekarang pertanyaan pembuka dapat dijawab. Untuk $\mathbf{w}=(1,1,1)$, misalkan
$U=\{\boldsymbol{\delta}:\mathbf{w}\cdot\boldsymbol{\delta}=0\}$ dan
$W=\operatorname{span}\{\mathbf{w}\}$. Irisannya hanya $\{\mathbf{0}\}$, sebab
$\mathbf{w}\cdot(t\mathbf{w})=3t$ yang nol hanya bila $t=0$; dengan $\dim U=2$ dan
$\dim W=1$ kita peroleh $U\oplus W=\mathbb{R}^{3}$. Setiap perturbasi terpecah tunggal menjadi
bagian yang tidak mengubah skor, yaitu anggota $U$, dan bagian yang mengubahnya, anggota $W$.

\kamus{jumlah langsung}{direct sum}{dua lorong yang hanya bertemu di titik asal, sehingga setiap tempat di gedung dapat dicapai dengan satu jalan yang tidak bisa diperselisihkan.}

\section{Contoh hitung}
\tingkat{menengah}

Empat hitungan berikut memakai seluruh isi bab; semuanya diperiksa ulang dengan Python di
bagian lab.

\subsection*{(a) Rank dan basis ruang kolom}

Kembali ke tabel pembuka. Eliminasi baris: $R_{2}\leftarrow R_{2}-2R_{1}$ mengubah baris
kedua menjadi $(0,0,1,1)$, dan $R_{3}\leftarrow R_{3}-3R_{1}$ mengubah baris ketiga menjadi
$(0,0,2,2)$. Lalu $R_{3}\leftarrow R_{3}-2R_{2}$ menghabisi baris ketiga, sehingga tersisa
$\mathbf{r}_{1}=(1,2,0,1)$ dan $(0,0,1,1)$.

Bentuk itu sudah eselon tereduksi, dengan kolom pivot $1$ dan $3$. Jadi
$\operatorname{rank}A=2$ dan ruang kolomnya berdimensi $2$ dengan basis
$\{\mathbf{c}_{1},\mathbf{c}_{3}\}=\{(1,2,3),(0,1,2)\}$.

Pemeriksaan langsung terhadap kolom yang dibuang:
\[
  \mathbf{c}_{2}=(2,4,6)=2\mathbf{c}_{1},
  \qquad
  \mathbf{c}_{4}=(1,3,5)=\mathbf{c}_{1}+\mathbf{c}_{3}.
\]
Keduanya eksak, tanpa pembulatan. Inilah jawaban ``dua'' untuk pertanyaan pembuka: kolom pivot
memberi himpunan bebas linear yang merentang ruang kolom, dan kolom tanpa pivot adalah
kombinasi kolom pivot di kirinya.

Ruang nol $A$ juga terbaca dari bentuk tereduksi. Dari $x_{1}+2x_{2}+x_{4}=0$ dan
$x_{3}+x_{4}=0$ kita peroleh $x_{1}=-2x_{2}-x_{4}$ dan $x_{3}=-x_{4}$, sehingga
\[
  \mathbf{x}=x_{2}\,(-2,1,0,0)+x_{4}\,(-1,0,-1,1),
\]
ditulis sebagai vektor baris agar ringkas. Jadi $\dim\ker A=2$, dan
$\operatorname{rank}A+\dim\ker A=2+2=4$ sama dengan banyak kolom; pernyataan umumnya adalah
rank-nullity, dibuktikan di Bab 7.

\subsection*{(b) Tiga vektor yang bergantung}

Tinjau $\mathbf{v}_{1}=(1,2,-1)$, $\mathbf{v}_{2}=(0,1,3)$, $\mathbf{v}_{3}=(2,5,1)$ sebagai
kolom $V$. Ekspansi kofaktor sepanjang baris pertama:
\[
  \det V=1(1\cdot1-3\cdot5)+0+2\bigl(2\cdot3-1\cdot(-1)\bigr)=(1-15)+2(6+1)=0.
\]
Determinan nol berarti kolomnya bergantung; relasinya ditebak lalu diperiksa:
$2\mathbf{v}_{1}+\mathbf{v}_{2}=(2,5,1)=\mathbf{v}_{3}$, jadi
$2\mathbf{v}_{1}+\mathbf{v}_{2}-\mathbf{v}_{3}=\mathbf{0}$ dan $\operatorname{rank}V=2$.

\subsection*{(c) Dimensi ruang polinomial}

$\dim P_{3}=4$, karena $\{1,x,x^{2},x^{3}\}$ basis. Himpunan
$B=\{1,\ 1+x,\ 1+x+x^{2},\ 1+x+x^{2}+x^{3}\}$ juga basis: matriks koordinatnya terhadap
basis standar berbentuk segitiga atas dengan diagonal $1$, jadi $B$ bebas linear, dan karena
jumlahnya $4=\dim P_{3}$, akibat (a) membuatnya basis.

Koordinat $p=2-x+3x^{3}$ terhadap $B$ dihitung dari derajat tertinggi ke bawah: $a_{4}=3$,
lalu $a_{3}+a_{4}=0$ memberi $a_{3}=-3$, $a_{2}+a_{3}+a_{4}=-1$ memberi $a_{2}=-1$, dan
konstanta memberi $a_{1}=3$. Jadi $p=3-(1+x)-3(1+x+x^{2})+3(1+x+x^{2}+x^{3})$, yang bila
diuraikan memberi $2-x+3x^{3}$.

\subsection*{(d) Penguraian jumlah langsung}

Ambil $U=\{(x,y,z):x+y+z=0\}$ dan $W=\operatorname{span}\{(1,1,1)\}$ di $\mathbb{R}^{3}$.
Basis $U$ adalah $\{(1,-1,0),(1,0,-1)\}$, jadi $\dim U=2$; dengan $\dim W=1$ dan
$U\cap W=\{\mathbf{0}\}$ kita peroleh $\mathbb{R}^{3}=U\oplus W$.

Uraikan $\mathbf{v}=(3,0,1)$ sebagai $\mathbf{u}+t(1,1,1)$. Syarat jumlah koordinat
$\mathbf{u}$ nol memberi $3+0+1-3t=0$, jadi $t=4/3$. Maka
\[
  \mathbf{v}=\underbrace{\Bigl(\tfrac53,-\tfrac43,-\tfrac13\Bigr)}_{\in U}
  +\underbrace{\Bigl(\tfrac43,\tfrac43,\tfrac43\Bigr)}_{\in W},
\]
dan $\tfrac53-\tfrac43-\tfrac13=0$ memastikan suku pertama memang di $U$. Menurut teorema
kriteria, tidak ada penguraian lain.

\kertas{Ulangi (a) tanpa melihat: kerjakan eliminasi pada $A$, tandai kolom pivot, lalu
tuliskan $\mathbf{c}_{2}$ dan $\mathbf{c}_{4}$ sebagai kombinasi $\mathbf{c}_{1}$ dan
$\mathbf{c}_{3}$. Kemudian hitung $\dim\ker A$ dari bentuk tereduksi dan periksa bahwa
jumlahnya dengan rank sama dengan banyak kolom. Terakhir, untuk matriks $V$ pada (b),
gantilah entri $(3,1)$ menjadi $0$ dan tentukan apakah kolomnya masih bergantung.}

\section{Lab komputasi}
\tingkat{lanjut}

Kode berikut dijalankan apa adanya dengan numpy. Fungsi \texttt{rref} ditulis sendiri supaya
kolom pivot terlihat; baris terakhirnya membersihkan sisa pembulatan dan tanda nol negatif.

\begin{lab}
\begin{verbatim}
import numpy as np

np.set_printoptions(precision=3, suppress=True)

def rref(M):
    M = M.astype(float).copy()
    r = 0
    for c in range(M.shape[1]):
        i = np.argmax(np.abs(M[r:, c])) + r
        if abs(M[i, c]) < 1e-12:
            continue
        M[[r, i]] = M[[i, r]]
        M[r] = M[r] / M[r, c]
        for j in range(M.shape[0]):
            if j != r:
                M[j] = M[j] - M[j, c] * M[r]
        r = r + 1
        if r == M.shape[0]:
            break
    M[np.abs(M) < 1e-12] = 0.0
    return M

A = np.array([[1, 2, 0, 1],
              [2, 4, 1, 3],
              [3, 6, 2, 5]], float)
print("RREF(A):"); print(rref(A))
print("rank(A) =", np.linalg.matrix_rank(A))
print("k2-2k1 =", A[:, 1] - 2 * A[:, 0])
print("k4-k1-k3 =", A[:, 3] - A[:, 0] - A[:, 2])

V = np.array([[1, 0, 2],
              [2, 1, 5],
              [-1, 3, 1]], float)
print("det(V) =", np.linalg.det(V))
print("2v1+v2-v3 =", 2 * V[:, 0] + V[:, 1] - V[:, 2])

x = np.array([0., 1, 2, 3])
M = np.vander(x, 4, increasing=True)
print("det(M) =", np.linalg.det(M))
print("rank(M) =", np.linalg.matrix_rank(M))

B = np.array([[1, 1, 1, 1],
              [0, 1, 1, 1],
              [0, 0, 1, 1],
              [0, 0, 0, 1]], float)
p = np.array([2, -1, 0, 3], float)
k = np.linalg.solve(B, p)
print("koordinat p =", k)
print("sisa =", B @ k - p)
\end{verbatim}
\end{lab}

Keluaran yang tercetak:
\begin{verbatim}
RREF(A):
[[1. 2. 0. 1.]
 [0. 0. 1. 1.]
 [0. 0. 0. 0.]]
rank(A) = 2
k2-2k1 = [0. 0. 0.]
k4-k1-k3 = [0. 0. 0.]
det(V) = 0.0
2v1+v2-v3 = [0. 0. 0.]
det(M) = 12.0
rank(M) = 4
koordinat p = [ 3. -1. -3.  3.]
sisa = [0. 0. 0. 0.]
\end{verbatim}

Empat hal terbaca dari keluaran itu. Kolom pivot $1$ dan $3$ serta
$\operatorname{rank}A=2$ cocok dengan hitungan tangan. Selisih
$\mathbf{c}_{2}-2\mathbf{c}_{1}$ dan $\mathbf{c}_{4}-\mathbf{c}_{1}-\mathbf{c}_{3}$ nol
eksak; pada data nyata selisih itu kecil tetapi tak nol, dan di situlah batas penerapan rank
secara numerik. $\det V=0$ menegaskan ketergantungan di (b), dan matriks Vandermonde titik $0,1,2,3$
berdeterminan $12$ serta berank $4$. Koordinat $p$ tercetak $(3,-1,-3,3)$, cocok dengan
hitungan tangan.

\begin{untukdosen}
Mulai dari tabel pembuka, bukan dari aksioma: minta mahasiswa menebak berapa kolom yang
mandiri, lalu tunjukkan bahwa tebakan ``empat'' salah. Aksioma masuk setelah itu, sebagai
jawaban atas kebutuhan. Uji subruang dan rentang pada pertemuan kedua, kebebasan linear dan
basis pada pertemuan ketiga, jumlah langsung pada pertemuan keempat.

Dua miskonsepsi khas. Pertama, memeriksa subruang hanya dengan ``apakah memuat nol'' lalu
berhenti; kuadran pertama memuat nol dan tertutup terhadap penjumlahan, tetapi gagal pada
kelipatan negatif. Kedua, menganggap basis berarti basis standar, sehingga $(5,7)$ dan
$(4,3)$ dianggap bukan basis $\mathbb{R}^{2}$; tekankan bahwa basis adalah sifat himpunan.

Aktivitas kelas: setiap kelompok menulis tiga vektor di $\mathbb{R}^{3}$ dan memutuskan bebas
atau bergantung lewat eliminasi; kelompok lain mengubah satu entri dan meminta keputusan
ulang. Mahasiswa siap lanjut bila dapat menjelaskan mengapa tiga vektor di $\mathbb{R}^{2}$
selalu bergantung.
\end{untukdosen}

\section{Ringkasan}
\tingkat{dasar}

\begin{ringkasan}
\begin{itemize}
  \item Ruang vektor: himpunan dengan penjumlahan dan perkalian skalar yang memenuhi delapan
        aksioma; anggotanya boleh vektor, matriks, polinomial, atau fungsi.
  \item Subruang cukup diperiksa dengan tiga hal: memuat $\mathbf{0}$ dan tertutup terhadap
        kedua operasi.
  \item $\operatorname{span}(S)$ adalah subruang terkecil yang memuat $S$; $\mathbf{b}$ di
        rentang kolom $A$ tepat ketika $A\mathbf{x}=\mathbf{b}$ konsisten.
  \item Bergantung linear: satu anggota berada di rentang anggota sebelumnya; tiga vektor di
        $\mathbb{R}^{2}$ selalu bergantung.
  \item Teorema pertukaran menutup $n+1$ vektor bebas linear di ruang berrentang $n$ vektor,
        sehingga semua basis sama ukuran dan $\dim V$ terdefinisi.
  \item $\dim(U+W)=\dim U+\dim W-\dim(U\cap W)$; $U+W$ langsung tepat ketika
        $U\cap W=\{\mathbf{0}\}$; setiap subruang berdimensi hingga punya komplemen.
\end{itemize}
\end{ringkasan}

\section{Latihan}
\tingkat{dasar}

\begin{latihan}
\begin{enumerate}
  \item Tentukan apakah $W=\{(x,y,z)\in\mathbb{R}^{3}:x-2y+z=0\ \text{dan}\ x+y=0\}$
        subruang $\mathbb{R}^{3}$. Bila ya, tentukan basis dan dimensinya.
  \item Buktikan bahwa himpunan semua matriks $2\times2$ yang entri diagonal utamanya nol
        adalah subruang $M_{2\times2}(\mathbb{R})$, lalu tentukan basis dan dimensinya.
  \item Tunjukkan bahwa $S=\{1+x,\ 1+x^{2},\ 2+x+x^{2}\}$ bergantung linear di $P_{2}$.
        Tentukan $\dim\operatorname{span}(S)$, lalu perluas himpunan bagian bebas linear
        terbesar dari $S$ menjadi basis $P_{2}$.
  \item Buktikan: jika $\{\mathbf{v}_{1},\mathbf{v}_{2},\mathbf{v}_{3}\}$ bebas linear di
        $V$, maka $\{\mathbf{v}_{1}+\mathbf{v}_{2},\ \mathbf{v}_{2}+\mathbf{v}_{3},\
        \mathbf{v}_{3}+\mathbf{v}_{1}\}$ juga bebas linear.
  \item Hitung $\dim(U\cap W)$ dan $\dim(U+W)$ untuk dua subruang
        \begin{align*}
          U &= \operatorname{span}\{(1,0,0),(0,1,0)\}, \\
          W &= \operatorname{span}\{(1,1,0),(0,0,1)\}
        \end{align*}
        di $\mathbb{R}^{3}$. Periksa rumus dimensi jumlah, lalu hitung juga dengan numpy.
  \item Untuk matriks $3\times4$ dengan baris $(1,2,3,4)$, $(2,4,6,9)$, dan $(0,1,1,1)$,
        tentukan rank, satu basis ruang kolom, dan satu basis ruang nol.
        Periksa dengan fungsi \texttt{matrix\_rank} dari numpy dan
        \texttt{null\_space} dari scipy.
  \item Nilai klaim ini: \emph{jika $U_{1},U_{2},U_{3}$ subruang $V$ dan setiap irisan
        berpasangan sama dengan $\{\mathbf{0}\}$, maka
        $U_{1}+U_{2}+U_{3}=U_{1}\oplus U_{2}\oplus U_{3}$.} Nyatakan benar, salah, atau
        belum cukup bukti, dan sertakan alasan atau contoh penyangkal.
  \item Buktikan bahwa jika $\dim V=n$ dan $U$ subruang dengan $\dim U=n-1$, maka ada
        vektor $\mathbf{v}\in V$ sehingga $V=U\oplus\operatorname{span}\{\mathbf{v}\}$.
  \item Tentukan dimensi ruang semua polinomial berderajat paling tinggi $4$ yang lenyap di
        $x=0$ dan $x=1$ sekaligus, lalu berikan satu basisnya. Periksa jawabanmu dengan
        menghitung rank matriks evaluasi berukuran $2\times5$ di dua titik itu.
  \item Diketahui $B=\{1,\ 1+x,\ 1+x+x^{2},\ 1+x+x^{2}+x^{3}\}$ basis $P_{3}$. Tentukan
        koordinat $q=4+2x-5x^{2}+x^{3}$ terhadap $B$, lalu periksa dengan menyelesaikan
        sistem linear di numpy.
\end{enumerate}

\noindent\textbf{Petunjuk tiga butir tersulit.} Butir 4: andaikan
$\sum c_{i}(\dots)=\mathbf{0}$, kumpulkan koefisien setiap $\mathbf{v}_{i}$, lalu pakai
kebebasan $\{\mathbf{v}_{1},\mathbf{v}_{2},\mathbf{v}_{3}\}$. Butir 8: ambil basis $U$,
perluas dengan satu vektor, lalu pakai akibat (d). Butir 9: polinomial berderajat paling
tinggi $4$ yang lenyap di $0$ dan $1$ habis dibagi $x(x-1)$; hitung parameter bebasnya.
\end{latihan}

\begin{untukpeneliti}
Yang benar-benar terbukti di bab ini ada tiga, dan ketiganya berlaku atas medan apa pun:
ketunggalan ukuran basis lewat teorema pertukaran (hasil Steinitz), rumus
$\dim(U+W)=\dim U+\dim W-\dim(U\cap W)$, dan keberadaan komplemen bagi setiap subruang
berdimensi hingga. Klaim populer bahwa ``rank sama dengan banyaknya fitur informatif'' hanya
bermakna bila ``informatif'' berarti bebas linear, dan pada data nyata syarat itu tidak pernah
terpenuhi secara eksak.

Sumber primer: Anton, \emph{Elementary Linear Algebra} edisi ke-11, Bab 4 bagian 4.1--4.5;
Axler, \emph{Linear Algebra Done Right} edisi ke-4, bagian 1B, 1C, 2A--2C; Hoffman dan Kunze,
\emph{Linear Algebra} edisi ke-2, Bab 2.

Pertanyaan terbuka: seberapa jauh $\dim(U\cap W)$ bergeser akibat perturbasi entri? Batas
keberlakuan: semua hasil mengandaikan aksioma ruang vektor; untuk ruang berdimensi tak hingga,
keberadaan basis memerlukan lemma Zorn. Rank yang dihitung komputer bergantung pada ambang.
\end{untukpeneliti}

\peta{Bab ini mengikuti Anton, \emph{Elementary Linear Algebra} edisi ke-11, Bab 4 (General
Vector Spaces) bagian 4.1 (Real Vector Spaces) untuk aksioma dan contoh, 4.2 (Subspaces) untuk
uji subruang dan rentang, 4.3 (Linear Independence) untuk kebebasan dan lemma ketergantungan,
serta 4.5 (Dimension) untuk teorema pertukaran dan akibatnya; pada edisi lain penomorannya
berbeda. Padanan yang lebih abstrak ada di Axler, \emph{Linear Algebra Done Right} edisi
ke-4, Bab 1 bagian 1B dan 1C (termasuk jumlah dan jumlah langsung subruang) serta Bab 2
bagian 2A, 2B, dan 2C, tempat teorema pertukaran muncul sebagai perbandingan panjang daftar
bebas linear dengan daftar merentang. Pembahasan koordinat, transformasi linear, dan
rank-nullity tidak ada di bab ini melainkan di Bab 7.}

\chapter{Basis, Dimensi, dan Transformasi Linear}

Sebuah detektor anomali tidak membaca lalu lintas jaringan mentah. Ia meringkasnya lebih
dulu. Empat pencacah, misalnya koneksi gagal, paket keluar, lama sesi, dan banyaknya porta
unik, diringkas menjadi tiga fitur oleh peta linear
$F\colon\mathbb{R}^{4}\to\mathbb{R}^{3}$.
Ambang keputusan hanya melihat $F(\mathbf{x})$.

Pertanyaan yang menentukan: adakah $\boldsymbol{\delta}\neq\mathbf{0}$ dengan
$F(\mathbf{x}+\boldsymbol{\delta})=F(\mathbf{x})$? Karena $F$ linear, syarat itu setara
dengan
$F(\boldsymbol{\delta})=\mathbf{0}$. Jawabannya subruang berdimensi
$4-\operatorname{rank}F$; bila $F$ berpangkat tiga, ada satu garis penuh arah buta, dan
penyerang dapat menggeser keempat pencacah sepanjang garis itu tanpa menggeser satu pun
fitur.

Untuk menjawabnya kita perlu tiga hal yang belum diberikan Bab 6: daftar koordinat terhadap
basis, aturan memindahkan daftar itu ketika basis berganti, dan terjemahan peta linear
menjadi matriks sehingga inti dapat dihitung dengan eliminasi. Bab 6 sudah menetapkan bahwa
basis ada dan ukurannya tunggal; aksiomanya tidak diulang di sini.

\begin{untukanak}
Bayangkan sebuah lemari buku dengan banyak rak. Setiap buku punya tempat tetap. Kalau kita
menata ulang lemarinya, buku-bukunya tidak berubah, tetapi alamat setiap buku berubah. Untuk
mengatakan di mana sebuah buku berada, kita harus menyebut penataan yang sedang dipakai.
Kawan kita yang memakai penataan lain akan menyebut alamat yang berbeda untuk buku yang
sama,
dan kalian berdua sama-sama benar.

Ada juga rak khusus di belakang lemari. Buku yang dipindahkan ke sana tidak terlihat dari
depan. Selama buku hanya berpindah di dalam rak itu, daftar buku di depan tidak berubah sama
sekali.

Pernahkah kamu memindahkan sesuatu ke tempat yang tidak diperiksa siapa pun? Kita akan tahu
mengapa daftar alamat bisa berubah sementara isi
lemari tetap, mengapa rak tersembunyi itu selalu punya ukuran tertentu, dan mengapa ada
perubahan yang tidak menggeser ringkasan sama sekali.
\end{untukanak}

\section{Koordinat dan matriks transisi}
\tingkat{menengah}

Vektor di $\mathbb{R}^{3}$ sudah berupa daftar tiga angka; vektor di $P_{2}$ atau
$M_{2\times2}(\mathbb{R})$ tidak. Supaya keduanya dapat dihitung dengan cara yang sama, kita
perlu terjemahan yang mengubah setiap vektor menjadi daftar angka. Terjemahan itu bergantung
pada basis, dan di situlah kesulitan bagian ini bermula.

Ketunggalan sudah kita punya: lemma ketunggalan koordinat di Bab 6 menyatakan penulisan
vektor dalam suatu basis selalu tunggal. Karena itu definisi berikut sah tanpa memilih di
antara beberapa penulisan.

\kamus{koordinat}{coordinate vector}{Alamat sebuah buku di lemari: isinya tetap, alamatnya
bergantung cara menata rak.}

\begin{definisi}[Vektor koordinat]
Misalkan $B=(\mathbf{b}_{1},\dots,\mathbf{b}_{n})$ adalah basis terurut dari ruang vektor
$V$
atas $\mathbb{R}$. Untuk setiap $\mathbf{v}\in V$ ada skalar tunggal $a_{1},\dots,a_{n}$
dengan $\mathbf{v}=a_{1}\mathbf{b}_{1}+\dots+a_{n}\mathbf{b}_{n}$. Kolom
\[
  [\mathbf{v}]_{B}=(a_{1},\dots,a_{n})^{\top}
\]
disebut \emph{vektor koordinat} $\mathbf{v}$ terhadap $B$.\index{koordinat}
\end{definisi}

Kata ``terurut'' bukan hiasan. Bila dua vektor basis ditukar, dua entri koordinat ikut
bertukar sedangkan vektornya tetap. Satu himpunan basis belum menentukan daftar koordinat;
urutannya juga menentukan.

Satu peringatan: $[\mathbf{v}]_{B}$ bukan $\mathbf{v}$. Angka-angka itu hidup di
$\mathbb{R}^{n}$, sedangkan $\mathbf{v}$ tinggal di $V$. Untuk $M_{2\times2}(\mathbb{R})$
bedanya terasa: vektornya berisi empat angka, dan koordinatnya juga empat angka, tetapi
tersusun sebagai kolom yang bergantung pada basis.

Sifat penting pertama dari terjemahan ini: ia menghormati kedua operasi ruang vektor.

\begin{proposisi}
Untuk setiap basis $B$ dari ruang $V$ berdimensi $n$, peta
$\varphi_{B}\colon V\to\mathbb{R}^{n}$ dengan $\varphi_{B}(\mathbf{v})=[\mathbf{v}]_{B}$
bersifat linear, injektif, dan surjektif.
\end{proposisi}

\begin{proof}
Ambil $\mathbf{v}=\sum_{i}a_{i}\mathbf{b}_{i}$ dan $\mathbf{w}=\sum_{i}c_{i}\mathbf{b}_{i}$.
Maka $\mathbf{v}+\mathbf{w}=\sum_{i}(a_{i}+c_{i})\mathbf{b}_{i}$. Ketunggalan koordinat
memaksa $[\mathbf{v}+\mathbf{w}]_{B}=(a_{i}+c_{i})=[\mathbf{v}]_{B}+[\mathbf{w}]_{B}$.
Dengan cara yang sama, $[c\mathbf{v}]_{B}=c[\mathbf{v}]_{B}$.

Untuk injektivitas, misalkan $[\mathbf{v}]_{B}=\mathbf{0}$. Maka semua koefisiennya nol,
sehingga $\mathbf{v}=\sum_{i}0\cdot\mathbf{b}_{i}=\mathbf{0}$. Untuk surjektivitas, ambil
sembarang $(a_{1},\dots,a_{n})\in\mathbb{R}^{n}$ dan bentuk
$\mathbf{v}=\sum_{i}a_{i}\mathbf{b}_{i}$; koordinatnya tepat daftar itu.
\end{proof}

Peta $\varphi_{B}$ memindahkan pertanyaan tentang $V$ menjadi pertanyaan tentang kolom di
$\mathbb{R}^{n}$, lalu memindahkannya kembali; pemindahan itu bergantung pada $B$. Koordinat
ditulis sebagai kolom agar matriks bekerja dari kiri: koordinat masuk, dikalikan matriks,
koordinat keluar. Dengan koordinat baris, setiap rumus harus ditranspos.

\begin{bacakeras}
Dua suara berdiri di depan kelas.

``Berapa panjang meja ini?''
``Tiga.''
``Tiga apa?''
``Tiga penggaris. Atau tiga jengkal. Atau tiga langkah. Pilih dulu alat ukurnya.''

Begitu juga ruang vektor. Sebuah vektor belum punya alamat sampai kita memilih basis. Basis
itu
penggaris: ia memberi skala, memberi arah, dan memberi urutan. Ganti penggarisnya, ganti
angkanya. Benda yang diukur tidak ikut berubah.

Satu penggaris, satu angka. Dua penggaris, dua angka. Tanpa penggaris, tidak ada angka sama
sekali.

Hati-hati dengan penggaris yang jelek: dua garis skalanya hampir sejajar. Angka tetap
keluar,
tetapi angkanya gemetar, dan dua orang yang mengukur dengan sabar akan mendapat angka
berbeda.

Ulangi bersama: basis bukan bendanya, basis alat ukurnya. Yang berubah alamat, yang tetap
isi.
\end{bacakeras}

Sekarang dua basis dipakai pada ruang yang sama. Setiap vektor punya dua daftar koordinat,
dan kita ingin rumus yang mengubah daftar pertama menjadi daftar kedua. Matriks yang
melakukan tugas itu disusun dari koordinat vektor basis lama dalam basis baru.

\begin{definisi}[Matriks transisi]
Misalkan $B=(\mathbf{b}_{1},\dots,\mathbf{b}_{n})$ dan
$C=(\mathbf{c}_{1},\dots,\mathbf{c}_{n})$
dua basis dari $V$. \emph{Matriks transisi} dari $B$ ke $C$ adalah matriks $n\times n$
\[
P_{C\leftarrow
B}=\Bigl(\,[\mathbf{b}_{1}]_{C}\;\Big|\;\cdots\;\Big|\;[\mathbf{b}_{n}]_{C}\,\Bigr),
\]
yaitu matriks yang kolom ke-$j$-nya adalah $[\mathbf{b}_{j}]_{C}$.\index{transisi}
\end{definisi}

Yang perlu dibuktikan: mengalikan daftar koordinat lama dengan matriks ini menghasilkan
daftar
koordinat baru.

\begin{teorema}
Untuk setiap $\mathbf{v}\in V$ berlaku $[\mathbf{v}]_{C}=P_{C\leftarrow B}[\mathbf{v}]_{B}$.
\end{teorema}

\begin{proof}
Tulis $\mathbf{v}=\sum_{j}a_{j}\mathbf{b}_{j}$, sehingga
$[\mathbf{v}]_{B}=(a_{1},\dots,a_{n})^{\top}$.
Kelinearan $\varphi_{C}$ memberi
$[\mathbf{v}]_{C}=\sum_{j}a_{j}[\mathbf{b}_{j}]_{C}$. Ruas kanan itu tepat hasil kali
matriks
$P_{C\leftarrow B}$ dengan kolom $(a_{j})$, menurut definisi perkalian matriks dengan vektor.
\end{proof}

Akibat yang membuat matriks itu dapat dipakai dua arah.

\begin{akibat}
Matriks transisi selalu invertible, dan $P_{B\leftarrow C}=P_{C\leftarrow B}^{-1}$. Selain itu
$P_{C\leftarrow B}=P_{C\leftarrow D}\,P_{D\leftarrow B}$ untuk tiga basis $B,C,D$ dari $V$.
\end{akibat}

\begin{proof}
Teorema sebelumnya memberi $P_{B\leftarrow C}P_{C\leftarrow B}[\mathbf{v}]_{B}=[\mathbf{v}]_{B}$
untuk semua $\mathbf{v}$. Karena $\varphi_{B}$ surjektif, hasil kali itu sama dengan matriks
identitas. Argumen yang sama berlaku pada urutan sebaliknya, jadi kedua matriks saling
invers. Untuk aturan rantai, terapkan teorema dua kali:
$[\mathbf{v}]_{C}=P_{C\leftarrow D}[\mathbf{v}]_{D}=P_{C\leftarrow D}P_{D\leftarrow B}[\mathbf{v}]_{B}$,
lalu bandingkan dengan $[\mathbf{v}]_{C}=P_{C\leftarrow B}[\mathbf{v}]_{B}$ pada $\mathbf{v}=\mathbf{b}_{j}$.
\end{proof}

Perhatikan bahwa $P_{C\leftarrow B}$ dan $P_{B\leftarrow C}$ bukan transpos satu sama lain; keduanya saling invers, dan itu berbeda. Arah subskrip dibaca dari kanan ke kiri, seperti komposisi peta: ambil vektor basis ke-$j$ dari $B$, lalu tanyakan alamatnya menurut $C$. Bila $B$ basis standar, kolom-kolom $P_{B\leftarrow C}$ adalah vektor-vektor $C$ itu sendiri.

Ambil $B=(\mathbf{e}_{1},\mathbf{e}_{2},\mathbf{e}_{3})$
dan $C=(\mathbf{b}'_{1},\mathbf{b}'_{2},\mathbf{b}'_{3})$ dengan
\[
  \mathbf{b}'_{1}=(1,1,0),\qquad \mathbf{b}'_{2}=(1,0,1),\qquad \mathbf{b}'_{3}=(0,1,1).
\]
Karena $B$ standar, matriks $P_{B\leftarrow C}$ memuat ketiga vektor itu sebagai kolom:
\[
  P_{B\leftarrow C}=\begin{pmatrix}1&1&0\\ 1&0&1\\ 0&1&1\end{pmatrix},
  \qquad \det P_{B\leftarrow C}=-2\neq0 .
\]
Inversnya
\[
  P_{C\leftarrow B}=\frac{1}{2}\begin{pmatrix}1&1&-1\\ 1&-1&1\\ -1&1&1\end{pmatrix}.
\]
Untuk $\mathbf{v}=(3,8,11)$ kita peroleh
\[
  [\mathbf{v}]_{C}=P_{C\leftarrow B}[\mathbf{v}]_{B}
  =\frac{1}{2}\begin{pmatrix}3+8-11\\ 3-8+11\\ -3+8+11\end{pmatrix}
  =\begin{pmatrix}0\\ 3\\ 8\end{pmatrix}.
\]
Pemeriksaannya langsung: $3\mathbf{b}'_{2}+8\mathbf{b}'_{3}=(3,8,11)$. Sebaliknya
$\mathbf{w}=(1,0,0)$ berkoordinat $(1/2,\,1/2,\,-1/2)$; entri pecahan wajar, karena
koordinat
hanyalah alamat.

\section{Transformasi linear, inti, dan jangkauan}
\tingkat{dasar}

Bab 6 membahas ruang dan subruang, tetapi belum membahas peta di antaranya. Padahal
pekerjaan utama aljabar linear adalah memindahkan vektor: memutar, memproyeksikan,
menurunkan, meringkas.

Aturan yang menghormati penjumlahan dan perkalian skalar disebut transformasi linear:
menjumlahkan dulu lalu memetakan sama hasilnya dengan memetakan dulu lalu menjumlahkan.

\begin{definisi}[Transformasi linear]
Peta $T\colon V\to W$ antara dua ruang vektor atas $\mathbb{R}$ disebut
\emph{transformasi linear} bila $T(\mathbf{u}+\mathbf{v})=T(\mathbf{u})+T(\mathbf{v})$ dan
$T(c\mathbf{v})=cT(\mathbf{v})$ untuk semua $\mathbf{u},\mathbf{v}\in V$ dan $c\in\mathbb{R}$.
\end{definisi}

Kedua syarat setara dengan satu syarat,
$T(c\mathbf{u}+\mathbf{v})=cT(\mathbf{u})+T(\mathbf{v})$. Konsekuensi pertama,
$T(\mathbf{0})=\mathbf{0}$, dipakai untuk menolak kelinearan dengan cepat: peta
$T(\mathbf{x})=\mathbf{x}+\mathbf{a}$ dengan $\mathbf{a}\neq\mathbf{0}$ langsung gugur.

Sebaliknya $T(\mathbf{0})=\mathbf{0}$ tidak cukup. Peta $T(x,y)=(x^{2},y)$ memenuhinya,
tetapi
$T(2,0)=(4,0)$ sedangkan $2T(1,0)=(2,0)$. Kuadrat memutus kelinearan.\index{transformasi
linear}

Rotasi $90^{\circ}$, $R(x,y)=(-y,x)$, linear: $R(3,1)=(-1,3)$. Turunan $D\colon P_{3}\to P_{2}$, $D(p)=p'$, linear karena $(p+q)'=p'+q'$. Peta $T(x,y)=(x+y,\,2x+2y)$ juga linear, tetapi seluruh $\mathbb{R}^{2}$ jatuh ke satu garis; ia linear sekaligus merusak informasi.

Ambil
$A=\begin{pmatrix}2&1\\1&2\end{pmatrix}$ dan $T(\mathbf{x})=A\mathbf{x}$. Vektor $(1,0)$ pergi
ke $(2,1)$ dan $(0,1)$ ke $(1,2)$, sehingga persegi satuan menjadi jajar genjang bertitik
sudut
$(0,0)$, $(2,1)$, $(3,3)$, $(1,2)$ dengan luas $3$, sebab $|2\cdot2-1\cdot1|=3$.
Gambar~\ref{fig:kisi-07} memperlihatkan kisi sebelum dan sesudah peta: garis tetap garis, kisi
tetap kisi, tetapi jarak dan luas tidak lagi sama.

\begin{figure}[htbp]
\centering
\begin{tikzpicture}[scale=0.38,>=stealth,font=\scriptsize]
  \foreach \k in {0,1,2,3}{
    \draw[layergrey!45] (\k,0)--(\k,3);
    \draw[layergrey!45] (0,\k)--(3,\k);
    \draw[layerteal,thick] (2*\k,\k)--(2*\k+3,\k+6);
    \draw[layerteal,thick] (\k,2*\k)--(\k+6,2*\k+3);}
  \fill[accent!15] (0,0)--(1,0)--(1,1)--(0,1)--cycle;
  \draw[accent,thick] (0,0)--(1,0)--(1,1)--(0,1)--cycle;
  \fill[layerorange!25] (0,0)--(2,1)--(3,3)--(1,2)--cycle;
  \draw[layerorange,very thick] (0,0)--(2,1)--(3,3)--(1,2)--cycle;
  \draw[->,ink,thick] (-0.6,0)--(9.8,0) node[below right]{$x$};
  \draw[->,ink,thick] (0,-0.6)--(0,9.8) node[left]{$y$};
  \foreach \x/\y in {2/1,1/2,3/3}{
    \fill[ink] (\x,\y) circle (1.8pt);
    \node[ink,anchor=north west] at (\x+0.05,\y-0.05) {$(\x,\y)$};}
  \node[layergrey,anchor=north west] at (0,-1.0) {tipis: kisi asal};
  \node[layerteal,anchor=north west] at (4.0,-1.0) {tebal: kisi setelah dipetakan};
\end{tikzpicture}
\caption{Kisi satuan dan bayangannya di bawah $T(\mathbf{x})=A\mathbf{x}$ dengan
$A=\begin{pmatrix}2&1\\1&2\end{pmatrix}$. Satu petak berubah menjadi jajar genjang bertitik
sudut $(0,0)$, $(2,1)$, $(3,3)$, dan $(1,2)$, sehingga luasnya naik dari $1$ menjadi $3$.
Kisi lama digambar tipis, kisi baru digambar tebal; keduanya tetap kisi.}
\label{fig:kisi-07}
\end{figure}

\begin{tangan}
Ambil selembar kertas berpetak, pensil, dan penggaris. Gambar sumbu mendatar dan sumbu
tegak,
lalu tandai perpotongannya sebagai pusat. Gambar sebuah rumah kecil: persegi lima petak,
atap
segitiga, pintu satu petak.

Sekarang geser. Pindahkan setiap titik rumah itu tiga petak ke kanan dan dua petak ke atas.
Jangan menebak; hitung. Titik yang tadinya satu petak ke kanan dan satu petak ke atas
sekarang
berada empat petak ke kanan dan tiga petak ke atas. Gambar ulang rumah itu. Bentuk dan
ukurannya
sama, hanya tempatnya berpindah.

Sekarang putar. Putar kertas seperempat putaran searah jarum jam, lalu gambar ulang rumah
yang
sama. Panjang setiap sisi tetap, arah setiap sisi berubah.

Yang harus terlihat: pada geseran, semua arah bertahan tetapi pusat ikut berpindah; pada
putaran, pusat bertahan dan semua arah berubah. Karena geseran memindahkan pusat, geseran
bukan
transformasi linear. Putaran termasuk.
\end{tangan}

Kabar baiknya: untuk mengenal peta linear, cukup mengenal nilainya pada vektor basis.

\begin{proposisi}
Misalkan $B=(\mathbf{b}_{1},\dots,\mathbf{b}_{n})$ basis $V$ dan
$\mathbf{w}_{1},\dots,\mathbf{w}_{n}\in W$ sebarang. Ada tepat satu peta linear
$T\colon V\to W$ dengan $T(\mathbf{b}_{j})=\mathbf{w}_{j}$ untuk setiap $j$.
\end{proposisi}

\begin{proof}
Untuk keberadaan, definisikan
$T\bigl(\sum_{j}a_{j}\mathbf{b}_{j}\bigr)=\sum_{j}a_{j}\mathbf{w}_{j}$.
Definisi ini tidak bergantung pada penulisan, karena koordinat terhadap basis tunggal. Untuk
kelinearan, pakai kelinearan koordinat dari proposisi di bagian sebelumnya:
$T(\mathbf{u}+\mathbf{v})=\sum_{j}(a_{j}+c_{j})\mathbf{w}_{j}=T(\mathbf{u})+T(\mathbf{v})$,
dan serupa untuk perkalian skalar.

Untuk ketunggalan, misalkan $S$ linear dan $S(\mathbf{b}_{j})=\mathbf{w}_{j}$. Untuk
$\mathbf{v}=\sum_{j}a_{j}\mathbf{b}_{j}$ berlaku
$S(\mathbf{v})=\sum_{j}a_{j}S(\mathbf{b}_{j})=\sum_{j}a_{j}\mathbf{w}_{j}=T(\mathbf{v})$.
\end{proof}

Akibat praktisnya: $n$ vektor di $W$ menentukan satu peta linear dari $V$ berdimensi $n$.
Bila $W=\mathbb{R}^{m}$, data itu tidak lain matriks $m\times n$.

Setiap peta linear membawa dua himpunan yang layak diperiksa: vektor yang dipetakan ke nol,
dan nilai yang benar-benar tercapai.

\kamus{inti}{kernel atau null space}{Daftar perubahan yang sama sekali tidak menggeser
ringkasan.}

\begin{definisi}[Inti, jangkauan, rank, dan nullity]
Untuk $T\colon V\to W$ linear, \emph{inti} adalah
$\ker T=\{\mathbf{v}\in V: T(\mathbf{v})=\mathbf{0}\}$ dan \emph{jangkauan} adalah
$\operatorname{im}T=\{T(\mathbf{v}):\mathbf{v}\in V\}$. Bila $V$ berdimensi hingga,
$\operatorname{nullity}(T)=\dim\ker T$ dan $\operatorname{rank}(T)=\dim\operatorname{im}T$.
\end{definisi}

Keduanya subruang, bukan sekadar himpunan.\index{inti}\index{jangkauan}

\begin{lemma}
$\ker T$ subruang $V$ dan $\operatorname{im}T$ subruang $W$.
\end{lemma}

\begin{proof}
Keduanya memuat $\mathbf{0}$ karena $T(\mathbf{0})=\mathbf{0}$. Bila
$\mathbf{u},\mathbf{v}\in\ker T$
dan $c\in\mathbb{R}$, maka $T(\mathbf{u}+\mathbf{v})=\mathbf{0}$ dan
$T(c\mathbf{u})=\mathbf{0}$,
jadi $\ker T$ tertutup. Bila $\mathbf{u}=T(\mathbf{x})$ dan $\mathbf{v}=T(\mathbf{y})$, maka
$\mathbf{u}+\mathbf{v}=T(\mathbf{x}+\mathbf{y})$ dan $c\mathbf{u}=T(c\mathbf{x})$ anggota
$\operatorname{im}T$.
\end{proof}

Inti mengukur kegagalan injektivitas, dan hanya itu yang perlu diperiksa.

\begin{proposisi}
$T$ injektif jika dan hanya jika $\ker T=\{\mathbf{0}\}$.
\end{proposisi}

\begin{proof}
Bila $T$ injektif, hanya $\mathbf{0}$ yang punya peta $\mathbf{0}$. Sebaliknya andaikan
$\ker T=\{\mathbf{0}\}$ dan $T(\mathbf{u})=T(\mathbf{v})$. Maka $T(\mathbf{u}-\mathbf{v})=\mathbf{0}$,
jadi $\mathbf{u}-\mathbf{v}\in\ker T$, sehingga $\mathbf{u}=\mathbf{v}$.
\end{proof}

Rank menghitung arah yang bertahan, nullity menghitung arah yang lenyap. Keduanya membagi
dimensi ruang asal.

\kamus{rank}{rank}{Banyaknya arah bebas yang masih bertahan setelah dipetakan.}

Contoh kecil memperlihatkan kedua nasib itu. Ambil
$B=\begin{pmatrix}1&2\\2&4\end{pmatrix}$. Maka $B\mathbf{e}_{1}=(1,2)$ dan
$B\mathbf{e}_{2}=(2,4)$, dua vektor segaris, jadi $\operatorname{rank}(B)=1$. Arah $(1,2)$
bertahan dan memanjang: $B(1,2)=(5,10)$. Arah $(2,-1)$ lenyap: $B(2,-1)=(0,0)$.
Gambar~\ref{fig:arah} menggambar keduanya, dan rank-nullity menutup hitungan: $2=1+1$.

\begin{figure}[htbp]
\centering
\begin{tikzpicture}[scale=0.60,>=stealth,font=\scriptsize]
  \draw[layerjade!60,thin] (-0.7,-1.4)--(2.8,5.6);
  \node[layerjade,anchor=west,rotate=63] at (2.0,4.3) {garis arah $(1,2)$};
  \draw[->,ink,thick] (-1.0,0)--(3.3,0) node[below right]{$x$};
  \draw[->,ink,thick] (0,-2.0)--(0,5.0) node[left]{$y$};
  \draw[->,layergrey,thick] (0,0)--(1,0) node[below]{$\mathbf{e}_{1}$};
  \draw[->,layergrey,thick] (0,0)--(0,1) node[left]{$\mathbf{e}_{2}$};
  \draw[->,layerorange,very thick] (0,0)--(1,2) node[above left]{$B\mathbf{e}_{1}=(1,2)$};
  \draw[->,layerbrown,very thick] (0,0)--(2,4) node[right]{$B\mathbf{e}_{2}=(2,4)$};
  \draw[->,layerplum,thick,dashed] (0,0)--(1.8,-0.9) node[below right]{$(2,-1)$};
  \node[layerplum,anchor=north] at (1.5,-1.6) {runtuh ke $\mathbf{0}$};
  \fill[ink] (0,0) circle (1.6pt) node[below left,ink]{$\mathbf{0}$};
\end{tikzpicture}
\caption{Dua arah di bawah $B=\begin{pmatrix}1&2\\2&4\end{pmatrix}$. Arah $(1,2)$ bertahan dan
memanjang menjadi $(5,10)$; arah $(2,-1)$ runtuh ke $\mathbf{0}$. Bayangan kedua vektor basis,
$B\mathbf{e}_{1}=(1,2)$ dan $B\mathbf{e}_{2}=(2,4)$, terletak pada satu garis, jadi
$\operatorname{rank}(B)=1$ dan $\operatorname{nullity}(B)=1$.}
\label{fig:arah}
\end{figure}

Sekarang teorema pusat bab ini. Isinya: tidak ada dimensi yang hilang di tengah jalan.

\begin{teorema}[Rank-nullity]
Bila $V$ berdimensi hingga dan $T\colon V\to W$ linear, maka
\[
  \dim V=\dim\ker T+\dim\operatorname{im}T .
\]
\end{teorema}

\begin{proof}
Misalkan $\dim V=n$ dan $\dim\ker T=k$. Ambil basis
$(\mathbf{u}_{1},\dots,\mathbf{u}_{k})$ dari $\ker T$. Menurut akibat perluasan basis di Bab 6,
basis itu dapat dilengkapi menjadi basis $V$ dengan menambahkan $n-k$ vektor
$\mathbf{w}_{1},\dots,\mathbf{w}_{n-k}$.

Pertama, himpunan $\{T(\mathbf{w}_{1}),\dots,T(\mathbf{w}_{n-k})\}$ merentang
$\operatorname{im}T$. Untuk sembarang $T(\mathbf{v})$, tulis
$\mathbf{v}=\sum_{i}\alpha_{i}\mathbf{u}_{i}+\sum_{j}\beta_{j}\mathbf{w}_{j}$. Suku-suku
$\mathbf{u}_{i}$ lenyap karena $T(\mathbf{u}_{i})=\mathbf{0}$, sehingga
$T(\mathbf{v})=\sum_{j}\beta_{j}T(\mathbf{w}_{j})$.

Kedua, himpunan itu bebas linear. Andaikan $\sum_{j}\beta_{j}T(\mathbf{w}_{j})=\mathbf{0}$.
Maka $T\bigl(\sum_{j}\beta_{j}\mathbf{w}_{j}\bigr)=\mathbf{0}$, jadi
$\sum_{j}\beta_{j}\mathbf{w}_{j}\in\ker T$, dan ada skalar $\alpha_{i}$ dengan
$\sum_{j}\beta_{j}\mathbf{w}_{j}=\sum_{i}\alpha_{i}\mathbf{u}_{i}$. Pindahkan semua suku ke
satu ruas:
\[
  \sum_{j}\beta_{j}\mathbf{w}_{j}-\sum_{i}\alpha_{i}\mathbf{u}_{i}=\mathbf{0}.
\]
Relasi itu berada di dalam basis $V$, sehingga semua koefisiennya nol; khususnya setiap
$\beta_{j}=0$. Jadi $\{T(\mathbf{w}_{j})\}$ basis $\operatorname{im}T$, dan
$\operatorname{rank}(T)=n-k$.
\end{proof}

Gambar~\ref{fig:batang} menyusun teorema itu sebagai dua batang. Untuk ruang berdimensi
$4$ panjang totalnya selalu $4$: identitas memberi $4+0$, proyeksi yang membuang satu
koordinat
memberi $3+1$, dan peta fitur $G$ nanti memberi $2+2$. Yang bergeser hanya letak pemisahnya.

\begin{figure}[htbp]
\centering
\begin{tikzpicture}[x=1cm,y=1cm,font=\scriptsize,>=stealth]
  \fill[layerjade!45] (0,2.4) rectangle (4,3.1);
  \draw[layergrey] (0,2.4) rectangle (4,3.1);
  \node at (2,2.75) {rank $4$};
  \node[layerjade,anchor=east] at (-0.12,2.75) {identitas};
  \node[layerjade,anchor=west] at (4.12,2.75) {nullity $0$};
  \fill[layerteal!45] (0,1.2) rectangle (3,1.9);
  \fill[layerorange!45] (3,1.2) rectangle (4,1.9);
  \draw[layergrey] (0,1.2) rectangle (4,1.9);
  \draw[layergrey] (3,1.2)--(3,1.9);
  \node at (1.5,1.55) {rank $3$};
  \node at (3.5,1.55) {$1$};
  \node[layerteal,anchor=east] at (-0.12,1.55) {proyeksi};
  \node[layerorange,anchor=west] at (4.12,1.55) {nullity $1$};
  \fill[layerteal!45] (0,0) rectangle (2,0.7);
  \fill[layerorange!45] (2,0) rectangle (4,0.7);
  \draw[layergrey] (0,0) rectangle (4,0.7);
  \draw[layergrey] (2,0)--(2,0.7);
  \node at (1,0.35) {rank $2$};
  \node at (3,0.35) {$2$};
  \node[layerteal,anchor=east] at (-0.12,0.35) {peta fitur $G$};
  \node[layerorange,anchor=west] at (4.12,0.35) {nullity $2$};
  \draw[layergrey,dashed] (4,-0.25)--(4,3.4);
  \node[layergrey,anchor=south] at (4,3.4) {$\dim V=4$};
  \fill[layerteal!45] (0,-0.95) rectangle (0.35,-0.6);
  \draw[layergrey] (0,-0.95) rectangle (0.35,-0.6);
  \node[anchor=west] at (0.45,-0.78) {arah yang bertahan};
  \fill[layerorange!45] (3.3,-0.95) rectangle (3.65,-0.6);
  \draw[layergrey] (3.3,-0.95) rectangle (3.65,-0.6);
  \node[anchor=west] at (3.75,-0.78) {arah yang lenyap};
\end{tikzpicture}
\caption{Rank-nullity sebagai dua batang yang jumlahnya tetap. Untuk $\dim V=4$: peta identitas
memberi $4+0$, proyeksi yang membuang satu koordinat memberi $3+1$, dan peta fitur $G$ memberi
$2+2$. Hanya letak pemisahnya yang bergeser; panjang totalnya tetap $4$.}
\label{fig:batang}
\end{figure}

Dua akibat yang paling sering dipakai.

\begin{akibat}
Misalkan $\dim V=\dim W$ berhingga dan $T\colon V\to W$ linear. Maka $T$ injektif jika dan
hanya jika $T$ surjektif. Selain itu
$\operatorname{rank}(T)\le\min(\dim V,\dim W)$.
\end{akibat}

\begin{proof}
Menurut rank-nullity, $\dim\ker T=n-\operatorname{rank}T$ dengan $n=\dim V$. Injektif tepat
ketika $\dim\ker T=0$, yaitu $\operatorname{rank}T=n=\dim W$, yang berarti
$\operatorname{im}T=W$ karena subruang berdimensi penuh memenuhi seluruh ruang. Implikasi
sebaliknya diperoleh dengan membalik urutan argumen. Batas atas rank jelas dari
$\operatorname{im}T\subseteq W$ dan dari rank-nullity.
\end{proof}

Sebagai gambaran angka, ambil $T\colon\mathbb{R}^{3}\to\mathbb{R}^{2}$ dengan
$T(x,y,z)=(x-y,\,y-z)$. Karena $T(t,t,t)=(0,0)$, intinya garis $\{(t,t,t)\}$ dan nullity
$1$.
Bayangan $T(1,0,0)=(1,0)$ dan $T(0,0,1)=(0,-1)$ merentang $\mathbb{R}^{2}$, jadi rank $2$;
rank-nullity: $3=1+2$.

Bagi pemelajaran mesin, akibat itu keras. Peta fitur $\mathbb{R}^{n}\to\mathbb{R}^{m}$
dengan $n>m$ tidak mungkin injektif, dan intinya berdimensi paling sedikit $n-m$. Setiap
pemampatan dimensi menyimpan arah buta; pertanyaannya bukan apakah arah itu ada, melainkan
apakah penyerang dapat menemukannya.

Prosedur baku untuk basis inti: bawa matriks ke bentuk eselon tereduksi, tandai kolom
tanpa pivot sebagai variabel bebas, lalu gilir setiap variabel bebas dengan nilai satu dan
variabel bebas lain nol. Variabel pivot dibaca dari baris-barisnya. Hasilnya satu vektor per
kolom tanpa pivot, bebas linear karena masing-masing punya posisi satu yang tidak dimiliki
yang
lain.

Untuk jangkauan, prosedurnya berbeda dan sering tertukar dengan prosedur inti di atas. Empat
langkah berikut cukup dihafal.

\begin{resep}
\begin{enumerate}
  \item Susun matriks $A$, lalu bawa ke bentuk eselon baris tereduksi dengan eliminasi.
  \item Tandai kolom yang memuat pivot. Kolom tanpa pivot bukan calon anggota basis.
  \item Kembali ke matriks asli $A$, bukan ke bentuk eselonnya. Ambil kolom asli yang posisinya
        sama dengan kolom pivot.
  \item Periksa: banyaknya vektor terpilih sama dengan banyaknya baris tak nol, dan vektor itu
        bebas linear.
\end{enumerate}
Peringatan urutan: langkah tiga tidak boleh mengambil kolom dari bentuk eselon. Operasi
baris
mengubah kolom, sedangkan ruang kolom tidak berubah. Mengambil kolom dari matriks eselon
memberi jawaban salah walaupun banyaknya benar.
\end{resep}

Inti juga menjawab pertanyaan yang lebih umum: bila satu nilai sudah tercapai, semua
prapetanya tersusun rapi. Di sini $T^{-1}(\mathbf{w})$ berarti himpunan prapeta, bukan peta
invers.

\begin{proposisi}
Misalkan $T\colon V\to W$ linear dan $T(\mathbf{v}_{0})=\mathbf{w}$. Maka
$T^{-1}(\mathbf{w})=\{\mathbf{v}_{0}+\mathbf{k}:\mathbf{k}\in\ker T\}$.
\end{proposisi}

\begin{proof}
Bila $\mathbf{k}\in\ker T$, maka $T(\mathbf{v}_{0}+\mathbf{k})=\mathbf{w}+\mathbf{0}=\mathbf{w}$,
jadi setiap $\mathbf{v}_{0}+\mathbf{k}$ termasuk prapeta $\mathbf{w}$. Sebaliknya, bila
$T(\mathbf{v})=\mathbf{w}$, maka $T(\mathbf{v}-\mathbf{v}_{0})=\mathbf{0}$, sehingga
$\mathbf{v}-\mathbf{v}_{0}\in\ker T$ dan $\mathbf{v}=\mathbf{v}_{0}+(\mathbf{v}-\mathbf{v}_{0})$.
\end{proof}

Himpunan masukan yang menghasilkan fitur sama adalah salinan inti yang digeser: bila intinya
berdimensi satu, keluarganya satu garis, yaitu arah buta pada pembuka bab; bila berdimensi
$k$,
penyerang punya $k$ arah.

\section{Matriks representasi dan komposisi}
\tingkat{menengah}

Sekarang dua alat sebelumnya digabung: koordinat mengubah vektor menjadi kolom, dan peta
linear ditentukan nilainya pada basis. Gabungan keduanya menghasilkan matriks, tempat
seluruh
perhitungan dilakukan.

Matriks representasi menyimpan bayangan setiap vektor basis sebagai satu kolom, ditulis
dalam koordinat basis tujuan.

\begin{definisi}[Matriks representasi]
Misalkan $B=(\mathbf{b}_{1},\dots,\mathbf{b}_{n})$ basis $V$,
$C=(\mathbf{c}_{1},\dots,\mathbf{c}_{m})$ basis $W$, dan $T\colon V\to W$ linear.
\emph{Matriks representasi} $T$ terhadap $B$ dan $C$ adalah matriks $m\times n$
\[
[T]_{C\leftarrow
B}=\Bigl(\,[T(\mathbf{b}_{1})]_{C}\;\Big|\;\cdots\;\Big|\;[T(\mathbf{b}_{n})]_{C}\,\Bigr).
\]
\end{definisi}

Kolom-kolom itu diperoleh dengan menyelesaikan sistem, bukan menebak: untuk setiap $j$,
selesaikan $\sum_{i}x_{i}\mathbf{c}_{i}=T(\mathbf{b}_{j})$, lalu tulis solusinya sebagai
kolom
ke-$j$. Bila $C$ basis standar, penyelesaian itu tidak perlu dikerjakan, karena koordinat
terhadap basis standar adalah vektornya sendiri.

Klaim berikut menjelaskan mengapa matriks itu berguna: ia bekerja dengan cara yang
diharapkan.

\begin{teorema}
Untuk setiap $\mathbf{v}\in V$ berlaku
$[T(\mathbf{v})]_{C}=[T]_{C\leftarrow B}\,[\mathbf{v}]_{B}$.
\end{teorema}

\begin{proof}
Tulis $\mathbf{v}=\sum_{j}a_{j}\mathbf{b}_{j}$, jadi $[\mathbf{v}]_{B}=(a_{j})$. Kelinearan
$T$ memberi $T(\mathbf{v})=\sum_{j}a_{j}T(\mathbf{b}_{j})$. Pengambilan koordinat juga
linear,
sehingga $[T(\mathbf{v})]_{C}=\sum_{j}a_{j}[T(\mathbf{b}_{j})]_{C}$, yaitu tepat hasil kali
matriks dengan kolom $(a_{j})$.
\end{proof}

Dua akibat langsung. Peta linear dan matriksnya saling menentukan, karena nilai pada
basis sudah cukup. Dan rank tidak berubah dalam penerjemahan itu.

\begin{proposisi}
$\operatorname{rank}(T)$ sama dengan rank matriks $[T]_{C\leftarrow B}$, dan $\ker T$
berkorespondensi dengan ruang nol matriks itu. Rank matriks yang dimaksud adalah dimensi
ruang
kolomnya, sebagaimana sudah dipakai di Bab 6.
\end{proposisi}

\begin{proof}
Kolom ke-$j$ matriks adalah koordinat $T(\mathbf{b}_{j})$, dan
$\{T(\mathbf{b}_{j})\}$ merentang $\operatorname{im}T$. Pengambilan koordinat mempertahankan
kombinasi linear, jadi dimensi jangkauan sama dengan dimensi ruang kolom. Teorema sebelumnya
menunjukkan $T(\mathbf{v})=\mathbf{0}$ tepat ketika $[T(\mathbf{v})]_{C}=[T]_{C\leftarrow B}[\mathbf{v}]_{B}=\mathbf{0}$,
yaitu ketika $[\mathbf{v}]_{B}$ berada di ruang nol matriks itu.
\end{proof}

Komposisi peta berpadanan dengan perkalian matriks. Itu bukan kebetulan, melainkan
alasan perkalian matriks didefinisikan seperti sekarang.

\begin{teorema}
Misalkan $T\colon V\to W$ dan $S\colon W\to U$ linear, dengan basis $B,C,E$ pada ketiga ruang.
Maka $[S\circ T]_{E\leftarrow B}=[S]_{E\leftarrow C}\,[T]_{C\leftarrow B}$.
\end{teorema}

\begin{proof}
Untuk setiap $\mathbf{v}\in V$,
\[
  [(S\circ T)(\mathbf{v})]_{E}
  =[S]_{E\leftarrow C}[T(\mathbf{v})]_{C}
  =[S]_{E\leftarrow C}[T]_{C\leftarrow B}[\mathbf{v}]_{B}.
\]
Jadi matriks $[S]_{E\leftarrow C}[T]_{C\leftarrow B}$ memenuhi definisi matriks representasi
untuk $S\circ T$. Matriks representasi tunggal: bila dua matriks memberi hasil sama pada
semua $\mathbf{v}$, hasilnya sama khususnya pada $\mathbf{v}=\mathbf{b}_{j}$, sehingga
kolom-kolomnya sama.
\end{proof}

Komposisi juga memberi uji invertibilitas lewat matriks.

\begin{proposisi}
Misalkan $B$ dan $C$ basis dari ruang $V$ berdimensi $n$, dan $T\colon V\to V$ linear. Maka
$T$ invertible jika dan hanya jika $[T]_{C\leftarrow B}$ invertible, dan dalam hal itu
$[T^{-1}]_{B\leftarrow C}=([T]_{C\leftarrow B})^{-1}$.
\end{proposisi}

\begin{proof}
Bila $T$ invertible, maka $T^{-1}\circ T=I$ dan $T\circ T^{-1}=I$. Teorema komposisi memberi
$[T^{-1}]_{B\leftarrow C}[T]_{C\leftarrow B}=I$ dan
$[T]_{C\leftarrow B}[T^{-1}]_{B\leftarrow C}=I$, jadi kedua matriks saling invers.

Sebaliknya, misalkan $A=[T]_{C\leftarrow B}$ invertible. Bila $T(\mathbf{v})=\mathbf{0}$, maka
$A[\mathbf{v}]_{B}=\mathbf{0}$, sehingga $[\mathbf{v}]_{B}=\mathbf{0}$ dan
$\mathbf{v}=\mathbf{0}$.
Jadi $T$ injektif. Karena ruang asal dan ruang tujuan sama-sama berdimensi $n$, akibat
rank-nullity membuat $T$ surjektif, sehingga $T$ invertible. Persamaan untuk matriks
$T^{-1}$
lalu diikuti dari bagian pertama.
\end{proof}

Dengan begitu pertanyaan tentang invers operator menjadi pertanyaan tentang invers
matriks; jawaban tuntasnya ada di Bab 8.

Kasus khusus yang paling sering dipakai adalah operator pada satu ruang, $T\colon V\to V$,
dengan dua basis berbeda. Di sini kedua matriks terhubung oleh matriks transisi.

\begin{teorema}[Perubahan basis]
Misalkan $T\colon V\to V$ linear, $B$ dan $B'$ dua basis $V$, dan $P=P_{B\leftarrow B'}$.
Maka
\[
  [T]_{B'}=P^{-1}[T]_{B}P .
\]
\end{teorema}

\begin{proof}
Tulis $T=I\circ T\circ I$ dengan $I$ peta identitas. Menurut teorema komposisi,
\[
  [T]_{B'\leftarrow B'}=[I]_{B'\leftarrow B}\,[T]_{B\leftarrow B}\,[I]_{B\leftarrow B'}.
\]
Matriks identitas terhadap dua basis adalah matriks transisi, jadi
$[I]_{B'\leftarrow B}=P^{-1}$ dan $[I]_{B\leftarrow B'}=P$.
\end{proof}\index{representasi}

Dua matriks persegi yang terhubung seperti itu mendeskripsikan operator yang sama dari
sudut pandang berbeda, jadi keduanya diberi nama sendiri.

\begin{definisi}[Matriks serupa]
Matriks $A$ dan $B$ berukuran $n\times n$ disebut \emph{serupa} bila ada matriks invertible
$P$ dengan $B=P^{-1}AP$.\index{similaritas}
\end{definisi}

Similaritas karena itu nama lain untuk peta yang sama dilihat dari basis berbeda.

\begin{proposisi}
Dua matriks $n\times n$ serupa jika dan hanya jika keduanya merupakan matriks representasi
satu operator linear yang sama terhadap dua basis.
\end{proposisi}

\begin{proof}
Bila $B=P^{-1}AP$, ambil basis standar $B_{0}$ dan basis $B'$ yang kolom-kolomnya adalah
kolom-kolom $P$. Maka $P=P_{B_{0}\leftarrow B'}$, dan teorema perubahan basis memberi
$B=[T]_{B'}$ untuk operator $T(\mathbf{x})=A\mathbf{x}$. Sebaliknya, teorema yang sama
memberi
hubungan itu untuk setiap pasangan basis.
\end{proof}

Rank termasuk besaran yang tidak berubah: rank $[T]_{B'}$ sama dengan rank $T$, dan
begitu pula rank $[T]_{B}$. Determinan dan trace juga invarian, tetapi keduanya baru
didefinisikan di Bab 8 dan dibahas di Bab 9.

Contoh kecil untuk operator turunan: $D\colon P_{2}\to P_{2}$, $D(p)=p'$, basis
$B=(1,x,x^{2})$. Karena $D(1)=0$, $D(x)=1$, dan $D(x^{2})=2x$,
\[
  [D]_{B}=\begin{pmatrix}0&1&0\\ 0&0&2\\ 0&0&0\end{pmatrix}.
\]
Untuk $p=x^{2}+3x$ koordinatnya $(0,3,1)^{\top}$, dan matriks memberi $(3,\,2,\,0)^{\top}$,
yaitu $3+2x$, tepat turunan $p$.

\section{Isomorfisme, invarian dimensi, dan ruang peta linear}
\tingkat{lanjut}

Dua ruang berbeda dapat berperilaku sama. $P_{2}$ dan $\mathbb{R}^{3}$ berisi objek
berbeda jenis, tetapi keduanya punya tiga arah bebas dan aturan penjumlahan yang sama
bentuknya.

Korespondensi yang mempertahankan seluruh struktur itu disebut isomorfisme: bijektif,
linear, dan karena itu mempertahankan basis, dimensi, serta inti.

\kamus{isomorfisme}{isomorphism}{Dua denah rumah yang gambarnya berbeda tetapi jumlah dan
hubungan ruangannya sama.}

\begin{definisi}[Isomorfisme]
Peta linear $T\colon V\to W$ disebut \emph{isomorfisme} bila $T$ bijektif. Dua ruang disebut
\emph{isomorf} bila ada isomorfisme antara keduanya, ditulis $V\cong W$.
\end{definisi}

Satu lemma sebelum teorema utamanya: isomorfisme memindahkan basis menjadi basis.

\begin{lemma}
Bila $T\colon V\to W$ isomorfisme dan $B$ basis $V$, maka $T(B)$ basis $W$.
\end{lemma}

\begin{proof}
Misalkan $B=(\mathbf{b}_{1},\dots,\mathbf{b}_{n})$. Untuk kebebasan linear, andaikan
$\sum_{j}c_{j}T(\mathbf{b}_{j})=\mathbf{0}$. Karena $T$ linear, vektor
$\sum_{j}c_{j}\mathbf{b}_{j}$ dipetakan ke $\mathbf{0}$, jadi vektor itu ada di inti.
Intinya
hanya $\{\mathbf{0}\}$ karena $T$ injektif, sehingga semua $c_{j}=0$.
Untuk merentang, ambil $\mathbf{w}\in W$. Ada $\mathbf{v}\in V$ dengan $T(\mathbf{v})=\mathbf{w}$
karena $T$ surjektif, dan $\mathbf{v}=\sum_{j}a_{j}\mathbf{b}_{j}$ karena $B$ merentang.
Maka
$\mathbf{w}=\sum_{j}a_{j}T(\mathbf{b}_{j})$.
\end{proof}

Dengan lemma itu, kesamaan dimensi ternyata satu-satunya penghalang.

\begin{teorema}
Dua ruang berdimensi hingga atas $\mathbb{R}$ isomorf jika dan hanya jika dimensinya sama.
\end{teorema}

\begin{proof}
Bila $V\cong W$ lewat $T$, lemma sebelumnya mengubah basis $V$ menjadi basis $W$ dengan
banyaknya anggota sama, jadi $\dim V=\dim W$.

Sebaliknya, misalkan $\dim V=\dim W=n$. Pilih basis $B$ pada $V$ dan $C$ pada $W$. Peta
koordinat $\varphi_{B}\colon V\to\mathbb{R}^{n}$ dan $\varphi_{C}\colon W\to\mathbb{R}^{n}$
bijektif dan linear. Maka $\varphi_{C}^{-1}\circ\varphi_{B}$ adalah peta linear bijektif
dari
$V$ ke $W$, karena komposisi peta linear bijektif juga bijektif dan linear.
\end{proof}

Kesimpulan itu langsung terasa pada ruang yang kita kenal.

\begin{akibat}
Setiap ruang berdimensi $n$ isomorf dengan $\mathbb{R}^{n}$. Secara khusus
$P_{n-1}\cong\mathbb{R}^{n}$ dan $M_{2\times2}(\mathbb{R})\cong\mathbb{R}^{4}$.
\end{akibat}

Isomorfisme mempertahankan segala sesuatu yang didefinisikan dari operasi: kebebasan
linear, rentang, inti, jangkauan, rank, dan nullity. Dua ruang isomorf tidak dapat dibedakan
oleh alat bab ini; dimensi menjadi satu-satunya bilangan yang perlu dihitung untuk mengenali
ruang berdimensi hingga.

Ada jebakan tafsir. Isomorfisme tadi bergantung pada pemilihan
basis; tidak ada isomorfisme kanonik antara $V$ dan $\mathbb{R}^{n}$. Karena itu ``$V$ sama
dengan $\mathbb{R}^{n}$'' tidak benar, sedangkan ``$V$ isomorf dengan $\mathbb{R}^{n}$''
benar.
Polinomial berderajat kurang dari tiga dan tripel bilangan real adalah objek berbeda jenis,
tetapi keduanya isomorf sebab setiap polinomial ditentukan tiga koefisiennya.

Sekarang kumpulan semua peta linear antara dua ruang; himpunan itu sendiri ruang vektor.

\begin{teorema}
Himpunan $\mathcal{L}(V,W)$ dari semua peta linear $V\to W$, dengan operasi
$(S+T)(\mathbf{v})=S(\mathbf{v})+T(\mathbf{v})$ dan $(cT)(\mathbf{v})=cT(\mathbf{v})$,
adalah
ruang vektor. Bila $\dim V=n$ dan $\dim W=m$, maka $\dim\mathcal{L}(V,W)=mn$. Pemetaan
$T\mapsto[T]_{C\leftarrow B}$ adalah isomorfisme $\mathcal{L}(V,W)\to M_{m\times n}(\mathbb{R})$.
\end{teorema}

\begin{proof}
$\mathcal{L}(V,W)$ termuat dalam ruang semua fungsi $V\to W$, yang merupakan ruang vektor
menurut contoh di Bab 6. Peta nol linear. Jumlah dua peta linear linear, dan kelipatan
skalar
peta linear juga linear; jadi $\mathcal{L}(V,W)$ subruang, dan karena itu ruang vektor.

Untuk klaim isomorfisme, misalkan $\Phi(T)=[T]_{C\leftarrow B}$. Kelinearan $\Phi$: kolom
ke-$j$ dari $\Phi(S+cT)$ adalah $[S(\mathbf{b}_{j})+cT(\mathbf{b}_{j})]_{C}$, yang sama
dengan
$[S(\mathbf{b}_{j})]_{C}+c[T(\mathbf{b}_{j})]_{C}$ karena koordinat linear. Jadi
$\Phi(S+cT)=\Phi(S)+c\Phi(T)$.

Injektivitas: bila $\Phi(T)$ matriks nol, semua $[T(\mathbf{b}_{j})]_{C}=\mathbf{0}$, jadi
$T(\mathbf{b}_{j})=\mathbf{0}$ untuk semua $j$, dan $T=\mathbf{0}$ karena peta linear
ditentukan oleh nilainya pada basis.

Surjektivitas: beri matriks $A=(a_{ij})$ berukuran $m\times n$. Definisikan
$T(\mathbf{b}_{j})=\sum_{i}a_{ij}\mathbf{c}_{i}$ dan perluas secara linear; menurut
proposisi
tentang nilai pada basis, perluasan itu ada dan tunggal. Kolom ke-$j$ dari $\Phi(T)$ adalah
$(a_{ij})_{i}$, jadi $\Phi(T)=A$.

Karena $\Phi$ isomorfisme dan $\dim M_{m\times n}(\mathbb{R})=mn$, maka
$\dim\mathcal{L}(V,W)=mn$.
\end{proof}

Hitungan itu memberi akibat langsung.

\begin{akibat}
$\dim\mathcal{L}(V,V)=n^{2}$ bila $\dim V=n$. Sebagai contoh,
$\dim\mathcal{L}(P_{2},P_{2})=9$ dan $\dim\mathcal{L}(P_{2},\mathbb{R}^{2})=6$.
\end{akibat}

Hasil terakhir itu menjelaskan mengapa matriks bukan sekadar alat hitung. Ruang operator
dan ruang matriks berdimensi sama dan terhubung oleh isomorfisme yang menghormati komposisi,
sehingga setiap pernyataan tentang operator punya padanan tentang matriks. Yang tidak
diwarisi
adalah kesederhanaan bentuk: operator yang sama dapat rumit pada satu basis dan sederhana
pada
basis lain; mencari basis terbaik adalah program Bab 9.

\begin{tantangan}
Cari satu peta linear $T\colon\mathbb{R}^{3}\to\mathbb{R}^{3}$ yang memetakan persegi satuan
pada bidang $xy$ ke persegi satuan lain dengan luas tetap $1$, tetapi
$\operatorname{rank}(T)=2$. Jelaskan mengapa luas yang bertahan pada satu bidang tidak
memaksa
rank penuh.

Petunjuk 1: periksa satu bidang, bukan seluruh ruang. Setelah itu tanyakan ke mana sumbu
ketiga
pergi.

Petunjuk 2: letakkan rotasi $90^{\circ}$ pada dua koordinat pertama dan nolkan koordinat
ketiga,
sehingga $T(\mathbf{e}_{3})=\mathbf{0}$. Hitung rank dari bayangan $\mathbf{e}_{1}$ dan
$\mathbf{e}_{2}$.

Kalau mentok, kembali ke bagian Transformasi linear, inti, dan jangkauan, khususnya definisi
rank dan akibat rank-nullity.
\end{tantangan}

\begin{untukanak}
Dua kelas di sekolah bisa berbeda ukuran, tetapi punya jumlah kursi yang sama. Kalau
jumlahnya
sama, setiap murid kelas pertama dapat dipasangkan dengan satu kursi di kelas kedua, tanpa
kursi
kosong dan tanpa murid kebagian dua kursi. Untuk urusan menghitung kursi, kedua kelas
setara.

Yang tidak ikut setara adalah nama murid dan wajah mereka. Pemetaan kursi tidak mengatakan
siapa duduk di mana; ia hanya memasangkan. Orang yang berbeda bisa membuat pemetaan berbeda,
dan semuanya sah.

Di bab ini kita akan tahu kapan dua ruang yang tampak berbeda boleh dianggap sama untuk
keperluan hitungan, dan kapan anggapan itu menyesatkan. Kesetaraan itu selalu bergantung
pada
pilihan kita, bukan pada sesuatu yang sudah ada di ruangnya.

Kalau dua kelas punya jumlah kursi sama, apakah kepala sekolah boleh menukar nama kelas
begitu
saja? Jawabannya tidak, dan alasannya ada di bab ini.
\end{untukanak}

\section{Contoh hitung}
\tingkat{menengah}

Tinjau matriks
\[
  A=\begin{pmatrix}1&2&1\\ 2&5&3\\ 3&7&4\end{pmatrix},
  \qquad \text{baris ketiga}=\text{baris pertama}+\text{baris kedua}.
\]
Kita tentukan rank, nullity, dan kedua basisnya, lalu periksa operator yang sama pada basis
lain.\index{rank}\index{nullity}

Eliminasi baris:
\[
  \begin{pmatrix}1&2&1\\ 2&5&3\\ 3&7&4\end{pmatrix}
  \xrightarrow[R_{3}\leftarrow R_{3}-3R_{1}]{R_{2}\leftarrow R_{2}-2R_{1}}
  \begin{pmatrix}1&2&1\\ 0&1&1\\ 0&1&1\end{pmatrix}
  \xrightarrow{R_{3}\leftarrow R_{3}-R_{2}}
  \begin{pmatrix}1&2&1\\ 0&1&1\\ 0&0&0\end{pmatrix}.
\]
Lalu $R_{1}\leftarrow R_{1}-2R_{2}$ memberi bentuk eselon tereduksi
\[
  \begin{pmatrix}1&0&-1\\ 0&1&1\\ 0&0&0\end{pmatrix}.
\]
Ada dua baris tak nol, jadi $\operatorname{rank}A=2$ dan $\det A=0$; kolom pivot adalah kolom
pertama dan kedua.

Inti: dari $x_{1}-x_{3}=0$ dan $x_{2}+x_{3}=0$ kita peroleh $x_{1}=x_{3}$ dan
$x_{2}=-x_{3}$.
Jadi
\[
  \ker A=\operatorname{span}\{(1,-1,1)^{\top}\},\qquad \operatorname{nullity}(A)=1 .
\]
Rank-nullity terpenuhi: $3=2+1$.

Jangkauan: kolom-kolom $A$ adalah $\mathbf{c}_{1}=(1,2,3)^{\top}$,
$\mathbf{c}_{2}=(2,5,7)^{\top}$, dan $\mathbf{c}_{3}=(1,3,4)^{\top}$. Kolom pivot pertama
dan
kedua memberi basis $\{\mathbf{c}_{1},\mathbf{c}_{2}\}$. Kolom ketiga tidak bebas:
$\mathbf{c}_{2}-\mathbf{c}_{1}=\mathbf{c}_{3}$, dan relasi itu persis vektor inti
$(1,-1,1)$: $\mathbf{c}_{1}-\mathbf{c}_{2}+\mathbf{c}_{3}=\mathbf{0}$.

Sekarang pandang $T(\mathbf{x})=A\mathbf{x}$ dan basis
$B'=(\mathbf{b}'_{1},\mathbf{b}'_{2},\mathbf{b}'_{3})$ dengan
$\mathbf{b}'_{1}=(1,1,0)$, $\mathbf{b}'_{2}=(1,0,1)$, $\mathbf{b}'_{3}=(0,1,1)$. Matriks
transisinya dari basis standar $B$ adalah
\[
  P=P_{B\leftarrow B'}=\begin{pmatrix}1&1&0\\ 1&0&1\\ 0&1&1\end{pmatrix},
  \qquad
  P^{-1}=\frac{1}{2}\begin{pmatrix}1&1&-1\\ 1&-1&1\\ -1&1&1\end{pmatrix}.
\]
Bayangan vektor basis baru dihitung lebih dulu:
\[
  A\mathbf{b}'_{1}=\begin{pmatrix}3\\7\\10\end{pmatrix},\quad
  A\mathbf{b}'_{2}=\begin{pmatrix}2\\5\\7\end{pmatrix},\quad
  A\mathbf{b}'_{3}=\begin{pmatrix}3\\8\\11\end{pmatrix}.
\]
Maka
\[
  AP=\begin{pmatrix}3&2&3\\ 7&5&8\\ 10&7&11\end{pmatrix},
  \qquad
  [T]_{B'}=P^{-1}AP=\begin{pmatrix}0&0&0\\ 3&2&3\\ 7&5&8\end{pmatrix}.
\]
Periksa satu entri: baris pertama kolom pertama $\tfrac{1}{2}(3+7-10)=0$; baris kedua
kolom ketiga $\tfrac{1}{2}(3-8+11)=3$. Sisanya mengikuti pola yang sama.

Baris pertama yang seluruhnya nol bukan kebetulan: ketiga bayangan berada di dalam
$\operatorname{span}(\mathbf{b}'_{2},\mathbf{b}'_{3})$, misalnya
$A\mathbf{b}'_{3}=(3,8,11)=3\mathbf{b}'_{2}+8\mathbf{b}'_{3}$. Karena itu koordinat pertama
setiap bayangan nol, dan jangkauan $T$ termuat dalam bidang koordinat pertama.

Pola itu umum. Baris nol menandakan jangkauan yang termuat dalam subruang koordinat
berdimensi lebih rendah; kolom nol menandakan vektor basis yang masuk ke inti. Kedua gejala
hanya tampak pada basis yang dipilih, sedangkan rank dan nullity tidak berubah.

Pemeriksaan menyeluruh:
\[
  \operatorname{rank}[T]_{B'}=2,\qquad \operatorname{trace}[T]_{B'}=0+2+8=10=\operatorname{trace}A,
  \qquad \det[T]_{B'}=0=\det A .
\]
Kesamaan trace bukan kebetulan aritmetika. Entri diagonal ke-$j$ adalah koefisien
$\mathbf{b}'_{j}$ pada bayangan $T(\mathbf{b}'_{j})$, yaitu seberapa besar $T$
mempertahankan
arah itu sendiri. Jumlahnya tidak bergantung pada arah yang dipilih, hanya pada operatornya,
sehingga trace dipakai sebagai sidik jari murah: trace berbeda berarti tidak serupa, trace
sama
belum menjamin apa pun.

Untuk $\mathbf{v}=(3,8,11)$ kita punya $[\mathbf{v}]_{B'}=(0,3,8)^{\top}$ dan
\[
  [T]_{B'}[\mathbf{v}]_{B'}=\begin{pmatrix}0\\30\\79\end{pmatrix},
\]
sedangkan $T(\mathbf{v})=A\mathbf{v}=(30,79,109)^{\top}$ dan koordinatnya
$(0,30,79)^{\top}$. Keduanya cocok.

Peta fitur $G\colon\mathbb{R}^{4}\to\mathbb{R}^{3}$
diberi matriks
\[
  G=\begin{pmatrix}1&2&0&1\\ 0&1&1&1\\ 1&3&1&2\end{pmatrix},
\]
dengan baris ketiga sama dengan jumlah dua baris pertama. Eliminasi memberi bentuk eselon
tereduksi
\[
  \begin{pmatrix}1&0&-2&-1\\ 0&1&1&1\\ 0&0&0&0\end{pmatrix},
\]
jadi $\operatorname{rank}G=2$ dan $\operatorname{nullity}(G)=4-2=2$. Dari
$x_{1}=2x_{2}+x_{4}$ dan $x_{3}=-x_{2}-x_{4}$:
\[
  \mathbf{k}_{1}=(-2,1,-1,0)^{\top},\qquad \mathbf{k}_{2}=(-1,0,-1,1)^{\top}.
\]
Keduanya memenuhi $G\mathbf{k}_{i}=\mathbf{0}$ dan bebas linear, jadi perubahan yang tidak
terlihat oleh ketiga fitur berbentuk bidang dua dimensi, bukan satu garis: penyerang dapat
menggeser masukan ke arah mana pun di bidang itu tanpa menggerakkan skor.

\kertas{Ulangi perhitungan dengan basis $B''=(\mathbf{b}'_{2},\mathbf{b}'_{3},\mathbf{b}'_{1})$,
yaitu $B'$ yang dipermutasi. Ramalkan bagaimana baris nol dan entri lain berpindah, lalu
hitung
$[T]_{B''}$ lewat $Q^{-1}[T]_{B'}Q$ dengan $Q$ matriks permutasi yang sesuai. Periksa bahwa
rank
dan trace tidak berubah, dan jelaskan mengapa $\det A=0$ dapat disimpulkan langsung dari ``baris
ketiga adalah jumlah dua baris pertama''.}

Semua angka di atas diperiksa dengan aritmetika eksak (SymPy); pemeriksaan numpy dan
scipy ada di bagian berikut, termasuk satu angka yang tidak persis nol karena pembulatan
mesin.

\section{Lab komputasi}
\tingkat{lanjut}

\begin{lab}
\begin{verbatim}
import numpy as np
from scipy.linalg import null_space

A = np.array([[1., 2., 1.],
              [2., 5., 3.],
              [3., 7., 4.]])
print("rank A  :", np.linalg.matrix_rank(A))
print("det A   :", np.linalg.det(A))
K = null_space(A)
print("dim ker :", K.shape[1])
print("basis ker:\n", np.round(K, 4))
print("A @ K   :\n", np.round(A @ K, 12))

P = np.array([[1., 1., 0.],
              [1., 0., 1.],
              [0., 1., 1.]])
B = np.linalg.inv(P) @ A @ P
print("[T]_B'  :\n", np.round(B, 10))
print("selisih :", np.linalg.norm(B - np.linalg.inv(P) @ A @ P))
print("trace   :", round(np.trace(B), 10), round(np.trace(A), 10))
print("rank B  :", np.linalg.matrix_rank(B))
print("cond P  :", np.linalg.cond(P))

v = np.array([3., 8., 11.])
print("[v]_B'  :", np.round(np.linalg.solve(P, v), 4))
print("[Av]_B' :", np.round(np.linalg.solve(P, A @ v), 4))
print("B[v]_B' :", np.round(B @ np.linalg.solve(P, v), 4))

x0 = np.array([1., 0., 0.])
d = 0.5 * K[:, 0]
print("A@d     :", np.round(A @ d, 12))
print("geser   :", np.linalg.norm(A @ (x0 + d) - A @ x0))
\end{verbatim}
Keluaran sebenarnya:
\begin{verbatim}
rank A  : 2
det A   : 6.66133814775094e-16
dim ker : 1
basis ker:
 [[ 0.5774]
 [-0.5774]
 [ 0.5774]]
A @ K   :
 [[-0.]
 [ 0.]
 [ 0.]]
[T]_B'  :
 [[0. 0. 0.]
 [3. 2. 3.]
 [7. 5. 8.]]
selisih : 0.0
trace   : 10.0 10.0
rank B  : 2
cond P  : 2.0000000000000004
[v]_B'  : [0. 3. 8.]
[Av]_B' : [ 0. 30. 79.]
B[v]_B' : [ 0. 30. 79.]
A@d     : [-0.  0.  0.]
geser   : 6.280369834735101e-16
\end{verbatim}
\end{lab}

Empat hal perlu dibaca dari keluaran itu.

Pertama, $\det A$ bukan nol melainkan $6{,}66\times10^{-16}$. Matriks ini singular
secara eksak, tetapi aritmetika pecahan biner tidak mengenal nol eksak.
\texttt{matrix\_rank} membandingkan nilai singular terbesar dan terkecil, bukan ambang
sederhana, sehingga rank tetap dilaporkan $2$; Bab 22 membahas mengapa ambang itu perlu.

Pilihan nilai singular, bukan determinan: determinan mengukur volume, dan volume dapat
sangat
kecil tanpa rank berubah. Yang menentukan rank adalah perbandingan nilai singular terkecil
dan
terbesar.

Kedua, \texttt{null\_space} mengembalikan vektor satuan dengan entri $0{,}5774$,
$-0{,}5774$, dan $0{,}5774$, yaitu $(1,-1,1)^{\top}$ dibagi $\sqrt{3}$, dan
$1/\sqrt{3}=0{,}57735$ sesuai. Arahnya sama dengan hasil tangan; panjangnya dinormalkan
karena
scipy memilih basis ortonormal. Basis kernel tidak tunggal, dimensinya tunggal.

Ketiga, $[T]_{B'}$ yang dilaporkan numpy sama persis dengan hasil tangan, dan trace
kedua matriks sama, $10$. Selisih $P^{-1}AP$ yang dihitung dua kali adalah $0{,}0$ karena
$P$
berisi bilangan bulat kecil. Bilangan kondisi $P$ bernilai $2{,}0000000000000004$, cukup
kecil;
basis buruk memperbesar galat, dan itu dibahas di Bab 10.

Keempat, arah inti benar-benar buta bagi $T$. Mendorong $\mathbf{x}_{0}=(1,0,0)$ sejauh
setengah vektor inti mengubah $\mathbf{x}_{0}$ secara nyata, tetapi $A\mathbf{x}$ hanya
bergeser
$6{,}28\times10^{-16}$, yaitu nol dalam presisi mesin.

Peringatan terakhir menyangkut pilihan basis. Basis yang nyaris bergantung membuat koordinat
meledak, dan angka yang dilaporkan kehilangan makna.

\begin{lab}
\begin{verbatim}
Pb = np.array([[1., 1., 0.],
               [1., 1.001, 0.],
               [0., 0., 1.]])
print("cond Pb:", np.linalg.cond(Pb))
print("koord  :", np.linalg.inv(Pb) @ v)
print("galat  :", np.linalg.norm(np.linalg.inv(Pb) @ (Pb @ v) - v))
\end{verbatim}
Keluaran sebenarnya:
\begin{verbatim}
cond Pb: 4002.00075012484
koord  : [-4997.  5000.    11.]
galat  : 4.823324033076557e-13
\end{verbatim}
\end{lab}

Dua vektor basis pertama berbeda hanya pada digit ketiga, dan bilangan kondisi $P_{b}$
mencapai $4002$. Koordinat $(3,8,11)$ meledak menjadi $(-4997,\,5000,\,11)$. Galat
rekonstruksi
masih $4{,}8\times10^{-13}$, tetapi galat relatif pada koordinat membesar sebesar faktor
bilangan
kondisi; dengan 16 digit presisi, sisanya tinggal sekitar 12. Pada basis yang lebih buruk
sisa
itu habis dan rank yang dilaporkan tidak dapat dipercaya.

\begin{untukdosen}
Urutan mengajar yang berhasil dipakai: koordinat lebih dulu, baru peta linear. Mahasiswa
yang
mengenal koordinat sebelum definisi abstrak lebih siap menerima matriks representasi, karena
matriks muncul sebagai daftar koordinat, bukan benda baru. Setelah rank-nullity, masuk
perubahan basis.

Dua miskonsepsi khas. Pertama, menyamakan vektor dengan daftar koordinatnya, lalu bingung
ketika dua basis memberi dua daftar untuk vektor yang sama. Kedua, arah matriks transisi
terbalik. Cara cepat memeriksa: kolom harus berisi alamat vektor basis, dan subskrip dibaca
kanan ke kiri.

Aktivitas papan tulis: beri satu operator pada $\mathbb{R}^{2}$ dan dua basis. Minta separuh
kelas menghitung $[T]_{B'}$ lewat definisi kolom, separuh lewat $P^{-1}AP$, lalu bandingkan
di
papan; kegaduhan yang muncul biasanya tepat pada arah $P$.

Tanda kesiapan lanjut: mahasiswa dapat menjelaskan mengapa sebuah baris nol muncul pada
matriks yang baru, tanpa menghitung ulang. Bila mereka menghubungkannya dengan jangkauan
yang
termuat di subruang koordinat, mereka siap masuk Bab 9.
\end{untukdosen}

\section{Ringkasan}
\tingkat{dasar}

\begin{ringkasan}
\begin{itemize}
  \item Koordinat $[\mathbf{v}]_{B}$ adalah daftar koefisien tunggal terhadap basis terurut;
        peta $\mathbf{v}\mapsto[\mathbf{v}]_{B}$ linear dan bijektif, tetapi bergantung pada
        basis.
  \item Matriks transisi memuat koordinat basis lama dalam basis baru:
        $[\mathbf{v}]_{C}=P_{C\leftarrow B}[\mathbf{v}]_{B}$, dengan
        $P_{B\leftarrow C}=P_{C\leftarrow B}^{-1}$.
  \item Peta linear ditentukan nilainya pada basis; $T$ injektif tepat ketika
        $\ker T=\{\mathbf{0}\}$.
  \item Rank-nullity: $\dim V=\dim\ker T+\dim\operatorname{im}T$; pada dimensi yang sama,
        injektif setara dengan surjektif.
  \item Matriks representasi memenuhi
$[T(\mathbf{v})]_{C}=[T]_{C\leftarrow B}[\mathbf{v}]_{B}$, komposisi menjadi perkalian
        matriks, dan $[T]_{B'}=P^{-1}[T]_{B}P$.
  \item $\dim\mathcal{L}(V,W)=mn$ bila $\dim V=n$, $\dim W=m$; dua ruang berdimensi hingga
        isomorf tepat ketika dimensinya sama.
\end{itemize}
\end{ringkasan}

\section{Latihan}
\tingkat{menengah}

\begin{latihan}
\begin{enumerate}
  \item Dengan $B'=(\mathbf{b}'_{1},\mathbf{b}'_{2},\mathbf{b}'_{3})$ seperti pada contoh
hitung, tentukan $[\mathbf{w}]_{B'}$ untuk $\mathbf{w}=(2,-1,3)$, lalu periksa hasilnya
        dengan mengalikan balik.

\item Buktikan bahwa setiap isomorfisme $V\to\mathbb{R}^{n}$ berbentuk $\varphi_{B}$ untuk
        suatu basis $B$, dan jelaskan mengapa isomorfisme ke $\mathbb{R}^{n}$ tidak tunggal.
        (Petunjuk: ambil prapeta basis standar, lalu tunjukkan himpunan itu basis.)

  \item Untuk $A=\begin{pmatrix}0&1&0\\ 0&0&2\\ 0&0&0\end{pmatrix}$, tentukan rank, nullity,
        basis inti, dan basis jangkauan. Jelaskan hubungannya dengan operator turunan pada
        $P_{2}$.

\item Buktikan $\operatorname{rank}(S\circ T)\le\min(\operatorname{rank}S,\operatorname{rank}T)$
        untuk peta linear yang terdefinisi. Gunakan
        $\operatorname{im}(S\circ T)\subseteq\operatorname{im}S$ dan bandingkan inti.
        (Petunjuk: untuk batas kedua, bandingkan $\ker T$ dengan $\ker(S\circ T)$ lewat
        rank-nullity.)

\item Nilai klaim ini: \emph{dua matriks persegi berukuran sama yang mempunyai rank sama
pasti serupa.} Tentukan benar, salah, atau belum cukup bukti, dengan alasan dan contoh
        penyangkal bila perlu.

\item Misalkan $T\colon\mathbb{R}^{3}\to\mathbb{R}^{2}$ dengan
        $T(x,y,z)=(x+2y,\;y+z)$. Hitung $[T]_{C\leftarrow B}$ untuk
$B=(\,(1,0,0),(1,1,0),(1,1,1)\,)$ dan $C=(\,(1,1),(1,0)\,)$, lalu periksa rumusnya pada
        $\mathbf{v}=(1,2,3)$.
        (Petunjuk: $[T]_{C\leftarrow B}=C^{-1}AB$; hitung $AB$ lebih dulu, lalu selesaikan
        $C\mathbf{x}=\mathbf{y}$ untuk setiap kolomnya.)

  \item Buktikan: bila $U$ subruang $V$ berdimensi hingga dan $\dim U=\dim V$, maka $U=V$.

  \item Buktikan bahwa bila $A$ dan $B$ serupa maka $\operatorname{nullity}(A)
        =\operatorname{nullity}(B)$. Jelaskan mengapa pernyataan itu setara dengan
        ``inti tidak bergantung pada basis''.

  \item Peta fitur $F\colon\mathbb{R}^{4}\to\mathbb{R}^{3}$ diberikan oleh matriks
        \[
          F=\begin{pmatrix}1&0&2&1\\ 0&1&1&3\\ 1&1&3&4\end{pmatrix}.
        \]
        Tentukan rank dan nullity $F$, beri satu vektor inti tak nol, dan jelaskan mengapa
        vektor itu merupakan arah perubahan yang tidak menggeser satu pun fitur.
\end{enumerate}
\end{latihan}

\begin{untukpeneliti}
Yang terbukti di bab ini adalah pernyataan berdimensi hingga atas medan $\mathbb{R}$:
rank-nullity, ketunggalan matriks representasi, kesetaraan isomorfisme dengan kesamaan
dimensi,
dan $\dim\mathcal{L}(V,W)=mn$. Semuanya teorema, bukan konvensi. Klaim ``matriks adalah
transformasinya'' hanya populer; tanpa pemilihan basis klaim itu tidak bermakna, dan contoh
baris
nol di atas menunjukkan bentuk matriks bergantung pada basis.

Sumber primer: Axler, \emph{Linear Algebra Done Right} edisi ke-4, bagian 3A sampai 3D,
memuat
ruang peta linear, inti dan jangkauan, matriks, isomorfisme, serta perubahan basis;
rank-nullity
di sana disebut teorema fundamental untuk peta linear. Hoffman dan Kunze, \emph{Linear
Algebra}
edisi ke-2, Bab 3, lebih rapat namun setara.

Pertanyaan terbuka yang jujur. Klasifikasi operator hingga similaritas tidak diberikan di
sini; Bab 9 hanya sampai diagonalisasi. Secara numerik, keputusan rank bergantung pada
ambang dan
bilangan kondisi; bab ini memakai hasil pustaka tanpa membahas stabilitasnya, utang yang
dibayar
di Bab 22. Batas keberlakuan: di ruang berdimensi tak hingga seluruh argumen runtuh,
rank-nullity
gagal, dan isomorfisme tidak lagi setara dengan kesamaan dimensi.
\end{untukpeneliti}

\peta{Bab ini mengikuti Anton, \emph{Elementary Linear Algebra} edisi ke-11, bagian 4.4, 4.6,
4.8, dan 8.1 sampai 8.5, dengan padanan abstrak di Axler, \emph{Linear Algebra Done Right}
edisi
ke-4, Bab 3, dan Hoffman dan Kunze, \emph{Linear Algebra} edisi ke-2, Bab 3.}
\chapter{Sistem Persamaan Linear, Determinan, dan Invers}

Sebuah penilai ancaman menjumlahkan tiga fitur dengan bobot yang belum diketahui,
$s=w_{1}f_{1}+w_{2}f_{2}+w_{3}f_{3}$. Kalibrasi pada tiga contoh berlabel memberi tiga
persamaan linear sekaligus:
\[
  2w_{1}+w_{2}-w_{3}=8,\qquad
  -3w_{1}-w_{2}+2w_{3}=-11,\qquad
  -2w_{1}+w_{2}+2w_{3}=-3 .
\]
Jawabannya $w=(2,3,-1)$, dan tidak ada yang mengejutkan.

Kemudian matriks koefisiennya diganti oleh hasil pengukuran lain,
\[
  S=\begin{pmatrix}1&2&3\\ 4&5&6\\ 7&8&9{,}001\end{pmatrix}.
\]
Baris kedua dikurangi baris pertama memberi $(3,3,3)$, sedangkan baris ketiga dikurangi
baris kedua memberi $(3,3,3{,}001)$: ketiga baris hampir bergantung. Bobot $w=(1,1,1)$
menghasilkan ruas kanan $(6,\,15,\,24{,}001)$, jadi sistem itu konsisten dan jawabannya tepat
$w=(1,1,1)$. Naikkan pengukuran ketiga sebesar $0{,}001$, dan jawabannya berpindah ke
$(2,-1,2)$. Perubahan relatif pada data $3{,}46\times10^{-5}$ menghasilkan perubahan relatif
pada jawaban $1{,}41$.

Ada empat pertanyaan, bukan satu: apakah sistem punya jawaban, berapa banyak, angka apa yang
memutuskannya tanpa menyelesaikan apa pun, dan seberapa jauh jawaban bergeser ketika data
bergeser. Bab 7 sudah memberi rank dan nullity untuk yang pertama. Yang belum ada tiga:
prosedur yang tercatat dan dapat diperiksa orang lain, satu bilangan yang memutuskan
invertibilitas sekaligus memberi rumus jawaban, dan ukuran kepekaan; ketiganya isi bab ini.

Determinan lazim diperkenalkan sebagai rumus kofaktor, lalu sifatnya dibuktikan satu per
satu. Kita balik: definisi dulu lewat tiga sifat, ketunggalan dibuktikan, rumus kofaktor
muncul sebagai akibat. Alasannya bukan gaya: rumus kofaktor menyembunyikan bahwa determinan
adalah satu-satunya fungsi bergantian yang ternormalisasi, dan fakta itulah yang membuat
bukti di bab ini pendek.

\begin{untukanak}
Tiga kawan menimbang tiga kotak dengan tiga timbangan berbeda. Setiap kawan menuliskan tiga
angka, dan dari angka-angka itu kita ingin tahu berat setiap kotak. Selama ketiga timbangan
mengukur hal yang berbeda, jawabannya mudah: setiap angka baru memberi satu keterangan baru.

Masalahnya begini. Timbangan kedua ternyata hampir menyalin timbangan pertama. Angkanya
tidak sama persis, hanya berbeda sedikit, mungkin karena debu di piringan. Timbangan ketiga
juga hampir menyalin keduanya. Sekarang ketiga keterangan itu hampir sama isinya. Kita
merasa punya tiga keterangan, padahal yang benar-benar baru hanya satu atau dua.

Di sinilah keanehan mulai. Kalau salah satu angka berubah seperseribu, berat kotak yang
kita hitung bisa berubah hampir setengah. Bukan karena hitungannya salah, tetapi karena
keterangannya memang hampir kembar. Ada satu bilangan yang dapat memberi tahu lebih dulu
apakah hal itu akan terjadi, sebelum kita menghitung apa pun.

Pernahkah kamu menanyakan hal yang sama kepada tiga orang, lalu mendapat tiga jawaban yang
isinya mirip, dan tetap tidak yakin? Di bab ini kita akan tahu mengapa kemiripan jawaban
justru melemahkan keterangan, dan bagaimana mengenalinya dari angkanya.
\end{untukanak}

\section{Eliminasi Gauss dan tiga nasib sebuah sistem}
\tingkat{dasar}

Sistem linear ditulis sebagai satu kesatuan, bukan sebagai daftar persamaan yang terpisah.
Caranya dengan menempelkan kolom ruas kanan ke matriks koefisien. Penyelesaian dikerjakan
pada susunan itu, dan setiap langkahnya adalah penggantian satu sistem oleh sistem lain yang
himpunan jawabannya sama.

\kamus{eliminasi}{elimination}{Mengganti pertanyaan dengan pertanyaan yang lebih sederhana
tanpa mengubah jawabannya, seperti menyederhanakan pecahan tanpa mengubah nilainya.}

\begin{definisi}[Matriks diperbesar dan operasi baris elementer]
Untuk sistem $A\mathbf{x}=\mathbf{b}$ dengan $A$ berukuran $m\times n$, \emph{matriks
diperbesar} adalah matriks $\bigl(A\mid\mathbf{b}\bigr)$ berukuran $m\times(n+1)$. Tiga
\emph{operasi baris elementer} adalah: menukar dua baris; mengalikan satu baris dengan
skalar tak nol; menambahkan kelipatan satu baris pada baris lain. Dua matriks disebut
\emph{ekuivalen baris} bila yang satu dapat dicapai dari yang lain dengan operasi-operasi
itu.\index{eliminasi}
\end{definisi}

Operasi itu dipilih karena masing-masing dapat dibatalkan. Menukar dua baris membatalkan
dirinya sendiri, mengalikan dengan $c\neq0$ dibatalkan dengan mengalikan $1/c$, dan
menambahkan $c$ kali baris $i$ pada baris $j$ dibatalkan dengan mengurangi $c$ kali baris
$i$. Karena itu tidak ada informasi yang hilang.

\begin{teorema}
Bila dua sistem linear mempunyai matriks diperbesar yang ekuivalen baris, kedua sistem itu
mempunyai himpunan jawaban yang sama.
\end{teorema}

\begin{proof}
Cukup diperiksa untuk satu operasi, karena operasi berulang mengikuti. Misalkan $S$ adalah
sistem semula dan $S'$ sistem setelah satu operasi. Setiap persamaan $S'$ adalah kombinasi
linear persamaan-persamaan $S$, sehingga setiap jawaban $S$ juga menjawab $S'$. Sebaliknya,
karena operasi itu dapat dibatalkan oleh operasi sejenis, setiap persamaan $S$ adalah
kombinasi linear persamaan $S'$, sehingga setiap jawaban $S'$ juga menjawab $S$. Kedua
himpunan jawaban karena itu saling memuat.
\end{proof}

Bentuk yang dituju adalah bentuk tangga. Setiap baris tak nol dimulai lebih ke kanan
daripada baris tak nol di atasnya, dan semua entri di bawah posisi pertama itu nol. Kolom
tempat angka pertama sebuah baris tak nol berada disebut kolom pivot, dan angka itu sendiri
disebut pivot.

\begin{definisi}[Bentuk eselon baris]
Matriks disebut dalam \emph{bentuk eselon baris} bila semua baris nol berada di bawah, dan
entri tak nol pertama setiap baris tak nol terletak lebih kanan daripada entri tak nol
pertama baris di atasnya. Bentuk itu disebut \emph{tereduksi} bila setiap pivot bernilai
$1$ dan merupakan satu-satunya entri tak nol pada kolomnya.\index{eselon}\index{pivot}
\end{definisi}

Bentuk tereduksi itu tunggal. Fakta ini dipakai setiap kali kita berbicara tentang ``rank
yang dihitung dengan eliminasi'', dan buktinya tidak sependek kelihatannya.

\begin{teorema}[Ketunggalan bentuk eselon tereduksi]
Setiap matriks ekuivalen baris dengan tepat satu matriks dalam bentuk eselon baris
tereduksi.
\end{teorema}

\begin{proof}
Bukti dilewati di sini. Pernyataan ini dibuktikan lengkap di Hoffman dan Kunze,
\emph{Linear Algebra} edisi ke-2, bagian 1.4, dengan menunjukkan bahwa ruang jawaban
$A\mathbf{x}=\mathbf{0}$ menentukan letak pivot secara tunggal. Yang kita pakai di bab ini
hanya akibatnya: banyaknya pivot dan letaknya tidak bergantung pada urutan operasi yang
dipilih.
\end{proof}

Dengan bentuk tereduksi, seluruh pertanyaan tentang jumlah jawaban dapat dijawab dengan
menghitung pivot. Inilah tiga nasib sebuah sistem.

\begin{teorema}[Tiga kasus]
Misalkan $\bigl(A\mid\mathbf{b}\bigr)$ berekuivalen baris dengan bentuk eselon tereduksi
$\bigl(R\mid\mathbf{c}\bigr)$, dengan $A$ berukuran $m\times n$.
\begin{enumerate}
  \item Sistem \emph{tidak konsisten} tepat ketika ada pivot pada kolom $(n+1)$.
  \item Bila konsisten, sistem punya \emph{jawaban tunggal} tepat ketika setiap kolom $A$
        memuat pivot.
  \item Bila konsisten, sistem punya \emph{banyak jawaban} tepat ketika ada kolom $A$
        tanpa pivot, dan banyaknya parameter bebas sama dengan banyaknya kolom tanpa pivot.
\end{enumerate}
\end{teorema}

\begin{proof}
Pivot pada kolom terakhir berarti ada baris berbentuk $0=\text{konstanta tak nol}$, yang
mustahil dipenuhi. Bila tidak ada pivot semacam itu, setiap baris tak nol memberi satu
persamaan yang menyatakan variabel pivot dalam variabel lain, sehingga sistem konsisten.

Bila setiap kolom memuat pivot, banyaknya pivot sama dengan $n$, sehingga tidak ada variabel
bebas dan setiap variabel pivot ditentukan tunggal. Bila ada kolom tanpa pivot, variabel
pada kolom itu bebas: ia boleh bernilai sebarang, dan setiap pilihan menentukan variabel
pivot secara tunggal. Jadi jawabannya tak berhingga banyak, dan banyaknya parameter bebas
tepat banyaknya kolom tanpa pivot.
\end{proof}

Gambar~\ref{fig:tiga-nasib} menaruh ketiga nasib itu pada satu bidang. Garis $x+y=2$
tidak berubah; yang berubah hanya garis kedua. Panel kiri: $x-y=0$ memotongnya di $(1,1)$,
dua pivot, jawaban tunggal. Panel tengah: $x+y=0$ sejajar, kolom terakhir matriks diperbesar
memuat pivot. Panel kanan: $2x+2y=4$ adalah garis yang sama dengan angka berbeda, satu pivot
dan satu parameter bebas.

\begin{figure}[htbp]
\centering
\begin{tikzpicture}[scale=0.66,>=Stealth]
  \begin{scope}
    \draw[layergrey!28,step=1] (-2.2,-2.2) grid (2.2,2.2);
    \draw[->,ink] (-2.3,0)--(2.4,0) node[right,font=\scriptsize]{$x$};
    \draw[->,ink] (0,-2.3)--(0,2.4) node[above,font=\scriptsize]{$y$};
    \draw[very thick,layerindigo] (-1.2,3.2)--(3.2,-1.2);
    \draw[very thick,layerteal] (-1.6,-1.6)--(2.2,2.2);
    \fill[accent] (1,1) circle (2.4pt);
    \node[font=\scriptsize,below right] at (1,1) {$(1,1)$};
    \node[font=\scriptsize] at (0,-2.9) {(a) berpotongan};
  \end{scope}
  \begin{scope}[xshift=6.4cm]
    \draw[layergrey!28,step=1] (-2.2,-2.2) grid (2.2,2.2);
    \draw[->,ink] (-2.3,0)--(2.4,0) node[right,font=\scriptsize]{$x$};
    \draw[->,ink] (0,-2.3)--(0,2.4) node[above,font=\scriptsize]{$y$};
    \draw[very thick,layerindigo] (-1.2,3.2)--(3.2,-1.2);
    \draw[very thick,accent,dashed] (-2.2,2.2)--(2.2,-2.2);
    \node[font=\scriptsize] at (0,-2.9) {(b) sejajar};
  \end{scope}
  \begin{scope}[xshift=12.8cm]
    \draw[layergrey!28,step=1] (-2.2,-2.2) grid (2.2,2.2);
    \draw[->,ink] (-2.3,0)--(2.4,0) node[right,font=\scriptsize]{$x$};
    \draw[->,ink] (0,-2.3)--(0,2.4) node[above,font=\scriptsize]{$y$};
    \draw[very thick,layerindigo] (-1.2,3.2)--(3.2,-1.2);
    \draw[line width=2.6pt,accent,densely dotted] (-1.2,3.2)--(3.2,-1.2);
    \node[font=\scriptsize] at (0,-2.9) {(c) berimpit};
  \end{scope}
\end{tikzpicture}
\caption{Tiga nasib sistem $2\times2$ dengan garis pertama $x+y=2$ tetap. Kiri: garis kedua
$x-y=0$ memotong di $(1,1)$, dua pivot, jawaban tunggal. Tengah: garis kedua $x+y=0$
sejajar, kolom ruas kanan memuat pivot, tidak ada jawaban. Kanan: garis kedua $2x+2y=4$
berimpit dengan garis pertama, satu pivot, tak berhingga jawaban.}
\label{fig:tiga-nasib}
\end{figure}

\begin{bacakeras}
Dua suara. Suara pertama menyebut keadaan, suara kedua menyebut akibatnya. Bacakan
bergantian, lalu bertukar peran.

Satu sistem. Dua garis. Tiga nasib.

Berpotongan: satu titik, satu jawaban.

Sejajar: tidak bertemu, tidak ada jawaban.

Berimpit: satu garis menyamar menjadi dua, jawaban tak berhingga.

Siapa yang memutuskan, sebelum kita menghitung apa pun? Pivot yang memutuskan.

Dua pivot, dua kolom: jawaban tunggal.

Pivot di kolom terakhir: ada baris yang berbunyi nol sama dengan bukan nol.

Ada kolom tanpa pivot: satu parameter bebas, jawaban berderet ke satu arah.

Ulangi baris terakhir bersama-sama: hitung pivot dulu, baru bicara jawaban.
\end{bacakeras}

Matriks dari Bab 7 memperlihatkan dua nasib sekaligus,
\[
  A=\begin{pmatrix}1&2&1\\ 2&5&3\\ 3&7&4\end{pmatrix},
  \qquad \text{baris ketiga}=\text{baris pertama}+\text{baris kedua}.
\]
Untuk $\mathbf{b}=(4,9,13)$ eliminasi memberi
\[
  \left(\begin{array}{ccc|c}1&2&1&4\\ 2&5&3&9\\ 3&7&4&13\end{array}\right)
  \longrightarrow
  \left(\begin{array}{ccc|c}1&0&-1&2\\ 0&1&1&1\\ 0&0&0&0\end{array}\right).
\]
Ada dua pivot dan tiga kolom, jadi ada satu variabel bebas. Jawabannya
$(2,1,0)+t\,(1,-1,1)$ untuk $t\in\mathbb{R}$, dan pemeriksaannya langsung:
$A(2,1,0)^{\top}=(4,9,13)^{\top}$ sedangkan $A(1,-1,1)^{\top}=\mathbf{0}$, jadi arah
pergeseran itu memang arah buta matriks dari Bab 7.

Untuk $\mathbf{b}=(4,9,14)$ baris terakhir menjadi $0=1$, dan sistem tidak konsisten.
Matriks koefisiennya sama persis; satu angka di ruas kanan memindahkan nasib sistem dari
``banyak jawaban'' ke ``tidak ada jawaban''. Karena itu pertanyaan tentang konsistensi tidak
dapat dijawab dari $A$ saja.

Kasus yang paling sering muncul dalam praktik adalah $m=n$. Di sana ketiga nasib itu
bertumpu pada satu bilangan, dan bilangan itulah yang kita cari di bagian berikut.

\begin{akibat}
Misalkan $A$ berukuran $n\times n$. Sistem $A\mathbf{x}=\mathbf{b}$ punya jawaban tunggal
untuk setiap $\mathbf{b}\in\mathbb{R}^{n}$ tepat ketika $\operatorname{rank}A=n$.
\end{akibat}

\begin{proof}
Bila $\operatorname{rank}A=n$, bentuk eselon tereduksi $A$ adalah $I_{n}$, sehingga
eliminasi pada $\bigl(A\mid\mathbf{b}\bigr)$ selalu berakhir dengan pivot pada setiap kolom
$A$ dan tanpa pivot pada kolom terakhir; teorema sebelumnya memberi jawaban tunggal.
Sebaliknya, bila $\operatorname{rank}A<n$, ada kolom tanpa pivot, dan untuk $\mathbf{b}$ yang
konsisten jawabannya tidak tunggal. Karena rank kurang dari $n$, ruang kolom $A$ tidak
memenuhi $\mathbb{R}^{n}$, sehingga ada $\mathbf{b}$ di luarnya; untuk $\mathbf{b}$ itu
sistemnya tidak konsisten.
\end{proof}

\section{Determinan sebagai fungsi bergantian}
\tingkat{menengah}

Kita ingin satu bilangan yang menjawab pertanyaan akibat sebelumnya tanpa menyelesaikan
sistem apa pun. Bilangan itu harus nol tepat ketika kolom-kolom $A$ tidak bebas linear,
berubah tanda ketika dua kolom bertukar, dan bernilai satu untuk matriks identitas. Ketiga
syarat itu ternyata sudah menentukan bilangannya secara tunggal.

\kamus{determinan}{determinant}{Faktor pengali volume: berapa kali ruang diregangkan atau
digepengkan oleh matriks itu.}

\begin{definisi}[Determinan]
Fungsi $D$ dari matriks $n\times n$ ke $\mathbb{R}$ disebut \emph{determinan} bila
\begin{enumerate}
  \item \emph{multilinear pada setiap baris}: untuk skalar $\alpha,\beta$ dan baris
        $\mathbf{u},\mathbf{v}$,
        $D(\dots,\alpha\mathbf{u}+\beta\mathbf{v},\dots)
         =\alpha D(\dots,\mathbf{u},\dots)+\beta D(\dots,\mathbf{v},\dots)$;
  \item \emph{bergantian}: $D(\dots,\mathbf{u},\dots,\mathbf{u},\dots)=0$ bila dua baris
        sama;
  \item \emph{ternormalisasi}: $D(I_{n})=1$.
\end{enumerate}
Nilai $D(A)$ ditulis $\det A$.\index{determinan}
\end{definisi}

Multilinearitas berarti determinan linear pada satu baris bila baris lain ditahan. Sifat
bergantian adalah syarat terkuat, dan akibatnya panjang.

\begin{lemma}
Menukar dua baris mengubah tanda determinan. Bila satu baris merupakan kelipatan skalar baris
lain, atau bila ada baris nol, determinannya nol.
\end{lemma}

\begin{proof}
Misalkan dua baris bersebelahan keduanya berisi $\mathbf{u}+\mathbf{v}$. Karena dua baris
sama, nilainya nol. Multilinearitas menguraikan nol itu menjadi empat suku,
\[
  0=D(\dots,\mathbf{u},\mathbf{u},\dots)+D(\dots,\mathbf{u},\mathbf{v},\dots)
   +D(\dots,\mathbf{v},\mathbf{u},\dots)+D(\dots,\mathbf{v},\mathbf{v},\dots),
\]
dan suku pertama serta keempat lenyap. Jadi menukar dua baris bersebelahan membalik tanda:
\[
  D(\dots,\mathbf{v},\mathbf{u},\dots)=-D(\dots,\mathbf{u},\mathbf{v},\dots).
\]
Penukaran yang tidak bersebelahan dicapai dengan beberapa penukaran bersebelahan, dan setiap
langkah membalik tanda.

Untuk kelipatan, $D(\dots,c\mathbf{u},\dots,\mathbf{u},\dots)
=c\,D(\dots,\mathbf{u},\dots,\mathbf{u},\dots)=0$ menurut multilinearitas dan sifat
bergantian. Baris nol adalah kasus $c=0$ dari bentuk yang sama.
\end{proof}

Sekarang ketunggalannya. Rumus yang muncul di tengah bukti adalah rumus Leibniz, dan rumus
itu nanti memberi ekspansi kofaktor sebagai akibat.

\begin{teorema}[Ketunggalan dan rumus Leibniz]
Hanya ada satu fungsi yang memenuhi ketiga syarat definisi, dan nilainya adalah
\[
  \det A=\sum_{\sigma}\operatorname{sgn}(\sigma)\,
  a_{\sigma(1)1}a_{\sigma(2)2}\cdots a_{\sigma(n)n},
\]
dengan jumlah diambil atas semua permutasi $\sigma$ dari $\{1,\dots,n\}$ dan
$\operatorname{sgn}(\sigma)=(-1)^{k}$ bila $\sigma$ dicapai dengan $k$ penukaran.
\end{teorema}

\begin{proof}
Tulis baris ke-$i$ sebagai $\mathbf{a}_{i}=\sum_{j}a_{ij}\mathbf{e}_{j}$. Multilinearitas
memungkinkan mengeluarkan setiap baris menjadi jumlah,
\[
  D(A)=\sum_{j_{1},\dots,j_{n}}a_{1j_{1}}a_{2j_{2}}\cdots a_{nj_{n}}\,
  D(\mathbf{e}_{j_{1}},\dots,\mathbf{e}_{j_{n}}).
\]
Suku dengan $j_{p}=j_{q}$ untuk $p\neq q$ lenyap karena $D$ bergantian, jadi hanya
$(j_{1},\dots,j_{n})$ yang merupakan permutasi $\sigma$ yang bertahan. Menurut lemma
sebelumnya, $D(\mathbf{e}_{\sigma(1)},\dots,\mathbf{e}_{\sigma(n)})
=\operatorname{sgn}(\sigma)D(I_{n})=\operatorname{sgn}(\sigma)$.

Untuk ketunggalan, misalkan $D$ memenuhi ketiga syarat. Perhitungan yang sama memberi
$D(A)=\bigl(\sum_{\sigma}\operatorname{sgn}(\sigma)a_{\sigma(1)1}\cdots
a_{\sigma(n)n}\bigr)D(I_{n})=\det A$. Jadi setiap fungsi yang memenuhi ketiga syarat sama
dengan rumus itu.
\end{proof}

Rumus Leibniz benar tetapi tidak praktis: ia menjumlahkan $n!$ suku, dan untuk $n=12$
banyaknya $479\,001\,600$. Yang dipakai orang adalah bentuk lain dari rumus yang sama.

\begin{akibat}[Ekspansi kofaktor]
Untuk setiap baris $i$ yang tetap berlaku
\[
  \det A=\sum_{j=1}^{n}(-1)^{i+j}a_{ij}M_{ij},
\]
dengan $M_{ij}$ adalah determinan submatriks $(n-1)\times(n-1)$ hasil menghapus baris $i$
dan kolom $j$. Angka $C_{ij}=(-1)^{i+j}M_{ij}$ disebut \emph{kofaktor} entri $a_{ij}$.
Rumus yang sama berlaku bila ekspansi dilakukan sepanjang kolom.
\end{akibat}

\begin{proof}
Dalam rumus Leibniz, kelompokkan suku-suku menurut nilai $j$ pada posisi baris $i$. Setiap
suku yang memuat $a_{ij}$ berbentuk $a_{ij}$ kali tanda kali hasilkali $n-1$ entri lain,
satu dari setiap baris selain $i$ dan satu dari setiap kolom selain $j$. Hasilkali itu tepat
suku Leibniz untuk minor $M_{ij}$, dan tandanya $(-1)^{i+j}$ karena memindahkan baris $i$
ke posisi terakhir menuntut $n-i$ penukaran dan memindahkan kolom $j$ ke posisi terakhir
menuntut $n-j$ penukaran, sehingga tandanya $(-1)^{2n-i-j}=(-1)^{i+j}$. Penjumlahan atas
$j$ memberi rumus yang ditulis.

Untuk ekspansi sepanjang kolom, dipakai $\det A=\det A^{\top}$. Kesamaan itu sendiri
mengikuti rumus Leibniz: setiap suku untuk $A^{\top}$ berbentuk
$a_{\sigma(1)1}\cdots a_{\sigma(n)n}$ dengan urutan indeks terbalik, dan pemetaan
$\sigma\mapsto\sigma^{-1}$ mempertahankan tanda karena banyaknya penukaran sama.
\end{proof}

\begin{contoh}
Untuk
$A=\begin{pmatrix}2&1&-1\\ -3&-1&2\\ -2&1&2\end{pmatrix}$ ekspansi sepanjang baris pertama
memberi kofaktor
\[
  C_{11}=\begin{vmatrix}-1&2\\ 1&2\end{vmatrix}=-4,\qquad
  C_{12}=-\begin{vmatrix}-3&2\\ -2&2\end{vmatrix}=2,\qquad
  C_{13}=\begin{vmatrix}-3&-1\\ -2&1\end{vmatrix}=-5,
\]
sehingga $\det A=2(-4)+1(2)+(-1)(-5)=-8+2+5=-1$. Tanda bergantian itu bukan hiasan: tanpa
tanda, hasilnya $-8+2-5=-11$, dan semua kesimpulan sesudahnya akan salah.
\end{contoh}

Determinan berguna untuk matriks hasil kali karena ia multiplikatif.

\begin{teorema}
Untuk matriks $n\times n$ berlaku $\det(AB)=\det A\,\det B$.
\end{teorema}

\begin{proof}
Tetapkan $A$ dan definisikan $f(B)=\det(AB)$. Baris ke-$i$ dari $AB$ adalah kombinasi linear
baris-baris $B$ dengan koefisien dari baris ke-$i$ $A$, sehingga $f$ multilinear terhadap
baris-baris $B$. Bila dua baris $B$ sama, dua baris $AB$ juga sama, jadi $f$ bergantian.
Menurut ketunggalan, setiap fungsi bergantian multilinear $f$ memenuhi
$f(B)=\det(B)\,f(I_{n})$. Di sini $f(I_{n})=\det A$, jadi
$f(B)=\det B\,\det A$.
\end{proof}

\begin{akibat}
$A$ invertibel jika dan hanya jika $\det A\neq0$. Bila invertibel,
$\det(A^{-1})=1/\det A$. Selain itu $\det(cA)=c^{n}\det A$, dan determinan matriks
segitiga adalah hasilkali entri diagonalnya.
\end{akibat}

\begin{proof}
Bila $A$ invertibel, $\det A\,\det(A^{-1})=\det I_{n}=1$, sehingga $\det A\neq0$ dan
inversnya $1/\det A$. Bila $A$ tidak invertibel, kolom-kolomnya bergantung linear menurut
Bab 7, sehingga salah satu kolom merupakan kombinasi linear kolom lain. Multilinearitas
menguraikan $\det A$ menjadi jumlah determinan dengan dua kolom sama, jadi semuanya nol.
Rumus $\det(cA)=c^{n}\det A$ diperoleh dengan mengeluarkan faktor $c$ dari setiap baris.
Untuk matriks segitiga, ekspansi kofaktor sepanjang baris yang hanya bersisi satu entri
diagonal, diulang $n$ kali, menyisakan hasilkali diagonal.
\end{proof}

Akibat itu membuat eliminasi dan determinan bertemu. Eliminasi membawa $A$ ke matriks
segitiga atas $U$ lewat penukaran dan penambahan kelipatan baris. Penambahan tidak mengubah
determinan; penukaran membalik tandanya. Karena itu $\det A=(-1)^{s}\det U$ dengan $s$
banyaknya penukaran, dan $\det U$ adalah hasilkali pivot: determinan gratis setelah eliminasi,
dan begitulah pustaka menghitungnya.

Sisi lain determinan yang sama pentingnya adalah luas. Kolom-kolom
$B=\begin{pmatrix}3&1\\ 1&2\end{pmatrix}$ adalah dua anak panah, $(3,1)$ dan $(1,2)$.
Persegi satuan yang sisinya dibawa oleh kedua panah itu berubah menjadi jajaran genjang,
dan luasnya lima, tepat $\det B=3\cdot2-1\cdot1=5$. Gambar~\ref{fig:luas-determinan}
memperlihatkan pemetaan itu. Kalau kedua kolom segaris, jajaran genjangnya gepeng menjadi
satu ruas, luasnya nol, dan determinannya juga nol.

\begin{figure}[htbp]
\centering
\begin{tikzpicture}[scale=0.95,>=Stealth]
  \draw[layergrey!28,step=1] (-0.4,-0.4) grid (4.4,3.4);
  \draw[->,ink] (-0.5,0)--(4.6,0) node[right,font=\scriptsize]{$x$};
  \draw[->,ink] (0,-0.5)--(0,3.5) node[above,font=\scriptsize]{$y$};
  \fill[layerteal!16] (0,0)--(3,1)--(4,3)--(1,2)--cycle;
  \draw[thick,layerteal] (0,0)--(3,1)--(4,3)--(1,2)--cycle;
  \fill[layerorange!22] (0,0)--(1,0)--(1,1)--(0,1)--cycle;
  \draw[thick,layerorange] (0,0)--(1,0)--(1,1)--(0,1)--cycle;
  \draw[very thick,accent,->] (0,0)--(3,1)
    node[below right,font=\scriptsize]{$(3,1)$};
  \draw[very thick,accent,->] (0,0)--(1,2)
    node[above left,font=\scriptsize]{$(1,2)$};
  \foreach \t in {1,2,3,4} {\node[font=\tiny,below] at (\t,-0.02) {\t};}
  \foreach \t in {1,2,3} {\node[font=\tiny,left] at (-0.02,\t) {\t};}
  \node[font=\tiny,fill=white,inner sep=0.5pt] at (0.5,0.5) {1};
  \node[font=\scriptsize,fill=white,inner sep=1.5pt] at (2,1.5) {luas $5$};
\end{tikzpicture}
\caption{Determinan sebagai luas. Persegi satuan (jingga, luas $1$) dipetakan oleh
$B=\begin{pmatrix}3&1\\ 1&2\end{pmatrix}$ menjadi jajaran genjang (teal) yang dibangun
$(3,1)$ dan $(1,2)$. Luasnya $5$, sama dengan $\det B=3\cdot2-1\cdot1=5$: determinan
mengalikan luas sebesar lima, dan tandanya mencatat apakah arah putaran dibalik.}
\label{fig:luas-determinan}
\end{figure}

\begin{tangan}
Bahan: kertas berpetak, gunting, penggaris, lem.
\begin{enumerate}
  \item Gambar persegi panjang tiga petak kali satu petak. Luasnya tiga petak.
  \item Gunting persegi panjang itu. Dari ujung kiri, potong satu segitiga siku-siku dengan alas
        dua petak dan tinggi satu petak, lalu tempelkan potongan itu ke ujung kanan.
  \item Bentuk barunya jajaran genjang yang miring. Luasnya tetap tiga petak, karena tidak
        ada kertas yang dibuang dan tidak ada yang ditambah.
  \item Ulangi dengan kemiringan satu petak, lalu dua, lalu tiga. Setiap kali hitung alas
        kali tinggi, jangan panjang sisi miring.
\end{enumerate}
Yang harus pembaca lihat: bentuk berubah, luas tidak. Kemiringan hanya menggeser, tidak
meregangkan, jadi determinan tidak bertambah.
\end{tangan}

\begin{untukanak}
Ambil selembar kertas persegi panjang, lalu miringkan dengan mendorong sisi atasnya ke samping
tanpa memendekkan atau memanjangkan apa pun. Tingginya tetap, lebarnya tetap, luasnya tetap,
walaupun bentuknya bukan persegi panjang lagi. Yang berubah hanya arah putaran: telusuri sisi
bawah lalu sisi kanan dari pojok kiri bawah, dan kertas miring masih berputar ke arah yang
sama; tukar kedua sisi, dan arahnya berbalik seperti sepatu kiri di kaki kanan.

Di bab ini kita akan tahu mengapa satu bilangan dapat memberitahu luas sekaligus arah, dan
mengapa bilangan itu menjadi nol tepat ketika kertas terlipat menjadi garis. Menurutmu, kalau
kertas itu dilipat dua, masih adakah cara menghitung luasnya dengan bilangan yang sama?
\end{untukanak}

\section{Aturan Cramer dan adjoin}
\tingkat{menengah}

Determinan menjawab apakah jawaban tunggal ada. Pertanyaan berikutnya: dapatkah jawaban itu
ditulis sebagai rumus, bukan sebagai hasil langkah? Untuk sistem kecil, dapat, dan rumusnya
rapi.

\kamus{adjoin}{adjugate}{Peta kofaktor yang disusun transpos, dipakai untuk menulis invers
sebagai satu rumus.}

\begin{definisi}[Adjoin]
Untuk $A$ berukuran $n\times n$ dengan kofaktor $C_{ij}$, \emph{matriks kofaktor} adalah
matriks dengan entri $C_{ij}$ pada posisi $(i,j)$, dan \emph{adjoin}
$\operatorname{adj}A$ adalah transposnya, yaitu matriks dengan entri $C_{ji}$ pada posisi
$(i,j)$.\index{kofaktor}\index{adjoin}
\end{definisi}

Perhatikan pembalikan indeksnya: kofaktor $C_{ij}$ berasal dari menghapus baris $i$ dan kolom
$j$, tetapi diletakkan pada posisi $(j,i)$. Inilah sumber kekeliruan paling sering saat
menghitung adjoin dengan tangan.

\begin{teorema}
Untuk setiap matriks $n\times n$ berlaku $A\operatorname{adj}A
=(\det A)I_{n}=\operatorname{adj}A\,A$.
\end{teorema}

\begin{proof}
Entri $(i,j)$ dari $A\operatorname{adj}A$ adalah
$\sum_{k}a_{ik}C_{jk}$. Bila $i=j$, jumlah itu tepat ekspansi kofaktor sepanjang baris $i$,
yaitu $\det A$. Bila $i\neq j$, jumlah itu adalah ekspansi sepanjang baris $j$ dari matriks
yang diperoleh dari $A$ dengan mengganti baris $j$ oleh baris $i$. Matriks itu punya dua
baris sama, jadi determinannya nol. Karena itu $A\operatorname{adj}A=(\det A)I_{n}$.
Kesamaan kedua diperoleh dengan argumen yang sama pada kolom, atau dari
$\operatorname{adj}(A^{\top})=(\operatorname{adj}A)^{\top}$.
\end{proof}

\begin{akibat}
Bila $\det A\neq0$, maka $A^{-1}=\dfrac{1}{\det A}\operatorname{adj}A$.
\end{akibat}

\begin{proof}
Dari teorema sebelumnya, $A\bigl(\operatorname{adj}A/\det A\bigr)=I_{n}$ dan
$\bigl(\operatorname{adj}A/\det A\bigr)A=I_{n}$. Invers matriks tunggal, jadi matriks itu
adalah $A^{-1}$.
\end{proof}

Rumus itu berlaku untuk semua $n$, tetapi harganya mahal: $\operatorname{adj}A$ memuat
$n^{2}$ kofaktor, dan setiap kofaktor adalah determinan $(n-1)\times(n-1)$.

Dengan rumus invers itu, jawaban sistem dapat ditulis sebagai nisbah dua determinan.

\begin{teorema}[Aturan Cramer]
Misalkan $A$ berukuran $n\times n$ dengan $\det A\neq0$, dan $\mathbf{b}\in\mathbb{R}^{n}$.
Jawaban tunggal $A\mathbf{x}=\mathbf{b}$ adalah
\[
  x_{j}=\frac{\det A_{j}}{\det A},\qquad j=1,\dots,n,
\]
dengan $A_{j}$ adalah matriks yang diperoleh dari $A$ dengan mengganti kolom ke-$j$ oleh
$\mathbf{b}$.\index{cramer}
\end{teorema}

\begin{proof}
Karena $\det A\neq0$, akibat sebelumnya memberi
$\mathbf{x}=A^{-1}\mathbf{b}=(\operatorname{adj}A)\,\mathbf{b}/\det A$. Entri ke-$j$ dari
hasil kali itu adalah $\sum_{i}b_{i}C_{ji}/\det A$. Jumlah $\sum_{i}b_{i}C_{ji}$ adalah
ekspansi kofaktor sepanjang kolom ke-$j$ dari matriks yang kolom ke-$j$-nya telah diganti
oleh $\mathbf{b}$, yaitu $A_{j}$. Jadi $x_{j}=\det A_{j}/\det A$.
\end{proof}

\begin{contoh}
Kembali ke sistem pembuka, $A\mathbf{x}=\mathbf{b}$ dengan
$\mathbf{b}=(8,-11,-3)^{\top}$. Kita sudah menghitung $\det A=-1$, dan
\[
  \det A_{1}=\begin{vmatrix}8&1&-1\\ -11&-1&2\\ -3&1&2\end{vmatrix}=-2,\qquad
  \det A_{2}=-3,\qquad \det A_{3}=1 .
\]
Jadi $x_{1}=(-2)/(-1)=2$, $x_{2}=(-3)/(-1)=3$, dan $x_{3}=1/(-1)=-1$. Pemeriksaan:
$A(2,3,-1)^{\top}=(4+3+1,\,-6-3-2,\,-4+3-2)^{\top}=(8,-11,-3)^{\top}$.

Adjoin matriks itu, dengan kofaktor baris pertamanya sudah kita hitung di atas, adalah
\[
  \operatorname{adj}A=\begin{pmatrix}-4&-3&1\\ 2&2&-1\\ -5&-4&1\end{pmatrix},
  \qquad
  A^{-1}=\frac{1}{-1}\operatorname{adj}A
  =\begin{pmatrix}4&3&-1\\ -2&-2&1\\ 5&4&-1\end{pmatrix}.
\]
Baris pertama $A^{-1}$ kali $\mathbf{b}$ memberi $4(8)+3(-11)-1(-3)=32-33+3=2$, cocok
dengan $x_{1}$.
\end{contoh}

Aturan Cramer indah dan jarang dipakai untuk menghitung. Nilainya terletak pada hal lain:
rumus itu menunjukkan jawaban bergantung secara mulus pada $\mathbf{b}$ dan pada entri $A$,
dan ketergantungan itulah yang kita ukur lewat biaya dan kepekaan.

\section{Matriks elementer, invers, dan syarat setara}
\tingkat{lanjut}

Eliminasi baris dapat dipandang sebagai perkalian matriks. Pandangan itu mengubah pertanyaan
``apakah $A$ invertibel'' menjadi pertanyaan tentang daftar faktor, dan daftar faktor itu
memberi algoritme invers.

\kamus{matriks elementer}{elementary matrix}{Matriks identitas yang satu langkahnya sudah
dikerjakan, sehingga menjadi catatan satu operasi baris.}

\begin{definisi}[Matriks elementer]
\emph{Matriks elementer} adalah matriks yang diperoleh dari $I_{n}$ dengan satu operasi
baris elementer. Ada tiga jenis, sesuai tiga operasi: penukaran dua baris, perkalian satu
baris dengan skalar tak nol, dan penambahan kelipatan satu baris pada baris lain.
\end{definisi}

\begin{lemma}
Bila $E$ matriks elementer $m\times m$ dan $A$ matriks $m\times n$, maka $EA$ adalah hasil
mengerjakan operasi baris yang sama pada $A$. Setiap matriks elementer invertibel, dan
inversnya matriks elementer sejenis.
\end{lemma}

\begin{proof}
Baris ke-$i$ dari $EA$ adalah kombinasi linear baris-baris $A$ dengan koefisien baris ke-$i$
$E$. Untuk $E$ jenis ketiga, baris itu adalah $\mathbf{a}_{i}+c\mathbf{a}_{j}$, tepat seperti
operasi barisnya; untuk dua jenis lain pemeriksaannya sama langsung. Inversnya diperoleh
dengan membatalkan operasi: menukar baris membatalkan dirinya sendiri, dan
$\operatorname{diag}(1,\dots,c,\dots,1)^{-1}=\operatorname{diag}(1,\dots,1/c,\dots,1)$.
\end{proof}

Gambar~\ref{fig:tiga-gerakan} menjalankan ketiga operasi pada sistem $x+2y=4$,
$3x+4y=10$, yang berpotongan di $(2,1)$. Menukar baris hanya menukar nama; mengalikan baris
pertama dengan $2$ menulis garis yang sama dengan angka dua kali; menambahkan baris kedua pada
baris pertama memberi $4x+6y=14$, yang masih melewati $(2,1)$. Titik potongnya tidak bergeser,
tetapi determinan matriks koefisiennya berganti tanda, berlipat dua, lalu tetap. Himpunan
jawaban kekal; bilangan yang mengukur luas tidak.

\begin{figure}[htbp]
\centering
\begin{tikzpicture}[scale=0.58,>=Stealth]
  \begin{scope}
    \draw[layergrey!28,step=1] (-0.4,-0.4) grid (4.4,3.2);
    \draw[->,ink] (-0.5,0)--(4.6,0) node[right,font=\scriptsize]{$x$};
    \draw[->,ink] (0,-0.5)--(0,3.4) node[above,font=\scriptsize]{$y$};
    \draw[very thick,layerindigo] (0,2)--(4,0) node[pos=0.15,above,font=\scriptsize]{$R_{1}$};
    \draw[very thick,layerteal] (0,2.5)--(3.33,0) node[pos=0.85,above left,font=\scriptsize]{$R_{2}$};
    \fill[accent] (2,1) circle (2.6pt);
    \node[font=\scriptsize,below right] at (2,1) {$(2,1)$};
    \node[font=\scriptsize] at (2,-1.1) {(i) tukar baris};
  \end{scope}
  \begin{scope}[xshift=5.6cm]
    \draw[layergrey!28,step=1] (-0.4,-0.4) grid (4.4,3.2);
    \draw[->,ink] (-0.5,0)--(4.6,0) node[right,font=\scriptsize]{$x$};
    \draw[->,ink] (0,-0.5)--(0,3.4) node[above,font=\scriptsize]{$y$};
    \draw[very thick,layerindigo] (0,2)--(4,0) node[pos=0.15,above,font=\scriptsize]{$R_{1}$};
    \draw[very thick,layerindigo,densely dashed] (0,2)--(4,0);
    \draw[very thick,layerteal] (0,2.5)--(3.33,0) node[pos=0.85,above left,font=\scriptsize]{$R_{2}$};
    \fill[accent] (2,1) circle (2.6pt);
    \node[font=\scriptsize,below right] at (2,1) {$(2,1)$};
    \node[font=\scriptsize] at (2,-1.1) {(ii) kalikan baris};
  \end{scope}
  \begin{scope}[xshift=11.2cm]
    \draw[layergrey!28,step=1] (-0.4,-0.4) grid (4.4,3.2);
    \draw[->,ink] (-0.5,0)--(4.6,0) node[right,font=\scriptsize]{$x$};
    \draw[->,ink] (0,-0.5)--(0,3.4) node[above,font=\scriptsize]{$y$};
    \draw[very thick,layergrey,dash pattern=on 3pt off 2pt] (0,2)--(4,0);
    \draw[very thick,layerteal] (0,2.5)--(3.33,0)
      node[pos=0.8,below right,font=\scriptsize]{$R_{2}$};
    \draw[very thick,layerplum] (0,2.33)--(3.5,0)
      node[pos=0.2,above,font=\scriptsize]{$R_{1}+R_{2}$};
    \fill[accent] (2,1) circle (2.6pt);
    \node[font=\scriptsize,below right] at (2,1) {$(2,1)$};
    \node[font=\scriptsize] at (2,-1.1) {(iii) tambah kelipatan};
  \end{scope}
\end{tikzpicture}
\caption{Tiga operasi baris elementer pada sistem $x+2y=4$, $3x+4y=10$. Menukar baris
hanya menukar nama $R_{1}$ dan $R_{2}$; mengalikan $R_{1}$ dengan $2$ menulis garis yang
sama; menambahkan $R_{2}$ pada $R_{1}$ memutar garis pertama menjadi $4x+6y=14$. Pada ketiga
panel, kedua garis melewati titik potong yang sama, $(2,1)$.}
\label{fig:tiga-gerakan}
\end{figure}

\begin{teorema}
Matriks $A$ berukuran $n\times n$ invertibel jika dan hanya jika $A$ merupakan hasilkali
matriks-matriks elementer.
\end{teorema}

\begin{proof}
Hasilkali matriks-matriks invertibel juga invertibel, jadi satu arah jelas.

Untuk arah lain, misalkan $A$ invertibel dan bawa $A$ ke bentuk eselon tereduksi lewat
operasi baris, yang sama dengan mengalikan dari kiri oleh matriks elementer
$E_{1},\dots,E_{k}$. Menurut teorema tiga kasus, $A$ invertibel berarti setiap kolom $A$
memuat pivot, sehingga bentuk tereduksi $A$ adalah $I_{n}$, sebab $A$ persegi dengan $n$
pivot. Jadi $E_{k}\cdots E_{1}A=I_{n}$, dan karena itu
$A=E_{1}^{-1}\cdots E_{k}^{-1}$, hasilkali matriks elementer.
\end{proof}

Bukti itu sekaligus algoritme: mengerjakan operasi yang sama pada $I_{n}$ memberi $A^{-1}$,
jadi $\bigl(A\mid I_{n}\bigr)\to\bigl(I_{n}\mid A^{-1}\bigr)$ menyelesaikan invers tanpa
satu pun determinan, jauh lebih murah daripada adjoin.

Kita kumpulkan yang sudah terbukti menjadi satu daftar. Daftar inilah yang dipakai orang
setiap hari, dan isinya wajah lain dari satu pernyataan.

\begin{teorema}[Syarat setara]
Untuk matriks $A$ berukuran $n\times n$, pernyataan berikut setara.
\begin{enumerate}
  \item $A$ invertibel.
  \item $A\mathbf{x}=\mathbf{0}$ hanya punya jawaban trivial $\mathbf{x}=\mathbf{0}$.
  \item Bentuk eselon tereduksi $A$ adalah $I_{n}$.
  \item $A$ adalah hasilkali matriks-matriks elementer.
  \item $A\mathbf{x}=\mathbf{b}$ konsisten untuk setiap $\mathbf{b}\in\mathbb{R}^{n}$.
  \item $A\mathbf{x}=\mathbf{b}$ punya tepat satu jawaban untuk setiap
        $\mathbf{b}\in\mathbb{R}^{n}$.
  \item $\det A\neq0$.
\end{enumerate}
\end{teorema}

\begin{proof}
Kita rangkaikan implikasi $(1)\Rightarrow(2)\Rightarrow(3)\Rightarrow(4)\Rightarrow(1)$
lalu $(1)\Rightarrow(5)\Rightarrow(6)\Rightarrow(1)$ dan $(1)\Leftrightarrow(7)$.

$(1)\Rightarrow(2)$: kalikan $A\mathbf{x}=\mathbf{0}$ dari kiri dengan $A^{-1}$.
$(2)\Rightarrow(3)$: inti $A$ hanya $\{\mathbf{0}\}$ berarti nullity nol, sehingga menurut
rank-nullity di Bab 7, $\operatorname{rank}A=n$; jadi ada $n$ pivot dan bentuk tereduksinya
$I_{n}$. $(3)\Rightarrow(4)$: langkah eliminasinya memberi
$E_{k}\cdots E_{1}A=I_{n}$. $(4)\Rightarrow(1)$: hasilkali matriks-matriks invertibel juga
invertibel.
$(1)\Rightarrow(5)$: ambil $\mathbf{x}=A^{-1}\mathbf{b}$. $(5)\Rightarrow(6)$: misalkan ada
$\mathbf{b}$ dengan dua jawaban berbeda $\mathbf{x}\neq\mathbf{y}$. Maka
$\mathbf{x}-\mathbf{y}\neq\mathbf{0}$ dan $A(\mathbf{x}-\mathbf{y})=\mathbf{0}$, sehingga
kolom-kolom $A$ bergantung linear dan $\operatorname{rank}A<n$. Ruang kolom $A$ karena itu
berdimensi kurang dari $n$, sehingga ada $\mathbf{b}'$ yang tidak terjangkau, bertentangan
dengan (5).
$(6)\Rightarrow(1)$: bila $A$ tidak invertibel, kolom-kolomnya bergantung, jadi ada
$\mathbf{x}\neq\mathbf{0}$ dengan $A\mathbf{x}=\mathbf{0}$, sehingga jawaban untuk
$\mathbf{b}=\mathbf{0}$ tidak tunggal. Terakhir $(1)\Leftrightarrow(7)$ adalah akibat di
bagian sebelumnya.
\end{proof}

Daftar itu menjelaskan mengapa empat pertanyaan yang tampak berbeda, tentang solusi, bentuk
eselon, faktorisasi, dan determinan, sebenarnya satu pertanyaan; yang berubah hanya alat
ukurnya. Tetapi daftar itu diam tentang biaya dan kepekaan, dan dua hal itulah yang menentukan
apakah jawaban yang diperoleh dapat dipercaya.

\begin{tantangan}
Sistem berikut tampak biasa: tiga persamaan untuk tiga bilangan, tidak ada baris kembar yang
mencolok, sehingga kelihatannya jawabannya tunggal.
\[
  x+y+z=3,\qquad x+2y+3z=6,\qquad 2x+3y+4z=9 .
\]
Tentukan semua jawabannya. Berapa banyak jawabannya, dan mengapa menghitung banyaknya
persamaan menyesatkan di sini?

Petunjuk 1: jumlahkan dua persamaan pertama, lalu bandingkan hasilnya dengan persamaan
ketiga.

Petunjuk 2: persamaan ketiga tidak menambah keterangan apa pun, jadi hanya dua keterangan
yang bebas. Setelah $z$ dipilih, $y$ lalu $x$ ikut tertentu; tuliskan jawaban sebagai satu
keluarga dengan satu parameter.

Kalau mentok, kembali ke bagian tiga nasib dan hitung pivotnya.
\end{tantangan}

\subsection*{Biaya dan kepekaan}

Eliminasi pada matriks $n\times n$ memakai sekitar $n^{3}/3$ perkalian: untuk $n=1000$
sekitar $3{,}33\times10^{8}$, yang masih cepat. Aturan Cramer dengan kofaktor memakai
$n\cdot n!$ perkalian, sekitar $5{,}75\times10^{9}$ untuk $n=12$ dan di atas $10^{19}$ untuk
$n=20$; karena itu ia hanya dipakai untuk $n$ kecil atau untuk keperluan teoretis.

Yang lebih halus adalah kepekaan. Misalkan $A$ invertibel dan $\mathbf{b}\neq\mathbf{0}$,
dan misalkan $\mathbf{b}$ berubah menjadi $\mathbf{b}+\delta\mathbf{b}$ sehingga jawabannya
berubah menjadi $\mathbf{x}+\delta\mathbf{x}$.

\begin{teorema}
Dengan notasi di atas, untuk sembarang norma matriks yang subordinate berlaku
\[
  \frac{\lVert\delta\mathbf{x}\rVert}{\lVert\mathbf{x}\rVert}
  \le\kappa(A)\,
  \frac{\lVert\delta\mathbf{b}\rVert}{\lVert\mathbf{b}\rVert},
  \qquad \kappa(A)=\lVert A\rVert\,\lVert A^{-1}\rVert .
\]
\end{teorema}

\begin{proof}
Dari $A\,\delta\mathbf{x}=\delta\mathbf{b}$ diperoleh
$\delta\mathbf{x}=A^{-1}\delta\mathbf{b}$, sehingga
$\lVert\delta\mathbf{x}\rVert\le\lVert A^{-1}\rVert\lVert\delta\mathbf{b}\rVert$.
Dari $\mathbf{b}=A\mathbf{x}$ diperoleh
$\lVert\mathbf{b}\rVert\le\lVert A\rVert\lVert\mathbf{x}\rVert$, yaitu
$1/\lVert\mathbf{x}\rVert\le\lVert A\rVert/\lVert\mathbf{b}\rVert$. Mengalikan kedua
ketaksamaan memberi hasil yang diminta.
\end{proof}

Bilangan $\kappa(A)$ disebut \emph{bilangan kondisi}. Ia tidak bergantung pada skala:
mengalikan $A$ dengan $c$ tidak mengubahnya, karena $\lVert cA\rVert\lVert (cA)^{-1}\rVert
=|c|\lVert A\rVert\cdot|c|^{-1}\lVert A^{-1}\rVert$. Determinan tidak punya sifat itu, dan
inilah alasan kita tidak memakai determinan sebagai ukuran kepekaan. Determinan kecil hanya
berarti volume kecil; ia dapat kecil karena seluruh matriks diperkecil, sedangkan kepekaan
tidak berubah.

Untuk sistem pembuka, $\kappa(S)\approx1{,}01\times10^{5}$. Dengan presisi mesin sekitar
$10^{-16}$, sisa angka yang dapat dipercaya pada jawaban sekitar $10^{-11}$ relatif: cukup
untuk contoh ini, habis untuk matriks yang lebih buruk. Bab 22 membahas ambang dan cara
mengenali matriks yang sudah tidak dapat dipercaya.
\index{kondisi}

\begin{untukanak}
Bayangkan papan panjang yang dipakai sebagai landasan untuk mengukur tinggi sebuah bukit.
Papan itu hampir mendatar, hanya miring sedikit. Kalau ujung yang di bawah diangkat sejari,
ujung yang lain naik jauh lebih banyak, karena papan itu bertumpu pada kemiringan yang
sangat kecil. Sejari di satu tempat menjadi sejengkal di tempat lain.

Sekarang bayangkan papannya dibuat sangat curam. Angkat ujung bawahnya sejari, dan ujung
atasnya hanya naik sedikit sekali. Pekerjaan yang sama menghasilkan dua akibat yang sangat
berbeda, hanya karena landasannya berbeda.

Hitungan kita bekerja seperti papan itu. Ada susunan angka yang membuat perubahan kecil pada
masukan menghasilkan perubahan kecil pada jawaban; ada pula susunan yang membuat perubahan
sebesar debu menghasilkan jawaban yang sama sekali lain. Yang menentukan bukan besar kecilnya
angka, melainkan bentuk susunannya.

Di bab ini kita akan tahu mengapa ada satu bilangan yang dapat meramalkan hal itu sebelum
kita menghitung, dan mengapa bilangan itu tidak berubah ketika semua angka diperbesar atau
diperkecil dengan perbandingan yang sama.

Kalau kamu menimbang dengan timbangan yang landasannya miring, apakah hasilnya salah, atau
hanya tidak dapat dipercaya?
\end{untukanak}

\section{Contoh hitung}
\tingkat{menengah}

Kita kerjakan satu sistem dengan dua cara, hitung satu determinan $4\times4$, lalu mengukur
kepekaan dua matriks.

\begin{resep}
\begin{enumerate}
  \item Tulis matriks diperbesar $\bigl(A\mid\mathbf{b}\bigr)$: kolom koefisien dulu, kolom
        ruas kanan paling akhir.
  \item Kerjakan kolom paling kiri. Pilih satu entri tak nol sebagai pivot; kalau entri
        teratas nol, tukar baris lebih dulu.
  \item Nolkan semua entri di bawah pivot: untuk tiap baris $j$ di bawahnya,
        $R_{j}\leftarrow R_{j}-\frac{a_{ji}}{a_{ii}}R_{i}$.
  \item Geser satu kolom ke kanan dan ulangi; berhenti saat tidak ada baris tak nol lagi,
        lalu baca nasib sistem dari letak pivot.
\end{enumerate}
Peringatan urutan: pivot dinolkan dari atas ke bawah, kolom demi kolom, sebelum substitusi
mundur dimulai. Membalik urutan itu menghapus pivot yang sudah diperoleh dan membuat
pekerjaan harus diulang.
\end{resep}

Sistemnya $A\mathbf{x}=\mathbf{b}$ dengan
\[
  A=\begin{pmatrix}2&1&-1\\ -3&-1&2\\ -2&1&2\end{pmatrix},
  \qquad \mathbf{b}=\begin{pmatrix}8\\ -11\\ -3\end{pmatrix}.
\]

Cara pertama, eliminasi, dengan langkah ditulis lengkap:
\[
  \left(\begin{array}{ccc|c}2&1&-1&8\\ -3&-1&2&-11\\ -2&1&2&-3\end{array}\right)
  \xrightarrow[R_{3}\leftarrow R_{3}+R_{1}]{R_{2}\leftarrow R_{2}+\tfrac{3}{2}R_{1}}
  \left(\begin{array}{ccc|c}2&1&-1&8\\ 0&\tfrac{1}{2}&\tfrac{1}{2}&1\\
  0&2&1&5\end{array}\right),
\]
lalu $R_{3}\leftarrow R_{3}-4R_{2}$ memberi
\[
  \left(\begin{array}{ccc|c}2&1&-1&8\\ 0&\tfrac{1}{2}&\tfrac{1}{2}&1\\
  0&0&-1&1\end{array}\right).
\]
Substitusi mundur: dari baris ketiga $-x_{3}=1$, jadi $x_{3}=-1$. Baris kedua
$\tfrac{1}{2}x_{2}+\tfrac{1}{2}(-1)=1$ memberi $x_{2}=3$. Baris pertama
$2x_{1}+3+1=8$ memberi $x_{1}=2$. Hasilnya $\mathbf{x}=(2,3,-1)$. Pivot yang muncul adalah
$2$, $\tfrac{1}{2}$, dan $-1$, jadi hasilkali pivot $\det A=2\cdot\tfrac{1}{2}\cdot(-1)=-1$,
cocok dengan ekspansi kofaktor tadi. Di sini tidak ada penukaran baris, sehingga tandanya
tidak berubah.

Cara kedua, aturan Cramer, sudah kita jalankan pada sistem yang sama di bagian sebelumnya.
Ketiga determinan pengganti memberi $\det A_{1}=-2$, $\det A_{2}=-3$, dan $\det A_{3}=1$,
sehingga $\mathbf{x}=(-2,-3,1)/(-1)=(2,3,-1)$. Dua jalan, satu jawaban; pemeriksaan semacam
ini murah dan sebaiknya selalu dikerjakan.

Sekarang determinan $4\times4$. Ambil
\[
  M=\begin{pmatrix}2&0&1&3\\ 1&4&0&2\\ 3&1&2&1\\ 0&2&1&5\end{pmatrix}.
\]
Ekspansi sepanjang baris pertama memakai empat minor, dan salah satunya kita hitung penuh:
\[
  M_{11}=\begin{vmatrix}4&0&2\\ 1&2&1\\ 2&1&5\end{vmatrix}
  =4(10-1)-0(5-2)+2(1-4)=36-6=30 .
\]
Dengan cara serupa $M_{12}=15$, $M_{13}=-45$, dan $M_{14}=-15$. Karena itu
\[
  \det M=2(30)-0(15)+1(-45)-3(-15)=60-45+45=60 .
\]
Banyaknya suku yang dijumlahkan adalah $4!=24$, dan untuk $5\times5$ sudah $120$:
pertumbuhannya faktorial. Inilah alasan pustaka tidak memakai cara ini.

Sekarang kepekaan. Matriks Hilbert $4\times4$ didefinisikan oleh $h_{ij}=1/(i+j-1)$.
Determinannya tepat $1/6\,048\,000\approx1{,}65\times10^{-7}$, sangat kecil, sehingga orang
mudah menyimpulkan bahwa matriks ini berbahaya. Bilangan kondisinya $15\,513{,}74$. Matriks
pembuka tadi punya $\det S=-0{,}003$ dan $\kappa(S)\approx1{,}01\times10^{5}$, jadi
determinannya delapan belas ribu kali lebih besar, sedangkan kepekaannya sekitar enam
setengah kali lebih buruk. Determinan bukan ukuran kepekaan; bilangan kondisi itu ukurannya.

Uji lompatan untuk $S$ memakai $\mathbf{x}=(1,1,1)$, sehingga
$\mathbf{b}=(6,\,15,\,24{,}001)$. Mengubah komponen ketiga menjadi $24{,}002$ mengubah
jawaban menjadi $(2,-1,2)$. Perubahan relatif data $3{,}46\times10^{-5}$, perubahan relatif
jawaban $\sqrt{2}\approx1{,}414$, nisbahnya $40\,916$: di bawah $\kappa(S)$ seperti
dijanjikan teorema, tetapi hanya dengan faktor dua setengah.

\kertas{Manakah yang lebih peka, $M$ atau matriks Hilbert $4\times4$? Ramalkan dengan membandingkan
determinannya, lalu hitung bilangan kondisi keduanya dengan \texttt{numpy.linalg.cond}.
Bandingkan juga dengan $10^{-3}I_{4}$, yang determinannya jauh lebih kecil, lalu jelaskan
mengapa aturan ``determinan kecil berarti peka'' gagal. Terakhir, kalikan $M$ dengan
$10^{-4}$: bilangan kondisinya tidak berubah, sedangkan determinannya mengecil sebesar
$10^{-16}$.}

\section{Lab komputasi}
\tingkat{lanjut}

\begin{lab}
\begin{verbatim}
import numpy as np, math

def det_kofaktor(A):
    n = A.shape[0]
    if n == 1:
        return A[0, 0]
    return sum((-1) ** j * A[0, j] *
               det_kofaktor(np.delete(np.delete(A, 0, 0), j, 1))
               for j in range(n))

A = np.array([[2., 1., -1.], [-3., -1., 2.], [-2., 1., 2.]])
b = np.array([8., -11., -3.])
x = np.linalg.solve(A, b)
print("det A    :", np.linalg.det(A))
print("det kof  :", det_kofaktor(A))
print("rank A   :", np.linalg.matrix_rank(A))
print("x        :", x)
print("sisa     :", np.linalg.norm(A @ x - b))
for j in range(3):
    Aj = A.copy(); Aj[:, j] = b
    print("det A%d   : %.1f  -> x%d = %.6f" % (j + 1, np.linalg.det(Aj),
          j + 1, np.linalg.det(Aj) / np.linalg.det(A)))

M = np.array([[2., 0., 1., 3.], [1., 4., 0., 2.],
              [3., 1., 2., 1.], [0., 2., 1., 5.]])
print("det M    :", det_kofaktor(M), np.linalg.det(M))

H = np.array([[1. / (i + j + 1) for j in range(4)] for i in range(4)])
print("det H4   :", np.linalg.det(H))
print("cond H4  : %.2f" % np.linalg.cond(H))

S = np.array([[1., 2., 3.], [4., 5., 6.], [7., 8., 9.001]])
bs = S @ np.array([1., 1., 1.])
print("det S    :", np.linalg.det(S), " cond S: %.1f" % np.linalg.cond(S))
print("x        :", np.linalg.solve(S, bs))
bs2 = bs.copy(); bs2[2] += 0.001
print("x'       :", np.linalg.solve(S, bs2))
print("biaya    :", math.factorial(12) * 12, round(1000 ** 3 / 3))
\end{verbatim}
Keluaran sebenarnya:
\begin{verbatim}
det A    : -1.0
det kof  : -1.0
rank A   : 3
x        : [ 2.  3. -1.]
sisa     : 1.9860273225978185e-15
det A1   : -2.0  -> x1 = 2.000000
det A2   : -3.0  -> x2 = 3.000000
det A3   : 1.0  -> x3 = -1.000000
det M    : 60.0 60.000000000000036
det H4   : 1.6534391534390412e-07
cond H4  : 15513.74
det S    : -0.0030000000000036655  cond S: 101072.3
x        : [1. 1. 1.]
x'       : [ 2. -1.  2.]
biaya    : 5748019200 333333333
\end{verbatim}
\end{lab}

Empat hal perlu dibaca dari keluaran itu.

Pertama, $x=(2,3,-1)$ muncul dengan sisa $1{,}99\times10^{-15}$, yaitu nol sampai presisi
mesin. Sisa itu sengaja dihitung: jawaban pustaka tidak pernah persis, dan sisa adalah cara
termurah memeriksa apakah jawaban itu memenuhi sistemnya. Pada matriks buruk, sisa kecil
tidak menjamin galat kecil; itu topik Bab 22.

Kedua, determinan hasil ekspansi kofaktor dan hasil pustaka sama-sama $-1$, dan untuk $M$
keduanya $60$. Pustaka tidak memakai kofaktor; ia memakai faktorisasi dan mengalikan pivot.
Selisih $3{,}6\times10^{-14}$ pada $\det M$ adalah pembulatan, bukan perbedaan definisi.

Ketiga, $\det H_{4}$ dilaporkan $1{,}65\times10^{-7}$ sedangkan nilainya tepat
$1/6\,048\,000$, jadi angka mesin itu benar sampai tujuh digit. Untuk orde lebih tinggi,
$1/\det$ melampaui $10^{12}$ dengan cepat, dan di situlah determinan berhenti berguna
sebagai penanda.

Keempat, lompatan jawaban pada $S$ terlihat utuh: dari $(1,1,1)$ ke $(2,-1,2)$ hanya karena
satu angka berubah seperseribu. Angka yang dilaporkan untuk $\det S$ pun tidak persis
$-0{,}003$; selisih $3{,}7\times10^{-15}$ muncul karena $9{,}001$ tidak terwakili tepat
dalam pecahan biner.

Terakhir, biaya: $5{,}748\times10^{9}$ perkalian untuk aturan Cramer pada $n=12$ melawan
$3{,}33\times10^{8}$ untuk eliminasi pada $n=1000$. Matriks yang jauh lebih besar dikerjakan
dengan pekerjaan yang jauh lebih sedikit; itu bedanya pangkat tiga dan faktorial.

\begin{untukdosen}
Urutan mengajar yang berhasil dipakai: eliminasi lebih dulu, dengan tiga kasus ditunjukkan
pada satu matriks koefisien tetap dan tiga ruas kanan berbeda. Mahasiswa yang melihat
konsistensi berubah karena satu angka di $\mathbf{b}$ biasanya langsung menerima bahwa
pertanyaan tentang $A$ dan pertanyaan tentang $\mathbf{b}$ harus dipisahkan. Baru setelah itu
determinan masuk sebagai jawaban atas pertanyaan tentang $A$ saja.

Dua miskonsepsi khas. Pertama, ``determinan kecil berarti sistem tidak stabil''. Matriks
$10^{-6}I$ punya determinan $10^{-18}$ dan bilangan kondisi satu; sebaliknya $H_{4}$ punya
determinan $1{,}65\times10^{-7}$ dan kondisi $1{,}5\times10^{4}$. Kekeliruan ini muncul
setiap tahun, dan cara membongkarnya adalah meminta mahasiswa mengalikan satu matriks
dengan skalar lalu membandingkan kedua bilangan. Kedua, ``aturan Cramer adalah cara
menyelesaikan sistem''. Secara teori benar; secara biaya salah. Beri satu sistem $12\times12$
dan minta mereka memperkirakan biaya kedua cara sebelum menghitung.

Aktivitas kelas yang berguna: bagikan tiga sistem dengan matriks koefisien sama dan ruas
kanan $(4,9,13)$, $(4,9,14)$, dan $(0,0,0)$. Minta setiap kelompok menentukan nasib
sistemnya hanya dari bentuk eselon, lalu menyimpulkan apa yang dapat dan tidak dapat
disimpulkan dari $A$ seorang diri.

Tanda kesiapan lanjut: mahasiswa dapat menyebutkan tujuh syarat setara tanpa menghafal
urutan, dan dapat menjelaskan mengapa syarat tentang determinan berada di daftar yang sama
dengan syarat tentang pivot.
\end{untukdosen}

\section{Ringkasan}
\tingkat{dasar}

\begin{ringkasan}
\begin{itemize}
  \item Operasi baris elementer tidak mengubah himpunan jawaban, karena setiap operasi dapat
        dibatalkan. Karena bentuk tereduksi tunggal, letak pivot menentukan tiga nasib: tidak
        konsisten bila pivot muncul di kolom terakhir, tunggal bila setiap kolom koefisien
        punya pivot, banyak bila ada kolom tanpa pivot.
  \item Determinan adalah satu-satunya fungsi multilinear, bergantian, dan ternormalisasi.
        Rumus Leibniz dan ekspansi kofaktor adalah akibatnya, bukan definisinya.
  \item $\det(AB)=\det A\det B$, sehingga $A$ invertibel tepat ketika $\det A\neq0$; dan
        $\det A$ sama dengan hasilkali pivot setelah eliminasi, kali tanda penukaran.
  \item Adjoin memberi $A\operatorname{adj}A=(\det A)I$ dan
        $A^{-1}=\operatorname{adj}A/\det A$; aturan Cramer menulis
        $x_{j}=\det A_{j}/\det A$.
  \item Tujuh syarat invertibilitas setara: inti trivial, bentuk tereduksi $I_{n}$, hasilkali
        matriks elementer, konsisten dan tunggal untuk semua $\mathbf{b}$, dan
        $\det A\neq0$.
  \item Biaya eliminasi sebanding $n^{3}/3$, biaya kofaktor sebanding $n\cdot n!$. Kepekaan
        diukur oleh $\kappa(A)=\lVert A\rVert\lVert A^{-1}\rVert$, bukan oleh $\det A$.
\end{itemize}
\end{ringkasan}

\section{Latihan}
\tingkat{menengah}

\begin{latihan}
\begin{enumerate}
  \item Selesaikan dengan eliminasi, lalu periksa dengan aturan Cramer:
        \[
          \begin{pmatrix}1&2&3\\ 2&5&3\\ 1&0&8\end{pmatrix}\mathbf{x}
          =\begin{pmatrix}9\\ 8\\ 26\end{pmatrix}.
        \]
        Tuliskan pivot yang muncul, catat berapa kali baris ditukar, lalu simpulkan
        $\det A$ dari hasilkali pivot.

  \item Hitung $\det M$ untuk
        $M=\begin{pmatrix}0&1&2&3\\ 1&0&3&2\\ 2&3&0&1\\ 3&2&1&0\end{pmatrix}$
        dengan ekspansi kofaktor sepanjang satu baris, lalu hitung dengan
        \texttt{numpy.linalg.det}. Kedua hasil tidak sama persis: yang pertama eksak, yang
        kedua $5{,}3\times10^{-15}$. Jelaskan mengapa. Sesudah itu cari vektor tak nol
        $\mathbf{v}$ dengan $M\mathbf{v}=\mathbf{0}$ dan simpulkan $\det M=0$ tanpa
        ekspansi apa pun.

  \item Buktikan bahwa bila dua kolom $A$ sama, maka $\det A=0$, tanpa memakai
        $\det A=\det A^{\top}$.
        (Petunjuk: pakai rumus Leibniz, kelompokkan suku-suku menurut pasangan permutasi yang
        bertukar pada dua kolom itu, dan tunjukkan tandanya berlawanan.)

  \item Buktikan bahwa $\det(\operatorname{adj}A)=(\det A)^{n-1}$ untuk $n\ge2$.
        (Petunjuk: ambil determinan pada $A\operatorname{adj}A=(\det A)I$ dan gunakan
        $\det(cI)=c^{n}$.)

  \item Nilai klaim ini: \emph{bila $\det A$ sangat kecil, maka jawaban $A\mathbf{x}
        =\mathbf{b}$ sangat peka terhadap perubahan $\mathbf{b}$}. Tentukan benar, salah,
        atau belum cukup bukti, dan sertakan contoh penyangkal atau alasan tambahan yang
        membuatnya benar.

  \item Untuk $A=\begin{pmatrix}1&1\\ 1&1+\varepsilon\end{pmatrix}$ dengan
        $\varepsilon>0$, hitung $\kappa(A)$ dalam norma $1$ sebagai fungsi
        $\varepsilon$. Tentukan nilai $\varepsilon$ yang membuat $\kappa(A)>10^{6}$, lalu
        periksa dengan \texttt{numpy.linalg.cond}.

  \item Buktikan: bila $A$ invertibel dan $B$ diperoleh dari $A$ dengan menambahkan
        kelipatan satu baris pada baris lain, maka $\det B=\det A$ dan
        $B^{-1}$ dapat dinyatakan lewat $A^{-1}$ dan satu matriks elementer.

  \item Misalkan $D=\operatorname{diag}(d_{1},\dots,d_{n})$ dan $A$ berukuran $n\times n$.
        Tunjukkan bahwa $\det(A+AD)=\det(A)\prod_{i=1}^{n}(1+d_{i})$, lalu tentukan syarat
        pada $d_{1},\dots,d_{n}$ agar $A+AD$ invertibel. Periksa jawabanmu dengan satu
        matriks $3\times3$ di Python.

  \item Pilih satu matriks $3\times3$ dengan entri bulat kecil dan $\det A=1$, hitung
        $A^{-1}$ lewat adjoin dan lewat \texttt{numpy.linalg.inv}, lalu bandingkan
        keduanya. Ulangi untuk matriks dengan $\det A=7$ dan jelaskan mengapa entri
        $A^{-1}$ menjadi pecahan.
\end{enumerate}

Petunjuk tiga butir tersulit. Butir 4: ambil determinan kedua ruas, ingat
$\det(cI)=c^{n}$, dan perhatikan bahwa $\det A$ boleh nol sehingga kesimpulan tentang
$\operatorname{adj}A$ harus dibaca sebagai matriks, bukan invers. Butir 6: dalam norma $1$,
$\lVert A\rVert_{1}$ adalah jumlah kolom terbesar, dan untuk matriks ini kolom keduanya
yang terbesar; hitung $\lVert A^{-1}\rVert_{1}$ dari rumus $2\times2$ dan perhatikan bahwa
jawabannya sebanding $1/\varepsilon$. Butir 8: tulis $A+AD$ sebagai $A(I+D)$ sebelum
mengeluarkan determinan, dan ingat bahwa $I+D$ bersifat diagonal sehingga determinannya
hasilkali entri diagonalnya.
\end{latihan}

\begin{untukpeneliti}
Yang benar-benar terbukti di bab ini, untuk medan $\mathbb{R}$ dan dimensi hingga: ketunggalan
determinan sebagai fungsi $n$-linear bergantian yang ternormalisasi beserta rumus permutasi,
di Hoffman dan Kunze, \emph{Linear Algebra} edisi ke-2, bagian 5.2 dan 5.3, berlaku lebih umum
atas gelanggang komutatif dengan identitas (bagian 5.1); multiplikativitas dan kesetaraan
invertibilitas dengan $\det\neq0$, di Axler, \emph{Linear Algebra Done Right} edisi ke-4,
bagian 9C, pernyataan 9.49 dan 9.50, dengan determinan operator yang bebas basis; serta daftar
tujuh syarat setara, di Anton edisi ke-11, Teorema 2.3.8.

Yang dipopulerkan: ``determinan adalah volume''. Yang terbukti adalah $|\det A|$ sama dengan
volume paralelepipedum kolom-kolomnya; tanpa nilai mutlak, klaim itu salah untuk determinan
negatif. Juga keliru secara numerik: aturan Cramer sebagai cara menyelesaikan sistem di atas
$n=11$, tempat biayanya faktorial.

Pertanyaan terbuka. Batas kepekaan pada teorema di atas bersifat normwise dan worst-case,
sering terlalu pesimistis dibanding analisis componentwise. Lebih mendasar: faktor pertumbuhan
maksimum eliminasi Gauss dengan pivot parsial belum diketahui untuk $n$ umum. Batas terbukti
$2^{n-1}$, sedangkan percobaan numerik memberi jauh lebih kecil, dan kesenjangan itu terbuka.

Batas keberlakuan. Argumen runtuh bila medan bukan lapangan, sebab modul tidak selalu punya
basis dan eliminasi tidak selalu dapat membagi. Untuk $n$ sangat besar, determinan meluap
atau mengalir habis jauh sebelum rank-nya hilang, sehingga pustaka melaporkan logaritma
determinan. Pernyataan kepekaan juga bergantung pada norma, dan tidak ada pilihan yang
kanonik.
\end{untukpeneliti}

\peta{Bab ini mengikuti Anton, \emph{Elementary Linear Algebra} edisi ke-11, bagian 1.1
sampai 1.6 (sistem linear, eliminasi Gauss-Jordan, invers, matriks elementer) dan 2.1 sampai
2.3 (kofaktor, determinan lewat reduksi baris, aturan Cramer, termasuk Teorema 2.3.8 tentang
syarat setara). Perlakuan aksiomatik mengikuti Hoffman dan Kunze, \emph{Linear Algebra} edisi
ke-2, bagian 5.2 dan 5.3, dengan bagian 1.4 untuk ketunggalan bentuk eselon tereduksi.
Pandangan bebas basis ada di Axler, \emph{Linear Algebra Done Right} edisi ke-4, bagian 9C.
Untuk eliminasi sebagai faktorisasi dan pandangan $A=CR$, lihat Strang, \emph{Introduction to
Linear Algebra} edisi ke-6, Bab 2 dan catatan ``Elimination and Factorization''. Sisi numerik,
yaitu bilangan kondisi, galat pembulatan, dan biaya, dibahas di Bab 22 buku ini.}

\chapter{Nilai Eigen dan Diagonalisasi}

Setiap menit, sebuah pemantau jaringan membagi seluruh permintaan yang masuk ke dalam tiga
keranjang: host A, host B, dan host C. Dari catatan lama ia menyusun aturan perpindahan. Dari
permintaan yang menuju A, 80 persen tetap di A, 10 persen pindah ke B, dan 10 persen pindah ke C.
Dari B, 30 persen pindah ke A, 60 persen bertahan, dan 10 persen pindah ke C. Dari C, 20 persen
pindah ke A, 30 persen ke B, dan 50 persen bertahan. Bila bagian permintaan di ketiga keranjang
pada menit ke-$k$ ditulis sebagai kolom $\mathbf{x}_{k}$, aturan itu berbentuk
$\mathbf{x}_{k+1}=M\mathbf{x}_{k}$ dengan
\[
  M=\begin{pmatrix}0.8&0.3&0.2\\ 0.1&0.6&0.3\\ 0.1&0.1&0.5\end{pmatrix}.
\]

Mulai dari $\mathbf{x}_{0}=(1,0,0)$, yaitu semua permintaan mula-mula menuju A, pemantau itu
menghitung sepuluh langkah dan memperoleh
$\mathbf{x}_{10}=(0.5672,\;0.2661,\;0.1666)$. Tiga puluh langkah kemudian angkanya
$(0.5667,\;0.2667,\;0.1667)$. Jumlah ketiga entri tetap satu, dan pada pembacaan kedua angka itu
berhenti bergerak sampai empat angka di belakang koma.

Pertanyaan yang layak diajukan bukan berapa hasilnya setelah empat puluh menit; perkalian berulang
sudah menjawabnya untuk satu titik awal. Pertanyaan yang lebih berguna: adakah vektor yang hanya
diperbesar oleh $M$ tanpa diubah arahnya? Di sini arah itu adalah
$(17/30,\;4/15,\;1/6)=(0.5667,\;0.2667,\;0.1667)$, tepat angka yang muncul setelah empat puluh
langkah, dan pengalinya $1$.

Bab 8 memberi determinan: $\det M=1/5$, jadi $M$ invertible dan setiap sistem $M\mathbf{y}=
\mathbf{b}$ punya solusi tunggal. Itu menjawab apakah proses dapat dibalik, bukan ke arah mana ia
menuju, dan bukan seberapa cepat. Angka yang mengatur laju di sini adalah $1/2$: nilai eigen
terbesar setelah $1$. Karena $(1/2)^{40}\approx9.1\times10^{-13}$, galat awal tinggal di belakang
enam angka pertama, dan itulah sebabnya pembacaan kedua terlihat diam.

Bab ini menyiapkan alat untuk membaca dua angka itu langsung dari matriksnya: cara menghitungnya,
syarat agar arah yang dipertahankan cukup banyak untuk menjadi basis, apa yang berubah ketika
basis diganti, lalu tiga pemakaiannya: pangkat matriks, sistem diferensial linear, dan rantai
Markov.

\begin{untukanak}
Ambil selembar karet yang diberi garis-garis kotak, lalu tarik dari dua sisi. Kotak-kotak berubah
menjadi belah ketupat, dan hampir semua garis menjadi miring. Tetapi ada satu garis yang tetap
lurus dan tetap sejajar dengan dirinya sendiri: garis yang searah tarikan. Panjangnya berubah,
arahnya tidak. Kadang ada juga garis lain yang tidak miring tetapi terbalik arahnya, seperti
panah yang diputar setengah lingkaran.

Kalau kita hanya diberi tahu berapa kali setiap garis istimewa itu mulur, kita sudah tahu seluruh
tarikan, tanpa perlu menggambar jaringannya lagi. Itu cara yang sangat hemat: tiga atau empat
angka menggantikan ribuan titik pada lembaran.

Masalahnya, tidak semua tarikan begitu rapi. Coba pelintir sedikit lembarnya sambil menarik.
Sekarang tidak ada garis yang benar-benar lurus lagi, dan cara hemat tadi gagal total. Ada
tarikan yang seluruhnya bisa dijelaskan oleh beberapa angka mulur, ada yang tidak bisa.

Pertanyaannya: dari gambar tarikan saja, bisakah kita menebak apakah ada garis istimewa yang
tidak miring, dan berapa banyak? Di bab ini kita akan tahu mengapa sebagian tarikan punya garis
seperti itu dan sebagian tidak, mengapa ada tarikan yang seluruhnya dapat diganti oleh sekadar
beberapa angka, dan mengapa mesin hitung kadang salah menebak pada tarikan yang hampir pelintir.
\end{untukanak}

\section{Arah yang tidak berputar}
\tingkat{dasar}

Sebuah matriks persegi bekerja pada vektor: ia memutar, meregangkan, atau memipihkan. Hampir
semua vektor keluar dengan arah yang berbeda dari arah masuknya. Sebagian kecil tidak. Vektor
yang keluar searah dengan masuknya, hanya berbeda panjang, menandai arah yang dipertahankan
matriks itu. Arah semacam itu yang kita beri nama lebih dulu, sebelum ada alat hitung apa pun.

\kamus{nilai eigen}{eigenvalue}{Pengali pada arah yang tidak diputar: arahnya tetap, panjangnya
dikali angka itu.}

\begin{definisi}[Nilai eigen dan vektor eigen]
Misalkan $A$ matriks $n\times n$ dengan entri real. Skalar $\lambda$ disebut \emph{nilai eigen}
matriks $A$ bila ada vektor tak nol $\mathbf{v}\in\mathbb{R}^{n}$ dengan
$A\mathbf{v}=\lambda\mathbf{v}$. Vektor $\mathbf{v}$ disebut \emph{vektor eigen} untuk
$\lambda$.\index{nilai eigen}\index{vektor eigen}
\end{definisi}

Syarat $\mathbf{v}\neq\mathbf{0}$ bukan hiasan. Tanpa syarat itu, $\mathbf{0}$ memenuhi
$A\mathbf{0}=\lambda\mathbf{0}$ untuk setiap $\lambda$, sehingga istilahnya kehilangan arti. Yang
kita cari justru arah, dan titik nol tidak punya arah.

Persamaan $A\mathbf{v}=\lambda\mathbf{v}$ dapat ditulis ulang sebagai
$(A-\lambda I)\mathbf{v}=\mathbf{0}$. Jadi $\lambda$ adalah nilai eigen tepat ketika matriks
$A-\lambda I$ mempunyai inti yang tidak trivial, yaitu tepat ketika $A-\lambda I$ singular.
Karena itu pertanyaan tentang nilai eigen selalu dapat diterjemahkan menjadi pertanyaan tentang
kekurangan pangkat, dan Bab 8 sudah menyediakan alat untuk itu.

Sebelum menghitung, satu peringatan penting. Tidak setiap matriks real punya nilai eigen real.
Matriks putar
\[
  R=\begin{pmatrix}0&-1\\ 1&0\end{pmatrix}
\]
mengirim $(1,0)$ ke $(0,1)$. Setiap vektor real berputar sembilan puluh derajat, sehingga tidak ada
garis yang dipertahankan; $R^{2}=-I$ memutar setengah lingkaran setelah dua langkah.
Ketidakadaan nilai eigen real di sini bukan kegagalan hitung, melainkan pernyataan geometris.

Untuk mencari nilai eigen tanpa menebak, kita pakai determinan: nilai eigen adalah bilangan yang
membuat matriks kehilangan inversnya setelah diagonalnya dikurangi bilangan itu.

\kamus{polinomial karakteristik}{characteristic polynomial}{Daftar bilangan yang membuat matriks
menjadi singular, ditulis sebagai polinomial dalam satu peubah.}

\begin{definisi}[Polinomial karakteristik]
Untuk matriks $n\times n$, \emph{polinomial karakteristik} didefinisikan sebagai
$p_{A}(\lambda)=\det(\lambda I-A)$.\index{polinomial karakteristik}
\end{definisi}

Kita memakai $\lambda I-A$, bukan $A-\lambda I$. Keduanya berbeda faktor $(-1)^{n}$, jadi akarnya
sama, tetapi pilihan pertama membuat polinomialnya monik, dan itu memudahkan pembacaan koefisien
lain seperti terlihat pada akibat berikutnya.

\begin{teorema}
Skalar $\lambda$ adalah nilai eigen $A$ jika dan hanya jika $p_{A}(\lambda)=0$.
\end{teorema}

\begin{proof}
Misalkan $A\mathbf{v}=\lambda\mathbf{v}$ dengan $\mathbf{v}\neq\mathbf{0}$. Maka
$(A-\lambda I)\mathbf{v}=\mathbf{0}$, sehingga $A-\lambda I$ mempunyai inti tidak trivial dan
karena itu singular. Menurut Bab 8, singularitas setara dengan determinan nol, jadi
$\det(A-\lambda I)=0$, yang berbeda faktor $(-1)^{n}$ dari $\det(\lambda I-A)=0$.

Sebaliknya, bila $\det(\lambda I-A)=0$ maka $\lambda I-A$ singular. Ruang intinya tidak trivial,
jadi ada $\mathbf{v}\neq\mathbf{0}$ dengan $(\lambda I-A)\mathbf{v}=\mathbf{0}$, yaitu
$A\mathbf{v}=\lambda\mathbf{v}$.
\end{proof}

Polinomial $p_{A}$ berderajat $n$ dan monik, jadi ia punya paling banyak $n$ akar berbeda. Di
atas $\mathbb{C}$ ia selalu terurai menjadi $n$ faktor linear, sehingga setiap matriks persegi
kompleks punya tepat $n$ nilai eigen bila dihitung dengan multiplisitas. Di atas $\mathbb{R}$ hal
itu bisa gagal: $p_{R}(\lambda)=\lambda^{2}+1$ tidak punya akar real, tetapi punya sepasang akar
kompleks $\pm i$.

Dua koefisien teratas polinomial karakteristik memuat dua besaran yang sudah kita kenal.

\begin{akibat}
Dengan nilai eigen dihitung beserta multiplisitasnya,
$\operatorname{tr}(A)=\sum_{i}\lambda_{i}$ dan $\det(A)=\prod_{i}\lambda_{i}$.
\index{jejak}
\end{akibat}

\begin{proof}
Koefisien $\lambda^{n-1}$ pada $\det(\lambda I-A)$ adalah $-\operatorname{tr}(A)$. Suku itu hanya
dihasilkan oleh permutasi identitas pada rumus Leibniz, karena setiap permutasi lain kehilangan
paling sedikit dua entri diagonal. Suku konstannya adalah $p_{A}(0)=\det(-A)=(-1)^{n}\det(A)$.
Membandingkan kedua koefisien itu dengan bentuk
$\prod_{i}(\lambda-\lambda_{i})=\lambda^{n}-(\sum_{i}\lambda_{i})\lambda^{n-1}+\dots+(-1)^{n}
\prod_{i}\lambda_{i}$ memberi kedua identitas.
\end{proof}

Akibat itu memberi pemeriksaan murah: setelah menghitung nilai eigen dengan susah payah, jumlah dan
hasil kalinya harus cocok dengan jejak dan determinan. Bila tidak cocok, ada kesalahan aritmetika.

\begin{bacakeras}
Dua suara, bergantian.

Suara pertama: apa yang dilakukan sebuah matriks?\\
Suara kedua: memutar, meregangkan, memipihkan.\\
Suara pertama: kepada siapa?\\
Suara kedua: kepada hampir semua vektor.\\
Suara pertama: lalu ada yang tidak berputar?\\
Suara kedua: ada. Satu arah. Kadang dua. Tidak pernah lebih dari banyaknya kolom.\\
Suara pertama: panjangnya berubah?\\
Suara kedua: berubah. Arahnya tetap.\\
Suara pertama: berapa kali?\\
Suara kedua: itu angkanya. Angka itu nilai eigen.\\
Suara pertama: dan arahnya?\\
Suara kedua: itu vektor eigen. Satu arah, satu angka.\\
Suara pertama: kalau arahnya berbalik?\\
Suara kedua: angkanya negatif. Garisnya tetap sama, hanya menghadap ke belakang.\\
Suara pertama: kalau angkanya nol?\\
Suara kedua: arah itu hilang, dan matriks kehilangan inversnya.

Bersama: arah yang tidak berubah menyimpan angka yang mengubah panjangnya.

Ulangi satu kali dengan suara biasa: arah yang tidak berubah menyimpan angka yang mengubah panjangnya.
\end{bacakeras}

\begin{resep}
Untuk $A=\begin{pmatrix}a&b\\c&d\end{pmatrix}$ yang entrinya sudah diketahui:
\begin{enumerate}
  \item Kurangi $\lambda$ dari kedua entri diagonal, sehingga tersisa $A-\lambda I$.
  \item Hitung $\det(A-\lambda I)$, atau langsung pakai $\lambda^{2}-(a+d)\lambda+(ad-bc)$.
  \item Selesaikan persamaan kuadrat itu; periksa jumlah akarnya $a+d$ dan hasil kalinya $ad-bc$.
  \item Untuk setiap akar $\lambda$, selesaikan $(A-\lambda I)\mathbf{v}=\mathbf{0}$ dan ambil satu
        solusi tak nol.
\end{enumerate}
Langkah 3 dan 4 tidak boleh ditukar: vektor eigen baru punya arti setelah $\lambda$ diketahui, dan
urutan kolom $P$ harus mengikuti urutan entri diagonal $D$.
\end{resep}

Sekarang satu contoh dengan tangan. Ambil
\[
  A=\begin{pmatrix}4&2\\ 1&3\end{pmatrix},\qquad
  p_{A}(\lambda)=\lambda^{2}-7\lambda+10=(\lambda-5)(\lambda-2).
\]
Jejaknya $7=5+2$ dan determinannya $10=5\cdot2$, cocok. Untuk $\lambda=5$ kita selesaikan
$(A-5I)\mathbf{v}=\mathbf{0}$, yaitu $-v_{1}+2v_{2}=0$ dan $v_{1}-2v_{2}=0$; kedua baris
mengatakan hal yang sama, jadi $\mathbf{v}=(2,1)$. Untuk $\lambda=2$ diperoleh
$2v_{1}+2v_{2}=0$, jadi $\mathbf{v}=(1,-1)$. Pemeriksaan langsung:
\[
  A\begin{pmatrix}2\\ 1\end{pmatrix}=\begin{pmatrix}10\\ 5\end{pmatrix}=5
  \begin{pmatrix}2\\ 1\end{pmatrix},\qquad
  A\begin{pmatrix}1\\ -1\end{pmatrix}=\begin{pmatrix}2\\ -2\end{pmatrix}=2
  \begin{pmatrix}1\\ -1\end{pmatrix}.
\]
Kedua pemeriksaan itu dapat dilihat pada bidang. Gambar~\ref{fig:panah-eigen} menggambar kedua
arah eigen beserta bayangannya, dan satu vektor yang bukan vektor eigen sebagai pembanding.

\begin{figure}[htbp]
\centering
\begin{tikzpicture}[x=0.4cm,y=0.4cm,>=latex]
  \draw[->,layergrey!60] (-0.8,0) -- (11.4,0) node[right,font=\scriptsize] {$x$};
  \draw[->,layergrey!60] (0,-2.9) -- (0,6.4) node[above,font=\scriptsize] {$y$};
  \draw[->,thick,layergrey] (0,0) -- (1,0);
  \draw[->,thick,dashed,layergrey] (0,0) -- (4,1)
        node[right,font=\scriptsize] {$(4,1)$};
  \draw[->,very thick,layerteal] (0,0) -- (2,1)
        node[above left,font=\scriptsize] {$(2,1)$};
  \fill[layerteal] (2,1) circle (2pt);
  \draw[->,very thick,accent] (0,0) -- (10,5)
        node[above,font=\scriptsize] {$(10,5)=5\,(2,1)$};
  \draw[->,very thick,layerindigo] (0,0) -- (1,-1)
        node[below left,font=\scriptsize] {$(1,-1)$};
  \fill[layerindigo] (1,-1) circle (2pt);
  \draw[->,very thick,layerindigo] (0,0) -- (2,-2)
        node[below right,font=\scriptsize] {$(2,-2)=2\,(1,-1)$};
\end{tikzpicture}
\caption{Arah eigen matriks $A=\begin{pmatrix}4&2\\1&3\end{pmatrix}$. Vektor $(2,1)$ keluar
sebagai $(10,5)$, yaitu $5$ kali lipat tanpa berputar; vektor $(1,-1)$ keluar sebagai $(2,-2)$,
yaitu $2$ kali lipat. Vektor $(1,0)$ keluar sebagai $(4,1)$ dan berputar, jadi ia bukan vektor
eigen.}
\label{fig:panah-eigen}
\end{figure}

Vektor eigen hanya ditentukan sampai pengali: $(4,2)$ juga vektor eigen untuk $5$. Yang tunggal
adalah arahnya, bukan panjangnya. Kebebasan itu akan kita pakai untuk merapikan matriks nanti.

Untuk matriks $2\times2$ akar itu dapat ditulis langsung. Karena
$p_{A}(\lambda)=\lambda^{2}-\operatorname{tr}(A)\lambda+\det(A)$, rumus akar memberi
\[
  \lambda_{1,2}=\frac{\operatorname{tr}(A)\pm\sqrt{\operatorname{tr}(A)^{2}-4\det(A)}}{2},
\]
sehingga tanda $\Delta=\operatorname{tr}(A)^{2}-4\det(A)$ menentukan seluruh spektrumnya. Untuk $A$
di atas, $\Delta=49-40=9$, dan akarnya $5$ serta $2$. Bila $\Delta=0$ kedua akar berimpit, dan bila
$\Delta<0$ tidak ada akar real: matriks putar tadi memberi $\Delta=0-4=-4$.

\begin{figure}[htbp]
\centering
\begin{tikzpicture}[x=0.85cm,y=0.85cm,>=latex,font=\scriptsize]
  \draw[->,layergrey!60] (-1.2,0.35) -- (6.6,0.35);
  \foreach \x in {0,1,2,3,4,5,6} \draw[layergrey!50] (\x,0.26) -- (\x,0.44);
  \fill[layerteal] (2,0.35) circle (2.2pt);
  \fill[layerteal] (5,0.35) circle (2.2pt);
  \node[above=2pt] at (2,0.35) {$2$};
  \node[above=2pt] at (5,0.35) {$5$};
  \node[left,align=right,text width=3.0cm] at (-1.3,0.35) {$\Delta=9$: dua akar real};

  \draw[->,layergrey!60] (-1.2,-1.25) -- (6.6,-1.25);
  \foreach \x in {0,1,2,3,4,5,6} \draw[layergrey!50] (\x,-1.34) -- (\x,-1.16);
  \fill[layerorange] (1,-1.25) circle (2.2pt);
  \node[above=2pt] at (1,-1.2) {$1$};
  \node[left,align=right,text width=3.0cm] at (-1.3,-1.25) {$\Delta=0$: satu akar ganda};

  \draw[->,layergrey!60] (-1.2,-2.85) -- (6.6,-2.85);
  \foreach \x in {0,1,2,3,4,5,6} \draw[layergrey!50] (\x,-2.94) -- (\x,-2.76);
  \draw[layergrey!80] (0,-3.55) -- (0,-2.15);
  \draw[accent,fill=white,thick] (0,-2.3) circle (2.4pt);
  \draw[accent,fill=white,thick] (0,-3.4) circle (2.4pt);
  \node[left=2pt,accent] at (0,-2.3) {$i$};
  \node[left=2pt,accent] at (0,-3.4) {$-i$};
  \node[left,align=right,text width=3.0cm] at (-1.3,-2.85) {$\Delta=-4$: tidak ada akar real};
\end{tikzpicture}
\caption{Spektrum matriks $2\times2$ pada garis bilangan, dipisahkan oleh tanda
$\Delta=\operatorname{tr}(A)^{2}-4\det(A)$. Baris atas: $A$ dengan $\Delta=49-40=9$ memberi akar
$2$ dan $5$. Baris tengah: matriks guntingan dengan $\Delta=0$ memberi akar $1$ yang dihitung dua
kali. Baris bawah: matriks putar dengan $\Delta=0-4=-4$ tidak memotong garis real, dan kedua
akarnya, $\pm i$, berada di luar garis itu.}
\label{fig:spektrum}
\end{figure}

Gambar~\ref{fig:spektrum} menyusun ketiga kemungkinan itu pada satu garis bilangan: dua titik,
satu titik yang dihitung dua kali, atau tidak ada titik sama sekali. Yang terakhir bukan kekurangan
gambar, melainkan pernyataan bahwa garis real tidak memuat arah yang dipertahankan.

\begin{tangan}
Bahan: satu karet gelang, dua pensil, selembar kertas, penggaris.

Letakkan kedua pensil sejajar, lalu bentangkan karet gelang memutari keduanya sehingga terbentuk
persegi panjang. Di kertas, gambar tiga panah dari sudut kiri bawah persegi itu, masing-masing pada
sudut 30, 45, dan 60 derajat terhadap sisi bawah. Beri nama ketiganya.

Geser pensil kanan dua kali lebih jauh, sehingga karet meregang. Bandingkan setiap panah dengan
panjang aslinya. Panah yang sejajar sisi tarikan tetap pada arahnya; dua panah lain miring, dan
yang paling miring adalah yang paling jauh dari sisi tarikan.

Sekarang putar seluruh rangka sembilan puluh derajat dan ulangi. Arah yang bertahan ikut berputar
bersama karet, karena yang istimewa adalah arahnya terhadap karet, bukan terhadap meja.

Yang harus terlihat: satu arah bertahan, dua arah berputar, dan angka regangan pada arah yang
bertahan itulah pengali yang kita sebut nilai eigen.
\end{tangan}

\section{Ruang eigen dan dua multiplisitas}
\tingkat{menengah}

Satu nilai eigen biasanya disertai lebih dari satu vektor eigen. Himpunan seluruh vektor eigen
untuk satu nilai eigen, ditambah vektor nol, bukan kumpulan acak: ia subruang. Dimensinya nanti
menentukan apakah matriksnya dapat disederhanakan.

\begin{definisi}[Ruang eigen]
Untuk nilai eigen $\lambda$ dari $A$, \emph{ruang eigen} adalah
$E_{\lambda}=\operatorname{null}(A-\lambda I)=\{\mathbf{v}:A\mathbf{v}=\lambda\mathbf{v}\}$.
\index{ruang eigen}
\end{definisi}

\kamus{ruang eigen}{eigenspace}{Seluruh arah yang tidak diputar oleh matriks, dikumpulkan menjadi
satu subruang.}

Ruang eigen selalu subruang karena ia inti sebuah matriks (Bab 7). Dimensinya paling sedikit satu
tepat ketika $\lambda$ benar-benar nilai eigen.

Ruang eigen untuk nilai eigen berbeda hanya beririsan di titik nol.

\begin{proposisi}
Bila $\lambda\neq\mu$, maka $E_{\lambda}\cap E_{\mu}=\{\mathbf{0}\}$.
\end{proposisi}

\begin{proof}
Ambil $\mathbf{v}$ di irisan itu. Maka $A\mathbf{v}=\lambda\mathbf{v}$ dan
$A\mathbf{v}=\mu\mathbf{v}$, sehingga $(\lambda-\mu)\mathbf{v}=\mathbf{0}$. Karena
$\lambda-\mu\neq0$, harus $\mathbf{v}=\mathbf{0}$.
\end{proof}

Sifat itu naik ke sejumlah berhingga vektor eigen, dengan induksi. Hasilnya lemma yang menjadi
tulang punggung bab ini.

\begin{lemma}
Vektor eigen yang bersesuaian dengan nilai eigen berbeda saling bebas linear.
\end{lemma}

\begin{proof}
Kita buktikan dengan induksi pada banyaknya vektor. Untuk satu vektor pernyataannya benar karena
vektor eigen tidak nol. Andaikan pernyataan berlaku untuk $m-1$ vektor, dan misalkan
\[
  c_{1}\mathbf{v}_{1}+\dots+c_{m}\mathbf{v}_{m}=\mathbf{0}
\]
dengan $\lambda_{1},\dots,\lambda_{m}$ berbeda. Kalikan persamaan itu dengan $A$, dan pakai
$A\mathbf{v}_{i}=\lambda_{i}\mathbf{v}_{i}$:
\[
  c_{1}\lambda_{1}\mathbf{v}_{1}+\dots+c_{m}\lambda_{m}\mathbf{v}_{m}=\mathbf{0}.
\]
Kurangi $\lambda_{m}$ kali persamaan pertama dari persamaan kedua:
\[
  c_{1}(\lambda_{1}-\lambda_{m})\mathbf{v}_{1}+\dots
  +c_{m-1}(\lambda_{m-1}-\lambda_{m})\mathbf{v}_{m-1}=\mathbf{0}.
\]
Menurut hipotesis induksi, $\mathbf{v}_{1},\dots,\mathbf{v}_{m-1}$ bebas, jadi
$c_{i}(\lambda_{i}-\lambda_{m})=0$ untuk $i<m$. Karena $\lambda_{i}\neq\lambda_{m}$, semua
$c_{i}=0$ untuk $i<m$. Persamaan semula tinggal $c_{m}\mathbf{v}_{m}=\mathbf{0}$, sehingga
$c_{m}=0$.
\end{proof}

Lemma itu memberi batas atas yang rapi: matriks $n\times n$ tidak mungkin punya lebih dari $n$
vektor eigen bebas. Kapan batas itu tercapai?

Untuk menjawabnya kita butuh dua cara menghitung banyaknya kemunculan satu nilai eigen: dari
polinomial dan dari dimensi ruang eigen. Keduanya sering berbeda, dan selisihnya yang menentukan.

\kamus{multiplisitas}{multiplicity}{Berapa kali sebuah nilai eigen muncul, dihitung dengan dua cara
yang tidak selalu sama.}

\begin{definisi}[Multiplisitas aljabar dan geometrik]
Misalkan $\lambda$ nilai eigen $A$. \emph{Multiplisitas aljabar} $\lambda$ adalah pangkat faktor
$(\lambda'-\lambda)$ pada $p_{A}(\lambda')$. \emph{Multiplisitas geometrik} $\lambda$ adalah
$\dim E_{\lambda}=\operatorname{nullity}(A-\lambda I)$.\index{multiplisitas}
\end{definisi}

Multiplisitas geometrik mudah dihitung: kurangi $\lambda$ dari diagonal, lalu hitung nullity
matriks hasilnya. Multiplisitas aljabar menuntut pemfaktoran polinomial. Hubungan keduanya selalu
satu arah.

\begin{lemma}
Untuk setiap nilai eigen $\lambda$ berlaku
$1\le\dim E_{\lambda}\le$ multiplisitas aljabar $\lambda$.
\end{lemma}

\begin{proof}
Pertidaksamaan kiri benar karena $\lambda$ nilai eigen, jadi $E_{\lambda}$ memuat vektor tak nol.

Untuk pertidaksamaan kanan, tulis $k=\dim E_{\lambda}$ dan pilih basis
$\mathbf{u}_{1},\dots,\mathbf{u}_{k}$ dari $E_{\lambda}$. Lengkapi menjadi basis
$\mathbf{u}_{1},\dots,\mathbf{u}_{n}$ dari $\mathbb{R}^{n}$, yang mungkin menurut Bab 7. Dalam
basis itu, kolom ke-$j$ matriks $A$ untuk $j\le k$ adalah $A\mathbf{u}_{j}=\lambda\mathbf{u}_{j}$,
yaitu $\lambda$ pada posisi diagonal dan nol di tempat lain. Jadi matriks $A$ dalam basis tersebut
berbentuk
\[
  \begin{pmatrix}\lambda I_{k}&C\\ 0&D\end{pmatrix}
\]
dengan $C$ berukuran $k\times(n-k)$ dan $D$ berukuran $(n-k)\times(n-k)$. Determinan
$\lambda' I-A$ dalam basis itu, lewat ekspansi kofaktor pada $k$ kolom pertama atau lewat rumus
determinan matriks blok segitiga, sama dengan $(\lambda'-\lambda)^{k}\det(\lambda' I-D)$. Jadi
$(\lambda'-\lambda)^{k}$ membagi $p_{A}$, sehingga multiplisitas aljabar $\lambda$ paling sedikit
$k$.
\end{proof}

Sekarang lihat kedua kemungkinan pada satu matriks. Ambil
\[
  A=\begin{pmatrix}2&1&1\\ 1&2&1\\ 1&1&2\end{pmatrix},\qquad
  p_{A}(\lambda)=(\lambda-1)^{2}(\lambda-4).
\]
Nilai eigennya $1$ dan $4$, dengan multiplisitas aljabar $2$ dan $1$. Jejaknya $6=4+1+1$ dan
determinannya $4=4\cdot1\cdot1$, cocok. Untuk $\lambda=4$, persamaan $(A-4I)\mathbf{v}=
\mathbf{0}$ memberi $v_{1}=v_{2}=v_{3}$, jadi $E_{4}$ direntang $(1,1,1)$. Untuk $\lambda=1$,
matriks $A-I$ adalah matriks yang seluruh entrinya $1$, sehingga intinya himpunan vektor yang
jumlah entrinya nol. Himpunan itu berdimensi dua, dengan basis $(1,-1,0)$ dan $(1,0,-1)$. Jadi
multiplisitas geometrik $1$ adalah $2$, sama dengan multiplisitas aljabarnya.

Nilai eigen berulang tidak otomatis menjadi masalah; yang menjadi masalah ruang eigen yang terlalu
kecil, dan itu terjadi pada matriks berikut.

\section{Kapan matriks dapat didiagonalkan}
\tingkat{menengah}

Kalau boleh memilih basis, kita ingin matriksnya sesederhana mungkin: diagonal, sehingga setiap
vektor basis hanya diskalakan dan tidak dicampur dengan yang lain. Apakah selalu ada basis semacam
itu?

\begin{definisi}[Dapat didiagonalkan]
Matriks $n\times n$ disebut \emph{dapat didiagonalkan} bila ada matriks $P$ yang invertible dan
matriks diagonal $D$ dengan $P^{-1}AP=D$, atau setara dengan itu $A=PDP^{-1}$.
\end{definisi}

\kamus{diagonalisasi}{diagonalization}{Mengganti alat ukur sampai seluruh pekerjaan matriks
tinggal beberapa perkalian angka.}

Kolom-kolom $P$ nanti akan berisi vektor eigen, dan $D$ berisi nilai eigen pada diagonalnya.
Karena itu pertanyaan tentang diagonalisasi sebenarnya pertanyaan tentang ketersediaan vektor
eigen yang cukup banyak.

\begin{teorema}
Matriks $n\times n$ dapat didiagonalkan jika dan hanya jika ia mempunyai $n$ vektor eigen yang
bebas linear.\index{diagonalisasi}
\end{teorema}

\begin{proof}
Misalkan $P^{-1}AP=D$ dengan $D=\operatorname{diag}(d_{1},\dots,d_{n})$. Kalikan dari kiri
dengan $P$: $AP=PD$. Kolom ke-$j$ kedua ruas memberi
$A\mathbf{p}_{j}=d_{j}\mathbf{p}_{j}$, dengan $\mathbf{p}_{j}$ kolom ke-$j$ dari $P$. Karena $P$
invertible, setiap kolomnya tak nol, jadi $\mathbf{p}_{j}$ vektor eigen. Kolom-kolom $P$ juga
bebas linear karena $P$ invertible. Jadi $A$ punya $n$ vektor eigen bebas.

Sebaliknya, misalkan $\mathbf{p}_{1},\dots,\mathbf{p}_{n}$ vektor eigen bebas dengan
$A\mathbf{p}_{j}=\lambda_{j}\mathbf{p}_{j}$. Susun $P$ dari vektor-vektor itu sebagai kolom dan
$D=\operatorname{diag}(\lambda_{1},\dots,\lambda_{n})$. Maka $AP=PD$, dan karena $P$ invertible
kita peroleh $P^{-1}AP=D$.
\end{proof}

Akibat yang paling sering dipakai dalam hitungan tangan: banyaknya nilai eigen berbeda sudah cukup.

\begin{akibat}
Bila matriks $n\times n$ mempunyai $n$ nilai eigen berbeda, ia dapat didiagonalkan.
\end{akibat}

\begin{proof}
Ambil satu vektor eigen untuk setiap nilai eigen. Menurut lemma sebelumnya, $n$ vektor eigen
dengan $n$ nilai eigen berbeda saling bebas. Teorema di atas lalu memberi diagonalisasi.
\end{proof}

Kebalikannya tidak berlaku: contoh $3\times3$ tadi hanya punya dua nilai eigen berbeda, tetapi
tetap dapat didiagonalkan karena ruang eigen untuk $\lambda=1$ berdimensi dua. Syarat yang tepat
bukan banyaknya nilai eigen, melainkan jumlah dimensi ruang eigennya.

\begin{teorema}[Syarat multiplisitas]
Matriks $n\times n$ dapat didiagonalkan jika dan hanya jika untuk setiap nilai eigennya,
multiplisitas aljabar sama dengan multiplisitas geometrik. Syarat itu setara dengan
$\sum_{\lambda}\dim E_{\lambda}=n$.
\end{teorema}

\begin{proof}
Pertama, anggap jumlah dimensi seluruh ruang eigen sama dengan $n$. Untuk setiap $\lambda$ pilih
basis $E_{\lambda}$, lalu gabungkan seluruh basis itu menjadi daftar $\mathbf{v}_{1},\dots,
\mathbf{v}_{n}$. Daftar itu bebas linear: bila suatu kombinasi linear bernilai nol, kelompokkan
suku menurut nilai eigennya menjadi $\mathbf{u}_{1}+\dots+\mathbf{u}_{m}=\mathbf{0}$ dengan
$\mathbf{u}_{k}\in E_{\lambda_{k}}$. Menurut lemma tentang vektor eigen berbeda, setiap
$\mathbf{u}_{k}=\mathbf{0}$; karena setiap $\mathbf{u}_{k}$ kombinasi basis $E_{\lambda_{k}}$,
semua koefisiennya nol. Daftar itu juga berisi $n$ vektor di ruang berdimensi $n$, jadi ia basis
vektor eigen, dan teorema sebelumnya menutup jalan.

Kedua, anggap $A$ dapat didiagonalkan, jadi ada basis vektor eigen
$\mathbf{p}_{1},\dots,\mathbf{p}_{n}$. Dalam basis itu matriks $A$ diagonal, sehingga
$p_{A}(\lambda')=\prod_{j}(\lambda'-\lambda_{j})$ dan multiplisitas aljabar setiap nilai eigen
sama dengan banyaknya $j$ dengan $\lambda_{j}=\lambda$. Banyaknya $j$ itu tidak melebihi
$\dim E_{\lambda}$, karena vektor-vektor tersebut berada di dalam $E_{\lambda}$ dan bebas. Jadi
multiplisitas aljabar tidak melebihi geometrik. Lemma sebelumnya memberi arah sebaliknya, sehingga
keduanya sama. Terakhir, $\sum_{\lambda}\dim E_{\lambda}=n$ karena basis vektor eigen mengisi
seluruh ruang.
\end{proof}

Sekarang contoh yang gagal. Ambil
\[
  S=\begin{pmatrix}1&1\\ 0&1\end{pmatrix},\qquad p_{S}(\lambda)=(\lambda-1)^{2}.
\]
Satu-satunya nilai eigen adalah $1$, dengan multiplisitas aljabar $2$. Tetapi
$S-I=\begin{pmatrix}0&1\\ 0&0\end{pmatrix}$ berpangkat satu, jadi
$\dim E_{1}=2-1=1$. Multiplisitas geometriknya kurang satu, dan menurut teorema di atas $S$ tidak
dapat didiagonalkan.

Secara geometris $S$ adalah guntingan: ia mempertahankan arah $(1,0)$ dan menggeser $(0,1)$ ke
$(1,1)$. Tidak ada basis yang membuat peta ini sekadar penskalaan, karena penskalaan tidak pernah
menggeser sesuatu yang tadinya sejajar. Kegagalan itu sifat petanya, bukan pilihan basis.

\begin{untukanak}
Pernahkah kamu mencoba merapikan sesuatu yang ternyata tidak bisa dirapikan? Ambil satu tumpukan
kartu, lalu geser bagian atasnya sedikit ke samping. Gambar pada setiap kartu tetap sama, tetapi
setiap titik bergeser sejauh jaraknya dari alas. Tidak ada cara memutar tumpukan itu supaya
geseran tersebut berubah menjadi sekadar memanjangkan atau memendekkan sesuatu. Geseran tetap
geseran.

Ada benda lain yang lebih ramah. Karet yang ditarik lurus dapat dijelaskan hanya dengan berapa
kali setiap arah mulur. Kalau kita tahu angka itu untuk beberapa arah penting, kita tahu seluruh
tarikan, termasuk untuk titik yang tidak pernah kita ukur.

Pertanyaannya: bagaimana membedakan keduanya hanya dari gambarnya, tanpa mencoba semua
kemungkinan? Di bab ini kita akan tahu mengapa geseran selalu menyimpan satu arah yang tidak bisa
dirapikan, mengapa ada tarikan yang punya arah istimewa cukup banyak, dan mengapa beda itu
menentukan apakah pangkat kesepuluh sebuah mesin hitung masih dapat dipercaya.
\end{untukanak}

Guntingan mudah dikenali karena segitiga. Kegagalan yang sama muncul pada matriks yang tidak tampak
segitiga. Ambil
\[
  B=\begin{pmatrix}-1&4&0\\ -1&3&0\\ 1&-2&2\end{pmatrix},\qquad
  p_{B}(\lambda)=(\lambda-1)^{2}(\lambda-2).
\]
Untuk $\lambda=1$, matriks $B-I$ berpangkat dua, jadi $\dim E_{1}=3-2=1$, padahal
multiplisitas aljabarnya $2$. Satu-satunya vektor eigen untuk $1$ adalah $(2,1,0)$. Untuk
$\lambda=2$ ruang eigennya direntang $(0,0,1)$ dan berdimensi satu, sesuai multiplisitas
aljabarnya. Jumlah dimensi seluruh ruang eigen hanya $2<3$, jadi $B$ tidak dapat didiagonalkan.
Matriks semacam ini disebut \emph{defective}: ia kekurangan vektor eigen karena ruang eigennya
sempit.\index{defective}

\kamus{defective}{defective}{Matriks yang nilai eigennya berulang tetapi ruang eigennya terlalu
sempit untuk menjadi basis.}

Satu catatan sebelum kita lanjut. Dapat atau tidaknya didiagonalkan bergantung pada medan
bilangannya: $R$ tidak dapat didiagonalkan atas $\mathbb{R}$, tetapi atas $\mathbb{C}$ ia punya dua
nilai eigen berbeda, $\pm i$, sehingga dapat didiagonalkan. Kalimat tanpa keterangan medan karena
itu tidak lengkap. Di sisa bab ini medan yang dipakai adalah $\mathbb{R}$ kecuali disebut lain.

\section{Matriks serupa dan invarian}
\tingkat{lanjut}

Diagonalisasi adalah perubahan basis yang menghasilkan bentuk paling sederhana. Tetapi tidak
setiap matriks dapat mencapai bentuk itu, dan pertanyaan berikutnya wajar: apa yang tetap sama
pada semua matriks yang dihubungkan oleh perubahan basis?\index{invarian}

\begin{definisi}[Matriks serupa]
Dua matriks $n\times n$, $A$ dan $B$, disebut \emph{serupa} bila ada matriks invertible $P$
dengan $B=P^{-1}AP$.\index{serupa}
\end{definisi}

\kamus{serupa}{similar}{Dua laporan tentang peta yang sama dari dua ruang ukur yang berbeda.}

Keserupaan adalah relasi ekuivalen: setiap matriks serupa dengan dirinya sendiri, relasinya
simetris karena $(P^{-1})^{-1}BP^{-1}=A$, dan transitif karena
$(PQ)^{-1}A(PQ)=Q^{-1}(P^{-1}AP)Q$. Kelas-kelas itu terbagi rapi, dan diagonalisasi berarti
menanyakan apakah suatu kelas memuat matriks diagonal.

\begin{proposisi}
Bila $A$ dan $B$ serupa, maka $p_{A}=p_{B}$. Akibatnya $A$ dan $B$ mempunyai nilai eigen yang
sama beserta multiplisitas aljabarnya, dan juga jejak, determinan, pangkat, serta nullity yang
sama.
\end{proposisi}

\begin{proof}
Misalkan $B=P^{-1}AP$. Maka
\[
  \lambda I-B=\lambda I-P^{-1}AP=P^{-1}(\lambda I-A)P,
\]
sehingga
$\det(\lambda I-B)=\det(P^{-1})\det(\lambda I-A)\det(P)=\det(\lambda I-A)$ karena
$\det(P^{-1})\det(P)=1$. Jadi polinomial karakteristiknya sama, dan semua data yang diturunkan
dari koefisiennya ikut sama. Kesamaan pangkat dan nullity diperoleh dari $B-\lambda I=
P^{-1}(A-\lambda I)P$: perkalian dengan matriks invertible tidak mengubah pangkat, dan peringkat
itulah yang menentukan nullity lewat rank-nullity.
\end{proof}

Kebalikannya tidak berlaku, dan contoh penyangkalnya sederhana: $I_{2}$ dan guntingan $S$ punya
polinomial karakteristik sama, $(\lambda-1)^{2}$, tetapi tidak serupa. Bila $P^{-1}I_{2}P=S$, ruas
kiri sama dengan $I_{2}$ untuk setiap $P$. Jadi polinomial karakteristik bukan invarian lengkap; ia
hanya menyaring.

Yang berubah dan yang tidak berubah dapat dipisahkan dengan rapi. Nilai eigen tidak berubah.
Vektor eigennya berubah, dan berubah dengan cara yang mudah diingat.

\begin{proposisi}
Bila $B=P^{-1}AP$ dan $A\mathbf{v}=\lambda\mathbf{v}$, maka
$B(P^{-1}\mathbf{v})=\lambda(P^{-1}\mathbf{v})$.
\end{proposisi}

\begin{proof}
Hitung $B(P^{-1}\mathbf{v})=P^{-1}APP^{-1}\mathbf{v}=P^{-1}A\mathbf{v}
=P^{-1}(\lambda\mathbf{v})=\lambda(P^{-1}\mathbf{v})$.
\end{proof}

Jadi vektor eigen ikut berpindah bersama koordinatnya, sedangkan pengalinya tinggal. Itu masuk
akal: arah yang dipertahankan adalah sifat peta, sedangkan daftar koordinatnya bergantung pada
basis. Karena itu nilai eigen dipakai sebagai sidik jari matriks, sedangkan vektor eigen tidak.

Di sisi lain, kebebasan memilih $P$ menjelaskan mengapa dua diagonalisasi dapat tampak berbeda
padahal keduanya benar: kolom $P$ boleh dikalikan skalar bukan nol atau dipermutasi asalkan $D$ ikut
berpermutasi. Yang harus sama adalah himpunan nilai eigen beserta multiplisitasnya.

\section{Contoh hitung}
\tingkat{menengah}

Semua angka di bagian ini dihitung dengan tangan lebih dulu, lalu diperiksa dengan Python pada
bagian lab. Kita mulai dari matriks $3\times3$ yang sudah muncul pada pembahasan ruang eigen,
\[
  A=\begin{pmatrix}2&1&1\\ 1&2&1\\ 1&1&2\end{pmatrix}.
\]
Polinomial karakteristiknya dihitung dari $\lambda I-A$, yang seluruh entri diagonalnya
$\lambda-2$ dan seluruh entri lainnya $-1$. Dengan menulis $t=\lambda-2$ dan memakai ekspansi
sepanjang baris pertama,
\[
  p_{A}(\lambda)=t\,(t^{2}-1)-(-1)\,(-t-1)+(-1)\,(1+t)=t^{3}-3t-2 .
\]
Bentuk itu terurai karena $t=-1$ akarnya: $t^{3}-3t-2=(t+1)^{2}(t-2)$. Kembali ke $\lambda$,
\[
  p_{A}(\lambda)=(\lambda-1)^{2}(\lambda-4)=\lambda^{3}-6\lambda^{2}+9\lambda-4 .
\]
Jejaknya $6$ dan determinannya $4$, cocok dengan $4+1+1$ dan $4\cdot1\cdot1$.

Vektor eigen untuk $\lambda=4$: matriks $A-4I$ mempunyai entri diagonal $-2$ dan entri lain $1$.
Baris-barisnya mengatakan $-2v_{1}+v_{2}+v_{3}=0$ dan seterusnya; ketiga persamaan itu terpenuhi
tepat ketika $v_{1}=v_{2}=v_{3}$. Jadi $E_{4}$ direntang $(1,1,1)$.

Vektor eigen untuk $\lambda=1$: matriks $A-I$ seluruhnya berisi $1$, jadi intinya adalah himpunan
vektor dengan $v_{1}+v_{2}+v_{3}=0$. Basisnya dua vektor, misalnya $(1,-1,0)$ dan $(1,0,-1)$.
Dengan demikian
\[
  P=\begin{pmatrix}1&1&1\\ 1&-1&0\\ 1&0&-1\end{pmatrix},\qquad
  P^{-1}=\frac{1}{3}\begin{pmatrix}1&1&1\\ 1&-2&1\\ 1&1&-2\end{pmatrix},\qquad
  D=\operatorname{diag}(4,1,1).
\]
Determinan $P$ adalah $3$, dan itulah penyebut pada $P^{-1}$. Satu entri kita periksa dengan
tangan: entri baris pertama kolom pertama $PDP^{-1}$ adalah
$\tfrac{1}{3}(4\cdot1\cdot1+1\cdot1\cdot1+1\cdot1\cdot1)=2$, sesuai entri $A$. Pemeriksaan
menyeluruh dengan SymPy memberi $PDP^{-1}-A=0$ persis.

Sekarang pangkat kesepuluh.\index{pangkat matriks} Mengalikan sembilan kali tidak memberi gambaran
apa pun; diagonalisasi memberinya. Karena $A=I+J$ dengan $J$ matriks yang seluruh entrinya $1$, dan
$J^{2}=3J$, setiap pangkat $J$ kembali menjadi kelipatan $J$:
\[
  J^{2}=3J,\quad J^{3}=9J,\quad\dots,\quad J^{k}=3^{k-1}J .
\]
Maka
\[
  A^{k}=(I+J)^{k}=I+\Bigl(\sum_{j=1}^{k}\binom{k}{j}3^{j-1}\Bigr)J
  =I+\frac{4^{k}-1}{3}\,J,
\]
dengan langkah terakhir memakai rumus deret geometri
$\sum_{j=1}^{k}\binom{k}{j}3^{j-1}=\bigl((1+3)^{k}-1\bigr)/3$. Untuk $k=10$ kita punya
$4^{10}=1048576$ dan $(1048576-1)/3=349525$, sehingga
\[
  A^{10}=I+349525\,J=
  \begin{pmatrix}349526&349525&349525\\ 349525&349526&349525\\ 349525&349525&349526\end{pmatrix}.
\]
Setiap entrinya dapat diperiksa: diagonal $1+349525$, luar diagonal $349525$. Perkalian langsung
dengan Python memberi matriks yang sama, dengan selisih terbesar $1.05\times10^{-9}$ karena
aritmetika floating-point.

Bentuk rumus itu menjelaskan apa yang tidak terlihat pada hasil akhir. Suku $4^{k}$ menguasai arah
$(1,1,1)$; suku konstan $1$ bertahan pada arah yang jumlah entrinya nol. Nilai eigen $4$ dan $1$
adalah dua laju pertumbuhan, dan pangkat matriks hanyalah campuran keduanya. Bila matriks tidak
dapat didiagonalkan, campuran itu berubah bentuk: muncul $k\lambda^{k}$ atau, untuk $\lambda=1$,
faktor $k$ murni.

Contoh defektif tadi memperlihatkan hal itu. Untuk $B$ yang gagal didiagonalkan tadi,
\[
  B^{k}=
  \begin{pmatrix}1-2k&4k&0\\ -k&1+2k&0\\ 2^{k}-1&-2(2^{k}-1)&2^{k}\end{pmatrix},
\]
yang dapat dibuktikan dengan induksi dan cocok untuk $k=1$. Pada $k=10$ rumus itu memberi
\[
  B^{10}=\begin{pmatrix}-19&40&0\\ -10&21&0\\ 1023&-2046&1024\end{pmatrix},
\]
sama dengan hasil perkalian berulang. Entri kiri atas, $-19=1-2\cdot10$, tumbuh linear, bukan
eksponensial: nilai eigen $1$ mengendalikan arah itu, dan kekurangan vektor eigen menambahkan satu
faktor $k$.

Sekarang penerapan pertama: proses yang menetap. Untuk matriks rantai Markov $M$ di pembuka,
\index{rantai Markov}
\[
  p_{M}(\lambda)=\lambda^{3}-\tfrac{19}{10}\lambda^{2}+\tfrac{11}{10}\lambda-\tfrac15
  =(\lambda-1)\bigl(\lambda-\tfrac12\bigr)\bigl(\lambda-\tfrac25\bigr).
\]
Jejaknya $19/10$ dan determinannya $1/5$, cocok dengan $1+\tfrac12+\tfrac25$ dan
$1\cdot\tfrac12\cdot\tfrac25$. Vektor eigen untuk $\lambda=1$ dicari dari $M\mathbf{v}=\mathbf{v}$,
yaitu $(M-I)\mathbf{v}=\mathbf{0}$, ditambah syarat jumlah entrinya satu. Penyelesaiannya memberi
$\boldsymbol{\pi}=(17/30,\;4/15,\;1/6)$. Pemeriksaan satu koordinat dengan tangan:
\[
  \tfrac{8}{10}\cdot\tfrac{17}{30}+\tfrac{3}{10}\cdot\tfrac{8}{30}
  +\tfrac{2}{10}\cdot\tfrac{5}{30}
  =\tfrac{13.6+2.4+1.0}{30}=\tfrac{17}{30}.
\]
Jadi $\boldsymbol{\pi}$ benar-benar titik tetap.

Karena $M$ punya tiga nilai eigen berbeda, ia dapat didiagonalkan. Setiap titik awal karena itu
terurai sebagai
\[
  \mathbf{x}_{k}=\boldsymbol{\pi}+c_{2}\bigl(\tfrac12\bigr)^{k}\mathbf{v}_{2}
  +c_{3}\bigl(\tfrac25\bigr)^{k}\mathbf{v}_{3},
\]
dengan $c_{2},c_{3}$ bergantung pada $\mathbf{x}_{0}$. Suku kedua dan ketiga mengecil, dan yang
menentukan laju adalah $1/2$. Untuk $\mathbf{x}_{0}=(1,0,0)$ kita peroleh $c_{2}\mathbf{v}_{2}$
berorde $10^{-1}$, sehingga galat setelah sepuluh langkah berorde
$10^{-1}\cdot(1/2)^{10}\approx10^{-4}$, dan setelah empat puluh langkah berorde
$(1/2)^{40}\approx9.1\times10^{-13}$. Hasil hitung sesuai: $\mathbf{x}_{10}$ berbeda dari
$\boldsymbol{\pi}$ pada angka keempat, sedangkan $\mathbf{x}_{40}$ sudah sama sampai enam angka.

Penerapan kedua: sistem diferensial. Bila $\mathbf{x}(t)$ memenuhi
$\mathbf{x}'=A\mathbf{x}$ dengan $A$ di atas, solusinya $\mathbf{x}(t)=e^{At}\mathbf{x}(0)$.
Karena $A=I+J$ dan $J^{2}=3J$,
\[
  e^{At}=e^{t}e^{Jt}=e^{t}\Bigl(I+\frac{e^{3t}-1}{3}J\Bigr)
  =e^{t}I+\frac{e^{4t}-e^{t}}{3}J .
\]
Pada $t=0$ rumus itu memberi $I$. Untuk $t=0.5$ entri luar diagonalnya $(e^{2}-e^{0.5})/3=1.9134$,
dan Python memberi angka yang sama sampai empat angka. Komponen sepanjang $(1,1,1)$ tumbuh seperti
$e^{4t}$; komponen pada bidang jumlah nol tumbuh seperti $e^{t}$, sehingga solusi cepat sejajar
dengan $(1,1,1)$. Nilai eigen real positif berarti pertumbuhan; nilai eigen imajiner murni seperti
$\pm i$ berarti osilasi tanpa pertumbuhan.

Penerapan ketiga: iterasi pangkat. Mengalikan berulang lalu menormalkan hasilnya,
$\mathbf{x}_{k+1}=A\mathbf{x}_{k}/\lVert A\mathbf{x}_{k}\rVert$, membuat arah berputar menuju vektor
eigen dominan, karena $\lambda_{1}^{k}$ mengalahkan $\lambda_{2}^{k}$ bila
$|\lambda_{1}|>|\lambda_{2}|$. Untuk $A$ di atas mulai dari $(1,0)$, sudutnya berjalan $0^\circ$,
$14.04^\circ$, $21.25^\circ$, $24.39^\circ$, $25.69^\circ$, $26.21^\circ$, dan menuju
$\arctan(1/2)=26.57^\circ$; nisbah Rayleigh $\mathbf{x}^{\top}A\mathbf{x}/\mathbf{x}^{\top}
\mathbf{x}$ berjalan $4.00$, $4.65$, $4.88$, $4.96$, $4.98$, $4.99$, menuju $5$.
Gambar~\ref{fig:iterasi-pangkat} memperlihatkan putaran itu.

\begin{figure}[htbp]
\centering
\begin{tikzpicture}[scale=3.2,>=latex]
  \draw[layergrey!35] (0,0) circle (1);
  \draw[->,layergrey!60] (-1.25,0) -- (1.55,0) node[right,font=\scriptsize] {$x$};
  \draw[->,layergrey!60] (0,-0.35) -- (0,1.2) node[above,font=\scriptsize] {$y$};
  \draw[->,semithick,layergrey] (0,0) -- (0:1);
  \draw[->,thick,layerindigo!45] (0,0) -- (14.04:1);
  \draw[->,thick,layerindigo!70] (0,0) -- (21.25:1);
  \draw[->,very thick,layerindigo] (0,0) -- (24.39:1);
  \foreach \a in {0,14.04,21.25,24.39} \fill[layerindigo!70] (\a:1) circle (1.3pt);
  \draw[->,very thick,dashed,accent] (0,0) -- (26.565:1.42);
  \node[right,font=\scriptsize,accent,align=left] at (26.565:1.42)
        {$(2,1)$\\[1pt]$26.57^\circ$};
\end{tikzpicture}
\caption{Iterasi pangkat pada $A=\begin{pmatrix}4&2\\1&3\end{pmatrix}$ mulai dari $(1,0)$.
Panah penuh adalah arah setelah $k=0,1,2,3$ langkah, dengan sudut $0^\circ$, $14.04^\circ$,
$21.25^\circ$, dan $24.39^\circ$. Panah putus-putus adalah arah eigen dominan $(2,1)$ pada
$26.57^\circ$, yang didekati tanpa pernah dilewati; dua panah terakhir berimpit karena sudutnya
makin rapat, dan itu justru inti gambarnya.}
\label{fig:iterasi-pangkat}
\end{figure}

\kertas{Untuk $J$ matriks yang seluruh entrinya $1$, ambil $B=J-I$. Tentukan nilai eigen $B$
lebih dulu, lalu ramalkan $B^{10}$ memakai pola $I+\frac{\mu^{10}-(-1)^{10}}{3}J$ dengan $\mu$
nilai eigen bukan nol dari $B$. Periksa ramalan itu dengan Python, lalu jelaskan mengapa $B$ dapat
didiagonalkan meskipun punya nilai eigen berulang. Sebagai latihan tambahan, hitung $e^{Bt}$ dan
tentukan limit $t\to\infty$ dari $e^{-2t}e^{Bt}$.}

\section{Lab komputasi}
\tingkat{lanjut}

\begin{lab}
\begin{verbatim}
import numpy as np
from scipy.linalg import expm

np.set_printoptions(precision=4, suppress=True)

A = np.array([[2., 1., 1.],
              [1., 2., 1.],
              [1., 1., 2.]])
lam, V = np.linalg.eig(A)
print("nilai eigen A :", np.round(np.sort(lam.real), 10))
for i in range(3):
    sisa = A @ V[:, i] - lam[i] * V[:, i]
    print("  sisa Av-lv  :", round(float(np.linalg.norm(sisa)), 16))
print("A^10 lewat diagonalisasi :")
print(np.round((V @ np.diag(lam ** 10) @ np.linalg.inv(V)).real).astype(int))
print("A^10 lewat pangkat biasa :")
print(np.linalg.matrix_power(A, 10).astype(int))
print("selisih maks :", np.abs((V @ np.diag(lam ** 10)
                              @ np.linalg.inv(V)).real
                             - np.linalg.matrix_power(A, 10)).max())
print("cond V       :", round(float(np.linalg.cond(V)), 4))

A2 = np.array([[4., 2.],
               [1., 3.]])
x = np.array([1., 0.])
print("iterasi pangkat pada A2:")
for k in range(6):
    sudut = np.degrees(np.arctan2(x[1], x[0]))
    rq = float(x @ A2 @ x / (x @ x))
    print("  k=%d sudut=%7.2f rayleigh=%.4f" % (k, sudut, rq))
    x = A2 @ x / np.linalg.norm(A2 @ x)
print("sudut batas  :", round(float(np.degrees(np.arctan2(1., 2.))), 2))

B = np.array([[-1., 4., 0.],
              [-1., 3., 0.],
              [1., -2., 2.]])
print("eigen B      :", np.round(np.sort(np.linalg.eigvals(B).real), 10))
print("rank(B - I)  :", np.linalg.matrix_rank(B - np.eye(3)))
print("B^10         :")
print(np.linalg.matrix_power(B, 10).astype(int))

M = np.array([[0.8, 0.3, 0.2],
              [0.1, 0.6, 0.3],
              [0.1, 0.1, 0.5]])
print("jumlah kolom :", M.sum(axis=0))
print("eigen M      :", np.round(np.sort(np.linalg.eigvals(M).real), 10))
x = np.array([1., 0., 0.])
for _ in range(10):
    x = M @ x
print("x setelah 10 :", np.round(x, 6))
for _ in range(30):
    x = M @ x
print("x setelah 40 :", np.round(x, 6))

C = np.array([[1., 1000.], [0., 2.]])
Cp = C.copy()
Cp[1, 0] = 1e-6
print("eigen C      :", np.round(np.sort(np.linalg.eigvals(C).real), 6))
print("eigen C+1e-6 :", np.round(np.sort(np.linalg.eigvals(Cp).real), 6))
print("cond V(C)    :", round(float(np.linalg.cond(np.linalg.eig(C)[1])), 1))
print("expm(0.5A)   :", round(float(expm(0.5 * A)[0, 1]), 4))
print("rumus        :", round(float((np.exp(2) - np.exp(0.5)) / 3), 4))
\end{verbatim}
Keluaran sebenarnya:
\begin{verbatim}
nilai eigen A : [1. 1. 4.]
  sisa Av-lv  : 8e-16
  sisa Av-lv  : 1.1e-15
  sisa Av-lv  : 3e-16
A^10 lewat diagonalisasi :
[[349526 349525 349525]
 [349525 349526 349525]
 [349525 349525 349526]]
A^10 lewat pangkat biasa :
[[349526 349525 349525]
 [349525 349526 349525]
 [349525 349525 349526]]
selisih maks : 1.0477378964424133e-09
cond V       : 1.5227
iterasi pangkat pada A2:
  k=0 sudut=   0.00 rayleigh=4.0000
  k=1 sudut=  14.04 rayleigh=4.6471
  k=2 sudut=  21.25 rayleigh=4.8820
  k=3 sudut=  24.39 rayleigh=4.9578
  k=4 sudut=  25.69 rayleigh=4.9840
  k=5 sudut=  26.21 rayleigh=4.9938
sudut batas  : 26.57
eigen B      : [1. 1. 2.]
rank(B - I)  : 2
B^10         :
[[  -19    40     0]
 [  -10    21     0]
 [ 1023 -2046  1024]]
jumlah kolom : [1. 1. 1.]
eigen M      : [0.4 0.5 1. ]
x setelah 10 : [0.5672 0.2661 0.1666]
x setelah 40 : [0.5667 0.2667 0.1667]
eigen C      : [1. 2.]
eigen C+1e-6 : [0.999 2.001]
cond V(C)    : 2000.0
expm(0.5A)   : 1.9134
rumus        : 1.9134
\end{verbatim}
\end{lab}

Baris pertama keluaran itu menyebut nilai eigen dalam urutan menaik, $1,1,4$. Tiga baris
berikutnya adalah sisa $A\mathbf{v}-\lambda\mathbf{v}$ untuk ketiga vektor eigen yang dikembalikan
pustaka. Sisa berorde $10^{-16}$ berarti persamaannya terpenuhi sampai batas presisi mesin; tidak ada sisa
yang persis nol karena vektor eigen dinormalkan dan dihitung secara iteratif. Vektor eigen pustaka
juga bukan $(1,1,1)$ dan $(1,-1,0)$, melainkan kelipatan ternormalkan: arahnya sama, panjangnya
berbeda, dan diagonalisasi hanya peduli arah.

Dua blok $A^{10}$ memberi angka bulat yang identik, dengan selisih terbesar $1.05\times10^{-9}$.
Selisih itu bukan kesalahan aljabar: entri sebesar $3.5\times10^{5}$ dihitung lewat $V$, $V^{-1}$,
dan pangkat bilangan real, sehingga galat pembulatan berorde $10^{-16}$ ikut membesar.

Nilai \texttt{cond V} sebesar $1.5227$ dekat dengan satu, dan itu kabar baik: $A$ simetris,
sehingga vektor eigennya dapat dipilih saling tegak lurus dan $P$ mendekati ortogonal.
Diagonalisasi untuk matriks semacam itu jauh lebih stabil.

Blok iterasi pangkat mengulang hitungan pada Gambar~\ref{fig:iterasi-pangkat}: sudutnya naik dari
$0$ ke $26.21$ derajat, tetap di bawah batas $26.57$, sedangkan nisbah Rayleigh naik dari $4.0000$
ke $4.9938$ menuju $5$.

Blok $B$ menunjukkan sisi lain: nilai eigennya $1,1,2$ dan \texttt{rank(B - I)} sama dengan $2$,
jadi nullity-nya $1$. Multiplisitas geometrik untuk $\lambda=1$ hanya satu, sedangkan aljabarnya
dua; matriksnya defektif. Blok $B^{10}$ memperlihatkan akibatnya: entri $-19$ dan $1024$
berdampingan, yang pertama tumbuh linear dan yang kedua eksponensial.

Blok $M$ memeriksa bahwa jumlah setiap kolomnya satu, sehingga $\mathbf{x}_{k}$ tetap stokastik.
Nilai eigennya $0.4$, $0.5$, dan $1$; setelah empat puluh langkah angkanya berhenti bergerak pada
enam angka, sesuai laju $(1/2)^{k}$ yang diturunkan dengan tangan.

Blok terakhir memperingatkan satu hal: matriks $C$ punya nilai eigen $1$ dan $2$, dan bila satu entri
nol di bawah diagonal diganti $10^{-6}$, nilai eigennya bergeser menjadi $0.999$ dan $2.001$;
perpindahan $10^{-3}$, seribu kali lebih besar daripada gangguannya. Angka itu dapat dihitung
persis: akar $\lambda^{2}-3\lambda+1.999$ adalah $1.5\mp\sqrt{2510}/100$, yaitu $0.999000998$ dan
$2.000999002$. Pergeseran itu sebanding dengan $10^{-6}\cdot1000/(2-1)$: gangguan dikali entri besar
dibagi celah kedua nilai eigen. Bilangan kondisi $2000$ menangkap pembesaran yang sama. Nilai eigen
tetap kontinu terhadap entri, tetapi tidak harus kontinu secara Lipschitz dengan konstanta satu.

\begin{tantangan}
Carilah satu matriks $2\times2$ yang kedua nilai eigennya real dan berbeda, tetapi kedua vektor
eigennya tidak saling tegak lurus. Tunjukkan pilihan kita benar dengan menghitung hasil kali
titiknya, lalu jelaskan mengapa matriks simetris tidak dapat dipakai (jawabannya menunggu Bab 10).

Petunjuk 1: dua vektor tegak lurus tepat ketika hasil kali titiknya nol. Mulailah dari matriks
segitiga atas yang sederhana, lalu hitung kedua vektor eigennya tanpa membulatkan.

Petunjuk 2: matriks $\begin{pmatrix}2&1\\0&3\end{pmatrix}$ sudah cukup. Vektor eigen untuk $3$
adalah $(1,1)$ dan untuk $2$ adalah $(1,0)$; hasil kali titiknya $1$, bukan nol. Bandingkan dengan
matriks simetris $A$ di lab, yang memberi $\texttt{cond V}\approx1.5$.

Kalau mentok, kembali ke bagian Matriks serupa dan invarian, lalu ke paragraf bilangan kondisi di
Lab komputasi.
\end{tantangan}

\begin{untukdosen}
Urutan mengajar yang berhasil dipakai: nilai eigen lebih dulu sebagai arah yang dipertahankan,
baru polinomial karakteristik sebagai alat hitung. Bila polinomial datang lebih dulu, mahasiswa
melihat akar polinomial dan lupa bahwa yang dicari arah. Setelah ruang eigen, masuk multiplisitas
geometrik lewat nullity, baru syarat diagonalisasi. Contoh guntingan sebaiknya diberikan sebelum
teorema syarat, supaya teorema menjawab pertanyaan yang sudah terasa.

Dua miskonsepsi khas. Pertama, mahasiswa menganggap nilai eigen berulang selalu berarti tidak
dapat didiagonalkan. Matriks $(2,1,1;1,2,1;1,1,2)$ di bagian contoh membantahnya dengan
multiplisitas geometrik dua. Kedua, mereka menuliskan vektor eigen sebagai baris pada $P$, lalu
mendapat $P^{-1}AP$ yang bukan diagonal dan menyalahkan rumusnya. Aturan yang perlu diulang:
vektor eigen adalah kolom, dan urutan kolom $P$ harus sama dengan urutan diagonal $D$.

Aktivitas kelas: bagi kelas menjadi dua kelompok untuk matriks $(2,1,1;1,2,1;1,1,2)$. Kelompok
pertama menghitung $A^{10}$ lewat diagonalisasi, kelompok kedua lewat pangkat berulang atau
perkalian berulang cepat. Bandingkan hasilnya di papan, lalu minta kelas menjelaskan mengapa kedua
jalan harus bertemu. Kegaduhan biasanya muncul pada penyebut $3$ di $P^{-1}$.

Tanda kesiapan lanjut: mahasiswa dapat mengatakan, tanpa menghitung, mengapa matriks guntingan
tidak dapat didiagonalkan, dan dapat menebak bentuk $S^{k}$ untuk guntingan $2\times2$. Bila
mereka menyebut munculnya faktor $k$, mereka siap masuk Bab 10.
\end{untukdosen}

\section{Ringkasan}
\tingkat{dasar}

\begin{ringkasan}
\begin{itemize}
  \item $\lambda$ nilai eigen $A$ bila ada $\mathbf{v}\neq\mathbf{0}$ dengan
        $A\mathbf{v}=\lambda\mathbf{v}$. Nilai eigen adalah akar $p_{A}(\lambda)=\det(\lambda I-A)$,
        dan $E_{\lambda}=\operatorname{null}(A-\lambda I)$ adalah subruangnya.
  \item Derajat $p_{A}$ adalah $n$, jadi ada paling banyak $n$ nilai eigen berbeda. Atas
        $\mathbb{C}$ selalu ada $n$ bila dihitung dengan multiplisitas; atas $\mathbb{R}$ tidak,
        seperti pada matriks putar.
  \item $\operatorname{tr}(A)=\sum_{i}\lambda_{i}$ dan $\det(A)=\prod_{i}\lambda_{i}$, keduanya
        dihitung dengan multiplisitas. Ini pemeriksaan murah untuk setiap hitungan tangan.
  \item Untuk setiap nilai eigen, $1\le\dim E_{\lambda}\le$ multiplisitas aljabar. Vektor eigen
        dengan nilai eigen berbeda selalu bebas linear.
  \item $A$ dapat didiagonalkan tepat ketika ada $n$ vektor eigen bebas, dan tepat ketika
        multiplisitas aljabar sama dengan geometrik untuk setiap nilai eigen. Bila tidak, $A$
        disebut defektif, dan $A^{k}$ pada umumnya memuat faktor $k$, bukan hanya $\lambda^{k}$.
  \item Matriks serupa mempunyai polinomial karakteristik, nilai eigen, jejak, dan determinan
        yang sama; kebalikannya tidak berlaku. Vektor eigen ikut berubah bersama basis.
\end{itemize}
\end{ringkasan}

\section{Latihan}
\tingkat{menengah}

\begin{latihan}
\begin{enumerate}
  \item Tentukan nilai eigen dan vektor eigen
        $A_{1}=\begin{pmatrix}3&1\\ 2&2\end{pmatrix}$, lalu susun $P$ dan $D$ serta periksa
        $P^{-1}A_{1}P=D$ dengan mengalikan balik.

  \item Hitung $G^{10}$ untuk $G=\begin{pmatrix}4&2\\ 1&3\end{pmatrix}$ lewat diagonalisasi,
        lalu periksa hasilnya dengan fungsi \texttt{matrix\_power} dari NumPy. Tuliskan pula
        bentuk umum $G^{k}$.

  \item Buktikan bahwa bila $\lambda$ nilai eigen $A$ maka $\lambda^{2}$ nilai eigen $A^{2}$,
        dan secara umum $\lambda^{k}$ nilai eigen $A^{k}$ untuk setiap bilangan bulat $k\ge1$.
        Jelaskan mengapa kebalikannya gagal: berikan $A$ dengan $A^{2}=-I$, lalu bandingkan nilai
        eigen $A^{2}$ dengan nilai eigen real $A$.
        (Petunjuk: pakai matriks putar di awal bab; ia tidak punya nilai eigen real.)

  \item Buktikan bahwa $A$ invertible jika dan hanya jika $0$ bukan nilai eigen $A$, dan bahwa
        dalam hal itu nilai eigen $A^{-1}$ adalah $1/\lambda$ untuk setiap nilai eigen $\lambda$
        dari $A$.
        (Petunjuk: mulai dari $A\mathbf{v}=\lambda\mathbf{v}$, kalikan dengan $A^{-1}$ dari kiri,
        dan pastikan $\lambda\neq0$ supaya pembagiannya sah.)

  \item Nilai klaim ini: \emph{dua matriks persegi berukuran sama yang mempunyai polinomial
        karakteristik sama pasti serupa.} Tentukan benar, salah, atau belum cukup bukti, dan
        sertakan alasan atau contoh penyangkal.

  \item Buktikan bahwa $A$ dan $A^{\top}$ mempunyai polinomial karakteristik yang sama, dengan
        memakai $\det(\lambda I-A)=\det((\lambda I-A)^{\top})$ dan sifat determinan pada Bab 8.
        Jelaskan mengapa ini berarti transpos tidak mengubah nilai eigen, tetapi biasanya
        mengubah vektor eigen.

  \item Matriks $K=\begin{pmatrix}2&1&0\\ 0&2&0\\ 0&0&1\end{pmatrix}$ mempunyai nilai eigen
        $2,2,1$. Tentukan pangkat $K-2I$, lalu putuskan apakah $K$ dapat didiagonalkan. Jelaskan
        jawabanmu dengan syarat multiplisitas.

  \item Rantai Markov dua keadaan mempunyai matriks
        $T=\begin{pmatrix}0.9&0.4\\ 0.1&0.6\end{pmatrix}$. Tentukan nilai eigen $T$, titik
        tetapnya, dan banyak langkah yang diperlukan agar galat awal menyusut di bawah
        $10^{-6}$. Periksa jawabanmu dengan Python.
        (Petunjuk: setelah nilai eigen diketahui, galat menyusut seperti $|\lambda_{2}|^{k}$;
        selesaikan $|\lambda_{2}|^{k}<10^{-6}$ dengan logaritma.)

  \item Buktikan: bila $A$ dapat didiagonalkan dengan $A=PDP^{-1}$, maka
        $\operatorname{tr}(A^{k})=\sum_{i}\lambda_{i}^{k}$ untuk setiap $k\ge1$, dengan memakai
        $\operatorname{tr}(XY)=\operatorname{tr}(YX)$. Untuk $B$ yang defektif tadi, hitung
        $\operatorname{tr}(B^{10})$ dari matriks yang tercetak, lalu bandingkan dengan
        $1^{10}+1^{10}+2^{10}$. Jelaskan mengapa kesamaan itu tetap berlaku walaupun $B$
        defektif, dan di titik mana bukti tadi memakai diagonalisasi.
        (Petunjuk: jejak tidak berubah oleh perubahan basis; pakai bentuk $B^{k}$ yang sudah
        diturunkan.)
\end{enumerate}
\end{latihan}

\begin{untukpeneliti}
Yang benar-benar terbukti di bab ini adalah pernyataan aljabar atas medan apa pun dan berdimensi
hingga: kesetaraan nilai eigen dengan akar polinomial karakteristik, kebebasan vektor eigen untuk
nilai eigen berbeda, ketaksamaan $1\le\dim E_{\lambda}\le$ multiplisitas aljabar, dan kesetaraan
diagonalisasi dengan dua syarat yang dipakai di sini. Rujukan bakunya Axler, \emph{Linear Algebra
Done Right} edisi ke-4, bagian 5A dan 5D (5.55, empat syarat setara, dan 5.54); versi matriks ada
di Anton, \emph{Elementary Linear Algebra} edisi ke-11, bagian 5.1 dan 5.2.

Yang lebih sering dipopulerkan daripada dibuktikan ada dua. Pertama, klaim bahwa nilai eigen
``mengukur kepentingan'', dasar penjelasan PageRank: vektor eigen utama matriks stokastik memberi
bobot tunak hanya bila rantainya tak tereduksi dan aperiodik, dan itu isi teorema Perron-Frobenius,
yang tidak dibuktikan di sini; bukti lengkapnya di Strang, \emph{Introduction to Linear Algebra}
edisi ke-6, Apendiks 8. Kedua, klaim bahwa polinomial karakteristik adalah cara praktis menghitung
nilai eigen: untuk $n\ge5$ tidak ada rumus akar umum, dan pustaka memakai algoritma QR pada
matriks, bukan pemfaktoran polinomial. Bab ini juga berhenti pada kasus dapat didiagonalkan; bentuk
Jordan dan polinomial minimal tidak dibahas, dan ada di Hoffman dan Kunze, \emph{Linear Algebra}
edisi ke-2, bagian 6.3 dan 7.3.

Pertanyaan terbuka: seberapa jauh nilai eigen bergeser bila entri diganggu, yang jawaban
kuantitatifnya menuntut teorema Bauer-Fike atau pseudospektrum; dan apakah kedefektifan itu sendiri
stabil secara numerik. Batas keberlakuan: dimensi hingga, medan real atau kompleks, tanpa klaim
tentang stabilitas algoritma.
\end{untukpeneliti}

\peta{Bab ini mengikuti Anton, \emph{Elementary Linear Algebra} edisi ke-11, Bab 5: bagian 5.1
(Eigenvalues and Eigenvectors), 5.2 (Diagonalization), 5.3 (Complex Vector Spaces), 5.4
(Application: Differential Equations), dan 5.5 (Application: Dynamical Systems and Markov
Chains). Padanan yang lebih abstrak ada di Axler, \emph{Linear Algebra Done Right} edisi ke-4,
Bab 5: 5A (nilai eigen dan ruang eigen), 5B (keberadaan nilai eigen atas ruang vektor kompleks),
dan 5D (operator yang dapat didiagonalkan, termasuk 5.55). Untuk rantai Markov dan teorema
Perron-Frobenius, lihat Strang, \emph{Introduction to Linear Algebra} edisi ke-6, bagian 6.1,
6.2, 6.5, dan Apendiks 8. Bentuk kanonik untuk matriks defektif ada di Hoffman dan Kunze,
\emph{Linear Algebra} edisi ke-2, Bab 6 (Elementary Canonical Forms), bagian 6.3 untuk polinomial
minimal, dan Bab 7 (The Rational and Jordan Forms), bagian 7.3 untuk bentuk Jordan.}

\chapter{Ortogonalitas, SVD, dan Kuadrat Terkecil}

Lima pembacaan dari satu probe suhu: $100$, $101$, $102$, $103$, dan $104$ derajat.
Bacaan yang tercatat berturut-turut $302{,}0$, $305{,}1$, $307{,}9$, $311{,}2$, dan
$313{,}8$. Model kalibrasinya satu garis lurus $y=a+bT$, jadi ada lima persamaan untuk
dua anu:
\[
  a+100b=302{,}0,\qquad a+101b=305{,}1,\qquad \dots,\qquad a+104b=313{,}8 .
\]
Tidak ada pasangan $(a,b)$ yang memenuhi kelimanya; kelima titik itu tidak pernah
segaris, dan selisih tiap persamaan berbeda tanda. Tetap ada satu pasangan yang
paling masuk akal. Dalam pengertian apa pasangan itu terbaik, dan seberapa jauh ia
bergeser bila satu bacaan berubah sepersepuluh?

Pertanyaan yang sama muncul dari sisi yang berlawanan. Sebuah detektor menghitung skor
$s=\mathbf{w}^{\top}\mathbf{x}$ dari tiga pencacah, dan penyerang ingin menggeser $s$
tepat sebesar $\Delta$ dengan perubahan $\mathbf{x}$ sesedikit mungkin, karena perubahan
besar mudah terlihat. Ada tiga anu dan satu syarat, jadi solusinya terlalu banyak, bukan
tidak ada: dari sekian perubahan yang memberi pergeseran sama, mana yang paling kecil?

Dua persoalan itu satu persoalan. Pada kasus pertama himpunan calon jawaban kosong dan
kita mencari yang terdekat; pada kasus kedua terlalu besar dan kita mencari yang
terpendek. Keduanya diselesaikan proyeksi tegak lurus. Bab 9 memberi nilai eigen, tetapi
hanya untuk matriks persegi yang vektornya cukup.
Matriks kalibrasi di atas berukuran $5\times2$, dan matriks detektor $2\times3$.
Memaksa keduanya menjadi persegi lewat $\mathbf{A}^{\top}\mathbf{A}$ memang bisa,
tetapi jalan itu mengkuadratkan kepekaan terhadap galat. Bagian keenam menunjukkan
harganya dengan angka.

\begin{untukanak}
Pernahkah kamu berdiri di bawah lampu dan melihat bayanganmu di lantai? Bayangan itu
bukan kamu: ia lebih pendek, gepeng, dan tidak bisa menunjukkan wajahmu. Tetapi ia hal
terbaik yang bisa dilakukan lantai: dari semua titik di lantai, bayanganmu yang paling
dekat dengan ujung kepalamu. Kalau lampu digeser, bayangan ikut bergeser, padahal kamu
tidak bergerak.

Sekarang bayangkan dua penggaris yang bedanya hanya seperseratus sentimeter. Kalau kamu
mengukur meja dengan keduanya lalu mengurangkan hasilnya, angka sisa itu hampir
seluruhnya berasal dari kekeliruan membaca. Selisih kecil mudah tenggelam.

Di bab ini kita akan tahu mengapa bayangan yang ``paling dekat'' itu bisa dihitung dengan
rapi, mengapa dari banyak jalan menuju tujuan yang sama ada satu jalan terpendek, dan
mengapa dua aturan ukur yang hampir kembar membuat jawaban kita goyah.
\end{untukanak}

\section{Hasilkali dalam, norma, dan ketaksamaan Cauchy-Schwarz}
\tingkat{dasar}

Panjang dan sudut adalah dua kata yang kita pakai sejak sekolah tanpa pernah
mendefinisikannya. Untuk vektor di $\mathbb{R}^{n}$ keduanya berasal dari satu operasi
yang memakan dua vektor dan mengembalikan satu bilangan; operasi itulah yang memberi
ruang vektor ukuran, karena aksioma ruang vektor di Bab 6 tidak menyebut jarak.

Aturannya: kalikan komponen demi komponen, lalu jumlahkan. Hasilnya nol tepat ketika
kedua vektor bersudut siku, sehingga operasi ini dipakai sebagai pengganti sudut bahkan
di ruang berdimensi seratus.

\kamus{hasilkali dalam}{inner product}{Dua daftar angka dilipat menjadi satu angka yang
mengukur seberapa searah keduanya.}

\begin{definisi}[Hasilkali dalam dan norma]
\emph{Hasilkali dalam} pada ruang vektor real $V$ adalah fungsi
$\langle\cdot,\cdot\rangle\colon V\times V\to\mathbb{R}$ yang memenuhi tiga syarat:
\begin{enumerate}
  \item simetris, $\langle\mathbf{u},\mathbf{v}\rangle=\langle\mathbf{v},\mathbf{u}\rangle$;
  \item linear pada argumen pertama,
        $\langle c\mathbf{u}+\mathbf{w},\mathbf{v}\rangle
         =c\langle\mathbf{u},\mathbf{v}\rangle+\langle\mathbf{w},\mathbf{v}\rangle$;
  \item definit positif, $\langle\mathbf{v},\mathbf{v}\rangle\ge0$ dengan kesamaan
        hanya bila $\mathbf{v}=\mathbf{0}$.
\end{enumerate}
Pada $\mathbb{R}^{n}$ hasilkali dalam bakunya
$\langle\mathbf{u},\mathbf{v}\rangle=\sum_{i}u_{i}v_{i}
=\mathbf{u}^{\top}\mathbf{v}$. \emph{Norma} yang ditimbulkannya adalah
$\|\mathbf{v}\|=\sqrt{\langle\mathbf{v},\mathbf{v}\rangle}$.\index{hasilkali dalam}
\index{norma}
\end{definisi}

Syarat ketiga yang membuat $\|\mathbf{v}\|=0$ hanya untuk vektor nol, sehingga norma
layak disebut panjang. Norma juga memenuhi $\|c\mathbf{v}\|=|c|\,\|\mathbf{v}\|$ dan
$\|\mathbf{u}+\mathbf{v}\|\le\|\mathbf{u}\|+\|\mathbf{v}\|$; yang kedua tidak otomatis,
melainkan akibat ketaksamaan berikut.

Dalam bahasa sehari-hari: hasilkali dalam dua vektor tidak pernah lebih besar daripada
hasil kali panjangnya. Dua vektor bisa sangat panjang, tetapi bila arahnya bersiku hasil
kalinya tetap kecil.

\begin{teorema}[Cauchy-Schwarz]
Untuk semua $\mathbf{u},\mathbf{v}\in V$ berlaku
$|\langle\mathbf{u},\mathbf{v}\rangle|\le\|\mathbf{u}\|\,\|\mathbf{v}\|$, dengan
kesamaan tepat ketika salah satu vektor kelipatan skalar vektor yang lain.
\index{Cauchy-Schwarz}
\end{teorema}

\begin{proof}
Bila $\mathbf{v}=\mathbf{0}$, kedua ruas nol dan tidak ada yang perlu dibuktikan.
Misalkan $\mathbf{v}\ne\mathbf{0}$ dan $t\in\mathbb{R}$ sebarang. Definit positif
memberi
\[
  0\le\|\mathbf{u}-t\mathbf{v}\|^{2}
   =\|\mathbf{u}\|^{2}-2t\langle\mathbf{u},\mathbf{v}\rangle+t^{2}\|\mathbf{v}\|^{2}.
\]
Ruas kanan adalah polinomial kuadrat dalam $t$ yang tidak pernah negatif. Koefisien
$t^{2}$ positif, jadi diskriminannya tidak boleh positif:
\[
  4\langle\mathbf{u},\mathbf{v}\rangle^{2}-4\|\mathbf{u}\|^{2}\|\mathbf{v}\|^{2}\le0 .
\]
Setelah dibagi empat dan diakarkan, diperoleh ketaksamaan yang diminta. Kesamaan
terjadi tepat ketika polinomial itu punya akar kembar, yaitu ketika
$\mathbf{u}-t\mathbf{v}=\mathbf{0}$ untuk suatu $t$.
\end{proof}

Bukti itu tidak memakai koordinat, jadi ketaksamaan berlaku di setiap ruang
berhasilkali dalam, termasuk ruang fungsi.

\begin{akibat}[Ketaksamaan segitiga]
$\|\mathbf{u}+\mathbf{v}\|\le\|\mathbf{u}\|+\|\mathbf{v}\|$.
\end{akibat}

\begin{proof}
Kembangkan $\|\mathbf{u}+\mathbf{v}\|^{2}
=\|\mathbf{u}\|^{2}+2\langle\mathbf{u},\mathbf{v}\rangle+\|\mathbf{v}\|^{2}$.
Cauchy-Schwarz membatasi suku tengah oleh $2\|\mathbf{u}\|\|\mathbf{v}\|$, sehingga
\[
  \|\mathbf{u}+\mathbf{v}\|^{2}\le\|\mathbf{u}\|^{2}+2\|\mathbf{u}\|\|\mathbf{v}\|
  +\|\mathbf{v}\|^{2}=(\|\mathbf{u}\|+\|\mathbf{v}\|)^{2}.
\]
Kedua ruas tidak negatif, jadi akar kuadratnya boleh diambil.
\end{proof}

Sekarang kita punya panjang. Sudut menyusul lewat definisi: dua vektor tegak lurus bila
hasilkali dalamnya nol. Itu masuk akal karena Cauchy-Schwarz menjamin
$\langle\mathbf{u},\mathbf{v}\rangle/(\|\mathbf{u}\|\|\mathbf{v}\|)$ terletak antara $-1$
dan $1$, sehingga boleh dibaca sebagai kosinus.

\kamus{ortogonal}{orthogonal}{Dua arah yang tidak saling menyumbang sama sekali; bayangan
satu arah pada arah lain panjangnya nol.}

\begin{definisi}[Ortogonal dan ortonormal]
$\mathbf{u}$ dan $\mathbf{v}$ \emph{ortogonal} bila
$\langle\mathbf{u},\mathbf{v}\rangle=0$, ditulis $\mathbf{u}\perp\mathbf{v}$. Himpunan
vektor disebut \emph{ortogonal} bila setiap dua anggotanya yang berbeda ortogonal, dan
\emph{ortonormal} bila semua anggotanya juga bernorma satu.\index{ortogonalitas}
\end{definisi}

Akibat kecil yang langsung dari pengembangan norma: bila $\mathbf{u}\perp\mathbf{v}$
maka $\|\mathbf{u}+\mathbf{v}\|^{2}=\|\mathbf{u}\|^{2}+\|\mathbf{v}\|^{2}$. Teorema
Pythagoras adalah kasus khusus definisi ini, bukan sebaliknya.

\section{Himpunan ortogonal, proyeksi, dan Gram-Schmidt}
\tingkat{menengah}

Untuk basis sembarang kita harus menyelesaikan sistem persamaan linear; untuk basis
ortonormal tidak, sebab tiap koefisien dibaca dengan satu hasilkali dalam. Satu fakta
membuatnya sah: himpunan ortogonal tanpa vektor nol selalu bebas linear.

\begin{lemma}
Bila $\mathbf{v}_{1},\dots,\mathbf{v}_{k}$ ortogonal dan semuanya bukan vektor nol,
maka himpunan itu bebas linear.
\end{lemma}

\begin{proof}
Misalkan $\sum_{i}c_{i}\mathbf{v}_{i}=\mathbf{0}$. Ambil hasilkali dalam dengan
$\mathbf{v}_{j}$. Semua suku dengan $i\ne j$ lenyap karena ortogonal, sehingga tinggal
$c_{j}\langle\mathbf{v}_{j},\mathbf{v}_{j}\rangle=0$. Karena $\mathbf{v}_{j}\ne\mathbf{0}$,
normanya bukan nol, jadi $c_{j}=0$. Ini berlaku untuk setiap $j$.
\end{proof}

Sekarang pernyataan tentang koefisien tadi, dalam bentuk yang paling sering dipakai.

\begin{proposisi}
Misalkan $(\mathbf{q}_{1},\dots,\mathbf{q}_{n})$ basis ortonormal dari $V$. Maka untuk
setiap $\mathbf{v}\in V$,
\[
  \mathbf{v}=\sum_{i=1}^{n}\langle\mathbf{v},\mathbf{q}_{i}\rangle\,\mathbf{q}_{i},
  \qquad
  \|\mathbf{v}\|^{2}=\sum_{i=1}^{n}\langle\mathbf{v},\mathbf{q}_{i}\rangle^{2}.
\]
\end{proposisi}

\begin{proof}
Tulis $\mathbf{v}=\sum_{i}a_{i}\mathbf{q}_{i}$; koefisien itu ada dan tunggal karena
basis. Hasilkali dalam dengan $\mathbf{q}_{j}$ memberi
$\langle\mathbf{v},\mathbf{q}_{j}\rangle=\sum_{i}a_{i}\langle\mathbf{q}_{i},\mathbf{q}_{j}\rangle
=a_{j}$, sebab $\langle\mathbf{q}_{i},\mathbf{q}_{j}\rangle$ bernilai satu bila $i=j$
dan nol bila tidak. Untuk rumus norma, pakai ortogonalitas pada pengembangan
$\|\mathbf{v}\|^{2}=\langle\sum_{i}a_{i}\mathbf{q}_{i},\sum_{j}a_{j}\mathbf{q}_{j}\rangle$;
suku silang semuanya nol.
\end{proof}

Kesamaan kedua disebut kesamaan Parseval: pada basis ortonormal, panjang total adalah
akar jumlah kuadrat koordinat; pada basis miring sifat itu gagal.

Pertanyaan berikutnya muncul dari masalah pembuka: $\mathbf{b}$ berada di luar subruang
$W$, dan kita ingin titik $W$ yang paling dekat dengannya. Titik itu disebut proyeksi,
dan cara menemukannya: jumlahkan bayangan $\mathbf{b}$ pada setiap arah ortonormal $W$.

\kamus{proyeksi}{projection}{Bayangan sebuah tongkat di lantai: titik di lantai yang
paling dekat dengan ujung tongkat.}

\begin{definisi}[Proyeksi ortogonal]
Misalkan $W\subseteq\mathbb{R}^{m}$ subruang dengan basis ortonormal
$(\mathbf{q}_{1},\dots,\mathbf{q}_{k})$. \emph{Proyeksi ortogonal} $\mathbf{b}$ pada $W$
adalah
\[
  \operatorname{proj}_{W}\mathbf{b}=\sum_{i=1}^{k}\langle\mathbf{b},\mathbf{q}_{i}\rangle\mathbf{q}_{i},
  \qquad
  \mathbf{r}=\mathbf{b}-\operatorname{proj}_{W}\mathbf{b}.
\]
Vektor $\mathbf{r}$ disebut \emph{sisa} (residu).\index{proyeksi}
\end{definisi}

Sisa itu ortogonal terhadap seluruh $W$, bukan hanya terhadap vektor basis: karena
$\langle\mathbf{r},\mathbf{q}_{j}\rangle=\langle\mathbf{b},\mathbf{q}_{j}\rangle
-\langle\mathbf{b},\mathbf{q}_{j}\rangle=0$ dan setiap vektor $W$ adalah kombinasi
$\mathbf{q}_{j}$.

Yang membuat proyeksi berguna: ia bukan sekadar satu titik di $W$, melainkan titik
terdekat. Itu klaim tentang semua titik lain di $W$, jadi perlu bukti.

\begin{bacakeras}
Bacalah keras, dua suara bergantian.

\textbf{Suara satu:} Ada satu titik menggantung di atas lantai, dan lantai diminta mengambilkan wakilnya.

\textbf{Suara dua:} Lantai tidak bisa naik. Ia hanya bisa menjatuhkan bayangan, dan bayangan itu titik lantai yang paling dekat dengan titik tadi.

\textbf{Suara satu:} Lalu ke mana perginya bagian yang tidak sampai ke lantai?

\textbf{Suara dua:} Bagian itu berdiri tegak. Ia tidak punya lantai untuk berpijak, jadi tidak ada jalan di lantai yang bisa memendekkannya.

Ulangi satu baris ini sampai terasa di lidah: \emph{yang jatuh ke lantai adalah yang terdekat, yang berdiri tegak adalah yang tak terpulangkan.}
\end{bacakeras}

\begin{teorema}[Hampiran terbaik]
Untuk setiap $\mathbf{w}\in W$ berlaku
$\|\mathbf{b}-\operatorname{proj}_{W}\mathbf{b}\|\le\|\mathbf{b}-\mathbf{w}\|$, dengan
kesamaan hanya bila $\mathbf{w}=\operatorname{proj}_{W}\mathbf{b}$.
\end{teorema}

\begin{proof}
Tulis $\mathbf{p}=\operatorname{proj}_{W}\mathbf{b}$. Untuk sembarang $\mathbf{w}\in W$,
\[
  \mathbf{b}-\mathbf{w}=(\mathbf{b}-\mathbf{p})+(\mathbf{p}-\mathbf{w}).
\]
Kedua suku ortogonal: yang pertama ortogonal terhadap seluruh $W$, dan
$\mathbf{p}-\mathbf{w}\in W$. Karena itu norma kuadratnya terbagi:
\[
  \|\mathbf{b}-\mathbf{w}\|^{2}=\|\mathbf{b}-\mathbf{p}\|^{2}+\|\mathbf{p}-\mathbf{w}\|^{2}.
\]
Suku kedua tidak negatif, sehingga ruas kanan paling kecil ketika $\mathbf{w}=\mathbf{p}$.
\end{proof}

Bacaan geometrinya ada pada Gambar~\ref{fig:proyeksi-sisa}.

\begin{figure}[htbp]
\centering
\begin{tikzpicture}[scale=1.18]
  \draw[->,layergrey] (-0.5,0)--(3.4,0) node[right,font=\small]{$x_{1}$};
  \draw[->,layergrey] (0,-0.5)--(0,3.9) node[above,font=\small]{$x_{2}$};
  \draw[very thick,accent] (-0.42,-0.21)--(2.75,1.375)
    node[right,font=\small]{$W$};
  \draw[densely dotted,thick,layergrey] (0.2,3.4)--(2.6,1.3);
  \fill[layergrey] (2.6,1.3) circle (1.2pt);
  \node[font=\small,layergrey,below right=-2pt] at (2.6,1.3) {$\mathbf{w}$};
  \node[font=\small,layergrey,above right=-3pt] at (1.05,2.62)
    {$\|\mathbf{b}-\mathbf{w}\|$};
  \draw[->,very thick,layerteal] (0,0)--(0.2,3.4)
    node[left,font=\small]{$\mathbf{b}$};
  \draw[->,thick,layerindigo] (0,0)--(1.52,0.76);
  \node[font=\small,layerindigo,below right=-2pt] at (1.34,0.6) {$\mathbf{p}$};
  \draw[-{Latex},very thick,layerplum] (1.52,0.76)--(0.2,3.4);
  \node[font=\small,layerplum,right=2pt] at (0.95,2.15) {$\mathbf{r}$};
  \draw[thick,layergrey] (1.717,0.858)--(1.619,1.055)--(1.422,0.957);
  \fill[layerplum] (1.52,0.76) circle (1.4pt);
\end{tikzpicture}
\caption{Proyeksi ortogonal $\mathbf{b}=(0{,}2;\,3{,}4)$ pada garis $W$ yang
direntang $\mathbf{q}=(2,1)/\sqrt{5}$. Kaki tegak lurusnya
$\mathbf{p}=(1{,}52;\,0{,}76)$, ditandai siku-siku dengan $W$, dan sisa
$\mathbf{r}=\mathbf{b}-\mathbf{p}$ tegak lurus terhadap seluruh $W$. Titik
$\mathbf{w}=(2{,}6;\,1{,}3)$ yang juga berada di $W$ memberi jarak
$\|\mathbf{b}-\mathbf{w}\|\approx3{,}19$, lebih panjang daripada
$\|\mathbf{b}-\mathbf{p}\|\approx2{,}95$.}
\label{fig:proyeksi-sisa}
\end{figure}

\begin{tangan}
Bahan: satu pensil, satu lampu meja, satu lembar kertas, dan penggaris.

Letakkan kertas di meja, lalu jauhkan lampu ke ujung ruangan sampai sinarnya hampir sejajar. Lampu yang terlalu dekat memanjangkan bayangan, dan yang terukur bukan lagi proyeksi ortogonal.

Pegang pensil tegak dengan ujungnya menyentuh kertas, dan tandai titik itu. Miringkan pensil sampai bersudut enam puluh derajat dengan kertas, tandai ujung bayangannya, lalu ulangi pada empat puluh lima dan tiga puluh derajat. Ukur ketiga bayangan dengan penggaris.

Yang perlu dilihat: bayangan paling panjang ketika pensil hampir menempel kertas, dan paling pendek ketika pensil tegak. Panjang bayangan tidak peduli seberapa tinggi ujung pensil di atas kertas; ia hanya mencatat bagian mendatar. Bagian tegak itulah sisa yang tidak bisa dijangkau lantai.
\end{tangan}

Bila $\mathbf{Q}$ kolomnya $\mathbf{q}_{1},\dots,\mathbf{q}_{k}$, rumus proyeksi menjadi
satu perkalian matriks: $\mathbf{p}=\mathbf{Q}\mathbf{Q}^{\top}\mathbf{b}$. Matriks
$\mathbf{P}=\mathbf{Q}\mathbf{Q}^{\top}$ bersifat $\mathbf{P}^{2}=\mathbf{P}$ dan
$\mathbf{P}^{\top}=\mathbf{P}$: memproyeksikan dua kali sama dengan sekali. Jejaknya $k$,
yaitu dimensi $W$; untuk contoh di bagian ketujuh jejaknya $2$.

Sekarang kita tahu memproyeksikan, tetapi belum punya cara memperoleh basis ortonormal.
Algoritme Gram-Schmidt membangunnya satu per satu: normalkan vektor pertama, lalu buang
dari vektor berikutnya bagian yang sejajar dengan yang sudah ada. Setiap langkahnya
hanya memakai proyeksi yang baru kita buktikan.

\begin{teorema}[Gram-Schmidt dan faktorisasi QR]
Misalkan $\mathbf{a}_{1},\dots,\mathbf{a}_{n}$ bebas linear di $\mathbb{R}^{m}$.
Definisikan berturut-turut
\[
  \mathbf{v}_{1}=\mathbf{a}_{1},\qquad
  \mathbf{v}_{j}=\mathbf{a}_{j}-\sum_{i<j}\langle\mathbf{a}_{j},\mathbf{q}_{i}\rangle\mathbf{q}_{i},
  \qquad
  \mathbf{q}_{j}=\frac{\mathbf{v}_{j}}{\|\mathbf{v}_{j}\|}
\]
untuk $j=2,\dots,n$. Maka $\mathbf{q}_{1},\dots,\mathbf{q}_{n}$ ortonormal dan
$\operatorname{span}(\mathbf{q}_{1},\dots,\mathbf{q}_{j})
=\operatorname{span}(\mathbf{a}_{1},\dots,\mathbf{a}_{j})$ untuk setiap $j$. Akibatnya
$\mathbf{A}=\mathbf{Q}\mathbf{R}$ dengan $\mathbf{Q}$ berkolom ortonormal dan
$\mathbf{R}$ segitiga atas dengan diagonal positif.\index{Gram-Schmidt}\index{QR}
\end{teorema}

\begin{proof}
Kita buktikan dengan induksi pada $j$. Untuk $j=1$, $\mathbf{q}_{1}$ adalah
$\mathbf{a}_{1}$ yang dinormalkan, jadi benar. Andaikan klaim benar sampai $j-1$.
Vektor $\mathbf{v}_{j}$ adalah $\mathbf{a}_{j}$ dikurangi kombinasinya dengan
$\mathbf{q}_{1},\dots,\mathbf{q}_{j-1}$, sehingga $\mathbf{v}_{j}$ tetap berada di
$\operatorname{span}(\mathbf{a}_{1},\dots,\mathbf{a}_{j})$. Untuk setiap $i<j$,
\[
  \langle\mathbf{v}_{j},\mathbf{q}_{i}\rangle
  =\langle\mathbf{a}_{j},\mathbf{q}_{i}\rangle-\langle\mathbf{a}_{j},\mathbf{q}_{i}\rangle=0,
\]
karena $\mathbf{q}_{i}$ ortonormal. Jadi $\mathbf{v}_{j}$ ortogonal terhadap semua
vektor sebelumnya. Vektor itu bukan nol: bila nol, $\mathbf{a}_{j}$ akan menjadi
kombinasi $\mathbf{a}_{1},\dots,\mathbf{a}_{j-1}$, bertentangan dengan kebebasan linear.
Karena itu $\mathbf{q}_{j}$ terdefinisi dan himpunan bertambah ortonormal.

Untuk faktorisasi, tulis setiap $\mathbf{a}_{j}$ sebagai
$\|\mathbf{v}_{j}\|\mathbf{q}_{j}+\sum_{i<j}\langle\mathbf{a}_{j},\mathbf{q}_{i}\rangle\mathbf{q}_{i}$.
Koefisien itu mengisi kolom ke-$j$ dari $\mathbf{R}$; yang berada di atas diagonal adalah
$\langle\mathbf{a}_{j},\mathbf{q}_{i}\rangle$ dan diagonalnya $\|\mathbf{v}_{j}\|>0$.
\end{proof}

Satu peringatan: Gram-Schmidt klasik menghitung setiap koefisien terhadap vektor asli
$\mathbf{a}_{j}$, bukan terhadap vektor yang sudah dikurangi. Bila kolom $\mathbf{A}$
hampir bergantung, hasilnya berhenti ortonormal; mengurangi $\mathbf{v}_{j}$ lebih dulu
memperbaikinya dengan biaya hampir nol. Bagian kedelapan menunjukkan angkanya: pada
matriks dengan $\kappa\approx6{,}8\times10^{6}$, versi klasik memberi
$\|\mathbf{Q}^{\top}\mathbf{Q}-\mathbf{I}\|=2{,}05\times10^{-2}$, versi modifikasi
$3{,}40\times10^{-10}$, Householder $1{,}19\times10^{-15}$; angka pertama berarti dua
kolom yang seharusnya bersiku berselisih sudut sekitar satu derajat.

\section{Kuadrat terkecil dan persamaan normal}
\tingkat{menengah}

Kembali ke lima bacaan probe. Sistem $\mathbf{A}\mathbf{x}=\mathbf{b}$ tidak punya solusi,
tetapi kita bisa menuntut sisa $\mathbf{b}-\mathbf{A}\mathbf{x}$ sekecil mungkin:
menuntut sisa nol berarti menuntut solusi, menuntut panjangnya minimum selalu punya
jawaban. Ukuran yang dipakai jumlah kuadrat selisih, bukan jumlah nilai mutlak, karena
jumlah kuadrat lahir dari hasilkali dalam sehingga seluruh mesin proyeksi bisa dipakai;
jumlah nilai mutlak tidak, dan jawabannya bisa tidak tunggal.

\kamus{kuadrat terkecil}{least squares}{Ketika tidak ada garis yang melewati semua titik,
pilih garis yang jumlah kuadrat kesalahannya paling kecil.}

\begin{definisi}[Penyelesaian kuadrat terkecil]
Untuk $\mathbf{A}\in\mathbb{R}^{m\times n}$ dan $\mathbf{b}\in\mathbb{R}^{m}$,
vektor $\hat{\mathbf{x}}\in\mathbb{R}^{n}$ disebut \emph{penyelesaian kuadrat terkecil}
bila $\|\mathbf{A}\hat{\mathbf{x}}-\mathbf{b}\|\le\|\mathbf{A}\mathbf{x}-\mathbf{b}\|$
untuk setiap $\mathbf{x}\in\mathbb{R}^{n}$.\index{kuadrat terkecil}
\end{definisi}

Definisi itu tidak menyebut turunan. Yang dipakai geometri: $\mathbf{A}\mathbf{x}$ selalu
di dalam $\operatorname{col}(\mathbf{A})$, jadi sisa terkecil adalah titik
$\operatorname{col}(\mathbf{A})$ terdekat dengan $\mathbf{b}$; itu tepat masalah proyeksi
tadi.

\begin{teorema}[Persamaan normal]
Vektor $\hat{\mathbf{x}}$ meminimumkan $\|\mathbf{A}\mathbf{x}-\mathbf{b}\|$ jika dan
hanya jika
\[
  \mathbf{A}^{\top}\mathbf{A}\hat{\mathbf{x}}=\mathbf{A}^{\top}\mathbf{b}.
\]
\end{teorema}

\begin{proof}
Ambil sembarang $\mathbf{x}$ dan sembarang $\mathbf{h}\in\mathbb{R}^{n}$, lalu tulis
$\mathbf{x}+\mathbf{h}$ sebagai tetangga $\mathbf{x}$. Maka
\[
  \|\mathbf{A}(\mathbf{x}+\mathbf{h})-\mathbf{b}\|^{2}
  =\|\mathbf{A}\mathbf{x}-\mathbf{b}\|^{2}
   +2\mathbf{h}^{\top}\mathbf{A}^{\top}(\mathbf{A}\mathbf{x}-\mathbf{b})
   +\|\mathbf{A}\mathbf{h}\|^{2}.
\]

Pertama, andaikan $\mathbf{A}^{\top}(\mathbf{A}\mathbf{x}-\mathbf{b})=\mathbf{0}$.
Suku tengah lenyap untuk setiap $\mathbf{h}$, dan suku terakhir tidak negatif, sehingga
$\|\mathbf{A}(\mathbf{x}+\mathbf{h})-\mathbf{b}\|^{2}\ge\|\mathbf{A}\mathbf{x}-\mathbf{b}\|^{2}$.
Jadi $\mathbf{x}$ meminimumkan.

Kedua, andaikan $\mathbf{A}^{\top}(\mathbf{A}\mathbf{x}-\mathbf{b})=\mathbf{z}\ne\mathbf{0}$.
Pilih $\mathbf{h}=-\epsilon\mathbf{z}$ dengan $\epsilon>0$ kecil. Suku tengah menjadi
$-2\epsilon\|\mathbf{z}\|^{2}$, sedangkan suku terakhir tumbuh sebesar
$\epsilon^{2}\|\mathbf{A}\mathbf{z}\|^{2}$. Untuk $\epsilon$ cukup kecil, yang negatif
mengalahkan yang positif, sehingga $\mathbf{x}$ bukan peminimum. Jadi syarat itu perlu.
\end{proof}

Sistem $\mathbf{A}^{\top}\mathbf{A}\mathbf{x}=\mathbf{A}^{\top}\mathbf{b}$ disebut
\emph{persamaan normal}. Syarat
$\mathbf{A}^{\top}(\mathbf{A}\hat{\mathbf{x}}-\mathbf{b})=\mathbf{0}$ mengatakan sisa
$\mathbf{r}=\mathbf{b}-\mathbf{A}\hat{\mathbf{x}}$ ortogonal terhadap setiap kolom
$\mathbf{A}$, jadi terhadap seluruh $\operatorname{col}(\mathbf{A})$. Dengan kata lain
$\mathbf{A}\hat{\mathbf{x}}$ adalah proyeksi $\mathbf{b}$ pada ruang kolom.\index{persamaan normal}

Dari sini satu pertanyaan praktis: kapan jawabannya tunggal? Matriks
$\mathbf{A}^{\top}\mathbf{A}$ selalu persegi dan simetris, tetapi belum tentu dapat
dibalik.

\begin{akibat}
Bila kolom-kolom $\mathbf{A}$ bebas linear, maka $\mathbf{A}^{\top}\mathbf{A}$ dapat
dibalik dan penyelesaian kuadrat terkecil tunggal, yaitu
$\hat{\mathbf{x}}=(\mathbf{A}^{\top}\mathbf{A})^{-1}\mathbf{A}^{\top}\mathbf{b}$.
\end{akibat}

\begin{proof}
Misalkan $\mathbf{A}^{\top}\mathbf{A}\mathbf{z}=\mathbf{0}$. Kalikan dari kiri dengan
$\mathbf{z}^{\top}$, diperoleh $\|\mathbf{A}\mathbf{z}\|^{2}
=\mathbf{z}^{\top}\mathbf{A}^{\top}\mathbf{A}\mathbf{z}=0$, jadi
$\mathbf{A}\mathbf{z}=\mathbf{0}$. Kebebasan linear kolom memaksa $\mathbf{z}=\mathbf{0}$.
Jadi inti $\mathbf{A}^{\top}\mathbf{A}$ hanya memuat vektor nol, dan matriks persegi
dengan inti trivial dapat dibalik. Ketunggalan $\hat{\mathbf{x}}$ lalu mengikuti
Teorema persamaan normal.
\end{proof}

\begin{resep}
\begin{enumerate}
  \item Susun $\mathbf{A}$ dan $\mathbf{b}$: satu baris $\mathbf{A}$ untuk tiap persamaan, kolomnya fungsi yang dipakai.
  \item Bentuk $\mathbf{A}^{\top}\mathbf{A}$, yang berukuran $n\times n$, dan $\mathbf{A}^{\top}\mathbf{b}$, yang berukuran $n\times1$.
  \item Selesaikan $\mathbf{A}^{\top}\mathbf{A}\hat{\mathbf{x}}=\mathbf{A}^{\top}\mathbf{b}$; periksa $\det(\mathbf{A}^{\top}\mathbf{A})\ne0$ lebih dulu, sebab nol berarti jawabannya tidak tunggal.
  \item Hitung $\mathbf{r}=\mathbf{b}-\mathbf{A}\hat{\mathbf{x}}$, periksa $\mathbf{A}^{\top}\mathbf{r}=\mathbf{0}$, lalu laporkan $\|\mathbf{r}\|$.
\end{enumerate}
Peringatan urutan: jangan membalik apa pun sebelum langkah ketiga. $\mathbf{A}$ tidak persegi, jadi $\mathbf{A}^{-1}$ tidak ada, dan menukar $\mathbf{A}^{\top}\mathbf{b}$ menjadi $\mathbf{b}^{\top}\mathbf{A}$ memberi vektor baris $1\times n$, bukan sistem persamaan.
\end{resep}

Sebaliknya, bila kolom $\mathbf{A}$ bergantung, $\mathbf{A}^{\top}\mathbf{A}$ singular
dan ada tak hingga penyelesaian; bagian kelima memilih satu di antaranya. Persamaan
normal mengubah masalah $m\times n$ menjadi $n\times n$, penghematan besar untuk
$m\gg n$, dan harganya terlihat di bagian keenam.

\section{Dekomposisi nilai singular}
\tingkat{lanjut}

Bab 9 mengubah matriks persegi menjadi bentuk diagonal dengan memilih satu basis.
Syaratnya berat: matriks harus punya cukup vektor eigen, dan untuk matriks real nilai
eigennya harus real. Matriks kalibrasi kita berukuran $5\times2$, jadi ia bahkan tidak
persegi, dan tidak ada yang namanya nilai eigen untuknya.

Jalan keluarnya: izinkan dua basis berbeda, satu di domain dan satu di kodomain,
keduanya ortonormal. Dengan kelonggaran itu setiap matriks bersifat diagonal; itulah
dekomposisi nilai singular. Bilangan di diagonal adalah panjang sumbu elips yang
dihasilkan dari bola satuan, selalu real dan tidak negatif meskipun matriksnya tidak
simetris.

\kamus{nilai singular}{singular value}{Seberapa panjang sebuah arah diregangkan oleh
matriks; nol berarti arah itu digepengkan sampai hilang.}

\begin{definisi}[Nilai singular]
Misalkan $\mathbf{A}\in\mathbb{R}^{m\times n}$. \emph{Nilai singular} $\mathbf{A}$
adalah akar kuadrat nilai eigen $\mathbf{A}^{\top}\mathbf{A}$, diurutkan menurun:
$\sigma_{1}\ge\sigma_{2}\ge\dots\ge\sigma_{n}\ge0$ dengan
$\sigma_{i}=\sqrt{\lambda_{i}(\mathbf{A}^{\top}\mathbf{A})}$.\index{nilai singular}
\end{definisi}

Definisi itu sah bila $\mathbf{A}^{\top}\mathbf{A}$ punya $n$ nilai eigen real dengan
vektor eigen ortonormal; untuk matriks simetris itu benar, dan kita nyatakan sebagai
lemma, dengan satu langkah yang jujur disebut belum dikerjakan.

\begin{lemma}[Spektral untuk matriks simetris]
Bila $\mathbf{S}\in\mathbb{R}^{n\times n}$ simetris, maka ada basis ortonormal
$\mathbf{v}_{1},\dots,\mathbf{v}_{n}$ dari $\mathbb{R}^{n}$ dan bilangan real
$\lambda_{1},\dots,\lambda_{n}$ dengan $\mathbf{S}\mathbf{v}_{i}=\lambda_{i}\mathbf{v}_{i}$.
\end{lemma}

\begin{proof}[Sketsa]
Fungsi $\mathbf{x}\mapsto\mathbf{x}^{\top}\mathbf{S}\mathbf{x}$ kontinu pada bola satuan
yang tertutup dan terbatas, jadi ia mencapai maksimum di suatu titik
$\mathbf{v}_{1}$ dengan norma satu. Dengan metode pengali Lagrange, titik itu memenuhi
$\mathbf{S}\mathbf{v}_{1}=\lambda_{1}\mathbf{v}_{1}$; nilai $\lambda_{1}$ adalah maksimum
tadi. Selanjutnya $\mathbf{v}_{1}^{\perp}$ juga invarian di bawah $\mathbf{S}$, karena
$\langle\mathbf{S}\mathbf{w},\mathbf{v}_{1}\rangle
=\langle\mathbf{w},\mathbf{S}\mathbf{v}_{1}\rangle
=\lambda_{1}\langle\mathbf{w},\mathbf{v}_{1}\rangle=0$ untuk $\mathbf{w}\perp\mathbf{v}_{1}$.
Batasan $\mathbf{S}$ pada $\mathbf{v}_{1}^{\perp}$ simetris dan berdimensi $n-1$, jadi
induksi memberi sisa basis ortonormalnya.

Langkah yang tidak dikerjakan di sini adalah bukti bahwa fungsi kontinu pada himpunan
tertutup dan terbatas mencapai maksimum. Itu teorema nilai ekstrem, dibuktikan di Bab 11,
dan bukti penuh lemma ini ada di Axler, \emph{Linear Algebra Done Right} edisi ke-4,
bagian 7B, atau di Hoffman dan Kunze bagian tentang operator simetris.
\end{proof}

Dengan lemma itu, $\mathbf{A}^{\top}\mathbf{A}$ simetris dan
$\mathbf{x}^{\top}\mathbf{A}^{\top}\mathbf{A}\mathbf{x}=\|\mathbf{A}\mathbf{x}\|^{2}\ge0$,
sehingga nilai eigennya tidak negatif dan akarnya real. Sekarang teorema utamanya.

\begin{teorema}[Dekomposisi nilai singular]
Setiap $\mathbf{A}\in\mathbb{R}^{m\times n}$ berpangkat $r$ dapat ditulis
\[
  \mathbf{A}=\mathbf{U}\boldsymbol{\Sigma}\mathbf{V}^{\top},
\]
dengan $\mathbf{U}\in\mathbb{R}^{m\times m}$ dan $\mathbf{V}\in\mathbb{R}^{n\times n}$
matriks berkolom ortonormal, serta $\boldsymbol{\Sigma}\in\mathbb{R}^{m\times n}$
bersifat diagonal dengan entri $\sigma_{1}\ge\dots\ge\sigma_{r}>0$ dan sisanya nol.
\index{SVD}
\end{teorema}

\begin{proof}
Karena $\mathbf{A}^{\top}\mathbf{A}$ simetris, lemma spektral memberi basis ortonormal
$\mathbf{v}_{1},\dots,\mathbf{v}_{n}$ dengan
$\mathbf{A}^{\top}\mathbf{A}\mathbf{v}_{i}=\sigma_{i}^{2}\mathbf{v}_{i}$ dan
$\sigma_{1}\ge\dots\ge\sigma_{n}\ge0$. Nilai eigen yang tidak nol berjumlah $r$;
sebut indeksnya $1,\dots,r$.

Untuk $i\le r$ definisikan $\mathbf{u}_{i}=\mathbf{A}\mathbf{v}_{i}/\sigma_{i}$, yang
bernorma satu karena $\|\mathbf{A}\mathbf{v}_{i}\|^{2}
=\mathbf{v}_{i}^{\top}\mathbf{A}^{\top}\mathbf{A}\mathbf{v}_{i}=\sigma_{i}^{2}$.
Untuk $i,j\le r$,
\[
  \langle\mathbf{u}_{i},\mathbf{u}_{j}\rangle
  =\frac{1}{\sigma_{i}\sigma_{j}}\mathbf{v}_{i}^{\top}\mathbf{A}^{\top}\mathbf{A}\mathbf{v}_{j}
  =\frac{\sigma_{j}^{2}}{\sigma_{i}\sigma_{j}}\langle\mathbf{v}_{i},\mathbf{v}_{j}\rangle
  =\delta_{ij},
\]
jadi $\mathbf{u}_{1},\dots,\mathbf{u}_{r}$ ortonormal. Untuk $i>r$ berlaku
$\|\mathbf{A}\mathbf{v}_{i}\|^{2}=\sigma_{i}^{2}=0$, sehingga $\mathbf{A}\mathbf{v}_{i}=\mathbf{0}$.

Vektor-vektor itu belum tentu mengisi $\mathbb{R}^{m}$; lengkapi menjadi basis ortonormal dengan
Gram-Schmidt pada sembarang basis komplemen ortogonalnya. Definisikan
$\mathbf{B}=\sum_{i\le r}\sigma_{i}\mathbf{u}_{i}\mathbf{v}_{i}^{\top}$. Untuk $j\le r$
berlaku $\mathbf{B}\mathbf{v}_{j}=\sigma_{j}\mathbf{u}_{j}=\mathbf{A}\mathbf{v}_{j}$, dan
untuk $j>r$ keduanya nol. Kedua matriks sepakat pada suatu basis, jadi $\mathbf{B}=\mathbf{A}$.
\end{proof}

Bacaan geometrinya: bola satuan dipetakan $\mathbf{A}$ menjadi elipsoid. Kolom
$\mathbf{v}_{i}$ menunjuk sumbu bola, $\mathbf{u}_{i}$ sumbu elipsoid, dan $\sigma_{i}$
panjang sumbu itu: $\mathbf{V}^{\top}$ memutar, $\boldsymbol{\Sigma}$ meregangkan,
$\mathbf{U}$ memutar lagi.

Gambar~\ref{fig:svd-lingkaran-elips} menjalankan bacaan itu pada
$\mathbf{A}=\begin{pmatrix}0&3\\4&1\end{pmatrix}$: sumbu $\mathbf{v}_{1}$ yang miring
$26{,}57^\circ$ berakhir sebagai sumbu elips $\mathbf{u}_{1}$ yang miring
$71{,}57^\circ$.

\begin{figure}[htbp]
\centering
\begin{tikzpicture}[scale=1.55]
  \draw[->,gray] (-1.3,0)--(1.3,0);
  \draw[->,gray] (0,-1.3)--(0,1.3);
  \draw[thick,accent] (0,0) circle (1);
  \draw[->,very thick,blue!55!black] (0,0)--(0.894,0.447)
    node[above right=-3pt] {$\mathbf{v}_{1}$};
  \draw[->,very thick,blue!55!black] (0,0)--(-0.447,0.894)
    node[above left=-3pt] {$\mathbf{v}_{2}$};
  \node[below,font=\small] at (0,-1.5) {lingkaran satuan};
\end{tikzpicture}
\hfill
\begin{tikzpicture}[scale=0.68]
  \draw[->,gray] (-3.5,0)--(3.5,0);
  \draw[->,gray] (0,-4.6)--(0,4.6);
  \draw[thick,accent] plot[domain=0:360,samples=160]
    ({3*sin(\x)},{4*cos(\x)+sin(\x)});
  \draw[->,very thick,blue!55!black] (0,0)--(1.342,4.025)
    node[right=3pt] {$\sigma_{1}\mathbf{u}_{1}$};
  \draw[->,very thick,blue!55!black] (0,0)--(2.683,-0.894)
    node[below right=-3pt] {$\sigma_{2}\mathbf{u}_{2}$};
  \node[below,font=\small] at (0,-5.1) {bayangannya};
\end{tikzpicture}
\caption{Aksi $\mathbf{A}=\begin{pmatrix}0&3\\4&1\end{pmatrix}$ pada lingkaran satuan.
Kiri: dua sumbu ortonormal lingkaran, $\mathbf{v}_{1}=(2,1)/\sqrt{5}$ dan
$\mathbf{v}_{2}=(-1,2)/\sqrt{5}$. Kanan: bayangannya elips dengan setengah-sumbu
$\sigma_{1}=3\sqrt{2}\approx4{,}243$ (numpy mencetak $4{,}242641$) sepanjang
$\mathbf{u}_{1}=(1,3)/\sqrt{10}$, dan $\sigma_{2}=2\sqrt{2}\approx2{,}828$ (numpy
$2{,}828427$) sepanjang $\mathbf{u}_{2}=(3,-1)/\sqrt{10}$.}
\label{fig:svd-lingkaran-elips}
\end{figure}

Bentuk penjumlahan $\mathbf{A}=\sum_{i}\sigma_{i}\mathbf{u}_{i}\mathbf{v}_{i}^{\top}$
adalah cara paling jujur membacanya: matriks diuraikan menjadi $r$ lapisan berpangkat
satu, diurutkan dari sumbangan terbesar. Memotong setelah $k$ suku memberi hampiran
berpangkat $k$ terbaik dalam norma Frobenius; itu teorema Eckart-Young, yang buktinya
tidak dikerjakan di sini dan dapat dibaca di Golub dan Van Loan bagian 2.4.

Gambar~\ref{fig:potret-singular} menampilkan potret nilai singular matriks yang sama:
memotong setelah suku pertama mematikan $\sigma_{2}$, dan galat yang tersisa tepat
$\sigma_{2}=2{,}828427$.

\begin{figure}[htbp]
\centering
\begin{tikzpicture}
\begin{axis}[
  width=9.6cm, height=5.4cm,
  ybar, bar width=15pt,
  xmin=0.35, xmax=2.65, ymin=0, ymax=6.4,
  ylabel={nilai singular},
  xtick={1,2}, xticklabels={$\sigma_{1}$,$\sigma_{2}$},
  tick label style={font=\small},
  label style={font=\small},
  axis lines=left,
  legend style={font=\small,draw=none,at={(0.03,0.97)},anchor=north west},
  legend cell align=left,
]
\addplot[fill=accent!55,draw=accent,bar shift=-8.5pt]
  coordinates {(1,4.243) (2,2.828)};
\addlegendentry{$\mathbf{A}$ penuh}
\addplot[fill=blue!25,draw=blue!55!black,bar shift=8.5pt]
  coordinates {(1,4.243) (2,0)};
\addlegendentry{$k=1$}
\draw[<->,gray,thick] (axis cs:2.30,0)--(axis cs:2.30,2.828)
  node[midway,right=2pt,font=\scriptsize] {$\sigma_{2}$};
\draw[gray,densely dotted] (axis cs:2.03,2.828)--(axis cs:2.30,2.828);
\end{axis}
\end{tikzpicture}
\caption{Potret nilai singular $\mathbf{A}=\begin{pmatrix}0&3\\4&1\end{pmatrix}$.
Bar merah: $\mathbf{A}$ penuh, $\sigma_{1}=4{,}243$ dan $\sigma_{2}=2{,}828$. Bar
biru: hampiran pangkat satu $\sigma_{1}\mathbf{u}_{1}\mathbf{v}_{1}^{\top}$, yang
menyisakan nol pada tempat kedua. Selisih ujung keduanya, yaitu $\sigma_{2}=2{,}828$,
adalah galat Frobenius hampiran itu.}
\label{fig:potret-singular}
\end{figure}

Beberapa akibat langsung, semuanya dipakai di bagian berikutnya.

\begin{akibat}
Untuk $\mathbf{A}\in\mathbb{R}^{m\times n}$ dengan nilai singular
$\sigma_{1}\ge\dots\ge\sigma_{r}>0$:
\begin{enumerate}
  \item $\operatorname{rank}\mathbf{A}=r$, yaitu banyaknya nilai singular tak nol;
  \item $\|\mathbf{A}\mathbf{x}\|\le\sigma_{1}\|\mathbf{x}\|$ untuk setiap $\mathbf{x}$,
        dan batas itu tercapai;
  \item $\|\mathbf{A}\|_{F}^{2}=\sum_{i}\sigma_{i}^{2}$ dan
        $\|\mathbf{A}\|_{2}=\sigma_{1}$;
  \item bila $\mathbf{A}$ persegi dan dapat dibalik,
        $|\det\mathbf{A}|=\prod_{i}\sigma_{i}$ dan
        $\kappa(\mathbf{A})=\sigma_{1}/\sigma_{n}$.
\end{enumerate}
\end{akibat}

\begin{proof}
Butir pertama: $\mathbf{A}\mathbf{x}=\sum_{i}\sigma_{i}\mathbf{u}_{i}
(\mathbf{v}_{i}^{\top}\mathbf{x})$, jadi jangkauan $\mathbf{A}$ direntang
$\mathbf{u}_{1},\dots,\mathbf{u}_{r}$ yang bebas linear. Butir kedua:
$\|\mathbf{A}\mathbf{x}\|^{2}=\sum_{i}\sigma_{i}^{2}(\mathbf{v}_{i}^{\top}\mathbf{x})^{2}
\le\sigma_{1}^{2}\|\mathbf{x}\|^{2}$, dengan kesamaan di $\mathbf{x}=\mathbf{v}_{1}$.
Butir ketiga: suku silang lenyap, jadi norma Frobenius kuadrat adalah jumlah
$\sigma_{i}^{2}$. Butir keempat: $|\det\mathbf{A}|=|\det\mathbf{U}\det\boldsymbol{\Sigma}
\det\mathbf{V}^{\top}|=\prod_{i}\sigma_{i}$, sebab determinan ortogonal bernilai $\pm1$.
\end{proof}

Bilangan $\kappa$ di butir keempat adalah \emph{bilangan kondisi}, nisbah nilai singular
terbesar dan terkecil; untuk matriks berpangkat kurang dari $n$ nisbahnya tak hingga.
$\kappa$ besar berarti satu arah hampir gepeng, sehingga galat kecil pada data menjadi
galat besar pada jawaban; bagian keenam memakainya untuk membandingkan dua cara
menyelesaikan kuadrat terkecil.\index{bilangan kondisi}

Satu peringatan: nilai singular bukan nilai eigen. Matriks
$\begin{pmatrix}0&1\\0&0\end{pmatrix}$ punya nilai eigen $0$ dan $0$, padahal nilai
singularnya $1$ dan $0$. Keduanya sama hanya bila matriksnya simetris dan semidefinit
positif.

\section{Pseudoinvers dan solusi norma minimum}
\tingkat{lanjut}

Sekarang kita kembali ke pertanyaan penyerang. Skornya satu bilangan, pencacahnya tiga,
dan himpunan perubahan yang memberi pergeseran sama berisi tak hingga jawaban. Semua
jawaban itu benar; yang membedakan mereka adalah besar perubahan, dan penyerang ingin
yang terkecil.

\begin{untukanak}
Bayangkan kamu harus pergi ke toko di ujung jalan. Ada banyak jalan: lewat pasar, lewat
lapangan, atau memutar lewat jembatan. Semuanya sampai di toko yang sama. Kalau yang
kamu hitung hanya ``sampai atau tidak'', semua jalan itu sama baiknya. Tetapi kalau kamu
menghitung langkah kaki, hanya satu jalan yang paling pendek, dan jalan itu biasanya
yang paling lurus.

Ingat dua penggaris yang bedanya seperseratus sentimeter? Selisih sekecil itu mudah
tenggelam dalam kekeliruan membaca. Karena itu dua jalan yang sama secara matematis bisa
memberi jawaban yang sangat berbeda ketika angkanya dihitung mesin.

Di bab ini kita akan tahu mengapa dari banyak jalan ada satu yang terpendek, mengapa
jalan terpendek itu selalu tegak lurus terhadap semua jalan lain yang tidak mengubah
tujuan, dan mengapa mesin kadang memilih jalan yang tidak paling langsung.
\end{untukanak}

Alat yang menjawabnya adalah invers yang diperumum. $\mathbf{A}$ belum tentu persegi, jadi
belum tentu punya invers. Dekomposisi nilai singular memberi cara membalik ``sebanyak yang
bisa dibalik'': balik setiap $\sigma_{i}$ yang tidak nol, biarkan arah yang sudah hilang
tetap hilang.

\kamus{pseudoinvers}{pseudoinverse}{Invers yang tetap bekerja ketika sebagian arah sudah
gepeng: yang bisa dibalik, dibalik; yang hilang dibiarkan hilang.}

\begin{definisi}[Pseudoinvers]
Misalkan $\mathbf{A}=\mathbf{U}\boldsymbol{\Sigma}\mathbf{V}^{\top}$ berpangkat $r$.
\emph{Pseudoinvers} $\mathbf{A}^{+}\in\mathbb{R}^{n\times m}$ didefinisikan sebagai
\[
  \mathbf{A}^{+}=\mathbf{V}\boldsymbol{\Sigma}^{+}\mathbf{U}^{\top},
  \qquad
  \boldsymbol{\Sigma}^{+}=\operatorname{diag}(1/\sigma_{1},\dots,1/\sigma_{r},0,\dots,0),
\]
yaitu matriks $\boldsymbol{\Sigma}$ yang ditranspos lalu setiap entri tak nolnya
dibalik.\index{pseudoinvers}
\end{definisi}

$\mathbf{A}^{+}$ selalu ada, bahkan untuk matriks nol. Bila $\mathbf{A}$ persegi dan
dapat dibalik, semua $\sigma_{i}>0$ dan $\mathbf{A}^{+}=\mathbf{A}^{-1}$; definisi ini
memperluas invers, bukan menggantinya.

Vektor $\mathbf{A}^{+}\mathbf{b}$ menjawab dua tuntutan sekaligus: sisanya terkecil, dan
di antara yang sisanya terkecil, panjangnya paling pendek.

\begin{teorema}
Untuk setiap $\mathbf{b}\in\mathbb{R}^{m}$, vektor
$\mathbf{x}^{+}=\mathbf{A}^{+}\mathbf{b}$ memenuhi
$\|\mathbf{A}\mathbf{x}^{+}-\mathbf{b}\|\le\|\mathbf{A}\mathbf{x}-\mathbf{b}\|$ untuk
semua $\mathbf{x}$, dan di antara semua vektor yang memenuhi ketaksamaan itu,
$\mathbf{x}^{+}$ yang normanya terkecil.
\end{teorema}

\begin{proof}
Tulis $\mathbf{A}=\mathbf{U}\boldsymbol{\Sigma}\mathbf{V}^{\top}$ dan misalkan
$\mathbf{y}=\mathbf{V}^{\top}\mathbf{x}$. Karena $\mathbf{V}$ persegi dan berkolom
ortonormal, $\|\mathbf{x}\|=\|\mathbf{y}\|$ dan pemetaan $\mathbf{x}\mapsto\mathbf{y}$
bijektif. Uraikan $\mathbf{b}$ menurut kolom $\mathbf{U}$:
$\mathbf{b}=\sum_{i=1}^{r}b_{i}\mathbf{u}_{i}+\mathbf{b}_{\perp}$ dengan
$b_{i}=\langle\mathbf{u}_{i},\mathbf{b}\rangle$ dan $\mathbf{b}_{\perp}$ ortogonal
terhadap semua $\mathbf{u}_{i}$. Maka
\[
  \mathbf{A}\mathbf{x}-\mathbf{b}
  =\sum_{i=1}^{r}(\sigma_{i}y_{i}-b_{i})\mathbf{u}_{i}+\sum_{i=r+1}^{m}0\cdot\mathbf{u}_{i}
   -\mathbf{b}_{\perp},
\]
dan ketiga kelompok suku itu saling ortogonal. Karena itu
\[
  \|\mathbf{A}\mathbf{x}-\mathbf{b}\|^{2}
  =\sum_{i=1}^{r}(\sigma_{i}y_{i}-b_{i})^{2}+\|\mathbf{b}_{\perp}\|^{2}.
\]
Suku terakhir tidak bergantung pada $\mathbf{x}$, jadi minimum dicapai dengan
$y_{i}=b_{i}/\sigma_{i}$ untuk $i\le r$. Koordinat $y_{r+1},\dots,y_{n}$ tidak muncul
di rumus mana pun: berapa pun nilainya, sisa kuadratnya sama. Tetapi
$\|\mathbf{y}\|^{2}=\sum_{i\le r}(b_{i}/\sigma_{i})^{2}+\sum_{i>r}y_{i}^{2}$, sehingga
norma minimum tercapai ketika $y_{i}=0$ untuk $i>r$. Vektor yang dihasilkan adalah
$\mathbf{x}^{+}=\sum_{i\le r}(b_{i}/\sigma_{i})\mathbf{v}_{i}
=\mathbf{V}\boldsymbol{\Sigma}^{+}\mathbf{U}^{\top}\mathbf{b}$.
\end{proof}

Bukti itu memuat jawaban pertanyaan penyerang. Pada sistem kurang tentsu, $n>r$, dan
koordinat $y_{r+1},\dots,y_{n}$ tidak mempengaruhi sisa sama sekali; koordinat itu
berada di sepanjang $\mathbf{v}_{r+1},\dots,\mathbf{v}_{n}$, tepat vektor yang merentang
inti $\mathbf{A}$. Jadi himpunan semua penyelesaian adalah $\mathbf{x}^{+}+\ker\mathbf{A}$,
ruang affine berdimensi $n-r$, dan pseudoinvers memilih titik dengan komponen inti nol:
satu-satunya yang ortogonal terhadap inti, yaitu di ruang baris $\mathbf{A}$.
Setiap penyelesaian lain adalah $\mathbf{x}^{+}$ ditambah vektor tak nol di inti, dan
penambahan itu menambah panjang menurut Pythagoras:
$\|\mathbf{x}^{+}+\mathbf{k}\|^{2}=\|\mathbf{x}^{+}\|^{2}+\|\mathbf{k}\|^{2}$.

Untuk dua kasus ekstrem ada rumus tertutup.

\begin{akibat}
Bila $\mathbf{A}$ berpangkat kolom penuh, $\mathbf{A}^{+}
=(\mathbf{A}^{\top}\mathbf{A})^{-1}\mathbf{A}^{\top}$. Bila $\mathbf{A}$ berpangkat
baris penuh, $\mathbf{A}^{+}=\mathbf{A}^{\top}(\mathbf{A}\mathbf{A}^{\top})^{-1}$.
\end{akibat}

\begin{proof}
Kasus pertama kita ulang alasannya dengan SVD. Ambil
$\mathbf{A}=\mathbf{U}\boldsymbol{\Sigma}\mathbf{V}^{\top}$ dengan $\mathbf{V}$ persegi
ortonormal; maka
\[
  (\mathbf{A}^{\top}\mathbf{A})^{-1}\mathbf{A}^{\top}
  =\mathbf{V}\boldsymbol{\Sigma}^{-2}\mathbf{V}^{\top}\mathbf{V}\boldsymbol{\Sigma}\mathbf{U}^{\top}
  =\mathbf{V}\boldsymbol{\Sigma}^{-1}\mathbf{U}^{\top}=\mathbf{A}^{+}.
\]
Kasus kedua diperoleh dengan menerapkan kasus pertama pada $\mathbf{A}^{\top}$, lalu
mentranspos: $(\mathbf{A}^{\top})^{+}=(\mathbf{A}\mathbf{A}^{\top})^{-1}\mathbf{A}$,
dan transposnya $\mathbf{A}^{\top}(\mathbf{A}\mathbf{A}^{\top})^{-1}$. Karena
$\mathbf{A}^{+}=(\mathbf{A}^{\top})^{+}{}^{\top}$ menurut definisi, hasilnya benar.
\end{proof}

Akibat kedua itu yang menjawab masalah penyerang. Matriks detektor $2\times3$ berpangkat
baris penuh, jadi perubahan minimum adalah
$\mathbf{A}^{\top}(\mathbf{A}\mathbf{A}^{\top})^{-1}\boldsymbol{\delta}$, dan rumus itu
hanya membutuhkan invers matriks $2\times2$. Perhatikan bahwa
$\mathbf{A}\mathbf{A}^{\top}$ dapat dibalik justru karena barisnya bebas linear, bukan
kolomnya.

\begin{tantangan}
Ambil
\[
  \mathbf{A}=\begin{pmatrix}1&0\\0&1\\1&1\end{pmatrix},\qquad \mathbf{b}=(1,1,3).
\]
Sistem $\mathbf{A}\mathbf{x}=\mathbf{b}$ tidak punya solusi: dua baris pertama memaksa
$\mathbf{x}=(1,1)$, sedangkan baris ketiga menuntut $x_{1}+x_{2}=3$.

Tetap ada satu penyelesaian kuadrat terkecil, karena kolomnya bebas linear. Tentukan $\hat{\mathbf{x}}$, lalu cari semua nilai $c$ sehingga untuk $\mathbf{b}=(1,1,c)$ panjang sisa kuadrat terkecilnya tepat $1$. Ada dua jawaban, dan keduanya simetris terhadap $c=2$.

Petunjuk 1: sisa selalu sejajar satu arah tetap; cari arah itu lebih dulu.

Petunjuk 2: jarak $\mathbf{b}$ ke bidang kolom adalah $|c-2|/\sqrt{3}$, jadi tinggal menyelesaikan satu persamaan.

Kalau mentok, kembali ke bagian Kuadrat terkecil dan persamaan normal.
\end{tantangan}

\section{Mengapa SVD lebih aman daripada persamaan normal}
\tingkat{lanjut}

Persamaan normal dan pseudoinvers memberi jawaban yang sama secara matematis. Secara
numerik keduanya berbeda jauh, dan alasannya satu baris: persamaan normal bekerja pada
$\mathbf{A}^{\top}\mathbf{A}$, yang nilai singularnya $\sigma_{i}^{2}$, bukan
$\sigma_{i}$. Karena bilangan kondisi adalah nisbah nilai singular, ia ikut dikuadratkan.

\begin{akibat}
Bila $\mathbf{A}$ berpangkat kolom penuh, $\kappa(\mathbf{A}^{\top}\mathbf{A})
=\kappa(\mathbf{A})^{2}$.
\end{akibat}

\begin{proof}
Dari $\mathbf{A}=\mathbf{U}\boldsymbol{\Sigma}\mathbf{V}^{\top}$ diperoleh
$\mathbf{A}^{\top}\mathbf{A}=\mathbf{V}\boldsymbol{\Sigma}^{2}\mathbf{V}^{\top}$, dengan
$\mathbf{V}$ ortogonal. Jadi nilai singular $\mathbf{A}^{\top}\mathbf{A}$ adalah
$\sigma_{i}^{2}$, dan nisbah terbesar terhadap terkecil menjadi
$\sigma_{1}^{2}/\sigma_{n}^{2}=(\sigma_{1}/\sigma_{n})^{2}$.
\end{proof}

Artinya begini. Dalam pecahan biner dengan sekitar enam belas digit, galat relatif
jawaban kira-kira $\kappa$ kali galat relatif data: pada $\kappa=10^{7}$ tersisa sembilan
digit, pada $\kappa^{2}=10^{14}$ tersisa dua, dan pada $\kappa=10^{8}$ seluruh digit
habis.

Angka pada tabel berikut dihitung dengan numpy dan scipy, dengan acuan 50 digit dari
mpmath. Baris pertama adalah kalibrasi probe dari pembuka bab; baris kedua adalah
pencocokan polinomial derajat sembilan pada dua belas titik, masalah yang tampak tidak
berbahaya.

\begin{center}
\begin{tabular}{lrrr}
\hline
kasus & $\kappa(\mathbf{A})$ & galat relatif persamaan normal & galat relatif SVD \\
\hline
kalibrasi $5\times2$ & $7{,}36\times10^{3}$ & $9{,}10\times10^{-11}$ & $4{,}72\times10^{-13}$ \\
polinomial $12\times10$ & $6{,}84\times10^{6}$ & $2{,}87\times10^{-4}$ & $4{,}90\times10^{-12}$ \\
\hline
\end{tabular}
\end{center}

Kedua baris berpola sama. Pada baris pertama, $\kappa=7{,}36\times10^{3}$ dan
$\kappa^{2}=5{,}42\times10^{7}$; persamaan normal berselisih $9{,}10\times10^{-11}$ dari
acuan, SVD $4{,}72\times10^{-13}$, rasionya $193$. Pada baris kedua,
$\kappa(\mathbf{A})=6{,}84\times10^{6}$ dan $\kappa(\mathbf{A}^{\top}\mathbf{A})
=4{,}68\times10^{13}$; persamaan normal berselisih $2{,}87\times10^{-4}$, SVD
$4{,}90\times10^{-12}$, rasionya $5{,}9\times10^{7}$. Selisih itu membesar bersama
$\kappa$, jauh lebih cepat pada persamaan normal; keduanya masih di bawah batas
$\kappa^{2}\epsilon$ karena batas itu kasus terburuk, bukan ramalan. Yang penting
urutannya, bukan angka mutlaknya.

Sisa kuadrat kedua jalan nyaris sama, $4{,}47019\times10^{-6}$ versus
$4{,}47018\times10^{-6}$, padahal koefisiennya berbeda jauh: sisa kecil tidak membuktikan
koefisien benar.

Kesimpulan yang jujur: SVD tidak menghapus kepekaan, ia hanya menghilangkan satu sumber
perbesaran, yaitu pengkuadratan, sehingga galat tumbuh sebesar $\kappa$ dan bukan
$\kappa^{2}$. Bila $\kappa$ mendekati $10^{16}$, tidak ada presisi ganda yang
menyelamatkan; yang tersisa adalah regularisasi, yaitu sengaja membuang arah dengan
$\sigma_{i}$ kecil, setara menghampiri $\mathbf{A}$ dengan matriks berpangkat lebih
rendah. Bab 26 membahas pilihan ambangnya.

\section{Contoh hitung}
\tingkat{menengah}

Empat hitungan berikut memakai pecahan sederhana, jadi dapat diperiksa ulang dengan
tangan tanpa mesin.

\emph{Pertama, proyeksi pada bidang.} Ambil
$\mathbf{a}_{1}=(1,1,0)$, $\mathbf{a}_{2}=(1,0,1)$, dan $\mathbf{b}=(3,1,3)$ di
$\mathbb{R}^{3}$. Subruang $W=\operatorname{span}(\mathbf{a}_{1},\mathbf{a}_{2})$
berdimensi dua. Gram-Schmidt memberi
\[
  \mathbf{q}_{1}=\frac{1}{\sqrt{2}}(1,1,0),\qquad
  \mathbf{u}_{2}=\mathbf{a}_{2}-\frac{1}{\sqrt{2}}\mathbf{q}_{1}
  =\Bigl(\tfrac12,-\tfrac12,1\Bigr),\qquad
  \mathbf{q}_{2}=\frac{1}{\sqrt{6}}(1,-1,2),
\]
karena $\|\mathbf{u}_{2}\|=\sqrt{3/2}=\sqrt{6}/2$. Koefisien proyeksi:
$\langle\mathbf{b},\mathbf{q}_{1}\rangle=4/\sqrt{2}=2\sqrt{2}$ dan
$\langle\mathbf{b},\mathbf{q}_{2}\rangle=8/\sqrt{6}$. Jadi
\[
  \mathbf{p}=2\sqrt{2}\,\mathbf{q}_{1}+\frac{8}{\sqrt{6}}\,\mathbf{q}_{2}
  =(2,2,0)+\frac{4}{3}(1,-1,2)
  =\Bigl(\tfrac{10}{3},\tfrac23,\tfrac83\Bigr).
\]
Sisanya $\mathbf{r}=\mathbf{b}-\mathbf{p}=(-\tfrac13,\tfrac13,\tfrac13)$, dengan
$\|\mathbf{r}\|=1/\sqrt{3}=0{,}577350$. Pemeriksaan: $\mathbf{a}_{1}^{\top}\mathbf{r}=0$
dan $\mathbf{a}_{2}^{\top}\mathbf{r}=0$, jadi sisa memang tegak lurus kedua pembangun.
Jarak $\mathbf{b}$ ke bidang itu $0{,}577$; tidak ada titik lain di $W$ yang lebih dekat.

\emph{Kedua, kuadrat terkecil $5\times2$.} Untuk data pembuka bab,
$\mathbf{A}$ berkolom $(1,1,1,1,1)^{\top}$ dan $(100,101,102,103,104)^{\top}$, sehingga
\[
  \mathbf{A}^{\top}\mathbf{A}=
  \begin{pmatrix}5&510\\ 510&52030\end{pmatrix},
  \qquad
  \mathbf{A}^{\top}\mathbf{y}=\begin{pmatrix}1540\\ 157109{,}7\end{pmatrix},
  \qquad
  \det(\mathbf{A}^{\top}\mathbf{A})=5\cdot52030-510^{2}=50 .
\]
Aturan Cramer memberi
\begin{align*}
  a&=\frac{1540\cdot52030-510\cdot157109{,}7}{50}=\frac{253}{50}=5{,}06,\\
  b&=\frac{5\cdot157109{,}7-510\cdot1540}{50}=\frac{148{,}5}{50}=2{,}97 .
\end{align*}
Ramalan bacaan menjadi $302{,}06$, $305{,}03$, $308{,}00$, $310{,}97$, dan $313{,}94$,
sedangkan data sebenarnya $302{,}0$, $305{,}1$, $307{,}9$, $311{,}2$, $313{,}8$. Sisanya
\[
  \mathbf{r}=(-0{,}06,\;0{,}07,\;-0{,}10,\;0{,}23,\;-0{,}14),
  \qquad
  \|\mathbf{r}\|^{2}=0{,}0910,\qquad \|\mathbf{r}\|=0{,}301662 .
\]
Dua pemeriksaan yang harus lolos bila hitungan benar: $\sum_{i}r_{i}=0$ dan
$\sum_{i}t_{i}r_{i}=0$. Yang pertama karena $\mathbf{r}$ tegak lurus kolom satuan, yang
kedua karena tegak lurus kolom $t$. Keduanya memang nol di sini.

\emph{Ketiga, SVD matriks $3\times2$.} Ambil
$\mathbf{A}=\begin{pmatrix}1&1\\ 1&1\\ 1&-1\end{pmatrix}$. Maka
\[
  \mathbf{A}^{\top}\mathbf{A}=\begin{pmatrix}3&1\\ 1&3\end{pmatrix},
  \qquad
  \lambda^{2}-6\lambda+8=0\;\Longrightarrow\;\lambda=4,\,2,
\]
jadi $\sigma_{1}=2$ dan $\sigma_{2}=\sqrt{2}=1{,}41421$, dengan
$\kappa(\mathbf{A})=\sqrt{2}$. Vektor singular kanannya
$\mathbf{v}_{1}=\frac{1}{\sqrt{2}}(1,1)$ dan
$\mathbf{v}_{2}=\frac{1}{\sqrt{2}}(1,-1)$; vektor singular kirinya dihitung dari
$\mathbf{u}_{i}=\mathbf{A}\mathbf{v}_{i}/\sigma_{i}$, yaitu
$\mathbf{u}_{1}=\frac{1}{\sqrt{2}}(1,1,0)$ dan $\mathbf{u}_{2}=(0,0,1)$. Periksa satu entri:
$(2\,\mathbf{u}_{1}\mathbf{v}_{1}^{\top}+\sqrt{2}\,\mathbf{u}_{2}\mathbf{v}_{2}^{\top})_{32}
=-\sqrt{2}/\sqrt{2}=-1$, cocok dengan $\mathbf{A}$.

\emph{Keempat, matriks berpangkat satu.} Ambil
$\mathbf{B}=\begin{pmatrix}1&2\\ 2&4\\ 3&6\end{pmatrix}$, yang kolom keduanya dua kali
kolom pertama. Maka $\mathbf{B}^{\top}\mathbf{B}=14\begin{pmatrix}1&2\\ 2&4\end{pmatrix}$
dengan nilai eigen $70$ dan $0$, jadi $\sigma_{1}=\sqrt{70}=8{,}36660$ dan
$\sigma_{2}=0$. Pangkatnya satu, dan $\kappa$ tak hingga. Untuk $\mathbf{b}=(1,0,0)$,
proyeksi pada ruang kolom adalah $\frac{1}{14}(1,2,3)$, sehingga $\mathbf{B}\mathbf{x}$
harus sama dengan vektor itu. Karena $\mathbf{B}(s,2s)=5s(1,2,3)$, kita perlu
$5s=1/14$, yaitu $s=1/70$. Solusi norma minimumnya
\[
  \mathbf{x}^{+}=\Bigl(\tfrac{1}{70},\tfrac{2}{70}\Bigr)=(0{,}0142857,\;0{,}0285714),
  \qquad
  \|\mathbf{x}^{+}\|=\frac{\sqrt{5}}{70}=0{,}031944 .
\]
Sisa kuadrat terkecil:
\[
  \mathbf{r}=\mathbf{b}-\mathbf{B}\mathbf{x}^{+}
  =\Bigl(\tfrac{13}{14},-\tfrac17,-\tfrac{3}{14}\Bigr),
  \qquad
  \|\mathbf{r}\|=\sqrt{13/14}=0{,}963624 .
\]
Sisa itu tidak mungkin nol: $\mathbf{b}$ di luar ruang kolom. Solusi lain tetap memberi
sisa yang sama, misalnya $\mathbf{x}^{+}+(-\tfrac27,\tfrac17)$, tetapi normanya lebih
besar.

\kertas{Ulangi hitungan kedua dengan bacaan terakhir $314{,}3$ dan bukan $313{,}8$.
Ramalkan lebih dulu apakah $a$ naik atau turun, baru hitung. Perhatikan bahwa titik
dengan $t$ terjauh menggeser garis lebih banyak daripada titik di tengah; jelaskan lewat
kolom $\mathbf{A}$ dan rumus $(\mathbf{A}^{\top}\mathbf{A})^{-1}\mathbf{A}^{\top}$.
Hitung juga $\|\mathbf{r}\|$ yang baru dan bandingkan dengan $0{,}301662$.}

\section{Lab komputasi}
\tingkat{lanjut}

Dua blok berikut dijalankan dengan numpy 2.5.2, scipy 1.18.0, Python 3.14.4; keluarannya
disalin apa adanya.

\begin{lab}
\begin{verbatim}
import numpy as np
from scipy.linalg import lstsq, qr

np.set_printoptions(precision=6, suppress=True, linewidth=72)

# 1. proyeksi pada bidang
a1 = np.array([1., 1., 0.])
a2 = np.array([1., 0., 1.])
b = np.array([3., 1., 3.])
Q, _ = qr(np.column_stack([a1, a2]), mode="economic")
P = Q @ Q.T
print("Q =\n", Q)
print("P b =", P @ b, " sisa =", b - P @ b)
nr = np.linalg.norm(b - P @ b)
print("norma sisa =", nr, " 1/sqrt3 =", 1/np.sqrt(3))
print("P^2 = P:", np.allclose(P @ P, P), " P^T = P:", np.allclose(P.T, P))

# 2. kuadrat terkecil 5x2
t = np.array([100., 101., 102., 103., 104.])
A = np.column_stack([np.ones(5), t])
y = np.array([302.0, 305.1, 307.9, 311.2, 313.8])
print("\nA^T A =\n", A.T @ A, " det =", np.linalg.det(A.T @ A))
print("A^T y =", A.T @ y)
xn = np.linalg.solve(A.T @ A, A.T @ y)
xs = lstsq(A, y)[0]
print("persamaan normal:", xn)
print("lstsq           :", xs)
print("selisih solusi  :", np.linalg.norm(xn - xs))
cA = np.linalg.cond(A)
print("cond A =", cA, " cond A^T A =", np.linalg.cond(A.T @ A))
r = y - A @ xn
print("sisa =", r, " norma =", np.linalg.norm(r))
print("A^T r =", A.T @ r)

# 3. SVD sebuah matriks 3x2
S = np.array([[1., 1.], [1., 1.], [1., -1.]])
U, s, Vt = np.linalg.svd(S, full_matrices=False)
print("\nsingular:", s)
print("rank", np.linalg.matrix_rank(S), " cond", np.linalg.cond(S))
print("U =\n", U, "\nV^T =\n", Vt)
print("galat rekonstruksi:", np.linalg.norm(U @ np.diag(s) @ Vt - S))
\end{verbatim}
Keluaran sebenarnya:
\begin{verbatim}
Q =
 [[-0.707107  0.408248]
 [-0.707107 -0.408248]
 [-0.        0.816497]]
P b = [3.333333 0.666667 2.666667]  sisa = [-0.333333  0.333333  0.333333]
norma sisa = 0.5773502691896251  1/sqrt3 = 0.5773502691896258
P^2 = P: True  P^T = P: True

A^T A =
 [[    5.   510.]
 [  510. 52030.]]  det = 49.999999999983494
A^T y = [  1540.  157109.7]
persamaan normal: [5.06 2.97]
lstsq           : [5.06 2.97]
selisih solusi  : 5.315035090945416e-10
cond A = 7358.860135917719  cond A^T A = 54152822.48728771
sisa = [-0.06  0.07 -0.1   0.23 -0.14]  norma = 0.30166206257998296
A^T r = [0. 0.]

singular: [2.       1.414214]
rank 2  cond 1.414213562373095
U =
 [[-0.707107 -0.      ]
 [-0.707107 -0.      ]
 [-0.        1.      ]] 
V^T =
 [[-0.707107 -0.707107]
 [ 0.707107 -0.707107]]
galat rekonstruksi: 6.568167990716596e-16
\end{verbatim}
\end{lab}

Tiga hal perlu dibaca. Pertama, $\mathbf{Q}$ dari scipy berbeda tanda dari hasil tangan:
kolomnya $-(1,1,0)/\sqrt{2}$ dan $-(1,-1,2)/\sqrt{6}$. Tanda kolom pada QR tidak tunggal,
dan $\mathbf{Q}\mathbf{Q}^{\top}$ tidak berubah karenanya; entri $\mathbf{P}$ tetap
$(10/3,2/3,8/3)$, dengan sisa $(-0.333333,$ $0.333333,$ $0.333333)$.

Kedua, nilai $\det(\mathbf{A}^{\top}\mathbf{A})$ dilaporkan $49{,}999999999983494$ dan bukan
$50$; selisih $1{,}7\times10^{-11}$ itu murni pembulatan. Yang penting adalah
$\kappa(\mathbf{A}^{\top}\mathbf{A})=5{,}415\times10^{7}$, kuadrat dari
$\kappa(\mathbf{A})=7{,}359\times10^{3}$.

Ketiga, $\mathbf{A}^{\top}\mathbf{r}$ dilaporkan $(0,0)$ sampai enam digit dan selisih
kedua solusi $5{,}32\times10^{-10}$: keduanya baru berselisih di digit kesembilan, dan itu
bukan jaminan.

\begin{lab}
\begin{verbatim}
import numpy as np
import mpmath as mp
from scipy.linalg import lstsq

mp.mp.dps = 50
np.set_printoptions(precision=4, suppress=True, linewidth=72)

# 4. peringkat satu: sisa bukan nol, solusi norma minimum
B = np.array([[1., 2.], [2., 4.], [3., 6.]])
b2 = np.array([1., 0., 0.])
sB = np.linalg.svd(B, compute_uv=False)
print("singular B:", sB, " rank", np.linalg.matrix_rank(B))
xb = np.linalg.pinv(B) @ b2
print("x+ =", xb, " B x+ =", B @ xb)
nr = np.linalg.norm(B @ xb - b2)
print("norma sisa =", nr, " sqrt(13/14) =", np.sqrt(13/14))

# 5. sistem kurang tentsu: norma minimum
C = np.array([[1., 1., 1.], [0., 1., 1.]])
b3 = np.array([3., 1.])
xp = np.linalg.pinv(C) @ b3
print("\nx+ =", xp, " C x+ =", C @ xp, " norma =", np.linalg.norm(xp))
print("norma (2,1,0) =", np.linalg.norm([2., 1., 0.]))

# 6. kuadrat terkecil: persamaan normal vs lstsq vs acuan 50 digit
x = np.linspace(0, 1, 12)
V = np.vander(x, 10, increasing=True)
b = V @ (0.1 * np.arange(1, 11)) + 0.001 * np.sin(np.arange(12))
Vm, bm = mp.matrix(V.tolist()), mp.matrix(b.tolist())
ref = np.array([float(v) for v in mp.lu_solve(Vm.T * Vm, Vm.T * bm)])
xne = np.linalg.solve(V.T @ V, V.T @ b)
xls = lstsq(V, b)[0]
rel = lambda u: np.linalg.norm(u - ref) / np.linalg.norm(ref)
print("\ncond V     :", np.linalg.cond(V))
print("cond V^T V :", np.linalg.cond(V.T @ V))
print("galat relatif persamaan normal:", rel(xne))
print("galat relatif lstsq           :", rel(xls))
print("sisa persamaan normal:", np.linalg.norm(V @ xne - b))
print("sisa lstsq           :", np.linalg.norm(V @ xls - b))

# 7. Gram-Schmidt klasik vs modifikasi vs Householder
def cgs(M):
    m, n = M.shape
    Q = np.zeros((m, n)); R = np.zeros((n, n))
    for j in range(n):
        v = M[:, j].astype(float).copy()
        for i in range(j):
            R[i, j] = Q[:, i] @ M[:, j]
            v -= R[i, j] * Q[:, i]
        R[j, j] = np.linalg.norm(v)
        Q[:, j] = v / R[j, j]
    return Q

def mgs(M):
    m, n = M.shape
    Q = np.zeros((m, n)); R = np.zeros((n, n)); W = M.astype(float).copy()
    for j in range(n):
        for i in range(j):
            R[i, j] = Q[:, i] @ W[:, j]
            W[:, j] -= R[i, j] * Q[:, i]
        R[j, j] = np.linalg.norm(W[:, j])
        Q[:, j] = W[:, j] / R[j, j]
    return Q

ort = lambda Q: np.linalg.norm(Q.T @ Q - np.eye(Q.shape[1]))
Qh, _ = np.linalg.qr(V)
print("\n||Q^T Q - I||  klasik     :", ort(cgs(V)))
print("||Q^T Q - I||  modifikasi :", ort(mgs(V)))
print("||Q^T Q - I||  Householder:", ort(Qh))
\end{verbatim}
Keluaran sebenarnya:
\begin{verbatim}
singular B: [8.3666 0.    ]  rank 1
x+ = [0.0143 0.0286]  B x+ = [0.0714 0.1429 0.2143]
norma sisa = 0.9636241116594315  sqrt(13/14) = 0.9636241116594315

x+ = [2.  0.5 0.5]  C x+ = [3. 1.]  norma = 2.121320343559642
norma (2,1,0) = 2.23606797749979

cond V     : 6841129.486331141
cond V^T V : 46788641042301.25
galat relatif persamaan normal: 0.0002869502788738492
galat relatif lstsq           : 4.899523677448492e-12
sisa persamaan normal: 4.4701872339775605e-06
sisa lstsq           : 4.47017846708843e-06

||Q^T Q - I||  klasik     : 0.020506481578936916
||Q^T Q - I||  modifikasi : 3.397900311236548e-10
||Q^T Q - I||  Householder: 1.1933579412809552e-15
\end{verbatim}
\end{lab}

Bacaan pertama: sisa kuadrat terkecil matriks berpangkat satu besarnya persis
$\sqrt{13/14}$; numpy memberi angka yang sama sampai enam belas digit, jadi rumus pecahan
di bagian ketujuh bukan hampiran.

Bacaan kedua: $\mathbf{C}\mathbf{x}=(3,1)$ memberi
$\mathbf{x}^{+}=(2,\,0{,}5,\,0{,}5)$ bernorma $2{,}12132$, sedangkan $(2,1,0)$ memberi
$2{,}23607$; selisih kuadrat normanya $0{,}5$ tepat $2t^{2}$ pada $t=0{,}5$ untuk
$\mathbf{k}=(0,1,-1)$. Pseudoinvers memilih $t=0$.

Bacaan ketiga dan keempat: pada Vandermonde $12\times10$, persamaan normal berselisih
$2{,}87\times10^{-4}$ dari acuan 50 digit dan \texttt{lstsq} $4{,}90\times10^{-12}$,
padahal sisanya praktis sama ($4{,}47019\times10^{-6}$ versus $4{,}47018\times10^{-6}$).
Koefisien salah, sisa benar. Pada Gram-Schmidt,
klasik memberi $2{,}05\times10^{-2}$, modifikasi $3{,}40\times10^{-10}$, Householder
$1{,}19\times10^{-15}$.

\begin{untukdosen}
Urutan yang berhasil dipakai: proyeksi dulu, definisi hasilkali dalam kemudian.
Mahasiswa yang melihat bayangan lebih dulu menerima $\mathbf{Q}\mathbf{Q}^{\top}$ sebagai
jawaban atas pertanyaan ``titik mana yang terdekat'', bukan rumus ajaib; sesudah itu
Gram-Schmidt menjadi kebutuhan, bukan hafalan.

Dua miskonsepsi khas. Pertama, $\mathbf{A}^{\top}\mathbf{A}$ dianggap selalu dapat dibalik
karena ``kuadrat matriks tidak pernah nol''; pakai akibat di bagian ketiga, lalu minta
inti $\mathbf{A}^{\top}\mathbf{A}$ untuk matriks $2\times2$ berpangkat satu. Kedua, SVD
dianggap memberi nilai eigen; matriks $\begin{pmatrix}0&1\\0&0\end{pmatrix}$ cukup untuk
mematahkannya.

Aktivitas papan tulis: beri lima titik yang jelas tidak segaris, minta separuh kelas
memproyeksikan $\mathbf{b}$ dengan tangan dan separuh memakai \texttt{lstsq}; bandingkan
sisa vektornya, bukan koefisiennya.

Tanda kesiapan lanjut: mahasiswa dapat menjelaskan mengapa menambah satu persamaan tidak
selalu memperbaiki jawaban, dan menyebut satu situasi di mana pseudoinvers lebih baik
daripada invers.
\end{untukdosen}

\section{Ringkasan}

\begin{ringkasan}
\begin{itemize}
  \item Hasilkali dalam memberi panjang dan sudut; Cauchy-Schwarz menopang ketaksamaan
        segitiga. Pada basis ortonormal
        $\mathbf{v}=\sum_{i}\langle\mathbf{v},\mathbf{q}_{i}\rangle\mathbf{q}_{i}$ dan
        $\|\mathbf{v}\|^{2}=\sum_{i}\langle\mathbf{v},\mathbf{q}_{i}\rangle^{2}$.
  \item Proyeksi $\mathbf{p}=\mathbf{Q}\mathbf{Q}^{\top}\mathbf{b}$ adalah titik terdekat
        di subruang, dan sisanya tegak lurus seluruh subruang. Gram-Schmidt membangun
        basis ortonormal dan memberi $\mathbf{A}=\mathbf{Q}\mathbf{R}$.
  \item Kuadrat terkecil memenuhi $\mathbf{A}^{\top}\mathbf{A}\hat{\mathbf{x}}
        =\mathbf{A}^{\top}\mathbf{b}$; bila kolom $\mathbf{A}$ bebas linear, jawabannya tunggal.
  \item Setiap matriks punya $\mathbf{A}=\mathbf{U}\boldsymbol{\Sigma}\mathbf{V}^{\top}$:
        pangkat adalah banyaknya $\sigma_{i}>0$, $\|\mathbf{A}\|_{2}=\sigma_{1}$, dan
        $\kappa(\mathbf{A})=\sigma_{1}/\sigma_{n}$.
  \item Pseudoinvers $\mathbf{A}^{+}=\mathbf{V}\boldsymbol{\Sigma}^{+}\mathbf{U}^{\top}$
        memberi sisa terkecil sekaligus norma terkecil, dan pada sistem kurang tentsu
        memilih jawaban tegak lurus inti.
  \item Persamaan normal bekerja pada $\kappa^{2}$, SVD pada $\kappa$; karena itu masalah
        berkondisi buruk kehilangan digit dua kali lebih cepat lewat $\mathbf{A}^{\top}\mathbf{A}$.
\end{itemize}
\end{ringkasan}

\section{Latihan}

\begin{latihan}
\begin{enumerate}
  \item Proyeksikan $\mathbf{b}=(1,2,2)$ pada $W=\operatorname{span}\{(1,0,0),(0,1,1)\}$
        dengan tangan, periksa dengan \texttt{numpy.linalg.qr}, lalu hitung jarak
        $\mathbf{b}$ ke $W$.

  \item Buktikan $\|\mathbf{b}\|^{2}=\|\operatorname{proj}_{W}\mathbf{b}\|^{2}
        +\|\mathbf{r}\|^{2}$ untuk sembarang $W$ dan $\mathbf{b}$, lalu jelaskan mengapa
        kesamaan itu gagal bila $\mathbf{p}$ bukan proyeksi ortogonal. (Petunjuk: ambil
        basis miring di $\mathbb{R}^{2}$ untuk bagian kedua.)

  \item Bila kolom $\mathbf{A}$ ortonormal, buktikan $\mathbf{A}^{\top}\mathbf{A}=\mathbf{I}$,
        $\hat{\mathbf{x}}=\mathbf{A}^{\top}\mathbf{b}$, dan $\kappa(\mathbf{A})=1$.

  \item Nilai klaim ini: \emph{bila $\mathbf{A}^{\top}\mathbf{A}$ dapat dibalik maka
        baris-baris $\mathbf{A}$ bebas linear.} Tentukan benar, salah, atau belum cukup
        bukti, dan sertakan contoh penyangkal bila perlu.

  \item Hitung SVD $\mathbf{A}=\begin{pmatrix}1&1\\ 1&1\end{pmatrix}$ dengan tangan: nilai
        singular, $\mathbf{U}$, $\mathbf{V}$, pangkat. Tentukan $\mathbf{A}^{+}\mathbf{b}$
        untuk $\mathbf{b}=(1,3)$, hitung sisanya, periksa dengan
        \texttt{numpy.linalg.pinv}.

  \item Buktikan sisa kuadrat terkecil selalu tegak lurus $\operatorname{col}(\mathbf{A})$
        dan sebaliknya, lalu sebutkan di mana kebebasan linear kolom dipakai.
        (Petunjuk: pakai penguraian norma seperti pada bukti persamaan normal; argumen
        ``perlu'' memilih arah $\mathbf{h}$.)

  \item Jalankan Gram-Schmidt pada $\{(1,1,1),(1,1,0),(1,0,0)\}$ dengan tangan, bandingkan
        dengan \texttt{numpy.linalg.qr}, hitung $\|\mathbf{Q}^{\top}\mathbf{Q}-\mathbf{I}\|$
        untuk keduanya, dan jelaskan mengapa tandanya boleh berbeda.

  \item Misalkan $m<n$ dan baris $\mathbf{A}$ bebas linear. Buktikan
        $\mathbf{A}\mathbf{A}^{\top}$ dapat dibalik, dan bahwa
        \[ \mathbf{A}^{\top}(\mathbf{A}\mathbf{A}^{\top})^{-1}\mathbf{b} \]
        adalah penyelesaian yang tegak lurus inti $\mathbf{A}$. (Petunjuk: tiru bukti
        akibat kasus kolom penuh dengan $\mathbf{A}^{\top}$ sebagai matriksnya.)

  \item Dengan data kalibrasi di bagian ketujuh, ganti bacaan terakhir menjadi
        $314{,}3$. Hitung ulang $a$ dan $b$ lewat aturan Cramer, bandingkan dengan
        \texttt{scipy.linalg.lstsq}, dan jelaskan mengapa titik di ujung menggeser garis
        lebih banyak daripada titik di tengah.

  \item Untuk $\mathbf{B}=\begin{pmatrix}1&2\\ 2&4\\ 3&6\end{pmatrix}$ dan
        $\mathbf{b}=(1,0,0)$, tuliskan semua penyelesaian kuadrat terkecil dalam bentuk
        $\mathbf{x}^{+}+t\mathbf{k}$, tentukan $\mathbf{k}$, lalu tunjukkan bahwa
        $\|\mathbf{x}^{+}+t\mathbf{k}\|^{2}$ minimum di $t=0$.
\end{enumerate}
\end{latihan}

\begin{untukpeneliti}
Yang terbukti di bab ini: Cauchy-Schwarz dan ketaksamaan segitiga untuk sebarang ruang
berhasilkali dalam; ketunggalan koefisien pada basis ortonormal; sifat terbaik proyeksi;
kebenaran Gram-Schmidt; kesetaraan kuadrat terkecil dengan persamaan normal; eksistensi
SVD; dan teorema pseudoinvers. Yang dipinjam hanya lemma spektral, yang buktinya memakai
teorema nilai ekstrem pada bola satuan (Bab 11); di sini hanya sketsanya.

Yang sering dipopulerkan ada dua. Pertama, klaim bahwa ``SVD selalu stabil''. Yang benar:
algoritme SVD stabil secara backward, artinya jawaban terhitung adalah jawaban eksak dari
matriks yang berselisih sedikit dari $\mathbf{A}$; jarak ke jawaban untuk $\mathbf{A}$
sendiri tetap dapat membesar sebesar $\kappa$. Kedua, klaim bahwa memotong nilai singular
kecil ``menghilangkan derau''. Pemotongan menghilangkan arah, dan apakah arah itu derau
atau sinyal tidak dapat diputuskan dari $\mathbf{A}$ saja.

Sumber primer: Axler edisi ke-4 bagian 6A sampai 6C dan 7E sampai 7F; Anton edisi ke-11
bagian 6.1 sampai 6.4 serta 9.4 dan 9.5; untuk sisi numerik Trefethen dan Bau,
\emph{Numerical Linear Algebra}, dan Golub dan Van Loan, \emph{Matrix Computations}
edisi ke-4, bagian 2.4 dan 5.3.

Pertanyaan terbuka. Ambang kapan nilai singular dianggap nol tidak punya jawaban bebas
parameter; pilihannya bergantung pada skala data dan galat yang sanggup diterima (Bab 22).
Untuk matriks besar, algoritme acak hanya menjamin secara probabilistik. Batas bab ini:
matriks real berdimensi hingga; untuk matriks kompleks transpos diganti transpos konjugat,
dan untuk ruang berdimensi tak hingga SVD diganti dekomposisi operator kompak, yang nilai
singularnya dapat menumpuk di nol.
\end{untukpeneliti}

\peta{Bab ini mengikuti Anton edisi ke-11 bagian 6.1 (Inner Products), 6.2 (Angle and
Orthogonality), 6.3 (Gram-Schmidt; QR-Decomposition), 6.4 (Best Approximation; Least
Squares), 9.4 (Singular Value Decomposition), dan 9.5 (Data Compression Using SVD).
Padanan abstraknya Axler edisi ke-4, Bab 6 bagian 6A sampai 6C dan Bab 7 bagian 7E
sampai 7F, dengan teorema spektral di 7B. Strang edisi ke-6 memberi jalan geometris,
tetapi salinan di direktori sumber hanya memuat sampul dan prakata. Sisi numerik:
Trefethen dan Bau; Golub dan Van Loan 2.4, 5.3.}


\part{Analisis dan Kalkulus}

Bagian ini membuat kata ``dekat'', ``halus'', dan ``menyusut'' menjadi pernyataan
yang bisa diperiksa dengan bilangan. Prasyaratnya Bab~3--5 dan~7. Setelah enam bab,
kamu bisa membuktikan kekontinuan, menurunkan dan mengintegralkan, menukar limit
dengan integral, dan menyelesaikan persamaan diferensial orde satu serta sistem
dua dimensi.
\chapter{Kekontinuan dan Kekontinuan Seragam}

Sebuah program memotret $f(x)=1/x$ pada kisi titik berjarak $10^{-4}$ di dalam
$(0,1]$. Di sekitar $x=10^{-4}$ dua titik bertetangga, yaitu $10^{-4}$ dan
$2\times 10^{-4}$, memberi keluaran $10000$ dan $5000$; selisihnya $5000$. Di sekitar
$x=1$ dua titik bertetangga memberi $1$ dan $1.0001$; selisihnya sekitar $10^{-4}$.
Fungsi yang sama, jarak kisi yang sama, dua selisih yang berbeda tujuh pangkat.

Pertanyaan yang muncul bukan ``apakah $1/x$ kontinu'', sebab jawabannya ya di setiap
titik $(0,1]$. Pertanyaannya lain: berapa jarak kisi yang harus dipakai supaya
\emph{seluruh} keluaran pada interval itu bertahan di dalam satu toleransi? Untuk
toleransi $10^{-2}$, titik $x=1$ hanya menuntut jarak di bawah
$9.901\times 10^{-3}$, sedangkan titik $x=10^{-4}$ menuntut jarak di bawah
$1.000\times 10^{-10}$. Nisbah kedua tuntutan itu $9.9\times 10^{7}$. Tidak ada satu
jarak pun yang melayani keduanya, dan itu bukan cacat program.

Bab 4 memberi limit barisan, Bab 3 memberi supremum, dan keduanya belum cukup: limit
fungsi membandingkan $f$ dengan satu bilangan tetap, sehingga ia tidak dapat
membandingkan dua keluaran dari dua masukan berdekatan. Yang dibutuhkan adalah satu
ambang untuk seluruh himpunan, beserta cara membuktikan bahwa ambang itu ada. Bab ini
menyusun alat tersebut, dari definisi $\varepsilon$-$\delta$ sampai himpunan kompak dan
kelas fungsi Lipschitz, lalu menutup dengan fungsi yang kontinu di satu titik saja.

\begin{untukanak}
Bayangkan tangga darurat dengan anak tangga yang makin ke atas makin tinggi. Di
bagian bawah, satu langkah selebar telapak kaki menaikkan kita sedikit. Di bagian
atas, langkah selebar itu menaikkan kita sejengkal, lalu sehasta, lalu hampir
setinggi badan. Kakak kita berdiri di samping dan menanyakan satu hal: berapa lebar
langkah yang aman untuk seluruh tangga? Jawaban jujur untuk tangga seperti itu tidak
ada. Lebar yang aman di bawah terlalu lebar di atas, dan lebar yang aman di atas
membuat kita merangkak di bawah.

Sekarang pertanyaan yang membuat penasaran. Apakah semua tangga begitu? Ada tangga
yang anak tangganya tetap, dan ada tangga yang anak tangganya berubah tetapi tidak
pernah melewati batas tertentu. Untuk tangga jenis kedua, satu lebar langkah berlaku
dari bawah sampai atas. Bedanya bukan pada tinggi tangga, melainkan pada seberapa
curam kemiringannya berubah.

Hal yang sama terjadi pada grafik fungsi. Ada grafik yang bisa dijaga dengan satu
aturan jarak di seluruh daerahnya, ada yang menuntut aturan berbeda di setiap titik.
Di bab ini kita akan tahu mengapa perbedaan itu bisa dibuktikan dengan pasti, dan
mengapa daerah yang tertutup dan terbatas selalu termasuk jenis yang pertama.
\end{untukanak}

\section{Limit fungsi dan ambang delta}
\tingkat{dasar}

Kata ``mendekati'' tidak dapat dipakai sebagai definisi, karena setiap orang boleh
mengartikannya berbeda. Yang dapat dipakai adalah taruhan: berapa pun ketelitian yang
diminta, keluaran akan masuk ke dalam batas ketelitian itu asalkan masukannya cukup
dekat. Kalimat sehari-harinya: kita memasang jendela sempit di sekitar $L$, dan
menuntut $f$ masuk ke jendela itu untuk semua $x$ yang cukup dekat ke $c$.

\kamus{limit fungsi}{functional limit}{janji bahwa keluaran bisa dipaksa masuk ke
jendela sempit asalkan masukan cukup dekat dengan pintunya}

Sebelum definisinya, satu istilah teknis diperlukan. Titik $c$ disebut \emph{titik
limit}\index{titik limit} himpunan $A$ bila setiap interval terbuka di sekitar $c$
memuat titik $A$ selain $c$. Tanpa syarat itu, definisi limit menjadi hampa di titik
terpencil.

\begin{definisi}[Limit fungsi]
\index{limit fungsi}
Misalkan $A\subseteq\mathbb{R}$, $c$ titik limit $A$, dan $f\colon A\to\mathbb{R}$.
Kita tulis $\lim_{x\to c}f(x)=L$ bila untuk setiap $\varepsilon>0$ terdapat
$\delta>0$ sehingga untuk setiap $x\in A$ dengan $0<|x-c|<\delta$ berlaku
$|f(x)-L|<\varepsilon$.
\end{definisi}

Dua hal mudah terlewat. Pertama, syarat $0<|x-c|$ berarti nilai $f$ di $c$ tidak ikut
dihitung; fungsi boleh tidak terdefinisi di $c$ dan limitnya tetap ada. Kedua, urutan
kuantor: $\delta$ boleh bergantung pada $\varepsilon$, tetapi tidak pada $x$. Pada
bagian tentang kekontinuan seragam, kebebasan itulah yang diperdebatkan.

Definisi ini alat memeriksa, bukan resep mencari $\delta$; $\delta$ dicari dengan
bekerja mundur dari $|f(x)-L|$. Untuk
\[
f(x)=\frac{x^{2}-1}{x-1}\quad (x\neq 1),\qquad L=2,
\]
aljabar memberi $f(x)=x+1$, sehingga $|f(x)-2|=|x-1|$. Jadi $\delta=\varepsilon$
bekerja. Untuk $\varepsilon=10^{-2}$ ambil $\delta=10^{-2}$: pada $x=1+10^{-2}$
selisihnya tepat $10^{-2}$, dan itu \emph{belum} memenuhi syarat ketat; pada
$x=1+5\times10^{-3}$ selisihnya $5\times10^{-3}$. Karena itu $\delta=\varepsilon/2$ adalah pilihan
yang aman, dan $\delta=\varepsilon$ adalah pilihan paling murah yang batasnya tidak
pernah dicapai dari dalam.

Satu hal yang wajar dituntut dari limit adalah ketunggalannya: sebuah fungsi tidak
dapat berhenti di dua tempat berbeda sekaligus.

\begin{teorema}[Ketunggalan limit]
\index{limit fungsi}
Jika $\lim_{x\to c}f(x)=L_{1}$ dan $\lim_{x\to c}f(x)=L_{2}$, maka $L_{1}=L_{2}$.
\end{teorema}

\begin{proof}[Bukti]
Andaikan $L_{1}\neq L_{2}$ dan tetapkan $\varepsilon=|L_{1}-L_{2}|/2>0$. Ada
$\delta_{1}>0$ sehingga $0<|x-c|<\delta_{1}$ memberi $|f(x)-L_{1}|<\varepsilon$, dan
ada $\delta_{2}>0$ sehingga $0<|x-c|<\delta_{2}$ memberi $|f(x)-L_{2}|<\varepsilon$.
Ambil $x$ dengan $0<|x-c|<\min(\delta_{1},\delta_{2})$; titik seperti itu ada karena
$c$ titik limit. Maka
\[
|L_{1}-L_{2}|\le|L_{1}-f(x)|+|f(x)-L_{2}|<2\varepsilon=|L_{1}-L_{2}|,
\]
suatu kemustahilan. Jadi $L_{1}=L_{2}$.
\end{proof}

Kita juga memakai satu jembatan yang tidak dibuktikan di sini: limit fungsi ada dan
bernilai $L$ tepat ketika setiap barisan $x_{n}\to c$ dengan $x_{n}\neq c$ memenuhi
$f(x_{n})\to L$. Namanya kriteria barisan; buktinya ada di Abbott, \emph{Understanding
Analysis}, Bab 4 bagian Functional Limits, dan di Rudin, \emph{Principles of
Mathematical Analysis}, Bab 4 bagian Continuous Functions. Jembatan itu yang membuat
barisan dari Bab 4 dapat dipakai membongkar limit yang tidak ada. Untuk $g(x)=x/|x|$,
barisan $1/n$ memberi $g\to 1$ sedangkan $-1/n$ memberi $g\to -1$; dua limit berbeda,
jadi $\lim_{x\to 0}g(x)$ tidak ada. Pola itu dipakai lagi pada kekontinuan seragam,
dengan dua barisan yang jaraknya menuju nol.

\section{Kekontinuan dan sifat aljabarnya}
\tingkat{menengah}

Limit membandingkan $f$ dengan bilangan $L$ yang berdiri sendiri. Kekontinuan meminta
sesuatu yang lebih ketat: bilangan itu haruslah nilai fungsi di titiknya. Kalimat
sehari-harinya: grafik tidak melompat di titik itu, dan nilainya di titik itu memang
nilai yang didekati dari kiri dan kanan.

\kamus{kontinu}{continuous}{grafik yang dapat digambar tanpa mengangkat pena di
titik itu}

\begin{definisi}[Kekontinuan di satu titik]
\index{kontinu}
Misalkan $f\colon A\to\mathbb{R}$ dan $c\in A$. Fungsi $f$ \emph{kontinu di $c$} bila
untuk setiap $\varepsilon>0$ terdapat $\delta>0$ sehingga untuk setiap $x\in A$
dengan $|x-c|<\delta$ berlaku $|f(x)-f(c)|<\varepsilon$. Fungsi disebut kontinu pada
$A$ bila ia kontinu di setiap titik $A$.
\end{definisi}

Perhatikan bahwa $x=c$ kini diizinkan, dan yang diminta bukan sekadar limitnya ada,
melainkan limitnya sama dengan nilai fungsinya. Karena itu kekontinuan bersifat
lokal: ia hanya berbicara tentang perilaku $f$ di sekitar satu titik.

Untuk $f(x)=x^{2}$ di $c=2$ dan $\varepsilon>0$ berlaku $|x^{2}-4|=|x-2|\,|x+2|$, dan
bila $|x-2|<1$ maka $|x+2|<5$. Jadi $\delta=\min(1,\varepsilon/5)$ bekerja, sedangkan
$\delta$ terbesar yang masih sah adalah akar positif dari $\delta(4+\delta)=\varepsilon$,
yaitu $2.49984376953\times10^{-4}$ untuk $\varepsilon=10^{-3}$. Selisih kedua angka itu
adalah harga keamanan yang dibayar karena tidak mau menghitung akar.
Gambar~\ref{fig:epsdelta} memperlihatkan kedua ambang itu untuk $\varepsilon=1$.

\begin{figure}[htbp]
\centering
\begin{tikzpicture}
\begin{axis}[width=0.82\linewidth,
  axis lines=left,xmin=1.5,xmax=2.6,ymin=2.6,ymax=5.8,
  xlabel={$x$},ylabel={$f(x)=x^{2}$},
  xtick={1.7639,2,2.2361},xticklabels={$2-\delta$,$2$,$2+\delta$},
  ytick={3,4,5},yticklabels={$4-\varepsilon$,$4$,$4+\varepsilon$},
  tick label style={font=\scriptsize}]
\fill[blue!8,draw=blue!45] (1.5,3) rectangle (2.6,5);
\fill[orange!14,draw=orange!60] (1.7639,2.6) rectangle (2.2361,5.8);
\draw[densely dashed,black!60] (1.7,2.6) rectangle (2.3,5.8);
\addplot[thick,domain=1.5:2.6] {x^2};
\draw[fill=white] (2.3,5.29) circle (2pt);
\node[font=\scriptsize,anchor=east] at (2.32,5.5) {$\delta=0.3$};
\end{axis}
\end{tikzpicture}
\caption{Pita $\varepsilon$ dan pita $\delta$ untuk $f(x)=x^{2}$ di $a=2$ dengan
$\varepsilon=1$. Batas $\delta$ memakai rumus $\delta=\sqrt{4+\varepsilon}-2$
pada bagian ini, jadi $\delta=\sqrt{5}-2\approx0.236$; pita putus-putus memakai
$\delta=0.3$, dan di $x=2.3$ nilai $f(2.3)=5.29$ sudah keluar pita.}
\label{fig:epsdelta}
\end{figure}

\begin{proposisi}
\index{kontinu}
Misalkan $c$ titik limit $A$. Fungsi $f$ kontinu di $c$ jika dan hanya jika
$\lim_{x\to c}f(x)=f(c)$.
\end{proposisi}

\begin{proof}[Bukti]
Bila $f$ kontinu di $c$ dan $\varepsilon>0$, ambil $\delta$ dari definisi
kekontinuan. Untuk $0<|x-c|<\delta$ berlaku $|f(x)-f(c)|<\varepsilon$, dan itu persis
definisi limit dengan $L=f(c)$. Sebaliknya, bila limitnya $f(c)$, ambil $\delta$ dari
definisi limit; syarat $0<|x-c|$ dapat dibuang karena untuk $x=c$ ketaksamaannya
menjadi $0<\varepsilon$, yang benar.
\end{proof}

Kekontinuan tertutup terhadap operasi yang biasa dipakai. Kalimat sehari-harinya:
menjumlah, mengalikan, dan membagi fungsi yang tidak melompat menghasilkan fungsi yang
juga tidak melompat, selama pembaginya tidak nol.

\begin{proposisi}[Sifat aljabar dan komposisi]
\index{kontinu}
Jika $f$ dan $g$ kontinu di $c$, maka $f+g$, $fg$, dan $cf$ kontinu di $c$; bila
$g(c)\neq 0$, maka $f/g$ juga kontinu di $c$. Jika $f$ kontinu di $c$ dan $g$ kontinu
di $f(c)$, maka $g\circ f$ kontinu di $c$.
\end{proposisi}

\begin{proof}[Bukti]
Bagian aljabar mengikuti aturan limit yang sama seperti pada barisan, dengan
mengganti $x_{n}\to c$ menjadi $x\to c$; karena $c$ titik limit, kriteria barisan
memindahkan setiap pernyataan limit fungsi menjadi pernyataan limit barisan, dan
aturan aljabar barisan sudah dibuktikan di Bab 4. Untuk komposisi, ambil
$\varepsilon>0$. Karena $g$ kontinu di $f(c)$, ada $\eta>0$ sehingga
$|y-f(c)|<\eta$ berakibat $|g(y)-g(f(c))|<\varepsilon$. Karena $f$ kontinu di $c$,
untuk $\eta$ itu ada $\delta>0$ sehingga $|x-c|<\delta$ berakibat
$|f(x)-f(c)|<\eta$. Rantai implikasi itu memberi $|g(f(x))-g(f(c))|<\varepsilon$
setiap kali $|x-c|<\delta$.
\end{proof}

Pola bukti komposisi layak diingat: dua $\varepsilon$ dipakai berurutan, dan $\eta$
berperan sebagai $\varepsilon$ bagi fungsi dalam. Akibatnya polinomial dan fungsi
rasional kontinu di seluruh daerah definisinya. Yang tidak dijamin adalah fungsi yang
didefinisikan bercabang, dan dari situ contoh penyangkal muncul di akhir bab.

Bentuk kegagalan yang paling kasar adalah lompatan. Fungsi pada
Gambar~\ref{fig:lompat} bernilai $x$ di sebelah kiri $1$, bernilai $2$ di titik $1$,
dan bernilai $x+1$ di sebelah kanannya. Di $x=1$ nilainya $2$, sedangkan nilai dari
kiri mendekati $1$. Selisih $1$ itu tidak dapat ditutup oleh $\delta$ sekecil apa pun,
sebab untuk setiap $\delta$ masih ada titik kiri dengan $|f(x)-f(1)|$ mendekati $1$.

\begin{figure}[htbp]
\centering
\begin{tikzpicture}
\begin{axis}[width=0.70\linewidth,height=5.0cm,
  axis lines=left,xmin=0,xmax=2.2,ymin=0,ymax=3.2,
  xlabel={$x$},ylabel={$f(x)$},xtick={1},xticklabels={$1$},
  ytick={1,2,3},tick label style={font=\scriptsize},
  label style={font=\scriptsize}]
\fill[layerplum!12,draw=layerplum!60] (0,1.5) rectangle (2.2,2.5);
\addplot[densely dashed,layerplum] coordinates {(0.9,0) (0.9,3.2)};
\addplot[densely dashed,layerplum] coordinates {(1.1,0) (1.1,3.2)};
\addplot[thick,ink,domain=0:1,samples=2] {x};
\addplot[thick,ink,domain=1:2.2,samples=2] {x+1};
\addplot[only marks,mark=*,mark size=1.6pt,accent] coordinates {(0.95,0.95)};
\addplot[only marks,mark=o,mark size=2.6pt,ink,mark options={fill=white,solid}]
  coordinates {(1,1)};
\addplot[only marks,mark=*,mark size=1.7pt,ink] coordinates {(1,2)};
\draw[accent,->,>=stealth,thin] (0.60,1.30) -- (0.93,0.99);
\node[font=\scriptsize,layerplum,anchor=west] at (1.13,2.75) {$\delta$};
\node[font=\scriptsize,accent,anchor=west] at (0.10,1.22) {$f(0.95)=0.95$};
\node[font=\scriptsize,ink,anchor=north east] at (1.02,1.88) {$f(1)=2$};
\end{axis}
\end{tikzpicture}
\caption{Lompatan di $x=1$: $f(x)=x$ untuk $x<1$, $f(1)=2$, $f(x)=x+1$ untuk $x>1$.
Pita $\varepsilon=\tfrac12$ di sekitar $f(1)=2$ adalah jalur mendatar antara $1.5$ dan
$2.5$, sedangkan kedua garis putus-putus menandai $|x-1|<0.1$. Titik $x=0.95$ sudah
masuk pita $\delta$ itu, tetapi nilainya $0.95$ masih $0.55$ di bawah batas bawah
pita $\varepsilon$. Lingkaran kosong di $(1,1)$ adalah nilai yang didekati dari kiri,
bukan nilai fungsinya. Mengecilkan $\delta$ tidak menolong: titik kiri yang lebih dekat
ke $1$ memberi nilai yang lebih dekat ke $1$, bukan ke $2$.}
\label{fig:lompat}
\end{figure}

\section{Teorema nilai antara}
\tingkat{menengah}

Kalimat sehari-harinya: untuk berpindah dari ruang bawah ke ruang atas, seorang
pendaki harus melewati setiap ketinggian di antaranya. Grafik fungsi kontinu tidak
dapat melompati nilai.

\begin{teorema}[Nilai antara]
\index{teorema nilai antara}
Misalkan $f$ kontinu pada $[a,b]$ dan $f(a)<0<f(b)$. Maka ada $c\in(a,b)$ dengan
$f(c)=0$. Lebih umum, untuk setiap $v$ di antara $f(a)$ dan $f(b)$ ada $c\in[a,b]$
dengan $f(c)=v$.
\end{teorema}

\begin{proof}[Bukti]
Cukup dibuktikan kasus $f(a)<0<f(b)$, karena kasus umum diperoleh dengan mengenakan
teorema itu pada $f-v$. Bentuk $S=\{x\in[a,b] : f(x)<0\}$. Himpunan itu tidak kosong
karena $a\in S$ dan terbatas di atas oleh $b$, jadi menurut sifat supremum di Bab 3,
$c=\sup S$ ada dan $a\le c\le b$.

Pertama, $f(c)\le 0$. Karena $c$ supremum, ada barisan $x_{n}\in S$ dengan $x_{n}\to c$,
dan $f(x_{n})<0$ untuk setiap $n$. Kekontinuan memberi $f(x_{n})\to f(c)$, dan limit
mempertahankan ketaksamaan lemah, jadi $f(c)\le 0$.

Kedua, $f(c)\ge 0$. Andaikan $f(c)<0$. Karena $f(b)>0$, titik $c$ bukan $b$, jadi
$c<b$. Kekontinuan di $c$ memberi $\delta>0$ sehingga $|x-c|<\delta$ berakibat
$|f(x)-f(c)|<-f(c)$, khususnya $f(x)<0$. Untuk $\delta<2(b-c)$ titik $c+\delta/2$
berada di $[a,b]$ dan memenuhi $f(c+\delta/2)<0$, jadi ia anggota $S$; padahal
$c+\delta/2>c=\sup S$. Kemustahilan itu memberi $f(c)\ge 0$.

Kedua ketaksamaan bersama memberi $f(c)=0$, dan karena $f(a)<0<f(b)$, titik itu di
$(a,b)$.
\end{proof}

Perhatikan bahwa bukti itu tidak memberi cara menemukan $c$. Ia hanya menunjukkan
bahwa himpunan titik negatif mempunyai atap, dan bahwa atap itu haruslah akar.
Kesenjangan antara ``ada'' dan ``dapat dihitung'' inilah yang membuat teorema nilai
antara dipakai sebagai jaminan keberadaan, bukan sebagai algoritme.

\begin{untukanak}
Seorang pendaki mulai dari lembah dan berhenti di puncak. Ia boleh beristirahat,
boleh berjalan mundur sebentar, boleh berhenti lama di satu tempat. Yang tidak ia
lakukan adalah melompat. Karena itu, kalau di lembah ada sungai dan di puncak tidak
ada sungai, pasti ada satu titik di jalan itu tempat ia pertama kali berada tepat di
ketinggian sungai.

Yang membuat penasaran: berapa banyak titik seperti itu? Bisa satu, bisa tiga, bisa
tak berhingga, dan cerita pendaki tadi tidak memberi tahu yang mana. Ia hanya
menjamin paling sedikit satu. Matematikawan memakai jaminan itu untuk hal lain:
membuktikan bahwa sebuah persamaan pasti punya jawaban, walaupun jawabannya tidak
selalu dapat ditulis dengan rumus sederhana.

Pertanyaan berikutnya lebih sulit. Bagaimana kalau jalannya tidak bersambung, misalnya
ada jurang di tengah? Jaminan tadi langsung hilang, dan itu bukan keisengan: ada
fungsi yang nilainya melompat dari negatif ke positif tanpa pernah menyentuh nol.
Di bab ini kita akan tahu mengapa sifat ``tidak melompat'' sudah cukup, dan mengapa
sifat itu harus disebut lebih dahulu sebelum teoremanya dipakai.
\end{untukanak}

Dua akibat langsung layak dicatat. Pertama, gambar interval di bawah fungsi kontinu
adalah interval: $f([a,b])$ memuat setiap nilai antara $f(a)$ dan $f(b)$. Kedua,
setiap fungsi kontinu $f\colon[0,1]\to[0,1]$ mempunyai titik tetap. Pandang
$h(x)=f(x)-x$; di $x=0$ berlaku $h(0)=f(0)\ge 0$ dan di $x=1$ berlaku
$h(1)=f(1)-1\le 0$. Bila salah satu bernilai nol, titik itu titik tetap. Bila tidak,
teorema nilai antara memberi $c$ dengan $h(c)=0$, yakni $f(c)=c$.

Angka pada kasus konkret dapat diperiksa. Persamaan $x^{3}+x-1=0$ mempunyai akar di
$(0,1)$ karena $f(0)=-1<0<1=f(1)$. Setelah dua puluh satu belah dua, akarnya terkurung
di antara $0.682327747$ dan $0.682328224$, dengan lebar selang $2^{-21}=4.768\times10^{-7}$;
nilai sebenarnya $0.6823278038280193$. Belah dua bukan bukti, tetapi angkanya dapat
dibandingkan dengan batas teoretis.

\section{Teorema nilai ekstrem}
\tingkat{menengah}

Kalimat sehari-harinya: pada jalan yang tertutup dan tidak putus, pasti ada titik
tertinggi dan titik terendah. Sebelum membuktikannya, kita perlu satu kendali: fungsi
seperti itu tidak dapat lari ke tak hingga.

\begin{lemma}
\index{teorema nilai ekstrem}
Jika $f$ kontinu pada $[a,b]$, maka $f$ terbatas pada $[a,b]$.
\end{lemma}

\begin{proof}[Bukti]
Andaikan $f$ tidak terbatas. Maka untuk setiap $n$ ada $x_{n}\in[a,b]$ dengan
$|f(x_{n})|>n$. Barisan $(x_{n})$ terbatas, jadi menurut teorema
Bolzano--Weierstrass di Bab 4 ada subbarisan $x_{n_{k}}\to x$. Karena
$a\le x_{n_{k}}\le b$ untuk semua $k$, limit mempertahankan ketaksamaan dan
$x\in[a,b]$. Kekontinuan di $x$ memberi $f(x_{n_{k}})\to f(x)$, sehingga barisan itu
terbatas. Padahal $|f(x_{n_{k}})|>n_{k}\to\infty$. Kemustahilan itu menutup
pengandaian.
\end{proof}

Sekarang pernyataan utamanya.

\begin{teorema}[Nilai ekstrem]
\index{teorema nilai ekstrem}
Jika $f$ kontinu pada $[a,b]$, maka ada $c,d\in[a,b]$ sehingga
$f(c)=\sup f([a,b])$ dan $f(d)=\inf f([a,b])$. Dengan kata lain, $f$ mencapai nilai
maksimum dan minimumnya.
\end{teorema}

\begin{proof}[Bukti]
Menurut lemma, himpunan $f([a,b])$ terbatas dan tidak kosong, jadi $M=\sup f([a,b])$
ada dan berhingga. Untuk setiap $n$ ada $x_{n}\in[a,b]$ dengan $f(x_{n})>M-1/n$;
barisan seperti itu ada karena $M$ batas atas terkecil. Menurut
Bolzano--Weierstrass ada subbarisan $x_{n_{k}}\to c$, dan karena
$a\le x_{n_{k}}\le b$ diperoleh $c\in[a,b]$. Kekontinuan di $c$ memberi
$f(x_{n_{k}})\to f(c)$, sedangkan pengurungan tadi memberi $f(x_{n_{k}})\to M$.
Ketunggalan limit memberi $f(c)=M$. Untuk minimum, terapkan hasil itu pada $-f$.
\end{proof}

Kedua bukti memakai jembatan yang sama: barisan terbatas di dalam interval tertutup
mempunyai subbarisan yang menuju titik di dalam interval itu. Ketertutupan yang
memastikan limitnya tidak keluar, dan itulah yang hilang pada interval terbuka.

Contoh penyangkalnya sederhana. Fungsi $f(x)=x$ kontinu pada $(0,1)$ tetapi tidak
mencapai maksimum maupun minimum, karena $0$ dan $1$ bukan anggota daerahnya. Fungsi
$g(x)=1/x$ kontinu pada $(0,1]$ tetapi tidak terbatas di sana: $g(1/n)=n$. Jadi
ketertutupan dan keterbatasan daerah bukan hiasan; keduanya adalah seluruh sebab
teorema bekerja.

Akibatnya, untuk $f$ kontinu pada $[a,b]$ dengan $m=\inf f$ dan $M=\sup f$, berlaku
$f([a,b])=[m,M]$. Kesamaan itu menggabungkan dua teorema: yang pertama mengatakan $m$
dan $M$ tercapai, yang kedua mengatakan semua nilai di antaranya terpakai.

\section{Kekontinuan seragam}
\tingkat{lanjut}

Pada definisi kekontinuan, $\delta$ boleh bergantung pada $\varepsilon$ \emph{dan}
pada titik $c$. Untuk $f(x)=1/x$ pada $(0,1]$, setiap titik punya $\delta$-nya
sendiri, dan $\delta$ itu mengecil tanpa batas ketika $c$ mendekati nol. Secara lokal
semuanya baik, tetapi tidak pernah ada satu $\delta$ untuk seluruh interval, dan itulah
yang dibutuhkan kisi pada pembuka bab.

\kamus{kekontinuan seragam}{uniform continuity}{satu lebar langkah yang berlaku dari
kaki tangga sampai puncaknya, bukan lebar berbeda di setiap anak tangga}

\begin{definisi}[Kekontinuan seragam]
\index{kekontinuan seragam}
Fungsi $f\colon A\to\mathbb{R}$ disebut \emph{seragam kontinu} pada $A$ bila untuk
setiap $\varepsilon>0$ terdapat $\delta>0$ sehingga untuk semua $x,y\in A$ dengan
$|x-y|<\delta$ berlaku $|f(x)-f(y)|<\varepsilon$.
\end{definisi}

Bandingkan kedua urutan kuantor secara berdampingan:
\[
\text{kontinu: }\ \forall c\in A\ \forall\varepsilon>0\ \exists\delta>0\
\forall x\in A,
\qquad
\text{seragam: }\ \forall\varepsilon>0\ \exists\delta>0\ \forall x,y\in A .
\]
Pada versi seragam, $\delta$ muncul sebelum $x$ dan $y$ diperkenalkan, jadi ia tidak
boleh memanfaatkan pengetahuan tentang titik mana pun. Itulah seluruh isi
perbedaannya, dan karena itu keseragaman selalu relatif terhadap himpunan $A$: fungsi
yang sama dapat seragam pada satu himpunan dan tidak pada himpunan lain.

\begin{bacakeras}
Bacalah keras, dua suara bergantian.

\textbf{Suara satu:} Dekat di satu titik. Berapa pun sempit jendela yang dipasang,
asal masukannya cukup dekat, keluarannya ada di dalam jendela.

\textbf{Suara dua:} Dekat di semua titik. Bukan pintu mana yang ditanyakan. Satu lebar
jendela dipasang lebih dahulu, lalu dituntut berlaku di seluruh lorong, dari pintu
pertama sampai pintu terakhir.

\textbf{Suara satu:} Kalau lorongnya menyempit di ujung, lebar itu harus dipotong
terus.

\textbf{Suara dua:} Dan potongan itu tidak boleh berhenti di satu angka. Berhenti
berarti satu lebar untuk seluruh lorong; terus mengecil berarti janji sepotong-sepotong.

Ulangi baris ini sampai terasa di lidah: \emph{satu titik cukup dekat, semua titik satu
lebar.}
\end{bacakeras}

Pernyataan bahwa keseragaman lebih kuat daripada kekontinuan dibuktikan langsung dari
definisinya: keseragaman memberi satu $\delta$ untuk semua pasangan titik, dan
kekontinuan hanya meminta $\delta$ untuk pasangan yang salah satu anggotanya titik
tetap.

\begin{proposisi}
\index{kekontinuan seragam}
Jika $f$ seragam kontinu pada $A$, maka $f$ kontinu di setiap titik $A$.
\end{proposisi}

\begin{proof}[Bukti]
Ambil $c\in A$ dan $\varepsilon>0$. Karena $f$ seragam kontinu, ada $\delta>0$
sehingga $|x-y|<\delta$ berakibat $|f(x)-f(y)|<\varepsilon$ untuk semua $x,y\in A$.
Dengan memilih $y=c$, syarat $|x-c|<\delta$ memberi $|f(x)-f(c)|<\varepsilon$. Itu
persis definisi kekontinuan di $c$.
\end{proof}

Cara kedua, yang dipakai berulang di sisa bab ini, bekerja dengan kontraposisi: bila
$f$ gagal kontinu di suatu titik, dua barisan yang jaraknya menuju nol dapat disusun
sehingga keluarannya tetap berjarak, dan itu menyangkal keseragaman. Bukti $1/x$ di
bawah memakai pola itu secara penuh.

Sekarang contoh yang memperlihatkan keseragaman bergantung pada daerah, bukan pada
rumus fungsinya.

\begin{proposisi}
\index{kekontinuan seragam}
$f(x)=x^{2}$ seragam kontinu pada $[0,1]$, dengan $\delta=\varepsilon/2$.
\end{proposisi}

\begin{proof}[Bukti]
Untuk $x,y\in[0,1]$ berlaku $|x+y|\le 2$, sehingga
\[
|x^{2}-y^{2}|=|x-y|\,|x+y|\le 2|x-y| .
\]
Bila $|x-y|<\varepsilon/2$, maka $|x^{2}-y^{2}|<\varepsilon$.
\end{proof}

Keseragaman itu tidak bergantung pada kedua ujungnya: pada $(0,1)$ bukti yang sama
tetap bekerja, karena batas $|x+y|\le 2$ tidak memakai titik ujung mana pun. Angka
$\varepsilon/2$ aman tetapi bukan yang terbaik. Untuk $|x-y|=\delta$ selisih
terbesar terjadi di ujung interval dan bernilai $\delta(2-\delta)$, jadi $\delta$ sah
asalkan $\delta\le 1-\sqrt{1-\varepsilon}$. Untuk $\varepsilon=10^{-2}$ batas itu
$5.012563\times10^{-3}$, dan dengan $\delta=\varepsilon/2$ selisih terbesarnya
$9.975\times10^{-3}$.

\begin{proposisi}
\index{kekontinuan seragam}
$f(x)=x^{2}$ tidak seragam kontinu pada $\mathbb{R}$.
\end{proposisi}

\begin{proof}[Bukti]
Ambil $\varepsilon=1$. Andaikan ada $\delta>0$ yang bekerja. Pilih bilangan asli
$n>1/\delta$ dan tetapkan $x_{n}=n$, $y_{n}=n+1/n$. Maka $|x_{n}-y_{n}|=1/n<\delta$,
tetapi
\[
|x_{n}^{2}-y_{n}^{2}|=\left|n^{2}-\left(n+\tfrac1n\right)^{2}\right|
=2+\frac{1}{n^{2}}>2>1=\varepsilon .
\]
Jadi tidak ada $\delta$ yang bekerja untuk $\varepsilon=1$.
\end{proof}

Angka $2$ pada perhitungan itu adalah intinya: selisih nilai tidak mengecil, justru
mendekati $2$, walaupun selisih masukannya menuju nol. Yang gagal adalah laju
perubahan $x^{2}$ yang tumbuh tanpa batas, bukan kekontinuannya. Dengan
$\delta=10^{-4}$ dan pasangan $x=10^{4}$, $y=10^{4}+5\times10^{-5}$, selisih
keluarannya $1.0000000149$; dengan $\delta=10^{-2}$ dan $x=100$, selisihnya
$1.000025$. Angka itu hampir tidak bergerak ketika $\delta$ mengecil sepuluh ribu
kali, dan justru ketidakbergerakan itu yang membuktikan kegagalan keseragaman.

Contoh yang sama pentingnya adalah $1/x$, karena di sini kekontinuan titik demi titik
tidak pernah gagal.

\begin{proposisi}
\index{kekontinuan seragam}
$f(x)=1/x$ tidak seragam kontinu pada $(0,1]$, tetapi seragam kontinu pada
$[1,\infty)$ dengan $\delta=\varepsilon$.
\end{proposisi}

\begin{proof}[Bukti bagian pertama, dengan dua barisan]
Ambil $x_{n}=1/n$ dan $y_{n}=1/(n+1)$. Maka
\[
|x_{n}-y_{n}|=\frac{1}{n(n+1)}\longrightarrow 0,
\qquad\text{tetapi}\qquad
\left|\frac{1}{x_{n}}-\frac{1}{y_{n}}\right|=|n-(n+1)|=1
\]
untuk setiap $n$. Bila $f$ seragam kontinu, ambil $\delta$ untuk $\varepsilon=1$;
untuk $n$ besar berlaku $|x_{n}-y_{n}|<\delta$, sehingga selisih keluarannya harus
kurang dari $1$. Itu bertentangan dengan hasil di atas.
\end{proof}

\begin{proof}[Bukti bagian kedua]
Untuk $x,y\ge 1$ berlaku $xy\ge 1$, sehingga
\[
\left|\frac{1}{x}-\frac{1}{y}\right|=\frac{|x-y|}{xy}\le|x-y| .
\]
Jadi $\delta=\varepsilon$ bekerja, dan keseragamannya berasal dari satu fakta saja:
hasil kali $xy$ tidak pernah lebih kecil dari $1$.
\end{proof}

Polanya jelas. Pada $[1,\infty)$ penyebut $xy$ dijaga jauh dari nol oleh batas bawah
daerahnya; pada $(0,1]$ penyebut itu dapat dibuat sekecil apa pun, dan di situlah
keseragaman hilang. Keseragaman adalah sifat pasangan fungsi dan himpunan, bukan sifat
fungsi saja: $1/x$ muncul di kedua proposisi dengan dua jawaban berbeda.

Ambang terbesar $\delta(c)$ di tiap titik memperlihatkan bedanya sebagai gambar.
Untuk $c\in(0,1]$ dan $\varepsilon>0$, syarat $1/(c-\delta)-1/c\le\varepsilon$ setara
dengan $\delta\le\varepsilon c^{2}/(1+\varepsilon c)$, jadi ambang terbesar di $c$
adalah $\delta(c)=\varepsilon c^{2}/(1+\varepsilon c)$, dan
Gambar~\ref{fig:delta-mengecil} memplotnya bersama satu nilai tetap
$\delta=\varepsilon/2$ yang bekerja untuk $x^{2}$ pada $[0,1]$. Kurva pertama turun
lurus dengan kemiringan dua pada skala ganda-log, dari $9.901\times10^{-3}$ di $c=1$
sampai $1.000\times10^{-10}$ di $c=10^{-4}$; garis kedua mendatar. Keseragaman pada
$[0,1]$ karena itu terbaca sebagai garis mendatar, dan kegagalannya pada $(0,1]$
sebagai kurva yang tidak pernah menyentuh lantai.

\begin{figure}[htbp]
\centering
\begin{tikzpicture}
\begin{axis}[width=0.80\linewidth,height=5.8cm,
  xmode=log,ymode=log,xmin=8e-5,xmax=1.4,ymin=8e-11,ymax=4e-2,
  xlabel={titik $c$},ylabel={ambang $\delta(c)$},
  grid=major,grid style={black!10},
  tick label style={font=\scriptsize},label style={font=\scriptsize},
  legend style={font=\scriptsize,at={(0.97,0.03)},anchor=south east,
    draw=black!15,fill=white,fill opacity=0.9,text opacity=1,row sep=1pt}]
\addplot[thick,accent,domain=1e-4:1,samples=80] {1e-2*x^2/(1+1e-2*x)};
\addlegendentry{$1/x$ pada $(0,1]$, $\varepsilon=10^{-2}$}
\addplot[thick,layerteal,densely dashed,domain=1e-4:1] {5e-3};
\addlegendentry{$x^{2}$ pada $[0,1]$, $\delta=\varepsilon/2$}
\addplot[only marks,mark=*,mark size=1.6pt,accent]
  coordinates {(1,9.901e-3) (1e-2,9.999e-7) (1e-4,1.0e-10)};
\node[font=\scriptsize,accent,anchor=south west] at (1.15e-4,1.7e-10)
  {$1.000\times10^{-10}$};
\node[font=\scriptsize,accent,anchor=east] at (0.92,1.05e-2) {$9.901\times10^{-3}$};
\end{axis}
\end{tikzpicture}
\caption{Ambang $\delta$ terbesar di tiap titik $c$ untuk $1/x$ pada $(0,1]$ dengan
$\varepsilon=10^{-2}$, yaitu $\delta(c)=\varepsilon c^{2}/(1+\varepsilon c)$, pada skala
ganda-log. Nilainya turun dari $9.901\times10^{-3}$ di $c=1$ menjadi
$9.999\times10^{-7}$ di $c=10^{-2}$ dan $1.000\times10^{-10}$ di $c=10^{-4}$,
sehingga grafiknya bergaris lurus dengan kemiringan dua. Garis putus-putus
$\delta=5\times10^{-3}$ adalah pilihan aman $\varepsilon/2$ dari bukti untuk $x^{2}$
pada $[0,1]$: satu angka yang sama berlaku di seluruh selang itu.}
\label{fig:delta-mengecil}
\end{figure}

\begin{tangan}
Bahan: lorong sepanjang dua meter, kapur atau selotip kertas, dan penggaris.

Tandai lantai di titik $0$, lalu berjalan dengan langkah yang terus diperkecil:
setengah meter, seperempat meter, seperdelapan meter, dan seterusnya. Tandai setiap
pendaratan kaki. Setelah enam atau tujuh langkah, berhenti dan lihat tanda itu dari
ujung lorong.

Yang perlu dilihat: tanda menggerombol di dekat ujung, dan sisa jarak mengecil tanpa
pernah menjadi nol. Aturan itu tidak pernah menyelesaikan perjalanan, walaupun setiap
langkah sah menurut aturannya sendiri.

Ulangi dengan aturan lain: langkah tetap setengah meter sampai ujung. Perjalanan
selesai dalam empat langkah, dan satu aturan berlaku dari tanda pertama sampai tanda
terakhir. Itulah bedanya aturan yang mengecil terus dengan satu aturan untuk seluruh
lorong.
\end{tangan}

\section{Kompak, Lipschitz, dan teorema Heine--Cantor}
\tingkat{lanjut}

Kalimat sehari-harinya: pada daerah yang tertutup dan tidak terbentang tak hingga,
kekontinuan biasa sudah cukup untuk memberi satu ambang bagi semua titik. Untuk
menyatakannya, kita butuh satu istilah tentang daerah.

\kamus{kompak}{compact}{kotak yang tertutup rapat dan tidak menjulur tak berujung}

\begin{definisi}[Kompak]
\index{kompak}
Himpunan $K\subseteq\mathbb{R}$ disebut \emph{kompak} bila setiap keluarga himpunan
terbuka yang gabungannya memuat $K$ mempunyai subkeluarga berhingga yang gabungannya
juga memuat $K$.
\end{definisi}

Definisi itu tidak praktis dipakai langsung. Yang menyelamatkan kita adalah
kesetaraannya dengan dua sifat yang mudah diperiksa: untuk himpunan di
$\mathbb{R}$, kompak setara dengan tertutup dan terbatas. Kesetaraan itu dikenal
sebagai teorema Heine--Borel; buktinya kita lewati, dan ada di Rudin,
\emph{Principles of Mathematical Analysis}, Bab 2 bagian Compact Sets.

\begin{teorema}[Heine--Cantor]
\index{heine-cantor}
Jika $f$ kontinu pada himpunan kompak $K$, maka $f$ seragam kontinu pada $K$.
\end{teorema}

\begin{proof}[Bukti]
Andaikan $f$ tidak seragam kontinu. Maka ada $\varepsilon>0$ sehingga untuk setiap $n$
terdapat $x_{n},y_{n}\in K$ dengan $|x_{n}-y_{n}|<1/n$ tetapi
$|f(x_{n})-f(y_{n})|\ge\varepsilon$. Karena $K$ terbatas, $(x_{n})$ terbatas, dan
menurut Bolzano--Weierstrass ada subbarisan $x_{n_{k}}\to x$. Karena $K$ tertutup,
$x\in K$. Dari $|x_{n_{k}}-y_{n_{k}}|<1/n_{k}$ diperoleh $y_{n_{k}}\to x$ juga, dan
$y_{n_{k}}\in K$. Kekontinuan $f$ di $x$ memberi $f(x_{n_{k}})\to f(x)$ dan
$f(y_{n_{k}})\to f(x)$, sehingga $|f(x_{n_{k}})-f(y_{n_{k}})|\to 0$. Itu bertentangan
dengan $|f(x_{n_{k}})-f(y_{n_{k}})|\ge\varepsilon$. Jadi pengandaian salah.
\end{proof}

Bukti itu memakai dua sifat kompak secara terpisah: keterbatasan untuk memicu
Bolzano--Weierstrass, dan ketertutupan untuk memastikan titik limitnya masih di dalam
$K$. Keduanya perlu. Pada $(0,1)$, fungsi $1/x$ kontinu tetapi tidak seragam, dan
daerahnya tidak tertutup. Pada $\mathbb{R}$, fungsi $x^{2}$ kontinu tetapi tidak
seragam, dan daerahnya tidak terbatas.

Kelas fungsi berikut memberi jalan pintas yang sering lebih cepat daripada memeriksa
definisi.

\begin{definisi}[Lipschitz]
\index{lipschitz}
Fungsi $f\colon A\to\mathbb{R}$ disebut \emph{Lipschitz} bila ada konstanta $L\ge 0$
sehingga $|f(x)-f(y)|\le L|x-y|$ untuk semua $x,y\in A$.
\end{definisi}

\begin{proposisi}
\index{lipschitz}
Jika $f$ Lipschitz pada $A$, maka $f$ seragam kontinu pada $A$.
\end{proposisi}

\begin{proof}[Bukti]
Ambil $\varepsilon>0$. Bila $L=0$, fungsi konstan dan pilihan $\delta$ apa pun bekerja.
Bila $L>0$, ambil $\delta=\varepsilon/L$. Untuk $|x-y|<\delta$ berlaku
$|f(x)-f(y)|\le L|x-y|<\varepsilon$.
\end{proof}

Kebalikannya tidak berlaku, dan contohnya memperlihatkan bahwa keseragaman lebih
lemah daripada Lipschitz. Fungsi $f(x)=\sqrt{x}$ kontinu pada $[0,1]$, jadi seragam
kontinu menurut teorema Heine--Cantor. Tetapi
\[
\frac{|\sqrt{x}-\sqrt{0}|}{|x-0|}=\frac{1}{\sqrt{x}}\longrightarrow\infty
\quad (x\to 0^{+}),
\]
sehingga tidak ada $L$ yang memenuhi definisi Lipschitz. Jadi Lipschitz adalah syarat
cukup, bukan syarat perlu. Sebaliknya $x^{2}$ pada $[0,1]$ Lipschitz dengan $L=2$
dan $1/x$ pada $[1,\infty)$ dengan $L=1$, persis angka di bukti sebelumnya.

\begin{tantangan}
Dua contoh kegagalan keseragaman di bab ini, $x^{2}$ pada $\mathbb{R}$ dan $1/x$ pada
$(0,1]$, sama-sama tidak terbatas pada daerahnya. Carilah satu fungsi yang
\emph{terbatas} pada $(0,1)$, kontinu di setiap titiknya, tetapi tidak seragam kontinu
di sana, lalu buktikan.

Petunjuk 1: kegagalan tidak menuntut nilai yang membesar; cukup nilai yang berayun,
dengan ayunan yang makin rapat ketika $x$ mendekati ujung selang.

Petunjuk 2: cobalah $f(x)=\sin(1/x)$ dengan $x_{n}=1/(2\pi n+\pi/2)$ dan
$y_{n}=1/(2\pi n-\pi/2)$. Nilai fungsinya $1$ dan $-1$, sedangkan
$|x_{n}-y_{n}|=\pi/((2\pi n)^{2}-(\pi/2)^{2})$ menuju nol.

Kalau mentok, kembali ke bagian Kekontinuan seragam, bukti $1/x$ dengan dua barisan.
\end{tantangan}

\section{Fungsi yang kontinu di satu titik saja}
\tingkat{lanjut}

Kekontinuan bukan bawaan semua fungsi; ia dapat gagal di himpunan titik yang sangat
besar. Contoh paling murah dibangun dari kepadatan bilangan rasional di Bab 2.

\begin{contoh}
\index{fungsi dirichlet}
Definisikan $f\colon\mathbb{R}\to\mathbb{R}$ dengan $f(x)=x$ bila $x$ rasional dan
$f(x)=0$ bila $x$ irasional. Maka $f$ kontinu di $0$ dan tidak kontinu di titik lain
mana pun.
\end{contoh}

\begin{proof}[Bukti]
Di $x=0$ berlaku $|f(x)-f(0)|=|f(x)|\le|x|$, karena $f(x)$ selalu bernilai $x$ atau
$0$. Jadi $\delta=\varepsilon$ bekerja dan $f$ kontinu di $0$.

Sekarang ambil $c\neq 0$. Karena $\mathbb{Q}$ dan $\mathbb{R}\setminus\mathbb{Q}$
keduanya rapat, ada barisan rasional $r_{n}\to c$ dan barisan irasional $s_{n}\to c$.
Maka $f(r_{n})=r_{n}\to c$ sedangkan $f(s_{n})=0\to 0$. Dua limit barisan berbeda,
jadi menurut kriteria barisan $\lim_{x\to c}f(x)$ tidak ada, dan $f$ tidak kontinu di
$c$.
\end{proof}

Kerabat dekat contoh itu adalah indikator bilangan rasional, yang bernilai $1$ di
rasional dan $0$ di irasional; fungsi itu tidak kontinu di satu titik pun. Contoh
berikut lebih menarik, karena himpunan titik kontinunya justru besar; ia dikenal
sebagai modifikasi fungsi Dirichlet.

\begin{contoh}[Fungsi Thomae, modifikasi Dirichlet]
\index{fungsi thomae}
Definisikan $t\colon[0,1]\to\mathbb{R}$ dengan $t(x)=0$ bila $x$ irasional, dan
$t(x)=1/q$ bila $x=p/q$ dalam bentuk paling sederhana dengan $q\ge 1$. Maka $t$
kontinu di setiap titik irasional dan tidak kontinu di setiap titik rasional.
\end{contoh}

\begin{proof}[Bukti]
Ambil $a$ irasional dan $\varepsilon>0$. Pilih $N$ dengan $1/N<\varepsilon$. Bilangan
rasional di $[0,1]$ dengan penyebut $q\le N$ ada berhingga banyak, jadi di antaranya
ada yang terdekat ke $a$; sebut jaraknya $\eta>0$. Karena $a$ irasional, jarak itu
tidak nol. Untuk $|x-a|<\eta$ ada dua kemungkinan. Bila $x$ irasional, maka
$|t(x)-t(a)|=0<\varepsilon$. Bila $x=p/q$ dengan $q>N$, maka
$|t(x)-t(a)|=1/q<1/N<\varepsilon$. Jadi $t$ kontinu di $a$.

Sekarang ambil $a=p/q$ rasional. Barisan irasional $s_{n}\to a$ memberi
$t(s_{n})=0\to 0$, sedangkan $t(a)=1/q>0$. Limit sepanjang barisan itu tidak sama
dengan nilai fungsinya, jadi $t$ tidak kontinu di $a$.
\end{proof}

Kedua contoh itu menunjukkan bahwa himpunan titik diskontinuitas dapat berupa seluruh
garis, himpunan rasional, atau satu titik. Klasifikasi lengkapnya, yang menyatakan
himpunan itu selalu gabungan terhitung himpunan tertutup, ada di Rudin, Bab 4 bagian
Discontinuities. Kekontinuan yang gagal di banyak titik bukan keanehan, melainkan
gejala terstruktur.

\section{Contoh hitung: delta yang bergantung titik}
\tingkat{menengah}

Kita kembali ke $f(x)=x^{2}$ di $x_{0}=2$ dengan $\varepsilon=10^{-3}$, kali ini
sampai angka terakhir.

\begin{resep}
Tiga langkah menulis bukti $\varepsilon$-$\delta$ untuk $|f(x)-f(c)|<\varepsilon$:
\begin{enumerate}
  \item Kerjakan mundur. Tulis $|f(x)-f(c)|$ sebagai $|x-c|$ dikali sisa, lalu cari
        batas atas sisa itu pada daerah yang dipakai.
  \item Batasi lebih dahulu $|x-c|<1$ (atau $<r$) supaya sisa itu terbatas. Sebut
        batas atasnya $M$, lalu tetapkan $\delta=\min(1,\varepsilon/M)$.
  \item Tulis bukti maju. Ambil $\varepsilon>0$, pilih $\delta$ itu, misalkan
        $|x-c|<\delta$, lalu turunkan $|f(x)-f(c)|<\varepsilon$ baris demi baris.
\end{enumerate}
Langkah 1 dan 3 tidak boleh ditukar. Kalimat bukti dibaca dari $\varepsilon$ ke
$\delta$, tetapi $\delta$ ditemukan dengan arah sebaliknya; menulis langsung dari
$\varepsilon$ biasanya menghasilkan $\delta$ yang masih memuat $x$, dan pilihan seperti
itu bukan $\delta$ yang sah.
\end{resep}

Syarat yang harus dipenuhi adalah
\[
|x^{2}-4|=|x-2|\,|x+2|<\varepsilon
\]
untuk semua $x$ dengan $|x-2|<\delta$. Bila $\delta<1$, maka $|x+2|<4+\delta$,
sehingga cukup menuntut $\delta(4+\delta)=\varepsilon$. Persamaan kuadrat itu
mempunyai akar positif
\[
\delta=\sqrt{4+\varepsilon}-2 .
\]
Untuk $\varepsilon=10^{-3}$ rumus itu memberi $\delta=2.49984376953\times10^{-4}$.
Pemeriksaan di tiga titik:
\[
\begin{array}{lll}
|x-2|=2.497343925761\times10^{-4} & \Rightarrow & |x^{2}-4|=9.989999375719\times10^{-4},\\
|x-2|=2.499843769530\times10^{-4} & \Rightarrow & |x^{2}-4|=1.000000000000\times10^{-3},\\
|x-2|=2.502343613300\times10^{-4} & \Rightarrow & |x^{2}-4|=1.001000062555\times10^{-3}.
\end{array}
\]
Baris kedua menunjukkan bahwa $\delta$ terbesar menghasilkan selisih tepat
$\varepsilon$; baris ketiga, bahwa melewatinya sedikit saja sudah melanggar. Pilihan
populer $\delta=\varepsilon/4=2.5\times10^{-4}$ ternyata \emph{gagal}: selisihnya
$1.0000625\times10^{-3}$, lebih besar dari $\varepsilon$ sebesar $6.25\times10^{-8}$.
Kegagalan sekecil itu mudah terlewat bila diperiksa dengan potongan tiga angka.

Sekarang perhatikan bagaimana $\delta$ bergantung pada titik. Untuk
$\varepsilon=10^{-2}$ dan $f(x)=x^{2}$, akar positif dari $\delta(2x_{0}+\delta)=\varepsilon$
adalah $\delta=\sqrt{x_{0}^{2}+\varepsilon}-x_{0}$.

\begin{center}
\small
\begin{tabular}{rrrr}
\hline
$x_{0}$ & $\delta$ terbesar & $\delta/x_{0}$ & $\varepsilon/(2x_{0})$ \\
\hline
$1$ & $4.987562\times10^{-3}$ & $4.987562\times10^{-3}$ & $5.000000\times10^{-3}$ \\
$2$ & $2.498439\times10^{-3}$ & $1.249220\times10^{-3}$ & $2.500000\times10^{-3}$ \\
$10$ & $4.999875\times10^{-4}$ & $4.999875\times10^{-5}$ & $5.000000\times10^{-4}$ \\
$100$ & $4.999999\times10^{-5}$ & $4.999999\times10^{-7}$ & $5.000000\times10^{-5}$ \\
$1000$ & $5.000000\times10^{-6}$ & $5.000000\times10^{-9}$ & $5.000000\times10^{-6}$ \\
$10^{4}$ & $5.000002\times10^{-7}$ & $5.000002\times10^{-11}$ & $5.000000\times10^{-7}$ \\
\hline
\end{tabular}
\end{center}

Kolom ketiga menyusut secepat $1/x_{0}^{2}$, dan untuk $x_{0}$ besar nilai $\delta$
praktis sama dengan $\varepsilon/(2x_{0})$. Jadi tidak ada batas bawah positif untuk
$\delta$ ketika $x_{0}$ membesar: laju perubahan $x^{2}$ di $x_{0}$ adalah $2x_{0}$, dan
laju itulah yang menentukan ambang, bukan pembulatan.

Kegagalan keseragaman $x^{2}$ pada $\mathbb{R}$ sudah dibuktikan di bagian Kekontinuan
seragam dan diulang mesin pada Lab komputasi bagian 3, jadi angkanya tidak diulang di
sini. Yang belum tampil sebagai gambar adalah ketergantungan pada daerah, dan itu
diperlihatkan oleh $1/x$. Pada $[1,\infty)$ dengan $\delta=\varepsilon$ selisih
terbesar terjadi di $x=1$ dan bernilai $\varepsilon/(1+\varepsilon)$, yaitu
$9.900990\times10^{-3}$ untuk $\varepsilon=10^{-2}$. Pada $(0,1]$ pasangan $x=1/n$,
$y=1/(n+1)$ memberi selisih keluaran tepat $1$ sementara jarak masukannya menyusut:

\begin{center}
\small
\begin{tabular}{rrr}
\hline
$n$ & $|x-y|$ & $|1/x-1/y|$ \\
\hline
$10$ & $9.091\times10^{-3}$ & $1$ \\
$100$ & $9.901\times10^{-5}$ & $1$ \\
$1000$ & $9.990\times10^{-7}$ & $1$ \\
$10^{4}$ & $9.999\times10^{-9}$ & $1$ \\
\hline
\end{tabular}
\end{center}

\kertas{Untuk $f(x)=x^{2}$ dan $\varepsilon=10^{-2}$, hitung $\delta$ terbesar di
$x_{0}=1000$ dari rumus $\sqrt{x_{0}^{2}+\varepsilon}-x_{0}$, lalu bandingkan dengan
$\varepsilon/(2x_{0})$. Periksa bahwa selisih di $x=x_{0}+\delta$ tepat
$\varepsilon$, dan jelaskan mengapa rumus itu tidak dapat dipakai untuk membuktikan
keseragaman pada $\mathbb{R}$.}

\section{Lab komputasi: mengukur keseragaman}
\tingkat{lanjut}

Delapan pemeriksaan berikut mengubah bagian sebelumnya menjadi angka yang dapat
diperiksa ulang: $\delta$ eksplisit dan batasnya, ketergantungannya pada titik,
$x^{2}$ dan $1/x$ pada dua daerah, ambang terbesar di tiap titik, dan belah dua untuk
teorema nilai antara.

\begin{lab}
\begin{verbatim}
import math

# 1. delta optimal untuk f(x)=x^2 di x0=2, eps=1e-3
def delta_sq(x0, eps):
    return math.sqrt(x0*x0 + eps) - x0

eps = 1e-3
d = delta_sq(2.0, eps)
print("== 1. delta eksplisit, x0=2, eps=1e-3 ==")
print("delta optimal = %.12f" % d)
for t in (d*0.999, d, d*1.001):
    print("  x-2=%.12e  |x^2-4|=%.12e" % (t, abs((2+t)**2-4)))
print("  delta=eps/4 -> |x^2-4|=%.12e" % abs((2+eps/4)**2-4))

# 2. delta bergantung titik
print("== 2. delta(x0) untuk eps=1e-2 ==")
print("      x0       delta     delta/x0")
for x0 in (1.0, 2.0, 10.0, 100.0, 1000.0, 1e4):
    d = delta_sq(x0, 1e-2)
    print("%8.0f  %.6e  %.6e" % (x0, d, d/x0))

# 3. kegagalan seragam x^2 pada R
print("== 3. x^2 pada R, eps=1 ==")
print("   delta   x=1/delta   |x^2-(x+d/2)^2|")
for delta in (1e-1, 1e-2, 1e-3, 1e-4):
    x = 1.0/delta
    print("%8.0e %9.0f   %.10f" % (delta, x, abs(x*x-(x+delta/2)**2)))

# 4. x^2 pada [0,1]
print("== 4. x^2 pada [0,1], |x-y|=delta ==")
print("   delta    2d-d^2       kisi")
for delta in (1e-1, 1e-2, 1e-3):
    n = int(round(1/delta))
    g = max(abs((i*delta/n)**2-(i*delta/n+delta)**2)
            for i in range(0, int(n*(1-delta)/delta)+1))
    print("%8.0e  %.10e %.10e" % (delta, 2*delta-delta*delta, g))

# 5. 1/x pada [1,inf) dan pada (0,1]
print("== 5. 1/x pada [1,inf), delta=eps ==")
for eps in (1e-1, 1e-2, 1e-3):
    s = max(abs(1/x-1/(x+eps)) for x in [1+i*eps/500 for i in range(501)])
    print("eps=%.0e  sup=%.10e  eps/(1+eps)=%.10e" % (eps, s, eps/(1+eps)))
print("== 6. 1/x pada (0,1]: x=1/n, y=1/(n+1) ==")
print("      n      |x-y|      |1/x-1/y|")
for n in (10, 100, 1000, 10000):
    print("%7d   %.3e       %.1f" % (n, 1/n-1/(n+1), 1.0))

# 7. delta terbesar yang sah di c, f=1/x, eps=1e-2
print("== 7. delta(c) untuk 1/x, eps=1e-2 ==")
print("       c      delta(c)    delta/c")
for c in (1.0, 1e-1, 1e-2, 1e-3, 1e-4):
    dd = 1e-2*c*c/(1+1e-2*c)
    print("%8.0e  %.4e  %.4e" % (c, dd, dd/c))

# 8. TVA numerik untuk x^3+x-1=0
print("== 8. akar x^3+x-1=0 dengan belah dua ==")
lo, hi = 0.0, 1.0
for k in range(1, 22):
    mid = (lo+hi)/2
    if mid**3+mid-1 < 0:
        lo = mid
    else:
        hi = mid
    if k in (1, 5, 10, 15, 20, 21):
        print("k=%2d  [%.12f, %.12f]  lebar=%.3e"
              % (k, lo, hi, hi-lo))
print("f(lo)=%.3e  f(hi)=%.3e" % (lo**3+lo-1, hi**3+hi-1))
\end{verbatim}
\end{lab}

\newpage
Keluarannya, apa adanya:

\begin{verbatim}
== 1. delta eksplisit, x0=2, eps=1e-3 ==
delta optimal = 0.000249984377
  x-2=2.497343925761e-04  |x^2-4|=9.989999375719e-04
  x-2=2.499843769530e-04  |x^2-4|=1.000000000000e-03
  x-2=2.502343613300e-04  |x^2-4|=1.001000062555e-03
  delta=eps/4 -> |x^2-4|=1.000062499999e-03
== 2. delta(x0) untuk eps=1e-2 ==
      x0       delta     delta/x0
       1  4.987562e-03  4.987562e-03
       2  2.498439e-03  1.249220e-03
      10  4.999875e-04  4.999875e-05
     100  4.999999e-05  4.999999e-07
    1000  5.000000e-06  5.000000e-09
   10000  5.000002e-07  5.000002e-11
== 3. x^2 pada R, eps=1 ==
   delta   x=1/delta   |x^2-(x+d/2)^2|
   1e-01        10   1.0025000000
   1e-02       100   1.0000250000
   1e-03      1000   1.0000002499
   1e-04     10000   1.0000000149
== 4. x^2 pada [0,1], |x-y|=delta ==
   delta    2d-d^2       kisi
   1e-01  1.9000000000e-01 1.9000000000e-01
   1e-02  1.9900000000e-02 1.9900000000e-02
   1e-03  1.9990000000e-03 1.9990000000e-03
== 5. 1/x pada [1,inf), delta=eps ==
eps=1e-01  sup=9.0909090909e-02  eps/(1+eps)=9.0909090909e-02
eps=1e-02  sup=9.9009900990e-03  eps/(1+eps)=9.9009900990e-03
eps=1e-03  sup=9.9900099900e-04  eps/(1+eps)=9.9900099900e-04
== 6. 1/x pada (0,1]: x=1/n, y=1/(n+1) ==
      n      |x-y|      |1/x-1/y|
     10   9.091e-03       1.0
    100   9.901e-05       1.0
   1000   9.990e-07       1.0
  10000   9.999e-09       1.0
== 7. delta(c) untuk 1/x, eps=1e-2 ==
       c      delta(c)    delta/c
   1e+00  9.9010e-03  9.9010e-03
   1e-01  9.9900e-05  9.9900e-04
   1e-02  9.9990e-07  9.9990e-05
   1e-03  9.9999e-09  9.9999e-06
   1e-04  1.0000e-10  1.0000e-06
== 8. akar x^3+x-1=0 dengan belah dua ==
k= 1  [0.500000000000, 1.000000000000]  lebar=5.000e-01
k= 5  [0.656250000000, 0.687500000000]  lebar=3.125e-02
k=10  [0.681640625000, 0.682617187500]  lebar=9.766e-04
k=15  [0.682312011719, 0.682342529297]  lebar=3.052e-05
k=20  [0.682327270508, 0.682328224182]  lebar=9.537e-07
k=21  [0.682327747345, 0.682328224182]  lebar=4.768e-07
f(lo)=-1.354e-07  f(hi)=1.007e-06
\end{verbatim}

Yang perlu dibaca dari keluaran itu ada empat, dan semuanya sudah diperkirakan sebelum
program dijalankan. Bagian 1 memastikan $\delta$ optimal menghasilkan selisih tepat
$10^{-3}$, sedangkan $\delta=\varepsilon/4$ menghasilkan $1.0000625\times10^{-3}$:
pilihan populer itu memang gagal, bukan sekadar kurang rapi. Bagian 2 memastikan
kolom $\delta/x_{0}$ turun secepat $1/x_{0}^{2}$ sampai $x_{0}=1000$, lalu berhenti
mengikuti pola di $x_{0}=10^{4}$. Di titik terakhir nilai sebenarnya
$4.99999999987\times10^{-7}$, sedangkan float64 mencetak $5.000002\times10^{-7}$; itu
galat pembatalan pada $\sqrt{x_{0}^{2}+\varepsilon}-x_{0}$, bukan sifat fungsinya, dan
itulah sebabnya rumus semacam itu tidak dipakai pada skala besar.

Bagian 3 sampai 7 memperlihatkan dua wajah fungsi yang sama. Pada $\mathbb{R}$ selisih
$x^{2}$ bertahan di atas $1$ untuk semua $\delta$ yang dicoba; pada $[0,1]$ selisih
terbesarnya $2\delta-\delta^{2}$, dan kisi mencetak angka yang sama sampai sepuluh
angka karena selisih terbesar terjadi di $x=1-\delta$, titik yang masuk daftar kisi.
Pada $[1,\infty)$ selisih $1/x$ persis $\varepsilon/(1+\varepsilon)$, sedangkan pada
$(0,1]$ ia tetap $1.0$ sementara jarak masukannya menyusut sampai
$9.999\times10^{-9}$. Bagian 8 menutup dengan lebar selang $2^{-21}$ setelah dua puluh
satu belah dua, dan akar $x^{3}+x-1$ terkurung di dalamnya.

Program tidak membuktikan satu pun dari semua itu. Ia memeriksa klaim yang sudah
dibuktikan, dan pada dua tempat menunjukkan batas dirinya: kegagalan
$\delta=\varepsilon/4$ baru terlihat pada angka kedelapan, dan galat pembatalan baru
terlihat pada $x_{0}=10^{4}$.

\begin{untukdosen}
Urutan mengajar yang saya sarankan: mulai dari kisi pada pembuka bab, bukan dari
definisi. Tulis $10000$ dan $5000$ di papan, tanyakan jarak kisi yang diperlukan di
titik itu, lalu bandingkan dengan jawaban di $x=1$; dari kesenjangan itu definisi
$\varepsilon$-$\delta$ muncul sebagai jawaban atas pertanyaan yang sudah terasa, bukan
sebagai kalimat hafalan. Sesudah itu kekontinuan, sifat aljabar, teorema nilai antara,
teorema nilai ekstrem, lalu kekontinuan seragam dan Heine--Cantor.

Dua miskonsepsi yang paling sering muncul. Pertama, mahasiswa menganggap $\delta$
boleh bergantung pada $x$ di dalam definisi kekontinuan, lalu ``membuktikan''
keseragaman dengan memilih $\delta$ berbeda di setiap titik; gejalanya tampak ketika
mereka tidak dapat menyebut satu $\delta$ untuk dua titik berjauhan. Kedua, mereka
menyangka teorema nilai antara memberi cara menghitung akar, padahal ia hanya
menjamin keberadaan.

Aktivitas papan tulis yang berhasil: minta dua mahasiswa memerankan penantang dan
penjawab untuk $f(x)=x^{2}$ di $x_{0}=2$ dan di $x_{0}=1000$ dengan $\varepsilon$
yang sama, lalu suruh kelas menebak di titik mana penjawab kehabisan angka. Tanda
bahwa kelas siap lanjut: mereka dapat menyebut satu $\delta$ eksplisit untuk $1/x$
pada $[1,\infty)$ dan menjelaskan mengapa $\delta$ yang sama gagal di $(0,1]$.
\end{untukdosen}

\section{Ringkasan}
\tingkat{dasar}

\begin{ringkasan}
\begin{itemize}
  \item $\lim_{x\to c}f(x)=L$ berarti untuk setiap $\varepsilon>0$ ada $\delta>0$
        sehingga $0<|x-c|<\delta$ memberi $|f(x)-L|<\varepsilon$. Nilai di $c$
        sendiri tidak ikut dihitung, dan limitnya tunggal.
  \item $f$ kontinu di $c$ bila limitnya ada dan sama dengan $f(c)$. Kekontinuan
        bersifat lokal, tertutup terhadap operasi aljabar, dan bertahan di bawah
        komposisi.
  \item Pada interval tertutup terbatas, fungsi kontinu terbatas, mencapai nilai
        maksimum dan minimum, dan memakai semua nilai di antaranya. Ketiganya dapat
        gagal pada interval terbuka.
  \item $f$ seragam kontinu bila satu $\delta$ berlaku untuk semua pasangan titik.
        Keseragaman lebih kuat daripada kekontinuan, dan keseragaman selalu relatif
        terhadap daerahnya.
  \item $x^{2}$ seragam pada $[0,1]$ tetapi tidak pada $\mathbb{R}$; $1/x$ seragam
        pada $[1,\infty)$ tetapi tidak pada $(0,1]$. Fungsi kontinu pada himpunan
        kompak selalu seragam, dan fungsi Lipschitz selalu seragam.
  \item Kekontinuan dapat gagal di himpunan titik yang besar: indikator rasional gagal
        di mana-mana, $x$ pada rasional dan $0$ pada irasional gagal di semua titik
        kecuali nol.
\end{itemize}
\end{ringkasan}

\section{Latihan}
\tingkat{menengah}

\begin{latihan}
\begin{enumerate}
  \item Buktikan dengan definisi bahwa $\lim_{x\to 3}(2x-1)=5$, dan berikan $\delta$
        eksplisit terkecil sebagai fungsi $\varepsilon$. Jelaskan mengapa
        $\delta=\varepsilon/2$ adalah pilihan terbesar yang sah, sedangkan
        $\delta=\varepsilon/3$ sah tetapi lebih kecil daripada yang diperlukan.
  \item Buktikan bahwa $f(x)=x^{2}$ kontinu di setiap titik $c\in\mathbb{R}$ dengan
        $\delta=\min(1,\varepsilon/(2|c|+1))$, lalu periksa rumus itu di $c=2$ dan
        $c=100$ untuk $\varepsilon=10^{-2}$.
  \item Nilai klaim ini: ``Karena $1/x$ kontinu di setiap titik $(0,1]$, fungsi itu
        seragam kontinu pada $(0,1]$.'' Benar, salah, atau belum cukup bukti? Sebutkan
        bagian yang benar dan bagian yang keliru.
  \item Buktikan bahwa bila $f$ dan $g$ seragam kontinu pada $A$, maka $f+g$ seragam
        kontinu pada $A$. Tunjukkan dengan contoh bahwa $fg$ belum tentu seragam
        kontinu pada $A$ yang tidak terbatas.
  \item Buktikan bahwa $f(x)=\sqrt{x}$ seragam kontinu pada $[0,\infty)$ dengan
        memakai $|\sqrt{x}-\sqrt{y}|\le\sqrt{|x-y|}$. Jelaskan mengapa bukti itu tidak
        memerlukan ketertutupan daerahnya.
  \item Hitung $\delta$ terbesar untuk $f(x)=x^{3}$ di $x_{0}=1$ dengan
        $\varepsilon=10^{-2}$ dengan menyelesaikan $(1+\delta)^{3}-1=\varepsilon$, lalu
        bandingkan dengan $\delta=\varepsilon/3$. Periksa kedua nilai dengan Python.
  \item Untuk $f(x)=x^{2}$ pada $[0,1]$ dan $\varepsilon=10^{-2}$, hitung selisih
        terbesar untuk $\delta=\varepsilon/2$ dan untuk $\delta=1-\sqrt{1-\varepsilon}$.
        Tunjukkan bahwa pilihan pertama aman tetapi tidak optimal, dan bahwa pilihan
        kedua menghasilkan selisih tepat $\varepsilon$.
  \item Buktikan bahwa jika $f$ kontinu pada $[a,b]$ dan $f(x)\neq 0$ untuk semua
        $x\in[a,b]$, maka ada $m>0$ dengan $|f(x)|\ge m$ untuk semua $x$. Gunakan
        teorema nilai ekstrem, dan jelaskan mengapa kesimpulannya gagal untuk
        $f(x)=x$ pada $(0,1)$.
  \item Definisikan $f(x)=x$ bila $x$ rasional dan $f(x)=-x$ bila $x$ irasional.
        Buktikan bahwa $f$ kontinu hanya di $0$, lalu tunjukkan bahwa $|f|$ kontinu di
        setiap titik. Jelaskan mengapa hasil kedua tidak melanggar hasil pertama.
\end{enumerate}

Petunjuk untuk tiga butir tersulit. Butir 4: untuk bagian kedua, ambil $f(x)=g(x)=x$
pada $\mathbb{R}$ dan pakai pasangan $x_{n}=n$, $y_{n}=n+1/n$. Butir 5: kuadratkan
$|\sqrt{x}-\sqrt{y}|\le\sqrt{|x-y|}$ untuk melihat asalnya. Butir 9: $|f|$ tidak lain
adalah $x\mapsto|x|$, yang kontinu di mana-mana; untuk $c\neq0$ pakai barisan rasional
dan irasional yang menuju $c$.
\end{latihan}

\begin{untukpeneliti}
Yang terbukti di bab ini: keseragaman lebih kuat daripada kekontinuan; kekontinuan
pada himpunan kompak memberi keseragaman; kekontinuan pada interval tertutup terbatas
memberi keterbatasan, nilai ekstrem, dan sifat nilai antara. Yang populer tetapi tidak
terbukti adalah kesan bahwa teorema nilai antara memberi algoritme: ia hanya memberi
keberadaan, dan laju $2^{-n}$ berasal dari prosedur, bukan dari teorema.

Sumber primer: Rudin, \emph{Principles of Mathematical Analysis}, Bab 4 bagian
Continuity and Compactness, yang membuktikan Heine--Cantor dalam kerangka ruang metrik
umum; dan Abbott, \emph{Understanding Analysis}, Bab 4 bagian Continuous Functions on
Compact Sets dan Sets of Discontinuity. Bukti Heine--Borel sengaja dilewati di sini dan
ada di Rudin, Bab 2 bagian Compact Sets.

Batas keberlakuan: seluruh bab bekerja di $\mathbb{R}$ dengan jarak biasa. Di ruang
metrik, definisi keseragaman tetap sama kata demi kata, tetapi teorema Heine--Cantor
menuntut kekompakan, bukan sekadar ketertutupan dan keterbatasan; di ruang berdimensi
tak hingga kedua sifat itu berpisah. Di ruang metrik lengkap, keseragaman pada himpunan
padat cukup untuk memperluas fungsi secara tunggal, dan itu pintu masuk ke teorema
ekstensi yang tidak dibahas di sini.

Pertanyaan yang masih terbuka untuk pembaca: seberapa cepat modulus kontinuitas
$\omega(\delta)=\sup_{|x-y|<\delta}|f(x)-f(y)|$ boleh turun. Untuk fungsi kontinu
pada interval tertutup, $\omega(\delta)\to 0$ dijamin, tetapi laju penurunannya tidak:
ada fungsi kontinu yang modulusnya turun lebih lambat daripada setiap pangkat
$\delta^{\alpha}$. Jaminan Lipschitz yang sering diandaikan pada batas serangan
adversarial di Bab 29 karena itu cukup, bukan perlu.
\end{untukpeneliti}

\peta{Bab 11 pada Rudin, \emph{Principles of Mathematical Analysis}: Bab 4 (Continuous
Functions, Continuity and Compactness, Continuity and Connectedness, Discontinuities),
dengan teorema Heine--Borel diambil dari Bab 2 bagian Compact Sets. Pada Abbott,
\emph{Understanding Analysis}: Bab 4 (Discussion: Examples of Dirichlet and Thomae,
Functional Limits, Combinations of Continuous Functions, Continuous Functions on
Compact Sets, The Intermediate Value Theorem, Sets of Discontinuity). Definisi
$\varepsilon$-$\delta$ dan kriteria barisan dibangun dari barisan di Bab 4 buku ini,
sedangkan supremum yang dipakai membuktikan teorema nilai antara berasal dari Bab 3.
Turunan di Bab 12 melanjutkan bab ini dengan memakai hasil bagi selisih sebagai limit,
dan Bab 27 memakai keseragaman sebagai salah satu jaminan yang sering diandaikan
tanpa diperiksa.}

\chapter{Turunan: Satu dan Banyak Variabel}

Sebuah pencatat posisi pada sepeda menulis satu titik setiap detik. Untuk memperoleh
kecepatan, kita membagi selisih jarak dengan selisih waktu. Dua bacaan berjarak satu detik
dan terpisah empat meter memberi empat meter per detik, meskipun itu hanya kecepatan
rata-rata. Bila jaraknya seperseratus detik, selisih jaraknya beberapa sentimeter sedangkan
derau alat beberapa meter; hasil baginya melompat ke ratusan meter per detik, bahkan bertanda
negatif. Pembagian yang sama, dua jawaban yang berbeda.

Turunan menyelesaikan ketegangan itu dengan memindahkan pertanyaan ke limit, tempat derau
tidak masuk hitungan. Tetapi limit tidak dapat diukur: yang terukur hanya hasil bagi pada
jarak berhingga, dan ia membawa dua galat sekaligus, pemotongan karena jaraknya tidak nol
dan pembulatan karena mesin hanya menyimpan sekian digit. Turunan adalah objek matematis
yang tepat; hasil bagi adalah aproksimasinya, dengan galat yang dapat dihitung dan diatur.

Kesulitan kedua muncul ketika masukan lebih dari satu. Sebuah model penilaian kredit
menerima dua belas masukan, dan menggeser satu masukan memberi satu angka laju perubahan.
Menggeser dua masukan sekaligus tidak sama dengan menjumlahkan kedua angka itu, kecuali
fungsinya berperilaku baik; nanti kita lihat contoh yang gagal, dan fungsi itu tetap
kontinu. Jawaban yang benar bukan daftar dua belas angka, melainkan satu peta linear pada vektor
pergeseran: itulah turunan sesungguhnya, dan matriksnya Jacobian. Kekontinuan Bab 11 tidak
mengatakan seberapa cepat sesuatu berubah.

\begin{untukanak}
Pernahkah kamu berdiri di atas atap yang melengkung? Kalau kamu melangkah ke depan, kamu
mendaki. Kalau ke samping, kamu bisa tetap pada ketinggian yang sama. Padahal kamu berdiri di
titik yang sama, di atap yang sama. Jadi pertanyaan ``seberapa curam tempat ini'' belum
lengkap sebelum kita menyebut arah langkah.

Ada masalah lain. Kalau kamu hanya melihat dua titik di atap, lalu membagi selisih ketinggian
dengan selisih jarak, kamu mendapat kemiringan rata-rata, bukan kemiringan di satu titik.
Untuk mendapat kemiringan tepat di bawah kakimu, kedua titik harus didekatkan terus-menerus.
Semakin dekat semakin tepat, tetapi kalau terlalu dekat angka pengukuranmu kacau karena alat
ukur tidak sempurna.

Di bab ini kita akan tahu mengapa kemiringan di satu titik butuh arah, mengapa arah paling
curam selalu ada, dan mengapa alat hitung yang terlalu rapat jaraknya justru salah.
\end{untukanak}

\section{Turunan sebagai limit dan dua artinya}
\tingkat{dasar}

Laju perubahan sesaat adalah bilangan yang didekati hasil bagi dua selisih ketika selisih
masukannya mengecil. Definisi berikut hanya menuliskan kalimat itu dengan hati-hati, karena
hasil baginya sendiri tidak terdefinisi ketika penyebutnya nol.

\kamus{turunan}{derivative}{Kemiringan garis yang paling rapat menempel pada grafik di satu
titik, yang sekaligus menjadi laju perubahan sesaat di titik itu.}

\begin{definisi}[Turunan di satu titik]
Misalkan $f$ terdefinisi pada interval terbuka yang memuat $c$. Turunan $f$ di $c$ adalah
\[
  f'(c)=\lim_{h\to0}\frac{f(c+h)-f(c)}{h},
\]
asalkan limit itu ada dan berhingga. Dalam hal ini $f$ disebut
\emph{terdiferensialkan} di $c$. Notasi lain yang dipakai di buku ini: $Df(c)$ dan
$\dfrac{df}{dx}(c)$.\index{turunan}
\end{definisi}

Hasil bagi itu terdefinisi untuk setiap $h\neq0$, jadi tidak ada pembagian dengan nol; yang
mungkin gagal adalah limitnya. Kegagalan itu nyata: untuk $f(x)=|x|$ di $c=0$ hasil baginya
$1$ untuk $h>0$ dan $-1$ untuk $h<0$, sehingga limit duanya tidak ada.

Secara fisis, turunan adalah laju perubahan sesaat: $s'(t)$ kecepatan bila $s(t)$ posisi,
$Q'(t)$ arus bila $Q(t)$ muatan yang mengalir, dan $C'(q)$ biaya marginal bila $C(q)$ biaya
produksi $q$ satuan.

Secara geometris, hasil bagi adalah kemiringan talibusur dari $(c,f(c))$ ke
$(c+h,f(c+h))$. Ketika $h\to0$, talibusur itu berputar mengelilingi titik pertama menuju satu
posisi; garis pada posisi limit disebut garis singgung, dan kemiringannya $f'(c)$. Persamaan
garis singgung di $c$ adalah
\[
  y=f(c)+f'(c)(x-c).
\]
Garis singgung bukan ``garis yang menyentuh grafik tepat satu kali'', sebab untuk
$f(x)=x^{3}$ di $c=0$ garis singgung $y=0$ memotong grafik di titik asal. Yang benar:
selisihnya terhadap grafik mengecil lebih cepat daripada jaraknya.

Putaran itu dapat diikuti angkanya. Untuk $f(x)=x^{2}$ di $c=2$, talibusur dengan $h=1$
berkemiringan $5$, dengan $h=0{,}5$ kemiringan $4{,}5$, dan garis singgungnya $4$.
Gambar~\ref{fig:sekan-singgung} menggambar ketiganya: dari $5$ ke $4{,}5$, menetap di $4$.

\begin{figure}[htbp]
\centering
\begin{tikzpicture}
\begin{axis}[
  width=\linewidth, height=6.4cm,
  axis lines=left,
  xlabel={$x$}, ylabel={$y$},
  xmin=1.55, xmax=3.25, ymin=2.2, ymax=10.2,
  xtick={2,3}, ytick={4,9},
  tick label style={font=\footnotesize},
  label style={font=\footnotesize},
  legend style={font=\scriptsize, draw=none, fill=none,
                at={(0.03,0.97)}, anchor=north west},
  clip=false,
]
\addplot[ink, thick, domain=1.6:3.15, samples=90] {x^2};
\addlegendentry{$f(x)=x^{2}$}
\addplot[accent, thick, domain=1.95:3.18] {5*(x-2)+4};
\addlegendentry{talibusur $h=1$, kemiringan $5$}
\addplot[layerteal, thick, dashed, domain=1.95:3.0] {4.5*(x-2)+4};
\addlegendentry{talibusur $h=0{,}5$, kemiringan $4{,}5$}
\addplot[layerindigo, thick, domain=1.68:3.05] {4*(x-2)+4};
\addlegendentry{garis singgung, kemiringan $4$}
\addplot[only marks, mark=*, mark size=1.5pt, ink]
  coordinates {(2,4) (2.5,6.25) (3,9)};
\node[font=\scriptsize, anchor=north west, text=ink] at (axis cs:2.04,3.75) {$(2,4)$};
\end{axis}
\end{tikzpicture}
\caption{Talibusur yang berputar menjadi garis singgung pada $f(x)=x^{2}$ di $(2,4)$.
Kemiringannya turun dari $5$ (untuk $h=1$) ke $4{,}5$ (untuk $h=0{,}5$) dan menetap di $4$
ketika $h\to0$.}
\label{fig:sekan-singgung}
\end{figure}

\begin{bacakeras}
Dua suara bergantian; jawabannya makin pendek setiap kali.

--- Berapa cepat dia bergerak? --- Dari titik mana ke titik mana?

--- Dari detik pertama ke detik kedua. --- Tiga meter.

--- Dari detik pertama ke detik satu setengah. --- Dua setengah meter.

--- Dari detik pertama ke detik satu sepersepuluh. --- Dua satu per lima.

--- Dari detik pertama ke detik satu seperseratus. --- Hampir dua.

--- Dari detik pertama ke detik pertama. --- Tidak bisa. Penyebutnya nol.

Berhenti di situ. Yang berubah bukan jawabannya, melainkan jarak antara dua titik yang
dibandingkan. Jarak mengecil, garis potong berputar, angka menetap: nilai yang tidak pernah
dicapai tetapi selalu didekati, itulah turunan.

Ulangi bersama: jarak mengecil, garis berputar, angka menetap.
\end{bacakeras}

Mari hitung satu turunan dari definisi. Untuk $f(x)=x^{2}$ di $c=2$,
\[
  \frac{f(2+h)-f(2)}{h}=\frac{(4+4h+h^{2})-4}{h}=\frac{h(4+h)}{h}=4+h .
\]
Karena $h\neq0$, pembagian itu sah, dan limitnya $4$. Untuk $h=0{,}01$ hasil baginya
$4{,}01$; selisih $0{,}01$ itu bukan kesalahan hitung, melainkan suku orde $h$ yang hilang
hanya di limit. Untuk $f(x)=1/x$ di $c=2$,
\[
  \frac{1/(2+h)-1/2}{h}
  =\frac{2-(2+h)}{2h(2+h)}=\frac{-1}{2(2+h)}\longrightarrow-\frac14 .
\]
Pada kedua hitungan aljabar menghapus faktor $h$ dari penyebut, sehingga bentuk $0/0$
berubah menjadi fungsi yang kontinu di $h=0$.

Bentuk hasil bagi itu penting, tetapi tidak bertahan ketika variabelnya menjadi banyak.
Bentuk yang bertahan adalah aproksimasi linear.

\begin{definisi}[Turunan sebagai koefisien aproksimasi linear]
Fungsi $f$ terdiferensialkan di $c$ dengan turunan $L$ bila
\[
  f(c+h)=f(c)+Lh+r(h),
  \qquad\text{dengan}\qquad
  \lim_{h\to0}\frac{r(h)}{h}=0 .
\]
\end{definisi}

Sisa $r(h)$ lebih kecil daripada $h$; itulah isi definisinya. Untuk satu variabel bentuk ini
terasa berlebihan, tetapi hanya bentuk inilah yang bertahan sampai bagian Jacobian.

\begin{proposisi}
Kedua definisi turunan di atas setara, dan bilangan $L$ pada definisi kedua tunggal.
\end{proposisi}

\begin{proof}
Bila limit hasil bagi sama dengan $L$, tetapkan $r(h)=f(c+h)-f(c)-Lh$. Maka $r(h)/h$ tepat
sama dengan hasil bagi dikurangi $L$, sehingga menuju $0$. Sebaliknya, bila $r$ dengan sifat
itu ada, hasil baginya sama dengan $L+r(h)/h$, yang menuju $L$.

Untuk ketunggalan, misalkan $L_{1}$ dan $L_{2}$ sama-sama memenuhi. Maka
$(L_{1}-L_{2})h=r_{2}(h)-r_{1}(h)$, dan setelah dibagi $h$ kita peroleh
$L_{1}-L_{2}=r_{2}(h)/h-r_{1}(h)/h$. Ruas kanan menuju $0$, sedangkan ruas kiri tidak
bergantung pada $h$. Jadi $L_{1}=L_{2}$.
\end{proof}

Sifat pertama yang perlu dipegang: diferensiabilitas lebih kuat daripada kekontinuan.

\begin{proposisi}
Bila $f$ terdiferensialkan di $c$, maka $f$ kontinu di $c$.
\end{proposisi}

\begin{proof}
Untuk $h\neq0$ berlaku
\[
  f(c+h)-f(c)=h\cdot\frac{f(c+h)-f(c)}{h}.
\]
Faktor kedua menuju $f'(c)$ dan faktor pertama menuju $0$, jadi hasil kalinya menuju $0$.
Dengan kata lain $\lim_{h\to0}f(c+h)=f(c)$.
\end{proof}

Kebalikannya tidak berlaku. Fungsi $|x|$ kontinu di $0$ tetapi tidak terdiferensialkan di
sana, karena hasil bagi duanya berbeda tanda. Fungsi $x\sin(1/x)$, bernilai $0$ di $x=0$,
juga kontinu tetapi tidak terdiferensialkan di sana: hasil baginya $\sin(1/x)$, yang berayun
tanpa menuju satu nilai.

\begin{tantangan}
Misalkan $f(0)=0$ dan, untuk $x\neq0$,
\[
  f(x)=|x|\bigl(2+\sin(1/x)\bigr).
\]
Tunjukkan bahwa $f$ kontinu di $0$ tetapi tidak punya turunan di sana. Kedua sifat tidak
bertentangan: kekontinuan menuntut nilai fungsi mendekati $f(0)$, turunan menuntut hasil
bagi selisih mendekati satu bilangan.

Petunjuk 1: untuk kekontinuan, cukup satu ketaksamaan yang mengapit $|f(x)|$ di antara $0$
dan sesuatu yang menuju $0$.

Petunjuk 2: untuk yang kedua, hitung $f(h)/h$ terpisah bagi $h>0$ dan $h<0$. Pada $h>0$
hasil baginya $2+\sin(1/h)$, yang berayun antara $1$ dan $3$; pada $h<0$ ia berayun antara
$-3$ dan $-1$. Kedua kumpulan nilai itu tidak beririsan.

Kalau mentok, kembali ke definisi turunan di awal bagian ini dan hitung hasil baginya untuk
$h>0$ lalu $h<0$ pada fungsi $|x|$.
\end{tantangan}

Aturan hitung berikut membebaskan kita dari definisi setiap kali.

\begin{teorema}[Aturan dasar]
Misalkan $f$ dan $g$ terdiferensialkan di $c$. Maka
\begin{enumerate}
  \item $(f+g)'(c)=f'(c)+g'(c)$ dan $(af)'(c)=a f'(c)$ untuk $a\in\mathbb{R}$;
  \item $(fg)'(c)=f'(c)g(c)+f(c)g'(c)$;
  \item $\left(\dfrac{f}{g}\right)'(c)=\dfrac{g(c)f'(c)-f(c)g'(c)}{g(c)^{2}}$ bila
        $g(c)\neq0$;
  \item bila $n$ bilangan bulat, turunan $x^{n}$ di $x\neq0$ adalah $nx^{n-1}$.
\end{enumerate}
\end{teorema}

\begin{proof}
Bagian (1) langsung dari limit jumlah. Untuk (2), tulis
\[
  f(x)g(x)-f(c)g(c)=f(x)\bigl[g(x)-g(c)\bigr]+g(c)\bigl[f(x)-f(c)\bigr].
\]
Setelah dibagi $x-c$ dan $x\to c$, faktor $f(x)$ menuju $f(c)$ karena $f$ kontinu di $c$
(proposisi sebelumnya), dan kedua hasil bagi menuju turunannya.

Untuk (3), cukup ditunjukkan bahwa $1/g$ terdiferensialkan dengan turunan $-g'(c)/g(c)^{2}$,
lalu pakai (2) pada $f\cdot(1/g)$. Hasil bagi untuk $1/g$ adalah
\[
  \frac{1/g(x)-1/g(c)}{x-c}
  =-\frac{1}{g(x)g(c)}\cdot\frac{g(x)-g(c)}{x-c},
\]
yang menuju $-g'(c)/g(c)^{2}$ karena $g$ kontinu dan $g(c)\neq0$. Bagian (4) diperoleh dari
(2) dengan induksi pada $n\ge0$, lalu dari (3) untuk $n<0$.
\end{proof}

\section{Aturan rantai}
\tingkat{menengah}

Bila satu besaran bergantung pada besaran kedua, dan besaran kedua bergantung pada waktu,
maka laju perubahan besaran pertama adalah hasil kali dua laju itu. Kalimat itu terdengar
seperti aritmetika sehari-hari, dan memang itulah isi aturan rantai.

Gambar~\ref{fig:roda-gigi} memakai dua roda gigi: roda dalam berputar tiga kali per putaran
roda tengah, dan roda tengah dua kali per putaran roda luar. Satu putaran roda luar memutar
roda dalam $3\times2=6$ kali; menjumlahkan kedua rasio memberi $5$, angka yang tidak muncul
pada putaran mana pun.

\begin{figure}[htbp]
\centering
\begin{tikzpicture}[scale=0.95]
% roda dalam, rasio 3
\draw[ink, thick] (0,0) circle (1.05);
\foreach \a in {0,20,...,340} { \draw[ink, thick] (\a:1.05) -- (\a:1.17); }
\fill[ink] (0,0) circle (1.5pt);
\draw[->, layerindigo, thick] (0,0) ++(150:0.55) arc (150:70:0.55);
% roda luar, rasio 2
\begin{scope}[xshift=2.9cm]
  \draw[ink, thick] (0,0) circle (1.8);
  \foreach \a in {10,30,...,350} { \draw[ink, thick] (\a:1.8) -- (\a:1.95); }
  \fill[ink] (0,0) circle (1.5pt);
  \draw[->, layerbrown, thick] (0,0) ++(-20:0.75) arc (-20:-100:0.75);
\end{scope}
\node[font=\footnotesize, text=ink, align=center] at (0,-2.35) {roda dalam:\\3 putaran};
\node[font=\footnotesize, text=ink, align=center] at (2.9,-2.35) {roda luar:\\2 putaran};
\node[font=\footnotesize, text=layerplum] at (1.45,2.15) {$3\times2=6$};
\end{tikzpicture}
\caption{Aturan rantai sebagai dua roda gigi. Roda dalam berputar $3$ kali untuk satu putaran
roda luar, dan roda luar $2$ kali untuk satu putaran keluarannya; rasio totalnya
$3\times2=6$, bukan $3+2=5$.}
\label{fig:roda-gigi}
\end{figure}

\kamus{aturan rantai}{chain rule}{Laju di ujung rantai sama dengan hasil kali laju pada
setiap mata rantai.}

\begin{teorema}[Aturan rantai]
Misalkan $g$ terdiferensialkan di $c$ dan $f$ terdiferensialkan di $g(c)$. Maka $f\circ g$
terdiferensialkan di $c$ dan
\[
  (f\circ g)'(c)=f'\bigl(g(c)\bigr)\,g'(c).
\]
\end{teorema}

\begin{proof}
Tulis $y=g(c)$. Karena $f$ terdiferensialkan di $y$, menurut definisi aproksimasi linear
ada fungsi $v$ dengan $v(s)\to0$ untuk $s\to y$ dan
\[
  f(s)-f(y)=(s-y)\bigl[f'(y)+v(s)\bigr].
\]
Dengan cara yang sama, karena $g$ terdiferensialkan di $c$, ada $u$ dengan $u(t)\to0$ untuk
$t\to c$ dan
\[
  g(t)-g(c)=(t-c)\bigl[g'(c)+u(t)\bigr].
\]
Substitusikan $s=g(t)$ pada persamaan pertama, lalu pakai persamaan kedua:
\[
  f\bigl(g(t)\bigr)-f\bigl(g(c)\bigr)
  =(t-c)\bigl[g'(c)+u(t)\bigr]\bigl[f'(y)+v(g(t))\bigr].
\]
Untuk $t\neq c$, bagi dengan $t-c$. Karena $g$ kontinu di $c$, $g(t)\to y$, sehingga
$v(g(t))\to0$, dan hasil baginya menuju $g'(c)f'(y)$.
\end{proof}

Pembuktian naif yang membagi dengan $g(t)-g(c)$ gagal ketika $g(t)=g(c)$ untuk $t$ dekat
$c$ tetapi $t\neq c$; fungsi konstan lokal, misalnya, membuat penyebut itu nol di semua
titik. Penyangga $v$ di atas menghindari pembagian tersebut sama sekali.

Contoh pendek: untuk $h(x)=(x^{2}+1)^{5}$, ambil $g(x)=x^{2}+1$ dan $f(u)=u^{5}$. Maka
$h'(x)=5(x^{2}+1)^{4}\cdot2x$ dan $h'(1)=5\cdot16\cdot2=160$; ekspansi langsung mungkin,
tetapi tidak menyenangkan.

\begin{resep}
\begin{enumerate}
  \item \textbf{Pisahkan dua lapis.} Tunjuk mana yang di dalam dan mana yang di luar. Untuk
        $k(x)=\sin(x^{3}+2x)$: lapis dalam $g(x)=x^{3}+2x$, lapis luar $f(u)=\sin u$.
  \item \textbf{Turunkan lapis luar, biarkan isinya utuh.} $f'(u)=\cos u$ menjadi
        $\cos(x^{3}+2x)$; jangan dulu menyederhanakan isinya.
  \item \textbf{Kalikan turunan lapis dalam.} $g'(x)=3x^{2}+2$, sehingga
        $k'(x)=(3x^{2}+2)\cos(x^{3}+2x)$.
  \item \textbf{Uji satu titik.} Di $x=0$ rumus memberi $2$, dan beda pusat dengan
        $h=10^{-3}$ memberi $1{,}9999997$.
\end{enumerate}
Peringatan urutan: lapis luar dinilai di $g(x)$, bukan di $x$. Menukar tempat penilaian
memberi $\cos(x)\,(3x^{2}+2)$, yang di $x=1$ bernilai $2{,}7015$, sedangkan nilai benarnya
$-4{,}9499$.
\end{resep}

Contoh laju: luas lingkaran $A=\pi r^{2}$ dengan $\dot r=0{,}5$ cm per detik. Di $r=10$ cm,
\[
  \frac{dA}{dt}=2\pi r\frac{dr}{dt}=2\pi\cdot10\cdot0{,}5=10\pi\approx31{,}4159
  \ \text{cm}^{2}\ \text{per detik}.
\]
Angka $10\pi$ eksak; di $r=20$ cm laju yang sama memberi $20\pi$: jawabannya bergantung
pada $r$.

\section{Teorema nilai rata-rata dan akibatnya}
\tingkat{menengah}

Di antara dua titik pada grafik yang mulus selalu ada satu titik yang kemiringannya sama
dengan kemiringan talibusur yang menghubungkan kedua titik itu. Pernyataan itu terlihat
seperti khayalan geometris, tetapi dialah jembatan antara turunan dan fungsi.

Kunci pembuktiannya adalah kasus khusus ketika kedua ujung bernilai sama.

\begin{teorema}[Rolle]
Bila $f$ kontinu pada $[a,b]$, terdiferensialkan pada $(a,b)$, dan $f(a)=f(b)$, maka ada
$c\in(a,b)$ dengan $f'(c)=0$.\index{teorema nilai rata-rata}
\end{teorema}

\begin{proof}
Bila $f$ konstan, setiap $c\in(a,b)$ memenuhi. Bila tidak, $f$ mencapai nilai maksimum atau
minimum yang berbeda dari nilai ujungnya; karena $f$ kontinu pada interval tertutup terbatas,
nilai itu dicapai di suatu titik, dan titik itu tidak boleh ujung, jadi $c\in(a,b)$ (teorema
nilai ekstrem, Bab 11). Di titik ekstrem interior, hasil bagi
$(f(c+h)-f(c))/h$ bernilai $\le0$ untuk $h>0$ kecil dan $\ge0$ untuk $h<0$ kecil. Kedua
ketaksamaan itu bersama memaksa limit duanya, yaitu $f'(c)$, bernilai $0$.
\end{proof}

\begin{teorema}[Nilai rata-rata]
Bila $f$ kontinu pada $[a,b]$ dan terdiferensialkan pada $(a,b)$, maka ada $c\in(a,b)$
dengan
\[
  f'(c)=\frac{f(b)-f(a)}{b-a}.
\]
\end{teorema}

\begin{proof}
Tetapkan $s=(f(b)-f(a))/(b-a)$ dan $g(x)=f(x)-f(a)-s(x-a)$. Maka $g$ kontinu pada $[a,b]$,
terdiferensialkan pada $(a,b)$, dan $g(a)=g(b)=0$. Menurut teorema Rolle ada $c\in(a,b)$
dengan $g'(c)=0$, yaitu $f'(c)=s$.
\end{proof}

Contoh angka: $f(x)=x^{3}$ pada $[0,3]$ berkemiringan talibusur $27/3=9$, dan $3c^{2}=9$
memberi $c=\sqrt{3}\approx1{,}7321$; tidak ada jaminan $c$ di tengah.

\begin{tangan}
Siapkan papan tipis 30 cm, dua buku, penggaris, dan pensil. Ganjal ujung kiri papan setinggi
4 cm, lalu tekan sedikit bagian tengahnya supaya melengkung ke atas.

Tegakkan penggaris pada 5 cm dan 25 cm dari ujung yang diangkat; selisih tinggi dibagi 20 cm
memberi satu angka kemiringan. Ulangi dengan 5 cm dan 10 cm: angkanya berubah, sebab yang
diukur adalah kemiringan rata-rata pada selang berbeda.

Sekarang letakkan penggaris rata di titik 15 cm dan putar sampai hanya satu titik yang
menempel. Angkat ujung yang lain perlahan: penggaris berputar pada titik itu, persis seperti
talibusur pada Gambar~\ref{fig:sekan-singgung}. Kemiringan penggaris yang menyentuh satu
titik itulah kemiringan di titik itu, dan di situlah teorema nilai rata-rata menjanjikan
adanya satu titik semacam itu.
\end{tangan}

\begin{akibat}[Uji monotonisitas]
Misalkan $f$ kontinu pada $[a,b]$ dan terdiferensialkan pada $(a,b)$. Bila $f'(x)\ge0$ untuk
semua $x\in(a,b)$, maka $f$ tidak turun pada $[a,b]$. Bila $f'(x)>0$ untuk semua $x$, maka
$f$ naik tegas. Bila $f'(x)=0$ untuk semua $x$, maka $f$ konstan.
\end{akibat}

\begin{proof}
Ambil $a\le x<y\le b$. Teorema nilai rata-rata pada $[x,y]$ memberi $c\in(x,y)$ dengan
$f(y)-f(x)=f'(c)(y-x)$. Karena $y-x>0$, tanda selisih nilai mengikuti tanda $f'(c)$. Bila
$f'$ nol di mana-mana, selisihnya nol untuk setiap pasangan, jadi $f$ konstan.
\end{proof}

Arah logikanya perlu diperhatikan: akibat ini mengubah informasi tentang turunan menjadi
informasi tentang fungsi. Kebalikannya tidak berlaku: fungsi yang naik tegas boleh punya
turunan nol di beberapa titik, seperti $x^{3}$ di $0$. Juga tidak berlaku kesimpulan balik
dari satu titik: $f'(c)=0$ tidak berarti $c$ titik ekstrem.

\begin{akibat}[Keterkendalian]
Bila $f$ terdiferensialkan pada interval $I$ dan $|f'(x)|\le M$ untuk semua $x\in I$, maka
$|f(x)-f(y)|\le M|x-y|$ untuk semua $x,y\in I$.
\end{akibat}

\begin{proof}
Untuk $x\neq y$, teorema nilai rata-rata pada interval di antara keduanya memberi $c$
dengan $f(x)-f(y)=f'(c)(x-y)$. Ambil nilai mutlak kedua ruas dan pakai $|f'(c)|\le M$.
\end{proof}

Karena itu fungsi dengan turunan terbatas otomatis seragam kontinu (ambil
$\delta=\varepsilon/M$), tetapi kebalikannya tidak berlaku: $\sqrt{x}$ seragam kontinu pada
$[0,1]$ dengan turunan tak terbatas di dekat nol.

\section{Polinomial Taylor dan sisanya}
\tingkat{lanjut}

Garis singgung adalah polinomial berderajat satu yang menempel pada grafik di satu titik.
Bila kita mau lebih lama menempel, kita tambah suku berderajat dua, tiga, dan seterusnya.
Pertanyaannya bukan apakah penambahan itu memperbaiki hampiran, tetapi seberapa besar
sisanya.

\begin{definisi}[Polinomial Taylor]
Misalkan $f$ mempunyai turunan sampai orde $n$ di $a$. Polinomial Taylor orde $n$ dari $f$
di $a$ adalah
\[
  P_{n}(x)=\sum_{k=0}^{n}\frac{f^{(k)}(a)}{k!}(x-a)^{k}.
\]
\end{definisi}

Syaratnya: $P_{n}$ dan $f$ sama nilai dan $n$ turunan pertamanya di $a$; karena $n+1$
koefisien menentukan satu polinomial berderajat $n$, $P_{n}$ tunggal.

\begin{teorema}[Taylor dengan sisa Lagrange]
Misalkan $f$ mempunyai turunan sampai orde $n-1$ yang kontinu pada $[a,b]$ dan turunan
ke-$n$ pada $(a,b)$. Untuk setiap $\beta\in(a,b]$ ada $\xi\in(a,\beta)$ dengan
\[
  f(\beta)=P_{n-1}(\beta)+\frac{f^{(n)}(\xi)}{n!}(\beta-a)^{n}.
\]
\end{teorema}

\begin{proof}
Tetapkan $M$ oleh $f(\beta)=P_{n-1}(\beta)+M(\beta-a)^{n}$, lalu
\[
  g(t)=f(t)-P_{n-1}(t)-M(t-a)^{n},\qquad a\le t\le b .
\]
$P_{n-1}$ dibangun dari turunan $f$ di $a$, jadi $g(a)=\cdots=g^{(n-1)}(a)=0$, dan pilihan
$M$ memberi $g(\beta)=0$. Rolle memberi $x_{1}\in(a,\beta)$ dengan $g'(x_{1})=0$; karena
$g'(a)=0$, Rolle pada $g'$ di $[a,x_{1}]$ memberi $x_{2}$ dengan $g''(x_{2})=0$. Ulangi sampai
$n$ langkah: ada $x_{n}\in(a,\beta)$ dengan $g^{(n)}(x_{n})=0$. Turunan ke-$n$ dari $P_{n-1}$
nol dan dari $(t-a)^{n}$ adalah $n!$, sehingga $0=f^{(n)}(x_{n})-n!M$, yaitu
$n!M=f^{(n)}(\xi)$ dengan $\xi=x_{n}$.
\end{proof}

Untuk $n=1$ teorema ini tepat teorema nilai rata-rata; yang baru adalah kemampuan menghitung
galat orde tinggi.

\begin{contoh}
Hampiri $e$ dengan $P_{3}(1)$ dari $f(x)=e^{x}$ di $a=0$, lalu bandingkan dengan batas
sisanya.
\end{contoh}

Di sini $f^{(k)}(0)=1$ untuk semua $k$, sehingga
\[
  P_{3}(1)=1+1+\frac{1}{2}+\frac{1}{6}=\frac{8}{3}=2{,}6666\ldots
\]
Nilai sebenarnya $e=2{,}71828\ldots$, jadi galatnya $0{,}0516152$. Sisa Lagrange memberi
$|e-P_{3}(1)|=e^{\xi}/24$ untuk suatu $\xi\in(0,1)$, dan karena $e^{\xi}\le e<3$, batasnya
$3/24=0{,}125$. Batas itu benar, tetapi longgar sekitar dua setengah kali. Memakai orde
lima, $P_{5}(1)=2{,}71666\ldots$ dengan galat $0{,}0016152$, sedangkan batasnya
$3/720=0{,}0041667$.

\begin{untukanak}
Bayangkan kamu menebak jalan kawanmu hanya dari beberapa petunjuk. Petunjuk pertama: di mana
dia sekarang. Tebakanmu satu titik, dan titik itu tidak bergerak. Petunjuk kedua: ke arah mana
dia melangkah. Tebakanmu garis lurus, sudah lebih mirip, tetapi jalan sungguhan biasanya
melengkung. Petunjuk ketiga: seberapa cepat arahnya berubah. Tebakanmu potongan parabola, yang
menempel lebih lama.

Semakin banyak petunjuk, semakin lama tebakanmu menempel, tetapi tidak pernah sama persis.
Selalu ada sisa. Sisa itu dapat dibatasi tanpa mengetahui jalan sebenarnya, hanya dari
seberapa besar perubahan pada tingkat berikutnya. Di bab ini kita akan tahu mengapa garis
lurus tebakan orde satu, mengapa parabola orde dua, dan mengapa batas sisa lebih berguna
daripada jawaban yang katanya tepat.
\end{untukanak}

\section{Turunan parsial, gradien, dan turunan berarah}
\tingkat{menengah}

Bila masukan lebih dari satu, ``laju perubahan'' belum tentu satu bilangan. Ia bergantung
pada arah kita melangkah. Cara paling sederhana mendapat satu bilangan adalah membekukan
semua koordinat kecuali satu, lalu memakai definisi yang sudah kita punya.

\begin{definisi}[Turunan parsial]
Misalkan $f\colon U\to\mathbb{R}$ terdefinisi pada himpunan terbuka
$U\subseteq\mathbb{R}^{n}$ dan $a\in U$. Turunan parsial $f$ terhadap $x_{j}$ di $a$ adalah
\[
  \frac{\partial f}{\partial x_{j}}(a)
  =\lim_{h\to0}\frac{f(a+h\,e_{j})-f(a)}{h},
\]
dengan $e_{j}$ vektor satuan arah $x_{j}$, asalkan limitnya ada.\index{turunan parsial}
\end{definisi}

Secara geometris, turunan parsial terhadap $x$ adalah kemiringan irisan grafik oleh bidang
$y$ tetap, jadi turunan biasa fungsi satu variabel.

\kamus{gradien}{gradient}{Panah yang menunjuk arah paling curam pada sebuah bukit, dengan
panjang sebesar kecuraman di arah itu.}

\begin{definisi}[Gradien]
Bila semua turunan parsial $f$ di $a$ ada, vektor
\[
  \nabla f(a)=\left(\frac{\partial f}{\partial x_{1}}(a),\dots,
  \frac{\partial f}{\partial x_{n}}(a)\right)
\]
disebut gradien $f$ di $a$.\index{gradien}
\end{definisi}

Untuk $f(x,y)=x^{2}y^{3}$ berlaku $\partial f/\partial x=2xy^{3}$ dan
$\partial f/\partial y=3x^{2}y^{2}$, sehingga $\nabla f(2,-1)=(-4,12)$.

\begin{definisi}[Turunan berarah]
Untuk vektor satuan $u\in\mathbb{R}^{n}$, turunan berarah $f$ di $a$ sepanjang $u$ adalah
\[
  D_{u}f(a)=\lim_{t\to0}\frac{f(a+tu)-f(a)}{t},
\]
asalkan limitnya ada.\index{turunan berarah}
\end{definisi}

Turunan parsial adalah kasus khusus $u=e_{j}$. Untuk arah lain $D_{u}f$ tidak dapat diduga
dari daftar parsial, kecuali dalam satu keadaan baik.

\begin{proposisi}
Bila $f$ terdiferensialkan di $a$ (dalam pengertian aproksimasi linear di bagian berikut),
maka untuk setiap vektor satuan $u$ berlaku $D_{u}f(a)=\nabla f(a)\cdot u$.
\end{proposisi}

\begin{proof}
Dengan $v=tu$ pada definisi aproksimasi linear,
$f(a+tu)=f(a)+t\,\nabla f(a)\cdot u+r(tu)$, sehingga
\[
  \frac{f(a+tu)-f(a)}{t}=\nabla f(a)\cdot u+\frac{r(tu)}{t}.
\]
Karena $\|r(tu)\|/\|tu\|\to0$ dan $\|tu\|=|t|$, suku terakhir menuju $0$.
\end{proof}

\begin{proposisi}[Arah naik tercepat]
Bila $\nabla f(a)\neq0$, maka di antara semua vektor satuan $u$, nilai $D_{u}f(a)$ terbesar
dicapai pada $u=\nabla f(a)/\|\nabla f(a)\|$, dan nilai terbesarnya $\|\nabla f(a)\|$.
\end{proposisi}

\begin{proof}
Menurut ketaksamaan Cauchy--Schwarz (Bab 10),
\[
  |D_{u}f(a)|=|\nabla f(a)\cdot u|\le\|\nabla f(a)\|\,\|u\|=\|\nabla f(a)\|.
\]
Kesamaan berlaku tepat ketika $u$ sejajar dengan $\nabla f(a)$. Karena $u$ vektor satuan,
hanya dua tanda yang mungkin pada arah itu, dan tanda positif memberi nilai terbesar.
\end{proof}

Untuk $f(x,y)=x^{2}y^{3}$ di $(2,-1)$, panjang gradiennya $\sqrt{160}\approx12{,}6491$, jadi
arah naik tercepat $(-4,12)/\sqrt{160}\approx(-0{,}3162,\,0{,}9487)$ dengan laju $12{,}6491$
per satuan jarak. Arah $x$ memberi $-4$, arah $y$ memberi $12$; keduanya lebih kecil.

Bila kurva $t\mapsto\gamma(t)$ tinggal pada himpunan nivel $f=\text{tetap}$, maka
$f(\gamma(t))$ konstan dan turunannya nol; aturan rantai memberi
$\nabla f(\gamma(t))\cdot\gamma'(t)=0$. Jadi gradien tegak lurus garis nivel, dan arah naik
tercepat tegak lurus garis kontur pada peta.

Contohnya $f(x,y)=x^{2}+2y^{2}$ di $(2,-1)$, titik dengan $f=6$, sehingga ia duduk pada
elips nivel $x^{2}+2y^{2}=6$. Gradiennya $(2x,4y)=(4,-4)$, panjangnya
$4\sqrt{2}\approx5{,}6569$, sedangkan arah singgung elips di sana $(1,1)$ dan
$(4,-4)\cdot(1,1)=0$. Gambar~\ref{fig:kontur-gradien} menggambar ketiganya: elips nivel,
panah gradien, dan garis singgung yang tegak lurus padanya. Panah gradien tidak menyusuri
kontur; ia menembusnya ke arah nilai $f$ bertambah.

\begin{figure}[htbp]
\centering
\begin{tikzpicture}[scale=0.92]
\draw[->, layergrey] (-3.0,0) -- (3.4,0) node[right, font=\scriptsize] {$x$};
\draw[->, layergrey] (0,-2.45) -- (0,2.5) node[above, font=\scriptsize] {$y$};
\draw[ink] (0,0) ellipse [x radius=1.414, y radius=1.000];
\draw[ink] (0,0) ellipse [x radius=2.000, y radius=1.414];
\draw[ink, thick] (0,0) ellipse [x radius=2.449, y radius=1.732];
\node[font=\scriptsize, anchor=west, text=ink] at (0.28,-0.98) {$c=2$};
\node[font=\scriptsize, anchor=west, text=ink] at (0.28,-1.40) {$c=4$};
\node[font=\scriptsize, anchor=west, text=ink] at (0.28,-1.98) {$c=6$};
\draw[accent, thick, dashed] (1.30,-1.70) -- (2.72,-0.28);
\draw[->, very thick, layerindigo] (2,-1) -- (2.707,-1.707);
\fill[ink] (2,-1) circle (1.7pt);
\node[font=\scriptsize, anchor=north west, text=ink] at (2.02,-1.05) {$(2,-1)$};
\node[font=\scriptsize, anchor=west, text=layerindigo] at (2.72,-1.86) {$\nabla f=(4,-4)$};
\node[font=\scriptsize, anchor=east, text=accent] at (1.32,-1.76) {arah $(1,1)$};
\end{tikzpicture}
\caption{Kurva nivel $x^{2}+2y^{2}=c$ untuk $c=2,4,6$ dan gradien di $(2,-1)$. Panah
gradien $(4,-4)$ tegak lurus arah singgung $(1,1)$ pada kurva nivel $c=6$, sebab
$(4,-4)\cdot(1,1)=0$.}
\label{fig:kontur-gradien}
\end{figure}

\section{Diferensiabilitas, Jacobian, dan aturan rantai multivariabel}
\tingkat{lanjut}

Turunan yang sesungguhnya bukan daftar parsial, melainkan satu peta linear yang meniru fungsi
di sekitar titik; daftar parsial dapat menyesatkan.

\begin{definisi}[Diferensiabilitas]
Misalkan $U\subseteq\mathbb{R}^{n}$ terbuka dan $F\colon U\to\mathbb{R}^{m}$. $F$
terdiferensialkan di $a\in U$ bila ada peta linear $L\colon\mathbb{R}^{n}\to\mathbb{R}^{m}$
dengan
\[
  F(a+v)=F(a)+Lv+r(v),
  \qquad
  \lim_{v\to0}\frac{\|r(v)\|}{\|v\|}=0 .
\]
Peta $L$ disebut turunan $F$ di $a$ dan ditulis $F'(a)$.\index{diferensiabilitas}
\end{definisi}

\begin{proposisi}
Peta $L$ pada definisi itu tunggal bila ada.
\end{proposisi}

\begin{proof}
Bila $L_{1}$ dan $L_{2}$ memenuhi, maka $(L_{1}-L_{2})v=r_{2}(v)-r_{1}(v)$. Ambil
$v=tw$ dengan $w$ tetap dan $t>0$, bagi dengan $t$, lalu $t\to0$: diperoleh
$(L_{1}-L_{2})w=0$. Karena berlaku untuk setiap $w$, kedua peta sama.
\end{proof}

\begin{teorema}
Bila $F$ terdiferensialkan di $a$, maka semua turunan parsialnya di $a$ ada, dan matriks
$L$ terhadap basis standar adalah
\[
  J_{F}(a)=\begin{pmatrix}
    \dfrac{\partial F_{1}}{\partial x_{1}}(a) & \cdots &
    \dfrac{\partial F_{1}}{\partial x_{n}}(a)\\
    \vdots & \ddots & \vdots\\
    \dfrac{\partial F_{m}}{\partial x_{1}}(a) & \cdots &
    \dfrac{\partial F_{m}}{\partial x_{n}}(a)
  \end{pmatrix}.
\]
\end{teorema}

\begin{proof}
Ambil $v=te_{j}$ dengan $t\neq0$. Definisi memberi
$F(a+te_{j})-F(a)=tLe_{j}+r(te_{j})$, dan setelah dibagi $t$ lalu $t\to0$, kolom ke-$j$
dari $L$ muncul sebagai limit hasil bagi, yaitu vektor turunan parsial terhadap $x_{j}$.
\end{proof}

\kamus{Jacobian}{Jacobian matrix}{Tabel semua kemiringan sebuah fungsi banyak variabel:
setiap baris satu keluaran, setiap kolom satu masukan.}

Matriks $m\times n$ itu disebut matriks Jacobian $F$ di $a$.\index{jacobian} Untuk $m=1$ ia
menjadi baris tunggal, yaitu transpos gradien. Contoh: $F(x,y)=(x^{2}y,\;x+y^{2})$ di $(1,2)$
memberi
\[
  J_{F}(1,2)=\begin{pmatrix}2xy & x^{2}\\ 1 & 2y\end{pmatrix}_{(1,2)}
  =\begin{pmatrix}4&1\\ 1&4\end{pmatrix},
\]
sehingga $F(1{,}01,\,1{,}98)\approx(2,5)+J_{F}(1,2)(0{,}01,-0{,}02)=(2{,}02,\,4{,}93)$,
sedangkan nilai eksaknya $(2{,}019798,\,4{,}9304)$; hampiran itu meleset $2\times10^{-4}$
pada komponen pertama.

Parsial yang ada tidak menjamin diferensiabilitas; dua contoh berikut menunjukkan dua tingkat
kegagalan.

\begin{contoh}[Parsial ada, tidak kontinu]
Misalkan $f(x,y)=\dfrac{xy}{x^{2}+y^{2}}$ untuk $(x,y)\neq(0,0)$ dan $f(0,0)=0$. Di titik
asal $f(h,0)=0$, sehingga kedua turunan parsialnya nol. Tetapi sepanjang garis $y=kx$ nilai
fungsinya $k/(1+k^{2})$: untuk $k=1$ nilainya $1/2$, untuk $k=0$ nilainya $0$. Limitnya tidak
ada, dan $f$ tidak kontinu di titik asal.
\end{contoh}

\begin{contoh}[Kontinu, parsial ada, tak terdiferensialkan]
Misalkan $g(x,y)=\dfrac{xy}{\sqrt{x^{2}+y^{2}}}$ untuk $(x,y)\neq(0,0)$ dan $g(0,0)=0$.
Karena $2|xy|\le x^{2}+y^{2}$, berlaku $|g(x,y)|\le\frac12\sqrt{x^{2}+y^{2}}$, jadi $g$
kontinu di titik asal, dan kedua parsialnya di sana nol. Bila $g$ terdiferensialkan, peta
linearnya harus nol, sehingga sisa $r(v)=g(v)$ mesti memenuhi $\|r(v)\|/\|v\|\to0$. Ambil
$v=t(1,1)/\sqrt{2}$ dengan $t>0$: maka $g(v)=(t^{2}/2)/t=t/2$, sehingga
$g(v)/\|v\|=1/2$ untuk setiap $t$. Nisbah itu menuju $1/2$, bukan $0$; jadi $g$ tidak
terdiferensialkan di titik asal.
\end{contoh}

Akibatnya penting: pada contoh kedua $\nabla g(0,0)=0$ tetapi turunan berarah diagonal tidak
ada, karena hasil baginya $\pm1/2$ bergantung tanda $t$. Jadi $D_{u}f=\nabla f\cdot u$
mensyaratkan diferensiabilitas.

Syarat yang mudah diperiksa dan cukup untuk diferensiabilitas adalah kekontinuan parsial.

\begin{teorema}
Bila turunan-turunan parsial $F$ ada dan kontinu di sekitar $a$, maka $F$ terdiferensialkan
di $a$.
\end{teorema}

\begin{proof}
Cukup dibuktikan untuk $m=1$ dan $n=2$; kasus lain mengikuti komponen demi komponen dengan
indeks yang bertambah. Tulis $a=(p,q)$ dan $v=(h,k)$. Maka
\[
  f(p+h,q+k)-f(p,q)
  =\bigl[f(p+h,q+k)-f(p,q+k)\bigr]+\bigl[f(p,q+k)-f(p,q)\bigr].
\]
Menurut teorema nilai rata-rata satu variabel, suku pertama sama dengan
$h\,\partial_{x}f(\xi,q+k)$ untuk suatu $\xi$ di antara $p$ dan $p+h$, dan suku kedua sama
dengan $k\,\partial_{y}f(p,\eta)$ untuk suatu $\eta$ di antara $q$ dan $q+k$. Karena itu
\[
  f(a+v)-f(a)-\nabla f(a)\cdot v
  =h\bigl[\partial_{x}f(\xi,q+k)-\partial_{x}f(a)\bigr]
  +k\bigl[\partial_{y}f(p,\eta)-\partial_{y}f(a)\bigr].
\]
Bagi dengan $\|v\|=\sqrt{h^{2}+k^{2}}$: karena $|h|/\|v\|\le1$ dan $|k|/\|v\|\le1$, dan
kedua selisih dalam kurung menuju $0$ (di sinilah kekontinuan parsial dipakai), ruas kanan
dibagi $\|v\|$ menuju $0$.
\end{proof}

Sekarang komposisi.

\begin{teorema}[Aturan rantai multivariabel]
Misalkan $H\colon\mathbb{R}^{n}\to\mathbb{R}^{k}$ terdiferensialkan di $a$ dan
$G\colon\mathbb{R}^{k}\to\mathbb{R}^{m}$ terdiferensialkan di $H(a)$. Maka $G\circ H$
terdiferensialkan di $a$ dan
\[
  (G\circ H)'(a)=G'(H(a))\,H'(a),
\]
yaitu komposisi dua peta linear; dalam matriks,
$J_{G\circ H}(a)=J_{G}(H(a))\,J_{H}(a)$.
\end{teorema}

\begin{proof}
Tulis $A=H'(a)$, $B=G'(y_{0})$ dengan $y_{0}=H(a)$, dan
\[
  H(a+v)=H(a)+Av+r(v),\qquad G(y_{0}+w)=G(y_{0})+Bw+q(w),
\]
dengan $\|r(v)\|\le\varepsilon(v)\|v\|$ dan $\|q(w)\|\le\eta(w)\|w\|$, di mana
$\varepsilon(v)\to0$ dan $\eta(w)\to0$. Untuk $w=H(a+v)-H(a)=Av+r(v)$ berlaku
$\|w\|\le(\|A\|+\varepsilon(v))\|v\|$, jadi $w\to0$ ketika $v\to0$. Maka
\[
  G\bigl(H(a+v)\bigr)-G(H(a))-BAv
  =B\bigl(H(a+v)-H(a)-Av\bigr)+q(w)
  =B\,r(v)+q(w).
\]
Bagi dengan $\|v\|$:
\[
  \frac{\|B\,r(v)+q(w)\|}{\|v\|}
  \le\|B\|\,\varepsilon(v)+\bigl(\|A\|+\varepsilon(v)\bigr)\eta(w),
\]
dan ruas kanan menuju $0$. Jadi $BA$ memenuhi definisi turunan $G\circ H$ di $a$.
\end{proof}

Bentuk matriksnya menjelaskan mengapa urutan perkalian tidak boleh dibalik: $J_{H}$
berukuran $k\times n$ dan $J_{G}$ berukuran $m\times k$, jadi hanya satu urutan yang dapat
dikalikan. Untuk fungsi skalar hasilnya $\frac{d}{dt}f(\gamma(t))=\nabla f(\gamma(t))\cdot
\gamma'(t)$.

Terakhir, turunan kedua untuk fungsi banyak variabel disimpan dalam satu matriks.

\kamus{Hessian}{Hessian matrix}{Tabel kelengkungan: apakah permukaan melengkung ke atas,
ke bawah, atau ke dua arah sekaligus.}

\begin{definisi}[Matriks Hessian]
Bila semua turunan parsial kedua $f\colon U\to\mathbb{R}$ ada di $a$, matriks
\[
  H_{f}(a)=\left(\frac{\partial^{2}f}{\partial x_{i}\,\partial x_{j}}(a)\right)_{i,j=1}^{n}
\]
disebut matriks Hessian $f$ di $a$.\index{hessian}
\end{definisi}

Untuk fungsi yang parsial keduanya kontinu, $H_{f}(a)$ simetris (teorema Clairaut, dibuktikan
di Stewart pada bagian turunan parsial). Tanpa kekontinuan, kesimpulan itu bisa gagal.

Hessian adalah turunan dari gradien. Bila parsial keduanya kontinu di sekitar titik kritis
$a$ ($\nabla f(a)=0$), maka
\[
  f(a+v)=f(a)+\tfrac12\,v^{\top}H_{f}(a)\,v+o\bigl(\|v\|^{2}\bigr),
\]
sehingga tanda bentuk kuadrat $v\mapsto v^{\top}Hv$ menentukan jenis titik kritis.

\begin{proposisi}[Uji turunan kedua]
Misalkan $f$ mempunyai turunan parsial kedua kontinu di sekitar titik kritis $a$
($\nabla f(a)=0$).
\begin{enumerate}
  \item Bila $H_{f}(a)$ definit positif, $a$ titik minimum lokal tegas.
  \item Bila $H_{f}(a)$ definit negatif, $a$ titik maksimum lokal tegas.
  \item Bila $H_{f}(a)$ mempunyai nilai eigen bertanda berbeda, $a$ titik sadel.
  \item Bila $H_{f}(a)$ semidefinit tetapi singular, uji ini tidak memberi kesimpulan.
\end{enumerate}
\end{proposisi}

Bagian pertama sampai ketiga mengikuti ekspansi kuadratik di atas: pada arah $v^{\top}Hv>0$
fungsi naik, pada arah $v^{\top}Hv<0$ fungsi turun. Bukti lengkapnya ada di Stewart pada
bagian nilai maksimum dan minimum. Bagian keempat benar-benar kosong: $x^{4}+y^{4}$ dan
$x^{4}-y^{4}$ sama-sama punya $H=0$ di titik asal, tetapi yang pertama minimum dan yang kedua
sadel. Untuk dua variabel, $AC-B^{2}>0$ dengan $A>0$ berarti definit positif, $AC-B^{2}<0$
berarti indefinit.

\section{Contoh hitung}
\tingkat{menengah}

Tiga hitungan berikut memakai angka yang sama dengan kode di bagian lab.

Pertama, beda hingga untuk $f(x)=\sin x$ di $x_{0}=1$; nilai eksaknya $\cos 1$, yang
angkanya muncul pada keluaran kode di bagian lab. Dua rumusnya:
\[
  \text{maju: }\frac{f(x_{0}+h)-f(x_{0})}{h},\qquad
  \text{pusat: }\frac{f(x_{0}+h)-f(x_{0}-h)}{2h}.
\]

\begin{center}
\begin{tabular}{rrrrr}
\hline
$h$ & beda maju & galat relatif & beda pusat & galat relatif\\
\hline
$10^{-1}$ & $0{,}497363752535$ & $7{,}947\times10^{-2}$ & $0{,}539402252170$ & $1{,}666\times10^{-3}$\\
$10^{-2}$ & $0{,}536085981012$ & $7{,}804\times10^{-3}$ & $0{,}540293300875$ & $1{,}667\times10^{-5}$\\
$10^{-4}$ & $0{,}540260231419$ & $7{,}787\times10^{-5}$ & $0{,}540302304968$ & $1{,}667\times10^{-9}$\\
$10^{-6}$ & $0{,}540301885121$ & $7{,}787\times10^{-7}$ & $0{,}540302305896$ & $5{,}130\times10^{-11}$\\
$10^{-8}$ & $0{,}540302302898$ & $5{,}497\times10^{-9}$ & $0{,}540302308449$ & $4{,}777\times10^{-9}$\\
$10^{-10}$ & $0{,}540302247387$ & $1{,}082\times10^{-7}$ & $0{,}540302247387$ & $1{,}082\times10^{-7}$\\
\hline
\end{tabular}
\end{center}

Pola galatnya bukan kebetulan. Untuk beda maju, ekspansi Taylor memberi
\[
  \frac{f(x_{0}+h)-f(x_{0})}{h}=f'(x_{0})+\frac{h}{2}f''(x_{0})+\mathcal{O}(h^{2}),
\]
jadi galat relatifnya berorde $h$ dan turun sepuluh kali setiap $h$ dibagi sepuluh:
$7{,}947\times10^{-2}$, $7{,}804\times10^{-3}$, $7{,}787\times10^{-5}$. Angkanya dapat
dicocokkan: $\frac{h}{2}f''(1)=-\frac{1}{2}\cdot10^{-2}\sin 1=-4{,}207\times10^{-3}$, dan
$4{,}207\times10^{-3}/0{,}5403=7{,}787\times10^{-3}$. Beda pusat memberi
\[
  \frac{f(x_{0}+h)-f(x_{0}-h)}{2h}=f'(x_{0})+\frac{h^{2}}{6}f'''(x_{0})+\mathcal{O}(h^{4}),
\]
sehingga galatnya berorde $h^{2}$ dan turun seratus kali: $1{,}666\times10^{-3}$,
$1{,}667\times10^{-5}$, $1{,}667\times10^{-9}$. Di sini $\frac{10^{-4}}{6}\cos 1
=9{,}005\times10^{-6}$, dan $9{,}005\times10^{-6}/0{,}5403$ kembali memberi $1{,}667\times10^{-5}$.

Pada $h=10^{-8}$ dan $10^{-10}$ beda pusat memburuk lagi: galat pemotongan sudah lebih kecil
daripada galat pembulatan yang tumbuh seperti $\varepsilon/h$, dengan
$\varepsilon\approx2{,}2\times10^{-16}$. Yang terbaik di tabel, $h=10^{-6}$ dengan galat
$5\times10^{-11}$, ordenya sesuai $\varepsilon^{1/3}\approx6\times10^{-6}$.

Kedua, gradien $f(x,y)=x^{2}y^{3}$ di $(2,-1)$ dengan nilai analitik $(-4,12)$. Beda pusat
memberi galat $3{,}333\times10^{-5}$ pada komponen $y$ untuk $h=10^{-2}$, turun ke
$3{,}333\times10^{-9}$ untuk $h=10^{-4}$: orde $h^{2}$. Komponen $x$ tepat sampai galat mesin,
sebab $f$ kuadratik dalam $x$.

Ketiga, aturan rantai pada komposisi berdimensi berbeda. Misalkan
\[
  H(x,y)=\bigl(x^{2}+y,\;xy,\;y^{3}\bigr),\qquad
  G(u,v,w)=\bigl(uv+w,\;u^{2}-vw\bigr),
\]
dan tinjau $F=G\circ H$ di $a=(1,2)$. Di sini $H(a)=(3,2,8)$, dan
\[
  J_{H}(1,2)=\begin{pmatrix}2&1\\ 2&1\\ 0&12\end{pmatrix}_{3\times2},
  \qquad
  J_{G}(3,2,8)=\begin{pmatrix}2&3&1\\ 6&-8&-2\end{pmatrix}_{2\times3}.
\]
Hasil kalinya berukuran $2\times2$:
\[
  J_{G}(3,2,8)\,J_{H}(1,2)=
  \begin{pmatrix}2&3&1\\ 6&-8&-2\end{pmatrix}
  \begin{pmatrix}2&1\\ 2&1\\ 0&12\end{pmatrix}
  =\begin{pmatrix}10&17\\ -4&-26\end{pmatrix}.
\]
\kertas{Untuk komposisi di atas, hitung $J_{F}$ di $(2,1)$ dengan dua jalan: hasil kali
$J_{G}(H(2,1))J_{H}(2,1)$ dan turunan langsung bentuk eksplisit. Lalu periksa arah $v=(1,3)$
dengan membandingkan $J_{F}(2,1)v$ terhadap $\bigl(F(a+sv)-F(a)\bigr)/s$ untuk $s=10^{-2}$
dan $s=10^{-4}$. Mengapa galatnya turun hanya sebesar faktor $10$, bukan $100$?}

\section{Lab komputasi}
\tingkat{lanjut}

Kode berikut dijalankan apa adanya dengan numpy 2.5.2 pada Python 3: bagian pertama mengulang
tiga hitungan di atas, lalu satu percobaan yang tidak dapat dikerjakan dengan tangan, beda
hingga di bawah derau pengukuran.

\begin{lab}
\begin{verbatim}
import numpy as np

# 1. beda hingga untuk f(x) = sin x di x0 = 1
x0 = 1.0
eksak = float(np.cos(x0))
print("cos(1) =", eksak)
for h in [1e-1, 1e-2, 1e-4, 1e-6, 1e-8, 1e-10]:
    maju = (np.sin(x0 + h) - np.sin(x0)) / h
    pusat = (np.sin(x0 + h) - np.sin(x0 - h)) / (2*h)
    em = abs(maju - eksak) / abs(eksak)
    ep = abs(pusat - eksak) / abs(eksak)
    print(f"{h:8.0e} fwd={maju:.12f} {em:.3e} "
          f"cen={pusat:.12f} {ep:.3e}")
\end{verbatim}
Keluaran sebenarnya:
\begin{verbatim}
cos(1) = 0.5403023058681398
   1e-01 fwd=0.497363752535 7.947e-02 cen=0.539402252170 1.666e-03
   1e-02 fwd=0.536085981012 7.804e-03 cen=0.540293300875 1.667e-05
   1e-04 fwd=0.540260231419 7.787e-05 cen=0.540302304968 1.667e-09
   1e-06 fwd=0.540301885121 7.787e-07 cen=0.540302305896 5.130e-11
   1e-08 fwd=0.540302302898 5.497e-09 cen=0.540302308449 4.777e-09
   1e-10 fwd=0.540302247387 1.082e-07 cen=0.540302247387 1.082e-07
\end{verbatim}
\end{lab}

Baris terakhir layak dibaca dua kali: pada $h=10^{-10}$ kedua rumus memberi angka identik,
karena keduanya menghitung galat pembulatan yang sama, bukan turunan.

\begin{lab}
\begin{verbatim}
def f(p):
    return p[0]**2 * p[1]**3

def grad_fd(fun, a, h):
    g = np.zeros_like(a)
    for j in range(a.size):
        e = np.zeros_like(a)
        e[j] = h
        g[j] = (fun(a + e) - fun(a - e)) / (2*h)
    return g

a = np.array([2.0, -1.0])
g_eks = np.array([2*a[0]*a[1]**3, 3*a[0]**2*a[1]**2])
print("gradien =", g_eks)
for h in [1e-2, 1e-4, 1e-6]:
    g = grad_fd(f, a, h)
    print(f"{h:8.0e} {g[0]:.12f} {g[1]:.12f} "
          f"{abs(g[0]+4)/4:.3e} {abs(g[1]-12)/12:.3e}")
\end{verbatim}
Keluaran sebenarnya:
\begin{verbatim}
gradien = [-4. 12.]
   1e-02 -4.000000000000 12.000400000000 1.577e-14 3.333e-05
   1e-04 -4.000000000004 12.000000039998 1.000e-12 3.333e-09
   1e-06 -4.000000000115 11.999999999679 2.876e-11 2.676e-11
\end{verbatim}
\end{lab}

\begin{lab}
\begin{verbatim}
def hh(p):
    x, y = p
    return np.array([x*x + y, x*y, y**3])

def gg(q):
    u, v, w = q
    return np.array([u*v + w, u*u - v*w])

def F(p):
    return gg(hh(p))

def jac(fun, p, step):
    n = p.size
    J = np.zeros((fun(p).size, n))
    for j in range(n):
        e = np.zeros(n)
        e[j] = step
        J[:, j] = (fun(p + e) - fun(p - e)) / (2*step)
    return J

a = np.array([1.0, 2.0])
q = hh(a)
Jh = np.array([[2*a[0], 1.0], [a[1], a[0]], [0.0, 3*a[1]**2]])
Jg = np.array([[q[1], q[0], 1.0], [2*q[0], -q[2], -q[1]]])
J_rantai = Jg @ Jh
J_num = jac(F, a, 1e-6)
print("h(1,2) =", q)
print("Jg @ Jh =\n", J_rantai)
print("selisih maks =", f"{np.max(np.abs(J_num - J_rantai)):.3e}")

v = np.array([0.7, -1.3])
d_rantai = J_rantai @ v
d_num = (F(a + 1e-6*v) - F(a - 1e-6*v)) / 2e-6
print("J v =", d_rantai)
print("beda hingga =", d_num)
print("galat relatif =",
      f"{np.max(np.abs(d_num - d_rantai)/np.abs(d_rantai)):.3e}")
\end{verbatim}
Keluaran sebenarnya:
\begin{verbatim}
h(1,2) = [3. 2. 8.]
Jg @ Jh =
 [[ 10.  17.]
 [ -4. -26.]]
selisih maks = 1.858e-09
J v = [-15.1  31. ]
beda hingga = [-15.1  31. ]
galat relatif = 9.860e-11
\end{verbatim}
\end{lab}

Selisih $1{,}858\times10^{-9}$ antara Jacobian aturan rantai dan Jacobian beda hingga adalah
galat pemotongan beda pusat, bukan kesalahan aljabar; pemeriksaan arah memberi
$9{,}86\times10^{-11}$, sebab keduanya memakai titik yang sama.

\begin{lab}
\begin{verbatim}
rng = np.random.default_rng(7)
sigma = 1e-3
for h in [1e-1, 1e-2, 1e-3, 1e-4, 1e-5]:
    galat = []
    for _ in range(2000):
        s1 = np.sin(x0 + h) + sigma*rng.standard_normal()
        s0 = np.sin(x0) + sigma*rng.standard_normal()
        galat.append(abs((s1 - s0)/h - eksak))
    print(f"{h:8.0e} galat rerata = {np.mean(galat):.3e}")
\end{verbatim}
Keluaran sebenarnya:
\begin{verbatim}
   1e-01 galat rerata = 4.203e-02
   1e-02 galat rerata = 1.112e-01
   1e-03 galat rerata = 1.160e+00
   1e-04 galat rerata = 1.117e+01
   1e-05 galat rerata = 1.088e+02
\end{verbatim}
\end{lab}

Percobaan terakhir adalah inti bab ini dalam angka. Dengan derau
$\sigma=10^{-3}$, hasil bagi membawa dua galat: bias pemotongan $\approx0{,}42h$ dan derau
yang menguat seperti $\sqrt{2}\sigma/h$. Pada $h=10^{-1}$ bias mendominasi ($0{,}042$), pada
$h=10^{-4}$ derau mendominasi ($11{,}2$), dan keduanya seimbang ketika
$0{,}42h=\sqrt{2}\cdot10^{-3}/h$, yaitu $h\approx0{,}058$: dengan alat seburuk itu, langkah
terbaik sekitar enam perseratus, bukan seperseratus juta.

Inilah jawaban jujur untuk pertanyaan pembuka. Turunan adalah limit, tetapi aproksimasi
terbaiknya bukan limit yang sedekat mungkin. Yang terbaik adalah kompromi antara dua galat
yang bergerak berlawanan arah, dan kompromi itu dapat dihitung sebelum satu angka pun
diukur.

\begin{untukdosen}
Urutan yang berhasil dipakai: definisi hasil bagi, lalu bentuk aproksimasi linear, baru
aturan hitung. Bentuk itu terasa berlebihan untuk satu variabel, tetapi hanya ia yang bertahan
sampai Jacobian; bila diperkenalkan belakangan, mahasiswa menghafal aturan rantai sebagai
perkalian matriks tanpa tahu mengapa urutannya begitu.

Dua miskonsepsi khas. Pertama, ``turunan parsial ada berarti fungsi mulus''; contoh
$xy/\sqrt{x^{2}+y^{2}}$ di papan mematahkannya asalkan dihitung bersama: parsial nol,
kontinu, tetapi arah diagonal memberi $1/2$. Kedua, arah perkalian pada aturan rantai:
mahasiswa menulis $J_{H}J_{G}$ karena urutan penulisan fungsi terbalik dari urutan
pengerjaan; cara memeriksanya adalah ukuran matriksnya.

Aktivitas papan tulis: minta kelas menghitung kedua parsial di titik asal untuk kedua contoh
penyangkal, lalu memutuskan apakah fungsi itu punya bidang singgung. Tanda kesiapan lanjut:
mahasiswa menjelaskan mengapa $D_{u}f=\nabla f\cdot u$ gagal pada contoh kedua, dan menyebut
kata ``sisa'', bukan ``rumus''.
\end{untukdosen}

\section{Ringkasan}
\tingkat{dasar}

\begin{ringkasan}
\begin{itemize}
  \item Turunan $f'(c)$ adalah limit hasil bagi selisih, sekaligus koefisien tunggal
        aproksimasi linear $f(c+h)=f(c)+Lh+r(h)$ dengan $r(h)/h\to0$. Diferensiabilitas
        menjamin kekontinuan, bukan sebaliknya.
  \item Aturan rantai $(f\circ g)'=f'(g(c))g'(c)$ dibuktikan lewat aproksimasi linear, dan
        buktinya bertahan ke dimensi banyak: $J_{G\circ H}(a)=J_{G}(H(a))J_{H}(a)$.
  \item Teorema nilai rata-rata menghubungkan tanda turunan dengan perilaku fungsi: $f'\ge0$
        berarti tidak turun, $|f'|\le M$ berarti Lipschitz; ia gagal untuk fungsi vektor.
  \item Taylor dengan sisa Lagrange memakai $\xi$ yang tidak diketahui, sehingga yang dipakai
        hanya batas atas $|f^{(n)}|$.
  \item Untuk banyak variabel, turunan adalah peta linear dengan matriks Jacobian. Parsial
        boleh ada tanpa kekontinuan, dan boleh ada bersama kekontinuan tanpa
        diferensiabilitas.
  \item Gradien menunjuk arah naik tercepat dengan laju $\|\nabla f\|$, tegak lurus kontur;
        Hessian yang definit menentukan jenis titik kritis.
  \item Beda hingga berorde $h$ (maju) atau $h^{2}$ (pusat), dan galat pembulatan tumbuh
        seperti $\varepsilon/h$. Langkah terkecil bukan langkah terbaik.
\end{itemize}
\end{ringkasan}

\section{Latihan}
\tingkat{menengah}

\begin{latihan}
\begin{enumerate}
  \item Hitung $f'(2)$ dari definisi untuk $f(x)=x^{3}$ dan $f(x)=\sqrt{x}$, sampai faktor
        $h$ hilang dari penyebut.

  \item Buktikan bahwa bila $f'(c)>0$, ada $\delta>0$ sehingga $f(c+h)>f(c)$ untuk setiap
        $h\in(0,\delta)$, dan jelaskan mengapa ini tidak berarti $f$ naik pada suatu interval.

  \item Buktikan aturan hasil bagi dengan menurunkan $1/g$ lebih dulu, lalu memakai aturan
        hasil kali. Sebutkan di langkah mana kekontinuan $g$ dipakai.

  \item Nilai klaim ini: \emph{bila $f$ mempunyai semua turunan parsial di suatu titik, maka
        $f$ kontinu di titik itu.} Nyatakan benar, salah, atau belum cukup bukti, dan
        sertakan contoh penyangkal bila perlu.

  \item Untuk $f(x,y)=x^{3}y-xy^{3}$, hitung gradien di $(2,1)$, tentukan arah naik tercepat,
        dan hitung laju maksimumnya. Periksa dengan beda pusat $h=10^{-4}$ dan bandingkan
        galat relatifnya dengan orde $h^{2}$.

  \item Buktikan bahwa bila $f'=0$ di mana-mana pada suatu interval, maka $f$ konstan.
        Jelaskan mengapa bukti itu gagal pada himpunan yang tidak terhubung, dan beri
        contohnya.

  \item Misalkan $H(x,y)=(e^{x}y,\;x+y^{2})$ dan $G(u,v)=(u^{2}v,\;u-v)$. Hitung
        $J_{G\circ H}$ di $(0,1)$ dengan dua jalan: hasil kali matriks dan turunan langsung.
        (Petunjuk: $H(0,1)=(1,1)$.)

  \item Tentukan semua titik kritis $f(x,y)=x^{3}+y^{3}-3xy$, lalu klasifikasikan dengan
        Hessian; jelaskan mengapa tanda determinan tidak cukup bila determinannya nol.

  \item Taksir $\sqrt{1{,}04}$ dengan polinomial Taylor orde dua di $a=1$, lalu batas
        galatnya dengan sisa Lagrange. Bandingkan batas itu dengan galat sebenarnya: apakah batasnya ketat? (Petunjuk:
        $f''(x)=-\frac14x^{-3/2}$ dan $f'''$ positif pada selang itu.)
\end{enumerate}
\end{latihan}

\begin{untukpeneliti}
Yang terbukti di bab ini adalah pernyataan klasik satu dan banyak variabel: aproksimasi
linear dan ketunggalan petanya, aturan rantai, teorema nilai rata-rata dan akibatnya, Taylor
dengan sisa Lagrange, kesetaraan diferensiabilitas dengan Jacobian, dan syarat cukup
kekontinuan parsial. Semua itu teorema, bukan konvensi, berlaku di ruang berdimensi hingga
atas $\mathbb{R}$.

Yang sering dipopulerkan adalah klaim ``gradien adalah turunan''. Tanpa diferensiabilitas
klaim itu salah, dan $xy/\sqrt{x^{2}+y^{2}}$ adalah penyangkal yang tetap kontinu. Klaim
``Hessian menentukan jenis titik kritis'' hanya berlaku bila Hessian nonsingular; kasus
singular memerlukan turunan orde lebih tinggi atau analisis langsung.

Sumber primer. Rudin, \emph{Principles of Mathematical Analysis}, bab
\emph{Differentiation} dan \emph{Functions of Several Variables}, memuat seluruh pernyataan
di atas. Spivak, \emph{Calculus on Manifolds}, bab \emph{Differentiation}, mengambil jalan
yang sama dan lebih singkat. Stewart, \emph{Calculus: Early Transcendentals} edisi ke-9,
bagian 14.3 sampai 14.7, memberi padanan komputasional, termasuk teorema Clairaut.

Batas keberlakuan: semua argumen menganggap himpunan definisi terbuka dan turunan yang
dipakai ada. Pada fungsi yang hanya terdiferensialkan hampir di mana-mana, kesimpulan
monotonisitas dan Taylor runtuh, dan itu wilayah teori ukur yang tidak dibahas buku ini.
Untuk fungsi vektor, teorema nilai rata-rata tidak berlaku; penggantinya ketaksamaan nilai
rata-rata yang menuntut norma operator. Orde beda hingga terbaik bergantung pada skala
fungsi, derau, dan aritmetika; bab ini hanya menunjukkan ordenya secara empiris.
\end{untukpeneliti}

\peta{Bab ini mengikuti Rudin, \emph{Principles of Mathematical Analysis}, bab
\emph{Differentiation} dan \emph{Functions of Several Variables}; versi komputasionalnya
Stewart, \emph{Calculus: Early Transcendentals} edisi ke-9, bagian 14.3 sampai 14.7; padanan
abstraknya Spivak, \emph{Calculus on Manifolds}, bab \emph{Differentiation}. Kekontinuan dan
teorema nilai ekstrem dipakai dari Bab 11, ketaksamaan Cauchy--Schwarz dari Bab 10, norma
Jacobian kembali di Bab 23 dan Bab 28.}

\chapter{Integral Riemann dan Teorema Dasar Kalkulus}

Sebuah program menghitung luas di bawah $y=x^2$ pada $[0,1]$ dengan 1000 potong sama
lebar. Bila tinggi persegi panjang diambil di ujung kiri, hasilnya $0.3328335$; bila di
ujung kanan, $0.3338335$. Selisihnya $0.001$, dan tidak ada alasan memilih salah satu.
Nilai yang benar adalah $1/3$, yang tidak muncul pada kedua hitungan itu. Menambah
potongan menjadi sejuta tidak menolong: kedua jawaban tetap mengapit, tidak pernah
menangkap.

Yang diperlukan bukan partisi lebih halus, melainkan cara membaca dua hampiran itu
sekaligus. Kalau selisihnya dapat dibuat lebih kecil dari bilangan positif mana pun,
maka ada satu bilangan yang terkurung di antaranya, dan itulah integralnya. Seluruh
bab ini berdiri di atas gagasan sederhana itu.

Bab 12 mengajari kita membongkar perubahan menjadi laju: dari posisi ke kecepatan,
dari fungsi ke turunan. Yang belum ada adalah jalan pulang. Kita tahu kecepatan pada
detik ke-$t$ adalah $v(t)$, tetapi belum tahu berapa jarak yang ditempuh selama satu
menit. Intuisinya jelas: jarak adalah jumlah kecepatan dikali potongan waktu kecil.
Yang belum jelas adalah cara menjumlahkan tak berhingga banyak suku itu dengan benar.
Bab ini menjawabnya, dan jawabannya mengejutkan: jalan pulang itu adalah antiturunan.

\begin{untukanak}
Bayangkan kita ingin tahu luas sebuah danau yang bentuknya tidak beraturan. Tidak ada
rumus untuk bentuk seperti itu. Yang kita punya hanya kertas berpetak bening. Kita
letakkan kertas itu di atas gambar danau, lalu hitung berapa petak yang seluruhnya
berada di dalam air. Hasilnya terlalu kecil. Kita hitung lagi berapa petak yang
tersentuh air, walau hanya sedikit. Hasilnya terlalu besar. Luas sebenarnya ada di
antara dua angka itu.

Sekarang potong petaknya menjadi sepuluh kali lebih kecil. Angka bawah naik, angka atas
turun, dan keduanya makin rapat. Kalau petaknya sangat kecil, kedua angka hampir sama,
dan kita bisa menyebut luas danau sampai ketelitian berapa pun.

Ada satu syarat yang mudah terlewat. Kalau danau punya pulau karang sangat tipis yang
tersebar di mana-mana, petak yang menyentuhnya bisa membuat taksiran atas tetap jauh
dari taksiran bawah. Matematikawan menyebut syarat itu keterintegralan. Di bab ini kita
akan tahu mengapa dua taksiran yang mengapit bisa dipercaya, dan kapan taksiran itu
gagal.
\end{untukanak}

\section{Partisi, jumlah atas, dan jumlah bawah}
\tingkat{dasar}

Kita ingin memberi arti pada luas di bawah grafik tanpa mengandaikan luas itu sudah
punya arti: tutup daerahnya dengan persegi panjang dari dalam dan dari luar, lalu
persempit selisihnya.

\begin{definisi}[Partisi dan jumlah atas-bawah]\index{partisi}
Misalkan $[a,b]$ selang tertutup. Sebuah \emph{partisi} $\mathcal{P}$ dari $[a,b]$
adalah himpunan berhingga titik
\[
a=x_0<x_1<\cdots<x_n=b .
\]
Lebar potongan ke-$i$ adalah $\Delta x_i=x_i-x_{i-1}$, dan \emph{mesh} partisi itu
adalah lebar potongan terlebar, $\|\mathcal{P}\|=\max_i\Delta x_i$.

Misalkan $f$ terbatas pada $[a,b]$. Pada potongan ke-$i$ tulis
\[
M_i=\sup\{f(x):x_{i-1}\le x\le x_i\},\qquad
m_i=\inf\{f(x):x_{i-1}\le x\le x_i\}.
\]
\emph{Jumlah atas} dan \emph{jumlah bawah} partisi $\mathcal{P}$ adalah
\begin{equation}
U(\mathcal{P},f)=\sum_{i=1}^{n}M_i\,\Delta x_i,
\qquad
L(\mathcal{P},f)=\sum_{i=1}^{n}m_i\,\Delta x_i .
\end{equation}
\end{definisi}

Kita memakai supremum dan infimum, bukan nilai terbesar dan terkecil, karena $f$ belum
tentu mencapai ekstremnya pada potongan: fungsi yang bernilai $x$ untuk $x<1$ dan $0$
di $x=1$ terbatas pada $[0,1]$, tetapi nilai terbesarnya pada potongan terakhir tidak
ada, sedangkan supremumnya ada, yaitu $1$.

Jumlah atas selalu paling sedikit sama dengan jumlah bawah, sebab $M_i\ge m_i$ untuk
setiap $i$. Untuk $f(x)=x^2$ pada $[0,1]$ dengan empat potongan sama lebar, jumlah
bawahnya
\[
L=\frac{1}{4}\Bigl(0+\frac{1}{16}+\frac{4}{16}+\frac{9}{16}\Bigr)=\frac{7}{32},
\]
dan jumlah atasnya $\frac{15}{32}$. Keduanya mengapit nilai yang dicari, dengan selisih
$\frac14$ yang masih besar. Bila $f$ bernilai negatif, jumlah ini menghitung luas
bertanda, dan tanda minusnya bagian dari definisi, bukan kekeliruan.

Gambar~\ref{fig:rect-halus} memperlihatkan dua nilai $n$ sekaligus, dengan ujung kanan
sebagai tinggi. Pada $n=4$ tangganya kasar dan $U_4=0.46875$; pada $n=16$ geriginya
halus dan $U_{16}=0.3652343750$. Jarak tangga ke garis $1/3$ menyusut dari $0.1354$
menjadi $0.0319$.

\begin{figure}[htbp]
\centering
\begin{tikzpicture}
\begin{axis}[width=0.47\linewidth,height=4.6cm,
  axis lines=left,xmin=0,xmax=1.02,ymin=0,ymax=1.12,
  xtick={0,0.5,1},ytick={0,0.5,1},
  tick label style={font=\scriptsize},label style={font=\scriptsize},
  xlabel={$x$},
  title={\scriptsize $n=4$: $U_4=0.46875$}]
\pgfplotsinvokeforeach{1,2,3,4}{
  \draw[fill=layerteal!18,draw=layerteal!65,line width=0.4pt]
    (axis cs:{(#1-1)/4},0) rectangle (axis cs:{#1/4},{#1*#1/16});
}
\draw[densely dashed,layergrey,line width=0.6pt]
  (axis cs:0,0.33333) -- (axis cs:1,0.33333);
\addplot[thick,accent,domain=0:1,samples=80] {x^2};
\node[font=\scriptsize,text=layergrey,anchor=west] at (axis cs:0.06,0.40) {$1/3$};
\end{axis}
\end{tikzpicture}\hfill
\begin{tikzpicture}
\begin{axis}[width=0.47\linewidth,height=4.6cm,
  axis lines=left,xmin=0,xmax=1.02,ymin=0,ymax=1.12,
  xtick={0,0.5,1},ytick={0,0.5,1},
  tick label style={font=\scriptsize},label style={font=\scriptsize},
  xlabel={$x$},
  title={\scriptsize $n=16$: $U_{16}=0.3652343750$}]
\pgfplotsinvokeforeach{1,...,16}{
  \draw[fill=layerteal!18,draw=layerteal!65,line width=0.3pt]
    (axis cs:{(#1-1)/16},0) rectangle (axis cs:{#1/16},{#1*#1/256});
}
\draw[densely dashed,layergrey,line width=0.6pt]
  (axis cs:0,0.33333) -- (axis cs:1,0.33333);
\addplot[thick,accent,domain=0:1,samples=80] {x^2};
\node[font=\scriptsize,text=layergrey,anchor=west] at (axis cs:0.06,0.40) {$1/3$};
\end{axis}
\end{tikzpicture}
\caption{Penghampiran $f(x)=x^{2}$ pada $[0,1]$ dengan persegi panjang berujung kanan.
Pada $n=4$ jumlah atas $U_4=0.46875$, pada $n=16$ jumlah atas $U_{16}=0.3652343750$,
dan keduanya masih di atas nilai sebenarnya $1/3=0.333\ldots$ (garis putus-putus).
Selisih $U_n-L_n$ tepat $1/n$, jadi naik dari $n=4$ ke $n=16$ memotong jarak tangga ke
garis itu menjadi seperempatnya.}
\label{fig:rect-halus}
\end{figure}


Ada dua kenyataan yang akan dipakai berulang kali. Yang pertama: memperhalus partisi
tidak memperburuk taksiran.

\begin{lemma}[Penghalusan]\index{penghalusan}
Jika $\mathcal{P}^{*}$ adalah penghalusan dari $\mathcal{P}$, artinya
$\mathcal{P}\subseteq\mathcal{P}^{*}$, maka
\[
L(\mathcal{P},f)\le L(\mathcal{P}^{*},f)\le U(\mathcal{P}^{*},f)\le U(\mathcal{P},f).
\]
\end{lemma}

\begin{proof}
Cukup dibuktikan untuk penghalusan yang menambah satu titik $c$ di dalam potongan
$[x_{i-1},x_i]$, sebab penghalusan umum diperoleh dengan mengulang langkah itu.
Misalkan $m'$ dan $m''$ infimum $f$ pada $[x_{i-1},c]$ dan pada $[c,x_i]$. Kedua
selang itu bagian dari $[x_{i-1},x_i]$, sehingga $m'\ge m_i$ dan $m''\ge m_i$.
Akibatnya
\[
m'(c-x_{i-1})+m''(x_i-c)\ \ge\ m_i(c-x_{i-1})+m_i(x_i-c)=m_i\,\Delta x_i ,
\]
jadi jumlah bawah tidak turun. Argumen yang sama dengan supremum memberi $M'\le M_i$
dan $M''\le M_i$, sehingga jumlah atas tidak naik.
\end{proof}

Kenyataan kedua lebih kuat: dua partisi berbeda tetap dapat dibandingkan, asal yang
dibandingkan adalah jumlah bawah dengan jumlah atas.

\begin{lemma}
Untuk sebarang dua partisi $\mathcal{P}_1$ dan $\mathcal{P}_2$ dari $[a,b]$ berlaku
\[
L(\mathcal{P}_1,f)\le U(\mathcal{P}_2,f).
\]
\end{lemma}

\begin{proof}
Ambil $\mathcal{P}$ penghalusan bersama, yaitu gabungan titik kedua partisi. Menurut
lemma sebelumnya, $L(\mathcal{P}_1,f)\le L(\mathcal{P},f)$ dan
$U(\mathcal{P},f)\le U(\mathcal{P}_2,f)$. Karena $L(\mathcal{P},f)\le U(\mathcal{P},f)$
selalu benar, rantai ketaksamaan itu memberi hasil yang diminta.
\end{proof}

Lemma itu memastikan setiap jumlah bawah merupakan batas bawah bagi seluruh jumlah
atas. Karena itu himpunan semua jumlah bawah terbatas di atas, dan semua jumlah atas
terbatas di bawah; kedua bilangan berikut selalu ada.

\begin{definisi}[Terintegralkan Riemann]\index{terintegralkan riemann}
Untuk $f$ terbatas pada $[a,b]$, \emph{integral bawah} dan \emph{integral atas}
didefinisikan sebagai
\[
\underline{I}(f)=\sup_{\mathcal{P}}L(\mathcal{P},f),
\qquad
\overline{I}(f)=\inf_{\mathcal{P}}U(\mathcal{P},f),
\]
dengan supremum dan infimum diambil atas seluruh partisi $[a,b]$. Karena
$\underline{I}(f)\le\overline{I}(f)$, fungsi $f$ disebut \emph{terintegralkan Riemann}
pada $[a,b]$ bila kedua bilangan itu sama. Nilai bersamanya ditulis
$\int_a^b f$ atau $\int_a^b f(x)\,dx$.
\end{definisi}

Perhatikan apa yang tidak ada pada definisi itu: tidak ada titik sampel dan tidak ada
limit atas $n$. Ukuran keberhasilan sebuah hampiran adalah selisih kedua kurungannya,
bukan kedekatannya pada tebakan.

\kamus{partisi}{partition}{Cara memotong selang menjadi beberapa potongan, seperti
membagi sebatang cokelat dengan pisau, bukan mengukur berapa panjang setiap potongan
dengan tepat.}

\kamus{jumlah atas}{upper sum}{Taksiran yang sengaja dibesar-besarkan, seperti
menimbang barang bersama kotaknya.}

\begin{bacakeras}
Dua suara. Sebut yang pertama \emph{Pemotong}, yang kedua \emph{Penjumlah}.

Pemotong: Aku punya sebatang pita. Aku potong jadi empat.
Penjumlah: Aku jumlah keempat potong itu. Hasilnya terlalu kecil.
Pemotong: Aku potong lagi. Jadi delapan.
Penjumlah: Hasilnya naik sedikit. Masih terlalu kecil.
Pemotong: Enam belas.
Penjumlah: Naik lagi. Makin dekat, belum sampai.
Pemotong: Seratus.
Penjumlah: Dekat sekali.
Pemotong: Seribu.
Penjumlah: Lebih dekat lagi, tetap belum sampai.

Pemotong: Kapan sampai?
Penjumlah: Tidak pernah, kalau kau hanya menambah potongan. Berhenti mengejar angka,
mulai mengejar kurungan.
Pemotong: Kurungan apa?
Penjumlah: Potongan yang kurus memberi batas bawah. Potongan yang gemuk memberi batas
atas. Kalau dua batas itu bertemu di satu titik, titik itulah jawabannya.
Pemotong: Jadi potonganku tidak pernah cukup?
Penjumlah: Tidak pernah. Potongan hanya membuat dua batas itu berdekatan. Yang menjawab
bukan potongan terakhir, melainkan pertemuan dua batas.

Sekarang ulangi satu baris ini: potong tipis, kurung rapat, batas bertemu.
\end{bacakeras}

\section{Kriteria keterintegralan}
\tingkat{menengah}

Definisi tadi menyebut supremum dan infimum atas seluruh partisi, dan partisi itu tak
berhingga banyaknya. Kriteria berikut menggantinya dengan satu pertidaksamaan yang
dapat diperiksa pada satu partisi saja.

\begin{teorema}[Kriteria keterintegralan]\index{kriteria keterintegralan}
Fungsi terbatas $f$ terintegralkan Riemann pada $[a,b]$ jika dan hanya jika untuk
setiap $\varepsilon>0$ terdapat partisi $\mathcal{P}$ dengan
\[
U(\mathcal{P},f)-L(\mathcal{P},f)<\varepsilon .
\]
\end{teorema}

\begin{proof}
Untuk setiap partisi berlaku
$L(\mathcal{P},f)\le\underline{I}(f)\le\overline{I}(f)\le U(\mathcal{P},f)$; ketaksamaan
pertama dari definisi supremum, yang terakhir dari definisi infimum.

$(\Rightarrow)$ Andaikan $f$ terintegralkan dengan nilai $I$, dan ambil
$\varepsilon>0$. Menurut definisi supremum, ada $\mathcal{P}_1$ dengan
$L(\mathcal{P}_1,f)>I-\varepsilon/2$. Menurut definisi infimum, ada $\mathcal{P}_2$
dengan $U(\mathcal{P}_2,f)<I+\varepsilon/2$. Ambil $\mathcal{P}$ penghalusan bersama
keduanya. Lemma penghalusan memberi
\[
U(\mathcal{P},f)-L(\mathcal{P},f)
<\Bigl(I+\tfrac{\varepsilon}{2}\Bigr)-\Bigl(I-\tfrac{\varepsilon}{2}\Bigr)
=\varepsilon .
\]

$(\Leftarrow)$ Andaikan syarat itu dipenuhi. Untuk setiap $\varepsilon>0$ ada
$\mathcal{P}$ dengan $U(\mathcal{P},f)-L(\mathcal{P},f)<\varepsilon$, sehingga
\[
0\le\overline{I}(f)-\underline{I}(f)\le U(\mathcal{P},f)-L(\mathcal{P},f)<\varepsilon .
\]
Bilangan $\overline{I}(f)-\underline{I}(f)$ tidak negatif dan lebih kecil dari setiap
$\varepsilon>0$, jadi bilangan itu nol. Kedua integral sama, dan $f$ terintegralkan.
\end{proof}

Kriteria itu mengubah pertanyaan tentang tak berhingga banyak partisi menjadi satu
pertidaksamaan. Akibat langsungnya: bila selisihnya kecil pada satu partisi, selisih
itu tetap kecil pada setiap penghalusannya, sebab lemma penghalusan hanya dapat
mempersempitnya.

Sebelum mencari kelas fungsi yang lolos kriteria, kita catat sifat integral yang akan
dipakai terus. Semuanya turunan dari definisi jumlah atas dan bawah, bukan asumsi
baru.

\begin{proposisi}[Sifat dasar integral]\index{integral, sifat dasar}
Misalkan $f$ dan $g$ terintegralkan Riemann pada $[a,b]$. Maka:
(a) $f+g$ terintegralkan dan $\int_a^b(f+g)=\int_a^b f+\int_a^b g$;
(b) jika $f\le g$ pada $[a,b]$ maka $\int_a^b f\le\int_a^b g$;
(c) jika $a<c<b$ maka $f$ terintegralkan pada $[a,c]$ dan pada $[c,b]$, dan
$\int_a^b f=\int_a^c f+\int_c^b f$;
(d) $|f|$ terintegralkan dan $\bigl|\int_a^b f\bigr|\le\int_a^b|f|$;
(e) untuk setiap konstanta $k$, $\int_a^b k\,dx=k(b-a)$.
\end{proposisi}

\begin{proof}
(a) Pada setiap potongan, $\sup(f+g)\le\sup f+\sup g$ dan $\inf(f+g)\ge\inf f+\inf g$.
Jadi $U(\mathcal{P},f+g)\le U(\mathcal{P},f)+U(\mathcal{P},g)$ dan
$L(\mathcal{P},f+g)\ge L(\mathcal{P},f)+L(\mathcal{P},g)$. Ambil $\varepsilon>0$ dan
pilih $\mathcal{P}$ yang membuat $U-L$ kurang dari $\varepsilon/2$ untuk $f$ dan $g$
sekaligus. Maka
\[
U(\mathcal{P},f+g)-L(\mathcal{P},f+g)
\le\bigl[U(\mathcal{P},f)-L(\mathcal{P},f)\bigr]
+\bigl[U(\mathcal{P},g)-L(\mathcal{P},g)\bigr]<\varepsilon ,
\]
sehingga $f+g$ terintegralkan. Nilainya diperoleh dengan mengurung dua arah, karena
$\varepsilon$ sebarang.

(b) Bila $f\le g$, maka $M_i(f)\le M_i(g)$ pada setiap potongan, sehingga
$U(\mathcal{P},f)\le U(\mathcal{P},g)$ untuk semua $\mathcal{P}$. Ambil infimum kedua
ruas.

(c) Bila $\mathcal{P}$ memuat $c$, jumlahnya terpecah menjadi bagian $[a,c]$ dan
$[c,b]$, dengan $L(\mathcal{P},f)=L(\mathcal{P}_1,f)+L(\mathcal{P}_2,f)$ dan hal sama
untuk $U$. Supremum memberi
$\underline{I}(f)\le\underline{I}_{[a,c]}(f)+\underline{I}_{[c,b]}(f)$, infimum
memberi arah sebaliknya. Untuk partisi tanpa $c$, tambahkan $c$ lebih dulu; lemma
penghalusan menjaga ketaksamaan.

(d) Karena $-|f|\le f\le|f|$, bagian (b) dan (a) memberi
$-\int|f|\le\int f\le\int|f|$. Keterintegralkan $|f|$ menyusul dari
$\bigl||f(x)|-|f(y)|\bigr|\le|f(x)-f(y)|$, yang memberi
$M_i(|f|)-m_i(|f|)\le M_i(f)-m_i(f)$.

(e) Pada setiap potongan, $M_i=m_i=k$, jadi $U=L=k(b-a)$ untuk semua partisi.
\end{proof}

Bagian (c) adalah alasan kita boleh memotong selang di titik mana pun; tanpanya fungsi
akumulasi tidak dapat didefinisikan dengan tenang.

\section{Kelas fungsi yang terintegralkan}
\tingkat{menengah}

Kriteria pada bagian sebelumnya adalah alat uji. Sekarang kita cari kelas fungsi yang
lolos uji itu tanpa harus mencari partisi khusus untuk setiap kasus. Dua kelas paling
berguna adalah fungsi kontinu dan fungsi monoton.

\begin{akibat}[Kekontinuan menjamin keterintegralan]
Jika $f$ kontinu pada $[a,b]$, maka $f$ terintegralkan Riemann pada $[a,b]$.
\end{akibat}

\begin{proof}
Ambil $\varepsilon>0$. Karena $[a,b]$ kompak, $f$ kontinu seragam di sana (Bab 11):
ada $\delta>0$ sehingga $|f(x)-f(t)|<\varepsilon/(b-a)$ bila $|x-t|<\delta$. Pilih
partisi dengan mesh kurang dari $\delta$. Pada setiap potongan, untuk semua $x,t$ di
potongan itu berlaku $|f(x)-f(t)|<\varepsilon/(b-a)$, sehingga
$M_i-m_i\le\varepsilon/(b-a)$. Maka
\[
U(\mathcal{P},f)-L(\mathcal{P},f)
=\sum_{i=1}^{n}(M_i-m_i)\,\Delta x_i
\le\frac{\varepsilon}{b-a}\sum_{i=1}^{n}\Delta x_i=\varepsilon .
\]
Kriteria keterintegralan menutup bukti.
\end{proof}

Perhatikan peran kekontinuan seragam: satu $\delta$ harus bekerja untuk semua potongan
sekaligus, sedangkan kekontinuan biasa hanya memberi $\delta$ yang bergantung pada
titik. Di sinilah Bab 11 dipakai, bukan sebagai hiasan.

Kelas kedua tidak memerlukan kekontinuan: fungsi monoton boleh melompat, asal
lompatannya berhingga banyak.

\begin{teorema}[Monoton menjamin keterintegralan]
Jika $f$ monoton pada $[a,b]$, maka $f$ terintegralkan Riemann pada $[a,b]$.
\end{teorema}

\begin{proof}
Andaikan $f$ naik. Bagi $[a,b]$ menjadi $n$ potongan sama lebar, $\Delta x=(b-a)/n$.
Pada potongan ke-$i$, nilai terkecil ada di ujung kiri dan nilai terbesar di ujung
kanan, sehingga $M_i-m_i=f(x_i)-f(x_{i-1})$. Maka
\[
U(\mathcal{P},f)-L(\mathcal{P},f)
=\Delta x\sum_{i=1}^{n}\bigl(f(x_i)-f(x_{i-1})\bigr)
=\frac{(b-a)\bigl(f(b)-f(a)\bigr)}{n},
\]
karena jumlah di tengah teleskopik. Pilih $n>(b-a)(f(b)-f(a))/\varepsilon$ dan kriteria
terpenuhi. Bila $f$ turun, terapkan hasil ini pada $-f$ dan pakai linearitas.
\end{proof}

Bukti itu tidak memakai $\delta$: yang dipakai hanya bahwa jumlah semua lompatan tidak
melebihi $f(b)-f(a)$, sehingga fungsi monoton mudah diintegralkan meski grafiknya penuh
sudut.

Kedua kelas itu dapat digabung: fungsi yang kontinu kecuali di berhingga banyak titik
tetap terintegralkan (Rudin, Teorema 6.10, tidak diulang di sini).

Sekarang sisi lain: fungsi terbatas yang gagal. Contoh paling bersih adalah fungsi
Dirichlet, bernilai $1$ pada bilangan rasional dan $0$ pada bilangan irasional.

\begin{contoh}[Fungsi Dirichlet tidak terintegralkan]\index{fungsi dirichlet}
Misalkan $D(x)=1$ bila $x$ rasional dan $D(x)=0$ bila $x$ irasional, pada $[0,1]$.
Setiap potongan memuat keduanya (Bab 3, kerapatan $\mathbb{Q}$), sehingga $M_i=1$ dan
$m_i=0$ untuk setiap $i$, dan
\[
U(\mathcal{P},D)=\sum_{i=1}^{n}1\cdot\Delta x_i=1,
\qquad
L(\mathcal{P},D)=0
\]
untuk \emph{setiap} partisi. Selisihnya selalu $1$, tidak pernah lebih kecil dari
$\varepsilon=1/2$, jadi menurut kriteria $D$ tidak terintegralkan meskipun terbatas.
\end{contoh}

Kegagalan itu bukan karena $D$ terputus-putus di banyak titik, melainkan karena
lompatannya tersebar di mana-mana. Fungsi berikut menunjukkan bahwa banyaknya titik
diskontinu bukan penghalang.

\begin{teorema}[Fungsi Thomae]\index{fungsi thomae}
Definisikan $T$ pada $[0,1]$ dengan $T(x)=1/q$ bila $x=p/q$ dalam bentuk paling
sederhana dengan $q\ge1$, dan $T(x)=0$ bila $x$ irasional. Maka $T$ terintegralkan
Riemann dan $\int_0^1 T=0$.
\end{teorema}

\begin{proof}
Ambil $\varepsilon>0$. Himpunan $S=\{x:T(x)\ge\varepsilon\}$ berisi pecahan dengan
penyebut $q\le1/\varepsilon$; untuk setiap $q$ hanya ada $q+1$ pecahan seperti itu di
dalam $[0,1]$, sehingga $S$ berhingga. Tutup $S$ dengan berhingga banyak selang
terbuka yang panjang totalnya kurang dari $\varepsilon$.

Bentuk partisi yang memuat semua ujung selang penutup itu. Potongan yang memuat $S$
panjangnya berjumlah kurang dari $\varepsilon$ dan di sana $M_i-m_i\le1$, jadi
sumbangannya kurang dari $\varepsilon$. Potongan lain memenuhi $M_i\le\varepsilon$ dan
$m_i\ge0$, sumbangannya paling banyak $\varepsilon$. Karena itu
$U(\mathcal{P},T)-L(\mathcal{P},T)<2\varepsilon$, dan kriteria menutup bukti. Nilainya
nol sebab setiap potongan memuat bilangan irasional, sehingga $L(\mathcal{P},T)=0$.
\end{proof}

Jadi $T$ terputus di setiap bilangan rasional, yang rapat di $[0,1]$, tetapi $T$ tetap
terintegralkan. Yang menentukan bukan jumlah titik diskontinu, melainkan besar ruang
yang ditempatinya: kriteria Lebesgue menyatakan fungsi terbatas terintegralkan Riemann
jika dan hanya jika titik diskontinunya berukuran nol (Rudin, Teorema 11.33(b); Abbott,
Bab 7). Kita mengutipnya tanpa bukti karena ukuran menuntut teori ukur.

\begin{tantangan}
Definisikan $h$ pada $[0,1]$ dengan $h(x)=1$ bila $x=p/q$ dalam bentuk paling sederhana
dan penyebut $q$ ganjil, serta $h(x)=0$ untuk semua $x$ lain, termasuk setiap bilangan
irasional. Apakah $h$ terintegralkan Riemann pada $[0,1]$? Bila ya, berapa nilai
integralnya? Bila tidak, buktikan memakai kriteria keterintegralan, bukan dengan
menghitung titik.

Petunjuk 1: jangan menghitung banyaknya titik tempat $h$ bernilai satu. Banyaknya tak
berhingga, dan kriteria tidak meminta hitungan itu.

Petunjuk 2: ambil sebarang potongan $[x_{i-1},x_i]$ dengan $x_{i-1}<x_i$. Potongan itu
memuat pecahan berpenyebut ganjil, jadi $M_i=1$, dan memuat bilangan irasional, jadi
$m_i=0$. Dengan begitu $U-L=1$ untuk setiap partisi.

Kalau mentok, kembali ke bagian Kelas fungsi yang terintegralkan, contoh Dirichlet.
\end{tantangan}

\kamus{keterintegralan}{Riemann integrability}{Sifat fungsi yang taksiran bawah dan
atasnya dapat dirapatkan sampai selisihnya hilang.}

\begin{untukanak}
Pernahkah kita melihat speedometer mobil, lalu bertanya berapa jarak yang sudah
ditempuh? Jarak tidak ditampilkan di situ, hanya kecepatan sesaat. Tetapi kalau kita
mencatat kecepatan setiap detik, kita bisa menjumlahkannya: satu detik dengan
kecepatan 10 meter per detik berarti 10 meter, detik berikutnya 12 meter, dan
seterusnya. Jumlah semua potongan itu adalah jarak.

Masalahnya, kecepatan berubah di tengah detik. Mobil bisa mempercepat lalu melambat
dalam satu detik yang sama. Kalau kita memakai kecepatan di awal detik, jarak yang
kita hitung terlalu kecil. Kalau memakai kecepatan di akhir detik, hasilnya terlalu
besar. Yang benar ada di antaranya.

Kalau kita mencatat sepuluh kali per detik, kedua taksiran makin dekat; seratus kali,
lebih dekat lagi. Matematikawan bertanya lebih jauh: kalau pencatatan bisa serapat apa
pun, apakah kedua taksiran itu menuju satu angka yang sama? Jawabannya tidak selalu ya,
dan di bab ini kita akan tahu kapan ya dan mengapa.
\end{untukanak}

\section{Integral sebagai fungsi akumulasi}
\tingkat{menengah}

Sampai sini integral hanyalah bilangan yang melekat pada selang. Kita dapat
memandangnya sebagai fungsi dengan membiarkan ujung atas bergerak: integralnya menjadi
laporan tentang berapa banyak yang sudah terkumpul sampai titik itu.

\begin{definisi}[Fungsi akumulasi]\index{fungsi akumulasi}
Misalkan $f$ terintegralkan pada $[a,b]$. Untuk $x\in[a,b]$ didefinisikan
\[
F(x)=\int_a^x f(t)\,dt .
\]
Untuk $x<a$ dipakai kesepakatan $\int_a^x f=-\int_x^a f$, dan untuk $x=a$ nilainya
nol. Fungsi $F$ disebut \emph{fungsi akumulasi} dari $f$ pada $[a,b]$.
\end{definisi}

Dengan kesepakatan itu sifat aditif berlaku untuk urutan titik apa pun, sehingga
\begin{equation}
F(y)-F(x)=\int_x^{y} f(t)\,dt
\end{equation}
untuk semua $x,y\in[a,b]$. Inilah bentuk yang dipakai pada semua bukti berikutnya:
pertambahan tabungan antara dua waktu adalah integral laju pemasukan di antaranya.

Fungsi akumulasi selalu lebih baik daripada fungsi asalnya: ia kontinu walaupun $f$
tidak, dan terdiferensialkan di setiap titik tempat $f$ kontinu.

\begin{proposisi}[Akumulasi selalu kontinu]
Jika $f$ terbatas pada $[a,b]$ dengan $|f|\le M$ dan terintegralkan, maka $F$ yang
didefinisikan di atas memenuhi
\[
|F(y)-F(x)|\le M\,|y-x|
\]
untuk semua $x,y\in[a,b]$. Khususnya $F$ kontinu pada $[a,b]$.
\end{proposisi}

\begin{proof}
Untuk $x<y$, sifat monoton dan bagian (d) Proposisi sifat dasar memberi
\[
|F(y)-F(x)|=\Bigl|\int_x^{y}f\Bigr|\le\int_x^{y}|f|\le M(y-x).
\]
Untuk $y<x$ hasil yang sama dengan tanda berlawanan. Jadi $F$ memenuhi syarat
Lipschitz dengan konstanta $M$, dan setiap fungsi Lipschitz kontinu.
\end{proof}

Sekarang arah pertama hubungan turunan-integral. Kalimat sehari-harinya: laju
pertumbuhan tabungan pada suatu saat sama dengan laju pemasukan pada saat itu.

\begin{teorema}[Teorema dasar kalkulus, bagian pertama]\index{teorema dasar kalkulus}
Misalkan $f$ terintegralkan pada $[a,b]$ dan kontinu di titik $x_0\in(a,b)$. Maka
fungsi akumulasi $F$ terdiferensialkan di $x_0$ dan
\[
F'(x_0)=f(x_0).
\]
\end{teorema}

\begin{proof}
Ambil $h\ne0$ dengan $x_0+h\in[a,b]$. Dari bentuk pertambahan,
\[
\frac{F(x_0+h)-F(x_0)}{h}-f(x_0)
=\frac{1}{h}\int_{x_0}^{x_0+h}\bigl(f(t)-f(x_0)\bigr)\,dt ,
\]
sebab integral fungsi konstan $f(x_0)$ pada selang panjang $|h|$ bernilai
$f(x_0)h$.

Ambil $\varepsilon>0$. Karena $f$ kontinu di $x_0$, ada $\delta>0$ sehingga
$|t-x_0|<\delta$ berakibat $|f(t)-f(x_0)|<\varepsilon$. Untuk $0<|h|<\delta$ seluruh
selang integrasi berada di dalam $(x_0-\delta,x_0+\delta)$, sehingga dengan sifat
monoton dan bagian (d) Proposisi sifat dasar,
\[
\left|\frac{1}{h}\int_{x_0}^{x_0+h}\bigl(f(t)-f(x_0)\bigr)dt\right|
\le\frac{1}{|h|}\left|\int_{x_0}^{x_0+h}\bigl|f(t)-f(x_0)\bigr|\,dt\right|
\le\frac{1}{|h|}\,\varepsilon\,|h|=\varepsilon .
\]
Karena $\varepsilon$ sebarang, hasil bagi selisih itu menuju $f(x_0)$.
\end{proof}

Teorema itu hanya menuntut kekontinuan di satu titik. Di tempat $f$ melompat, turunan
$F$ boleh gagal ada, dan itu bukan kelemahan bukti: untuk $f(t)=\operatorname{sgn} t$
pada $[-1,1]$ berlaku $F(x)=\int_0^x\operatorname{sgn} t\,dt=|x|$, yang tidak
terdiferensialkan di $0$, persis di tempat tanda $f$ berubah.

Satu akibat penting: setiap fungsi kontinu punya antiturunan, dan antiturunannya dapat
ditulis sebagai integral. Bab 12 menurunkan aturan turunan tetapi tidak menjawab apakah
setiap fungsi kontinu merupakan turunan dari sesuatu. Sekarang pertanyaan itu terjawab,
dan jawabannya konstruktif.

\begin{akibat}[Antiturunan selalu ada]
Jika $f$ kontinu pada $[a,b]$, maka $F(x)=\int_a^x f$ terdiferensialkan pada $[a,b]$
dengan $F'=f$. Jadi $F$ adalah antiturunan $f$ pada $[a,b]$.
\end{akibat}

\begin{proof}
Kekontinuan $f$ memberi keterintegralan (Akibat sebelumnya) dan kekontinuan di setiap
titik, sehingga teorema bagian pertama berlaku di setiap $x_0\in(a,b)$. Di titik ujung
dipakai turunan satu sisi.
\end{proof}

Aturan rantai memberi satu keluwesan tambahan: bila batas atas bergantung pada $x$,
\[
\frac{d}{dx}\int_a^{g(x)}f(t)\,dt=f\bigl(g(x)\bigr)\,g'(x),
\]
asalkan $g$ terdiferensialkan. Contohnya,
$\frac{d}{dx}\int_0^{x^2}e^{-t^2}dt=2x\,e^{-x^4}$.

Kedua arah itu mudah tertukar, jadi ada baiknya digambar sekali.
Gambar~\ref{fig:tdk-dua-arah} menaruh keduanya pada satu diagram: panah atas
mengumpulkan potongan tipis, panah bawah menurunkan. Yang dijamin teorema adalah bahwa
menurunkan hasil pengumpulan mengembalikan fungsi semula, bukan kebalikan aljabar yang
rapi, sebab antiturunan hanya pasti sampai konstanta.

\begin{figure}[htbp]
\centering
\begin{tikzpicture}[>=Stealth,every node/.style={font=\footnotesize},
  kotak/.style={draw,rounded corners=3pt,align=center,inner sep=5pt,
    minimum height=10mm,text width=3.3cm}]
\node[kotak,fill=layerindigo!8,draw=layerindigo] (f) at (0,0) {fungsi $f$};
\node[kotak,fill=layerbrown!8,draw=layerbrown] (F) at (7.6,0)
  {fungsi akumulasi $F$};
\draw[->,line width=0.9pt,draw=layerindigo] ([yshift=2mm]f.north) to[bend left=24]
  node[above,text=layerindigo,align=center]
  {integral tentu: kumpulkan potongan tipis} ([yshift=2mm]F.north);
\draw[->,line width=0.9pt,draw=layerbrown] ([yshift=-2mm]F.south) to[bend left=24]
  node[below,text=layerbrown,align=center]
  {turunan: $F'=f$, inilah bagian 1} ([yshift=-2mm]f.south);
\node[draw=layergrey,dashed,rounded corners=3pt,align=center,inner sep=4pt]
  at (3.8,-2.7)
  {bagian 2 menutup lingkaran: $\int_a^b F'(x)\,dx=F(b)-F(a)$};
\end{tikzpicture}
\caption{Dua arah teorema dasar kalkulus sebagai panah bolak-balik. Panah atas:
mengumpulkan potongan tipis menghasilkan $F(x)=\int_a^x f$. Panah bawah: menurunkan $F$
mengembalikan $f$, dan itulah bagian pertama. Kotak putus-putus di bawah adalah bagian
kedua, yang memakai antiturunan untuk menghitung integral. Untuk $f(x)=x^2$ pada
$[0,1]$, $F(x)=x^3/3$, dan bagian kedua memberi $F(1)-F(0)=1/3$, angka yang sama dengan
pengurungan pada Gambar~\ref{fig:rect-halus}.}
\label{fig:tdk-dua-arah}
\end{figure}

\section{Teorema dasar kalkulus, bagian kedua}
\tingkat{lanjut}

Bagian pertama berjalan dari integral ke turunan. Bagian kedua berjalan sebaliknya:
bila kita sudah memegang antiturunan, integral dapat dihitung dengan satu pengurangan.
Kalimat sehari-harinya: total perubahan sama dengan nilai akhir dikurangi nilai awal.

\begin{teorema}[Teorema dasar kalkulus, bagian kedua]
Misalkan $F$ kontinu pada $[a,b]$, terdiferensialkan pada $(a,b)$, dan $F'$
terintegralkan Riemann pada $[a,b]$. Maka
\[
\int_a^b F'(x)\,dx=F(b)-F(a).
\]
\end{teorema}

\begin{proof}
Ambil partisi $\mathcal{P}=\{x_0,\dots,x_n\}$ dari $[a,b]$. Pada setiap subselang, $F$
kontinu pada $[x_{i-1},x_i]$ dan terdiferensialkan di dalamnya, sehingga teorema nilai
rata-rata (Bab 12) memberi $t_i\in(x_{i-1},x_i)$ dengan
\[
F(x_i)-F(x_{i-1})=F'(t_i)\,\Delta x_i .
\]
Karena $m_i\le F'(t_i)\le M_i$, penjumlahan atas $i$ memberi
\[
L(\mathcal{P},F')\le\sum_{i=1}^{n}\bigl(F(x_i)-F(x_{i-1})\bigr)\le U(\mathcal{P},F').
\]
Jumlah di tengah teleskopik dan sama dengan $F(b)-F(a)$, jadi bilangan itu terkurung di
antara jumlah bawah dan jumlah atas setiap partisi. Karena $F'$ terintegralkan, kedua
kurungan menuju nilai yang sama, yaitu $\int_a^b F'$.
\end{proof}

\begin{resep}
\begin{enumerate}
\item Cari $F$ dengan $F'=f$, lalu periksa dengan menurunkan $F$. Bila turunannya tidak
      kembali ke $f$, calonnya salah dan langkah berikutnya sia-sia.
\item Tulis $F$ pada batas atas, tulis $F$ pada batas bawah, lalu kurangkan:
      $F(b)-F(a)$. Jangan sisipkan $+C$; konstanta itu saling menghapus.
\item Rapikan tandanya bila selangnya terbalik: $\int_b^a f=-\int_a^b f$, dan
      $\int_a^a f=0$.
\item Uji dengan satu potongan kasar, trapesium atau \texttt{quad}. Bila hasilnya jauh,
      galatnya ada di langkah 1 atau 2, bukan di langkah 3.
\end{enumerate}
Peringatan urutan: memeriksa antiturunan harus terjadi sebelum mengurangkan. Dua nilai
fungsi sembarang tetap memberi selisih yang rapi dan masuk akal, sehingga keliru di
langkah 1 tidak terlihat dari bentuk jawabannya.
\end{resep}

Begitu pula kedua bagian itu tidak saling menggantikan: bagian pertama menurunkan
fungsi akumulasi tanpa memerlukan antiturunan, sedangkan bagian kedua memerlukan
antiturunan dan keterintegralan turunannya.

Bagian kedua menjelaskan mengapa hasilnya disebut perubahan bersih, bukan luas. Untuk
$F(x)=\frac{x^4}{4}-\frac{x^2}{2}$ pada $[-1,1]$, $F'(x)=x^3-x$, sehingga
\[
\int_{-1}^{1}(x^3-x)\,dx=F(1)-F(-1)=\Bigl(\frac14-\frac12\Bigr)-\Bigl(\frac14-\frac12\Bigr)=0 .
\]
Daerah di atas sumbu seluas $1/4$, daerah di bawah sumbu juga $1/4$, jadi luas totalnya
$1/2$ sedangkan integralnya nol karena tandanya berlawanan.
Gambar~\ref{fig:luas-bertanda} memperlihatkan keduanya.

\begin{figure}[htbp]
\centering
\begin{tikzpicture}
\begin{axis}[width=0.74\linewidth,height=5.4cm,
  axis lines=middle,xmin=-1.3,xmax=1.3,ymin=-0.6,ymax=0.6,
  xtick={-1,0,1},ytick=\empty,
  tick label style={font=\scriptsize},
  xlabel={$x$},ylabel={$y=x^3-x$},
  xlabel style={font=\scriptsize},ylabel style={font=\scriptsize,rotate=-90}]
\addplot[fill=layerjade!22,draw=layerjade!70,domain=-1:0,samples=60]
  {x^3-x} \closedcycle;
\addplot[fill=layerplum!22,draw=layerplum!70,domain=0:1,samples=60]
  {x^3-x} \closedcycle;
\addplot[thick,ink,domain=-1.2:1.2,samples=120] {x^3-x};
\node[font=\scriptsize,text=layerjade!55!black] at (axis cs:-0.52,0.36)
  {$+\frac14$};
\node[font=\scriptsize,text=layerplum!55!black] at (axis cs:0.52,-0.36)
  {$-\frac14$};
\end{axis}
\end{tikzpicture}
\caption{Luas bertanda untuk $F'(x)=x^3-x$ pada $[-1,1]$. Daerah di atas sumbu-$x$
menyumbang $+1/4$, daerah di bawah sumbu menyumbang $-1/4$, sehingga
$\int_{-1}^{1}(x^3-x)\,dx=F(1)-F(-1)=0$. Luas geometrinya tetap $1/2$; yang nol adalah
jumlah bertandanya, dan itulah yang dihitung teorema bagian kedua.}
\label{fig:luas-bertanda}
\end{figure}

Hipotesis keterintegralan pada bagian kedua tidak berlebihan. Ada fungsi
terdiferensialkan yang turunannya terbatas tetapi tidak Riemann-terintegralkan; contoh
klasiknya ditemukan Volterra pada abad kesembilan belas, dan kita tidak
membuktikannya di sini. Teorema bagian kedua bukan pernyataan tanpa syarat.

\section{Substitusi dan integral parsial}
\tingkat{lanjut}

Dua aturan hitung integral yang paling sering dipakai bukan aksioma baru: keduanya
aturan turunan yang dibalik, dan pembuktiannya hanya memakai teorema dasar kalkulus.
Yang pertama membalik aturan rantai.

\begin{proposisi}[Substitusi]\index{substitusi}
Misalkan $\varphi$ terdiferensialkan dengan $\varphi'$ kontinu pada $[a,b]$, dan $f$
kontinu pada selang yang memuat $\varphi([a,b])$. Maka
\[
\int_a^b f\bigl(\varphi(t)\bigr)\,\varphi'(t)\,dt
=\int_{\varphi(a)}^{\varphi(b)}f(u)\,du .
\]
\end{proposisi}

\begin{proof}
Karena $f$ kontinu, ia mempunyai antiturunan $G(u)=\int_c^u f$, dengan $c$ di dalam
selang nilai $\varphi$; menurut teorema bagian pertama, $G'=f$. Aturan rantai (Bab 12)
memberi
\[
\frac{d}{dt}G\bigl(\varphi(t)\bigr)=G'\bigl(\varphi(t)\bigr)\varphi'(t)
=f\bigl(\varphi(t)\bigr)\varphi'(t).
\]
Ruas kanan kontinu, jadi terintegralkan. Teorema bagian kedua pada fungsi
$G\circ\varphi$ memberi
\[
\int_a^b f\bigl(\varphi(t)\bigr)\varphi'(t)\,dt
=G\bigl(\varphi(b)\bigr)-G\bigl(\varphi(a)\bigr)
=\int_{\varphi(a)}^{\varphi(b)}f(u)\,du .
\]
\end{proof}

Aturan kedua membalik aturan perkalian.

\begin{proposisi}[Integral parsial]\index{integral parsial}
Misalkan $F$ dan $G$ terdiferensialkan pada $[a,b]$ dengan $F'$ dan $G'$ kontinu. Maka
\[
\int_a^b F(x)\,G'(x)\,dx=F(b)G(b)-F(a)G(a)-\int_a^b F'(x)\,G(x)\,dx .
\]
\end{proposisi}

\begin{proof}
Misalkan $H=FG$. Aturan perkalian memberi $H'=F'G+FG'$, dan ruas kanan kontinu
sehingga $H'$ terintegralkan. Teorema bagian kedua pada $H$ memberi
\[
F(b)G(b)-F(a)G(a)=\int_a^b H'=\int_a^b F'G+\int_a^b FG' .
\]
Pindahkan suku pertama ke ruas kiri.
\end{proof}

\kamus{antiturunan}{antiderivative}{Fungsi yang turunannya adalah fungsi yang sedang
kita integralkan, seperti arah pulang dari sebuah jalan: ia hanya berguna kalau arah
berangkatnya sudah diketahui.}

Sebagai pemeriksaan angka, ambil $\int_0^1 x\,e^{x^2}dx$. Substitusi $u=x^2$,
$du=2x\,dx$ memberi
\[
\int_0^1 x\,e^{x^2}dx=\frac12\int_0^1 e^{u}du=\frac{e-1}{2}=0.8591409142295225 .
\]
\texttt{scipy.integrate.quad} memberi $0.8591409142295227$, selisih $2\times10^{-16}$
dari pembulatan mesin. Untuk aturan kedua, $F(x)=x$ dan $G'(x)=e^x$ memberi
$\int_0^1 xe^xdx=e-(e-1)=1$, dan \texttt{quad} mengembalikan $1.0$ dengan perkiraan
galat $1.1\times10^{-14}$.

\section{Teorema nilai rata-rata untuk integral}
\tingkat{menengah}

Rata-rata biasanya dihitung dengan menjumlah lalu membagi banyaknya data. Untuk fungsi
yang berubah terus, penjumlahan diganti integral dan pembaginya panjang selang: rata-rata
adalah tinggi persegi panjang yang luasnya sama dengan luas di bawah kurva.

\begin{definisi}[Nilai rata-rata]
Untuk $f$ terintegralkan pada $[a,b]$, \emph{nilai rata-rata} $f$ pada selang itu
adalah
\[
\bar{f}=\frac{1}{b-a}\int_a^b f(x)\,dx .
\]
\end{definisi}

\begin{teorema}[Nilai rata-rata untuk integral]\index{nilai rata-rata integral}
Jika $f$ kontinu pada $[a,b]$, maka terdapat $c\in[a,b]$ dengan
\[
f(c)=\frac{1}{b-a}\int_a^b f(x)\,dx .
\]
\end{teorema}

\begin{proof}
Karena $f$ kontinu pada selang tertutup terbatas, ia mencapai nilai terkecil $m$ dan
nilai terbesar $M$ di sana (Bab 11). Sifat monoton memberi
\[
m(b-a)\le\int_a^b f\le M(b-a),
\]
sehingga $\bar{f}\in[m,M]$. Menurut teorema nilai antara (Bab 11), fungsi kontinu $f$
mengambil setiap nilai di antara $m$ dan $M$, jadi ada $c\in[a,b]$ dengan
$f(c)=\bar{f}$.
\end{proof}

Sebagai angka, untuk $f(x)=x^2$ pada $[0,1]$ rata-ratanya $1/3$, dan titik yang
dicapai teorema adalah $c=1/\sqrt3=0.5773502691896258$, sebab
$c^2=0.3333333333333334$. Titik itu bukan titik tengah selang: teorema menjamin
keberadaan, bukan lokasi.

\section{Contoh hitung}
\tingkat{menengah}

Sekarang kita kerjakan tiga hal: jumlah atas dan bawah $f(x)=x^2$ pada $[0,1]$,
perbandingan lima aturan kuadratur, dan satu integral tanpa antiturunan elementer.
Pecahan dihitung dengan \texttt{fractions.Fraction}, tanpa pembulatan.

Mulai dari $n$ potong sama lebar dengan $x_i=i/n$ dan $\Delta x=1/n$, sehingga
\[
U_n=\sum_{i=1}^{n}\Bigl(\frac{i}{n}\Bigr)^2\frac1n=\frac{1}{n^3}\sum_{i=1}^{n}i^2
=\frac{(n+1)(2n+1)}{6n^2}=\frac13+\frac{1}{2n}+\frac{1}{6n^2},
\]
memakai rumus jumlah kuadrat dari Bab 1. Dengan cara yang sama
$L_n=\frac13-\frac{1}{2n}+\frac{1}{6n^2}$, sehingga
\begin{equation}
U_n-L_n=\frac1n .
\end{equation}
Hasil itu bukan hampiran. Untuk setiap $n$ selisihnya persis $1/n$; karena $1/n$ dapat
dibuat lebih kecil dari $\varepsilon$ mana pun dengan $n>1/\varepsilon$, kriteria
keterintegralan terpenuhi. Jadi tabel berikut sekaligus bukti bahwa
$\int_0^1x^2dx=1/3$: kedua jumlah mengapit $1/3$, dan jaraknya menuju nol.

\begin{center}\small
\begin{tabular}{rrrr}
\hline
$n$ & $U_n$ & $L_n$ & $U_n-L_n$ \\
\hline
1 & $1.0000000000$ & $0.0000000000$ & $1$ \\
4 & $0.4687500000$ & $0.2187500000$ & $1/4$ \\
16 & $0.3652343750$ & $0.3027343750$ & $1/16$ \\
64 & $0.3411865234$ & $0.3255615234$ & $1/64$ \\
256 & $0.3352890015$ & $0.3313827515$ & $1/256$ \\
1024 & $0.3338217735$ & $0.3328452110$ & $1/1024$ \\
\hline
\end{tabular}
\end{center}

Pada $n=1024$ selisihnya $9.77\times10^{-4}$ dan jarak $U_n$ ke $1/3$ separuh dari itu.
Setiap kali $n$ dilipatgandakan selisihnya tepat separuh, dan itulah sebabnya jumlah
atas dan bawah bukan cara efisien untuk mencapai ketelitian tinggi.

\begin{tangan}
Bahan: satu lembar kertas berpetak, gunting, penggaris, selotip, pensil.

Langkah. (1) Di kertas berpetak gambar kurva yang naik makin cepat: satu petak dari
tepi kiri tingginya satu petak, dua petak tingginya empat petak, tiga petak tingginya
sembilan petak. (2) Gunting daerah di bawah kurva menjadi jalur tegak selebar satu
petak. (3) Susun jalur itu berdampingan di meja, dari terpendek ke terpanjang, alas
rata. (4) Perhatikan sisi atasnya: yang muncul tangga bergerigi, bukan garis mulus.
(5) Ulangi dengan jalur selebar setengah petak, lalu seperempat petak.

Yang harus kita lihat: gerigi di sisi atas makin halus ketika jalur makin tipis,
sedangkan luas totalnya tidak berubah, sebab memotong lalu menyusun ulang tidak
menambah atau mengurangi kertas. Jadi luas daerah melengkung itu sama dengan luas
bangun bergerigi, dan bangun bergerigi itu makin mirip satu persegi panjang. Tinggi
persegi panjang itulah yang nanti kita sebut nilai rata-rata fungsi pada selang itu.
\end{tangan}

Sekarang bandingkan lima aturan kuadratur pada integral yang sama: ujung kiri, ujung
kanan, titik tengah $x_i=(i-\frac12)/n$, trapesium, dan Simpson dengan bobot
$1,4,2,4,\dots,4,1$ dibagi $3n$. Galat dihitung terhadap nilai eksak $1/3$.

\begin{center}\small
\begin{tabular}{rrrrrr}
\hline
$n$ & kiri & kanan & titik tengah & trapesium & Simpson \\
\hline
10 & $4.833\times10^{-2}$ & $5.167\times10^{-2}$ & $8.333\times10^{-4}$
   & $1.667\times10^{-3}$ & $1.1\times10^{-16}$ \\
100 & $4.983\times10^{-3}$ & $5.017\times10^{-3}$ & $8.333\times10^{-6}$
   & $1.667\times10^{-5}$ & $0$ \\
1000 & $4.998\times10^{-4}$ & $5.002\times10^{-4}$ & $8.333\times10^{-8}$
   & $1.667\times10^{-7}$ & $0$ \\
\hline
\end{tabular}
\end{center}

Tiga pola terbaca. Galat aturan kiri dan kanan turun sepuluh kali setiap $n$ naik
sepuluh kali: orde $\mathcal{O}(1/n)$, sekitar $1/(2n)$ karena $f$ naik. Galat titik
tengah dan trapesium turun seratus kali, orde $\mathcal{O}(1/n^2)$, dengan nilai
$1/(12n^2)$ dan $1/(6n^2)$; keduanya berlawanan tanda dan trapesium dua kali lipat,
sehingga $(2M+T)/3$ menghapus sukunya. Simpson tepat sampai pembulatan mesin, sebab
aturannya eksak untuk polinomial berderajat paling tinggi tiga.

Integral ketiga tidak punya antiturunan elementer:
\[
\int_0^1 e^{-x^2}dx .
\]
Ketiadaan itu akibat teorema Liouville, kita kutip tanpa bukti. Acuannya fungsi galat
$\frac{\sqrt{\pi}}{2}\operatorname{erf}(1)=0.746824132812427$; \texttt{quad} memberi
angka yang sama dengan perkiraan galat $8.3\times10^{-15}$. Aturan Simpson memberi galat
$3.563\times10^{-4}$ pada $n=2$, $1.988\times10^{-6}$ pada $n=8$, dan
$7.795\times10^{-9}$ pada $n=32$.

\kertas{Untuk $f(x)=x^2$ pada $[0,1]$, turunkan sendiri rumus
$U_n-L_n=1/n$ dari rumus jumlah kuadrat, lalu hitung $U_{10}$ dan $L_{10}$ sebagai
pecahan. Periksa bahwa rata-ratanya $(U_{10}+L_{10})/2$ sama dengan
$1/3+1/(6\cdot100)$, bukan $1/3$. Jelaskan mengapa rata-rata itu tetap bukan taksiran
terbaik, padahal galatnya jauh lebih kecil daripada galat $U_{10}$ sendiri.}

\section{Integral tak wajar dan kaitannya dengan deret}
\tingkat{lanjut}

Selang integrasi sejauh ini selalu terbatas dan fungsinya terbatas. Banyak perhitungan
nyata tidak begitu: kita ingin luas daerah yang memanjang tanpa ujung, atau integral
fungsi yang meledak di ujung selang. Caranya bukan memperluas definisi Riemann,
melainkan menulis integral itu sebagai limit integral biasa.

\begin{definisi}[Integral tak wajar]\index{integral tak wajar}
Bila $f$ terintegralkan Riemann pada $[a,R]$ untuk setiap $R>a$, didefinisikan
\[
\int_a^\infty f=\lim_{R\to\infty}\int_a^R f ,
\]
asalkan limit itu ada dan berhingga. Bila $f$ terintegralkan pada $[c,b]$ untuk setiap
$c\in(a,b]$ tetapi tidak terbatas di dekat $a$, didefinisikan
\[
\int_a^b f=\lim_{c\to a^{+}}\int_c^b f .
\]
Dalam kedua kasus, integral disebut \emph{konvergen} bila limitnya ada, dan
\emph{berdivergen} bila tidak.
\end{definisi}

Contoh terpenting adalah keluarga $f(x)=x^{-p}$ pada $[1,\infty)$. Untuk $p\ne1$,
\[
\int_1^R x^{-p}dx=\frac{R^{1-p}-1}{1-p},
\]
sehingga limitnya ada hanya bila $p>1$, dengan nilai $1/(p-1)$; untuk $p=1$ hasilnya
$\ln R$, yang menuju tak berhingga. Jadi $\int_1^\infty x^{-2}dx=1$, dan integral
parsialnya menuju nilai itu dengan laju $1/R$: $0.9$, $0.99$, $0.999$ untuk
$R=10,100,1000$.

Kaitan dengan deret Bab 5 bukan kebetulan: untuk $f$ positif dan turun, luas persegi
panjang setinggi $f(n)$ mengurung luas di bawah kurva, sehingga
\[
\sum_{n=N+1}^{\infty}f(n)\ \le\ \int_N^\infty f(x)\,dx\ \le\ \sum_{n=N}^{\infty}f(n) .
\]
Karena itu uji integral berlaku: $\sum f(n)$ konvergen jika dan hanya jika
$\int_1^\infty f$ konvergen. Untuk $f(x)=1/x^2$ keduanya konvergen, dengan
$\sum_{n=1}^{\infty}1/n^2=\pi^2/6=1.6449340668482264$. Ketaksamaan kiri juga memberi
batas galat, dan makin besar $N$ makin dekat galat ke batas itu:

\begin{center}\small
\begin{tabular}{rrrr}
\hline
$N$ & galat jumlah parsial & $1/N$ \\
\hline
10 & $9.516634\times10^{-2}$ & $0.1$ \\
100 & $9.950167\times10^{-3}$ & $0.01$ \\
1000 & $9.995002\times10^{-4}$ & $0.001$ \\
\hline
\end{tabular}
\end{center}

Perlu ditegaskan: integral tak wajar bukan integral Riemann, melainkan limitnya. Karena
itu substitusi dan integral parsial hanya berlaku bila limitnya ada, dan untuk fungsi
berganti tanda kekonvergenannya belum tentu mutlak.

\section{Lab komputasi}
\tingkat{lanjut}

Kode berikut dijalankan dengan Python 3, \texttt{scipy} 1.18.0, dan modul
\texttt{fractions} untuk aritmetika eksak. Semua partisi sama lebar pada $[0,1]$.

\begin{lab}
\begin{verbatim}
import math
from fractions import Fraction as F
from scipy.integrate import quad

def jumlah(n):
    U = sum(F(i * i, n ** 3) for i in range(1, n + 1))
    L = sum(F(i * i, n ** 3) for i in range(n))
    return U, L

print("n        U_n          L_n          U_n - L_n")
for n in (1, 2, 4, 8, 16, 32, 64, 128, 256, 512, 1024):
    U, L = jumlah(n)
    print(f"{n:5d}  {float(U):.10f}  {float(L):.10f}  {float(U - L):.2e}")

def aturan(f, n):
    h = 1 / n
    kiri = sum(f(i * h) for i in range(n)) * h
    kanan = sum(f(i * h) for i in range(1, n + 1)) * h
    tengah = sum(f((i + 0.5) * h) for i in range(n)) * h
    trap = (f(0) / 2 + sum(f(i * h) for i in range(1, n)) + f(1) / 2) * h
    simp = (f(0) + f(1) + sum((4 if i % 2 else 2) * f(i * h)
                              for i in range(1, n))) * h / 3
    return kiri, kanan, tengah, trap, simp

f = lambda x: x * x
print()
print("galat |hampiran - 1/3|, integral x^2 pada [0,1]")
print("   n      kiri      kanan     tengah   trapesium    Simpson")
for n in (10, 100, 1000):
    galat = [f"{abs(v - 1 / 3):.3e}" for v in aturan(f, n)]
    print(f"{n:4d}  " + "  ".join(galat))

g = lambda x: math.exp(-x * x)
v, e = quad(g, 0, 1)
acuan = math.sqrt(math.pi) / 2 * math.erf(1)
print()
print(f"quad exp(-x^2) = {v:.15f}   galat perkiraan {e:.1e}")
print(f"acuan erf      = {acuan:.15f}   selisih {v - acuan:+.1e}")
print(f"quad x^2       = {quad(f, 0, 1)[0]:.17f}")
print()
print("galat |Simpson - acuan| untuk exp(-x^2)")
for n in (2, 4, 8, 16, 32):
    print(f"  n={n:3d}  {abs(aturan(g, n)[4] - acuan):.3e}")
print()
print("galat |trapesium - acuan| untuk exp(-x^2)")
for n in (10, 100, 1000):
    print(f"  n={n:4d}  {abs(aturan(g, n)[3] - acuan):.3e}")
\end{verbatim}
\end{lab}

Keluaran yang diperoleh, apa adanya:

\begin{verbatim}
n        U_n          L_n          U_n - L_n
    1  1.0000000000  0.0000000000  1.00e+00
    2  0.6250000000  0.1250000000  5.00e-01
    4  0.4687500000  0.2187500000  2.50e-01
    8  0.3984375000  0.2734375000  1.25e-01
   16  0.3652343750  0.3027343750  6.25e-02
   32  0.3491210938  0.3178710938  3.12e-02
   64  0.3411865234  0.3255615234  1.56e-02
  128  0.3372497559  0.3294372559  7.81e-03
  256  0.3352890015  0.3313827515  3.91e-03
  512  0.3343105316  0.3323574066  1.95e-03
 1024  0.3338217735  0.3328452110  9.77e-04

galat |hampiran - 1/3|, integral x^2 pada [0,1]
   n      kiri      kanan     tengah   trapesium    Simpson
  10  4.833e-02  5.167e-02  8.333e-04  1.667e-03  1.110e-16
 100  4.983e-03  5.017e-03  8.333e-06  1.667e-05  0.000e+00
1000  4.998e-04  5.002e-04  8.333e-08  1.667e-07  0.000e+00

quad exp(-x^2) = 0.746824132812427   galat perkiraan 8.3e-15
acuan erf      = 0.746824132812427   selisih +1.1e-16
quad x^2       = 0.33333333333333337

galat |Simpson - acuan| untuk exp(-x^2)
  n=  2  3.563e-04
  n=  4  3.125e-05
  n=  8  1.988e-06
  n= 16  1.246e-07
  n= 32  7.795e-09

galat |trapesium - acuan| untuk exp(-x^2)
  n=  10  6.133e-04
  n= 100  6.131e-06
  n=1000  6.131e-08
\end{verbatim}

Empat hal terbaca. Selisih terbagi dua setiap kali $n$ berlipat dua, sesuai $1/1024$
pada baris terakhir. Galat kiri dan kanan turun sepuluh kali lipat, titik tengah dan
trapesium seratus kali lipat. Kolom Simpson berisi $1.110\times10^{-16}$ lalu nol, sisa
pembulatan, dan galat Simpson untuk $e^{-x^2}$ turun sekitar enam belas kali lipat
setiap $n$ berlipat dua.

Kode itu tidak membuktikan satu pun teorema. Ia memeriksa bahwa rumus bab ini memberi
angka yang disebutkan. Bukti keterintegralan $x^2$ ada pada perhitungan $U_n-L_n=1/n$
di bagian sebelumnya: keluaran program adalah data, bukan argumen.

\begin{untukdosen}
Urutan mengajar yang saya sarankan: mulai dari dua jawaban yang berbeda di pembuka, lalu
definisi partisi dan hitung $U$, $L$ untuk $x^2$ dengan empat potongan di papan. Setelah
itu lemma penghalusan, kriteria, baru kelas fungsi. Jangan menunda teorema dasar
kalkulus: mahasiswa perlu melihat kerja keras kriteria berbuah menjadi alat hitung.

Dua miskonsepsi yang sering muncul. Pertama, integral dianggap sama dengan antiturunan,
sehingga $F'=f$ dianggap berlaku karena definisi, padahal itu teorema. Kedua, fungsi
dengan banyak titik diskontinu dianggap pasti tidak terintegralkan, karena contoh
Dirichlet terlalu sering dipakai sebagai satu-satunya contoh. Fungsi Thomae menghapus
anggapan itu dalam satu halaman.

Aktivitas papan tulis: minta mahasiswa menebak $\int_0^1x^2dx$ hanya dari $U_n$ dan
$L_n$ untuk $n=1,2,4$, lalu menyebutkan berapa $n$ yang diperlukan agar galat di bawah
$10^{-3}$ tanpa menghitung integralnya. Jawaban $n>1000$ muncul dari $1/n$, dan momen
itu menjelaskan mengapa antiturunan lebih berharga daripada kuadratur. Tanda siap
lanjut: mahasiswa dapat menjelaskan perbedaan arah kedua bagian teorema dasar kalkulus,
dan menyebutkan satu fungsi terintegralkan yang tidak kontinu.
\end{untukdosen}

\section{Ringkasan}
\tingkat{dasar}

\begin{ringkasan}
\begin{itemize}
\item Integral Riemann adalah satu-satunya bilangan yang terkurung di antara semua
      jumlah bawah dan semua jumlah atas; bila kedua kurungan bertemu, nilainya disebut
      integral.
\item Kriteria keterintegralan menggantikan pemeriksaan tak berhingga banyak partisi
      dengan satu pertidaksamaan $U(\mathcal{P},f)-L(\mathcal{P},f)<\varepsilon$.
\item Fungsi kontinu dan fungsi monoton terintegralkan; fungsi Dirichlet tidak, karena
      setiap potongannya memuat rasional dan irasional, sedangkan fungsi Thomae
      terintegralkan walaupun terputus di setiap bilangan rasional.
\item Fungsi akumulasi $F(x)=\int_a^xf$ selalu kontinu dan turunannya $f$ di setiap
      titik tempat $f$ kontinu, jadi setiap fungsi kontinu punya antiturunan.
\item Bagian kedua menyatakan $\int_a^bF'=F(b)-F(a)$; buktinya teorema nilai rata-rata
      ditambah teleskopik, dengan syarat $F'$ terintegralkan.
\item Substitusi dan integral parsial adalah aturan rantai dan aturan perkalian yang
      dibaca terbalik. Integral tak wajar adalah limit integral Riemann, dan uji
      integral menghubungkannya dengan deret Bab 5.
\end{itemize}
\end{ringkasan}

\section{Latihan}
\tingkat{menengah}

\begin{latihan}
\begin{enumerate}
  \item Hitung $U_n$ dan $L_n$ untuk $f(x)=x$ pada $[0,1]$ dengan partisi seragam,
        nyatakan keduanya sebagai pecahan dalam $n$, lalu tentukan $U_n-L_n$. Simpulkan
        bahwa $f$ terintegralkan dan hitung integralnya.
  \item Buktikan bahwa $f$ terintegralkan pada $[a,b]$ berakibat $|f|$ terintegralkan
        dan $\bigl|\int_a^bf\bigr|\le\int_a^b|f|$. Tunjukkan kebalikannya salah dengan
        satu contoh.
  \item Nilai klaim ini: ``Setiap fungsi terbatas pada $[a,b]$ terintegralkan, sebab
        jumlah atas dan jumlah bawahnya selalu berhingga.'' Benar, salah, atau belum
        cukup bukti? Sebutkan bagian mana yang didukung bukti dan bagian mana yang
        tidak.
  \item Buktikan bahwa fungsi naik pada $[a,b]$ terintegralkan dengan menunjukkan
        $U(\mathcal{P},f)-L(\mathcal{P},f)\le(b-a)\bigl(f(b)-f(a)\bigr)/n$ untuk
        partisi seragam $n$ bagian. Jelaskan mengapa batas itu tidak berlaku untuk
        fungsi yang tidak monoton.
  \item Hitung $\int_0^1(3x^2-2x+1)\,dx$ dengan teorema dasar kalkulus, lalu periksa
        hasilnya memakai aturan Simpson dengan $n=100$ dan fungsi
        \texttt{scipy.integrate.quad}. Laporkan ketiga angka dan selisihnya.
  \item Hitung $\int_1^{10}x^{-2}dx$ dan $\int_1^{\infty}x^{-2}dx$. Laporkan selisih
        keduanya, lalu jelaskan mengapa selisih itu sama dengan $1/10$ tanpa
        menghitung ulang.
  \item Buktikan bahwa bila $f$ kontinu, $f\ge0$ pada $[a,b]$, dan $\int_a^bf=0$, maka
        $f(x)=0$ untuk setiap $x\in[a,b]$. Tunjukkan juga bahwa pernyataan itu gagal
        bila $f$ hanya terintegralkan.
  \item Hitung $\int_0^{\pi/2}\sin x\cos x\,dx$ dengan dua cara: substitusi
        $u=\sin x$, dan identitas $\sin2x=2\sin x\cos x$. Bandingkan hasilnya dengan
        nilai numerik dari \texttt{quad}.
  \item Tunjukkan bahwa fungsi Dirichlet pada $[0,1]$ memenuhi
        $U(\mathcal{P},D)-L(\mathcal{P},D)=1$ untuk setiap partisi, lalu jelaskan
        mengapa fakta itu cukup untuk menyimpulkan ketidakterintegralannya.
  \item Taksir $\sum_{n=101}^{\infty}1/n^2$ dengan batas integral dan bandingkan dengan
        nilai sebenarnya $9.950167\times10^{-3}$; berapa galat relatifnya?
\end{enumerate}

\noindent\textbf{Petunjuk.} Butir 3: klaim memuat dua pernyataan berbeda, tentang
keterbatasan jumlah dan tentang kesamaan kedua integral. Butir 7: gunakan
kekontinuan di satu titik untuk membuat integral kecil pada selang sempit di sekitar
titik itu. Butir 10: pakai $\sum_{n=N+1}^\infty1/n^2\le\int_N^\infty x^{-2}dx=1/N$
dengan $N=100$.
\end{latihan}

\begin{untukpeneliti}
Yang benar-benar terbukti di bab ini: kriteria keterintegralan dengan kedua arahnya,
keterintegralan fungsi kontinu (memakai kekontinuan seragam dari Bab 11), fungsi
monoton, fungsi Thomae, kedua bagian teorema dasar kalkulus, teorema nilai rata-rata
untuk integral, dan status substitusi serta integral parsial sebagai akibat, bukan
aksioma.

Yang dikutip tanpa bukti dan sebaiknya disebut begitu di kelas: kriteria Lebesgue
(Rudin, Teorema 11.33(b); Abbott, Bab 7), contoh Volterra tentang turunan terbatas yang
tidak terintegralkan, dan ketiadaan antiturunan elementer untuk $e^{-x^2}$ (teorema
Liouville).

Sumber primer: Rudin, \emph{Principles of Mathematical Analysis}, Bab 6, bagian
Definition and Existence of the Integral (Teorema 6.4, 6.6, 6.8, 6.9, 6.10) dan bagian
Integration and Differentiation (6.19--6.22). Rudin menyusunnya untuk integral
Riemann-Stieltjes; integral Riemann adalah kasus $\alpha(x)=x$. Abbott,
\emph{Understanding Analysis}, Bab 7, lebih lambat tetapi lebih terbaca; Stewart hanya
untuk aturan hitung.

Pertanyaan terbuka: seberapa tajam batas galat kuadratur bila hanya $|f''|$ yang
diketahui, dan bagaimana memperkirakan galat tanpa turunan sama sekali; jawaban
standarnya kuadratur adaptif dan selisih dua tingkat penghalusan. Batas keberlakuan:
fungsi terbatas pada selang tertutup terbatas, integral tak wajar sebagai limit, tanpa
teori ukur, dan integral bernilai vektor tidak disinggung.
\end{untukpeneliti}

\peta{Bab 13 pada Rudin, Bab 6: bagian Definition and Existence of the Integral
(partisi, jumlah atas-bawah, Lemma penghalusan sebagai Teorema 6.4, kriteria sebagai
6.6, fungsi kontinu dan monoton sebagai 6.8 dan 6.9) dan bagian Integration and
Differentiation (6.19 perubahan variabel, 6.21 teorema dasar kalkulus, 6.22 integral
parsial). Pada Abbott, Bab 7, dari The Definition of the Riemann Integral sampai
Lebesgue's Criterion. Pada Stewart: The Substitution Rule dan Average Value of a
Function (Bab 5), Integration by Parts dan Improper Integrals (Bab 7), serta 11.3
The Integral Test, jembatan ke Bab 5 buku ini.}
\chapter{Deret Fungsi dan Pertukaran Limit}

Sebuah pustaka numerik kecil menghitung $\ln(1+x)$ dengan deret
$\sum_{k\ge 1}(-1)^{k+1}x^{k}/k$. Penulisnya menguji fungsi itu pada $x=0.5$
dengan dua belas suku dan membandingkannya dengan pustaka standar: galatnya
$6.4\times 10^{-6}$. Uji itu lolos, dokumentasi menyebut deretnya akurat sampai
enam angka, dan kodenya masuk produksi. Pemakai pertama memanggilnya dengan
$x=0.99$. Dua belas suku memberi $0.6529204447$, sedangkan
$\ln 1.99=0.6881346387$. Galatnya $3.5\times10^{-2}$, sekitar lima ribu kali
lebih besar, pada fungsi yang sama dengan jumlah suku yang sama.

Kerusakan itu bukan kesalahan ketik. Untuk setiap $x$ di $(-1,1)$ deretnya
konvergen ke $\ln(1+x)$; yang berubah hanya berapa suku yang dibutuhkan: $16$
suku di $x=0.5$, $89$ di $x=0.9$, $721$ di $x=0.99$, $5248$ di $x=0.999$,
semuanya untuk galat di bawah $10^{-6}$. Tidak ada satu bilangan suku yang cukup
untuk semua $x$, sebab $x$ boleh sedekat apa pun ke $1$.

Bandingkan dengan potongan tertutup $[-r,r]$, $0<r<1$. Di sana satu anggaran
berlaku serentak untuk semua $x$: sisa deret berganti tanda dibatasi
$r^{N+1}/(N+1)$, dan batas itu tidak menyebut $x$. Untuk $r=0.5$ dan $N=12$
batasnya $9.4\times10^{-6}$, galat sebenarnya $6.4\times10^{-6}$. Inilah
perbedaan yang seluruh bab ini bicarakan.

Bab 5 mengurus deret bilangan, dan setiap $x$ memberi deret bilangan yang
berbeda. Bab 11 dan 13 mendefinisikan kekontinuan dan integral titik demi titik,
jadi keduanya tidak punya cara memindahkan limit melewati tanda integral tanpa
syarat tambahan. Alat yang belum ada adalah jarak antara dua fungsi; begitu
jarak itu ada, seluruh keluarga deret menjadi satu barisan di ruang metrik.

\begin{untukanak}
Bayangkan sebuah kelas dengan tiga puluh murid, dan guru menulis satu bilangan
di papan: seratus. Tugas setiap murid menebak bilangan itu, dan setelah tiap
tebakan guru mengumumkan apakah tebakan itu makin dekat. Murid pertama cepat
sekali: dalam lima tebakan ia sudah meleset kurang dari satu. Murid kedua
lambat: butuh dua puluh tebakan. Murid ketiga sangat lambat: butuh seratus.
Tetapi ketiganya, kalau diberi waktu cukup, akhirnya meleset kurang dari satu.

Sekarang guru mengganti pertanyaannya. Bukan "apakah setiap murid akhirnya
dekat", melainkan "adakah satu saat ketika semua murid sudah dekat pada saat
yang sama". Untuk kelas tadi jawabannya tidak: murid paling lambat selalu
tertinggal, dan selalu bisa ditambah murid yang lebih lambat lagi.

Pertanyaan yang sama muncul pada rumus yang dipakai komputer. Sebuah rumus
mungkin benar untuk setiap nilai yang dimasukkan, tetapi jumlah suku yang
dibutuhkan bisa membengkak tanpa batas begitu nilainya mendekati ujung. Kalau
program menganggap "benar untuk setiap nilai" sama dengan "satu rumus cukup
untuk semua nilai", program itu akan meleset jauh pada nilai yang paling
merepotkan, dan melesetnya tidak kelihatan pada uji yang dipilih penulisnya.

Di bab ini kita akan tahu mengapa dua pertanyaan tadi bukan soal kata-kata,
melainkan soal apakah dua operasi boleh ditukar urutannya: limit dengan
integral, limit dengan turunan, dan penjumlahan dengan keduanya.
\end{untukanak}

\section{Konvergensi titik demi titik dan konvergensi seragam}
\tingkat{dasar}

Ambil keluarga fungsi $f_1,f_2,\ldots$ yang semuanya terdefinisi pada satu
himpunan $E\subseteq\mathbb{R}$. Untuk satu titik $x$ yang kita pegang tetap,
nilai $f_1(x),f_2(x),\ldots$ hanyalah barisan bilangan, dan Bab 4 sudah
mengatakan apa artinya barisan itu konvergen. Karena itu definisi pertama hampir
menulis dirinya sendiri.

\begin{definisi}
Misalkan $f_n$ dan $f$ terdefinisi pada $E\subseteq\mathbb{R}$.
\begin{enumerate}
\item Kita katakan $f_n\to f$ \emph{titik demi titik}\index{konvergensi titik
      demi titik} pada $E$ bila untuk setiap $x\in E$ dan setiap
      $\varepsilon>0$ terdapat $N$ sehingga $n\ge N$ berakibat
      $|f_n(x)-f(x)|<\varepsilon$. Di sini $N$ boleh bergantung pada $x$ dan
      pada $\varepsilon$.
\item Kita katakan $f_n\to f$ \emph{seragam}\index{konvergensi seragam} pada
      $E$ bila untuk setiap $\varepsilon>0$ terdapat $N$ sehingga $n\ge N$
      berakibat $|f_n(x)-f(x)|<\varepsilon$ untuk \emph{semua} $x\in E$
      sekaligus. Di sini $N$ hanya bergantung pada $\varepsilon$.
\item Bila $h$ terbatas pada $E$, norma supremum\index{norma supremum} $h$
      adalah $\|h\|_\infty=\sup_{x\in E}|h(x)|$.
\end{enumerate}
\end{definisi}

\kamus{konvergensi seragam}{uniform convergence}{semua titik harus sudah dekat
pada saat yang sama, bukan satu per satu menurut jadwalnya masing-masing.}

Satu kuantor berpindah tempat, dan seluruh isi bab ini bergantung pada
perpindahan itu. Pada definisi pertama urutannya "untuk setiap titik, ada
ambang". Pada definisi kedua urutannya "ada satu ambang untuk semua titik".
Kalimat pertama benar untuk deret pada pembuka bab, kalimat kedua salah, dan
selisih keduanya adalah selisih antara $16$ suku dan $5248$ suku.

\begin{bacakeras}
Dua suara bergantian.

Suara pertama: ``Apakah deret ini sampai?''
Suara kedua: ``Sampai. Perlahan, tetapi sampai.''
Suara pertama: ``Sampai untuk titik ini?''
Suara kedua: ``Untuk titik ini. Untuk titik sebelah, lain kali, dengan tempo sendiri.''

Itulah janji yang pertama. Setiap titik boleh lambat, asal ia punya satu tempo
untuk dirinya sendiri. Titik yang jauh dari ujung sampai dalam belasan langkah;
titik yang menempel di ujung butuh ribuan.

Sekarang suara pertama bertanya dengan kata lain: ``Adakah satu tempo untuk
semua titik sekaligus?''

Suara kedua diam lebih lama. ``Di selang tertutup yang berhenti sebelum ujung,
ada. Di seluruh selang terbuka, tidak. Dan yang tidak ada itu bukan kegagalan
hitungan; ia hanya menuntut tempo yang bergerak bersama titiknya.''

Ulangi baris ini sampai terasa di lidah: semua titik boleh lambat, asal satu
tempo. Bila temponya harus berbeda untuk tiap titik, namanya bukan seragam.
\end{bacakeras}

\begin{contoh}
Misalkan $f_n(x)=x^n$ pada $E=[0,1]$. Untuk $0\le x<1$ berlaku $x^n\to0$,
sedangkan $x^n=1$ untuk semua $n$ bila $x=1$. Jadi $f_n$ konvergen titik demi
titik ke
\[
  f(x)=\begin{cases}0, & 0\le x<1,\\ 1, & x=1.\end{cases}
\]
Kekonvergenan itu tidak seragam, sebab $\sup_{x\in[0,1]}|x^n-f(x)|=1$ untuk
setiap $n$: di $x=1$ selisihnya tepat $1$, dan untuk $x$ yang cukup dekat ke $1$
dari kiri $x^n$ mendekati $1$ sementara $f(x)=0$.
\end{contoh}

Contoh itu memperlihatkan dua hal. Pertama, limit titik demi titik fungsi kontinu
tidak harus kontinu: setiap $x^n$ kontinu, limitnya melompat di $x=1$. Kedua,
keseragaman bergantung pada himpunannya, bukan hanya pada rumus fungsinya.

Ketergantungan itu bisa diukur. Pada $[0,r]$ dengan $r<1$ kekonvergenannya
seragam, sebab $\sup_{[0,r]}|x^n|=r^n\to0$. Yang menarik adalah bila batas
kanannya ikut bergerak bersama $n$. Pada $[0,1-1/n]$ supremumnya
$(1-1/n)^n\to1/e=0.3678794412$; jadi menyempit sebesar $1/n$ belum
menyelamatkan keseragaman. Pada $[0,1-1/\sqrt{n}]$ supremumnya
$(1-1/\sqrt n)^n$, yang runtuh seperti $e^{-\sqrt n}$: $2.2\times10^{-2}$ pada
$n=10$, $2.7\times10^{-5}$ pada $n=100$, $2.2\times10^{-44}$ pada $n=10^4$. Laju
penyempitan, bukan rumus fungsinya, yang menentukan.

Uji diagonal memberi keterangan tambahan yang lebih tajam. Ambil titik
$x_n=1-1/n$ yang mengejar titik putus. Maka $f_n(x_n)=(1-1/n)^n\to1/e$,
sedangkan $f(x_n)=0$ untuk semua $n$. Jadi bukan hanya supremumnya yang tidak
mengecil: nilai di titik-titik yang dipilih dengan tepat pun menolak menuju
limit titik demi titiknya. Pola ini dipakai berulang kali di bab ini. Untuk
menyangkal keseragaman, cari satu barisan titik yang membuat selisihnya tetap
besar.

Dua keluarga sederhana memisahkan kedua gejala itu. Yang pertama, $g_n$ adalah
segitiga bertinggi $1$ dengan alas $1/n$ dan puncak di $1-1/n$: puncaknya
bergerak ke kanan, tingginya tidak pernah turun, dan luasnya $1/(2n)$. Titik
$x<1$ hanya tersentuh puncak selama $1/(2(1-x))<n<3/(2(1-x))$, jadi di $x=0.9$
hanya sembilan nilai $n$ memberi nilai positif, di $x=0.99$ sembilan puluh
sembilan, dan setelah puncak terakhir lewat nilainya nol selamanya. Limitnya nol
di setiap titik, sedangkan $\|g_n\|_\infty=1$. Yang kedua, $h_n$ bertinggi $1/n$
dengan alas $1$ di tempat yang sama: puncaknya mengecil, bukan berpindah, dan
$\|h_n\|_\infty=1/n\to0$ (Gambar~\ref{fig:puncak}).

\begin{figure}[htbp]
\centering
\begin{tikzpicture}[x=5.6cm,y=3.1cm,font=\small]
  \draw[->,layergrey] (0,0)--(1.15,0) node[right,font=\scriptsize]{$x$};
  \draw[->,layergrey] (0,0)--(0,1.32) node[above,font=\scriptsize]{$y$};
  \node[below,font=\scriptsize,layergrey] at (1,0) {$1$};
  \foreach \n/\c in {16/layerindigo!15,8/layerindigo!30,4/layerindigo!50,2/layerindigo!75}{
    \fill[\c] ({1-1/\n-1/(2*\n)},0) -- ({1-1/\n},1) -- ({1-1/\n+1/(2*\n)},0) -- cycle;
  }
  \draw[dashed,layergrey] (0,1)--(1.06,1);
  \node[above,font=\scriptsize,layergrey] at (0.72,1) {tinggi $1$};
  \begin{scope}[xshift=7.0cm]
    \draw[->,layergrey] (0,0)--(1.15,0) node[right,font=\scriptsize]{$x$};
    \draw[->,layergrey] (0,0)--(0,1.32) node[above,font=\scriptsize]{$y$};
    \node[below,font=\scriptsize,layergrey] at (1,0) {$1$};
    \foreach \n/\c in {8/layerteal!18,4/layerteal!32,2/layerteal!52,1/layerteal!72}{
      \fill[\c] (0,0) -- (0.5,{1/\n}) -- (1,0) -- cycle;
    }
    \node[right,font=\scriptsize,layergrey] at (0.52,0.5) {$n=2$};
    \node[right,font=\scriptsize,layergrey] at (0.52,0.25) {$n=4$};
    \node[right,font=\scriptsize,layergrey] at (0.52,0.125) {$n=8$};
  \end{scope}
\end{tikzpicture}
\caption{Puncak yang bergerak dan puncak yang mengecil. Kiri: $g_n$ bertinggi
tetap $1$ dengan alas $1/n$ dan puncak di $1-1/n$; tingginya tidak turun,
sehingga $\|g_n\|_\infty=1$ walaupun $g_n(x)\to0$ di setiap titik. Kanan: $h_n$
bertinggi $1/n$ dengan alas $1$ di tempat yang sama; seluruh grafiknya masuk
pita setinggi $1/n$, sehingga konvergen seragam.}
\label{fig:puncak}
\end{figure}

\begin{lemma}
$f_n\to f$ seragam pada $E$ jika dan hanya jika $\|f_n-f\|_\infty\to0$.
\end{lemma}

\begin{proof}
Kekonvergenan seragam berarti: untuk setiap $\varepsilon>0$ ada $N$ sehingga
untuk $n\ge N$ berlaku $|f_n(x)-f(x)|\le\varepsilon$ pada semua $x\in E$,
yaitu $\sup_x|f_n(x)-f(x)|\le\varepsilon$. Itu persis pernyataan bahwa barisan
bilangan real $\|f_n-f\|_\infty$ menuju nol. Satu-satunya selisih teknis adalah
bahwa definisi memakai $<\varepsilon$ sedangkan norma memakai $\le\varepsilon$;
keduanya setara karena $\varepsilon$ sebarang.
\end{proof}

Lemma itu mengubah pertanyaan tentang fungsi menjadi pertanyaan tentang satu
barisan bilangan: memeriksa keseragaman cukup dengan satu supremum, menolaknya
cukup dengan menunjukkan supremum itu tidak menuju nol, seperti pada $x^n$.

Keseragaman punya gambar yang tepat, yaitu pita. Ambil $\varepsilon=0.1$. Bila
seluruh grafik $f_n$ untuk $n\ge N$ termuat di antara $f-\varepsilon$ dan
$f+\varepsilon$, maka satu $N$ itu bekerja untuk semua titik. Yang diperiksa
bukan satu titik, melainkan apakah ada bagian grafik yang keluar dari pita
selebar $0.2$. Pada $g_n$ pita itu selalu ditembus puncaknya; keadaan yang
diharapkan terlihat di Gambar~\ref{fig:pita}.

\begin{figure}[htbp]
\centering
\begin{tikzpicture}[x=2.55cm,y=4.0cm,font=\small]
  \draw[->,layergrey] (0,0)--(4.5,0) node[right,font=\scriptsize]{$x$};
  \draw[->,layergrey] (0,0)--(0,1.18) node[above,font=\scriptsize]{$y$};
  \fill[layergrey!12] plot[domain=0:4,samples=90] (\x,{0.55+0.22*sin(2*\x r)+0.1})
    -- plot[domain=4:0,samples=90] (\x,{0.55+0.22*sin(2*\x r)-0.1}) -- cycle;
  \draw[dashed,layergrey] plot[domain=0:4,samples=90] (\x,{0.55+0.22*sin(2*\x r)+0.1});
  \draw[dashed,layergrey] plot[domain=0:4,samples=90] (\x,{0.55+0.22*sin(2*\x r)-0.1});
  \draw[thick,accent] plot[domain=0:4,samples=90] (\x,{0.55+0.22*sin(2*\x r)});
  \draw[thick,layerindigo] plot[domain=0:4,samples=180]
    (\x,{0.55+0.22*sin(2*\x r)+0.06*cos(9*\x r)});
  \draw[<->,layergrey] (3.78,{0.55+0.22*sin(2*3.78 r)-0.1})
    -- (3.78,{0.55+0.22*sin(2*3.78 r)+0.1});
  \node[right,font=\scriptsize,layergrey] at (3.8,{0.55+0.22*sin(2*3.78 r)}) {$0.2$};
  \node[accent,font=\scriptsize,above left] at (0.42,{0.55+0.22*sin(2*0.42 r)}) {$f$};
  \node[layerindigo,font=\scriptsize,below right] at (0.95,{0.55+0.22*sin(2*0.95 r)+0.06*cos(9*0.95 r)}) {$f_n$};
\end{tikzpicture}
\caption{Pita seragam selebar $0.2$ di sekeliling $f$, yaitu di antara $f-0.1$
dan $f+0.1$. Grafik $f_n$ yang bergelombang masih termuat seluruhnya di dalam
pita, jadi $\|f_n-f\|_\infty\le0.1$ dan satu $N$ yang sama bekerja untuk semua
titik sekaligus.}
\label{fig:pita}
\end{figure}

\begin{tangan}
Ambil tali sepanjang dua meter, selotip, dan dua orang. Bentangkan tali lurus di
lantai, lalu tempel lima tanda berjarak sama. Orang pertama memegang ujung kiri
dan mengangkatnya $5$ cm saja. Lihat: hanya ujung itu yang naik, empat tanda
lain masih menempel di lantai. Itu kekonvergenan titik demi titik.

Sekarang kedua orang mengangkat kedua ujung bersama-sama dengan hitungan yang
sama: $5$ cm, lalu $3$ cm, lalu $1$ cm, lalu setengah sentimeter. Amati seluruh
panjang tali. Tidak ada titik yang tertinggal di lantai, dan jarak terjauh tali
ke lantai sama di semua titik, sebab keduanya diangkat serentak. Itu
kekonvergenan seragam.

Terakhir, orang pertama mengangkat ujungnya sambil berjalan ke tengah, sementara
orang kedua menurunkan ujungnya. Talinya tetap miring, dan selalu ada bagian
yang $5$ cm di atas lantai: setiap titik akhirnya turun, tetapi jarak terjauh
tidak pernah mengecil. Itulah $g_n$ pada Gambar~\ref{fig:puncak}.
\end{tangan}

\begin{lemma}[Kriteria Cauchy seragam]
Barisan $f_n$ konvergen seragam pada $E$ jika dan hanya jika untuk setiap
$\varepsilon>0$ ada $N$ sehingga $n,m\ge N$ berakibat
$|f_n(x)-f_m(x)|<\varepsilon$ untuk semua $x\in E$. Syarat terakhir ini disebut
kriteria Cauchy seragam\index{kriteria Cauchy seragam}.
\end{lemma}

\begin{proof}
$(\Rightarrow)$ Bila $f_n\to f$ seragam, ambil $N$ dari definisi kekonvergenan
untuk $\varepsilon/2$. Untuk $n,m\ge N$ dan setiap $x$,
\[
  |f_n(x)-f_m(x)|\le|f_n(x)-f(x)|+|f(x)-f_m(x)|<\varepsilon .
\]
$(\Leftarrow)$ Untuk setiap $x$ tetap, barisan $f_n(x)$ adalah barisan Cauchy
di $\mathbb{R}$, jadi konvergen ke suatu bilangan yang kita sebut $f(x)$. Ini
mendefinisikan sebuah fungsi $f$ titik demi titik. Untuk memperoleh
keseragamannya, tetapkan $n\ge N$ pada kondisi Cauchy dan biarkan
$m\to\infty$. Karena $f_m(x)\to f(x)$, diperoleh $|f_n(x)-f(x)|\le\varepsilon$
untuk semua $x\in E$ sekaligus, dan $N$ yang dipakai tidak bergantung pada $x$.
\end{proof}

Bukti arah kedua hanya memakai kelengkapan $\mathbb{R}$: ruang fungsi lengkap
terhadap norma supremum, dan kelengkapan itulah yang membuat deret fungsi bisa
dijumlahkan tanpa takut limitnya keluar dari kelas fungsi yang kita pakai.

\begin{catatan}
Kekonvergenan seragam adalah sifat pasangan: barisan dan himpunannya. Kalimat
"$x^n$ konvergen seragam" tidak lengkap tanpa menyebut pada himpunan mana. Pada
$[0,1]$ jawabannya tidak, pada $[0,0.99]$ ya: norma supremum galatnya
$0.99^n$, yaitu $0.37$ pada $n=100$ dan $4.3\times10^{-5}$ pada $n=1000$.
Karena itu setiap klaim keseragaman di bab ini menyebut himpunannya.
\end{catatan}

\section{Uji Weierstrass: keseragaman dari satu deret bilangan}
\tingkat{menengah}

Deret fungsi adalah deret yang sukunya fungsi. Jumlah parsialnya
$S_N=\sum_{n=1}^{N}f_n$ adalah fungsi, jadi deret fungsi tidak lain adalah
barisan fungsi yang khusus, dan semua definisi bagian sebelumnya berlaku begitu
kita mengganti $f_n$ dengan $S_N$.

\begin{definisi}
Deret $\sum_{n\ge1}f_n$ \emph{konvergen titik demi titik} pada $E$ bila untuk
setiap $x\in E$ deret bilangan $\sum_n f_n(x)$ konvergen. Deret itu
\emph{konvergen seragam}\index{deret fungsi} pada $E$ bila barisan jumlah
parsialnya konvergen seragam pada $E$.
\end{definisi}

Uji berikut adalah cara termurah memperoleh keseragaman, dan hampir semua
contoh di buku ini memakainya. Idenya: bila setiap suku dibatasi oleh bilangan
yang deretnya konvergen, maka ekor deret fungsi dibatasi oleh ekor deret
bilangan itu, dan ekor bilangan tidak tahu-menahu soal $x$.

\begin{teorema}[Uji-$M$ Weierstrass]\index{uji Weierstrass}
Misalkan $f_n$ terdefinisi pada $E$ dan ada bilangan $M_n\ge0$ dengan
$|f_n(x)|\le M_n$ untuk semua $x\in E$ dan semua $n$. Jika $\sum M_n$
konvergen, maka $\sum f_n$ konvergen seragam pada $E$. Bila $E$ adalah selang
tertutup terbatas dan setiap $f_n$ kontinu, maka jumlahnya kontinu pada $E$.
\end{teorema}

\begin{proof}
Tulis $S_N=\sum_{i=1}^{N}f_i$. Untuk $n>m$ dan setiap $x\in E$,
\[
  |S_n(x)-S_m(x)|\le\sum_{i=m+1}^{n}|f_i(x)|\le\sum_{i=m+1}^{n}M_i .
\]
Karena $\sum M_n$ konvergen, ekor $\sum_{i=m+1}^{n}M_i$ dapat dibuat lebih
kecil dari $\varepsilon$ asalkan $m,n$ cukup besar, dan pilihan itu tidak
menyebut $x$. Kriteria Cauchy seragam menutup bukti keseragamannya. Untuk
kekontinuan, setiap $S_N$ kontinu sebagai jumlah berhingga fungsi kontinu, dan
limit seragam fungsi kontinu kontinu; pernyataan terakhir itu dibuktikan di
bagian berikut.
\end{proof}

\begin{contoh}
Pada $E=[-1,1]$ deret $\sum_{n\ge1}x^n/n^2$ konvergen seragam, sebab
$|x^n/n^2|\le1/n^2$ dan $\sum1/n^2=\pi^2/6$ konvergen. Jumlahnya kontinu pada
$[-1,1]$, dan nilainya di $x=1$ adalah $1.6449340668$. Jumlah dua juta suku
memberi $1.6449335668$; selisihnya $5.0\times10^{-7}$, sebanding dengan batas
ekor $1/N=5\times10^{-7}$. Di $x=-1$ deretnya konvergen ke $-\pi^2/6$, sebab
suku-sukunya berganti tanda dengan nilai mutlak yang menurun menuju nol.
\end{contoh}

\begin{contoh}
Pada $E=\mathbb{R}$ deret $\sum_{n\ge0}2^{-n}\sin(2^n x)$ konvergen seragam,
sebab $|2^{-n}\sin(2^n x)|\le2^{-n}$ dan $\sum2^{-n}=1$. Setiap suku kontinu,
jadi jumlahnya kontinu di setiap titik. Yang mengejutkan, jumlah itu tidak
terdiferensialkan di titik mana pun: hasil klasik Weierstrass (1872), diperluas
Hardy sampai nisbah $1$, kita kutip tanpa bukti. Keseragaman menjamin
kekontinuan, tetapi tidak menjamin apa pun tentang turunan.
\end{contoh}

\begin{catatan}
Uji-$M$ memberi syarat cukup, bukan syarat perlu, dan kegagalannya bukan bukti
apa-apa. Deret $\sum_{n\ge1}(-1)^n(1-x)x^n$ pada $[0,1]$ konvergen seragam:
sisa setelah suku ke-$N$ dibatasi $(1-x)x^{N+1}$, yang maksimumnya pada $[0,1]$
adalah $\tfrac{1}{N+2}\left(\tfrac{N+1}{N+2}\right)^{N+1}\to0$, yaitu
$3.20\times10^{-2}$ pada $N=10$ dan $3.67\times10^{-4}$ pada $N=1000$. Tetapi
$\sup_{[0,1]}|f_n|=n^n/(n+1)^{n+1}$, sebanding dengan $1/(en)$, sehingga
$\sum_n\sup|f_n|$ divergen: $0.88$ sampai $N=10$ dan $2.53$ sampai $N=1000$.
Tidak ada deret pembanding yang konvergen, jadi uji-$M$ tidak terpakai walaupun
keseragamannya benar.
\end{catatan}

\begin{untukanak}
Ada dua cara mengerjakan sederet tugas. Cara pertama: selesaikan semua tugas
satu per satu, baru kemudian hitung totalnya. Cara kedua: hitung totalnya sambil
mengerjakan, lalu selesaikan sisanya. Untuk penjumlahan biasa kedua cara memberi
hasil yang sama, dan kita terbiasa menukarnya tanpa berpikir.

Tetapi ada tugas yang tidak suka ditukar. Bayangkan kamu mengukur panjang
bayangan sebuah tiang setiap jam selama satu hari, menghitung rata-ratanya, lalu
mencari jam ketika rata-rata itu paling panjang. Balik urutannya: cari dulu jam
terpanjang pada tiap hari, lalu rata-ratakan jam itu sepanjang bulan. Dua
hasilnya bisa berbeda jauh, bukan karena hitungannya salah, tetapi karena
"paling panjang" dan "rata-rata" memang tidak saling bertukar.

Hal yang sama terjadi pada rumus yang dipakai untuk menjumlahkan tak hingga
banyak suku, lalu mengintegralkan atau menurunkannya. Kadang urutannya boleh
ditukar dan hasilnya sama persis; kadang tidak, dan selisihnya besar. Di bab ini
kita akan tahu tanda bahayanya sebelum hasilnya meleset.
\end{untukanak}

\section{Kapan limit boleh ditukar}
\tingkat{menengah}

Tiga operasi yang biasa kita pakai pada fungsi adalah menghitung limit di
sebuah titik, mengintegralkan, dan menurunkan. Ketiganya sendiri adalah limit:
kekontinuan di $x_0$ berarti $\lim_{x\to x_0}f(x)=f(x_0)$, integral adalah
limit jumlah Riemann, turunan adalah limit hasil bagi selisih. Jadi memindahkan
limit barisan melewati salah satu operasi itu berarti menukar dua limit, dan
pertukaran itu tidak gratis: harganya satu hipotesis, yaitu
keseragaman\index{pertukaran limit}.

\begin{teorema}
Bila $f_n$ kontinu pada $E$ dan $f_n\to f$ seragam pada $E$, maka $f$ kontinu
pada $E$.
\end{teorema}

\begin{proof}
Ambil $x_0\in E$ dan $\varepsilon>0$. Karena kekonvergenannya seragam, ada $n$
dengan $\|f_n-f\|_\infty<\varepsilon/3$. Karena $f_n$ kontinu di $x_0$, ada
$\delta>0$ sehingga $|x-x_0|<\delta$ berakibat $|f_n(x)-f_n(x_0)|<\varepsilon/3$.
Maka untuk $x\in E$ dengan $|x-x_0|<\delta$,
\[
  |f(x)-f(x_0)|\le|f(x)-f_n(x)|+|f_n(x)-f_n(x_0)|+|f_n(x_0)-f(x_0)|
  <\varepsilon .
\]
\end{proof}

Perhatikan urutan logisnya: $n$ dipilih lebih dulu, baru $\delta$, dan
$\delta$ boleh bergantung pada $n$. Yang tidak boleh adalah memilih $\delta$
yang harus bekerja untuk semua $n$ sekaligus.

\begin{tantangan}
Deret $\sum_{n\ge0}\bigl(x^{2n}-x^{2n+1}\bigr)$ bersuku kontinu pada $[0,1]$
dan konvergen di setiap titiknya. Tentukan jumlahnya, lalu tunjukkan jumlah itu
tidak kontinu di $x=1$. Mengapa hasil ini tidak membantah teorema di atas?
Jawaban yang benar harus menyebut satu bilangan yang boleh bergantung pada
$\varepsilon$ tetapi tidak boleh bergantung pada $x$.

Petunjuk 1: jumlah parsialnya punya bentuk tertutup, dan bentuk itu memuat
$x^{2N+2}$.

Petunjuk 2: untuk $x<1$ jumlahnya $1/(1+x)$, sedangkan di $x=1$ setiap suku
bernilai nol; jadi nilainya melompat dari $1/2$ ke $0$.

Kalau mentok, kembali ke bagian uji diagonal dan tanyakan apa yang terjadi pada
$x_N=1-1/N$.
\end{tantangan}

\begin{teorema}[Integrasi suku demi suku]
Bila $f_n$ terintegralkan Riemann pada $[a,b]$ dan $f_n\to f$ seragam pada
$[a,b]$, maka $f$ terintegralkan pada $[a,b]$ dan
\[
  \int_a^b f=\lim_{n\to\infty}\int_a^b f_n .
\]
\end{teorema}

\begin{proof}
Untuk keterintegralannya, ambil $\varepsilon>0$ dan pilih $n$ dengan
$\|f_n-f\|_\infty<\varepsilon$. Karena $f_n$ terintegralkan, ada partisi
$\mathcal{P}$ dengan $U(\mathcal{P},f_n)-L(\mathcal{P},f_n)<\varepsilon(b-a)$.
Pada setiap subselang, $M_i(f)\le M_i(f_n)+\varepsilon$ dan
$m_i(f)\ge m_i(f_n)-\varepsilon$, sehingga
$U(\mathcal{P},f)\le U(\mathcal{P},f_n)+\varepsilon(b-a)$ dan
$L(\mathcal{P},f)\ge L(\mathcal{P},f_n)-\varepsilon(b-a)$. Ketiga ketaksamaan
itu memberi $U(\mathcal{P},f)-L(\mathcal{P},f)<3\varepsilon(b-a)$, dan kriteria
keterintegralan Bab 13 menutup argumennya. Untuk nilai integralnya, cukup
dicatat bahwa
\[
  \left|\int_a^b f_n-\int_a^b f\right|
  =\left|\int_a^b (f_n-f)\right|
  \le\int_a^b|f_n-f|\le(b-a)\,\|f_n-f\|_\infty ,
\]
dan ruas kanan menuju nol.
\end{proof}

\begin{akibat}[Integrasi deret]
Bila $\sum_{n\ge1}f_n$ konvergen seragam pada $[a,b]$ dan setiap $f_n$
terintegralkan di sana, maka
\[
  \int_a^b\sum_{n=1}^{\infty}f_n=\sum_{n=1}^{\infty}\int_a^b f_n .
\]
\end{akibat}

Akibat ini diterapkan pada jumlah parsial, dengan keseragaman biasanya
diperiksa lewat uji-$M$. Contohnya, deret $\sum_{n\ge1}x^n/n^2$ konvergen seragam
pada $[0,1]$, sehingga
\[
  \int_0^1\sum_{n=1}^{\infty}\frac{x^n}{n^2}\,dx
  =\sum_{n=1}^{\infty}\frac{1}{n^2(n+1)}=\frac{\pi^2}{6}-1=0.6449340668 .
\]
Angka terakhir juga diperoleh dengan menjumlahkan dua juta suku pertama deret di
ruas kanan, yang memberi $0.6449340668$ pada presisi mesin.

\begin{contoh}[Kegagalan tanpa keseragaman]
Misalkan $f_n(x)=nx(1-x^2)^n$ pada $[0,1]$. Untuk setiap $x\in(0,1]$ faktor
$(1-x^2)^n$ mengecil secara geometri, jadi $n(1-x^2)^n\to0$ dan $f_n(x)\to0$;
di $x=0$ semua $f_n$ bernilai $0$. Jadi limit titik demi titiknya fungsi nol.
Integralnya dapat dihitung tepat dengan substitusi $u=1-x^2$:
\[
  \int_0^1 nx(1-x^2)^n\,dx=\frac{n}{2}\int_0^1 u^n\,du=\frac{n}{2(n+1)} .
\]
Barisan itu menuju $1/2$, bukan $0$. Kuadratur numerik memberi
$0.4545454545$ untuk $n=10$ dan $0.4950495050$ untuk $n=100$, sama dengan
$n/(2(n+1))$ sampai digit terakhir. Jadi $\lim_n\int_0^1f_n=1/2$ sedangkan
$\int_0^1\lim_nf_n=0$.
\end{contoh}

Kegagalan itu punya bentuk geometris. Fungsi $f_n$ bernilai nol di kedua ujung
dan memuncak di $x=1/\sqrt{2n+1}$, dengan tinggi $1.3396685$ pada $n=10$ dan
$4.2834707$ pada $n=100$: puncaknya makin tinggi dan makin sempit, sehingga
luasnya tetap sekitar $1/2$ sementara tingginya di setiap titik menuju nol.

\begin{teorema}[Turunan suku demi suku]
Misalkan $f_n$ terdiferensialkan pada $[a,b]$, barisan $f_n'$ konvergen seragam
pada $[a,b]$, dan ada satu titik $x_0\in[a,b]$ tempat $f_n(x_0)$ konvergen.
Maka $f_n$ konvergen seragam ke suatu $f$ yang terdiferensialkan pada $[a,b]$,
dan
\[
  f'(x)=\lim_{n\to\infty}f_n'(x)
\]
untuk setiap $x\in[a,b]$.
\end{teorema}

\begin{proof}[Rangka bukti]
Intinya teorema nilai rata-rata pada selisih dua suku: untuk $m,n$ dan setiap
$x,y\in[a,b]$,
\[
  \bigl|(f_n-f_m)(x)-(f_n-f_m)(y)\bigr|\le\|f_n'-f_m'\|_\infty\,|x-y| .
\]
Karena $f_n'$ Cauchy seragam, ruas kanan dapat dibuat kecil serentak untuk semua
pasangan $x,y$ begitu $m,n$ besar. Bersama kekonvergenan di satu titik $x_0$,
ketaksamaan itu membuat $f_n$ Cauchy seragam, jadi $f_n\to f$ seragam. Untuk
turunan, bagi dengan $|x-y|$, biarkan $y\to x$, lalu tukar dua limit yang
muncul; ketaksamaan yang sama menjamin pertukaran itu sah. Rincian teknisnya ada
pada Rudin, Teorema 7.17.
\end{proof}

Syarat "ada satu titik tempat $f_n(x_0)$ konvergen" tidak dapat dibuang: $f_n$
dan $f_n+1$ punya barisan turunan yang sama, tetapi hanya satu di antaranya yang
konvergen di $x_0$. Sebaliknya, keseragaman $f_n$ saja juga tidak cukup untuk
menyimpulkan apa pun tentang turunan.

\begin{contoh}[Turunan yang tidak konvergen]
Misalkan $f_n(x)=\sin(nx)/n$ pada $[0,\pi]$. Karena $\|f_n\|_\infty=1/n$,
barisan itu konvergen seragam ke fungsi nol; nilainya $0.1$ pada $n=10$ dan
$0.01$ pada $n=100$. Tetapi $f_n'(x)=\cos(nx)$: di $x=0$ barisan itu bernilai
$1$ untuk semua $n$, sedangkan di $x=\pi/2$ nilainya $0,-1,0,1,0,-1,\ldots$
menurut paritas $n$, yang tidak konvergen. Jadi $f_n'$ tidak konvergen titik
demi titik di mana pun kecuali titik-titik khusus, apalagi seragam. Turunan
tidak memaafkan osilasi kecil berfrekuensi tinggi: $f_n$ mengecil menuju nol,
tetapi kemiringannya tidak.
\end{contoh}

Gambar~\ref{fig:turunan} memuat tiga kurva turunan itu pada $[0,\pi]$. Di
$x=\pi/2$ ketiganya bernilai $1$, $0$, dan $-1$: tiga nilai berbeda dari satu
barisan di satu titik, dan itu sudah cukup untuk mengatakan limitnya tidak ada.

\begin{figure}[htbp]
\centering
\begin{tikzpicture}[x=3.4cm,y=1.5cm,font=\small]
  \draw[->,layergrey] (0,0)--(3.55,0) node[right,font=\scriptsize]{$x$};
  \draw[->,layergrey] (0,-1.28)--(0,1.28) node[above,font=\scriptsize]{$y$};
  \draw[dashed,layergrey] (1.571,-1.25)--(1.571,1.25);
  \node[below,font=\scriptsize,layergrey] at (1.571,-1.25) {$\pi/2$};
  \draw[thick,layerindigo] plot[domain=0:3.1416,samples=160] (\x,{cos(4*\x r)});
  \draw[thick,layerteal] plot[domain=0:3.1416,samples=160] (\x,{cos(5*\x r)});
  \draw[thick,layerorange] plot[domain=0:3.1416,samples=160] (\x,{cos(6*\x r)});
  \fill[ink] (1.571,1) circle (1.5pt);
  \fill[ink] (1.571,0) circle (1.5pt) node[right,font=\scriptsize] {$0$};
  \fill[ink] (1.571,-1) circle (1.5pt) node[right,font=\scriptsize] {$-1$};
  \node[right,font=\scriptsize] at (1.585,1.06) {$1$};
\end{tikzpicture}
\caption{Turunan $f_n'=\cos(nx)$ pada $[0,\pi]$ untuk $n=4$ (indigo), $n=5$
(teal), dan $n=6$ (oranye). Di $x=\pi/2$ nilainya $1$, $0$, dan $-1$: tiga nilai
berbeda dari satu barisan di satu titik, jadi limit turunannya tidak ada,
walaupun $f_n=\sin(nx)/n$ konvergen seragam ke nol dengan amplitudo $1/4$,
$1/5$, dan $1/6$.}
\label{fig:turunan}
\end{figure}

\begin{catatan}
Dalam pembelajaran mesin, fungsi risiko sering ditulis sebagai harapan atas
distribusi data, $\mathbb{E}_{(x,y)\sim D}[\ell(f(x),y)]$, lalu diturunkan
terhadap parameter model. Bila distribusinya diskret berhingga, seperti pada
himpunan latih, harapan itu jumlah berhingga dan pertukaran turunan dengan
penjumlahan selalu sah tanpa syarat. Bila kontinu, harapan itu integral, dan
memindahkan turunan ke dalamnya menuntut syarat dominasi yang lebih halus
daripada keseragaman; alatnya adalah teorema kekonvergenan dominan dari teori
ukur, yang tidak dibahas di sini. "Turunkan di dalam integral" adalah operasi
berizin, bukan operasi yang selalu benar.
\end{catatan}

\section{Deret pangkat dan jari-jari konvergensi}
\tingkat{lanjut}

Deret pangkat adalah deret fungsi dengan bentuk paling kaku: suku ke-$n$ adalah
koefisien tetap dikali $x^n$. Kekakuan itu berbuah. Himpunan tempat deretnya
konvergen selalu sebuah selang yang simetris terhadap pusatnya, sehingga seluruh
informasi tentang di mana deret itu berguna termuat dalam satu bilangan.

\begin{definisi}
Deret berbentuk $\sum_{n\ge0}a_nx^n$ dengan $a_n\in\mathbb{R}$ disebut deret
pangkat\index{deret pangkat} berpusat di $0$. Bilangan
\[
  R=\sup\Bigl\{r\ge0:\ \sum_{n\ge0}a_nr^n\ \text{konvergen}\Bigr\}
\]
disebut jari-jari konvergensi\index{jari-jari konvergensi} deret itu.
\end{definisi}

\kamus{deret pangkat}{power series}{deret yang sukunya koefisien tetap dikali
pangkat naik dari satu variabel, seperti resep yang sama untuk semua nilai.}

\kamus{jari-jari konvergensi}{radius of convergence}{jarak terjauh dari pusat
tempat resep itu masih menghasilkan bilangan berhingga.}

\begin{teorema}
Misalkan $\sum a_nx^n$ deret pangkat dengan jari-jari $R$. Maka deret itu
konvergen mutlak untuk setiap $|x|<R$ dan divergen untuk setiap $|x|>R$.
Selain itu
\[
  \frac{1}{R}=\limsup_{n\to\infty}|a_n|^{1/n},
\]
dengan kesepakatan $1/0=\infty$ dan $1/\infty=0$.
\end{teorema}

\begin{proof}
Tetapkan $x\ne0$ dan tulis $|a_nx^n|^{1/n}=|x|\cdot|a_n|^{1/n}$. Menurut uji
akar Bab 5, deret $\sum a_nx^n$ konvergen mutlak bila
$\limsup_n|a_nx^n|^{1/n}<1$, yaitu bila $|x|<1/\limsup_n|a_n|^{1/n}$, dan
divergen bila limit superior itu lebih besar dari $1$, sebab sukunya tidak
menuju nol. Karena $x$ sebarang, batas antara kedua rejim itu tepat
$R=1/\limsup_n|a_n|^{1/n}$, dan himpunan kekonvergenannya termuat di $[-R,R]$.
\end{proof}

Pada $|x|=R$ uji akar tidak memberi keputusan, dan memang tidak ada keputusan
umum: $\sum x^n/n^2$ konvergen di kedua ujung, $\sum x^n/n$ hanya di $x=-1$, dan
$\sum x^n$ di keduanya divergen.

\begin{akibat}
Bila $a_n\ne0$ untuk $n$ cukup besar dan $|a_{n+1}/a_n|\to L$, maka $R=1/L$
dengan kesepakatan yang sama seperti di atas.
\end{akibat}

Alasannya adalah rantai ketaksamaan
\[
  \liminf|a_{n+1}/a_n|\le\liminf|a_n|^{1/n}\le\limsup|a_n|^{1/n}
  \le\limsup|a_{n+1}/a_n|,
\]
yang dibuktikan sebagai latihan di Bab 5. Bila nisbah punya limit, keempat
bilangan itu sama, jadi rumus nisbah dan rumus akar memberi $R$ yang sama. Dalam
hitungan tangan, rumus nisbah hampir selalu lebih cepat.

Sekarang bagian yang membuat deret pangkat berguna: ia boleh diturunkan dan
diintegralkan suku demi suku, dan hasilnya berlaku di seluruh selang
kekonvergenannya, bukan hanya di potongan kecilnya.

\begin{teorema}[Turunan dan integral suku demi suku]
Misalkan $\sum a_nx^n$ berjari-jari $R>0$. Maka untuk setiap $x$ dengan
$|x|<R$,
\[
  \frac{d}{dx}\sum_{n\ge0}a_nx^n=\sum_{n\ge1}n\,a_nx^{n-1},
  \qquad
  \int_0^x\sum_{n\ge0}a_nt^n\,dt=\sum_{n\ge0}\frac{a_n}{n+1}x^{n+1},
\]
dan kedua deret baru itu juga berjari-jari $R$.
\end{teorema}

\begin{proof}
Ambil $r$ dengan $0<r<R$ dan pilih $\rho$ dengan $r<\rho<R$. Deret
$\sum|a_n|\rho^n$ konvergen, jadi ada $C$ dengan $|a_n|\rho^n\le C$ untuk semua
$n$. Untuk $|x|\le r$,
\[
  |n\,a_nx^{n-1}|\le n\,\frac{C}{\rho}\left(\frac{r}{\rho}\right)^{n-1},
\]
dan deret $\sum_n n(r/\rho)^{n-1}$ konvergen sebab $r/\rho<1$. Uji-$M$ karena
itu memberi keseragaman deret turunan pada $[-r,r]$, dan teorema turunan suku
demi suku memberi kesamaan di sana. Setiap titik $|x|<R$ termuat pada suatu
$[-r,r]$ dengan $r<R$, jadi kesamaannya berlaku di seluruh $(-R,R)$. Untuk
integral, jalan yang sama dipakai pada $\sum|a_n|r^{n-1}$, yang juga konvergen.
Jari-jari tidak berubah karena koefisien deret turunan adalah $na_n$, dan
\[
  \limsup_n(n|a_n|)^{1/n}
  =\limsup_n|a_n|^{1/n}\cdot\lim_n n^{1/n}
  =\limsup_n|a_n|^{1/n},
\]
sebab $n^{1/n}\to1$. Argumen yang sama berlaku untuk koefisien $a_n/(n+1)$,
karena $(n+1)^{-1/n}\to1$.
\end{proof}

Bukti itu memakai seluruh isi bab: uji-$M$ untuk keseragaman, teorema turunan
untuk memindahkan limit, dan rumus jari-jari untuk memastikan deret baru tidak
menyempit. Keseragaman hanya diklaim pada $[-r,r]$, $r<R$; pada contoh pembuka,
$\ln(1+x)$ seragam pada setiap $[-r,r]$ dengan $r<1$ tetapi tidak pada $(-1,1)$.

\begin{contoh}
Untuk $|x|<1$, deret geometri memberi $\sum_{n\ge0}x^n=1/(1-x)$. Menurunkan
suku demi suku,
\[
  \sum_{n\ge1}n\,x^{n-1}=\frac{1}{(1-x)^2},
\]
dan pada $x=1/2$ ruas kanannya $4$. Jumlah dua ratus suku pertama deret di ruas
kiri memberi $4.000000$ pada presisi mesin. Mengintegralkan suku demi suku dari
$0$ ke $x$ memberi
\[
  \sum_{n\ge1}\frac{(-1)^{n+1}x^n}{n}=\ln(1+x),\qquad |x|<1,
\]
yaitu deret pada pembuka bab. Kedua hasil itu bukan kebetulan numerik: keduanya
konsekuensi teorema di atas, karena $1/2$ dan $x$ yang dipakai terletak di dalam
jari-jari.
\end{contoh}

Dengan teorema itu, pertanyaan pembuka bab terjawab. Deret $\ln(1+x)$ konvergen
seragam pada setiap $[-r,r]$ dengan $r<1$, sebab $|(-1)^{n+1}x^n/n|\le r^n/n$
dan $\sum r^n/n$ konvergen. Pada $(-1,1)$ keseragamannya gagal, dan harganya
terlihat pada jumlah suku: $16$ di $x=0.5$, $5248$ di $x=0.999$. Di titik ujung
$x=1$ deretnya tetap konvergen, sebab menjadi deret harmonik berganti tanda. Itu
bukan akibat jari-jari melainkan teorema Abel: bila $\sum a_nR^n$ konvergen,
maka $\sum a_nx^n$ konvergen seragam pada $[0,R]$. Kita kutip; buktinya ada di
Rudin, bab "Some Special Functions", Teorema 8.2.

Satu peringatan praktis: batas seragam dirancang untuk kasus terburuk. Pada
deret $\ln(1+x)$ dengan $N=12$ dan $r=0.5$, batas $r^{13}/13=9.4\times10^{-6}$
hanya tiga kali galat sebenarnya di $x=0.5$, tetapi di $x=0.1$ ia longgar sekitar
sepuluh ribu kali.

\begin{resep}
\begin{enumerate}
\item Tulis selisihnya, $M_N=\sup_{x\in E}|f_N(x)-f(x)|$, atau taksir dengan
      nilai di titik yang paling buruk.
\item Cari titik $x_N$ yang memaksimalkan selisih itu; biasanya puncak, atau
      titik yang mengejar ujung selang.
\item Hitung $M_N$ sebagai ungkapan yang hanya memuat $N$, lalu periksa
      $M_N\to0$.
\end{enumerate}
Peringatan urutan: cari $x_N$ lebih dulu untuk setiap $N$, baru hitung
supremumnya. Membalik urutan itu, yaitu memilih $N$ lalu mencari satu $x$ yang
mewakili semua titik, memberi kesimpulan seragam yang salah pada $g_n$ di
Gambar~\ref{fig:puncak}, yang puncaknya justru bergeser bersama $N$.
\end{resep}

\section{Contoh hitung}
\tingkat{menengah}

Tiga kelompok angka di bawah ini dihitung ulang di lab komputasi. Yang pertama
mengukur kegagalan keseragaman $x^n$ pada $[0,1]$.

\begin{center}
\small
\begin{tabular}{rrrrr}
\hline
$n$ & $\sup_{[0,1]}x^n$ & $(1-\tfrac1n)^n$ & $(1-\tfrac1{\sqrt n})^n$ & $0.9^{\,n}$\\
\hline
$10$    & $1$ & $0.3486784401$ & $2.234\times10^{-2}$ & $0.3486784401$\\
$100$   & $1$ & $0.3660323413$ & $2.656\times10^{-5}$ & $2.656\times10^{-5}$\\
$1000$  & $1$ & $0.3676954248$ & $1.108\times10^{-14}$ & $1.748\times10^{-46}$\\
$10000$ & $1$ & $0.3678610464$ & $2.249\times10^{-44}$ & $<10^{-450}$\\
\hline
\end{tabular}
\end{center}

Kolom kedua dan kelima adalah dua laju yang berbeda. Di titik tetap $x=0.9$
galatnya $0.9^n$, turun dengan faktor $0.9$ per suku: perlu $132$ suku untuk
mencapai $10^{-6}$, sedangkan di $x=0.5$ cukup $20$ suku dan di $x=0.99$ perlu
$1375$. Setiap titik punya laju eksponensialnya sendiri, dengan basis $x$. Kolom
kedua tidak bergerak: galat seragam pada $[0,1]$ tetap $1$ untuk semua $n$, jadi
laju penurunannya nol. Yang dibandingkan karena itu bukan dua bilangan,
melainkan sesuatu yang menuju nol dan sesuatu yang tidak.

Kelompok angka kedua adalah galat aproksimasi $e^{x}$ dengan deret Taylor di
$x=1$. Sisa deretnya dibatasi $\frac{e}{(N+1)!}$ untuk $x\in[0,1]$, sebab
turunan ke-$(N+1)$ dari $e^x$ adalah $e^x\le e$.

\begin{center}
\small
\begin{tabular}{rrrr}
\hline
$N$ & $S_N(1)=\sum_{k=0}^{N}\frac{1}{k!}$ & $e-S_N(1)$ & batas $\frac{e}{(N+1)!}$\\
\hline
$5$  & $2.716666666667$ & $1.615\times10^{-3}$ & $3.775\times10^{-3}$\\
$10$ & $2.718281801146$ & $2.731\times10^{-8}$ & $6.810\times10^{-8}$\\
$15$ & $2.718281828459$ & $5.063\times10^{-14}$ & $1.299\times10^{-13}$\\
$20$ & $2.718281828459$ & $0$ (di bawah presisi) & $5.320\times10^{-20}$\\
\hline
\end{tabular}
\end{center}

Batas itu selalu benar dan lebih besar daripada galat sebenarnya, dengan faktor
sekitar $2.3$ sampai $2.6$: sisanya kira-kira
$\frac{1}{(N+1)!}\left(1+\frac{1}{N+2}+\cdots\right)$, dan $e$ muncul karena
sukunya dijumlahkan di $x=1$. Pada $N=20$ galat sebenarnya tidak terukur lagi,
sebab $5.3\times10^{-20}$ jauh di bawah satu unit pembulatan terakhir pada
besaran $2.718$, yaitu sekitar $4.4\times10^{-16}$.

Kelompok angka ketiga adalah tiga jari-jari konvergensi, dihitung dengan dua
cara berbeda pada suku ke-$1000$.

\begin{center}
\small
\begin{tabular}{llrr}
\hline
deret & $R$ dari rumus & $|a_{1000}/a_{1001}|$ & $|a_{1000}|^{1/1000}$\\
\hline
$\sum x^n/n^2$          & $1$              & $1.0020010000$ & $0.9862794856$\\
$\sum \frac{n!}{n^n}x^n$ & $e=2.7182818285$ & $2.7169239322$ & $0.3694916635$\\
$\sum \frac{2^n}{n}x^n$  & $0.5$            & $0.5005000000$ & $1.9862320968$\\
\hline
\end{tabular}
\end{center}

Kolom ketiga menaksir $R$ dari nisbah dua suku berurutan, kolom keempat dari
kebalikan taksiran akar. Untuk deret kedua, $|a_n/a_{n+1}|=(1+1/n)^n$ naik ke
$e$; pada $n=1000$ nilainya $2.7169239$, meleset $0.0014$. Untuk deret pertama,
$(1+1/n)^2$ turun ke $1$ dari atas, jadi taksirannya $1.0020$, dan kekeliruan
$0.0020$ itu tanda kita baru di suku ke-$1000$. Kolom keempat lebih lambat menuju
limitnya, sebab $|a_n|^{1/n}$ mendekati $1$ dengan laju $|\ln a_n|/n$; pada
$n=1000$ selisihnya $0.0016$, sebanding dengan $\ln(2\pi n)/(2n)$ dikali $1/e$.

\kertas{Ringkasan angka. Galat seragam $x^n$ pada $[0,1]$ adalah $1$ untuk setiap
$n$, sementara galat di titik tetap $x$ adalah $x^n$: $132$ suku untuk $10^{-6}$
di $x=0.9$, $1375$ suku di $x=0.99$. Deret $e^x$ di $x=1$ mencapai galat
$2.7\times10^{-8}$ dengan $10$ suku, di bawah batas $6.8\times10^{-8}$.
Jari-jari tiga deret adalah $1$, $e$, dan $0.5$; integral
$\int_0^1 nx(1-x^2)^n dx=n/(2(n+1))$ bernilai $0.4545$ pada $n=10$, sementara
limit fungsinya nol.}

\section{Lab komputasi}
\tingkat{lanjut}

Kode berikut dijalankan apa adanya pada Python 3.14.4 dengan NumPy 2.5.2 dan
SciPy 1.18.0. Ia menghitung ulang ketiga kelompok angka di atas, ditambah
kuadratur untuk contoh tandingan integral.

\begin{lab}
\begin{verbatim}
import math
import numpy as np
from scipy.integrate import quad

print("n      sup[0,1]   (1-1/n)^n   (1-1/sqrt n)^n     0.9^n")
for n in (10, 100, 1000, 10000):
    a = (1 - 1/n)**n
    b = (1 - 1/math.sqrt(n))**n
    print(f"{n:6d} {1.0:9.6f} {a:11.6f} {b:17.3e} {0.9**n:10.3e}")

print()
print("suku ln(1+x) agar galat < 1e-6")
for x in (0.5, 0.9, 0.99, 0.999):
    N = 1
    while x**N / N > 1e-6:
        N += 1
    s = sum(((-1)**(k+1)) * x**k / k for k in range(1, N + 1))
    print(f"x={x:6.3f}  N={N:5d}  galat={abs(math.log(1+x) - s):.3e}")

print()
print("N      S_N              e - S_N        e/(N+1)!")
for N in (5, 10, 15, 20):
    s = math.fsum(1/math.factorial(k) for k in range(N + 1))
    b = math.e / math.factorial(N + 1)
    print(f"{N:3d} {s:.12f} {math.e - s:14.3e} {b:13.3e}")

print()
print("jari-jari dari nisbah suku ke-1000 dan ke-1001 (log-space)")
la = lambda n: -2 * math.log(n)
lb = lambda n: math.lgamma(n + 1) - n * math.log(n)
lc = lambda n: n * math.log(2) - math.log(n)
for nama, g in (("a", la), ("b", lb), ("c", lc)):
    R = math.exp(g(1000) - g(1001))
    H = math.exp(g(1000) / 1000)
    print(f"deret {nama}: R = {R!r}   |a_n|^(1/n) = {H!r}")
print("e =", repr(math.e))

print()
print("f_n(x) = n x (1-x^2)^n di [0,1]")
for n in (10, 100):
    val, _ = quad(lambda x, n=n: n * x * (1 - x*x)**n, 0, 1)
    print(f"n={n:4d}  integral={val!r}  n/(2(n+1))={n/(2*(n+1))!r}")
\end{verbatim}
\end{lab}

Keluaran yang diperoleh:

\begin{verbatim}
n      sup[0,1]   (1-1/n)^n   (1-1/sqrt n)^n     0.9^n
    10  1.000000    0.348678         2.234e-02  3.487e-01
   100  1.000000    0.366032         2.656e-05  2.656e-05
  1000  1.000000    0.367695         1.108e-14  1.748e-46
 10000  1.000000    0.367861         2.249e-44  0.000e+00

suku ln(1+x) agar galat < 1e-6
x= 0.500  N=   16  galat=3.049e-07
x= 0.900  N=   89  galat=4.478e-07
x= 0.990  N=  721  galat=4.915e-07
x= 0.999  N= 5248  galat=4.993e-07

N      S_N              e - S_N        e/(N+1)!
  5 2.716666666667      1.615e-03     3.775e-03
 10 2.718281801146      2.731e-08     6.810e-08
 15 2.718281828459      5.063e-14     1.299e-13
 20 2.718281828459      0.000e+00     5.320e-20

jari-jari dari nisbah suku ke-1000 dan ke-1001 (log-space)
deret a: R = 1.0020010000000001   |a_n|^(1/n) = 0.9862794856312105
deret b: R = 2.7169239322350776   |a_n|^(1/n) = 0.36949166347196
deret c: R = 0.5004999999999973   |a_n|^(1/n) = 1.9862320968418674
e = 2.718281828459045

f_n(x) = n x (1-x^2)^n di [0,1]
n=  10  integral=0.4545454545454544  n/(2(n+1))=0.45454545454545453
n= 100  integral=0.4950495049504951  n/(2(n+1))=0.49504950495049505
\end{verbatim}

Empat hal perlu dibaca dari keluaran ini. Pertama, kolom kedua tetap $1.000000$
untuk keempat nilai $n$, sementara kolom kelima sudah $3.5\times10^{-1}$ pada
$n=10$: kegagalan keseragaman terlihat sebagai angka. Kedua, baris terakhir kolom
kelima tercetak $0.000e+00$ bukan karena galatnya nol, melainkan karena
$0.9^{10000}=10^{-457.6}$ di bawah batas \texttt{float64} terkecil yang masih
normal: underflow.

Ketiga, kolom keempat memberi bukti numerik untuk dua limit yang sudah
dibuktikan: $(1-1/n)^n\to1/e=0.3678794$ dan $(1-1/\sqrt n)^n\to0$, yang runtuh
jauh lebih cepat.

Keempat, kuadratur SciPy memberi $0.4545454545454544$ untuk $n=10$, sedangkan
rumus $n/(2(n+1))$ memberi $0.45454545454545453$; selisih satu unit pembulatan
terakhir. Tanpa bentuk tertutup, kita tidak punya cara memastikan kuadratur tidak
mengukur hal lain.

Satu hal yang tidak terlihat di keluaran: semua batas di bab ini batas atas.
Baris $N$ yang dilaporkan kode dipilih dengan memeriksa $x^N/N$, jadi $N$ itu
dijamin cukup, bukan terkecil; galat sebenarnya $3.0\times10^{-7}$ sampai
$5.0\times10^{-7}$ di bawah $10^{-6}$.

\begin{untukdosen}
Urutan mengajar yang saya sarankan: mulai dari kegagalan pustaka $\ln(1+x)$ pada
pembuka, lalu definisi titik demi titik dan seragam dengan $x^n$ sebagai contoh
tunggal. Jangan buru-buru ke uji-$M$. Habiskan satu pertemuan penuh untuk
memastikan mahasiswa bisa mengatakan dengan tepat mengapa $x^n$ gagal seragam
dan mengapa $(1-1/\sqrt n)^n$ berhasil. Setelah itu uji-$M$, lalu tiga teorema
pertukaran, masing-masing dengan contoh tandingannya sendiri.

Dua miskonsepsi yang paling sering muncul. Pertama, mahasiswa mengira "setiap
$f_n$ kontinu dan $f_n\to f$" sudah cukup untuk menyimpulkan $f$ kontinu; mereka
lupa bahwa yang dibutuhkan keseragaman, dan $x^n$ adalah obatnya. Kedua,
mahasiswa mengira keseragaman turun-menurun: bila $f_n\to f$ seragam maka
$f_n'\to f'$. Contoh $\sin(nx)/n$ mematahkan itu dalam satu baris; sebaiknya
ditulis di papan sebelum teorema turunan muncul.

Aktivitas papan tulis yang berhasil: tulis empat barisan fungsi tanpa nama,
yaitu $x^n$, $x^n/n$, $nx(1-x^2)^n$, dan $\sin(nx)/n$, lalu minta kelas
menentukan untuk masing-masing apakah konvergen titik demi titik, seragam, dan
apakah integral serta turunannya bertukar. Jawabannya membentuk tabel $4\times4$
yang bisa diisi dalam dua puluh menit. Mahasiswa siap lanjut bila ia bisa
menjelaskan mengapa deret pangkat hanya seragam pada $[-r,r]$.
\end{untukdosen}

\section{Ringkasan}
\tingkat{dasar}

\begin{ringkasan}
\begin{itemize}
\item Kekonvergenan titik demi titik memberi setiap titik ambangnya sendiri;
      kekonvergenan seragam menuntut satu ambang untuk semua titik sekaligus,
      dan perbedaan itu setara dengan $\|f_n-f\|_\infty\to0$.
\item Uji-$M$ Weierstrass menyimpulkan keseragaman dari satu deret bilangan
      pembanding. Ia syarat cukup, bukan perlu; kegagalannya tidak membuktikan
      apa pun.
\item Limit seragam fungsi kontinu kontinu, dan integralnya boleh ditukar dengan
      limit. Turunan menuntut keseragaman barisan turunan ditambah kekonvergenan
      di satu titik.
\item Ketiga hipotesis itu tidak dapat dihapus: $x^n$, $nx(1-x^2)^n$, dan
      $\sin(nx)/n$ masing-masing mematahkan satu di antaranya.
\item Deret pangkat konvergen mutlak di dalam jari-jarinya dan divergen di luar;
      di dalam jari-jari ia boleh diturunkan dan diintegralkan suku demi suku
      tanpa mengubah jari-jarinya.
\item Keseragaman selalu lokal: $\ln(1+x)$ seragam pada setiap $[-r,r]$, $r<1$,
      tetapi tidak pada $(-1,1)$, dan harganya terlihat pada jumlah suku, dari
      $16$ di $x=0.5$ sampai $5248$ di $x=0.999$.
\end{itemize}
\end{ringkasan}

\section{Latihan}
\tingkat{menengah}

\begin{latihan}
\begin{enumerate}
\item Hitung dengan Python, untuk $n=10,100,1000$, nilai
      $\sup_{x\in[0,1]}|x^n|$ dan nilai $x^n$ di $x=0.9$. Jelaskan mengapa
      kolom pertama datar dan kolom kedua tidak.

\item Buktikan: bila $f_n\to f$ seragam pada $E$ dan setiap $f_n$ terbatas pada
      $E$, maka $f$ terbatas pada $E$ dan barisan $\|f_n\|_\infty$ terbatas.
      Tunjukkan dengan contoh bahwa kekonvergenan titik demi titik tidak cukup.

\item Misalkan $f_n(x)=x/(1+nx^2)$ pada $\mathbb{R}$. Tunjukkan bahwa
      $f_n\to0$ titik demi titik, tentukan $\|f_n\|_\infty$, dan buktikan bahwa
      kekonvergenannya seragam. Berapa $n$ terkecil sehingga
      $\|f_n\|_\infty<0.01$?

\item Nilai klaim ini: benar, salah, atau belum cukup bukti. "Karena $x^n\to0$
      untuk setiap $x\in[0,1)$, barisan $f_n(x)=x^n$ konvergen seragam pada
      $[0,1)$." Sebutkan tepat di mana penalaran itu patah, dan berikan satu
      barisan titik yang membuktikannya.

\item Hitung galat $|e-S_N(1)|$ untuk $N=5,10,15,20$ dan bandingkan dengan
      batas $e/(N+1)!$. Berapa $N$ terkecil agar batas itu di bawah
      $10^{-16}$? Apakah galat sebenarnya masih bisa diukur pada $N$ itu?

\item Buktikan bahwa deret $\sum_{n\ge0}x^n(1-x)$ konvergen titik demi titik
      pada $[0,1]$ tetapi tidak seragam, dengan menghitung jumlahnya lebih dulu
      dan memeriksa kekontinuannya.

\item Hitung jari-jari konvergensi $\sum x^n/n^2$, $\sum(n!/n^n)x^n$, dan
      $\sum 2^nx^n/n$ dari nisbah suku ke-$1000$ dan ke-$1001$, lalu bandingkan
      dengan nilai eksaknya. Deret mana yang taksirannya paling lambat, dan
      mengapa?

\item Buktikan bahwa bila $f_n\to f$ seragam pada $E$ dan $g$ terbatas pada $E$,
      maka $f_ng\to fg$ seragam pada $E$. Tunjukkan dengan contoh bahwa
      keterbatasan $g$ tidak dapat dihapus.

\item Untuk $f_n(x)=nx(1-x^2)^n$ pada $[0,1]$, hitung titik puncak, tinggi
      puncak, dan $\int_0^1f_n$ untuk $n=10,100$. Bandingkan tinggi puncak
      dengan $\sqrt{n/(2e)}$ dan jelaskan mengapa luasnya tidak menuju nol.

\item Turunkan $\sum_{n\ge0}x^n=1/(1-x)$ suku demi suku dan pakai hasilnya
      untuk menghitung $\sum_{n\ge1}n/2^n$. Berapa suku yang diperlukan agar
      jumlah parsialnya berada dalam $10^{-10}$ dari nilai eksaknya?
\end{enumerate}

\noindent\emph{Petunjuk untuk tiga butir tersulit.} Butir 4: sifat seragam
menuntut satu $N$ untuk semua $x$, jadi cukup tunjukkan bahwa untuk setiap $n$
ada $x$ dengan $x^n\ge\tfrac12$; ambil $x=2^{-1/n}$. Butir 6: jumlahkan deret
geometrinya lebih dulu; hasilnya fungsi yang berbeda dari limitnya di satu titik,
dan limit seragam fungsi kontinu harus kontinu. Butir 8: tulis
$|f_ng-fg|=|f_n-f|\,|g|$ dan pakai $\|g\|_\infty$; untuk contoh tandingannya
ambil $g(x)=1/x$ pada $(0,1)$ dan $f_n$ yang mengecil seragam tetapi tidak
cukup cepat.
\end{latihan}

\begin{untukpeneliti}
Yang terbukti dengan bukti lengkap di bab ini: kesetaraan keseragaman dengan
norma supremum, kriteria Cauchy seragam, uji-$M$ Weierstrass, ketiga teorema
pertukaran beserta contoh tandingannya, rumus jari-jari lewat limit superior, dan
ketetapan jari-jari pada penurunan serta pengintegralan suku demi suku. Dua hasil
dikutip tanpa bukti: ketidakdiferensialan fungsi Weierstrass (1872, diperluas
Hardy 1916) dan teorema Abel. Sumber primernya Rudin, \emph{Principles of
Mathematical Analysis}, bab "Sequences and Series of Functions" (7.1--7.17) dan
"Some Special Functions" (8.1--8.2); Abbott, \emph{Understanding Analysis}, bab
"Sequences and Series of Functions"; Stewart, \emph{Calculus}, bagian "Power
Series".

Yang sering dipopulerkan berlebihan adalah gagasan bahwa keseragaman
menghilangkan semua masalah. Ia menghilangkan tiga masalah tertentu, pada selang
tempat ia berlaku, dan tidak lebih. Pada deret Fourier keseragaman justru jarang,
dan pertukaran integral tetap sah karena alasan lain. Pada aproksimasi jaringan
saraf, yang dipakai adalah kekonvergenan pada himpunan hingga.

Pertanyaan terbuka: bila $f_n$ hanya diketahui lewat nilai pada grid berhingga,
seberapa besar kekeliruan taksiran $\|f_n\|_\infty$ oleh maksimum di grid? Tanpa
asumsi keteraturan tidak ada jawaban umum, dan untuk puncak sempit grid yang
tampak rapat bisa meleset seluruhnya. Batas keberlakuan: semua fungsi di sini
bernilai real pada himpunan bagian $\mathbb{R}$, dengan aritmetika 64-bit pada
Python 3.14.4, NumPy 2.5.2, dan SciPy 1.18.0.
\end{untukpeneliti}

\peta{Bab 14 pada Rudin, \emph{Principles of Mathematical Analysis}, bab
"Sequences and Series of Functions" (7.1--7.17) dan "Some Special Functions"
(8.1--8.2); pada Abbott, \emph{Understanding Analysis}, bab "Sequences and
Series of Functions"; pada Stewart, \emph{Calculus}, bagian "Power Series" dan
"Taylor and Maclaurin Series". Basisnya paruh kedua bab \emph{Integral, Deret
Fungsi, dan Pertukaran Limit} edisi pertama buku ini. Bab ini memakai definisi
limit dan kriteria Cauchy dari Bab 4, uji akar Bab 5, kekontinuan seragam Bab 11,
serta kriteria keterintegralan Bab 13, semuanya dirujuk, tidak diulang.}
\chapter[Integral Lipat dan Teorema Lapangan Vektor]{Integral Lipat, Perubahan Variabel, dan Teorema Lapangan Vektor}

Sebuah pelat pendingin berbentuk cakram berjari-jari satu meter. Tebalnya
seragam, tetapi kerapatannya menurun dari pusat: $e^{-(x^{2}+y^{2})}$ kilogram
per meter persegi. Massanya adalah jumlah kerapatan di seluruh cakram: pustaka numerik memberi
$1.9858653038$ kilogram setelah $11907$ evaluasi fungsi.

Rekan kita menulis versi lain dari perhitungan yang sama. Ia pindah ke
koordinat polar, menuliskan
\[
\int_{0}^{2\pi}\!\!\int_{0}^{1} e^{-r^{2}}\,dr\,d\theta ,
\]
dan mendapat $4.6924344183$ kilogram. Programnya berjalan tanpa peringatan.
Angka itu $2.36$ kali lebih besar daripada yang pertama, dan tidak satu pun
program memperingatkan. Kesalahannya satu huruf: faktor $r$ hilang, sehingga luas
tiap petak kecil dihitung terlalu besar.

Kerusakan sekelas itu muncul juga di sisi lain bab ini. Hitung kerja
$F(x,y)=(-y,x)/(x^{2}+y^{2})$ mengelilingi lingkaran satuan: hasilnya
$6.2831853072$, padahal curl medan itu nol di setiap titik tempat ia
terdefinisi. Kalau curl nol berarti kerja tertutup nol, salah satu hitungan
itu bohong. Yang bohong bukan hitungannya, melainkan syarat yang tidak
diperiksa: medan itu tidak terdefinisi di titik asal, dan titik asal ada di
dalam lingkaran.

Bab 13 membangun integral pada selang, dan Bab 12 memberi Jacobian sebagai
matriks turunan. Keduanya belum cukup: belum ada cara menjumlahkan pada daerah
berdimensi lebih dari satu, aturan tentang apa yang terjadi pada luas petak
ketika koordinat diganti, dan arti tanda pada batas. Bab ini mengisi
ketiganya, lalu menunjukkan bahwa Green, Gauss, dan Stokes sebenarnya satu
pernyataan.

\begin{untukanak}
Bayangkan penjahit dengan pola baju di atas kertas berpetak. Polanya persegi
dan rapi; kain yang akan dipotong melengkung mengikuti badan. Kalau pola persegi
itu ditempelkan begitu saja, bajunya sempit di satu sisi dan longgar di sisi
lain, walaupun jumlah petaknya tetap sama.

Penjahit berpengalaman melakukan hal lain: setiap petak kecil akan meregang
atau mengkerut dengan perbandingan tertentu saat dipindahkan ke kain. Petak di
dekat pinggang meregang sedikit; petak di dekat bahu mengkerut. Ia mencatat
perbandingan itu, mengalikan luas petak dengannya, baru menjumlahkan.

Pekerjaan yang sama muncul pada bentuk lengkung: banyak benda, luas, atau berat.
Kalau perbandingan regangan lupa dicatat, jumlahnya tetap bilangan rapi, dan
justru karena rapi ia sulit dicurigai. Di bab ini kita akan tahu mengapa
perbandingan itu harus ada, dan bagaimana menghitungnya untuk bentuk bulat.
\end{untukanak}

\section{Dari persegi panjang ke daerah umum}
\tingkat{dasar}

Kita mulai dari benda paling sederhana: sebuah persegi panjang, dan sebuah
bilangan yang hendak dijumlahkan di seluruh permukaannya. Caranya sama dengan
Bab 13, hanya potongannya sekarang berbentuk petak, bukan potongan garis.

\begin{definisi}[Integral lipat dua]\index{integral lipat dua}
Misalkan $R=[a,b]\times[c,d]$ sebuah persegi panjang tertutup. Sebuah
\emph{partisi} $\mathcal{P}$ dari $R$ adalah pembagian $[a,b]$ menjadi
selang-selang dan $[c,d]$ menjadi selang-selang; hasilnya adalah $mn$ persegi
panjang kecil $R_{ij}$. \emph{Mesh} partisi adalah diagonal terpanjang di
antara petak-petak itu, dan ditulis $\|\mathcal{P}\|$.

Untuk $f$ terbatas pada $R$, tulis $M_{ij}=\sup_{R_{ij}}f$ dan
$m_{ij}=\inf_{R_{ij}}f$, lalu
\[
U(\mathcal{P},f)=\sum_{i,j}M_{ij}\,\mathrm{luas}(R_{ij}),
\qquad
L(\mathcal{P},f)=\sum_{i,j}m_{ij}\,\mathrm{luas}(R_{ij}).
\]
Fungsi $f$ disebut \emph{terintegralkan} pada $R$ bila supremum semua $L$ sama
dengan infimum semua $U$; nilai bersama itu adalah $\iint_{R}f\,dA$.
\end{definisi}

Definisi ini tidak menyebut urutan pengintegralan. Bila $f$ kontinu, jumlah
atas dan jumlah bawah bertemu karena beda terbesar $f$ pada satu petak
mengecil bersama mesh: pada himpunan kompak, fungsi kontinu kontinu seragam
(Bab 11).

Pertanyaan berikutnya praktis. Menjumlahkan pada petak berarti menghitung
sesuatu berdimensi dua, sedangkan tangan kita hanya bisa menghitung integral satu
variabel. Jembatannya: memotong menurut satu arah lebih dulu, lalu menjumlahkan
hasil potongan itu.

\begin{teorema}[Fubini]\index{teorema Fubini}
Jika $f$ kontinu pada $R=[a,b]\times[c,d]$, maka
\[
\iint_{R} f\,dA=\int_{a}^{b}\!\!\int_{c}^{d} f(x,y)\,dy\,dx
=\int_{c}^{d}\!\!\int_{a}^{b} f(x,y)\,dx\,dy .
\]
\end{teorema}

\begin{proof}
Bukti penuh menuntut pencocokan jumlah atas-bawah dua dimensi dengan integral
satu dimensi secara seragam, dan ada di Spivak, bab \emph{Integration},
Teorema 3-10. Gagasannya begini.

Untuk $f\ge 0$, integral lipat dua adalah volume benda di bawah grafik $f$.
Iris benda itu dengan bidang $x$ tetap. Irisannya berupa daerah di bawah kurva
$y\mapsto f(x,y)$, sehingga luasnya
\[
A(x)=\int_{c}^{d}f(x,y)\,dy .
\]
Volume benda adalah integral dari luas irisan itu:
\[
\iint_{R}f\,dA=\int_{a}^{b}A(x)\,dx=\int_{a}^{b}\!\!\int_{c}^{d}f(x,y)\,dy\,dx .
\]
Karena $f$ kontinu pada himpunan kompak, $A$ kontinu. Untuk $f$ bertanda apa
pun, pisahkan $f=f^{+}-f^{-}$; arah sebaliknya diberikan oleh irisan $y$
tetap.
\end{proof}

Yang penting dari teorema ini bukan rumusnya, melainkan izinnya: urutan boleh
ditukar. Izin itu gugur pada fungsi yang tidak terintegralkan, dan itu dibahas
nanti.

Benda yang diiris dalam bukti itu tampak pada Gambar~\ref{fig:kolom-tipis}:
tiap kolom menyumbang $f(x,y)\,\Delta A$.

\begin{bacakeras}
--- Berapa isi gunungan pasir ini?
--- Jumlahkan kolom-kolom tipisnya. Berdiri di sisi utara, iris gunungan menjadi
jalur, ukur tiap jalur, lalu jumlahkan.
--- Berapa hasilnya?
--- Delapan belas ember.
--- Sekarang berdiri di sisi timur, iris dengan arah lain, ukur, jumlahkan.
--- Berapa hasilnya?
--- Delapan belas ember juga.
--- Jalur utara-selatan dan jalur timur-barat tidak satu pun sama. Mengapa
hasilnya sama?
--- Karena gunungannya satu. Kita tidak menjumlahkan dua gunungan; kita mengiris
benda yang sama dengan dua pisau berbeda arah.
--- Jadi arah iris boleh ditukar?
--- Selama bendanya terukur. Kalau gunungannya berlubang di tengah, dua arah itu
masih sama; yang gagal nanti lubangnya, bukan urutannya.

Ulangi sekali, dengan suara biasa: arah iris berganti, isi tetap.
\end{bacakeras}

\begin{figure}[htbp]
\centering
\begin{tikzpicture}[scale=1.55,font=\small]
  \fill[layergrey!10] (0,0) -- (6.1,0) -- (7.1,0.90) -- (1.0,0.90) -- cycle;
  \draw[layergrey,thick] (0,0) -- (6.1,0) -- (7.1,0.90) -- (1.0,0.90) -- cycle;
  \fill[layerteal!12] (0.91,0.82) rectangle (1.43,1.38);
  \draw[layerteal!70,line width=0.4pt] (0.91,0.82) rectangle (1.43,1.38);
  \fill[layerteal!30] (0.91,1.38) -- (1.08,1.47) -- (1.60,1.47) -- (1.43,1.38) -- cycle;
  \draw[layerteal!70,line width=0.4pt] (0.91,1.38) -- (1.08,1.47) -- (1.60,1.47);
  \fill[layerteal!12] (1.67,0.82) rectangle (2.19,1.60);
  \draw[layerteal!70,line width=0.4pt] (1.67,0.82) rectangle (2.19,1.60);
  \fill[layerteal!30] (1.67,1.60) -- (1.83,1.69) -- (2.35,1.69) -- (2.19,1.60) -- cycle;
  \draw[layerteal!70,line width=0.4pt] (1.67,1.60) -- (1.83,1.69) -- (2.35,1.69);
  \fill[layerteal!12] (2.42,0.82) rectangle (2.93,1.97);
  \draw[layerteal!70,line width=0.4pt] (2.42,0.82) rectangle (2.93,1.97);
  \fill[layerteal!30] (2.42,1.97) -- (2.58,2.06) -- (3.10,2.06) -- (2.93,1.97) -- cycle;
  \draw[layerteal!70,line width=0.4pt] (2.42,1.97) -- (2.58,2.06) -- (3.10,2.06);
  \fill[layerteal!12] (3.17,0.82) rectangle (3.68,2.30);
  \draw[layerteal!70,line width=0.4pt] (3.17,0.82) rectangle (3.68,2.30);
  \fill[layerteal!30] (3.17,2.30) -- (3.33,2.39) -- (3.85,2.39) -- (3.68,2.30) -- cycle;
  \draw[layerteal!70,line width=0.4pt] (3.17,2.30) -- (3.33,2.39) -- (3.85,2.39);
  \fill[layerteal!12] (3.92,0.82) rectangle (4.43,2.32);
  \draw[layerteal!70,line width=0.4pt] (3.92,0.82) rectangle (4.43,2.32);
  \fill[layerteal!30] (3.92,2.32) -- (4.08,2.41) -- (4.60,2.41) -- (4.43,2.32) -- cycle;
  \draw[layerteal!70,line width=0.4pt] (3.92,2.32) -- (4.08,2.41) -- (4.60,2.41);
  \fill[layerteal!12] (4.67,0.82) rectangle (5.18,2.00);
  \draw[layerteal!70,line width=0.4pt] (4.67,0.82) rectangle (5.18,2.00);
  \fill[layerteal!30] (4.67,2.00) -- (4.83,2.09) -- (5.35,2.09) -- (5.18,2.00) -- cycle;
  \draw[layerteal!70,line width=0.4pt] (4.67,2.00) -- (4.83,2.09) -- (5.35,2.09);
  \fill[layerteal!12] (5.42,0.82) rectangle (5.93,1.62);
  \draw[layerteal!70,line width=0.4pt] (5.42,0.82) rectangle (5.93,1.62);
  \fill[layerteal!30] (5.42,1.62) -- (5.58,1.71) -- (6.10,1.71) -- (5.93,1.62) -- cycle;
  \draw[layerteal!70,line width=0.4pt] (5.42,1.62) -- (5.58,1.71) -- (6.10,1.71);
  \fill[layerteal!12] (6.17,0.82) rectangle (6.68,1.39);
  \draw[layerteal!70,line width=0.4pt] (6.17,0.82) rectangle (6.68,1.39);
  \fill[layerteal!30] (6.17,1.39) -- (6.33,1.48) -- (6.85,1.48) -- (6.68,1.39) -- cycle;
  \draw[layerteal!70,line width=0.4pt] (6.17,1.39) -- (6.33,1.48) -- (6.85,1.48);
  \fill[layerteal!12] (0.59,0.50) rectangle (1.11,1.07);
  \draw[layerteal!70,line width=0.4pt] (0.59,0.50) rectangle (1.11,1.07);
  \fill[layerteal!30] (0.59,1.07) -- (0.76,1.16) -- (1.28,1.16) -- (1.11,1.07) -- cycle;
  \draw[layerteal!70,line width=0.4pt] (0.59,1.07) -- (0.76,1.16) -- (1.28,1.16);
  \fill[layerteal!12] (1.34,0.50) rectangle (1.86,1.33);
  \draw[layerteal!70,line width=0.4pt] (1.34,0.50) rectangle (1.86,1.33);
  \fill[layerteal!30] (1.34,1.33) -- (1.51,1.42) -- (2.03,1.42) -- (1.86,1.33) -- cycle;
  \draw[layerteal!70,line width=0.4pt] (1.34,1.33) -- (1.51,1.42) -- (2.03,1.42);
  \fill[layerteal!12] (2.09,0.50) rectangle (2.61,1.75);
  \draw[layerteal!70,line width=0.4pt] (2.09,0.50) rectangle (2.61,1.75);
  \fill[layerteal!30] (2.09,1.75) -- (2.26,1.84) -- (2.78,1.84) -- (2.61,1.75) -- cycle;
  \draw[layerteal!70,line width=0.4pt] (2.09,1.75) -- (2.26,1.84) -- (2.78,1.84);
  \fill[layerteal!12] (2.84,0.50) rectangle (3.36,2.14);
  \draw[layerteal!70,line width=0.4pt] (2.84,0.50) rectangle (3.36,2.14);
  \fill[layerteal!30] (2.84,2.14) -- (3.01,2.23) -- (3.53,2.23) -- (3.36,2.14) -- cycle;
  \draw[layerteal!70,line width=0.4pt] (2.84,2.14) -- (3.01,2.23) -- (3.53,2.23);
  \fill[layerteal!12] (3.59,0.50) rectangle (4.11,2.15);
  \draw[layerteal!70,line width=0.4pt] (3.59,0.50) rectangle (4.11,2.15);
  \fill[layerteal!30] (3.59,2.15) -- (3.76,2.24) -- (4.28,2.24) -- (4.11,2.15) -- cycle;
  \draw[layerteal!70,line width=0.4pt] (3.59,2.15) -- (3.76,2.24) -- (4.28,2.24);
  \fill[layerteal!12] (4.34,0.50) rectangle (4.86,1.79);
  \draw[layerteal!70,line width=0.4pt] (4.34,0.50) rectangle (4.86,1.79);
  \fill[layerteal!30] (4.34,1.79) -- (4.51,1.88) -- (5.03,1.88) -- (4.86,1.79) -- cycle;
  \draw[layerteal!70,line width=0.4pt] (4.34,1.79) -- (4.51,1.88) -- (5.03,1.88);
  \fill[layerteal!12] (5.09,0.50) rectangle (5.61,1.35);
  \draw[layerteal!70,line width=0.4pt] (5.09,0.50) rectangle (5.61,1.35);
  \fill[layerteal!30] (5.09,1.35) -- (5.26,1.44) -- (5.78,1.44) -- (5.61,1.35) -- cycle;
  \draw[layerteal!70,line width=0.4pt] (5.09,1.35) -- (5.26,1.44) -- (5.78,1.44);
  \fill[layerteal!12] (5.84,0.50) rectangle (6.36,1.08);
  \draw[layerteal!70,line width=0.4pt] (5.84,0.50) rectangle (6.36,1.08);
  \fill[layerteal!30] (5.84,1.08) -- (6.01,1.17) -- (6.53,1.17) -- (6.36,1.08) -- cycle;
  \draw[layerteal!70,line width=0.4pt] (5.84,1.08) -- (6.01,1.17) -- (6.53,1.17);
  \fill[layerteal!12] (0.26,0.17) rectangle (0.78,0.73);
  \draw[layerteal!70,line width=0.4pt] (0.26,0.17) rectangle (0.78,0.73);
  \fill[layerteal!30] (0.26,0.73) -- (0.43,0.82) -- (0.95,0.82) -- (0.78,0.73) -- cycle;
  \draw[layerteal!70,line width=0.4pt] (0.26,0.73) -- (0.43,0.82) -- (0.95,0.82);
  \fill[layerteal!12] (1.02,0.17) rectangle (1.54,0.95);
  \draw[layerteal!70,line width=0.4pt] (1.02,0.17) rectangle (1.54,0.95);
  \fill[layerteal!30] (1.02,0.95) -- (1.19,1.04) -- (1.71,1.04) -- (1.54,0.95) -- cycle;
  \draw[layerteal!70,line width=0.4pt] (1.02,0.95) -- (1.19,1.04) -- (1.71,1.04);
  \fill[layerteal!12] (1.76,0.17) rectangle (2.29,1.32);
  \draw[layerteal!70,line width=0.4pt] (1.76,0.17) rectangle (2.29,1.32);
  \fill[layerteal!30] (1.76,1.32) -- (1.93,1.41) -- (2.46,1.41) -- (2.29,1.32) -- cycle;
  \draw[layerteal!70,line width=0.4pt] (1.76,1.32) -- (1.93,1.41) -- (2.46,1.41);
  \fill[layerteal!12] (2.51,0.17) rectangle (3.04,1.65);
  \draw[layerteal!70,line width=0.4pt] (2.51,0.17) rectangle (3.04,1.65);
  \fill[layerteal!30] (2.51,1.65) -- (2.69,1.74) -- (3.21,1.74) -- (3.04,1.65) -- cycle;
  \draw[layerteal!70,line width=0.4pt] (2.51,1.65) -- (2.69,1.74) -- (3.21,1.74);
  \fill[layerteal!12] (3.26,0.17) rectangle (3.79,1.67);
  \draw[layerteal!70,line width=0.4pt] (3.26,0.17) rectangle (3.79,1.67);
  \fill[layerteal!30] (3.26,1.67) -- (3.44,1.76) -- (3.96,1.76) -- (3.79,1.67) -- cycle;
  \draw[layerteal!70,line width=0.4pt] (3.26,1.67) -- (3.44,1.76) -- (3.96,1.76);
  \fill[layerteal!12] (4.01,0.17) rectangle (4.53,1.35);
  \draw[layerteal!70,line width=0.4pt] (4.01,0.17) rectangle (4.53,1.35);
  \fill[layerteal!30] (4.01,1.35) -- (4.18,1.44) -- (4.70,1.44) -- (4.53,1.35) -- cycle;
  \draw[layerteal!70,line width=0.4pt] (4.01,1.35) -- (4.18,1.44) -- (4.70,1.44);
  \fill[layerteal!12] (4.76,0.17) rectangle (5.28,0.97);
  \draw[layerteal!70,line width=0.4pt] (4.76,0.17) rectangle (5.28,0.97);
  \fill[layerteal!30] (4.76,0.97) -- (4.93,1.06) -- (5.45,1.06) -- (5.28,0.97) -- cycle;
  \draw[layerteal!70,line width=0.4pt] (4.76,0.97) -- (4.93,1.06) -- (5.45,1.06);
  \fill[layerteal!12] (5.51,0.17) rectangle (6.03,0.74);
  \draw[layerteal!70,line width=0.4pt] (5.51,0.17) rectangle (6.03,0.74);
  \fill[layerteal!30] (5.51,0.74) -- (5.68,0.83) -- (6.20,0.83) -- (6.03,0.74) -- cycle;
  \draw[layerteal!70,line width=0.4pt] (5.51,0.74) -- (5.68,0.83) -- (6.20,0.83);
  \fill[accent!42] (2.69,0.17) rectangle (2.86,1.65);
  \draw[accent,line width=1pt] (2.69,0.17) rectangle (2.86,1.65);
  \fill[accent!55] (2.69,1.65) -- (2.78,1.70) -- (2.96,1.70) -- (2.86,1.65) -- cycle;
  \draw[accent,line width=1pt] (2.69,1.65) -- (2.78,1.70) -- (2.96,1.70);
  \draw[accent,densely dashed] (2.86,0.16) -- (3.42,-0.26);
  \node[accent,anchor=north west,align=left]
    at (3.50,-0.30) {satu kolom tipis: alas $\Delta A$,\\ tinggi $f(x,y)$, isi $\approx f(x,y)\,\Delta A$};
  \node[layergrey,anchor=north west] at (-0.05,-0.08) {alas $R$};
  \node[layergrey,anchor=north] at (5.95,-0.02) {$x$};
  \node[layergrey,anchor=west] at (7.15,0.55) {$y$};
\end{tikzpicture}
\caption{Benda di bawah grafik $f$ di atas persegi panjang $R$, dipotong menjadi
kolom-kolom tipis. Setiap kolom hampir berupa balok dengan alas $\Delta A$ dan tinggi
$f(x,y)$, sehingga isinya $f(x,y)\,\Delta A$; menjumlahkan seluruh kolom memberi
$\iint_R f\,dA$. Kolom yang ditandai merah adalah satu suku dalam jumlah itu.
Mengiris dari arah $x$ maupun dari arah $y$ menghasilkan kolom yang berbeda, tetapi
jumlah isinya sama.}
\label{fig:kolom-tipis}
\end{figure}

Daerah yang muncul di soal jarang berupa persegi panjang. Kita menanganinya
dengan satu langkah tambahan: perluas fungsinya dengan nol di luar daerah.

\begin{lemma}[Daerah di antara dua grafik]\index{daerah bertipe I}
Misalkan $g_{1},g_{2}\colon[a,b]\to\mathbb{R}$ kontinu dengan $g_{1}\le
g_{2}$, dan
\[
D=\{(x,y): a\le x\le b,\ g_{1}(x)\le y\le g_{2}(x)\}.
\]
Jika $f$ kontinu pada $D$, maka
\[
\iint_{D}f\,dA=\int_{a}^{b}\!\!\int_{g_{1}(x)}^{g_{2}(x)}f(x,y)\,dy\,dx .
\]
\end{lemma}

\begin{proof}
Ambil persegi panjang $R$ yang memuat $D$, dan perluas $f$ dengan nilai $0$ di
$R\setminus D$. Batas $D$ dapat ditutup oleh berhingga persegi panjang dengan
luas total kurang dari $\varepsilon$: pada grafik, kekontinuan seragam pada
selang kompak memberi partisi yang goyangan $g$ kecil di tiap potongan.
Himpunan yang dapat ditutup begitu tidak mengubah nilai integral, sehingga
\[
\iint_{D}f\,dA=\iint_{R} f\cdot\mathbf{1}_{D}\,dA .
\]
Menurut Fubini, integral kanan dihitung dengan mengintegralkan $x$ lebih dulu;
untuk $x$ tetap, integran $\mathbf{1}_{D}f$ bernilai $f$ hanya pada
$g_{1}(x)\le y\le g_{2}(x)$.
\end{proof}

Daerah bertipe I cukup untuk hampir semua soal di bab ini; batas yang memotong
garis tegak lebih dari dua kali dipecah lalu dijumlahkan. Cakram paling penting:
di Cartesius ia menuntut dua potongan, di polar cukup satu.

\begin{untukanak}
Sekarang ambil sebuah ruangan tertutup dengan beberapa tungku di dalamnya.
Setiap tungku mengeluarkan panas. Pertanyaannya: berapa panas yang keluar dari
ruangan itu setiap detik, dijumlahkan dari seluruh dinding, lantai, dan atap?

Cara pertama menghitungnya sulit: berkeliling ke setiap dinding, mengukur aliran
panas di tiap titik, lalu menjumlahkan. Cara kedua jauh lebih mudah: hitung
berapa panas yang dihasilkan semua tungku di dalam. Kalau tidak ada panas yang
hilang lewat celah lain, kedua angka harus sama: jumlah di dalam sama dengan
jumlah yang menembus kulit.

Kalimat itu terdengar terlalu bagus untuk benar. Kesamaan itu hanya berlaku
kalau kulit ruangan benar-benar tertutup dan arah "keluar" ditentukan dengan tegas. Kalau ada satu lubang
kecil yang tidak dihitung, atau kalau kita salah menandai mana sisi luar dan
mana sisi dalam, angkanya meleset walaupun semua penjumlahan kita benar.

Di bab ini kita akan tahu mengapa jumlah di dalam bisa dibaca dari kulitnya,
untuk bentuk lengkung apa pun, dan bagaimana satu lubang kecil dapat merusak
seluruh kesimpulan.
\end{untukanak}

\section{Perubahan variabel dan determinan Jacobian}
\tingkat{menengah}

Mengganti koordinat berarti mengganti petak: petak persegi di bidang $(u,v)$
menjadi petak melengkung di bidang $(x,y)$, dan luasnya berubah. Kita butuh
satu bilangan di tiap titik.

\begin{definisi}[Perubahan variabel]
\index{perubahan variabel}\index{Jacobian} Misalkan $D^{*}$ daerah terbuka di
$\mathbb{R}^{2}$ dan $T\colon D^{*}\to\mathbb{R}^{2}$,
$T(u,v)=(x(u,v),y(u,v))$, fungsi dengan turunan parsial kontinu. $T$ disebut
\emph{perubahan variabel} yang sah bila $T$ satu-satu pada bagian dalam
$D^{*}$ dan determinan Jacobiannya
\[
\det J_{T}(u,v)=\frac{\partial x}{\partial u}\frac{\partial y}{\partial v}
-\frac{\partial x}{\partial v}\frac{\partial y}{\partial u}
\]
tidak nol di sana. Matriks $J_{T}$ sendiri sudah kita kenal di Bab 12.
\end{definisi}

Kalimat sehari-harinya: ia angka yang memberi tahu berapa kali sebuah petak
kecil membesar ketika dipindahkan oleh $T$. Nilai $3$ berarti petak menjadi tiga kali lebih luas, $1/2$ berarti menyusut
setengah, dan nol berarti petak gepeng menjadi garis dan informasinya hilang.

Gambar~\ref{fig:juring-jacobian} memperlihatkan satu petak dan bayangannya:
luas $0.04$ menjadi $0.064$, dan angka $1.6$ itu adalah $|{\det J}|$ di tengah
petak.

\begin{figure}[htbp]
\centering
\begin{tikzpicture}[font=\footnotesize]
  \begin{scope}[x=1.9cm,y=1.9cm]
    \fill[layergrey!7] (1,0) rectangle (2.4,1.2);
    \draw[layergrey!40,very thin] (1,0) grid[xstep=0.2,ystep=0.2] (2.4,1.2);
    \draw[->,ink] (0.88,0) -- (2.62,0) node[right] {$u$};
    \draw[->,ink] (1,-0.08) -- (1,1.42) node[above] {$v$};
    \fill[layerindigo!30] (1.5,0.6) rectangle (1.7,0.8);
    \draw[layerindigo,thick] (1.5,0.6) rectangle (1.7,0.8);
    \node[layerindigo!85!black,anchor=north,align=center,fill=white,inner sep=1.2pt]
      at (1.60,0.55) {petak $0.2\times0.2$\\ luas $0.04$};
    \node[layergrey,anchor=north] at (1.7,-0.10) {bidang $(u,v)$};
  \end{scope}
  \draw[-{Latex[length=2.4mm]},ink,thick] (4.62,1.14) -- (5.85,1.14);
  \node[ink,above] at (5.23,1.18) {$T$};
  \node[ink,below,font=\scriptsize] at (5.23,1.06) {$\times\,1.6$};
  \begin{scope}[xshift=5.9cm,x=1.9cm,y=1.9cm]
    \foreach \r in {1,1.2,1.4,1.6,1.8,2,2.2,2.4}{
      \draw[layergrey!40,very thin] (\r,0) arc (0:69:\r);}
    \foreach \a in {0,11.5,23,34.5,46,57.5,69}{
      \draw[layergrey!40,very thin] (0,0) -- (\a:2.4);}
    \draw[->,ink] (-0.08,0) -- (2.62,0) node[right] {$x$};
    \draw[->,ink] (0,-0.08) -- (0,2.52) node[above] {$y$};
    \fill[accent!30] (34.4:1.5) arc (34.4:45.8:1.5) -- (45.8:1.7)
      arc (45.8:34.4:1.7) -- cycle;
    \draw[accent,thick] (34.4:1.5) arc (34.4:45.8:1.5) -- (45.8:1.7)
      arc (45.8:34.4:1.7) -- cycle;
    \node[accent,anchor=west,align=left] at (1.62,1.42)
      {juring melengkung\\ luas $0.064=1.6\times0.04$};
    \node[layergrey,anchor=north] at (1.6,-0.10) {bidang $(x,y)$};
  \end{scope}
\end{tikzpicture}
\caption{Satu petak kecil di bidang $(u,v)$ dan bayangannya di bidang $(x,y)$ di bawah
$T(u,v)=(u\cos v,\;u\sin v)$. Petak $0.2\times0.2$ berluas $0.04$ menjadi juring dengan
luas $0.064$, yaitu $1.6$ kali lipat, persis nilai $|{\det J}|=u$ di tengah petak
($u=1.6$). Perhatikan bahwa sisinya ikut melengkung: yang berubah bukan hanya ukuran,
melainkan juga bentuk, dan satu angka di tiap titik cukup untuk mencatat perubahannya.}
\label{fig:juring-jacobian}
\end{figure}

\begin{tangan}
Ambil kertas persegi $12$ cm kali $12$ cm, penggaris, pensil, gunting, dan
selotip. Gambar petak $1$ cm di seluruh kertas; ada $144$ petak. Potong kertas
menjadi dua belas jalur tegak, masing-masing $1$ cm kali $12$ cm. Pada setiap
jalur, tandai titik $3$ cm ke kanan di tepi atas, tarik garis dari sudut kiri
bawah ke titik itu, lalu potong sepanjang garis. Setiap jalur kini dua potongan
miring. Susun semua potongan menjadi satu jajar genjang yang rapat, lalu
rekatkan.

Lihat yang terjadi: bingkai luarnya $15$ cm kali $12$ cm, lebih lebar $3$ cm
daripada kertas asal, tetapi jajar genjang di dalamnya tetap berukuran alas
$12$ cm dan tinggi $12$ cm. Ukur keduanya dengan penggaris, tegak lurus, bukan
sepanjang sisi miring: hasilnya $144$ cm persegi, sama dengan luas kertas
sebelum dipotong. Kalau angkamu lain, yang salah hampir pasti tinggi.
\end{tangan}

\begin{teorema}[Perubahan variabel pada integral lipat dua]
\index{teorema perubahan variabel} Misalkan $T$ perubahan variabel yang sah
dari $D^{*}$ ke $D=T(D^{*})$, dan $f$ kontinu pada $D$. Maka
\[
\iint_{D}f(x,y)\,dx\,dy
=\iint_{D^{*}} f\bigl(T(u,v)\bigr)\,\bigl|\det J_{T}(u,v)\bigr|\,du\,dv .
\]
\end{teorema}

\begin{proof}
Kita buktikan gagasannya; bukti penuh ada di Spivak, bab \emph{Integration},
bagian perubahan variabel, dan di Rudin, bab \emph{Integration of Differential
Forms}.

Partisi $D^{*}$ menjadi petak kecil. Pada petak di sekitar $(u_{0},v_{0})$,
peta $T$ hampir linear: $T(u_{0}+h,v_{0}+k)\approx
T(u_{0},v_{0})+J_{T}(u_{0},v_{0})(h,k)$. Petak itu dipetakan menjadi jajar
genjang bersisi $h\,\partial T/\partial u$ dan $k\,\partial T/\partial v$,
yang luasnya dihitung dengan determinan (Bab 7):
\[
\mathrm{luas}
=\Bigl|h\,k\,\det\Bigl(\frac{\partial T}{\partial u},
\frac{\partial T}{\partial v}\Bigr)\Bigr|
=|h\,k|\,|\det J_{T}(u_{0},v_{0})| .
\]
Setiap petak menyumbang $f(T(u_{0},v_{0}))$ kali luas asalnya kali
$|\det J_{T}|$, dan jumlahnya adalah jumlah Riemann untuk integral di ruas
kanan; karena turunan $T$ kontinu pada himpunan kompak, galat aproksimasi linear
mengecil seragam bersama mesh.
\end{proof}

Nilai mutlak pada determinan bukan hiasan: determinan bisa negatif, luas tidak.
Tanda negatifnya menyimpan arah; pada integral lipat dua arah itu dibuang, tetapi
pada integral garis dan permukaan nanti tanda itu menentukan jawaban.

Sekarang determinan untuk dua sistem koordinat yang paling sering dipakai,
mulai dari polar.

\begin{proposisi}[Koordinat polar]\index{koordinat polar}
Dengan $x=r\cos\theta$, $y=r\sin\theta$, $r>0$, berlaku
\[
\det J=\begin{vmatrix}\cos\theta & -r\sin\theta\\ \sin\theta & r\cos\theta
\end{vmatrix}=r\cos^{2}\theta+r\sin^{2}\theta=r ,
\]
sehingga
$\iint_{D}f\,dx\,dy=\iint_{D^{*}}f(r\cos\theta,r\sin\theta)\,r\,dr\,d\theta$.
\end{proposisi}

Faktor $r$ itulah yang hilang pada program di pembukaan bab. Pada $r=1$
nilainya $1$, tetapi pada $r=0.1$ nilainya $0.1$: petak di dekat pusat
menyumbang sepuluh kali lebih kecil daripada dugaan yang mengabaikannya, dan
kesalahan totalnya bisa lebih dari dua kali.

\begin{proposisi}[Koordinat silinder dan bola]
\index{koordinat silinder}\index{koordinat bola} Untuk koordinat silinder
$x=r\cos\theta$, $y=r\sin\theta$, $z=z$, determinan Jacobiannya $r$. Untuk
koordinat bola
\[
x=\rho\sin\varphi\cos\theta,\quad
y=\rho\sin\varphi\sin\theta,\quad
z=\rho\cos\varphi ,
\]
dengan $\rho\ge0$, $0\le\varphi\le\pi$, $0\le\theta<2\pi$, determinan
Jacobiannya $\rho^{2}\sin\varphi$.
\end{proposisi}

\begin{proof}
Untuk silinder, kolom ketiga $(0,0,1)$ membuat determinan tinggal determinan
blok polar berukuran $2\times2$, yaitu $r$.

Untuk bola, hitung determinan $3\times3$ lewat ekspansi kofaktor. Dengan
$s_{\varphi}=\sin\varphi$, $c_{\varphi}=\cos\varphi$,
$s_{\theta}=\sin\theta$, $c_{\theta}=\cos\theta$, matriksnya
\[
J=\begin{pmatrix}
s_{\varphi}c_{\theta} & \rho c_{\varphi}c_{\theta} & -\rho s_{\varphi}s_{\theta}\\
s_{\varphi}s_{\theta} & \rho c_{\varphi}s_{\theta} & \rho s_{\varphi}c_{\theta}\\
c_{\varphi} & -\rho s_{\varphi} & 0
\end{pmatrix}.
\]
Entri baris ketiga kolom ketiga nol, jadi ekspansi kolom ketiga menyisakan dua
suku:
\[
\det J=\rho^{2}s_{\varphi}s_{\theta}^{2}+\rho^{2}s_{\varphi}c_{\theta}^{2}
=\rho^{2}s_{\varphi}\bigl(s_{\theta}^{2}+c_{\theta}^{2}\bigr)=\rho^{2}s_{\varphi}.
\]
\end{proof}

Faktor $\rho^{2}\sin\varphi$ nol di kutub $\varphi=0$ dan $\varphi=\pi$,
karena di sana seluruh bola $\rho$ runtuh menjadi satu titik dan peta berhenti
satu-satu; itu sebabnya kutub dikecualikan dari $D^{*}$. Ia paling besar di
ekuator, tempat perubahan kecil $\varphi$ memindahkan titik paling jauh.

\begin{resep}
\begin{enumerate}
\item Tulis $T$ eksplisit: $x=x(u,v)$, $y=y(u,v)$, dan periksa bahwa $T$
  satu-satu pada bagian dalam $D^{*}$.
\item Hitung empat turunan parsial, susun matriks $2\times2$, lalu hitung
  $\det J(u,v)$. Simpan tandanya; nilai mutlak dipasang nanti.
\item Terjemahkan batas: gambar $D$ di bidang $(x,y)$, tanyakan nilai $(u,v)$
  mana yang mendarat di sana, lalu tulis $D^{*}$.
\item Ganti seluruh $x,y$ pada $f$ dengan $x(u,v),y(u,v)$, kalikan
  $|\det J(u,v)|$, lalu tutup dengan $du\,dv$ sesuai urutan batas.
\end{enumerate}
Peringatan urutan: langkah 2 dan 3 tidak boleh ditukar. Determinan harus sudah
berbentuk $(u,v)$ sebelum batas dipasang. Kalau batas ditulis lebih dulu, faktor itu tertinggal dalam $x$ dan $y$ dan
integralnya tidak dapat dihitung.
\end{resep}

\kamus{determinan Jacobian}{Jacobian determinant}{Angka yang mengatakan berapa kali petak kecil meregang ketika koordinat diganti.}

\kamus{integral lipat dua}{double integral}{Menjumlahkan nilai pada petak kecil di sebuah permukaan, lalu mempersempitnya.}

\kamus{orientasi}{orientation}{Kesepakatan tentang arah putar mana yang dihitung positif.}

\section{Integral garis dan medan konservatif}
\tingkat{menengah}

Sekarang kita berhenti menjumlahkan pada permukaan dan mulai menjumlahkan
sepanjang jalan. Bayangkan sebuah benda kecil diseret sepanjang lintasan
melengkung oleh gaya yang arah dan besarnya berubah di setiap titik. Kerja
gaya itu adalah jumlah komponen gaya searah jalan dikali potongan panjang
jalan.

\begin{definisi}[Integral garis]\index{integral garis}
Misalkan $\gamma\colon[a,b]\to\mathbb{R}^{2}$ kurva dengan turunan kontinu
sepotong-potong, dan $F=(P,Q)$ medan vektor kontinu pada lintasannya.
\emph{Integral garis} $F$ sepanjang $\gamma$ adalah
\[
\int_{\gamma}F\cdot d\mathbf{r}
=\int_{a}^{b}\Bigl(P(\gamma(t))\,\gamma_{1}'(t)
+Q(\gamma(t))\,\gamma_{2}'(t)\Bigr)\,dt .
\]
Bila $\gamma$ tertutup, artinya $\gamma(a)=\gamma(b)$, integral itu disebut
\emph{sirkulasi} $F$ mengelilingi $\gamma$.
\end{definisi}

Rumus itu tidak memilih titik awal, hanya arah. Mengganti parameter dengan
laju berbeda tidak mengubah nilainya, karena substitusi (Bab 13) menyerap
faktor laju; membalik arah mengubah tanda. Jadi integral garis berorientasi,
dan itu perbedaan pertamanya dengan integral lipat dua.

Gambar~\ref{fig:sirkulasi} menunjukkan medan berputar seperti itu, dengan satu
kurva tertutup untuk mengukur berapa banyak putaran yang tertangkap.

\begin{figure}[htbp]
\centering
\begin{tikzpicture}[scale=1.45,font=\footnotesize]
  \draw[->,layergrey] (-2.85,0) -- (2.85,0) node[right] {$x$};
  \draw[->,layergrey] (0,-2.85) -- (0,2.85) node[above] {$y$};
  \foreach \x in {-2.4,-1.6,-0.8,0,0.8,1.6,2.4}{
    \foreach \y in {-2.4,-1.6,-0.8,0,0.8,1.6,2.4}{
      \draw[-{Latex[length=1.2mm]},layerteal!75]
        (\x,\y) -- ({\x-0.30*\y},{\y+0.30*\x});}}
  \draw[accent,very thick] (0,0) circle (1.75);
  \draw[-{Latex[length=2.8mm]},accent,very thick]
    (1.126,-1.341) -- (1.340,-1.126);
  \node[accent,anchor=north west,align=left] at (0.28,-1.86)
    {kurva tertutup, arah positif};
  \node[ink,align=center,fill=white,fill opacity=0.85,text opacity=1,inner sep=2pt] at (0,0.95)
    {$\displaystyle\oint_{\gamma}F\cdot d\mathbf{r}=2\pi$};
  \node[ink,font=\scriptsize,fill=white,fill opacity=0.85,text opacity=1,inner sep=1.5pt] at (0,0.42)
    {sedangkan $\operatorname{curl}F=2$ di tiap titik};
\end{tikzpicture}
\caption{Medan $F(x,y)=(-y,x)$ dengan panah di setiap titik kisi. Panahnya berputar
mengelilingi titik asal, dan besarnya bertambah sejauh titik itu dari pusat. Sirkulasi
mengelilingi lingkaran berjari-jari $1$ adalah $2\pi$, sama dengan luas $\pi$ dikali
$\operatorname{curl}F=2$: inilah Green untuk medan ini. Arah panah pada kurva adalah arah
positif, berlawanan arah jarum jam.}
\label{fig:sirkulasi}
\end{figure}

Kalau medan gaya berasal dari sebuah fungsi ketinggian, kerja hanya bergantung
pada titik awal dan titik akhir. Ini versi kurva dari teorema dasar kalkulus,
dan namanya mengikuti hal itu.

\begin{teorema}[Teorema dasar untuk integral garis]
\index{medan konservatif} Misalkan $f\colon U\to\mathbb{R}$ mempunyai turunan
parsial kontinu pada himpunan terbuka $U\subseteq\mathbb{R}^{2}$, dan $\gamma$
kurva dari $A$ ke $B$ di dalam $U$. Maka
\[
\int_{\gamma}\nabla f\cdot d\mathbf{r}=f(B)-f(A).
\]
\end{teorema}

\begin{proof}
Tulis $g(t)=f(\gamma(t))$. Menurut aturan rantai multivariabel (Bab 12),
\[
g'(t)=\nabla f(\gamma(t))\cdot\gamma'(t),
\]
asalkan $\gamma$ terdiferensialkan di $t$. Jadi integran pada integral garis
adalah turunan $g$, dan teorema dasar kalkulus (Bab 13) memberi
$\int_{a}^{b}g'(t)\,dt=g(b)-g(a)=f(B)-f(A)$. Untuk kurva yang hanya
terdiferensialkan sepotong-potong, jumlahkan pada tiap potongan; suku di titik
sambung saling menghapus.
\end{proof}

Medan berbentuk $\nabla f$ disebut \emph{konservatif}, dan $f$ potensialnya.
Akibat langsung teoremanya: pada medan konservatif, sirkulasi mengelilingi
lintasan tertutup selalu nol.

Kebalikannya juga benar.

\begin{teorema}[Kebalikan]
Misalkan $F$ kontinu pada himpunan terbuka terhubung
$U\subseteq\mathbb{R}^{2}$, dan integral garis $F$ tidak bergantung pada
lintasan di dalam $U$. Maka $F=\nabla f$ untuk suatu $f$, yaitu
\[
f(\mathbf{x})=\int_{\mathbf{x}_{0}}^{\mathbf{x}}F\cdot d\mathbf{r},
\]
dengan $\mathbf{x}_{0}\in U$ tetap.
\end{teorema}

\begin{proof}
Ambil $\mathbf{x}\in U$ dan $h$ kecil sehingga segmen dari $\mathbf{x}$ ke
$\mathbf{x}+h\mathbf{e}_{1}$ tinggal di dalam $U$. Karena lintasan tidak
penting, kita boleh menempuh jalan dari $\mathbf{x}_{0}$ ke $\mathbf{x}$ lalu
menambah segmen pendek itu, sehingga
\[
f(\mathbf{x}+h\mathbf{e}_{1})-f(\mathbf{x})
=\int_{0}^{1}F_{1}(\mathbf{x}+th\mathbf{e}_{1})\,h\,dt .
\]
Bagi dengan $h$ dan pakai kekontinuan: nilai rata-rata integran menuju
$F_{1}(\mathbf{x})$ ketika $h\to0$, jadi $\partial f/\partial x_{1}=F_{1}$,
dan argumen sama memberi $\partial f/\partial x_{2}=F_{2}$. Kedua turunan
parsial itu kontinu, maka $F=\nabla f$.
\end{proof}

Kalau diberi $F$, bagaimana tahu ia konservatif tanpa mencari potensial? Di
bidang ada uji cepat: $F=(P,Q)$ berpotensial hanya jika $\partial P/\partial
y=\partial Q/\partial x$, sebab kedua ruas adalah turunan kedua $f$ dan Clairaut
(Bab 12) menyamakannya. Arah sebaliknya menuntut domain tanpa lubang dan
dibuktikan lewat Green nanti. Medan pembukaan bab tandingannya: curl-nya nol,
tetapi ia bukan gradien apa pun.

\begin{contoh}[Dua cara menghitung kerja]
\index{potensial} Ambil $F(x,y)=(2xy,\ x^{2}+3y^{2})$ dan lintasan
$\gamma(t)=(t,2t^{2})$ untuk $0\le t\le1$, dari $(0,0)$ ke $(1,2)$.

Cara pertama, langsung. Di sini $dx=dt$ dan $dy=4t\,dt$, sehingga
\[
\int_{\gamma}F\cdot d\mathbf{r}
=\int_{0}^{1}\Bigl(2t\cdot 2t^{2}+(t^{2}+12t^{4})\cdot 4t\Bigr)dt
=\int_{0}^{1}\bigl(8t^{3}+48t^{5}\bigr)dt=2+8=10 .
\]
Cara kedua, lewat potensial. Periksa dulu bahwa $\partial P/\partial
y=2x=\partial Q/\partial x$, lalu tebak $f(x,y)=x^{2}y+y^{3}$. Turunannya
memang $\partial f/\partial x=2xy$ dan $\partial f/\partial y=x^{2}+3y^{2}$.
Karena itu
\[
\int_{\gamma}F\cdot d\mathbf{r}=f(1,2)-f(0,0)=(2+8)-0=10 .
\]
Kedua cara memberi $10$; kuadratur adaptif berselisih $1.78\times10^{-15}$.
\end{contoh}

\section{Teorema Green}
\tingkat{menengah}

Gagasan di kotak tentang ruangan dan tungku sekarang menjadi teorema: apa yang
terjadi di dalam sebuah daerah dapat dibaca dari tepinya.

Sebelum itu, satu kesepakatan arah.

\begin{definisi}[Orientasi positif]
\index{orientasi positif} Kurva tertutup sederhana $\gamma$ berorientasi
\emph{positif} bila daerah yang dibatasinya selalu berada di sebelah kiri
ketika kita menelusurinya. Untuk lingkaran, arah positif adalah berlawanan
arah jarum jam.
\end{definisi}

\begin{teorema}[Green]\index{teorema Green}
Misalkan $D$ daerah tertutup terbatas di $\mathbb{R}^{2}$ yang batasnya
$\partial D$ terdiri dari berhingga kurva tertutup sederhana yang mulus
sepotong-potong dan berorientasi positif. Misalkan $P,Q$ mempunyai turunan
parsial kontinu pada himpunan terbuka yang memuat $D$. Maka
\[
\oint_{\partial D}P\,dx+Q\,dy=\iint_{D}\Bigl(\frac{\partial Q}{\partial x}
-\frac{\partial P}{\partial y}\Bigr)dA .
\]
\end{teorema}

\begin{proof}
Cukup dibuktikan untuk daerah bertipe I sekaligus bertipe II, karena daerah
umum dipecah menjadi berhingga potongan semacam itu dan integralnya
dijumlahkan; sisi potongan yang bersebelahan menempuh arah berlawanan sehingga
saling menghapus.

Andaikan $D$ bertipe I dengan $g_{1},g_{2}$ mulus. Pada sisi tegak $dx=0$,
jadi sumbangannya nol dan
\[
\oint_{\partial D}P\,dx=\int_{a}^{b}P(x,g_{1}(x))\,dx+\int_{b}^{a}P(x,g_{2}(x))\,dx
=\int_{a}^{b}\bigl[P(x,g_{1})-P(x,g_{2})\bigr]dx .
\]
Menurut teorema dasar kalkulus,
$P(x,g_{1})-P(x,g_{2})=-\int_{g_{1}(x)}^{g_{2}(x)}\partial P/\partial y\,dy$,
sehingga
\[
\oint_{\partial D}P\,dx=-\int_{a}^{b}\!\!\int_{g_{1}(x)}^{g_{2}(x)}
\frac{\partial P}{\partial y}\,dy\,dx=-\iint_{D}\frac{\partial P}{\partial y}\,dA .
\]
Dengan cara yang sama pada $D$ bertipe II, tempat $dy=0$ pada sisi datar,
$\oint_{\partial D}Q\,dy=\iint_{D}\partial Q/\partial x\,dA$.
\end{proof}

\begin{contoh}[Green pada persegi satuan]
Ambil $P=x^{2}-y$, $Q=x+2y$, dan $D=[0,1]\times[0,1]$. Di sini $\partial
Q/\partial x-\partial P/\partial y=1-(-1)=2$, sehingga ruas kanan Green adalah
$2$ kali luas persegi, yaitu $2$.

Ruas kiri kita hitung pada empat sisi, berlawanan arah jarum jam:
\begin{align*}
\text{bawah }(y=0,\,dy=0):&\quad \int_{0}^{1}x^{2}\,dx=\tfrac13 ,\\
\text{kanan }(x=1,\,dx=0):&\quad \int_{0}^{1}(1+2y)\,dy=2 ,\\
\text{atas }(y=1,\,dy=0):&\quad \int_{1}^{0}(x^{2}-1)\,dx=\tfrac23 ,\\
\text{kiri }(x=0,\,dx=0):&\quad \int_{1}^{0}2y\,dy=-1 .
\end{align*}
Jumlahnya $\tfrac13+2+\tfrac23-1=2$, sama dengan ruas kanan. Kuadratur adaptif
untuk keempat sisi memberi $2.000000000000000$, dan integral lipat duanya juga
$2.000000000000000$.
\end{contoh}

Sekarang kesalahan tanda yang diminta. Ulangi hitungan yang sama searah jarum
jam. Setiap sisi ditempuh terbalik, jadi setiap suku berganti tanda dan
jumlahnya menjadi $-2$; ruas kanan Green tidak berubah karena daerahnya sama.
Jawaban yang seharusnya $2$ menjadi $-2$: bukan meleset sedikit, melainkan salah
tanda sepenuhnya, dengan selisih $4$ yang tidak dapat disalahkan pada
pembulatan.

Jadi orientasi bukan kesepakatan sopan: rumus Green hanya benar untuk satu
arah putar, dan arah itu harus diperiksa sebelum menghitung.

Dengan Green di tangan, arah sebaliknya dari uji curl di bagian sebelumnya
dapat dibuktikan.

\begin{akibat}[Uji curl di domain tanpa lubang]\index{uji curl}
Misalkan $U\subseteq\mathbb{R}^{2}$ terbuka dan \emph{terhubung sederhana},
artinya setiap kurva tertutup di $U$ dapat disusutkan ke satu titik di dalam
$U$. Misalkan $P,Q$ mempunyai turunan parsial kontinu pada $U$ dan $\partial
P/\partial y=\partial Q/\partial x$ di seluruh $U$. Maka $F=(P,Q)$ konservatif
pada $U$.
\end{akibat}

\begin{proof}
Ambil kurva tertutup sederhana $\gamma$ di $U$ dan daerah $D$ yang
dibatasinya. Karena $U$ terhubung sederhana, $D$ tinggal di dalam $U$, jadi
Green berlaku dan
\begin{align*}
\oint_{\gamma}P\,dx+Q\,dy
&=\iint_{D}\Bigl(\frac{\partial Q}{\partial x}-\frac{\partial P}{\partial y}\Bigr)dA\\
&=0 .
\end{align*}
Kurva tertutup yang memotong dirinya sendiri dipecah menjadi lengkung-lengkung
sederhana, dan tiap potongan memberi nol. Semua sirkulasi nol, maka integral
garis tidak bergantung lintasan, dan teorema kebalikan tadi memberi potensial.
\end{proof}

Syarat itu tidak bisa dibuang. Pada daerah berlubang, Green menuntut batas
dalam ikut dihitung; di situlah medan pembukaan bab lolos dari uji curl.

\section{Teorema Gauss dan Stokes}
\tingkat{lanjut}

Green berbicara tentang daerah di bidang dan tepinya. Di ruang berdimensi tiga
ada dua versi, keduanya memakai divergensi dan curl dari Bab 12.

\begin{definisi}[Divergensi dan curl]
\index{divergensi}\index{curl} Untuk medan $F=(F_{1},F_{2},F_{3})$ dengan
turunan parsial kontinu,
\[
\operatorname{div}F=\frac{\partial F_{1}}{\partial x}
+\frac{\partial F_{2}}{\partial y}+\frac{\partial F_{3}}{\partial z},
\qquad
\operatorname{curl}F=\Bigl(\frac{\partial F_{3}}{\partial y}
-\frac{\partial F_{2}}{\partial z},\
\frac{\partial F_{1}}{\partial z}-\frac{\partial F_{3}}{\partial x},\
\frac{\partial F_{2}}{\partial x}-\frac{\partial F_{1}}{\partial y}\Bigr).
\]
\end{definisi}

Divergensi mengukur berapa banyak medan lahir di sebuah titik: positif berarti
sumber, negatif penyerap. Curl mengukur putaran medan di sekitar titik itu.

\begin{teorema}[Gauss, atau teorema divergensi]\index{teorema Gauss}
Misalkan $E\subseteq\mathbb{R}^{3}$ daerah tertutup terbatas yang batasnya
$\partial E$ berupa permukaan mulus sepotong-potong, dan $\mathbf{n}$ vektor
normal satuan yang menunjuk keluar. Jika $F$ mempunyai turunan parsial kontinu
pada himpunan terbuka yang memuat $E$, maka
\[
\iiint_{E}\operatorname{div}F\,dV=\iint_{\partial E}F\cdot\mathbf{n}\,dA .
\]
\end{teorema}

\begin{proof}
Kita buktikan untuk kotak dulu.

Ambil $E=[a,b]\times[c,d]\times[e,f]$. Untuk suku ketiga,
\[
\iiint_{E}\frac{\partial F_{3}}{\partial z}\,dV
=\iint_{[a,b]\times[c,d]}\bigl[F_{3}(x,y,f)-F_{3}(x,y,e)\bigr]dx\,dy ,
\]
menurut Fubini dan teorema dasar kalkulus pada arah $z$. Ruas kanan itu tepat
fluks lewat sisi atas (normal $+\mathbf{k}$) dan sisi bawah (normal
$-\mathbf{k}$), karena di sana $\mathbf{n}\,dA=\pm\mathbf{k}\,dx\,dy$. Suku
lain memberi fluks lewat dua pasang sisi lainnya.

Untuk daerah umum, potong menjadi kotak-kotak kecil. Sisi dalam muncul dua kali
dengan normal berlawanan dan saling menghapus, sehingga yang tersisa hanya sisi
yang menyentuh batas luar; limit mesh menuju nol memberi pernyataan untuk batas
lengkung. Bukti penuh ada di Spivak, Teorema 5-8.
\end{proof}

\begin{contoh}[Gauss pada bola satuan]
Ambil $F(x,y,z)=(x^{3},y^{3},z^{3})$ dan $E$ bola satuan. Di sini
$\operatorname{div}F=3(x^{2}+y^{2}+z^{2})=3\rho^{2}$, sehingga dengan
determinan Jacobian bola
\[
\iiint_{E}3\rho^{2}\,dV=\int_{0}^{2\pi}\!\!\int_{0}^{\pi}\!\!\int_{0}^{1}
3\rho^{4}\sin\varphi\,d\rho\,d\varphi\,d\theta
=3\cdot2\pi\cdot2\cdot\frac15=\frac{12\pi}{5}.
\]
Ruas permukaan: pada bola satuan $\mathbf{n}=(x,y,z)$, jadi
$F\cdot\mathbf{n}=x^{4}+y^{4}+z^{4}$ dan
\[
\iint_{\partial E}(x^{4}+y^{4}+z^{4})\,dA=3\iint_{\partial E}z^{4}\,dA
=3\int_{0}^{2\pi}\!\!\int_{0}^{\pi}\cos^{4}\varphi\sin\varphi\,d\varphi\,d\theta
=3\cdot2\pi\cdot\frac25=\frac{12\pi}{5},
\]
dengan faktor $3$ dari simetri. Nilainya $7.5398223686155$, dan kuadratur pada
permukaan memberi angka yang sama sampai digit terakhir.
\end{contoh}

Teorema kedua berbicara tentang permukaan tidak tertutup yang batasnya kurva.
Rumusnya memerlukan aturan tangan kanan.

\begin{teorema}[Stokes]\index{teorema Stokes}
Misalkan $S$ permukaan mulus sepotong-potong berorientasi dengan normal
$\mathbf{n}$, dan $\partial S$ kurva batasnya dengan orientasi yang diberikan
aturan tangan kanan. Jika $F$ mempunyai turunan parsial kontinu pada himpunan
terbuka yang memuat $S$, maka
\[
\iint_{S}(\operatorname{curl}F)\cdot\mathbf{n}\,dA
=\oint_{\partial S}F\cdot d\mathbf{r}.
\]
\end{teorema}

\begin{proof}
Parametrisasi $S$ oleh $\mathbf{r}(u,v)$ atas daerah $D$ di bidang $(u,v)$,
dengan orientasi yang sesuai. Aturan rantai mengubah integrannya menjadi
turunan komponen $F$ terhadap $u$ dan $v$, dan setelah disusun ulang integral
permukaan itu menjadi
\[
\iint_{D}\Bigl(\frac{\partial \tilde{Q}}{\partial u}
-\frac{\partial \tilde{P}}{\partial v}\Bigr)du\,dv ,
\]
untuk dua fungsi $\tilde{P},\tilde{Q}$ yang dibentuk dari $F$ dan
$\mathbf{r}$. Itu tepat ruas kanan Green, jadi Green memberi integral garis
pada $\partial D$, yang tidak lain integral garis $F$ pada $\partial S$.
Aljabar penuhnya di Spivak, Teorema 5-9.
\end{proof}

Kalau $S$ datar di bidang $xy$, Stokes menjadi Green; kalau $S$ hanya sebuah
selang, ia menjadi teorema dasar kalkulus.

\begin{contoh}[Stokes pada belahan bola]
Ambil $F=(z,x,y)$, yang curl-nya $(1,1,1)$, dan $S$ belahan atas bola satuan
dengan normal keluar, sehingga batasnya lingkaran satuan di bidang $xy$
berlawanan jarum jam.

Ruas permukaan:
\[
\iint_{S}(1,1,1)\cdot\mathbf{n}\,dA=\iint_{S}(x+y+z)\,dA=\iint_{S}z\,dA
=\int_{0}^{2\pi}\!\!\int_{0}^{\pi/2}\cos\varphi\sin\varphi\,d\varphi\,d\theta
=2\pi\cdot\frac12=\pi ,
\]
karena suku $x$ dan $y$ saling menghapus oleh simetri belahan. Ruas batas, dengan
$\mathbf{r}(t)=(\cos t,\sin t,0)$:
\[
\oint_{\partial S}F\cdot d\mathbf{r}=\int_{0}^{2\pi}\cos^{2}t\,dt=\pi .
\]
Keduanya $3.1415926535898$; kuadratur memberi angka yang sama pada kedua ruas.
\end{contoh}

Sekarang kesalahan tanda pada Stokes. Balik normal permukaan menjadi ke dalam:
ruas permukaan menjadi $-\pi$, ruas batas tidak berubah selama arah putarnya
tetap, dan kesamaan gagal dengan selisih $2\pi$. Aturan tangan kanan yang
mengikat kedua ruas; membalik keduanya bersama memulihkan kesamaan.

\section{Bentuk diferensial: satu kalimat untuk semua}
\tingkat{lanjut}

Kita sudah melihat tiga teorema dengan tiga wajah. Ada bahasa yang
menuliskannya sebagai satu pernyataan, dan bahasanya bukan sekadar kerapian: ia
menjelaskan mengapa tanda orientasi muncul di mana-mana.

\begin{definisi}[Bentuk diferensial]
\index{bentuk diferensial} Sebuah \emph{bentuk-$k$} pada himpunan terbuka
$U\subseteq\mathbb{R}^{n}$ adalah jumlah suku berbentuk
$a(x)\,dx_{i_{1}}\wedge\cdots\wedge dx_{i_{k}}$, dengan $a$ fungsi mulus.
Perkalian baji $\wedge$ bersifat antisimetris:
\[
dx\wedge dy=-dy\wedge dx,\qquad dx\wedge dx=0 .
\]
Turunan luar $d$ bekerja pada koefisien dan menambah satu baji:
\[
d(a\,dx_{I})=da\wedge dx_{I},\qquad da=\sum_{j}\frac{\partial a}{\partial x_{j}}dx_{j}.
\]
\end{definisi}

Kalimat sehari-harinya: bentuk diferensial adalah benda yang siap
diintegralkan pada permukaan berdimensi tertentu, dan $d$ menaikkan dimensi satu
tingkat. Antisimetri membuat membalik orientasi otomatis membalik tanda;
kesulitan tanda sebelumnya tertanam di satu aturan itu.

\begin{teorema}[Stokes umum]
Misalkan $M$ manifold berorientasi berdimensi $k$ dengan batas, dan $\omega$
bentuk-$(k-1)$ dengan koefisien mulus. Maka
\[
\int_{\partial M}\omega=\int_{M}d\omega ,
\]
dengan orientasi $\partial M$ yang diinduksi oleh orientasi $M$.
\end{teorema}

Bukti lengkapnya dilewati: Spivak, \emph{Integration on Chains}, Teorema 4-13
untuk rantai, dan Teorema 5-5 pada \emph{Integration of Manifolds} untuk
manifold; juga Rudin, bab \emph{Integration of Differential Forms}. Di sana
Green, Gauss, dan Stokes klasik diturunkan sebagai kasus khusus.

Terjemahannya di ruang berdimensi tiga:

\begin{center}
\begin{tabular}{lll}
\hline
Bahasa bentuk & Medan vektor & Teorema \\
\hline
bentuk-0 $f$ & fungsi skalar & $df$ memberi gradien \\
bentuk-1 $P\,dx+Q\,dy+R\,dz$ & $F=(P,Q,R)$ & $d\omega$ memberi $\operatorname{curl}F$ \\
bentuk-2 $P\,dy\wedge dz+\cdots$ & $F=(P,Q,R)$ & $d\eta$ memberi $\operatorname{div}F$ \\
$\int_{\partial M}\omega=\int_{M}d\omega$ & & Green, Gauss, Stokes, FTC \\
\hline
\end{tabular}
\end{center}

Satu akibat kecil: karena $d(d\omega)=0$, turunan luar dua kali selalu nol.
Diterjemahkan, $\operatorname{curl}\nabla f=0$ dan
$\operatorname{div}\operatorname{curl}F=0$. Pemeriksaan simbolik untuk fungsi
dan medan umum memberi nol pada keenam komponen, jadi dua identitas yang tampak
kebetulan itu satu pernyataan.

\section{Kapan teorema ini gagal}
\tingkat{lanjut}

Semua teorema di bab ini kesamaan bersyarat. Bagian ini mengumpulkan empat
cara kesamaan itu rusak, dengan angka.

Kegagalan pertama adalah singularitas. Ambil $F(x,y)=(-y,x)/(x^{2}+y^{2})$
pada bidang tanpa titik asal. Hitung curl-nya:
\[
\frac{\partial Q}{\partial x}-\frac{\partial P}{\partial y}
=\frac{(x^{2}+y^{2})-2x^{2}}{(x^{2}+y^{2})^{2}}
+\frac{(x^{2}+y^{2})-2y^{2}}{(x^{2}+y^{2})^{2}}=0 .
\]
Beda hingga $h=10^{-5}$ di $(1,0)$ memberi $-2.01\times10^{-10}$, nol sampai
galat pemotongan. Tetapi sirkulasi
mengelilingi lingkaran satuan $2\pi=6.2831853072$, karena di sana
$F\cdot\gamma'=1$ di setiap titik. Jadi curl nol tidak berarti sirkulasi nol.

Yang dilanggar adalah syarat Green: $P$ dan $Q$ harus punya turunan parsial
kontinu pada himpunan terbuka yang memuat $D$, sedangkan keduanya tidak
terdefinisi di titik asal di dalam $D$, dan $D$ tidak terhubung sederhana. Pada cincin $1\le r\le2$ Green tetap berlaku dengan dua batas
berlawanan arah, dan integralnya $2\pi-2\pi=0$; nol itu benar karena di sana
medan berpotensial $\theta$, sudut polar.

Kegagalan kedua adalah pertukaran urutan tanpa keterintegralan. Ambil
\[
f(x,y)=\frac{x^{2}-y^{2}}{(x^{2}+y^{2})^{2}}
\]
pada persegi $(0,1)\times(0,1)$. Untuk $y>0$ tetap,
\[
\int_{0}^{1}f(x,y)\,dx=\Bigl[-\frac{x}{x^{2}+y^{2}}\Bigr]_{x=0}^{x=1}
=-\frac{1}{1+y^{2}},
\]
sehingga integral luar memberi $-\pi/4=-0.7853981634$; dengan urutan terbalik
hasilnya $+\pi/4$. Kedua integral berulang ada tetapi berbeda, jadi menurut
Fubini integral lipat duanya tidak ada. Di koordinat polar
$f=\cos2\theta/r^{2}$, dan pada seperempat cakram satuan
\[
\iint|f|\,dA\ge\int_{0}^{\pi/2}|\cos2\theta|\,d\theta
\int_{0}^{1}\frac{dr}{r}=\infty ,
\]
karena $\int_{0}^{1}dr/r$ menyimpang di $r=0$. Fubini tidak salah: syaratnya
tidak dipenuhi.

Kegagalan ketiga adalah batas yang tidak terurus. Gauss menuntut $E$ tertutup
dan terbatas dengan batas mulus sepotong-potong. Pada himpunan tak terbatas,
batasnya dapat lari ke tak hingga dan integral permukaannya berhenti berarti;
Spivak menaruh hal ini sebagai latihan setelah Teorema 5-5, yang gagal tanpa
kekompakan. Berhingga sudut masih aman, karena normalnya hanya melompat di
berhingga titik; batas bergerigi tanpa akhir, misalnya garis pantai, tidak.

Kegagalan keempat adalah permukaan yang tidak dapat diorientasikan. Pada pita
Mobius, berjalan mengelilingi pita membawa normal kembali ke sisi yang
berlawanan, sehingga $\mathbf{n}$ pada rumus Stokes tidak terdefinisi.
Teoremanya bukan salah; pernyataannya belum lengkap untuk benda semacam itu.

\begin{tantangan}
Ambil $T(u,v)=(u^{2},v)$ pada $D^{*}=(-2,2)\times(0,1)$. Peta ini tampak
melebarkan: titik $u=2$ mendarat di $x=4$, dan turunannya $2u$ tumbuh tanpa
batas ketika $u$ menjauh dari nol. Hitung luas $T(D^{*})$ dari gambarnya, lalu
hitung dengan rumus perubahan variabel. Kedua angka berbeda dengan faktor tepat
$2$. Jelaskan mengapa, dan sebutkan garis tempat determinan Jacobiannya nol.

Petunjuk 1: periksa apakah $T$ satu-satu. Ada dua nilai $u$ yang mendarat di
$x=1$.

Petunjuk 2: garis $u=0$ dilipat menjadi satu titik, sehingga tiap titik bayangan
terhitung dua kali oleh integral; $|\det J|=|2u|$ tidak tahu bahwa lipatan itu
hanya satu lembar.

Kalau mentok, kembali ke Definisi perubahan variabel yang sah dan periksa syarat
satu-satu.
\end{tantangan}

Satu peringatan penutup, yang tidak ada di buku teks. Kuadratur numerik tidak
pernah memeriksa syarat. Ia memberi $6.2831853072$ dengan tujuh digit benar
untuk medan yang melanggar syarat Green, dan tidak ada pesan galat. Bilangan yang rapi bukan bukti syaratnya dipenuhi; itu hanya bukti
aritmetikanya berjalan.

\section{Contoh hitung}
\tingkat{menengah}

Sekarang tiga hitungan yang menjadi tulang bab ini. Ketiganya bisa diperiksa
dengan Python, dan ketiganya pernah salah di tangan penulis sendiri.

\textbf{Satu integral, dua sistem koordinat.}
\nopagebreak[4]

Hitung
\[
I=\iint_{D}e^{-(x^{2}+y^{2})}\,dA,\qquad D=\{(x,y):x^{2}+y^{2}\le1\}.
\]
Di koordinat polar, $D$ menjadi $0\le r\le1$, $0\le\theta\le2\pi$, dan
$dA=r\,dr\,d\theta$. Integralnya terpisah menjadi dua faktor:
\[
I=\int_{0}^{2\pi}\!\!\int_{0}^{1}e^{-r^{2}}r\,dr\,d\theta
=2\pi\cdot\Bigl[-\frac12e^{-r^{2}}\Bigr]_{0}^{1}
=2\pi\cdot\frac{1-e^{-1}}{2}=\pi\bigl(1-e^{-1}\bigr).
\]
Nilainya $1.9858653038$. Antiturunan terhadap $r$ ada justru karena faktor
$r$: tanpa faktor itu, antiturunannya melibatkan fungsi galat.

Di koordinat Cartesius, daerahnya menuntut dua potongan:
\[
I=\int_{-1}^{1}\!\!\int_{-\sqrt{1-x^{2}}}^{\sqrt{1-x^{2}}}
e^{-(x^{2}+y^{2})}\,dy\,dx .
\]
Integral dalam tidak punya antiturunan elementer, jadi kita harus memakai
kuadratur. Hasil kedua bentuk, dengan toleransi yang sama:
\begin{align*}
\text{kartesius}&: 1.9858653037988705 \quad (11907\ \text{evaluasi}),\\
\text{polar}&: 1.9858653037988712 \quad (441).
\end{align*}
Galatnya $8.9\times10^{-16}$ dan $2.2\times10^{-16}$; selisih ketelitiannya
tidak penting, tetapi selisih biayanya $27$ kali.

Aturan titik tengah pada grid Cartesius yang sama rata menurunkan galat lebih
lambat: $1.141\times10^{-3}$ pada $50$ petak per sisi, dan
$4.667\times10^{-5}$ pada $400$. Setiap kali jumlah petak digandakan, galatnya
turun sekitar $2.9$ kali, bukan $4$ kali, karena tinggi kolom di dekat tepi
berperilaku seperti $\sqrt{1-x^{2}}$ sehingga galatnya seorde $h^{3/2}$.

\textbf{Volume bola dengan determinan Jacobian.}

Volume bola berjari-jari $R$ adalah
\[
V=\int_{0}^{2\pi}\!\!\int_{0}^{\pi}\!\!\int_{0}^{R}
\rho^{2}\sin\varphi\,d\rho\,d\varphi\,d\theta
=2\pi\cdot2\cdot\frac{R^{3}}{3}=\frac{4\pi R^{3}}{3}.
\]
Untuk $R=1$ nilainya $4.1887902048$. Kuadratur tiga dimensi pada bola memberi
\[
4.188790204786397 \quad(\text{selisih }6.2\times10^{-15}),
\]
dan simulasi Monte Carlo dengan $10^{6}$ titik pada kubus $[-1,1]^{3}$ memberi
$4.192824\pm0.004$.

Bandingkan dengan menghapus determinan Jacobian. Tanpa faktor
$\rho^{2}\sin\varphi$,
\[
\int_{0}^{2\pi}\!\!\int_{0}^{\pi}\!\!\int_{0}^{1}1\,d\rho\,d\varphi\,d\theta
=2\pi^{2}=19.7392088022,
\]
yaitu $3\pi/2$ kali terlalu besar, dan hasilnya tetap bilangan yang tampak
wajar untuk sebuah volume.



Tabel berikut mengumpulkan semua angka di bab ini.

\begin{center}
{\small
\begin{tabular}{llll}
\hline
Besaran & Eksak & Numerik & Galat \\
\hline
$\iint_{D}e^{-(x^{2}+y^{2})}dA$, polar & $1.985865$ & $1.985865$ & $2.2\times10^{-16}$ \\
$\iint_{D}e^{-(x^{2}+y^{2})}dA$, Cartesius & $1.985865$ & $1.985865$ & $8.9\times10^{-16}$ \\

$\int_{\gamma}F\cdot d\mathbf{r}$ & $10$ & $10.000000$ & $1.8\times10^{-15}$ \\
$\oint_{\partial D}$ Green, persegi & $2$ & $2.000000$ & $<10^{-15}$ \\
$\iint_{S}(\operatorname{curl}F)\cdot\mathbf{n}\,dA$ & $3.141593$ & $3.141593$ & $<10^{-13}$ \\
\hline
\end{tabular}}
\end{center}

\kertas{
Uji sendiri dalam tiga langkah. (1) Hitung ulang $I$ dengan tangan memakai
antiturunan $-\tfrac12e^{-r^{2}}$, lalu bandingkan dengan $1.9858653038$;
selisih di digit terakhir adalah galat pembulatan, bukan galat rumus. (2)
Jalankan \texttt{scipy.integrate.dblquad} pada kedua bentuk dan hitung berapa
kali fungsinya dipanggil; angka $11907$ dan $441$ akan muncul kembali, dengan
syarat toleransi keduanya sama. (3) Ganti $D$ menjadi cakram berjari-jari $2$
dan periksa bahwa jawaban eksaknya menjadi $\pi(1-e^{-4})$, yaitu $3.0841$.
Kalau langkah ketiga memberi angka lain, yang salah hampir pasti batas $r$,
bukan rumusnya. }

\section{Lab komputasi}
\tingkat{lanjut}

Kode berikut memeriksa lima hal sekaligus: biaya dua sistem koordinat, Green,
Gauss, Stokes, dan medan bersingularitas. Angka di tabel sebelumnya berasal
dari program ini.

\begin{lab}
\begin{verbatim}
import numpy as np
from scipy import integrate

# 1. Satu integral lipat dua, dua sistem koordinat
def fc(y, x):
    return np.exp(-(x*x + y*y))

cnt = [0]
def fc_count(y, x):
    cnt[0] += 1
    return fc(y, x)

kart = integrate.dblquad(fc_count, -1, 1,
                         lambda x: -np.sqrt(1-x*x),
                         lambda x: np.sqrt(1-x*x),
                         epsabs=1e-12, epsrel=1e-12)
nkart = cnt[0]
cnt[0] = 0

def fp(r, th):
    cnt[0] += 1
    return np.exp(-r*r)*r

pol = integrate.dblquad(fp, 0, 2*np.pi, 0, 1,
                        epsabs=1e-12, epsrel=1e-12)
nilai = np.pi*(1 - np.exp(-1))
print("kartesius %.16f galat %.2e eval %d"
      % (kart[0], abs(kart[0]-nilai), nkart))
print("polar     %.16f galat %.2e eval %d"
      % (pol[0], abs(pol[0]-nilai), cnt[0]))
print("eksak     %.16f" % nilai)

# 2. Green: (x^2-y)dx + (x+2y)dy pada persegi satuan
P = lambda x, y: x*x - y
Q = lambda x, y: x + 2*y
tepi = [integrate.quad(lambda t: P(t, 0.0), 0, 1)[0],
        integrate.quad(lambda t: Q(1.0, t), 0, 1)[0],
        integrate.quad(lambda t: P(t, 1.0), 1, 0)[0],
        integrate.quad(lambda t: Q(0.0, t), 1, 0)[0]]
ganda = integrate.dblquad(lambda y, x: 2.0, 0, 1, 0, 1)[0]
print("Green tepi %.15f  ganda %.15f" % (sum(tepi), ganda))

# 3. Gauss: F=(x^3,y^3,z^3) pada bola satuan
def fluks(phi, th):
    sp, cp = np.sin(phi), np.cos(phi)
    x, y, z = sp*np.cos(th), sp*np.sin(th), cp
    return (x**4 + y**4 + z**4)*sp

F = integrate.dblquad(fluks, 0, 2*np.pi, 0, np.pi,
                      epsabs=1e-11, epsrel=1e-11)[0]
print("Gauss fluks %.13f  12pi/5 %.13f" % (F, 12*np.pi/5))

# 4. Stokes: F=(z,x,y) pada belahan atas bola satuan
def curl_n(phi, th):
    sp, cp = np.sin(phi), np.cos(phi)
    x, y, z = sp*np.cos(th), sp*np.sin(th), cp
    return (x + y + z)*sp

S = integrate.dblquad(curl_n, 0, 2*np.pi, 0, np.pi/2)[0]
L = integrate.quad(lambda t: np.cos(t)**2, 0, 2*np.pi)[0]
print("Stokes permukaan %.13f  batas %.13f" % (S, L))

# 5. Medan dengan singularitas: F=(-y,x)/(x^2+y^2)
def lingkar(t):
    x, y = np.cos(t), np.sin(t)
    return (-y*(-np.sin(t)) + x*np.cos(t))/(x*x + y*y)

def Pg(x, y): return -y/(x*x + y*y)
def Qg(x, y): return x/(x*x + y*y)

h = 1e-5
curl_01 = ((Qg(1+h, 0) - Qg(1-h, 0))/(2*h)
           - (Pg(1, h) - Pg(1, -h))/(2*h))
C = integrate.quad(lingkar, 0, 2*np.pi)[0]
print("lingkaran %.13f  2pi %.13f" % (C, 2*np.pi))
print("curl(1,0) %.2e" % curl_01)
\end{verbatim}
\end{lab}

Keluaran sebenarnya, apa adanya:

\begin{verbatim}
kartesius 1.9858653037988705 galat 8.88e-16 eval 11907
polar     1.9858653037988712 galat 2.22e-16 eval 441
eksak     1.9858653037988714
Green tepi 2.000000000000000  ganda 2.000000000000000
Gauss fluks 7.5398223686155  12pi/5 7.5398223686155
Stokes permukaan 3.1415926535898  batas 3.1415926535898
lingkaran 6.2831853071796  2pi 6.2831853071796
curl(1,0) -2.01e-10
\end{verbatim}

Empat hal patut dibaca. Pertama, dua baris awal sama sampai digit ke-15 dengan
biaya berbeda $27$ kali: koordinat yang tepat tidak mengubah jawaban, hanya
harganya. Kedua, Green cocok sampai digit terakhir, tetapi itu
hanya bukti implementasinya konsisten dengan rumusnya. Ketiga, Gauss dan Stokes
dihitung dari sisi berbeda dan tetap bertemu; di situlah verifikasi punya gigi.
Keempat, kuadratur memberi $6.2831853071796$ untuk medan yang melanggar syarat
Green, tanpa satu pun pesan galat.

\begin{untukdosen}
Urutan mengajar yang saya pakai: kerapatan pelat dan program yang lupa faktor
$r$, lalu Fubini, lalu determinan Jacobian sebagai luas petak yang meregang,
lalu integral garis dan teorema dasarnya, baru Green. Green sebaiknya dibuktikan
penuh untuk daerah bertipe I di papan, karena bukti itu hanya memakai dua alat
yang sudah mereka punya. Gauss dan Stokes cukup dibuktikan untuk kotak dan
permukaan berparameter.

Dua miskonsepsi yang paling sering muncul. Pertama, mahasiswa menganggap
$\partial P/\partial y=\partial Q/\partial x$ cukup untuk konservatif, karena
mereka hanya menguji medan pada domain tanpa lubang. Medan
$(-y,x)/(x^{2}+y^{2})$ harus muncul di papan, bukan di daftar bacaan. Kedua,
mereka memperlakukan determinan Jacobian sebagai formalitas notasi dan menulis
$r$ atau $\rho^{2}\sin\varphi$ sebagai hafalan tanpa memeriksa batas.
Bandingkan $4\pi/3$ dengan $2\pi^{2}$; rasio $4.712$ kali itu lebih mengesankan
daripada seratus kali peringatan.

Satu aktivitas kelas: empat kelompok, masing-masing menghitung satu ruas dari
Green, Gauss, dan Stokes untuk medan yang sama, lalu mencocokkan angka. Kelompok
yang salah tanda ketahuan dalam lima menit, dan itu lebih baik daripada diberi
tahu. Tanda kesiapan lanjut: mahasiswa mampu menjelaskan mengapa orientasi
mengubah tanda tanpa melihat rumus.
\end{untukdosen}

\section{Ringkasan}
\begin{ringkasan}
\begin{itemize}
\item Integral lipat dua didefinisikan lewat jumlah atas dan bawah pada petak;
Fubini mengizinkan urutan pengintegralan ditukar bila fungsinya terintegralkan.
\item Mengganti koordinat mengalikan luas petak dengan $|\det J|$; polar memberi
$r$, bola memberi $\rho^{2}\sin\varphi$. Tanpa faktor itu hasilnya rapi tetapi
salah.
\item Pada medan konservatif, integral garis hanya bergantung pada titik ujung.
\item Green, Gauss, dan Stokes menyatakan satu hal: integral turunan pada daerah
sama dengan integral medan pada batasnya, dan orientasi menentukan tanda.
\item Semua teorema itu bersyarat: singularitas, batas tak terurus, permukaan tak
terorientasi, dan fungsi tak terintegralkan membuatnya gagal tanpa peringatan.
\end{itemize}
\end{ringkasan}

\section{Latihan}
\begin{latihan}
\begin{enumerate}
\item Hitung $\iint_{D}(x^{2}+y^{2})\,dA$ pada cakram satuan dengan Cartesius dan polar, periksa dengan \texttt{dblquad}, dan sebutkan jumlah evaluasi tiap cara.

\item Buktikan bahwa bila $T(u,v)=A(u,v)$ linear dan $\det A\ne0$, rumus
perubahan variabel menjadi $|\det A|\iint_{D}f\circ T\,dA$.

\item Ulangi contoh Green pada persegi dengan arah jarum jam: tunjukkan hasilnya
$-2$ dan sebutkan sisi mana menyumbang tanda mana.

\item Hitung volume elipsoid $x^{2}/a^{2}+y^{2}/b^{2}+z^{2}/c^{2}\le1$ dengan
perubahan variabel $x=au$, $y=bv$, $z=cw$, lalu periksa rumusnya pada
$a=b=c=1$.

\item Nilai klaim ini: benar, salah, atau belum cukup bukti. ``Karena
$\operatorname{curl}F=0$ di setiap titik tempat $F$ terdefinisi, medan
$F=(-y,x)/(x^{2}+y^{2})$ konservatif pada
$\mathbb{R}^{2}\setminus\{(0,0)\}$.'' Sertakan alasan dan satu angka
pendukung.

\item Buktikan bahwa integral garis tidak bergantung lintasan jika dan hanya jika
sirkulasi setiap kurva tertutup nol, dengan memakai dua lintasan dari $A$ ke
$B$ yang membentuk satu kurva tertutup.

\item Hitung kedua ruas Stokes untuk $F=(-y,x,z)$ pada paraboloid
$z=1-x^{2}-y^{2}$, $z\ge0$, dengan normal mengarah ke atas, lalu periksa
dengan kuadratur.

\item Hitung kedua ruas Gauss untuk $F=(x,y,z)/(x^{2}+y^{2}+z^{2})^{3/2}$ pada bola
satuan. Jelaskan mengapa hasilnya tidak bertentangan dengan teorema, dan
syarat mana yang dilanggar. (Petunjuk: hitung $\operatorname{div}F$ lebih dulu.)

\item Tunjukkan bahwa determinan Jacobian koordinat bola nol pada $\varphi=0$ dan
$\varphi=\pi$, lalu jelaskan mengapa peta itu tetap boleh dipakai pada
$0<\varphi<\pi$. (Petunjuk: periksa di mana peta berhenti satu-satu.)
\end{enumerate}
\end{latihan}

\begin{untukpeneliti}
Yang benar-benar terbukti di bab ini, dengan bukti yang ditulis: lemma daerah
di antara dua grafik, teorema dasar untuk integral garis, kebalikannya lewat
potensial integral, Green untuk daerah bertipe I dan II, akibat uji curl pada
domain terhubung sederhana, dan sketsa bukti Gauss lewat kotak serta Stokes
lewat parametrisasi. Yang saya kutip tanpa bukti: Fubini (Spivak, bab
\emph{Integration}, Teorema 3-10), perubahan variabel penuh (Spivak bab yang
sama; Rudin, bab \emph{Integration of Differential Forms}), Gauss (Spivak,
Teorema 5-8), dan Stokes (Spivak, Teorema 5-9); Stokes umum adalah Teorema
4-13 dan 5-5.

Yang dipopulerkan secara berlebihan adalah ungkapan ``curl nol berarti
konservatif''. Yang benar: curl nol berarti bentuk tertutup, dan bentuk tertutup
eksak hanya bila grup kohomologi pertama domainnya nol. Untuk $\mathbb{R}^{3}$
itu berlaku, untuk $\mathbb{R}^{2}\setminus\{0\}$ tidak, dan selisihnya
diukur oleh integral pada lingkaran, yaitu $2\pi$ di sini.

Pertanyaan terbuka yang sungguh terbuka bagi saya: dapatkah kuadratur memberi
sertifikat bahwa syarat teoremanya dipenuhi, bukan hanya taksiran nilainya?
Keluaran di lab menunjukkan bahwa jawaban numerik yang benar sampai lima belas
digit dapat berasal dari medan yang melanggar syarat, tanpa tanda bahaya. Untuk
medan yang hanya diketahui lewat sampel pada grid, tidak ada cara umum
memastikan tidak ada singularitas di antara titik sampel.

Batas keberlakuan: semua fungsi di sini bernilai real dan mulus
sepotong-sepotong; semua angka dari Python 3.14.4, NumPy 2.5.2, SciPy 1.18.0
dengan toleransi $10^{-12}$ kecuali disebut lain, dan permukaan bola
diparametrisasi dengan $\mathbf{n}$ mengarah keluar.
\end{untukpeneliti}

\peta{Bab 15 pada Stewart, \emph{Calculus: Early Transcendentals}, bab ``Multiple Integrals'' (15.1--15.3 dan 15.6--15.9) dan bab ``Vector Calculus'' (16.2
integral garis, 16.3 teorema dasar, 16.4 Green, 16.5 curl
dan divergensi, 16.8 Stokes, 16.9 Gauss); pada Spivak, \emph{Calculus on
Manifolds}, bab ``Integration'' (Fubini, Teorema 3-10, perubahan variabel),
bab ``Integration on Chains'' dan bab ``Integration of Manifolds'' (Teorema
5-5, 5-8, 5-9); pada Rudin, \emph{Principles of Mathematical Analysis}, bab
``Functions of Several Variables'' dan ``Integration of Differential Forms''; pada Active Calculus, bab 11 dan 12 untuk urutan yang lebih bertahap.
Basisnya bab \emph{Integral, Deret Fungsi, dan Pertukaran Limit} (bagian integral)
dan bab \emph{Kekontinuan dan Turunan} (turunan parsial, Jacobian) pada edisi pertama
buku ini. Bab ini memakai integral Riemann dan teorema dasar kalkulus
dari Bab 13, kekontinuan seragam dan himpunan kompak dari Bab 11, determinan
dari Bab 7, serta gradien dan matriks Jacobian dari Bab 12, semuanya dirujuk
dan tidak diulang.}

\chapter{Persamaan Diferensial dan Sistem Dinamik}

Jam bandul di ruang tamu berayun dengan amplitudo enam puluh derajat, dan panjang
lengannya satu meter. Buku pelajaran memberi periode
\[
T=2\pi\sqrt{L/g}=2.006409\ \text{detik},
\]
dan rumus itu memang benar untuk ayunan yang sangat kecil. Diukur pada amplitudo
enam puluh derajat, bandul yang sama memerlukan $2{,}153242$ detik untuk satu ayunan
penuh: $7{,}32$ persen lebih lama. Jam yang disetel memakai rumus kecil itu
ketinggalan $5892$ detik, sekitar sembilan puluh delapan menit, setiap hari.

Rumus itu bukan salah hitung. Ia jawaban tepat untuk persamaan yang berbeda.
Ayunan sebenarnya memenuhi
\[
\frac{d^{2}\theta}{dt^{2}}=-\frac{g}{L}\sin\theta ,
\]
sedangkan rumus periode kecil berasal dari penggantian $\sin\theta$ dengan
$\theta$. Pada $\theta=60^{\circ}$ penggantian itu mengubah $0{,}866025$ menjadi
$1{,}047198$, selisih tujuh belas persen di titik terjauh, dan kesalahan periode
yang tertinggal $7{,}32$ persen. 

Bab 12 mengajari kita menurunkan, Bab 14 deret fungsi, Bab 15 medan vektor.
Tidak satu pun menjawab pertanyaan di sini: bila laju perubahan diketahui,
bagaimana mencari fungsinya, apakah ia tunggal, dan apa yang terjadi dalam jangka
panjang. Bab ini mengisi keempat kekosongan itu: cara menyelesaikan, syarat
keberadaan dan ketunggalan, bidang fase, dan uji kestabilan.

\begin{untukanak}
Bayangkan bak air di kamar mandi. Keran di atas mengalir tetap, lubang pembuangan
di bawah mengalir keluar, dan semakin tinggi air, semakin deras yang keluar.
Kalau kita membuka keran lalu pergi, tinggi air besok tidak dapat dihitung hanya
dari besar keran: yang menentukan adalah tinggi air sekarang. Air yang tinggi
mendorong keluar lebih cepat, kenaikannya melambat, dan akhirnya berhenti ketika
air masuk sama dengan air keluar.

Benda lain bekerja dengan cara yang sama. Kopi panas mendingin lebih cepat ketika
masih sangat panas. Kuman berkembang biak lebih cepat ketika jumlahnya banyak,
lalu melambat ketika makanan habis. Yang kita tahu adalah aturan perubahan, bukan
hasil akhirnya.

Sekarang pertanyaan yang membuat pekerjaan ini sulit. Dua bak sama, keran sama,
lubang sama, tetapi satu mulai kosong dan satu mulai setengah penuh: apakah
akhirnya tingginya sama? Kadang ya, kadang tidak. Dan ada lubang yang jahat,
tersumbat sebagian lalu terbuka tiba-tiba: dua cerita berbeda bisa keluar dari
keadaan awal yang sama.

Di bab ini kita akan tahu mengapa aturan perubahan kadang cukup menentukan
seluruh masa depan, dan kadang tidak.
\end{untukanak}

\section{Laju yang bergantung pada dirinya sendiri}
\tingkat{dasar}

Persamaan diferensial menyamakan laju perubahan sebuah fungsi dengan sesuatu
yang memuat fungsi itu sendiri. Menyelesaikannya berarti mencari fungsi yang
memenuhi kalimat itu pada seluruh sebuah selang. Turunan di Bab 12 dihitung di
satu titik; solusi di sini adalah fungsi pada selang, dan perbedaan itu yang
paling sering terlewat.

\begin{definisi}[Persamaan diferensial biasa]\index{persamaan diferensial}
Persamaan
\[
F\bigl(t,y,y',y'',\dots,y^{(n)}\bigr)=0
\]
disebut \emph{persamaan diferensial biasa} (PDB) orde $n$ untuk fungsi $y$ dari
peubah $t$ bila $F$ memuat turunan ke-$n$ dan tidak memuat turunan yang lebih
tinggi. PDB itu disebut \emph{linear} bila $F$ berbentuk
\[
a_{n}(t)\,y^{(n)}+\cdots+a_{1}(t)\,y'+a_{0}(t)\,y=b(t),
\]
dan disebut \emph{otonom} bila $F$ tidak memuat $t$ secara eksplisit.
\end{definisi}

 Semua contoh di
bab ini orde satu atau orde dua, karena keduanya sudah memuat seluruh gagasan
penting.

\begin{definisi}[Solusi dan masalah nilai awal]\index{masalah nilai awal}
Fungsi $y$ disebut \emph{solusi} PDB pada selang $I$ bila $y$ mempunyai turunan
sampai orde $n$ pada $I$ dan memenuhi persamaan di setiap titik $I$. Bila
ditambahkan syarat $y(t_{0})=y_{0}$, gabungan keduanya disebut \emph{masalah
nilai awal}, dan $y_{0}$ disebut nilai awal.
\end{definisi}

Satu contoh untuk memastikan definisinya bekerja. Ambil model populasi
\[
y'=y(1-y),\qquad y(0)=\tfrac12 ,
\]
dengan $y$ menyatakan bagian populasi terhadap daya tampung. Calon solusinya
\[
y(t)=\frac{1}{1+e^{-t}} .
\]
Turunannya $y'(t)=e^{-t}/(1+e^{-t})^{2}$, dan ruas kanan
$y(1-y)=e^{-t}/(1+e^{-t})^{2}$ sama persis. Nilai awalnya $y(0)=1/2$ juga cocok.
Jadi fungsinya solusi, dan pemeriksaan itu hanya memakai aturan turunan hasil
bagi dari Bab 12.

Yang perlu diperhatikan: $y(t)=1$ dan $y(t)=0$ juga solusi, keduanya konstan,
dan keduanya muncul saat kita menanyakan kapan ruas kanan nol. Jadi satu
persamaan dapat punya banyak solusi; nilai awal memilih satu di antaranya.
Koleksi semua solusi disebut \emph{solusi umum}, dan ia memuat satu konstanta
untuk setiap orde.

Selang tempat solusi hidup adalah bagian dari jawaban. Ambil
\[
y'=y^{2},\qquad y(0)=1 .
\]
Pemisahan variabel memberi $y(t)=1/(1-t)$, dan rumus itu hanya masuk akal pada
$t<1$. Di $t=0{,}9$ nilainya $10$; di $t=0{,}99$ nilainya $100$; di $t=0{,}999$ nilainya
$1000$. Tidak ada bilangan berhingga di $t=1$, dan tidak ada solusi yang
melanjutkan fungsi itu melewatinya. Solusinya \emph{meledak} dalam waktu
berhingga, dan program yang melaporkan nilai di $t=2$ tidak menemukan solusi; ia
menemukan pembulatan.

\kamus{persamaan diferensial}{differential equation}{Kalimat yang mengikat laju perubahan sebuah besaran dengan besar besaran itu sendiri.}

\kamus{solusi umum}{general solution}{Keluarga semua fungsi yang memenuhi persamaan, biasanya memuat satu konstanta bebas per orde.}

\begin{bacakeras}
--- Kalau bola ini kulepas di lereng, ia akan bergerak ke bawah. Berapa cepat?
--- Tergantung di mana bola itu sekarang. Di dekat puncak, pelan. Di tengah
lereng, cepat.
--- Jadi kecepatannya bukan satu angka tetap?
--- Bukan. Ia aturan: di tiap titik, kecepatan ditentukan oleh titik itu.
--- Kalau aturannya sudah tahu segalanya, apakah jalannya sudah tertentu?
--- Belum. Aturan memberi arah di setiap titik; jalan lahir dari titik awal yang
kau pilih.
--- Dua bola, dua titik awal, satu aturan?
--- Satu aturan, dua jalan.
--- Kalau dua jalan bertemu di titik yang sama, bagaimana kita tahu mana yang
datang dari mana?
--- Kita tidak bisa. Karena itu bab ini menuntut satu titik awal untuk satu
jalan.

Ulangi sekali, dengan suara biasa: aturan sama, titik awal berbeda, jalan berbeda.
\end{bacakeras}

\section{Dua keluarga yang dapat diselesaikan dengan tangan}
\tingkat{dasar}

Tidak semua persamaan diferensial punya rumus penyelesaian, tetapi dua keluarga
besar punya, dan keduanya muncul di hampir setiap penerapan tingkat pertama.
Keluarga pertama memisahkan peubah; keluarga kedua linear orde satu.

\subsection{Pemisahan variabel}

Bila persamaan dapat ditulis $y'=f(t)g(y)$, kita memindahkan semua $y$ ke satu
sisi dan semua $t$ ke sisi lain:
\[
\frac{dy}{g(y)}=f(t)\,dt,\qquad \int\frac{dy}{g(y)}=\int f(t)\,dt .
\]
Langkah itu bukan trik notasi, melainkan aturan rantai yang dibalik: bila $H$
antiturunan $1/g$, maka $\frac{d}{dt}H(y(t))=H'(y)y'=y'/g(y)$, sehingga
persamaan asal menjadi $\frac{d}{dt}H(y(t))=f(t)$ dan Bab 13 memberi
$H(y)=\int f\,dt+C$. Hasilnya biasanya implisit, dan itu jawaban yang sah.

Model logistik di atas dipakai sepanjang bab. Dengan
$g(y)=y(1-y)$ dan pecahan parsial,
\[
\int\frac{dy}{y(1-y)}=\int\Bigl(\frac{1}{y}+\frac{1}{1-y}\Bigr)dy
=\ln\Bigl|\frac{y}{1-y}\Bigr| ,
\]
maka $\ln|y/(1-y)|=t+C$ dan
\[
y(t)=\frac{1}{1+Ae^{-t}},\qquad A=\frac{1-y_{0}}{y_{0}} .
\]
Untuk $y_{0}=1/2$ diperoleh $A=1$, yaitu rumus yang kita periksa tadi. Perhatikan
bahwa $A$ ditentukan oleh nilai awal, dan bahwa $A=0$ memberi solusi konstan $y=1$
yang tidak tercakup pembagian di atas: saat memisahkan variabel kita membagi
dengan $g(y_{0})$, dan pembagian itu hanya sah bila $g(y_{0})\ne0$.

\subsection{Persamaan linear dan faktor integrasi}

Bentuk baku persamaan linear orde satu adalah
\[
y'+p(t)\,y=q(t).
\]
Ruas kiri hampir merupakan turunan sebuah hasil kali, hanya kurang satu faktor.
Faktor yang hilang itulah yang dicari teorema berikut.

\begin{teorema}[Faktor integrasi]\index{faktor integrasi}
Misalkan $p$ dan $q$ kontinu pada selang buka $I$ yang memuat $t_{0}$, dan
misalkan
\[
\mu(t)=\exp\Bigl(\int_{t_{0}}^{t}p(s)\,ds\Bigr).
\]
Maka masalah nilai awal $y'+p(t)y=q(t)$, $y(t_{0})=y_{0}$ mempunyai tepat satu
solusi pada $I$, yaitu
\[
y(t)=\frac{1}{\mu(t)}\Bigl(y_{0}+\int_{t_{0}}^{t}\mu(s)\,q(s)\,ds\Bigr).
\]
\end{teorema}

\begin{proof}
Karena $p$ kontinu, $\mu$ terdefinisi, positif, dan terdiferensialkan pada $I$,
dengan $\mu'=p\mu$ menurut teorema dasar kalkulus (Bab 13). Kalikan persamaan
dengan $\mu$:
\[
\mu y'+\mu p\,y=\mu q
\quad\Longleftrightarrow\quad
(\mu y)'=\mu q ,
\]
karena aturan turunan hasil kali memberi $(\mu y)'=\mu'y+\mu y'=\mu(py+y')$.
Integralkan kedua ruas dari $t_{0}$ sampai $t$ dan pakai teorema dasar kalkulus:
\[
\mu(t)y(t)-\mu(t_{0})y(t_{0})=\int_{t_{0}}^{t}\mu(s)q(s)\,ds .
\]
Karena $\mu(t_{0})=1$, pembagian dengan $\mu(t)>0$ memberi rumus yang diminta.
Ketunggalan: bila $y_{1},y_{2}$ dua solusi, selisihnya $w=y_{1}-y_{2}$ memenuhi
$w'+pw=0$ dengan $w(t_{0})=0$; mengalikan dengan $\mu$ memberi $(\mu w)'=0$,
jadi $\mu w$ konstan, dan karena nilainya nol di $t_{0}$ ia nol di mana-mana.
\end{proof}

Bukti itu tidak menebak apa pun: $\mu$ dicari dengan menuntut $\mu'=p\mu$,
persamaan yang lebih mudah daripada asalnya. Pola ini berulang: satu persoalan
diselesaikan dengan mengubahnya menjadi persoalan yang lebih sederhana.

\begin{contoh}[Satu persamaan linear, dua pemeriksaan]
Selesaikan $y'+2y=4t$ dengan $y(0)=1$. Di sini $p=2$ dan $q=4t$, sehingga
$\mu=e^{2t}$ dan
\[
(e^{2t}y)'=4t\,e^{2t}.
\]
Integral parsial memberi
\[
\int 4t\,e^{2t}\,dt=\Bigl(2t-1\Bigr)e^{2t}+C ,
\]
sebab turunan $(2t-1)e^{2t}$ adalah $2e^{2t}+2(2t-1)e^{2t}=(4t)e^{2t}$. Jadi
\[
y(t)=2t-1+Ce^{-2t},
\]
dan syarat $y(0)=1$ memberi $C=2$, yaitu $y(t)=2t-1+2e^{-2t}$.

Pemeriksaan: substitusi memberi $y'+2y=4t$, dan $y(1)=1+2e^{-2}=1{,}2706705665$;
integrator adaptif dengan toleransi $10^{-12}$ memberi $1{,}270670566473$.
\end{contoh}

\subsection{Rangkaian RC: contoh nyata yang paling sederhana}

Sebuah sumber tegangan $V_{s}$ dihubungkan seri dengan resistor $R$ dan
kapasitor $C$. Hukum Kirchhoff tentang tegangan memberi $V_{s}=Ri+V_{C}$, dan
hukum kapasitor memberi $i=C\,dV_{C}/dt$. Keduanya memberi
\[
RC\,\frac{dV_{C}}{dt}+V_{C}=V_{s},
\qquad\text{atau}\qquad
\frac{dV_{C}}{dt}+\frac{1}{RC}V_{C}=\frac{V_{s}}{RC}.
\]
Ini persamaan linear dengan $p=1/(RC)$ konstan. Menurut teorema tadi, dengan
$V_{C}(0)=0$,
\[
V_{C}(t)=V_{s}\bigl(1-e^{-t/\tau}\bigr),\qquad \tau=RC .
\]
Angka $\tau$ disebut tetapan waktu. Untuk $R=1\,\mathrm{k}\Omega$ dan
$C=1\,\mu\mathrm{F}$, $\tau=10^{-3}$ detik. Dengan $V_{s}=5$ volt, tegangan
pada $t=\tau$ adalah $3{,}160603$ volt dan pada $t=5\tau$ adalah $4{,}966310$ volt,
yaitu $99{,}33$ persen dari $5$ volt. Untuk mencapai $90$ persen dibutuhkan
$t=\tau\ln10=2{,}302585$ milidetik.

\begin{resep}
\begin{enumerate}
\item Tulis persamaan dalam bentuk baku $y'+p(t)y=q(t)$: bagi dengan koefisien
  $y'$ sehingga koefisiennya tepat satu.
\item Hitung $\mu(t)=\exp\bigl(\int p(t)\,dt\bigr)$, tanpa konstanta integrasi.
\item Kalikan seluruh persamaan dengan $\mu$, lalu tulis ruas kiri sebagai
  $(\mu y)'$. Bila ruas kiri tidak mau menjadi satu turunan, langkah 1 salah.
\item Integralkan kedua ruas terhadap $t$, lalu bagi dengan $\mu$.
\item Pakai nilai awal untuk menentukan konstanta, dan tulis selang berlakunya.
\end{enumerate}
Peringatan: langkah 1 tidak boleh dilewati, dan langkah 3 tidak boleh dikerjakan
sebelum langkah 2. Menghitung $\mu$ dari koefisien yang belum dinormalkan memberi
faktor integrasi yang salah, dan kesalahannya baru kelihatan saat solusinya
diuji.
\end{resep}

\kamus{faktor integrasi}{integrating factor}{Pengali yang mengubah ruas kiri persamaan linear menjadi satu turunan hasil kali.}

\section{Keberadaan, ketunggalan, dan iterasi Picard}
\tingkat{menengah}

Contoh $y'=y^{2}$ menunjukkan solusi bisa berhenti ada. Diberi $f$ dan satu titik
awal, apakah solusinya ada dan apakah hanya satu? Jawabannya bergantung pada satu
sifat $f$ terhadap $y$.

\begin{definisi}[Syarat Lipschitz]\index{syarat Lipschitz}
Fungsi $f(t,y)$ memenuhi syarat Lipschitz terhadap $y$ pada daerah
$D\subseteq\mathbb{R}^{2}$ bila ada konstanta $L$ sehingga
\[
|f(t,y_{1})-f(t,y_{2})|\le L\,|y_{1}-y_{2}|
\]
untuk setiap $(t,y_{1}),(t,y_{2})\in D$. Konstanta terkecil yang mungkin disebut
konstanta Lipschitz.
\end{definisi}

Kalimat sehari-harinya: $f$ tidak boleh berubah terlalu mendadak ketika $y$
bergeser sedikit. Bila $f$ punya turunan parsial terhadap $y$ yang kontinu dan
terbatas oleh $L$ pada $D$, syarat itu otomatis berlaku, karena teorema nilai
rata-rata (Bab 12) memberi
$f(t,y_{1})-f(t,y_{2})=\partial_{y}f(t,\xi)(y_{1}-y_{2})$.

\begin{teorema}[Picard--Lindel\"of]\index{teorema Picard-Lindelof}
Misalkan $f$ kontinu pada persegi panjang tertutup
$D=[t_{0}-a,t_{0}+a]\times[y_{0}-b,y_{0}+b]$ dan memenuhi syarat Lipschitz
terhadap $y$ pada $D$ dengan konstanta $L$. Misalkan $M=\max_{D}|f|$ dan
$h=\min(a,\,b/M)$. Maka masalah nilai awal $y'=f(t,y)$, $y(t_{0})=y_{0}$
mempunyai tepat satu solusi pada selang $|t-t_{0}|\le h$.
\end{teorema}

Bukti lengkapnya memakai teorema titik tetap Banach pada ruang fungsi kontinu
dengan norma maksimum, yang lengkap karena kekontinuan seragam dibahas di Bab 11
dan Bab 14. Saya berikan sketsanya, karena sketsa itu memuat seluruh gagasan dan
dapat dijadikan algoritma.

Tulis persamaan sebagai persamaan integral, yang setara karena kedua sisinya
kontinu:
\[
y(t)=y_{0}+\int_{t_{0}}^{t}f\bigl(s,y(s)\bigr)\,ds .
\]
Definisikan pemetaan $T$ dengan rumus itu, mulai dari $y_{0}(t)=y_{0}$ konstan,
lalu ulangi:
\[
y_{n+1}(t)=y_{0}+\int_{t_{0}}^{t}f\bigl(s,y_{n}(s)\bigr)\,ds .
\]
Setiap $y_{n}$ terletak di dalam $D$ selama $|t-t_{0}|\le h$, karena
$|y_{n+1}-y_{0}|\le M|t-t_{0}|\le Mh\le b$. Selisih dua iterasi berturut-turut
memenuhi
\[
|y_{n+1}(t)-y_{n}(t)|\le L\Bigl|\int_{t_{0}}^{t}|y_{n}(s)-y_{n-1}(s)|\,ds\Bigr|
\le \frac{(L|t-t_{0}|)^{n}}{n!}\,\|y_{1}-y_{0}\| ,
\]
dan induksi pada $n$ memberi batas itu. Karena $\sum_{n}(Lh)^{n}/n!$ konvergen
(uji rasio, Bab 5), deret $y_{0}+\sum(y_{n+1}-y_{n})$ konvergen seragam, dan
batasnya adalah solusi. Ketunggalan datang dari pertidaksamaan yang sama: dua
solusi berbeda akan memenuhi $|y_{1}-y_{2}|\le L|t-t_{0}|\,\|y_{1}-y_{2}\|$,
yang untuk $Lh<1$ hanya mungkin bila selisihnya nol; kasus $Lh\ge1$ ditangani
dengan memperkecil $h$, dan itulah sebab teorema hanya menjanjikan selang
pendek.

Jalankan iterasi itu pada $y'=y$, $y(0)=1$. Mulai dari $y_{0}=1$:
\[
y_{1}=1+t,\quad y_{2}=1+t+\frac{t^{2}}{2},\quad
y_{3}=1+t+\frac{t^{2}}{2}+\frac{t^{3}}{6},
\]
sehingga $y_{n}$ adalah jumlah parsial deret eksponensial dari Bab 5. Di $t=1$,
iterasi kedelapan memberi $2{,}718278769841$, sedangkan $e=2{,}718281828459$;
galatnya $3{,}06\times10^{-6}$.

Syarat Lipschitz tidak boleh dibuang. Ambil
\[
y'=\sqrt{|y|},\qquad y(0)=0 .
\]
Ruas kanannya kontinu, jadi solusi ada; tetapi ia tidak memenuhi syarat Lipschitz
di sekitar $y=0$, karena kemiringannya tak terbatas di sana. Dan memang ada dua
solusi: fungsi nol, $y(t)\equiv0$, dan
\[
y(t)=\frac{t^{2}}{4}\quad (t\ge0).
\]
Pemeriksaannya satu baris: $y'=t/2$ dan $\sqrt{|y|}=\sqrt{t^{2}/4}=t/2$ untuk
$t\ge0$, sedangkan $y(0)=0$. Dua riwayat berbeda dari satu keadaan awal yang
sama, persis seperti lubang bak yang terbuka tiba-tiba pada kotak pembuka bab.

Lipschitz adalah syarat \emph{cukup}, bukan perlu. Keberadaan tanpa
ketunggalan dijamin teorema Peano, yang memakai kekompakan (Bab 11) dan
Arzel\`a--Ascoli; rujukannya Coddington dan Levinson, bab 1--2.

\kamus{ketunggalan}{uniqueness}{Jaminan bahwa satu keadaan awal hanya melahirkan satu riwayat, tanpa cabang.}
\section{Sistem, bidang fase, dan osilator}
\tingkat{menengah}

Satu persamaan orde dua bukan binatang baru; ia dua persamaan orde satu yang
berjalan bersama. Dengan menamai $v=\theta'$, persamaan pendulum menjadi
\[
\theta'=v,\qquad v'=-\frac{g}{L}\sin\theta .
\]
Pasangan ini otonom: ruas kanannya tidak memuat $t$, sehingga keadaan awal
menentukan seluruh riwayat (bila syarat Lipschitz dipenuhi) dan kita boleh
melupakan waktu sebentar untuk bertanya apa yang terjadi pada bidang
$(\theta,v)$.

\begin{definisi}[Sistem otonom dan bidang fase]\index{bidang fase}
Sistem
\[
\mathbf{x}'=\mathbf{f}(\mathbf{x}),\qquad \mathbf{x}\in\mathbb{R}^{n},
\]
disebut \emph{otonom}. Untuk $n=2$, bidang tempat titik $\mathbf{x}$ hidup
disebut \emph{bidang fase}; kurva $\{\mathbf{x}(t)\}$ disebut \emph{lintasan},
dan gambar seluruh panah $\mathbf{f}(\mathbf{x})$ disebut \emph{medan arah}.
Titik $\mathbf{x}^{*}$ dengan $\mathbf{f}(\mathbf{x}^{*})=0$ disebut
\emph{titik setimbang} atau titik tetap.
\end{definisi}

Sumbu mendatar bidang fase menyatakan posisi, sumbu tegak menyatakan laju, dan
satu titik berarti "di sini, dengan kecepatan ini". Waktu tidak digambar; ia
hanya parameter yang menyusuri kurva. Itulah keuntungannya: seluruh riwayat, untuk
semua titik awal sekaligus, tampak sebagai satu gambar.

Titik setimbang pendulum mudah dicari. $\theta'=0$ menuntut $v=0$, dan $v'=0$
menuntut $\sin\theta=0$, jadi $\theta=k\pi$. Di $\theta=0$ bandul tergantung
lurus ke bawah; di $\theta=\pi$ ia berdiri tegak ke atas. Yang pertama stabil
dalam pengertian sehari-hari, yang kedua tidak: sentuhan sekecil apa pun
menjatuhkannya.

\begin{untukanak}
Pergi ke taman, cari ayunan, dan duduklah. Minta seorang teman mendorongmu
pelan, hanya sedikit. Hitung sepuluh ayunan dengan menghitung di dalam kepala:
satu, dua, tiga, sampai sepuluh. Sekarang minta dorongan yang jauh lebih tinggi,
sampai ayunan naik hampir setinggi mistar atas, lalu hitung sepuluh ayunan lagi.

Yang akan kau lihat: ayunan tinggi memerlukan waktu lebih lama untuk satu
putaran penuh, walaupun kau merasa seolah-olah ayunan tinggi lebih cepat. Rasa
itu menipu karena jarak yang ditempuh jauh lebih panjang, sedangkan waktu satu
putaran hampir sama.

Sekarang lihat bagian atas ayunan. Saat kau berada di titik tertinggi, ada satu
saat pendek ketika kau hampir berhenti dan menggantung. Di titik terendah, kau
melaju paling cepat. Ayunan tidak pernah berhenti di tengah, dan tidak pernah
paling cepat di ujung. Dua hal itu selalu bertukar tempat.

Di bab ini kita akan tahu mengapa ayunan tinggi lebih lambat, dan mengapa
bandul yang berdiri tegak di atas hampir tidak mungkin dipertahankan.
\end{untukanak}

Sekarang rumus periode yang tepat. Kalikan $v'=-(g/L)\sin\theta$ dengan $v=\theta'$
dan integrasikan:
\[
\frac{d}{dt}\Bigl(\frac{v^{2}}{2}+\frac{g}{L}(1-\cos\theta)\Bigr)=0 .
\]
Kurung itu adalah energi per satuan massa, dan ia konstan sepanjang lintasan.
Dengan $E$ menyatakan nilainya,
\[
\frac{v^{2}}{2}+\frac{g}{L}\bigl(1-\cos\theta\bigr)=E ,
\]
sehingga $v=\pm\sqrt{2E-2(g/L)(1-\cos\theta)}$. Lintasan di bidang fase adalah
kurva-kurva energi itu; pada Gambar~\ref{fig:bidang-fase-pendulum} mereka
digambar untuk satuan $g/L=1$.

\begin{figure}[htbp]
\centering
\begin{tikzpicture}[x=1.05cm,y=1.0cm,font=\footnotesize]
  \draw[->,layergrey] (-4.15,0) -- (4.15,0) node[right] {$\theta$};
  \draw[->,layergrey] (0,-2.95) -- (0,2.95) node[above] {$v$};
  \foreach \tick/\lab in {-3.1416/{$-\pi$},-1.5708/{$-\pi/2$},
                          1.5708/{$\pi/2$},3.1416/{$\pi$}}{
    \draw[layergrey!55] (\tick,0.07) -- (\tick,-0.07);
    \node[layergrey,font=\scriptsize,below] at (\tick,-0.07) {\lab};}
  \node[layergrey,font=\scriptsize,below left] at (-0.06,0.02) {$0$};
  \foreach \th in {-3.5,-2.6,-1.7,-0.85,0.85,1.7,2.6,3.5}{
    \foreach \w in {-2.4,-1.6,-0.8,0,0.8,1.6,2.4}{
      \pgfmathsetmacro{\nn}{sqrt(\w*\w+sin(\th r)*sin(\th r))}
      \pgfmathsetmacro{\dx}{0.32*\w/\nn}
      \pgfmathsetmacro{\dy}{-0.32*sin(\th r)/\nn}
      \draw[-{Latex[length=0.9mm]},layerteal!60]
        (\th,\w) -- ({\th+\dx},{\w+\dy});}}
  \foreach \En/\tm in {0.4/0.9273,1.0/1.5708,1.6/2.2143}{
    \draw[layerindigo,thick]
      plot[domain=-\tm:\tm,samples=121,variable=\t]
        ({\t},{sqrt(max(0,2*(\En-1+cos(\t r))))});
    \draw[layerindigo,thick]
      plot[domain=-\tm:\tm,samples=121,variable=\t]
        ({\t},{-sqrt(max(0,2*(\En-1+cos(\t r))))});}
  \draw[accent,very thick]
    plot[domain=-3.1416:3.1416,samples=161,variable=\t]
      ({\t},{2*cos(\t r/2)});
  \draw[accent,very thick]
    plot[domain=-3.1416:3.1416,samples=161,variable=\t]
      ({\t},{-2*cos(\t r/2)});
  \draw[layerplum,thick,densely dashed]
    plot[domain=-3.6:3.6,samples=161,variable=\t]
      ({\t},{sqrt(2*(1.6+cos(\t r)))});
  \fill[ink] (0,0) circle (1.6pt);
  \node[ink,anchor=north east,font=\scriptsize] at (-0.12,-0.52) {pusat};
  \fill[accent] (3.1416,0) circle (1.6pt);
  \fill[accent] (-3.1416,0) circle (1.6pt);
  \node[accent,anchor=north west,font=\scriptsize,align=left]
    at (2.35,-0.18) {pelana};
  \node[accent,anchor=west,font=\scriptsize,align=left]
    at (2.30,1.30) {separatriks $E=2$};
  \node[layerplum,anchor=west,font=\scriptsize,align=left]
    at (0.55,2.72) {lintasan terbuka $E=2.6$};
\end{tikzpicture}
\caption{Bidang fase pendulum dalam satuan $g/L=1$, dengan sumbu mendatar sudut
$\theta$ dan sumbu tegak kecepatan $v$. Panah kecil menunjukkan medan arah
$(\theta',v')=(v,-\sin\theta)$. Kurva biru tua adalah tiga lintasan tertutup dengan
energi $E=0{,}4$, $1{,}0$, dan $1{,}6$; makin besar $E$, makin lebar putarannya.
Kurva merah adalah separatriks $E=2$, yaitu lintasan yang tepat mencapai
$\theta=\pm\pi$ dengan kecepatan nol; ia memisahkan ayunan tertutup dari putaran
terbuka. Kurva putus-putus adalah lintasan terbuka $E=2{,}6$, tempat bandul
berputar penuh tanpa pernah berhenti. Titik di $\theta=0$ adalah pusat, titik di
$\theta=\pm\pi$ adalah pelana.}
\label{fig:bidang-fase-pendulum}
\end{figure}

Gambar itu memuat seluruh perilaku kualitatif pendulum tanpa satu persamaan
diselesaikan; yang dibutuhkan hanya tanda $\sin\theta$ dan simetri energi.

Untuk amplitudo kecil, $\sin\theta\approx\theta$ dan sistemnya menjadi
osilator harmonik
\[
\theta'=v,\qquad v'=-\omega_{0}^{2}\theta,\qquad \omega_{0}=\sqrt{g/L},
\]
dengan solusi $\theta(t)=\theta_{0}\cos(\omega_{0}t)$ dan periode
$2\pi/\omega_{0}=2{,}006409$ detik untuk $L=1$ meter. Lintasannya elips
$v^{2}+\omega_{0}^{2}\theta^{2}=\text{konstan}$, dan itulah tiga lingkaran
kecil di dekat pusat pada gambar tadi.

Periode yang tepat diperoleh dari kekekalan energi, bukan dari linearisasi.
Dengan $E=(g/L)(1-\cos\theta_{0})$ untuk bandul yang dilepas dari diam di
$\theta_{0}$, waktu turun dari $\theta_{0}$ ke nol adalah
\[
\frac{T}{4}=\sqrt{\frac{L}{2g}}\int_{0}^{\theta_{0}}
\frac{d\theta}{\sqrt{\cos\theta-\cos\theta_{0}}}
=\sqrt{\frac{L}{g}}\,K(k),\qquad k=\sin\frac{\theta_{0}}{2},
\]
dengan $K$ integral eliptik lengkap jenis pertama. Untuk $\theta_{0}=60^{\circ}$
diperoleh $k=0{,}5$ dan $K$ pada parameter $m=k^{2}=0{,}25$ bernilai $1{,}6857503548$,
sehingga $T=2{,}153242$ detik. Kuadratur pada integral itu memberi angka yang sama
sampai $9\times10^{-13}$; dua jalan berbeda, satu hasil. Untuk $\theta_{0}=5^{\circ}$
rasionya $1{,}00047617$, yaitu kesalahan kurang dari seperdua ribu: di situlah
rumus kecil pantas dipakai, dan di situlah batas keberlakuannya.

\begin{tangan}
Ambil tali sepanjang satu meter, sebuah mur logam sebagai beban, dan stopwatch
di telepon. Ikat mur di ujung tali, tempelkan ujung lain ke kusen pintu dengan
selotip, dan ukur panjang tali dari titik gantung ke pusat mur: usahakan tepat
satu meter. Tarik mur ke samping hanya sekitar lima sentimeter, lepaskan, dan
catat waktu untuk sepuluh ayunan penuh. Bagi sepuluh; hasilnya sekitar dua
detik.

Sekarang tarik jauh lebih tinggi, kira-kira setinggi setengah meter dari titik
terendah, dan ulangi. Bandingkan kedua angka. Yang akan kau lihat: angka kedua
lebih besar, tetapi tidak dua kali lebih besar. Selisihnya sekitar tujuh persen,
dan itu cocok dengan hitungan di atas bila sudutnya sekitar enam puluh derajat.

Ulangi dengan tali seperempat meter: periodenya menjadi setengah, karena
periode bergantung pada akar panjang.
\end{tangan}

\section{Kestabilan: nilai eigen, linearisasi, dan Lyapunov}
\tingkat{lanjut}

Pertanyaan yang paling sering diajukan tentang sebuah sistem bukan "berapa
nilainya pada $t=7$", melainkan "apakah ia kembali setelah diganggu". Untuk
menjawabnya, kita tidak perlu solusi; cukup turunan di sekitar titik setimbang.

\begin{definisi}[Kestabilan]\index{kestabilan}
Misalkan $\mathbf{x}^{*}$ titik setimbang $\mathbf{x}'=\mathbf{f}(\mathbf{x})$.
Titik itu disebut \emph{stabil} bila untuk setiap $\varepsilon>0$ ada
$\delta>0$ sehingga $\|\mathbf{x}(0)-\mathbf{x}^{*}\|<\delta$ menjamin
$\|\mathbf{x}(t)-\mathbf{x}^{*}\|<\varepsilon$ untuk setiap $t\ge0$. Titik itu
disebut \emph{stabil asimtotik} bila stabil dan ada $\delta>0$ sehingga
$\mathbf{x}(t)\to\mathbf{x}^{*}$ untuk $t\to\infty$. Bila tidak stabil, ia
disebut \emph{tidak stabil}.
\end{definisi}

\begin{teorema}[Linearisasi]\index{linearisasi}
Misalkan $\mathbf{f}$ mempunyai turunan parsial kontinu di sekitar titik
setimbang $\mathbf{x}^{*}$ dan misalkan $A=J_{\mathbf{f}}(\mathbf{x}^{*})$
matriks Jacobian dari Bab 12. Bila tidak ada nilai eigen $A$ yang bagian
realnya nol (titik setimbang \emph{hiperbolik}), maka perilaku kualitatif sistem
nonlinear di sekitar $\mathbf{x}^{*}$ sama dengan perilaku sistem linear
$\mathbf{u}'=A\mathbf{u}$: jenis kestabilannya sama, dan lintasan di dekatnya
dapat dipetakan satu-satu dan kontinu ke lintasan linear.
\end{teorema}

Ini teorema Hartman--Grobman; buktinya panjang dan memakai teorema titik tetap
pada ruang fungsi, jadi saya kutip: Hartman, bab 9, Perko, bagian 2.8. Akibat
praktisnya: kestabilan lokal dibaca dari nilai eigen.

Untuk sistem dua dimensi, seluruh kemungkinan dibaca dari jejak
$\tau=\operatorname{tr}A$ dan determinan $\Delta=\det A$; nilai eigennya akar
$\lambda^{2}-\tau\lambda+\Delta=0$. Bila $\Delta<0$, akarnya real dan
berlawanan tanda: \emph{pelana}. Bila $\Delta>0$ dan $\tau^{2}\ge4\Delta$,
keduanya real dan bertanda sama: \emph{simpul}, stabil bila $\tau<0$. Bila
$\tau^{2}<4\Delta$, akarnya konjugat dengan bagian real $\tau/2$:
\emph{spiral}, stabil bila $\tau<0$, dan pusat bila $\tau=0$.

Contoh: matriks $A=\begin{pmatrix}0&1\\-2&-3\end{pmatrix}$ punya $\tau=-3$
dan $\Delta=2$, sehingga akarnya $\lambda^{2}+3\lambda+2=0$, yaitu $-1$ dan
$-2$: simpul stabil. Matriks $B=\begin{pmatrix}0&1\\-2&0\end{pmatrix}$ punya
$\tau=0$ dan $\Delta=2>0$, akarnya $\pm i\sqrt{2}$: pusat.

Gambar~\ref{fig:potret-linear} memperlihatkan dua potret yang paling berbeda
sifatnya: pelana, tempat hampir semua titik awal menjauh, dan spiral stabil,
tempat semuanya berputar masuk.

\begin{figure}[htbp]
\centering
\begin{tikzpicture}[font=\footnotesize]
  \begin{scope}[x=1.35cm,y=1.35cm]
    \begin{scope}
      \clip (-1.62,-1.52) rectangle (1.62,1.52);
      \foreach \c in {0.05,0.30,0.70,1.20}{
        \draw[layerteal,thick] plot[domain=-1.55:1.55,samples=90,variable=\s]
          ({sqrt(\c+\s*\s)},{\s});
        \draw[layerteal,thick] plot[domain=-1.55:1.55,samples=90,variable=\s]
          ({-sqrt(\c+\s*\s)},{\s});}
      \draw[layerindigo,densely dashed] (-1.5,-1.5) -- (1.5,1.5);
      \draw[layerindigo,densely dashed] (-1.5,1.5) -- (1.5,-1.5);
    \end{scope}
    \draw[->,layergrey] (-1.75,0) -- (1.75,0) node[right] {$x$};
    \draw[->,layergrey] (0,-1.55) -- (0,1.55) node[above] {$y$};
    \fill[ink] (0,0) circle (1.5pt);
    \node[layerindigo,font=\scriptsize,anchor=south east] at (-1.52,1.30)
      {arah eigen};
    \node[layergrey,anchor=north] at (0,-1.62)
      {(a) pelana: $\lambda=1,-1$};
  \end{scope}
  \begin{scope}[xshift=6.6cm,x=1.35cm,y=1.35cm]
    \draw[->,layergrey] (-1.75,0) -- (1.75,0) node[right] {$x$};
    \draw[->,layergrey] (0,-1.55) -- (0,1.55) node[above] {$y$};
    \foreach \r in {0.35,0.7,1.1,1.5}{
      \draw[accent,thick]
        plot[domain=0:12.6,samples=200,variable=\t]
          ({\r*exp(-0.3*\t)*cos(\t r)},{\r*exp(-0.3*\t)*(-sin(\t r))});}
    \fill[ink] (0,0) circle (1.5pt);
    \node[layergrey,anchor=north] at (0,-1.62)
      {(b) spiral stabil: $\lambda=-0.3\pm i$};
  \end{scope}
\end{tikzpicture}
\caption{Dua potret fase linear yang berbeda jenis. (a) Sistem
$x'=y$, $y'=x$: lintasan memenuhi $x^{2}-y^{2}=\text{konstan}$, hiperbola
dengan asimtot $y=\pm x$, dan garis putus-putus adalah arah eigen, satu masuk
dan satu keluar, sehingga titik asal adalah pelana dan tidak stabil. (b) Sistem $x'=-0{,}3x+y$, $y'=-x-0{,}3y$:
nilai eigen $-0{,}3\pm i$, sehingga lintasan berputar sambil menyusut dengan
faktor $e^{-0{,}3t}$ dan titik asal adalah spiral stabil. Empat lintasan di (b)
berbeda hanya pada jari-jari awal, karena sistemnya linear.}
\label{fig:potret-linear}
\end{figure}

Satu peringatan yang mahal. Bila $\tau=0$ dan $\Delta>0$, linearisasi memberi
pusat, tetapi sistem nonlinear di titik itu belum tentu pusat. Ambil
\[
x'=-y+x(x^{2}+y^{2}),\qquad y'=x+y(x^{2}+y^{2}).
\]
Jacobiannya di titik asal adalah matriks rotasi, nilai eigennya $\pm i$:
linearisasi berkata pusat. Tetapi dalam koordinat polar $r^{2}=x^{2}+y^{2}$,
perkalian silang memberi
\[
\frac{1}{2}\frac{d}{dt}r^{2}=xx'+yy'=r^{2}(x^{2}+y^{2})=r^{4},
\]
jadi $r'=r^{3}>0$ untuk $r>0$: lintasan berputar keluar, dan titik asal tidak
stabil. Linearisasi tidak salah; ia hanya tidak cukup ketika ada nilai eigen di
sumbu imajiner. Di situlah teorema tadi meminta titik setimbang hiperbolik.

Untuk kasus yang tidak hiperbolik, alat kedua dipakai: fungsi energi yang
menurun, yang kita karang sendiri.

\begin{definisi}[Fungsi Lyapunov]\index{fungsi Lyapunov}
Fungsi $V$ yang kontinu dan mempunyai turunan kontinu pada lingkungan $U$ dari
$\mathbf{x}^{*}$ disebut \emph{fungsi Lyapunov} bila $V(\mathbf{x}^{*})=0$ dan
$V(\mathbf{x})>0$ untuk $\mathbf{x}\in U\setminus\{\mathbf{x}^{*}\}$. Turunannya
sepanjang lintasan ditulis
\[
\dot V(\mathbf{x})=\nabla V(\mathbf{x})\cdot\mathbf{f}(\mathbf{x}).
\]
\end{definisi}

Bukit buatan itu adalah $V$, dan $\dot V$ kecuraman turunnya.

\begin{teorema}[Uji Lyapunov]
Misalkan $V$ fungsi Lyapunov di sekitar titik setimbang $\mathbf{x}^{*}$.
\begin{enumerate}
\item Bila $\dot V\le0$ pada $U$, maka $\mathbf{x}^{*}$ stabil.
\item Bila $\dot V<0$ pada $U\setminus\{\mathbf{x}^{*}\}$, maka
  $\mathbf{x}^{*}$ stabil asimtotik.
\end{enumerate}
\end{teorema}

\begin{proof}
Untuk pernyataan pertama, ambil $\varepsilon>0$ kecil sehingga bola
$\|\mathbf{x}-\mathbf{x}^{*}\|\le\varepsilon$ tinggal di dalam $U$. Di bola
itu $V$ positif dan bola itu kompak, jadi
$m=\min_{\|\mathbf{x}-\mathbf{x}^{*}\|=\varepsilon}V>0$. Pilih $\delta>0$
sehingga $V(\mathbf{x})<m$ untuk $\|\mathbf{x}-\mathbf{x}^{*}\|<\delta$;
mungkin karena $V$ kontinu dan $V(\mathbf{x}^{*})=0$. Lintasan dari bola $\delta$
memenuhi $V(\mathbf{x}(t))\le V(\mathbf{x}(0))<m$, sebab $\dot V\le0$,
sehingga ia tidak dapat menyeberangi bola $\varepsilon$; jadi stabil.

Untuk pernyataan kedua, $V(\mathbf{x}(t))$ menurun dan terbatas di bawah oleh
nol, jadi punya limit $L\ge0$. Andaikan $L>0$; lintasan itu tinggal di himpunan
kompak $\{V\ge L\}$, dan di sana $-\dot V$ punya minimum positif, sehingga $V$
harus turun melewati $L$ dalam waktu berhingga. Kontradiksi. Jadi $L=0$, dan
karena $V$ hanya nol di $\mathbf{x}^{*}$, lintasan menuju ke sana.
\end{proof}

Contohnya pendulum teredam, $\theta'=v$, $v'=-(g/L)\sin\theta-bv$ dengan
$b>0$, dan $V=v^{2}/2+(g/L)(1-\cos\theta)$, yaitu energi mekanik, positif
kecuali di titik setimbang. Turunannya
\[
\dot V=v\Bigl(-\frac{g}{L}\sin\theta-bv\Bigr)+\frac{g}{L}\sin\theta\,v=-bv^{2}\le0 .
\]
Jadi titik setimbang bawah stabil dan energi tidak pernah bertambah. Uji ini
tidak memberi kestabilan asimtotik, karena $\dot V$ nol di seluruh garis $v=0$;
versi LaSalle menambahkan bahwa satu-satunya himpunan invarian
di dalam $\{\dot V=0\}$ adalah titik setimbang, sehingga semua lintasan menuju
ke sana.

Linearisasi memberi jawaban yang sama. Pada $\theta=0$ Jacobiannya
$\begin{pmatrix}0&1\\-g/L&-b\end{pmatrix}$, dengan $\tau=-b$ dan
$\Delta=g/L>0$; untuk $b=0{,}5$ dan $g/L=9{,}80665$ akarnya $-0{,}25\pm3{,}121562\,i$,
yaitu spiral stabil. Pada $\theta=\pi$ muncul $+g/L$ di pojok kanan bawah,
sehingga $\Delta=-g/L<0$: pelana dengan nilai eigen $\pm3{,}131557$. Bandul yang
berdiri tegak memang pelana, bukan kiasan: ada satu arah tepat yang
mempertahankannya di atas, dan setiap gangguan membawanya pergi.

\kamus{stabil asimtotik}{asymptotically stable}{Tetap dekat setelah diganggu, lalu benar-benar tiba kembali ke titik setimbang.}

\kamus{fungsi Lyapunov}{Lyapunov function}{Bukit buatan yang dipakai sebagai penggaris: selama sistem menuruninya, ia tidak bisa lari jauh.}

\section{Euler, Runge--Kutta, dan kekakuan}
\tingkat{lanjut}

Hampir tidak ada persamaan yang dapat diselesaikan dengan rumus, jadi metode
numerik bukan pelengkap melainkan cara utama. Gagasannya kembali ke Bab 12: bila
kita tahu nilai dan turunan di satu titik, kita dapat menaksir nilai di sedikit
sebelah kanannya.

Metode Euler melakukan hal paling sederhana. Dari $y'=f(t,y)$ dan nilai awal,
\[
y_{n+1}=y_{n}+h\,f(t_{n},y_{n}),\qquad t_{n+1}=t_{n}+h .
\]
Ia tepat merupakan dua suku pertama deret Taylor di sekitar $t_{n}$, sehingga
galat satu langkahnya adalah suku $h^{2}y''(\xi)/2$, berorde $h^{2}$. Galat
lokal itu menumpuk menjadi galat global berorde $h$: menggandakan ketelitian
menuntut dua kali lipat jumlah langkah, bukan empat kali.

Angka untuk $y'=y$, $y(0)=1$, pada $t=1$ dengan $h=0{,}1$: Euler memberi
$(1{,}1)^{10}=2{,}5937424601$, sedangkan $e=2{,}7182818285$. Galatnya $0{,}1245$, yaitu
$4{,}58$ persen, dan galat satu langkah pertamanya $5{,}171\times10^{-3}$.

Metode Runge--Kutta orde empat memperbaiki ini dengan mengambil empat contoh
kemiringan di setiap langkah:
\[
y_{n+1}=y_{n}+\frac{h}{6}\bigl(k_{1}+2k_{2}+2k_{3}+k_{4}\bigr),
\]
dengan $k_{1}=f(t_{n},y_{n})$, $k_{2}=f(t_{n}+h/2,\,y_{n}+hk_{1}/2)$,
$k_{3}=f(t_{n}+h/2,\,y_{n}+hk_{2}/2)$, dan $k_{4}=f(t_{n}+h,\,y_{n}+hk_{3})$.
Galat lokalnya berorde $h^{5}$ dan galat globalnya berorde $h^{4}$: separuh
langkah berarti galat turun sekitar enam belas kali.

Tabel berikut berasal dari program di bagian lab, pada model logistik dengan
$y(0)=1/2$ dan $t=2$; nilai tepatnya $0{,}8807970780$.

\begin{center}
{\small
\begin{tabular}{lrrr}
\hline
$h$ & galat Euler & galat RK4 & rasio Euler \\
\hline
$0{,}400$ & $1{,}958\times10^{-2}$ & $1{,}245\times10^{-5}$ & --- \\
$0{,}200$ & $9{,}423\times10^{-3}$ & $7{,}237\times10^{-7}$ & $2{,}08$ \\
$0{,}100$ & $4{,}630\times10^{-3}$ & $4{,}359\times10^{-8}$ & $2{,}04$ \\
$0{,}050$ & $2{,}296\times10^{-3}$ & $2{,}674\times10^{-9}$ & $2{,}02$ \\
$0{,}025$ & $1{,}143\times10^{-3}$ & $1{,}656\times10^{-10}$ & $2{,}01$ \\
\hline
\end{tabular}}
\end{center}

Kolom terakhir adalah bukti orde pertama: setiap kali $h$ dibagi dua, galat
Euler turun hampir tepat dua kali; untuk RK4 rasionya sekitar $16$, sesuai orde
keempat. Gambar~\ref{fig:galat-numerik} menampilkan angka yang sama pada skala
logaritma kedua sumbu.

\begin{figure}[htbp]
\centering
\begin{tikzpicture}
\begin{axis}[width=0.80\linewidth,height=5.6cm,
  xlabel={langkah $h$},ylabel={$|y(2)-y_{h}(2)|$},
  xmode=log,ymode=log,xmin=0.02,xmax=0.5,ymin=1e-11,ymax=1e-1,
  grid=both,minor tick num=1,
  legend style={font=\scriptsize,at={(0.5,-0.22)},anchor=north,
    legend columns=2,draw=none,fill=none},
  tick label style={font=\scriptsize},label style={font=\small},
  axis line style={layergrey},tick style={layergrey}]
  \addplot[layerindigo,thick,mark=*,mark size=1.7pt] coordinates {
    (0.4,0.01958) (0.2,0.009423) (0.1,0.004630)
    (0.05,0.002296) (0.025,0.001143)};
  \addplot[accent,thick,mark=square*,mark size=1.6pt] coordinates {
    (0.4,0.00001245) (0.2,0.0000007237) (0.1,0.00000004359)
    (0.05,0.000000002674) (0.025,0.0000000001656)};
  \legend{Euler (kemiringan $1$), Runge--Kutta 4 (kemiringan $4$)}
\end{axis}
\end{tikzpicture}
\caption{Galat di $t=2$ untuk model logistik sebagai fungsi langkah $h$, pada
skala logaritma kedua sumbu. Garis Euler menurun dengan kemiringan satu: galat
berbanding lurus dengan $h$. Garis Runge--Kutta menurun dengan kemiringan empat:
galat berbanding dengan $h^{4}$. Pada $h=0{,}4$ kedua metode berjarak sekitar tiga
orde, dan jarak itu melebar ketika $h$ diperkecil. Angka pada gambar adalah
angka yang sama dengan tabel di teks.}
\label{fig:galat-numerik}
\end{figure}

Sekarang sifat yang membuat banyak orang tertipu: kekakuan. Sebuah sistem
disebut \emph{kaku} (\emph{stiff}) bila di dalamnya ada proses yang bergerak
jauh lebih cepat daripada proses yang kita minati. Ciri praktisnya: metode
eksplisit menuntut langkah sangat kecil untuk tetap \emph{stabil}, walaupun
ketelitian yang kita minta hanya menuntut langkah besar.

Contohnya kecil dan jujur: $y'=-50(y-\cos t)-\sin t$, $y(0)=1$, yang solusi
tepatnya $y=\cos t$. Di sini $y$ mengikuti $\cos t$ dengan lambat, tetapi setiap
penyimpangan diseret kembali dengan laju $50$ per detik. Euler stabil hanya bila
$|1-50h|<1$, yaitu $h<0{,}04$.

Hasil pada $t=3$, dengan nilai tepat $\cos3=-0{,}989992$: Euler dengan $h=0{,}02$
memberi galat $1{,}98\times10^{-4}$, dengan $h=0{,}039$ memberi
$7{,}97\times10^{-4}$, dan dengan $h=0{,}05$ memberi $-1{,}839\times10^{7}$. Satu
kenaikan dari $0{,}039$ ke $0{,}05$ mengubah jawaban dari hampir benar menjadi tidak
berarti. RK4 bertahan lebih lama, karena daerah stabilnya lebih lebar: pada
$h=0{,}05$ galatnya $1{,}11\times10^{-3}$, tetapi pada $h=0{,}1$ ia memberi
$-1{,}32\times10^{31}$; batas praktisnya $2{,}785/50=0{,}0557$.

Pelajaran yang dapat dibawa: galat dan kestabilan dua hal berbeda. Galat lokal
pada langkah pertama kecil dan tidak memperingatkan apa pun; yang meledak adalah
kestabilan, dan itu urusan nilai eigen.

Penawarnya metode implisit, yang paling sederhana Euler mundur:
\[
y_{n+1}=y_{n}+h\,f(t_{n+1},y_{n+1}),
\]
yang menuntut penyelesaian persamaan di setiap langkah, tetapi stabil untuk $h$
besar pada persamaan kaku; untuk contoh di atas, Euler mundur dengan $h=0{,}5$
memberi jawaban yang mendekati $\cos t$. Rujukan: Hairer dan Wanner, bab IV.

\begin{tantangan}
Linearisasi memberi pusat, tetapi sistem nonlinear tidak harus pusat. Susun satu
sistem dua dimensi yang titik asalnya \emph{bukan} pusat walaupun linearisasinya
pusat, tentukan sifat titik itu (stabil, stabil asimtotik, atau tidak stabil),
dan sertakan rumus eksplisit, nilai eigen Jacobian di titik asal, serta bukti
singkat klaimmu.

Petunjuk 1: linearisasi hanya melihat suku berderajat satu. Tambahkan suku
berderajat tiga yang simetris terhadap putaran, lalu tulis sistemnya dalam
koordinat polar.

Petunjuk 2: dalam koordinat polar, matriks rotasi menjadi $\theta'=1$, dan
gangguan radial menjadi $r'=cr^{3}$ untuk suatu konstanta $c$. Tanda $c$
menentukan segalanya, dan $V=r^{2}$ adalah fungsi Lyapunov yang bekerja untuk
salah satu tanda.

Kalau mentok, lihat kembali contoh $r'=r^{3}$ di bagian Kestabilan.
\end{tantangan}
\section{Contoh hitung}
\tingkat{menengah}

Empat hitungan berikut adalah tulang bab ini; semuanya dapat diperiksa ulang
dengan kalkulator atau program di bagian lab, dan semuanya pernah salah di tangan
saya sendiri.

\textbf{Model logistik: satu rumus, dua metode.}
Ambil $y'=y(1-y)$ dengan $y(0)=1/2$ dan tanyakan nilai di $t=2$. Rumus dari
bagian pemisahan variabel memberi $y(2)=1/(1+e^{-2})=0{,}8807970780$. Dengan
$h=0{,}1$ ada dua puluh langkah: Euler memberi $0{,}8854273272$, galat
$4{,}630\times10^{-3}$ atau $0{,}53$ persen, sedangkan Runge--Kutta dengan langkah
yang sama memberi $0{,}8807970344$, galat $4{,}359\times10^{-8}$. Biayanya empat kali
lipat, ketelitiannya sekitar seratus ribu kali.

Perhatikan juga arah galatnya. Euler di sini terlalu tinggi, karena
$y(1-y)$ cekung ke bawah pada selang itu sehingga garis lurus selalu berada di
atas kurva. Memeriksa tanda galat adalah cara murah untuk mengetahui apakah program kita
bekerja seperti seharusnya.

\textbf{Koefisien variabel: $y'+y/t=t^{2}$, $y(1)=0$.}
Di sini $p(t)=1/t$, sehingga
\[
\mu(t)=\exp\Bigl(\int\frac{dt}{t}\Bigr)=t
\]
pada $t>0$. Mengalikan persamaan dengan $t$ memberi
\[
t\,y'+y=t^{3}
\quad\Longleftrightarrow\quad
\bigl(t\,y\bigr)'=t^{3},
\]
jadi $t\,y=t^{4}/4+C$ dan
\[
y(t)=\frac{t^{3}}{4}+\frac{C}{t}.
\]
Syarat $y(1)=0$ memberi $C=-1/4$, sehingga
\[
y(t)=\frac{t^{3}}{4}-\frac{1}{4t}=\frac{t^{4}-1}{4t},
\qquad y(2)=\frac{15}{8}=1.875 .
\]
Pemeriksaan: $y'=3t^{2}/4+1/(4t^{2})$ dan $y/t=t^{2}/4-1/(4t^{2})$, jumlahnya
$t^{2}$, tepat ruas kanan. Integrator adaptif memberi $1{,}875000000000$.

Tanpa $\mu=t$, ruas kiri $y'+y/t$ tidak tampak seperti turunan apa pun; setelah
dikali $t$, ia menjadi turunan $ty$.

\textbf{Berapa besar ayunan sebelum rumus kecil meleset satu persen.}
Rasio periode tepat terhadap periode kecil adalah
\[
\frac{T(\theta_{0})}{T_{0}}
=\frac{2}{\pi}K\Bigl(\sin^{2}\frac{\theta_{0}}{2}\Bigr),
\]
dengan $K$ integral eliptik lengkap. Beberapa nilainya:

\begin{center}
{\small
\begin{tabular}{rrrr}
\hline
$\theta_{0}$ & $T/T_{0}$ & galat & keterangan \\
\hline
$5^{\circ}$ & $1{,}000476$ & $0{,}048$ persen & jam bandul \\
$10^{\circ}$ & $1{,}001907$ & $0{,}191$ persen & masih aman \\
$20^{\circ}$ & $1{,}007669$ & $0{,}767$ persen & mendekati batas \\
$22{,}81^{\circ}$ & $1{,}010000$ & $1{,}000$ persen & batas satu persen \\
$30^{\circ}$ & $1{,}017409$ & $1{,}741$ persen & perlu koreksi \\
$60^{\circ}$ & $1{,}073182$ & $7{,}318$ persen & jam meleset sejam \\
\hline
\end{tabular}}
\end{center}

Jadi rumus kecil itu layak dipakai sampai sekitar dua puluh derajat, dan
$22{,}81^{\circ}$ adalah ambang galat satu persen; galatnya tumbuh seperti kuadrat
amplitudo, bukan amplitudo.
Menggandakan amplitudo dari $10^{\circ}$ ke $20^{\circ}$ menaikkan galat dari
$0{,}19$ menjadi $0{,}77$ persen, yaitu empat kali, cocok dengan
$1+(\theta_{0}^{2}/16)$ sebagai dua suku pertama taksiran itu.

\kertas{
Uji sendiri dalam tiga langkah. (1) Kerjakan ulang contoh logistik dengan
$h=0{,}2$; engkau harus mendapat galat Euler $9{,}423\times10^{-3}$ dan galat RK4
$7{,}237\times10^{-7}$. (2) Turunkan sendiri rumus periode eliptik dari kekekalan
energi, lalu periksa bahwa $K$ pada parameter $0{,}25$ bernilai $1{,}6857503548$.
(3) Cari ambang $h$ untuk Euler pada $y'=-100y$; jawabannya $h<0{,}02$, dan pada
$h=0{,}02$ tepat, metode itu berayun antara $+1$ dan $-1$ tanpa meluruh maupun
meledak.}

\section{Lab komputasi}
\tingkat{lanjut}

Kode berikut menghitung semua angka di atas: galat Euler dan RK4 pada model
logistik, periode pendulum dengan dua jalan berbeda, dan ambang kekakuan. Ia
hanya memakai NumPy dan SciPy.

\begin{lab}
\begin{verbatim}
import numpy as np
from scipy import integrate, special

def euler(f, y0, T, h):
    y, t, n = y0, 0.0, int(round(T / h))
    for _ in range(n):
        y += h * f(t, y); t += h
    return y

def rk4(f, y0, T, h):
    y, t, n = y0, 0.0, int(round(T / h))
    for _ in range(n):
        k1 = f(t, y); k2 = f(t + h/2, y + h*k1/2)
        k3 = f(t + h/2, y + h*k2/2); k4 = f(t + h, y + h*k3)
        y += h * (k1 + 2*k2 + 2*k3 + k4) / 6; t += h
    return y

logi = lambda t, y: y * (1 - y)
ex = lambda t: 1 / (1 + np.exp(-t))
print("logistik y(2) eksak %.10f" % ex(2.0))
for h in (0.4, 0.2, 0.1, 0.05, 0.025):
    ee = abs(euler(logi, 0.5, 2, h) - ex(2))
    er = abs(rk4(logi, 0.5, 2, h) - ex(2))
    print("h=%.3f euler %.3e rk4 %.3e" % (h, ee, er))

g, L = 9.80665, 1.0
T0 = 2 * np.pi * np.sqrt(L / g)
th0 = np.radians(60.0)
k = np.sin(th0 / 2)
Tq = 4 * np.sqrt(L / (2*g)) * integrate.quad(
        lambda th: 1/np.sqrt(np.cos(th) - np.cos(th0)), 0, th0, limit=200)[0]
Te = 4 * np.sqrt(L / g) * special.ellipk(k**2)
print("T0 %.6f T60kuad %.6f T60eliptik %.6f" % (T0, Tq, Te))
print("rasio %.6f persen %.4f" % (Te / T0, 100 * (Te / T0 - 1)))

stiff = lambda t, y: -50 * (y - np.cos(t)) - np.sin(t)
for h in (0.02, 0.039, 0.05):
    print("kaku euler h=%.3f y(3)=%.3e" % (h, euler(stiff, 1.0, 3, h)))
for h in (0.05, 0.1):
    print("kaku rk4 h=%.3f y(3)=%.3e" % (h, rk4(stiff, 1.0, 3, h)))
print("eksak cos(3) %.10f" % np.cos(3))
\end{verbatim}
\end{lab}

Keluaran sebenarnya, apa adanya:

\begin{verbatim}
logistik y(2) eksak 0.8807970780
h=0.400 euler 1.958e-02 rk4 1.245e-05
h=0.200 euler 9.423e-03 rk4 7.237e-07
h=0.100 euler 4.630e-03 rk4 4.359e-08
h=0.050 euler 2.296e-03 rk4 2.674e-09
h=0.025 euler 1.143e-03 rk4 1.656e-10
T0 2.006409 T60kuad 2.153242 T60eliptik 2.153242
rasio 1.073182 persen 7.3182
kaku euler h=0.020 y(3)=-9.902e-01
kaku euler h=0.039 y(3)=-9.908e-01
kaku euler h=0.050 y(3)=-1.839e+07
kaku rk4 h=0.050 y(3)=-9.883e-01
kaku rk4 h=0.100 y(3)=-1.317e+31
eksak cos(3) -0.9899924966
\end{verbatim}

Tiga hal patut dibaca. Pertama, rasio galat Euler $2{,}08$, $2{,}04$, $2{,}02$, $2{,}01$
dan RK4 menuju $16=2^{4}$: orde pertama dan keempat terlihat sebagai angka. Kedua,
dua jalan untuk periode pendulum bertemu sampai enam digit, padahal kuadratur dan
integral eliptik tidak memakai rumus yang sama. Ketiga, baris kaku: $h=0{,}039$
masih wajar, $h=0{,}05$ tidak berarti, dan tidak satu pun baris memperingatkan
kita.

Program itu tidak memeriksa syarat teorema: apakah $f$ memenuhi Lipschitz,
apakah titik setimbangnya hiperbolik. Itu tetap pekerjaan kita.

\begin{untukdosen}
Urutan mengajar yang saya pakai: mulai dari jam bandul dan selisih $7{,}32$
persen, lalu definisi solusi beserta contoh ledakan $1/(1-t)$, baru pemisahan
variabel dan faktor integrasi. Faktor integrasi sebaiknya dibuktikan penuh di
papan, karena buktinya hanya memakai aturan hasil kali dan teorema dasar
kalkulus, dua alat yang sudah mereka punya sejak Bab 13. Picard--Lindelof cukup
sebagai sketsa dengan satu iterasi dijalankan; yang harus tinggal adalah arti
syarat Lipschitz, bukan teknik titik tetapnya. Bidang fase sebelum nilai eigen:
gambar dulu, hitung kemudian.

Dua miskonsepsi yang hampir selalu muncul. Pertama, mahasiswa menganggap
"solusi" berarti rumus tertutup, lalu menyimpulkan persamaan tanpa rumus tidak
punya solusi; contoh $y'=\sqrt{|y|}$ membantu, karena di sana ada dua solusi dan
tidak ada rumus tunggal. Kedua, mereka memperlakukan galat kecil sebagai bukti
kebenaran, dan terkejut ketika Euler dengan $h=0{,}05$ memberi $-1{,}8\times10^{7}$.
Tunjukkan baris keluaran itu, jangan hanya diperingatkan.

Satu aktivitas papan tulis: gambar bidang fase pendulum kosong, lalu minta
mahasiswa menandai di mana kecepatan nol, di mana percepatan nol, dan di mana
energi tertinggi. Kelas akan berdebat tentang daerah di antara dua pelana, dan
perdebatan itu persis isi teorema. Tanda kesiapan lanjut: mahasiswa dapat
menebak jenis titik setimbang dari jejak dan determinan tanpa menghitung akar.

\end{untukdosen}

\section{Ringkasan}
\begin{ringkasan}
\begin{itemize}
\item Persamaan diferensial menyamakan laju perubahan dengan keadaan; solusinya
fungsi pada sebuah selang, dan selang itu bagian dari jawaban.
\item Dua keluarga terpecahkan dengan tangan: pemisahan variabel dan persamaan
linear orde satu dengan faktor integrasi $\mu=\exp(\int p\,dt)$.
\item Picard--Lindel\"of menjamin solusi tunggal pada selang pendek bila $f$
kontinu dan memenuhi syarat Lipschitz terhadap $y$; tanpa syarat itu, satu
keadaan awal dapat melahirkan dua riwayat.
\item Bidang fase mengganti waktu dengan gambar: titik setimbang, lintasan
tertutup, separatriks, dan lintasan terbuka terbaca tanpa rumus solusi.
\item Kestabilan lokal dibaca dari nilai eigen Jacobian, kecuali bila ada nilai
eigen di sumbu imajiner; di situ linearisasi tidak memutuskan, dan fungsi
Lyapunov dipakai.
\item Euler berorde satu secara global dan RK4 berorde empat; keduanya eksplisit
dan karena itu gagal pada sistem kaku, tempat kestabilan, bukan ketelitian,
yang menentukan besar langkah.
\end{itemize}
\end{ringkasan}

\section{Latihan}
\begin{latihan}
\begin{enumerate}
\item Selesaikan $y'=2ty$, $y(0)=1$, lalu hitung $y(1)$ dan bandingkan dengan
  $e$. Hitung pula Euler dengan $h=0{,}1$ dan sebutkan galatnya. (Hasil saya:
  $2{,}33463336$, galat $3{,}84\times10^{-1}$.)
\item Buktikan bahwa $y'+p(t)y=q(t)$ dengan $p,q$ kontinu punya paling banyak
  satu solusi, tanpa rumus eksplisit: selisih dua solusi memenuhi
  $w'=-p(t)w$ dan $w(t_{0})=0$.
\item Benar, salah, atau belum cukup: ``karena linearisasi di titik asal memberi
  nilai eigen $\pm i$, titik asal adalah pusat yang stabil''. Beri satu contoh
  pendukung atau satu penyangkal.
\item Buktikan $E=v^{2}/2+\omega_{0}^{2}\theta^{2}/2$ konstan pada osilator
  harmonik, lalu jelaskan mengapa lintasannya elips dan periodenya
  $2\pi/\omega_{0}$.
\item Rangkaian RC dengan $R=2{,}2\,\mathrm{k}\Omega$, $C=10\,\mu\mathrm{F}$,
  $V_{s}=12$ volt, mulai dari nol. Hitung $\tau$, tegangan pada $t=\tau$, dan
  waktu sampai $6$ volt. ($22$ milidetik; $7{,}585447$ volt; $15{,}249$ milidetik.)
\item Tunjukkan bahwa $v=\pm2\cos(\theta/2)$ memenuhi $E=2$, lalu jelaskan
  mengapa kurva itu menghubungkan $(-\pi,0)$ dengan $(\pi,0)$ dan disebut
  separatriks.
\item Bandul $L=0{,}25$ meter dilepas dari $30^{\circ}$. Hitung periode kecil dan
  periode tepatnya, lalu sebutkan selisih persennya. ($1{,}003205$ detik;
  $1{,}020669$ detik; $1{,}74$ persen.)
\item Untuk $y'=-100y$ dengan $y(0)=1$, tentukan ambang $h$ yang membuat Euler
  eksplisit tetap stabil, lalu jelaskan apa yang terjadi tepat pada ambang itu.
  Periksa dengan program.
\item Buktikan bahwa persamaan otonom skalar $y'=f(y)$ dengan $f$ kontinu dan
  memenuhi Lipschitz tidak dapat mempunyai solusi periodik tidak konstan.
\end{enumerate}

\textbf{Petunjuk tiga butir tersulit.} Butir 3: linearisasi hanya melihat suku
berderajat satu; tambahkan suku kubik simetris seperti pada contoh $r'=r^{3}$
dan bandingkan dengan $r'=-r^{3}$. Butir 6: pakai
$\cos\theta=2\cos^{2}(\theta/2)-1$ dan periksa nilai $v$ di $\theta=\pm\pi$.
Butir 9: andaikan $y$ periodik tidak konstan; di titik maksimum $y'=0$, jadi
$f(y_{\max})=0$, dan ketunggalan memaksa lintasan itu konstan.
\end{latihan}

\begin{untukpeneliti}
Yang benar-benar terbukti di bab ini, dengan bukti yang ditulis: rumus faktor
integrasi beserta ketunggalannya, sketsa lengkap Picard--Lindelof dengan
batas galat iterasi, ketunggalan solusi linear lewat $w'=-pw$, uji Lyapunov
untuk kedua pernyataannya, dan seluruh identitas energi pada osilator serta
pendulum. Yang saya kutip tanpa bukti: teorema Hartman--Grobman (Hartman,
\emph{Ordinary Differential Equations}, bab 9; Perko, \emph{Differential
Equations and Dynamical Systems}, bagian 2.8), teorema Peano dan
Arzel\`a--Ascoli, serta prinsip invarians LaSalle. 

Yang dipopulerkan secara berlebihan adalah ungkapan "nilai eigen menentukan
segalanya". Yang benar: nilai eigen menentukan perilaku lokal hanya bila tidak
ada yang bagian realnya nol; di sumbu imajiner suku berderajat tiga dapat
mengubah pusat menjadi spiral, seperti contoh $r'=r^{3}$ di sini. Untuk sistem
berdimensi tinggi, "stabil asimtotik" lokal bahkan tidak berarti aman dari
gangguan besar, karena daerah tariknya dapat sangat kecil.

Pertanyaan terbuka: seberapa jauh sifat kualitatif dapat disertifikasi secara
numerik? Galat kecil pada satu lintasan tidak memberi tahu apa pun tentang
kestabilan; Euler dengan $h=0{,}039$ benar sampai empat digit, sedangkan $h=0{,}05$
memberi $10^{7}$. Untuk sistem yang hanya diketahui lewat sampel medan, tidak ada
cara umum memastikan tidak ada daerah kaku yang tersembunyi di antara dua titik
sampel.

Batas keberlakuan: semua fungsi di sini real dan cukup mulus; semua angka dari
Python 3.14.4, NumPy 2.5.2, dan SciPy 1.18.0, dengan kuadratur bertoleransi
$10^{-13}$ kecuali disebut lain; $g=9{,}80665\ \mathrm{m/s^{2}}$ dan gesekan
diabaikan pada hitungan periode.
\end{untukpeneliti}

\peta{Bab 16 pada Boyce dan DiPrima, \emph{Elementary Differential Equations
and Boundary Value Problems}, bab 1 (pengantar), 2 (persamaan orde satu), 3
(persamaan linear orde dua), 4 (sistem orde satu dan bidang fase), dan 9
(persamaan nonlinear dan kestabilan); pada Strogatz, \emph{Nonlinear Dynamics
and Chaos}, bab 2 (aliran pada garis), 5 (sistem dua dimensi), 6 (bidang fase),
7 (osilator), dan 8 (bifurkasi); pada Arnold, \emph{Ordinary Differential
Equations}, bab 1--4 untuk keberadaan, ketunggalan, dan klasifikasi titik
setimbang; pada Coddington dan Levinson, bab 1--2 untuk Picard--Lindelof dan
Peano; pada Hartman, bab 9, dan Perko, bagian 2.8, untuk Hartman--Grobman;
pada Hairer dan Wanner, \emph{Solving Ordinary Differential Equations II}, bab
IV, untuk kekakuan dan metode implisit. Dasar hitungannya bab \emph{Turunan} (turunan parsial dan Jacobian), bab
\emph{Deret Fungsi dan Pertukaran Limit} (limit seragam, yang membuat iterasi
Picard konvergen), dan bab \emph{Integral Lipat} (medan vektor) pada edisi
pertama buku ini; nilai eigen dari Bab 9 dipakai langsung tanpa diulang, dan
analisis galat floating-point dari Bab 22 menjadi lanjutan bagian numerik bab
ini.}


\part{Peluang, Statistika, dan Aritmetika Komputer}

Bagian ini menyiapkan dua hal yang tampak tidak berhubungan: cara mengukur
ketidakpastian dan cara komputer menyimpan bilangan. Prasyaratnya Bab~2, 5, dan~10.
Setelah enam bab, kamu bisa menghitung peluang, menaksir parameter, menguji klaim,
memasang regresi, dan menjelaskan mengapa $0{,}1+0{,}2$ tidak sama dengan $0{,}3$.
\chapter{Peluang dan Peluang Bersyarat}

Sebuah klinik memasang iklan uji saring untuk penyakit yang jarang: hanya satu dari
seribu orang di daerah itu mengidapnya. Ujinya memang bagus. Dari setiap seratus orang
sakit, sembilan puluh sembilan memberi hasil positif. Dari setiap seratus orang sehat,
hanya satu yang salah dituduh positif. Iklan menuliskan kedua angka itu sebagai
``ketepatan 99 persen''.

Sekarang hasil positif itu ada di tangan kita. Berapa peluang kita benar-benar sakit?
Pada populasi $100{,}000$ orang, sekitar $100$ mengidap penyakit itu dan $99$ di
antaranya positif. Sisanya, $99{,}900$ orang sehat, menyumbang $999$ positif palsu.
Jadi ada $1{,}098$ hasil positif dan hanya $99$ yang benar: $99/1{,}098=11/122$,
yaitu $0.0902$. Sekitar satu dari sebelas, bukan sembilan puluh sembilan dari seratus.

Kedua angka itu benar dan keduanya menjawab pertanyaan yang berbeda. Iklan memberi
peluang hasil positif bila orangnya sakit. Pasien ingin tahu peluang orangnya sakit
bila hasilnya positif. Menukar dua arah itu bukan kesalahan aritmetika; itu kesalahan
menamai besaran, dan akibatnya sering sepuluh kali lipat.

Bab 2 sudah memberi bahasa untuk kejadian: gabungan, irisan, komplemen. Bab 5 memberi
cara menjumlahkan tak hingga banyak suku. Keduanya belum cukup, karena himpunan tidak
membawa bobot: $A\cup B$ tetap himpunan yang sama, entah $A$ hampir pasti terjadi atau
hampir pasti tidak. Yang dibutuhkan tiga hal: fungsi bobot dengan aturan yang tidak boleh
dilanggar, cara mencacah bobot, dan aturan memperbarui bobot ketika informasi baru masuk.
Bab ini menyusun ketiganya, lalu memeriksanya dengan simulasi berbenih tetap.

\begin{untukanak}
Bayangkan penjaga gerbang sekolah memeriksa tas dengan alat yang berbunyi. Alat itu
bagus: dari seratus tas yang benar-benar membawa barang terlarang, sembilan puluh
sembilan membuatnya berbunyi. Tetapi alat itu juga berbunyi untuk satu dari seratus
tas yang isinya biasa saja. Sekarang misalkan di sekolah itu hanya satu dari seribu
tas yang benar-benar membawa barang terlarang.

Pertanyaannya sederhana dan mengejutkan. Kalau alat berbunyi, apakah tas itu pasti
berbahaya? Ternyata tidak. Tas berbahaya memang jarang, sedangkan tas biasa sangat
banyak, sehingga bunyi palsu dari tas biasa jauh lebih banyak daripada bunyi benar.
Penjaga yang hanya melihat lampu merah akan menahan banyak anak yang tidak bersalah.

Pertanyaan yang membuat penasaran: berapa bagian dari semua bunyi itu yang benar?
Jawabannya dapat dihitung tanpa menebak, asalkan kita tahu berapa banyak tas dari
masing-masing jenis. Di bab ini kita akan tahu mengapa bunyi alat yang ``sembilan puluh
sembilan persen akurat'' bisa berubah menjadi tuduhan yang salah dalam sembilan dari
sepuluh kasus, dan mengapa arah pertanyaannya menentukan jawabannya.
\end{untukanak}

\section{Aksioma peluang dan akibatnya}
\tingkat{dasar}

Sebelum satu angka pun dihitung, kita harus menyepakati benda apa yang diberi bobot.
Kalimat sehari-harinya: peluang adalah bobot yang kita taruh pada setiap kemungkinan,
dengan aturan bahwa total bobot seluruh kemungkinan adalah satu, dan bobot tidak pernah
negatif. Aturan itu tidak menjelaskan arti filosofis peluang; ia menetapkan aritmetika yang
harus dipatuhi setiap perhitungan.

\kamus{peluang}{probability}{bobot pada setiap kemungkinan, dengan total bobot seluruh
kemungkinan sama dengan satu}

\index{peluang}
\index{ruang sampel}

\begin{definisi}[Ruang peluang]
Misalkan $\Omega$ adalah himpunan hasil, $\mathcal{F}$ kumpulan kejadian, yaitu
himpunan bagian $\Omega$ yang memuat $\Omega$ dan tertutup terhadap komplemen serta
gabungan terbilang, dan $P\colon\mathcal{F}\to\mathbb{R}$. Fungsi $P$ disebut
\emph{peluang} bila
\begin{enumerate}
  \item $P(A)\ge 0$ untuk setiap $A\in\mathcal{F}$;
  \item $P(\Omega)=1$;
  \item bila $A_{1},A_{2},\dots$ saling lepas, maka
        $P\bigl(\bigcup_{i\ge 1}A_{i}\bigr)=\sum_{i\ge 1}P(A_{i})$.
\end{enumerate}
\end{definisi}

Aksioma ketiga, aditif terbilang, membedakan teori peluang modern dari perhitungan
untung-untungan abad ketujuh belas. Tanpanya, deret peluang seperti
$\sum_{n\ge 1}2^{-n}=1$ tidak punya dasar, padahal deret itulah yang muncul ketika kita
bertanya pada lemparan ke berapa keberhasilan pertama terjadi. Aksioma ini juga yang
membuat peluang dapat dipakai pada limit, persoalan Bab 4 dan Bab 5.

Tiga akibat langsung berikut dipakai terus-menerus, jadi layak dibuktikan sekali.

\begin{proposisi}[Akibat aksioma]
\index{aksioma peluang}
Untuk kejadian $A$ dan $B$ dalam ruang peluang berlaku
$P(A^{c})=1-P(A)$, $P(\varnothing)=0$, dan bila $A\subseteq B$ maka
$P(A)\le P(B)$. Selain itu
$P(A\cup B)=P(A)+P(B)-P(A\cap B)$.
\end{proposisi}

\begin{proof}[Bukti]
Karena $A$ dan $A^{c}$ saling lepas dan gabungannya $\Omega$, aksioma kedua dan ketiga
memberi $1=P(\Omega)=P(A)+P(A^{c})$. Untuk yang kedua, ambil $A=\varnothing$ pada hasil
pertama. Untuk yang ketiga, tulis $B$ sebagai gabungan saling lepas $A$ dan
$B\setminus A$; maka $P(B)=P(A)+P(B\setminus A)\ge P(A)$ karena peluang tidak negatif.
Untuk yang terakhir, pecah $A\cup B$ menjadi tiga bagian saling lepas $A\setminus B$,
$A\cap B$, dan $B\setminus A$. Dengan cara yang sama $A$ dan $B$ masing-masing terpecah
menjadi dua bagian, sehingga
\[
P(A\cup B)=P(A\setminus B)+P(A\cap B)+P(B\setminus A)
=P(A)+P(B)-P(A\cap B).\qedhere
\]
\end{proof}

Sifat terakhir itu, inklusi--eksklusi, adalah alasan mengapa ``atau'' tidak boleh
dijumlahkan begitu saja. Contohnya, di satu kelas $30$ mahasiswa, $18$ menyukai kopi
dan $12$ menyukai teh, serta $8$ menyukai keduanya. Maka peluang seorang mahasiswa yang
dipilih acak menyukai kopi atau teh adalah
\[
\frac{18+12-8}{30}=\frac{22}{30}=0.7333,
\]
bukan $30/30=1$ seperti yang diberikan penjumlahan naif $18+12=30$. Delapan orang
dihitung dua kali, dan inklusi--eksklusi menguranginya sekali.

Untuk tiga kejadian, rumusnya memperpanjang dengan tanda berganti:
\[
P(A\cup B\cup C)=P(A)+P(B)+P(C)-P(A\cap B)-P(A\cap C)-P(B\cap C)+P(A\cap B\cap C).
\]
Buktinya menerapkan versi dua kejadian dua kali, dan itu latihan pertama bab ini. Pola
tandanya bukan hiasan: bagian yang terlalu banyak dikurangi harus ditambahkan kembali.

Satu peringatan untuk Bab 18. Pada ruang berhingga yang setiap hasilnya berpeluang
positif, $P(A)=0$ berarti $A$ kosong; pada ruang kontinu tidak, sebab satu titik dapat
berpeluang nol tanpa mustahil. Bab ini bekerja di ruang berhingga.

\begin{contoh}
Dua dadu dilempar. Berapa peluang jumlah matanya $8$ atau kedua matanya sama? Dengan
$36$ hasil yang setara, kejadian pertama memuat $5$ hasil, kejadian kedua memuat $6$,
dan irisannya satu hasil, yaitu $(4,4)$. Jadi peluangnya
$(5+6-1)/36=10/36=0.2778$. Tanpa mengurangkan irisan, jawabannya $11/36=0.3056$, terlalu
besar sebesar $1/36$.
\end{contoh}

\section{Pencacahan dan peluang yang setara}
\tingkat{menengah}

Bila semua hasil sama mungkin, peluang sebuah kejadian tinggal nisbah banyak hasil yang
menguntungkan terhadap banyak hasil seluruhnya. Kalimat sehari-harinya: membagi jumlah
cara yang kita inginkan dengan jumlah cara yang ada. Definisi ini bukan hukum alam; ia
sebuah asumsi tentang simetri, dan asumsi itu harus diperiksa, bukan diandaikan.

\index{pencacahan}
\index{kombinasi}

\kamus{kombinasi}{combination}{pilihan beberapa benda tanpa memperhatikan urutannya,
seperti mengambil lima kartu dari satu tumpukan}

\begin{proposisi}[Aturan perkalian dan pencacahan]
\index{aturan perkalian}
Bila $|A|=m$ dan $|B|=n$, maka $|A\times B|=mn$. Banyak barisan panjang $k$ dari $n$
benda berbeda tanpa pengulangan adalah $n!/(n-k)!$, dan banyak himpunan bagian
berukuran $k$ adalah
\[
\binom{n}{k}=\frac{n!}{k!\,(n-k)!}.
\]
Khususnya $\sum_{k=0}^{n}\binom{n}{k}=2^{n}$.
\end{proposisi}

\begin{proof}[Bukti]
Untuk hasilkali, setiap pasangan $(a,b)$ ditentukan oleh satu pilihan dari $m$ dan satu
pilihan dari $n$. Untuk barisan, ada $n$ pilihan untuk tempat pertama, $n-1$ untuk
tempat kedua, dan seterusnya sampai $n-k+1$. Untuk himpunan bagian, setiap himpunan
berukuran $k$ dapat diurutkan dengan $k!$ cara, sehingga banyak barisan dibagi $k!$
memberi banyak himpunan. Identitas terakhir menghitung semua himpunan bagian dengan dua
cara: menjumlahkan menurut ukuran, atau memilih masuk atau tidak untuk setiap benda.
\end{proof}

Pencacahan yang rapi sering menyingkat perhitungan panjang. Tangan penuh dalam poker,
tiga kartu seperingkat ditambah sepasang kartu seperingkat lain, memberi contoh bagus:
pilih peringkat trio ($13$ cara), pilih tiga dari empat jenisnya ($\binom{4}{3}=4$),
pilih peringkat pasangan ($12$ cara), pilih dua dari empat jenisnya ($\binom{4}{2}=6$).
Jumlahnya $13\cdot 4\cdot 12\cdot 6=3{,}744$ dari $\binom{52}{5}=2{,}598{,}960$ tangan,
sehingga peluangnya $0.001441$, yaitu satu dari $694.2$ tangan. Dalam $694$ tangan,
banyaknya tangan penuh yang diharapkan $0.9998$, jadi kita menunggu sekitar tujuh ratus
tangan untuk melihat yang pertama.

Pencacahan juga menjelaskan kejadian yang terasa ajaib. Ambil masalah ulang tahun:
$k$ orang di satu ruangan, setiap hari lahir setara di antara $365$ hari, dan hari
lahir antarorang saling bebas. Peluang tidak ada dua orang berbagi hari lahir adalah
\[
P(\text{tidak ada pasangan})=\frac{365\cdot 364\cdots (365-k+1)}{365^{k}},
\]
sehingga peluang ada pasangan adalah pelengkapnya.

\index{paradoks ulang tahun}

\begin{center}
\small
\begin{tabular}{rrr}
\hline
$k$ & $P(\text{tidak ada pasangan})$ & $P(\text{ada pasangan})$ \\
\hline
$10$ & $0.8831$ & $0.1169$ \\
$23$ & $0.4927$ & $0.5073$ \\
$41$ & $0.0968$ & $0.9032$ \\
$50$ & $0.0296$ & $0.9704$ \\
$57$ & $0.0099$ & $0.9901$ \\
\hline
\end{tabular}
\end{center}

Pada $k=23$ peluangnya melewati setengah, dan pada $k=50$ sudah $0.9704$. Sumber
kebingungannya satuan yang salah: yang bertabrakan bukan orang, melainkan pasangan orang,
dan $23$ orang membentuk $\binom{23}{2}=253$ pasangan. Setiap pasangan cocok dengan
peluang $1/365$, sehingga banyak pasangan yang diharapkan $253/365=0.69315$. Bila
tabrakan dianggap jarang dan saling bebas, peluang tidak ada tabrakan mendekati
$e^{-0.69315}=0.5000$, berselisih $0.0073$ dari nilai eksak $0.4927$. Pendekatan itu
bukan bukti; ia menjelaskan mengapa ambangnya jatuh di sekitar dua puluh tiga.

Pencacahan di atas mengandaikan setiap hari lahir setara. Pada populasi nyata
ketidaksetaraan itu justru menaikkan peluang tabrakan; asumsi simetri di sini adalah
pendekatan yang perlu diperiksa.

Pencacahan tidak dapat menjawab pertanyaan klinik tadi. Hasil ``sakit'' dan ``sehat''
tidak sama mungkin: nisbahnya satu berbanding sembilan ratus sembilan puluh sembilan.
Yang dibutuhkan bukan cara menghitung hasil, melainkan cara memakai informasi bahwa
seseorang sudah dinyatakan positif. Itulah peluang bersyarat.

\section{Peluang bersyarat dan aturan perkalian}
\tingkat{menengah}

Informasi baru mempersempit dunia. Kalau kita tahu $B$ terjadi, dunia yang relevan bukan
lagi $\Omega$ melainkan $B$, dan semua bobot harus dibaca ulang di dalam dunia yang lebih
sempit itu. Kalimat sehari-harinya: peluang bersyarat adalah peluang yang dihitung ulang
setelah sebagian kemungkinan dicoret.

\kamus{peluang bersyarat}{conditional probability}{peluang yang dihitung ulang setelah
kita tahu sebagian kemungkinan sudah tidak berlaku}

\index{peluang bersyarat}

\begin{definisi}[Peluang bersyarat]
Bila $A$ dan $B$ kejadian dengan $P(B)>0$, peluang $A$ dengan syarat $B$ adalah
\[
P(A\mid B)=\frac{P(A\cap B)}{P(B)}.
\]
\end{definisi}

Pembagian dengan $P(B)$ adalah pilihan yang membuat peluang bersyarat berjumlah satu di
atas $B$. Bila kita juga menuntut $P(B\mid B)=1$ dan nisbah antar-kejadian di dalam $B$
tidak berubah, pembagian itulah satu-satunya pilihan; karena itu definisi ini bukan
konvensi sembarangan.

Penyempitan itu dapat dilihat langsung pada tabel silang. Gambar~\ref{fig:silang}
menyusun populasi klinik menjadi empat sel. Sebelum ada hasil, penyebutnya
seluruh $100{,}000$ orang dan $P(S)=100/100{,}000=0.001$. Setelah hasil positif
diketahui, yang masih dihitung hanya kolom $+$ yang berisi $1{,}098$ orang, dan
$P(S\mid+)=99/1{,}098=0.0902$. Orangnya tidak berubah; yang berubah hanya himpunan yang
masih boleh dihitung.

    \begin{figure}[htbp]
\centering
\begin{tikzpicture}[font=\small, x=1cm, y=1cm]
  \fill[accent!10] (1.95,1.05) rectangle (3.25,-1.1);
  \draw[layergrey, thick] (1.9,1.05)--(7.3,1.05);
  \draw[layergrey] (1.9,0.35)--(7.3,0.35);
  \draw[layergrey] (1.9,-0.35)--(7.3,-0.35);
  \draw[layergrey, thick] (1.9,-1.1)--(7.3,-1.1);
  \node[anchor=east] at (1.75,1.4) {hasil uji};
  \node at (2.6,1.4) {$+$};
  \node at (4.3,1.4) {$-$};
  \node at (6.5,1.4) {jumlah};
  \node[anchor=east] at (1.75,0.7) {sakit $S$};
  \node at (2.6,0.7) {99};
  \node at (4.3,0.7) {1};
  \node at (6.5,0.7) {100};
  \node[anchor=east] at (1.75,0) {sehat $S^{c}$};
  \node at (2.6,0) {999};
  \node at (4.3,0) {98\,901};
  \node at (6.5,0) {99\,900};
  \node[anchor=east] at (1.75,-0.7) {jumlah};
  \node at (2.6,-0.7) {1\,098};
  \node at (4.3,-0.7) {98\,902};
  \node at (6.5,-0.7) {100\,000};
  \draw[->, accent, thick] (2.6,-1.65)--(2.6,-1.15);
  \node[anchor=north, align=center, text width=10.4cm, font=\footnotesize] at (4.6,-1.95)
    {setelah $+$ diketahui, hanya kolom ini yang dihitung:
     $99/1{,}098=11/122=0.0902$};
\end{tikzpicture}
\caption{Tabel silang uji saring pada $100{,}000$ orang. Baris adalah keadaan
sebenarnya, kolom adalah hasil uji. Setelah hasil positif diketahui, ruang sampel
menyempit dari seluruh tabel ke kolom $+$ yang berisi $1{,}098$ orang, dan $99$
di antaranya benar-benar sakit.}
\label{fig:silang}
\end{figure}    

\begin{proposisi}[Peluang bersyarat adalah peluang]
\index{peluang bersyarat}
Untuk $B$ tetap dengan $P(B)>0$, fungsi $A\mapsto P(A\mid B)$ memenuhi ketiga aksioma
peluang pada $\mathcal{F}$.
\end{proposisi}

\begin{proof}[Bukti]
Ketidaknegatifan berasal dari $P(A\cap B)\ge 0$. Norma berasal dari $P(\Omega\mid B)=1$,
sebab $P(B)/P(B)=1$. Untuk aditivitas, bila $A_{1},A_{2},\dots$ saling lepas
maka $A_{i}\cap B$ juga saling lepas, sehingga
\[
P\Bigl(\bigcup_{i}A_{i}\mid B\Bigr)
=\frac{P\bigl(\bigcup_{i}(A_{i}\cap B)\bigr)}{P(B)}
=\frac{\sum_{i}P(A_{i}\cap B)}{P(B)}
=\sum_{i}P(A_{i}\mid B).\qedhere
\]
\end{proof}

Karena itu semua akibat aksioma berlaku untuk peluang bersyarat: komplemen, monotonisitas,
dan inklusi--eksklusi tetap sah bila setiap $P$ diganti $P(\cdot\mid B)$. Yang tidak
berlaku adalah menukar arah syaratnya.

Dari definisi langsung lahir aturan perkalian
\[
P(A\cap B)=P(A\mid B)\,P(B)=P(B\mid A)\,P(A),
\]
yang memperpanjang menjadi
$P(A_{1}\cap\cdots\cap A_{n})=P(A_{1})P(A_{2}\mid A_{1})\cdots
P(A_{n}\mid A_{1}\cap\cdots\cap A_{n-1})$. Aturan ini membaca percobaan bertahap
sebagai perkalian peluang di setiap tahap.

Satu tumpukan kartu $52$ memberi contoh arah yang mudah tertukar. Misalkan $A$ kejadian
kartu pertama hati dan $B$ kejadian kartu kedua merah. Maka
\[
P(A\cap B)=\frac{13\cdot 25}{52\cdot 51}=\frac{25}{204},
\]
sehingga
\[
P(A\mid B)=\frac{25/204}{1/2}=\frac{25}{102}=0.2451,
\qquad
P(B\mid A)=\frac{25/204}{1/4}=\frac{25}{51}=0.4902 .
\]
Dua angka yang berbeda dua kali lipat, dari data yang sama. Kesalahannya bukan pada
aritmetika, melainkan pada penempatan kejadian di kiri dan kanan garis tegak.

\begin{bacakeras}
Dua suara. Yang pertama belum tahu, yang kedua sudah tahu. Bacakan bergantian.

Suara pertama: ``Duniaku seluas lapangan. Semua pintu mungkin. Aku menebak dengan mata
tertutup.''

Suara kedua: ``Duniaku satu kamar. Satu pintu sudah terbuka, dan aku hanya menghitung
yang ada di dalamnya.''

Suara pertama: ``Kalau aku tahu sesuatu, apakah peluangku berubah?''

Suara kedua: ``Bukan peluangmu yang berubah. Dunia yang kamu hitung yang berubah. Dunia
mengecil setelah kita tahu sesuatu: bobot yang tadinya tersebar sekarang berkumpul di
tempat yang masih mungkin. Penyebutnya menyempit, dan angka yang keluar terdengar baru
walaupun aritmetikanya tidak pernah berubah sejak awal.''

Suara pertama: ``Jadi yang bergerak penyebutnya?''

Suara kedua: ``Yang bergerak penyebutnya, dan yang tinggal dihitung ulang adalah sisa
yang masih boleh ikut.''

Ulangi bersama, tiga kali, makin pelan: ``Dunia mengecil, bukan peluangnya yang
berkhianat. Yang berubah hanya yang masih boleh dihitung.''
\end{bacakeras}

Alat kedua memecah populasi menjadi beberapa kelompok yang tidak bertumpuk. Bila
$B_{1},\dots,B_{k}$ saling lepas dan gabungannya $\Omega$, maka
\[
P(A)=\sum_{i=1}^{k}P(A\mid B_{i})\,P(B_{i}).
\]

\begin{proposisi}[Hukum peluang total]
\index{peluang total}
Bila $B_{1},\dots,B_{k}$ saling lepas, gabungannya $\Omega$, dan setiap $P(B_{i})>0$,
maka $P(A)=\sum_{i}P(A\mid B_{i})P(B_{i})$ untuk setiap kejadian $A$.
\end{proposisi}

\begin{proof}[Bukti]
Kejadian $A\cap B_{1},\dots,A\cap B_{k}$ saling lepas dan gabungannya $A$. Jadi
$P(A)=\sum_{i}P(A\cap B_{i})$, dan setiap suku ditulis ulang dengan aturan perkalian
sebagai $P(A\mid B_{i})P(B_{i})$.
\end{proof}

Hukum ini tampak seperti formalitas sampai ia dipakai pada contoh klinik: ia yang
mengubah dua angka iklan menjadi satu angka yang berguna, dan itulah isi bagian
berikutnya.

Sebelum itu, satu contoh bahwa kalimat Indonesia dan kejadian matematis tidak selalu
sepadan. Sebuah keluarga punya dua anak, setiap anak laki-laki atau perempuan dengan
peluang sama, dan jenis kelamin keduanya saling bebas. Bila kita tahu \emph{paling sedikit
satu} anak laki-laki, peluang keduanya laki-laki $1/3=0.3333$: dari empat susunan setara
tersisa tiga, dan hanya satu yang cocok. Bila yang kita tahu \emph{anak yang lebih tua}
laki-laki, peluangnya $1/2=0.5$. Dua kalimat yang mirip memberi jawaban berbeda, sebab
yang pertama mengecualikan satu susunan dan yang kedua dua.

\section{Aturan Bayes dan tafsir hasil uji}
\tingkat{menengah}

Aturan Bayes membalik arah syarat: dari peluang hasil bila sakit, ia menghitung peluang
sakit bila hasil. Kalimat sehari-harinya: dari semua cara sebuah hasil dapat muncul,
berapa bagian yang lewat jalan yang kita tanyakan.

\index{aturan bayes}
\index{kelaziman}

\begin{teorema}[Aturan Bayes]
Bila $B_{1},\dots,B_{k}$ saling lepas dengan gabungan $\Omega$ dan $P(B_{i})>0$, maka
untuk setiap kejadian $A$ dengan $P(A)>0$,
\[
P(B_{j}\mid A)=\frac{P(A\mid B_{j})\,P(B_{j})}
{\sum_{i=1}^{k}P(A\mid B_{i})\,P(B_{i})}.
\]
\end{teorema}

\begin{proof}[Bukti]
Menurut definisi, $P(B_{j}\mid A)=P(A\cap B_{j})/P(A)$. Pembilangnya adalah
$P(A\mid B_{j})P(B_{j})$ menurut aturan perkalian, dan penyebutnya $P(A)$ menurut hukum
peluang total. Jadi rumus itu tidak memuat asumsi baru; ia definisi peluang bersyarat
yang ditulis dua kali dengan arah berbeda.
\end{proof}

Karena itu tidak ada ``metode Bayes'' yang berdiri di atas aksioma; yang ada adalah
pembukuan. Pembilang adalah jalur yang kita tanyakan, penyebut adalah jumlah semua jalur
yang menghasilkan hasil yang sama.

\begin{resep}
\begin{enumerate}
  \item Namai dulu dua hal: hipotesis $H$ yang ditanyakan dan bukti $E$ yang sudah
        terlihat. Tulis pertanyaannya sebagai $P(H\mid E)$, bukan $P(E\mid H)$.
  \item Tulis pembilang sebagai satu jalur saja: $P(E\mid H)\,P(H)$. Itu bobot dunia
        tempat $H$ benar.
  \item Tulis penyebut sebagai jumlah semua jalur yang menghasilkan $E$, termasuk jalur
        lewat $H^{c}$: $P(E\mid H)P(H)+P(E\mid H^{c})P(H^{c})$.
  \item Bagi, lalu periksa akal: hasilnya antara $0$ dan $1$, dan lebih kecil daripada
        $P(E\mid H)$ bila $H$ jarang.
\end{enumerate}
Peringatan urutan: pembilang dan penyebut sama-sama memakai $P(E\mid\cdot)$. Menulis
$P(\cdot\mid E)$ di penyebut adalah kesalahan yang paling sering terjadi, dan besarnya
angka tidak memberi tahu bahwa arahnya sudah tertukar.
\end{resep}

Pohon pada Gambar~\ref{fig:pohon} menuliskan seluruh jalan menuju hasil positif. Setiap
daun adalah darab peluang di cabangnya, dan dua daun bertanda $+$ berjumlah
$P(+)=0.01098$. Pembilang Bayes hanyalah daun pertama, sedangkan penyebutnya jumlah
kedua daun itu.

    \begin{figure}[htbp]
\centering
\begin{tikzpicture}[font=\small, every node/.style={inner sep=2pt}]
  \node (r) at (0,0) {awal};
  \node (s) at (2.2,1.7) {$S$};
  \node (sc) at (2.2,-1.7) {$S^{c}$};
  \node (sp) at (4.6,2.5) {$+$};
  \node (sm) at (4.6,0.9) {$-$};
  \node (scp) at (4.6,-0.9) {$+$};
  \node (scm) at (4.6,-2.5) {$-$};
  \draw[->] (r)--(s) node[midway, above left=-2pt] {$0.001$};
  \draw[->] (r)--(sc) node[midway, below left=-2pt] {$0.999$};
  \draw[->, accent, thick] (s)--(sp) node[midway, above left=-2pt] {$0.99$};
  \draw[->, layergrey] (s)--(sm) node[midway, below left=-2pt] {$0.01$};
  \draw[->, accent, thick] (sc)--(scp) node[midway, above left=-2pt] {$0.01$};
  \draw[->, layergrey] (sc)--(scm) node[midway, below left=-2pt] {$0.99$};
  \node[anchor=west, font=\footnotesize] at (5.0,2.5) {$0.001\cdot 0.99=0.00099$};
  \node[anchor=west, font=\footnotesize] at (5.0,0.9) {$0.001\cdot 0.01=0.00001$};
  \node[anchor=west, font=\footnotesize] at (5.0,-0.9) {$0.999\cdot 0.01=0.00999$};
  \node[anchor=west, font=\footnotesize] at (5.0,-2.5) {$0.999\cdot 0.99=0.98901$};
  \node[anchor=north, align=center, text width=10.6cm, font=\footnotesize] at (4.4,-3.0)
    {$P(+)=0.00099+0.00999=0.01098$, \quad $P(S\mid+)=0.00099/0.01098=0.0902$};
\end{tikzpicture}
\caption{Pohon dua tahap uji saring. Cabang pertama memisahkan sakit dari sehat,
cabang kedua memisahkan hasil positif dari negatif; darab di tiap daun adalah peluang
jalur itu, dan dua daun positif berjumlah $P(+)=0.01098$.}
\label{fig:pohon}
\end{figure}    

Kembali ke klinik. Tulis $S$ untuk kejadian sakit dan $+$ untuk hasil positif. Dari
iklan dan kelaziman penyakit kita punya $P(S)=0.001$, $P(+\mid S)=0.99$, dan
$P(+\mid S^{c})=0.01$. Hukum peluang total memberi
\[
P(+)=0.99\cdot 0.001+0.01\cdot 0.999=0.00099+0.00999=0.01098,
\]
sehingga
\[
P(S\mid +)=\frac{0.00099}{0.01098}=\frac{11}{122}=0.0902 .
\]
Suku positif palsu, $0.00999$, sepuluh kali lebih besar daripada suku deteksi benar,
$0.00099$; nisbahnya $999/99=10.09$. Bukan karena ujinya buruk: orang sehat $999$ kali
lebih banyak, $99{,}900$ berbanding $100$.

Perubahan kelaziman mengubah jawaban jauh lebih kuat daripada perubahan mutu uji.
Dengan sensitivitas dan spesifisitas tetap $0.99$, tabel berikut menunjukkan
$P(S\mid +)$ untuk beberapa kelaziman.

\begin{center}
\small
\begin{tabular}{rrrr}
\hline
$P(S)$ & $P(S\mid +)$ & $1/P(S\mid +)$ & tafsiran \\
\hline
$0.0001$ & $0.0098$ & $102.0$ & satu dari seratus dua \\
$0.001$ & $0.0902$ & $11.09$ & satu dari sebelas \\
$0.010$ & $0.5000$ & $2.00$ & satu dari dua \\
$0.100$ & $0.9167$ & $1.09$ & sebelas dari dua belas \\
\hline
\end{tabular}
\end{center}

Pada kelaziman $1$ persen, uji yang sama memberi $P(S\mid+)=1/2$ tepat, bukan $0.99$.
Pada kelaziman $10$ persen, $P(S\mid+)=11/12=0.9167$. Angka-angka ini menjelaskan
kebijakan kesehatan yang sering dikritik: uji saring massal untuk penyakit langka
menghasilkan banyak positif palsu, dan karena itu hasil positif biasanya diikuti uji
kedua, bukan pengobatan.

Uji kedua memakai aturan yang sama, dengan prior yang sudah diperbarui. Bila uji kedua
bebas dari yang pertama dan bermutu sama, maka
\[
P(S\mid +,+)=\frac{0.99\cdot (11/122)}{0.99\cdot (11/122)+0.01\cdot (111/122)}
=\frac{10.89}{10.89+1.11}=\frac{363}{400}=0.9075 .
\]
Perhatikan bahwa $0.9075$ bukan $0.99$ dan bukan $0.9999$. Dua uji positif menaikkan
peluang dari $0.0902$ ke $0.9075$, hampir sepuluh kali, tetapi belum memberi kepastian.
Yang berubah hanya kelaziman di dalam kelompok yang sudah positif; mutu uji tetap. Dengan
mutu $0.99$ ini, posterior melampaui $0.99$ hanya bila kelaziman awal sudah melewati
$1/2$.

Satu hal teknis: $P(S\mid+)=0.0902$ berlaku untuk populasi berkelaziman $1/1000$, bukan
untuk pasien yang datang dengan gejala. Gejala mengubah kelaziman lebih dahulu, dan
itulah sebabnya dokter menanyakan riwayat sebelum menafsirkan hasil laboratorium.

Contoh terakhir bagian ini melepaskan Bayes dari uji klinik. Tiga kantong berisi sepuluh
kelereng: kantong $A$ satu merah, $B$ lima merah, $C$ sembilan merah. Kita memilih
kantong secara acak, mengambil satu kelereng, dan yang keluar merah. Sebelum pengambilan
ketiga kantong berpeluang $1/3$; sesudahnya $0.0667$, $0.3333$, dan $0.6$, seperti
Gambar~\ref{fig:bayes-batang}. Tidak ada bobot yang hilang; bobot hanya berpindah ke
kantong yang lebih mudah memberi bukti itu.

    \begin{figure}[htbp]
\centering
\begin{tikzpicture}[font=\small, x=1cm, y=1cm]
  \draw[->, layergrey] (0,0)--(7.4,0);
  \draw[->, layergrey] (0,0)--(0,4.3);
  \foreach \y/\lab in {0/0, 1.2/0.2, 2.4/0.4, 3.6/0.6} {
    \draw[layergrey] (0,\y)--(-0.08,\y);
    \node[anchor=east, font=\footnotesize] at (-0.12,\y) {\lab};
  }
  \fill[layerindigo] (0.65,0) rectangle (1.20,2.0);
  \fill[layerteal] (1.40,0) rectangle (1.95,2.0);
  \fill[accent] (2.15,0) rectangle (2.70,2.0);
  \fill[layerindigo] (4.35,0) rectangle (4.90,0.4);
  \fill[layerteal] (5.10,0) rectangle (5.65,2.0);
  \fill[accent] (5.85,0) rectangle (6.40,3.6);
  \draw[dashed, layergrey] (0,2.0)--(6.9,2.0);
  \node[font=\footnotesize, anchor=south] at (0.92,2.02) {$1/3$};
  \node[font=\footnotesize, anchor=south] at (1.67,2.02) {$1/3$};
  \node[font=\footnotesize, anchor=south] at (2.42,2.02) {$1/3$};
  \node[font=\footnotesize, anchor=south] at (4.62,0.44) {$0.0667$};
  \node[font=\footnotesize, anchor=south] at (5.37,2.02) {$0.3333$};
  \node[font=\footnotesize, anchor=south] at (6.12,3.62) {$0.6$};
  \node[font=\footnotesize, anchor=north] at (0.92,-0.08) {$A$};
  \node[font=\footnotesize, anchor=north] at (1.67,-0.08) {$B$};
  \node[font=\footnotesize, anchor=north] at (2.42,-0.08) {$C$};
  \node[font=\footnotesize, anchor=north] at (4.62,-0.08) {$A$};
  \node[font=\footnotesize, anchor=north] at (5.37,-0.08) {$B$};
  \node[font=\footnotesize, anchor=north] at (6.12,-0.08) {$C$};
  \node[font=\footnotesize] at (1.67,-0.85) {sebelum};
  \node[font=\footnotesize] at (5.37,-0.85) {sesudah};
\end{tikzpicture}
\caption{Pembaruan Bayes pada tiga kantong. Sebelum ada bukti ketiga kantong berbagi
bobot $1/3$; setelah satu kelereng merah terlihat, bobot berpindah menjadi $0.0667$,
$0.3333$, dan $0.6$. Kantong yang paling mudah menghasilkan bukti itu naik paling
tinggi.}
\label{fig:bayes-batang}
\end{figure}    

\section{Independensi, dan mengapa saling lepas bukan independen}
\tingkat{lanjut}

Dua kejadian disebut saling lepas bila keduanya tidak dapat terjadi bersamaan. Dua
kejadian disebut independen bila mengetahui yang satu tidak mengubah peluang yang lain.
Kedua kalimat itu kedengarannya serupa, dan menukarnya adalah kekeliruan paling mahal
di bab ini. Kalimat sehari-harinya: saling lepas berarti tidak pernah bersama;
independen berarti tidak saling memberi kabar.

\kamus{independen}{independent}{kejadian yang tidak saling memberi kabar: mengetahui
satu tidak mengubah peluang yang lain}

\begin{untukanak}
Ada dua hal yang sering tertukar, dan bedanya dapat dilihat tanpa angka. Di satu
sekolah, dua siswa tidak pernah masuk kelas bersamaan: yang satu belajar pagi, yang satu
belajar siang. Kalau kita tahu Andi sudah masuk, kita tahu Budi belum masuk. Mereka tidak
pernah bersama, dan justru karena itu kabar tentang yang satu memberi kabar tentang yang
lain.

Sekarang dua siswa lain. Keduanya masuk kelas yang sama, tetapi duduk berjauhan dan
mengerjakan ujian yang berbeda. Nilai mereka tidak saling memberi kabar: tahu nilai Sari
tidak menolong kita menebak nilai Tono.

Kelompok pertama disebut tidak pernah bersama. Kelompok kedua disebut tidak saling
memberi kabar. Keduanya kedengaran seperti ``tidak ada hubungan'', padahal yang pertama
justru hubungan yang paling kuat: kabar tentang yang satu menentukan yang lain. Di bab
ini kita akan tahu mengapa dua hal yang tidak pernah bersama tidak mungkin saling memberi
kabar, dan mengapa menukar keduanya membuat hitungan meleset bukan sedikit.
\end{untukanak}

\index{independensi}
\index{kejadian saling lepas}

\begin{definisi}[Independensi dua kejadian]
Kejadian $A$ dan $B$ \emph{independen} bila
\[
P(A\cap B)=P(A)\,P(B).
\]
\end{definisi}

Definisi ini simetris, dan bila $P(B)>0$ ia setara dengan $P(A\mid B)=P(A)$. Perhatikan
bahwa definisi itu tidak menyebut sebab akibat: dua kejadian dapat independen tanpa satu
menyebabkan yang lain, dan dapat bergantung tanpa satu menyebabkan yang lain.

\begin{proposisi}[Saling lepas bukan independen]
\index{kejadian saling lepas}
Bila $A$ dan $B$ saling lepas dan $P(A)>0$ serta $P(B)>0$, maka $A$ dan $B$ tidak
independen. Lebih tepatnya $P(A\cap B)=0<P(A)P(B)$, dan $P(A\mid B)=0<P(A)$.
\end{proposisi}

\begin{proof}[Bukti]
Saling lepas berarti $A\cap B=\varnothing$, sehingga $P(A\cap B)=P(\varnothing)=0$.
Karena $P(A)P(B)>0$, kesamaan definisi independensi gagal. Untuk pernyataan bersyarat,
$P(A\mid B)=0/P(B)=0$, yang lebih kecil daripada $P(A)>0$.
\end{proof}

Kekeliruannya bergerak dua arah, dan kedua arah itu dapat dihitung.

Ambil satu dadu dan dua pasang kejadian dengan marginal sama tetapi hubungan berbeda.
Pasangan pertama: $A=\{1,2\}$ dan $B=\{3,4\}$, saling lepas, dengan $P(A)=P(B)=1/3$.
Pasangan kedua: $C=\{1,2,3\}$ dan $D=\{1,4\}$, dengan $P(C)=1/2$ dan $P(D)=1/3$. Untuk
pasangan pertama,
\[
P(A)P(B)=\frac{1}{9}=0.1111,\qquad P(A\cap B)=0,\qquad P(A\mid B)=0,
\]
jadi selisih definisi independensi $0.1111$ dan selisih peluang bersyaratnya $0.3333$,
sebab $P(A)=1/3$. Untuk pasangan kedua,
\[
P(C)P(D)=\frac{1}{6}=0.1667=P(C\cap D),\qquad P(C\mid D)=\frac{1}{2}=P(C),
\]
jadi kedua selisihnya nol.

Yang penting bukan angka $1/6$ itu, melainkan bahwa dua pasangan dapat bermarginal sama
dengan hubungan berlawanan total. Pasangan $E=\{1,2,3\}$ dan $F=\{4,6\}$ juga bermarginal
$1/2$ dan $1/3$ tetapi saling lepas. Marginal tidak menentukan hubungan; hanya irisan
yang menentukannya, sehingga setiap klaim independensi adalah klaim tambahan tentang
dunia.

Untuk tiga kejadian, independensi berpasangan tidak cukup. Kejadian $A,B,C$ disebut
saling independen bila $P(A\cap B)=P(A)P(B)$, $P(A\cap C)=P(A)P(C)$,
$P(B\cap C)=P(B)P(C)$, dan $P(A\cap B\cap C)=P(A)P(B)P(C)$.

\begin{contoh}[Berpasangan tetapi tidak saling independen]
\index{independensi berpasangan}
Dua koin dilempar. Misalkan $A$ koin pertama gambar, $B$ koin kedua gambar, dan $C$
kedua koin sama; masing-masing berpeluang $1/2$. Setiap
pasangan beririsan satu hasil, jadi $P(A\cap B)=P(A\cap C)=P(B\cap C)=1/4$, yang sama
dengan $1/2\cdot 1/2$. Tetapi $A\cap B\cap C$ terjadi hanya bila keduanya gambar,
sehingga peluangnya $1/4$, sedangkan $P(A)P(B)P(C)=1/8$. Jadi ketiganya berpasangan
independen dan tidak saling independen; selisihnya $0.125$.
\end{contoh}

Peringatan terakhir: independensi tidak bertahan di bawah pengkondisian. Bila $A$ dan
$B$ independen berpeluang $1/2$, maka dengan syarat $A\cup B$, $P(A\mid A\cup B)=2/3$
sedangkan $P(A\cap B\mid A\cup B)=1/3\neq 4/9$. Menggabungkan dua kejadian independen
lalu mengkondisikan pada gabungannya justru menciptakan ketergantungan, sehingga kalimat
``$A$ dan $B$ independen'' selalu relatif terhadap ruang peluang yang dipakai.

\begin{tangan}
Bahan: dua belas kartu, enam pasang berperingkat sama (As sampai enam), setiap pasangan
satu merah dan satu hitam.

Langkah pertama, buat pasangan. Susun dua belas kartu terbuka dalam dua baris; letakkan
setiap pasangan seperingkat bersebelahan, merah di kiri, hitam di kanan. Periksa dengan
mata: peringkat tidak memberi kabar tentang warna, sebab setiap peringkat punya merah
dan hitam. Dua sifat itu tidak saling memberi kabar.

Langkah kedua, saring. Kocok, tutup semua, lalu buka satu kartu: misalnya empat hati.
Beri satu petunjuk, ``pasangan kartu ini berwarna hitam'', lalu sisihkan semua kartu
merah yang masih tertutup. Dari sebelas kartu tertutup, tinggal enam yang hitam.

Yang harus kaulihat: penyebutnya berubah dari sebelas menjadi enam, sedangkan kartunya
tidak berubah sama sekali. Peluang menebak pasangannya naik dari satu per sebelas menjadi
satu per enam. Dua sifat yang tadi tidak saling memberi kabar sekarang saling memberi
kabar, dan itu yang dilakukan satu petunjuk.
\end{tangan}

\section{Dua paradoks yang menuntut hitungan}
\tingkat{lanjut}

Paradoks dalam peluang hampir selalu tabrakan antara intuisi dan definisi, dan
tabrakan itu hanya selesai bila dihitung. Dua contoh berikut menunjukkan hal berbeda:
jawaban bergantung pada aturan main, dan pada satuan yang dihitung.

\index{monty hall}

Monty Hall. Tiga pintu tertutup; satu berisi mobil, dua berisi kambing. Kita memilih
satu pintu. Pembawa acara, yang tahu isi semua pintu, membuka satu pintu lain yang
berisi kambing, lalu menawarkan kesempatan berpindah ke pintu yang masih tertutup.
Apakah berpindah menaikkan peluang?

Hitung dengan hukum peluang total terhadap tiga kemungkinan letak mobil, masing-masing
berpeluang $1/3$, dan andaikan kita mula-mula memilih pintu $1$. Bila mobil di pintu
$1$, berpindah kalah. Bila mobil di pintu $2$, pembawa acara wajib membuka pintu $3$,
sehingga berpindah menang. Bila mobil di pintu $3$, ia membuka pintu $2$, dan berpindah
menang. Jadi berpindah menang pada dua dari tiga kemungkinan:
\[
P(\text{menang dengan berpindah})=\frac{2}{3},\qquad
P(\text{menang tanpa berpindah})=\frac{1}{3}.
\]
Selisihnya $1/3$, bukan nol.

Argumen ``dua pintu tersisa, jadi $1/2$'' keliru karena mengabaikan bahwa pembawa acara
tidak memilih pintu secara bebas. Ia dilarang membuka pintu berisi mobil dan dilarang
membuka pintu pilihan kita, sehingga tindakannya membawa informasi tentang pintu yang
tidak ia buka.

Informasi itu hilang bila aturannya berubah. Andaikan pembawa acara membuka satu dari
dua pintu lain secara acak, tanpa melihat isinya, dan permainan dibatalkan bila mobil
terlihat. Dengan syarat permainan berjalan sampai tawaran,
\[
P(\text{mobil di pintu }1\mid \text{kambing terlihat})
=\frac{1/3}{1/3+1/6+1/6}=\frac{1}{2},
\]
sebab bila mobil di pintu $2$ atau $3$, ia hanya punya peluang $1/2$ menampilkan kambing.
Pada aturan itu berpindah tidak menolong: ``peluang berpindah $2/3$'' bukan sifat pintu,
melainkan sifat protokol.

\index{paradoks ulang tahun}

Ulang tahun. Bagian sebelumnya sudah menghitung $0.5073$ untuk $23$ orang dan $0.9704$
untuk $50$ orang. Paradoksnya bukan pada angkanya, melainkan pada satuan: kita
membandingkannya dengan peluang satu orang tertentu cocok dengan kita, padahal yang
dihitung $253$ pasangan sekaligus. Kesalahan yang sama muncul pada tabel hash dan uji
kecocokan DNA.

Kekeliruan jaksa adalah kesalahan arah yang sama pada perkara pidana. Misalkan peluang
sebuah ciri muncul secara kebetulan satu per satu juta, dan ciri itu ditemukan pada
tersangka yang diambil dari basis data sejuta orang tanpa bukti lain. Maka
$P(\text{cocok}\mid\text{tidak bersalah})=10^{-6}$, tetapi
\[
P(\text{bersalah}\mid\text{cocok})
=\frac{10^{-6}\cdot 1}{10^{-6}\cdot 1+(1-10^{-6})\cdot 10^{-6}}=0.5000003,
\]
yaitu praktis setengah, bukan hampir pasti. Dua bilangan $10^{-6}$ itu besaran berbeda:
yang pertama peluang ciri secara kebetulan, yang kedua kelaziman tersangka di dalam basis
data. Menukarnya adalah wajah lain dari kekeliruan yang sama, kali ini di ruang
pengadilan.

\begin{tantangan}
Tiga kartu di atas meja. Kartu pertama merah di kedua sisinya, kartu kedua hitam di kedua
sisinya, kartu ketiga merah di satu sisi dan hitam di sisi lain. Kita mengocok ketiganya
tertutup, mengambil satu, dan meletakkannya di meja. Sisi yang terlihat merah. Berapa
peluang sisi di baliknya juga merah?

Hampir semua orang menjawab setengah: katanya tinggal dua kartu yang punya sisi merah,
jadi dua kemungkinan. Jawaban yang benar bukan setengah.

Petunjuk 1: jangan menghitung kartu, hitung sisi. Ada berapa sisi merah seluruhnya, dan
berapa di antaranya bersebelahan dengan sisi merah?

Petunjuk 2: sisi merah yang terlihat dapat berasal dari tiga sisi merah yang berbeda,
dan dua di antaranya berada pada kartu merah-merah.

Kalau mentok, kembali ke bagian peluang bersyarat dan tulis ulang ruang sampelnya sebagai
enam sisi, bukan tiga kartu.
\end{tantangan}

\section{Contoh hitung}
\tingkat{menengah}

Tiga hitungan berikut dikerjakan dengan pecahan eksak, supaya tidak ada pembulatan yang
menyembunyikan langkah. Semuanya sudah diperiksa dengan Python, dan angka yang muncul
adalah angka yang sama dengan yang dipakai di bagian sebelumnya.

Pertama, Bayes dua langkah. Dengan $P(S)=1/1000$, $P(+\mid S)=99/100$, dan
$P(+\mid S^{c})=1/100$, hukum peluang total memberi
\[
P(+)=\frac{99}{100}\cdot\frac{1}{1000}+\frac{1}{100}\cdot\frac{999}{1000}
=\frac{99+999}{100{,}000}=\frac{1098}{100{,}000}.
\]
Pembilangnya dua kelompok terpisah: $99$ orang sakit yang terdeteksi dan $999$ orang
sehat yang salah dituduh. Maka
\[
P(S\mid +)=\frac{99/100{,}000}{1098/100{,}000}=\frac{99}{1098}=\frac{11}{122}=0.090164 .
\]
Uji kedua bermutu sama memakai $11/122$ sebagai kelaziman baru dan memberi $0.9075$,
sama seperti hitungan di bagian sebelumnya. Penyebutnya $1200$, bukan angka berjuta-juta:
faktor $1/122$ muncul di kedua suku dan saling menghapus, tanda bahwa yang berubah hanya
kelaziman, bukan mutu uji.

Kedua, saling lepas versus independen, dengan marginal yang disamakan.

\begin{center}
\small
\begin{tabular}{lrrrr}
\hline
pasangan & $P(\text{pertama})$ & $P(\text{kedua})$ & $P(\text{irisan})$ & $P(\text{pertama}\mid\text{kedua})$ \\
\hline
$C=\{1,2,3\}$, $D=\{1,4\}$ & $0.5$ & $0.3333$ & $0.1667$ & $0.5$ \\
$E=\{1,2,3\}$, $F=\{4,6\}$ & $0.5$ & $0.3333$ & $0$ & $0$ \\
\hline
\end{tabular}
\end{center}

Pada baris pertama $P(C)P(D)=1/6=0.1667$ sama dengan irisan, jadi selisihnya $0$. Pada
baris kedua $P(E)P(F)=1/6$ sedangkan irisannya $0$, jadi selisihnya $0.1667$; peluang
bersyaratnya $0$ sedangkan $P(E)=0.5$. Siapa pun yang menyangka ``tidak bisa terjadi
bersama'' berarti ``tidak saling memengaruhi'' akan melaporkan $0.1667$ di tempat yang
seharusnya $0$.

Ketiga, Monty Hall dengan protokol yang dinyatakan. Bila pembawa acara selalu membuka
pintu kambing, peluang menang dengan berpindah $2/3=0.6667$ dan tanpa berpindah
$1/3=0.3333$. Bila ia membuka secara acak dan permainan dibatalkan ketika mobil
terlihat, permainan berjalan dengan peluang $2/3$, dan dengan syarat itu kedua strategi
bernilai $1/2=0.5$. Dua protokol, satu himpunan pintu, dua jawaban.

\kertas{Untuk uji dengan kelaziman $1/2000$, sensitivitas $0.95$, dan spesifisitas
$0.90$, hitung $P(S\mid+)$ dengan tabel $100{,}000$ orang seperti di atas, lalu hitung
lagi dengan spesifisitas $0.99$. Tulis kedua jawaban sebagai pecahan dan sebagai desimal
empat angka, dan jelaskan bagian mana yang berubah paling banyak: pembilang atau
penyebut. Setelah itu hitung $P(S\mid +,+)$ untuk uji kedua yang bermutu sama, dan
periksa apakah hasilnya masih di bawah $0.95$.}

\section{Lab komputasi: Monte Carlo sebagai alat verifikasi}
\tingkat{lanjut}

Semua angka di bab ini diturunkan, bukan diukur. Simulasi tidak menggantikan penurunan,
tetapi ia menangkap kesalahan yang paling sering terjadi: arah syarat yang tertukar,
penyebut yang keliru, dan asumsi independensi yang dipakai tanpa dasar. Bangkitkan
percobaan yang sama banyak kali, hitung nisbah keberhasilan, lalu bandingkan dengan nilai
eksak. Menurut hukum bilangan besar (Bab 19), nisbah itu menuju peluang sebenarnya; yang
tersisa hanya seberapa cepat, diukur oleh galat baku $\sqrt{p(1-p)/m}$ dengan $m$ banyak
percobaan yang menghasilkan kejadian itu.

\index{monte carlo}
\index{galat baku}

Benih tetap dipakai supaya keluarannya dapat diperiksa ulang kata demi kata. Benih tetap
tidak membuat hasilnya lebih benar; ia hanya membuatnya dapat diulang.

\begin{lab}
    \begin{verbatim}
import numpy as np

rng = np.random.default_rng(20261024)

print("== 1. uji diagnostik: prevalensi 0,1%, sens 99%, spek 99% ==")
n = 1_000_000
sakit = rng.random(n) < 0.001
positif = np.where(sakit, rng.random(n) < 0.99, rng.random(n) < 0.01)
m = int(positif.sum())
ppv = sakit[positif].mean()
galat = np.sqrt(ppv * (1 - ppv) / m)
eksak = 0.99 * 0.001 / (0.99 * 0.001 + 0.01 * 0.999)
print("sakit=%d  positif=%d  benar_positif=%d"
      % (sakit.sum(), m, int((sakit & positif).sum())))
print("PPV empiris=%.6f  eksak=%.6f  galat baku=%.6f  z=%.3f"
      % (ppv, eksak, galat, (ppv - eksak) / galat))

print("== 2. Monty Hall, 10^5 permainan ==")
N = 100_000
mobil = rng.integers(0, 3, N)
pilih = rng.integers(0, 3, N)
buka = np.where(pilih == mobil,
                (pilih + 1 + rng.integers(0, 2, N)) % 3,
                3 - pilih - mobil)
ganti = 3 - pilih - buka
menang_tetap = float((pilih == mobil).mean())
menang_ganti = float((ganti == mobil).mean())
se_ganti = np.sqrt(menang_ganti * (1 - menang_ganti) / N)
print("menang tanpa ganti=%.4f  menang ganti=%.4f  galat baku=%.5f"
      % (menang_tetap, menang_ganti, se_ganti))
print("eksak 1/3=%.6f  2/3=%.6f" % (1/3, 2/3))

print("== 3. pembawa acara membuka pintu secara acak ==")
buka_acak = (pilih + 1 + rng.integers(0, 2, N)) % 3
sah = buka_acak != mobil
ganti_acak = 3 - pilih - buka_acak
print("permainan sah=%d dari %d (%.4f)" % (sah.sum(), N, sah.mean()))
print("menang ganti (sah)=%.4f  tanpa ganti (sah)=%.4f"
      % (float((ganti_acak[sah] == mobil[sah]).mean()),
         float((pilih[sah] == mobil[sah]).mean())))

print("== 4. peluang ulang tahun, 20.000 kamar ==")
for k in (23, 50):
    hari = rng.integers(0, 365, size=(20000, k))
    hari.sort(axis=1)
    ada = (np.diff(hari, axis=1) == 0).any(axis=1)
    q = 1.0
    for i in range(k):
        q *= (365 - i) / 365
    f = float(ada.mean())
    se = np.sqrt(f * (1 - f) / 20000)
    print("k=%2d  empiris=%.4f  eksak=%.4f  galat baku=%.5f  z=%+.2f"
          % (k, f, 1 - q, se, (f - (1 - q)) / se))

print("== 5. saling lepas versus independen, 10^6 lemparan dadu ==")
N = 1_000_000
dadu = rng.integers(1, 7, N)
A = np.isin(dadu, [1, 2])
B = np.isin(dadu, [3, 4])
C = np.isin(dadu, [1, 2, 3])
D = np.isin(dadu, [1, 4])
print("lepas:  P(A)=%.4f P(B)=%.4f P(AnB)=%.4f P(A|B)=%.4f selisih=%.4f"
      % (A.mean(), B.mean(), (A & B).mean(), (A & B).sum() / B.sum(),
         A.mean() * B.mean() - (A & B).sum() / B.sum()))
print("bebas:  P(C)=%.4f P(D)=%.4f P(CnD)=%.4f P(C|D)=%.4f selisih=%.4f"
      % (C.mean(), D.mean(), (C & D).mean(), (C & D).sum() / D.sum(),
         abs(C.mean() - (C & D).sum() / D.sum())))
print("eksak lepas: 1/3, 1/3, 0, 0, 1/9=%.6f" % (1 / 9))
print("eksak bebas: 1/2, 1/3, 1/6, 1/2, 0")
\end{verbatim}    
\end{lab}

Keluaran sesungguhnya, apa adanya:

    \begin{verbatim}
== 1. uji diagnostik: prevalensi 0,1%, sens 99%, spek 99% ==
sakit=963  positif=11094  benar_positif=952
PPV empiris=0.085812  eksak=0.090164  galat baku=0.002659  z=-1.637
== 2. Monty Hall, 10^5 permainan ==
menang tanpa ganti=0.3325  menang ganti=0.6675  galat baku=0.00149
eksak 1/3=0.333333  2/3=0.666667
== 3. pembawa acara membuka pintu secara acak ==
permainan sah=66472 dari 100000 (0.6647)
menang ganti (sah)=0.4998  tanpa ganti (sah)=0.5002
== 4. peluang ulang tahun, 20.000 kamar ==
k=23  empiris=0.5019  eksak=0.5073  galat baku=0.00354  z=-1.54
k=50  empiris=0.9701  eksak=0.9704  galat baku=0.00120  z=-0.19
== 5. saling lepas versus independen, 10^6 lemparan dadu ==
lepas:  P(A)=0.3333 P(B)=0.3335 P(AnB)=0.0000 P(A|B)=0.0000 selisih=0.1112
bebas:  P(C)=0.5001 P(D)=0.3334 P(CnD)=0.1667 P(C|D)=0.5000 selisih=0.0001
eksak lepas: 1/3, 1/3, 0, 0, 1/9=0.111111
eksak bebas: 1/2, 1/3, 1/6, 1/2, 0
\end{verbatim}    

Empat hal perlu dibaca dari keluaran itu, dan dua di antaranya justru menunjukkan batas
simulasi.

Pertama, baris pertama bagian 1 menyebut $963$ orang sakit dari sejuta, sedangkan
harapannya $1{,}000$ dengan simpangan baku $\sqrt{10^{6}(0.001)(0.999)}=31.6$. Jadi
$963$ berada $1.17$ simpangan baku di bawah harapan: ayunan wajar, bukan kelalaian
program. Harapan hasil positif menjadi $963\cdot 0.99+999{,}037\cdot 0.01=10{,}943.7$,
dan yang teramati $11{,}094$ berjarak $1.4$ simpangan baku. Peluang bersyarat empiris
$0.0858$ berbeda dari eksak $0.0902$ sebesar $0.0044$, yaitu $1.64$ galat baku. Selisih
itu jauh lebih kecil daripada jarak $0.09$ ke $1$, sehingga tafsirnya tidak berubah.

Kedua, Monty Hall memberi $0.6675$ dan $0.3325$ dengan galat baku $0.0015$, sedangkan
selisih yang diharapkan antarstrategi $1/3$. Kesepakatan sampai tiga angka itu bukan
kebetulan; bentuk verifikasi yang paling berguna memang bukan memastikan angka ketiga,
melainkan memastikan selisih yang kita klaim ada dan bertanda benar.

Ketiga, bagian 3 dan 5 menunjukkan bahwa protokol dan hubungan antar-kejadian masuk ke
angka. Pembawa acara yang membuka pintu secara acak hanya membiarkan $66{,}472$ dari
$100{,}000$ permainan berjalan sampai tawaran, dan di antaranya berpindah menang
$0.4998$ sedangkan bertahan $0.5002$; selisih $0.0004$ itu seperempat galat baku.
Pasangan saling lepas memberi $P(A)P(B)-P(A\cap B)=0.1112$, cocok dengan $1/9$ sampai
tiga angka, sedangkan pasangan independen memberi $0.0001$, yaitu nol dalam ayunan acak.
Angka $0.1112$ adalah besar kesalahan orang yang menyangka dua kejadian saling lepas
pasti independen.

Yang tidak dibuktikan program: tidak satu pun dari kelima bagian itu membuktikan rumus.
Program hanya memeriksa rumus yang sudah dibuktikan, dengan sampel jauh lebih besar
daripada jangkauan tangan; Bab 19 menjelaskan mengapa itu sah.

\begin{untukdosen}
Urutan mengajar yang saya sarankan: mulai dari tabel $100{,}000$ orang, bukan dari
aksioma. Suruh mahasiswa mengisi empat sel: sakit dan positif, sakit dan negatif, sehat
dan positif, sehat dan negatif. Angka $99$, $1$, $999$, dan $98{,}901$ muncul sendiri,
dan dari situ $99/1{,}098$ menjadi kesimpulan kelas, bukan hasil ceramah. Sesudah itu
aksioma, akibatnya, pencacahan, peluang bersyarat, Bayes, lalu independensi. Monty Hall
paling baik ditaruh paling akhir, setelah mahasiswa mampu menuliskan protokol pembawa
acara sebagai kejadian bersyarat.

Dua miskonsepsi yang paling sering muncul. Pertama, ``saling lepas berarti independen''.
Gejalanya terlihat di aritmetika: mahasiswa menjumlahkan peluang dua kejadian yang
seharusnya dikalikan, atau sebaliknya, lalu tidak dapat menjelaskan mengapa hasilnya
melebihi satu. Kedua, ``sensitivitas sama dengan peluang sakit bila hasil positif''.
Gejalanya adalah jawaban $0.99$ untuk soal klinik, tanpa memakai kelaziman yang sudah
diberikan di soal.

Aktivitas papan tulis yang berhasil: dua mahasiswa memerankan pembawa acara dengan
protokol tertulis berbeda, satu selalu membuka pintu kambing dan satu membuka acak; kelas
mencatat dua puluh permainan lalu membandingkannya dengan $2/3$ dan $1/2$. Mahasiswa siap
lanjut bila dapat menyebut satu contoh dua kejadian independen dan satu contoh saling
lepas dengan marginal sama, serta menghitung peluang bersyarat dari tabel dua arah tanpa
membalik arah.
\end{untukdosen}

\section{Ringkasan}

\begin{ringkasan}
\begin{itemize}
  \item Peluang adalah fungsi pada kejadian yang tidak negatif, bernilai satu pada
        seluruh ruang, dan aditif terbilang pada kejadian saling lepas; komplemen,
        monotonisitas, dan inklusi--eksklusi diturunkan dari situ.
  \item Bila hasil sama mungkin, peluang adalah nisbah pencacahan, dengan aturan
        perkalian dan $\binom{n}{k}$; asumsi kesetaraan itu harus diperiksa.
  \item Peluang bersyarat $P(A\mid B)=P(A\cap B)/P(B)$ menghitung ulang di dalam $B$
        dan memenuhi ketiga aksioma; aturan perkalian dan hukum peluang total adalah
        akibatnya.
  \item Aturan Bayes bukan asumsi baru, melainkan definisi yang dibaca dua arah.
        Kelaziman menentukan jawaban lebih kuat daripada mutu uji: dengan sensitivitas
        dan spesifisitas $0.99$, kelaziman $0.1$ persen memberi $P(S\mid+)=0.0902$.
  \item Saling lepas dan independen berbeda; bila kedua peluang positif, kejadian
        saling lepas tidak mungkin independen, dan hanya irisan yang menentukan
        hubungan.
  \item Jawaban paradoks bergantung pada protokol dan satuan hitung: Monty Hall $2/3$
        pada protokol tetap dan $1/2$ pada protokol acak; ulang tahun $0.5073$ untuk
        $23$ orang karena yang dihitung $253$ pasangan. Simulasi berbenih tetap
        memeriksa, bukan membuktikan.
\end{itemize}
\end{ringkasan}

\section{Latihan}

\begin{latihan}
\begin{enumerate}
  \item Buktikan rumus inklusi--eksklusi untuk tiga kejadian dengan menerapkan versi
        dua kejadian pada $A\cup B$ dan $C$. Tunjukkan pada langkah mana suku
        $P(A\cap B\cap C)$ muncul, dan mengapa tandanya positif.
  \item Nilai klaim ini: ``Karena $P(A\mid B)=0$, $A$ dan $B$ saling lepas.'' Benar,
        salah, atau belum cukup bukti? Beri contoh dengan $P(A\mid B)=0$ tetapi
        $A\cap B\neq\varnothing$, atau jelaskan mengapa itu tidak mungkin.
  \item Sebuah uji berkelaziman $2$ persen, sensitivitas $0.95$, spesifisitas $0.90$.
        Hitung $P(S\mid+)$ dari tabel $10{,}000$ orang, lalu ulangi dengan
        spesifisitas $0.99$; nyatakan keduanya sebagai pecahan.
  \item Buktikan bahwa bila $A$ dan $B$ independen, maka $A$ dan $B^{c}$ juga
        independen. Gunakan $P(A)=P(A\cap B)+P(A\cap B^{c})$ dan jelaskan mengapa
        hasil yang sama berlaku untuk $A^{c}$ dan $B^{c}$.
  \item Hitung peluang ada pasangan hari lahir sama untuk $30$ orang, bandingkan
        dengan $1-e^{-\binom{30}{2}/365}$, lalu jelaskan sumber selisihnya.
  \item Buktikan hukum peluang total untuk $k=2$ langsung dari aksioma, tanpa memakai
        hukum untuk $k$ umum: tulis $A=(A\cap B)\cup(A\cap B^{c})$ dan periksa
        kesalinglepasannya.
  \item Empat pintu, satu mobil, tiga kambing. Pembawa acara yang tahu isi membuka dua
        pintu kambing di antara tiga pintu lain dan menawarkan perpindahan ke satu pintu
        tersisa. Hitung peluang menang bila berpindah, lalu periksa dengan simulasi
        $10^{4}$ permainan.
  \item Nilai klaim ini: ``Dua uji bebas bersensitivitas dan spesifisitas $0.99$
        memberi kepastian $0.9999$ bila keduanya positif, asalkan kelazimannya $0.1$
        persen.'' Hitung angka yang benar dan sebutkan asumsinya.
  \item Apakah independensi berpasangan untuk tiga kejadian cukup untuk menyimpulkan
        $P(A\cap B\cap C)=P(A)P(B)P(C)$? Bantah dengan contoh dari dua lemparan koin,
        dan hitung kedua nilainya.
\end{enumerate}

Petunjuk untuk tiga butir tersulit. Butir 2: pada ruang berhingga yang setiap hasilnya
berpeluang positif, $P(A\mid B)=0$ memaksa $A\cap B=\varnothing$, jadi pada lemparan dadu
klaim itu benar; contoh penyangkal harus datang dari ruang kontinu, misalnya
$A=\{1/2\}$ dan $B=(0,1)$ dengan peluang seragam, tempat $P(A)=0$ walaupun $A$ berisi
titik. Butir 7: mula-mula mobil di pintu pilihan berpeluang $1/4$; bila tidak, satu-satunya
pintu yang tidak dibuka pembawa acara pasti berisi mobil, jadi peluang berpindah $3/4$.
Butir 9: ambil $A$ koin pertama gambar, $B$ koin kedua gambar, dan $C$ kedua koin sama;
ketiga pasangan memberi $1/4$ sedangkan irisannya $1/4$, bukan $1/8$.
\end{latihan}

\begin{untukpeneliti}
Yang benar-benar terbukti di bab ini: ketiga aksioma beserta akibatnya; aturan Bayes
sebagai teorema yang diturunkan dari definisi peluang bersyarat; kemustahilan dua
kejadian saling lepas berpeluang positif untuk juga independen; dan ketergantungan
jawaban Monty Hall pada protokol, yang analisisnya dirumuskan tegas oleh Morgan,
Chaganty, Dahiya, dan Doviak, \emph{The American Statistician} 1991, pada artikel
``Let's Make a Deal: The Player's Dilemma''. Yang populer tetapi tidak terbukti di sini
adalah klaim bahwa format frekuensi alami memperbaiki penalaran kelaziman pada manusia;
itu klaim empiris tentang pembaca, bukan teorema, dan buktinya ada di literatur
psikologi pengambilan keputusan.
\end{untukpeneliti}

Sumber primer: Blitzstein dan Hwang, \emph{Introduction to Probability}, Bab 1 (Why
probability, Naive definition of probability, How to count, Story proofs, Non-naive
definition of probability) dan Bab 2 (The importance of thinking conditionally,
Definition and intuition, Bayes' rule and the law of total probability, Conditional
probabilities are probabilities, Independence of events, Coherency of Bayes' rule);
serta Wasserman, \emph{All of Statistics}, Bab 1 bagian Probability untuk perumusan
aksioma dan Bab 2 bagian Independent Events.

Batas keberlakuan: seluruh bab bekerja pada ruang berhingga yang setiap hasilnya
berpeluang positif, tempat $P(A)=0$ setara dengan $A=\varnothing$. Pada ruang tak
tercacah, pengkondisian pada kejadian berpeluang nol tidak didefinisikan oleh rumus
nisbah dan memerlukan distribusi bersyarat reguler; paradoks Borel menunjukkan jawabannya
dapat bergantung pada cara mengambil limit. Independensi juga tidak stabil: ia dapat
hilang atau muncul di bawah pengkondisian, dan pada model grafis muncul sebagai bias
seleksi.

Pertanyaan terbuka: seberapa jauh kelaziman harus diketahui sebelum aturan Bayes dipakai
secara jujur, dan apa yang dapat dikatakan bila kelaziman hanya diketahui dari sampel
kecil. Jawaban formalnya inferensi Bayesian dengan prior, yang membawa persoalan
pemilihan prior; jawaban frekuentisnya uji hipotesis di Bab 20, yang tidak memakai prior
tetapi menjawab pertanyaan berbeda. Bab ini tidak memilih salah satu.

\peta{Bab 17 pada Blitzstein dan Hwang, \emph{Introduction to Probability}: Bab 1
(Probability and counting, termasuk How to count dan contoh masalah ulang tahun) untuk
aksioma, pencacahan, dan peluang kombinatorial; Bab 2 (Conditional probability, Bayes'
rule and the law of total probability, Independence of events) untuk peluang bersyarat,
Bayes, dan independensi. Pada Wasserman, \emph{All of Statistics}: Bab 1 bagian
Probability dan Bab 2 bagian Independent Events. Himpunan dan komplemen dipakai dari
Bab 2 buku ini; deret yang mendasari aditif terbilang dari Bab 5; hukum bilangan besar
yang menopang simulasi dibuktikan di Bab 18. Variabel acak di Bab 17 memindahkan bobot
dari kejadian ke bilangan, sedangkan uji hipotesis di Bab 19 memakai aturan Bayes dan hukum
peluang total untuk pertanyaan berbeda.}
\chapter{Variabel Acak dan Distribusi}

Sebuah layanan pesan menerima rata-rata tiga permintaan per detik pada jam sibuk.
Angka itu tercetak di dasbor, dan ia tidak menjawab pertanyaan yang mendesak: berapa
sering satu detik berisi delapan permintaan atau lebih, sehingga antrean meluap? Dua sistem bisa punya rata-rata sama persis dan berperilaku
sangat berbeda: yang satu hampir tidak pernah meluap, yang lain beberapa kali sehari.
Rata-rata saja tidak cukup, dan bukan karena hitungannya kurang teliti.

Bab 17 memberi kita bahasa untuk kejadian: peluang satu hasil, peluang gabungan,
peluang bersyarat. Bahasa itu berhenti bekerja begitu yang kita ukur berupa bilangan.
``Jumlah permintaan dalam satu detik'' bukan satu kejadian, melainkan keluarga
kejadian yang tak berhingga banyaknya, satu untuk setiap nilai yang mungkin. Untuk
cacahan, peluang tiap nilai masih bisa dijumlahkan satu per satu. Untuk besaran yang
berlanjut, cara itu runtuh: waktu tunggu bisa bernilai $0{,}317$ detik atau
$0{,}3170001$ detik, dan tidak ada cara menyebut ``peluang tepat $0{,}317$ detik''
tanpa memberi bobot nol pada setiap nilai.

Jadi dibutuhkan dua hal: cara memindahkan peluang ke garis bilangan, dan cara
mengganti penjumlahan dengan integral ketika nilainya berlanjut.
Setelah keduanya ada, rata-rata, sebaran, dan bentuk ekor bisa dihitung.

\begin{untukanak}
Bayangkan sebuah toko roti. Setiap jam, penjual mencatat berapa banyak roti yang
terjual: kadang 3, kadang 11, kadang tidak ada sama sekali. Catatan itu berupa
bilangan bulat, dan tidak bisa bernilai setengah. Kalau kita ingin tahu berapa peluang
terjual tepat 7 roti, kita bisa menghitungnya: hitung berapa jam penjualan 7 roti,
bagi dengan jumlah jam yang dicatat. Semua peluang bisa dijumlahkan satu per satu,
karena pilihannya sedikit dan bisa didaftar.

Sekarang timbang satu roti di atas timbangan digital. Hasilnya bukan bilangan bulat:
250,3 gram, atau 249,8 gram. Mau sampai berapa angka di belakang koma pun, selalu ada
kemungkinan lain yang lebih dekat lagi. Karena itu pertanyaan ``berapa peluang beratnya
tepat 250,0000 gram'' tidak berguna, sedangkan ``berapa peluang beratnya antara 249
dan 251 gram'' sangat berguna. Peluang di sini bukan milik satu titik, melainkan milik
selang.

Di bab ini kita akan tahu mengapa kedua cara itu berbeda, dan mengapa keduanya tetap
bisa memberi rata-rata yang masuk akal.
\end{untukanak}

\section{Dari hasil ke bilangan}
\tingkat{dasar}

Alat ukur tidak mencatat kejadian; ia mencatat bilangan. Jembatannya adalah fungsi
yang mengambil setiap hasil percobaan dan menetapkan satu bilangan real. Nama
``variabel acak'' menyesatkan: ia bukan variabel, dan ia tidak acak. Keacakannya
berasal dari peluang pada ruang hasil, sedangkan fungsinya biasa saja.

\begin{definisi}[Variabel acak]\index{variabel acak}
Misalkan $\Omega$ ruang hasil dan $\mathcal{F}$ kumpulan kejadian. Pemetaan
$X:\Omega\to\mathbb{R}$ disebut \emph{variabel acak} bila untuk setiap $x\in\mathbb{R}$
himpunan
\[
\{X\le x\}=\{\omega\in\Omega : X(\omega)\le x\}
\]
merupakan kejadian, yaitu anggota $\mathcal{F}$.
\end{definisi}

Syarat terakhir itu menjamin $P(X\le x)$ punya arti. Untuk hampir semua fungsi yang
muncul dalam praktik ia terpenuhi otomatis, dan hanya gagal pada penanda himpunan yang
tidak terukur.

Sebagai contoh, lempar tiga koin setimbang dan misalkan $X$ jumlah gambar yang muncul.
Nilainya $0,1,2,3$ dengan peluang $1/8,3/8,3/8,1/8$. Perhatikan bahwa $X$ memetakan
delapan hasil ke empat bilangan, jadi beberapa hasil berbagi nilai yang sama. Itulah
sebabnya peluang sebuah nilai adalah peluang gabungan kejadian.

Kita perlu satu alat yang menjawab ``sejauh mana nilainya sudah bergerak'' untuk semua
ambang sekaligus, karena ambang dipilih pembaca, bukan penulis buku.

\begin{definisi}[Fungsi distribusi kumulatif]\index{fungsi distribusi kumulatif}
Fungsi distribusi kumulatif dari variabel acak $X$ adalah
$F_X(x)=P(X\le x)$ untuk $x\in\mathbb{R}$.
\end{definisi}

\begin{proposisi}[Sifat fungsi distribusi]
Untuk setiap variabel acak $X$ berlaku:
\begin{enumerate}
  \item $F_X$ tidak turun: bila $a\le b$ maka $F_X(a)\le F_X(b)$;
  \item $\lim_{x\to-\infty}F_X(x)=0$ dan $\lim_{x\to+\infty}F_X(x)=1$;
  \item $F_X$ kontinu kanan, yaitu $F_X(x)=F_X(x^+)$ untuk setiap $x$;
  \item $P(a<X\le b)=F_X(b)-F_X(a)$ untuk setiap $a<b$.
\end{enumerate}
\end{proposisi}

\begin{proof}
Untuk (1), kejadian $\{X\le a\}$ termuat dalam $\{X\le b\}$ bila $a\le b$, dan
peluang bersifat monoton terhadap penyertaan.

Untuk (2), misalkan $A_n=\{X\le -n\}$. Kejadian ini turun, dan
$\bigcap_n A_n=\varnothing$ karena setiap hasil memberi nilai berhingga. Kekontinuan
peluang dari atas, yang menyusul dari aditif terbilang, memberi $F_X(-n)\to0$. Arah
sebaliknya memakai $B_n=\{X\le n\}$ yang naik menuju $\Omega$, sehingga
$F_X(n)\to1$.

Untuk (3), ambil barisan turun $y_n\downarrow x$ dan $C_n=\{X\le y_n\}$. Maka
$\bigcap_n C_n=\{X\le x\}$, jadi $F_X(y_n)\to F_X(x)$ menurut kekontinuan dari atas.
Karena barisan itu sebarang, limit kanan sama dengan nilai di $x$.

Untuk (4), kejadian $\{X\le b\}$ terpecah menjadi $\{X\le a\}$ dan $\{a<X\le b\}$,
sehingga $F_X(b)=F_X(a)+P(a<X\le b)$.
\end{proof}

Untuk dua lemparan koin setimbang dengan $X$ jumlah gambar,
\[
F_X(x)=\begin{cases}
0, & x<0,\\
1/4, & 0\le x<1,\\
3/4, & 1\le x<2,\\
1, & x\ge 2.
\end{cases}
\]
Fungsi ini terdefinisi di seluruh garis bilangan, walaupun $X$ hanya bernilai $0,1,2$.
Karena itu $F_X(1{,}4)=3/4$: yang ditanyakan bukan ``apakah $X$ sama dengan $1{,}4$'',
melainkan ``apakah $X$ paling besar $1{,}4$''. Setiap lompatan grafiknya adalah massa
peluang pada satu nilai. Gambar~\ref{fig:tangga-cdf} menggambar tangga itu: tiga anak
tangga mendatar, tiga lompatan tegak, dan tinggi setiap lompatan sama dengan massa
peluang di titik lompatnya.

\begin{figure}[htbp]
\centering
\begin{tikzpicture}[font=\footnotesize,>=Stealth,x=1.75cm,y=4.2cm]
  \foreach \y/\lab in {0.25/{1/4},0.75/{3/4},1/{1}}{
    \draw[layergrey!35,dashed] (-1.35,\y) -- (3.55,\y);
    \node[left,font=\scriptsize] at (0,\y) {$\lab$};
  }
  \draw[->,layergrey] (-1.35,0) -- (3.6,0) node[below] {$x$};
  \draw[->,layergrey] (0,-0.1) -- (0,1.18) node[left] {$F_X(x)$};
  \draw[very thick,layerindigo] (-1.25,0) -- (0,0);
  \draw[very thick,layerindigo] (0,0.25) -- (1,0.25);
  \draw[very thick,layerindigo] (1,0.75) -- (2,0.75);
  \draw[very thick,layerindigo] (2,1) -- (3.5,1);
  \draw[very thick,layerplum] (0,0) -- (0,0.25);
  \draw[very thick,layerplum] (1,0.25) -- (1,0.75);
  \draw[very thick,layerplum] (2,0.75) -- (2,1);
  \fill[layerindigo] (0,0.25) circle (1.8pt);
  \fill[layerindigo] (1,0.75) circle (1.8pt);
  \fill[layerindigo] (2,1) circle (1.8pt);
  \draw[layerindigo,fill=white,thick] (1,0.25) circle (1.8pt);
  \draw[layerindigo,fill=white,thick] (2,0.75) circle (1.8pt);
  \node[below,font=\scriptsize] at (0,-0.02) {$0$};
  \node[below,font=\scriptsize] at (1,-0.02) {$1$};
  \node[below,font=\scriptsize] at (2,-0.02) {$2$};
  \node[right,font=\scriptsize,text=layerplum] at (1.06,0.5) {lompatan $1/2$};
  \node[right,font=\scriptsize,text=layerplum] at (0.06,0.12) {lompatan $1/4$};
\end{tikzpicture}
\caption{Fungsi distribusi kumulatif jumlah gambar pada dua lemparan koin. Grafiknya
tangga datar di antara nilai yang mungkin, lalu melompat tepat di nilai itu. Lompatan
$1/4$ di $x=0$, $1/2$ di $x=1$, dan $1/4$ di $x=2$ adalah massa peluang $p_X(0)$,
$p_X(1)$, dan $p_X(2)$.
Titik penuh di ujung kiri setiap anak tangga menandakan $F_X$ kontinu kanan, sehingga
$F_X(1)=3/4$ dan bukan $1/4$.}
\label{fig:tangga-cdf}
\end{figure}

\kamus{variabel acak}{random variable}{fungsi yang menempelkan satu bilangan pada tiap
hasil percobaan, seperti juru tulis yang mencatat satu angka untuk setiap kejadian}

\kamus{fungsi distribusi kumulatif}{cumulative distribution function}{tangga yang
menanjak: tingginya di titik $x$ menyatakan berapa banyak peluang yang sudah
terkumpul sampai ambang itu}

\section{Fungsi massa dan fungsi kerapatan}
\tingkat{dasar}

Ada dua cara memberi bobot pada nilai variabel acak. Pada variabel yang nilainya bisa
didaftar, bobot menempel pada titik; pada yang nilainya mengisi selang, bobot menempel
pada panjang. Kekeliruan paling umum di topik ini adalah memakai cara pertama untuk
membaca hasil cara kedua.

\begin{definisi}[Variabel diskret dan fungsi massa peluang]\index{fungsi massa peluang}
$X$ disebut \emph{diskret} bila himpunan nilainya terbilang, katakanlah
$\{x_1,x_2,\dots\}$. \emph{Fungsi massa peluang} dari $X$ adalah $p_X(x)=P(X=x)$.
Fungsi ini memenuhi $p_X(x)\ge0$, $p_X(x)\le1$, dan $\sum_i p_X(x_i)=1$, serta
$F_X(x)=\sum_{x_i\le x}p_X(x_i)$.
\end{definisi}

Untuk $X$ binomial dengan $n=2$ dan $p=1/2$, massa peluangnya $1/4,1/2,1/4$ pada
$x=0,1,2$ dan nol di tempat lain. Massa peluang tidak pernah lebih besar dari satu,
karena ia memang sebuah peluang.

\begin{definisi}[Variabel kontinu dan fungsi kerapatan peluang]\index{fungsi kerapatan}
$X$ disebut \emph{kontinu} bila terdapat fungsi $f_X\ge0$ sehingga
\[
P(a\le X\le b)=\int_a^b f_X(x)\,dx
\]
untuk setiap $a\le b$, dengan $\int_{-\infty}^{\infty}f_X(x)\,dx=1$. Fungsi $f_X$
disebut \emph{fungsi kerapatan peluang}.
\end{definisi}

Kerapatan bukan peluang, dan perbedaannya bukan soal istilah. Ambil $X$ dengan
kerapatan $f_X(x)=5$ pada $[0,1/5]$ dan nol di luar itu. Maka $f_X(0{,}1)=5$, angka
yang mustahil bagi sebuah peluang, sedangkan
$P(X\le0{,}1)=\int_0^{0{,}1}5\,dx=0{,}5$. Kerapatan juga boleh tak terbatas:
$f_X(x)=\tfrac23x^{-1/3}$ pada $(0,1)$ berintegral satu walaupun meledak di dekat nol.
Satuannya berbeda: peluang per satuan panjang, seperti kepadatan penduduk adalah jiwa
per kilometer persegi. Gambar~\ref{fig:batang-dan-kurva} menaruh
kedua cara itu berdampingan: batang di kiri dijumlahkan menurut tinggi, segitiga di
kanan menurut luas, dan keduanya berjumlah satu.

\begin{figure}[htbp]
\centering
\begin{tikzpicture}[font=\footnotesize,>=Stealth]
  \begin{scope}[x=1.05cm,y=5.4cm]
    \draw[->,layergrey] (0,0) -- (5.4,0) node[right] {$x$};
    \draw[->,layergrey] (0,0) -- (0,0.62) node[left] {$p_X(x)$};
    \foreach \x/\p in {1/0.1,2/0.2,3/0.3,4/0.4}{
      \fill[layerorange!40] (\x-0.3,0) rectangle (\x+0.3,\p);
      \draw[layerbrown,thick] (\x-0.3,0) rectangle (\x+0.3,\p);
      \node[below,font=\scriptsize] at (\x,0) {$\x$};
    }
    \node[font=\scriptsize,text=layerbrown,align=left,anchor=west] at (0.45,0.44)
      {empat batang,\\ jumlah tinggi $=1$};
  \end{scope}
  \begin{scope}[xshift=7.7cm,x=1.05cm,y=5.4cm]
    \draw[->,layergrey] (0,0) -- (5.4,0) node[right] {$x$};
    \draw[->,layergrey] (0,0) -- (0,0.62) node[left] {$f_X(x)$};
    \fill[layerteal!25] (0.4,0) -- (2.4,0.5) -- (4.4,0) -- cycle;
    \draw[layerteal,thick] (0.4,0) -- (2.4,0.5) -- (4.4,0) -- cycle;
    \draw[layergrey,dashed] (2.4,0) -- (2.4,0.5);
    \node[above,font=\scriptsize] at (2.4,0.5) {$0{,}5$};
    \node[below,font=\scriptsize] at (0.4,0) {$0$};
    \node[below,font=\scriptsize] at (4.4,0) {$4$};
    \node[font=\scriptsize,text=layerteal,align=center] at (2.4,0.14)
      {luas $=\tfrac12\cdot4\cdot0{,}5=1$};
  \end{scope}
\end{tikzpicture}
\caption{Dua cara memberi bobot yang sama-sama berjumlah satu. Kiri: empat batang
setinggi $0{,}1$, $0{,}2$, $0{,}3$, $0{,}4$ pada nilai $1,2,3,4$, dijumlahkan menurut
tinggi. Kanan: kerapatan berbentuk segitiga dengan puncak $0{,}5$ pada $[0,4]$,
dijumlahkan menurut luas, dan $\tfrac12\cdot4\cdot0{,}5=1$. Tinggi $0{,}5$ di panel
kanan bukan peluang apa pun; yang berpeluang adalah luas potongan di bawahnya.}
\label{fig:batang-dan-kurva}
\end{figure}

\begin{bacakeras}
Bacakan berpasangan, satu orang memilih suara Tinggi dan satu orang suara Berat. Mulai
pelan, lalu percepat.

Tinggi: Saya tinggi kurva. Saya menjulang di satu titik, dan semua orang menunjuk saya.
Berat: Saya berat di bawah kurva. Saya menumpuk sepanjang selang, dan semua orang
menimbang saya.
Tinggi: Angka saya boleh lima. Boleh sepuluh. Boleh tak terbatas.
Berat: Angka saya tidak pernah melebihi satu, karena saya adalah seluruh peluang.
Tinggi: Tanpa saya, kau tidak tahu di mana nilai paling sering muncul.
Berat: Tanpa saya, kau tidak tahu apakah puncak itu berisi apa pun.
Bersama: Tinggi menunjuk tempat. Berat menimbang isi.

Ulangi bersama tiga kali, makin pelan: peluang itu berat, bukan tinggi.
\end{bacakeras}

\begin{teorema}\index{kerapatan}
Bila $X$ kontinu dengan kerapatan $f_X$, maka $F_X(x)=\int_{-\infty}^{x}f_X(t)\,dt$,
dan $f_X(x)=F_X'(x)$ di setiap titik $x$ tempat $F_X$ terturunkan.
\end{teorema}

\begin{proof}
Persamaan pertama adalah definisi kerapatan dengan $a=-\infty$ dan $b=x$. Persamaan
kedua adalah teorema dasar kalkulus bagian pertama: turunan fungsi akumulasi sama
dengan $f_X$ di setiap titik kekontinuan $f_X$.
\end{proof}

Pertanyaan yang paling sering menimbulkan kebingungan: bila $X$ kontinu maka
$P(X=x)=0$ untuk setiap $x$, sebab $\int_x^x f_X=0$. Apakah itu berarti kejadian $X=x$
mustahil? Tidak. Himpunan $\{\omega:X(\omega)=x\}$ tetap ada dan tetap bisa berisi
hasil; yang nol adalah peluangnya, bukan keberadaannya.

Kesulitan muncul saat mencoba menjumlahkan semua titik. Selang $[0,1]$ adalah gabungan
tak terbilang dari titik berpeluang nol, tetapi $P(0\le X\le1)=1$. Tidak ada
kontradiksi: aksioma ketiga Bab 17 hanya menjamin aditif untuk gabungan
\emph{terbilang}, dan integral mengukur panjang, bukan menghitung titik.

Ada lapisan praktis. Pembangkit bawaan NumPy membentuk nilai
sebagai kelipatan $2^{-53}$, jadi hanya ada $2^{53}$ keluaran mungkin pada $[0,1)$, dan
peluang satu nilai tertentu muncul dalam $10^7$ tarikan kira-kira
$10^7/2^{53}=1{,}1\times10^{-9}$. Pada komputer ``peluang nol'' sebenarnya
$2^{-53}\approx1{,}1\times10^{-16}$; nol yang kita tulis adalah pengidealan.

\begin{contoh}
Misalkan $X$ seragam pada $(0,1)$, dengan $f_X(x)=1$ untuk $0<x<1$. Maka
\[
P(0{,}49<X<0{,}51)=\int_{0{,}49}^{0{,}51}1\,dx=0{,}02,
\qquad P(X=0{,}5)=0 .
\]
Kerapatan di sekitar $0{,}5$ sama besar dengan di sekitar $0{,}01$; yang membedakan
hanya panjang selangnya. Karena itu pertanyaan yang benar untuk variabel kontinu
selalu berbentuk selang.
\end{contoh}

\begin{tangan}
Bahan: meteran badan, gunting, kertas karton, dan sepuluh kartu kosong.

Ukur tinggi sepuluh teman sampai sentimeter terdekat, lalu tulis satu tinggi pada tiap
kartu. Buat lima kotak di lantai dari potongan karton selebar lima sentimeter:
$150$--$154$, $155$--$159$, sampai $170$--$174$. Letakkan setiap kartu di kotak yang
cocok, bertumpuk ke atas. Yang harus kita lihat: tumpukan itu punya bentuk, dan jumlah
kartu di semua kotak tetap sepuluh.

Potong ulang karton menjadi kotak selebar sepuluh sentimeter, susun kembali kartunya,
dan lihat apa yang berubah. Tumpukan menjadi hampir separuh tingginya, tetapi tidak ada
satu kartu pun yang hilang. Lebar kotak mengubah tinggi tumpukan; jumlah kartu tidak
berubah. Itulah bedanya tinggi dan berat, dan itulah sebabnya kerapatan harus dikalikan
lebar selang sebelum boleh disebut peluang.
\end{tangan}

\begin{untukanak}
Bayangkan lapangan sepak bola yang penuh orang. Kita ingin tahu di mana orangnya
paling banyak. Kita tidak bisa menjawabnya dengan bertanya ``berapa orang tepat di
titik ini''. Jawabannya selalu nol, karena satu titik tidak punya luas. Yang bisa
dijawab adalah ``berapa orang di setiap meter persegi''. Angka itu disebut kepadatan,
dan nilainya berbeda-beda: di depan pintu masuk padat, di pojok lapangan lengang.

Kepadatan tinggi bukan berarti di situ ada banyak orang di satu titik. Kepadatan
tinggi berarti di sekitar situ, dalam wilayah kecil, orangnya menumpuk. Kalau kita
mengalikan kepadatan dengan luas wilayah, barulah keluar jumlah orang. Hal yang sama
berlaku untuk peluang: kerapatan dikali panjang selang memberi peluang, dan kerapatan
sendiri bukan peluang.

Inilah sebabnya dua kurva yang bentuknya sangat berbeda bisa memberi total yang sama,
yaitu satu. Yang penting bukan tinggi kurva di satu titik, melainkan luas di bawahnya.
Di bab ini kita akan tahu mengapa pertanyaan tentang titik harus diganti pertanyaan
tentang selang.
\end{untukanak}

\section{Harapan dan variansi}
\tingkat{menengah}

Kalau seluruh distribusi harus diringkas menjadi satu bilangan, bilangan yang dipilih
adalah titik tumpunya: letak pusat massa bila peluang dianggap berat.

\begin{definisi}[Harapan]\index{harapan}
Harapan dari variabel acak $X$ adalah
\[
\mathbb{E}[X]=\sum_i x_i\,p_X(x_i)
\qquad\text{atau}\qquad
\mathbb{E}[X]=\int_{-\infty}^{\infty}x\,f_X(x)\,dx,
\]
bergantung pada apakah $X$ diskret atau kontinu. Harapan ada hanya bila
$\mathbb{E}[|X|]<\infty$; bila tidak, dikatakan harapan $X$ tidak ada.
\end{definisi}

Syarat nilai mutlak itu bukan kehati-hatian berlebihan. Untuk distribusi Cauchy dengan
kerapatan $f(x)=1/(\pi(1+x^2))$, $\int|x|f(x)dx$ berdivergen, sehingga $\mathbb{E}[X]$
tidak ada walaupun distribusinya simetris terhadap nol; rata-rata sampelnya tidak
memusat ke mana pun, seberapa banyak pun data ditambah. Bab 19 memakai contoh ini
sebagai batas hukum bilangan besar.

Sebagai hitungan pendek, untuk $X$ Bernoulli dengan $P(X=1)=p$,
$\mathbb{E}[X]=0\cdot(1-p)+1\cdot p=p$. Untuk $X$ seragam pada $(0,1)$,
$\mathbb{E}[X]=\int_0^1x\,dx=1/2$.

Seringkali yang dibutuhkan bukan harapan $X$, melainkan harapan fungsi $X$, misalnya
$\mathbb{E}[X^2]$. Jalan lurusnya adalah mencari kerapatan $X^2$ lebih dahulu, lalu
mengintegralkan. Ada jalan yang jauh lebih pendek.

\begin{teorema}[Harapan fungsi variabel acak]
Untuk fungsi $r$ dan variabel acak $X$,
\[
\mathbb{E}[r(X)]=\sum_i r(x_i)\,p_X(x_i)
\qquad\text{atau}\qquad
\mathbb{E}[r(X)]=\int_{-\infty}^{\infty}r(x)\,f_X(x)\,dx .
\]
\end{teorema}

\begin{proof}
Untuk kasus diskret, kelompokkan nilai menurut $y=r(x)$:
\[
\sum_i r(x_i)p_X(x_i)=\sum_y y\sum_{i:\,r(x_i)=y}p_X(x_i)=\sum_y y\,P(r(X)=y),
\]
dan bentuk terakhir adalah definisi $\mathbb{E}[r(X)]$ karena $r(X)$ diskret dengan
massa peluang $P(r(X)=y)$. Untuk kasus kontinu, buktinya adalah rumus perubahan
variabel pada integral di Bab 15; pengelompokan yang sama berlaku dengan integral
menggantikan jumlah.
\end{proof}

\begin{proposisi}[Linearitas harapan]
Untuk variabel acak $X,Y$ dan konstanta $a,b$ berlaku
$\mathbb{E}[aX+bY]=a\,\mathbb{E}[X]+b\,\mathbb{E}[Y]$.
\end{proposisi}

\begin{proof}
Untuk kasus diskret dengan massa bersama $p(x,y)$,
\[
\mathbb{E}[aX+bY]=\sum_{x,y}(ax+by)p(x,y)
=a\sum_x x\sum_y p(x,y)+b\sum_y y\sum_x p(x,y)=a\,\mathbb{E}[X]+b\,\mathbb{E}[Y].
\]
Kasus kontinu identik dengan integral menggantikan jumlah.
\end{proof}

Yang tidak dipakai di sini adalah kebebasan: linearitas berlaku untuk variabel apa pun,
termasuk yang bergantung kuat. Karena itu ia aman untuk jumlah penanda, tetapi tidak
boleh dipakai untuk perkalian: $\mathbb{E}[XY]=\mathbb{E}[X]\mathbb{E}[Y]$ menuntut
kebebasan.

Harapan memberi pusat, tetapi tidak memberi lebar; ukuran lebar yang paling dipakai
adalah rata-rata kuadrat jarak ke pusat.

\begin{definisi}[Variansi]\index{variansi}
Variansi dari $X$ dengan $\mu=\mathbb{E}[X]$ adalah
$\operatorname{Var}(X)=\mathbb{E}\bigl[(X-\mu)^2\bigr]$, dan simpangan bakunya
$\sigma_X=\sqrt{\operatorname{Var}(X)}$.
\end{definisi}

\begin{proposisi}
Bila variansi ada, maka $\operatorname{Var}(X)=\mathbb{E}[X^2]-\mu^2$ dan
$\operatorname{Var}(aX+b)=a^2\operatorname{Var}(X)$.
\end{proposisi}

\begin{proof}
Kembangkan kuadratnya dan pakai linearitas:
\[
\mathbb{E}[(X-\mu)^2]=\mathbb{E}[X^2]-2\mu\,\mathbb{E}[X]+\mu^2
=\mathbb{E}[X^2]-\mu^2 .
\]
Untuk sifat kedua, $\mathbb{E}[aX+b]=a\mu+b$, sehingga
$\operatorname{Var}(aX+b)=\mathbb{E}[(aX+b-a\mu-b)^2]=a^2\mathbb{E}[(X-\mu)^2]$.
\end{proof}

Untuk $X$ Bernoulli, $\mathbb{E}[X^2]=p$ karena $X^2=X$ untuk nilai $0$ dan $1$,
sehingga $\operatorname{Var}(X)=p-p^2=p(1-p)$. Untuk $X$ seragam pada $(0,1)$,
$\mathbb{E}[X^2]=1/3$, sehingga $\operatorname{Var}(X)=1/3-1/4=1/12\approx0{,}083333$.
Variansi bersatuan kuadrat, sedangkan simpangan baku bersatuan sama dengan $X$. Satu
peringatan yang mahal: $\mathbb{E}[X^2]\ne\mathbb{E}[X]^2$, dan selisihnya justru
variansi; untuk seragam di atas, $0{,}25$ berbanding $0{,}333333$.

\begin{resep}
Bahan: tabel nilai $x_i$ dan peluang $p_X(x_i)$.
\begin{enumerate}
  \item Jumlahkan seluruh $p_X(x_i)$. Kalau hasilnya bukan satu, tabelnya belum sah.
  \item Kalikan setiap $x_i$ dengan $p_X(x_i)$, lalu jumlahkan: itulah $\mathbb{E}[X]$.
  \item Kuadratkan setiap $x_i$ lebih dahulu, kalikan dengan $p_X(x_i)$, lalu jumlahkan:
        itulah $\mathbb{E}[X^2]$.
  \item Kurangkan kuadrat hasil langkah 2 dari hasil langkah 3, yaitu
        $\operatorname{Var}(X)=\mathbb{E}[X^2]-(\mathbb{E}[X])^2$; akar kuadratnya
        adalah simpangan baku.
\end{enumerate}
Contoh: untuk $x=0,1,2,3$ dengan $p=0{,}1;\,0{,}2;\,0{,}3;\,0{,}4$, langkah 2 memberi
$2$, langkah 3 memberi $5$, dan langkah 4 memberi $1$.

Peringatan urutan: kuadratkan $x_i$ sebelum mengalikan dengan peluang, dan kurangkan
$\mu^2$ sesudah $\mathbb{E}[X^2]$ selesai dijumlahkan. Menukar urutan itu menghasilkan
angka yang tidak mengukur lebar apa pun.
\end{resep}

\kamus{harapan}{expectation}{titik tumpu penggaris berbeban: satu bilangan yang
menandai di mana nilai-nilai menumpuk}

\kamus{variansi}{variance}{ukuran seberapa jauh nilai biasa menyimpang dari titik
tumpunya, dihitung dari kuadrat jarak agar tanda plus dan minus tidak saling menghapus}

\section{Empat distribusi diskret}
\tingkat{menengah}

Empat keluarga distribusi diskret cukup untuk sebagian besar pekerjaan praktis: satu
percobaan biner, jumlah percobaan biner, waktu tunggu sampai sukses pertama, dan
kejadian langka pada selang waktu.

\paragraph{Bernoulli.} $X$ Bernoulli dengan parameter $p\in[0,1]$ bernilai $1$ dengan
peluang $p$ dan $0$ dengan peluang $1-p$, sehingga $p_X(x)=p^x(1-p)^{1-x}$ untuk
$x\in\{0,1\}$. Dari dua titik itu, $\mathbb{E}[X]=p$ dan
$\operatorname{Var}(X)=p(1-p)$; variansi terbesar pada $p=1/2$, yaitu $1/4$, tepat
ketika kedua kemungkinan seimbang.

\paragraph{Binomial.} Bila $n$ percobaan Bernoulli bebas dengan peluang sukses $p$
dijalankan, jumlah keberhasilan $X$ binomial dengan
\[
p_X(k)=\binom{n}{k}p^k(1-p)^{n-k},\qquad k=0,1,\dots,n .
\]
Rata-rata dan variansinya paling mudah diturunkan dengan memecah $X$ menjadi penanda,
bukan dari rumus itu.

\begin{teorema}\index{distribusi binomial}
Bila $X$ binomial dengan parameter $n$ dan $p$, maka
$\mathbb{E}[X]=np$ dan $\operatorname{Var}(X)=np(1-p)$.
\end{teorema}

\begin{proof}
Tulis $X=\sum_{i=1}^{n}X_i$ dengan $X_i=1$ bila percobaan ke-$i$ berhasil dan $0$ bila
gagal. Setiap $X_i$ Bernoulli dengan parameter $p$, jadi $\mathbb{E}[X_i]=p$ dan
$\operatorname{Var}(X_i)=p(1-p)$. Linearitas memberi
$\mathbb{E}[X]=\sum_i\mathbb{E}[X_i]=np$ tanpa memakai kebebasan sama sekali.

Untuk variansi, kebebasan baru diperlukan. Karena percobaannya bebas, pasangan
$X_i,X_j$ dengan $i\ne j$ juga bebas, sehingga
$\operatorname{Cov}(X_i,X_j)=p^2-p^2=0$. Maka
\[
\operatorname{Var}(X)=\sum_{i=1}^{n}\operatorname{Var}(X_i)
+2\sum_{i<j}\operatorname{Cov}(X_i,X_j)=np(1-p).
\]
\end{proof}

Rumus variansi lebih rapuh daripada rumus rata-rata: yang pertama menuntut kebebasan,
yang kedua tidak. Untuk $n=20$ dan $p=0{,}35$, $\mathbb{E}[X]=7$ dan
$\operatorname{Var}(X)=4{,}55=91/20$, dengan simpangan baku $2{,}133073$; massa di
sekitar pusatnya $p_X(7)=0{,}184401$, $p_X(10)=0{,}068614$, $p_X(12)=0{,}013564$.
Ekornya kecil tetapi bukan nol: $P(X\ge12)=0{,}019579$ dan $P(X\le4)=0{,}118197$, jadi
aturan ``tolak kotak bila ada dua belas cacat atau lebih'' menolak sekitar dua persen
kotak.

\paragraph{Geometrik.} Bila percobaan Bernoulli bebas diulang sampai sukses pertama,
banyaknya percobaan yang diperlukan adalah $X$ geometrik dengan
$p_X(k)=p(1-p)^{k-1}$ untuk $k=1,2,\dots$ Jumlah massanya satu karena deret geometri:
$\sum_{k\ge1}p(1-p)^{k-1}=p/(1-(1-p))=1$. Dengan $q=1-p$,
\[
\sum_{k\ge1}kq^{k-1}=\frac{d}{dq}\frac{1}{1-q}=\frac{1}{(1-q)^2}=\frac{1}{p^2},
\]
sehingga $\mathbb{E}[X]=p\cdot p^{-2}=1/p$ dan, dengan cara serupa,
$\operatorname{Var}(X)=(1-p)/p^2$ dari $\mathbb{E}[X^2]=(2-p)/p^2$. Untuk $p=0{,}2$:
$\mathbb{E}[X]=5$, $\operatorname{Var}(X)=20$, $P(X=5)=0{,}081920$ dan
$P(X>10)=0{,}10737418$; simulasi $10^5$ tarikan memberi rata-rata $4{,}99176$ dengan
galat baku $0{,}01414$.

\paragraph{Poisson.} Untuk $X$ Poisson dengan parameter $\lambda>0$,
\[
p_X(k)=e^{-\lambda}\frac{\lambda^k}{k!},\qquad k=0,1,2,\dots
\]
Jumlah massanya satu karena deret eksponensial:
$e^{-\lambda}\sum_k\lambda^k/k!=1$. Rata-rata dan variansinya sama, dan keduanya
keluar dari satu trik. Karena $k/k!=1/(k-1)!$,
\[
\mathbb{E}[X]=\sum_{k\ge1}k\,e^{-\lambda}\frac{\lambda^k}{k!}
=\lambda e^{-\lambda}\sum_{k\ge1}\frac{\lambda^{k-1}}{(k-1)!}=\lambda .
\]
Dengan menulis $k(k-1)/k!=1/(k-2)!$ diperoleh $\mathbb{E}[X(X-1)]=\lambda^2$, sehingga
$\operatorname{Var}(X)=\lambda$. Satu parameter mengatur pusat sekaligus lebar, dan itu
membatasi Poisson: data yang variansinya jauh melebihi rata-ratanya (disebut
\emph{overdispersi}) tidak cocok dimodelkan Poisson.

Untuk $\lambda=3$: $p_X(0)=0{,}049787$, $p_X(3)=0{,}224042$,
$P(X\ge8)=0{,}011905$. Angka terakhir menjawab pertanyaan antrean di pembuka bab:
dengan tiga permintaan per detik, sekitar $1{,}2$ persen detik berisi delapan
permintaan atau lebih.

Poisson adalah limit binomial ketika $n$ besar, $p$ kecil, dan $np\to\lambda$. Untuk
$n=1000$ dan $p=0{,}003$, binomial memberi $P(0)=0{,}049563083$ sedangkan Poisson
$e^{-3}=0{,}049787068$, selisih relatif $-4{,}5\times10^{-3}$. Kesalahannya kecil tetapi
tidak nol dan tandanya berganti menurut $k$; Poisson dipakai sebagai hampiran, bukan
kesamaan.

\section{Tiga distribusi kontinu}
\tingkat{menengah}

\paragraph{Seragam.} $X$ seragam pada $(a,b)$ bila kerapatannya $1/(b-a)$ di dalam
selang itu dan nol di luar. Fungsi distribusinya naik linear dari $0$ ke $1$, dan
\[
\mathbb{E}[X]=\frac{a+b}{2},\qquad
\operatorname{Var}(X)=\frac{(b-a)^2}{12}.
\]
Faktor $12$ muncul dari $\int_a^b(x-\tfrac{a+b}{2})^2\frac{dx}{b-a}=(b-a)^2/12$; untuk
$(0,1)$ hasilnya $1/12\approx0{,}083333$, sama dengan hitungan lewat
$\mathbb{E}[X^2]$ di bagian sebelumnya.

\paragraph{Eksponensial.} $X$ eksponensial dengan parameter $\beta>0$ memiliki
kerapatan $f_X(x)=\beta^{-1}e^{-x/\beta}$ untuk $x>0$, dengan
$F_X(x)=1-e^{-x/\beta}$, $\mathbb{E}[X]=\beta$, dan $\operatorname{Var}(X)=\beta^2$.
Parameter $\beta$ adalah rata-rata waktu tunggu. Sifat yang membedakannya dari semua
distribusi kontinu lain adalah ketiadaan memori.

\begin{proposisi}[Tanpa memori]\index{distribusi eksponensial}
Untuk $s,t\ge0$ berlaku $P(X>s+t\mid X>s)=P(X>t)$.
\end{proposisi}

\begin{proof}
$P(X>u)=e^{-u/\beta}$, sehingga
\[
P(X>s+t\mid X>s)=\frac{P(X>s+t)}{P(X>s)}
=\frac{e^{-(s+t)/\beta}}{e^{-s/\beta}}=e^{-t/\beta}=P(X>t).
\]
\end{proof}

Komponen yang sudah bertahan lima tahun tidak lebih mungkin bertahan satu tahun lagi
dibanding komponen baru: masuk akal untuk proses tanpa penuaan, tidak masuk akal untuk
lampu yang aus. Untuk $\beta=2$, $P(X>3)=0{,}22313016$ dan simulasi $10^5$ tarikan
memberi $0{,}22237$;
dengan syarat $X>2$, peluang empiris $X>5$ adalah $0{,}22028$ berbanding teori
$0{,}22313$. Eksponensial adalah kembaran kontinu geometrik: bila kedatangan mengikuti
Poisson dengan laju $\lambda$, jarak antar kedatangan eksponensial dengan
$\beta=1/\lambda$.

\paragraph{Normal.} $X$ normal dengan parameter $\mu$ dan $\sigma^2>0$ memiliki
kerapatan
\[
f_X(x)=\frac{1}{\sigma\sqrt{2\pi}}
\exp\!\left(-\frac{(x-\mu)^2}{2\sigma^2}\right),\qquad x\in\mathbb{R}.
\]
Konstanta di depan bukan pilihan: ia satu-satunya nilai yang membuat luas di bawah
kurva sama dengan satu. Bahwa luasnya satu dibuktikan dengan koordinat polar di Bab 15:
tulis $I^2$ sebagai integral ganda $e^{-(x^2+y^2)/2}$, ubah ke polar, dan dapatkan
$I^2=2\pi$.

\begin{proposisi}
Bila $X\sim N(\mu,\sigma^2)$, maka $Z=(X-\mu)/\sigma\sim N(0,1)$.
\end{proposisi}

\begin{proof}
$F_Z(z)=P(X\le\mu+\sigma z)=F_X(\mu+\sigma z)$. Turunkan terhadap $z$ dengan aturan
rantai: $f_Z(z)=\sigma f_X(\mu+\sigma z)=\varphi(z)$, yaitu kerapatan normal baku.
\end{proof}

Fungsi $\Phi(z)=P(Z\le z)$ tidak punya bentuk elementer; ia dihitung numerik atau
dibaca dari tabel. Tiga angka yang layak dihafal: $P(|Z|\le1)=0{,}682689492137$,
$P(|Z|\le2)=0{,}954499736104$, $P(|Z|\le3)=0{,}997300203937$. Untuk peluang tak
simetris, kurangkan dua nilai $\Phi$: $P(1<Z<2)=0{,}135905122$. Kuantilnya,
$\Phi^{-1}(0{,}975)=1{,}959964$, muncul di setiap interval kepercayaan 95 persen pada
Bab 20.

Normal bukan distribusi yang lebih benar daripada yang lain. Ia penting karena dua
alasan berbeda: banyak besaran alam mendekatinya, dan jumlah banyak variabel bebas
cenderung mendekatinya apa pun bentuk aslinya. Alasan kedua adalah teorema limit
pusat, dibuktikan pada Bab 19.

\kamus{distribusi normal}{normal distribution}{kurva lonceng yang simetris: hampir
semua nilai menumpuk dekat pusat, dan ekornya menipis sangat cepat}

\kamus{distribusi eksponensial}{exponential distribution}{waktu tunggu pada proses
tanpa penuaan: yang sudah lama menunggu tidak lebih dekat ke kejadian berikutnya}

\section{Distribusi bersama, kovarians, dan korelasi}
\tingkat{lanjut}

Dua distribusi terpisah tidak menyimpan informasi tentang apakah kedua besaran
cenderung naik bersama. Alat untuk itu adalah distribusi bersama.

\begin{contoh}
Misalkan $X,Y\in\{0,1\}$ dengan massa bersama $P(0,0)=0{,}1$, $P(0,1)=0{,}2$,
$P(1,0)=0{,}3$, $P(1,1)=0{,}4$. Menjumlahkan baris dan kolom memberi marginal
$P(X=1)=0{,}7$ dan $P(Y=1)=0{,}6$. Bila $X$ dan $Y$ bebas, massa pada $(1,1)$
seharusnya $0{,}42$, sedangkan yang sebenarnya $0{,}4$; jadi keduanya tidak bebas,
walaupun selisihnya kecil.
\end{contoh}

Pada kasus kontinu, distribusi bersama digambarkan oleh kerapatan bersama
$f_{X,Y}\ge0$ dengan $\iint f_{X,Y}=1$, dan marginalnya diperoleh dengan
mengintegralkan satu peubah keluar: $f_X(x)=\int f_{X,Y}(x,y)\,dy$. Kebebasan berarti
$f_{X,Y}(x,y)=f_X(x)f_Y(y)$ untuk semua $x,y$.

\begin{definisi}[Kovarians dan korelasi]\index{kovarians}\index{korelasi}
Dengan $\mu_X=\mathbb{E}[X]$ dan $\mu_Y=\mathbb{E}[Y]$,
\[
\operatorname{Cov}(X,Y)=\mathbb{E}\bigl[(X-\mu_X)(Y-\mu_Y)\bigr],
\qquad
\rho(X,Y)=\frac{\operatorname{Cov}(X,Y)}{\sigma_X\sigma_Y}.
\]
\end{definisi}

Kovarians positif berarti kedua variabel cenderung menyimpang ke arah yang sama;
negatif berarti berlawanan. Korelasi adalah versi tak bersatuan dari kovarians.

\begin{proposisi}
$\operatorname{Cov}(X,Y)=\mathbb{E}[XY]-\mathbb{E}[X]\mathbb{E}[Y]$ dan
$-1\le\rho(X,Y)\le1$.
\end{proposisi}

\begin{proof}
Kembangkan perkaliannya dan pakai linearitas. Untuk batas korelasi, pandang hasil kali
dalam $\langle U,V\rangle=\mathbb{E}[UV]$ pada peubah terpusat; ketaksamaan
Cauchy-Schwarz Bab 10 memberi
$|\langle U,V\rangle|^2\le\langle U,U\rangle\langle V,V\rangle$, yaitu
$\operatorname{Cov}(X,Y)^2\le\sigma_X^2\sigma_Y^2$.
\end{proof}

\begin{teorema}
$\operatorname{Var}(X+Y)=\operatorname{Var}(X)+\operatorname{Var}(Y)
+2\operatorname{Cov}(X,Y)$.
\end{teorema}

\begin{proof}
Tulis $X+Y-\mu_X-\mu_Y=(X-\mu_X)+(Y-\mu_Y)$, kuadratkan, lalu pakai linearitas pada
tiga suku yang muncul; suku silangnya tepat dua kali kovarians.
\end{proof}

Kembali ke tabel di atas: $\mathbb{E}[X]=7/10$, $\mathbb{E}[Y]=3/5$,
$\mathbb{E}[XY]=2/5$, sehingga $\operatorname{Cov}(X,Y)=-1/50=-0{,}02$. Dengan
$\operatorname{Var}(X)=21/100$ dan $\operatorname{Var}(Y)=6/25$, korelasinya
$-0{,}089087$ dan $\operatorname{Var}(X+Y)=0{,}41$.

Kovarians nol sering disalahartikan sebagai kebebasan, padahal arahnya hanya satu:
bebas berakibat kovarians nol, sebaliknya tidak. Ambil $X$ bernilai $-1,0,1$ dengan
peluang sama dan $Y=X^2$; maka $\mathbb{E}[X]=0$ dan
$\mathbb{E}[XY]=\mathbb{E}[X^3]=0$, sehingga korelasinya nol, padahal $Y$ ditentukan
sepenuhnya oleh $X$. Kovarians hanya menangkap ketergantungan yang berbentuk garis
lurus.

\kamus{kovarians}{covariance}{ukuran apakah dua besaran cenderung naik dan turun
bersama, dihitung dari hasil kali simpangan keduanya}

\section{Transformasi variabel acak}
\tingkat{lanjut}

Pekerjaan ini muncul terus: harga diubah ke logaritma, panjang ke luas, waktu tunggu
ke laju. Pertanyaannya selalu sama, yaitu bagaimana distribusi berubah ketika
$Y=r(X)$.

Pada kasus diskret jawabannya sederhana: massa dipindahkan dari nilai lama ke nilai
baru. Misalkan $P(X=-1)=1/4$, $P(X=0)=1/2$, $P(X=1)=1/4$, dan $Y=X^2$. Maka
$P(Y=0)=1/2$ dan $P(Y=1)=P(X=-1)+P(X=1)=1/2$; nilai $Y$ hanya dua, lebih sedikit
daripada $X$, karena $r(x)=x^2$ tidak satu-satu.

Pada kasus kontinu massanya nol di setiap titik, jadi yang dipindahkan adalah panjang,
dan panjang berubah skala. Tiga langkah baku menjawab semuanya: tentukan
$A_y=\{x:r(x)\le y\}$, hitung $F_Y(y)=\int_{A_y}f_X(x)\,dx$, lalu turunkan
$f_Y=F_Y'$.

\begin{teorema}[Transformasi monoton]\index{transformasi variabel acak}
Bila $r$ monoton tegas dan terturunkan, dengan invers $s=r^{-1}$, maka
\[
f_Y(y)=f_X\bigl(s(y)\bigr)\left|\frac{ds(y)}{dy}\right| .
\]
\end{teorema}

\begin{proof}
Misalkan $r$ naik tegas. Maka $r(x)\le y$ setara dengan $x\le s(y)$, sehingga
$F_Y(y)=F_X(s(y))$, dan aturan rantai memberi $f_Y(y)=f_X(s(y))s'(y)$ dengan
$s'(y)>0$. Bila $r$ turun tegas, ketaksamaan berbalik: $F_Y(y)=1-F_X(s(y))$, sehingga
$f_Y(y)=-f_X(s(y))s'(y)=f_X(s(y))|s'(y)|$.
\end{proof}

Faktor $|ds/dy|$ adalah peregangan lokal: bila $r$ merenggangkan skala, kerapatan
menurun, karena peluang yang sama tersebar pada selang lebih panjang. Di dimensi lebih
dari satu faktor itu menjadi determinan Jacobian, bentuk umumnya di Bab 15.
Gambar~\ref{fig:transformasi-sumbu} memperlihatkan peregangan itu pada $Y=X^2$: empat
selang sama panjang mendarat pada selang yang makin panjang ketika $y$ mendekati satu,
sementara kerapatan $Y$ meninggi di dekat nol, tempat selangnya memendek.

\begin{figure}[htbp]
\centering
\begin{tikzpicture}[font=\footnotesize,>=Stealth]
  \draw[->,layergrey] (-0.25,0) -- (7.3,0) node[right] {$x$};
  \foreach \px/\lab in {0/0,1.75/0{,}25,3.5/0{,}5,5.25/0{,}75,7/1}{
    \draw[layergrey] (\px,0.07) -- (\px,-0.07);
    \node[above,font=\scriptsize] at (\px,0.07) {$\lab$};
  }
  \node[font=\scriptsize,text=layergrey,anchor=west] at (0.15,0.55)
    {selang sama panjang};
  \draw[->,layergrey] (-0.25,-2.6) -- (7.3,-2.6) node[right] {$y=x^2$};
  \foreach \px/\lab/\anc in {0/0/east,0.44/0{,}0625/east,1.75/0{,}25/north,
                             3.94/0{,}5625/west,7/1/west}{
    \draw[layergrey] (\px,-2.53) -- (\px,-2.67);
    \node[below,font=\scriptsize,anchor=\anc] at (\px,-2.67) {$\lab$};
  }
  \foreach \px/\py in {0/0,1.75/0.44,3.5/1.75,5.25/3.94,7/7}{
    \draw[->,layerindigo] (\px,-0.05) -- (\py,-2.55);
  }
  \node[font=\scriptsize,text=layerindigo,align=center,anchor=west] at (1.0,-1.85)
    {selang memendek,\\ kerapatan meninggi};
  \node[font=\scriptsize,text=layerindigo,align=center,anchor=east] at (6.6,-1.85)
    {selang memanjang,\\ kerapatan menurun};
  \begin{scope}[xshift=8.7cm,yshift=-2.6cm]
    \draw[->,layergrey] (0,0) -- (5.4,0) node[right] {$y$};
    \draw[->,layergrey] (0,0) -- (0,3.95) node[left] {$f_Y(y)$};
    \draw[layerplum,very thick,domain=0.12:1,samples=100,smooth,variable=\t]
      plot ({5.0*\t},{1.175/sqrt(\t)});
    \draw[layerplum,dashed] (0,0) -- (0,3.8);
    \node[below,font=\scriptsize,anchor=east] at (0,0) {$0$};
    \node[below,font=\scriptsize] at (1.25,0) {$0{,}25$};
    \node[below,font=\scriptsize,anchor=west] at (5.0,0) {$1$};
    \draw[layergrey,dotted] (5.0,0) -- (5.0,1.175);
    \node[font=\scriptsize,text=layerplum,anchor=west] at (0.25,3.45) {meledak di dekat $0$};
  \end{scope}
\end{tikzpicture}
\caption{Transformasi $Y=X^2$ untuk $X$ seragam pada $(0,1)$, digambar pada dua sumbu.
Titik $x=0{,}25$ mendarat di $y=0{,}0625$, $x=0{,}5$ di $y=0{,}25$, dan $x=0{,}75$ di
$y=0{,}5625$; selang $0{,}25$ di dekat nol memendek menjadi $0{,}0625$, sedangkan selang
yang sama di dekat satu memanjang menjadi $0{,}4375$. Panel kanan adalah kerapatan
$f_Y(y)=1/(2\sqrt{y})$: nilainya $2$ di $y=0{,}0625$, $1$ di $y=0{,}25$, dan $0{,}5$ di
$y=1$. Kerapatan meninggi tepat di tempat selang memendek, sehingga luas totalnya tetap
satu.}
\label{fig:transformasi-sumbu}
\end{figure}

Bila $r$ tidak monoton, teorema di atas tidak berlaku langsung, tetapi tiga langkah
baku tetap bekerja dengan $A_y$ berupa gabungan selang.

\begin{contoh}
Misalkan $X$ seragam pada $(-1,3)$, jadi $f_X(x)=1/4$ di dalamnya, dan $Y=X^2$
bernilai di $(0,9)$. Untuk $0<y<1$ himpunan $A_y=[-\sqrt y,\sqrt y]$ seluruhnya di
dalam $(-1,3)$, jadi $F_Y(y)=\tfrac14\cdot2\sqrt y=\sqrt y/2$; untuk $1\le y<9$
himpunan itu menjadi $[-1,\sqrt y]$, jadi $F_Y(y)=(\sqrt y+1)/4$. Turunannya
\[
f_Y(y)=\begin{cases}
\dfrac{1}{4\sqrt y}, & 0<y<1,\\[4pt]
\dfrac{1}{8\sqrt y}, & 1<y<9,\\[4pt]
0, & \text{di luar itu}.
\end{cases}
\]
Kedua cabang berjumlah satu, $\tfrac12$ dan $\tfrac12$. Kerapatan $Y$ meledak di
dekat nol: $x^2$ memampatkan selang di sana.
\end{contoh}

Bila $X$ seragam pada $(0,1)$ dan $Y=X^2$, hasilnya lebih rapi: $F_Y(y)=\sqrt y$ dan
$f_Y(y)=1/(2\sqrt y)$, yang diperiksa numerik di Lab. Bila $Y=-\ln X$, maka
$F_Y(y)=1-e^{-y}$, eksponensial dengan $\beta=1$: membuang logaritma dari bilangan
seragam menghasilkan waktu tunggu eksponensial, dan kebalikannya dipakai komputer.
Bila $U$ seragam dan $F$ fungsi distribusi yang diinginkan, maka $F^{-1}(U)$
berdistribusi $F$.

\begin{tantangan}
Sebuah catatan kuliah mengusulkan kerapatan waktu tunggu
\[
f(x)=\frac{1}{2\sqrt{x}},\qquad 0<x<4,
\]
dengan alasan ``positif, turun, dan bentuknya mirip contoh di depan''. Dari kerapatan
itu penulisnya menghitung $P(X\le4)=1$ lalu menyimpulkan bahwa distribusinya sudah sah.

Periksa klaim itu. Tentukan luas total di bawah $f$, perbaiki kerapatannya dengan dua
cara sah yang berbeda, dan jelaskan mengapa perbaikannya tidak tunggal.

Petunjuk 1: yang menentukan bukan nilai $f$ di satu titik, melainkan seluruh daerah di
bawah kurvanya.

Petunjuk 2: hitung $\int_0^4\tfrac12 x^{-1/2}\,dx$; hasilnya $2$, bukan $1$. Supaya
totalnya satu, konstanta di depan harus dibagi dua menjadi $1/4$, atau selangnya
dipersempit menjadi $(0,1)$ dengan konstanta tetap.

Kalau mentok, kembali ke bagian ``Fungsi massa dan fungsi kerapatan''.
\end{tantangan}

\begin{untukanak}
Bayangkan sebuah lampu di lantai dan sebuah bola di antara lampu dan dinding.
Bayangan bola di dinding jauh lebih besar daripada bolanya, dan makin jauh bola dari
lampu, makin besar bayangannya. Tidak ada yang bertambah banyak: jumlah cat pada bola
tetap sama, hanya tersebar pada luas yang lebih besar.

Sekarang geser bola sedikit ketika ia sudah jauh dari lampu. Bayangannya bergeser
jauh. Ketika bola masih dekat lampu, geseran yang sama hanya menggeser bayangan
sedikit. Artinya, satu sentimeter di dekat lampu menjadi banyak sentimeter di dinding,
sedangkan satu sentimeter jauh dari lampu menjadi sangat banyak sentimeter.

Mengubah variabel acak bekerja seperti itu. Ada daerah yang dipampatkan dan ada
daerah yang direnggangkan, dan kerapatan hasil harus menyesuaikan diri agar totalnya
tetap satu. Kalau tidak disesuaikan, peluang akan muncul dari udara atau hilang entah
ke mana.
\end{untukanak}

\section{Contoh hitung}
\tingkat{menengah}

Sebuah lini produksi mengisi kotak berisi 20 sekering. Setiap sekering cacat dengan
peluang $0{,}35$, saling bebas, dan kotak ditolak bila berisi dua belas cacat atau
lebih.

Dengan $n=20$ dan $p=0{,}35$, aturan penolakan memberi $P(X\ge12)=0{,}019579$ dari
ekor binomial, bukan dari hampiran normal. Nilai paling mungkin, $X=7$, pun hanya
berpeluang $18{,}4$ persen, sehingga satu kotak tidak bercerita banyak.

Simulasi $10^5$ kotak dengan \texttt{default\_rng(20261024)} memberi rata-rata
$7{,}005260$ dan variansi $4{,}552918$. Galat baku rata-ratanya $0{,}006745$, jadi
simpangannya $+0{,}78$ galat baku; galat baku variansi sekitar $0{,}020348$, jadi
simpangannya $+0{,}14$. Keduanya dekat, dan itu memang diharapkan: simulasi tidak
membuktikan rumus, ia hanya gagal menyangkalnya.

Pemeriksaan kedua menyangkut total massa peluang Poisson. Jumlah
$e^{-\lambda}\lambda^k/k!$ sampai $K$ yang cukup besar harus sama dengan satu. Dihitung
dalam ruang log agar tidak meluap:

\begin{center}\small
\begin{tabular}{rrrr}
\hline
$\lambda$ & $K$ & jumlah & $1-\text{jumlah}$ \\
\hline
$0{,}5$ & 68 & $1{,}00000000000000000$ & $0$ \\
$1$ & 73 & $1{,}00000000000000000$ & $0$ \\
$3$ & 83 & $1{,}00000000000000022$ & $-2{,}2\times10^{-16}$ \\
$10$ & 107 & $1{,}00000000000000311$ & $-3{,}1\times10^{-15}$ \\
$50$ & 194 & $0{,}99999999999999778$ & $+2{,}2\times10^{-15}$ \\
$100$ & 280 & $1{,}00000000000003886$ & $-3{,}9\times10^{-14}$ \\
\hline
\end{tabular}
\end{center}

Penyimpangan terbesar muncul pada $\lambda=100$, yaitu $3{,}9\times10^{-14}$, sekitar
175 kali epsilon mesin $2{,}2\times10^{-16}$. Itu bukan tanda rumusnya salah, melainkan
tanda penjumlahan $281$ suku berukuran sangat berbeda kehilangan beberapa digit
terakhir; empat belas angka pertama masih benar.

Pemeriksaan ketiga menyangkut total kerapatan normal, karena konstanta
$1/\sqrt{2\pi}$ mudah salah ketik. Kuadratur adaptif pada $[-8,8]$ memberi
$0{,}99999999999999878$, sehingga kekurangannya $1{,}22\times10^{-15}$. Teori
meramalkan kekurangan itu tepat sebesar ekor di luar selang, yaitu
$2\Phi(-8)=1{,}24\times10^{-15}$. Kedua angka sepakat sampai dua angka penting: yang
hilang memang ekor, bukan kesalahan konstanta. Pada seluruh garis bilangan hasilnya
$0{,}99999999999999978$ dengan perkiraan galat tercetak $1{,}0\times10^{-8}$, taksiran
yang konservatif.

Pemeriksaan keempat, transformasi: kerapatan $1/(2\sqrt y)$ berintegral
$0{,}9999999999999992$, dengan sisa $8\times10^{-16}$ dari ujung $0$ tempat kerapatan
meledak. Simulasi memberi $F_Y(0{,}09)=0{,}29944$, $F_Y(0{,}25)=0{,}50003$,
$F_Y(0{,}49)=0{,}69944$, $F_Y(0{,}81)=0{,}90057$ berbanding teori $0{,}3$, $0{,}5$,
$0{,}7$, $0{,}9$. Galat baku tiap angka sekitar $0{,}00145$ pada $p=0{,}3$ dan
$0{,}00095$ pada $p=0{,}9$; keempat simpangan di bawah satu galat baku.

\kertas{Periksa dengan tangan. Hitung $\binom{20}{7}0{,}35^7 0{,}65^{13}$ dan
bandingkan dengan $0{,}184401$; jelaskan mengapa $\sum_{k=0}^{20}p_X(k)=1$ mengikuti
teorema binomial, bukan hasil pengukuran. Terakhir: rata-rata sampel $7{,}005260$
berbeda dari $7$ sebesar $0{,}00526$. Apakah itu kesalahan program? Jawab dengan
galat baku $0{,}006745$.}

\section{Lab komputasi}
\tingkat{lanjut}

Kode berikut dijalankan dengan Python 3, NumPy 2.5.2, dan SciPy 1.18.0. Semua angka
pada tabel dan paragraf sebelumnya berasal dari eksekusi ini; benihnya tetap
$20261024$, sehingga hasilnya dapat diulang.

\begin{lab}
\begin{verbatim}
import numpy as np
from math import exp, lgamma, log, sqrt, pi
from scipy import stats
from scipy.integrate import quad

rng = np.random.default_rng(20261024)
R = 100_000

# 1. binomial: rumus eksak vs Monte Carlo
n, p = 20, 0.35
mu, var = n * p, n * p * (1 - p)
sd = sqrt(var)
x = rng.binomial(n, p, size=R)
se = sd / sqrt(R)
print("Binomial n=20 p=0.35")
print(f"  teori E={mu:.6f} Var={var:.6f} sd={sd:.6f}")
print(f"  MC    rata2={x.mean():.6f} Var={x.var(ddof=1):.6f}")
print(f"  galat baku rata2={se:.6f} z={(x.mean()-mu)/se:+.4f}")
print(f"  galat baku Var  ={var*sqrt(2/(R-1)):.6f}")
print(f"  P(X>=12)={stats.binom.sf(11, n, p):.6f}")

# 2. jumlah pmf Poisson, dihitung dalam ruang log
print("\nJumlah pmf Poisson, k=0..K")
for lam in (0.5, 1.0, 3.0, 10.0, 50.0, 100.0):
    K = int(lam + 12 * sqrt(lam) + 60)
    s = sum(exp(-lam + k*log(lam) - lgamma(k+1)) for k in range(K+1))
    print(f"  lam={lam:6.1f} K={K:4d} jumlah={s:.17f}"
          f" 1-jumlah={1-s:+.2e}")

# 3. total kerapatan normal harus 1
phi = lambda t: exp(-t*t/2) / sqrt(2*pi)
print("\nKerapatan normal baku")
print(f"  phi(0)={phi(0):.15f} 1/sqrt(2pi)={1/sqrt(2*pi):.15f}")
print(f"  quad [-8,8]  ={quad(phi, -8, 8)[0]:.17f}")
print(f"  sisa={1-quad(phi, -8, 8)[0]:.2e}"
      f"  2*Phi(-8)={2*stats.norm.cdf(-8):.2e}")
v, e = quad(phi, -np.inf, np.inf)
print(f"  quad (-inf,inf)={v:.17f} galat lapor={e:.1e}")

# 4. transformasi X~Unif(0,1)
u = rng.uniform(0, 1, R)
y = u * u
print("\nY=X^2, X~Unif(0,1)")
print(f"  integral f_Y={quad(lambda t: 1/(2*sqrt(t)), 0, 1)[0]:.16f}")
for yy in (0.09, 0.25, 0.49, 0.81):
    print(f"  F_Y({yy:.2f}) MC={np.mean(y<=yy):.5f} teori={sqrt(yy):.5f}")
z = -np.log(u)
print("Y=-ln(X), harusnya Exp(1)")
print(f"  rata2={z.mean():.5f} galat baku={1/sqrt(R):.5f}"
      f" Var={z.var(ddof=1):.5f}")
print(f"  P(Y>2) MC={np.mean(z>2):.5f} teori={exp(-2):.5f}")

# 5. geometrik dan eksponensial
g = rng.geometric(0.2, R)
print("\nGeometrik p=0.2: E=5 Var=20")
print(f"  rata2={g.mean():.5f} galat baku={sqrt(20/R):.5f}"
      f" Var={g.var(ddof=1):.5f}")
w = rng.exponential(2.0, R)
print("Eksponensial beta=2: E=2 Var=4")
print(f"  rata2={w.mean():.5f} galat baku={2/sqrt(R):.5f}"
      f" Var={w.var(ddof=1):.5f}")
print(f"  P(X>3) MC={np.mean(w>3):.5f} teori={exp(-1.5):.5f}")
print(f"  P(X>5|X>2) MC={np.mean(w[w>2]>5):.5f} teori={exp(-1.5):.5f}")
\end{verbatim}
\end{lab}

Keluaran yang diperoleh, apa adanya:

\begin{verbatim}
Binomial n=20 p=0.35
  teori E=7.000000 Var=4.550000 sd=2.133073
  MC    rata2=7.005260 Var=4.552918
  galat baku rata2=0.006745 z=+0.7798
  galat baku Var  =0.020348
  P(X>=12)=0.019579

Jumlah pmf Poisson, k=0..K
  lam=   0.5 K=  68 jumlah=1.00000000000000000 1-jumlah=+0.00e+00
  lam=   1.0 K=  73 jumlah=1.00000000000000000 1-jumlah=+0.00e+00
  lam=   3.0 K=  83 jumlah=1.00000000000000022 1-jumlah=-2.22e-16
  lam=  10.0 K= 107 jumlah=1.00000000000000311 1-jumlah=-3.11e-15
  lam=  50.0 K= 194 jumlah=0.99999999999999778 1-jumlah=+2.22e-15
  lam= 100.0 K= 280 jumlah=1.00000000000003886 1-jumlah=-3.89e-14

Kerapatan normal baku
  phi(0)=0.398942280401433 1/sqrt(2pi)=0.398942280401433
  quad [-8,8]  =0.99999999999999878
  sisa=1.22e-15  2*Phi(-8)=1.24e-15
  quad (-inf,inf)=0.99999999999999978 galat lapor=1.0e-08

Y=X^2, X~Unif(0,1)
  integral f_Y=0.9999999999999992
  F_Y(0.09) MC=0.29944 teori=0.30000
  F_Y(0.25) MC=0.50003 teori=0.50000
  F_Y(0.49) MC=0.69944 teori=0.70000
  F_Y(0.81) MC=0.90057 teori=0.90000
Y=-ln(X), harusnya Exp(1)
  rata2=0.99811 galat baku=0.00316 Var=0.99623
  P(Y>2) MC=0.13478 teori=0.13534

Geometrik p=0.2: E=5 Var=20
  rata2=4.99176 galat baku=0.01414 Var=19.92571
Eksponensial beta=2: E=2 Var=4
  rata2=1.99895 galat baku=0.00632 Var=3.98400
  P(X>3) MC=0.22237 teori=0.22313
  P(X>5|X>2) MC=0.22028 teori=0.22313
\end{verbatim}

Empat hal terbaca dari keluaran itu. Pertama, setiap rata-rata sampel menyimpang kurang
dari satu galat baku: $z=+0{,}78$ untuk binomial, sedangkan $0{,}99811$, $4{,}99176$,
dan $1{,}99895$ berbanding $1$, $5$, dan $2$ untuk tiga kasus lain. Penyimpangan
sebesar itu harus muncul; kalau nol semua, justru ada yang salah.

Kedua, variansi sampel sedikit di bawah teori pada kasus berekor panjang,
$19{,}92571$ berbanding $20$, karena satu tarikan sangat besar mendominasi jumlah
kuadrat. Ketiga, galat penjumlahan Poisson tumbuh bersama $\lambda$, dari nol menjadi
$10^{-14}$ ketika jumlah sukunya bertambah dari $74$ menjadi $281$. Keempat, kerapatan
hasil berintegral satu sampai $16$ angka, dan $F_Y$ cocok dengan $\sqrt y$ di empat
titik uji; yang kedua lebih penting, karena integral satu hanya memeriksa konstanta,
sedangkan kecocokan di empat titik memeriksa bentuk fungsi.

Kode itu tidak membuktikan satu pun teorema di bab ini; ia memeriksa bahwa rumus yang
ditulis memberi angka yang disebutkan. Bukti $\mathbb{E}[X]=np$ ada pada argumen
penanda, bukti $f_Y(y)=1/(2\sqrt y)$ pada perhitungan fungsi distribusi; keluaran
program adalah data, bukan argumen.

\begin{untukdosen}
Urutan yang saya sarankan: dua sistem dengan rata-rata sama tetapi ekor berbeda, lalu
definisi variabel acak, lalu langsung fungsi distribusi kumulatif. Kerjakan satu tabel
kecil di papan, gambar tangganya, tanyakan $F(1{,}4)$, baru setelah itu kerapatan.
Tahan diri untuk tidak menyebut integral sebelum mahasiswa melihat mengapa penjumlahan
tidak lagi cukup.

Dua miskonsepsi yang paling sering muncul. Pertama, $f_X(x)$ dibaca sebagai peluang,
sehingga kerapatan bernilai $5$ terasa mustahil dan $P(X=x)$ dianggap sama dengan
$f_X(x)$. Kedua, $P(X=x)=0$ disimpulkan berarti mustahil, sehingga
$P(0\le X\le1)=1$ menjadi paradoks. Jawabannya aditif terbilang: gabungan tak
terbilang tidak dijumlah.

Tanda siap lanjut: mahasiswa dapat menjelaskan perbedaan massa dan kerapatan tanpa
membaca catatan, dan menghitung $\mathbb{E}[X^2]$ tanpa mencari kerapatan $X^2$.
\end{untukdosen}

\section{Ringkasan}
\tingkat{dasar}

\begin{ringkasan}
\begin{itemize}
\item Variabel acak adalah fungsi dari ruang hasil ke bilangan real, bukan bilangan
      acak. Fungsi distribusi $F_X(x)=P(X\le x)$ memuat seluruh informasi distribusi:
      tidak turun, kontinu kanan, berlimit $0$ dan $1$.
\item Fungsi massa membobot titik dan nilainya paling besar satu; fungsi kerapatan
      membobot panjang, boleh melebihi satu, dan $P(X=x)=0$ di setiap titik. Peluang
      nol bukan kemustahilan, karena yang menjumlah hanya gabungan terbilang.
\item Harapan adalah titik tumpu dan linear tanpa syarat kebebasan;
      $\operatorname{Var}(X)=\mathbb{E}[X^2]-\mathbb{E}[X]^2$ adalah rata-rata kuadrat
      simpangan.
\item Binomial $\mathbb{E}=np$, $\operatorname{Var}=np(1-p)$, diturunkan lewat
      penanda; geometrik $\mathbb{E}=1/p$, $\operatorname{Var}=(1-p)/p^2$; Poisson
      $\mathbb{E}=\operatorname{Var}=\lambda$, limit binomial untuk kejadian langka.
\item Normal distandardisasi menjadi $N(0,1)$ dan dibaca lewat $\Phi$; eksponensial
      tanpa memori, kembaran kontinu geometrik.
\item Kovarians mengukur ketergantungan linear saja: nol tidak berarti bebas, dan
      transformasi monoton mengubah kerapatan dengan faktor $|ds/dy|$.
\end{itemize}
\end{ringkasan}

\section{Latihan}
\tingkat{menengah}

\begin{latihan}
\begin{enumerate}
  \item Misalkan $p_X(-2)=0{,}1$, $p_X(0)=0{,}4$, $p_X(3)=c$. Tentukan $c$, lalu
        hitung $\mathbb{E}[X]$, $\mathbb{E}[X^2]$, dan $\operatorname{Var}(X)$.
  \item Buktikan $P(a<X\le b)=F_X(b)-F_X(a)$, lalu jelaskan mengapa
        $P(a\le X\le b)$ bisa berbeda dari $P(a<X\le b)$ untuk variabel diskret dan
        mengapa keduanya sama untuk variabel kontinu.
  \item Nilai klaim ini: ``karena $P(X=x)=0$ untuk setiap $x$, kejadian $X=x$
        mustahil.'' Sebutkan bagian yang benar, bagian yang keliru, dan peran aditif
        terbilang.
  \item Turunkan $\mathbb{E}[X^2]=(2-p)/p^2$ untuk geometrik dengan menurunkan
        $\sum k^2q^{k-1}$ dua kali, lalu simpulkan $\operatorname{Var}(X)$ dan periksa
        untuk $p=0{,}2$.
  \item Misalkan $X\sim N(3,4)$. Hitung $P(X>1)$, $P(1<X<5)$, dan kuantil 0,2 untuk
        $X$; bandingkan dengan \texttt{scipy.stats.norm}.
  \item Hitung $\sum_{k=0}^{40}e^{-2{,}5}2{,}5^k/k!$ dan $P(X\ge8)$ untuk Poisson
        $\lambda=2{,}5$, lalu bandingkan dengan binomial $n=1000$, $p=0{,}0025$;
        jelaskan arah selisihnya.
  \item Buktikan eksponensial tanpa memori, lalu jelaskan mengapa geometrik juga tanpa
        memori dan mengapa seragam pada bilangan bulat tidak.
  \item Untuk $X$ seragam pada $(-1,3)$ dan $Y=X^2$, turunkan $f_Y$, periksa
        $\int f_Y=1$, lalu hitung $\mathbb{E}[Y]$ lewat $f_Y$ dan lewat
        $\mathbb{E}[X^2]$.
  \item Konstruksikan dua variabel yang kovariansnya nol tetapi tidak bebas. Hitung
        $\operatorname{Cov}$ dan $\rho$ secara eksak, dan tunjukkan peluang bersama
        yang membuktikan ketidakbebasan.
  \item Berapa ulangan $R$ agar galat baku rata-rata binomial $n=20$, $p=0{,}35$
        paling banyak $0{,}005$? Berapa pula agar galat baku taksiran variansi paling
        banyak $0{,}02$?
\end{enumerate}

\noindent\textbf{Petunjuk.} Butir 3: pisahkan pernyataan tentang peluang dari
pernyataan tentang keberadaan himpunan. Butir 4: turunkan $\sum_{k\ge0}q^k=1/(1-q)$
dua kali terhadap $q$, dan ingat $\sum k^2q^{k-1}$ memuat $\sum k(k-1)q^{k-1}$.
Butir 8: untuk $0<y<1$ himpunan $A_y$ memuat dua cabang, untuk $y\ge1$ hanya satu.
\end{latihan}

\begin{untukpeneliti}
Yang benar-benar terbukti di bab ini: sifat fungsi distribusi beserta buktinya,
hubungan kerapatan dengan fungsi distribusi lewat teorema dasar kalkulus, rumus
harapan fungsi variabel acak untuk kasus diskret, linearitas harapan, rumus variansi,
rata-rata dan variansi binomial lewat penanda, rata-rata geometrik dan Poisson lewat
deret, ketiadaan memori eksponensial, dan rumus transformasi monoton.

Yang dikutip, bukan dibuktikan: bahwa fungsi yang memenuhi tiga sifat fungsi distribusi
selalu berasal dari suatu ruang peluang (Wasserman, Teorema 2.8, yang arahnya memakai
teorema Caratheodory), bahwa fungsi distribusi menentukan seluruh peluang (Wasserman,
Teorema 2.7), normalisasi normal lewat koordinat polar (Bab 15), dan batas korelasi
lewat Cauchy-Schwarz (Bab 10).

Sumber primer: Wasserman, \emph{All of Statistics}, Bagian 2.1--2.4 (termasuk Lemma
2.15 dan peringatan bahwa kerapatan boleh melebihi satu) dan 2.11 (Contoh 2.47 untuk
$Y=X^2$ pada $U(-1,3)$), serta 3.1--3.4. Blitzstein dan Hwang, \emph{Introduction to
Probability}, Bab 3--5 (LOTUS di 4.5, Poisson di 4.7, eksponensial di 5.5) dan 7.3
untuk kovariansi.

Batas keberlakuan: kerapatan tidak tunggal, sehingga setiap pernyataan tentang $f_X$
adalah pernyataan tentang kelas fungsi yang sama hampir di mana-mana; rumus Jacobian
menuntut monoton dan terturunkan; harapan Cauchy tidak ada, sehingga setiap
``rata-rata'' pada data berekor sangat tebal harus dicurigai lebih dahulu. Pertanyaan
terbuka: seberapa baik kerapatan dapat ditaksir dari data tanpa asumsi bentuk, dan
apakah ekor distribusi nyata cukup tipis untuk membenarkan pemakaian variansi.
\end{untukpeneliti}

\peta{Bab 18 pada Wasserman, \emph{All of Statistics}, Bab 2 (2.1--2.4, Lemma 2.15,
dan 2.11 untuk transformasi) dan Bab 3 (3.1--3.4). Pada Blitzstein dan Hwang,
\emph{Introduction to Probability}, Bab 3 (3.6--3.7), Bab 4 (4.2, 4.5, 4.7), Bab 5
(5.1--5.5) dan 7.3. Integral yang dipakai adalah integral Riemann Bab 13; perubahan
variabel dimensi lebih dari satu ada di Bab 15.}

\chapter{Hukum Bilangan Besar dan Teorema Limit Pusat}

Satu program pendek menghitung rata-rata dari $\exp(U)$ untuk sepuluh ribu bilangan
acak $U$ yang seragam pada $(0,1)$. Hasil cetaknya $1{,}722206$. Nilai sebenarnya
adalah $e-1=1{,}7182818285$, jadi galat jalankan itu $+0{,}003924$. Program
dijalankan ulang dengan empat kali lipat bilangan: $1{,}717502$, galat
$-0{,}000780$. Dengan sepuluh kali lipat: $1{,}717767$, galat $-0{,}000515$.
Polanya tidak rapi. Empat kali data memangkas galat lebih dari lima kali, sedangkan
sepuluh kali data hanya memangkasnya $7{,}6$ kali. Tebakan naif "sepuluh kali data,
sepersepuluh galat" tidak berlaku pada satu jalankan. Satu jalankan program tidak
memuat laju apa pun; yang punya laju adalah sebaran dari banyak jalankan.

Sebaran itu bisa diukur. Dengan $n=10^4$ simpangan baku penaksir ini $0{,}00492$;
dengan $n=4\cdot10^4$ menjadi $0{,}00246$; dengan $n=10^5$ menjadi $0{,}00156$.
Setiap kali $n$ naik empat kali, simpangan itu tepat separuh. Di sini ada dua
pernyataan berbeda yang sering tertukar. Pertama, rata-rata memang mendekat ke
nilai sebenarnya: itu kecenderungan, dan namanya hukum bilangan besar. Kedua, sisa
galat punya ukuran dan bentuk tertentu: itu fluktuasi, dan namanya teorema limit
pusat. Bab ini memberi nama, bukti, dan batas keberlakuan untuk keduanya.

Bab 18 sudah menyiapkan bahan: harapan, variansi, dan distribusi jumlah lewat
konvolusi. Konvolusi hanya berguna bila $n$ kecil dan bentuk asalnya sederhana;
untuk sebagian besar distribusi tidak ada bentuk tertutup. Yang kita butuhkan
bukan distribusi tepat pada satu $n$, melainkan pernyataan yang berlaku untuk
semua $n$ yang besar. Pertanyaan praktisnya dua: seberapa dekat rata-rata boleh
dipercaya, dan seberapa besar galat yang wajar. Jawabannya datang dari
memeriksa bagaimana barisan distribusi itu berperilaku, bukan dari menghitung
satu distribusi.

\begin{untukanak}
Bayangkan satu kelas diminta menebak banyaknya kelereng di dalam stoples besar.
Tebakan anak-anak jauh meleset ke dua arah. Ada yang menebak tiga ratus, ada yang
menebak dua ribu, padahal isinya delapan ratus. Guru lalu menuliskan semua tebakan
di papan dan menghitung rata-ratanya. Angkanya ternyata dekat sekali dengan
delapan ratus.

Sekarang gurunya menambah murid. Dengan tiga puluh penebak, rata-ratanya masih
goyang. Dengan tiga ratus penebak, rata-ratanya lebih tenang. Dengan tiga ribu
penebak, rata-ratanya hampir tidak bergeser lagi. Ada dua hal aneh di sini.
Pertama, tidak seorang pun penebak yang menjadi lebih jago; yang berubah hanya
banyaknya penebak. Kedua, menambah penebak dari tiga puluh ke tiga ratus menolong
jauh lebih banyak daripada menambah dari tiga ratus ke tiga ribu. Pertolongannya
terus bertambah, tetapi makin lama makin sedikit.

Pertanyaan yang membuat penasaran: kalau setiap tebakan bisa meleset sangat jauh,
mengapa rata-ratanya justru bisa dipercaya? Dan mengapa pertolongannya melambat
dengan pola yang bisa ditebak, bukan sembarangan? Di bab ini kita akan tahu
mengapa, dan sekaligus tahu kapan kebiasaan mempercayai rata-rata itu menipu.
\end{untukanak}

\section{Tiga arti kata menuju pada variabel acak}
\tingkat{dasar}

Kalimat sehari-harinya: barisan variabel acak bisa "menuju" dalam tiga arti yang
berbeda, dan arti yang paling longgar sama sekali tidak mengatakan bahwa nilainya
mendekat.

Barisan bilangan real hanya punya satu arti menuju, yaitu jarak $|x_n-L|$ menjadi
kecil. Untuk variabel acak, jarak itu sendiri adalah variabel acak, sehingga kita
harus memilih besaran mana yang dibuat kecil. Pilihan pertama adalah peluangnya.
Ini pilihan yang paling lemah secara intuitif tetapi paling berguna secara praktis.

\kamus{konvergensi dalam peluang}{convergence in probability}{janji bahwa peluang
meleset lebih dari toleransi mengecil menuju nol, bukan janji bahwa nilainya
mendekat pada satu jalankan}

\begin{definisi}[Konvergensi dalam peluang]
Barisan variabel acak $X_1,X_2,\dots$ \emph{konvergen dalam peluang}
\index{konvergensi dalam peluang} ke variabel acak $X$ jika untuk setiap
$\varepsilon>0$ berlaku
\[
P\bigl(|X_n-X|>\varepsilon\bigr)\to 0 \qquad (n\to\infty).
\]
\end{definisi}

Perhatikan letak kuantornya. Nilai $\varepsilon$ dipilih lebih dahulu, lalu
peluangnya boleh mengecil sesudah itu. Definisi ini tidak mengatakan bahwa
$|X_n-X|$ mengecil untuk satu hasil percobaan tertentu; ia mengatakan bahwa
kejadian "meleset lebih dari $\varepsilon$" menjadi makin jarang. Pada percobaan
yang sama, $X_n$ masih boleh meleset jauh, asal kejadian itu jarang.

\kamus{konvergensi dalam distribusi}{convergence in distribution}{yang menyerupai
adalah bentuk sebarannya, bukan nilainya; dua variabel bisa sebentuk tanpa pernah
berdekatan}

\begin{definisi}[Konvergensi dalam distribusi]
Misalkan $F_n$ fungsi sebaran $X_n$ dan $F$ fungsi sebaran $X$. Kita tulis
$X_n\rightsquigarrow X$ bila $F_n(t)\to F(t)$ untuk setiap titik $t$ tempat $F$
kontinu. Batasnya disebut \emph{limit dalam distribusi}\index{konvergensi dalam
distribusi}.
\end{definisi}

Syarat "di titik kontinu" bukan keisengan teknis. Ambil $X_n$ yang bernilai $1/n$
dengan peluang satu, dan $X$ yang bernilai $0$ dengan peluang satu. Maka $F_n(0)=0$
untuk setiap $n$, sedangkan $F(0)=1$, jadi barisan itu tidak konvergen di $t=0$.
Di setiap $t\neq0$ barisan itu konvergen, dan $t=0$ adalah satu-satunya titik
ketakkontinuan $F$. Tanpa syarat tadi, kesimpulan yang benar akan ditolak hanya
karena satu titik.

Arti ketiga lebih kuat dan lebih harfiah: hampir setiap hasil percobaan, sebagai
satu titik $\omega$ di ruang sampel, menghasilkan barisan bilangan real yang
konvergen.

\begin{definisi}[Konvergensi hampir pasti]
$X_n$ \emph{konvergen hampir pasti}\index{konvergensi hampir pasti} ke $X$ bila
\[
P\bigl(\{\omega : X_n(\omega)\to X(\omega)\}\bigr)=1 .
\]
\end{definisi}

Di sini barisan bilangan real $X_1(\omega),X_2(\omega),\dots$ benar-benar
konvergen dalam arti Bab 4, dan himpunan hasil percobaan yang gagal konvergen
hanya boleh berpeluang nol. Karena itu konvergensi hampir pasti berbicara tentang
satu jalankan yang panjang, sedangkan konvergensi dalam peluang berbicara tentang
satu $n$ pada banyak jalankan sekaligus.

Ketiga arti itu tidak setara, tetapi ada arah implikasi yang tetap.

\begin{teorema}[Hubungan tiga modus]
Untuk variabel acak $X_n$ dan $X$ berlaku:
(a) $X_n\to X$ hampir pasti mengakibatkan $X_n\to X$ dalam peluang;
(b) $X_n\to X$ dalam peluang mengakibatkan $X_n\rightsquigarrow X$;
(c) bila $X_n\rightsquigarrow c$ dan $c$ konstanta, maka $X_n\to c$ dalam peluang.
\end{teorema}

\begin{proof}[Bukti]
Untuk (a), tetapkan $\varepsilon>0$ dan misalkan
$A_n=\{\sup_{k\ge n}|X_k-X|>\varepsilon\}$. Himpunan $A_n$ menurun, dan setiap
$\omega$ di $\bigcap_n A_n$ memenuhi $|X_k(\omega)-X(\omega)|>\varepsilon$ untuk
tak hingga banyak $k$, sehingga barisan itu tidak konvergen ke $X(\omega)$.
Jadi $\bigcap_n A_n$ termuat di himpunan berpeluang nol. Karena peluang bersifat
kontinu dari atas pada himpunan berpeluang hingga,
$P(A_n)\to P(\bigcap_n A_n)=0$. Kejadian $\{|X_n-X|>\varepsilon\}$ termuat di
$A_n$, maka peluangnya juga menuju nol.

Untuk (b), ambil $t$ titik kontinu $F$ dan $\varepsilon>0$. Dari
$\{X_n\le t\}\subseteq\{X\le t+\varepsilon\}\cup\{|X_n-X|>\varepsilon\}$ kita
peroleh
\[
F_n(t)\le F(t+\varepsilon)+P\bigl(|X_n-X|>\varepsilon\bigr).
\]
Dengan cara yang sama, dari $\{X\le t-\varepsilon\}\subseteq\{X_n\le t\}\cup
\{|X_n-X|>\varepsilon\}$ kita peroleh
$F(t-\varepsilon)\le F_n(t)+P(|X_n-X|>\varepsilon)$. Kedua ketaksamaan itu
memberi
\[
F(t-\varepsilon)-\delta_n\le F_n(t)\le F(t+\varepsilon)+\delta_n,
\qquad \delta_n=P\bigl(|X_n-X|>\varepsilon\bigr).
\]
Karena $\delta_n\to0$, limit inferior dan limit superior $F_n(t)$ terkurung di
antara $F(t-\varepsilon)$ dan $F(t+\varepsilon)$. Biarkan $\varepsilon\downarrow0$
dan pakai kekontinuan $F$ di $t$; keduanya menuju $F(t)$.

Untuk (c), ambil $\varepsilon>0$. Karena $F$ adalah sebaran titik massa di $c$,
titik $c-\varepsilon$ dan $c+\varepsilon$ keduanya titik kontinu $F$, sehingga
\[
P\bigl(|X_n-c|>\varepsilon\bigr)
\le P(X_n<c-\varepsilon)+P(X_n>c+\varepsilon)
\le F_n(c-\varepsilon)+1-F_n(c+\varepsilon)
\to 0+1-1=0 .
\]
\end{proof}

Tiga arah sebaliknya tidak berlaku, dan contoh penyangkalnya perlu diingat karena
masing-masing menunjukkan apa yang sesungguhnya hilang.

\begin{contoh}[Dalam peluang tetapi tidak hampir pasti]
Pecah $(0,1)$ menjadi $k$ bagian sama panjang, untuk $k=1,2,3,\dots$, lalu
urutkan semua bagian itu menjadi satu daftar. Definisikan $X_n(\omega)=1$ bila
$\omega$ berada di bagian ke-$n$ dan $0$ bila tidak. Bagian ke-$n$ berasal dari
pemecahan dengan $k_n$ bagian, jadi $P(X_n=1)=1/k_n\to0$ dan $X_n\to0$ dalam
peluang. Tetapi untuk setiap $\omega\in(0,1)$, dalam setiap blok pemecahan ada
tepat satu $n$ dengan $X_n(\omega)=1$. Nilai satu itu muncul tak hingga sering,
sehingga $X_n(\omega)$ tidak konvergen, untuk setiap $\omega$. Jadi tidak ada
konvergensi hampir pasti. Barisan ini bahkan konvergen dalam $L^1$, karena
$\mathbb{E}|X_n|=1/k_n\to0$, maka kekuatan $L^1$ pun tidak menutup celah itu.
\end{contoh}

\begin{contoh}[Dalam distribusi tetapi tidak dalam peluang]
Ambil $X\sim N(0,1)$ dan tetapkan $X_n=-X$ untuk setiap $n$. Sebaran $X_n$ sama
persis dengan sebaran $X$, jadi $F_n(t)=F(t)$ dan $X_n\rightsquigarrow X$ secara
hampa. Tetapi $P(|X_n-X|>\varepsilon)=P(2|X|>\varepsilon)>0$ untuk setiap
$\varepsilon$, dan peluang itu tidak bergantung pada $n$. Dua variabel yang
sebentuk belum tentu berdekatan: $X$ dan $-X$ bisa berjarak jauh pada jalankan
yang sama.
\end{contoh}

Pelajaran dari dua contoh itu: konvergensi dalam distribusi hanya menyentuh
sebaran satu per satu, sedangkan dua modus lainnya menyentuh hubungan
antar-variabel pada ruang sampel yang sama. Karena itu teorema (c) penting secara
praktis: dalam statistika limit kita hampir selalu konstanta, dan di situ ketiga
modus bertemu.

Satu peringatan lain harus dicatat sejak awal, karena kesalahannya mahal.

\begin{catatan}
Konvergensi dalam peluang tidak membawa serta konvergensi harapan. Misalkan
$P(X_n=n^2)=1/n$ dan $P(X_n=0)=1-1/n$. Untuk setiap $\varepsilon>0$,
$P(|X_n|>\varepsilon)=1/n\to0$, jadi $X_n\to0$ dalam peluang. Tetapi
$\mathbb{E}[X_n]=n$, sehingga harapannya menuju tak hingga. Rata-rata yang
memusat bukan jaminan bahwa harapannya memusat.
\end{catatan}

\section{Ketaksamaan Chebyshev dan hukum bilangan besar lemah}
\tingkat{menengah}

Kalimat sehari-harinya: peluang menyimpang jauh dari rata-rata dibatasi oleh
variansi dibagi kuadrat jarak penyimpangan itu.

Untuk membuktikan hukum bilangan besar kita hanya butuh satu ketaksamaan, dan
ketaksamaan itu butuh satu pengamatan sederhana tentang harapan.

\begin{lemma}[Markov]
Jika $Y\ge0$ dan $a>0$, maka $P(Y\ge a)\le \mathbb{E}[Y]/a$.
\index{ketaksamaan Markov}
\end{lemma}

\begin{proof}[Bukti]
Tulis $Y=Y\mathbf{1}_{\{Y<a\}}+Y\mathbf{1}_{\{Y\ge a\}}$. Suku pertama tak
negatif. Pada kejadian $Y\ge a$ berlaku $Y\ge a$, sehingga harapan suku kedua
paling sedikit $a\,P(Y\ge a)$. Jadi
$\mathbb{E}[Y]\ge 0+a\,P(Y\ge a)$.
\end{proof}

\begin{akibat}[Ketaksamaan Chebyshev]
Jika $\mathbb{E}[X]=\mu$ dan $\operatorname{Var}(X)=\sigma^2<\infty$, maka untuk
setiap $a>0$
\[
P\bigl(|X-\mu|\ge a\bigr)\le\frac{\sigma^2}{a^2}.
\]
\index{ketaksamaan Chebyshev}
\end{akibat}

\begin{proof}[Bukti]
Terapkan Lemma Markov pada $Y=(X-\mu)^2$ dengan ambang $a^2$. Kejadian
$\{Y\ge a^2\}$ sama dengan $\{|X-\mu|\ge a\}$, dan
$\mathbb{E}[Y]=\sigma^2$.
\end{proof}

Ketaksamaan ini lemah, dan kelemahannya bukan cacat tersembunyi melainkan
kekuatannya. Ia tidak mengetahui bentuk distribusi apa pun, hanya dua momen.
Karena itu ia berlaku untuk semua distribusi bervariansi hingga, dan justru
karena itu ia cukup untuk membuktikan teorema yang berlaku umum.

\kamus{hukum bilangan besar}{law of large numbers}{janji bahwa rata-rata banyak
pengamatan menempel ke rata-rata sebenarnya, dengan peluang meleset yang mengecil}

\begin{teorema}[Hukum bilangan besar lemah]
Misalkan $X_1,X_2,\dots$ bebas dan berdistribusi sama dengan
$\mathbb{E}[X_i]=\mu$ dan $\operatorname{Var}(X_i)=\sigma^2<\infty$. Maka untuk
setiap $\varepsilon>0$\index{hukum bilangan besar}
\[
P\bigl(|\bar X_n-\mu|\ge\varepsilon\bigr)\to0,
\qquad \bar X_n=\frac1n\sum_{i=1}^n X_i .
\]
\end{teorema}

\begin{proof}[Bukti]
Karena harapan linear, $\mathbb{E}[\bar X_n]=\mu$. Karena variabelnya bebas,
variansi jumlah adalah jumlah variansi, sehingga
$\operatorname{Var}(\bar X_n)=\sigma^2/n$. Ketaksamaan Chebyshev dengan
$a=\varepsilon$ memberi
\[
P\bigl(|\bar X_n-\mu|\ge\varepsilon\bigr)\le\frac{\sigma^2}{n\varepsilon^2}.
\]
Ruas kanan menuju nol untuk setiap $\varepsilon>0$ tetap.
\end{proof}

Perhatikan apa yang baru saja dibuktikan dan apa yang tidak. Yang dibuktikan:
untuk toleransi $\varepsilon$ yang kita pilih sendiri, peluang rata-rata meleset
lebih dari $\varepsilon$ menuju nol. Yang tidak dibuktikan: seberapa cepat, dan
bagaimana bentuk sisa galatnya. Keduanya pertanyaan berbeda.

Gambar~\ref{fig:tiga-lebar} menggambar isi pernyataan itu untuk dadu. Yang
berubah dari $n=4$ ke $n=64$ hanya dua: lebarnya menyusut, puncaknya meninggi,
sedangkan pusatnya tetap di $3{,}5$. Hukum bilangan besar bukan janji bahwa satu
jalankan lebih tepat, melainkan bahwa sebaran kemungkinannya mengerut ke satu
titik.

\begin{figure}[htbp]
\centering
\begin{tikzpicture}
\begin{axis}[
  width=0.95\linewidth, height=5.4cm,
  xmin=1, xmax=6, ymin=0, ymax=2.05,
  xlabel={nilai rata-rata $\bar x_n$}, ylabel={kerapatan},
  label style={font=\small}, tick label style={font=\footnotesize},
  axis lines=left,
  legend style={font=\footnotesize, draw=layergrey!35, fill=white,
    fill opacity=0.9, text opacity=1, at={(0.015,0.985)},
    anchor=north west, cells={anchor=west}},
  legend cell align=left,
]
\addplot[draw=layerteal!60!black, fill=layerteal!18, fill opacity=0.85,
  line width=0.4pt] coordinates{
  (1,0.0031) (1.25,0.0123) (1.5,0.0309) (1.75,0.0617) (2,0.108) (2.25,0.1728)
  (2.5,0.2469) (2.75,0.321) (3,0.3858) (3.25,0.4321) (3.5,0.4506) (3.75,0.4321)
  (4,0.3858) (4.25,0.321) (4.5,0.2469) (4.75,0.1728) (5,0.108) (5.25,0.0617)
  (5.5,0.0309) (5.75,0.0123) (6,0.0031)};
\addlegendentry{$n=4$}
\addplot[draw=layerindigo!70!black, fill=layerindigo!22, fill opacity=0.8,
  line width=0.4pt] coordinates{
  (1.6875,0) (1.75,0.0001) (1.8125,0.0002) (1.875,0.0004) (1.9375,0.0008) (2,0.0014)
  (2.0625,0.0025) (2.125,0.0043) (2.1875,0.0072) (2.25,0.0117) (2.3125,0.0184) (2.375,0.028)
  (2.4375,0.0416) (2.5,0.0601) (2.5625,0.0847) (2.625,0.1163) (2.6875,0.1558) (2.75,0.2039)
  (2.8125,0.2607) (2.875,0.3258) (2.9375,0.3981) (3,0.4759) (3.0625,0.5567) (3.125,0.6374)
  (3.1875,0.7145) (3.25,0.7843) (3.3125,0.8431) (3.375,0.8878) (3.4375,0.9156) (3.5,0.9251)
  (3.5625,0.9156) (3.625,0.8878) (3.6875,0.8431) (3.75,0.7843) (3.8125,0.7145) (3.875,0.6374)
  (3.9375,0.5567) (4,0.4759) (4.0625,0.3981) (4.125,0.3258) (4.1875,0.2607) (4.25,0.2039)
  (4.3125,0.1558) (4.375,0.1163) (4.4375,0.0847) (4.5,0.0601) (4.5625,0.0416) (4.625,0.028)
  (4.6875,0.0184) (4.75,0.0117) (4.8125,0.0072) (4.875,0.0043) (4.9375,0.0025) (5,0.0014)
  (5.0625,0.0008) (5.125,0.0004) (5.1875,0.0002) (5.25,0.0001) (5.3125,0)};
\addlegendentry{$n=16$}
\addplot[draw=accent!80!black, fill=accent!30, fill opacity=0.8,
  line width=0.5pt] coordinates{
  (2.5312,0) (2.5469,0.0001) (2.5625,0.0001) (2.5781,0.0001) (2.5938,0.0002) (2.6094,0.0003)
  (2.625,0.0004) (2.6406,0.0005) (2.6562,0.0007) (2.6719,0.0009) (2.6875,0.0012) (2.7031,0.0016)
  (2.7188,0.0021) (2.7344,0.0028) (2.75,0.0036) (2.7656,0.0047) (2.7812,0.0061) (2.7969,0.0079)
  (2.8125,0.01) (2.8281,0.0127) (2.8438,0.0161) (2.8594,0.0202) (2.875,0.0252) (2.8906,0.0312)
  (2.9062,0.0385) (2.9219,0.0472) (2.9375,0.0576) (2.9531,0.0698) (2.9688,0.0842) (2.9844,0.1009)
  (3,0.1203) (3.0156,0.1426) (3.0312,0.1681) (3.0469,0.1971) (3.0625,0.2297) (3.0781,0.2663)
  (3.0938,0.307) (3.1094,0.3521) (3.125,0.4015) (3.1406,0.4553) (3.1562,0.5136) (3.1719,0.5762)
  (3.1875,0.6429) (3.2031,0.7134) (3.2188,0.7875) (3.2344,0.8645) (3.25,0.944) (3.2656,1.0252)
  (3.2812,1.1075) (3.2969,1.19) (3.3125,1.2718) (3.3281,1.352) (3.3438,1.4295) (3.3594,1.5035)
  (3.375,1.573) (3.3906,1.6369) (3.4062,1.6943) (3.4219,1.7445) (3.4375,1.7867) (3.4531,1.8202)
  (3.4688,1.8445) (3.4844,1.8592) (3.5,1.8641) (3.5156,1.8592) (3.5312,1.8445) (3.5469,1.8202)
  (3.5625,1.7867) (3.5781,1.7445) (3.5938,1.6943) (3.6094,1.6369) (3.625,1.573) (3.6406,1.5035)
  (3.6562,1.4295) (3.6719,1.352) (3.6875,1.2718) (3.7031,1.19) (3.7188,1.1075) (3.7344,1.0252)
  (3.75,0.944) (3.7656,0.8645) (3.7812,0.7875) (3.7969,0.7134) (3.8125,0.6429) (3.8281,0.5762)
  (3.8438,0.5136) (3.8594,0.4553) (3.875,0.4015) (3.8906,0.3521) (3.9062,0.307) (3.9219,0.2663)
  (3.9375,0.2297) (3.9531,0.1971) (3.9688,0.1681) (3.9844,0.1426) (4,0.1203) (4.0156,0.1009)
  (4.0312,0.0842) (4.0469,0.0698) (4.0625,0.0576) (4.0781,0.0472) (4.0938,0.0385) (4.1094,0.0312)
  (4.125,0.0252) (4.1406,0.0202) (4.1562,0.0161) (4.1719,0.0127) (4.1875,0.01) (4.2031,0.0079)
  (4.2188,0.0061) (4.2344,0.0047) (4.25,0.0036) (4.2656,0.0028) (4.2812,0.0021) (4.2969,0.0016)
  (4.3125,0.0012) (4.3281,0.0009) (4.3438,0.0007) (4.3594,0.0005) (4.375,0.0004) (4.3906,0.0003)
  (4.4062,0.0002) (4.4219,0.0001) (4.4375,0.0001) (4.4531,0.0001) (4.4688,0)};
\addlegendentry{$n=64$}
\draw[layergrey, densely dotted, line width=0.6pt] (axis cs:3.5,0) -- (axis cs:3.5,1.95);
\node[font=\footnotesize, layergrey, anchor=south east] at (axis cs:3.45,1.95) {$3{,}5$};
\end{axis}
\end{tikzpicture}
\caption{Tiga sebaran rata-rata $n$ lemparan dadu, digambar pada satu skala yang
sama: $n=4$, $16$, dan $64$. Lebarnya $0{,}853913$, $0{,}426956$, dan
$0{,}213478$, jadi setiap kali $n$ dilipat empatkan lebarnya tinggal separuh.
Puncaknya meninggi dari $0{,}4506$ ke $0{,}9251$ dan $1{,}8641$, sedangkan
pusatnya tetap di $3{,}5$.}
\label{fig:tiga-lebar}
\end{figure}

Angka-angka membuat perbedaan itu terasa. Untuk lemparan koin dengan $p=1/2$ dan
toleransi $\varepsilon=0{,}1$, Chebyshev memberi
$\sigma^2/(n\varepsilon^2)=0{,}25/(0{,}01n)=25/n$. Agar peluang meleset paling
besar $0{,}3$, cukup $n\ge25/0{,}3=83{,}33$, yaitu $n=84$. Pada $n=84$ batas itu
bernilai $0{,}297619$, jadi Chebyshev menjamin peluang berada dalam
$0{,}4$ sampai $0{,}6$ paling sedikit $0{,}702381$. Peluang sebenarnya, dihitung
dari sebaran binomial, adalah $0{,}937028$, dan aproksimasi normal memberi
$0{,}933197$. Jaminan itu benar tetapi terlalu rendah sekitar $23$ poin
persentase.

Kelemahan yang sama menjadi parah bila $\sigma$ besar. Untuk dadu adil,
$\operatorname{Var}(X)=35/12=2{,}916667$ dan $\sigma=1{,}707825$. Pada $n=100$
dengan $\varepsilon=0{,}1$, Chebyshev memberi batas $2{,}916667$, angka yang lebih
besar dari satu dan karena itu tidak berarti apa pun. Untuk mendapatkan galat
$0{,}01$ dengan keyakinan $95$ persen, rumus normal menuntut paling sedikit
$112043$ pengamatan, sedangkan Chebyshev menuntut jauh lebih banyak lagi.

Di sinilah letak persoalan yang menyiapkan bagian berikutnya. Chebyshev menjawab
"apakah rata-rata boleh dipercaya" dengan benar tetapi tumpul. Ia tidak mampu
menjawab "seberapa dekat, dengan peluang berapa" secara layak dipakai. Jawaban
tajam menuntut pengetahuan tentang bentuk sebaran galat, bukan hanya ukurannya.

\begin{bacakeras}
Dua suara, bergantian. Yang pertama menyebut banyaknya lemparan; yang kedua
menyebut jarak meleset yang masih wajar.

--- Sepuluh lemparan. Melesetnya masih lebar.
--- Seratus lemparan. Sudah lebih sempit.
--- Seribu lemparan. Makin sempit lagi.

Sekarang lipat empatkan, jangan tambah sepuluh: sepuluh jadi empat puluh,
seratus jadi empat ratus, seribu jadi empat ribu. Melesetnya tinggal separuh.
Bukan seperempat. Separuh. Setiap kali empat kali, tepat separuh.

Lalu balik arahnya, dan bacakan lebih pelan. Berapa ongkos satu angka tambahan
di belakang koma? Seratus kali lemparan. Seratus kali kerja, satu angka.
Rata-rata tidak sabar: ia menuntut empat kali, bukan dua kali, dan tidak
menawarkan diskon.

Baris yang diulang tiga kali: empat kali data, separuh galat.
\end{bacakeras}

\begin{untukanak}
Bayangkan mesin pengisi kantong beras. Setiap kantong meleset sedikit dari berat
yang diminta, kadang lebih kadang kurang, dan besar melesetnya berbeda-beda;
menimbang satu kantong tidak memberi tahu arah melesetnya.

Sekarang mesin itu mengisi seratus kantong, lalu kita menimbang semuanya
sekaligus. Meleset ke atas dan ke bawah saling memotong, sehingga berat totalnya
jauh lebih dekat ke seratus kali berat yang diminta. Kalau kantongnya sepuluh
ribu, jumlahnya hampir tepat.

Yang mengejutkan, bentuk sebaran meleset itu tidak perlu kita ketahui. Entah
melesetnya berupa lonceng, berupa dua nilai saja, atau berupa apa pun yang aneh,
sebaran jumlahnya tetap menyerupai lonceng yang sama, asalkan kantongnya cukup
banyak dan melesetnya tidak terlalu sering menjadi raksasa. Inilah yang membuat
para penjudi, pembuat jembatan, dan peneliti cuaca bisa memakai rumus yang sama.

Pertanyaan yang membuat penasaran: mengapa bentuk asli melesetnya tidak lagi
penting setelah kantongnya banyak? Dan mengapa satu kantong yang melesetnya
raksasa bisa merusak seluruh kesimpulan? Di bab ini kita akan tahu mengapa.
\end{untukanak}

\section{Hukum yang kuat: apa yang sesungguhnya ditambahkan}
\tingkat{lanjut}

Kalimat sehari-harinya: dengan keyakinan satu, rata-rata pada satu jalankan
panjang akhirnya menetap, bukan sekadar sering dekat.

Hukum lemah berbicara tentang setiap $n$ secara terpisah: pada $n$ yang cukup
besar, peluang meleset kecil. Ia tidak melarang rata-rata meleset jauh pada
$n=10^6$, lalu kembali tenang, lalu meleset lagi pada $n=10^9$, asalkan kejadian
semacam itu jarang pada setiap $n$. Hukum kuat menutup celah itu.

\begin{teorema}[Hukum bilangan besar kuat]
Misalkan $X_1,X_2,\dots$ bebas dan berdistribusi sama dengan
$\mathbb{E}|X_1|<\infty$ dan $\mathbb{E}[X_1]=\mu$. Maka
$\bar X_n\to\mu$ hampir pasti.\index{hukum bilangan besar kuat}
\end{teorema}

Bedanya halus tetapi tajam. Barisan $X_n$ pada bagian sebelumnya konvergen dalam
peluang ke nol, tetapi pada setiap $\omega$ nilainya satu tak hingga sering.
Hukum lemah tidak dapat menolaknya; hukum kuat menolaknya, karena ia menuntut
konvergensi pada setiap $\omega$ di luar himpunan berpeluang nol.

Perhatikan juga syaratnya. Hukum kuat hanya menuntut $\mathbb{E}|X_1|<\infty$,
tanpa variansi. Hukum lemah yang kita buktikan di atas menuntut $\sigma^2<\infty$
hanya karena buktinya memakai Chebyshev; syarat itu bukan syarat alami teoremanya.
Untuk hukum lemah, syarat perlu dan cukup bagi barisan iid adalah
$n\,P(|X_1|>n)\to0$, yang lebih lemah daripada $\mathbb{E}|X_1|<\infty$. Rincian
itu kita tinggalkan, dan pembaca dapat menemukannya pada Feller, \emph{An
Introduction to Probability Theory and Its Applications}, jilid 2, bab tentang
hukum bilangan besar.

Bukti hukum kuat tidak kita kerjakan di sini. Bentuk klasiknya memotong variabel
pada ambang $n$, menerapkan Borel--Cantelli pada ekor potongan, lalu memakai
kriteria Cauchy untuk kekonvergenan hampir pasti; Etemadi memberi bukti yang
menghindari asumsi variansi dan hanya memakai pasangan variabel yang saling
bebas. Lokasi buktinya: Blitzstein dan Hwang, \emph{Introduction to Probability},
bab 10 bagian 10.2; Wasserman, \emph{All of Statistics}, lampiran bab 5 bagian
5.7.1. Yang perlu diingat dari pernyataan itu: hukum kuat tidak memberi laju.
Ia mengatakan bahwa galat akhirnya kecil dan tinggal kecil, bukan seberapa cepat.

\begin{tangan}
Ambil satu koin dan selembar kertas bergaris. Angka bernilai satu, gambar
bernilai nol.

Lempar empat puluh kali, tulis 1 atau 0 berderet. Di baris kedua, tulis
rata-rata berjalan di bawah setiap lemparan: jumlah sampai situ dibagi
banyaknya lemparan sampai situ.

Yang harus kita lihat: baris kedua bergoyang paling liar pada lima lemparan
pertama, lalu makin tenang. Pada lemparan keempat puluh angkanya belum tentu
$0{,}5$; masih wajar $0{,}475$ atau $0{,}525$.

Sekarang ulangi percobaan itu empat kali dan gambar keempat jalur di satu
grafik dengan pensil berbeda. Keempat jalur menyempit ke $0{,}5$, dan lebar
sebarannya tinggal separuh setiap kali lemparannya dilipat empatkan.

Tempel kertasnya di dinding. Kalau ada yang berkata hasilnya pasti seimbang,
tunjuk lemparan keempat puluh, bukan yang kesepuluh.
\end{tangan}

\section{Teorema limit pusat}
\tingkat{lanjut}

Kalimat sehari-harinya: galat rata-rata, setelah dibesarkan dengan faktor
$\sqrt n$, selalu berbentuk lonceng yang sama, apa pun bentuk asli datanya.

Hukum bilangan besar memberi tahu bahwa $\bar X_n$ menempel ke $\mu$. Ia tidak
memberi tahu seberapa lebar tempelan itu. Teorema limit pusat menjawabnya, dan
jawabannya mengejutkan karena tidak bergantung pada distribusi asal.

\kamus{teorema limit pusat}{central limit theorem}{bentuk lonceng yang muncul pada
galat rata-rata, bukan pada datanya, dan hanya setelah galat itu dibesarkan
sebesar akar banyaknya data}

\begin{teorema}[Limit pusat, bentuk Lindeberg--Levy]\index{teorema limit pusat}
Misalkan $X_1,X_2,\dots$ bebas dan berdistribusi sama dengan
$\mathbb{E}[X_i]=\mu$ dan $0<\operatorname{Var}(X_i)=\sigma^2<\infty$. Maka
\[
Z_n=\frac{\sqrt n\,(\bar X_n-\mu)}{\sigma}
\;\rightsquigarrow\; Z\sim N(0,1),
\]
artinya $P(Z_n\le z)\to\Phi(z)$ untuk setiap $z\in\mathbb{R}$.
\end{teorema}

Perhatikan letak $\sqrt n$. Simpangan baku $\bar X_n$ adalah
$\sigma/\sqrt n$, bukan $\sigma/n$. Karena itu untuk memangkas galat menjadi
separuh kita perlu empat kali data, dan untuk menambah satu digit ketelitian kita
perlu seratus kali data. Ini kabar buruk yang tidak bisa diakali: ketelitian
datang dengan harga kuadratik.

Gambar~\ref{fig:lonceng-dari-kotak} memperlihatkan dari mana bentuk lonceng itu
datang. Panel kiri adalah sebaran satu dadu: datar, tanpa lonceng sama sekali.
Panel kanan adalah sebaran rata-rata dua belas dadu yang sudah distandarkan.
Bahannya bukan lonceng, hasilnya lonceng.

\begin{figure}[htbp]
\centering
\begin{tikzpicture}
\begin{axis}[
  name=satu, width=4.3cm, height=4.5cm,
  ybar, bar width=11pt,
  xmin=0.5, xmax=6.5, ymin=0, ymax=0.21,
  xtick={1,2,3,4,5,6}, ytick={0,0.0833,0.1667},
  yticklabels={$0$,$\frac{1}{12}$,$\frac16$},
  xlabel={mata dadu}, ylabel={peluang},
  title={\footnotesize satu dadu},
  label style={font=\footnotesize}, tick label style={font=\scriptsize},
  axis lines=left,
]
\addplot[fill=layerindigo!50, draw=layerindigo!80!black, line width=0.3pt]
  coordinates{
  (1,0.1667) (2,0.1667) (3,0.1667) (4,0.1667) (5,0.1667) (6,0.1667)};
\end{axis}
\begin{axis}[
  name=duabelas, at={(satu.east)}, anchor=west, xshift=7pt,
  width=9.3cm, height=4.5cm,
  ybar, bar width=3.4pt,
  xmin=-5.3, xmax=5.3, ymin=0, ymax=0.45,
  xtick={-4,-2,0,2,4}, ytick={0,0.2,0.4},
  xlabel={skor $z$ dari rata-rata}, ylabel={kerapatan},
  title={\footnotesize dua belas dadu, sudah distandarkan},
  label style={font=\footnotesize}, tick label style={font=\scriptsize},
  axis lines=left,
]
\addplot[fill=accent!45, draw=accent!70!black, line width=0.2pt]
  coordinates{
  (-4.3948,0) (-4.2258,0.00001) (-4.0567,0.00003) (-3.8877,0.00009) (-3.7187,0.0002) (-3.5496,0.00044)
  (-3.3806,0.00091) (-3.2116,0.00177) (-3.0426,0.00327) (-2.8735,0.00575) (-2.7045,0.00966) (-2.5355,0.01559)
  (-2.3664,0.02418) (-2.1974,0.03614) (-2.0284,0.05215) (-1.8593,0.07276) (-1.6903,0.09826) (-1.5213,0.1286)
  (-1.3522,0.16323) (-1.1832,0.20113) (-1.0142,0.24073) (-0.8452,0.28001) (-0.6761,0.31669) (-0.5071,0.34838)
  (-0.3381,0.37287) (-0.169,0.38835) (0,0.39365) (0.169,0.38835) (0.3381,0.37287) (0.5071,0.34838)
  (0.6761,0.31669) (0.8452,0.28001) (1.0142,0.24073) (1.1832,0.20113) (1.3522,0.16323) (1.5213,0.1286)
  (1.6903,0.09826) (1.8593,0.07276) (2.0284,0.05215) (2.1974,0.03614) (2.3664,0.02418) (2.5355,0.01559)
  (2.7045,0.00966) (2.8735,0.00575) (3.0426,0.00327) (3.2116,0.00177) (3.3806,0.00091) (3.5496,0.00044)
  (3.7187,0.0002) (3.8877,0.00009) (4.0567,0.00003) (4.2258,0.00001) (4.3948,0)};
\addplot[domain=-4.6:4.6, samples=120, layerindigo, line width=0.9pt]
  {0.39894228*exp(-x^2/2)};
\end{axis}
\end{tikzpicture}
\caption{Dari datar ke lonceng. Kiri: sebaran satu dadu, rata di enam nilai.
Kanan: sebaran rata-rata dua belas dadu setelah distandarkan, ditumpuk dengan
lengkung normal baku. Jarak supremum keduanya $0{,}03384$, kurang dari separuh
jarak pada $n=2$ yang $0{,}08333$. Bahan asalnya tidak lonceng, hasilnya
lonceng.}
\label{fig:lonceng-dari-kotak}
\end{figure}

Bukti berikut memakai fungsi pembangkit momen. Ia lengkap untuk kasus yang momen
ketiganya hingga, dan kita nyatakan batasnya secara jujur di akhir.

\begin{lemma}[Kekonvergenan lewat fungsi pembangkit momen]\index{fungsi pembangkit momen}
Misalkan $\psi_n(t)=\mathbb{E}[e^{tZ_n}]$ dan $\psi(t)=\mathbb{E}[e^{tZ}]$ hingga
pada suatu selang terbuka di sekitar $t=0$. Jika $\psi_n(t)\to\psi(t)$ untuk
setiap $t$ di selang itu, maka $Z_n\rightsquigarrow Z$.
\end{lemma}

Lemma itu adalah teorema kontinuitas untuk fungsi pembangkit momen. Buktinya
memakai kenyataan bahwa fungsi sebaran ditentukan oleh fungsi pembangkit momennya
pada selang di sekitar nol, dan kita tidak mengulanginya di sini; lihat Wasserman,
\emph{All of Statistics}, bab 5 bagian 5.7.2, Lemma 5.20.

\begin{proof}[Bukti teorema limit pusat]
Tetapkan $Y_i=(X_i-\mu)/\sigma$, sehingga $\mathbb{E}[Y_i]=0$,
$\operatorname{Var}(Y_i)=1$, dan $Z_n=n^{-1/2}\sum_{i=1}^n Y_i$. Misalkan
$\psi(t)=\mathbb{E}[e^{tY_1}]$ fungsi pembangkit momen $Y_1$, yang hingga pada
selang $|t|\le\delta$. Karena variabelnya bebas, fungsi pembangkit momen $Z_n$
adalah
\[
\xi_n(t)=\Bigl[\psi\bigl(t/\sqrt n\bigr)\Bigr]^n .
\]
Turunan di nol memberi $\psi(0)=1$, $\psi'(0)=\mathbb{E}[Y_1]=0$, dan
$\psi''(0)=\mathbb{E}[Y_1^2]=1$. Ekspansi Taylor orde dua dengan sisa orde tiga,
yang sah bila $\mathbb{E}|Y_1|^3<\infty$, berbunyi
\[
\psi(s)=1+\frac{s^2}{2}+R(s),\qquad |R(s)|\le C|s|^3 \quad (|s|\le\delta).
\]
Untuk $n$ cukup besar berlaku $|t|/\sqrt n\le\delta$, sehingga
\[
\xi_n(t)=\Bigl(1+\frac{t^2}{2n}+R\bigl(t/\sqrt n\bigr)\Bigr)^n
=\Bigl(1+\frac{t^2}{2n}+O\bigl(n^{-3/2}\bigr)\Bigr)^n .
\]
Ambil logaritma. Dengan $\log(1+a)=a-a^2/2+\dots$ dan
$a_n=t^2/(2n)+O(n^{-3/2})$,
\[
\log\xi_n(t)=n\Bigl(a_n-\frac{a_n^2}{2}+\dots\Bigr)
=\frac{t^2}{2}+O\bigl(n^{-1/2}\bigr)-\frac{t^4}{8n}+\dots
\longrightarrow \frac{t^2}{2}.
\]
Jadi $\xi_n(t)\to e^{t^2/2}$ untuk setiap $t$ tetap. Fungsi terakhir adalah fungsi
pembangkit momen $N(0,1)$, sehingga Lemma di atas memberi
$Z_n\rightsquigarrow N(0,1)$.
\end{proof}

Ada satu syarat tambahan yang menyelinap masuk: sisa Taylor orde tiga menuntut
$\mathbb{E}|Y_1|^3<\infty$. Untuk distribusi berekor sangat berat syarat itu bisa
gagal walaupun variansinya hingga, dan bukti di atas harus diganti. Jalan yang
umum adalah memakai fungsi karakteristik $\varphi(t)=\mathbb{E}[e^{itX}]$, yang
selalu ada karena $|e^{itX}|=1$, lalu memakai teorema kontinuitas Levy. Rangka
buktinya sama: logaritma fungsi karakteristik $Z_n$ menuju $-t^2/2$. Versi tanpa
asumsi momen ketiga ada di Wasserman, bab 5 bagian 5.7.2, dan di Blitzstein dan
Hwang, bab 10 bagian 10.3.

Teorema ini tidak mengatakan bahwa galatnya kecil. Ia mengatakan bahwa galat itu
berbentuk lonceng setelah dibagi $\sigma/\sqrt n$. Ukurannya tetap diberikan oleh
$\sigma/\sqrt n$, dan $\sigma$ adalah besaran yang biasanya tidak kita ketahui.
Dalam praktik $\sigma$ diganti oleh simpangan baku sampel $s$, dan penggantian itu
tidak merusak limit; pernyataannya menjadi $\sqrt n(\bar X_n-\mu)/s_n
\rightsquigarrow N(0,1)$. Besaran $\sigma/\sqrt n$ disebut galat baku
\index{galat baku} (\emph{standard error}), dan inilah dasar interval kepercayaan
di Bab 20.

Seberapa cepat aproksimasi itu membaik juga bisa dibatasi.

\begin{teorema}[Ketaksamaan Berry--Esseen]
Jika $\rho=\mathbb{E}|X_1-\mu|^3<\infty$, maka
\[
\sup_z\bigl|P(Z_n\le z)-\Phi(z)\bigr|\le \frac{C\rho}{\sigma^3\sqrt n},
\]
dengan $C$ konstanta universal; nilai $C=0{,}4748$ cukup untuk semua distribusi.
\end{teorema}

Bentuk ketaksamaan ini penting: jaraknya mengecil seperti $1/\sqrt n$, dengan
konstanta yang bergantung pada kemencengan distribusi asal lewat $\rho/\sigma^3$.
Distribusi yang sangat menceng akan butuh $n$ jauh lebih besar daripada
distribusi yang simetris. Angka pada bagian lab menunjukkan hal ini secara
kuantitatif untuk distribusi eksponensial, yang momen ketiganya $\rho=2{,}414553$
dan karenanya batasnya $1{,}1464/\sqrt n$.

\begin{resep}
Tiga langkah menormalkan rata-rata sampel. Urutannya tidak boleh ditukar.
\begin{enumerate}
\item \textbf{Kurangi pusat.} Hitung $\bar x-\mu$. Lonceng bergeser sehingga
      puncaknya jatuh di nol.
\item \textbf{Bagi lebar.} Bagi hasilnya dengan galat baku $\sigma/\sqrt n$.
      Bentuk lonceng tidak berubah; hanya skalanya yang menjadi satu satuan.
\item \textbf{Baca tabel.} Sisa angka $z$ dibaca pada $\Phi(z)$; untuk ekor
      kanan pakai $1-\Phi(z)$.
\end{enumerate}
Peringatan urutan: membagi sebelum mengurangi pusat menggeser hasil sebesar
$\mu\sqrt n/\sigma$. Pada dadu dengan $\mu=3{,}5$, $\sigma=1{,}707825$, dan
$n=12$, pergeseran itu $7{,}10$; angkanya sudah di luar tabel normal mana pun.
\end{resep}

\section{Aproksimasi normal dan koreksi kontinuitas}
\tingkat{menengah}

Kalimat sehari-harinya: untuk menghitung peluang jumlah, kita boleh memakai
tabel normal, tetapi bila jumlahnya hanya bisa bernilai bulat, ambangnya digeser
setengah.

Aproksimasi normal untuk jumlah bekerja paling sederhana pada distribusi diskret,
dan di situ muncul satu koreksi kecil yang dampaknya besar. Namanya koreksi
kontinuitas\index{koreksi kontinuitas}. Ambil $S\sim
\text{Binomial}(100;\,0{,}5)$, sehingga $\mu=50$ dan $\sigma=\sqrt{100\cdot0{,}5
\cdot0{,}5}=5$. Kejadian $\{S\le40\}$ adalah kejadian yang sama dengan
$\{S<40{,}5\}$, karena $S$ hanya bernilai bulat. Karena itu ambang yang benar
untuk distandarkan adalah $40{,}5$, bukan $40$.

\begin{center}
\begin{tabular}{lrrr}
\hline
Kejadian & Eksak & Tanpa koreksi & Dengan koreksi \\
\hline
$P(S\le40)$ & $0{,}028444$ & $0{,}022750$ & $0{,}028717$ \\
$P(S\le55)$ & $0{,}864373$ & $0{,}841345$ & $0{,}864334$ \\
$P(S=60)$   & $0{,}010844$ & $0{,}010798$ & $0{,}010852$ \\
\hline
\end{tabular}
\end{center}

Bacalah tabel itu baris demi baris. Pada baris pertama, aproksimasi tanpa koreksi
meleset $0{,}005694$, sedangkan dengan koreksi meleset $0{,}000273$, hampir dua
puluh satu kali lebih dekat. Pada baris ketiga keduanya sudah dekat, karena di
sekitar pusat sebaran lebar satu satuan sudah kecil dibandingkan $\sigma=5$.
Pelajaran umumnya: koreksi kontinuitas paling menolong di ekor dan ketika
$\sigma$ kecil. Besar pergeseran $z$ yang ditimbulkannya adalah $1/(2\sigma)$,
yaitu $0{,}1$ pada contoh ini.

\kertas{Hitung ulang baris pertama tabel ini dengan tangan. $z=(40-50)/5=-2$,
$\Phi(-2)=0{,}022750$; lalu $z=(40{,}5-50)/5=-1{,}9$, $\Phi(-1{,}9)=0{,}028717$.
Nilai eksaknya $\sum_{k=0}^{40}\binom{100}{k}2^{-100}=0{,}028444$. Periksa juga
bahwa galat tanpa koreksi $0{,}005694$ sekitar $21$ kali galat dengan koreksi
$0{,}000273$.}

Koreksi ini bukan rumus baru, hanya pembacaan ulang kejadian. Untuk variabel
diskret dengan jarak antar-nilai $h$, ambang digeser $h/2$ ke arah dalam. Untuk
variabel kontinu tidak ada yang perlu digeser, karena tidak ada lompatan.

Satu kebiasaan buruk perlu disebut di sini. Banyak buku menulis bahwa
aproksimasi normal "boleh dipakai bila $n\ge30$". Itu bukan teorema melainkan
kebiasaan. Angka pada bagian lab memperlihatkan mengapa: untuk rata-rata
eksponensial, jarak supremum ke normal pada $n=30$ adalah $0{,}02429$, dan
kuantil $95$ persen dari $Z_{30}$ adalah $1{,}7419$ sedangkan normal memberi
$1{,}6449$. Selisih $0{,}097$ itu berarti interval kepercayaan yang terlalu
sempit sekitar enam persen. Yang menentukan bukan $n$ saja, melainkan kemencengan
dan kurtosis distribusi asal.

\section{Bila teorema limit pusat menyesatkan: ekor berat}
\tingkat{lanjut}

Kalimat sehari-harinya: bila satu pengamatan sesekali bisa jauh lebih besar
daripada semua yang lain, rata-rata tidak lagi menetap, dan lonceng tidak muncul.

\kamus{ekor berat}{heavy tail}{peluang nilai raksasa yang tidak cukup cepat
mengecil, sehingga satu pengamatan dapat mengalahkan seluruh sisanya}

Semua teorema di atas memakai satu syarat yang mudah dilewati: variansi hingga.
Distribusi Cauchy\index{distribusi Cauchy}, dengan kepadatan
\[
f(x)=\frac{1}{\pi(1+x^2)},\qquad x\in\mathbb{R},
\]
tidak memilikinya. Ekornya mengecil seperti $1/x$, terlalu lambat. Peluang
melewati sepuluh adalah $1-(2/\pi)\arctan 10=0{,}063451$, dan peluang melewati
seratus masih $0{,}006366$. Bandingkan dengan normal baku, yang peluang melewati
sepuluh praktis nol. Bahkan $\mathbb{E}|X|$ tidak hingga: harapan terpangkas pada
ambang $M$ bernilai $(2/\pi)\ln(1+M^2)$, yaitu $2{,}9381$ untuk $M=10$,
$8{,}7952$ untuk $M=10^3$, dan $17{,}5905$ untuk $M=10^6$. Angka itu tumbuh
seperti logaritma, jadi tidak ada batas.

Akibatnya tegas, dan bukan sekadar "CLT kurang akurat".

\begin{proposisi}
Jika $X_1,\dots,X_n$ iid Cauchy, maka $\bar X_n$ berdistribusi Cauchy untuk
setiap $n$.
\end{proposisi}

\begin{proof}[Bukti]
Fungsi karakteristik Cauchy adalah $\varphi(t)=e^{-|t|}$. Karena variabelnya
bebas, fungsi karakteristik $\bar X_n$ adalah
$\bigl[e^{-|t|/n}\bigr]^n=e^{-|t|}$, yaitu fungsi karakteristik Cauchy lagi.
Karena fungsi karakteristik menentukan sebaran, $\bar X_n$ berdistribusi Cauchy.
\end{proof}

Jadi rata-rata $10^5$ pengamatan Cauchy sama tidak informatifnya dengan satu
pengamatan. Hukum bilangan besar gagal, teorema limit pusat gagal, dan "galat
baku" tidak berarti apa pun. Angka pada bagian lab menunjukkan tiga gejala itu
sekaligus. Median $|\bar X_n|$ tinggal di sekitar satu untuk semua $n$
($1{,}004$, $1{,}008$, $1{,}040$, $0{,}973$), padahal untuk distribusi
bervariansi hingga ia harus mengecil seperti $1/\sqrt n$. Peluang
$|\bar X_n|>10$ tinggal di sekitar $0{,}0635$, nilai teoretisnya, untuk semua $n$.
Simpangan baku empirisnya melompat-lompat: $33{,}04$, $863{,}61$, $23{,}87$,
$379{,}53$ untuk $n=10$ sampai $10^4$. Lompatan itu bukan kesalahan program.
Simpangan baku sampel dari data Cauchy sendiri adalah variabel acak yang tidak
punya nilai khas.

Gambar~\ref{fig:cauchy-berjalan} memperlihatkan bedanya pada satu jalankan.
Rata-rata berjalan data seragam menempel ke $0{,}5$ sesudah beberapa puluh
tarikan. Rata-rata berjalan data Cauchy melompat dari $-0{,}0866$ ke $4{,}4572$
pada tarikan ke-20, karena satu nilai $90{,}79$ kebetulan masuk, lalu mengembara
lagi. Satu nilai raksasa cukup untuk membatalkan seluruh riwayat sebelumnya.

\begin{figure}[htbp]
\centering
\begin{tikzpicture}
\begin{axis}[
  width=0.95\linewidth, height=6.0cm,
  xmode=log, xmin=1, xmax=4000, ymin=-1.1, ymax=4.9,
  xlabel={banyak tarikan $n$ (skala logaritma)}, ylabel={rata-rata berjalan},
  label style={font=\small}, tick label style={font=\footnotesize},
  axis lines=left,
  legend style={font=\footnotesize, draw=layergrey!35, fill=white,
    fill opacity=0.9, text opacity=1, at={(0.015,0.985)},
    anchor=north west, cells={anchor=west}},
  legend cell align=left,
]
\addplot[accent, line width=0.7pt] coordinates{
  (1,-0.253) (2,1.045) (3,0.582) (4,0.564) (5,0.553) (6,1.064)
  (7,-0.579) (8,-0.517) (9,-0.409) (10,-0.325) (11,-0.409) (12,-0.361)
  (13,-0.344) (14,-0.262) (15,0.008) (16,-0.038) (17,0.017) (18,-0.08)
  (19,-0.087) (20,4.457) (21,4.29) (22,4.153) (23,4.046) (24,3.762)
  (25,3.656) (26,3.551) (27,3.425) (28,3.257) (29,3.186) (30,3.083)
  (31,2.934) (32,2.786) (33,2.837) (34,2.743) (35,2.633) (36,2.528)
  (37,2.464) (38,2.381) (39,2.279) (40,2.218) (41,2.157) (44,1.732)
  (47,1.387) (50,1.367) (53,1.191) (57,1.113) (61,1.077) (65,0.942)
  (70,0.796) (75,0.773) (80,0.725) (85,0.492) (91,0.79) (97,0.655)
  (104,0.54) (111,0.478) (119,0.494) (127,0.401) (135,0.981) (145,0.868)
  (155,2.5) (165,2.351) (177,2.207) (189,2.168) (202,1.935) (216,1.854)
  (230,1.776) (246,1.665) (263,1.573) (281,0.848) (300,0.794) (321,0.791)
  (343,0.511) (367,0.922) (392,0.986) (419,0.984) (447,1.03) (478,0.951)
  (511,0.827) (546,0.685) (583,0.37) (623,0.548) (666,0.953) (712,1.021)
  (761,0.945) (813,2.066) (869,2.135) (929,1.855) (992,1.814) (1060,1.649)
  (1133,1.31) (1211,0.557) (1294,0.422) (1383,0.434) (1478,0.466) (1579,0.324)
  (1688,0.406) (1803,0.386) (1927,0.444) (2059,0.403) (2201,0.585) (2352,0.419)
  (2513,0.528) (2686,0.397) (2870,0.356) (3067,0.326) (3278,0.276) (3503,-0.096)
  (3743,-0.159) (4000,-0.203)};
\addlegendentry{Cauchy}
\addplot[layerteal, dashed, line width=1.0pt] coordinates{
  (1,0.221) (2,0.272) (3,0.323) (4,0.472) (5,0.462) (6,0.519)
  (7,0.524) (8,0.573) (9,0.55) (10,0.574) (11,0.605) (12,0.601)
  (13,0.57) (14,0.581) (15,0.543) (16,0.532) (17,0.529) (18,0.528)
  (19,0.532) (20,0.544) (21,0.545) (22,0.547) (23,0.539) (24,0.531)
  (25,0.54) (26,0.555) (27,0.561) (28,0.541) (29,0.545) (30,0.547)
  (31,0.562) (32,0.549) (33,0.533) (34,0.542) (35,0.531) (36,0.535)
  (37,0.537) (38,0.526) (39,0.519) (40,0.516) (41,0.517) (44,0.518)
  (47,0.519) (50,0.503) (53,0.518) (57,0.515) (61,0.525) (65,0.53)
  (70,0.529) (75,0.53) (80,0.535) (85,0.538) (91,0.539) (97,0.538)
  (104,0.541) (111,0.535) (119,0.532) (127,0.533) (135,0.528) (145,0.526)
  (155,0.518) (165,0.523) (177,0.513) (189,0.516) (202,0.513) (216,0.507)
  (230,0.5) (246,0.498) (263,0.495) (281,0.496) (300,0.496) (321,0.495)
  (343,0.499) (367,0.499) (392,0.501) (419,0.503) (447,0.497) (478,0.505)
  (511,0.508) (546,0.509) (583,0.504) (623,0.502) (666,0.497) (712,0.494)
  (761,0.488) (813,0.484) (869,0.485) (929,0.484) (992,0.485) (1060,0.484)
  (1133,0.483) (1211,0.485) (1294,0.487) (1383,0.489) (1478,0.488) (1579,0.491)
  (1688,0.489) (1803,0.489) (1927,0.49) (2059,0.491) (2201,0.49) (2352,0.492)
  (2513,0.491) (2686,0.493) (2870,0.492) (3067,0.49) (3278,0.492) (3503,0.491)
  (3743,0.493) (4000,0.493)};
\addlegendentry{seragam $(0,1)$}
\draw[layergrey, densely dotted, line width=0.6pt]
  (axis cs:1,0.5) -- (axis cs:4000,0.5);
\node[font=\scriptsize, layergrey, anchor=south west] at (axis cs:1.3,0.53) {$0{,}5$};
\end{axis}
\end{tikzpicture}
\caption{Rata-rata berjalan dari 4000 tarikan. Garis putus-putus: seragam pada
$(0,1)$, menetap dekat $0{,}5$ dan berakhir di $0{,}4930$. Garis penuh: Cauchy,
melompat dari $-0{,}0866$ ke $4{,}4572$ pada tarikan ke-20 ketika satu nilai
$90{,}79$ masuk, lalu mengembara lagi dan berakhir di $-0{,}203$. Lebar sumbu
mendatar logaritma, supaya lompatan awal tidak tertutup.}
\label{fig:cauchy-berjalan}
\end{figure}

Yang menggantikan teorema limit pusat adalah keluarga sebaran stabil
\index{sebaran stabil}. Untuk
variabel iid dengan ekor yang mengecil seperti $x^{-\alpha}$ dengan
$0<\alpha<2$, ada konstanta $a_n$ sehingga
\[
\frac{S_n-a_n}{n^{1/\alpha}}\;\rightsquigarrow\;Y_\alpha,
\]
dengan $Y_\alpha$ sebaran stabil berindeks $\alpha$. Untuk $\alpha=1$ kita
mendapat Cauchy dan $n^{1/1}=n$, persis seperti proposisi di atas. Untuk
$\alpha\in(1,2)$ rata-rata masih konvergen ke $\mu$, tetapi galatnya mengecil
seperti $n^{1-1/\alpha}$, lebih lambat daripada $1/\sqrt n$, dan bentuknya bukan
lonceng. Pernyataan umum ini adalah hasil Gnedenko dan Kolmogorov; rinciannya ada
pada Feller jilid 2, bab tentang sebaran stabil.

Bila $\alpha>2$ tetapi momen ketiga tidak ada, teorema limit pusat tetap benar
karena variansinya hingga. Yang hilang adalah jaminan lajunya: batas
Berry--Esseen menjadi tak berhingga, sehingga tidak ada lagi pegangan bahwa galat
mengecil seperti $1/\sqrt n$ dengan konstanta yang terkendali. Dalam praktik ini
yang paling sering menggigit: sebarannya tidak ekstrem seperti Cauchy, tetapi
cukup berekor untuk membuat interval kepercayaan terlalu sempit pada ukuran
sampel yang realistis.

\begin{tantangan}
Distribusi tanpa rata-rata. Ambil $X$ dengan
$P(X=2^k)=P(X=-2^k)=2^{-k-1}$ untuk $k=1,2,3,\dots$

Tunjukkan dua hal. Pertama, $\mathbb{E}|X|=\infty$, sehingga tidak ada $\mu$
yang boleh dituju. Kedua, rata-rata sampelnya tidak konvergen dalam peluang ke
konstanta mana pun: hukum bilangan besar gagal, bukan sekadar melambat.

Petunjuk 1: yang menentukan bukan rata-rata semua suku, melainkan suku
terbesar. Tanyakan seberapa sering satu suku melewati ambang $n$.

Petunjuk 2: hitung $n\,P(|X|>n)$ untuk $n$ di antara dua pangkat dua, lalu
pakai syarat perlu $n\,P(|X_1|>n)\to0$ yang disebut pada bagian hukum kuat.

Kalau mentok, kembali ke bagian tentang ekor berat dan ke paragraf syarat perlu
dan cukup di bagian hukum kuat.
\end{tantangan}

\begin{untukanak}
Di sebuah desa ada seratus orang dengan penghasilan biasa-biasa saja, dan satu
orang yang sangat kaya. Kalau kita menghitung penghasilan rata-rata desa itu,
angkanya ditentukan hampir sepenuhnya oleh orang kaya tadi. Kalau orang kaya itu
kebetulan sedang pergi, rata-ratanya jatuh drastis. Kalau ia pulang, rata-ratanya
melompat naik.

Sekarang tambahkan penduduk. Seratus orang lagi, seribu orang lagi; selama
pendatangnya orang biasa, rata-ratanya berubah sedikit. Tetapi kalau di antara
pendatang itu ada satu orang raksasa lagi, rata-ratanya melompat lagi. Menambah
banyak orang tidak menolong, karena yang menentukan bukan banyaknya orang,
melainkan siapa yang terbesar.

Inilah yang terjadi pada data yang ekornya berat: satu pengamatan bisa
mengalahkan semua yang lain. Rata-rata menjadi tidak stabil, dan lonceng yang
biasa dipakai untuk menghitung "seberapa yakin" tidak muncul sama sekali.

Pertanyaan yang membuat penasaran: bagaimana kita bisa tahu, sebelum menghitung,
apakah data kita termasuk jenis yang tenang atau jenis yang dikuasai raksasa? Di
bab ini kita akan tahu mengapa pertanyaan itu tidak bisa dijawab hanya dengan
melihat rata-ratanya.
\end{untukanak}

\section{Lab komputasi: lima pemeriksaan}
\tingkat{lanjut}

Lima klaim di bab ini dapat diperiksa dengan satu skrip: laju $1/\sqrt n$ pada
hukum bilangan besar, kemencengan dan kurtosis rata-rata eksponensial, jarak
supremum ke normal, galat baku Monte Carlo, dan kegagalan total pada Cauchy.
Semua angka di bawah ini adalah keluaran nyata dari eksekusi dengan NumPy 2.5.2
dan SciPy 1.18.0, pembangkit \texttt{default\_rng(20261024)}.

\begin{lab}
\begin{verbatim}
import numpy as np
from scipy import stats

rng = np.random.default_rng(20261024)

print("== 1. hukum bilangan besar: rata-rata dadu ==")
for k in range(1, 6):
    n = 10**k
    m = rng.integers(1, 7, size=(2000, n)).mean(axis=1)
    print("n=%6d  rata2=%.4f  rms=%.4f  sd/akar n=%.4f"
          % (n, m.mean(), np.sqrt(np.mean((m-3.5)**2)),
             np.sqrt(35/12/n)))

print("== 2. CLT: rata-rata eksponensial ==")
for n in (2, 5, 10, 30, 100, 500):
    z = (rng.exponential(size=(50000, n)).mean(axis=1)
         - 1)*np.sqrt(n)
    ks = stats.kstest(z, "norm")
    print("n=%4d skew=%+.3f (%+.3f) kurt=%+.3f (%+.3f)"
          % (n, stats.skew(z), 2/np.sqrt(n),
             stats.kurtosis(z), 6.0/n))
    print("      KS D=%.4f p=%.3f" % (ks.statistic, ks.pvalue))

print("== 3. jarak sup ke normal, eksak lewat gamma ==")
rho = 2 + 2*(6/np.e - 2)          # E|X-1|^3 untuk X ~ Exp(1)
for n in (2, 5, 10, 30, 100, 500):
    zg = np.linspace(-4, 4, 40001)
    Fn = stats.gamma.cdf(n*(1+zg/np.sqrt(n)), a=n)
    D = np.max(np.abs(Fn - stats.norm.cdf(zg)))
    print("n=%4d sup|Fn-Phi|=%.5f  D*akar n=%.3f  BE=%.3f"
          % (n, D, D*np.sqrt(n), 0.4748*rho/np.sqrt(n)))

print("== 4. galat baku Monte Carlo untuk E[exp(U)] ==")
for k in (2, 3, 4, 5):
    n = 10**k
    est = np.exp(rng.random(size=(2000, n))).mean(axis=1)
    print("n=%6d  sd=%.5f  sd*akar n=%.5f"
          % (n, est.std(), est.std()*np.sqrt(n)))

print("== 5. Cauchy: rata-rata tidak memusat ==")
for k in (1, 2, 3, 4):
    n = 10**k
    m = rng.standard_cauchy(size=(4000, n)).mean(axis=1)
    print("n=%5d  median|m|=%.3f P(|m|>10)=%.4f sd=%.2f"
          % (n, np.median(np.abs(m)), np.mean(np.abs(m) > 10),
             m.std()))
print("teori P(|rata2|>10) = %.4f"
      % (1 - 2/np.pi*np.arctan(10)))
\end{verbatim}
\end{lab}

Keluarannya, apa adanya:

\begin{verbatim}
== 1. hukum bilangan besar: rata-rata dadu ==
n=    10  rata2=3.5014  rms=0.5440  sd/akar n=0.5401
n=   100  rata2=3.4992  rms=0.1698  sd/akar n=0.1708
n=  1000  rata2=3.4986  rms=0.0527  sd/akar n=0.0540
n= 10000  rata2=3.5001  rms=0.0171  sd/akar n=0.0171
n=100000  rata2=3.4997  rms=0.0053  sd/akar n=0.0054
== 2. CLT: rata-rata eksponensial ==
n=   2 skew=+1.440 (+1.414) kurt=+3.130 (+3.000)
      KS D=0.0929 p=0.000
n=   5 skew=+0.919 (+0.894) kurt=+1.393 (+1.200)
      KS D=0.0602 p=0.000
n=  10 skew=+0.631 (+0.632) kurt=+0.589 (+0.600)
      KS D=0.0480 p=0.000
n=  30 skew=+0.355 (+0.365) kurt=+0.181 (+0.200)
      KS D=0.0219 p=0.000
n= 100 skew=+0.196 (+0.200) kurt=+0.038 (+0.060)
      KS D=0.0167 p=0.000
n= 500 skew=+0.112 (+0.089) kurt=-0.003 (+0.012)
      KS D=0.0062 p=0.043
== 3. jarak sup ke normal, eksak lewat gamma ==
n=   2 sup|Fn-Phi|=0.09446  D*akar n=0.134  BE=0.811
n=   5 sup|Fn-Phi|=0.05963  D*akar n=0.133  BE=0.513
n=  10 sup|Fn-Phi|=0.04211  D*akar n=0.133  BE=0.363
n=  30 sup|Fn-Phi|=0.02429  D*akar n=0.133  BE=0.209
n= 100 sup|Fn-Phi|=0.01330  D*akar n=0.133  BE=0.115
n= 500 sup|Fn-Phi|=0.00595  D*akar n=0.133  BE=0.051
== 4. galat baku Monte Carlo untuk E[exp(U)] ==
n=   100  sd=0.04928  sd*akar n=0.49281
n=  1000  sd=0.01586  sd*akar n=0.50156
n= 10000  sd=0.00495  sd*akar n=0.49538
n=100000  sd=0.00155  sd*akar n=0.49056
== 5. Cauchy: rata-rata tidak memusat ==
n=   10  median|m|=1.004 P(|m|>10)=0.0630 sd=33.04
n=  100  median|m|=1.008 P(|m|>10)=0.0680 sd=863.61
n= 1000  median|m|=1.040 P(|m|>10)=0.0653 sd=23.87
n=10000  median|m|=0.973 P(|m|>10)=0.0580 sd=379.53
teori P(|rata2|>10) = 0.0635
\end{verbatim}

Bagian 1 memperlihatkan hukum bilangan besar bekerja, dan memperlihatkan
harganya. Kolom rms cocok dengan $\sigma/\sqrt n$ sampai dua angka di belakang
koma untuk $n\ge100$: $0{,}1698$ lawan $0{,}1708$, $0{,}0171$ lawan $0{,}0171$,
$0{,}0053$ lawan $0{,}0054$. Tetapi lihat besarnya: pada $n=10^5$ rata-ratanya
masih meleset sekitar $0{,}005$, yaitu $0{,}15$ persen dari $3{,}5$. Rata-rata
$10^5$ lemparan bukan "praktis tepat"; ia tepat sampai dua angka di belakang
koma.

Bagian 2 dan 3 memperlihatkan dua cara mengukur kedekatan ke normal. Kemencengan
dan kurtosis mengikuti ramalan teori dengan rapi mulai $n=10$: $0{,}631$ lawan
$2/\sqrt{10}=0{,}632$ dan $0{,}589$ lawan $6/10=0{,}600$. Pada $n=2$ dan $n=5$
penaksir kurtosis masih berisik, karena simpangannya bergantung pada momen
kedelapan distribusi itu, bukan hanya pada banyaknya ulangan. Bagian 3 menghitung
jarak supremum secara eksak, tanpa simulasi, dengan memakai sebaran Gamma untuk
rata-rata eksponensial: $0{,}04211$ pada $n=10$ dan $0{,}00595$ pada $n=500$.
Kolom $D\sqrt n$ hampir tetap di $0{,}133$, bukti langsung bahwa lajunya
$1/\sqrt n$. Batas Berry--Esseen $1{,}1464/\sqrt n$ berlaku, dengan angka
$0{,}209$ lawan $0{,}02429$ pada $n=30$; batas itu sekitar $8{,}6$ kali lebih
longgar daripada kenyataan, dan itu wajar untuk batas yang seragam atas semua
distribusi.

Uji Kolmogorov--Smirnov pada lima puluh ribu tarikan menolak kenormalan dengan
$p=0{,}000$ untuk semua $n$ sampai $100$, lalu memberi $p=0{,}043$ pada $n=500$.
Itu bukan tanda program salah: aproksimasi normal tidak pernah tepat untuk $n$
berhingga, dan jaraknya hanya mengecil sampai uji sebesar itu tidak mampu
melihatnya lagi. Pertanyaan yang benar bukan "apakah normal" melainkan
"seberapa besar galat yang masih boleh kita terima", dan bagian 3 menjawabnya
dengan angka.

Bagian 4 memisahkan dua besaran yang sering tertukar. Simpangan baku penaksir
turun dari $0{,}04928$ ke $0{,}00155$ ketika $n$ naik seribu kali, sedangkan
$\text{sd}\cdot\sqrt n$ hampir tetap di $0{,}49$. Nilai tetap itu adalah
$\sigma$ dari $\exp(U)$, yaitu $\sqrt{(e^2-1)/2-(e-1)^2}=0{,}491971$. Inilah
hukum $1/\sqrt n$ dalam bentuk paling telanjang: untuk ketelitian $0{,}001$ pada
$95$ persen keyakinan kita perlu $(1{,}96\cdot0{,}491971/0{,}001)^2=929803$
tarikan.

Bagian 5 adalah tandingannya. Median $|\bar X_n|$ tidak mengecil; peluang
melewati sepuluh tidak mengecil; dan simpangan baku empirisnya melompat dari
$33{,}04$ ke $863{,}61$ lalu turun ke $23{,}87$. Teori menyebut $0{,}0635$, dan
kelima angka simulasi berkisar di situ. Tidak ada hukum bilangan besar di sini,
dan tidak ada teorema limit pusat. Satu-satunya cara menyelamatkan analisis adalah
mengubah pertanyaannya: menaksir indeks ekor, bukan rata-rata.

\begin{untukdosen}
Urutan mengajar yang saya sarankan: mulai dari tabel simulasi, bukan dari definisi.
Tunjukkan pada kelas bahwa galat satu jalankan tidak mengikuti pola apa pun,
sedangkan sebaran galatnya mengikuti. Sesudah itu tiga modus konvergensi, dengan
contoh penyangkal untuk setiap arah yang gagal; jangan menunda contoh-contoh itu,
karena di situlah mahasiswa paham bahwa "menuju" perlu diperjelas. Baru kemudian
Markov, Chebyshev, hukum lemah, hukum kuat sebagai pernyataan, lalu limit pusat
dengan bukti fungsi pembangkit momen.

Dua miskonsepsi yang paling sering muncul. Pertama, mahasiswa membaca "konvergen
dalam peluang" sebagai "nilainya menjadi tetap", lalu heran ketika simulasi
menunjukkan lompatan besar pada $n$ besar; contoh $(0,1)$ yang dipecah pada bab
ini adalah penawarnya. Kedua, mereka membaca teorema limit pusat sebagai "rata-rata
menjadi normal", padahal yang menjadi normal adalah galat yang dibagi
$\sigma/\sqrt n$; akibatnya mereka menuliskan $\bar X_n\approx N(\mu,\sigma^2)$
dan lupa $\sqrt n$.

Aktivitas papan tulis yang berhasil: minta dua mahasiswa memerankan penjudi dan
bandar. Penjudi menyebut toleransi, bandar menyebut banyak lemparan yang
menjaminnya, dan kelas menghitung dengan Chebyshev lalu membandingkannya dengan
angka eksak dari SciPy. Selisihnya selalu besar, dan itulah pintu masuk ke
aproksimasi normal. Tanda kelas siap lanjut: mereka dapat menjelaskan mengapa
Chebyshev tetap berguna walaupun longgar, dan dapat menyebut satu distribusi yang
membuat teorema limit pusat gagal.
\end{untukdosen}

\section{Ringkasan}

\begin{ringkasan}
\begin{itemize}
  \item Ada tiga modus konvergensi yang berbeda: hampir pasti, dalam peluang, dan
        dalam distribusi. Hampir pasti mengakibatkan dalam peluang, dalam peluang
        mengakibatkan dalam distribusi, dan arah sebaliknya gagal kecuali bila
        limitnya konstanta.
  \item Ketaksamaan Markov dan Chebyshev hanya memakai momen, karena itu berlaku
        untuk semua distribusi bervariansi hingga, dan karena itu pula longgar.
  \item Hukum bilangan besar lemah: $P(|\bar X_n-\mu|\ge\varepsilon)\le
        \sigma^2/(n\varepsilon^2)\to0$. Hukum kuat menambahkan bahwa
        $\bar X_n\to\mu$ hampir pasti bila $\mathbb{E}|X_1|<\infty$, tanpa
        menuntut variansi.
  \item Teorema limit pusat: $\sqrt n(\bar X_n-\mu)/\sigma\rightsquigarrow N(0,1)$
        bila variansi hingga. Lajunya $1/\sqrt n$, maka empat kali data untuk
        separuh galat, dan seratus kali data untuk satu digit tambahan.
  \item Aproksimasi normal untuk jumlah diskret memerlukan koreksi kontinuitas
        setengah satuan; koreksi itu memangkas galat di ekor sekitar dua puluh
        kali pada contoh binomial di bab ini.
  \item Bila variansi tak hingga, teorema limit pusat gagal, bukan hanya melambat.
        Rata-rata $n$ pengamatan Cauchy tetap Cauchy, dan yang menggantikan
        lonceng adalah sebaran stabil berindeks $\alpha<2$.
\end{itemize}
\end{ringkasan}

\section{Latihan}

\begin{latihan}
\begin{enumerate}
  \item Buktikan: bila $X_n\to X$ dan $Y_n\to Y$ dalam peluang, maka
        $X_n+Y_n\to X+Y$ dalam peluang. Tuliskan pertidaksamaan himpunan yang
        dipakai, lalu jelaskan mengapa $\varepsilon/2$ muncul.
  \item Buktikan bahwa ketaksamaan Chebyshev tajam: untuk $0<\sigma^2\le a^2$,
        ambil $X$ dengan $P(X=a)=P(X=-a)=\sigma^2/(2a^2)$ dan
        $P(X=0)=1-\sigma^2/a^2$, lalu hitung $\operatorname{Var}(X)$ dan
        $P(|X|\ge a)$.
  \item Nilai klaim ini: "Karena teorema limit pusat berlaku untuk semua
        distribusi, rata-rata tiga puluh pengamatan selalu boleh dianggap
        normal." Benar, salah, atau belum cukup bukti? Sebutkan syarat yang
        dilanggar atau yang belum diperiksa, dan pakai angka dari bagian lab.
  \item Untuk koin dengan $p=1/2$ dan toleransi $\varepsilon=0{,}05$, cari $n$
        terkecil yang dijamin Chebyshev untuk peluang meleset paling besar
        $0{,}05$. Lalu cari $n$ terkecil yang benar-benar memenuhi syarat itu,
        dengan menghitung sebaran binomial di Python. Bandingkan kedua angka.
  \item Buktikan bahwa konvergensi dalam peluang tidak menjamin konvergensi
        harapan, dengan memakai $P(X_n=n^2)=1/n$ dan $P(X_n=0)=1-1/n$. Hitung
        $\mathbb{E}[X_n]$ dan $P(|X_n|>\varepsilon)$ untuk $\varepsilon$ tetap.
  \item Untuk $S\sim\text{Binomial}(200;\,0{,}3)$, aproksimasi $P(S\le50)$ dengan
        normal tanpa dan dengan koreksi kontinuitas, lalu bandingkan keduanya
        dengan nilai eksak dari SciPy. Laporkan kedua galat.
  \item Bangkitkan $4000$ rata-rata dari $n$ bilangan Cauchy untuk
        $n=10,10^2,10^3,10^4$, lalu laporkan median $|\bar X_n|$, peluang
        $|\bar X_n|>10$, dan simpangan bakunya. Bandingkan dengan
        $1-(2/\pi)\arctan10$ dan jelaskan mengapa kolom simpangan baku tidak
        mengecil.
  \item Buktikan dengan fungsi karakteristik bahwa rata-rata $n$ variabel iid
        Cauchy berdistribusi Cauchy. Lalu jelaskan mengapa laporan berbentuk
        "rata-rata $\pm$ dua galat baku" pada data semacam itu tidak dapat
        ditafsirkan.
  \item Berapa tarikan Monte Carlo yang diperlukan agar penaksir
        $\mathbb{E}[\exp(U)]$ memiliki setengah lebar interval $95$ persen
        sebesar $0{,}0005$? Pakai $\sigma=0{,}491971$ dari bagian lab, lalu
        bandingkan dengan $n$ yang diperlukan untuk setengah lebar $0{,}001$.
        Nyatakan rasionya dan jelaskan asalnya.
  \item Buktikan bahwa bila $X_n\rightsquigarrow c$ dengan $c$ konstanta, maka
        $X_n\to c$ dalam peluang, tanpa memakai bagian (c) teorema hubungan tiga
        modus di awal bab, lalu tunjukkan dengan contoh bahwa syarat "konstanta"
        tidak dapat dibuang.
\end{enumerate}

Petunjuk untuk tiga butir tersulit. Butir 1: tulis
$\{|(X_n+Y_n)-(X+Y)|>\varepsilon\}$ termuat dalam gabungan dua kejadian, lalu
pakai batas gabungan. Butir 8: fungsi karakteristik jumlah adalah hasilkali
fungsi karakteristik, dan fungsi karakteristik Cauchy $e^{-|t|}$ bersifat
istimewa karena bentuknya tidak berubah saat dipangkatkan. Butir 10: tulis
$P(X_n\le t)\le P(X_n\le c-\varepsilon)+P(|X_n-c|>\varepsilon)$ untuk $t$ di
kiri $c$, dan ingat bahwa sebaran titik massa melompat tepat di $c$.
\end{latihan}

\begin{untukpeneliti}
Yang benar-benar terbukti di bab ini: implikasi hampir pasti $\Rightarrow$ dalam
peluang $\Rightarrow$ dalam distribusi, hukum lemah lewat Chebyshev, dan limit
pusat lewat fungsi pembangkit momen dengan asumsi $\mathbb{E}|Y_1|^3<\infty$.
Yang populer tetapi tidak terbukti di sini: hukum kuat (pernyataan saja, bukti
Etemadi dilewati), versi limit pusat tanpa momen ketiga (butuh fungsi
karakteristik dan teorema kontinuitas Levy), dan seluruh teori sebaran stabil
(pernyataan Gnedenko--Kolmogorov, tanpa bukti).

Sumber primer yang saya pakai: Wasserman, \emph{All of Statistics}, bab 5 bagian
5.2 (tipe konvergensi dan diagram implikasinya), 5.3 (hukum lemah), 5.4 (limit
pusat dan Berry--Esseen), serta lampiran 5.7.1 dan 5.7.2 untuk hukum kuat dan
bukti limit pusat. Blitzstein dan Hwang, \emph{Introduction to Probability}, bab
10 bagian 10.1 (ketaksamaan), 10.2 (hukum bilangan besar), 10.3 (limit pusat).
Untuk sisi ekor berat: Feller jilid 2.

Batas keberlakuan yang perlu disebut. Pertama, semua teorema di sini mengandaikan
kebebasan dan distribusi sama; untuk data dependen, misalnya deret waktu atau
iterasi rantai Markov, yang berlaku adalah versi martingale atau mixing, dan
lajunya bisa jauh lebih lambat. Kedua, konstanta optimal pada Berry--Esseen belum
diketahui; $0{,}4748$ cukup, dan batas bawah yang terbaik bagi konstanta itu
masih di bawah $0{,}5$. Ketiga, laju $1/\sqrt n$ adalah laju paling buruk yang
dijamin secara seragam; untuk distribusi tertentu laju sebenarnya lebih baik,
dan untuk ekor berat dengan $\alpha\in(1,2)$ laju sebenarnya lebih buruk.
Pertanyaan terbuka yang masih hidup: seberapa tajam batas Berry--Esseen untuk
kelas distribusi yang dibatasi momennya, dan bagaimana menaksir indeks ekor
dengan interval kepercayaan yang sah pada $n$ kecil.
\end{untukpeneliti}

\peta{Bab 19 pada Wasserman, \emph{All of Statistics}: bab 5 (Convergence of
Random Variables), bagian 5.2 Types of Convergence, 5.3 The Law of Large Numbers,
5.4 The Central Limit Theorem, dan lampiran 5.7.1 Almost Sure and L1 Convergence
serta 5.7.2 Proof of the Central Limit Theorem. Pada Blitzstein dan Hwang,
\emph{Introduction to Probability}: bab 10 (Inequalities and limit theorems),
bagian 10.1 Inequalities, 10.2 Law of large numbers, 10.3 Central limit theorem.
Untuk sebaran stabil: Feller, \emph{An Introduction to Probability Theory and Its
Applications}, jilid 2. Distribusi dasar yang dipakai di sini dibangun di Bab 17,
dan interval kepercayaan yang bersandar pada bab ini dibahas di Bab 19.}

\chapter{Estimasi, Interval Kepercayaan, dan Uji Hipotesis}

Sebuah tim menguji pendeteksi serangan pada $50$ contoh berlabel, dan pendeteksi itu
menangkap $33$ di antaranya. Laporan yang ditulis malam itu berbunyi: recall $66$
persen. Pendeteksi kedua, diuji pada $100$ contoh, menangkap $53$: recall $53$ persen.
Rapat pagi memutuskan memakai yang pertama, karena $66$ lebih besar daripada $53$, dan
keputusan itu selesai dalam satu menit.

Aritmetikanya benar: $33/50=0{,}66$. Yang keliru adalah anggapan bahwa angka itu
menunjuk satu bilangan. Bila lima puluh contoh itu diganti lima puluh contoh lain dari
kumpulan yang sama, hasilnya bukan $33$ lagi; ia bergoyang, dan goyangannya tidak
kecil. Sesudah dua jam menulis program, tim itu sebenarnya memegang dua kalimat yang
berbeda jauh kekuatannya. Kalimat pertama, ``pada data ini terlihat $66$ persen'',
hampir tidak berisi apa pun. Kalimat kedua, ``laju tangkapan sebenarnya ada di suatu
tempat dalam selang ini'', berisi sesuatu yang bisa dipakai, dengan syarat kita tahu
seberapa lebar selang itu dan apa arti lebar tersebut.

Bab 19 sudah memberi dua alat: hukum bilangan besar, yang mengatakan rata-rata sampel
memusat ke nilai sebenarnya, dan teorema limit pusat, yang mengatakan galatnya menyebar
normal dengan simpangan baku $\sigma/\sqrt n$. Keduanya pernyataan tentang banyak
percobaan sekaligus, dan tidak satu pun menjawab pertanyaan rapat pagi tadi: dari satu
percobaan yang sudah selesai, kalimat apa yang boleh diucapkan tentang parameter yang
tidak pernah dilihat, dan seberapa sering kalimat itu salah bila prosedurnya diulang.
Bab ini menyusun jawabannya, dari penaksir sampai nilai-p.

\begin{untukanak}
Seorang petani ingin tahu berapa banyak rumput liar di sawahnya. Seluruh sawah terlalu
luas untuk diperiksa, jadi ia mengambil dua puluh petak kecil yang tersebar, lalu
menghitung rumput di setiap petak. Dari dua puluh angka itu ia menebak keadaan seluruh
sawah. Masalahnya, tebakan selalu berupa satu angka, sedangkan keadaan sawah tidak pernah
sesederhana itu.

Sekarang bagian yang membuat penasaran. Di dua puluh petak itu ia tidak menemukan
rumput sama sekali. Apakah ia boleh pulang dan berkata bahwa sawahnya bersih? Tentu
tidak. Bisa jadi rumputnya memang tidak ada, bisa jadi ia hanya kebetulan memilih
petak-petak yang terdiri dari tanah kosong. Yang dapat ia katakan hanyalah: kalau
rumputnya cukup banyak, dua puluh petak kosong itu hampir mustahil. Dari situ ia bisa
menghitung batas atas yang jujur, misalnya ``paling banyak tiga dari seratus petak''.

Hal yang sama berlaku ketika guru menilai ulangan beberapa murid dan ketika program
komputer diuji pada beberapa contoh.
Di bab ini kita akan tahu mengapa satu angka tidak pernah cukup, mengapa jawaban jujur
selalu berupa selang, dan mengapa selang itu melebar ketika buktinya sedikit.
\end{untukanak}

\section{Penaksir, bias, dan galat baku}
\tingkat{dasar}

Setiap kesimpulan statistika dimulai dari satu keputusan kecil: besaran apa yang kita
hitung dari data. Rata-rata, proporsi, simpangan baku sampel, semuanya adalah fungsi
dari data dan tidak memerlukan pengetahuan tentang parameter yang dicari. Kalimat
sehari-harinya: penaksir adalah juru ukur yang menebak tinggi pohon dari panjang
bayangannya, dan kita perlu tahu apakah juru ukur itu sistematis meleset atau hanya
goyah.

\kamus{penaksir}{estimator}{juru ukur yang menebak tinggi pohon dari panjang
bayangannya, bukan dari pohon itu sendiri}

\begin{definisi}[Statistik dan penaksir]
\index{penaksir}
Misalkan $X_1,\dots,X_n$ sampel acak dari distribusi yang bergantung pada parameter
$\theta$. \emph{Statistik} adalah fungsi dari $X_1,\dots,X_n$ yang tidak memuat
$\theta$. Statistik yang dipakai menebak $\theta$ disebut \emph{penaksir}
(\emph{estimator}) dan ditulis $\hat\theta$; nilai yang diperoleh pada satu sampel
disebut \emph{taksiran} (\emph{estimate}).
\end{definisi}

Penaksir adalah variabel acak, sebab ia fungsi dari variabel acak. Karena itu ia punya
harapan dan variansi, dan kedua besaran itulah yang menentukan mutunya. Harapannya
menjawab pertanyaan ``apakah tebakannya meleset ke satu arah terus'', sedangkan
variansinya menjawab ``seberapa jauh tebakan itu bergoyang dari sampel ke sampel''.

\begin{definisi}[Bias dan penaksir tak bias]
\index{tak bias}
Bias penaksir $\hat\theta$ untuk $\theta$ adalah $\operatorname{bias}(\hat\theta)
=\mathbb{E}[\hat\theta]-\theta$. Penaksir disebut \emph{tak bias} bila biasnya nol,
yakni $\mathbb{E}[\hat\theta]=\theta$ untuk setiap nilai $\theta$ yang mungkin.
\end{definisi}

Rata-rata sampel tak bias untuk $\mu$, dan itu sudah dipakai pada pembuktian hukum
bilangan besar di Bab 19, jadi tidak diulang di sini. Yang lebih menarik adalah
penaksir yang dipakai untuk variansi. Dengan pembagi $n$, penaksir
$S_n^2=\frac1n\sum(X_i-\bar X_n)^2$ bias ke bawah; dengan pembagi $n-1$, penaksir
$S^2=\frac1{n-1}\sum(X_i-\bar X_n)^2$ tak bias. Sebabnya, $\sum(X_i-\bar X_n)^2$
kehilangan satu derajat kebebasan karena $\bar X_n$ sendiri dihitung dari data yang
sama, sehingga harapannya $(n-1)\sigma^2$ dan bukan $n\sigma^2$. Pembagi $n-1$
adalah koreksi untuk kehilangan itu, bukan kerapian belaka.

Kualitas penaksir biasanya diringkas oleh satu angka, yaitu galat kuadrat rata-rata.
Kalimat sehari-harinya: seberapa besar kesalahan yang kita harapkan, dihitung dengan
cara yang menghukum kesalahan besar lebih berat.

\begin{definisi}[Galat kuadrat rata-rata]
\index{galat kuadrat rata-rata}
$\operatorname{MSE}(\hat\theta)=\mathbb{E}[(\hat\theta-\theta)^2]$.
\end{definisi}

\begin{teorema}[Dekomposisi MSE]
Untuk penaksir $\hat\theta$ dengan variansi berhingga,
\[
\operatorname{MSE}(\hat\theta)=\operatorname{bias}(\hat\theta)^2
+\operatorname{Var}(\hat\theta).
\]
\end{teorema}

\begin{proof}[Bukti]
Tulis $\hat\theta-\theta=(\hat\theta-\mathbb{E}\hat\theta)
+(\mathbb{E}\hat\theta-\theta)$. Kuadratkan dan ambil harapannya:
\[
\mathbb{E}[(\hat\theta-\theta)^2]
=\mathbb{E}[(\hat\theta-\mathbb{E}\hat\theta)^2]
+2(\mathbb{E}\hat\theta-\theta)\,\mathbb{E}[\hat\theta-\mathbb{E}\hat\theta]
+(\mathbb{E}\hat\theta-\theta)^2 .
\]
Suku tengah lenyap karena $\mathbb{E}[\hat\theta-\mathbb{E}\hat\theta]=0$, sedangkan
suku pertama adalah variansi dan suku ketiga adalah bias dikuadratkan.
\end{proof}

Identitas itu menjelaskan mengapa ``tak bias'' bukan gelar juara. Penaksir tak bias
boleh bervariansi besar, dan penaksir berbias sedikit boleh bervariansi jauh lebih
kecil sehingga MSE-nya lebih rendah. Untuk menaksir $\mu$ dari $n$ pengamatan,
statistik $X_1$ tak bias dengan variansi $\sigma^2$, sedangkan $\bar X_n$ tak bias
dengan variansi $\sigma^2/n$; pada $n=20$ selisihnya dua puluh kali lipat.

Contoh kedua menunjukkan bias yang harus dihitung, karena ia dipakai setiap hari
untuk proporsi. Bila $K\sim\mathrm{Bin}(n,p)$ dan $\hat p=K/n$, maka
$\hat p(1-\hat p)$ adalah penaksir alami bagi $p(1-p)$, yaitu variansi satu percobaan.
Harapannya
\[
\mathbb{E}[\hat p(1-\hat p)]
=\mathbb{E}[\hat p]-\mathbb{E}[\hat p^2]
=p-\Bigl(p^2+\frac{p(1-p)}{n}\Bigr)
=p(1-p)\Bigl(1-\frac1n\Bigr),
\]
sehingga biasnya $-p(1-p)/n$. Untuk $p=0{,}3$ dan $n=20$ nilainya $p(1-p)=0{,}21$
dan biasnya $-0{,}0105$, yakni lima persen ke bawah. Pada $n=200$ biasnya tinggal
$-0{,}00105$. Bias yang mengecil sebanding $1/n$ ini khas: ia hilang di tak hingga,
tetapi hadir pada setiap sampel berhingga.

Sifat kedua yang dituntut dari penaksir adalah bahwa ia benar-benar menuju sasaran
ketika data bertambah. Kalimat sehari-harinya: dengan cukup banyak bukti, tebakan yang
salah menjadi semakin mustahil, bukan hanya semakin jarang.

\begin{definisi}[Konsisten]
\index{konsisten}
Barisan penaksir $\hat\theta_n$ \emph{konsisten} untuk $\theta$ bila untuk setiap
$\varepsilon>0$ berlaku $P(|\hat\theta_n-\theta|\ge\varepsilon)\to0$ ketika
$n\to\infty$.
\end{definisi}

\begin{proposisi}[MSE menuju nol cukup untuk konsistensi]
\index{konsisten}
Bila $\operatorname{MSE}(\hat\theta_n)\to0$, maka $\hat\theta_n$ konsisten untuk
$\theta$.
\end{proposisi}

\begin{proof}[Bukti]
Menurut ketaksamaan Chebyshev, yang sudah dibuktikan di Bab 19,
\[
P(|\hat\theta_n-\theta|\ge\varepsilon)
\le\frac{\mathbb{E}[(\hat\theta_n-\theta)^2]}{\varepsilon^2}
=\frac{\operatorname{MSE}(\hat\theta_n)}{\varepsilon^2}.
\]
Untuk $\varepsilon$ tetap, ruas kanan menuju nol, dan ketaksamaan itu menutup
pembuktian.
\end{proof}

Perhatikan bahwa Chebyshev memakai MSE, bukan variansi. Suku bias dikuadratkan, jadi
penaksir berbias tetap konsisten asalkan biasnya menuju nol. Sebaliknya penaksir $X_1$
tak bias untuk setiap $n$ tetapi tidak pernah memakai sisa data, sehingga tidak
konsisten.

Besaran berikutnya dipakai terus sampai akhir bab. Variansi penaksir mengatakan
seberapa jauh tebakan bergoyang, dan akarnya punya nama sendiri.

\kamus{galat baku}{standard error}{kisaran goyangan satu tebakan, bukan lebar seluruh
data}

\begin{definisi}[Galat baku]
\index{galat baku}
\emph{Galat baku} penaksir $\hat\theta$ adalah $\operatorname{se}(\hat\theta)
=\sqrt{\operatorname{Var}(\hat\theta)}$. Bila rumusnya memuat parameter yang tidak
diketahui, galat baku \emph{tertaksir} diperoleh dengan mengganti parameter itu oleh
taksirannya.
\end{definisi}

Untuk proporsi, $\operatorname{se}(\hat p)=\sqrt{p(1-p)/n}$, dan versi tertaksirnya
$\sqrt{\hat p(1-\hat p)/n}$. Untuk data $33$ dari $50$, $\hat p=0{,}66$ dan
\[
\sqrt{\frac{0{,}66\cdot0{,}34}{50}}=\sqrt{0{,}004488}=0{,}066993 .
\]
Untuk dua belas pengukuran waktu tunda (dalam milidetik),
\begin{align*}
&18{,}2,\ 19{,}1,\ 17{,}8,\ 20{,}3,\ 18{,}9,\ 19{,}7,\\
&18{,}4,\ 21{,}0,\ 17{,}5,\ 19{,}9,\ 18{,}6,\ 20{,}1,
\end{align*}
rata-ratanya $19{,}125$ dan simpangan baku sampelnya $1{,}082191$. Galat bakunya
$1{,}082191/\sqrt{12}=0{,}312401$.
Angka itu tidak mengatakan bahwa rata-rata sebenarnya berjarak $0{,}312$ dari
$19{,}125$; ia mengatakan bahwa tebakan seperti ini, bila diulang, bergoyang dengan
skala sekitar $0{,}31$. Selang kepercayaan adalah cara mengubah kalimat terakhir
menjadi kalimat yang bisa dipakai.

\section{Interval kepercayaan untuk rata-rata}
\tingkat{menengah}

Satu angka tanpa ukuran goyangan mudah disalahgunakan. Yang dibutuhkan sepasang angka
yang, dengan peluang besar, mengurung parameter. Kalimat sehari-harinya: kita tidak
menembak satu titik, kita memasang jaring dan berjanji jaring itu menangkap ikan yang
kita cari pada sebagian besar lemparan.

\kamus{interval kepercayaan}{confidence interval}{jaring yang dipasang dengan lebar
tertentu, dan janji bahwa jaring itu menangkap sasaran pada sebagian besar lemparan}

\begin{definisi}[Interval kepercayaan dan cakupan]
\index{interval kepercayaan}
Misalkan $C(X_1,\dots,X_n)$ adalah selang yang dihitung dari data. $C$ disebut
\emph{interval kepercayaan} bertaraf $1-\alpha$ untuk $\theta$ bila
$P_\theta(\theta\in C)\ge1-\alpha$ untuk setiap $\theta$. Peluang itu disebut
\emph{cakupan} (\emph{coverage}) prosedurnya.
\end{definisi}

Perhatikan letak peluangnya: $P_\theta$ dihitung sebelum data diambil, dan yang acak
adalah $C$, bukan $\theta$. Ini tempat kesalahan tafsir paling sering terjadi, dan kita
kembali ke sana sesudah rumusnya ada.

\begin{teorema}[Interval-z]
\index{interval kepercayaan}
Misalkan $X_1,\dots,X_n$ bebas dan berdistribusi $N(\mu,\sigma^2)$ dengan $\sigma$
diketahui. Untuk $z=z_{1-\alpha/2}$, selang
\[
C=\Bigl[\bar X_n-z\frac{\sigma}{\sqrt n},\ \bar X_n+z\frac{\sigma}{\sqrt n}\Bigr]
\]
mempunyai cakupan tepat $1-\alpha$. Untuk distribusi lain dengan variansi berhingga,
cakupannya menuju $1-\alpha$ ketika $n\to\infty$.
\end{teorema}

\begin{proof}[Bukti]
Misalkan $Z=\sqrt n(\bar X_n-\mu)/\sigma$. Untuk sampel normal, $Z\sim N(0,1)$ tepat;
untuk distribusi lain, $Z$ menuju $N(0,1)$ menurut teorema limit pusat di Bab 19.
Ketaksamaan $\mu\in C$ setara dengan $|Z|\le z$, karena
\[
|\bar X_n-\mu|\le z\frac{\sigma}{\sqrt n}
\iff \frac{\sqrt n\,|\bar X_n-\mu|}{\sigma}\le z .
\]
Jadi $P(\mu\in C)=P(|Z|\le z)=1-\alpha$ pada kasus normal, dan menuju $1-\alpha$
pada kasus umum.
\end{proof}

Bila $\sigma$ tidak diketahui, seperti pada hampir semua data nyata, penggantian
$\sigma$ oleh $s$ menambah ketidakpastian dan menebalkan ekor. Hasilnya memakai
distribusi $t$ dengan $n-1$ derajat kebebasan, diterbitkan Gosset pada 1908 dengan nama
samaran Student:
\[
C=\Bigl[\bar X_n-t_{n-1,1-\alpha/2}\frac{s}{\sqrt n},\
\bar X_n+t_{n-1,1-\alpha/2}\frac{s}{\sqrt n}\Bigr],
\]
dengan cakupan tepat $1-\alpha$ bila sampelnya normal. Buktinya menuntut distribusi
bersama $\bar X_n$ dan $s$ yang bebas, dan ada di Wasserman, \emph{All of Statistics},
Bab 6. Untuk $n=12$ data waktu tunda tadi, $t_{11,0{,}975}=2{,}200985$, sehingga
setengah lebarnya $2{,}200985\cdot0{,}312401=0{,}687591$ dan selangnya
$(18{,}4374;\ 19{,}8126)$. Dengan $z=1{,}96$ setengah lebarnya $0{,}612296$: mengakui
$\sigma$ tidak diketahui melebarkan selang dua belas persen pada $n=12$, dan satu
persen pada $n=100$.

Sekarang tafsirannya, dan di sini kita harus lambat. Sesudah data diambil,
$C$ bukan lagi variabel acak: ia selang tetap $(18{,}4374;\ 19{,}8126)$. Rata-rata
sebenarnya $\mu$ juga bilangan tetap, hanya saja tidak diketahui. Kalimat
``peluang $\mu$ berada di dalam $(18{,}4374;\ 19{,}8126)$ sama dengan $0{,}95$''
karena itu tidak bermakna dalam kerangka ini: tidak ada peluang yang tersisa, sebab
tidak ada yang acak. Yang berpeluang $0{,}95$ adalah \emph{prosedurnya}, dan itu
pernyataan tentang percobaan yang belum dijalankan.

Gambar~\ref{fig:dua-puluh-selang} menjalankan percobaan itu dua puluh kali, dengan
$\mu=19{,}125$ dan $\sigma=1$ yang diketahui: sembilan belas selang memuat $\mu$, satu
tidak.

\begin{figure}[htbp]
\centering
\begin{tikzpicture}[x=4.15cm,y=0.235cm,font=\scriptsize]
  \draw[->,layergrey] (17.80,0)--(20.22,0);
  \draw[layergrey,densely dashed,thick] (19.125,0)--(19.125,21.4);
  \node[above,font=\scriptsize] at (19.125,21.4) {$\mu=19{,}125$};
  \foreach \y/\a/\b in {
    1/18.5955/19.7271,
    2/18.5211/19.6527,
    3/18.7441/19.8757,
    4/18.5895/19.7211,
    5/18.4046/19.5362,
    6/18.6636/19.7952,
    7/18.9356/20.0672,
    8/18.8326/19.9642,
    9/18.3561/19.4876,
    10/18.1939/19.3255,
    11/18.3793/19.5109,
    12/18.5711/19.7027,
    14/18.4960/19.6276,
    15/18.1995/19.3311,
    16/18.3478/19.4794,
    17/18.4021/19.5337,
    18/18.4679/19.5995,
    19/18.6780/19.8096,
    20/18.8602/19.9917
  } {
    \draw[very thick,layerindigo] (\a,\y)--(\b,\y);
    \draw[layerindigo] (\a,\y-0.22)--(\a,\y+0.22);
    \draw[layerindigo] (\b,\y-0.22)--(\b,\y+0.22);
  }
  \draw[very thick,accent] (17.8880,13)--(19.0196,13);
  \draw[accent] (17.8880,12.7)--(17.8880,13.3);
  \draw[accent] (19.0196,12.7)--(19.0196,13.3);
  \node[accent,anchor=west,font=\scriptsize] at (17.82,0.62) {satu selang meleset};
  \foreach \t in {18,19,20} {\draw[layergrey] (\t,-0.2)--(\t,0.2);}
  \foreach \t in {18,19,20} {\node[below,font=\scriptsize] at (\t,-0.45) {\t};}
\end{tikzpicture}
\caption{Dua puluh selang kepercayaan $95$ persen untuk $\mu$, masing-masing dari satu
sampel berukuran $n=12$ dengan $\sigma=1$ diketahui, sehingga setengah lebarnya
$1{,}96/\sqrt{12}=0{,}5658$ dan lebarnya $1{,}1316$. Garis putus-putus tegak adalah
$\mu=19{,}125$ yang sebenarnya. Sembilan belas selang memuatnya; selang pada baris
ke-13, yaitu $(17{,}8880;\,19{,}0196)$, tidak. Satu kegagalan dalam dua puluh adalah
angka yang wajar untuk janji $95$ persen, bukan tanda prosedur yang rusak.}
\label{fig:dua-puluh-selang}
\end{figure}

\begin{catatan}
Tiga kalimat yang benar, dan satu yang salah. Benar: ``prosedur ini menangkap $\mu$
pada $95$ persen percobaan yang diulang''. Benar: ``kita $95$ persen yakin $\mu$
berada di antara $18{,}44$ dan $19{,}81$'', asalkan ``yakin'' dibaca sebagai ringkasan
percobaan berulang. Benar: ``prosedur ini salah pada sekitar lima dari seratus
kasus''. Salah: ``peluang $\mu$ ada di selang ini $0{,}95$''. Kalimat terakhir
berpura-pura mengetahui sesuatu tentang $\mu$ yang tidak diberikan data.
\end{catatan}

Lebar selang menyusut sebanding $1/\sqrt n$, bukan $1/n$: memangkasnya menjadi separuh
menuntut empat kali data, menjadi sepersepuluh menuntut seratus kali. Aturan ini dapat
dihitung sebelum percobaan dimulai, sebab setengah lebar $z\sigma/\sqrt n\le d$ setara
dengan $n\ge(z\sigma/d)^2$; dengan $\sigma=1$, $z=1{,}96$, dan $d=0{,}1$, dibutuhkan
$n\ge385$.

\begin{tangan}
Siapkan dua puluh kartu, tulis angka $1$ sampai $20$ pada kartu, satu angka per kartu.
Kocok, ambil lima kartu tanpa melihat, jumlahkan, lalu bagi lima. Tulis hasilnya pada
satu baris kertas. Kembalikan kelima kartu, kocok lagi, ulangi. Lakukan tiga puluh kali,
selalu lima kartu, dan susun angka-angka itu sebagai titik pada satu garis bilangan.

Yang harus terlihat: titik-titik itu tidak menumpuk di satu angka. Rata-rata seluruh
kartu sebenarnya $10{,}5$, dan titikmu menyebar di sekitarnya, ada yang $8$, ada yang
$13$. Sekarang ulangi seluruh pekerjaan dengan sepuluh kartu per pengambilan, tiga puluh
kali lagi. Sebaran kedua lebih sempit, tetapi tidak setengahnya: ia menyempit kira-kira
sebanding akar dua, sebab banyak kartu berlipat dua. Tempel kedua baris titik itu
berdampingan, dan lihat sendiri bahwa menggandakan pekerjaan hanya memangkas goyangan
sekitar tiga puluh persen.
\end{tangan}

\section{Proporsi: Wald, Wilson, dan aturan tiga}
\tingkat{menengah}

Untuk proporsi, rumus yang paling banyak dipakai adalah versi paling sederhana dari
teorema sebelumnya: ganti $\bar X_n$ oleh $\hat p$ dan $\sigma/\sqrt n$ oleh galat
baku proporsi. Kalimat sehari-harinya: kita pasang selang simetris di sekitar proporsi
yang teramati, selebar dua kali galat bakunya.

\begin{definisi}[Interval Wald]
\index{interval wald}
Interval Wald bertaraf $1-\alpha$ untuk proporsi adalah
\[
\hat p\pm z\sqrt{\frac{\hat p(1-\hat p)}{n}},\qquad z=z_{1-\alpha/2}.
\]
\end{definisi}

Rumus itu wajar selama $\hat p$ jauh dari nol dan satu. Kesulitannya muncul tepat di
tempat data paling sedikit, dan bentuknya tidak halus. Bila $\hat p=0$, rumus itu
memberi selang $[0;\,0]$ selebar nol: sebuah klaim percaya diri penuh bahwa laju
kejadiannya nol, padahal yang kita punya hanya ketiadaan kejadian pada $n$ percobaan.
Bila $\hat p=1/20$, batas bawahnya $-0{,}0455$, sebuah bilangan yang bukan peluang.

Penyebabnya bukan pembulatan. Wald memakai $\hat p$ untuk dua tugas sekaligus:
sebagai titik tengah dan sebagai penaksir variansi. Ketika $\hat p$ mendekati nol,
variansi taksirannya ikut mengecil, padahal ketidakpastian sebenarnya justru besar pada
skala itu.

Jalan keluarnya adalah membalik pertanyaannya: alih-alih menganggap $\hat p$ menyebar
normal di sekitar $p$ yang benar, kita tanyakan nilai $p$ mana saja yang masih masuk
akal dengan data teramati. Kriterianya, jarak ternormalisasi antara $\hat p$ dan $p$
tidak melampaui $z$, dengan variansi di penyebut memakai $p$, bukan $\hat p$.

\kamus{interval wilson}{Wilson interval}{selang yang dicari dengan bertanya nilai mana
yang masih masuk akal, bukan dengan memasang selang simetris di sekitar tebakan}

\begin{teorema}[Interval Wilson]
\index{interval wilson}
Untuk $0\le k\le n$, $\hat p=k/n$, dan $z>0$, himpunan $p\in[0,1]$ yang memenuhi
\[
n(\hat p-p)^2\le z^2p(1-p)
\]
adalah selang tertutup dengan ujung
\[
\frac{\hat p+\dfrac{z^2}{2n}\ \pm\
z\sqrt{\dfrac{\hat p(1-\hat p)}{n}+\dfrac{z^2}{4n^2}}}{1+\dfrac{z^2}{n}} .
\]
Kedua ujungnya selalu terletak di $[0,1]$, dan selangnya tidak pernah kosong.
\end{teorema}

\begin{proof}[Bukti]
Susun ulang ketaksamaan menjadi $(n+z^2)p^2-(2n\hat p+z^2)p+n\hat p^2\le0$.
Diskriminannya
\[
(2n\hat p+z^2)^2-4(n+z^2)n\hat p^2
=4z^2\Bigl(n\hat p(1-\hat p)+\frac{z^2}{4}\Bigr)\ge0,
\]
jadi akarnya real dan rumus akar persamaan kuadrat memberi bentuk pada pernyataan
teorema setelah pembilang dan penyebut dibagi $2n$.

Untuk letaknya, misalkan $f(p)=n(\hat p-p)^2-z^2p(1-p)$. Koefisien $p^2$ adalah
$n+z^2>0$, jadi $\{f\le0\}$ adalah selang $[r_-,r_+]$. Karena
$f(0)=n\hat p^2\ge0$, titik $0$ tidak berada di bagian dalam selang itu, sehingga
$r_-\ge0$; karena $f(1)=n(1-\hat p)^2\ge0$, titik $1$ juga tidak, sehingga
$r_+\le1$. Terakhir $f(\hat p)=-z^2\hat p(1-\hat p)\le0$ menempatkan $\hat p$ di dalam
selang, jadi selang itu tidak kosong.
\end{proof}

Sekarang bandingkan keduanya pada data yang sama. Tabel berikut memakai $z=1{,}96$ dan
selang kepercayaan $95$ persen; semua angka dihitung dengan rumus di atas dan diperiksa
ulang dengan Python.

\begin{center}
\small
\begin{tabular}{lrrrr}
\hline
data & $\hat p$ & Wald $95\%$ & Wilson $95\%$ & lebar Wilson \\
\hline
$33/50$  & 0,660 & $(0{,}5287;\ 0{,}7913)$ & $(0{,}5215;\ 0{,}7756)$ & 0,254 \\
$53/100$ & 0,530 & $(0{,}4322;\ 0{,}6278)$ & $(0{,}4329;\ 0{,}6249)$ & 0,192 \\
$13/20$  & 0,650 & $(0{,}4410;\ 0{,}8590)$ & $(0{,}4329;\ 0{,}8188)$ & 0,386 \\
$1/20$   & 0,050 & $(-0{,}0455;\ 0{,}1455)$ & $(0{,}0089;\ 0{,}2361)$ & 0,227 \\
$0/20$   & 0,000 & $(0;\ 0)$ & $(0;\ 0{,}1611)$ & 0,161 \\
\hline
\end{tabular}
\end{center}

Dua baris pertama hampir tidak berbeda: ketika $\hat p$ dekat setengah, keduanya sah
dan sepakat sampai dua angka desimal. Baris keempat dan kelima memperlihatkan
pemisahan. Pada $1/20$ selang Wald \emph{lebih sempit} daripada Wilson, $0{,}191$
melawan $0{,}227$, dan kesempitan itulah tandanya salah: batas bawahnya keluar dari
$[0,1]$. Pada $0/20$ Wald memberi lebar nol, sedangkan Wilson menyisakan ruang sampai
$0{,}1611$. Kesimpulannya bukan bahwa Wilson selalu lebih baik, melainkan bahwa lebar
selang tidak dapat dinilai tanpa cakupannya.

Untuk kasus nol kejadian ada hasil khusus yang lebih tua daripada Wilson dan lebih
mudah diingat. Kalimat sehari-harinya: tidak ada kejadian sama sekali bukan bukti
bahwa laju kejadiannya nol, dan batas atasnya kira-kira tiga dibagi banyak percobaan.

\kamus{aturan tiga}{rule of three}{jika tidak ada kejadian dalam sekian percobaan,
laju kejadiannya masih boleh setinggi tiga dibagi banyak percobaan}

\begin{teorema}[Aturan tiga]
\index{aturan tiga}
Misalkan $K\sim\mathrm{Bin}(n,p)$ dan teramati $K=0$. Batas atas satu sisi pada taraf
$95\%$ adalah $p_{\mathrm{hi}}=1-0{,}05^{1/n}$, dan berlaku
\[
p_{\mathrm{hi}}\le\frac{\ln 20}{n}<\frac{3}{n}.
\]
\end{teorema}

\begin{proof}[Bukti]
Peluang mengamati nol kejadian adalah $(1-p)^n$. Menetapkan peluang itu sama dengan
$0{,}05$ memberi $(1-p_{\mathrm{hi}})^n=0{,}05$, yakni
$p_{\mathrm{hi}}=1-0{,}05^{1/n}=1-e^{-c/n}$ dengan $c=\ln 20=2{,}9957\ldots$ Karena
$1-e^{-x}\le x$ untuk $x\ge0$, diperoleh $p_{\mathrm{hi}}\le c/n<3/n$.
\end{proof}

Untuk $n=100$, aturan tiga memberi $3/100=0{,}03$, batas eksaknya $0{,}0295130$, dan
batas atas Wilson dua sisi pada $k=0$ adalah $z^2/(n+z^2)=0{,}0369935$. Ketiga angka
menjawab pertanyaan yang sama dengan tiga janji berbeda; aturan tiga adalah batas satu
sisi yang eksak, Wilson aproksimasi dua sisi. Yang tidak boleh adalah menyebut angka
tanpa menyebut prosedurnya.

\section{Uji hipotesis dan nilai-p}
\tingkat{menengah}

Sampai di sini pertanyaannya berbentuk ``berapa nilainya''. Pertanyaan berikutnya
berbentuk ``apakah nilainya sama dengan angka tertentu''. Dalam bahasa sehari-hari,
kita memasang tuduhan bahwa tidak ada efek, lalu memeriksa apakah data cukup
mengejutkan untuk membatalkan tuduhan itu.

\begin{definisi}[Hipotesis nol, alternatif, dan taraf]
\index{hipotesis nol}
\emph{Hipotesis nol} $H_0$ adalah pernyataan tentang parameter yang ingin diuji;
\emph{hipotesis alternatif} $H_1$ adalah himpunan nilai yang dianggap benar bila $H_0$
ditolak. \emph{Taraf} $\alpha$ adalah peluang menolak $H_0$ ketika $H_0$ sebenarnya
benar. Aturan penolakan ditentukan oleh satu statistik uji $T$ dan satu daerah tolak
$R$: kita menolak $H_0$ bila $T\in R$.
\end{definisi}

Taraf $\alpha$ dipilih sebelum data dilihat, dan di situlah letak seluruh
kredibilitas prosedur: kita bersedia salah menuduh pada sebagian kecil percobaan, dan
kita menyebutkan berapa besar bagian itu. Nilai-p memindahkan gagasan yang sama ke
satu dataset.

\begin{definisi}[Nilai-p]
\index{nilai-p}
Untuk data $x$ dan statistik uji $T$, nilai-p adalah peluang memperoleh hasil yang
sekurang-kurangnya sejauh hasil teramati bila $H_0$ benar,
\[
\text{nilai-}p=\sup_{\theta\in H_0}P_\theta\bigl(T\text{ sejauh }T(x)\bigr).
\]
Kita menolak $H_0$ pada taraf $\alpha$ bila nilai-p $\le\alpha$.
\end{definisi}

Kalimat ``sekurang-kurangnya sejauh'' perlu ditentukan untuk uji dua sisi, dan
ketentuan yang lazim adalah menjumlahkan kedua ekor: untuk statistik $z$ yang
simetris, nilai-p $=2(1-\Phi(|z|))$. Definisi itu membuat nilai-p menjadi variabel
acak, dan sifatnya di bawah $H_0$ dapat dinyatakan tepat.

\begin{teorema}[Nilai-p seragam di bawah $H_0$]
\index{nilai-p}
Bila $H_0$ benar dan $T$ berdistribusi kontinu dengan fungsi sebaran $F_0$, maka
$P_{H_0}(\text{nilai-}p\le\alpha)=\alpha$ untuk setiap $\alpha\in(0,1)$. Akibatnya
$P_{H_0}(\text{nilai-}p\le0{,}05)=0{,}05$: uji dengan aturan ``tolak bila nilai-p
$\le\alpha$'' mempunyai taraf tepat $\alpha$.
\end{teorema}

\begin{proof}[Bukti]
Untuk uji satu sisi kanan, nilai-p $=1-F_0(T)$, sehingga
\[
P_{H_0}(1-F_0(T)\le\alpha)
=P_{H_0}\bigl(F_0(T)\ge1-\alpha\bigr)
=P_{H_0}\bigl(T\ge F_0^{-1}(1-\alpha)\bigr)=\alpha,
\]
karena $F_0(T)$ berdistribusi seragam pada $(0,1)$ bila $F_0$ kontinu. Untuk uji dua
sisi, nilai-p bergantung pada $|T|$ dan pembuktiannya memakai langkah yang sama pada
sebaran $|T|$.
\end{proof}

Keseragaman itu punya akibat yang sering terlewat: di bawah $H_0$ nilai-p tersebar
rata, jadi nilai-p $=0{,}04$ muncul pada empat persen percobaan yang hipotesis nolnya
benar. Yang membedakan dua percobaan bukan nilai-p tunggal, melainkan daya
masing-masing.

\begin{bacakeras}
Dua suara, bergantian.

``Nilai p nol koma nol empat. Berarti peluang tidak ada efek tinggal empat persen.''

``Bukan. Nilai p menghitung data, bukan tuduhan. Seandainya tidak ada efek, dan
percobaan ini diulang seratus kali, sekitar empat kali data akan tampak sejauh ini.''

``Jadi yang dihitung data.''

``Ya, data. Tuduhan hanya tempat berpijak.''

``Dan kalau nilai p besar?''

``Berarti data tidak mengejutkan. Diamnya timbangan bukan bukti bahwa barangnya ringan.''

Ucapkan sekali lagi, pelan: nilai p adalah peluang data, bukan peluang tuduhan. Ulangi
baris itu sampai terasa wajar di mulut: peluang data, bukan peluang tuduhan.
\end{bacakeras}


\begin{resep}
\begin{enumerate}
  \item Tulis $H_0$ dan $H_1$ sebagai pernyataan tentang satu parameter, sebelum
        menghitung apa pun.
  \item Tetapkan $\alpha$ dan daerah tolaknya, lalu berhenti sejenak. Jangan melihat
        data dulu.
  \item Hitung statistik uji dengan galat baku di bawah $H_0$, lalu nilai-p.
  \item Bandingkan nilai-p dengan $\alpha$ yang sudah ditetapkan, dan laporkan
        selangnya bersama keputusan itu.
\end{enumerate}
Peringatan urutan: langkah 2 harus selesai sebelum langkah 3. Bila $\alpha$ dipilih
sesudah data terlihat, $\alpha$ bukan lagi janji melainkan catatan hasil, dan langkah 4
kehilangan artinya. Yang tidak boleh ditukar: langkah 2 dan langkah 3.
\end{resep}

Sekarang contoh dengan angka. Untuk data $33$ dari $50$ dan $H_0\colon p=0{,}5$,
statistik $z$ memakai galat baku di bawah $H_0$, yaitu $\sqrt{0{,}25/50}=0{,}070711$:
\[
z=\frac{0{,}66-0{,}5}{0{,}070711}=2{,}2627,\qquad
\text{nilai-}p=2(1-\Phi(2{,}2627))=0{,}0237 .
\]
Uji binomial eksak atas aturan yang sama, $P(X\ge33)+P(X\le17)$ untuk
$X\sim\mathrm{Bin}(50;0{,}5)$, memberi $0{,}0328$. Kedua angka menolak $H_0$ pada
taraf lima persen, tetapi aproksimasi normal melaporkan nilai-p hampir tiga puluh
persen lebih kecil. Pada $n=50$ selisih sebesar itu cukup membalik keputusan bila
tarafnya $0{,}03$; laporan yang baik menyebutkan uji mana yang dipakai, bukan hanya
nilai-p-nya.

Gambar~\ref{fig:dua-ekor} menaruh kedua batas itu pada sebaran $z$ di bawah $H_0$: dua
ekor seluas $0{,}0250$ masing-masing, dan statistik teramati $2{,}2627$ di dalam ekor
kanan.

\begin{figure}[htbp]
\centering
\begin{tikzpicture}
\begin{axis}[
  width=11.2cm, height=5.4cm,
  xmin=-4, xmax=4, ymin=0, ymax=0.46,
  axis lines=left,
  xtick={-1.96,0,1.96},
  xticklabels={$-1{,}96$,$0$,$1{,}96$},
  ytick=\empty,
  tick label style={font=\small},
  xlabel={nilai $z$ bila $H_0$ benar},
  xlabel style={font=\small},
]
\addplot[fill=accent!25,draw=none,domain=-3.95:-1.96,samples=100]
  {0.39894228*exp(-0.5*x^2)} \closedcycle;
\addplot[fill=accent!25,draw=none,domain=1.96:3.95,samples=100]
  {0.39894228*exp(-0.5*x^2)} \closedcycle;
\addplot[very thick,layerindigo,domain=-3.95:3.95,samples=200]
  {0.39894228*exp(-0.5*x^2)};
\draw[densely dashed,layergrey] (axis cs:-1.96,0)--(axis cs:-1.96,0.0584);
\draw[densely dashed,layergrey] (axis cs:1.96,0)--(axis cs:1.96,0.0584);
\draw[very thick,accent] (axis cs:2.2627,0)--(axis cs:2.2627,0.0308);
\node[font=\scriptsize] (tl) at (axis cs:-2.55,0.135) {$0{,}025$};
\draw[-{Stealth[length=4pt]},layergrey] (tl)--(axis cs:-2.55,0.012);
\node[accent,font=\scriptsize] (tz) at (axis cs:2.05,0.215) {$z=2{,}2627$};
\draw[-{Stealth[length=4pt]},accent] (tz)--(axis cs:2.2627,0.055);
\node[font=\scriptsize] (tr) at (axis cs:3.05,0.135) {$0{,}025$};
\draw[-{Stealth[length=4pt]},layergrey] (tr)--(axis cs:3.05,0.004);
\end{axis}
\end{tikzpicture}
\caption{Sebaran statistik $z=\sqrt n(\hat p-p_0)/\sqrt{p_0(1-p_0)}$ bila $H_0$ benar,
untuk data $33$ dari $50$ dengan $H_0\colon p=0{,}5$. Dua garis putus-putus di
$\pm1{,}96$ memotong dua ekor yang diwarnai, masing-masing seluas
$1-\Phi(1{,}96)=0{,}0250$, sehingga taraf dua sisi $0{,}05$. Garis tebal di
$2{,}2627$ adalah statistik yang teramati pada data itu; ekor kanan dari titik itu
seluas $0{,}0118$, dan dua kali luas itu adalah nilai-p $0{,}0237$ yang dilaporkan uji
normal.}
\label{fig:dua-ekor}
\end{figure}

Interval kepercayaan dan uji hipotesis bukan dua prosedur berbeda; keduanya dua cara
membaca himpunan yang sama.

\begin{proposisi}[Kesetaraan selang dan uji]
\index{interval kepercayaan}
Untuk selang Wald dan uji-$z$ dua sisi yang memakai galat baku $\hat p$, berlaku
$p_0$ termuat dalam selang Wald bertaraf $1-\alpha$ jika dan hanya jika $H_0\colon
p=p_0$ tidak ditolak pada taraf $\alpha$. Secara umum, selang Wilson adalah himpunan
$p_0$ yang tidak ditolak oleh uji skor pada taraf $\alpha$.
\end{proposisi}

\begin{proof}[Bukti]
Untuk pasangan Wald, kedua pernyataan memakai ketaksamaan yang sama, yaitu
$|\hat p-p_0|\le z\sqrt{\hat p(1-\hat p)/n}$: yang pertama membacanya sebagai
``$p_0$ berada di dalam selang'', yang kedua sebagai ``statistik uji tidak melewati
$z$''. Untuk Wilson, himpunan $\{p_0:n(\hat p-p_0)^2\le z^2p_0(1-p_0)\}$ adalah
definisi uji skor yang dibalik, dan teorema Wilson sudah menghitung ujungnya.
\end{proof}

\section{Galat tipe I, galat tipe II, dan daya}
\tingkat{lanjut}

Setiap keputusan biner dapat salah dengan dua cara, dan keduanya tidak simetris.
Kalimat sehari-harinya: alarm kebakaran dapat berbunyi padahal tidak ada api, dan
dapat diam padahal ada api. Kesalahan pertama membuat kita keluar rumah tanpa sebab;
kesalahan kedua jauh lebih mahal.

\begin{definisi}[Galat tipe I, galat tipe II, dan daya]
\index{galat tipe i}
\emph{Galat tipe I} adalah menolak $H_0$ ketika $H_0$ benar; peluangnya adalah taraf
$\alpha$. \emph{Galat tipe II} adalah tidak menolak $H_0$ ketika $H_1$ benar;
peluangnya ditulis $\beta$. \emph{Daya} uji adalah $1-\beta$, yaitu peluang menolak
$H_0$ ketika efek yang dicari benar-benar ada.
\end{definisi}

Gambar~\ref{fig:dua-galat} menyusun keempat kemungkinan itu sebagai kotak keputusan,
dengan peluang dari hitungan eksak pada $n=20$, $p_0=0{,}5$, dan $p_1=0{,}75$ yang
muncul di Lab komputasi: $\alpha=0{,}0414$, daya $0{,}6172$, jadi $\beta=0{,}3828$.

\begin{figure}[htbp]
\centering
\begin{tikzpicture}[font=\small,
  keputusan/.style={draw=layergrey,thick,rounded corners=2pt,
    minimum width=4.3cm,minimum height=1.6cm,align=center,inner sep=4pt}]
\node[keputusan,fill=layerjade!10] (a) at (0,0)
  {tidak menolak $H_0$\\[3pt] benar, peluang $1-\alpha=0{,}9586$};
\node[keputusan,fill=accent!12] (b) at (5.1,0)
  {menolak $H_0$\\[3pt] galat tipe I, peluang $\alpha=0{,}0414$};
\node[keputusan,fill=accent!12] (c) at (0,-2.4)
  {tidak menolak $H_0$\\[3pt] galat tipe II, peluang $\beta=0{,}3828$};
\node[keputusan,fill=layerjade!10] (d) at (5.1,-2.4)
  {menolak $H_0$\\[3pt] benar, daya $1-\beta=0{,}6172$};
\node[font=\scriptsize,above=2pt] at (0,1.0) {$H_0$ benar: $p=0{,}5$};
\node[font=\scriptsize,above=2pt] at (5.1,1.0) {$H_1$ benar: $p=0{,}75$};
\end{tikzpicture}
\caption{Dua jenis kesalahan sebagai empat kotak keputusan. Baris atas: keadaan
sebenarnya $H_0$ benar; baris bawah: $H_1$ benar dengan $p=0{,}75$. Kolom kiri: uji
tidak menolak; kolom kanan: uji menolak. Kotak berwarna merah adalah dua kesalahan yang
tidak dapat ditekan bersamaan: mengecilkan $\alpha=0{,}0414$ akan memperbesar
$\beta=0{,}3828$, dan sebaliknya. Angka dihitung dari sebaran binomial eksak pada
$n=20$ dengan aturan tolak $k\le5$ atau $k\ge15$.}
\label{fig:dua-galat}
\end{figure}

Daya bukan sifat uji saja, melainkan sifat uji bersama ukuran sampel dan besar efek.
Karena itu ia harus dihitung untuk satu nilai $p_1$ yang dianggap penting, bukan
secara umum. Untuk uji dua sisi dengan $H_0\colon p=0{,}5$, $n=100$, dan aturan tolak
bila $|z|>1{,}96$, daerah tolaknya adalah $k\le40$ atau $k\ge60$. Hitungan eksak atas
binomial memberi ukuran $0{,}0569$, bukan $0{,}05$, dan daya $0{,}8750$ pada
$p_1=0{,}65$ sedangkan aproksimasi normal memberi $0{,}8508$. Aproksimasi normal,
yakni
\[
1-\beta\approx\Phi\Bigl(\frac{(p_1-p_0)\sqrt n}{\sqrt{p_0(1-p_0)}}-z\Bigr),
\]
cukup dekat untuk perencanaan, tetapi selisih dua setengah poin itu nyata dan berasal
dari kediskretan data.

Pelajaran pertama: taraf nominal jarang tercapai tepat. Pada $n=20$ aturan yang sama
menolak bila $k\le5$ atau $k\ge15$, dengan ukuran eksak $0{,}0414$; pada $n=100$
ukurannya $0{,}0569$. Uji diskret tidak dapat memilih ambang di antara dua bilangan
bulat, jadi taraf sebenarnya menggantung pada $n$ dan $p_0$. Pelajaran kedua: pada
$n=20$ daya terhadap $p_1=0{,}75$ hanya $0{,}6172$, sehingga hampir empat dari sepuluh
percobaan berukuran dua puluh gagal menemukannya.

Angka terakhir itu membuka salah tafsir yang paling mahal: membaca ``nilai-p
$>0{,}05$'' sebagai ``tidak ada efek''. Pertama, nilai-p adalah
$P(\text{data}\mid H_0)$, bukan $P(H_0\mid\text{data})$; keduanya peluang bersyarat
dengan syarat tertukar, dan pertukaran itu sudah kita temui pada uji diagnostik di
Bab 17. Kedua, nilai-p besar adalah pernyataan tentang data, bukan tentang dunia: data
yang sedikit tidak dapat mengejutkan siapa pun. Ketiga, kegagalan menolak dapat terjadi
justru karena ujinya lemah, seperti pada $n=20$ tadi.

Contoh kecil membuat ketiganya konkret. Sepuluh keberhasilan dari dua puluh percobaan
terhadap $H_0\colon p=0{,}5$ memberi nilai-p dua sisi tepat $1$: data itu berada
persis di tengah. Apakah itu bukti bahwa $p=0{,}5$? Tidak. Selang Wilson $95$ persen
untuk data yang sama adalah $(0{,}2993;\ 0{,}7007)$. Jadi laju sebenarnya boleh
$0{,}30$, boleh $0{,}70$, dan data ini tidak dapat membedakannya dari $0{,}50$.
Nilai-p $=1$ hanya berarti ``tidak ada arah penyimpangan'', dan itu jauh lebih sedikit
daripada yang biasa dibaca orang.

\begin{tantangan}
Dua versi pendeteksi diuji pada \mbox{$100\,000$} contoh masing-masing. Versi A menangkap
\mbox{$66\,000$}, versi B menangkap \mbox{$65\,000$}. Uji dua proporsi memberi nilai-p sekitar
$0{,}0000026$, dan tim menulis di laporan: ``A lebih baik, $p<0{,}001$''. Apakah
kalimat itu didukung data? Jawab dengan satu selang, bukan dengan nilai-p.

Petunjuk 1: nilai-p mengukur kejutan, bukan besar efek. Hitung selang $95$ persen untuk
selisih kedua proporsi, dan perhatikan satuannya.

Petunjuk 2: galat baku selisihnya $0{,}002126$, jadi selangnya terbentang dari
$0{,}0058$ sampai $0{,}0142$, yaitu nol koma enam sampai satu koma empat poin persen.
Bandingkan dengan ambang praktis yang tim itu sendiri pakai: tiga poin persen.

Kalau mentok, kembali ke bagian Galat tipe I, galat tipe II, dan daya.
\end{tantangan}

\begin{untukanak}
Seorang wasit garis bertugas mengangkat bendera bila bola keluar lapangan. Ia hanya
mengangkat bendera kalau melihat sendiri bola itu keluar, dengan cukup jelas. Bila ia
tidak mengangkat bendera, apakah itu berarti bola selalu di dalam? Tentu tidak. Bisa
jadi bola keluar, tetapi pandangannya terhalang pemain lain atau bolanya sangat cepat.

Cara ia bekerja mirip dengan cara sebuah uji hipotesis bekerja. Ia diberi satu aturan
yang harus diputuskan sebelum pertandingan: seberapa yakin ia harus sebelum
mengangkat bendera. Kalau aturannya terlalu longgar, ia sering mengangkat bendera
padahal bola di dalam, dan penonton marah. Kalau aturannya terlalu ketat, ia hampir
tidak pernah mengangkat bendera, dan pelanggaran besar lolos begitu saja. Kedua
kesalahan itu berbeda, dan ia tidak bisa menekan keduanya sekaligus hanya dengan
kehati-hatian.

Di bab ini kita akan tahu mengapa ``tidak ada bukti'' bukan sama dengan ``tidak ada
pelanggaran'', mengapa jumlah pengamatan berperan sebesar aturan yang dipilih, dan
mengapa wasit yang tidak pernah mengangkat bendera bukan wasit yang lebih teliti.
\end{untukanak}

\section{Contoh hitung: 33 dari 50 dan 53 dari 100}
\tingkat{menengah}

Kita kembali ke dua pendeteksi pada pembuka bab dan menghitung kedua selangnya sampai
ujung. Semua angka di bawah memakai $z=1{,}96$, sehingga $z^2=3{,}8416$.

Untuk A, $\hat p=0{,}66$ dan $n=50$. Suku pertama di dalam akar Wilson adalah
$0{,}66\cdot0{,}34/50=0{,}2244/50=0{,}004488$; suku kedua
$z^2/(4n^2)=3{,}8416/10000=0{,}00038416$; jumlahnya $0{,}00487216$ dan akarnya
$0{,}069801$. Titik tengahnya
\[
\frac{0{,}66+3{,}8416/100}{1+3{,}8416/50}
=\frac{0{,}698416}{1{,}076832}=0{,}648584,
\]
dan setengah lebarnya $1{,}96\cdot0{,}069801/1{,}076832=0{,}127048$. Jadi selang
Wilson untuk A adalah $(0{,}5215;\ 0{,}7756)$, selebar $0{,}2541$. Selang Wald untuk
data yang sama memakai galat baku $0{,}066993$ dan setengah lebar
$1{,}96\cdot0{,}066993=0{,}131303$, sehingga $(0{,}5287;\ 0{,}7913)$ selebar
$0{,}2626$.

Untuk B, $\hat p=0{,}53$ dan $n=100$. Suku pertama
$0{,}53\cdot0{,}47/100=0{,}002491$, suku kedua $3{,}8416/40000=0{,}00009604$,
jumlahnya $0{,}00258704$ dan akarnya $0{,}050863$. Titik tengahnya
$(0{,}53+0{,}019208)/1{,}038416=0{,}528890$, setengah lebarnya
$1{,}96\cdot0{,}050863/1{,}038416=0{,}096003$, jadi $(0{,}4329;\ 0{,}6249)$ selebar
$0{,}1920$. Wald memberi $(0{,}4322;\ 0{,}6278)$ selebar $0{,}1956$.

Sekarang bandingkan. Selang A dan B bertumpuk pada $(0{,}5215;\ 0{,}6249)$, seluas
$0{,}1034$: sekitar $41$ persen panjang A dan $54$ persen panjang B. Uji dua proporsi
berbasis normal memberi $z=1{,}5175$ dan nilai-p dua sisi $0{,}1291$. Meskipun angka
mentahnya berbeda tiga belas poin, data ini tidak cukup menyimpulkan bahwa A lebih
baik; selang A juga lebih lebar, sebab dibangun dari lima puluh contoh.

\kertas{Hitung ulang bagian A tanpa melihat teks: $0{,}66\cdot0{,}34$, bagi $50$,
tambahkan $3{,}8416/10000$, tarik akarnya, lalu bagi dengan $1{,}076832$ sesudah
dikali $1{,}96$. Cocokkan bahwa hasilnya $0{,}127048$. Lalu jawab: berapa banyak
contoh yang diperlukan agar lebar selang A turun menjadi $0{,}15$, bila $\hat p$
tetap $0{,}66$? Petunjuk: lebar sebanding $1/\sqrt n$, dan Anda sudah tahu lebarnya
pada $n=50$.}

\section{Lab komputasi: cakupan, galat tipe I, dan daya}
\tingkat{lanjut}

Seluruh klaim di atas dapat diperiksa dengan komputer, dan ada satu klaim yang hanya
dapat diperiksa dengan cara itu: cakupan sebenarnya dari sebuah prosedur. Kuncinya
adalah mengetahui $p$ yang benar, sesuatu yang tidak pernah terjadi pada data nyata.
Karena itu kita bekerja dengan simulasi dan dengan hitungan eksak atas binomial.

\begin{lab}
\begin{verbatim}
import numpy as np
from math import sqrt
from scipy.stats import norm, binom

z = norm.ppf(0.975)


def wald(k, n):
    p = k / n
    se = sqrt(p * (1 - p) / n)
    return p - z * se, p + z * se


def wilson(k, n):
    p = k / n
    z2 = z * z
    mid = (p + z2 / (2 * n)) / (1 + z2 / n)
    half = z * sqrt(p * (1 - p) / n + z2 / (4 * n * n)) / (1 + z2 / n)
    return mid - half, mid + half


# 1. lebar interval 95% untuk beberapa proporsi
print("== 1. lebar interval 95% ==")
print("   n    k   phat   lebar Wald   lebar Wilson")
for n in (20, 100):
    for p in (0.00, 0.05, 0.25, 0.50, 0.66):
        k = round(p * n)
        a = wald(k, n)
        b = wilson(k, n)
        print("%4d %4d  %5.2f   %9.4f    %9.4f"
              % (n, k, k / n, a[1] - a[0], b[1] - b[0]))

# 2. cakupan eksak pada p = 0,05
print("== 2. cakupan eksak, p=0.05 ==")
for n in (20, 50, 100):
    for nama, f in (("Wald", wald), ("Wilson", wilson)):
        c = 0.0
        for k in range(n + 1):
            lo, hi = f(k, n)
            if lo <= 0.05 <= hi:
                c += binom.pmf(k, n, 0.05)
        print("n=%3d %-6s cakupan=%.4f" % (n, nama, c))

# 3. cakupan simulasi dengan seed tetap
rng = np.random.default_rng(20261024)
R = 20000
print("== 3. cakupan simulasi, R=%d ==" % R)
for n in (20, 100):
    X = rng.binomial(n, 0.05, R)
    for nama, f in (("Wald", wald), ("Wilson", wilson)):
        ok = 0
        w = 0.0
        for k in X:
            lo, hi = f(int(k), n)
            ok += (lo <= 0.05 <= hi)
            w += hi - lo
        print("n=%3d %-6s cakupan=%.4f  lebar=%.4f" % (n, nama, ok / R, w / R))

# 4. galat tipe I dan daya untuk H0: p=0.5
print("== 4. galat tipe I dan daya, H0 p=0.5 ==")


def tolak(k, n, p0=0.5):
    return abs((k / n - p0) / sqrt(p0 * (1 - p0) / n)) > z


for n in (20, 100):
    rej = [k for k in range(n + 1) if tolak(k, n)]
    lo = max(k for k in rej if k < n / 2)
    hi = min(k for k in rej if k > n / 2)
    print("n=%3d tolak bila k<=%d atau k>=%d" % (n, lo, hi))
    print("      ukuran eksak di p=0.5 : %.4f"
          % sum(binom.pmf(k, n, 0.5) for k in rej))
    for p1 in (0.65, 0.75):
        ex = sum(binom.pmf(k, n, p1) for k in rej)
        zc = (p1 - 0.5) / sqrt(0.25 / n)
        ap = norm.cdf(zc - z) + norm.cdf(-zc - z)
        print("      daya di p=%.2f: eksak %.4f  aproksimasi %.4f"
              % (p1, ex, ap))

# 4b. simulasi tingkat galat tipe I ketika H0 benar
print("== 4b. simulasi galat tipe I, R=%d ==" % R)
for n in (20, 100):
    X0 = rng.binomial(n, 0.5, R)
    tolak0 = np.abs((X0 / n - 0.5) / sqrt(0.25 / n)) > z
    print("n=%3d p=0.5: tolak=%.4f" % (n, tolak0.mean()))

# 5. interval-t untuk 12 pengukuran
x = np.array([18.2, 19.1, 17.8, 20.3, 18.9, 19.7,
              18.4, 21.0, 17.5, 19.9, 18.6, 20.1])
n = len(x)
m = x.mean()
s = x.std(ddof=1)
from scipy.stats import t
tq = t.ppf(0.975, n - 1)
h = tq * s / sqrt(n)
print("== 5. interval-t, n=12 ==")
print("rata-rata=%.6f  s=%.6f  galat baku=%.6f" % (m, s, s / sqrt(n)))
print("t=%.6f  setengah lebar=%.6f  IK=(%.6f, %.6f)" % (tq, h, m - h, m + h))
print("setengah lebar dengan z=1.96: %.6f" % (z * s / sqrt(n)))
\end{verbatim}
\end{lab}

Keluaran yang sebenarnya, apa adanya:

\begin{verbatim}
== 1. lebar interval 95% ==
   n    k   phat   lebar Wald   lebar Wilson
  20    0   0.00      0.0000       0.1611
  20    1   0.05      0.1910       0.2272
  20    5   0.25      0.3795       0.3568
  20   10   0.50      0.4383       0.4014
  20   13   0.65      0.4181       0.3860
 100    0   0.00      0.0000       0.0370
 100    5   0.05      0.0854       0.0902
 100   25   0.25      0.1697       0.1676
 100   50   0.50      0.1960       0.1923
 100   66   0.66      0.1857       0.1826
== 2. cakupan eksak, p=0.05 ==
n= 20 Wald   cakupan=0.6389
n= 20 Wilson cakupan=0.9245
n= 50 Wald   cakupan=0.9199
n= 50 Wilson cakupan=0.9622
n=100 Wald   cakupan=0.8775
n=100 Wilson cakupan=0.9659
== 3. cakupan simulasi, R=20000 ==
n= 20 Wald   cakupan=0.6430  lebar=0.1472
n= 20 Wilson cakupan=0.9224  lebar=0.2193
n=100 Wald   cakupan=0.8783  lebar=0.0829
n=100 Wilson cakupan=0.9667  lebar=0.0884
== 4. galat tipe I dan daya, H0 p=0.5 ==
n= 20 tolak bila k<=5 atau k>=15
      ukuran eksak di p=0.5 : 0.0414
      daya di p=0.65: eksak 0.2457  aproksimasi 0.2687
      daya di p=0.75: eksak 0.6172  aproksimasi 0.6088
n=100 tolak bila k<=40 atau k>=60
      ukuran eksak di p=0.5 : 0.0569
      daya di p=0.65: eksak 0.8750  aproksimasi 0.8508
      daya di p=0.75: eksak 0.9997  aproksimasi 0.9988
== 4b. simulasi galat tipe I, R=20000 ==
n= 20 p=0.5: tolak=0.0430
n=100 p=0.5: tolak=0.0554
== 5. interval-t, n=12 ==
rata-rata=19.125000  s=1.082191  galat baku=0.312401
t=2.200985  setengah lebar=0.687591  IK=(18.437409, 19.812591)
setengah lebar dengan z=1.96: 0.612296
\end{verbatim}

Empat hal perlu dibaca dari keluaran itu. Pertama, bagian 2 dan 3 hampir sama sampai
tiga angka, dan kesamaan itu sendiri adalah pemeriksaan, sebab simulasi dengan
$R=20000$ ulangan hanya dapat mendekati hitungan eksak sampai sekitar
$\sqrt{0{,}25/20000}=0{,}0035$. Pada $p=0{,}05$ dan $n=20$ selisihnya $0{,}0041$ untuk
Wald dan $0{,}0021$ untuk Wilson. Simulasi yang jujur harus melaporkan $R$.

Kedua, cakupan Wald pada $n=20$ dan $p=0{,}05$ adalah $0{,}6389$, bukan $0{,}95$.
Penyebabnya dapat dihitung dengan tangan: peluang mengamati nol kejadian $0{,}3585$,
dan pada $k=0$ selang Wald adalah $[0;0]$, yang tidak memuat $0{,}05$; satu titik itu
saja menghabiskan hampir tiga puluh enam poin cakupan. Wilson hanya gagal bila
$k\ge3$, peluangnya $0{,}0755$, dan kegagalannya tidak total karena selangnya tetap
lebar. Itulah sebabnya ia mencapai $0{,}9245$. Perhatikan juga kolom lebar: rata-rata
Wald $0{,}1472$ lebih kecil daripada Wilson $0{,}2193$, dan justru di situ intinya.
Selang yang lebih sempit bukan selang yang lebih baik bila ia sempit karena salah
tempat.

Ketiga, bagian 4 memisahkan ukuran uji dari taraf nominal: $0{,}0414$ pada $n=20$ dan
$0{,}0569$ pada $n=100$, keduanya untuk aturan $|z|>1{,}96$; bagian 4b mengulangnya
dengan simulasi dan mendarat dekat, $0{,}0430$ dan $0{,}0554$. Daya bergerak jauh lebih
drastis: $0{,}2457$ pada $n=20$ dan $0{,}8750$ pada $n=100$ untuk $p_1=0{,}65$.
Aproksimasi normal tidak selalu di sisi yang aman: pada $n=20$ ia melaporkan $0{,}2687$
sedangkan kebenarannya $0{,}2457$.

Keempat, bagian 5 hanya mengulang interval-$t$ yang sudah dihitung di teks, dengan
selisih setengah lebar $t$ dan $z$ sebesar $0{,}075295$.

Program tidak membuktikan satu pun dari semua itu. Ia memeriksa klaim yang sudah
dibuktikan, dan pada dua tempat ia menunjukkan batas dirinya sendiri: cakupan eksak
memakai distribusi yang kita andaikan benar, sehingga bila modelnya salah, angka
$0{,}9245$ itu tidak berarti apa-apa; dan simulasi hanya mengukur seberapa sering
prosedurnya berhasil pada model itu, bukan seberapa benar modelnya.

\begin{untukdosen}
Urutan mengajar yang saya sarankan: mulai dari rapat pagi pada pembuka bab, bukan dari
definisi penaksir. Tulis $33/50$ dan $53/100$ di papan, minta kelas memilih, lalu
tanyakan apa yang terjadi bila lima puluh contoh itu diganti. Dari kegelisahan itu
definisi penaksir, bias, dan galat baku muncul sebagai jawaban. Dekomposisi MSE
dibuktikan di papan dengan satu baris aljabar. Tafsiran selang kepercayaan masuk
sesudahnya, sebagai bahan kuliah tersendiri, bukan catatan kaki.

Dua miskonsepsi yang paling sering muncul. Pertama, mahasiswa membaca ``$95$ persen''
sebagai peluang bersyarat pada data yang sudah ada, lalu menulis ``peluang $\mu$ ada di
selang ini $0{,}95$''; gejalanya tampak ketika mereka tidak dapat menjelaskan apa yang
acak pada selang yang sudah dihitung. Kedua, mereka menganggap nilai-p besar sebagai
bukti $H_0$ benar, dan akan menghitung nilai-p untuk $10/20$ lalu menyimpulkan bahwa
$p=0{,}5$ terbukti.

Aktivitas papan tulis yang berhasil: bagi kelas menjadi dua kelompok, satu menghitung
selang Wald dan satu Wilson untuk $1/20$, lalu suruh keduanya mempertahankan hasilnya
di depan kelas. Kelompok Wald akan menemukan batas bawah negatif dan harus
menjelaskannya. Tanda bahwa kelas siap lanjut: mereka dapat menyebut satu contoh
ukuran sampel yang membuat nilai-p besar padahal efeknya besar, dan dapat menjelaskan
mengapa hal itu terjadi.
\end{untukdosen}

\section{Ringkasan}
\tingkat{dasar}

\begin{ringkasan}
\begin{itemize}
  \item Penaksir adalah statistik yang menebak parameter, jadi ia variabel acak dengan
        harapan dan variansi. Mutunya diukur oleh
        $\operatorname{MSE}=\operatorname{bias}^2+\operatorname{Var}$; tak bias saja
        tidak menjamin apa pun, dan MSE menuju nol menjamin konsistensi lewat
        Chebyshev.
  \item Galat baku adalah akar variansi penaksir, bukan lebar data. Untuk proporsi
        nilainya $\sqrt{\hat p(1-\hat p)/n}$, mengecil sebanding $1/\sqrt n$,
        sehingga empat kali data hanya memangkasnya menjadi separuh.
  \item Interval kepercayaan $95$ persen adalah pernyataan tentang prosedur, bukan
        tentang parameter: prosedur itu menangkap nilai sebenarnya pada $95$ persen
        percobaan yang diulang. Sesudah data diambil, selang itu tetap dan tidak lagi
        berpeluang.
  \item Interval Wald runtuh di dekat $0$ dan $1$: pada $0/20$ lebarnya nol, dan pada
        $p=0{,}05$ dengan $n=20$ cakupannya hanya $0{,}6389$. Interval Wilson, hasil
        membalik uji skor, mencapai $0{,}9245$ pada setelan yang sama.
  \item Bila tidak ada kejadian dalam $n$ percobaan, batas atas $95$ persen kira-kira
        $3/n$; untuk $n=100$ batas eksaknya $0{,}0295130$ dan Wilson dua sisi
        $0{,}0369935$.
  \item Nilai-p adalah peluang data sejauh itu bila $H_0$ benar, dan di bawah $H_0$
        ia tersebar seragam. Nilai-p besar bukan bukti $H_0$ benar: pada $n=20$ daya
        terhadap $p=0{,}75$ hanya $0{,}6172$.
\end{itemize}
\end{ringkasan}

\section{Latihan}
\tingkat{menengah}

\begin{latihan}
\begin{enumerate}
  \item Nilai klaim ini: ``Karena nilai-p sama dengan $0{,}42$, hipotesis nol
        $p=0{,}5$ terbukti benar.'' Benar, salah, atau belum cukup bukti? Sebutkan
        bagian yang benar, bagian yang keliru, dan satu perhitungan tambahan yang
        akan memperkuat jawaban Anda.
  \item Buktikan bahwa $\mathbb{E}[\hat p(1-\hat p)]=p(1-p)(n-1)/n$ untuk
        $\hat p=K/n$ dengan $K\sim\mathrm{Bin}(n,p)$, lalu tentukan penaksir tak bias
        bagi $p(1-p)$ yang diperoleh dari situ. Periksa hasilnya pada $p=0{,}3$ dan
        $n=20$.
  \item Buktikan bahwa bila $\operatorname{bias}(\hat\theta_n)\to0$ dan
        $\operatorname{Var}(\hat\theta_n)\to0$, maka $\hat\theta_n$ konsisten untuk
        $\theta$. Jelaskan mengapa syarat itu lebih lemah daripada
        $\operatorname{MSE}\to0$ tetapi cukup untuk kesimpulan yang sama.
  \item Hitung selang Wilson dan Wald $95$ persen untuk $7/40$, lalu bandingkan
        lebarnya dengan $13/20$. Jelaskan mengapa lebar yang lebih kecil pada $7/40$
        bukan tanda ketelitian yang lebih tinggi, dan periksa kedua selang dengan
        Python.
  \item Untuk $n=30$ dan $H_0\colon p=0{,}5$ dengan uji dua sisi berbasis normal,
        tentukan daerah tolaknya, hitung ukuran eksaknya, lalu hitung dayanya pada
        $p_1=0{,}7$. Bandingkan daya itu dengan $n=100$ pada nilai $p_1$ yang sama.
  \item Buktikan bahwa selang Wald dan uji-$z$ dua sisi yang memakai galat baku
        $\hat p$ memang setara, dengan menuliskan kedua ketaksamaannya berdampingan.
        Jelaskan mengapa kesetaraan itu tidak berlaku bila uji memakai galat baku di
        bawah $H_0$.
  \item Sebuah laporan menulis: ``tidak ada perbedaan yang signifikan, $p=0{,}31$,
        jadi kedua metode sama baik''. Tulis ulang kalimat itu sehingga jujur, dengan
        menyebut selang, taraf, dan daya. Anda boleh mengarang angkanya asalkan
        konsisten.
  \item Untuk $n=100$ tanpa satu pun kejadian, hitung aturan tiga, batas eksak satu
        sisi, dan batas atas Wilson. Jelaskan mengapa ketiga angka berbeda, lalu
        sebutkan mana yang boleh dipakai untuk kalimat ``paling tinggi $5$ persen''.
  \item Hitung selang-$t$ $90$ persen untuk dua belas data waktu tunda pada bab ini,
        lalu bandingkan dengan selang $95$ persen. Berapa banyak data yang diperlukan
        agar lebar selang $95$ persen turun menjadi separuh, bila $s$ tidak berubah?
\end{enumerate}

Petunjuk untuk tiga butir tersulit. Butir 3: tulis
$\hat\theta_n-\theta=\operatorname{bias}_n+(\hat\theta_n-\mathbb{E}\hat\theta_n)$,
lalu pakai ketaksamaan segitiga dan Chebyshev pada suku kedua. Butir 6: ketaksamaan
pertama adalah $|\hat p-p_0|\le z\sqrt{\hat p(1-\hat p)/n}$, dan uji yang memakai
$\sqrt{p_0(1-p_0)/n}$ memakai penyebut berbeda, sehingga himpunan $p_0$ yang tidak
ditolak bukan lagi selang Wald melainkan himpunan yang lebih rumit. Butir 5: daerah
tolaknya adalah $|k-15|>1{,}96\sqrt{30\cdot0{,}25}$, jadi $k\le9$ atau $k\ge21$, dan
dayanya dihitung dengan menjumlahkan pmf binomial pada $p=0{,}7$ di kedua ekor itu.
\end{latihan}

\begin{untukpeneliti}
Yang benar-benar terbukti di bab ini: dekomposisi MSE; kekonsistenan dari MSE menuju
nol; cakupan tepat interval-$z$ pada model normal; bentuk tertutup interval Wilson dan
letaknya di $[0,1]$; aturan tiga sebagai batas eksak satu sisi beserta
ketaksamaan $p_{\mathrm{hi}}<\ln 20/n$; keseragaman nilai-p di bawah $H_0$ untuk
statistik kontinu; dan kesetaraan selang dengan uji pada pasangan yang disebutkan.
Yang hanya populer, bukan teorema: bahwa $95$ persen berarti peluang bersyarat pada
data yang sudah teramati; bahwa nilai-p besar mendukung $H_0$; dan bahwa lebar selang
yang lebih kecil selalu lebih baik, padahal perbandingan lebar tanpa cakupan tidak
berarti. Cakupan Wilson di bawah nominal pada beberapa $(n,p)$ juga fakta hitung, bukan
teorema.

Sumber primer: Wasserman, \emph{All of Statistics}, Bab 6 (Models, Statistical
Inference, and Learning) untuk bias, MSE, dan tafsiran selang; Bab 9 (Parametric
Inference) untuk interval dan metode plug-in; Bab 10 (Hypothesis Testing and p-values)
untuk nilai-p, galat tipe I dan II, daya, serta hubungan uji dengan selang. Interval
Wilson berasal dari E.~B. Wilson, ``Probable Inference, the Law of Succession, and
Statistical Inference'' (1927); aturan tiga dari B.~D. Jovanovic dan P.~S. Levy,
``A Look at the Rule of Three'' (1997). Keduanya tidak dibahas Wasserman, jadi
rujukannya disebut langsung, bukan lewat buku teks perantara.

Batas keberlakuan: seluruh bab bekerja dengan sampel bebas dan berdistribusi sama,
satu parameter, dan aproksimasi normal atau binomial yang diketahui. Interval-$z$
hanya tepat pada model normal dengan $\sigma$ diketahui; di luar itu ia pernyataan
asimtotik, dan pada $n$ kecil dengan distribusi berekor berat cakupannya dapat jauh di
bawah nominal. Interval Wilson adalah aproksimasi skor, bukan selang eksak; selang
eksak untuk proporsi dihitung dari sebaran binomial dan lebih lebar. Semua cakupan yang
dilaporkan bersyarat pada model binomial yang benar.

Pertanyaan terbuka yang layak dikerjakan pembaca: seberapa baik interval proporsi
berperilaku ketika $n$ bukan bilangan tetap melainkan berhenti secara adaptif, seperti
pada pengambilan sampel berurutan; dan bagaimana menyusun selang untuk metrik
klasifikasi yang dihitung dari data dependen, misalnya ketika satu contoh muncul di
beberapa lipatan validasi silang. Keduanya masih menjadi bahan penelitian.
\end{untukpeneliti}

\peta{Bab 20 pada Wasserman, \emph{All of Statistics}: Bab 6 (Models, Statistical
Inference, and Learning: penaksir, bias, MSE, galat baku, tafsiran selang), Bab 9
(Parametric Inference: interval kepercayaan dan metode plug-in), Bab 10 (Hypothesis
Testing and p-values: nilai-p, galat tipe I dan II, daya, kesetaraan uji dan selang).
Interval Wilson dan aturan tiga diambil dari makalah aslinya, Wilson (1927) dan
Jovanovic--Levy (1997), karena Wasserman tidak membahas keduanya. Hukum bilangan besar,
teorema limit pusat, dan ketaksamaan Chebyshev yang dipakai di sini dibuktikan di
Bab 18; peluang bersyarat dan kekeliruan pembalik syarat dibahas di Bab 16. Bab 25
melanjutkan bab ini untuk metrik klasifikasi, termasuk aturan tiga pada pelaporan
evaluasi model.}

\chapter{Regresi dan Statistika Terapan}

Sebuah kafe kecil mencatat sembilan hari dagang. Dua angka ditulis setiap malam:
banyak pesanan daring dalam puluhan, dan jam lembur dapur. Pemilik kafe membutuhkan
satu angka untuk satu pertanyaan: bila pesanan bertambah sepuluh, berapa tambahan jam
lembur yang harus ia siapkan besok?

\begin{center}
\small
\begin{tabular}{l*{9}{r}}
\hline
hari               & 1 & 2 & 3 & 4 & 5 & 6 & 7 & 8 & 9 \\
pesanan (puluhan)  & 1 & 2 & 3 & 4 & 5 & 6 & 7 & 8 & 9 \\
lembur (jam)       & 10 & 12 & 13 & 14 & 16 & 18 & 19 & 20 & 22 \\
\hline
\end{tabular}
\end{center}

Jawaban cepat tersedia tanpa alat apa pun: dari hari pertama ke hari terakhir pesanan
naik delapan puluhan dan lembur naik dua belas jam, jadi kira-kira $1{,}5$ jam per
sepuluh pesanan. Kuadrat terkecil, yang akan kita susun di bab ini, memberi
$1{,}466667$. Kedua angka itu dekat, dan justru karena dekat kita tergoda menganggap
persoalannya selesai. Tiga pertanyaan penting tidak terjawab oleh keduanya. Seberapa
jauh angka itu bergeser bila sembilan hari itu diganti sembilan hari lain? Bila besok
pesanannya delapan puluh, berapa jam yang harus disiapkan? Dan bila satu hari ternyata
salah catat, berapa besar kerusakannya?

Bab 10 sudah memberi gambarannya: kuadrat terkecil adalah proyeksi ortogonal vektor
pengamatan ke ruang kolom matriks rancangan, dan persamaan normal menyatakan bahwa
sisaannya tegak lurus ruang itu. Yang belum ada di sana adalah ketidakpastian. Bab 20
sudah memberi selang kepercayaan dan uji hipotesis, tetapi untuk satu parameter dari
pengamatan yang saling bebas dan berbagi satu sebaran. Di sini parameter yang dicari
adalah kemiringan, nisbah dua jumlah yang keduanya acak, dan pengamatan tidak lagi
sederajat: titik di ujung rentang lebih berpengaruh daripada titik di tengah. Bab ini
menyambung keduanya dan berhenti di tempat yang jujur: seberapa jauh kesimpulan boleh
dibawa.

\begin{untukanak}
Seorang nenek menjual es kelapa di depan rumah. Ia tidak pernah mencatat apa pun
dengan rapi, tetapi ingatannya cukup: hari yang panas, dagangannya habis sebelum
sore. Hari yang mendung, separuh es mencair di termos. Suatu hari cucunya datang
dengan buku tulis dan mencatat Sembilan hari berturut-turut: berapa kira-kira panasnya,
dan berapa gelas yang terjual.

Dari sembilan catatan itu cucunya menarik satu garis lurus: makin panas, makin banyak
yang terjual. Neneknya mengangguk, tetapi lalu bertanya, ``kalau besok lebih panas
sedikit, berapa gelas yang harus saya siapkan?'' Cucunya terdiam. Garis itu mengatakan
arahnya, bukan angkanya. Lagi pula, satu hari pernah hujan deras padahal panas, dan
penjualannya jelek sekali; titik itu jauh dari garis.

Di bab ini kita akan tahu mengapa garis yang paling pas tetap bukan ramalan yang pasti,
mengapa satu hari yang aneh dapat menggeser garis lebih jauh daripada lima hari biasa,
dan mengapa orang yang menjual angka harus menyebutkan seberapa lebar ketidakpastiannya.
\end{untukanak}

\section{Sembilan hari dan satu garis}
\tingkat{dasar}

Sembilan pasang angka, dua bilangan yang dicari, sembilan persamaan. Sistem seperti ini
sudah kita temui di Bab 8: lebih banyak persamaan daripada bilangan yang tidak diketahui,
sehingga solusi tepat hanya ada bila kesembilan titik itu kebetulan segaris. Tidak ada
garis yang lewat persis di semua titik di atas, dan itu bukan kegagalan pengukuran,
melainkan keadaan biasa. Yang bisa dicari hanyalah garis yang paling dekat, dan
``paling dekat'' harus didefinisikan lebih dahulu.

Kalimat sehari-harinya: kita mencari satu garis yang jumlah kesalahannya sekecil mungkin,
dengan kesalahan dihitung sebagai jarak tegak dari titik ke garis.

\kamus{regresi}{regression}{menarik satu garis yang paling dekat ke semua titik, bukan
garis yang lewat di semua titik}

\begin{definisi}[Model regresi linear sederhana]
\index{regresi}
Untuk pengamatan $i=1,\dots,n$,
\[
y_i = a + b\,x_i + \varepsilon_i ,
\]
dengan $x_i$ peubah penjelas yang teramati, $y_i$ peubah yang dijelaskan, $a$ titik
potong, $b$ kemiringan, dan $\varepsilon_i$ galat yang tidak teramati. Model disebut
\emph{linear} karena linear dalam parameter $a$ dan $b$, bukan karena garisnya lurus
dalam $x$.
\end{definisi}

Satu perbedaan kecil yang menentukan seluruh bab: $\varepsilon_i$ dan sisaan bukan
benda yang sama. Galat $\varepsilon_i$ adalah selisih antara nilai sebenarnya dan garis
populasi yang tidak pernah kita lihat; sisaan $e_i=y_i-\hat y_i$ adalah selisih antara
nilai teramati dan garis yang kita pasang dari data itu sendiri. Sisaan dapat dihitung
dan digambar, galat hanya dapat diasumsikan. Setiap kali sebuah buku menyebut asumsi
galat lalu memeriksanya lewat gambar sisaan, ada satu lompatan yang tidak boleh
dilupakan: yang diperiksa adalah pengganti, bukan barang aslinya.

\section{Kuadrat terkecil sebagai proyeksi}
\tingkat{menengah}

Susun $y\in\mathbb{R}^n$ sebagai vektor pengamatan dan $X$ matriks $n\times2$ yang
kolomnya satu dan $x$. Mencari garis terbaik berarti memilih $\hat y$ di ruang kolom $X$
yang paling dekat ke $y$: persoalan Bab 10, dengan tambahan bahwa kita butuh bentuk
eksplisit koefisiennya untuk menghitung ketidakpastiannya.

\begin{definisi}[Penaksir kuadrat terkecil]
\index{kuadrat terkecil}
$\hat\beta=(\hat a,\hat b)$ adalah penaksir kuadrat terkecil bila ia meminimumkan
\[
\operatorname{RSS}(\beta)=\|y-X\beta\|^2=\sum_{i=1}^{n}\bigl(y_i-a-bx_i\bigr)^2 .
\]
\end{definisi}

\kamus{kuadrat terkecil}{least squares}{menjumlahkan semua kesalahan setelah
dikuadratkan, lalu memilih garis yang jumlahnya paling kecil}

\begin{teorema}[Persamaan normal]
\index{persamaan normal}
$\hat\beta$ meminimumkan RSS jika dan hanya jika
\[
X^{\top}(y-X\hat\beta)=0 ,
\qquad\text{yaitu}\qquad X^{\top}X\hat\beta=X^{\top}y .
\]
Bila kolom $X$ bebas linear, $\hat\beta=(X^{\top}X)^{-1}X^{\top}y$ tunggal.
\end{teorema}

\begin{proof}[Bukti]
Tulis $f(\beta)=\|y-X\beta\|^2=y^{\top}y-2\beta^{\top}X^{\top}y+\beta^{\top}X^{\top}X\beta$.
Gradiennya $-2X^{\top}y+2X^{\top}X\beta$, jadi syarat stasioner tepat persamaan normal.
Karena $X^{\top}X$ definit positif bila kolomnya bebas linear, $f$ konveks ketat dan
titik stasioner itu minimum global.

Arah sebaliknya lebih informatif. Andaikan $X^{\top}(y-X\hat\beta)=0$. Untuk sebarang
$\beta$,
\[
\|y-X\beta\|^2=\|y-X\hat\beta\|^2+\|X(\hat\beta-\beta)\|^2
+2(\hat\beta-\beta)^{\top}X^{\top}(y-X\hat\beta),
\]
dan suku silang lenyap. Jadi $\operatorname{RSS}(\beta)\ge\operatorname{RSS}(\hat\beta)$
untuk semua $\beta$, dengan kesamaan hanya bila $X\beta=X\hat\beta$.
\end{proof}

Pembuktian itu memuat seluruh isi geometrisnya: sisaan tegak lurus setiap kolom $X$.
Karena kolom pertama vektor satu dan kolom kedua $x$, berlaku $\sum_i e_i=0$ dan
$\sum_i x_ie_i=0$, dan kedua kesamaan itulah yang menopang dekomposisi jumlah kuadrat
di bagian berikutnya.

\begin{proposisi}[Matriks topi]
\index{matriks topi}
Dengan $H=X(X^{\top}X)^{-1}X^{\top}$, berlaku $\hat y=Hy$, $H^{\top}=H$, $H^2=H$,
$\operatorname{tr}H=2$, dan $0\le h_{ii}\le1$ dengan $\sum_i h_{ii}=2$. Nilai $h_{ii}$
disebut \emph{leverage} titik $i$.
\end{proposisi}

\begin{proof}[Bukti]
Simetri dan idempotensi diperoleh langsung dari definisi:
\[
H^2=X(X^{\top}X)^{-1}X^{\top}X(X^{\top}X)^{-1}X^{\top}=H .
\]
Karena itu $H$ adalah matriks proyeksi ortogonal,
sehingga $h_{ii}=e_i^{\top}He_i=\|He_i\|^2\ge0$ dan $h_{ii}\le\|e_i\|^2=1$. Terakhir
$\operatorname{tr}H=\operatorname{tr}\bigl((X^{\top}X)^{-1}X^{\top}X\bigr)
=\operatorname{tr}I_2=2$.
\end{proof}

Sekarang angkanya, karena seluruh bab bergantung pada satu perhitungan yang dapat
diperiksa ulang. Untuk data kafe,
\[
X^{\top}X=\begin{pmatrix}9&45\\45&285\end{pmatrix},
\qquad X^{\top}y=\begin{pmatrix}144\\808\end{pmatrix},
\qquad \det(X^{\top}X)=9\cdot285-45^2=540 .
\]
Maka
\[
\hat\beta=\frac1{540}\begin{pmatrix}285&-45\\-45&9\end{pmatrix}
\begin{pmatrix}144\\808\end{pmatrix}
=\frac1{540}\begin{pmatrix}4680\\792\end{pmatrix}
=\begin{pmatrix}8{,}666667\\1{,}466667\end{pmatrix}.
\]
Angka $808$ layak diperiksa dengan tangan: $\sum x_iy_i=10+24+39+56+80+108+133+160+198$.
Angka $4680$ juga: $285\cdot144-45\cdot808=41040-36360$.

Satu teorema lagi diperlukan sebelum kita boleh bicara tentang mutu penaksir. Teorema
itu menjelaskan mengapa kuadrat terkecil dipakai di hampir semua bidang, dan sekaligus
mengapa ia bukan pilihan suci.

\begin{teorema}[Gauss--Markov]
\index{gauss-markov}
Andaikan (A1) $y=X\beta+\varepsilon$; (A2) $\mathbb{E}[\varepsilon\mid X]=0$;
(A3) $\operatorname{Var}(\varepsilon\mid X)=\sigma^2I$; (A4) kolom $X$ bebas linear.
Maka $\hat\beta$ tak bias dengan $\operatorname{Var}(\hat\beta)=\sigma^2(X^{\top}X)^{-1}$,
dan di antara semua penaksir linear tak bias, $\hat\beta$ yang bervariansi terkecil.
Penaksir seperti itu disebut \emph{BLUE} (\emph{best linear unbiased estimator}).
\end{teorema}

\begin{proof}[Bukti]
Ketakbiasan: $\mathbb{E}[\hat\beta\mid X]=(X^{\top}X)^{-1}X^{\top}X\beta=\beta$ memakai
(A2). Variansi: $\hat\beta-\beta=(X^{\top}X)^{-1}X^{\top}\varepsilon$, jadi
$\operatorname{Var}(\hat\beta)=(X^{\top}X)^{-1}X^{\top}(\sigma^2I)X(X^{\top}X)^{-1}
=\sigma^2(X^{\top}X)^{-1}$.

Bagian ketiga tidak dibuktikan penuh di sini. Sketsanya: setiap penaksir linear tak
bias berbentuk $\tilde\beta=Cy$ dengan $CX=I$, sehingga $\tilde\beta-\beta=C\varepsilon$
dan $\operatorname{Var}(\tilde\beta)=\sigma^2CC^{\top}$. Menulis $C=(X^{\top}X)^{-1}X^{\top}+D$
memberi $DX=0$ dan
\[
CC^{\top}-(X^{\top}X)^{-1}=DD^{\top},
\]
yang matriks semi-definit positif, sehingga setiap komponen $\tilde\beta$ bervariansi
tak lebih kecil daripada komponen yang sepadan dari $\hat\beta$. Bukti lengkap ada di
Seber dan Lee, \emph{Linear Regression Analysis}, Bab 3, dan di Greene,
\emph{Econometric Analysis}, Bab 3.
\end{proof}

\section{Koefisien dan satuannya}
\tingkat{dasar}

Kemiringan $1{,}466667$ tidak berarti apa pun sebelum satuannya disebut. Satuannya jam
per puluhan pesanan: setiap tambahan sepuluh pesanan menaikkan ramalan lembur sebesar
$1{,}466667$ jam, yaitu $88$ menit, atau $8{,}8$ menit per pesanan. Kalimat itu dapat
dipakai pemilik kafe; kalimat ``koefisiennya $1{,}47$'' tidak.

Titik potong $8{,}666667$ jam adalah ramalan untuk $x=0$, yaitu hari tanpa satu pun
pesanan daring. Tidak ada hari seperti itu dalam data, dan di luar rentang $1$ sampai
$9$ garisnya tidak dijamin berlaku. Menyebut angka itu ``biaya tetap lembur'' adalah
\index{ekstrapolasi}ekstrapolasi yang dibungkus nama: kita memakai garis di luar
tempat ia dipasang. Yang aman: titik potong adalah bilangan yang membuat garis lewat
$(\bar x,\bar y)$, sebab $\hat a=\bar y-\hat b\bar x$; ia berguna untuk perhitungan,
bukan untuk tafsiran.

Perubahan satuan tidak mengubah kesimpulan. Bila $x$ ditulis dalam pesanan, bukan
puluhan, kemiringannya $0{,}146667$ jam per pesanan dan galat bakunya $0{,}004714$,
keduanya terbagi sepuluh; $t$ dan $R^2$ tidak bergerak. Mengalikan satu peubah dengan
konstanta $c$ mengalikan koefisien dan galat bakunya dengan $c$ atau $1/c$, sedangkan
nisbah keduanya tetap, sehingga laporan yang menyebut $t$ tanpa satuan tidak terbaca.

\begin{tangan}
Siapkan kertas berpetak, sembilan koin, dan seutas benang sepanjang sisi kertas. Letakkan
koin pada petak sesuai sembilan pasang angka di awal bab: mendatar pesanan, tegak lembur,
satu petak satu satuan. Ikat ujung benang pada dua ujung kertas sehingga benang dapat
digeser sebagai garis lurus. Geser dan putar benang sampai jumlah jarak tegak semua koin
ke benang terasa paling kecil; jangan mengukur, cukup bandingkan dengan mata, lalu
kunci dengan selotip.

Yang harus terlihat: dua orang akan menaruh benang di tempat yang sedikit berbeda, dan
selisihnya paling besar ketika engkau menggeser ujung benang. Sekarang angkat koin
paling kiri satu petak ke atas, dan perhatikan betapa besar benang harus berputar.
Kembalikan, lalu angkat koin tengah satu petak: hampir tidak ada yang berubah. Itulah
leverage, dan engkau baru saja melihatnya dengan tangan sebelum bertemu rumusnya.
\end{tangan}

\section{R kuadrat dan korelasinya}
\tingkat{menengah}

Sesudah garisnya ada, orang bertanya ``seberapa bagus''. Jawaban yang paling sering
dipakai adalah satu angka yang membandingkan sisaan dengan sebaran asli $y$. Kalimat
sehari-harinya: berapa bagian dari goyangan $y$ yang berhasil dijelaskan garis, dan
berapa bagian yang tertinggal.

\begin{definisi}[$R^2$]
\index{koefisien determinasi}
Dengan $\mathrm{SST}=\sum_i(y_i-\bar y)^2$, $\mathrm{SSE}=\sum_i e_i^2$, dan
$\mathrm{SSR}=\sum_i(\hat y_i-\bar y)^2$, koefisien determinasi adalah
\[
R^2=1-\frac{\mathrm{SSE}}{\mathrm{SST}}=\frac{\mathrm{SSR}}{\mathrm{SST}} .
\]
\end{definisi}

\kamus{koefisien determinasi}{coefficient of determination}{bagian dari goyangan yang
berhasil dijelaskan garis, dinyatakan sebagai pecahan dari keseluruhan}

\begin{teorema}[Dekomposisi jumlah kuadrat]
Bila $X$ memuat kolom satu, maka $\mathrm{SST}=\mathrm{SSR}+\mathrm{SSE}$ dan
$0\le R^2\le1$.
\end{teorema}

\begin{proof}[Bukti]
Tulis $y-\bar y\mathbf{1}=(\hat y-\bar y\mathbf{1})+e$. Karena $\sum e_i=0$, vektor $e$
tegak lurus $\bar y\mathbf{1}$, dan menurut persamaan normal $e$ juga tegak lurus
$\hat y$. Jadi kedua suku di ruas kanan saling tegak lurus, dan teorema Pythagoras
memberi $\|y-\bar y\mathbf{1}\|^2=\|\hat y-\bar y\mathbf{1}\|^2+\|e\|^2$.
\end{proof}

Untuk data kafe, $\mathrm{SST}=130$, $\mathrm{SSE}=0{,}933333$, sehingga
$R^2=1-0{,}933333/130=0{,}992821$. Pada regresi sederhana, $R^2$ sama dengan kuadrat
korelasi: $r=88/\sqrt{60\cdot130}=0{,}996404$ dan $r^2=0{,}992821$. Kesamaan itu
khusus untuk satu peubah penjelas; begitu penjelasnya dua, $R^2$ bukan lagi kuadrat
korelasi apa pun.

Angka $0{,}992821$ mudah membuat orang berhenti berpikir, jadi mari kita jinakkan. Yang
dijelaskan adalah $99{,}28$ persen \emph{goyangan} lembur, dan goyangan itu diukur
terhadap rata-ratanya sendiri. Sisaannya masih bervariansi $s^2=0{,}933333/7=0{,}133333$,
yaitu simpangan baku $s=0{,}365148$ jam atau $21{,}9$ menit. Pemilik kafe yang menyiapkan
lembur atas dasar garis ini masih bisa meleset lebih dari dua puluh menit. $R^2$ bersifat
nisbi; $s$ bersifat mutlak, dan yang dirasakan orang adalah $s$.

\begin{untukanak}
Bayangkan dua siswa. Yang pertama ikut lomba lari seratus meter dan finis $0{,}2$
detik lebih lambat daripada juara. Yang kedua ikut lomba maraton dan finis dua puluh
menit lebih lambat daripada juara. Keduanya kalah, tetapi tidak dalam arti yang sama.
Dua persepuluh detik itu sangat dekat; dua puluh menit itu sangat jauh. Kalau gurunya
hanya menuliskan ``peringkat dua'' untuk keduanya, kita kehilangan seluruh ceritanya.

Angka yang mirip terjadi pada bab ini. Ada angka yang mengukur \emph{bagian} yang
berhasil ditebak, dan ada angka yang mengukur \emph{besar} kesalahan yang masih
tersisa. Yang pertama bisa tampak hampir sempurna sementara yang kedua tetap besar.
Orang yang hanya membaca angka pertama akan menyiapkan terlalu sedikit, lalu bingung
mengapa ramalannya meleset.

Di bab ini kita akan tahu mengapa dua angka harus disebut berpasangan, mengapa angka
bagus tidak sama dengan ramalan yang sempit, dan kapan garis yang rapi sebaiknya
ditinggalkan.
\end{untukanak}

\begin{proposisi}[Menambah peubah tidak pernah menurunkan $R^2$]
Bila kolom $X_2$ ditambahkan pada $X_1$ sehingga ruang kolomnya memuat ruang kolom
$X_1$, maka $\mathrm{SSE}$ tidak bertambah dan $R^2$ tidak berkurang.
\end{proposisi}

\begin{proof}[Bukti]
Himpunan vektor $X_1\beta_1$ termuat dalam himpunan $X_2\beta_2$, jadi minimum jarak
ke $y$ dihimpun atas himpunan yang lebih besar tidak dapat lebih tinggi. Karena itu
$R^2$ tidak dapat turun, meskipun peubah yang ditambahkan tidak berhubungan dengan $y$.
\end{proof}

Akibat praktisnya penting: membandingkan model dengan $R^2$ mentah selalu memihak model
yang lebih gemuk. Koreksi yang lazim adalah
$R^2_{\text{adj}}=1-(1-R^2)(n-1)/(n-p)$ dengan $p$ banyak parameter. Untuk model
berganda di bagian berikutnya, $R^2=0{,}997951$ turun menjadi $0{,}997496$ ketika $p=3$
dan $n=12$; pada $n$ kecil dengan banyak peubah koreksi itu dapat membalik urutan.

\section{Regresi berganda dan kolinearitas}
\tingkat{lanjut}

Sampai sini satu peubah penjelas. Kenyataannya hampir setiap persoalan menuntut lebih
dari satu, dan di situlah muncul penyakit khasnya: dua penjelas yang membawa informasi
hampir sama.

Dua belas cabang sebuah kedai dicatat: penjualan harian $y$ dalam juta rupiah, banyak
karyawan $x_1$, dan luas ruang duduk $x_2$ dalam meter persegi.

\begin{center}
\small
\begin{tabular}{l*{12}{r}}
\hline
cabang   & 1 & 2 & 3 & 4 & 5 & 6 & 7 & 8 & 9 & 10 & 11 & 12 \\
$x_1$    & 6 & 7 & 7 & 8 & 8 & 9 & 9 & 10 & 10 & 11 & 12 & 12 \\
$x_2$    & 18 & 22 & 20 & 26 & 24 & 28 & 26 & 32 & 30 & 34 & 36 & 38 \\
$y$      & 11,2 & 12,4 & 12,0 & 14,1 & 13,6 & 15,2 & 14,8 & 17,0 & 16,4 & 18,1 & 19,4 & 19,9 \\
\hline
\end{tabular}
\end{center}

Kuadrat terkecil memberi
\[
\hat y=2{,}442991+0{,}686916\,x_1+0{,}239252\,x_2,
\]
dengan galat baku $0{,}207583$, $0{,}135366$, dan $0{,}042101$; nilai $t$ berturut-turut
$11{,}7687$, $5{,}0745$, dan $5{,}6828$ pada $9$ derajat kebebasan; $s=0{,}143607$ dan
$R^2=0{,}997951$. Semua tampak sehat. Lalu kita pasang satu penjelas saja: $y$ atas $x_1$
memberi $1{,}446019$, $y$ atas $x_2$ memberi $0{,}450075$, keduanya hampir dua kali
koefisien bergandanya.

Bukan karena salah hitung. Karyawan dan luas ruang hampir saling menentukan: korelasi
keduanya $r_{12}=0{,}986801$. Setelah $x_2$ masuk, $x_1$ hampir tidak membawa
keterangan baru, jadi koefisiennya menyusut menjadi efek tambahannya saja, yaitu
$0{,}687$ juta per karyawan pada luas yang tetap. Koefisien tunggal $1{,}446$ mencampur
efek karyawan dengan efek luas yang biasanya menyertainya, dan jurang antara keduanya
adalah ukuran seberapa jauh kedua peubah itu bertumpang.

\begin{definisi}[Faktor inflasi varians]
\index{varians inflasi}
Untuk peubah penjelas ke-$j$, misalkan $R_j^2$ adalah koefisien determinasi regresi
$x_j$ atas seluruh penjelas lain. Faktor inflasi variansnya
\[
\mathrm{VIF}_j=\frac{1}{1-R_j^2}.
\]
\end{definisi}

\kamus{varians inflasi}{variance inflation factor}{seberapa banyak galat baku sebuah
koefisien digelembungkan oleh penjelas lain yang membawa informasi serupa}

\begin{proposisi}[VIF dua penjelas]
\index{kolinearitas}
Bila hanya ada dua penjelas, $\operatorname{Var}(\hat b_1)=\sigma^2/\bigl(S_{11}
(1-r_{12}^2)\bigr)$, dengan $S_{11}=\sum_i(x_{i1}-\bar x_1)^2$. Jadi galat baku
$\hat b_1$ tergelembung tepat sebesar $\sqrt{\mathrm{VIF}_1}$ dibanding keadaan penjelas
yang saling ortogonal.
\end{proposisi}

\begin{proof}[Bukti]
Untuk dua penjelas, invers matriks $2\times2$ dari bagian terpusat $X^{\top}X$ memberi
unsur pertama $S_{22}/(S_{11}S_{22}-S_{12}^2)$, dan $S_{12}=r_{12}\sqrt{S_{11}S_{22}}$.
Membagi pembilang dan penyebut dengan $S_{22}$ menghasilkan
$\sigma^2/\bigl(S_{11}(1-r_{12}^2)\bigr)$.
\end{proof}

Angkanya dapat diperiksa: $S_{11}=42{,}9167$, $r_{12}^2=0{,}973775$, sehingga
$\mathrm{VIF}=38{,}132232$ dan $\sqrt{\mathrm{VIF}}=6{,}1751$. Galat baku $\hat b_1$
memang $0{,}135366$, sedangkan $\sigma/\sqrt{S_{11}}$ hanya $0{,}021921$; nisbahnya
$6{,}1751$, cocok dengan akar VIF sampai empat angka.

Kolinearitas tidak merusak ramalan; ia merusak tafsiran. Nilai $R^2=0{,}997951$
menunjukkan garisnya masih dekat ke titik, dan ramalan pada $x_1=9$, $x_2=28$ masih
berguna. Yang tidak dapat dipercaya adalah kalimat ``satu karyawan tambahan menaikkan
penjualan $0{,}687$ juta'': selang $95$ persennya $(0{,}380697;\ 0{,}993135)$ lebar,
dan bila data bergeser sedikit koefisien itu dapat berubah banyak tanpa garisnya
berubah banyak.

\section{Selang, uji $t$, dan selang prediksi}
\tingkat{lanjut}

Pertanyaan berikutnya sama dengan Bab 20, hanya bendanya berganti: bukan rata-rata,
melainkan kemiringan. Kita memerlukan satu besaran yang sebarannya diketahui bila
modelnya benar, lalu membaliknya menjadi selang.

Kalimat sehari-harinya: kita mengukur seberapa jauh kemiringan yang teramati dari nol,
dalam satuan goyangan khas kemiringan itu sendiri.

\kamus{selang prediksi}{prediction interval}{selang untuk satu hari yang belum terjadi,
selalu lebih lebar daripada selang untuk rata-rata karena ia harus menanggung kejutan
hari itu juga}

\begin{teorema}[Uji $t$ dan selang untuk satu koefisien]
Andaikan (A1) sampai (A4) dan galatnya normal. Dengan $p$ banyak parameter,
$s^2=\mathrm{SSE}/(n-p)$, dan $\operatorname{se}(\hat b_j)=s\sqrt{[(X^{\top}X)^{-1}]_{jj}}$,
maka
\[
\frac{\hat b_j-\beta_j}{\operatorname{se}(\hat b_j)}\sim t_{n-p},
\]
sehingga selang $1-\alpha$ untuk $\beta_j$ adalah
$\hat b_j\pm t_{n-p,1-\alpha/2}\operatorname{se}(\hat b_j)$, dan uji $H_0\colon\beta_j=0$
menolak pada taraf $\alpha$ bila $|t|>t_{n-p,1-\alpha/2}$.
\end{teorema}

\begin{proof}[Bukti]
$\hat b_j$ adalah kombinasi linear $y$ dengan koefisien tetap, jadi ia normal dengan
harapan $\beta_j$ dan variansi $\sigma^2[(X^{\top}X)^{-1}]_{jj}$ menurut Gauss--Markov.
Selanjutnya $\mathrm{SSE}/\sigma^2\sim\chi^2_{n-p}$ dan bebas dari $\hat\beta$, sebab
sisaan tegak lurus ruang kolom sedangkan $\hat\beta$ seluruhnya di dalam ruang itu.
Nisbah peubah normal baku dengan akar peubah khi-kuadrat dibagi derajat kebebasannya
mengikuti sebaran $t$. Bukti lengkapnya ada di Wasserman, \emph{All of Statistics},
Bab 13.
\end{proof}

Untuk data kafe, $\hat b=1{,}466667$, $s=0{,}365148$, $S_{xx}=60$, jadi
$\operatorname{se}(\hat b)=0{,}047140$. Dengan $t_{7,0{,}975}=2{,}364624$,
\[
\hat b\pm t\,\operatorname{se}(\hat b)=1{,}466667\pm0{,}111470
=(1{,}355197;\ 1{,}578136),
\]
dan $t=31{,}1127$ dengan nilai-p di bawah $10^{-6}$. Menolak $\beta=0$ di sini bukan
prestasi: dengan sisaan sesempit $22$ menit dan rentang $x$ selebar ini, hampir semua
data akan menolaknya. Yang layak dilaporkan justru selangnya, karena ia mengatakan
bahwa tambahan lembur per sepuluh pesanan ada di antara $1{,}36$ dan $1{,}58$ jam.

Sekarang bedakan dua pertanyaan yang sering tertukar. Berapa rata-rata lembur pada hari
dengan delapan puluh pesanan? Dan berapa lembur pada satu hari tertentu dengan delapan
puluh pesanan? Jawaban pertama menyempit sebanding $1/\sqrt n$; jawaban kedua tidak
pernah menyempit di bawah simpangan baku sisaan, sebab satu hari membawa kejutannya
sendiri.

\begin{proposisi}[Dua ragam ramalan]
Untuk $x_0$ tetap dan pengamatan baru $y_0=x_0^{\top}\beta+\varepsilon_0$ yang bebas
dari data lama,
\[
\operatorname{Var}(\hat y_0)=\sigma^2x_0^{\top}(X^{\top}X)^{-1}x_0,
\qquad
\operatorname{Var}(y_0-\hat y_0)=\sigma^2\bigl(1+x_0^{\top}(X^{\top}X)^{-1}x_0\bigr).
\]
\end{proposisi}

\begin{proof}[Bukti]
Yang pertama sudah muncul di Gauss--Markov. Untuk yang kedua, tulis
$y_0-\hat y_0=\varepsilon_0-x_0^{\top}(\hat\beta-\beta)$; kedua suku bebas karena
$\varepsilon_0$ bebas dari data lama, sehingga variansinya jumlah kedua variansi.
\end{proof}

Untuk $x_0=8$ pada data kafe, $x_0^{\top}(X^{\top}X)^{-1}x_0=1/9+9/60=0{,}261111$.
Setengah lebar selang rata-rata $2{,}364624\sqrt{0{,}133333\cdot0{,}261111}=0{,}441209$,
jadi $(19{,}9588;\ 20{,}8412)$ selebar $0{,}882$. Setengah lebar selang prediksi
$2{,}364624\sqrt{0{,}133333\cdot1{,}261111}=0{,}969635$, jadi
$(19{,}4304;\ 21{,}3696)$ selebar $1{,}939$. Nisbah lebarnya
$\sqrt{1{,}261111/0{,}261111}=2{,}198$: selang prediksi lebih dari dua kali lebih lebar,
dan tidak ada jumlah data yang dapat menghapus kelebihan itu. Gambar~\ref{fig:dua-selang}
menunjukkan kedua pita untuk seluruh rentang $x$.

\begin{figure}[htbp]
\centering
\begin{tikzpicture}
\begin{axis}[
  width=11.4cm, height=6.0cm,
  xmin=0.6, xmax=9.4, ymin=8.2, ymax=25.4,
  axis lines=left,
  xtick={1,2,3,4,5,6,7,8,9},
  tick label style={font=\scriptsize},
  xlabel={pesanan daring (puluhan)}, ylabel={jam lembur},
  xlabel style={font=\small}, ylabel style={font=\small},
  legend style={font=\scriptsize,draw=layergrey,at={(0.02,0.98)},
    anchor=north west,row sep=0.5pt},
]
\addplot[accent,very thick,domain=0.7:9.3,samples=60]
  {8.666667+1.466667*x+0.863439*sqrt(1.111111+(x-5)^2/60)};
\addplot[accent,very thick,domain=0.7:9.3,samples=60]
  {8.666667+1.466667*x-0.863439*sqrt(1.111111+(x-5)^2/60)};
\addplot[accent,very thick,domain=0.7:9.3,samples=60,forget plot]
  {8.666667+1.466667*x};
\addplot[layerindigo,densely dashed,thick,domain=0.7:9.3,samples=60]
  {8.666667+1.466667*x+0.863439*sqrt(0.111111+(x-5)^2/60)};
\addplot[layerindigo,densely dashed,thick,domain=0.7:9.3,samples=60]
  {8.666667+1.466667*x-0.863439*sqrt(0.111111+(x-5)^2/60)};
\addplot[only marks,mark=*,mark size=1.8pt,layerindigo] coordinates {
  (1,10) (2,12) (3,13) (4,14) (5,16) (6,18) (7,19) (8,20) (9,22)};
\legend{batas prediksi, , garis pas, batas rata-rata, }
\end{axis}
\end{tikzpicture}
\caption{Selang kepercayaan $95$ persen untuk rata-rata lembur (garis putus-putus) dan
selang prediksi $95$ persen untuk satu hari baru (garis tebal), keduanya di sekitar
garis pas $\hat y=8{,}666667+1{,}466667x$. Pada $x=8$ selang rata-rata
$(19{,}9588;\ 20{,}8412)$ selebar $0{,}882$, sedangkan selang prediksi
$(19{,}4304;\ 21{,}3696)$ selebar $1{,}939$. Kedua pita melebar di ujung rentang,
tempat garis kurang ditopang data; nisbah lebar keduanya $2{,}198$ tidak bergantung
pada $n$.}
\label{fig:dua-selang}
\end{figure}

\section{Diagnostik residual}
\tingkat{menengah}

Semua selang di atas bersyarat pada model yang benar. Karena itu langkah pertama sesudah
memasang garis bukan menghitung nilai-p, melainkan memeriksa sisaan. Kalimat
sehari-harinya: kita melihat apakah yang tertinggal dari garis masih berpola; sisaan
yang berpola berarti garisnya salah bentuk, bukan kebetulan.

\index{residual}
Empat pola yang perlu dikenali: lengkungan terhadap $x$, tanda hubungannya tidak linear;
corong yang melebar bersama $x$, tanda (A3) gagal; satu titik jauh sendirian, tersangka
salah catat; dan pola menurut urutan waktu, tanda pengamatan tidak bebas.

Untuk data kafe, sisaannya berurut
\[
-0{,}1333,\ 0{,}4,\ -0{,}0667,\ -0{,}5333,\ 0,\ 0{,}5333,\ 0{,}0667,\ -0{,}4,\ 0{,}1333 .
\]
Tidak ada lengkungan yang mencolok, dan jumlahnya nol seperti seharusnya. Yang menarik
adalah ukurannya: sisaan terbesar $0{,}5333$ jam, sekitar $32$ menit, pada hari keempat
dan keenam, keduanya di dekat pusat data. Titik-titik itu justru \emph{tidak}
berpengaruh besar, karena leverage-nya hanya $0{,}1278$.

Leverage titik $i$ pada regresi sederhana punya bentuk tertutup
\[
h_{ii}=\frac1n+\frac{(x_i-\bar x)^2}{S_{xx}},
\]
sehingga pada data ini $h_{11}=h_{99}=1/9+16/60=0{,}3778$, sedangkan $h_{55}=0{,}1111$.
Titik ujung memegang $3{,}4$ kali bobot titik tengah, dan itu bukan pilihan perancang,
melainkan akibat geometri: garis yang diputar di sekitar pusat data bergerak paling
jauh di ujung. Pengaruh sebenarnya menggabungkan besarnya sisaan dengan leverage,
dan ukuran yang lazim adalah jarak Cook
\[
D_i=\frac{e_i^2}{p\,s^2}\cdot\frac{h_{ii}}{(1-h_{ii})^2},
\]
yang di sini bernilai terbesar $0{,}2870$ pada hari kedua dan kedelapan, keduanya
sisaan $0{,}4$ dengan leverage $0{,}2611$. Nilai di bawah satu umumnya dianggap tidak
mencemaskan, tetapi ambang itu konvensi, bukan hukum.

Satu peringatan yang jarang ditulis: dengan $n=9$ diagnostik hampir tidak berdaya.
Pola pada sembilan titik mudah berasal dari kebetulan, dan pola yang sungguh ada mudah
tidak terlihat. Diagnostik memberi tuduhan, bukan putusan.

\begin{bacakeras}
Dua suara.

``$R$ kuadratnya nol koma sembilan sembilan. Berarti modelnya hampir sempurna.''

``Sempurna untuk apa? Sisaannya masih dua puluh dua menit.''

``Tetapi sembilan puluh sembilan persen goyangannya sudah dijelaskan.''

``Dijelaskan oleh apa? Oleh pesanan. Bukan oleh dapur yang tiba-tiba lambat.''

``Jadi $R$ kuadrat tinggi bukan jaminan.''

``$R$ kuadrat tinggi hanya berarti garisnya rapat ke titik yang ada. Ia tidak tahu
apa pun tentang hari yang belum datang.''

Ucapkan sekali lagi, pelan: rapat ke titik lama bukan janji tentang hari baru. Ulangi
sampai terasa wajar di mulut: rapat ke titik lama, bukan janji.
\end{bacakeras}

\begin{resep}
\begin{enumerate}
  \item Gambar $y$ terhadap $x$ sebelum menghitung apa pun, dan tanyakan apakah garis
        lurus masuk akal.
  \item Pasang garis dengan kuadrat terkecil, lalu tulis persamaannya lengkap dengan
        satuan setiap koefisien.
  \item Periksa sisaan terhadap $x$ dan terhadap urutan waktu; cari lengkungan,
        corong, dan titik menyendiri.
  \item Laporkan $s$ bersama $R^2$, lalu selang untuk koefisien yang menjadi
        kesimpulan utama.
  \item Bila tujuannya meramalkan satu kejadian, pakai selang prediksi; bila tujuannya
        rata-rata, pakai selang rata-rata.
\end{enumerate}
Peringatan urutan: langkah 1 tidak boleh ditukar dengan langkah 2. Menggambar sesudah
memasang garis membuat mata mencari pembenaran, bukan mencari kesalahan.
\end{resep}

\section{Ketika jawabannya ya atau tidak}
\tingkat{lanjut}

Sampai sini $y$ berupa bilangan yang dapat diukur. Banyak pertanyaan penting berbentuk
biner: pelanggan kembali atau tidak, transaksi curang atau bukan, pasien sembuh atau
tidak. Masih menggoda memasang garis lurus biasa pada $0$ dan $1$, lalu membacanya
sebagai peluang. Dua hal salah sekaligus, dan keduanya dapat dihitung.

Data berikut: $26$ pelanggan, jarak rumah dari kedai $x$ dalam kilometer, dan $y=1$ bila
mereka kembali dalam tiga puluh hari. Tercatat $15$ dari $26$ kembali. Regresi linear
pada $y$ biner, yang disebut \emph{model peluang linear}, memberi
$\hat y=1{,}346154-0{,}096154x$. Pada $x=2$ ramalannya $1{,}153846$: peluang lebih besar
dari satu, dan itu bukan kesalahan pembulatan melainkan cacat bentuk. Pada $x=20$,
masih di luar rentang data tetapi hanya tujuh kilometer di luar, ramalannya $-0{,}576923$.

Kerusakan kedua lebih halus. Untuk $y$ biner, variansi sisaan sebenarnya $p(1-p)$ dan
karenanya bergantung pada $x$: pada $x=2$ nilainya $0{,}007225$, pada $x=14$ nilainya
$0{,}028890$. Model peluang linear mengandaikan satu ragam untuk semua titik, yaitu
$s^2=0{,}124199$ pada data ini, sehingga galat bakunya salah di tempat yang paling
penting.
Kuadrat terkecil biasa, yang bekerja baik ketika ragamnya tetap, di sini memakai
asumsi yang jelas tidak berlaku.

\index{regresi logistik}
\begin{definisi}[Model logistik]
Untuk $p_i=P(y_i=1\mid x_i)$,
\[
\log\frac{p_i}{1-p_i}=a+b\,x_i,
\qquad\text{yaitu}\qquad
p_i=\frac{1}{1+e^{-(a+bx_i)}} .
\]
Parameter ditaksir dengan kemungkinan maksimum, bukan dengan kuadrat terkecil, sehingga
tidak ada bentuk tertutup dan perhitungannya memakai iterasi.
\end{definisi}

\kamus{rasio peluang}{odds ratio}{berapa kali peluang berlipat untuk satu satuan
tambahan, bukan berapa poin peluangnya bertambah}

Untuk data itu, kemungkinan maksimum memberi $\hat a=6{,}315432$ dan $\hat b=-0{,}699945$
dengan galat baku $2{,}383413$ dan $0{,}258594$. Rasio peluangnya
$e^{-0{,}699945}=0{,}496613$: setiap kilometer tambahan memangkas peluang kembali
menjadi sekitar separuh, dalam arti nisbah, bukan selisih poin. Taksiran peluangnya
$0{,}992722$ pada $x=2$ dan $0{,}029777$ pada $x=14$, keduanya sah sebagai peluang.
Selang $95$ persen untuk $b$ adalah $(-1{,}233657;\ -0{,}166233)$; seluruhnya negatif,
jadi arahnya jelas meskipun besarnya tidak.

Satu kesalahan tafsir yang mahal: koefisien logistik bukan efek marginal pada peluang.
Nilai $b=-0{,}699945$ bekerja pada skala log nisbah, dan efeknya pada peluang paling
besar di sekitar $p=0{,}5$ serta paling kecil di ujung. Laporan yang baik menyebutkan
peluang pada beberapa nilai $x$ yang bermakna, bukan hanya koefisiennya.

\section{Korelasi bukan kausalitas}
\tingkat{dasar}

Sekarang bagian yang paling sering dilanggar di laporan yang tampak rapi. Dua kelompok
karyawan mengikuti pelatihan dengan lama berbeda, dan dicatat banyak kesalahan per
seratus transaksi.

\begin{center}
\small
\begin{tabular}{l*{6}{r}}
\hline
shift pagi  & jam latihan & 1 & 2 & 3 & 4 & 5 \\
            & kesalahan   & 4,0 & 3,6 & 3,2 & 2,9 & 2,5 \\
\hline
shift malam & jam latihan & 5 & 6 & 7 & 8 & 9 \\
            & kesalahan   & 8,0 & 7,5 & 7,0 & 6,4 & 6,0 \\
\hline
\end{tabular}
\end{center}

Dalam shift pagi, kemiringannya $-0{,}36$ kesalahan per jam latihan. Dalam shift malam,
$-0{,}474286$. Keduanya negatif: latihan lebih lama berpasangan dengan kesalahan lebih
sedikit. Tetapi bila kedua kelompok digabung, kemiringannya $+0{,}359036$ dengan
$R^2=0{,}226121$. Tanda berbalik, dan tidak ada satu pun angka yang salah.

Sebabnya bukan sihir, melainkan penjumlahan dua sumber:
\[
b_{\text{gab}}=\frac{\sum_g S_{xy,g}+\sum_g n_g(\bar x_g-\bar x)(\bar y_g-\bar y)}
{\sum_g S_{xx,g}+\sum_g n_g(\bar x_g-\bar x)^2}.
\]
Kelompok pagi menyumbang $S_{xy}=-6{,}3$ dan $S_{xx}=17{,}5$; kelompok malam
$S_{xy}=-8{,}3$ dan $S_{xx}=17{,}5$; suku antarkelompok
$6\cdot(-2)(-1{,}85)+6\cdot(2)(1{,}85)=44{,}4$; penyebutnya $83$. Maka
$b_{\text{gab}}=(-6{,}3-8{,}3+44{,}4)/83=29{,}8/83=0{,}359036$. Suku antarkelompok
$44{,}4$ mengalahkan jumlah suku dalam kelompok $-14{,}6$, dan itulah seluruh rahasianya.

Tafsirannya: jam latihan tidak diberikan secara acak. Karyawan shift malam, yang
kesalahannya lebih banyak karena jam kerja dan jenis transaksi, juga menerima latihan
lebih banyak, sehingga yang terukur sebagai efek latihan sebagian besar adalah
perbedaan shift. Gambar~\ref{fig:simpson} menaruh semuanya pada satu bidang.

\begin{figure}[htbp]
\centering
\begin{tikzpicture}
\begin{axis}[
  width=11.4cm, height=6.2cm,
  xmin=0.5, xmax=10.5, ymin=1.6, ymax=8.9,
  axis lines=left,
  xtick={1,2,3,4,5,6,7,8,9,10},
  tick label style={font=\scriptsize},
  xlabel={jam latihan}, ylabel={kesalahan per 100 transaksi},
  xlabel style={font=\small}, ylabel style={font=\small},
  legend style={font=\scriptsize,draw=layergrey,at={(0.98,0.98)},
    anchor=north east,row sep=0.5pt},
]
\addplot[only marks,mark=*,mark size=2pt,layerteal] coordinates {
  (1,4.0) (2,3.6) (3,3.2) (4,2.9) (5,2.5)};
\addplot[only marks,mark=square*,mark size=1.9pt,layerplum] coordinates {
  (6,7.5) (7,7.0) (8,6.4) (9,6.0) (10,5.7)};
\addplot[layerteal,very thick,domain=0.8:5.4,samples=2] {4.326667-0.36*x};
\addplot[layerplum,very thick,domain=5.6:10.2,samples=2]
  {10.32381-0.474286*x};
\addplot[accent,thick,densely dashed,domain=0.8:10.2,samples=2]
  {2.941968+0.359036*x};
\legend{pagi, malam, gabungan}
\end{axis}
\end{tikzpicture}
\caption{Paradoks Simpson pada data pelatihan. Di dalam setiap shift, garisnya menurun:
$-0{,}36$ kesalahan per jam untuk pagi dan $-0{,}4743$ untuk malam. Garis gabungan,
yang dihitung dari kedua belas titik tanpa memandang shift, justru menanjak dengan
kemiringan $+0{,}3590$. Perubahan tanda itu berasal dari selisih antar kelompok, bukan
dari satu titik yang aneh; suku antarkelompok $44{,}4$ mengalahkan suku dalam kelompok
$-14{,}6$.}
\label{fig:simpson}
\end{figure}

\begin{untukanak}
Dua kelas mengikuti lomba membaca. Kelas A membaca buku tipis, banyak halaman, dan
nilainya bagus. Kelas B membaca buku tebal, lebih sedikit halaman, dan nilainya jelek.
Kalau semua murid digabung, terlihat seolah-olah makin sedikit halaman yang dibaca,
makin tinggi nilainya. Padahal di dalam setiap kelas, makin banyak halaman, makin tinggi
nilainya.

Yang salah bukan angkanya, melainkan pencampurannya. Kelas B memang diberi buku yang
lebih tebal, dan kelas B juga memang lebih lemah membaca. Dua hal itu bercampur menjadi
satu garis yang menipu.

Di bab ini kita akan tahu mengapa mencampur dua kelompok dapat membalik arah kesimpulan,
mengapa pertanyaan ``apakah latihan mengurangi kesalahan'' tidak dapat dijawab hanya
dengan menggambar semua titik, dan mengapa orang yang mengubah kebijakan harus tahu
kelompok mana yang sebenarnya ia bandingkan.
\end{untukanak}

\begin{tantangan}
Susun sebuah tabel dengan dua kelompok, masing-masing lima titik, sedemikian sehingga
kedua kemiringan dalam kelompok positif tetapi kemiringan gabungan negatif, dan
selisih rata-rata $y$ antarkelompok paling kecil yang masih memungkinkan. Berapa
selisih terkecil itu, dan mengapa tidak bisa lebih kecil?

Petunjuk 1: tanda kemiringan gabungan ditentukan oleh suku antarkelompok. Buat
kelompok pertama memiliki $x$ kecil dan $y$ kecil, kelompok kedua $x$ besar dan $y$
besar, lalu balik arah di dalam masing-masing kelompok.

Petunjuk 2: dengan lima titik dan kemiringan dalam kelompok $+1$ pada
$S_{xx}=10$, setiap kelompok menyumbang $S_{xy}=+10$ dan suku antarkelompok harus
lebih kecil dari $-20$. Karena $n_g(\bar x_g-\bar x)(\bar y_g-\bar y)$ bernilai
negatif hanya bila salah satu selisih rata-rata berlawanan tanda dengan yang lain,
selisih rata-rata harus melewati titik di mana kedua suku itu saling meniadakan.

Kalau mentok, kembali ke bagian Korelasi bukan kausalitas dan hitung ulang
$29{,}8/83$ langkah demi langkah.
\end{tantangan}

\section{Pelaporan yang jujur}
\tingkat{menengah}

Sebuah analisis regresi dapat dikerjakan dengan benar dan dilaporkan sehingga menyesatkan
tanpa satu pun angka dipalsukan. Empat kebiasaan menutup sebagian besar celah itu.

Pertama, sebutkan berapa pengamatan dan berapa parameter, karena pembaca memerlukan
keduanya untuk menghitung derajat kebebasan dan menilai apakah $s$ dapat dipercaya.
Pada data kafe, $n=9$ dan $p=2$, jadi $s$ dibangun dari tujuh derajat kebebasan;
menyembunyikan angka itu membuat $s$ tampak lebih kokoh daripada kenyataannya.

Kedua, pisahkan data yang dipakai memilih model dari data yang dipakai menilai.
Memilih bentuk model, memilih peubah, dan menghitung galat pada himpunan yang sama
selalu menghasilkan angka yang terlalu bagus. Pada $n=9$ pemisahan itu tidak mungkin,
dan justru karena itu kesimpulannya harus disebut lemah, bukan diperkuat dengan
$R^2=0{,}992821$.

Ketiga, sebutkan berapa banyak model yang dicoba. Bila sepuluh kemungkinan diuji dan
yang terbaik dilaporkan, peluang sekurangnya satu di antaranya tampak berarti adalah
$1-0{,}95^{10}=0{,}401263$ meskipun tidak satu pun efeknya nyata. Angka itu bukan
peluang salah, melainkan peluang terpilihnya sesuatu yang kebetulan.

Keempat, laporkan besaran dengan satuan dan selang, bukan hanya tanda bintang. Kalimat
``$\hat b=0{,}687$ ($95\%$: $0{,}381$ sampai $0{,}993$ juta per karyawan)'' dapat
diperiksa; kalimat ``berpengaruh signifikan'' tidak. Bila kesimpulannya tentang kausalitas,
sebutkan mengapa: pengacakan, atau penyesuaian terhadap pencampur yang disebut namanya.

\section{Contoh hitung dari ujung ke ujung}
\tingkat{menengah}

Kita ulang seluruh rantai pada data kafe, dengan angka antara ditulis, supaya setiap
langkah dapat diperiksa dengan kalkulator.

Rata-ratanya $\bar x=5$ dan $\bar y=16$. Simpangan kuadratnya
$S_{xx}=\sum(x_i-5)^2=60$ dan $S_{yy}=130$; hasil kalinya
$S_{xy}=\sum(x_i-5)(y_i-16)=808-9\cdot5\cdot16=88$. Maka
$\hat b=88/60=1{,}466667$ dan $\hat a=16-1{,}466667\cdot5=8{,}666667$.

Ramalan pada sembilan titik itu $10{,}1333$; $11{,}6$; $13{,}0667$; $14{,}5333$;
$16$; $17{,}4667$; $18{,}9333$; $20{,}4$; $21{,}8667$. Sisaannya
$-0{,}1333$; $0{,}4$; $-0{,}0667$; $-0{,}5333$; $0$; $0{,}5333$; $0{,}0667$;
$-0{,}4$; $0{,}1333$, dan jumlah kuadratnya
$\mathrm{SSE}=0{,}933333$. Dengan $n-p=7$, $s^2=0{,}133333$ dan $s=0{,}365148$ jam.

Galat baku kemiringan $s/\sqrt{S_{xx}}=\sqrt{0{,}133333/60}=\sqrt{0{,}00222222}
=0{,}047140$. Statistik $t=1{,}466667/0{,}047140=31{,}1127$, dan
$t_{7,0{,}975}=2{,}364624$, sehingga setengah lebar selangnya $0{,}111470$ dan
selangnya $(1{,}355197;\ 1{,}578136)$.

\kertas{Ulangi perhitungan yang sama untuk satu perubahan data. Hari kesembilan
tercatat $22$ jam lembur; andaikan catatan yang benar adalah $26$ jam. Hitung ulang
$\hat b$, $s$, $\operatorname{se}(\hat b)$, dan selang $95$ persennya, lalu bandingkan
dengan tabel di atas. Jawaban yang benar: $\hat b=1{,}733333$, $s=1{,}306880$,
$\operatorname{se}(\hat b)=0{,}168718$, selang $(1{,}3344;\ 2{,}1323)$. Lebarnya
$0{,}798$ melawan $0{,}223$, yaitu $3{,}6$ kali. Jelaskan mengapa satu titik di ujung
dapat melakukan itu, padahal ia hanya satu dari sembilan.}

\section{Lab komputasi}
\tingkat{lanjut}

Semua angka di bab ini dapat dihitung ulang dalam satu berkas pendek. Program berikut
memakai $X^{\top}X$ secara langsung, jadi ia sekaligus memeriksa aljabar di bagian
kedua.

\begin{lab}
\begin{verbatim}
import numpy as np
from math import sqrt
from scipy import stats

x = np.arange(1, 10, dtype=float)
y = np.array([10, 12, 13, 14, 16, 18, 19, 20, 22], dtype=float)
n = len(y)
X = np.column_stack([np.ones(n), x])

beta = np.linalg.solve(X.T @ X, X.T @ y)
fit = X @ beta
e = y - fit
s2 = (e ** 2).sum() / (n - 2)
sst = ((y - y.mean()) ** 2).sum()

print("X'X =", (X.T @ X).tolist(), " X'y =", (X.T @ y).tolist())
print("beta = %.6f %.6f" % tuple(beta))
sse = (e ** 2).sum()
print("SSE = %.6f  R2 = %.6f  s = %.6f"
      % (sse, 1 - sse / sst, sqrt(s2)))
print("residu =", np.round(e, 4).tolist())

cov = s2 * np.linalg.inv(X.T @ X)
se = np.sqrt(np.diag(cov))
tq = stats.t.ppf(0.975, n - 2)
print("se = %.6f %.6f  t = %.4f %.4f"
      % (se[0], se[1], beta[0] / se[0], beta[1] / se[1]))
lo, hi = beta[1] - tq * se[1], beta[1] + tq * se[1]
print("IK95 kemiringan = (%.6f, %.6f)" % (lo, hi))

for x0 in (5.0, 8.0):
    c = np.array([1.0, x0])
    y0 = c @ beta
    v = s2 * (c @ np.linalg.inv(X.T @ X) @ c)
    print("x0=%.0f  yhat=%.4f  rata2 +/-%.4f  prediksi +/-%.4f"
          % (x0, y0, tq * sqrt(v), tq * sqrt(v + s2)))

h = np.einsum("ij,jk,ik->i", X, np.linalg.inv(X.T @ X), X)
cook = e ** 2 / (2 * s2) * h / (1 - h) ** 2
print("h_ii =", np.round(h, 4).tolist())
print("D_Cook =", np.round(cook, 4).tolist())

x1 = np.array([6, 7, 7, 8, 8, 9, 9, 10, 10, 11, 12, 12], dtype=float)
x2 = np.array([18, 22, 20, 26, 24, 28, 26, 32, 30, 34, 36, 38], dtype=float)
yb = np.array([11.2, 12.4, 12.0, 14.1, 13.6, 15.2, 14.8, 17.0, 16.4, 18.1,
               19.4, 19.9])
Xb = np.column_stack([np.ones(12), x1, x2])
bb = np.linalg.solve(Xb.T @ Xb, Xb.T @ yb)
eb = yb - Xb @ bb
s2b = (eb ** 2).sum() / 9
seb = np.sqrt(np.diag(s2b * np.linalg.inv(Xb.T @ Xb)))
r12 = np.corrcoef(x1, x2)[0, 1]
print("berganda beta = %.4f %.4f %.4f  s = %.4f"
      % (bb[0], bb[1], bb[2], sqrt(s2b)))
print("berganda se   = %.4f %.4f %.4f" % tuple(seb))
print("r12 = %.4f  VIF = %.2f  akar VIF = %.2f"
      % (r12, 1 / (1 - r12 ** 2), 1 / sqrt(1 - r12 ** 2)))


def ols(xx, yy):
    A = np.column_stack([np.ones(len(xx)), xx])
    return np.linalg.solve(A.T @ A, A.T @ yy)[1]


xp = np.array([1, 2, 3, 4, 5, 6], dtype=float)
yp = np.array([4.0, 3.6, 3.2, 2.9, 2.5, 2.2])
xm = np.array([5, 6, 7, 8, 9, 10], dtype=float)
ym = np.array([8.0, 7.5, 7.0, 6.4, 6.0, 5.7])
print("Simpson: pagi %.4f  malam %.4f  gabung %.4f"
      % (ols(xp, yp), ols(xm, ym), ols(np.r_[xp, xm], np.r_[yp, ym])))

xd = np.array([2, 2, 3, 3, 4, 4, 5, 5, 6, 6, 7, 7, 8, 8, 9, 9, 10, 10,
               11, 11, 12, 12, 13, 13, 14, 14], dtype=float)
yd = np.array([1, 1, 1, 1, 1, 1, 1, 1, 1, 1, 1, 0, 1, 0, 1, 0, 1, 0, 0,
               1, 0, 0, 0, 0, 0, 0], dtype=float)
A = np.column_stack([np.ones(len(yd)), xd])
lpm = np.linalg.solve(A.T @ A, A.T @ yd)
print("LPM: %.4f %+.4f  prediksi di x=2: %.4f  di x=20: %.4f"
      % (lpm[0], lpm[1], lpm @ [1, 2], lpm @ [1, 20]))


def nll(th):
    z = A @ th
    return float(np.sum(np.log1p(np.exp(-z)) + (1 - yd) * z))


g = np.array([0.0, 0.0])
for _ in range(40):
    p = 1 / (1 + np.exp(-(A @ g)))
    grad = A.T @ (p - yd)
    Hm = A.T @ (A * (p * (1 - p))[:, None])
    g = g - np.linalg.solve(Hm, grad)
se_g = np.sqrt(np.diag(np.linalg.inv(Hm)))
p = 1 / (1 + np.exp(-(A @ g)))
print("logit: %.4f %+.4f  se %.4f %.4f" % (g[0], g[1], se_g[0], se_g[1]))
print("devians = %.4f  p(x=2) = %.4f  p(x=14) = %.4f"
      % (2 * nll(g), p[0], p[-1]))
print("rasio peluang per km = %.4f" % np.exp(g[1]))
\end{verbatim}
\end{lab}

Keluarannya, apa adanya:

\begin{verbatim}
X'X = [[9.0, 45.0], [45.0, 285.0]]  X'y = [144.0, 808.0]
beta = 8.666667 1.466667
SSE = 0.933333  R2 = 0.992821  s = 0.365148
residu = [-0.1333, 0.4, -0.0667, -0.5333, 0.0, 0.5333, 0.0667, -0.4, 0.1333]
se = 0.265274 0.047140  t = 32.6706 31.1127
IK95 kemiringan = (1.355197, 1.578136)
x0=5  yhat=16.0000  rata2 +/-0.2878  prediksi +/-0.9101
x0=8  yhat=20.4000  rata2 +/-0.4412  prediksi +/-0.9696
h_ii = [0.3778, 0.2611, 0.1778, 0.1278, 0.1111, 0.1278, 0.1778, 0.2611, 0.3778]
D_Cook = [0.0651, 0.287, 0.0044, 0.1792, 0.0, 0.1792, 0.0044, 0.287, 0.0651]
berganda beta = 2.4430 0.6869 0.2393  s = 0.1436
berganda se   = 0.2076 0.1354 0.0421
r12 = 0.9868  VIF = 38.13  akar VIF = 6.18
Simpson: pagi -0.3600  malam -0.4743  gabung 0.3590
LPM: 1.3462 -0.0962  prediksi di x=2: 1.1538  di x=20: -0.5769
logit: 6.3154 -0.6999  se 2.3834 0.2586
devians = 17.8843  p(x=2) = 0.9927  p(x=14) = 0.0298
rasio peluang per km = 0.4966
\end{verbatim}

Dua hal perlu dibaca dari keluaran itu. Pertama, $\mathrm{X}^{\top}\mathrm{X}$ dan
$\mathrm{X}^{\top}y$ tercetak persis seperti dihitung tangan, dan $\hat\beta$ keluar
dari penyelesaian sistem; kesepakatan semacam ini bukan bukti kebenaran teori, tetapi ia
menyingkirkan salah hitung, yang paling sering menjadi penyebab. Kedua, tiga bagian
terakhir memperlihatkan mengapa tafsiran lebih rapuh daripada hitungan: $0{,}6869$
berbanding $1{,}4460$ pada regresi tunggal, VIF $38{,}13$, dan kemiringan gabungan yang
tandanya berbalik. Tidak satu pun dari ketiganya kesalahan program.

\begin{untukdosen}
Urutan mengajar yang saya sarankan: mulai dari tabel sembilan hari, minta kelas menebak
kemiringannya dengan dua cara, yaitu selisih ujung dan ``garis yang paling dekat''.
Selisih kedua jawaban kecil, dan justru itu pintu masuknya: pertanyaan sebenarnya bukan
angkanya, melainkan goyangannya. Baru sesudah kelas gelisah, masuk ke persamaan normal
dan tunjukkan bahwa $\sum x_ie_i=0$ adalah pernyataan geometris, bukan trik aljabar.

Dua miskonsepsi yang paling sering muncul. Pertama, mahasiswa membaca $R^2$ sebagai
``peluang ramalan benar'', lalu terkejut ketika selang prediksi selebar dua jam pada
$R^2=0{,}99$; gejalanya tampak ketika mereka tidak dapat menjelaskan mengapa $s$ tidak
ikut mengecil bersama $R^2$. Kedua, mereka menafsirkan koefisien berganda seperti
koefisien tunggal, padahal pada data dua belas cabang angkanya berubah dari $1{,}446$
menjadi $0{,}687$ hanya karena penjelas lain masuk.

Aktivitas papan tulis yang berhasil: suruh dua kelompok menghitung kemiringan pada data
pelatihan, satu kelompok memakai semua dua belas titik, satu kelompok memisahkan shift.
Tulis kedua hasil berdampingan, lalu minta masing-masing mempertahankannya. Tanda bahwa
kelas siap lanjut: mereka dapat menyebut satu keadaan yang membuat koefisien berganda
berubah tanda, dan dapat menjelaskan mengapa ramalannya tetap baik.
\end{untukdosen}

\section{Ringkasan}
\tingkat{dasar}

\begin{ringkasan}
\begin{itemize}
  \item Regresi memasang garis yang meminimumkan jumlah kuadrat sisaan. Menurut
        persamaan normal sisaannya tegak lurus setiap kolom $X$; pada data kafe
        $\hat b=88/60=1{,}466667$ dan $\hat a=8{,}666667$.
  \item Koefisien hanya bermakna bersama satuannya: $1{,}466667$ jam per puluhan pesanan
        sama dengan $8{,}8$ menit per pesanan. Mengubah skala $x$ mengubah koefisien dan
        galat bakunya, tidak mengubah $t$ maupun $R^2$.
  \item $R^2$ nisbi, $s$ mutlak. Pada data kafe $R^2=0{,}992821$ tetapi
        $s=0{,}365148$ jam, yaitu $21{,}9$ menit.
  \item Di bawah (A1)--(A4) Gauss--Markov menjamin kuadrat terkecil adalah BLUE dengan
        $\operatorname{Var}(\hat\beta)=\sigma^2(X^{\top}X)^{-1}$; jaminan itu gugur bila
        ragam galat tidak tetap atau penjelasnya kolinear.
  \item Kolinearitas tidak merusak ramalan tetapi merusak tafsiran: pada dua belas cabang,
        $r_{12}=0{,}986801$ memberi $\mathrm{VIF}=38{,}13$ dan galat baku $6{,}18$ kali
        lebih besar, sementara $R^2$ tetap $0{,}997951$.
  \item Selang prediksi selalu lebih lebar daripada selang rata-rata; pada $x_0=8$
        nisbahnya $2{,}198$ dan tidak mengecil bersama $n$. Tanda kemiringan dapat
        berbalik ketika dua kelompok dicampur: $-0{,}36$ dan $-0{,}474286$ di dalam
        kelompok, $+0{,}359036$ setelah digabung.
\end{itemize}
\end{ringkasan}

\section{Latihan}
\tingkat{menengah}

\begin{latihan}
\begin{enumerate}
  \item Nilai klaim ini: ``Karena $R^2=0{,}99$, garis ini dapat meramalkan satu hari
        dengan galat paling banyak dua puluh menit.'' Benar, salah, atau belum cukup
        bukti? Sebutkan bagian yang benar, bagian yang keliru, dan angka tambahan yang
        menyelesaikan penilaian itu.
  \item Buktikan bahwa dengan kolom satu pada $X$ berlaku $\sum_i e_i=0$ dan
        $\sum_i x_ie_i=0$, lalu jelaskan mengapa keduanya gagal bila garis dipaksa lewat
        titik asal.
  \item Buktikan $\operatorname{Var}(\hat b)=\sigma^2/S_{xx}$ dari
        $\hat b=\sum_i(x_i-\bar x)y_i/S_{xx}$ tanpa rumus matriks.
  \item Hitung $\hat a$, $\hat b$, $s$, $\operatorname{se}(\hat b)$, dan selang $95$
        persen untuk $(1,2)$, $(2,2{,}5)$, $(3,4)$, $(4,3{,}5)$, $(5,5)$, lalu periksa
        dengan Python. Mengapa selangnya jauh lebih lebar daripada selang data kafe?
  \item Hitung selang prediksi $95$ persen untuk $x_0=2$ dan $x_0=9$ pada data kafe.
        Mengapa lebarnya berbeda padahal $n$ sama, dan mengapa selang itu tidak pernah
        menyempit ke nol ketika $n$ diperbesar?
  \item Tunjukkan bahwa menambah penjelas tidak dapat menaikkan SSE, lalu jelaskan
        mengapa $R^2_{\text{adj}}$ dapat turun. Hitung $R^2_{\text{adj}}$ untuk data kafe
        bila $p$ naik dari $2$ menjadi $3$ pada $R^2$ yang sama.
  \item Untuk $S_{11}=50$, $S_{22}=200$, $r_{12}=0{,}9$, dan $\sigma^2=2$, hitung VIF
        serta galat baku $\hat b_1$ dengan dan tanpa $x_2$ dalam model.
  \item Susun tabel dua kelompok dengan kemiringan dalam kelompok $-0{,}2$ dan kemiringan
        gabungan $+0{,}4$, lalu periksa hasilnya dengan rumus dekomposisi.
  \item Tulis ulang kalimat ini sehingga jujur: ``model terbaik dari dua belas model yang
        kami coba memberi $p=0{,}03$, jadi efeknya nyata''.
\end{enumerate}

Petunjuk untuk tiga butir tersulit. Butir 3: tulis $\hat b-\beta=\sum_i(x_i-\bar x)
\varepsilon_i/S_{xx}$, ambil variansi, pakai ketakbebasan dan
$\operatorname{Var}(\varepsilon_i)=\sigma^2$. Butir 6: bandingkan ruang kolom kedua model,
lalu ingat pembagi $n-p$ yang mengecil. Butir 7: $\mathrm{VIF}=1/(1-0{,}81)=5{,}263158$,
jadi galat bakunya $\sqrt{5{,}263158}$ kali lebih besar.
\end{latihan}

\begin{untukpeneliti}
Yang benar-benar terbukti di sini: persamaan normal sebagai syarat perlu dan cukup bagi
minimum beserta ketunggalannya saat $X$ berpangkat kolom penuh; sifat matriks topi dan
$\operatorname{tr}H=p$; ketakbiasan serta
$\operatorname{Var}(\hat\beta)=\sigma^2(X^{\top}X)^{-1}$; ketidaksamaan Gauss--Markov
dalam kelas penaksir linear tak bias; dekomposisi $\mathrm{SST}=\mathrm{SSR}+\mathrm{SSE}$
bila ada kolom satu; ketakmonotonan $R^2$ terhadap penambahan kolom; sebaran $t$ untuk
nisbah koefisien bila galat normal; dan variansi galat prediksi.

Yang hanya populer: bahwa $R^2$ mengukur mutu model, bahwa koefisien berganda selalu
merupakan efek, dan bahwa ambang $\mathrm{VIF}>10$ atau $D_i>1$ adalah hasil turunan.
Sumber primer: Seber dan Lee, \emph{Linear Regression Analysis}, Bab 3; Greene,
\emph{Econometric Analysis}, Bab 3; Wasserman, \emph{All of Statistics}, Bab 13;
Simpson (1951) untuk pembalikan arah pada tabel kontingensi; Hosmer, Lemeshow, dan
Sturdivant, \emph{Applied Logistic Regression}, Bab 1--3.

Batas keberlakuan: semua selang bersyarat pada model yang benar, galat bebas, dan ragam
tetap. Pada data berkelompok, deret waktu, atau pengamatan berulang, ketakbebasan gagal
dan galat baku yang dilaporkan terlalu kecil. Pertanyaan terbuka: seberapa jauh selang
pascapemilihan model menyimpang dari nominalnya, dan bagaimana diagnostik pengaruh
disusun ketika $n$ jauh lebih kecil daripada $p$.
\end{untukpeneliti}

\peta{Bab 21 pada Wasserman, \emph{All of Statistics}: Bab 13 (Linear and Logistic
Regression: kuadrat terkecil, sebaran $t$ untuk koefisien, dan regresi logistik) serta
Bab 6 (Models, Statistical Inference, and Learning: tafsiran selang dan galat baku).
Geometri kuadrat terkecil sebagai proyeksi ortogonal disusun di Bab 10 buku ini, tempat
persamaan normal pertama kali muncul sebagai akibat ortogonalitas, dan dipakai lagi di
Bab 24 untuk regularisasi. Selang kepercayaan, nilai-p, dan tafsiran taraf diambil dari
Bab 20 tanpa diulang di sini. Teorema Gauss--Markov beserta syaratnya dirujuk ke Seber
dan Lee Bab 3 dan Greene Bab 3, karena Wasserman tidak membuktikannya. Paradoks Simpson
dirujuk ke makalah Simpson (1951) dan dipakai lagi di Bab 25 ketika membahas pemisahan
data dan kebocoran evaluasi, serta di Bab 26 ketika membahas klaim kausal yang harus
dibuktikan, bukan disimpulkan dari korelasi.}

\chapter{Aritmetika Floating-Point dan Analisis Galat}

Sebuah program menyelesaikan $x^{2}+10^{8}x+1=0$ dengan rumus kuadrat yang kita pakai
sejak sekolah, $x=(-b\pm\sqrt{b^{2}-4ac})/(2a)$. Untuk $a=1$, $b=10^{8}$, $c=1$ ia
melaporkan akar $-10^{8}$ dan $-7{,}450580596923828\times10^{-9}$. Akar pertama benar
sampai digit terakhir. Akar kedua salah seperempat: nilai sebenarnya $-10^{-8}$, dan
galat relatifnya $0{,}255$. Rumusnya tidak salah tulis. Mesinnya juga tidak rusak.
Koefisiennya tepat, diskriminannya tepat, dan pengurangannya bahkan tepat menurut
teorema yang akan kita buktikan di bawah.

Yang runtuh adalah bentuk aljabarnya. Bilangan $b=10^{8}$ dan $\sqrt{b^{2}-4ac}$ berbeda
hanya satu ulp, yaitu satu langkah terkecil yang masih bisa dibedakan mesin pada skala
itu: selisihnya $1{,}49\times10^{-8}$, atau $1{,}5\times10^{-16}$ kali besar bilangannya
sendiri. Selisih sekecil itu tidak berarti apa-apa selama kita memandang bilangan
berukuran $10^{8}$. Begitu keduanya dikurangkan, hasilnya berukuran $10^{-8}$ dan
selisih tadi menjadi seluruh jawabannya: yang tersisa hanyalah galat, bukan angka.

Bab 4 dan Bab 5 membangun limit, kekontinuan, dan turunan di atas $\mathbb{R}$ yang
setiap titiknya eksak. Bab 20 memakai rumus variansi dan akar kuadrat dengan anggapan
yang sama. Komputer tidak bekerja di $\mathbb{R}$. Ia bekerja di himpunan berhingga
bilangan yang punya aturan sendiri: setiap operasi dibulatkan, penjumlahan tidak
asosiatif, dan dua cara aljabar yang setara bisa memberi jawaban berbeda jauh. Bab ini
membangun cara membaca aturan itu, mengukur kerusakannya sebelum program dijalankan,
dan memilih bentuk rumus yang tidak runtuh.

\begin{untukanak}
Bayangkan sebuah timbangan dapur digital yang layarnya hanya muat lima angka. Kita
menimbang satu karung beras bersama wadahnya, lalu menimbang wadah kosongnya, lalu
mengurangkan kedua angka itu untuk mendapat berat berasnya. Kalau wadah dan beras
sama-sama beratnya sekitar dua puluh kilogram, layar hanya bisa menunjukkan angka
sampai puluhan gram terdekat. Selisih dua angka besar yang hampir sama itu memang
benar secara hitungan, tetapi angka terakhirnya sudah tidak berarti apa-apa. Beras
yang sebenarnya delapan ratus gram bisa muncul sebagai tujuh ratus gram, atau bahkan
angka yang tidak masuk akal.

Ada masalah kedua yang lebih halus. Timbangan yang sama tidak bisa menimbang satu
butir beras: jarumnya tidak bergerak sama sekali. Kalau kita menuang sejuta butir satu
per satu, tidak ada satu pun butir yang tercatat, padahal sejuta butir itu berat.
Mesin hitung bekerja persis seperti itu: ada bilangan yang terlalu kecil untuk
mengubah apa pun, dan ada pengurangan yang membuang hampir semua angka.

Pertanyaannya: kalau mesinnya memang begitu, bagaimana kita bisa tahu kapan jawabannya
masih bisa dipercaya? Di bab ini kita akan tahu mengapa kesalahan itu bukan kejadian
acak, melainkan sesuatu yang bisa dihitung batasnya sebelum satu baris program pun
dijalankan.
\end{untukanak}

\section{Bilangan yang tidak punya wakil}
\tingkat{dasar}

Di antara $0$ dan $1$ ada tak berhingga banyak bilangan real. Di dalam komputer, semua
bilangan disimpan dalam 32 atau 64 bit, jadi yang tersedia hanya berhingga banyak.
Kompromi yang diambil sejak mesin pertama sampai standar IEEE-754 adalah notasi ilmiah
biner dengan lebar bidang tetap: satu tanda, satu eksponen, satu mantissa.

\kamus{sistem floating-point}{floating-point system}{Notasi ilmiah biner dengan jumlah
digit tetap, seperti menulis $1{,}01\times2^{7}$ dengan tempat terbatas.}

\begin{definisi}[Sistem bilangan floating-point]
Misalkan $\beta=2$ basis, $p$ presisi, dan $e_{\min}\le e\le e_{\max}$ batas eksponen.
Bilangan normal adalah
\[
  x=\pm\left(1.d_{1}d_{2}\dots d_{p-1}\right)_{2}\times2^{e},
\]
dengan $d_{i}\in\{0,1\}$. Bilangan subnormal memakai $e=e_{\min}$ dengan digit
terdepan $0$.\index{subnormal}\index{aritmetika floating-point} Himpunan seluruh
bilangan yang terwakili, ditambah nol, kedua tak hingga, dan NaN, kita sebut
$\mathbb{F}$.
\end{definisi}

Digit terdepan selalu $1$ untuk bilangan normal, jadi digit itu tidak disimpan;
itulah sebabnya bidang mantissa punya $p-1$ bit. Satu bit yang dihemat menggandakan
presisi: 23 bit tersimpan memberi 24 bit
signifikan.



\begin{teorema}[Banyaknya bilangan mesin]
Untuk setiap eksponen $e$, banyaknya bilangan normal dalam selang
$\left[2^{e},2^{e+1}\right)$ adalah $2^{p-1}$. Himpunan $\mathbb{F}$ berhingga
banyaknya, jadi terhitung.
\end{teorema}

\begin{proof}
Pada selang itu setiap bilangan berbentuk $M\cdot2^{e-p+1}$ dengan $M$ bilangan bulat
dan $2^{p-1}\le M<2^{p}$. Banyaknya $M$ semacam itu tepat $2^{p}$, dikurangi
$2^{p-1}$, yaitu $2^{p-1}$. Ada berhingga banyak eksponen yang diizinkan, jadi
$\mathbb{F}$ adalah gabungan berhingga himpunan berhingga.
\end{proof}

Untuk binary64, $2^{p-1}=2^{52}=4\,503\,599\,627\,370\,496$. Itu banyak, tetapi Bab 2
sudah menunjukkan $\mathbb{R}$ tak terhitung. Jadi hampir setiap bilangan real tidak
punya wakil: representasi selalu berarti pembulatan, dan pembulatan selalu berarti
galat. Pertanyaannya bukan apakah ada galat, melainkan seberapa besar dan ke mana ia
pergi.

\begin{bacakeras}
Dua suara, dibaca bergantian dengan teman. Ketuk meja setiap kali kata ``potongan''
diucapkan.

``Satu dibagi tiga, berapa?'' ``Nol koma tiga tiga tiga, terus, tanpa henti.''
``Mesin menulisnya sampai habis?'' ``Tidak. Ia berhenti di tempat yang muat.''
``Jadi yang disimpan apa?'' ``Potongan. Bukan nilainya, potongannya.''
``Kalau potongannya dijumlahkan?'' ``Yang bertambah potongan, bukan nilai.''
``Kalau potongannya dikurangkan?'' ``Angka depannya habis, yang tinggal debu.''

Ulangi tiga baris tengah, makin cepat, dengan ketukan yang sama: potongan bukan nilai,
potongan bukan nilai, potongan bukan nilai.

Suara kedua menutup pelan: mesin tidak berdusta, ia hanya memotong. Yang keliru adalah
kita kalau membaca potongan sebagai nilai.

Baris yang harus diulang pembaca, sekali, dengan suara jelas: komputer menyimpan
potongan, bukan nilai, jadi setiap hasil hitung harus diperiksa, bukan dipercaya.
\end{bacakeras}

Standar IEEE-754 menetapkan dua lebar yang dipakai di mana-mana. Tabel ini memuat
seluruh angkanya; tidak ada yang perlu dihafal kecuali $p$ dan biasnya.

\begin{center}
\begin{tabular}{lcc}
\hline
 & binary32 & binary64 \\
\hline
bit tanda / eksponen / mantissa & $1/8/23$ & $1/11/52$ \\
presisi $p$ & $24$ & $53$ \\
bias eksponen & $127$ & $1023$ \\
$e_{\min},\ e_{\max}$ & $-126,\ 127$ & $-1022,\ 1023$ \\
$\mathrm{eps}=2^{-p+1}$ & $1{,}1920929\times10^{-7}$ & $2{,}2204460\times10^{-16}$ \\
$u=2^{-p}$ & $5{,}9604645\times10^{-8}$ & $1{,}1102230\times10^{-16}$ \\
terkecil normal & $1{,}1754944\times10^{-38}$ & $2{,}2250739\times10^{-308}$ \\
terkecil subnormal & $1{,}4012985\times10^{-45}$ & $4{,}9406565\times10^{-324}$ \\
terbesar & $3{,}4028235\times10^{38}$ & $1{,}7976931\times10^{308}$ \\
\hline
\end{tabular}
\end{center}

Eksponen disimpan dengan bias supaya tidak perlu tanda sendiri: nilai tersimpan $123$
berarti eksponen sebenarnya $123-127=-4$. Trik ini membuat urutan bit dari nol sampai
tak hingga positif naik monoton, sehingga dua bilangan positif dapat dibandingkan
sebagai bilangan bulat.

Mari bongkar tiga pola bit yang benar-benar dihasilkan Python. Pada binary32,
$0{,}1$ tercetak sebagai \texttt{0 01111011 10011001100110011001101} dengan
$e=123-127=-4$; $6$ menjadi \texttt{0 10000001 10000000000000000000000}, yaitu
$1{,}5\times2^{2}$; dan $-12{,}375$ menjadi
\texttt{1 10000010 10001100000000000000000}, yaitu $-1{,}546875\times2^{3}$.
Dua kasus batas melengkapi gambarannya:

\begin{center}
\begin{tabular}{ll}
\hline
subnormal terkecil & \texttt{0 00000000 00000000000000000000001} \\
terbesar & \texttt{0 11111110 11111111111111111111111} \\
\hline
\end{tabular}
\end{center}

Yang pertama bernilai $2^{-149}$: seluruh bit eksponennya nol dan hanya bit terakhir
mantissa menyala. Yang kedua bernilai $3{,}4028235\times10^{38}$, dengan eksponen
tersimpan $254$ dan mantissa penuh.

Bentuk yang lebih mudah dibaca untuk binary64 disediakan oleh \texttt{float.hex()}. Hasilnya untuk beberapa bilangan:

\begin{center}
\begin{tabular}{ll}
\hline
$0{,}1$ & \texttt{0x1.999999999999ap-4} \\
$6{,}0$ & \texttt{0x1.8000000000000p+2} \\
$10^{-8}$ & \texttt{0x1.5798ee2308c3ap-27} \\
$5\times10^{-324}$ & \texttt{0x0.0000000000001p-1022} \\
\hline
\end{tabular}
\end{center}

Bentuk itu memperlihatkan tiga hal. Digit di belakang titik heksadesimal selalu tiga
belas, dan $13\times4=52$, yaitu lebar mantissa binary64. Digit terakhir $0{,}1$
adalah \texttt{a}, yaitu $1010$, hasil pembulatan ke atas dari pola $1001$ yang
berulang. Baris terakhir punya nol di depan titik: tanda bilangan subnormal.

Nilai khusus melengkapi sistem ini: $\pm\infty$ untuk hasil yang melampaui
$e_{\max}$, nol bertanda, bilangan subnormal di bawah $2^{e_{\min}}$, dan NaN
(\emph{not a number}) untuk bentuk tak tentu seperti $0/0$. Satu sifat NaN layak
diingat karena sering mengejutkan: ia tidak sama dengan dirinya sendiri, sehingga
perbandingan \texttt{x == x} bernilai salah justru ketika $x$ adalah NaN.\index{NaN}

\section{Satu-satunya jembatan: pembulatan}
\tingkat{dasar}

Tidak ada operasi lain yang menghubungkan $\mathbb{R}$ dengan $\mathbb{F}$ selain
pembulatan. Karena itu seluruh analisis galat berdiri di atas satu pertanyaan: seberapa
jauh bilangan mesin terdekat dari bilangan aslinya, dan bagaimana jarak itu berskala.

IEEE-754 menetapkan aturan \emph{round to nearest, ties to even}: pilih bilangan mesin
terdekat, dan bila dua kandidat sama jauhnya pilih yang digit terakhir mantissanya
genap. Aturan kedua bukan kemewahan teoretis. Tanpa aturan itu, pembulatan pada titik
seri selalu condong ke satu arah, dan galat sejuta operasi akan menumpuk ke arah yang
sama, bukan saling meniadakan.

Titik seri itu benar-benar muncul, dan kita bisa membuatnya dengan sengaja. Dalam
binary64, jarak antara $1$ dan bilangan sesudahnya adalah $2^{-52}$. Maka
$1+2^{-53}$ jatuh tepat di tengah, dan aturan genap memilih $1$ karena mantissanya
berakhiran nol. Python memberi

\begin{center}
\texttt{1 + 2.0**-53 = 1.0}, sedangkan \texttt{1 + 3*2.0**-53 = 1.0000000000000004}.
\end{center}

Kasus pertama membulat ke bawah, kasus kedua ke atas, padahal keduanya sama-sama tepat
setengah jalan. Yang membedakan hanya paritas digit terakhir. Inilah sebabnya
$1+10^{-16}=1$ dalam binary64: bukan karena $10^{-16}$ diabaikan dengan sembarangan,
tetapi karena tidak ada tempat untuk menyimpannya.

\kamus{epsilon mesin}{machine epsilon}{Jarak antara $1$ dan bilangan mesin
berikutnya: ukuran satu langkah di skala satu.}\index{epsilon mesin}

Dua angka sering tertukar dan pertukarannya mengubah batas galat dengan faktor dua.
Epsilon mesin $\mathrm{eps}=2^{-p+1}$ adalah jarak antara $1$ dan bilangan
sesudahnya, yaitu $2^{-52}$ untuk binary64. Satuan pembulatan
$u=2^{-p}=2^{-53}$ adalah setengah jarak itu, dan hanya $u$ yang muncul dalam batas
galat. Python menyebut yang pertama \texttt{numpy.finfo(float).eps}; yang kedua harus
kita hitung sendiri.

Jarak antar bilangan tidak seragam. Di dekat $1$ jaraknya $2^{-52}$, di dekat
$10^{10}$ jaraknya $1{,}9\times10^{-6}$. Yang seragam adalah jarak itu relatif
terhadap besar bilangannya, dan itulah isi lemma berikut.

\begin{lemma}[Ulp]
Jika $2^{e}\le|x|<2^{e+1}$, maka jarak antara $x$ dan bilangan mesin berikutnya adalah
$\operatorname{ulp}(x)=2^{e-p+1}$.
\end{lemma}

\begin{proof}
Pada selang itu seluruh bilangan normal berbentuk $M\cdot2^{e-p+1}$ dengan $M$ bulat
dan $2^{p-1}\le M<2^{p}$; banyaknya $M$ adalah $2^{p-1}$ dan selisih dua $M$
berurutan adalah satu satuan $2^{e-p+1}$. Selang berikutnya memakai $e+1$ sehingga
jaraknya berlipat dua, dan itulah satu-satunya titik tempat kerapatan berubah.
\end{proof}

\kamus{ulp}{unit in the last place}{Satu langkah terkecil yang bisa dibedakan mesin
pada skala sebuah bilangan.}\index{ulp}

Lemma itu menjelaskan dua hal sekaligus.
$\operatorname{ulp}(1)=2^{-52}=2{,}220446049250313\times10^{-16}$, sedangkan
$\operatorname{ulp}(10^{10})=2^{-52}\cdot2^{33}\approx1{,}9\times10^{-6}$. Jadi
menambahkan $1$ kepada $10^{16}$ tidak mengubah apa pun dalam binary64, karena
$\operatorname{ulp}(10^{16})=2$: memang tidak ada bilangan mesin di antara keduanya.

Gambar~\ref{fig:titikmesin} memperlihatkan akibat lemma itu pada empat oktaf
pertama: setiap oktaf memuat $2^{23}$ titik, tetapi lebarnya berlipat dua, sehingga
jarak antar titik berlipat dua pula.

\begin{figure}[htbp]
\centering
\begin{tikzpicture}[x=2.1cm,y=1cm,every node/.style={font=\footnotesize}]
  \draw[ink,thick,->] (-0.55,0) -- (4.30,0);
  \draw[ink,thick,fill=white] (-0.42,0) circle (1.7pt);
  \draw[ink,thick] (-0.48,-0.09) -- (-0.36,0.09);
  \foreach \x in {0,0.125,...,0.875} {\draw[ink] (\x,0.055) -- (\x,-0.055);}
  \foreach \x in {1,1.25,...,1.875} {\draw[ink] (\x,0.055) -- (\x,-0.055);}
  \foreach \x in {2,2.5,...,3.0} {\draw[ink] (\x,0.055) -- (\x,-0.055);}
  \foreach \x in {3,4} {\draw[ink] (\x,0.055) -- (\x,-0.055);}
  \foreach \x/\lab in {0/{$2^{0}$},1/{$2^{1}$},2/{$2^{2}$},3/{$2^{3}$},4/{$2^{4}$}} {
    \draw[ink,thick] (\x,0.14) -- (\x,-0.14);
    \node[below=1pt] at (\x,-0.14) {\lab};
  }
  \draw[{Latex}-{Latex},accent] (0.125,-0.62) -- (0.25,-0.62) node[midway,below] {$d$};
  \draw[{Latex}-{Latex},accent] (1.25,-0.62) -- (1.50,-0.62) node[midway,below] {$2d$};
  \draw[{Latex}-{Latex},accent] (2.50,-0.62) -- (3.00,-0.62) node[midway,below] {$4d$};
  \draw[{Latex}-{Latex},accent] (3.00,-0.62) -- (4.00,-0.62) node[midway,below] {$8d$};
  \node[align=center,text width=3.0cm] at (1.95,0.92)
    {tiap oktaf berisi $2^{23}$ titik, lebarnya dua kali oktaf sebelumnya};
  \node[align=center,text width=3.0cm] at (0.40,-1.32)
    {ke kiri menuju nol: jarak mengecil};
\end{tikzpicture}
\caption{Empat oktaf pertama bilangan mesin binary32. Setiap oktaf memuat $2^{23}$
titik, tetapi lebarnya dua kali oktaf sebelumnya, jadi jarak antar titik berlipat dua:
$d=2^{-23}$ pada $[1,2)$, $2d$ pada $[2,4)$, $4d$ pada $[4,8)$, dan $8d$ pada
$[8,16)$. Titik di ekor tidak lebih sedikit, hanya lebih jarang; ke arah nol jaraknya
mengecil sampai batas subnormal.}
\label{fig:titikmesin}
\end{figure}

Ada satu kehalusan yang perlu disebut sebelum kita memakai ulp sebagai satuan galat.
Galat pembulatan paling besar setengah ulp, tetapi bila dinyatakan sebagai galat
relatif, angka itu bukan satu bilangan melainkan satu selang. Untuk $x$ yang tepat
sama dengan pangkat dua, setengah ulp relatifnya $2^{-53}$; untuk $x$ yang hampir
mencapai pangkat dua berikutnya, nilainya mendekati $2^{-54}$. Jadi galat relatif satu
pembulatan bisa berayun dengan faktor dua. Goldberg menamai ayunan itu \emph{wobble}, dan
karena itu batas $u$ selalu dipakai sebagai batas terburuk, bukan sebagai nilai khas.

\begin{teorema}[Galat pembulatan relatif]
Untuk setiap $x\in\mathbb{R}$ dalam jangkauan normal, pembulatan benar memenuhi
$\operatorname{fl}(x)=x(1+\delta)$ dengan $|\delta|\le u=2^{-p}$.
\end{teorema}

\begin{proof}
Ambil $2^{e}\le|x|<2^{e+1}$. Bilangan mesin terdekat berjarak paling jauh setengah
ulp, yaitu $2^{e-p}$, sehingga
$|\operatorname{fl}(x)-x|\le2^{e-p}=u\cdot2^{e}\le u|x|$. Membagi dengan $|x|$
memberi $|\delta|\le u$, dan tanda $\delta$ mengikuti arah pembulatan.
\end{proof}

Bentuk $x(1+\delta)$ adalah awal dari segalanya: ia mengubah pertanyaan ``seberapa
salah'' menjadi ``berapa banyak faktor $(1+\delta)$ yang menumpuk''. Operasi
$+,-,\times,\div,\sqrt{\phantom{x}}$ dalam IEEE-754 wajib mengembalikan pembulatan
benar dari hasil eksaknya, jadi setiap operasi menyumbang satu faktor
semacam itu dan tidak lebih. Jaminan inilah yang membuat analisis galat mungkin.


\section{Galat, bilangan kondisi, dan pembatalan}
\tingkat{menengah}

Sekarang kita punya satuan. Yang belum ada adalah ukuran seberapa besar galat sebuah
perhitungan dibandingkan galat masukannya. Dua masalah bisa sama-sama menerima
masukan yang salah $10^{-16}$ dan keluar dengan jawaban yang salah $10^{-16}$ atau
$10^{-1}$, bergantung pada masalahnya sendiri.

\kamus{bilangan kondisi}{condition number}{Pengali yang menunjukkan berapa kali galat
masukan diperbesar oleh sebuah perhitungan.}\index{bilangan kondisi}

\begin{definisi}[Galat absolut, relatif, dan ke belakang]
Misalkan $\tilde{x}$ penghampiran bagi $x$. Galat absolutnya $|\tilde{x}-x|$, galat
relatifnya $|\tilde{x}-x|/|x|$ untuk $x\neq0$. Jika $\tilde{x}$ adalah jawaban eksak
bagi masalah dengan data $x+\Delta x$, maka $\|\Delta x\|/\|x\|$ disebut galat ke
belakang, dan algoritma yang selalu memberi galat ke belakang kecil disebut
\emph{stabil ke belakang}.\index{galat relatif}\index{stabil ke belakang}
\end{definisi}

Pembedaan terakhir paling berguna dalam praktik. Algoritma stabil ke belakang boleh
salah jauh, asalkan jawabannya benar untuk masalah yang sedikit berbeda:
kesalahannya milik masalah. Galat ke belakang yang besar adalah kesalahan penulis
kode.

Untuk fungsi satu peubah, pembesaran itu dihitung langsung dengan turunan: bilangan
kondisi adalah kemiringan fungsi di titik itu, dikali besar titiknya, dibagi besar
nilainya.

\begin{teorema}[Bilangan kondisi]
Misalkan $f$ terturunkan pada selang yang memuat $x$ dan $f(x)\neq0$. Bila
$\tilde{x}=x(1+\delta)$ dengan $|\delta|\le u$, maka
\[
  \frac{|f(\tilde{x})-f(x)|}{|f(x)|}\;\le\;\kappa_{f}(x)\,u+O(u^{2}),
  \qquad \kappa_{f}(x)=\frac{|x\,f'(x)|}{|f(x)|}.
\]
\end{teorema}

\begin{proof}
Menurut teorema nilai rata-rata, ada $\xi$ di antara $x$ dan $\tilde{x}$ dengan
$f(\tilde{x})-f(x)=f'(\xi)(\tilde{x}-x)$. Karena $|\tilde{x}-x|\le u|x|$ dan $f'$
kontinu di $x$, berlaku $|f'(\xi)|=|f'(x)|+O(u)$. Maka
$|f(\tilde{x})-f(x)|\le u|x|\,|f'(x)|+O(u^{2}|x|)$. Membagi dengan $|f(x)|$ memberi
ketaksamaan yang diklaim.
\end{proof}

Bilangan $\kappa_{f}$ tidak bergantung pada mesin, hanya pada masalahnya. Untuk
$f(x)=\sqrt{x}$ ia bernilai $1/2$: akar kuadrat tidak pernah memperburuk galat relatif.
Untuk $f(x)=\log x$ ia bernilai $1/|\log x|$, yang besar ketika $x$ dekat $1$ karena
di situ $\log x$ hampir nol. Untuk $f(x)=x-c$ dengan $c$ konstanta besar, ia bernilai
$|x|/|x-c|$, yang membesar tanpa batas ketika $x$ mendekati $c$. Pola terakhir itulah
yang menjadi sumber seluruh persoalan bab ini.

Pengurangan dua bilangan yang berdekatan adalah operasi paling berbahaya dalam
aritmetika mesin, dan bahayanya sering salah alamat. Kesalahpahaman lazim berbunyi:
pengurangan kehilangan digit, jadi operasi pengurangan itu sendiri tidak akurat. Yang
benar lebih halus dan lebih menenteramkan: bila kedua bilangan berdekatan,
pengurangannya justru eksak.

\begin{teorema}[Sterbenz]
Jika $a,b\in\mathbb{F}$ dan $b/2\le a\le2b$ dengan $b>0$, maka $a-b$ terwakili persis
di $\mathbb{F}$.
\end{teorema}

\begin{proof}
Misalkan $2^{e}\le b<2^{e+1}$. Setiap bilangan mesin pada selang itu kelipatan bulat
$2^{e-p+1}$, jadi $a$ juga, sehingga $a-b$ kelipatan $2^{e-p+1}$. Dari
$b/2\le a\le2b$ kita peroleh $0\le a-b\le b/2<2^{e}$. Kelipatan $2^{e-p+1}$ yang
terletak di $[0,2^{e})$ berbentuk $M\cdot2^{e-p+1}$ dengan $0\le M<2^{p}$, dan
semuanya terwakili dengan $p$ digit. Jadi hasilnya eksak.
\end{proof}

Jadi bila hasil pengurangan tampak buruk, yang buruk adalah masukannya. Angka yang
``hilang'' sudah hilang sebelum pengurangan dilakukan: $a$ dan $b$ masing-masing hanya
diketahui sampai galat relatif $u$, dan galat mutlak sebesar $u|a|$ itu tidak ikut
menyusut ketika kedua bilangan saling meniadakan. Yang menyusut hanya besar hasilnya,
sehingga galat relatifnya meledak.

\begin{proposisi}[Bilangan kondisi pengurangan]
Misalkan $x=a-b\neq0$ dan mesin bekerja dengan $\tilde{a}=a(1+\delta_{a})$,
$\tilde{b}=b(1+\delta_{b})$, $|\delta_{a}|,|\delta_{b}|\le u$. Maka
\[
  \frac{|\operatorname{fl}(\tilde{a}-\tilde{b})-x|}{|x|}
  \;\le\;\kappa\,u+u,
  \qquad \kappa=\frac{|a|+|b|}{|a-b|}.
\]
\end{proposisi}

\begin{proof}
Karena $\operatorname{fl}(\tilde{a}-\tilde{b})$ adalah pembulatan benar dari
$\tilde{a}-\tilde{b}$, galat tambahannya paling banyak $u|\tilde{a}-\tilde{b}|$.
Selanjutnya
$|\tilde{a}-\tilde{b}-(a-b)|\le|a||\delta_{a}|+|b||\delta_{b}|\le u(|a|+|b|)$,
dan membagi dengan $|x|=|a-b|$ memberi $u\kappa$. Menjumlahkan kedua galat dengan
ketaksamaan segitiga menyelesaikan bukti.
\end{proof}

Bilangan $\kappa$ tidak pernah lebih kecil dari satu dan membesar tanpa batas ketika
$|a-b|$ mengecil. Untuk $a=12345{,}678$ dan $b=12345{,}671$ kita peroleh $a-b=0{,}007$
dan $\kappa\approx3{,}53\times10^{6}$, sehingga galat relatif hasil bisa mencapai
$3{,}9\times10^{-10}$ walaupun setiap masukan hanya menyimpang $10^{-16}$. Karena setiap
faktor sepuluh dalam $\kappa$ menghapus satu digit signifikan,
$\log_{10}\kappa\approx6{,}5$ berarti enam setengah dari enam belas digit mesin hilang,
dan hasil $0{,}007$ hanya menyisakan sekitar sepuluh digit yang bisa dipercaya.

Satu catatan yang menyelamatkan: aturan Sterbenz berlaku untuk selisih, bukan untuk
penjumlahan. Untuk penjumlahan $a+b$ bilangan kondisi adalah
$(|a|+|b|)/|a+b|$, yang mendekati satu ketika $a$ dan $b$ bertanda sama. Menjumlahkan
bilangan besar tidak berbahaya. Mengurangkan dua bilangan besar yang hampir sama
berbahaya, dan itulah satu-satunya kebiasaan yang perlu diwaspadai.

\section{Pembatalan katastrofik: dua wajah}
\tingkat{menengah}

Kita sudah tahu mekanismenya. Sekarang kita beri namanya, karena dua jenis pembatalan
yang mirip sering dianggap sama padahal akibatnya berbeda jauh.

\begin{untukanak}
Kalau kita ingin tahu tebal selembar kertas, penggaris biasa tidak akan menolong: satu
lembar jauh lebih tipis daripada satu milimeter, jadi jarumnya seolah tidak bergerak.
Caranya bukan mencari penggaris yang lebih teliti, melainkan mengubah apa yang diukur:
tumpuk lima ratus lembar, ukur tebal tumpukannya, lalu bagi lima ratus. Hasilnya jauh
lebih dapat dipercaya, walaupun penggarisnya tetap penggaris yang sama.

Ternyata rumus juga bisa diperlakukan seperti itu. Ada rumus yang bekerja dengan
mengurangkan dua angka besar yang hampir sama, dan itu sama saja dengan menimbang
satu lembar kertas dengan penggaris. Ada rumus lain yang secara matematika sama
persis, tetapi menyusun langkahnya supaya angka kecil diukur langsung, bukan dicari
dari selisih. Hasilnya bisa berbeda jauh, padahal kertasnya tidak berubah.

Pertanyaannya: kalau dua rumus yang sama-sama benar memberi jawaban berbeda, mana yang
harus kita percaya? Di bab ini kita akan tahu mengapa rumus yang salah bukan yang
hitungannya keliru, melainkan yang memilih langkah yang membuang keterangan, dan
bagaimana mengenalinya sebelum program dijalankan.
\end{untukanak}

\kamus{pembatalan katastrofik}{catastrophic cancellation}{Hilangnya digit yang berarti
ketika dua bilangan hampir sama dikurangkan, sehingga yang tersisa hanya galat.}
\index{pembatalan katastrofik}

\begin{definisi}[Pembatalan katastrofik dan jinak]
Misalkan $x$ dan $y$ adalah bilangan mesin yang dihitung dari data sebelumnya.
Pengurangan $x-y$ disebut \emph{pembatalan jinak} bila $x$ dan $y$ tepat, sehingga
hasilnya eksak menurut teorema Sterbenz. Ia disebut \emph{pembatalan katastrofik} bila
$x$ dan $y$ sendiri sudah membawa galat pembulatan, sehingga digit yang sama-sama
mereka miliki ikut hilang dan yang tersisa didominasi galat itu.
\end{definisi}

Contoh jinak: $1{,}0000000000000002-1$ dalam binary64. Keduanya bilangan mesin,
selisihnya $2^{-52}$ dan terwakili persis. Hasilnya kecil, tetapi tidak ada satu digit
pun yang hilang: kita memang tahu kedua masukannya sampai bit terakhir.

Contoh katastrofik adalah rumus kuadrat di awal bab. Di situ $b^{2}$ dan $4ac$
masing-masing hasil perkalian, jadi keduanya dibulatkan. Untuk $a=1$, $b=10^{8}$,
$c=1$ keduanya kebetulan tepat, dan justru kekhususan itulah yang membuat contoh ini
layak diperiksa sampai bit: keruntuhannya bukan berasal dari perkalian, melainkan dari
pembulatan akar kuadratnya. Python menghitung

\begin{center}
\begin{tabular}{ll}
\hline
diskriminan $b^{2}-4ac$ & \texttt{9999999999999996.0} \\
$\sqrt{b^{2}-4ac}$ & \texttt{99999999.99999999} \\
$b-\sqrt{b^{2}-4ac}$ & \texttt{1.4901161193847656e-08} \\
$\operatorname{ulp}(b)$ & \texttt{1.4901161193847656e-08} \\
\hline
\end{tabular}
\end{center}

Nilai sebenarnya $b-\sqrt{b^{2}-4ac}$ adalah $2\times10^{-8}$, yaitu $1{,}34$ ulp.
Yang terhitung tepat $1$ ulp, dan galat setengah ulp pada akar kuadrat sudah cukup
untuk itu: akar kuadratnya benar sampai $0{,}46$ ulp pada skala $10^{8}$, tetapi
selisih yang seharusnya $1{,}34$ ulp hanya bisa mendarat pada $1$ atau $2$ ulp, jadi
informasinya tinggal satu pilihan biner. Hasil akhirnya $-7{,}45\times10^{-9}$, yaitu
$-2^{-27}$, sedangkan akar sebenarnya $-10^{-8}$. Galat relatif $0{,}255$.

Perhatikan pembagian tugasnya. Akar kuadrat bekerja dengan galat relatif
$5\times10^{-17}$, jauh di bawah $u$. Masalahnya sendiri berkondisi baik: menggeser
$b$ sebesar satu ulp hanya menggeser akar kecil sebesar satu ulp, jadi
$\kappa\approx|b|/\sqrt{b^{2}-4ac}\approx1$. Yang buruk adalah \emph{bentuk} yang kita
pilih: bilangan kondisi pengurangan pada langkah itu adalah
\[
  \kappa=\frac{|b|+\sqrt{b^{2}-4ac}}{\left|\sqrt{b^{2}-4ac}-|b|\right|}
  =\frac{2\times10^{8}}{1{,}49\times10^{-8}}\approx1{,}34\times10^{16},
\]
sehingga $\kappa u\approx1{,}49$: batas galat relatifnya melewati satu, artinya tidak
ada satu digit pun yang dijamin.

Gambar~\ref{fig:katastrofik} memisahkan dua hal yang sering tercampur pada langkah itu:
besar bilangannya, dan banyaknya digit yang masih berarti setelah pengurangan.

\begin{figure}[htbp]
\centering
\begin{tikzpicture}[every node/.style={font=\footnotesize}]
  \draw[ink,thick,->] (1.2,3.20) -- (10.6,3.20);
  \foreach \dec/\tk in {-9/{$10^{-9}$},-6/{$10^{-6}$},-3/{$10^{-3}$},0/{$10^{0}$},
                        3/{$10^{3}$},6/{$10^{6}$},9/{$10^{9}$}} {
    \draw[ink] ({1.4+(\dec+9)*0.5},3.26) -- ({1.4+(\dec+9)*0.5},3.14);
    \node[below] at ({1.4+(\dec+9)*0.5},3.12) {\tk};
  }
  \fill[accent] ({1.4+17*0.5},3.20) circle (2.4pt);
  \node[above=2pt] at ({1.4+17*0.5},3.24) {$b=10^{8}$};
  \fill[layerteal] ({1.4+1.173*0.5},3.20) circle (2.4pt);
  \node[above=2pt] at ({1.4+1.173*0.5},3.24) {$b-\sqrt{\Delta}$};
  \draw[<->,thick,dashed,layergrey] ({1.4+1.173*0.5},3.72) -- ({1.4+17*0.5},3.72);
  \node[above] at ({1.4+9.1*0.5},3.74) {enam belas orde besaran};
  \foreach \i in {0,...,16} {
    \fill[accent!28] ({1.4+\i*0.45},1.60) rectangle ++(0.45,0.40);
    \draw[ink!45] ({1.4+\i*0.45},1.60) rectangle ++(0.45,0.40);
  }
  \foreach \i in {0,...,15} {
    \fill[layergrey!16] ({1.4+\i*0.45},0.85) rectangle ++(0.45,0.40);
    \draw[ink!45] ({1.4+\i*0.45},0.85) rectangle ++(0.45,0.40);
  }
  \fill[accent!70] ({1.4+16*0.45},0.85) rectangle ++(0.45,0.40);
  \draw[ink!45] ({1.4+16*0.45},0.85) rectangle ++(0.45,0.40);
  \node[left] at (1.35,1.80) {$b$};
  \node[left] at (1.35,1.05) {$b-\sqrt{\Delta}$};
  \draw[{Latex}-{Latex},layergrey] (1.40,0.62) -- ({1.4+16*0.45},0.62);
  \node[below] at ({1.4+8*0.45},0.60) {enam belas digit sama, saling meniadakan};
  \node[right] at ({1.4+17*0.45},1.05) {satu digit};
\end{tikzpicture}
\caption{Atas: $b$ dan $\sqrt{b^{2}-4ac}$ pada skala log, berjarak enam belas orde
besaran dari selisihnya. Bawah: kotak adalah digit berarti, tujuh belas masuk dan enam
belas di antaranya sama sehingga saling meniadakan. Yang tersisa satu digit, yaitu
selisih $1{,}49\times10^{-8}$ itu sendiri; karena setengah ulp galat pada akar kuadrat
cukup untuk menggeser digit terakhir, akar kecil keluar dengan galat relatif $0{,}255$.}
\label{fig:katastrofik}
\end{figure}

Perbaikannya bukan menambah presisi, melainkan memilih akar yang dihitung lebih dulu.
Hasil kali kedua akar adalah $c/a$, jadi setelah akar besar $r_{1}$ diketahui, akar
kecil diperoleh dari pembagian, bukan dari pengurangan:
\[
  r_{1}=\frac{-b-\sqrt{b^{2}-4ac}}{2a},\qquad r_{2}=\frac{c}{a\,r_{1}}
  =\frac{2c}{-b-\sqrt{b^{2}-4ac}}.
\]
Bentuk kedua tidak pernah mengurangkan dua bilangan berdekatan. Untuk contoh di atas
ia memberi $-10^{-8}$ tepat, yaitu akar yang dibulatkan benar. Rumus yang sama dengan
tanda diatur menurut tanda $b$ dipakai di hampir setiap pustaka yang menyelesaikan
persamaan kuadrat, dan alasannya bukan kerapian melainkan kestabilan.

Wajah kedua pembatalan katastrofik muncul di tempat yang tampaknya tidak berbahaya:
selisih $1-\cos x$ untuk $x$ kecil. Untuk $x$ kecil,
\[
  \frac{1-\cos x}{x^{2}}\longrightarrow\frac{1}{2},
\]
karena $\cos x=1-x^{2}/2+x^{4}/24-\cdots$. Rumus itu benar, dan rumus itu juga yang
paling wajar ditulis. Tetapi $\cos x$ disimpan dengan galat relatif $u$, dan karena
$|\cos x|\le1$, galat mutlaknya sekitar $u$. Setelah dibagi $x^{2}$, galat itu menjadi
galat relatif sekitar $2u/x^{2}$. Untuk $x=10^{-8}$ faktor itu bernilai $2{,}2$:
seluruh digit lenyap.

Berikut hasil eksekusi nyata untuk sepuluh nilai $x$, dibandingkan dengan acuan yang
dihitung pada 50 digit dan dengan deret Taylor $0{,}5-x^{2}/24+x^{4}/720$; untuk
$x\le10^{-3}$ kedua pembanding itu sepakat sampai digit terakhir. Kolom ``lugu'' adalah
$(1-\cos x)/x^{2}$, kolom ``identitas'' adalah
$\tfrac12\left(\sin(x/2)/(x/2)\right)^{2}$ yang berasal dari
$1-\cos x=2\sin^{2}(x/2)$.

\begin{center}
\begin{tabular}{rccc}
\hline
$x$ & lugu & galat relatif & identitas \\
\hline
$10^{-1}$ & $0{,}499583472197418$ & $1{,}1\times10^{-14}$ & $0{,}499583472197423$ \\
$10^{-2}$ & $0{,}499995833347366$ & $2{,}9\times10^{-13}$ & $0{,}499995833347222$ \\
$10^{-3}$ & $0{,}499999958325503$ & $1{,}6\times10^{-11}$ & $0{,}499999958333335$ \\
$10^{-4}$ & $0{,}499999996961265$ & $5{,}2\times10^{-9}$ & $0{,}499999999583333$ \\
$10^{-5}$ & $0{,}500000041370185$ & $8{,}3\times10^{-8}$ & $0{,}499999999995833$ \\
$10^{-6}$ & $0{,}500044450291171$ & $8{,}9\times10^{-5}$ & $0{,}499999999999958$ \\
$10^{-7}$ & $0{,}499600361081320$ & $8{,}0\times10^{-4}$ & $0{,}500000000000000$ \\
$10^{-8}$ & $0{,}000000000000000$ & $1{,}0$ & $0{,}500000000000000$ \\
$10^{-9}$ & $0{,}000000000000000$ & $1{,}0$ & $0{,}500000000000000$ \\
$10^{-10}$ & $0{,}000000000000000$ & $1{,}0$ & $0{,}500000000000000$ \\
\hline
\end{tabular}
\end{center}

Polanya khas dan layak dihafal. Selama $x^{2}/2$ jauh di atas $u$, rumus lugu hampir
benar. Ketika $x^{2}/2$ menyusut ke skala $u$, digit berartinya habis dan galat
membesar mengikuti $2u/x^{2}$: setiap penurunan $x$ dengan faktor sepuluh menambah dua
digit kehilangan. Di $x=10^{-8}$, $1-\cos x$ lebih kecil daripada setengah ulp di
sekitar $1$, sehingga $\cos x$ dibulatkan tepat ke $1$ dan hasilnya nol. Titik runtuh
itu bukan misteri: $x^{2}/2\approx u$, yaitu $x\approx\sqrt{2u}\approx1{,}5\times10^{-8}$.

Kolom terakhir menunjukkan bahwa masalahnya berkondisi baik. Identitas sudut ganda
memberi jawaban benar sampai digit terakhir untuk setiap $x$ di tabel, karena ia tidak
pernah mengurangkan dua bilangan berdekatan. Galat bukan sifat masalah. Ia sifat rumus
yang dipilih.

\begin{resep}
Empat langkah memilih bentuk rumus yang aman dari pemusnahan:
\begin{enumerate}
  \item Tandai setiap pengurangan $a-b$ di dalam rumus, lalu hitung
        $\kappa=(|a|+|b|)/|a-b|$ pada data terburuk yang masuk akal.
  \item Terjemahkan $\kappa$ menjadi digit yang hilang: $\kappa u\approx10^{-d}$
        berarti $d$ digit terakhir tidak dapat dipercaya.
  \item Cari bentuk setara yang tidak memuat selisih itu: kalikan dengan sekawan,
        pakai identitas sudut ganda, faktorkan, atau panggil \texttt{log1p},
        \texttt{expm1}, dan \texttt{hypot}.
  \item Ulangi hitungan pada presisi lebih tinggi. Selisih kedua hasil harus sepadan
        dengan $\kappa u$; bila jauh lebih besar, bentuk barunya masih menyembunyikan
        pengurangan.
\end{enumerate}
Peringatan urutan: langkah 1 tidak boleh ditukar dengan langkah 3. Menulis ulang rumus
sebelum tahu pengurangan mana yang berbahaya hanya memindahkan masalah, seperti pada
bentuk satu-lintasan variansi yang menukar satu pengurangan berbahaya dengan
pengurangan berbahaya yang lain.
\end{resep}

\section{Penjumlahan: galat yang tumbuh bersama $n$}
\tingkat{lanjut}

Penjumlahan adalah operasi yang paling sering dianggap remeh, padahal ia tempat galat
paling mudah dihitung. Sebelum sampai ke batasnya, satu demonstrasi yang membongkar
anggapan bahwa urutan tidak penting:

\begin{center}
\texttt{(0.1 + 0.2) + 0.3 = 0.6000000000000001}, \quad
\texttt{0.1 + (0.2 + 0.3) = 0.6}.
\end{center}

Selisihnya $1{,}11\times10^{-16}$, tepat satu ulp pada skala $0{,}6$. Tidak ada
kesalahan tulis; hanya tiga bilangan yang masing-masing tidak punya wakil biner
berhingga, dan pembulatan yang berbeda urutan. Dalam matematika, penjumlahan bersifat
asosiatif. Di dalam mesin, ia asosiatif hanya bila kebetulan tidak ada pembulatan yang
berbeda.

\begin{teorema}[Galat penjumlahan berurutan]
Misalkan $s_{1}=x_{1}$ dan $s_{k}=\operatorname{fl}(s_{k-1}+x_{k})$ untuk
$k=2,\dots,n$ dengan $x_{i}\in\mathbb{F}$. Maka
\[
  \left|s_{n}-\sum_{i=1}^{n}x_{i}\right|
  \;\le\;(n-1)\,u\sum_{i=1}^{n}|x_{i}|+O(u^{2}).
\]
\end{teorema}

\begin{proof}
Setiap langkah memakai satu pembulatan, jadi $s_{k}=(s_{k-1}+x_{k})(1+\delta_{k})$
dengan $|\delta_{k}|\le u$. Menguraikan hasilnya secara rekursif memberi
\[
  s_{n}=\sum_{i=1}^{n}x_{i}\prod_{k=i}^{n}(1+\delta_{k}),
\]
dengan suku pertama tidak memuat $\delta_{1}$ karena $s_{1}=x_{1}$ belum dibulatkan.
Karena $\prod(1+\delta_{k})=1+\theta$ dengan
$|\theta|\le(1+u)^{n-1}-1\le(n-1)u+O(u^{2})$, ketaksamaan segitiga memberi batas yang
diklaim.
\end{proof}

Batas itu jujur tetapi longgar: ia menganggap setiap pembulatan mencapai galat
terburuknya dan semuanya searah, yang praktis tidak pernah terjadi. Tabel berikut
menunjukkan jurang antara batas dan kenyataan untuk deret harmonik
$H_{n}=\sum_{i=1}^{n}1/i$ dalam binary32, dengan nilai eksak dihitung oleh
\texttt{math.fsum} yang menjumlahkan secara eksak.

\begin{center}
\begin{tabular}{rccc}
\hline
$n$ & berurutan & eksak & batas $(n-1)u\sum|x_{i}|$ \\
\hline
$10^{2}$  & $5{,}1873779297$ & $5{,}1873775176$ & $3{,}1\times10^{-5}$ \\
$10^{3}$  & $7{,}4854784012$ & $7{,}4854708606$ & $4{,}5\times10^{-4}$ \\
$10^{4}$  & $9{,}7876129150$ & $9{,}7876060360$ & $5{,}8\times10^{-3}$ \\
$10^{5}$  & $12{,}0908508301$ & $12{,}0901461299$ & $7{,}2\times10^{-2}$ \\
$10^{6}$  & $14{,}3573579788$ & $14{,}3927267229$ & $8{,}6\times10^{-1}$ \\
\hline
\end{tabular}
\end{center}

Galat sesungguhnya pada $n=10^{6}$ adalah $3{,}54\times10^{-2}$, sekitar dua puluh lima
kali lebih kecil daripada batas, tetapi sudah menghapus dua setengah digit signifikan.
Galatnya juga tidak naik mulus: ia sempat turun pada $n=10^{4}$ karena pembulatan
searah dan berlawanan kebetulan saling meniadakan. Yang bisa diandalkan bukan besar
galat pada satu $n$, melainkan kecenderungannya membesar, dengan laju terburuk
sebanding $n$.

Penyebabnya bukan pembulatan yang besar, melainkan akumulator yang menjauh dari nol.
Begitu $s$ bernilai sekitar $14$, suku $1/i$ untuk $i$ besar berada di bawah setengah
ulp $s$ dan setiap penambahan kehilangan seluruh informasinya. Kasus ekstremnya mudah
dibuat: jumlahkan $1$ dengan sejuta suku bernilai $10^{-16}$; secara eksak hasilnya
$1{,}0000000001$, sedangkan penjumlahan berurutan tetap $1$ untuk berapa pun banyaknya
suku. Yang hilang bukan satu suku, melainkan seluruh kontribusi.

Gambar~\ref{fig:galatjumlah} merangkum pola yang sama untuk kelima nilai $n$ pada tabel
tadi, kali ini dengan penjumlahan Kahan sebagai pembanding.

\begin{figure}[htbp]
\centering
\begin{tikzpicture}
\begin{axis}[
  width=0.88\linewidth, height=6.0cm,
  xmode=log, ymode=log,
  xlabel={$n$}, ylabel={$|s_{n}-H_{n}|$},
  xmin=70, xmax=1.6e6, ymin=1e-8, ymax=1,
  grid=both, grid style={layergrey!15},
  tick label style={font=\footnotesize},
  label style={font=\footnotesize},
  legend style={font=\footnotesize,draw=layergrey!40,
                at={(0.03,0.97)},anchor=north west},
  legend cell align=left,
]
\addplot[ink,thick,mark=*,mark size=1.7pt] coordinates {
  (100,4.12e-7) (1000,7.54e-6) (10000,6.88e-6) (100000,7.05e-4) (1000000,3.54e-2)};
\addlegendentry{berurutan}
\addplot[accent,thick,mark=square*,mark size=1.6pt] coordinates {
  (100,6.48e-8) (1000,8.88e-8) (10000,2.03e-7) (100000,6.51e-8) (1000000,1.75e-7)};
\addlegendentry{Kahan}
\end{axis}
\end{tikzpicture}
\caption{Galat mutlak jumlah harmonik dalam binary32 pada lima ukuran $n$, skala
ganda-log. Penjumlahan berurutan naik dari $4{,}1\times10^{-7}$ pada $n=10^{2}$ ke
$3{,}54\times10^{-2}$ pada $n=10^{6}$; Kahan bertahan di orde $10^{-7}$ sepanjang
rentang, sehingga pada $n=10^{6}$ keduanya terpaut $2\times10^{5}$ kali. Data, presisi,
dan banyaknya operasi sama; yang berbeda hanya urutan penjumlahannya.}
\label{fig:galatjumlah}
\end{figure}

\begin{tangan}
Bahan: sebelas kartu kosong, spidol, dan satu aturan main. Setiap kali dua kartu
dijumlahkan, tulis hasilnya pada kartu baru dengan paling banyak empat angka penting,
lalu buang kartu lamanya.

Tulis $1{,}000$ pada satu kartu dan $0{,}0001$ pada sepuluh kartu lain. Jalur pertama:
letakkan kartu besar, lalu tambahkan sepuluh kartu kecil satu per satu. Setiap kali
$1{,}000+0{,}0001=1{,}0001$ dipotong menjadi $1{,}000$, sehingga kartu kecil itu hilang
tanpa jejak dan tumpukan tetap $1{,}000$.

Kocok ulang kartunya, jalur kedua: jumlahkan dulu sepuluh kartu kecil menjadi
$0{,}0010$, baru tambahkan kartu besar. Tumpukan akhir $1{,}001$.

Yang harus terlihat: sebelas kartu yang sama, aturan pemotongan yang sama, dua hasil
yang berbeda pada digit terakhir. Yang berubah hanya urutan, dan mesin mengulang urutan
itu sejuta kali.
\end{tangan}

Penjumlahan Kahan menyerang masalah itu secara langsung.\index{penjumlahan Kahan}
Idenya sederhana: setiap kali
sebuah suku hilang karena pembulatan, simpan sisa yang hilang di variabel kompensasi,
lalu masukkan kembali pada langkah berikutnya.

\begin{quote}
\texttt{s = 0.0; c = 0.0}\\
\texttt{for x in data:}\\
\hspace*{1.5em}\texttt{y = x - c}\\
\hspace*{1.5em}\texttt{t = s + y}\\
\hspace*{1.5em}\texttt{c = (t - s) - y}\\
\hspace*{1.5em}\texttt{s = t}
\end{quote}

Baris \texttt{c = (t - s) - y} adalah intinya, dan ia bekerja justru karena Sterbenz:
$t-s$ adalah selisih dua bilangan berdekatan sehingga eksak, dan $(t-s)-y$ karena itu
tepat sama dengan galat yang baru saja dibuang. Hasilnya tidak lagi menumpuk galat
sebanding $n$.

\begin{teorema}[Batas penjumlahan Kahan]
Dengan algoritme Kahan di atas, jumlah yang terhitung memenuhi ketaksamaan
\[
  s=\sum_{i=1}^{n}x_{i}(1+\delta_{i})+O(nu^{2})\sum_{i=1}^{n}|x_{i}|,
  \qquad |\delta_{i}|\le2u .
\]
\end{teorema}

Bukti teorema ini tidak kita ulang di sini. Ia adalah Teorema 8 pada kitab Goldberg,
dan buktinya ada di lampiran kitab itu, diambil dari Knuth. Yang perlu dibaca dari
pernyataannya: galat relatif tiap suku dibatasi $2u$, dua kali lipat galat satu
pembulatan, tetapi faktor $n$ hilang dari suku utama. Untuk $n$ kecil batas ini justru
lebih longgar daripada batas lugu, karena $2u>u$. Keunggulannya baru muncul ketika
$n$ besar, dan di situlah seluruh gunanya.

Angka nyatanya pada deret harmonik dengan $n=10^{6}$: penjumlahan berurutan dalam
binary32 memberi galat $3{,}54\times10^{-2}$, sedangkan Kahan memberi
$1{,}75\times10^{-7}$, yaitu $0{,}18$ ulp dari jawaban benar. Dalam binary64
selisihnya sama tajam: berurutan meleset $7{,}35\times10^{-13}$, atau $414$ ulp,
sedangkan Kahan menghasilkan angka yang sama persis dengan \texttt{math.fsum}.

Dua perbaikan lain layak disebut. Penjumlahan berpasangan, yang dipakai
\texttt{numpy.sum},\index{penjumlahan berpasangan} membagi data secara rekursif
sehingga galatnya tumbuh seperti $\log n$; ia hampir gratis dan cocok untuk data di
memori. \texttt{math.fsum} memakai akumulator eksak dan memberi hasil yang dibulatkan
benar untuk seluruh daftar, dengan biaya waktu lebih tinggi. Menjumlahkan dari suku
terkecil ke terbesar adalah perbaikan termurah, tetapi ia hanya mengurangi galat.

\section{Panduan praktis}
\tingkat{dasar}

Semua yang sudah kita bangun bisa diringkas menjadi beberapa kebiasaan yang murah
dilakukan dan mahal kalau diabaikan.

Pertama, pilih presisi berdasarkan bilangan kondisi, bukan berdasarkan kebiasaan.
Galat relatif hasil kira-kira $\kappa u$. Dengan binary64, sebuah perhitungan
berkondisi $10^{6}$ masih menyisakan sepuluh digit; dengan binary32, perhitungan yang
sama hanya menyisakan satu digit. Untuk $10^{8}$ operasi berurutan batas kasarnya
$nu$, dan dalam binary32 itu sudah melewati satu, artinya jawabannya tidak bermakna.

Kedua, jangan menguji kesamaan bilangan pecahan dengan \texttt{==}. Dua lintasan yang
secara matematis identik boleh menghasilkan bit yang berbeda, seperti
$0{,}1+0{,}2$ tadi. Uji yang benar memakai toleransi berskala, misalnya
$|x-y|\le10^{-12}\max(1,|x|,|y|)$, dan toleransi itu harus berasal dari analisis
galat, bukan dari percobaan sampai uji lulus.

Ketiga, curigai setiap rumus yang memuat selisih dua besaran hampir sama. Bentuk baku
variansi, $(1-\cos x)/x^{2}$, dan $\log(1+x)$ untuk $x$ kecil semuanya bisa ditulis
ulang. Pustaka standar sudah melakukannya: \texttt{log1p}, \texttt{expm1}, dan
\texttt{hypot} ada justru karena bentuk lugunya tidak stabil.

Keempat, untuk jumlah yang panjang atau bersuku sangat berbeda besarnya, gunakan
\texttt{math.fsum} atau penjumlahan berpasangan, dan simpan akumulator dalam presisi
lebih tinggi bila tersedia. Urutkan dari suku terkecil bila tidak ada pilihan lain.

Kelima, jangan menyimpulkan kebenaran dari kesepakatan angka. Dua program yang sepakat
sampai dua belas digit cukup meyakinkan; yang sepakat sampai tiga digit hanya memberi
tahu bahwa keduanya memakai rumus yang sama. Uji yang lebih baik adalah mengulang
hitungan pada presisi berbeda: bila hasilnya bergeser jauh melebihi ramalan
$\kappa u$, ada yang salah dalam rumusnya, bukan dalam mesinnya.


\pagebreak[4]
\begin{tantangan}
Hitung $10^{8}\left(\sqrt{10^{16}+1}-10^{8}\right)$ dalam binary64. Nilai sebenarnya
$0{,}5$, tetapi program yang menuliskan rumus itu apa adanya melaporkan $0{,}0$.

Jelaskan mengapa hasilnya nol, lalu hitung bentuk yang memberi $0{,}5$ tanpa menaikkan
presisi sama sekali.

Petunjuk 1: tanyakan berapa $\operatorname{ulp}(10^{16})$, lalu apakah $10^{16}+1$
masih punya tempat di grid bilangan mesin.

Petunjuk 2: kalikan pembilang dan penyebut dengan sekawan
$\sqrt{10^{16}+1}+10^{8}$; pengurangan berbahaya itu menjadi pembagian.

Kalau mentok, kembali ke bagian Pembatalan katastrofik: dua wajah.
\end{tantangan}

\section{Contoh hitung}
\tingkat{menengah}

Semua angka di bagian ini bisa diperiksa dengan kalkulator, \texttt{struct}, dan
pecahan biasa. Kita mulai dari membongkar satu pola bit, yaitu $0{,}1$ dalam binary32:

\[
  \underbrace{0}_{1}\;\underbrace{01111011}_{8}\;
  \underbrace{10011001100110011001101}_{23}.
\]

Eksponen tersimpan $123$, jadi setelah dikurangi bias $127$ diperoleh $e=-4$. Bilangan
bulat 23 bit pada mantissa adalah $5033165$, sehingga signifikannya
$1+5033165/2^{23}$. Maka
\[
  x=\left(1+\frac{5033165}{2^{23}}\right)2^{-4}=\frac{13421773}{134217728}.
\]
Pembilangnya berasal dari $2^{23}+5033165=13421773$. Sementara itu penyebutnya
$134217728$, dan sepersepuluhnya $13421772{,}8$, sehingga galatnya tepat
$0{,}2/134217728$, yaitu $1{,}4901161\times10^{-9}$,
dan galat relatifnya $1{,}49\times10^{-8}$, yaitu $u_{32}/4$. Angka itu masuk akal:
bilangan $1/10$ sama dengan $1{,}1001100110011\dots_{2}\times2^{-4}$, yang memerlukan
mantissa periodik dan tidak berhenti, sehingga pembulatan benar memotongnya pada bit
ke-24. Pola yang tersimpan
berakhir \texttt{\dots 1101}, bukan \texttt{\dots 1100}, karena sisa di belakangnya
lebih besar daripada setengah ulp dan pembulatan naik.

\kertas{Bongkar pola bit \texttt{0 01111100 01000000000000000000000} dengan tangan:
tentukan tanda, eksponen setelah dikurangi bias, nilai mantissa, lalu tulis bilangannya
sebagai pecahan dan sebagai desimal. Hitung $\operatorname{ulp}$ bilangan itu, lalu
putuskan hasil dari menambahkan $2^{-27}$ dan $3\cdot2^{-27}$ dengan aturan
\emph{tie to even}. Periksa jawabanmu dengan \texttt{struct.unpack} dan
\texttt{numpy.float32}. Bilangan ini dipilih agar mantissanya berhenti setelah dua
digit, sehingga seluruh aritmetikanya bisa dikerjakan tanpa kalkulator.}

Sekarang kasus $0{,}1+0{,}2$ dalam binary64, yang seluruhnya bisa dihitung dengan
pecahan. Ketiga bilangan yang terlibat berpenyebut pangkat dua:
\[
  \operatorname{fl}(0{,}1)=\frac{3602879701896397}{2^{55}},\qquad
  \operatorname{fl}(0{,}2)=\frac{3602879701896397}{2^{54}},
\]
\[
  \operatorname{fl}(0{,}3)=\frac{5404319552844595}{2^{54}}.
\]
Jumlah dua yang pertama sebelum dibulatkan adalah $10808639105689191/2^{55}$,
sedangkan $\operatorname{fl}(0{,}3)=10808639105689190/2^{55}$. Jadi jumlah eksak itu
terletak tepat setengah ulp di atas $\operatorname{fl}(0{,}3)$: titik seri, bukan di
dekatnya. Digit terakhir mantissanya $1$, yaitu ganjil, sehingga \emph{tie to even}
membulatkan ke atas: $0{,}300000000000000044408920985006\dots$

Selisihnya terhadap $\operatorname{fl}(0{,}3)$ karena itu tepat satu ulp,
$2^{-54}=5{,}551115123125783\times10^{-17}$, bukan galat acak. Pengurangan sesudahnya
eksak menurut Sterbenz, jadi angka yang tercetak program adalah angka itu persis.

Angka yang sama menjelaskan mengapa kesimpulannya bergantung pada presisi: dalam
binary32, $0{,}1+0{,}2$ dan $0{,}3$ menghasilkan pola bit yang identik karena kedua
galat pembulatannya kebetulan saling meniadakan. Jadi \texttt{0.1+0.2 != 0.3} bukan
hukum alam, melainkan akibat kebetulan pada satu presisi.

Sekarang akar kuadrat tadi, dengan aritmetika ulp. Diskriminan $b^{2}-4ac$ tepat
sama dengan $10^{16}-4$, karena $10^{16}$ masih di bawah $2^{54}$ dan bilangan itu
kelipatan $4$. Akar kuadratnya benar sampai $0{,}46$ ulp, jadi ia jatuh satu ulp di
bawah $b$ karena $\operatorname{ulp}(b)=2^{-26}$. Selisih sebenarnya $2\times10^{-8}=1{,}34$ ulp,
selisih yang terhitung $1$ ulp; kesalahan $0{,}34$ ulp pada skala $10^{8}$ itu, dibagi
dua, menjadi seluruh galat akar kecil:
\[
  r_{2}^{\text{naif}}=-\frac{2^{-26}}{2}=-2^{-27}
  =-7{,}450580596923828\times10^{-9},
  \qquad r_{2}=-10^{-8}.
\]
Bentuk stabil $2c/(-b-\sqrt{b^{2}-4ac})$ tidak mengalami pengurangan berbahaya, dan
hasilnya $-10^{-8}$ tepat, akar yang dibulatkan benar.

\section{Lab komputasi}
\tingkat{lanjut}

\begin{lab}
\begin{verbatim}
import math, struct
import numpy as np

def bit32(x):
    return ''.join(f'{b:08b}' for b in struct.pack('>f', float(x)))

print("0.1     :", bit32(0.1))
print("6.0     :", bit32(6.0))
print("-12.375 :", bit32(-12.375))
print("hex     :", (0.1).hex(), (6.0).hex(), (1e-8).hex())
print("eps     :", np.finfo(np.float32).eps, np.finfo(np.float64).eps)
print("ulp     :", math.ulp(1.0), math.ulp(1e10))
print("seri    :", 1 + 2.0**-53, 1 + 3*2.0**-53)

s = (0.1 + 0.2) - 0.3
print("0.1+0.2-0.3 :", repr(s), "=", s / math.ulp(0.3), "ulp")
print("asosiatif   :", repr((0.1 + 0.2) + 0.3), repr(0.1 + (0.2 + 0.3)))

n = 10**6
exact = math.fsum(1.0 / i for i in range(1, n + 1))

def jumlah(prec, kahan):
    s, c = prec(0), prec(0)
    for i in range(1, n + 1):
        y = prec(1) / prec(i)
        if kahan:
            y = y - c
            t = s + y
            c = (t - s) - y
            s = t
        else:
            s = s + y
    return s

print("H_n eksak   :", f"{exact:.12f}")
for prec, nama in ((np.float32, "float32"), (np.float64, "float64")):
    for kahan in (False, True):
        v = float(jumlah(prec, kahan))
        print(f"{nama} {'Kahan' if kahan else 'naif ':7} {v:.12f}"
              f"  galat {abs(v - exact):.2e}")

print(f"{'x':>8} {'lugu':>19} {'galat':>9} {'identitas':>19} {'galat':>9}")
for k in range(1, 11):
    x = 10.0**-k
    lugu = (1 - math.cos(x)) / x**2
    stabil = 2 * math.sin(x / 2)**2 / x**2
    acuan = 0.5 - x**2/24 + x**4/720 - x**6/40320
    print(f"{x:8.0e} {lugu:19.15f} {abs(lugu-acuan)/acuan:9.1e}"
          f" {stabil:19.15f} {abs(stabil-acuan)/acuan:9.1e}")

a, b, c = 1.0, 1e8, 1.0
d = b*b - 4*a*c
sq = math.sqrt(d)
print("diskriminan :", d, " akar:", sq)
print("b - akar    :", b - sq, " ulp(b):", math.ulp(b))
print("akar naif   :", (-b + sq) / (2*a))
print("akar stabil :", 2*c / (-b - sq))
\end{verbatim}
Keluaran sebenarnya:
\begin{verbatim}
0.1     : 00111101110011001100110011001101
6.0     : 01000000110000000000000000000000
-12.375 : 11000001010001100000000000000000
hex     : 0x1.999999999999ap-4 0x1.8000000000000p+2 0x1.5798ee2308c3ap-27
eps     : 1.1920929e-07 2.220446049250313e-16
ulp     : 2.220446049250313e-16 1.9073486328125e-06
seri    : 1.0 1.0000000000000004
0.1+0.2-0.3 : 5.551115123125783e-17 = 1.0 ulp
asosiatif   : 0.6000000000000001 0.6
H_n eksak   : 14.392726722866
float32 naif    14.357357978821  galat 3.54e-02
float32 Kahan   14.392726898193  galat 1.75e-07
float64 naif    14.392726722865  galat 7.35e-13
float64 Kahan   14.392726722866  galat 0.00e+00
       x                lugu     galat           identitas     galat
   1e-01   0.499583472197418   5.6e-15   0.499583472197423   5.4e-15
   1e-02   0.499995833347366   2.9e-13   0.499995833347222   1.1e-16
   1e-03   0.499999958325503   1.6e-11   0.499999958333335   0.0e+00
   1e-04   0.499999996961265   5.2e-09   0.499999999583333   1.1e-16
   1e-05   0.500000041370185   8.3e-08   0.499999999995833   1.1e-16
   1e-06   0.500044450291171   8.9e-05   0.499999999999958   0.0e+00
   1e-07   0.499600361081320   8.0e-04   0.500000000000000   1.1e-16
   1e-08   0.000000000000000   1.0e+00   0.500000000000000   0.0e+00
   1e-09   0.000000000000000   1.0e+00   0.500000000000000   0.0e+00
   1e-10   0.000000000000000   1.0e+00   0.500000000000000   0.0e+00
diskriminan : 9999999999999996.0  akar: 99999999.99999999
b - akar    : 1.4901161193847656e-08  ulp(b): 1.4901161193847656e-08
akar naif   : -7.450580596923828e-09
akar stabil : -1e-08
\end{verbatim}
\end{lab}

Tiga baris pertama memastikan tabel di bagian awal bab bukan karangan: pola bit
$0{,}1$, $6$, dan $-12{,}375$ keluar persis seperti hasil bongkar tangan, termasuk bit
tanda. Baris \texttt{hex} memperlihatkan tiga belas digit heksadesimal. Baris
\texttt{seri} adalah demonstrasi aturan genap yang paling telanjang: menambah
setengah ulp kepada $1$ tidak mengubah apa pun, menambah satu setengah ulp memberi
$1+2^{-51}$. Baris \texttt{0.1+0.2-0.3} mencetak
$5{,}551115123125783\times10^{-17}$ dan menyatakannya sendiri sebagai $1{,}0$ ulp;
baris \texttt{asosiatif} memberi dua hasil berbeda, berselisih
$1{,}11\times10^{-16}$.

Empat baris berikutnya adalah inti bab ini dalam bentuk paling ringkas: pada
$n=10^{6}$, binary32 berurutan kehilangan $3{,}54\times10^{-2}$ sedangkan Kahan hanya
$1{,}75\times10^{-7}$; binary64 berurutan $7{,}35\times10^{-13}$ sedangkan Kahan nol,
bit yang sama dengan \texttt{math.fsum}. Galat $7{,}35\times10^{-13}$ itu setara
$414$ ulp, jadi presisi ganda hanya memindahkan masalah ke skala lebih kecil.

Tabel $(1-\cos x)/x^{2}$ memperlihatkan pola yang sudah dihitung di kertas. Pada
$x=10^{-1}$ kedua cara berbeda dari acuan sekitar $5\times10^{-15}$ karena acuannya
sendiri terpotong pada suku $x^{6}$; mulai $x=10^{-2}$ acuan itu jauh lebih akurat
daripada mesin, dan galat rumus lugu naik dari $2{,}9\times10^{-13}$ ke $1{,}0$
ketika $x$ turun ke $10^{-8}$, sedangkan kolom identitas bertahan pada $10^{-16}$.
Nol pada $x=10^{-8}$ bukan pembulatan malas: $\cos(10^{-8})$ memang sama dengan $1$.

Blok terakhir memeriksa akar kuadrat: diskriminan $10^{16}-4$, akar $10^{8}-2^{-26}$,
selisih satu ulp $b$ yaitu $2^{-26}$, akar naif $-2^{-27}$, dan akar stabil $10^{-8}$
yang dibulatkan benar. Dua baris
terakhir berjarak $25{,}5$ persen, padahal rumusnya setara secara aljabar dan datanya
sama persis.

\begin{untukdosen}
Urutan mengajar yang berhasil dipakai: mulai dari satu kegagalan yang bisa dihitung
ulang, yaitu rumus kuadrat di awal bab, baru setelah itu masuk ke representasi. Bila
representasi datang lebih dulu, mahasiswa menghafal $1/8/23$ dan $1/11/52$ tanpa tahu
untuk apa. Setelah kegagalan itu, urutan yang bekerja adalah: banyaknya bilangan mesin
(agar jelas bahwa pembulatan tak terhindarkan), aturan pembulatan, lemma ulp, batas
$(1+\delta)$, lalu bilangan kondisi. Penjumlahan sebaiknya masuk setelah bilangan
kondisi, karena batas $(n-1)u\sum|x_{i}|$ baru terbaca sebagai ``galat yang tumbuh''
bila $\kappa$ sudah dikenal.

Dua miskonsepsi khas. Pertama, mahasiswa menganggap galat pembulatan bersifat acak dan
karena itu saling meniadakan pada penjumlahan panjang. Tabel harmonik membantahnya
dengan keras: galatnya bergerak searah, dan pada $n=10^{6}$ ia mencapai tiga puluh
lima ribu ulp. Kedua, mereka menganggap pengurangan selalu merusak digit, lalu
menghindari semua pengurangan. Teorema Sterbenz menunjukkan sebaliknya: pengurangan
dua bilangan mesin yang berdekatan justru eksak, dan yang merusak adalah galat yang
sudah ada pada masukannya. Aturan yang perlu diulang: yang berbahaya bukan tanda
minus, melainkan selisih kecil dari dua bilangan yang masing-masing sudah dibulatkan.

Aktivitas papan tulis yang murah dan efektif: minta seluruh kelas menghitung
$1+2^{-53}$ dan $1+3\cdot2^{-53}$ dengan aturan genap, lalu bandingkan dengan Python.
Setelah itu, minta dua sukarelawan mengerjakan
$(0{,}1+0{,}2)+0{,}3$ dan $0{,}1+(0{,}2+0{,}3)$ dengan menuliskan mantissa 52 bit
sebagai pecahan. Kegaduhan biasanya muncul saat mereka menemukan bahwa keduanya benar.

Tanda kesiapan lanjut: mahasiswa dapat menebak, tanpa menjalankan program, bahwa
$\operatorname{fl}(1+10^{-16})=1$, dapat menjelaskan mengapa bentuk
$2c/(-b-\sqrt{\Delta})$ lebih baik daripada $(-b+\sqrt{\Delta})/(2a)$, dan dapat
memperkirakan berapa digit yang hilang dari sebuah $\kappa$ yang diberikan. Bila
mereka menyebut kata ``galat ke belakang'' tanpa dipancing, mereka siap masuk Bab 25,
tempat gradien jaringan dihitung dengan aturan serupa.
\end{untukdosen}

\section{Ringkasan}
\tingkat{dasar}

\begin{ringkasan}
\begin{itemize}
  \item $\mathbb{F}$ berhingga: $2^{p-1}$ bilangan normal per pangkat dua, jadi
        hampir setiap bilangan real tidak punya wakil dan setiap penyimpanan berarti
        pembulatan.
  \item Pembulatan IEEE adalah \emph{round to nearest, ties to even}. Akibatnya
        $\operatorname{fl}(x)=x(1+\delta)$ dengan $|\delta|\le u=2^{-p}$, yaitu
        $1{,}11\times10^{-16}$ untuk binary64 dan $5{,}96\times10^{-8}$ untuk
        binary32.
  \item $\operatorname{eps}=2^{-p+1}$ adalah jarak $1$ ke bilangan sesudahnya;
        $\operatorname{ulp}(x)=2^{e-p+1}$ untuk $2^{e}\le|x|<2^{e+1}$. Galat satu
        pembulatan adalah setengah ulp, dan bila dinyatakan relatif ia berayun dengan
        faktor dua.
  \item Bilangan kondisi $\kappa_{f}=|xf'(x)|/|f(x)|$ mengukur pembesaran galat
        masukan. Untuk pengurangan, $\kappa=(|a|+|b|)/|a-b|$, yang membesar tanpa
        batas ketika $|a-b|$ mengecil.
  \item Pengurangan dua bilangan mesin yang berdekatan eksak (Sterbenz). Pembatalan
        katastrofik terjadi karena masukannya sudah membawa galat, bukan karena
        pengurangannya.
  \item Rumus kuadrat lugu menghitung akar kecil lewat pengurangan dua bilangan
        berukuran $|b|$; bentuk $2c/(-b-\sqrt{\Delta})$ menghindarinya dan memberi akar
        yang dibulatkan benar.
  \item Penjumlahan berurutan punya galat terburuk $(n-1)u\sum|x_{i}|$ dan tidak
        asosiatif. Kahan memindahkan galat ke variabel kompensasi sehingga faktor $n$
        hilang dari suku utama.
\end{itemize}
\end{ringkasan}

\section{Latihan}
\tingkat{menengah}

\begin{latihan}
\begin{enumerate}
  \item Tentukan $\operatorname{ulp}$ dari $1$; $1024$; $0{,}1$; dan $3\times10^{8}$
        dalam binary64, lalu bandingkan dengan $|x|\,u$. Berapa banyak bilangan
        binary64 yang terletak di $[1,2)$?

  \item Tunjukkan bahwa $1+2^{-53}=1$ dan $1+2^{-52}>1$ dalam binary64, lalu tentukan
        bilangan positif terkecil $x$ dengan $1+x>1$ dalam binary32. Nyatakan
        jawabanmu sebagai $2^{-k}$ dan bandingkan dengan tabel di bagian awal.

  \item Buktikan bahwa untuk setiap $x$ normal berlaku
        $|\operatorname{fl}(x)-x|\le\tfrac12\operatorname{ulp}(x)$, lalu jelaskan
        mengapa ketaksamaan itu tidak dapat diganti dengan kesamaan untuk semua $x$.

  \item Nilai klaim ini: \emph{karena setiap operasi IEEE-754 dibulatkan dengan benar,
        galat relatif hasil akhir program yang hanya memakai $+,-,\times,\div$ selalu
        paling besar beberapa $u$.} Putuskan benar, salah, atau belum cukup bukti, dan
        sebutkan satu contoh dari bab ini yang mendukung keputusanmu.

  \item Hitung dengan tangan $(0{,}1+0{,}2)+0{,}3$ dan $0{,}1+(0{,}2+0{,}3)$ sebagai
        pecahan berpenyebut pangkat dua, lalu tunjukkan bahwa selisihnya satu ulp.
        Periksa dengan \texttt{math.fsum} dan dengan \texttt{float.hex()}.
        (Petunjuk: nyatakan setiap bilangan sebagai $m/2^{55}$ atau $m/2^{54}$ dan
        hitung pembilangnya sebagai bilangan bulat.)

  \item Ulangi perhitungan $(1-\cos x)/x^{2}$ untuk $x=10^{-7}$ dengan identitas sudut
        ganda, lalu bandingkan dengan $0{,}5-x^{2}/24+x^{4}/720$. Berapa galat relatif
        sisanya, dan dari mana asalnya?

  \item Untuk $a=1$, $b=10^{4}$, $c=1$, hitung kedua akar dengan rumus lugu dan dengan
        bentuk stabil dalam binary64. Tentukan bilangan kondisi pengurangannya, lalu
        ramalkan berapa digit yang hilang sebelum menjalankan program. Periksa
        ramalanmu dengan Python.
        (Petunjuk: $\sqrt{10^{8}-4}\approx10^{4}-2\times10^{-4}$; bandingkan dengan
        $\operatorname{ulp}(10^{4})=2^{-39}$.)

  \item Buktikan bahwa penjumlahan berpasangan atas $n=2^{k}$ suku mempunyai galat
        terburuk sebanding $ku\sum|x_{i}|$, bukan $nu\sum|x_{i}|$, dengan menguraikan
        pohon penjumlahannya dan menghitung kedalaman pohon.

  \item Misalkan $x_{i}=10^{-16}$ untuk $i=1,\dots,10^{6}$ dan $x_{0}=1$. Hitung
        penjumlahan berurutan dan penjumlahan Kahan keduanya dengan tangan (boleh
        dibantu kalkulator), lalu jelaskan mengapa jawaban Kahan berbeda padahal
        setiap langkahnya juga dibulatkan.
        (Petunjuk: pada penjumlahan berurutan, setiap suku di bawah setengah ulp
        akumulator; pada Kahan, kompensasi $c$ menyimpan sisa itu dan menambahkannya
        kembali pada langkah berikutnya.)

  \item Sebuah program menghitung variansi dengan
        $\frac1n\sum x_{i}^{2}-\left(\frac1n\sum x_{i}\right)^{2}$ pada data yang
        seluruhnya sekitar $999$. Tentukan bilangan kondisi bentuk itu, ramalkan
        berapa digit yang hilang dalam binary32, dan bandingkan dengan bentuk
        dua-lintasan $\frac1n\sum(x_{i}-\bar{x})^{2}$. Jelaskan mengapa bentuk kedua
        lebih baik meskipun memerlukan dua kali lintasan data.
\end{enumerate}
\end{latihan}

\begin{untukpeneliti}
Yang benar-benar terbukti di bab ini adalah pernyataan-pernyataan baku dan tidak
bergantung pada mesin tertentu: batas $\operatorname{fl}(x)=x(1+\delta)$ dengan
$|\delta|\le u$, lemma ulp, teorema Sterbenz, batas galat penjumlahan berurutan, dan
ketaksamaan bilangan kondisi lewat teorema nilai rata-rata. Rujukan primernya Goldberg,
\emph{What Every Computer Scientist Should Know About Floating-Point Arithmetic}
(ACM Computing Surveys 23(1), 1991), bagian ``Relative Error and Ulps'',
``Cancellation'', ``Exactly Rounded Operations'', dan ``Errors in Summation''. Batas
Kahan adalah Teorema 8 kitab itu; buktinya di lampiran, diambil dari Knuth, dan tidak
diulang di sini. Untuk perlakuan sistematis mengenai kestabilan, rujukan lanjutannya
Higham, \emph{Accuracy and Stability of Numerical Algorithms}, bab tentang analisis
galat dan penjumlahan.

Yang lebih sering dipopulerkan daripada dibuktikan ada tiga. Pertama, klaim bahwa
``galat pembulatan bersifat acak dan saling meniadakan''. Itu benar hanya sebagai
model statistik untuk pembulatan yang tak berkorelasi, dan justru gagal pada kasus
yang paling penting: akumulasi pada penjumlahan panjang bersifat sistematis, seperti
terlihat pada tabel harmonik. Kedua, klaim bahwa galat selalu bisa dikurangi dengan
presisi lebih tinggi. Untuk pembatalan katastrofik, menaikkan presisi hanya
memindahkan titik runtuh, bukan menghapusnya: dalam binary64 titik itu ada di
$x\approx1{,}5\times10^{-8}$, dan dalam binary128 ia ada di sekitar $10^{-17}$. Ketiga,
klaim bahwa ulp adalah satuan galat yang bebas ambiguitas. Goldberg menunjukkan bahwa
ulp dapat didefinisikan sebagai jarak antar bilangan atau setengah jarak itu, dan
selisih kedua definisi itu muncul sebagai faktor dua dalam setiap batas. Angka
``$0{,}46$ ulp'' pada pembahasan akar kuadrat di bab ini memakai definisi pertama.

Batas keberlakuan hasil di sini perlu dinyatakan. Seluruh analisis menganggap model
pembulatan benar dan tidak ada \emph{fused multiply-add} yang mengubah urutan operasi,
tidak ada pengaturan ulang oleh pengompilasi, dan tidak ada mode pembulatan selain
\emph{to nearest}. Di perangkat nyata, vektorisasi dan instruksi FMA dapat mengubah
hasil bit demi bit; itu bukan pelanggaran standar, tetapi membuat perbandingan antar
mesin menjadi tidak deterministik. Analisis galat juga tidak menyentuh galat model,
galat data, atau galat pemotongan: bab ini hanya membahas galat pembulatan.

Pertanyaan terbuka yang jujur. Pertama, analisis galat klasik memberi batas terburuk,
sedangkan yang dibutuhkan praktisi adalah perkiraan yang ketat untuk kasus khas;
probabilistik pembulatan masih menjadi alat yang aktif diteliti dan belum menjadi
teori yang selesai. Kedua, untuk algoritme iteratif dan penyelesaian sistem linear,
galat pembulatan berinteraksi dengan bilangan kondisi matriks sedemikian rupa sehingga
batas naif terlalu pesimistis; analisis ke belakang, yang menjadi gagasan utama
Higham dan Wilkinson, memberi kerangka yang lebih tajam tetapi tidak dibahas di sini.
Ketiga, keputusan praktis seperti ``apakah binary32 cukup'' bergantung pada distribusi
data, dan tidak ada jawaban yang berlaku umum.
\end{untukpeneliti}

\peta{Bab ini mengikuti Goldberg, \emph{What Every Computer Scientist Should Know About
Floating-Point Arithmetic} (ACM Computing Surveys, 1991), bagian ``Introduction'',
``Relative Error and Ulps'', ``Cancellation'', ``Exactly Rounded Operations'',
``Formats and Operations'', ``Special Quantities'', dan ``Errors in Summation'', dengan
Teorema 8 untuk batas penjumlahan Kahan. Untuk sisi analisis real dari gagasan limit
dan teorema nilai rata-rata yang dipakai pada bukti bilangan kondisi, lihat Abbott,
\emph{Understanding Analysis}, bagian 2.2 (Limit of a Sequence), 5.3 (The Mean Value
Theorem), dan 6.6 (Taylor Series). Contoh $0{,}1+0{,}2$, pembongkaran pola bit,
analisis $(1-\cos x)/x^{2}$, dan seluruh angka pada tabel serta lab adalah perhitungan
sendiri yang dijalankan dengan Python 3.14 dan NumPy 2.5 dan diverifikasi ulang dengan
pecahan biasa serta aritmetika 50 digit. Untuk perlakuan lanjutan mengenai kestabilan
dan galat ke belakang, rujukan yang biasa dipakai adalah Higham, \emph{Accuracy and
Stability of Numerical Algorithms}.}


\part{Optimasi}

Bagian ini mengubah ``cari yang terbaik'' menjadi soal yang punya syarat
selesai. Prasyaratnya Bab~6 dan~12. Setelah dua bab, kamu bisa membedakan minimum
lokal dari global, memakai dualitas untuk membuktikan batas, dan mengenali kapan
sebuah soal linear selesai dalam waktu polinomial dan kapan tidak.
\chapter{Optimasi Konveks dan Dualitas}

Sebuah pencampur warna bekerja dengan dua pigmen. Tiap pigmen punya takaran yang paling
disukai mesin: dua bagian untuk pigmen pertama, dua bagian untuk pigmen kedua. Tetapi
resep menuntut jumlah keduanya tepat dua bagian, dan kelarutan pigmen pertama tidak
mengizinkan lebih dari setengah bagian. Kita ingin campuran yang sedekat mungkin dengan
kesukaan mesin, dengan dua tuntutan yang tidak bisa dilanggar. Tanpa tuntutan, jawabannya
dua dan dua; dengan tuntutan jumlah saja, satu dan satu. Tuntutan kelarutan menggeser
jawaban lagi, dan tidak jelas lagi ke arah mana.

Bentuk soal itu muncul di tempat yang tidak berbau pigmen. Dua bobot model punya nilai
terbaik dari pencocokan sebelumnya, sebuah syarat kalibrasi menuntut jumlah keduanya
tetap, dan satu pembatas perangkat membatasi salah satunya. Berapa bobot terbaik yang
masih memenuhi keduanya, dan berapa harga tiap tuntutan dalam satuan fungsi sasaran?
Pertanyaan kedua itu yang dipakai insinyur untuk memutuskan tuntutan mana yang layak
dilonggarkan.

Bab 10 memberi kuadrat terkecil, tetapi hanya tanpa kendala. Bab 12 memberi gradien, tetapi
gradien hanya berbicara tentang tetangga: titik dengan gradien nol bisa minimum, maksimum,
atau pelana. Bab 9 memberi nilai eigen, yang dipakai di sini, tetapi nilai eigen tidak
menjawab pertanyaan terpenting: kalau algoritme berhenti di suatu titik, apakah itu
jawaban terbaik di seluruh daerah layak? Untuk fungsi konveks ya. Untuk soal umum tidak,
dan bagian keempat menunjukkan bagaimana dualitas memberi batas bawah yang sah tanpa
memeriksa seluruh daerah layak.

\begin{untukanak}
Pernahkah kamu menaruh kelereng di dalam mangkuk? Kelereng itu selalu berhenti di dasar, di
tempat yang sama, ke mana pun ia dilepas, karena dasarnya cekung ke atas seperti huruf U.
Sekarang bayangkan kelereng yang sama dilepas di lereng bukit berbatu. Ia berhenti di
lekukan pertama yang cukup dalam, dan lekukan itu belum tentu yang paling rendah. Dari
tempat berhenti, kamu tidak bisa tahu apakah ada lembah lain yang lebih dalam tanpa
berjalan jauh dan memeriksa.

Di bab ini kita akan tahu mengapa mangkuk begitu bersahabat dan bukit begitu menipu, dan
mengapa kita bisa membedakan keduanya sebelum menghitung apa pun. Ada gagasan kedua yang
lebih halus: kadang kita tidak perlu tahu letak dasar lembah, cukup tahu batas terendah
yang mungkin dicapai. Kalau batas itu sama dengan tempat kita berdiri, kita yakin tidak
ada yang lebih baik, tanpa memeriksa satu titik pun. Batas semacam itu dijual oleh daftar
harga untuk setiap tuntutan, dan daftar harga itulah yang membuat banyak soal optimasi
bisa dipercaya tanpa diselesaikan sampai habis.
\end{untukanak}

\section{Himpunan konveks dan kombinasi konveks}
\tingkat{dasar}

Kata ``konveks'' dalam bahasa sehari-hari berarti cembung ke luar, seperti punggung
kura-kura. Dalam optimasi artinya hampir kebalikannya: himpunan disebut konveks bila ia
tidak punya lekuk ke dalam, sehingga ruas garis antara dua anggotanya tetap di dalam.
Mangkuk tadi konveks; daerah di antara dua tanduk bulan tidak.

Aturan mainnya satu: campur dua titik dengan bobot tak negatif yang jumlahnya satu, dan
hasil campuran itu harus tetap anggota himpunan. Bobot satu berarti memilih titik pertama, bobot
nol titik kedua, bobot setengah tepat di tengah.

\kamus{kombinasi konveks}{convex combination}{Resep mencampur beberapa titik dengan
takaran yang jumlahnya tepat satu, seperti mencampur cat: hasilnya selalu di antara
bahan-bahannya.}

\begin{definisi}[Himpunan konveks dan kombinasi konveks]
Himpunan $C\subseteq\mathbb{R}^{n}$ disebut \emph{konveks} bila untuk setiap
$\mathbf{x},\mathbf{y}\in C$ dan setiap $t\in[0,1]$ berlaku
$t\mathbf{x}+(1-t)\mathbf{y}\in C$. Titik itu disebut \emph{kombinasi konveks} dari
$\mathbf{x}$ dan $\mathbf{y}$. Secara umum, $\sum_{i=1}^{k}\theta_{i}\mathbf{x}_{i}$
disebut kombinasi konveks bila $\theta_{i}\ge0$ dan $\sum_{i}\theta_{i}=1$.
\emph{Selubung konveks} $\operatorname{conv}S$ adalah himpunan semua kombinasi konveks
berhingga titik di $S$.\index{himpunan konveks}\index{kombinasi konveks}
\end{definisi}

Definisi dengan dua titik cukup untuk semua titik: bila $C$ memuat setiap campuran dua
anggota, ia memuat setiap campuran berhingga anggotanya, dengan induksi pada banyak
titik.

Empat contoh penting. \emph{Hiperbidang}
$\{\mathbf{x}:\mathbf{a}^{\top}\mathbf{x}=b\}$ konveks, sebab
$\mathbf{a}^{\top}(t\mathbf{x}+(1-t)\mathbf{y})=tb+(1-t)b=b$. \emph{Setengah ruang}
$\{\mathbf{x}:\mathbf{a}^{\top}\mathbf{x}\le b\}$ konveks dengan alasan yang sama memakai
ketaksamaan. \emph{Bola} $\{\mathbf{x}:\|\mathbf{x}-\mathbf{c}\|\le r\}$ konveks karena
norma memenuhi ketaksamaan segitiga (Bab 10):
\[
  \|t\mathbf{x}+(1-t)\mathbf{y}-\mathbf{c}\|
  \le t\|\mathbf{x}-\mathbf{c}\|+(1-t)\|\mathbf{y}-\mathbf{c}\|\le tr+(1-t)r=r .
\]
\emph{Polihedron}, yaitu irisan berhingga hiperbidang dan setengah ruang, konveks karena
alasan berikut.

\begin{proposisi}[Irisan mempertahankan konveksitas]
Bila $C_{1},\dots,C_{k}$ konveks, maka $\bigcap_{i=1}^{k}C_{i}$ konveks.
\end{proposisi}

\begin{proof}
Ambil $\mathbf{x},\mathbf{y}$ di irisan dan $t\in[0,1]$. Untuk setiap $i$, kedua titik
itu anggota $C_{i}$, jadi $t\mathbf{x}+(1-t)\mathbf{y}\in C_{i}$ karena $C_{i}$ konveks.
Karena itu titik campuran berada di setiap $C_{i}$, yaitu di irisan.
\end{proof}

Irisan murah hati; gabungan tidak. Ambil $C_{1}=\{(x_{1},0):x_{1}\in[0,1]\}$ dan
$C_{2}=\{(x_{1},1):x_{1}\in[0,1]\}$, keduanya konveks. Titik $(0,0)$ dan $(0,1)$ ada di
gabungan, tetapi campuran setengahnya $(0,\tfrac12)$ tidak. Hal yang sama terjadi pada
$\{(x_{1},x_{2}):x_{1}x_{2}\ge1\}$: titik $(-1,-1)$ dan $(1,1)$ keduanya anggota,
sedangkan titik asal bukan, sebab $0<1$. Himpunan itu dua daerah terpisah.

Selubung konveks adalah tambalan termurah untuk lubang semacam itu: himpunan konveks
terkecil yang memuat $S$. Ia selalu ada, karena irisan semua himpunan konveks yang memuat
$S$ konveks menurut proposisi tadi dan tidak kosong, sebab $\mathbb{R}^{n}$ sendiri
konveks.

\section{Fungsi konveks dan dua ujinya}
\tingkat{menengah}

Himpunan konveks berbicara tentang daerah; fungsi konveks berbicara tentang bentuk
permukaan di atas daerah itu. Yang dituntut: nilai di titik campuran tidak lebih besar
daripada campuran nilai. Pada Gambar~\ref{fig:fungsi-konveks} kiri, tali dari $(-1,1)$ ke
$(1,1)$ di atas grafik $f(t)=t^{2}$; pada panel kanan talinya terpotong.

\kamus{fungsi konveks}{convex function}{Permukaan yang tidak pernah melengkung ke atas
seperti atap; tali antar dua titiknya selalu di atas permukaan.}

\begin{definisi}[Fungsi konveks]
Misalkan $C\subseteq\mathbb{R}^{n}$ konveks. Fungsi $f\colon C\to\mathbb{R}$ disebut
\emph{konveks} bila
\[
  f\bigl(t\mathbf{x}+(1-t)\mathbf{y}\bigr)
  \le t\,f(\mathbf{x})+(1-t)\,f(\mathbf{y})
\]
untuk semua $\mathbf{x},\mathbf{y}\in C$ dan $t\in[0,1]$. Ia \emph{konveks ketat} bila
ketaksamaan itu ketat untuk $\mathbf{x}\ne\mathbf{y}$ dan $t\in(0,1)$, dan \emph{cekung}
bila $-f$ konveks.\index{fungsi konveks}
\end{definisi}

\begin{figure}[!ht]
\centering
\begin{tikzpicture}
\begin{axis}[width=6cm,height=4.8cm,xlabel={$t$},ylabel={$f(t)$}]
\addplot[domain=-1.3:1.3,samples=50,thick] {x^2};
\addplot[dashed,thick,red] coordinates {(-1,1) (1,1)};
\addplot[only marks,mark=*] coordinates {(-1,1) (1,1)};
\end{axis}
\end{tikzpicture}\quad
\begin{tikzpicture}
\begin{axis}[width=6cm,height=4.8cm,xlabel={$t$},ylabel={$g(t)$}]
\addplot[domain=-1.3:1.3,samples=80,thick] {x^3-x};
\addplot[dashed,thick,red] coordinates {(-1,0) (1,0)};
\addplot[only marks,mark=*] coordinates {(-1,0) (1,0)};
\end{axis}
\end{tikzpicture}
\caption{Tali antara dua titik grafik. Kiri: $f(t)=t^{2}$, tali dari $(-1,1)$ ke $(1,1)$
bernilai $1$ di $t=0$ sedangkan grafiknya $0$, jadi tali di atas grafik. Kanan:
$g(t)=t^{3}-t$, tali dari $(-1,0)$ ke $(1,0)$ memotong grafik di $t=-1$, $0$, $1$; di
$t=-0{,}5$ grafik bernilai $0{,}375$ sedangkan talinya $0$.}
\label{fig:fungsi-konveks}
\end{figure}

Definisi itu setara dengan pernyataan geometris tentang \emph{epigraf}, yaitu himpunan
titik di atas grafik: $f$ konveks tepat ketika
$\operatorname{epi}f=\{(\mathbf{x},t):f(\mathbf{x})\le t\}$ konveks. Kesetaraannya murni
aljabar, dan memindahkan pertanyaan tentang fungsi menjadi pertanyaan tentang himpunan.

Untuk fungsi yang dapat diturunkan ada dua uji yang jauh lebih cepat, dan keduanya
perlu-dan-cukup, bukan hanya cukup.

\begin{teorema}[Uji orde pertama]
Misalkan $f$ terdiferensialkan pada himpunan konveks terbuka $C$. Fungsi $f$ konveks jika
dan hanya jika
\[
  f(\mathbf{y})\ge f(\mathbf{x})+\nabla f(\mathbf{x})^{\top}(\mathbf{y}-\mathbf{x})
\]
untuk semua $\mathbf{x},\mathbf{y}\in C$.
\end{teorema}

\begin{proof}
Andaikan $f$ konveks, ambil $\mathbf{x}\ne\mathbf{y}$ dan $t\in(0,1)$. Dari definisi,
\[
  \frac{f\bigl(\mathbf{x}+t(\mathbf{y}-\mathbf{x})\bigr)-f(\mathbf{x})}{t}
  \le f(\mathbf{y})-f(\mathbf{x}).
\]
Saat $t\downarrow0$ ruas kiri menuju turunan berarah
$\nabla f(\mathbf{x})^{\top}(\mathbf{y}-\mathbf{x})$, sehingga ketaksamaan yang diminta
tercapai di limit. Untuk arah sebaliknya, andaikan ketaksamaan itu berlaku dan ambil
$\mathbf{z}=t\mathbf{x}+(1-t)\mathbf{y}$ dengan $t\in[0,1]$. Terapkan ketaksamaan di
$\mathbf{z}$ terhadap $\mathbf{x}$ dan terhadap $\mathbf{y}$:
\[
  f(\mathbf{x})\ge f(\mathbf{z})+\nabla f(\mathbf{z})^{\top}(\mathbf{x}-\mathbf{z}),
  \qquad
  f(\mathbf{y})\ge f(\mathbf{z})+\nabla f(\mathbf{z})^{\top}(\mathbf{y}-\mathbf{z}).
\]
Kalikan yang pertama dengan $t$, yang kedua dengan $1-t$, lalu jumlahkan. Suku gradien
lenyap karena $t(\mathbf{x}-\mathbf{z})+(1-t)(\mathbf{y}-\mathbf{z})=\mathbf{0}$, dan
tersisa $t\,f(\mathbf{x})+(1-t)f(\mathbf{y})\ge f(\mathbf{z})$.
\end{proof}

Bacaan praktisnya: grafiknya di atas setiap bidang singgung. Satu bidang singgung sudah
memberi batas bawah global, dan itulah sebabnya gradien di titik mana pun menjadi alat
yang jujur, bukan ramalan lokal.

\begin{teorema}[Uji orde kedua]
Misalkan $f$ terdiferensialkan dua kali pada himpunan konveks terbuka $C$. Fungsi $f$
konveks jika dan hanya jika Hessian $\nabla^{2}f(\mathbf{x})$ semidefinit positif untuk
setiap $\mathbf{x}\in C$, yaitu
$\mathbf{v}^{\top}\nabla^{2}f(\mathbf{x})\mathbf{v}\ge0$ untuk semua $\mathbf{v}$.
\index{Hessian}
\end{teorema}

\begin{proof}
Bila $f$ konveks, untuk setiap $\mathbf{x}$ dan $\mathbf{v}$ fungsi satu variabel
$g(s)=f(\mathbf{x}+s\mathbf{v})$ konveks pada selang di sekitar nol, sebab komposisi
dengan peta afin mempertahankan konveksitas. Uji orde pertama pada $g$ menunjukkan $g'$
naik, sehingga $g''(0)\ge0$; dan $g''(0)=\mathbf{v}^{\top}\nabla^{2}f(\mathbf{x})
\mathbf{v}$.

Sebaliknya, bila Hessian semidefinit positif di seluruh $C$, teorema Taylor dengan sisa
integral (Bab 12) memberi
\[
  f(\mathbf{y})=f(\mathbf{x})+\nabla f(\mathbf{x})^{\top}(\mathbf{y}-\mathbf{x})
  +\int_{0}^{1}(1-s)\,(\mathbf{y}-\mathbf{x})^{\top}
   \nabla^{2}f\bigl(\mathbf{x}+s(\mathbf{y}-\mathbf{x})\bigr)(\mathbf{y}-\mathbf{x})\,ds .
\]
Integran tidak negatif karena Hessian semidefinit positif dan $(1-s)\ge0$, jadi suku
integralnya tidak negatif. Ketaksamaan orde pertama terpenuhi, dan $f$ konveks.
\end{proof}

Dua peringatan tentang uji orde kedua. Pertama, Hessian \emph{definit} positif cukup
untuk konveks ketat, tetapi bukan syarat perlu: $f(x)=x^{4}$ konveks ketat padahal
$f''(0)=0$. Kedua, Hessian semidefinit positif hanya memberi konveksitas, bukan ketatnya;
$f(x_{1},x_{2})=x_{1}^{2}$ punya Hessian $\operatorname{diag}(2,0)$ dan minimumnya
seluruh garis $x_{1}=0$.

\begin{bacakeras}
Dua suara, dibaca bergantian dengan teman. Suara pertama bertanya, suara kedua menjawab.

``Aku menjatuhkan kelereng di mangkuk. Ke mana ia pergi?'' ``Ke dasar. Selalu ke dasar
yang sama.'' ``Kalau aku melepasnya dari sisi yang lain?'' ``Ke dasar itu juga.'' ``Jadi
hanya ada satu tempat terendah?'' ``Hanya satu, dan setiap jalan menurun menemukannya.''

``Sekarang aku melepasnya di lereng batu. Ke mana ia pergi?'' ``Ke lekukan pertama yang
cukup dalam.'' ``Apakah itu yang terdalam?'' ``Belum tentu. Lekukan itu menahannya, bukan
meyakinkannya.''

Ulangi bergantian, makin cepat, dengan ketukan meja: mangkuk tidak punya jebakan, bukit
punya jebakan.

Baris yang harus diulang pembaca, sekali, dengan suara jelas: kalau setiap jalan menurun
berakhir di titik yang sama, di situ tidak ada jebakan, dan itulah arti konveks.
\end{bacakeras}

Beberapa operasi mempertahankan konveksitas. Jumlah dua fungsi konveks tetap konveks,
karena ketaksamaannya dijumlahkan; kelipatan positif juga. Komposisi
$f(\mathbf{A}\mathbf{x}+\mathbf{b})$ konveks bila $f$ konveks, sebab peta afin memetakan
campuran ke campuran. Maksimum titik demi titik juga konveks: bila $f=\max_{i}f_{i}$,
ketaksamaan bagi setiap $f_{i}$ berlaku bagi maksimumnya. Itu menjelaskan mengapa
$\max\{|x_{1}|,|x_{2}|\}$ berpatahan tetapi tetap konveks.

\begin{resep}
Empat langkah memeriksa kecembungan, dari yang paling murah.
\begin{enumerate}
  \item Periksa daerahnya lebih dulu: ambil dua titik di $C$ dan pastikan ruas di
        antaranya tetap di $C$. Bila gagal, berhenti, karena uji berikutnya mengandaikan
        $C$ konveks.
  \item Uji lewat definisi dengan pasangan umum $\mathbf{x},\mathbf{y}$: tulis
        $f(t\mathbf{x}+(1-t)\mathbf{y})$ dan bandingkan dengan
        $t\,f(\mathbf{x})+(1-t)f(\mathbf{y})$ untuk $t\in[0,1]$. Satu pasangan yang
        membalik ketaksamaan sudah cukup untuk menolak.
  \item Hitung Hessian $\nabla^{2}f(\mathbf{x})$ sebagai fungsi $\mathbf{x}$, bukan satu
        matriks angka di satu titik.
  \item Putuskan tandanya: nilai eigen tak negatif, atau semua minor utama
        $1\times1$ sampai $n\times n$ tak negatif. Satu nilai eigen negatif sudah
        menutup perkara.
\end{enumerate}
Peringatan urutan: langkah 4 tidak boleh dikerjakan sebelum langkah 1. Hessian yang
definit positif di satu titik hanya berbicara tentang titik itu, sedangkan definisi
berbicara tentang seluruh $C$.
\end{resep}

\section{Bentuk baku dan mengapa optimum lokal menjadi global}
\tingkat{menengah}

Sampai sini kita punya dua gagasan terpisah: daerah yang tidak berlubang dan permukaan
yang tidak melengkung ke atas. Optimasi menyatukan keduanya lewat satu bentuk tulisan.
Setiap masalah optimasi dapat ditulis dengan cara yang sama: satu fungsi yang hendak
dikecilkan, sekumpulan syarat ketaksamaan, dan sekumpulan syarat kesamaan.

\kamus{bentuk baku}{standard form}{Cara menulis semua soal optimasi dengan satu pola:
satu sasaran, syarat ``tidak lebih dari nol'', dan syarat ``sama dengan''.}

\begin{definisi}[Masalah optimasi bentuk baku]
\emph{Masalah optimasi bentuk baku} adalah
\[
  \begin{array}{ll}
  \text{minimumkan} & f_{0}(\mathbf{x})\\
  \text{dengan} & f_{i}(\mathbf{x})\le0,\qquad i=1,\dots,m,\\
                & h_{j}(\mathbf{x})=0,\qquad j=1,\dots,p,
  \end{array}
\]
dengan variabel $\mathbf{x}\in\mathbb{R}^{n}$, fungsi sasaran $f_{0}$, fungsi kendala
ketaksamaan $f_{i}$, dan fungsi kendala kesamaan $h_{j}$. Daerah asalnya
$D=\bigcap_{i}\operatorname{dom}f_{i}\cap\bigcap_{j}\operatorname{dom}h_{j}$, daerah
layaknya $X=\{\mathbf{x}\in D:f_{i}(\mathbf{x})\le0,\ h_{j}(\mathbf{x})=0\}$, dan
\emph{nilai optimum} $p^{*}=\inf_{\mathbf{x}\in X}f_{0}(\mathbf{x})$, boleh
$-\infty$ bila masalahnya tak terbatas ke bawah. Masalahnya disebut \emph{masalah
konveks} bila $f_{0},\dots,f_{m}$ konveks dan $h_{1},\dots,h_{p}$ afin.
\index{bentuk baku}
\end{definisi}

Bahwa setiap kendala dapat dibawa ke pola ini mudah diperiksa. Syarat
$f(\mathbf{x})\ge b$ menjadi $b-f(\mathbf{x})\le0$; kesamaan linear tetap kesamaan;
memaksimumkan $f$ sama dengan meminimumkan $-f$. Yang tidak gratis hanyalah konveksitas:
membalik tanda ketaksamaan membalik konveksitas fungsinya, jadi arah ketaksamaan tidak
boleh diubah tanpa memeriksa ulang.

\begin{akibat}[Daerah layak konveks]
Pada masalah konveks, daerah layak $X$ konveks.
\end{akibat}

\begin{proof}
Himpunan bawah $S_{i}=\{\mathbf{x}:f_{i}(\mathbf{x})\le0\}$ konveks untuk setiap $i$:
bila $f_{i}(\mathbf{x})\le0$ dan $f_{i}(\mathbf{y})\le0$, konveksitas memberi
$f_{i}(t\mathbf{x}+(1-t)\mathbf{y})\le t\,f_{i}(\mathbf{x})+(1-t)f_{i}(\mathbf{y})\le0$.
Himpunan $\{\mathbf{x}:h_{j}(\mathbf{x})=0\}$ adalah hiperbidang karena $h_{j}$ afin,
jadi konveks. Daerah layak adalah irisan semuanya, dan irisan himpunan konveks juga konveks
menurut proposisi di bagian pertama.
\end{proof}

Akibat itu yang membuat seluruh bab ini masuk akal: dua titik layak selalu bisa
dihubungkan tanpa keluar dari daerah layak. Bentuk paling sederhana polihedron adalah
daerah layak program linear dua variabel pada Gambar~\ref{fig:lp-verteks}: segi empat,
dan sasaran linearnya optimal di verteks.

\begin{figure}[!ht]
\centering
\begin{tikzpicture}
\begin{axis}[width=9cm,height=6cm,xlabel={$x_{1}$},ylabel={$x_{2}$}]
\addplot[forget plot,fill=blue!8,draw=blue!50!black,thick]
  coordinates {(0,0) (4,0) (3,1.5) (0,3)} \closedcycle;
\addplot[domain=0:2.1,densely dotted,thick] {2.5-1.25*x};
\addlegendentry{laba $10$}
\addplot[domain=0:2.9,densely dotted,thick] {3.75-1.25*x};
\addlegendentry{laba $15$}
\addplot[domain=0:4.1,thick,red] {5.25-1.25*x};
\addlegendentry{laba $21$}
\addplot[forget plot,only marks,mark=*] coordinates {(3,1.5)};
\node[font=\scriptsize,below right] at (axis cs:3,1.5) {$(3;\,1{,}5)$};
\end{axis}
\end{tikzpicture}
\caption{Daerah layak program linear bengkel Bab 24: $6x_{1}+4x_{2}\le24$ (kayu lapis),
$x_{1}+2x_{2}\le6$ (jam), $x_{1},x_{2}\ge0$, sasaran $5x_{1}+4x_{2}$. Verteks
$(0;0)$, $(4;0)$, $(3;1{,}5)$, $(0;3)$ bernilai $0$, $20$, $21$, $12$. Garis level $21$
hanya menyentuh verteks $(3;1{,}5)$; garis $20$ masih memotong daerah layak.}
\label{fig:lp-verteks}
\end{figure}


Sekarang pernyataan yang dijanjikan pembuka bab, dan buktinya pendek justru karena
daerah layak konveks.

\begin{teorema}[Minimum lokal adalah minimum global]
Misalkan $f_{0}$ konveks pada himpunan konveks $X$, dan $\mathbf{x}^{*}\in X$ memenuhi
$f_{0}(\mathbf{x}^{*})\le f_{0}(\mathbf{x})$ untuk semua $\mathbf{x}\in X$ yang cukup
dekat dengan $\mathbf{x}^{*}$. Maka $f_{0}(\mathbf{x}^{*})\le f_{0}(\mathbf{x})$ untuk
semua $\mathbf{x}\in X$. Himpunan semua titik minimum juga konveks.
\index{optimum global}
\end{teorema}

\begin{proof}
Andaikan ada $\mathbf{y}\in X$ dengan $f_{0}(\mathbf{y})<f_{0}(\mathbf{x}^{*})$. Untuk
$t\in(0,1]$ titik $\mathbf{z}_{t}=(1-t)\mathbf{x}^{*}+t\mathbf{y}$ berada di $X$ dan
memusat ke $\mathbf{x}^{*}$ saat $t\to0$, jadi untuk $t$ cukup kecil ia masuk ke
tetangga tempat $\mathbf{x}^{*}$ minimum lokal. Tetapi konveksitas memberi
\[
  f_{0}(\mathbf{z}_{t})\le(1-t)f_{0}(\mathbf{x}^{*})+t\,f_{0}(\mathbf{y})
  <f_{0}(\mathbf{x}^{*}),
\]
bertentangan dengan keminimuman lokal. Untuk bagian kedua, bila
$\mathbf{x}^{*},\mathbf{y}^{*}$ keduanya minimum bernilai $p^{*}$, maka
$f_{0}(\mathbf{z}_{t})\le p^{*}$, sedangkan $p^{*}$ adalah nilai terkecil yang mungkin,
jadi $f_{0}(\mathbf{z}_{t})=p^{*}$ dan $\mathbf{z}_{t}$ juga minimum.
\end{proof}

\begin{akibat}[Syarat stasioner tanpa kendala]
Bila $f$ konveks dan terdiferensialkan pada himpunan terbuka, maka
$\nabla f(\mathbf{x}^{*})=\mathbf{0}$ jika dan hanya jika $\mathbf{x}^{*}$ minimum
global.
\end{akibat}

\begin{proof}
Bila $\nabla f(\mathbf{x}^{*})=\mathbf{0}$, uji orde pertama memberi
$f(\mathbf{y})\ge f(\mathbf{x}^{*})$ untuk setiap $\mathbf{y}$. Sebaliknya, bila
$\mathbf{x}^{*}$ minimum, maka untuk setiap arah $\mathbf{v}$ dan $t>0$ kecil berlaku
$f(\mathbf{x}^{*}+t\mathbf{v})\ge f(\mathbf{x}^{*})$, sehingga turunan berarahnya
$\nabla f(\mathbf{x}^{*})^{\top}\mathbf{v}\ge0$. Karena $\mathbf{v}$ juga boleh diganti
$-\mathbf{v}$, haruslah $\nabla f(\mathbf{x}^{*})=\mathbf{0}$.
\end{proof}

\begin{tangan}
Bahan: tali dua meter, satu cincin kecil, penutup mata, dan seorang teman.

Ikat tali pada dua tiang setinggi sama, lalu regangkan sampai lurus. Minta teman menutup
mata, menggantungkan cincin di tengah tali, dan menggesernya perlahan ke kiri dan ke
kanan dengan satu jari. Cincin meluncur ke titik terendah, ke mana pun ia dilepas. Minta
teman menandai titik berhentinya tanpa membuka mata; ulangi lima kali dari posisi awal
berbeda. Semua tanda harus jatuh pada satu titik.

Sekarang ikat satu simpul kecil di tali, agak ke kiri, sehingga talinya membentuk dua
lekuk. Ulangi percobaan. Cincin berhenti di lekuk kiri bila dilepas di dekat situ, dan
di tengah bila dilepas di kanan simpul: dua tempat berhenti, dua lembah.

Yang harus terlihat: tali tanpa simpul hanya punya satu titik terendah, tali bersimpul
punya dua, dan jari yang berhenti di lekuk kiri tidak tahu apa yang terjadi di kanan,
persis seperti algoritme yang berhenti di titik stasioner pertama.
\end{tangan}

Tanpa konveksitas, akibat itu salah. Ambil $f(x_{1},x_{2})=x_{1}x_{2}$ pada bujur
sangkar $[-1,1]^{2}$. Gradiennya $(x_{2},x_{1})$ lenyap di titik asal, padahal
$f(\mathbf{0})=0$ sedangkan $f(1,-1)=-1$. Hessiannya
$\begin{pmatrix}0&1\\1&0\end{pmatrix}$ dengan nilai eigen $1$ dan $-1$, jadi uji orde
kedua menolak konveksitas: titik asal pelana, bukan lembah. Bab 24 memberi contoh yang
lebih menipu: pembulatan jawaban bersambung memberi laba 19, padahal jawaban bulat
terbaik memberi 20.

\section{Lagrangian dan batas bawah dari dual}
\tingkat{lanjut}

Kendala selalu berbentuk ``sesuatu tidak boleh positif'', jadi kita bisa memberi harga
pada pelanggaran: setiap kali $f_{i}(\mathbf{x})$ positif, kita membayar denda
$\lambda_{i}$ per satuan, dan setiap penyimpangan $h_{j}(\mathbf{x})$ dari nol dibayar
atau diterima sebesar $\nu_{j}$ per satuan. Fungsi yang dihasilkan tidak lebih besar
daripada sasaran asli, asalkan harganya tak negatif. Harga buruk memberi batas longgar;
harga baik memberi batas tepat.

\begin{definisi}[Lagrangian dan fungsi dual]
Untuk masalah bentuk baku, \emph{Lagrangian}-nya adalah
\[
  L(\mathbf{x},\boldsymbol{\lambda},\boldsymbol{\nu})
  =f_{0}(\mathbf{x})+\sum_{i=1}^{m}\lambda_{i}f_{i}(\mathbf{x})
   +\sum_{j=1}^{p}\nu_{j}h_{j}(\mathbf{x}),
\]
dengan $\boldsymbol{\lambda}\in\mathbb{R}^{m}$ dan $\boldsymbol{\nu}\in\mathbb{R}^{p}$
disebut \emph{pengali} atau \emph{variabel dual}. \emph{Fungsi dual} adalah infimumnya
terhadap variabel asal,
\[
  g(\boldsymbol{\lambda},\boldsymbol{\nu})
  =\inf_{\mathbf{x}\in D}L(\mathbf{x},\boldsymbol{\lambda},\boldsymbol{\nu}),
\]
dan \emph{masalah dual} adalah memaksimumkan $g$ dengan syarat
$\boldsymbol{\lambda}\ge\mathbf{0}$. Pasangan $(\boldsymbol{\lambda},\boldsymbol{\nu})$
disebut \emph{layak dual} bila $\boldsymbol{\lambda}\ge\mathbf{0}$ dan
$g(\boldsymbol{\lambda},\boldsymbol{\nu})>-\infty$. Nilai optimumnya ditulis $d^{*}$.
\index{Lagrangian}\index{dualitas}
\end{definisi}

\begin{proposisi}[Fungsi dual selalu cekung]
$g$ adalah infimum fungsi-fungsi afin dalam $(\boldsymbol{\lambda},\boldsymbol{\nu})$,
karena untuk setiap $\mathbf{x}$ tetap, $L(\mathbf{x},\cdot,\cdot)$ afin. Karena itu $g$
cekung, dan masalah dual selalu masalah konveks, apa pun bentuk masalah asalnya.
\end{proposisi}

\begin{proof}
Ambil dua titik dual $(\boldsymbol{\lambda},\boldsymbol{\nu})$ dan
$(\boldsymbol{\lambda}',\boldsymbol{\nu}')$ serta $t\in[0,1]$. Untuk setiap $\mathbf{x}$,
kelinearan memberi
\[
  L\bigl(\mathbf{x},\,t\boldsymbol{\lambda}+(1-t)\boldsymbol{\lambda}',\,
  t\boldsymbol{\nu}+(1-t)\boldsymbol{\nu}'\bigr)
  =t\,L(\mathbf{x},\boldsymbol{\lambda},\boldsymbol{\nu})
   +(1-t)\,L(\mathbf{x},\boldsymbol{\lambda}',\boldsymbol{\nu}').
\]
Ambil infimum kedua ruas terhadap $\mathbf{x}$; infimum jumlah tidak kurang daripada
jumlah infimum, sehingga
$g\ge t\,g(\boldsymbol{\lambda},\boldsymbol{\nu})
+(1-t)g(\boldsymbol{\lambda}',\boldsymbol{\nu}')$, yaitu $g$ cekung.
\end{proof}

Akibat yang menyenangkan: memaksimumkan fungsi cekung dengan kendala
$\boldsymbol{\lambda}\ge\mathbf{0}$ selalu masalah konveks. Setiap soal, sekalipun penuh
lubang, punya kembaran konveks; kembaran itu tidak selalu memberi jawaban yang sama,
tetapi selalu memberi batas.

\begin{teorema}[Dualitas lemah]
Untuk setiap $(\boldsymbol{\lambda},\boldsymbol{\nu})$ layak dual dan setiap $\mathbf{x}$
layak asal berlaku $L(\mathbf{x},\boldsymbol{\lambda},\boldsymbol{\nu})\le
f_{0}(\mathbf{x})$. Karena itu $g(\boldsymbol{\lambda},\boldsymbol{\nu})\le p^{*}$,
yaitu $d^{*}\le p^{*}$.
\end{teorema}

\begin{proof}
Ambil $\mathbf{x}\in X$. Setiap $f_{i}(\mathbf{x})\le0$ dan $\lambda_{i}\ge0$, jadi
$\lambda_{i}f_{i}(\mathbf{x})\le0$. Setiap $h_{j}(\mathbf{x})=0$, jadi
$\nu_{j}h_{j}(\mathbf{x})=0$ berapa pun $\nu_{j}$. Karena itu
$L(\mathbf{x},\boldsymbol{\lambda},\boldsymbol{\nu})\le f_{0}(\mathbf{x})$. Ketaksamaan
itu berlaku untuk setiap $\mathbf{x}$ layak, sehingga juga berlaku bagi infimum atas
semua $\mathbf{x}$ layak; infimum atas seluruh $D$ tidak lebih besar, jadi
$g(\boldsymbol{\lambda},\boldsymbol{\nu})\le p^{*}$. Ketaksamaan terakhir diambil atas
semua pasangan layak dual.
\end{proof}

Bukti itu tidak memakai konveksitas dan tidak memakai keterdiferensialan, sehingga
dualitas lemah berlaku untuk soal integer, kombinatorial, dan soal apa pun dalam bentuk
baku. Itulah sebabnya batas bawah dari dual dipakai di percabangan dan pembatasan pada
Bab 24: batas itu sah meskipun relaksasinya bukan jawaban.

Selisih $p^{*}-d^{*}$ disebut \emph{celah dualitas}. Karena
$d^{*}\le p^{*}$, celah itu tidak pernah negatif. Bila celahnya nol, dikatakan
\emph{dualias kuat} berlaku; bagian berikutnya membahas kapan itu terjadi.
\index{dualitas lemah}\index{celah dualitas}

Sebelum itu, satu akibat yang dipakai memeriksa hasil hitungan. Bila ada titik layak
$\mathbf{x}$ dan pasangan layak dual $(\boldsymbol{\lambda},\boldsymbol{\nu})$ dengan
$f_{0}(\mathbf{x})=g(\boldsymbol{\lambda},\boldsymbol{\nu})$, maka
$d^{*}\ge g=f_{0}(\mathbf{x})\ge p^{*}\ge d^{*}$, sehingga semua ketaksamaan menjadi
kesamaan: titik itu optimal, pasangan itu optimal, dan celahnya nol. Kesamaan nilai adalah
sertifikat optimalitas yang bisa diperiksa tanpa mengetahui jawabannya lebih dulu.

\begin{untukanak}
Pernahkah kamu menawar di pasar? Kamu menyebut satu harga, penjual menyebut harga lain,
dan kalau tawar-menawar berhenti di tengah, kalian berdua belum tahu harga sebenarnya.
Tetapi kalian sudah tahu satu hal: penjual tidak akan mau di bawah angka tertentu, dan
kamu tidak akan mau di atas angka tertentu. Dua angka itu mengurung harga sebenarnya,
walaupun harga itu belum pernah disebut siapa pun.

Gagasan dualitas bekerja persis begitu. Alih-alih mencari jawaban, kita memasang harga
untuk setiap tuntutan, lalu bertanya: dengan daftar harga ini, berapa hasil terburuk yang
masih mungkin? Setiap daftar harga yang masuk akal memberi satu angka yang tidak mungkin
dilampaui jawaban sebenarnya. Kalau ada daftar harga yang angkanya sama dengan jawaban
yang kita temukan, urusannya selesai.

Di bab ini kita akan tahu mengapa daftar harga semacam itu selalu ada, mengapa ia kadang
tidak cukup menutup jarak, dan mengapa jarak yang tersisa bisa diukur dengan angka.
\end{untukanak}

\section{Dualitas kuat dan syarat Slater}
\tingkat{lanjut}

Konveksitas membuat minimum lokal menjadi global, tetapi belum menjamin celah dualitas
nol. Yang hilang adalah syarat tentang \emph{geometri kelayakan}: bila daerah layak hanya
menyentuh batas kendala tanpa masuk ke dalamnya, tidak ada bidang pemisah yang cukup rapi.
Syarat itu dinamai Slater.

\begin{teorema}[Syarat Slater]
Misalkan masalahnya konveks, $f_{0},\dots,f_{m}$ terdiferensialkan, dan
$h_{j}(\mathbf{x})=\mathbf{a}_{j}^{\top}\mathbf{x}-b_{j}$. Bila terdapat titik layak
$\mathbf{x}\in\operatorname{ri}D$ dengan $f_{i}(\mathbf{x})<0$ untuk setiap $i$ dan
$h_{j}(\mathbf{x})=0$ untuk setiap $j$, maka dualitas kuat berlaku, $d^{*}=p^{*}$, dan
nilai optimum dual tercapai. Bila sebagian kendala ketaksamaan ternyata afin, cukup
kendala yang tidak afin yang dipenuhi secara ketat.
\index{Slater}
\end{teorema}

Titik dengan $f_{i}(\mathbf{x})<0$ disebut \emph{layak ketat}. Adanya satu titik semacam
itu sudah cukup, dan syaratnya tidak menuntut apa pun tentang fungsi sasarannya.

Buktinya tidak dikerjakan di sini. Jalannya begini: himpunan
$\mathcal{A}=\{(u,t):\exists\mathbf{x},\ f_{i}(\mathbf{x})\le u_{i},\
h_{j}(\mathbf{x})=0,\ f_{0}(\mathbf{x})\le t\}$ konveks, dan titik $(0,p^{*})$ berada di
batasnya; teorema pemisah hiperbidang (Bab 10 memakai gagasan serupa untuk proyeksi)
memberi bidang pendukung dengan koefisien $(\boldsymbol{\lambda},\boldsymbol{\nu})$ dan
nilai $d^{*}$. Bukti lengkapnya di Boyd dan Vandenberghe bagian 5.3.2, dan teorema itulah
yang membenarkan langkah ``selesaikan KKT lalu percaya hasilnya''.

Slater cukup, bukan perlu, dan bedanya terlihat pada contoh kecil berikut. Ambil
$\min x$ dengan kendala $x^{2}\le0$. Daerah layaknya satu titik, $x=0$, jadi tidak ada
titik layak ketat dan Slater gagal. Nilai optimumnya $p^{*}=0$. Lagrangian
$L=x+\lambda x^{2}$ memberi $g(\lambda)=-1/(4\lambda)$ untuk $\lambda>0$, sebab minimum
terjadi di $x=-1/(2\lambda)$, sehingga $\sup_{\lambda\ge0}g(\lambda)=0=d^{*}$. Dualitas
kuat tetap berlaku, tetapi nilai itu hanya didekati, tidak dicapai: pada setiap
$\lambda$ berhingga, $g(\lambda)<0$. Jadi yang hilang ketika Slater gagal bukan selalu
kesamaan nilai, melainkan jaminan bahwa optimum dual tercapai.

Sebaliknya, Slater menjamin sesuatu yang tidak diberikan dualitas kuat sendirian: pengali
optimal ada, sehingga persamaan KKT punya solusi untuk dicari.

\section{Kondisi KKT: perlu, dan cukup bila konveks}
\tingkat{lanjut}

Kalau pengali sudah optimal, titik asal yang optimal harus meminimumkan Lagrangian tanpa
kendala, dan di situ gradiennya nol. Menguraikan syarat itu komponen demi komponen
memberi empat kelompok persamaan: kondisi Karush-Kuhn-Tucker.

\kamus{pengali Lagrange}{Lagrange multiplier}{Harga per satuan pelanggaran sebuah
kendala; nol berarti kendala itu tidak sedang menekan, positif berarti ia menekan.}

\begin{definisi}[Kondisi KKT]
Titik $(\mathbf{x},\boldsymbol{\lambda},\boldsymbol{\nu})$ memenuhi \emph{kondisi KKT}
bila
\[
  \begin{array}{ll}
  f_{i}(\mathbf{x})\le0,\quad h_{j}(\mathbf{x})=0, & \text{(kelayakan asal)}\\
  \lambda_{i}\ge0, & \text{(kelayakan dual)}\\
  \lambda_{i}f_{i}(\mathbf{x})=0,\quad i=1,\dots,m, & \text{(kelengkapan komplementer)}\\
  \nabla f_{0}(\mathbf{x})+\sum_{i}\lambda_{i}\nabla f_{i}(\mathbf{x})
  +\sum_{j}\nu_{j}\nabla h_{j}(\mathbf{x})=\mathbf{0}. & \text{(kestasioneran)}
  \end{array}
\]
\index{KKT}\index{pengali Lagrange}\index{kelengkapan komplementer}
\end{definisi}

Kelompok ketiga paling banyak menimbulkan salah paham, jadi mari diterjemahkan: pengali
hanya boleh bukan nol untuk kendala yang benar-benar menekan. Kendala yang longgar tidak
dibayar apa pun, karena melonggarkannya tidak mengubah apa-apa. Pengali nol pada kendala
yang menekan juga mungkin, dan itu tanda kemerosotan.

Sisi geometri kelompok terakhir terlihat pada Gambar~\ref{fig:kkt-singgung}.

\begin{figure}[!ht]
\centering
\begin{tikzpicture}[scale=1.42]
  \draw[->,layergrey] (-0.3,0)--(4.4,0) node[right,font=\small]{$x_{1}$};
  \draw[->,layergrey] (0,-0.3)--(0,4.4) node[above,font=\small]{$x_{2}$};
  \fill[layerteal!10] (0,0) circle (2);
  \begin{scope}
    \clip (-0.3,-0.3) rectangle (4.3,4.3);
    \draw[layergrey,densely dashed] (2.6,2.6) circle (1.25);
    \draw[layergrey,densely dashed] (2.6,2.6) circle (2.15);
    \draw[thick,accent] (2.6,2.6) circle (1.677);
  \end{scope}
  \draw[very thick,layerteal] (0,0) circle (2);
  \node[font=\scriptsize,layerteal,fill=white,inner sep=1pt,anchor=west]
    at (-0.4,1.9) {$f_{1}(\mathbf{x})=0$};
  \node[font=\scriptsize,accent,fill=white,inner sep=1pt,anchor=south east]
    at (1.1,2.02) {$f_{0}=2{,}81$};
  \node[font=\scriptsize,layergrey,fill=white,inner sep=1pt]
    at (2.6,3.9) {nilai lebih kecil, tidak layak};
  \fill[layergrey] (2.6,2.6) circle (1.4pt);
  \node[font=\scriptsize,layergrey,anchor=west] at (2.72,2.6) {$\mathbf{c}$};
  \fill[accent] (1.414,1.414) circle (1.9pt);
  \node[font=\scriptsize,accent,below left=-2pt] at (1.414,1.414)
    {$\mathbf{x}^{*}$};
  \draw[-{Latex},very thick,accent] (1.414,1.414)--(0.8,0.8)
    node[below right=-1pt,font=\scriptsize]{$\nabla f_{0}$};
  \draw[-{Latex},very thick,layerteal] (1.414,1.414)--(2.03,2.03)
    node[above left=-1pt,font=\scriptsize]{$\lambda\nabla f_{1}$};
\end{tikzpicture}
\caption{Kontur sasaran $f_{0}(\mathbf{x})=(x_{1}-2{,}6)^{2}+(x_{2}-2{,}6)^{2}$
sepusat di titik $\mathbf{c}=(2{,}6;\,2{,}6)$, yaitu minimum tanpa kendala yang
jatuh di luar daerah layak
$f_{1}(\mathbf{x})=x_{1}^{2}+x_{2}^{2}-4\le0$. Kontur $f_{0}=2{,}81$ hanya
menyinggung lingkaran kendala di $\mathbf{x}^{*}=(\sqrt{2};\sqrt{2})$; kontur
putus-putus yang lebih dalam bernilai lebih kecil tetapi seluruhnya di luar
daerah layak, sedangkan yang lebih luar masih memotong daerah layak dengan nilai
lebih besar. Di titik singgung itu $\nabla f_{0}$ dan $\nabla f_{1}$ sejajar
berlawanan arah, dengan pengali positif.}
\label{fig:kkt-singgung}
\end{figure}

Sekarang pernyataan yang menjadi judul bagian ini, dan yang harus dibaca dengan hati-hati
karena dua arahnya punya syarat berbeda.

\begin{teorema}[KKT cukup untuk optimalitas]
Misalkan masalahnya konveks, $f_{0},\dots,f_{m}$ terdiferensialkan, dan $h_{j}$ afin.
Bila $(\mathbf{x},\boldsymbol{\lambda},\boldsymbol{\nu})$ memenuhi kondisi KKT, maka
$\mathbf{x}$ optimal asal, $(\boldsymbol{\lambda},\boldsymbol{\nu})$ optimal dual, dan
celah dualitasnya nol.
\end{teorema}

\begin{proof}
Karena $\lambda_{i}\ge0$ dan $f_{i}$ konveks, setiap suku $\lambda_{i}f_{i}$ konveks;
karena $h_{j}$ afin, setiap suku $\nu_{j}h_{j}$ afin. Jadi
$L(\cdot,\boldsymbol{\lambda},\boldsymbol{\nu})$ konveks dalam $\mathbf{x}$. Syarat
kestasioneran mengatakan gradiennya lenyap di $\mathbf{x}$, jadi menurut akibat di bagian
ketiga titik itu meminimumkan $L$ secara global:
$g(\boldsymbol{\lambda},\boldsymbol{\nu})=L(\mathbf{x},\boldsymbol{\lambda},
\boldsymbol{\nu})$. Kelayakan asal dan kelengkapan komplementer menghapus semua suku
kendala:
\[
  g(\boldsymbol{\lambda},\boldsymbol{\nu})
  =f_{0}(\mathbf{x})+\sum_{i}\lambda_{i}f_{i}(\mathbf{x})
   +\sum_{j}\nu_{j}h_{j}(\mathbf{x})=f_{0}(\mathbf{x}).
\]
Karena $g\le d^{*}\le p^{*}\le f_{0}(\mathbf{x})$ menurut dualitas lemah, semua
ketaksamaan itu menjadi kesamaan. Karena itu $\mathbf{x}$ optimal asal, pasangan itu
optimal dual, dan $p^{*}-d^{*}=0$.
\end{proof}

Arah sebaliknya berbeda sifatnya, dan perbedaannya bukan teknis.

\begin{teorema}[KKT perlu di bawah syarat kualifikasi]
Misalkan masalahnya konveks dan terdiferensialkan, memenuhi syarat Slater, dan
$\mathbf{x}^{*}$ optimal asal dengan $p^{*}$ berhingga. Maka terdapat
$(\boldsymbol{\lambda}^{*},\boldsymbol{\nu}^{*})$ sehingga
$(\mathbf{x}^{*},\boldsymbol{\lambda}^{*},\boldsymbol{\nu}^{*})$ memenuhi kondisi KKT.
\end{teorema}

Buktinya memakai teorema pemisah hiperbidang pada himpunan capaian, dan tidak diulang di
sini; lihat Boyd dan Vandenberghe bagian 5.5.3. Yang perlu diingat adalah peran Slater di
dalamnya: tanpa syarat kualifikasi, titik optimal bisa tidak memenuhi KKT sama sekali,
sehingga ``selesaikan KKT'' bukan jalan yang sah untuk mencari jawaban.

Contoh penolak yang paling kecil adalah $\min x$ dengan kendala $x^{2}\le0$. Titik
optimalnya $x^{*}=0$ dan $\nabla f_{0}=1$, sedangkan $\nabla f_{1}=2x=0$ di titik itu.
Syarat kestasioneran menuntut $1+\lambda\cdot0=0$, yang mustahil. Jadi titik optimal itu
tidak memenuhi KKT, dan yang gagal hanya kualifikasi kendala: daerah layaknya satu titik,
tidak ada arah untuk bergerak, tidak ada bidang pendukung.

Kesimpulan logisnya, dan ini sering dikutip salah: pada masalah konveks, KKT
\emph{cukup} tanpa syarat tambahan, tetapi \emph{perlu} hanya bila ada kualifikasi
kendala. Pernyataan ``$\mathbf{x}$ optimal jika dan hanya jika ada pengali yang memenuhi
KKT'' menuntut kedua arah itu, jadi Slater termasuk di dalamnya.

Satu tafsiran untuk mengingat pengali. Misalkan kendala ke-$i$ dilonggarkan
menjadi $f_{i}(\mathbf{x})\le u_{i}$, dan nilai optimumnya $p^{*}(u)$. Bila $p^{*}$
terdiferensialkan di $u=0$, maka
\[
  \lambda_{i}^{*}=-\frac{\partial p^{*}}{\partial u_{i}}(0).
\]
Pengali adalah harga bayangan: berapa nilai sasaran turun bila kendala itu dilonggarkan
satu satuan. Pada contoh hitung nanti, melonggarkan pembatas kelarutan dari setengah
menjadi $0{,}6$ menurunkan penyimpangan dari $2{,}5$ ke $2{,}32$, sedangkan harga bayangan
memprediksi $2{,}5-2(0{,}1)=2{,}30$. Selisih $0{,}02$ adalah suku orde
kedua, dan arahnya selalu sama: ramalan linear itu batas bawah, sehingga perbaikan yang
dijanjikannya tidak pernah terlampaui.

Tafsiran mekanisnya, yang sudah dipakai Lagrange sendiri: dua balok dihubungkan pegas dan
tidak boleh saling menembus atau menembus dinding. Energi potensial adalah fungsi sasaran;
pada kesetimbangan, kestasioneran adalah keseimbangan gaya, dan pengali adalah gaya kontak
dinding. Gaya kontak hanya menekan, tidak menarik, itulah $\lambda_{i}\ge0$; ia muncul
hanya bila terjadi sentuhan, itulah kelengkapan komplementer (Boyd dan Vandenberghe bagian
5.5.4).

\begin{catatan}[Subgradien]
Bila $f$ konveks tetapi tidak terdiferensialkan, gradien diganti \emph{subgradien}:
$\mathbf{s}$ adalah subgradien $f$ di $\mathbf{x}$ bila
$f(\mathbf{y})\ge f(\mathbf{x})+\mathbf{s}^{\top}(\mathbf{y}-\mathbf{x})$ untuk semua
$\mathbf{y}$; himpunan semua subgradien ditulis $\partial f(\mathbf{x})$. Untuk
$f(x)=|x|$ di $x=0$ himpunan itu selang $[-1,1]$, sebab setiap kemiringan antara $-1$
dan $1$ memberi garis di bawah grafik. Syarat optimal tanpa kendala menjadi
$\mathbf{0}\in\partial f(\mathbf{x}^{*})$, dan KKT tetap berlaku dengan $\nabla f_{i}$
diganti subgradien (Rockafellar bagian 23).\index{subgradien}
\end{catatan}

\section{Contoh hitung}
\tingkat{menengah}

Tiga hitungan berikut dikerjakan sampai angka terakhir, dengan pecahan sederhana sehingga
dapat diperiksa dengan tangan.

\emph{Pertama, uji konveksitas lewat nilai eigen.} Ambil
$f(x_{1},x_{2})=x_{1}^{2}+4x_{1}x_{2}+5x_{2}^{2}$. Hessiannya konstan,
\[
  \nabla^{2}f=\begin{pmatrix}2&4\\4&10\end{pmatrix},
  \qquad
  \operatorname{tr}=12,\quad \det=20-16=4 .
\]
Nilai eigennya $\lambda_{1,2}=6\pm\sqrt{36-4}=6\pm4\sqrt{2}$, yaitu
$0{,}343146$ dan $11{,}656854$. Keduanya positif, jadi Hessian definit positif dan $f$
konveks ketat. Pembandingnya $g(x_{1},x_{2})=x_{1}^{2}-x_{2}^{2}$ punya Hessian
$\operatorname{diag}(2,-2)$ dengan nilai eigen $2$ dan $-2$: tidak konveks, dan titik
kritisnya di titik asal adalah pelana. Fungsi $h(x_{1},x_{2})=x_{1}^{2}$ punya nilai
eigen $0$ dan $2$: konveks, tidak ketat, dan minimumnya seluruh garis $x_{1}=0$,
sehingga jawabannya tidak tunggal.

\begin{tantangan}
Ambil $f(x_{1},x_{2})=x_{1}^{2}x_{2}^{2}$. Tahan $x_{2}=c$ tetap: grafiknya parabola
yang membuka ke atas, jadi setiap irisan mendatar tampak konveks, dan hal yang sama
berlaku untuk setiap irisan tegak. Apakah $f$ konveks pada $\mathbb{R}^{2}$? Bila tidak,
tunjukkan dua titik yang membuktikannya, lalu jelaskan mengapa pemeriksaan per irisan
tidak pernah cukup.

Petunjuk 1: jangan bergerak sejajar sumbu. Gerakkan kedua koordinat sekaligus, lalu
periksa titik tengahnya.

Petunjuk 2: hitung $f(1,0)$, $f(0,1)$, dan nilai $f$ di titik tengah potongan itu; atau
hitung $\det\nabla^{2}f$ di $(1,1)$, yang sama dengan $-12x_{1}^{2}x_{2}^{2}$.

Kalau mentok, kembali ke bagian Fungsi konveks dan dua ujinya.
\end{tantangan}

\emph{Kedua, masalah pencampur warna lewat KKT.} Dengan $x_{1},x_{2}$ takaran kedua
pigmen, masalahnya
\[
  \text{minimumkan } (x_{1}-2)^{2}+(x_{2}-2)^{2}
  \quad\text{dengan}\quad x_{1}+x_{2}=2,\quad x_{1}\le\tfrac12 .
\]
Tulis $g(\mathbf{x})=x_{1}-\tfrac12\le0$ dan $h(\mathbf{x})=x_{1}+x_{2}-2=0$.
Kestasioneran memberi
\[
  2(x_{1}-2)+\lambda+\nu=0,\qquad 2(x_{2}-2)+\nu=0 ,
\]
dan kelengkapan komplementer memberi $\lambda(x_{1}-\tfrac12)=0$. Bila $\lambda=0$, kedua
persamaan itu memberi $x_{1}=x_{2}$, dan dengan $x_{1}+x_{2}=2$ diperoleh $x_{1}=x_{2}=1$,
yang melanggar $x_{1}\le\tfrac12$. Jadi $\lambda>0$: $x_{1}=\tfrac12$ dan $x_{2}=\tfrac32$.
Persamaan kedua memberi $\nu=2(2-\tfrac32)=1$, persamaan pertama memberi
$\lambda=4-2x_{1}-\nu=4-1-1=2$. Karena $\lambda=2\ge0$, titik itu optimal global menurut
teorema kecukupan, dengan
\[
  p^{*}=\bigl(\tfrac12-2\bigr)^{2}+\bigl(\tfrac32-2\bigr)^{2}
  =\tfrac94+\tfrac14=\tfrac52=2{,}5 .
\]

Pemeriksaan lewat dual, bukan lewat kepercayaan. Fungsi dualnya dihitung dengan
melengkapkan kuadrat pada tiap variabel:
\[
  \min_{x_{1}}\Bigl[(x_{1}-2)^{2}+(\lambda+\nu)x_{1}\Bigr]
  =4-\frac{(\lambda+\nu-4)^{2}}{4},
  \qquad
  \min_{x_{2}}\Bigl[(x_{2}-2)^{2}+\nu x_{2}\Bigr]
  =4-\frac{(\nu-4)^{2}}{4},
\]
sehingga
\[
  g(\lambda,\nu)=8-\frac{(\lambda+\nu-4)^{2}}{4}-\frac{(\nu-4)^{2}}{4}
  -\frac{\lambda}{2}-2\nu .
\]
Di $(\lambda,\nu)=(2,1)$ diperoleh $g=8-\tfrac14-\tfrac94-1-2=\tfrac52=p^{*}$: celahnya
nol, sesuai dualitas kuat yang dijamin Slater, sebab $(0{,}4;\,1{,}6)$ layak ketat. Turunan
parsialnya memberi $-\tfrac12(\lambda+\nu-4)-\tfrac12(\nu-4)-2=0$ dan
$-\tfrac12(\lambda+\nu-4)-\tfrac12=0$; dari yang kedua $\lambda+\nu=3$, lalu $\nu=1$ dan
$\lambda=2$. Pasangan itu satu-satunya pemaksimum.

Harga bayangannya bisa diperiksa dengan angka. Bila pembatas kelarutan dilonggarkan
menjadi $0{,}6$, jawabannya $x_{1}=0{,}6$, $x_{2}=1{,}4$, dan
$p^{*}=1{,}96+0{,}36=2{,}32$. Ramalan linear $p^{*}-\lambda\cdot0{,}1=2{,}5-0{,}2=2{,}30$
meleset $0{,}02$, dan $0{,}02=\tfrac12\cdot4\cdot(0{,}1)^{2}$ tepat, sebab turunan kedua
$p^{*}$ terhadap batas itu $4$. Bila jumlah yang dituntut dinaikkan menjadi $2{,}1$ dengan
pembatas tetap, jawabannya $x_{1}=\tfrac12$, $x_{2}=1{,}6$, dan
$p^{*}=2{,}25+0{,}16=2{,}41$; ramalan $2{,}5-\nu\cdot0{,}1=2{,}40$ meleset
$0{,}01=\tfrac12\cdot2\cdot(0{,}1)^{2}$. Kedua harga itu kemiringan yang bisa dipakai
memutuskan, dengan kesalahan yang bisa dihitung.

\emph{Ketiga, celah dualitas pada masalah nonkonveks.} Ambil satu variabel,
\[
  \text{minimumkan } (x-2)^{2}\quad\text{dengan}\quad x^{3}-x\le0 .
\]
Kendala itu setara dengan $x(x-1)(x+1)\le0$, jadi daerah layaknya
$(-\infty,-1]\cup[0,1]$: dua potong terpisah, tidak konveks. Pada $[0,1]$ sasaran menurun,
jadi nilai terkecilnya $1$ di $x=1$. Pada $(-\infty,-1]$ sasaran naik ke kiri dari $x=-1$
yang bernilai $9$. Jadi $p^{*}=1$, dicapai di $x=1$.

Sisi dual sangat berbeda. Lagrangiannya $L=(x-2)^{2}+\lambda(x^{3}-x)$. Untuk
$\lambda>0$ suku kubik mendominasi saat $x\to-\infty$ dan bertanda negatif, sehingga
$L\to-\infty$ dan $g(\lambda)=-\infty$: tidak ada batas bawah yang bisa dipungut. Hanya
$\lambda=0$ yang menyisakan nilai berhingga, dan di situ $g(0)=\min_{x}(x-2)^{2}=0$ di
$x=2$. Maka $d^{*}=0$ dan
\[
  p^{*}-d^{*}=1-0=1 ,
\]
celah sebesar $100\%$ dari nilai optimumnya sendiri. Batas bawah dari dual di sini bukan
batas yang longgar, melainkan batas yang salah tempat: $x=2$ jawaban tanpa kendala, dan
kendala itulah yang menentukan jawaban.

Contoh yang sama menunjukkan mengapa konveksitas tidak boleh dihapus dari teorema
kecukupan. Periksa titik $x=-1$, yang layak karena $(-1)^{3}-(-1)=0$. Di situ
$\nabla f_{0}=2(x-2)=-6$ dan $\nabla g=3x^{2}-1=2$, sehingga kestasioneran menuntut
$-6+2\lambda=0$, yaitu $\lambda=3$; kelengkapan komplementer terpenuhi karena kendalanya
aktif. Jadi $(x,\lambda)=(-1,3)$ memenuhi keempat kelompok kondisi KKT, padahal
$f_{0}(-1)=9$ sedangkan $p^{*}=1$.

\kertas{Ulangi hitungan kedua dengan pembatas $x_{1}\le1$. Sebelum menghitung, putuskan
apakah kendala itu masih menekan, lalu periksa apakah pengalinya masih
$2$. Setelah itu hitung nilai optimumnya dan bandingkan dengan $2{,}5$. Sebagai latihan
tambahan, cari nilai pembatas $c$ tempat kendala berhenti menekan, dan jelaskan mengapa
pengali harus nol tepat di $c$ itu.}

\section{Lab komputasi}
\tingkat{lanjut}

Dua blok berikut dijalankan dengan Python 3.14.4, numpy 2.5.2, dan scipy 1.18.0.
Keluaran disalin apa adanya.

\begin{lab}
\begin{verbatim}
import numpy as np
from scipy.optimize import minimize

np.set_printoptions(precision=8, suppress=True, linewidth=78)

# masalah pencampur: min (x1-2)^2+(x2-2)^2  s.t. x1+x2=2, x1<=0.5
f  = lambda x: (x[0]-2)**2 + (x[1]-2)**2
df = lambda x: np.array([2*(x[0]-2), 2*(x[1]-2)])
cons = [{"type": "eq",   "fun": lambda x: x[0]+x[1]-2,
         "jac": lambda x: np.array([1., 1.])},
        {"type": "ineq", "fun": lambda x: 0.5-x[0],
         "jac": lambda x: np.array([-1., 0.])}]
r = minimize(f, [0., 0.], jac=df, constraints=cons, method="SLSQP",
             options={"ftol": 1e-15, "maxiter": 500})
xs = r.x
print("x* =", xs, " f(x*) =", f(xs))
print("multipliers (scipy):", r.multipliers)

# pengali dari kestasioneran, memakai himpunan aktif {eq, ineq}
A = np.column_stack([[1., 0.], [1., 1.]])   # kolom: grad g, grad h
lam, nu = np.linalg.solve(A, -df(xs))
print("lambda, nu =", lam, nu)

# sisa KKT
print("h(x*)          =", xs[0]+xs[1]-2)
print("g(x*)          =", 0.5-xs[0])
print("lambda*g(x*)   =", lam*(0.5-xs[0]))
print("||grad L||     =", np.linalg.norm(df(xs) + lam*np.array([1.,0.])
                                         + nu*np.array([1.,1.])))

# fungsi dual dan pemaksimumannya
g = lambda lam, nu: (8 - (lam+nu-4)**2/4 - (nu-4)**2/4
                     - 0.5*lam - 2*nu)
print("g(2,1)         =", g(2., 1.))
rd = minimize(lambda z: -g(z[0], z[1]), [0.5, 0.5], method="SLSQP",
              bounds=[(0, None), (None, None)], options={"ftol": 1e-15})
print("dual argmax    =", rd.x, " d* =", -rd.fun)
print("celah p*-d*    =", f(xs) + rd.fun)

# uji nilai eigen Hessian
for name, H in [("x1^2+4x1x2+5x2^2", np.array([[2.,4.],[4.,10.]])),
                ("x1^2-x2^2",        np.array([[2.,0.],[0.,-2.]])),
                ("x1^2",             np.array([[2.,0.],[0.,0.]]))]:
    print(name, "eig =", np.linalg.eigvalsh(H))
\end{verbatim}
Keluaran sebenarnya:
\begin{verbatim}
x* = [0.5 1.5]  f(x*) = 2.500000000000001
multipliers (scipy): [-1.  2.]
lambda, nu = 2.0000000000000027 0.9999999999999987
h(x*)          = 0.0
g(x*)          = 6.106226635438361e-16
lambda*g(x*)   = 1.2212453270876738e-15
||grad L||     = 0.0
g(2,1)         = 2.5
dual argmax    = [1.99999997 0.99999999]  d* = 2.499999999999999
celah p*-d*    = 1.7763568394002505e-15
x1^2+4x1x2+5x2^2 eig = [ 0.34314575 11.65685425]
x1^2-x2^2 eig = [-2.  2.]
x1^2 eig = [0. 2.]
\end{verbatim}
\end{lab}

\begin{lab}
\begin{verbatim}
import numpy as np
from scipy.optimize import minimize

# nonkonveks: min (x-2)^2  s.t. x^3 - x <= 0
f3 = lambda x: (x[0]-2)**2
r3 = minimize(f3, [0.5], constraints=[{"type": "ineq",
              "fun": lambda x: -(x[0]**3 - x[0])}],
              method="SLSQP", options={"ftol": 1e-14, "maxiter": 500})
print("x* =", r3.x, " p* =", f3(r3.x))

# Lagrangian untuk beberapa lambda, dievaluasi jauh di kiri
for lam in (0.0, 1e-3, 1e-2, 0.1):
    L = lambda x: (x[0]-2)**2 + lam*(x[0]**3 - x[0])
    print(f"lam={lam:<6} L(-1e4)={L(np.array([-1e4])):.3e}",
          f" L(-1e6)={L(np.array([-1e6])):.3e}")

# titik yang memenuhi KKT tetapi bukan jawaban
for xv in (-1.0, 1.0, 0.0):
    f0p = 2*(xv-2); gp = 3*xv*xv - 1
    lam = -f0p/gp if gp != 0 else float("nan")
    print(f"x={xv:>4}  f0'={f0p:>5}  g'={gp:>4}  lam={lam:>6}"
          f"  f0={ (xv-2)**2 }  layak={xv**3-xv <= 0}")
\end{verbatim}
Keluaran sebenarnya:
\begin{verbatim}
x* = [1.]  p* = 0.999999999995369
lam=0.0    L(-1e4)=1.000e+08  L(-1e6)=1.000e+12
lam=0.001  L(-1e4)=-9.000e+08  L(-1e6)=-9.990e+14
lam=0.01   L(-1e4)=-9.900e+09  L(-1e6)=-9.999e+15
lam=0.1    L(-1e4)=-9.990e+10  L(-1e6)=-1.000e+17
x=-1.0  f0'= -6.0  g'= 2.0  lam=   3.0  f0=9.0  layak=True
x= 1.0  f0'= -2.0  g'= 2.0  lam=   1.0  f0=1.0  layak=True
x= 0.0  f0'= -4.0  g'=-1.0  lam=  -4.0  f0=4.0  layak=True
\end{verbatim}
\end{lab}

Empat hal perlu dibaca. Pertama, $x^{*}$ yang dilaporkan adalah $(0{,}5;\,1{,}5)$ dengan
nilai $2{,}500000000000001$; selisih $10^{-15}$ dari $2{,}5$ itu pembulatan pecahan biner.
Kedua, sisa KKT yang dihitung terpisah sangat kecil: kendala ketaksamaan
$6{,}1\times10^{-16}$, hasil kali $\lambda g$ sekitar $1{,}2\times10^{-15}$, dan norma
gradien Lagrangian nol, sehingga titik itu memang memenuhi syarat, bukan hampir
memenuhi.

Ketiga, yang paling mudah menjerat, \texttt{scipy} melaporkan pengali $(-1,\,2)$,
sedangkan kestasioneran memberi $\lambda=2{,}0000000000000027$ dan
$\nu=0{,}9999999999999987$, yaitu $(\nu,\lambda)=(1,\,2)$ sampai empat belas digit. Entri
kedua cocok; entri pertama berlawanan tanda. Konvensi tanda pengali pada rutin SQP berbeda
antar pustaka, jadi tanda tidak boleh ditebak: periksa ulang lewat persamaan kestasioneran.
Angka $2$ pada kendala ketaksamaan kebetulan cocok, dan itu menyamarkan kekeliruan
tanda.

Keempat, $d^{*}=2{,}499999999999999$ dan celah $p^{*}-d^{*}=1{,}78\times10^{-15}$:
nol sampai lima belas digit. Pemaksimum dual mengembalikan
$(\lambda,\nu)=(1{,}99999997,\;0{,}99999999)$, yaitu $(2,1)$ sampai delapan digit, dengan
selisih dari toleransi \texttt{ftol}. Kesepakatan tiga jalan, KKT tangan, KKT numerik, dan
pemaksimuman dual, yang membuat angka $2{,}5$ layak dipercaya.

Keluaran kedua berlawanan arah. Primalnya memberi $x^{*}=1$ dengan $p^{*}\approx1$,
sesuai hitungan tangan. Lagrangian dengan $\lambda=0{,}001$ sudah bernilai
$-9\times10^{8}$ di $x=-10^{4}$ dan $-9{,}99\times10^{14}$ di $x=-10^{6}$; menambah
$\lambda$ memperburuknya sebanding pangkat tiga, sehingga tidak ada $\lambda>0$ yang layak
dual dan $d^{*}=0$ dari $\lambda=0$ satu-satunya batas yang tersedia. Tiga titik
kritisnya: $x=-1$ dengan $\lambda=3$ memenuhi KKT dan bernilai $9$; $x=1$ jawaban sejati
bernilai $1$; $x=0$ gagal karena $\lambda=-4$.

\begin{untukdosen}
Urutan yang berhasil dipakai: mulai dari mangkuk dan kelereng, lalu definisi himpunan
konveks, lalu langsung ke teorema minimum lokal. Uji orde pertama dan kedua sebaiknya
datang setelah mahasiswa melihat mengapa konveksitas dipakai; kalau tidak, keduanya
terasa seperti dua rumus hafalan.

Dua miskonsepsi khas. Pertama, ``konveks berarti cembung'', sehingga mahasiswa menolak
mangkuk sebagai konveks; yang menolong adalah ruas garis antara dua titik di dalam
mangkuk. Kedua, KKT dianggap perlu-dan-cukup tanpa syarat, sehingga mahasiswa kecewa
ketika contoh $x^{3}-x\le0$ memberi titik KKT yang bukan jawaban. Contoh itu sebaiknya
ditulis di papan lebih awal, supaya kecukupan terlihat sebagai akibat konveksitas.

Aktivitas papan tulis: beri satu QP dua variabel dengan satu kendala ketaksamaan, minta
kelas menebak apakah kendalanya menekan, lalu selesaikan KKT dengan pencabangan
$\lambda=0$ dan $\lambda>0$, lalu hitung fungsi dualnya dan tunjukkan nilainya sama.
Mahasiswa yang bisa menjelaskan mengapa cabang $\lambda=0$ gagal di sini sudah siap untuk
dualitas.

Tanda kesiapan lanjut: mahasiswa dapat menyebut satu soal konveks dan satu nonkonveks
dari bacaan mereka sendiri, dan menyebut apa yang berubah pada jaminan optimalitasnya.
\end{untukdosen}

\section{Ringkasan}

\begin{ringkasan}
\begin{itemize}
  \item Himpunan konveks memuat setiap ruas antara dua anggotanya. Irisan mempertahankan
        konveksitas, gabungan tidak. Bola, hiperbidang, setengah ruang, dan polihedron
        konveks; himpunan $\{x_{1}x_{2}\ge1\}$ bukan.
  \item Fungsi konveks berada di bawah talinya. Uji orde pertama: grafik di atas setiap
        bidang singgung. Uji orde kedua: Hessian semidefinit positif, perlu dan cukup.
        Hessian definit positif cukup untuk konveks ketat, tetapi tidak perlu.
  \item Pada bentuk baku dengan $f_{i}$ konveks dan $h_{j}$ afin, daerah layak konveks,
        sehingga setiap minimum lokal adalah minimum global. Tanpa konveksitas, titik
        stasioner bisa pelana, seperti $x_{1}x_{2}$ di titik asal.
  \item Lagrangian memberi harga pada kendala, dan fungsi dual adalah infimumnya terhadap
        variabel asal. Fungsi dual selalu cekung, jadi masalah dual selalu konveks, dan
        dualitas lemah $d^{*}\le p^{*}$ berlaku tanpa syarat apa pun.
  \item Dualitas kuat $d^{*}=p^{*}$ butuh syarat Slater, yaitu satu titik layak ketat, yang
        sekaligus menjamin pengali optimal ada. Slater cukup, tidak perlu: pada $\min x$
        dengan $x^{2}\le0$ dualitas kuat berlaku tanpa tercapainya optimum dual.
  \item KKT cukup untuk optimalitas pada masalah konveks tanpa syarat tambahan, tetapi
        perlu hanya bila ada kualifikasi kendala. Pengali adalah harga bayangan, dan celah
        dualitas bisa positif, bahkan tak hingga, pada masalah nonkonveks: $\min (x-2)^{2}$
        dengan $x^{3}-x\le0$ memberi $p^{*}=1$ dan $d^{*}=0$.
\end{itemize}
\end{ringkasan}

\section{Latihan}

\begin{latihan}
\begin{enumerate}
  \item Tentukan apakah $f=x_{1}^{2}+x_{1}x_{2}+x_{2}^{2}$, $g=x_{1}x_{2}$, dan
        $h=2x_{1}^{2}-4x_{1}x_{2}+3x_{2}^{2}$ konveks, cekung, atau tidak keduanya.
        Sebutkan Hessian dan nilai eigennya; untuk yang semidefinit positif tetapi tidak
        definit, jelaskan akibatnya bagi ketunggalan minimum.

  \item Buktikan bahwa irisan dua himpunan konveks konveks, lalu berikan dua himpunan
        konveks yang gabungannya tidak konveks, dengan dua titik konkret dan satu titik
        campuran yang gagal.

  \item Buktikan bahwa bila $f_{1},f_{2}$ konveks maka $\max\{f_{1},f_{2}\}$ konveks.
        Periksa $f(x_{1},x_{2})=\max\{x_{1}^{2},x_{2}\}$ di $(0,0)$ dan $(1,1)$: apakah
        gradiennya ada, dan apakah subdiferensialnya di titik patahan memuat nol?
        (Petunjuk: pakai definisi subgradien dan uji orde pertama untuk fungsi tak mulus.)

  \item Selesaikan dengan KKT: minimumkan $x_{1}^{2}+x_{2}^{2}$ dengan $x_{1}+2x_{2}=4$
        dan $x_{1}\ge1$. Tentukan titik optimum, kedua pengali, dan nilai optimum, lalu
        cocokkan dengan \texttt{scipy.optimize.minimize}. (Petunjuk: selesaikan dulu tanpa
        $x_{1}\ge1$; bila jawaban itu melanggar, kendala menekan dan $x_{1}=1$.)

  \item Hitung fungsi dual dari $\min x_{1}^{2}+x_{2}^{2}$ dengan $x_{1}+x_{2}\ge2$,
        tentukan pemaksimumnya, dan bandingkan dengan nilai optimum asal $2$. Mengapa
        celahnya nol meskipun tidak ada kendala kesamaan?

  \item Nilai klaim ini: \emph{dualitas kuat menjamin optimum primal tercapai.} Benar,
        salah, atau belum cukup bukti? Sertakan contoh penyangkal bila perlu.

  \item Nilai klaim ini: \emph{kelengkapan komplementer menyatakan pengali kendala yang
        aktif selalu positif.} Benar, salah, atau belum cukup bukti? Tunjukkan apa yang
        sebenarnya dinyatakan syarat itu.

  \item Pada masalah pencampur warna, ganti pembatasnya menjadi $x_{1}\le1$ dan hitung
        ulang. Tentukan nilai pembatas $c$ tempat kendala berhenti menekan, dan tunjukkan
        pengalinya nol di $c$ itu. (Petunjuk: pakai syarat bahwa jawaban tanpa pembatas
        sudah memenuhi $x_{1}\le c$.)

  \item Untuk masalah nonkonveks $\min (x-2)^{2}$ dengan $x^{3}-x\le0$, buktikan
        $g(\lambda)=-\infty$ untuk setiap $\lambda>0$ dengan memilih barisan
        $x_{k}\to-\infty$, lalu jelaskan mengapa $d^{*}=0$. Hitung juga $p^{*}$ dengan
        memeriksa kedua potongan daerah layak.

  \item Buktikan bahwa bila $f$ konveks dan terdiferensialkan serta
        $\nabla f(\mathbf{x}^{*})=\mathbf{0}$, maka $\mathbf{x}^{*}$ minimum global,
        tanpa memakai ketatnya konveksitas. Mengapa bukti itu gagal untuk $f(x)=x^{3}$,
        dan di titik mana syaratnya dilanggar?
\end{enumerate}
\end{latihan}

\begin{untukpeneliti}
Yang benar-benar terbukti di bab ini: ketertutupan konveksitas di bawah irisan; kesetaraan
konveksitas dengan kedua uji turunan, dua arah; konveksitas daerah layak; teorema minimum
lokal sama dengan global; kecekungan fungsi dual; dualitas lemah tanpa syarat apa pun;
kecukupan KKT tanpa kualifikasi kendala. Yang dipinjam dua: teorema Slater dan keniscayaan
KKT, keduanya lewat teorema pemisah hiperbidang, dengan bukti di Boyd dan Vandenberghe
bagian 5.3.2 dan 5.5.3.

Yang dipopulerkan melampaui buktinya. ``Konveks berarti mudah'': konveksitas menjamin
setiap titik KKT adalah global, tetapi tidak memberi laju maupun jaminan bahwa algoritme
mencapainya. ``Pengali selalu harga'': tafsiran itu sah hanya bila $p^{*}$
terdiferensialkan terhadap ruas kanan kendala. ``Celah mengukur nonkonveksitas'': contoh
$x^{3}-x$ menunjukkan dual bisa lebih lemah daripada nilai minimum di selubung konveks
daerah layak.

Sumber primer: Boyd dan Vandenberghe bab 2 sampai 5; Nesterov, \emph{Lectures on Convex
Optimization}, bab 1 sampai 3; Rockafellar, \emph{Convex Analysis}, bagian 2 dan 23.

Batas keberlakuan: dimensi berhingga, fungsi terdiferensialkan kecuali pada catatan
subgradien, dan kendala dinyatakan eksplisit. Untuk kendala implisit, misalnya yang
didefinisikan simulasi, KKT tidak bisa ditulis dan yang tersisa metode tanpa turunan.
Pertanyaan terbuka: laju metode orde pertama pada masalah konveks besar masih diselidiki.
\end{untukpeneliti}

\peta{Bab ini mengikuti Boyd dan Vandenberghe, \emph{Convex Optimization}: bagian 2.1
sampai 2.3 (himpunan konveks, contoh penting, operasi yang mempertahankannya), 3.1 (sifat
dasar fungsi konveks, termasuk uji orde satu dan dua), 4.1 dan 4.2 (masalah optimasi dan
masalah konveks), 5.1 dan 5.2 (Lagrangian, dualitas lemah dan kuat, syarat Slater), 5.5
(KKT dan tafsiran mekanis). Sudut pandang sertifikat dual diambil dari Bertsimas dan
Tsitsiklis, \emph{Introduction to Linear Optimization}, bab 1 dan 4; untuk subgradien,
Rockafellar, \emph{Convex Analysis}, bagian 2 dan 23.}
\chapter{Pemrograman Linear dan Integer}

Sebuah bengkel kayu menerima 24 lembar kayu lapis dan 6 jam mesin perakitan untuk hari
itu. Satu rak buku memakai 6 lembar kayu lapis dan 1 jam perakitan. Satu bangku memakai
4 lembar dan 2 jam. Rak memberi laba 5 satuan, bangku 4 satuan. Berapa banyak rak dan
bangku yang sebaiknya dibuat?

Naluri pertama memilih rak, sebab labanya per unit lebih tinggi. Naluri itu salah: satu rak
menghabiskan kayu lapis sebanyak satu setengah bangku. Jawaban terbaik model bersambungnya
3 rak dan 1,5 bangku, berlaba 21, dan setengah bangku tidak bisa dijual. Pembulatan ke bawah
memberi 19; 4 rak tanpa bangku memberi 20. Jadi jawaban bulat terbaik bukan pembulatan
jawaban pecahan terbaik, dan selisihnya kesalahan cara berpikir, bukan kesalahan hitung.

Alat dari bab sebelumnya belum cukup. Bab 8 mengajari kita menyelesaikan
$A\mathbf{x}=\mathbf{b}$, tetapi di sini persamaannya lebih sedikit daripada variabelnya,
jadi ada tak berhingga jawaban dan kita harus memilih yang terbaik. Bab 23 memberi kondisi
KKT, tetapi pada soal ini gradien sasaran konstan dan tidak pernah lenyap: jawabannya
menempel di batas daerah layak, bukan di titik stasioner. Jaminan bab 21 bahwa optimum lokal
juga global tidak menunjukkan di mana optimumnya, dan hilang begitu variabel dituntut
bulat.

Karena itu bab ini menempuh tiga langkah: menulis soal sebagai program linear dan
mengenali bahwa jawabannya ada di sudut daerah layak; memberi harga pada setiap kendala
lewat dualitas, sehingga keoptimalan dapat diperiksa tanpa menyisir seluruh daerah; lalu
menghadapi tuntutan kebulatan, saat model bersambung memberi batas yang sah tetapi belum
tentu tercapai.

\begin{untukanak}
Bayangkan kamu punya 24 lembar kertas warna dan 6 jam untuk membuat dua jenis hiasan
kelas. Hiasan pertama memakai 6 lembar dan 1 jam, hiasan kedua memakai 4 lembar dan 2 jam.
Hiasan pertama lebih banyak dipesan, jadi harganya lebih tinggi. Berapa banyak masing-masing
yang sebaiknya kamu buat?

Kalau kamu coba-coba, ada banyak campuran yang muat. Yang mengejutkan, campuran paling
menguntungkan ternyata memakai setengah hiasan. Setengah hiasan tidak bisa dijual, dan di
sinilah akal sehat mulai menipu: kalau kamu membulatkan ke bawah, hasilnya bukan yang
terbaik. Ada campuran lain yang seluruhnya bulat dan justru lebih menguntungkan, padahal
tidak kelihatan seperti pembulatan apa pun.

Pertanyaan yang membuat penasaran begini. Kalau ada ribuan campuran yang mungkin, bagaimana
kita tahu satu campuran itu yang terbaik tanpa mencoba semuanya? Dan kalau temanmu
menawarkan satu lembar kertas tambahan, berapa harga wajar untuk lembar itu? Di bab ini
kita akan tahu mengapa jawaban terbaik selalu ada di pojok daerah yang mungkin, mengapa
setiap bahan punya harga yang bisa dihitung, dan mengapa jawaban pecahan harus dibulatkan
dengan hati-hati, bukan dengan naluri.
\end{untukanak}

\section{Dari rencana produksi ke bentuk baku}
\tingkat{dasar}

Satu rencana produksi adalah sepasang bilangan: banyak rak dan banyak bangku. Rencana
disebut layak bila pemakaian bahannya tidak melebihi persediaan. Untuk rak $x_{1}$ dan
bangku $x_{2}$, syaratnya
\[
  6x_{1}+4x_{2}\le24,\qquad x_{1}+2x_{2}\le6,\qquad x_{1}\ge0,\quad x_{2}\ge0 .
\]
Tiap syarat memotong bidang dengan satu garis dan menyisakan satu sisi; irisannya segiempat
dengan empat sudut. Labanya $5x_{1}+4x_{2}$, dan kita ingin angka itu sebesar mungkin.

Pola soal semacam ini disebut program linear. Namanya tepat: semua yang terlibat berbentuk
garis lurus, tidak ada lengkungan, tidak ada pangkat.

\kamus{program linear}{linear program}{Soal mencari nilai terbaik dari sesuatu yang dihitung
dengan penjumlahan berbobot, ketika semua aturannya berbentuk garis lurus.}

\begin{definisi}[Program linear]
Program linear adalah masalah memaksimumkan atau meminimumkan fungsi linear
$\mathbf{c}^{\top}\mathbf{x}$ terhadap variabel $\mathbf{x}\in\mathbb{R}^{n}$ dengan
kendala berbentuk ketaksamaan atau kesamaan linear. Himpunan titik yang memenuhi semua
kendala disebut \emph{daerah layak}. Bila daerah layaknya kosong, program disebut
\emph{tidak layak}; bila nilai sasarannya dapat dibuat sebesar apa pun, program disebut
\emph{tak terbatas}.\index{program linear}
\end{definisi}

Perangkat lunak dan buku teks memakai satu pola tulisan yang seragam, dan pola itulah yang
dipakai sepanjang bab ini.

\begin{definisi}[Bentuk baku]
Program linear dalam \emph{bentuk baku} adalah
\[
  \text{maksimumkan } \mathbf{c}^{\top}\mathbf{x}
  \quad\text{dengan}\quad A\mathbf{x}=\mathbf{b},\qquad \mathbf{x}\ge\mathbf{0},
\]
dengan $A$ berukuran $m\times n$ berperingkat $m$, $\mathbf{b}\in\mathbb{R}^{m}$, dan
$\mathbf{c}\in\mathbb{R}^{n}$. Setiap program linear dapat dibawa ke bentuk ini tanpa
mengubah nilai optimumnya.\index{bentuk baku}
\end{definisi}

Empat terjemahan melakukan seluruh pekerjaan itu. Kendala
$\mathbf{a}^{\top}\mathbf{x}\le b$ menjadi kesamaan dengan menambah variabel yang menampung
sisa pemakaian; kendala $\mathbf{a}^{\top}\mathbf{x}\ge b$ dengan mengurangkannya. Variabel
yang boleh bernilai negatif diganti selisih dua variabel tak negatif, dan memaksimumkan
$\mathbf{c}^{\top}\mathbf{x}$ sama dengan meminimumkan $-\mathbf{c}^{\top}\mathbf{x}$.

\kamus{variabel kendur}{slack variable}{Bilangan yang dicatat sebagai sisa pemakaian, supaya
``tidak boleh lebih dari'' berubah menjadi ``sama dengan''.}

Untuk bengkel tadi, bentuk bakunya memakai $s_{1}$ sebagai sisa kayu lapis dan $s_{2}$
sebagai sisa jam perakitan:
\[
  \begin{array}{rcl}
    6x_{1}+4x_{2}+s_{1} &=& 24,\\
    x_{1}+2x_{2}+s_{2} &=& 6,
  \end{array}
  \qquad x_{1},x_{2},s_{1},s_{2}\ge0 .
\]
Di sini $m=2$ dan $n=4$, sehingga $A$ berukuran $2\times4$. Kedua variabel kendur itu bukan
hiasan: nanti keduanya muncul sebagai kolom di tabel simpleks dan sebagai petunjuk kendala
mana yang masih longgar.

Bentuk baku juga menuntut $\mathbf{b}\ge\mathbf{0}$, dan itu selalu bisa diatur
dengan mengalikan satu baris dengan $-1$. Yang tidak bisa diatur adalah peringkat: baris $A$
yang bergantung linear harus dibuang lebih dulu, sebab kalau tidak, matriks basis yang dipilih
bisa singular.

\section{Geometri: mengapa jawaban ada di sudut}
\tingkat{menengah}

Empat sudut segiempat tadi adalah $(0;0)$, $(4;0)$, $(3;1{,}5)$, dan $(0;3)$, dengan laba
$0$, $20$, $21$, dan $12$. Yang terbaik ada di sudut, dan itu bukan kebetulan: pada program
linear nilai terbaik selalu dapat dicari di sudut, sebab itulah metode simpleks hanya
berjalan dari sudut ke sudut.

\begin{figure}[!ht]
\centering
\begin{tikzpicture}
\begin{axis}[width=9.3cm,height=6.3cm,
  xlabel={$x_{1}$ (rak)},ylabel={$x_{2}$ (bangku)},
  xmin=-0.35,xmax=5.1,ymin=-0.35,ymax=3.9,
  axis lines=left,xtick={0,1,2,3,4,5},ytick={0,1,2,3},
  tick label style={font=\scriptsize},label style={font=\scriptsize}]
\addplot[forget plot,fill=layerteal!10,draw=layerteal,thick]
  coordinates {(0,0) (4,0) (3,1.5) (0,3)} \closedcycle;
\addplot[forget plot,domain=0:2.4,densely dotted,thick,layerindigo] {3-1.25*x};
\addplot[forget plot,domain=0:3.6,densely dotted,thick,layerindigo] {4.5-1.25*x};
\addplot[forget plot,domain=0.2:4.2,thick,accent] {5.25-1.25*x};
\addplot[forget plot,only marks,mark=*,mark size=1.7pt,draw=ink]
  coordinates {(0,0) (4,0) (0,3)};
\addplot[forget plot,only marks,mark=*,mark size=2.6pt,draw=accent]
  coordinates {(3,1.5)};
\node[font=\scriptsize,anchor=north west] at (axis cs:0,0) {$(0;0)$};
\node[font=\scriptsize,anchor=north east] at (axis cs:4,0) {$(4;0)$};
\node[font=\scriptsize,anchor=south west] at (axis cs:0,3) {$(0;3)$};
\node[font=\scriptsize,anchor=west] at (axis cs:3.05,1.45) {$(3;1{,}5)$};
\node[font=\scriptsize,anchor=west,layerindigo] at (axis cs:0.15,1.15) {$z=12$};
\node[font=\scriptsize,anchor=west,layerindigo] at (axis cs:0.15,2.15) {$z=18$};
\node[font=\scriptsize,anchor=west,accent] at (axis cs:0.15,3.45) {$z=21$};
\draw[-{Stealth},thick,accent] (axis cs:2.3,2.75) -- (axis cs:3.25,3.35)
  node[font=\scriptsize,above,accent,pos=0.9] {laba naik};
\end{axis}
\end{tikzpicture}
\caption{Daerah layak bengkel dan tiga garis laba sama besar. Poligon punya empat verteks,
$(0;0)$, $(4;0)$, $(3;1{,}5)$, $(0;3)$, bernilai $0$, $20$, $21$, $12$. Garis $z=12$ dan
$z=18$ masih memotong daerah layak; garis $z=21$ hanya menyentuh verteks $(3;1{,}5)$,
sehingga tidak ada titik layak yang bernilai lebih tinggi.}
\label{fig:ch22-daerah-optimum}
\end{figure}

Gambar~\ref{fig:ch22-daerah-optimum} mempertajam kebetulan itu. Tiga garis di sana adalah
garis laba sama besar, $5x_{1}+4x_{2}=12$, $18$, dan $21$. Garis laba yang lebih tinggi
lebih jauh dari titik nol, dan menggesernya ke arah itu memindahkannya keluar daerah layak;
garis $21$ hanya menyentuh verteks $(3;1{,}5)$.

Sudut punya nama yang lebih tepat, dan namanya tidak mengacu pada gambar.

\begin{definisi}[Verteks]
Misalkan $P\subseteq\mathbb{R}^{n}$ himpunan konveks. Titik $\mathbf{x}\in P$ disebut
\emph{verteks} atau \emph{titik ekstrem} bila tidak ada $\mathbf{y},\mathbf{z}\in P$ dengan
$\mathbf{y}\ne\mathbf{z}$ dan $t\in(0,1)$ sehingga
$\mathbf{x}=t\mathbf{y}+(1-t)\mathbf{z}$. Dengan kata lain, $\mathbf{x}$ bukan titik
tengah ruas antara dua titik lain di $P$.\index{verteks}
\end{definisi}

Pada segiempat tadi, titik $(3;1{,}5)$ adalah verteks: setiap ruas di dalam segiempat yang
melewatinya harus berakhir di titik itu sendiri. Sebaliknya $(2;1)$ bukan verteks, sebab ia
titik tengah antara $(1;1)$ dan $(3;1)$.

\kamus{verteks}{vertex}{Sudut daerah layak, yaitu titik yang tidak pernah menjadi titik
tengah antara dua titik layak lain.}

Langkah pertama pembuktian adalah mengubah penyangkalan menjadi arah yang bisa dipakai.

\begin{lemma}[Titik bukan verteks memberi garis]
Bila $\mathbf{x}\in P$ bukan verteks, maka ada $\mathbf{d}\ne\mathbf{0}$ sehingga
$\mathbf{x}+\mathbf{d}\in P$ dan $\mathbf{x}-\mathbf{d}\in P$.
\end{lemma}

\begin{proof}
Karena $\mathbf{x}$ bukan verteks, ada $\mathbf{y}\ne\mathbf{z}\in P$ dan $t\in(0,1)$
dengan $\mathbf{x}=t\mathbf{y}+(1-t)\mathbf{z}$. Tetapkan
$s=\tfrac12\min\{t,1-t\}>0$ dan $\mathbf{d}=s(\mathbf{y}-\mathbf{z})$. Maka
$\mathbf{x}+\mathbf{d}=(t+s)\mathbf{y}+(1-t-s)\mathbf{z}$ dan
$\mathbf{x}-\mathbf{d}=(t-s)\mathbf{y}+(1-t+s)\mathbf{z}$. Karena $0<t\pm s<1$, kedua
titik itu kombinasi konveks $\mathbf{y}$ dan $\mathbf{z}$, jadi anggota $P$ sebab $P$
konveks. Juga $\mathbf{d}\ne\mathbf{0}$, sebab $\mathbf{y}\ne\mathbf{z}$ dan $s>0$.
\end{proof}

Sekarang pernyataan yang menjadi alasan metode simpleks masuk akal. Kalimatnya dalam
bahasa sehari-hari: kalau daerah layaknya terkurung dan sasarannya tidak tumbuh ke arah
mana pun, maka puncak tertinggi ada di salah satu sudut, bukan di tengah bidang.

\begin{teorema}[Optimum tercapai di verteks]
Misalkan $P=\{\mathbf{x}:A\mathbf{x}\le\mathbf{b}\}$ tak kosong dan terbatas, dan misalkan
$\mathbf{c}^{\top}\mathbf{x}$ terbatas di atas pada $P$. Maka ada verteks
$\mathbf{x}^{*}$ dari $P$ dengan $\mathbf{c}^{\top}\mathbf{x}^{*}
=\max_{\mathbf{x}\in P}\mathbf{c}^{\top}\mathbf{x}$.
\index{optimum di verteks}
\end{teorema}

\begin{proof}
Himpunan optimal $F=\{\mathbf{x}\in P:\mathbf{c}^{\top}\mathbf{x}=p^{*}\}$ tak kosong,
tertutup, terbatas, dan konveks, jadi $F$ juga polytope. Tulis kendala $P$ sebagai
$\mathbf{a}_{i}^{\top}\mathbf{x}\le b_{i}$, $i=1,\dots,k$. Pilih $\mathbf{x}^{*}\in F$
yang memaksimumkan banyaknya kendala yang aktif, yaitu kendala dengan
$\mathbf{a}_{i}^{\top}\mathbf{x}^{*}=b_{i}$.

Andaikan $\mathbf{x}^{*}$ bukan verteks. Menurut lemma, ada $\mathbf{d}\ne\mathbf{0}$
dengan $\mathbf{x}^{*}\pm\mathbf{d}\in F$. Setiap kendala yang aktif di
$\mathbf{x}^{*}$ juga aktif di kedua titik itu, sehingga
$\mathbf{a}_{i}^{\top}\mathbf{d}=0$ bagi semua $i$ aktif. Tinjau garis
$\mathbf{x}^{*}+s\mathbf{d}$ untuk $s\in\mathbb{R}$. Irisannya dengan $P$ adalah selang
tertutup $[s_{-},s_{+}]$ yang memuat $[-1,1]$, dan karena $P$ terbatas kedua ujungnya
berhingga. Ujung $s_{+}$ ditentukan oleh kendala dengan $\mathbf{a}_{i}^{\top}\mathbf{d}>0$ yang
paling cepat tercapai; kendala itu aktif di $s_{+}$ tetapi tidak di $\mathbf{x}^{*}$, jadi
titik baru itu layak dan aktif di lebih banyak kendala. Nilainya tetap $p^{*}$ sebab
$\mathbf{c}^{\top}(\mathbf{x}^{*}+s_{+}\mathbf{d})=p^{*}+s_{+}\mathbf{c}^{\top}
\mathbf{d}=p^{*}$, dan $\mathbf{c}^{\top}\mathbf{d}=0$ karena kedua titik
$\mathbf{x}^{*}\pm\mathbf{d}$ bernilai $p^{*}$. Titik baru itu anggota $F$ dengan lebih banyak kendala
aktif, bertentangan dengan pilihan $\mathbf{x}^{*}$; maka $\mathbf{x}^{*}$ verteks.
\end{proof}

\begin{bacakeras}
Dua suara, dibaca bergantian dengan teman. Ketuk meja sekali setiap kali kata ``puncak''
diucapkan.

``Di mana jawaban terbaik?'' ``Di puncak.''
``Kenapa bukan di tengah?'' ``Karena di tengah ada tetangga di dua sisi.''
``Kalau ada tetangga, kenapa salah?'' ``Karena masih ada arah untuk bergeser.''
``Bergeser ke arah mana?'' ``Ke arah yang menanjak.''
``Kalau tidak ada arah menanjak?'' ``Maka kita sudah berdiri di puncak.''

Ulangi tiga baris tengah, makin pelan, dengan ketukan yang sama: di tengah ada tetangga, di
tengah masih bisa bergeser, puncak tidak punya arah menanjak.

Baris yang harus diulang pembaca, sekali, dengan suara jelas: jawaban terbaik ada di puncak,
bukan di tengah, dan puncak dikenali dari tidak adanya arah yang menanjak.
\end{bacakeras}

Bukti itu memakai keterbatasan $P$. Bila $P$ tak terbatas tetapi sasarannya
terbatas di atas, kesimpulannya masih benar; jalannya lewat dekomposisi $P$ menjadi selubung
konveks verteks ditambah kerucut arah resesi, dan bukti lengkapnya ada pada
\texttt{WSPrinted.pdf}, bagian struktur himpunan polihedral.

Banyaknya verteks berhingga, tetapi tidak selalu sedikit. Dalam bentuk baku,
verteks bersesuaian dengan pemilihan $m$ kolom $A$ yang bebas linear: kolom terpilih
menentukan $\mathbf{x}_{B}=B^{-1}\mathbf{b}$, kolom lain nol, dan titik itu layak bila
$\mathbf{x}_{B}\ge\mathbf{0}$. Calonnya sebanyak $\binom{n}{m}$, yang untuk $n=100$ dan
$m=50$ melampaui $10^{29}$; memeriksa semuanya bukan rencana masuk akal. Di sinilah simpleks
berperan: ia berjalan dari verteks ke verteks tetangga yang lebih baik.

\begin{tangan}
Bahan: lantai berpetak, dua utas tali atau selotip kertas, dan segenggam kancing.

Bentangkan dua tali lurus sebagai dua batas; satu untuk kayu lapis, satu untuk jam
perakitan. Berdirilah di petak pojok dan sebut petak itu titik nol: setiap petak ke kanan
menambah satu rak, setiap petak ke depan menambah satu bangku.

Berjalanlah menyusuri tali dan berhenti hanya di petak tempat dua tali bertemu. Letakkan
kancing di situ, dan di situ saja. Yang harus terlihat: semua kancing menempel di tikungan
pagar, tidak satu pun berdiri di tengah halaman. Sesudah itu hitung laba di setiap kancing
dengan jari, lalu pindahkan satu kancing ke tikungan yang paling menguntungkan.

Geser satu tali satu petak ke dalam, lalu ulangi: kancingnya berpindah, pemenangnya tetap
kancing di tikungan.
\end{tangan}

\section{Metode simpleks: satu tabel, dua pivot}
\tingkat{menengah}

Gagasannya sederhana. Berdiri di satu verteks, periksa apakah ada arah sisi yang menaikkan
laba; bila ada, berjalan ke verteks di ujungnya. Bila tidak, berhenti: tidak ada tetangga
yang lebih baik, yang pada program linear berarti tidak ada titik lain yang lebih baik.

Arah sisi itu ditentukan oleh \emph{biaya tereduksi}: seberapa banyak laba bertambah bila
satu variabel yang sekarang nol dinaikkan satu satuan, dengan mengorbankan variabel basis
secara optimal.

\kamus{solusi layak basis}{basic feasible solution}{Titik yang diperoleh dengan menyetel
sebagian variabel ke nol lalu menyelesaikan sisanya, dan kebetulan masih memenuhi semua
batas; titik ini selalu verteks.}

Untuk bengkel tadi, mulailah dari verteks $(0;0)$, yaitu $s_{1}=24$ dan $s_{2}=6$. Tabel
pertamanya memuat baris laba $z$ dan satu baris untuk tiap kendala:
\begin{center}
\begin{tabular}{c|cccc|c}
\hline
basis & $x_{1}$ & $x_{2}$ & $s_{1}$ & $s_{2}$ & ruas kanan \\
\hline
$z$   & $-5$ & $-4$ & $0$ & $0$ & $0$ \\
$s_{1}$ & $6$ & $4$ & $1$ & $0$ & $24$ \\
$s_{2}$ & $1$ & $2$ & $0$ & $1$ & $6$ \\
\hline
\end{tabular}
\end{center}
Baris $z$ ditulis sebagai $z-5x_{1}-4x_{2}=0$; koefisien negatif berarti menaikkan variabel
itu menaikkan $z$, jadi kolom dengan koefisien paling negatif masuk, di sini $x_{1}$ dengan
$-5$.

Uji nisbah menentukan siapa yang keluar: bagi ruas kanan dengan entri positif di kolom
masuk, lalu ambil nisbah terkecil. $24/6=4$ dan $6/1=6$, jadi $s_{1}$ keluar dan pivotnya
$6$. Membagi baris pivot dengan $6$ lalu menghapus kolom itu dari baris lain memberi
\[
  \begin{array}{rcl}
    x_{1}+\tfrac{2}{3}x_{2}+\tfrac{1}{6}s_{1} &=& 4,\\
    \tfrac{4}{3}x_{2}-\tfrac{1}{6}s_{1}+s_{2} &=& 2,\\
    z-\tfrac{2}{3}x_{2}+\tfrac{5}{6}s_{1} &=& 20 .
  \end{array}
\]
Titiknya kini $(4;0)$ berlaba $20$, dan masih ada koefisien negatif: $x_{2}$ dengan
$-2/3$. Uji nisbah memberi $4/(2/3)=6$ dan $2/(4/3)=1{,}5$, jadi $s_{2}$ keluar dan
pivotnya $4/3$. Sesudah satu putaran lagi,
\[
  \begin{array}{rcl}
    x_{1}+\tfrac{1}{4}s_{1}-\tfrac{1}{2}s_{2} &=& 3,\\
    x_{2}-\tfrac{1}{8}s_{1}+\tfrac{3}{4}s_{2} &=& \tfrac{3}{2},\\
    z+\tfrac{3}{4}s_{1}+\tfrac{1}{2}s_{2} &=& 21 .
  \end{array}
\]
Baris $z$ sekarang tidak punya koefisien negatif, jadi tidak ada arah yang menaikkan laba:
$(3;1{,}5)$ optimal dengan laba $21$. Seluruh aritmetika di atas diperiksa ulang dengan pecahan eksak, supaya tidak ada
pembulatan desimal yang menyelinap.

Baris terakhir itu lebih dari sekadar tanda berhenti. Koefisien $s_{1}$ dan
$s_{2}$ di sana, $3/4$ dan $1/2$, adalah harga dua bahan tadi. Angka itu muncul dari matriks
basis, $\mathbf{y}^{\top}=\mathbf{c}_{B}^{\top}B^{-1}$, dan
$\mathbf{b}^{\top}\mathbf{y}=24\cdot\tfrac34+6\cdot\tfrac12=21$, sama dengan laba.
Kesamaan itu adalah dualitas, dan bagian berikutnya menjelaskan mengapa ia sendiri sudah
merupakan bukti optimalitas.

Dua catatan. Aturan koefisien paling negatif tidak menjamin berhenti, sebab ada contoh kecil
dengan putaran tanpa akhir; aturan Bland yang memilih indeks terkecil menjamin berhenti.
Perangkat lunak modern memakai simpleks dual atau metode titik dalam, tetapi melaporkan
verteks dan harga yang sama.

\section{Dualitas: harga untuk setiap kendala}
\tingkat{lanjut}

Setiap bahan yang terbatas punya harga tersembunyi. Satu lembar kayu lapis tambahan
menaikkan laba, dan besar kenaikannya adalah harga lembar itu bagi bengkel. Harga itu tidak
tertulis pada soal aslinya, tetapi bisa dihitung, dan dualitas adalah cara menghitungnya.

Kalikan kendala kayu lapis dengan $y_{1}\ge0$ dan kendala jam dengan $y_{2}\ge0$, lalu
jumlahkan. Hasilnya
\[
  (6y_{1}+y_{2})x_{1}+(4y_{1}+2y_{2})x_{2}\le24y_{1}+6y_{2}.
\]
Bila $6y_{1}+y_{2}\ge5$ dan $4y_{1}+2y_{2}\ge4$, ruas kiri paling sedikit
$5x_{1}+4x_{2}$, yaitu laba. Jadi $24y_{1}+6y_{2}$ menjadi batas atas laba bagi setiap harga yang memenuhi dua syarat
itu, dan batas terbaiknya dibuat sekecil mungkin.

\begin{definisi}[Pasangan primal dan dual]
Untuk program linear
\[
  \text{(P)}\quad \text{maksimumkan } \mathbf{c}^{\top}\mathbf{x}
  \ \text{ dengan } A\mathbf{x}\le\mathbf{b},\ \mathbf{x}\ge\mathbf{0},
\]
\emph{program dualnya} adalah
\[
  \text{(D)}\quad \text{minimumkan } \mathbf{b}^{\top}\mathbf{y}
  \ \text{ dengan } A^{\top}\mathbf{y}\ge\mathbf{c},\ \mathbf{y}\ge\mathbf{0}.
\]
Variabel $\mathbf{y}$ disebut \emph{variabel dual} atau \emph{harga bayangan}.
\index{dualitas}\index{harga bayangan}
\end{definisi}

Untuk bengkel tadi, dualnya meminimumkan $24y_{1}+6y_{2}$ dengan kendala
$6y_{1}+y_{2}\ge5$, $4y_{1}+2y_{2}\ge4$, dan $y_{1},y_{2}\ge0$. Penyelesaiannya
$y_{1}=3/4$, $y_{2}=1/2$, bernilai $21$, sama dengan laba terbaik.

Bahwa batas atas dari dual tidak pernah di bawah nilai primal adalah hal pertama yang perlu
dipegang, dan buktinya tiga baris.

\begin{teorema}[Dualitas lemah]
Untuk setiap $\mathbf{x}$ layak (P) dan setiap $\mathbf{y}$ layak (D) berlaku
$\mathbf{c}^{\top}\mathbf{x}\le\mathbf{b}^{\top}\mathbf{y}$.
\index{dualitas lemah}
\end{teorema}

\begin{proof}
Karena $\mathbf{y}\ge\mathbf{0}$ dan $\mathbf{b}-A\mathbf{x}\ge\mathbf{0}$, berlaku
$\mathbf{b}^{\top}\mathbf{y}\ge\mathbf{y}^{\top}A\mathbf{x}$. Karena
$\mathbf{x}\ge\mathbf{0}$ dan $A^{\top}\mathbf{y}-\mathbf{c}\ge\mathbf{0}$, berlaku
$\mathbf{y}^{\top}A\mathbf{x}\ge\mathbf{c}^{\top}\mathbf{x}$. Rantai keduanya memberi
hasil yang diminta.
\end{proof}

Akibat praktisnya besar: sepasang titik layak, satu primal dan satu dual, dengan nilai sama
sudah membuktikan keduanya optimal tanpa memeriksa titik lain. Pada bengkel tadi,
$(3;1{,}5)$ dan $(0{,}75;\,0{,}5)$ memberi nilai $21$, jadi $21$ tidak bisa diperbaiki
maupun diturunkan.

Yang tidak diberikan dualitas lemah adalah jaminan bahwa kedua nilai itu selalu bertemu.
Jaminan itu isi teorema berikut.

\begin{teorema}[Dualitas kuat]
Bila (P) atau (D) mempunyai nilai optimum berhingga, maka keduanya mempunyai nilai optimum
dan nilainya sama.
\end{teorema}

Buktinya tidak dikerjakan di sini. Jalannya ada dua: lewat lemma Farkas tentang sistem
ketaksamaan linear, atau lewat kenyataan bahwa simpleks berhenti pada tabel tanpa koefisien
negatif, dan tabel itu sekaligus memberi $\mathbf{y}$ layak dual bernilai sama. Bukti
lengkapnya ada pada \texttt{WSPrinted.pdf} bagian dualitas program linear dan pada Boyd dan
Vandenberghe bagian 5.2.

Dari kesamaan itu lahir syarat yang dipakai untuk memeriksa hasil hitungan dengan cepat.
Kalimat sehari-harinya: kendala yang masih longgar tidak perlu dibayar, dan variabel yang
dipakai harus membayar harganya dengan pas.

\begin{akibat}[Kelengkapan komplementer]
Bila $\mathbf{x}$ dan $\mathbf{y}$ keduanya optimal, maka
\[
  y_{i}\bigl(b_{i}-\mathbf{a}_{i}^{\top}\mathbf{x}\bigr)=0 \ \text{ untuk setiap } i,
  \qquad
  x_{j}\bigl((A^{\top}\mathbf{y})_{j}-c_{j}\bigr)=0 \ \text{ untuk setiap } j .
\]
Sebaliknya, bila $\mathbf{x}$ layak (P), $\mathbf{y}$ layak (D), dan kedua kelompok
kesamaan itu berlaku, maka keduanya optimal.
\index{kelengkapan komplementer}
\end{akibat}

\begin{proof}
Karena keduanya optimal, $\mathbf{c}^{\top}\mathbf{x}=\mathbf{b}^{\top}\mathbf{y}$, dan
rantai pada bukti dualitas lemah menjadi rantai kesamaan. Setiap suku yang dijumlahkan di
dalamnya tak negatif, jadi setiap suku nol; suku-suku itu tepat
$y_{i}(b_{i}-\mathbf{a}_{i}^{\top}\mathbf{x})$ dan
$x_{j}((A^{\top}\mathbf{y})_{j}-c_{j})$. Arah sebaliknya: bila semua suku nol, rantai yang
sama memberi kesamaan nilai, dan menurut dualitas lemah nilai itu optimum kedua masalah.
\end{proof}

Sekarang harga bayangan itu sendiri. Bila ruas kanan diganggu, nilai optimum
berubah, dan laju perubahannya adalah komponen $\mathbf{y}$. Menambah satu lembar kayu lapis
menaikkan laba dari $21$ menjadi $21{,}75$; menambah satu jam perakitan menaikkannya menjadi
$21{,}5$. Jadi selembar kayu lapis bernilai $0{,}75$ satuan laba, satu jam perakitan
$0{,}5$.

Harga itu tidak berlaku selamanya, dan batas berlakunya bisa dihitung. Selama susunan
kendala aktif tidak berubah, $\mathbf{x}_{B}=B^{-1}\mathbf{b}$ linear dalam $\mathbf{b}$,
sehingga nilainya juga linear. Untuk bengkel tadi,
\[
  B^{-1}=\begin{pmatrix}\tfrac14&-\tfrac12\\[2pt]-\tfrac18&\tfrac34\end{pmatrix},
  \qquad
  \mathbf{x}_{B}=\begin{pmatrix}\tfrac14b_{1}-\tfrac12b_{2}\\[2pt]
  -\tfrac18b_{1}+\tfrac34b_{2}\end{pmatrix}.
\]
Dengan $b_{2}=6$ tetap, syarat $\mathbf{x}_{B}\ge\mathbf{0}$ memberi $b_{1}\in[12,36]$; di
luar selang itu satu variabel basis menjadi negatif dan kemiringannya berganti. Dengan $b_{1}=24$ tetap, selangnya $b_{2}\in[4,12]$; harga bayangan karena itu selalu
dilaporkan bersama selangnya, bukan sebagai angka telanjang.

Satu peringatan berguna sejak awal: harga bayangan adalah gagasan model bersambung. Pada
model integer, satu lembar kayu lapis tambahan bisa menaikkan laba sebesar $1$ atau tidak
sama sekali, dan tidak ada harga tunggal yang merangkum keduanya.

\section{Relaksasi LP: batas yang sah, pembulatan yang menipu}
\tingkat{menengah}

Kalau rak dan bangku boleh pecahan, soal tadi selesai: $3$ rak dan $1{,}5$ bangku.
Kenyataannya tidak, dan hubungan antara soal bersambung dan soal bulat perlu dinyatakan
dengan hati-hati.

\begin{definisi}[Program integer, relaksasi, dan celah integralitas]
\emph{Program integer} adalah program linear yang sebagian atau seluruh variabelnya
dituntut bulat. Bila seluruh variabel dituntut bulat, ia disebut program integer murni.
\emph{Relaksasi LP}-nya adalah program linear yang sama tanpa tuntutan kebulatan. Untuk
masalah maksimisasi dengan nilai optimum $z^{*}_{\mathrm{LP}}$ dan $z^{*}_{\mathrm{IP}}$,
\emph{celah integralitas} adalah $z^{*}_{\mathrm{LP}}-z^{*}_{\mathrm{IP}}$.
\index{program integer}\index{relaksasi}\index{celah integralitas}
\end{definisi}

\kamus{celah integralitas}{integrality gap}{Jarak antara batas yang dijanjikan model
bersambung dan jawaban bulat sebenarnya; jarak itu mengukur seberapa jauh model bersambung
menyesatkan.}

Alasan relaksasi berguna adalah satu kalimat: himpunan titik bulat yang layak adalah
bagian dari himpunan titik bersambung yang layak, jadi memaksimumkan di himpunan yang lebih
besar tidak mungkin memberi nilai yang lebih kecil.

\begin{proposisi}[Relaksasi memberi batas yang sah]
Untuk masalah maksimisasi, $z^{*}_{\mathrm{IP}}\le z^{*}_{\mathrm{LP}}$. Untuk masalah
minimisasi, ketaksamaannya berbalik.
\index{batas relaksasi}
\end{proposisi}

\begin{proof}
Setiap titik bulat yang layak bagi program integer juga layak bagi relaksasinya, sebab
relaksasi hanya menghapus syarat, tidak menambah syarat. Jadi daerah layak program integer
termasuk di dalam daerah layak relaksasi. Maksimum atas himpunan yang lebih besar tidak
lebih kecil daripada maksimum atas himpunan bagiannya.
\end{proof}

Batas itu sah, tetapi tidak selalu rapat, dan selisihnya bisa sangat menyesatkan bila
dipakai untuk membulatkan. Ada dua cara pembulatan gagal, dan keduanya muncul pada contoh
kecil kita.

Cara pertama: hasil pembulatan tidak layak. Titik $(3;1{,}5)$ dibulatkan ke atas menjadi
$(3;2)$. Labanya $23$, lebih tinggi daripada $21$, tetapi kayu lapisnya
$6\cdot3+4\cdot2=26$ lembar, melebihi $24$: itu bukan rencana produksi, hanya deretan
angka.

Cara kedua, lebih halus: hasil pembulatan layak tetapi bukan yang terbaik. Pembulatan ke
bawah memberi $(3;1)$ berlaba $19$, padahal $(4;0)$ berlaba $20$ dan seluruhnya layak.
Membulatkan koordinat satu demi satu ke arah terdekat karena itu memberi jawaban lebih buruk
daripada titik bulat lain yang jaraknya lebih jauh dari $(3;1{,}5)$.

Kedua kegagalan itu akibat bentuk daerah layaknya. Himpunan titik bulat layak
tidak konveks: $(3;1)$ dan $(4;0)$ anggotanya, tetapi titik tengahnya $(3{,}5;\,0{,}5)$
bukan, sebab tidak bulat. Pada himpunan tak konveks, teorema optimum di verteks tidak
berlaku: tidak ada verteks yang dijamin memuat jawaban, dan tidak ada tetangga yang bisa
diandalkan.

\begin{figure}[!ht]
\centering
\begin{tikzpicture}
\begin{axis}[width=9.3cm,height=6.3cm,
  xlabel={$x_{1}$ (rak)},ylabel={$x_{2}$ (bangku)},
  xmin=-0.35,xmax=5.1,ymin=-0.35,ymax=3.9,
  axis lines=left,xtick={0,1,2,3,4,5},ytick={0,1,2,3},
  tick label style={font=\scriptsize},label style={font=\scriptsize}]
\addplot[forget plot,fill=layerteal!8,draw=layerteal,thick]
  coordinates {(0,0) (4,0) (3,1.5) (0,3)} \closedcycle;
\addplot[forget plot,only marks,mark=*,mark size=1.5pt,draw=layerindigo]
  coordinates {(0,0) (1,0) (2,0) (3,0) (4,0)
               (0,1) (1,1) (2,1) (3,1)
               (0,2) (1,2) (2,2) (0,3)};
\addplot[forget plot,only marks,mark=square*,mark size=2.4pt,draw=layerjade]
  coordinates {(4,0)};
\addplot[forget plot,only marks,mark=*,mark size=1.5pt,draw=layerplum]
  coordinates {(3,1)};
\addplot[forget plot,only marks,mark=o,mark size=3pt,thick,accent]
  coordinates {(3,1.5)};
\draw[densely dashed,layergrey] (axis cs:3,1) -- (axis cs:3,2.35);
\node[font=\scriptsize,anchor=west,layerjade] at (axis cs:4.1,-0.05) {$(4;0)$, $z=20$};
\node[font=\scriptsize,anchor=east,layerplum] at (axis cs:2.9,0.85) {$(3;1)$, $z=19$};
\node[font=\scriptsize,anchor=west,accent] at (axis cs:3.15,1.6) {$(3;1{,}5)$, $z=21$};
\node[font=\scriptsize,anchor=north,layergrey] at (axis cs:3,-0.05) {};
\end{axis}
\end{tikzpicture}
\caption{Kisi titik bulat versus titik relaksasi. Ada $13$ titik bulat layak, ditandai titik
penuh. Relaksasi memberi $(3;1{,}5)$ bernilai $21$, yang berdiri di tengah antara $(3;1)$
bernilai $19$ dan $(3;2)$ yang tidak layak sebab $6\cdot3+4\cdot2=26>24$. Jawaban bulat
terbaik $(4;0)$ bernilai $20$ justru berjarak lebih jauh dari titik relaksasi.}
\label{fig:ch22-kisi-relaksasi}
\end{figure}

Gambar~\ref{fig:ch22-kisi-relaksasi} memuat seluruhnya: tiga belas titik bulat layak, dan
satu titik relaksasi yang berdiri di antara dua titik. Yang di antara itu bukan titik bulat,
dan yang terdekat belum tentu yang terbaik.

Pembulatan juga tidak punya arti tunggal di dimensi banyak. Dengan dua variabel,
``titik bulat terdekat'' masih bisa dibayangkan; dengan seratus variabel, yang terdekat belum
tentu layak, dan yang layak belum tentu terbaik. Praktiknya bukan membulatkan, melainkan
memakai batas relaksasi untuk memangkas pencarian.

\begin{tantangan}
Maksimumkan $6x_{1}+5x_{2}$ dengan $2x_{1}+3x_{2}\le12$, $4x_{1}+x_{2}\le13$, dan
$x_{1},x_{2}\ge0$ bulat. Relaksasinya memberi $(2{,}7;\,2{,}2)$ bernilai $27{,}2$.
Pembulatan ke bilangan bulat terdekat, $(3;2)$, bernilai $28$: lebih tinggi daripada nilai
relaksasi mana pun yang masih layak. Tunjukkan mengapa $28$ tidak boleh dipercaya, lalu
tentukan jawaban bulat yang benar beserta nilainya.

Petunjuk 1: periksa kelayakan, bukan hanya nilai. Angka yang lebih tinggi bisa datang dari
titik yang bukan rencana.

Petunjuk 2: jepit $x_{1}$. Bilangan bulat terbesar yang masih layak adalah $3$; untuk
$x_{1}=3$, hitung batas $x_{2}$ dari kedua kendala dan ambil yang lebih kecil.

Kalau mentok, kembali ke bagian Relaksasi LP: batas yang sah, pembulatan yang menipu.
\end{tantangan}

\section{Mengapa program integer lebih sulit}
\tingkat{menengah}

Sampai di sini pertanyaannya wajar: kalau relaksasi hanya menambah satu syarat, mengapa
program integer tidak diselesaikan dengan menyelesaikan relaksasinya lalu membulatkan?
Jawaban singkatnya, kesulitannya bukan pada mencari jawaban, melainkan pada memastikan
bahwa jawaban itu yang terbaik.

Tiga perbedaan menjelaskannya.

Pertama, tidak ada sertifikat optimalitas yang sederhana. Pada program linear,
sepasang $\mathbf{x}$ dan $\mathbf{y}$ bernilai sama adalah bukti selesai yang bisa diperiksa
dalam waktu sebanding ukuran data. Pada program integer, harga dari relaksasi hanya memberi
batas, bukan titik bulat yang mencapainya: dual dari relaksasi memberi batas bawah, bukan
jawaban.

Kedua, selisih antara batas dan jawaban bisa besar. Pada selimut verteks graf
lengkap $n$ simpul, relaksasinya membolehkan $x_{v}=1/2$ untuk semua $v$ dengan nilai $n/2$,
sedangkan jawaban bulatnya $n-1$; nisbahnya $2(n-1)/n$ mendekati $2$. Perhitungan lengkapnya
ada pada Laboratorium komputasi dan pada Latihan 10.

Ketiga, biaya komputasinya berbeda kelas, dan jumlah kandidatnya tumbuh eksponensial. Program
linear selesai dalam waktu polinomial terhadap ukuran masukan; hasil Khachiyan lewat metode
elipsoid dan Karmarkar lewat metode titik dalam adalah dua tonggaknya. Versi keputusan
program integer, ``apakah ada titik bulat layak bernilai paling sedikit $K$'', termasuk
NP-lengkap, dan daftar dua puluh satu masalah Karp memuat knapsack, yang sendirinya punya
$2^{n}$ himpunan bagian. Tidak ada algoritme polinomial yang diketahui untuk program integer,
dan bila ada, banyak hal lain akan runtuh bersamanya.

\begin{untukanak}
Pernahkah kamu memasang ubin di lantai? Kamu mengukur panjang ruangan, lalu membagi dengan
ukuran ubin. Hasilnya jarang bilangan bulat. Kalau kamu bulatkan ke bawah, ada celah di
ujung lantai. Kalau kamu bulatkan ke atas, kamu memotong ubin dan membuang sisanya. Dan
kadang jawaban terbaik bukan salah satu dari keduanya, melainkan pola pemasangan yang sama
sekali lain: mulai dari tengah, atau memakai ubin potongan di kedua sisi.

Orang yang hanya melihat satu angka hasil pembagian tidak akan menemukan pola
itu. Ia butuh dua batas: paling banyak berapa ubin terpakai, dan paling sedikit berapa.
Dengan dua angka itu ia berhenti memikirkan pola yang jelas terlalu boros, dan memusatkan
perhatian pada yang masih mungkin menang.

Di bab ini kita akan tahu mengapa batas dari model bersambung selalu sah, mengapa
pembulatan bisa memberi jawaban yang layak tetapi bukan yang terbaik, dan bagaimana
komputer memakai batas itu untuk memangkas jutaan kemungkinan tanpa memeriksanya satu per
satu.
\end{untukanak}

\section{Percabangan dan pembatasan}
\tingkat{lanjut}

Kalau jawaban relaksasi pecahan, kita tidak membulatkannya, melainkan memecah soalnya
menjadi dua dan bertanya lagi dengan relaksasi pada setiap potongan. Setiap jawaban bulat
memenuhi salah satu dari dua kemungkinan yang saling meniadakan, $x_{2}\le1$ atau
$x_{2}\ge2$, jadi menyelesaikan keduanya sama dengan menyelesaikan soal aslinya.

Alurnya empat langkah. Selesaikan relaksasi pada satu simpul. Bila hasilnya
bulat, catat sebagai calon jawaban dan jangan lanjutkan simpul itu. Bila pecahan, pilih satu
variabel yang pecah, $x_{j}=v$, lalu buat dua anak: $x_{j}\le\lfloor v\rfloor$ dan
$x_{j}\ge\lceil v\rceil$. Sebelum memeriksa sebuah simpul, bandingkan batas relaksasinya
dengan calon terbaik; bila tidak lebih baik, buang simpul itu beserta seluruh keturunannya.
Berhenti bila tidak ada simpul tersisa.

Aturan pembuangan itu sah, dan alasannya satu baris.

\begin{proposisi}[Pemangkasan tidak membuang jawaban]
Misalkan $z_{\mathrm{LP}}(N)$ adalah nilai relaksasi pada simpul $N$ dan $z_{\mathrm{best}}$
calon terbaik yang sudah tercatat. Bila $z_{\mathrm{LP}}(N)\le z_{\mathrm{best}}$, maka
tidak ada titik bulat layak di $N$ yang lebih baik daripada $z_{\mathrm{best}}$.
\end{proposisi}

\begin{proof}
Setiap titik bulat layak di $N$ juga layak bagi relaksasi $N$, jadi nilainya paling banyak
$z_{\mathrm{LP}}(N)$ menurut proposisi batas relaksasi. Bila
$z_{\mathrm{LP}}(N)\le z_{\mathrm{best}}$, semua titik itu bernilai paling banyak
$z_{\mathrm{best}}$, sehingga tidak ada yang perlu diperiksa.
\end{proof}

Pohon untuk bengkel tadi kecil, dan seluruhnya bisa ditulis sebagai tabel.
\begin{center}
\begin{tabular}{llcc}
\hline
simpul & tambahan kendala & $\mathbf{x}$ relaksasi & $z_{\mathrm{LP}}$ \\
\hline
akar & -- & $(3;\,1{,}5)$ & $21$ \\
kiri & $x_{2}\le1$ & $(10/3;\,1)$ & $20{,}\overline{6}$ \\
kanan & $x_{2}\ge2$ & $(2;\,2)$ & $18$ \\
kiri-kiri & $x_{2}\le1,\ x_{1}\le3$ & $(3;\,1)$ & $19$ \\
kiri-kanan & $x_{2}\le1,\ x_{1}\ge4$ & $(4;\,0)$ & $20$ \\
\hline
\end{tabular}
\end{center}
Akar memberi $21$ dengan $x_{2}$ pecah. Percabangan pada $x_{2}$ memberi dua anak: kiri
$20{,}\overline{6}$ dengan $x_{1}=10/3$ masih pecah, kanan $18$ dan sudah bulat. Simpul kanan
dicatat sebagai calon. Pada simpul kiri, percabangan pada $x_{1}$ memberi $(3;1)$ bernilai
$19$ dan $(4;0)$ bernilai $20$; yang kedua memperbarui calon. Tidak ada simpul tersisa, dan
jawabannya $20$.

\begin{figure}[!ht]
\centering
\begin{tikzpicture}[
  simpul/.style={draw=layergrey,fill=layergrey!6,rounded corners=1.5pt,
    inner sep=3pt,font=\scriptsize,align=center,text width=3.0cm},
  pangkas/.style={draw=accent,fill=accent!7,rounded corners=1.5pt,
    inner sep=3pt,font=\scriptsize,align=center,text width=3.0cm},
  panah/.style={-{Stealth},thick,draw=layergrey}]
\node[simpul] (akar) at (0,0) {akar\\ $z_{\mathrm{LP}}=21$\\ $(3;\,1{,}5)$};
\node[simpul] (kiri) at (-3.6,-1.75) {$x_{2}\le1$\\ $z_{\mathrm{LP}}=20{,}\overline{6}$\\
  $(10/3;\,1)$};
\node[pangkas] (kanan) at (3.6,-1.75) {$x_{2}\ge2$\\ $z_{\mathrm{LP}}=18$\\ $(2;\,2)$ bulat};
\node[pangkas] (kk) at (-5.1,-3.6) {$x_{1}\le3$\\ $z_{\mathrm{LP}}=19$\\ $(3;\,1)$ bulat};
\node[simpul] (kn) at (-2.1,-3.6) {$x_{1}\ge4$\\ $z_{\mathrm{LP}}=20$\\ $(4;\,0)$ bulat};
\draw[panah] (akar.south) -- (kiri.north);
\draw[panah] (akar.south) -- (kanan.north);
\draw[panah] (kiri.south) -- (kk.north);
\draw[panah] (kiri.south) -- (kn.north);
\node[font=\scriptsize,accent,anchor=north] at (3.6,-2.55) {dipangkas: $18<20$};
\node[font=\scriptsize,accent,anchor=north] at (-5.1,-4.45) {dipangkas: $19<20$};
\end{tikzpicture}
\caption{Pohon percabangan dan pembatasan untuk bengkel: lima simpul. Akar memberi $z=21$
dengan $x_{2}$ pecah. Cabang $x_{2}\ge2$ berakhir di $(2;\,2)$ bernilai $18$; cabang
$x_{2}\le1$ bercabang lagi menjadi $(3;\,1)$ bernilai $19$ dan $(4;\,0)$ bernilai $20$.
Kedua simpul berlabel bulat yang bernilai di bawah $20$ berbingkai merah: batasnya sudah
tidak lebih baik daripada calon terbaik, jadi keduanya dipangkas tanpa diperiksa lagi.}
\label{fig:ch22-pohon-percabangan}
\end{figure}

Gambar~\ref{fig:ch22-pohon-percabangan} menggambar pohon yang sama. Dua simpul di sana
berbingkai merah, yaitu $(2;\,2)$ bernilai $18$ dan $(3;\,1)$ bernilai $19$: keduanya bulat,
tetapi batasnya di bawah calon terbaik $20$, jadi keduanya dibuang.

Simpul kanan yang bernilai $18$ tidak perlu diperiksa lebih jauh seandainya
ia pecahan: batas $18$ sudah di bawah calon $20$, jadi seluruh cabangnya bisa dibuang. Pada
contoh sekecil ini penghematannya tidak terasa; pada soal dengan ratusan variabel, aturan
itulah yang membedakan selesai dalam semenit dan tidak selesai sama sekali, sebab tanpa
batas yang baik percabangan dan pembatasan merosot menjadi $2^{n}$ simpul.

Dua hal perlu jujur disebut. Waktu terburuknya tetap eksponensial, dan tidak ada
jaminan pohonnya kecil. Pemilihan variabel bercabang dan urutan simpul juga sangat
memengaruhi ukuran pohon; perangkat lunak memakai heuristik, misalnya variabel yang paling
dekat setengah. Keduanya menjelaskan mengapa program integer bisa jauh lebih lambat daripada
program linear berukuran sama.

\section{Contoh hitung}
\tingkat{menengah}

Sekarang seluruh angka dikerjakan sampai akhir, sehingga setiap klaim di atas bisa
diperiksa ulang dengan tangan atau dengan Python.

\begin{resep}
Empat langkah menyelesaikan program linear dua variabel secara grafik:
\begin{enumerate}
  \item Ubah tiap kendala menjadi garis dengan mengganti tanda pertidaksamaan menjadi tanda
        sama dengan, lalu gambar garisnya lewat dua titik potong sumbu.
  \item Tentukan sisi yang layak dengan menguji satu titik, biasanya titik nol; arsir sisi
        yang salah supaya daerah layak tetap bersih.
  \item Potongkan setiap pasangan garis, buang titik potong yang tidak memenuhi seluruh
        kendala, dan yang tinggal itulah verteks.
  \item Hitung nilai sasaran di setiap verteks, lalu ambil yang terbesar untuk maksimisasi
        atau terkecil untuk minimisasi.
\end{enumerate}
Peringatan urutan: langkah 3 tidak boleh dikerjakan sebelum langkah 2. Titik potong yang
dihitung lebih dulu sering terletak di luar daerah layak dan nilainya justru lebih tinggi
daripada jawaban sebenarnya; angka bagus yang muncul sebelum kelayakan diperiksa sulit
dibuang.
\end{resep}

\emph{Pertama, program linear bengkel.} Daerah layaknya segiempat dengan empat verteks,
dan keempatnya diperiksa:
\begin{center}
\begin{tabular}{lcc}
\hline
verteks & kendala aktif & laba $5x_{1}+4x_{2}$ \\
\hline
$(0;\,0)$ & $x_{1}=0,\ x_{2}=0$ & $0$ \\
$(4;\,0)$ & $6x_{1}+4x_{2}=24,\ x_{2}=0$ & $20$ \\
$(3;\,1{,}5)$ & $6x_{1}+4x_{2}=24,\ x_{1}+2x_{2}=6$ & $21$ \\
$(0;\,3)$ & $x_{1}+2x_{2}=6,\ x_{1}=0$ & $12$ \\
\hline
\end{tabular}
\end{center}
Titik $(3;1{,}5)$ dihitung dari kedua persamaan: $x_{1}=6-2x_{2}$ disubstitusi ke
$6x_{1}+4x_{2}=24$ memberi $36-8x_{2}=24$, jadi $x_{2}=3/2$, $x_{1}=3$, dan labanya
$15+6=21$.

\emph{Kedua, dual dan dualitas kuat.} Dualnya meminimumkan $24y_{1}+6y_{2}$ dengan
$6y_{1}+y_{2}\ge5$ dan $4y_{1}+2y_{2}\ge4$. Karena $x_{1},x_{2}>0$, kelengkapan
komplementer menuntut kedua kendala dual aktif; menyelesaikan keduanya sebagai kesamaan
memberi $y_{1}=3/4$ dan $y_{2}=1/2$. Nilainya $18+3=21$, sama dengan laba. Pemeriksaan silang:
$6\cdot\tfrac34+\tfrac12=5$ dan $4\cdot\tfrac34+2\cdot\tfrac12=4$, tepat sama dengan ruas
kanan, sehingga celahnya nol.

\emph{Ketiga, harga bayangan dan batas berlakunya.} Mengganggu ruas kanan memberi
tabel berikut, dengan $z$ dihitung ulang dari nol pada setiap baris.
\begin{center}
\begin{tabular}{lccc}
\hline
gangguan & $\mathbf{x}$ optimum & $z$ & perubahan \\
\hline
$b=(25;\,6)$ & $(3{,}25;\,1{,}375)$ & $21{,}75$ & $+0{,}75$ \\
$b=(24;\,7)$ & $(2{,}5;\,2{,}25)$ & $21{,}50$ & $+0{,}50$ \\
$b=(12;\,6)$ & $(0;\,3)$ & $12{,}00$ & $-9{,}00$ \\
$b=(24;\,4)$ & $(4;\,0)$ & $20{,}00$ & $-1{,}00$ \\
\hline
\end{tabular}
\end{center}
Baris ketiga mengurangi kayu lapis menjadi $12$, dengan perubahan
$-9=0{,}75\cdot(-12)$; baris keempat mengurangi $2$ jam, dengan perubahan
$-1=0{,}5\cdot(-2)$. Keduanya masih di dalam selang $b_{1}\in[12,36]$ dan
$b_{2}\in[4,12]$, jadi kelinierannya berlaku sampai ujung. Di luarnya kemiringan berganti:
dengan $b_{1}=40$, relaksasinya memberi $(0;\,3)$ bernilai $12$, bukan
$21+0{,}75\cdot16=33$.

\emph{Keempat, program integer bengkel.} Enumerasi dua arah tadi memberi jawaban $20$ di
$(4;0)$, sedangkan relaksasinya $21$ di $(3;1{,}5)$. Celahnya $1$, yaitu
$1/21=4{,}76\%$ dari nilai relaksasi. Pembulatan ke bawah memberi $19$, jadi kehilangan
$1$ satuan laba, atau $5\%$ dari jawaban sebenarnya, dan itu terjadi pada soal dengan dua
variabel saja.

\emph{Kelima, knapsack empat barang.} Empat barang dengan berat $5,4,3,2$ dan
nilai $10,7,5,3$ dimuat ke wadah berkapasitas $8$. Kerapatan nilainya $2$, $1{,}75$,
$1{,}667$, $1{,}5$, jadi aturan rakus mengambil barang pertama lalu $\tfrac34$ barang kedua:
berat $5+3=8$, nilai $15{,}25$. Jawaban bulat terbaiknya barang pertama dan ketiga, berat $8$
tepat penuh, nilai $15$. Celah integralitasnya $0{,}25$, hanya $1{,}64\%$ dari batas. Tetapi
pembulatan ke bawah membuang barang kedua sepenuhnya dan menyisakan nilai $10$, yaitu
$33{,}3\%$ di bawah jawaban sebenarnya: celah kecil tidak berarti pembulatan yang aman, sebab
keduanya mengukur hal berbeda.

\kertas{Untuk bengkel tadi, andaikan mesin perakitan hanya tersedia $5$ jam, bukan $6$.
Hitung ulang verteksnya, tentukan laba optimum bersambung, lalu tentukan laba optimum
bulatnya dengan enumerasi kecil. Sebelum menghitung, tulis dugaanmu: apakah celah
integralitasnya membesar atau mengecil? Sesudah itu periksa harga bayangan jam perakitan
pada $b_{2}=5$ dan bandingkan dengan $0{,}5$. Sebagai tambahan, cari $b_{2}$ tempat jawaban bulat terbaik berpindah dari $(4;0)$.}

\section{Lab komputasi}
\tingkat{lanjut}

Dua blok berikut dijalankan dengan Python 3.14.4, numpy 2.5.2, dan scipy 1.18.0. Blok
pertama menyelesaikan program linear, dualnya, dan gangguan ruas kanan. Blok kedua
menyelesaikan pohon percabangan, knapsack, dan selimut verteks pada graf lengkap lima
simpul. Keluaran disalin apa adanya.

\begin{lab}
\begin{verbatim}
import numpy as np
from scipy.optimize import linprog

np.set_printoptions(precision=6, suppress=True, linewidth=78)

c = np.array([5., 4.])              # laba per rak, per bangku
A = np.array([[6., 4.], [1., 2.]])  # kayu lapis, jam perakitan
b = np.array([24., 6.])

r = linprog(-c, A_ub=A, b_ub=b, bounds=[(0, None)]*2, method="highs")
print("LP   x* =", r.x, "  z* =", -r.fun)
print("sisa b - Ax =", b - A @ r.x)
print("marginals (scipy):", r.ineqlin.marginals)

rd = linprog(b, A_ub=-A.T, b_ub=-c, bounds=[(0, None)]*2, method="highs")
print("dual y* =", rd.x, "  nilai dual =", rd.fun)
print("celah dualitas =", rd.fun + r.fun)

print("B^-1 =", np.linalg.inv(A), "  x_B =", np.linalg.inv(A) @ b)

for nama, db in [("b=(25, 6)", (1., 0.)), ("b=(24, 7)", (0., 1.)),
                 ("b=(12, 6)", (-12., 0.)), ("b=(24, 4)", (0., -2.))]:
    rr = linprog(-c, A_ub=A, b_ub=b + np.array(db), bounds=[(0, None)]*2,
                 method="highs")
    print(f"{nama:<10} x = {np.round(rr.x, 6) + 0.0}  z = {-rr.fun:7.4f}"
          f"  delta z = {-rr.fun + r.fun:+.4f}")
\end{verbatim}
Keluaran sebenarnya:
\begin{verbatim}
LP   x* = [3.  1.5]   z* = 21.0
sisa b - Ax = [0. 0.]
marginals (scipy): [-0.75 -0.5 ]
dual y* = [0.75 0.5 ]   nilai dual = 21.0
celah dualitas = 0.0
B^-1 = [[ 0.25  -0.5  ]
 [-0.125  0.75 ]]   x_B = [3.  1.5]
b=(25, 6)  x = [3.25  1.375]  z = 21.7500  delta z = +0.7500
b=(24, 7)  x = [2.5  2.25]  z = 21.5000  delta z = +0.5000
b=(12, 6)  x = [0. 3.]  z = 12.0000  delta z = -9.0000
b=(24, 4)  x = [4. 0.]  z = 20.0000  delta z = -1.0000
\end{verbatim}
\end{lab}

\begin{lab}
\begin{verbatim}
import itertools
import numpy as np
from scipy.optimize import linprog, milp, LinearConstraint, Bounds

c = np.array([5., 4.])
A = np.array([[6., 4.], [1., 2.]])
b = np.array([24., 6.])


def lp_bengkel(tambahan_A=None, tambahan_b=None):
    Aub = A if tambahan_A is None else np.vstack([A, tambahan_A])
    bub = b if tambahan_b is None else np.concatenate([b, tambahan_b])
    r = linprog(-c, A_ub=Aub, b_ub=bub, bounds=[(0, None)]*2, method="highs")
    if r.status != 0:
        return (None, "tidak layak")
    return (-r.fun, np.round(r.x, 6) + 0.0)


print("-- pohon percabangan dan pembatasan --")
simpul = [("akar",        None, None),
          ("x2<=1",       [[0., 1.]],  [1.]),
          ("x2>=2",       [[0., -1.]], [-2.]),
          ("x2<=1,x1<=3", [[0., 1.], [1., 0.]],  [1., 3.]),
          ("x2<=1,x1>=4", [[0., 1.], [-1., 0.]], [1., -4.])]
for nama, ta, tb in simpul:
    z, x = lp_bengkel(None if ta is None else np.array(ta),
                      None if tb is None else np.array(tb))
    print(f"{nama:<14} z_LP = {z:8.4f}  x = {x}")

rb = milp(-c, constraints=LinearConstraint(A, -np.inf, b),
          integrality=[1, 1], bounds=Bounds(0, np.inf))
print("IP optimum   z_IP =", -rb.fun, " x =", np.round(rb.x, 6) + 0.0,
      " celah =", 21.0 + rb.fun)

print("-- knapsack 4 barang, kapasitas 8 --")
w = np.array([5., 4., 3., 2.])
v = np.array([10., 7., 5., 3.])
for mask in range(16):
    S = [i + 1 for i in range(4) if mask >> i & 1]
    bw, bv = w[[i - 1 for i in S]].sum(), v[[i - 1 for i in S]].sum()
    print(f"{str(S):<12} berat {bw:4.0f}  nilai {bv:4.0f}"
          f"  {'layak' if bw <= 8 else 'tidak layak'}")
rk = milp(-v, constraints=LinearConstraint(w, -np.inf, 8.),
          integrality=[1]*4, bounds=Bounds(0, 1))
print("enumerasi/MILP  z_IP =", -rk.fun, " x =", np.round(rk.x, 6) + 0.0)
rl = linprog(-v, A_ub=[w], b_ub=[8.], bounds=[(0, 1)]*4, method="highs")
print("relaksasi       z_LP =", -rl.fun, " x =", np.round(rl.x, 6) + 0.0)
print("celah integralitas =", -rl.fun + rk.fun)

print("-- selimut verteks K5: LP vs IP --")
n = 5
E = list(itertools.combinations(range(n), 2))
M = np.zeros((len(E), n))
for k, (i, j) in enumerate(E):
    M[k, i] = M[k, j] = 1.
rv = linprog(np.ones(n), A_ub=-M, b_ub=-np.ones(len(E)),
             bounds=[(0, 1)]*n, method="highs")
rc = milp(np.ones(n), constraints=LinearConstraint(M, 1, np.inf),
          integrality=[1]*n, bounds=Bounds(0, 1))
print("LP =", rv.fun, " x =", np.round(rv.x, 6) + 0.0)
print("IP =", rc.fun, " x =", np.round(rc.x, 6) + 0.0)
\end{verbatim}
Keluaran sebenarnya:
\begin{verbatim}
-- pohon percabangan dan pembatasan --
akar           z_LP =  21.0000  x = [3.  1.5]
x2<=1          z_LP =  20.6667  x = [3.333333 1.      ]
x2>=2          z_LP =  18.0000  x = [2. 2.]
x2<=1,x1<=3    z_LP =  19.0000  x = [3. 1.]
x2<=1,x1>=4    z_LP =  20.0000  x = [4. 0.]
IP optimum   z_IP = 20.0  x = [4. 0.]  celah = 1.0
-- knapsack 4 barang, kapasitas 8 --
[]           berat    0  nilai    0  layak
[1]          berat    5  nilai   10  layak
[2]          berat    4  nilai    7  layak
[1, 2]       berat    9  nilai   17  tidak layak
[3]          berat    3  nilai    5  layak
[1, 3]       berat    8  nilai   15  layak
[2, 3]       berat    7  nilai   12  layak
[1, 2, 3]    berat   12  nilai   22  tidak layak
[4]          berat    2  nilai    3  layak
[1, 4]       berat    7  nilai   13  layak
[2, 4]       berat    6  nilai   10  layak
[1, 2, 4]    berat   11  nilai   20  tidak layak
[3, 4]       berat    5  nilai    8  layak
[1, 3, 4]    berat   10  nilai   18  tidak layak
[2, 3, 4]    berat    9  nilai   15  tidak layak
[1, 2, 3, 4] berat   14  nilai   25  tidak layak
enumerasi/MILP  z_IP = 15.0  x = [1. 0. 1. 0.]
relaksasi       z_LP = 15.25  x = [1.   0.75 0.   0.  ]
celah integralitas = 0.25
-- selimut verteks K5: LP vs IP --
LP = 2.5  x = [0.5 0.5 0.5 0.5 0.5]
IP = 4.0  x = [1. 1. 1. 1. 0.]
\end{verbatim}
\end{lab}

Tiga hal perlu dibaca dari keluaran pertama. $\mathbf{x}^{*}$ yang dilaporkan
$(3;\,1{,}5)$ dengan $z^{*}$ tepat $21$, tanpa sisa pada kedua kendala, sesuai kelengkapan
komplementer. \texttt{scipy} melaporkan pengali $(-0{,}75,\,-0{,}5)$, negatif dari harga
bayangan, sebab rutin meminimumkan $-\mathbf{c}^{\top}\mathbf{x}$; angka benarnya $0{,}75$
dan $0{,}5$, dan cara memastikannya adalah menghitung ulang lewat $B^{-1}$. Celah dualitas
keluar $0$ tepat, bukan $10^{-15}$, karena $21$ terwakili tepat dalam biner; pada data
desimal, sisa orde $10^{-12}$ adalah hal biasa.

Keluaran kedua cocok dengan hitungan tangan: lima simpul dengan nilai yang sama
sampai empat desimal, jawaban bulat $20$ di $(4;\,0)$, relaksasi knapsack $15{,}25$, dan
jawaban bulatnya $15$ dengan celah $0{,}25$. Enumerasi keenam belas himpunan bagian
menegaskan $15$ yang terbaik: semua himpunan bernilai lebih tinggi melampaui kapasitas.

Bagian terakhir keluaran kedua memeriksa klaim faktor dua tadi. Untuk graf
lengkap lima simpul, relaksasinya $2{,}5$ sedangkan jawaban bulatnya $4$; nisbahnya $1{,}6$,
dan untuk $n$ simpul nisbahnya $2(n-1)/n$ menuju $2$ dari bawah. Batas relaksasi di sini bisa
meleset hampir dua kali.

\begin{untukdosen}
Urutan yang berhasil dipakai: mulai dari bengkel dan jawaban setengah bangku, lalu bentuk
baku, lalu verteks, baru simpleks. Jangan mulai dari tabel. Mahasiswa yang melihat tabel
lebih dulu menghafal langkah pivot tanpa tahu mengapa langkah itu boleh berhenti; yang sudah
melihat jawaban ada di sudut akan menganggap pivot sebagai cara berpindah sudut, dan uji
keoptimalan sebagai cara tahu tidak ada sudut yang lebih baik.

Dua miskonsepsi khas. Pertama, ``verteks adalah semua perpotongan garis
kendala''. Banyak perpotongan terletak di luar daerah layak, misalnya $(4;3)$ pada contoh
kita, yang memotong kedua garis tetapi melanggar keduanya; memeriksa kelayakan setiap kandidat
perlu dikerjakan di papan. Kedua, ``jawaban bulat adalah pembulatan jawaban pecahan''. Bengkel
memberi $19$ melawan $20$, knapsack memberi $10$ melawan $15$; keduanya cukup untuk
mematahkan dugaan itu.

Aktivitas kelas: tulis tabel awal di papan, minta kelas memilih kolom masuk dan menghitung
uji nisbah, lalu memprediksi laba sebelum pivot. Sesudah tabel akhir, tanyakan laba bila
kayu lapis menjadi $25$ lembar; $21{,}75$ membuka jalan ke selang $[12,36]$.

Tanda kesiapan lanjut: mahasiswa dapat menjelaskan mengapa metode simpleks tidak perlu
memeriksa seluruh verteks, dan dapat menyebut satu alasan mengapa relaksasi tidak boleh
dibulatkan.
\end{untukdosen}

\section{Ringkasan}

\begin{ringkasan}
\begin{itemize}
  \item Program linear menuntut sasaran dan kendala berbentuk garis. Bentuk baku
        $\max\mathbf{c}^{\top}\mathbf{x}$ dengan $A\mathbf{x}=\mathbf{b}$ dan
        $\mathbf{x}\ge\mathbf{0}$ berlaku untuk semua kasus, dengan variabel kendur untuk
        mengubah ketaksamaan menjadi kesamaan.
  \item Bila daerah layak terbatas dan sasarannya terbatas di atas, nilai terbaik dicapai di
        verteks. Pada bentuk baku, verteks sama dengan solusi layak basis, dan jumlahnya
        $\binom{n}{m}$, sehingga memeriksa semuanya bukan rencana masuk akal.
  \item Metode simpleks berjalan dari verteks ke verteks lewat kolom dengan biaya tereduksi
        negatif dan uji nisbah, dan berhenti ketika tidak ada biaya tereduksi negatif; saat
        itu baris $z$ sudah memuat harga dual.
  \item Dualitas lemah berlaku tanpa syarat; dualitas kuat menyatakan kedua nilai sama bila
        keduanya berhingga; kelengkapan komplementer menyatakan kendala yang longgar
        berharga nol.
  \item Harga bayangan $y_{i}$ adalah laju perubahan nilai optimum terhadap ruas kanan, dan
        hanya berlaku selama susunan kendala aktif tidak berubah; pada bengkel $0{,}75$ per
        lembar dan $0{,}5$ per jam, sah untuk $b_{1}\in[12,36]$ dan $b_{2}\in[4,12]$.
  \item Relaksasi LP memberi batas yang sah, bukan jawaban. Pembulatan bisa memberi titik
        tidak layak, seperti $(3;2)$, atau titik layak yang bukan terbaik, seperti $(3;1)$
        bernilai $19$ melawan $20$.
  \item Percabangan dan pembatasan memakai batas relaksasi untuk membuang cabang. Pohon
        bengkel berisi lima simpul, jawaban bulatnya $20$ dengan celah $1$. Waktu terburuknya
        tetap eksponensial, dan versi keputusannya NP-lengkap.
\end{itemize}
\end{ringkasan}

\section{Latihan}

\begin{latihan}
\begin{enumerate}
  \item Sebuah toko kue memakai $10$ kg tepung dan $8$ jam oven. Kue A memakai $2$ kg dan
        $1$ jam dengan laba $30$; kue B memakai $1$ kg dan $2$ jam dengan laba $20$.
        Tuliskan program linearnya, gambar daerah layaknya, dan hitung laba di setiap
        verteks. Lalu tentukan jawaban bulat terbaik dengan enumerasi, dan bandingkan
        dengan hasil pembulatan ke bawah dari jawaban bersambung.

  \item Buktikan bahwa himpunan $\{\mathbf{x}:A\mathbf{x}\le\mathbf{b}\}$ konveks, lalu
        jelaskan mengapa bukti itu tidak berlaku untuk himpunan
        $\{\mathbf{x}\in\mathbb{Z}^{n}:A\mathbf{x}\le\mathbf{b}\}$.

  \item Ganti laba bengkel menjadi $x_{1}+2x_{2}$. Hitung nilainya di keempat verteks,
        sebutkan verteks mana yang optimal, dan tunjukkan bahwa seluruh ruas di antara
        dua verteks optimal juga optimal. Terangkan mengapa titik di tengah ruas itu bukan
        verteks, lalu hubungkan hasilnya dengan harga bayangan kendala jam.

  \item Selesaikan dengan tabel simpleks: maksimumkan $3x_{1}+2x_{2}$ dengan
        $x_{1}+x_{2}\le4$, $x_{1}\le3$, $x_{2}\le3$, $x\ge0$. Tuliskan setiap tabel antara,
        sebutkan kolom masuk dan alasannya, lalu bacakan harga dualnya dari tabel akhir.
        (Petunjuk: dua pivot cukup; bandingkan hasil akhir dengan verteks yang diperoleh
        dari perpotongan $x_{1}+x_{2}=4$ dan $x_{1}=3$.)

  \item Tuliskan dual dari program pada latihan 1, selesaikan, dan periksa dualitas kuat
        secara numerik. Jelaskan arti setiap harga bayangan dalam satu kalimat.

  \item Nilai klaim ini: \emph{pembulatan ke bawah dari jawaban relaksasi selalu memberi
        titik bulat yang layak.} Tentukan benar, salah, atau belum cukup bukti, dan
        tunjukkan satu contoh penyangkal dengan angka.

  \item Nilai klaim ini: \emph{harga bayangan yang positif berarti kendala itu pasti
        menekan.} Tentukan benar, salah, atau belum cukup bukti, lalu jelaskan hubungan
        yang sebenarnya antara harga bayangan dan kelengkapan komplementer.
        (Petunjuk: periksa apa yang terjadi pada kendala dengan $y_{i}>0$ menurut
        $y_{i}(b_{i}-\mathbf{a}_{i}^{\top}\mathbf{x})=0$.)

  \item Untuk knapsack dengan berat $(5,4,3,2)$ dan nilai $(10,7,5,3)$ serta kapasitas
        $8$, hitung relaksasi LP-nya dengan tangan memakai urutan kerapatan, lalu bandingkan
        dengan hasil enumerasi. Hitung celah integralitasnya dan nyatakan sebagai persen
        dari batas. (Petunjuk: urutkan menurut nilai per satuan berat, lalu isi wadah
        sampai penuh; barang terakhir masuk sebagian.)

  \item Ulangi latihan 8 dengan kapasitas $7$. Tentukan relaksasi, jawaban bulat, dan
        celahnya, lalu jelaskan mengapa celahnya bisa berubah walaupun hanya satu angka
        data yang berubah.

  \item Pada selimut verteks graf lengkap $n$ simpul, buktikan bahwa relaksasi bernilai
        $n/2$ dan jawaban bulatnya $n-1$, lalu hitung nisbah keduanya untuk $n=4,5,10$.
        Jelaskan mengapa hasil itu menunjukkan bahwa pembulatan bukan cara umum untuk
        menyelesaikan program integer.
\end{enumerate}
\end{latihan}

\begin{untukpeneliti}
Yang benar-benar terbukti di bab ini: dualitas lemah untuk pasangan primal-dual dalam bentuk
baku, tanpa syarat apa pun; arah balik kelengkapan komplementer sebagai uji keoptimalan;
keoptimalan di verteks untuk daerah layak terbatas, dengan bukti lewat argumen kendala aktif;
dan kesahan pemangkasan pada percabangan dan pembatasan. Yang dipinjam: dualitas kuat, yang bukti
lazimnya lewat lemma Farkas atau lewat berhentinya simpleks; keduanya tercatat pada
\texttt{WSPrinted.pdf} bagian dualitas program linear dan pada Boyd dan Vandenberghe
bagian 5.2.

Yang dipopulerkan melampaui buktinya cukup banyak, dan tiga perlu disebut.
``Relaksasi LP sudah hampir benar'': faktor aproksimasi terburuknya bisa tumbuh tanpa batas
pada kelas masalah tertentu, dan bahkan pada selimut verteks mencapai $2$. ``Harga bayangan
adalah turunan nilai optimum'': itu benar hanya bila nilai optimum terdiferensialkan terhadap
ruas kanan; pada titik kemerosotan, saat dua basis optimal berbagi verteks yang sama,
turunannya tidak ada dan yang tersisa hanya subgradien. ``Cabang dan batas selalu cepat'':
waktu terburuknya eksponensial, dan kelambatan pada instans tertentu bukan tanda implementasi
yang buruk.

Pertanyaan terbuka: seberapa kuat hierarki pemotongan Chvatal-Gomory untuk kelas polytope
tertentu, dan apakah ada batas bawah superpolinomial untuk ukuran formulasi diperluas
polytope traveling salesman; hasil Fiorini dan kawan-kawan pada 2012 menunjukkan batas
eksponensial untuk kasus itu, sementara pencarian batas serupa untuk masalah lain
berjalan. Batas keberlakuan bab ini: data rasional, aritmetika eksak
untuk pembuktian, toleransi titik mengambang untuk praktik; seluruh pernyataan tentang harga
bayangan berlaku pada model bersambung, bukan model integer.

\peta{Bab ini mengikuti \emph{Introduction to Linear Optimization} pada berkas
\texttt{WSPrinted.pdf}: bagian formulasi, geometri himpunan polihedral (verteks dan titik
ekstrem), sistem ketaksamaan linear dan dualitas (kelengkapan komplementer, analisis
sensitivitas), serta simpleks (solusi layak basis, biaya tereduksi, uji keoptimalan). Sisi
program integer mengikuti Wolsey, \emph{Integer Programming}, bab 1 (formulasi) dan bab 2
(optimalitas, relaksasi, batas), yang menempatkan relaksasi sebagai sumber batas dan
percabangan sebagai cara memakainya. Sudut pandang konveks pada bagian kedua dan keempat
adalah Boyd dan Vandenberghe, \emph{Convex Optimization}, bab 4 dan 5. Hubungan dengan bab
lain: bentuk baku dan daerah layak melanjutkan Bab 21, penyelesaian sistem linear memakai
Bab 8, dan celah integralitas dipakai lagi pada Bab 25.}

\end{untukpeneliti}


\part{Pembelajaran Mesin, Keamanan, dan Privasi}

Bagian ini menutup buku dengan pertanyaan yang paling sering dihindari: kapan
sebuah mesin yang tampak pintar sebenarnya rapuh, dan bagaimana kerapuhannya
dibuktikan dengan angka. Prasyaratnya Bab~10, 12, 18, dan~23. Setelah delapan bab,
kamu bisa melatih dan mengevaluasi model secara jujur, menyerangnya, membuktikan
batas serangan, serta menjelaskan kriptografi dan privasi diferensial pada tingkat
yang cukup untuk membaca makalah aslinya.
\chapter{Jaringan Saraf dan Turunan Balik}

Sebuah sistem pencegahan intrusi mengamati empat besaran pada setiap koneksi: banyak
permintaan per detik, lama koneksi, banyak paket yang ditolak, dan ukuran muatan rata-rata.
Analis yang mencoba memisahkan lalu lintas berbahaya dari lalu lintas wajar menemukan pola
yang tidak enak dibaca. Koneksi dengan permintaan tinggi tetapi tanpa paket ditolak biasanya
sah, sebab begitulah pengambilan berkas besar. Koneksi dengan banyak paket ditolak tetapi
permintaan rendah juga biasanya sah, sebab begitulah pemindaian yang gagal. Gabungan
keduanya hampir selalu serangan. Setiap ciri sendiri tidak memisahkan apa pun, sedangkan
gabungannya memisahkan dengan tajam.

Ambang linear tidak dapat menggambar pemisah semacam itu, dan bagian pertama bab ini
menunjukkan bahwa batas itu bukan soal jumlah data. Yang harus berubah adalah bentuk
fungsinya. Satu lapisan tersembunyi sudah cukup untuk pola tadi. Setelah bentuknya berubah,
muncul pertanyaan berikutnya yang lebih gawat: bagaimana kita menyesuaikan ketiga belas
parameter jaringan kecil itu tanpa menurunkan rumus baru untuk setiap arsitektur?

Jawabannya adalah aturan rantai Bab 12, dijalankan mundur pada graf komputasi. Bab ini
menurunkan aturan itu sampai ke angka: satu langkah maju dengan bobot yang ditentukan,
ketiga belas gradien dihitung dengan tangan lalu diperiksa dengan NumPy, satu langkah
penurunan gradien, dan gradien terhadap masukan yang nanti menjadi senjata di Bab 29.
Yang dijanjikan bab ini bukan jaringan yang pintar, melainkan mesin turunan yang dapat
diperiksa baris demi baris.

\begin{untukanak}
Bayangkan sebuah lini perakitan dengan beberapa tukang. Tukang pertama mengerjakan bahan
mentah, tukang kedua mengerjakan hasil tukang pertama, dan begitu seterusnya sampai tukang
terakhir menyerahkan barang jadi kepada pemeriksa. Pemeriksa menimbang barang itu, lalu
berkata: terlalu berat, kurang dua ratus gram.

Sekarang pertanyaannya: siapa yang harus mengubah cara kerjanya? Tukang pertama? Tukang
terakhir? Kalau mandor hanya bisa mengubah satu setelan lalu menjalankan ulang seluruh lini, ia butuh dua kali percobaan per setelan. Dengan tiga belas setelan, itu dua puluh enam kali menjalankan lini untuk satu perbaikan.

Ternyata ada cara yang jauh lebih pintar. Sesudah satu kali jalan, kita bisa menelusuri
kembali dari pemeriksa ke tukang terakhir, lalu ke tukang sebelumnya, sampai ke tukang
pertama. Di setiap perhentian kita hitung satu angka: seberapa besar kesalahan akhir berubah
kalau orang ini mengubah sedikit hasil kerjanya. Sekali telusur mundur, semua tukang tahu
bagiannya masing-masing, tanpa ada yang perlu dijalankan ulang.

Di bab ini kita akan tahu mengapa informasi bisa mengalir mundur padahal barangnya hanya
mengalir maju, dan mengapa telusur mundur itu memberi jawaban yang sama tepatnya dengan
menjalankan lini dua puluh enam kali.
\end{untukanak}

\section{Jembatan pemahaman: dari cerita ke simbol}
\tingkat{dasar}

Rumus di bab ini lebih mudah dibaca kalau urutannya dibalik: cerita tanpa simbol, lalu
hitungan di kertas, lalu simbol sebagai singkatan dari langkah yang sudah kita jalankan.

\subsection{Analogi nyata: seberapa besar, bukan siapa}

Analogi tukang tadi menjawab pertanyaan ``siapa yang harus diubah''. Yang belum dijawab
adalah ``seberapa besar''. Bayangkan tiga keran berderet di atas satu ember dengan satu
garis sasaran. Memutar keran ketiga paling terasa, tetapi keran itu hanya mengalirkan
kiriman keran kedua; membuka keran pertama paling menentukan, tetapi besar pengaruhnya
bergantung pada dua keran di antaranya. Yang kita butuhkan bukan memutar semua keran lalu
mengukur ulang, melainkan satu angka per keran: berapa milimeter tinggi air berubah kalau
keran ini diputar sedikit? Satu kali jalan, telusuri mundur dari ember ke keran pertama.

\subsection{Hitung manual dulu}

\begin{contoh}
Satu neuron, empat angka: bobot $w = 0.5$, bias $b = 0.1$, masukan $x = 2.0$, jawaban benar
$y = 1$, dan tingkat belajar $\eta = 0.1$. Ikuti langkahnya berurutan; setiap angka di sini
bisa diperiksa dengan kalkulator biasa.

\begin{enumerate}
  \item Skor mentah: $w x = 0.5 \times 2.0 = 1.0$, ditambah bias menjadi $z = 1.1$.
  \item Peluang: $e^{-1.1} = 0.33287$, jadi $\sigma = 1/(1+0.33287) = 0.7502601$.
  \item Kerugian: karena $y = 1$, cukup satu suku, $L = -\ln 0.7502601 = 0.2873353$.
  \item Selisih peluang dan kenyataan: $\sigma - y = 0.7502601 - 1 = -0.2497399$.
  \item Aturan rantai ke bobot: kalikan selisih itu dengan masukan, $-0.2497399 \times 2.0
        = -0.4994798$. Ke bias tidak dikalikan apa pun, jadi tetap $-0.2497399$.
  \item Geser: $w' = 0.5 - 0.1(-0.4994798) = 0.5499480$ dan $b' = 0.1 - 0.1(-0.2497399)
        = 0.1249740$. Skor baru $z' = 1.2248699$, peluang baru $0.7729194$, kerugian baru
        $0.2575805$. Turun $0.0297549$ dari $0.2873353$.
\end{enumerate}

Perhatikan tanda di Langkah 5: gradien negatif, dan karena kita \emph{mengurangi} bobot
dengan gradien, bobot justru naik. Tanda minus di rumus geser bukan hiasan; ia yang
membalik arah.
\end{contoh}

Pemeriksaan yang tidak boleh dilewat: hitung kerugian pada $w + 0.000001$ dan
$w - 0.000001$, bagi selisihnya dengan $0.000002$. Hasilnya $-0.4994798$, sama sampai tujuh
angka. Kalau manual dan numerik berbeda, yang salah hampir selalu tanda.

\subsection{Kamus baca simbol}

Simbol berikut muncul di buku mana pun tentang jaringan saraf. Tidak ada yang perlu
dihafal: bacanya cukup satu kalimat, dan kalimat itu sudah kita hitung sendiri di atas.

\begin{description}
  \item[$\partial L/\partial w$] dibaca ``del L del w'', artinya: kalau bobot $w$ digeser
        sedikit, berapa perubahan kerugian $L$? Di contoh kita jawabannya $-0.4994798$
        (Langkah 5).
  \item[$z = Wx + b$] dibaca ``z sama dengan W kali x tambah b'', artinya skor mentah
        sebelum ditekan menjadi peluang. Versi satu angka dari rumus ini adalah Langkah 1.
  \item[$\sigma(z)$] dibaca ``sigma dari z'', artinya fungsi kurva S yang menekan skor agar
        selalu di antara nol dan satu. Di Langkah 2 hasilnya $0.7502601$.
  \item[$\nabla_w L$] dibaca ``nabla w L'', artinya kumpulan semua turunan kerugian
        terhadap seluruh bobot sekaligus; tanda panahnya menjawab ``ke arah mana kerugian
        paling cepat turun''. Isinya angka-angka seperti Langkah 5, hanya banyak.
\end{description}

\subsection{Kenapa ini penting untuk keamanan?}

Turunan yang sama yang dipakai memperbaiki model juga dipakai merusaknya. Penyerang hanya
butuh satu angka: kalau satu ciri masukan digeser sedikit, berapa kerugian naik? Di contoh
kita $x$ punya turunan $-0.1248699$, jadi mengecilkan $x$ menaikkan kerugian. Menggeser
$1.9$ membuat $x$ menjadi $0.1$, peluang jatuh ke $0.5374298$, kerugian naik dari $0.2873353$ ke
$0.6209570$, yaitu $2.16$ kali lipat. Satu ciri, satu arah, satu langkah.

Bagi pihak bertahan, pertanyaan yang sama dipakai terbalik: ciri mana yang turunannya
paling besar? Ciri itulah yang paling murah dipalsukan, sehingga paling perlu dibulatkan,
dipangkas, atau diperiksa keanehannya. Bab 29 memakai resep ini secara utuh, dengan nama
gangguan $\epsilon$ dan rumus
$x_{\mathrm{adv}} = x + \epsilon \cdot \operatorname{sign}(\nabla_x L)$.

\begin{catatan}
Peringatan yang sering terlewat: pada sistem nyata penyerang jarang bisa menghitung gradien
eksak; ia menaksirnya dengan mengubah masukan berkali-kali. Taksiran itu lebih mahal dan
lebih berisik, tetapi tetap bekerja. Karena itu ``modelnya tertutup'' bukan pertahanan.
\end{catatan}

\subsection{Latihan menghitung dengan tangan}

Kerjakan dengan kalkulator, tulis di kertas, baru cocokkan. Jawaban ada di tanda kurung.

\begin{catatan}
\begin{enumerate}
  \item Ulangi seluruh contoh dengan $b = 0.3$ dan angka lain tetap. Hitung $z$, $\sigma$,
        $L$, $\partial L/\partial w$, $\partial L/\partial b$.
        ($z = 1.3$; $\sigma = 0.7858350$; $L = 0.2410085$; $-0.4283300$; $-0.2141650$.)
  \item Ulangi contoh asli dengan $\eta = 0.5$. Berapa $w'$, $b'$, dan kerugian barunya?
        ($w' = 0.7497399$; $b' = 0.2248699$; $L' = 0.1640633$, turun $0.1232720$.)
  \item Ulangi contoh asli dengan $y = 0$, angka lain tetap. Berapa $\partial L/\partial w$
        sekarang, dan ke arah mana bobot bergerak?
        ($0.7502601$ dan $1.5005202$; positif, jadi bobot turun.)
  \item Pada contoh asli, geser $x$ sebesar $\epsilon = 0.5$ ke arah yang menaikkan
        kerugian. Berapa $x$ baru dan kerugiannya? ($x = 1.5$; $L = 0.3558651$.)
\end{enumerate}
\end{catatan}

\subsection{Dari kertas ke kode}

Kode berikut menjalankan enam langkah yang sama, dengan nama variabel yang menunjuk ke
langkahnya. Tidak ada yang baru di sini: setiap baris adalah satu langkah di atas.

\begin{lab}
\begin{verbatim}
import numpy as np

# Data satu contoh, sama dengan angka di "Hitung manual dulu".
x, y = 2.0, 1.0
w, b, eta = 0.5, 0.1, 0.1

# Langkah 1: z = w*x + b  (skor mentah)
z = w*x + b
print("Langkah 1  z        =", round(z, 7))

# Langkah 2: sigma(z) = 1/(1+e^-z)  (tekan skor jadi peluang 0..1)
s = 1/(1+np.exp(-z))
print("Langkah 2  sigma    =", round(s, 7))

# Langkah 3: L = -ln(s)  (kerugian; y=1 jadi cukup satu suku)
L = -np.log(s)
print("Langkah 3  L        =", round(L, 7))

# Langkah 4: dL/dz = sigma - y  (selisih peluang dan kenyataan)
dLdz = s - y
print("Langkah 4  dL/dz    =", round(dLdz, 7))

# Langkah 5: aturan rantai ke bobot dan bias
dLdw, dLdb = dLdz*x, dLdz
print("Langkah 5  dL/dw    =", round(dLdw, 7), " dL/db =", round(dLdb, 7))

# Langkah 6: satu langkah penurunan gradien
w2, b2 = w - eta*dLdw, b - eta*dLdb
s2 = 1/(1+np.exp(-(w2*x + b2)))
print("Langkah 6  w', b'   =", round(w2, 7), round(b2, 7),
      " L' =", round(-np.log(s2), 7))

# Pemeriksa: gradien hasil aturan rantai vs beda hingga numerik.
def loss(ww, bb):
    return -np.log(1/(1+np.exp(-(ww*x + bb))))
h = 1e-6
print("cek beda hingga dL/dw =", round((loss(w+h, b)-loss(w-h, b))/(2*h), 7),
      " dL/db =", round((loss(w, b+h)-loss(w, b-h))/(2*h), 7))

# Gradien terhadap MASUKAN: inilah yang dipakai menyerang model (Bab 29).
dLdx = dLdz*w
xa = x + 1.9*np.sign(dLdx)          # tanda gradien menentukan arah geser
sa = 1/(1+np.exp(-(w*xa + b)))
print("dL/dx =", round(dLdx, 7), " x_adv =", round(xa, 4),
      " sigma =", round(sa, 7), " L =", round(-np.log(sa), 7))
\end{verbatim}

Keluarannya harus sama dengan hitungan kertas: $1.1$, $0.7502601$, $0.2873353$,
$-0.2497399$, $-0.4994798$, lalu $0.2575805$. Kalau satu angka meleset, jangan menambal
kodenya: cari dulu langkah mana di kertas yang berbeda.
\end{lab}

\section{Ambang linear dan perceptron}
\tingkat{dasar}

Bentuk pengklasifikasi paling tua adalah ambang linear. Untuk masukan
$\mathbf{x}\in\mathbb{R}^{d}$, bobot $\mathbf{w}\in\mathbb{R}^{d}$, dan bias
$b\in\mathbb{R}$, kita hitung satu skor dan memutuskan kelas dari tanda skor itu.

\kamus{perceptron}{perceptron}{Pengklasifikasi yang memutuskan hanya dari sisi mana sebuah
titik jatuh terhadap satu garis atau satu bidang pemisah.}

\begin{definisi}[Ambang linear]
Fungsi $f:\mathbb{R}^{d}\to\{-1,+1\}$ disebut ambang linear bila ada $\mathbf{w}$ dan $b$
sehingga $f(\mathbf{x})=\operatorname{sign}(\mathbf{w}^{\top}\mathbf{x}+b)$. Himpunan
$\{\mathbf{x}:\mathbf{w}^{\top}\mathbf{x}+b=0\}$ adalah hiperbidang pemisah, dan hiperbidang
itu satu-satunya bentuk pemisah yang dapat dipelajari model ini.\index{perceptron}
\end{definisi}

Aturan pelatihan perceptron bekerja dengan memperbaiki bobot hanya ketika ada kesalahan.
Untuk label $y\in\{-1,+1\}$ dan contoh $\mathbf{x}$, bila
$y\,\mathbf{w}^{\top}\mathbf{x}\le 0$ maka kita perbarui
\[
  \mathbf{w}\leftarrow\mathbf{w}+y\,\mathbf{x}.
\]
Perhatikan bahwa tidak ada laju belajar di sini: besarnya langkah ditentukan oleh besarnya
contoh, dan pelatihan berhenti sendiri begitu semua contoh benar. Apakah ia selalu berhenti?

\begin{teorema}[Kekonvergenan perceptron]
Misalkan setiap contoh memenuhi $\lVert\mathbf{x}_{i}\rVert\le1$ dan terdapat vektor satuan
$\mathbf{w}^{*}$ serta $\gamma>0$ sehingga $y_{i}\,\mathbf{w}^{*\top}\mathbf{x}_{i}\ge\gamma$
untuk setiap $i$. Maka algoritme perceptron yang mulai dari $\mathbf{w}_{0}=\mathbf{0}$
melakukan paling banyak $\lfloor1/\gamma^{2}\rfloor$ pembaruan.\index{perceptron}
\end{teorema}

\begin{proof}
Misalkan pembaruan terjadi $k$ kali, pada contoh $\mathbf{x}_{1},\dots,\mathbf{x}_{k}$.
Setiap pembaruan hanya terjadi ketika contoh itu salah klasifikasi, jadi
$y_{j}\,\mathbf{w}_{j-1}^{\top}\mathbf{x}_{j}\le0$. Maka
\[
  \lVert\mathbf{w}_{j}\rVert^{2}
  =\lVert\mathbf{w}_{j-1}\rVert^{2}+2y_{j}\mathbf{w}_{j-1}^{\top}\mathbf{x}_{j}
   +\lVert\mathbf{x}_{j}\rVert^{2}
  \le\lVert\mathbf{w}_{j-1}\rVert^{2}+1 .
\]
Dengan induksi, $\lVert\mathbf{w}_{k}\rVert^{2}\le k$.

Sisi lain, hasil kali dengan $\mathbf{w}^{*}$ bertambah paling sedikit $\gamma$ setiap kali:
\[
  \mathbf{w}^{*\top}\mathbf{w}_{j}
  =\mathbf{w}^{*\top}\mathbf{w}_{j-1}+y_{j}\mathbf{w}^{*\top}\mathbf{x}_{j}
  \ge\mathbf{w}^{*\top}\mathbf{w}_{j-1}+\gamma ,
\]
sehingga $\mathbf{w}^{*\top}\mathbf{w}_{k}\ge k\gamma$. Karena $\mathbf{w}^{*}$ vektor
satuan, ketaksamaan Cauchy-Schwarz memberi
\[
  k\gamma\le\mathbf{w}^{*\top}\mathbf{w}_{k}\le\lVert\mathbf{w}_{k}\rVert\le\sqrt{k}.
\]
Setelah dibagi $\sqrt{k}$ dan dikuadratkan, $k\gamma^{2}\le1$, yaitu
$k\le1/\gamma^{2}$.
\end{proof}

Batas itu tidak bergantung pada dimensi maupun banyaknya contoh, dan itulah daya tariknya. Tetapi bila $\gamma$ tidak ada, teorema ini diam, dan perceptron memang berayun selamanya; ia juga bukan pernyataan generalisasi.

Sekarang batas yang lebih mendasar. Empat titik pada bidang sudah cukup untuk menunjukkan
bahwa tidak setiap pelabelan dapat dipisahkan oleh satu hiperbidang. Titik $(0,0)$ dan
$(1,1)$ berlabel $0$, titik $(1,0)$ dan $(0,1)$ berlabel $1$. Pola ini biasa disebut XOR.

\begin{proposisi}
Keempat titik XOR tidak terpisah secara linear.
\end{proposisi}

\begin{proof}
Andaikan ada $w_{1},w_{2},b$ yang memisahkan, dengan label $1$ pada sisi positif. Label $0$
di $(0,0)$ memberi $b\le0$. Label $1$ di $(1,0)$ memberi $w_{1}+b>0$, sehingga
$w_{1}>-b\ge0$. Dengan jalan yang sama, $w_{2}>-b\ge0$. Karena itu $w_{1}+w_{2}>-2b$, dan
\[
  w_{1}+w_{2}+b>-b\ge0 .
\]
Tetapi label $0$ di $(1,1)$ menuntut $w_{1}+w_{2}+b\le0$. Kontradiksi.
\end{proof}

Kontradiksi ini berlaku untuk setiap pilihan $\mathbf{w}$ dan $b$, bukan hanya untuk hasil
pelatihan tertentu. Menambah contoh, mengganti laju belajar, atau menormalkan fitur tidak
akan menolong. Yang harus berubah adalah bentuk fungsinya, dan perubahan terkecil yang
tersedia adalah menyisipkan satu lapisan di antaranya.

\kamus{lapisan tersembunyi}{hidden layer}{Lapisan unit yang tidak dilihat dunia luar:
masukannya dari lapisan sebelumnya, keluarannya ke lapisan berikutnya.}

Lapisan tersembunyi bekerja dengan dua langkah. Pertama, hitung
$\mathbf{z}_{1}=\mathbf{W}_{1}\mathbf{x}+\mathbf{b}_{1}$ dengan
$\mathbf{W}_{1}\in\mathbb{R}^{2\times4}$, lalu terapkan
$\operatorname{ReLU}(t)=\max\{t,0\}$ pada setiap komponen, dan tulis hasilnya $\mathbf{h}$.
Kedua, hitung skor linear dari $\mathbf{h}$. Untuk XOR, dengan
\[
  \mathbf{W}_{1}=\begin{pmatrix}1&1\\1&1\end{pmatrix},\qquad
  \mathbf{b}_{1}=\begin{pmatrix}0\\-1\end{pmatrix},\qquad
  \mathbf{w}_{2}=\begin{pmatrix}1\\-2\end{pmatrix},\qquad b_{2}=0,
\]
keempat titik memberi: $(0,0)$ menghasilkan $\mathbf{z}_{1}=(0,-1)$ dan
$\mathbf{h}=(0,0)$ dengan skor $0$; titik $(1,0)$ dan $(0,1)$ keduanya menghasilkan
$\mathbf{z}_{1}=(1,0)$ dan $\mathbf{h}=(1,0)$ dengan skor $1$; titik $(1,1)$ menghasilkan
$\mathbf{z}_{1}=(2,1)$ dan $\mathbf{h}=(2,1)$ dengan skor $1(2)-2(1)=0$. Dengan ambang
$0{,}5$ keempat titik terklasifikasi benar.

Yang membuat ini bekerja bukan sihir, melainkan aritmetika sederhana. Unit pertama menghitung
$\max\{x_{1}+x_{2},0\}$ dan unit kedua menghitung $\max\{x_{1}+x_{2}-1,0\}$. Keduanya
linear, tetapi selisih keduanya adalah fungsi potongan yang berbelok di dua tempat.
Kombinasi linear fungsi linear tetap linear; kombinasi linear fungsi potongan tidak.

Batas seperti itu tersusun dari ruas garis. Setiap unit di lapis pertama menyumbang satu
garis lurus, dan lapis kedua memilih sisi mana yang dipakai, sehingga batas akhirnya
berupa poligon; Gambar~\ref{fig:dua-lapis} menunjukkan bentuknya.

\begin{figure}[htbp]
\centering
\begin{tikzpicture}[>=latex, font=\small]
  % panel kiri: bagan dua lapis tersembunyi
  \foreach \i/\yy in {1/1.5,2/-0.5}
    \node[circle,draw=layergrey,fill=layergrey!10,inner sep=1.7pt] (I\i) at (0.4,\yy) {};
  \foreach \i/\yy in {1/2.4,2/1.65,3/0.9,4/0.15,5/-0.6}
    \node[circle,draw=layerteal,fill=layerteal!12,inner sep=1.7pt] (H1\i) at (2.1,\yy) {};
  \foreach \i/\yy in {1/1.6,2/0.7,3/-0.2}
    \node[circle,draw=layerteal,fill=layerteal!12,inner sep=1.7pt] (H2\i) at (3.7,\yy) {};
  \node[circle,draw=layerindigo,fill=layerindigo!12,inner sep=1.7pt] (O) at (5.1,0.7) {};
  \foreach \i in {1,2} \foreach \j in {1,...,5} \draw[layergrey!45] (I\i) -- (H1\j);
  \foreach \i in {1,...,5} \foreach \j in {1,2,3} \draw[layergrey!45] (H1\i) -- (H2\j);
  \foreach \j in {1,2,3} \draw[layergrey!45] (H2\j) -- (O);
  \node[font=\scriptsize] at (0.4,-1.35) {masukan};
  \node[font=\scriptsize] at (2.1,-1.35) {lima unit};
  \node[font=\scriptsize] at (3.7,-1.35) {tiga unit};
  \node[font=\scriptsize] at (5.1,-1.35) {keluaran};
  % panel kanan: bidang masukan
  \draw[layergrey] (6.7,-0.9) rectangle (14.4,2.9);
  \node[font=\scriptsize] at (14.1,-1.35) {$x_{1}$};
  \node[font=\scriptsize] at (6.4,2.75) {$x_{2}$};
  \fill[layerorange!18] (8.3,0.3) -- (9.7,2.4) -- (12.6,2.55) -- (13.3,1.1) -- (11.2,0.0) -- cycle;
  \draw[layerorange,thick] (8.3,0.3) -- (9.7,2.4) -- (12.6,2.55) -- (13.3,1.1) -- (11.2,0.0) -- cycle;
  \draw[dashed,layergrey!70] (7.50,-0.9) -- (11.43,2.9);
  \draw[dashed,layergrey!70] (6.70,2.23) -- (14.40,2.64);
  \draw[dashed,layergrey!70] (12.43,2.9) -- (14.27,-0.9);
  \draw[dashed,layergrey!70] (14.40,1.68) -- (9.48,-0.9);
  \draw[dashed,layergrey!70] (6.70,0.50) -- (14.40,-0.33);
  \foreach \xx/\yy in {9.9/1.6, 11.4/1.3, 10.6/0.6}
    \draw[accent,thick] (\xx-0.09,\yy-0.09) -- (\xx+0.09,\yy+0.09)
      (\xx-0.09,\yy+0.09) -- (\xx+0.09,\yy-0.09);
  \foreach \xx/\yy in {7.7/1.4, 8.5/-0.3, 12.3/-0.5, 13.6/2.2}
    \draw[layergrey,fill=white] (\xx,\yy) circle (2.1pt);
  \draw[accent,thick] (11.0,-1.30) -- (11.15,-1.15) (11.0,-1.15) -- (11.15,-1.30);
  \node[font=\scriptsize,anchor=west] at (11.2,-1.23) {kelas 1};
  \draw[layergrey,fill=white] (13.0,-1.23) circle (2.1pt);
  \node[font=\scriptsize,anchor=west] at (13.2,-1.23) {kelas 0};
\end{tikzpicture}
\caption{Skema bagaimana dua lapis tersembunyi mengubah satu garis menjadi lipatan.
Kiri: setiap unit di lapis pertama menarik satu garis lurus pada bidang masukan, dan
lapis kedua menggabungkan sisi-sisi garis itu. Kanan: batas keputusannya berupa poligon
lima sisi, digambar tebal; garis putus-putus adalah perpanjangan tiap sisi, yaitu garis
milik satu unit. Tanda silang adalah kelas 1 di dalam poligon, bulatan kelas 0 di luar.
Gambar ini skema, bukan hasil latihan: banyaknya sisi mengikuti banyaknya unit di lapis
pertama.}
\label{fig:dua-lapis}
\end{figure}


\section{Fungsi aktivasi dan dua penyakit lapisan dalam}
\tingkat{menengah}

Tanpa fungsi taklinear di antara lapisan, seluruh jaringan runtuh menjadi satu matriks. Bila
ReLU diganti identitas, keluarannya $\mathbf{W}_{2}\mathbf{W}_{1}\mathbf{x}$, yaitu
perkalian matriks tunggal $\mathbf{W}\mathbf{x}$ dengan
$\mathbf{W}=\mathbf{W}_{2}\mathbf{W}_{1}$. Karena itu setiap lapisan tersembunyi memerlukan
aktivasi taklinear, dan pilihan aktivasinya menentukan penyakit apa yang akan muncul saat
jaringan dibuat dalam.

\kamus{fungsi aktivasi}{activation function}{Pintu yang menentukan seberapa besar sinyal
sebuah unit diteruskan, dan sekaligus penentu bentuk turunannya.}

\begin{definisi}[ReLU, sigmoid, dan tanh]
Untuk $t\in\mathbb{R}$,
\[
  \operatorname{ReLU}(t)=\max\{t,0\},\qquad
  \sigma(t)=\frac{1}{1+e^{-t}},\qquad
  \tanh(t)=\frac{e^{t}-e^{-t}}{e^{t}+e^{-t}} .
\]
Turunan ReLU bernilai $1$ untuk $t>0$ dan $0$ untuk $t<0$; di $t=0$ turunannya tidak ada,
dan buku ini memakai nilai $0$ sebagai konvensi. Turunan tanh memenuhi
$\tanh'(t)=1-\tanh^{2}(t)$.\index{fungsi aktivasi}\index{relu}
\end{definisi}

Nilai $0$ di titik patah ReLU adalah konvensi: himpunan tempat turunan gagal hanya $\{0\}$, berukuran nol, sehingga tidak memengaruhi turunan di hampir semua tempat. Akibat praktisnya: unit yang masukannya selalu negatif menerima gradien nol dan tidak akan pernah berubah.

Sigmoid punya sifat yang lebih rapi, dan sifat itulah sumber masalahnya.

\begin{lemma}[Turunan sigmoid]
Untuk setiap $t$ berlaku $\sigma'(t)=\sigma(t)\bigl(1-\sigma(t)\bigr)$, sehingga
$0<\sigma'(t)\le\tfrac14$.\index{saturasi}
\end{lemma}

\begin{proof}
Tulis $\sigma(t)=(1+e^{-t})^{-1}$. Maka
\[
  \sigma'(t)=-(1+e^{-t})^{-2}\cdot(-e^{-t})
  =\frac{e^{-t}}{(1+e^{-t})^{2}}
  =\frac{1}{1+e^{-t}}\cdot\frac{e^{-t}}{1+e^{-t}}
  =\sigma(t)\bigl(1-\sigma(t)\bigr).
\]
Karena $0<\sigma(t)<1$, hasil kalinya positif. Dengan menulis $s=\sigma(t)$ diperoleh
$s(1-s)=\tfrac14-(s-\tfrac12)^{2}\le\tfrac14$, dan kesamaan tercapai ketika $s=\tfrac12$.
\end{proof}

Batas $\tfrac14$ itu yang harus dipegang. Tabel berikut menghitung $\sigma$, $\sigma'$, dan
$\tanh$ pada beberapa nilai; angkanya dihasilkan NumPy dan ditulis sampai enam desimal.

\begin{center}
\small
\begin{tabular}{rrrrr}
\hline
$t$ & $\sigma(t)$ & $\sigma'(t)$ & $\tanh(t)$ & $1-\tanh^{2}(t)$ \\
\hline
$0$  & $0{,}500000$ & $0{,}250000$ & $0{,}000000$ & $1{,}000000$ \\
$1$  & $0{,}731059$ & $0{,}196612$ & $0{,}761594$ & $0{,}419974$ \\
$2$  & $0{,}880797$ & $0{,}104994$ & $0{,}964028$ & $0{,}070650$ \\
$5$  & $0{,}993307$ & $0{,}006648$ & $0{,}999909$ & $0{,}000181$ \\
$10$ & $0{,}999955$ & $0{,}000045$ & $0{,}999999$ & $0{,}000002$ \\
\hline
\end{tabular}
\end{center}

Kolom keempat dan kelima menunjukkan apa yang disebut saturasi: begitu $t$ menjauh dari nol,
keluaran hampir tidak berubah terhadap masukan, dan turunannya menuju nol. Pada $t=10$ turunannya tinggal $4{,}54\times10^{-5}$, sekitar seperdua puluh ribu dari maksimumnya.

Akibatnya bisa dihitung. Bila galat merambat mundur melalui $L$ lapisan berturut-turut yang
turunannya paling besar $\tfrac14$, sisa sinyalnya paling besar $4^{-L}$: untuk $L=5$ sekitar
$9{,}77\times10^{-4}$, dan untuk $L=10$ sekitar $9{,}54\times10^{-7}$. Lapisan yang jauh dari
keluaran menerima sinyal yang tinggal sepersejuta, sehingga praktis tidak belajar. Gejala ini
disebut gradien hilang. ReLU menghindarinya dengan alasan yang berbeda: turunannya $1$ di sisi
positif, jadi tidak ada penyusutan, tetapi turunannya $0$ di sisi negatif, jadi ada
kematian unit. Kedua penyakit itu bukan kesalahan implementasi; keduanya konsekuensi bentuk
fungsi.\index{gradien hilang}

Satu pernyataan lagi yang sering dikutip melampaui isinya: satu lapisan tersembunyi cukup untuk semua fungsi kontinu.

\begin{teorema}[Aproksimasi universal, pernyataan]
Misalkan $\varphi$ terbatas, naik monoton, dan memiliki limit berbeda di $-\infty$ dan
$+\infty$. Untuk setiap himpunan kompak $K\subset\mathbb{R}^{d}$, setiap fungsi kontinu
$f:K\to\mathbb{R}$, dan setiap $\varepsilon>0$, terdapat $N$ dan parameter
$\mathbf{w}_{i},b_{i},v_{i}$ sehingga
\[
  \sup_{\mathbf{x}\in K}\Bigl|f(\mathbf{x})-\sum_{i=1}^{N}v_{i}\,
  \varphi(\mathbf{w}_{i}^{\top}\mathbf{x}+b_{i})\Bigr|<\varepsilon .
\]
\end{teorema}

Bukti pernyataan ini tidak dimuat di sini; bukti aslinya ada pada Cybenko (1989) dan
versi yang lebih umum pada Hornik (1991), dan perluasan ke ReLU ada pada Leshno dkk. (1993).
Yang penting dibaca dengan hati-hati adalah apa yang \emph{tidak} dikatakan teorema itu. Ia
menyatakan bahwa himpunan fungsi semacam itu padat, bukan bahwa jaringan tertentu sama
dengan $f$. Ia tidak mengatakan berapa $N$ yang diperlukan untuk galat $\varepsilon$ pada dimensi $d$; dalam kasus terburuk $N$ tumbuh secara eksponensial terhadap $d$ (Goodfellow dkk., bagian 6.4.1). Ia juga tidak mengatakan bahwa parameter itu dapat ditemukan oleh penurunan gradien, dan tidak mengatakan apa pun tentang data yang belum dilihat.
Kalimat ``jaringan bisa menghampiri apa saja'' karena itu benar sebagai pernyataan
kepadatan dan menyesatkan sebagai janji praktis.

\section{Dari skor ke kerugian}
\tingkat{menengah}

Keluaran jaringan harus dibandingkan dengan label melalui satu bilangan yang dapat
diminimumkan, dan pilihan bilangan itu menentukan seberapa besar sinyal yang sampai ke
lapisan bawah. Untuk klasifikasi biner dengan label $y\in\{0,1\}$ dan keluaran
$a=\sigma(z_{2})$, kerugian yang dipakai adalah entropi silang biner
\[
  L(a,y)=-\bigl[y\ln a+(1-y)\ln(1-a)\bigr].
\]
Entropi silang adalah negatif logaritma kepercayaan yang diberikan model kepada kelas yang
benar, sehingga kesalahan yang percaya diri dihukum jauh lebih berat daripada kesalahan yang
ragu.

\kamus{entropi silang}{cross-entropy}{Ukuran kejutan: makin yakin model pada jawaban yang
salah, makin besar dendanya, dan dendanya tumbuh tanpa batas.}\index{entropi silang}

\begin{proposisi}[Gradien entropi silang]
Untuk $a=\sigma(z)$ dan $L$ entropi silang biner berlaku
$\dfrac{\partial L}{\partial z}=a-y$.
\end{proposisi}

\begin{proof}
Turunkan $L$ terhadap $a$:
\[
  \frac{\partial L}{\partial a}=-\frac{y}{a}+\frac{1-y}{1-a}
  =\frac{a-y}{a(1-a)} .
\]
Dengan aturan rantai dan lemma turunan sigmoid,
$\dfrac{\partial a}{\partial z}=a(1-a)$, sehingga
\[
  \frac{\partial L}{\partial z}=\frac{a-y}{a(1-a)}\cdot a(1-a)=a-y .
\]
Pembilang dan penyebut saling menghapus, dan itulah intinya.
\end{proof}

Bandingkan dengan kerugian kuadrat $L=\tfrac12(a-y)^{2}$, yang memberi
$\partial L/\partial z=(a-y)\,a(1-a)$. Faktor $a(1-a)$ tidak hilang; ia mengecil persis ketika model paling salah dengan paling percaya diri. Untuk $z=10$ dan $y=0$, $a(1-a)=4{,}5396\times10^{-5}$, jadi sinyal ke bobot hanya $4{,}5394\times10^{-5}$ walaupun galatnya hampir satu. Dengan entropi silang sinyalnya
tetap $a-y=0{,}999955$. Perbandingan kedua angka itu sekitar $2{,}2\times10^{4}$; inilah
sebabnya keluaran sigmoid dipasangkan dengan entropi silang, bukan dengan kuadrat.

Untuk $K$ kelas, pasangan yang sama berlaku, dan kali ini kita buktikan.

\kamus{softmax}{softmax}{Mengubah deretan skor menjadi pembagian peluang yang jumlahnya
satu, dengan cara membandingkan setiap skor terhadap jumlah semuanya.}\index{softmax}

\begin{proposisi}[Softmax dan entropi silang]
Misalkan $\mathbf{p}=\operatorname{softmax}(\mathbf{z})$ dengan
$p_{k}=e^{z_{k}}/\sum_{j}e^{z_{j}}$, dan $L=-\sum_{k}y_{k}\ln p_{k}$ dengan
$\sum_{k}y_{k}=1$. Maka $\dfrac{\partial L}{\partial z_{i}}=p_{i}-y_{i}$.
\end{proposisi}

\begin{proof}
Turunan komponen softmax adalah
$\dfrac{\partial p_{k}}{\partial z_{i}}=p_{k}(\delta_{ki}-p_{i})$, dengan $\delta_{ki}$
delta Kronecker. Maka
\[
  \frac{\partial L}{\partial z_{i}}
  =-\sum_{k}\frac{y_{k}}{p_{k}}\cdot p_{k}(\delta_{ki}-p_{i})
  =-\sum_{k}y_{k}\delta_{ki}+p_{i}\sum_{k}y_{k}
  =-y_{i}+p_{i},
\]
sebab $\sum_{k}y_{k}=1$.
\end{proof}

Untuk dua kelas, $\operatorname{softmax}((t,0))_{1}=\sigma(t)$, sehingga rumus ini kembali
menjadi $a-y$ seperti sebelumnya. Perhatikan pula bahwa menambahkan tetapan $c$ pada semua $z_{k}$ tidak mengubah
$\mathbf{p}$, karena faktor $e^{c}$ saling menghapus; hanya selisih skor yang bermakna,
sekaligus sebabnya penambahan tetapan besar pada semua logit merusak ketelitian floating-point
(Bab 22).

\begin{untukanak}
Pernahkah kamu main permainan menebak angka, dan setiap kali salah kamu hanya diberi tahu
``terlalu besar'' atau ``terlalu kecil'', tanpa diberi tahu seberapa jauh? Kalau tebakanmu
sembilan ratus dan jawabannya seratus, keterangan ``terlalu besar'' sama saja dengan kalau
tebakanmu seratus satu.

Jaringan saraf punya masalah yang sama, tetapi lebih rumit. Ia punya tiga belas tombol, dan
setiap tombol memengaruhi jawaban lewat beberapa perantara. Kalau kita hanya tahu jawabannya
salah, kita tidak tahu tombol mana yang harus diputar, ke arah mana, dan sejauh apa.

Untungnya, jawaban salah itu punya derajat: seberapa jauh, dan ke arah mana. Dengan satu ukuran yang disebut kerugian, kita bisa bertanya pada setiap tombol: kalau kamu diputar sedikit, apakah kerugiannya naik atau turun, dan seberapa cepat? Jawabannya satu angka per tombol.

Di bab ini kita akan tahu mengapa satu ukuran kerugian saja cukup untuk memberi tiga belas
jawaban sekaligus, dan mengapa jawaban itu bisa dipercaya.
\end{untukanak}

\section{Turunan balik}
\tingkat{lanjut}

Jaringan adalah komposisi fungsi, jadi turunannya dihitung dengan aturan rantai Bab 12. Yang
membuatnya murah adalah urutan pengerjaannya: dari keluaran mundur ke masukan, bukan dari
masukan maju ke keluaran. Untuk setiap besaran antara $u$ dalam graf komputasi, kita
tetapkan satu bilangan yang menyatakan seberapa besar kerugian berubah bila $u$ berubah.

\kamus{adjoint}{adjoint}{Bayangan seberapa besar kesalahan akhir terpengaruh oleh satu
besaran di tengah jaringan; satu angka per besaran.}

\begin{definisi}[Adjoint]
Misalkan $C$ kerugian pada satu contoh dan $u$ salah satu besaran antara. Adjoint $u$
adalah $\bar{u}=\partial C/\partial u$.\index{adjoint}
\end{definisi}

\begin{proposisi}[Rekursi adjoint]
Bila $u$ memengaruhi $C$ hanya melalui $v_{1},\dots,v_{m}$, maka
$\bar{u}=\sum_{i=1}^{m}\bar{v}_{i}\,\dfrac{\partial v_{i}}{\partial u}$.
\end{proposisi}

\begin{proof}
Ini aturan rantai untuk fungsi beberapa variabel (Bab 12) yang diterapkan pada $C$ sebagai
fungsi dari $v_{1},\dots,v_{m}$, dengan setiap $v_{i}$ sebagai fungsi dari $u$. Karena $u$
tidak memengaruhi $C$ dengan jalan lain, semua suku lain dalam aturan rantai itu nol.
\end{proof}

Rekursi itu mengubah pekerjaan menjadi satu kali telusur. Mulai dari $\bar{C}=1$, kita
mengunjungi besaran antara dalam urutan terbalik dari hitung maju. Setiap besaran menerima
sumbangan dari semua besaran yang memakainya, jadi tidak ada yang terlewat atau dihitung dua kali. Inilah turunan balik, dan ia bukan algoritme baru: ia hanya aturan rantai
yang dijalankan pada graf dengan urutan yang tepat.\index{turunan balik}

Gambar~\ref{fig:graf-balik} menaruh kedua arah itu pada satu graf. Panah penuh membawa
nilai maju; panah putus-putus membawa adjoint, dan angka pada panah itu diambil dari
contoh hitung di bagian berikutnya.

\begin{figure}[htbp]
\centering
\begin{tikzpicture}[>=latex, font=\small]
  \node[draw=layerteal,rounded corners=2pt,fill=layerteal!10,inner sep=4pt] (X)  at (0,0)    {$\mathbf{x}$};
  \node[draw=layerteal,rounded corners=2pt,fill=layerteal!10,inner sep=4pt] (Z1) at (2.7,0)  {$\mathbf{z}_{1}$};
  \node[draw=layerteal,rounded corners=2pt,fill=layerteal!10,inner sep=4pt] (H)  at (5.3,0)  {$\mathbf{h}$};
  \node[draw=layerteal,rounded corners=2pt,fill=layerteal!10,inner sep=4pt] (Z2) at (7.5,0)  {$z_{2}$};
  \node[draw=layerteal,rounded corners=2pt,fill=layerteal!10,inner sep=4pt] (A)  at (9.3,0)  {$a$};
  \node[draw=layerteal,rounded corners=2pt,fill=layerteal!10,inner sep=4pt] (L)  at (11.4,0) {$L$};
  \node (W1) at (2.7,1.45) {$\mathbf{W}_{1},\mathbf{b}_{1}$};
  \node (W2) at (7.5,1.45) {$\mathbf{w}_{2},b_{2}$};
  \foreach \a/\b in {X/Z1,Z1/H,H/Z2,Z2/A,A/L} \draw[->,thick,ink] (\a) -- (\b);
  \draw[->,thick,ink] (W1) -- (Z1);
  \draw[->,thick,ink] (W2) -- (Z2);
  \draw[->,dashed,accent] (L.south)  .. controls +(0,-1.0) and +(0,-1.0) .. (Z2.south);
  \draw[->,dashed,accent] (Z2.south) .. controls +(0,-1.0) and +(0,-1.0) .. (H.south);
  \draw[->,dashed,accent] (H.south)  .. controls +(0,-1.0) and +(0,-1.0) .. (Z1.south);
  \draw[->,dashed,accent] (Z1.south) .. controls +(0,-1.0) and +(0,-1.0) .. (X.south);
  \draw[->,dashed,accent] (3.05,0.25) -- (3.05,1.25);
  \draw[->,dashed,accent] (7.85,0.25) -- (7.85,1.25);
  \node[font=\scriptsize] at (9.9,-1.4)  {$\bar{z}_{2}=a-y=-0{,}2554$};
  \node[font=\scriptsize] at (6.5,-2.0)  {$\bar{\mathbf{h}}=(-0{,}2043;\,0{,}1532)$};
  \node[font=\scriptsize] at (4.0,-1.4)  {$\bar{\mathbf{z}}_{1}=\bar{\mathbf{h}}$};
  \node[font=\scriptsize,align=center] at (1.1,-1.95)
    {$\nabla_{\mathbf{x}}L=(-0{,}0868;\,0{,}1328;$\\$(-0{,}0971;\,-0{,}1430)$};
  \node[font=\scriptsize] at (2.7,2.15) {$\nabla_{\mathbf{W}_{1}}L=\bar{\mathbf{z}}_{1}\mathbf{x}^{\top}$};
  \node[font=\scriptsize] at (9.0,2.15) {$\nabla_{\mathbf{w}_{2}}L=\bar{z}_{2}\mathbf{h}$};
\end{tikzpicture}
\caption{Graf komputasi jaringan 4--2--1 pada contoh hitung. Panah tebal penuh adalah
arah maju: masukan dan bobot menjadi $\mathbf{z}_{1}$, gerbang ReLU menjadi
$\mathbf{h}$, lalu skor $z_{2}$, keluaran $a$, dan kerugian $L$. Panah putus-putus
adalah arah mundur yang membawa adjoint: $\bar{z}_{2}=a-y=-0{,}2554$ masuk ke bobot
keluaran, $\bar{\mathbf{h}}=(-0{,}2043;\,0{,}1532)$ masuk ke gerbang, dan
$\nabla_{\mathbf{x}}L=(-0{,}0868;\,0{,}1328;\,-0{,}0971;\,-0{,}1430)$ keluar di
masukan. Satu telusur mundur menghasilkan ketiga belas gradien sekaligus.}
\label{fig:graf-balik}
\end{figure}

\begin{bacakeras}
Dua suara. Suara pertama: ``Kerugiannya terlalu besar.'' Suara kedua: ``Siapa yang
salah?'' Suara pertama: ``Bukan satu orang. Yang salah di ujung, dan kesalahannya dibagi
ke belakang, sedikit demi sedikit.''

Dengar sekali lagi, pelan. Kesalahan lahir di ujung, pada angka terakhir. Dari sana ia
berjalan mundur, berhenti di setiap persimpangan, dan bertanya: lewat sini, berapa
bagianmu? Yang paling dekat ujung menjawab paling besar. Yang jauh di belakang menjawab
paling kecil, sebab jawabannya sudah disaring dua kali di jalan pulang.

Suara kedua: ``Jadi yang di depan tidak bersalah?'' Suara pertama: ``Ia bersalah sebesar
bagian yang sampai kepadanya. Tidak lebih. Karena itu kita tidak menuduh, kita
menghitung.''

Ulangi tiga kali sampai enak di telinga: salah di ujung, dibagi ke belakang, sedikit
demi sedikit, sesuai jalannya.
\end{bacakeras}


Sekarang turunkan seluruhnya untuk jaringan 4--2--1 dengan ReLU di lapisan tersembunyi dan
sigmoid di keluaran. Hitung majunya
\[
  \mathbf{z}_{1}=\mathbf{W}_{1}\mathbf{x}+\mathbf{b}_{1},\qquad
  \mathbf{h}=\operatorname{ReLU}(\mathbf{z}_{1}),\qquad
  z_{2}=\mathbf{w}_{2}^{\top}\mathbf{h}+b_{2},\qquad
  a=\sigma(z_{2}),
\]
lalu $L$ entropi silang biner dari $a$ dan $y$. Setiap mata rantai punya turunan lokalnya
sendiri, dan semuanya sudah kita punya:
\[
  \frac{\partial L}{\partial z_{2}}=a-y\ \text{(proposisi sebelumnya)},\qquad
  \frac{\partial z_{2}}{\partial h_{i}}=w_{2}[i],
\]
\[
  \frac{\partial h_{i}}{\partial z_{1}[i]}=\mathbb{1}\bigl[z_{1}[i]>0\bigr],\qquad
  \frac{\partial z_{1}[i]}{\partial W_{1}[i,j]}=x_{j}.
\]
Mengalikan mundur sepanjang rantai itu memberi, satu per satu,
\[
  \bar{z}_{2}=a-y,\qquad
  \nabla_{\mathbf{w}_{2}}L=\bar{z}_{2}\,\mathbf{h},\qquad
  \frac{\partial L}{\partial b_{2}}=\bar{z}_{2},
\]
lalu sinyal yang sampai ke lapisan tersembunyi
\[
  \bar{\mathbf{h}}=\bar{z}_{2}\,\mathbf{w}_{2},
  \qquad
  \bar{\mathbf{z}}_{1}=\bar{\mathbf{h}}\odot\mathbb{1}\bigl[\mathbf{z}_{1}>0\bigr],
\]
dengan $\odot$ hasil kali per komponen. Perhatikan bahwa $\bar{\mathbf{z}}_{1}$ adalah
$\bar{\mathbf{h}}$ yang dicukur oleh gerbang ReLU: komponen yang gerbangnya tertutup menjadi
tepat nol, bukan kecil. Setelah itu
\[
  \nabla_{\mathbf{W}_{1}}L=\bar{\mathbf{z}}_{1}\mathbf{x}^{\top},\qquad
  \nabla_{\mathbf{b}_{1}}L=\bar{\mathbf{z}}_{1},\qquad
  \nabla_{\mathbf{x}}L=\mathbf{W}_{1}^{\top}\bar{\mathbf{z}}_{1}.
\]
Bentuk hasil kali luar pada $\nabla_{\mathbf{W}_{1}}L$ bukan kebetulan: turunan
$\partial z_{1}[i]/\partial W_{1}[i,j]=x_{j}$ tidak bergantung pada $i$, sehingga seluruh
baris ke-$i$ adalah $\bar{z}_{1}[i]$ dikali vektor $\mathbf{x}$. Setiap bobot menerima
gradien yang besarnya sebanding dengan masukan yang dikalikannya, dan itu sebabnya fitur
berskala besar otomatis menerima langkah besar.

Tidak satu pun rumus itu bergantung pada $d=4$ atau pada dua unit tersembunyi; polanya sama untuk lapisan mana pun.

Biayanya sebanding: satu langkah maju memerlukan satu perkalian matriks-vektor per lapisan, satu langkah mundur satu perkalian transpos dan satu hasil kali luar per lapisan. Selisih hingga terpusat untuk $P$ parameter memerlukan $2P$ langkah maju, sedangkan turunan balik hanya sekitar dua. Untuk $P=13$ di bab ini, dua puluh enam lawan dua; untuk sejuta parameter, dua juta lawan dua, dan di situlah seluruh bedanya. Nielsen menamakan pengamatan ini sebagai alasan mengapa turunan balik dianggap
penemuan besar, bukan sekadar trik aljabar.

\begin{tangan}
Ambil tiga kertas dan sebatang pensil, lalu beri label di atasnya: MASUKAN, BOBOT,
KELUARAN. Di kertas MASUKAN tulis empat angka: $1$; $0{,}5$; $-1$; $2$. Di kertas BOBOT
tulis dua angka bobot ke keluaran: $0{,}8$ dan $-0{,}6$. Di kertas KELUARAN tulis
$0{,}7446$, peluang yang diberikan jaringan, dan ingat label benarnya $1$.

Letakkan ketiga kertas berjajar dari kiri ke kanan: MASUKAN, BOBOT, KELUARAN. Sekarang
bekerja mundur, selalu dari kanan ke kiri. Di kertas KELUARAN, kurangkan $0{,}7446$ dari
$1$ dan tulis $-0{,}2554$. Geser ke kertas BOBOT: kalikan $-0{,}2554$ dengan $0{,}8$,
tulis $-0{,}2043$; kalikan dengan $-0{,}6$, tulis $0{,}1532$. Geser ke kertas MASUKAN:
kalikan $-0{,}2043$ dengan tiap angka di kertas itu, lalu tulis barisnya: $-0{,}2043$;
$-0{,}1022$; $0{,}2043$; $-0{,}4086$.

Yang harus kamu lihat: setiap kali kertas bergeser satu langkah ke kiri, tanda angkanya
bisa berbalik, dan besarnya mengikuti isi kertas yang sedang kamu pegang. Kalau arahnya
kamu balik, dari MASUKAN ke KELUARAN, urutan pengerjaannya salah dan hasilnya tidak
berarti apa-apa.
\end{tangan}


\section{Inisialisasi dan penurunan gradien}
\tingkat{menengah}

Setelah gradien ada, langkah berikutnya sederhana: bergerak sedikit ke arah yang menurunkan
kerugian. Untuk seluruh parameter $\boldsymbol{\theta}$,
\[
  \boldsymbol{\theta}\leftarrow\boldsymbol{\theta}-\eta\,\nabla_{\boldsymbol{\theta}}L,
\]
dengan $\eta$ laju belajar. Untuk $B$ contoh dan kerugian rata-rata
$L=\frac{1}{B}\sum_{i}L_{i}$, gradiennya juga rata-rata gradien per contoh; rumus
turunan balik tidak berubah, hanya berapa kali telusur mundur.
Syarat dan batas $\eta$ ada di Bab~23; karena jaringan tak
konveks, jaminan di sana pembanding.

\kamus{penurunan gradien stokastik}{stochastic gradient descent}{Menuruni lereng dengan
memakai tebakan lereng dari sebagian kecil data, bukan dari seluruh data.}\index{penurunan gradien}

Tiga varian perlu disebut. SGD mengambil $B$ kecil dan
memperbarui parameter setiap kali; gradiennya berderau, dan simpangan derau itu mengecil
seperti $1/\sqrt{B}$ karena variansi rata-rata $B$ peubah bebas adalah $\sigma^{2}/B$.
Momentum menyimpan kecepatan,
\[
  \mathbf{v}\leftarrow\beta\mathbf{v}+\mathbf{g},\qquad
  \boldsymbol{\theta}\leftarrow\boldsymbol{\theta}-\eta\mathbf{v},
\]
sehingga arah yang konsisten menumpuk dan arah yang berayun saling menghapus. Adam
menormalkan setiap koordinat dengan taksiran momen keduanya,
\[
  \mathbf{m}\leftarrow\beta_{1}\mathbf{m}+(1-\beta_{1})\mathbf{g},\qquad
  \mathbf{v}\leftarrow\beta_{2}\mathbf{v}+(1-\beta_{2})\mathbf{g}^{2},
\]
\[
  \hat{\mathbf{m}}=\frac{\mathbf{m}}{1-\beta_{1}^{t}},\qquad
  \hat{\mathbf{v}}=\frac{\mathbf{v}}{1-\beta_{2}^{t}},\qquad
  \boldsymbol{\theta}\leftarrow\boldsymbol{\theta}
  -\eta\,\frac{\hat{\mathbf{m}}}{\sqrt{\hat{\mathbf{v}}}+\epsilon},
\]
dengan $t$ nomor langkah. Faktor $1-\beta_{1}^{t}$ dan $1-\beta_{2}^{t}$ disebut koreksi
bias, dan gunanya menghapus kecenderungan momen menuju nol pada langkah-langkah awal. Pada
langkah pertama, $\hat{m}_{1}=g$ dan $\hat{v}_{1}=g^{2}$, sehingga pembaruannya tepat
$-\eta\,g/(\lvert g\rvert+\epsilon)$: besarnya hampir tepat $\eta$, berapa pun besar
gradiennya. Sifat itu membuat Adam tidak peka skala, dan sekaligus menjelaskan kelemahannya:
karena langkahnya tidak mengecil ketika gradien mengecil, Adam tidak dijamin konvergen pada
semua masalah konveks, dan contoh penyangkalnya bisa dibaca pada Reddi, Kale, dan Kumar
(2018).

Yang tersisa, dan yang paling sering diremehkan, adalah titik awal.

\begin{proposisi}[Simetri yang tidak pernah pecah]
Bila dua baris $\mathbf{W}_{1}$ identik dan $b_{1}[1]=b_{1}[2]$ pada awal pelatihan, maka
keduanya tetap identik setelah setiap pembaruan gradien, untuk setiap himpunan data dan
setiap laju belajar.\index{inisialisasi}
\end{proposisi}

\begin{proof}
Kita buktikan dengan induksi atas langkah pembaruan. Andaikan baris pertama dan kedua
$\mathbf{W}_{1}$ sama dan $b_{1}[1]=b_{1}[2]$. Maka untuk setiap masukan
$z_{1}[1]=z_{1}[2]$, sehingga kedua gerbang ReLU memberi keputusan yang sama dan
$\bar{z}_{1}[1]=\bar{z}_{1}[2]$. Karena
$\nabla_{\mathbf{W}_{1}}L=\bar{\mathbf{z}}_{1}\mathbf{x}^{\top}$, kedua baris gradien itu
identik, begitu pula kedua komponen $\nabla_{\mathbf{b}_{1}}L$. Pembaruan dengan laju yang
sama mempertahankan kesamaan. Untuk sekumpulan contoh, gradien total adalah jumlah gradien
per contoh, jadi kesamaan juga bertahan.
\end{proof}

Akibatnya tegas: delapan unit yang diinisialisasi seragam berperilaku seperti satu unit, seberapa pun lama ia dilatih. 

Inisialisasi nol lebih buruk lagi. Bila seluruh $\mathbf{W}_{1}$ dan $\mathbf{b}_{1}$ nol,
maka $\mathbf{z}_{1}=\mathbf{0}$, dan karena gerbangnya $\mathbb{1}[z>0]$, semua gerbang
tertutup: gradien terhadap $\mathbf{W}_{1}$ tepat nol, dan bobotnya tetap nol selamanya.
Yang masih bisa belajar hanya $b_{2}$, dan hasilnya adalah prediksi tetap sebesar proporsi
kelas. Pada data lab, kerugian akhirnya $0{,}690923$, tepat sama dengan entropi label. Delapan unit, nol manfaat.

\begin{tantangan}
Rancang jaringan 2--2--1 dengan ReLU di lapis tersembunyi yang gradiennya tepat nol pada
keenam bobot dan kedua bias lapis pertama sejak langkah pertama, sedangkan
$\mathbf{W}_{1}$ tidak nol, dan jaringannya tetap salah pada sedikitnya satu titik.
Hanya $b_{2}$ yang boleh bergerak. Tunjukkan mengapa pilihanmu bekerja.

Tulis jawabanmu sebagai tabel: unit, $z_{1}[i]$ untuk setiap contoh, keadaan gerbangnya, dan gradien yang masuk.

Petunjuk 1: yang membekukan bobot bukan nilai bobotnya, melainkan gerbang ReLU yang menutup untuk semua contoh.

Petunjuk 2: pakai data di kuadran tak negatif, lalu ambil
$\mathbf{W}_{1}=-\mathbf{1}\mathbf{1}^{\top}$ dan $\mathbf{b}_{1}=\mathbf{0}$. Maka
$z_{1}[i]=-(x_{1}+x_{2})\le0$, sehingga $\mathbf{h}=\mathbf{0}$,
$\bar{\mathbf{z}}_{1}=\mathbf{0}$, dan $\nabla_{\mathbf{W}_{1}}L=\mathbf{0}$ walaupun
$\mathbf{W}_{1}\neq\mathbf{0}$.

Kalau mentok, kembali ke bagian Inisialisasi dan penurunan gradien.
\end{tantangan}


Karena itu inisialisasi acak bukan hiasan, melainkan syarat agar unit-unit itu menjelajahi
fungsi yang berbeda. Pertanyaan berikutnya: seberapa besar. Jawabannya berasal dari
menghitung variansi. Bila komponen $\mathbf{x}$ saling bebas dengan rata-rata nol dan
variansi $\sigma_{x}^{2}$, dan bobot saling bebas dengan variansi $\sigma_{w}^{2}$, maka
\[
  \operatorname{Var}\bigl(z_{1}[i]\bigr)
  =\sum_{j=1}^{n_{\mathrm{in}}}\operatorname{Var}\bigl(W_{1}[i,j]x_{j}\bigr)
  =n_{\mathrm{in}}\,\sigma_{w}^{2}\sigma_{x}^{2}.
\]
Agar variansi itu tidak bergantung pada $n_{\mathrm{in}}$, kita butuh
$\sigma_{w}^{2}\propto1/n_{\mathrm{in}}$. Inisialisasi He memakai $\sigma_{w}^{2}=2/n_{\mathrm{in}}$,
dan angka $2$ itu bukan hiasan: untuk $z$ yang simetris, ReLU memotong separuh sebaran,
sehingga $\mathbb{E}[h^{2}]=\operatorname{Var}(z)/2=\sigma_{x}^{2}$, dan variansi itu
bertahan lapisan demi lapisan. Inisialisasi Xavier memakai
$\sigma_{w}^{2}=2/(n_{\mathrm{in}}+n_{\mathrm{out}})$, kompromi antara menjaga sinyal maju
dan sinyal mundur. Kompromi itu tepat hanya bila $n_{\mathrm{in}}=n_{\mathrm{out}}$.

Untuk lapisan sigmoid, inisialisasi yang terlalu besar langsung menghukum: pada jaringan dua lapis dengan simpangan $1{,}0$, kerugian awal mencapai $2{,}223899$, sedangkan dengan simpangan $0{,}5$ hanya $1{,}220409$. Kebalikannya lebih gawat: dengan simpangan $0{,}5$ dan delapan lapisan, pelatihan berhenti pada akurasi $0{,}5333$.
Inisialisasi bukan detail administratif; ia menentukan apakah sinyal masih ada ketika
sampai ke lapisan pertama.

\section{Contoh hitung}
\tingkat{menengah}

\begin{resep}
\begin{enumerate}
  \item Mulai dari ujung: $\bar{z}_{2}=a-y$ untuk sigmoid dengan entropi silang, atau
        $\mathbf{p}-\mathbf{y}$ untuk softmax. Tanpa angka ini tidak ada yang bisa dibagi.
  \item Bobot keluaran: $\nabla_{\mathbf{w}_{2}}L=\bar{z}_{2}\mathbf{h}$ dan
        $\partial L/\partial b_{2}=\bar{z}_{2}$.
  \item Sinyal ke lapis sebelumnya: $\bar{\mathbf{h}}=\bar{z}_{2}\mathbf{w}_{2}$, lalu
        cukur dengan gerbang
        $\bar{\mathbf{z}}_{1}=\bar{\mathbf{h}}\odot\mathbb{1}[\mathbf{z}_{1}>0]$.
  \item Bobot lapis pertama: $\nabla_{\mathbf{W}_{1}}L=\bar{\mathbf{z}}_{1}\mathbf{x}^{\top}$
        dan $\nabla_{\mathbf{b}_{1}}L=\bar{\mathbf{z}}_{1}$; bila perlu gradien masukan
        $\nabla_{\mathbf{x}}L=\mathbf{W}_{1}^{\top}\bar{\mathbf{z}}_{1}$.
\end{enumerate}
Peringatan urutan: langkah 3 tidak boleh mendahului langkah 1. $\bar{\mathbf{z}}_{1}$
dihitung dari $\bar{z}_{2}$, dan gerbangnya memakai $\mathbf{z}_{1}$ dari hitung maju
yang sudah tersimpan. Menukar kedua langkah itu mematikan seluruh gradien bobot lapis
pertama.
\end{resep}


Kita pakai satu contoh dan menghitung semuanya dengan bilangan eksak. Masukannya
$\mathbf{x}=(1{,}0;\,0{,}5;\,-1{,}0;\,2{,}0)$, label $y=1$, dan parameter awalnya
\[
  \mathbf{W}_{1}=\begin{pmatrix}
  0{,}5&-0{,}5&0{,}25&1{,}0\\
  0{,}1&0{,}2&-0{,}3&0{,}4\end{pmatrix},\quad
  \mathbf{b}_{1}=\begin{pmatrix}0{,}1\\-0{,}2\end{pmatrix},\quad
  \mathbf{w}_{2}=\begin{pmatrix}0{,}8\\-0{,}6\end{pmatrix},\quad b_{2}=0{,}05 .
\]
Itu seluruh parameternya.

\emph{Langkah maju.} Unit pertama menerima
$0{,}5(1)-0{,}5(0{,}5)+0{,}25(-1)+1(2)+0{,}1=0{,}5-0{,}25-0{,}25+2+0{,}1=2{,}1$.
Unit kedua menerima
$0{,}1(1)+0{,}2(0{,}5)-0{,}3(-1)+0{,}4(2)-0{,}2=0{,}1+0{,}1+0{,}3+0{,}8-0{,}2=1{,}1$.
Keduanya positif, jadi $\mathbf{h}=(2{,}1;\,1{,}1)$ tanpa pemotongan. Skor keluarannya
$z_{2}=0{,}8(2{,}1)-0{,}6(1{,}1)+0{,}05=1{,}68-0{,}66+0{,}05=1{,}07$, sehingga
$a=\sigma(1{,}07)=0{,}7445969158$ dan kerugiannya $L=-\ln a=0{,}2949122596$.

\emph{Langkah mundur.} Adjoint keluaran $\bar{z}_{2}=a-y=-0{,}2554030842$. Dari situ
\[
  \nabla_{\mathbf{w}_{2}}L=\bar{z}_{2}\mathbf{h}
  =(-0{,}5363464768;\,-0{,}2809433926),\qquad
  \frac{\partial L}{\partial b_{2}}=-0{,}2554030842 .
\]
Sinyal ke lapisan tersembunyi adalah
$\bar{\mathbf{h}}=\bar{z}_{2}\mathbf{w}_{2}=(-0{,}2043224673;\,0{,}1532418505)$; kedua
gerbang terbuka, jadi $\bar{\mathbf{z}}_{1}=\bar{\mathbf{h}}$. Hasil kali luar memberi
\[
  \nabla_{\mathbf{W}_{1}}L=
  \begin{pmatrix}
  -0{,}2043224673&-0{,}1021612337&0{,}2043224673&-0{,}4086449347\\
  0{,}1532418505&0{,}0766209253&-0{,}1532418505&0{,}3064837010
  \end{pmatrix},
\]
dan $\nabla_{\mathbf{b}_{1}}L=(-0{,}2043224673;\,0{,}1532418505)$. Dua pemeriksaan cepat: baris kedua tepat $-0{,}75$ kali baris pertama, karena $-0{,}6/0{,}8=-0{,}75$; dan setiap kolom sebanding dengan $x_{j}$, sehingga kolom keempat dua kali kolom pertama, kolom kedua setengahnya, dan kolom ketiga negatifnya. Norma Euclidean gradien seluruhnya
$0{,}9796964023$.

Satu langkah penurunan gradien dengan $\eta=0{,}5$ memperbarui setiap parameter menurut
$\theta\leftarrow\theta-\eta\,\partial L/\partial\theta$, memberi
\[
  \mathbf{W}_{1}'=\begin{pmatrix}
  0{,}6021612337&-0{,}4489193832&0{,}1478387663&1{,}2043224673\\
  0{,}0233790747&0{,}1616895374&-0{,}2233790747&0{,}2467581495
  \end{pmatrix},
\]
\[
  \mathbf{b}_{1}'=(0{,}2021612337;\,-0{,}2766209253),\qquad
  \mathbf{w}_{2}'=(1{,}0681732384;\,-0{,}4595283037),
\]
\[
  b_{2}'=0{,}1777015421 .
\]
Setelah pembaruan itu $\mathbf{z}_{1}'=(2{,}8406689;\,0{,}5444983)$ dan
$a'=0{,}9508189702$, sehingga kerugiannya turun menjadi $0{,}0504315919$; selisihnya
$-0{,}2444806677$, atau sekitar $83\%$ dari kerugian awal.

Angka $83\%$ itu mengundang salah baca. Yang diukur hanya satu contoh, tanpa data uji, pada lanskap yang tidak konveks. Kesimpulan yang sah bukan ``laju $0{,}5$ bagus'', melainkan ``arah gradien ini menurunkan kerugian pada contoh ini''.

\kertas{Periksa satu baris tabel lab dengan tangan. Ambil $W_{1}[2,1]$ dan hitung
$\bar{z}_{1}[2]\,x_{1}$ untuk mendapat $0{,}1532418505$. Lalu hitung sendiri selisih hingga
terpusat untuk parameter itu: jalankan satu langkah maju dengan $W_{1}[2,1]+10^{-6}$ dan satu
langkah maju dengan $W_{1}[2,1]-10^{-6}$, bagi selisih kerugiannya dengan $2\times10^{-6}$,
dan bandingkan dengan angka analitik tadi. Setelah itu jelaskan mengapa hasilnya tidak sama
persis, dan mengapa tanda selisihnya bisa berubah bila $h$ diperkecil menjadi $10^{-12}$.}

\section{Lab komputasi}
\tingkat{lanjut}

Ketiga blok dijalankan dengan Python 3.14.4 dan NumPy 2.5.2, tanpa pustaka pembelajaran
mesin; keluaran disalin apa adanya.

Blok pertama memeriksa contoh tangan di atas, lalu membandingkan setiap gradien analitik
dengan selisih hingga terpusat pada $h=10^{-6}$.

\begin{lab}
\begin{verbatim}
import numpy as np
np.set_printoptions(precision=10, suppress=True, linewidth=76)

# jaringan 4-2-1: ReLU tersembunyi, sigmoid keluaran
x  = np.array([1.0, 0.5, -1.0, 2.0])
y  = 1.0
W1 = np.array([[0.5, -0.5, 0.25, 1.0],
               [0.1,  0.2, -0.3,  0.4]])
b1 = np.array([0.1, -0.2])
w2 = np.array([0.8, -0.6])
b2 = 0.05

def forward(x, W1, b1, w2, b2):
    z1 = W1 @ x + b1
    h  = np.maximum(z1, 0.0)
    o  = w2 @ h + b2
    p  = 1.0/(1.0+np.exp(-o))
    return z1, h, o, p, -(y*np.log(p) + (1-y)*np.log(1-p))

def backward(x, W1, b1, w2, b2):
    z1, h, o, p, L = forward(x, W1, b1, w2, b2)
    d   = p - y                    # dL/dz2 untuk sigmoid + entropi silang
    dz1 = d*w2*(z1 > 0)            # gerbang ReLU
    return dict(W1=np.outer(dz1, x), b1=dz1, w2=d*h, b2=d,
                x=W1.T @ dz1), L, p, z1, h, dz1

g, L, p, z1, h, dz1 = backward(x, W1, b1, w2, b2)
print("z1 =", z1, "  h =", h)
print("p = %.12f  L = %.12f" % (p, L))
print("dz1 =", dz1)
print("g_w2 =", g["w2"], " g_b2 = %.12f" % g["b2"])
for i in range(2):
    print("g_W1[%d] =" % i, np.array2string(g["W1"][i], precision=10,
                                            suppress_small=False))
print("g_x =", g["x"])

# pemeriksaan beda hingga terpusat, satu parameter sekali jalan
eps = 1e-6
def loss_at(nama, idx, delta):
    W, B, V, c = W1.copy(), b1.copy(), w2.copy(), b2
    if nama == "W1":  W[idx] += delta
    if nama == "b1":  B[idx] += delta
    if nama == "w2":  V[idx] += delta
    if nama == "b2":  c += delta
    return forward(x, W, B, V, c)[4]

daftar = [("W1", i, j) for i in range(2) for j in range(4)]
daftar += [("b1", i) for i in range(2)] + [("w2", i) for i in range(2)]
print("\nparameter     analitik        numerik     galat relatif")
for nama, *idx in daftar + [("b2",)]:
    idx = tuple(idx)
    ana = {"W1": g["W1"], "b1": g["b1"], "w2": g["w2"], "b2": np.array(g["b2"])}
    a = ana[nama][idx] if idx else ana[nama]
    n = (loss_at(nama, idx, eps) - loss_at(nama, idx, -eps))/(2*eps)
    rel = abs(a-n)/max(abs(a), abs(n), 1e-300)
    print("%-10s %15.10f %15.10f %12.2e" %
          (nama + "".join("[%d]" % i for i in idx), a, n, rel))

# satu langkah penurunan gradien dengan laju 0.5
eta = 0.5
W1n, b1n = W1 - eta*g["W1"], b1 - eta*g["b1"]
w2n, b2n = w2 - eta*g["w2"], b2 - eta*g["b2"]
z1b, hb, ob, pb, Lb = forward(x, W1n, b1n, w2n, b2n)
print("\nsatu langkah: p = %.12f  L = %.12f  selisih = %.12f"
      % (pb, Lb, L - Lb))
for i in range(2):
    print("W1'[%d] =" % i, np.array2string(W1n[i], precision=10))
print("w2' =", w2n, " b2' = %.12f" % b2n)

# gradien terhadap masukan: analitik vs beda hingga
E = np.eye(4)
gx = np.array([(forward(x + eps*E[j], W1, b1, w2, b2)[4]
              - forward(x - eps*E[j], W1, b1, w2, b2)[4])/(2*eps)
               for j in range(4)])
print("\ng_x analitik    =", g["x"])
print("g_x beda hingga =", gx)
print("galat relatif   =", ["%.1e" % v for v in np.abs(gx-g["x"])/np.abs(g["x"])])
\end{verbatim}
Keluaran sebenarnya:
\begin{verbatim}
z1 = [2.1 1.1]   h = [2.1 1.1]
p = 0.744596915827  L = 0.294912259605
dz1 = [-0.2043224673  0.1532418505]
g_w2 = [-0.5363464768 -0.2809433926]  g_b2 = -0.255403084173
g_W1[0] = [-0.2043224673 -0.1021612337  0.2043224673 -0.4086449347]
g_W1[1] = [ 0.1532418505  0.0766209253 -0.1532418505  0.306483701 ]
g_x = [-0.0868370486  0.1328096038 -0.097053172  -0.1430257271]

parameter     analitik        numerik     galat relatif
W1[0][0]     -0.2043224673   -0.2043224674     1.17e-10
W1[0][1]     -0.1021612337   -0.1021612337     2.53e-10
W1[0][2]      0.2043224673    0.2043224674     1.17e-10
W1[0][3]     -0.4086449347   -0.4086449347     4.95e-11
W1[1][0]      0.1532418505    0.1532418505     7.22e-11
W1[1][1]      0.0766209253    0.0766209252     8.33e-10
W1[1][2]     -0.1532418505   -0.1532418505     7.22e-11
W1[1][3]      0.3064837010    0.3064837010     1.09e-10
b1[0]        -0.2043224673   -0.2043224674     1.17e-10
b1[1]         0.1532418505    0.1532418506     6.16e-10
w2[0]        -0.5363464768   -0.5363464768     9.80e-11
w2[1]        -0.2809433926   -0.2809433925     2.41e-10
b2           -0.2554030842   -0.2554030842     3.59e-11

satu langkah: p = 0.950818970185  L = 0.050431591881  selisih = 0.244480667724
W1'[0] = [ 0.6021612337 -0.4489193832  0.1478387663  1.2043224673]
W1'[1] = [ 0.0233790747  0.1616895374 -0.2233790747  0.2467581495]
w2' = [ 1.0681732384 -0.4595283037]  b2' = 0.177701542086

g_x analitik    = [-0.0868370486  0.1328096038 -0.097053172  -0.1430257271]
g_x beda hingga = [-0.0868370486  0.1328096037 -0.097053172  -0.1430257272]
galat relatif   = ['7.7e-10', '5.0e-10', '1.4e-10', '9.8e-11']
\end{verbatim}
\end{lab}

Tabel itu adalah bukti yang diminta. Semua tiga belas parameter cocok sampai sembilan atau sepuluh digit, dengan galat relatif terbesar $8{,}33\times10^{-10}$ pada $W_{1}[1,2]$. Angka itu tidak seharusnya nol: selisih hingga terpusat membawa galat pemotongan orde $h^{2}$ dan galat pembulatan orde $\varepsilon_{\mathrm{mach}}/h$ (Bab 22), sehingga ada $h$ terbaik yang tidak dapat membuat galatnya nol. Untuk $h=10^{-6}$ orde $10^{-10}$ memang yang diharapkan.
 

Gradien terhadap masukan cocok sampai $7{,}7\times10^{-10}$ pada komponen pertama. Angka itu berarti: menaikkan fitur kedua sedikit menaikkan kerugian paling cepat, sedangkan menaikkan fitur keempat menurunkannya paling cepat. Di Bab 29 arah ini dipakai terbalik, untuk menaikkan
kerugian secepat mungkin pada masukan yang boleh diubah penyerang.

Blok kedua melatih jaringan 2--8--1 dari nol pada data dua dimensi berpola XOR, lalu
membandingkannya dengan garis dasar linear dan dengan dua inisialisasi yang cacat.

\begin{lab}
\begin{verbatim}
import numpy as np

rng = np.random.default_rng(7)
n = 240
X = rng.uniform(-1.0, 1.0, size=(n, 2))
y = (X[:, 0]*X[:, 1] > 0).astype(float)        # label 1 bila tanda sama
print("contoh =", n, " kelas 1 =", int(y.sum()))

def sig(o):
    return 1.0/(1.0+np.exp(-o))

def bce(p, t):
    p = np.clip(p, 1e-12, 1-1e-12)
    return float(-np.mean(t*np.log(p) + (1-t)*np.log(1-p)))

def init(mode, rs, nh=8):
    if mode == "nol":
        return np.zeros((nh, 2)), np.zeros(nh), np.zeros(nh), 0.0
    if mode == "seragam":
        return np.full((nh, 2), 0.5), np.zeros(nh), np.full(nh, 0.5), 0.0
    return (rs.normal(0, np.sqrt(2/2), (nh, 2)), np.zeros(nh),
            rs.normal(0, np.sqrt(2/nh), nh), 0.0)          # He

def fwd(X, W1, b1, w2, b2):
    Z1 = X @ W1.T + b1
    H  = np.maximum(Z1, 0.0)
    return Z1, H, sig(H @ w2 + b2)

def grads(X, t, W1, b1, w2, b2):
    Z1, H, P = fwd(X, W1, b1, w2, b2)
    d  = (P - t)/len(t)
    dZ = np.outer(d, w2)*(Z1 > 0)
    return dZ.T @ X, dZ.sum(0), H.T @ d, d.sum()

def latih(mode, langkah=3000, lr=0.5, seed=0):
    W1, b1, w2, b2 = init(mode, np.random.default_rng(seed))
    jejak = []
    for t in range(langkah + 1):
        gW1, gb1, gw2, gb2 = grads(X, y, W1, b1, w2, b2)
        W1, b1 = W1 - lr*gW1, b1 - lr*gb1
        w2, b2 = w2 - lr*gw2, b2 - lr*gb2
        if t % 500 == 0:
            jejak.append((t, bce(fwd(X, W1, b1, w2, b2)[2], y)))
    Z1, H, P = fwd(X, W1, b1, w2, b2)
    mati = int(np.sum(np.all(H <= 0, axis=0)))
    return jejak, float(np.mean((P > 0.5) == (y > 0.5))), bce(P, y), mati

jejak, akur, L, mati = latih("he")
for t, v in jejak:
    print("langkah %4d   L = %.6f" % (t, v))
print("He      : akurasi = %.4f  unit mati = %d/8" % (akur, mati))
for mode in ("nol", "seragam"):
    j, a, l, m = latih(mode)
    print("%-8s: L akhir = %.6f  akurasi = %.4f  unit mati = %d/8"
          % (mode, l, a, m))

w, b = np.zeros(2), 0.0             # garis dasar linear, tanpa lapisan
for t in range(4000):
    p = sig(X @ w + b)
    w -= 0.5*(X.T @ (p - y))/n
    b -= 0.5*np.mean(p - y)
pl = sig(X @ w + b)
print("linear  : L = %.6f  akurasi = %.4f"
      % (bce(pl, y), float(np.mean((pl > 0.5) == (y > 0.5)))))
\end{verbatim}
Keluaran sebenarnya:
\begin{verbatim}
contoh = 240  kelas 1 = 112
langkah    0   L = 0.648815
langkah  500   L = 0.060245
langkah 1000   L = 0.036445
langkah 1500   L = 0.027164
langkah 2000   L = 0.022039
langkah 2500   L = 0.018717
langkah 3000   L = 0.016349
He      : akurasi = 1.0000  unit mati = 0/8
nol     : L akhir = 0.690923  akurasi = 0.5333  unit mati = 8/8
seragam : L akhir = 0.539475  akurasi = 0.7167  unit mati = 0/8
linear  : L = 0.686130  akurasi = 0.7208
\end{verbatim}
\end{lab}

Kurva kerugiannya menurun pada setiap pencuplikan, dari $0{,}648815$ menjadi $0{,}016349$. Akurasi latih mencapai $1{,}0000$ dengan delapan unit dan nol unit mati. Garis dasar linear, yang
memakai optimasi dan data yang sama tetapi tanpa lapisan tersembunyi, berhenti pada
$0{,}7208$. Selisih $0{,}28$ itu harga satu-satunya bidang pemisah, dan ia sejalan dengan proposisi XOR: pola ini tidak terpisah secara linear.

Gambar~\ref{fig:kurva-kerugian} memplot ketujuh titik itu.

\begin{figure}[htbp]
\centering
\begin{tikzpicture}
\begin{axis}[
  width=12.0cm, height=5.6cm,
  xlabel={langkah pelatihan}, ylabel={kerugian $L$},
  xmin=-80, xmax=3080, ymin=0, ymax=0.75,
  axis lines=left,
  grid=major, grid style={layergrey!15},
  tick label style={font=\scriptsize},
  ytick={0,0.2,0.4,0.6}, yticklabels={$0$,$0{,}2$,$0{,}4$,$0{,}6$},
  label style={font=\small},
  legend style={font=\scriptsize, draw=none, fill=none,
                at={(0.98,0.98)}, anchor=north east},
]
\addplot[thick, accent, mark=*, mark size=1.7pt] coordinates {
  (0,0.648815) (500,0.060245) (1000,0.036445) (1500,0.027164)
  (2000,0.022039) (2500,0.018717) (3000,0.016349) };
\addplot[dashed, layergrey, domain=0:3000, samples=2] {0.693147};
\legend{kerugian latih, $\ln 2=0{,}693147$}
\end{axis}
\end{tikzpicture}
\caption{Kurva kerugian latih jaringan 2--8--1 pada blok kedua, dari $0{,}648815$ pada
langkah $0$ menjadi $0{,}016349$ pada langkah $3000$. Garis putus-putus adalah
$\ln 2=0{,}693147$, kerugian tebakan seragam; titik pertama sudah berada di bawahnya
karena diukur sesudah pembaruan pertama.}
\label{fig:kurva-kerugian}
\end{figure}


Inisialisasi nol memberi pelajaran yang lebih keras. Kerugian akhirnya $0{,}690923$, sama dengan entropi label sampai enam desimal, dan akurasinya $0{,}5333$, yaitu proporsi kelas mayoritas; kedelapan unit mati. Inisialisasi seragam memberi $0{,}7167$, praktis sama dengan garis dasar linear.

Blok ketiga mengukur dua hal: sampai berapa lapisan sigmoid masih belajar, dan seberapa cepat sinyal menyusut dari keluaran ke lapisan pertama.

\begin{lab}
\begin{verbatim}
import numpy as np

rng = np.random.default_rng(7)
n = 240
X = rng.uniform(-1.0, 1.0, size=(n, 2))
y = (X[:, 0]*X[:, 1] > 0).astype(float)

def sig(o):
    return 1.0/(1.0+np.exp(-o))

def bce(p, t):
    p = np.clip(np.ravel(p), 1e-12, 1-1e-12)
    t = np.ravel(t)
    return float(-np.mean(t*np.log(p) + (1-t)*np.log(1-p)))

def bangun(dalam, std, seed=0):
    rs = np.random.default_rng(seed)
    ukuran = [2] + [16]*(dalam-1) + [1]
    W = [rs.normal(0, std, (ukuran[i+1], ukuran[i])) for i in range(dalam)]
    return W, [np.zeros(ukuran[i+1]) for i in range(dalam)]

def maju(W, b, X):
    A = [X]
    for Wl, bl in zip(W[:-1], b[:-1]):
        A.append(sig(A[-1] @ Wl.T + bl))
    return A, sig(A[-1] @ W[-1].T + b[-1])

def mundur(W, b, X, t):
    A, P = maju(W, b, X)
    d = (P - t[:, None])/len(t)
    gW, gb = [None]*len(W), [None]*len(b)
    gW[-1], gb[-1] = d.T @ A[-1], d.sum(0)
    delta = d @ W[-1]
    for l in range(len(W)-1, 0, -1):
        delta = delta*A[l]*(1 - A[l])
        gW[l-1], gb[l-1] = delta.T @ A[l-1], delta.sum(0)
        if l > 1:
            delta = delta @ W[l-1]
    return gW, gb, bce(P, t), P

print("dalam  std   L(awal)   L(akhir)  akurasi  ||gW1||/||gWakhir||")
for dalam in (1, 2, 4, 8):
    for std in (0.5, 1.0):
        W, b = bangun(dalam, std)
        gW, gb, L0, P = mundur(W, b, X, y)
        rasio = np.linalg.norm(gW[0])/np.linalg.norm(gW[-1])
        for t in range(2000):
            gW, gb, L, P = mundur(W, b, X, y)
            for l in range(dalam):
                W[l] -= 0.5*gW[l]
                b[l] -= 0.5*gb[l]
        akur = float(np.mean((np.ravel(P) > 0.5) == (y > 0.5)))
        print("%5d %5.1f %9.6f %9.6f %8.4f %14.2e"
              % (dalam, std, L0, bce(P, y), akur, rasio))
\end{verbatim}
Keluaran sebenarnya:
\begin{verbatim}
dalam  std   L(awal)   L(akhir)  akurasi  ||gW1||/||gWakhir||
    1   0.5  0.693190  0.686130   0.7208       1.00e+00
    1   1.0  0.693951  0.686130   0.7208       1.00e+00
    2   0.5  1.220409  0.332913   0.8417       1.66e-02
    2   1.0  2.223899  0.299519   0.8542       2.16e-02
    4   0.5  0.691597  0.147887   0.9458       5.77e-02
    4   1.0  0.701977  0.023491   0.9958       2.05e-01
    8   0.5  0.739920  0.690867   0.5333       1.26e-04
    8   1.0  0.875314  0.002383   1.0000       3.76e-03
\end{verbatim}
\end{lab}

Tiga baris pertama tabel itu adalah janji dan peringatan sekaligus. Kedalaman satu memberi $0{,}7208$, sama dengan garis dasar linear di blok sebelumnya; kedalaman dua memberi $0{,}8417$ dan kedalaman empat $0{,}9458$. Tetapi baris terakhir dengan simpangan $0{,}5$ memberi $0{,}5333$ dengan kerugian $0{,}690867$:
delapan lapisan sigmoid dengan inisialisasi kecil praktis tidak belajar. Rasio norma gradien di kolom terakhir menjelaskan mengapa:
$1{,}26\times10^{-4}$, empat orde lebih kecil daripada pada kedalaman empat.
Baris yang sama dengan simpangan $1{,}0$ berhasil sempurna. Jadi kegagalan itu bukan soal kedalaman semata, melainkan soal ukuran awal.

Dua batas hasil lab ini harus disebut. Data hanya 240 contoh dua dimensi buatan sendiri, tanpa himpunan uji, jadi seluruhnya akurasi latih; generalisasi adalah urusan Bab 26 dan Bab 27. Semua percobaan memakai satu benih, sehingga variasi antar benih tidak terlihat. Yang boleh
disimpulkan hanya yang diperiksa langsung: turunan balik cocok dengan beda hingga, kerugian
menurun, inisialisasi nol dan seragam gagal dengan cara yang dapat dijelaskan, dan jaringan
dalam dengan inisialisasi kecil kehilangan sinyal.

\begin{untukdosen}
Urutan yang berhasil dipakai: mulai dari satu garis pemisah dan empat titik XOR, lalu tunjukkan
konstruksi 2--2--1 di papan, baru setelah itu definisi umum. Bukti ketidakmungkinan XOR biasanya diterima dengan mudah; yang menghentikan mahasiswa adalah langkah mundur, jadi sediakan waktu paling banyak di turunan balik, bukan di fungsi aktivasi.

Dua miskonsepsi khas. Pertama, ``turunan balik adalah algoritme lain, bukan aturan rantai'',
sehingga mahasiswa menghafal rumus tanpa bisa menurunkannya untuk arsitektur lain; obatnya,
minta mereka menurunkan ulang jaringan 4--2--1 dengan tiga unit tersembunyi di papan. Kedua,
``galat relatif $10^{-10}$ berarti implementasinya terbukti benar''; galat kecil hanya
memeriksa kecocokan pada satu titik, untuk satu himpunan parameter, dan tidak memeriksa tanda,
penskalaan, maupun kasus gerbang yang tertutup.

Aktivitas papan tulis yang paling menolong: beri satu jaringan dengan satu unit tersembunyi
yang $z_{1}<0$, minta kelas menghitung gradien seluruh baris $\mathbf{W}_{1}$, lalu tanyakan
apa yang terjadi pada unit itu setelah seratus langkah. Jawaban ``tidak terjadi apa-apa,
selamanya'' biasanya membuka diskusi tentang kematian unit dan tentang mengapa laju belajar
besar berbahaya.

Tanda kesiapan lanjut: mahasiswa dapat menuliskan turunan balik untuk jaringan dengan dua
lapisan tersembunyi tanpa melihat catatan, dan dapat menjelaskan mengapa inisialisasi nol
gagal bukan karena ``kurang data'' melainkan karena simetri.
\end{untukdosen}

\section{Ringkasan}

\begin{ringkasan}
\begin{itemize}
  \item Ambang linear hanya dapat memisahkan himpunan yang terpisah secara linear, dan empat
        titik XOR bukan salah satunya. Buktinya tidak bergantung pada algoritme, jadi
        menambah data tidak menolong. Satu lapisan tersembunyi dengan ReLU sudah cukup untuk
        pola itu.
  \item Tanpa aktivasi taklinear, jaringan berlapis runtuh menjadi satu matriks. Sigmoid dan
        tanh menyusutkan sinyal karena turunannya paling besar $\tfrac14$ dan $1$; melalui
        $L$ lapisan sigmoid, sisa sinyalnya paling besar $4^{-L}$, yaitu
        $9{,}54\times10^{-7}$ untuk $L=10$. ReLU tidak menyusutkan, tetapi mematikan unit
        yang masukannya selalu negatif.
  \item Entropi silang berpasangan dengan sigmoid atau softmax sehingga adjoint keluarannya
        tepat $a-y$ atau $\mathbf{p}-\mathbf{y}$. Dengan kerugian kuadrat, faktor $a(1-a)$
        tertinggal dan sinyalnya mengecil justru ketika model paling salah dengan yakin;
        pada $z=10$ perbandingannya sekitar $2{,}2\times10^{4}$.
  \item Turunan balik adalah aturan rantai Bab 12 yang dijalankan pada graf komputasi dengan
        urutan terbalik. Untuk 4--2--1, seluruh gradiennya:
        $\bar{z}_{2}=a-y$, $\nabla_{\mathbf{w}_{2}}L=\bar{z}_{2}\mathbf{h}$,
        $\bar{\mathbf{z}}_{1}=\bar{z}_{2}\mathbf{w}_{2}\odot\mathbb{1}[\mathbf{z}_{1}>0]$,
        $\nabla_{\mathbf{W}_{1}}L=\bar{\mathbf{z}}_{1}\mathbf{x}^{\top}$,
        $\nabla_{\mathbf{x}}L=\mathbf{W}_{1}^{\top}\bar{\mathbf{z}}_{1}$. Biayanya sekitar dua langkah maju, bukan $2P$. Pada contoh 4--2--1, ketiga belas gradien itu cocok dengan selisih hingga terpusat sampai galat relatif $8{,}33\times10^{-10}$.
  \item Inisialisasi menentukan apakah jaringan dapat belajar. Nol mematikan semua gerbang
        dan meninggalkan prediksi sebesar proporsi kelas; seragam mempertahankan simetri
        sehingga kapasitas efektifnya satu unit; terlalu kecil menyusutkan sinyal sampai
        lapisan pertama; terlalu besar menyaturasi sigmoid. He memakai
        $\sigma_{w}^{2}=2/n_{\mathrm{in}}$, Xavier $2/(n_{\mathrm{in}}+n_{\mathrm{out}})$.
  \item Penurunan gradien, SGD, momentum, dan Adam hanya berbeda pada cara memakai gradien
        yang sama. Langkah pertama Adam besarnya hampir tepat $\eta$ berapa pun gradiennya,
        dan sifat itu menjelaskan baik ketangguhannya maupun kegagalannya pada sebagian
        masalah.
  
\end{itemize}
\end{ringkasan}

\section{Latihan}

\begin{latihan}
\begin{enumerate}
  \item Untuk data dengan margin $\gamma=0{,}2$ dan $\lVert\mathbf{x}_{i}\rVert\le1$, berapa batas atas banyaknya pembaruan perceptron menurut teorema kekonvergenan, dan berapa pula untuk $\gamma=0{,}05$? Jelaskan mengapa batas itu tidak bergantung pada banyaknya contoh dan mengapa ia tidak berguna ketika $\gamma$ kecil.

  \item Buktikan bahwa jaringan dua lapis dengan aktivasi identitas di lapisan tersembunyi
        tidak lebih kuat daripada ambang linear, yaitu bahwa
        $\mathbf{W}_{2}\operatorname{id}(\mathbf{W}_{1}\mathbf{x}+\mathbf{b}_{1})+\mathbf{b}_{2}$
        dapat ditulis sebagai satu ambang linear. Jelaskan akibatnya untuk lapisan identitas berikutnya.

  

  \item Hitung $\nabla_{\mathbf{W}_{1}}L$ untuk jaringan pada contoh hitung bila $\mathbf{x}=(0{,}5;\,-1{,}0;\,0{,}25;\,2{,}0)$ dan $y=0$, dengan parameter yang sama. Tentukan unit mana yang gerbangnya tertutup, dan mengapa baris itu bergradien nol tanpa berarti kerugiannya minimum.

  \item Nilai klaim ini: \emph{karena galat relatif terbesar hanya
        $8{,}33\times10^{-10}$, implementasi turunan balik pada bab ini terbukti benar.}
        Tentukan benar, salah, atau belum cukup bukti, dan sebutkan pemeriksaan tambahan
        yang diperlukan.

  \item Nilai klaim ini: \emph{teorema aproksimasi universal menjamin bahwa jaringan dengan
        satu lapisan tersembunyi dapat mencapai galat sekecil apa pun pada data uji.}
        Tentukan benar, salah, atau belum cukup bukti, dan sebutkan dua hal yang tidak
        dinyatakan teorema itu.

  \item Dengan $\bar{z}_{2}=-0{,}2554030842$, hitung satu langkah Adam untuk $b_{2}=0{,}05$
        dengan $\eta=0{,}01$, $\beta_{1}=0{,}9$, $\beta_{2}=0{,}999$, $\epsilon=10^{-8}$,
        pada $t=1$. Tunjukkan bahwa besar langkahnya praktis $\eta$, lalu hitung langkah yang sama bila gradiennya $100$ kali lebih besar.

  \item Tunjukkan bahwa untuk dua unit ReLU di $\mathbb{R}^{2}$ keluarannya adalah fungsi
        linear sepotong-sepotong, dan bahwa bidang masukan terbagi paling banyak empat
        daerah dengan rumus linear masing-masing. Berapa daerah untuk $k$ unit ReLU di
        $\mathbb{R}^{2}$ dalam susunan umum? (Petunjuk: setiap unit memberi satu garis
        $z_{1}[i]=0$, dan di dalam setiap sel semua gerbang tetap.)

  \item Hitung gradien terhadap masukan untuk contoh hitung, lalu tentukan gangguan
        $\boldsymbol{\delta}$ dengan $\lVert\boldsymbol{\delta}\rVert_{\infty}\le0{,}05$
        yang paling menaikkan kerugian menurut linearisasi, yaitu
        $\boldsymbol{\delta}=0{,}05\operatorname{sign}(\nabla_{\mathbf{x}}L)$. Hitung kerugian pada $\mathbf{x}+\boldsymbol{\delta}$ dan bandingkan dengan ramalan linear $L+\nabla_{\mathbf{x}}L^{\top}\boldsymbol{\delta}$.

  \item Tanpa menjalankan program, ramalkan apa yang terjadi bila seluruh $\mathbf{b}_{1}$ diinisialisasi $+10$ sementara $\mathbf{W}_{1}$ tetap acak kecil, dengan tabel saturasi; lalu bila $\mathbf{b}_{1}=-10$. Periksa ramalanmu dengan menjalankan blok lab kedua setelah mengubah satu baris.
\end{enumerate}
\end{latihan}

Petunjuk untuk tiga butir tersulit. Butir 2: tulis $\mathbf{W}=\mathbf{W}_{2}\mathbf{W}_{1}$ dan $\mathbf{b}=\mathbf{W}_{2}\mathbf{b}_{1}+\mathbf{b}_{2}$, lalu bandingkan dengan definisi ambang linear. Butir 7: hitung jumlah daerah untuk $k=1,2,3$ (hasilnya $2,4,7$), lalu cari pola $1+k+k(k-1)/2$. Butir 9: blok lab kedua memakai ReLU, jadi $\mathbf{b}_{1}=+10$ tidak menyaturasi apa pun; yang jenuh adalah $\mathbf{b}_{1}=-10$, ketika semua unit mati dan turunannya nol. Pada blok ketiga yang memakai sigmoid, $z\approx10$ membuat turunan lapis tersembunyi sekitar $4{,}5\times10^{-5}$, jadi gradien ke $\mathbf{W}_{1}$ mengecil sebanding itu, sedangkan $\mathbf{b}_{1}$ masih menerima gradien penuh.

\begin{untukpeneliti}
Yang benar-benar terbukti di bab ini: ketidakmungkinan pemisahan linear untuk XOR; batas
$\lfloor1/\gamma^{2}\rfloor$ pada perceptron, dengan bukti lengkap; turunan sigmoid dan
batas $\tfrac14$; kesetaraan $\partial L/\partial z=a-y$ untuk sigmoid dengan entropi silang,
dan $\mathbf{p}-\mathbf{y}$ untuk softmax, keduanya dengan bukti; rekursi adjoint sebagai
aturan rantai; dan proposisi simetri inisialisasi. Yang dipinjam tanpa bukti: teorema
aproksimasi universal, dengan bukti aslinya pada Cybenko (1989) dan Hornik (1991), keduanya
memakai teorema representasi fungsi linear kontinu. Yang dipopulerkan melampaui buktinya:
``jaringan bisa menghampiri apa saja'' sebagai janji praktis, padahal itu pernyataan tentang kepadatan fungsi, bukan keterjangkauan parameter oleh optimasi; dan ``Adam selalu lebih baik'', padahal contoh penyangkalnya ada pada Reddi, Kale, dan Kumar (2018).

Bacaan primer yang membentuk bab ini: Nielsen bagian 2.4 dan 2.8 (empat persamaan turunan balik dan mengapa ia murah) serta bagian 3.1 (entropi silang, softmax); Goodfellow dkk. bagian 6.4.1 (aproksimasi universal) dan 6.5 (turunan balik). Untuk inisialisasi, Glorot dan Bengio (2010) serta He dkk. (2015); untuk Adam, Kingma dan Ba (2015).

Batas keberlakuan: seluruh hitungan di bab ini untuk satu contoh pada jaringan dengan satu lapisan tersembunyi, dengan konvensi turunan ReLU nol di titik patah. Tidak ada pernyataan generalisasi di sini; itu isi Bab 26. Pertanyaan terbuka: mengapa jaringan lebar lebih mudah dilatih daripada yang sempit pada galat yang sama, mengapa jaringan dengan parameter jauh melebihi banyaknya contoh masih menggeneralisasi, dan apakah ada inisialisasi yang optimal untuk kedalaman berapa pun.
\end{untukpeneliti}

\peta{Bab ini mengikuti Nielsen, \emph{Neural Networks and Deep Learning}: bagian 1.1
(perceptron), 1.2 (neuron sigmoid), 1.5 (penurunan gradien), 2.1 sampai 2.5 (notasi matriks, dua asumsi tentang fungsi kerugian, hasil kali Hadamard,
empat persamaan turunan balik beserta buktinya), 2.8 (mengapa turunan balik cepat), 3.1
(entropi silang, softmax), dan 3.3 (inisialisasi bobot). Untuk perumusan umum dan untuk sisi yang tidak dibuktikan Nielsen,
dipakai Goodfellow, Bengio, dan Courville, \emph{Deep Learning}: bagian 6.2 (pembelajaran
berbasis gradien, fungsi kerugian, dan unit keluaran), 6.3 (unit tersembunyi dan pilihan
aktivasi), 6.4.1 (aproksimasi universal dan kedalaman), 6.5 (turunan balik dan graf
komputasi), 8.3 (algoritme dasar: SGD, momentum), 8.4 (strategi inisialisasi parameter), dan
8.5 (algoritme dengan laju belajar adaptif, termasuk Adam). Untuk gradien terhadap masukan sebagai dasar
serangan, lihat pembahasan FGSM di Bab 27.}

\chapter{Generalisasi, Regularisasi, dan Pelatihan}

Sebuah tim keamanan membangun pendeteksi intrusi dari dua besaran mentah: banyak
permintaan per detik dan ukuran muatan rata-rata. Dari dua besaran itu mereka
menurunkan dua puluh satu ciri dengan mengalikan dan memangkatkan keduanya sampai
derajat lima, lalu melatih pengklasifikasi logistik pada enam puluh koneksi berlabel.
Pada data latihan, pendeteksi itu salah pada nol dari enam puluh koneksi. Pada dua
puluh ribu koneksi baru yang tidak pernah dilihatnya, ia salah pada tiga puluh dua
persen. Model yang sama, bobot yang sama, dua angka yang berbeda jauh.

Angka nol itu bukan kebetulan dan bukan prestasi. Dengan dua puluh satu ciri dan enam
puluh contoh, sebuah bidang pemisah selalu dapat ditemukan yang memotong setiap contoh
latih dengan benar, termasuk contoh yang labelnya keliru. Yang dipelajari model itu
bukan bentuk pemisah yang benar, melainkan posisi enam puluh titik itu. Begitu titiknya
berganti, bidangnya salah tempat.

Bab 25 mengajari kita menghitung gradien dan menurunkannya sampai galat latih kecil;
Bab 20 mengajari kita membaca satu angka beserta ketidakpastiannya. Tidak satu pun
menjawab pertanyaan yang menentukan di sini: besaran apa yang harus dikecilkan supaya
angka di lapangan ikut mengecil, dan bagaimana memilih setelan latihan tanpa menipu diri
sendiri.

\begin{untukanak}
Seorang murid diberi enam puluh soal latihan menjelang ujian. Ia belajar dengan cara
menghafal: setiap soal beserta jawabannya, sampai huruf per huruf. Pada soal latihan ia
selalu benar, seratus dari seratus, dan ia pulang dengan bangga. Pada ujian, soalnya
berbeda: angkanya lain, kalimatnya lain, urutannya lain. Nilainya jatuh, dan gurunya
bingung karena murid ini tampak paling rajin.

Murid kedua belajar lebih lambat. Ia tidak menghafal jawaban, ia mencari polanya: kapan
harus dibagi, kapan harus dikali, kapan jawabannya masuk akal. Pada soal latihan ia
salah pada beberapa nomor, dan ia tidak bangga. Pada ujian ia justru lebih baik.

Sekarang bagian yang membuat penasaran. Andaikan murid pertama diberi dua ribu soal
latihan, bukan enam puluh. Apakah menghafal masih buruk? Ternyata tidak selalu: kalau
soalnya cukup banyak dan bervariasi, menghafal menjadi cara belajar yang mahal tetapi
bekerja. Jadi masalahnya bukan menghafal, melainkan menghafal dari contoh yang sedikit.
Di bab ini kita akan tahu mengapa jumlah contoh dan kerumitan cara belajar harus
seimbang, mengapa menambah kerumitan justru memperburuk nilai ujian, dan bagaimana
mengukur keseimbangan itu tanpa harus menunggu ujian.
\end{untukanak}

\section{Dua galat dan satu jurang}
\tingkat{dasar}

Setiap model yang dilatih punya dua angka, dan keduanya dihitung dengan cara sama
persis: proporsi kesalahan pada sekumpulan contoh. Bedanya hanya himpunan mana yang
dipakai. Kalimat sehari-harinya: nilai ujian di rumah berbeda dari nilai ujian di
sekolah, walaupun soalnya sama jenisnya.

\kamus{generalisasi}{generalization}{kemampuan mengerjakan soal baru, bukan soal yang
sudah dihafal}

\begin{definisi}[Galat latih, galat populasi, dan jurang]
\index{generalisasi}
Misalkan contoh $(x_i,y_i)$ diambil dari distribusi $\mathcal{D}$ yang tidak diketahui.
Untuk hipotesis $h$ dan galat 0-1, \emph{galat populasi} adalah
$R(h)=\mathbb{P}_{(x,y)\sim\mathcal{D}}\bigl(h(x)\neq y\bigr)$, sedangkan \emph{galat
latih} pada $n$ contoh adalah
$\hat R_n(h)=\frac1n\sum_{i=1}^n \mathbf{1}\{h(x_i)\neq y_i\}$.
Selisih $R(h)-\hat R_n(h)$ disebut \emph{jurang generalisasi} untuk $h$ itu.
\end{definisi}

Galat latih dapat dihitung; galat populasi tidak, sebab distribusinya tidak diketahui.
Yang dapat kita lakukan hanyalah menaksirnya pada himpunan lain yang tidak dipakai
melatih, dan penaksir itu sendiri bergoyang seperti setiap proporsi pada Bab 20. Jadi
setiap angka "galat validasi" dalam bab ini adalah taksiran, bukan kebenaran.

Sekarang dua penyakit yang berlawanan. Kalimat sehari-harinya: model yang terlalu
sederhana tidak mampu menangkap polanya, sedangkan model yang terlalu rumit menghafal
kebetulan.

\begin{definisi}[Underfitting dan overfitting]
\index{underfitting}
\index{overfitting}
Sebuah model \emph{underfitting} bila galat latih dan galat populasi keduanya tinggi:
modelnya terlalu kaku untuk bentuk yang sebenarnya. Sebuah model \emph{overfitting}
bila galat latih rendah tetapi jurang generalisasinya besar: modelnya menyesuaikan diri
pada kebetulan data latih, bukan pada bentuk yang sebenarnya.
\end{definisi}

Kedua penyakit itu tidak simetris dalam praktik. Underfitting terlihat jelas: galat latih
tinggi, dan menambah ciri menolong. Overfitting menipu, sebab satu-satunya angka yang
tersedia dengan mudah, yaitu galat latih, justru tampak sempurna. Model pada pembuka bab
adalah contoh paling murni: bobot terbesarnya $406{,}204$, tanda bahwa skornya bergantung
pada perbedaan kecil di antara titik, bukan pada skala yang masuk akal.

Yang mengatur keduanya adalah \emph{kapasitas}: seberapa kaya himpunan fungsi yang boleh
dipilih model. Untuk pengklasifikasi linear atas $d$ ciri, kapasitas naik dengan $d$;
untuk jaringan saraf, dengan jumlah unit, kedalaman, dan jumlah langkah latihan.
Kapasitas bukan sekadar cacah parameter: model besar yang diregularisasi kuat dapat
berperilaku seperti model kecil.

\begin{bacakeras}
Dua suara, bergantian.

``Seratus dari seratus. Berarti ia paham, bukan?''

``Belum tentu. Seratus dari seratus hanya berarti tidak ada soal yang baru.''

``Jadi menghafal itu salah?''

``Menghafal bukan salah. Menghafal dari sedikit contoh yang salah.''

``Bedanya di mana?''

``Di soal yang tidak ada di catatan. Hafalan menjawab yang sudah pernah dilihat.
Pemahaman menjawab yang belum pernah dilihat.''

``Dan kalau soalnya dua ribu, bukan enam puluh?''

``Maka hafalan mulai menolong. Bukan karena hafalan berubah menjadi pemahaman, tetapi
karena kebetulan yang tersimpan semakin banyak yang terpakai lagi.''

Ucapkan sekali lagi, pelan: hafalan menjawab yang sudah ada, pemahaman menjawab yang
belum ada. Ulangi baris itu sampai terasa wajar di mulut: sudah ada, belum ada.
\end{bacakeras}

Jurang generalisasi menyusut ketika contoh bertambah. Pada eksperimen kapasitas dengan
anggaran langkah tetap, jurangnya $0{,}1338$ pada $n=60$, $0{,}0274$ pada $n=150$,
$0{,}0193$ pada $n=600$, dan $0{,}0097$ pada $n=2400$, sedangkan galat latihnya hampir
tidak bergerak di kisaran $0{,}17$ sampai $0{,}21$.

\section{Dekomposisi bias-variansi}
\tingkat{menengah}

Pertanyaan yang lebih tajam: dari mana jurang itu berasal. Kalimat sehari-harinya:
seorang penembak dapat meleset karena bidikannya bergeser ke satu arah, atau karena
tangannya gemetar, dan dua sebab itu menuntut obat yang berbeda.

\kamus{bias-variansi}{bias-variance}{meleset sistematis lawan gemetar; keduanya
menyumbang galat, dan menekan yang satu kadang menaikkan yang lain}

\begin{teorema}[Dekomposisi bias-variansi]
\index{bias-variansi}
Misalkan $y=f(x)+\varepsilon$ dengan $\mathbb{E}[\varepsilon]=0$ dan
$\operatorname{Var}(\varepsilon)=\sigma^2$, dan misalkan $\hat f$ dilatih pada
himpunan latih acak $D$ yang bebas dari $\varepsilon$. Untuk setiap $x$ tetap,
\[
\mathbb{E}_{D,\varepsilon}\bigl[(y-\hat f(x))^2\bigr]
=\sigma^2
+\bigl(\mathbb{E}_D\hat f(x)-f(x)\bigr)^2
+\operatorname{Var}_D\bigl(\hat f(x)\bigr).
\]
\end{teorema}

\begin{proof}[Bukti]
Tulis $y-\hat f(x)=\varepsilon+\bigl(f(x)-\mathbb{E}_D\hat f(x)\bigr)
+\bigl(\mathbb{E}_D\hat f(x)-\hat f(x)\bigr)$, yaitu tiga suku: derau yang tidak dapat
dikurangi, bias, dan simpangan dari rata-rata. Kuadratkan dan ambil harapannya. Suku
silang yang memuat $\varepsilon$ lenyap karena $\mathbb{E}[\varepsilon]=0$ dan
$\varepsilon$ bebas dari $D$; suku silang antara bias dan simpangan lenyap karena
$\mathbb{E}_D[\hat f(x)-\mathbb{E}_D\hat f(x)]=0$. Yang tersisa tepat tiga suku pada
pernyataan teorema.
\end{proof}

Suku pertama tidak dapat dikurangi model apa pun; ia milik datanya, disebut galat tak
tereduksikan, dan pada data pembuka bab nilainya sekitar $0{,}21$ karena labelnya acak.
Suku kedua dan ketiga dapat diperdagangkan: menambah kapasitas menurunkan bias dan
menaikkan variansi, mengurangi kapasitas melakukan sebaliknya. Regularisasi adalah nama
untuk perdagangan itu.

Contoh paling murah yang memuat kedua suku sekaligus adalah penaksir menyusut. Untuk
mengukur $\mu$ dari $n$ pengamatan berderau $\sigma^2$, orang biasa memakai
$\bar X_n$, yang tak bias. Kita boleh memakai $c\bar X_n$ untuk $c\in(0,1)$, dan
penaksir itu berbias dengan sengaja. Galat kuadrat rata-ratanya
\[
\operatorname{MSE}(c)=(c-1)^2\mu^2+c^2\frac{\sigma^2}{n},
\]
diminimalkan di $c^\ast=\mu^2/(\mu^2+\sigma^2/n)$. Untuk $\mu=2$, $\sigma=1$, dan
$n=25$ diperoleh $c^\ast=0{,}990099$: penaksir terbaik justru sedikit menyusut.

\begin{center}
\small
\begin{tabular}{rrrr}
\hline
$c$ & bias$^2$ & variansi & MSE \\
\hline
0,70 & 0,360000 & 0,019600 & 0,379600 \\
0,90 & 0,040000 & 0,032400 & 0,072400 \\
0,99 & 0,000400 & 0,039204 & 0,039604 \\
1,00 & 0,000000 & 0,040000 & 0,040000 \\
1,10 & 0,040000 & 0,048400 & 0,088400 \\
1,30 & 0,360000 & 0,067600 & 0,427600 \\
\hline
\end{tabular}
\end{center}

Tiga hal perlu dibaca dari tabel itu. Keuntungan menyusut kecil: MSE turun dari
$0{,}040000$ ke $0{,}039604$, sekitar satu persen. Kerugian menyimpang ke arah yang
salah jauh lebih besar: $c=1{,}3$ memberi $0{,}4276$, sedangkan $c=0{,}7$ hanya
$0{,}3796$. Dan minimumnya landai, sehingga memilih $c$ yang sedikit keliru bukan
bencana.

Teorema itu tidak memberi rumus untuk $\hat f$ yang rumit seperti jaringan saraf.
Dekomposisinya berlaku formal, tetapi kedua sukunya tidak dapat dihitung dari data saja;
yang teramati hanyalah jumlahnya, yaitu galat validasi. Karena itu sisa bab ini bekerja
dengan angka teramati, bukan dengan bias dan variansi yang sebenarnya.

\section{Regularisasi L2 dan L1}
\tingkat{menengah}

Cara paling langsung menekan variansi adalah melarang bobot tumbuh besar. Kalimat
sehari-harinya: kita menambahkan denda ke dalam sasaran, sehingga model harus membayar
untuk setiap bobot besar, dan hanya mau membayar bila keuntungannya di data latih sepadan.

\kamus{regularisasi}{regularization}{denda yang dibayar model untuk bobot besar, supaya
ia tidak menghafal data latih}

\begin{definisi}[Sasaran ternormalisasi]
\index{regularisasi}
Untuk sasaran asal $J(\mathbf{w})$ dan penalti $\Omega$, sasaran yang diminimalkan
adalah $\tilde J(\mathbf{w})=J(\mathbf{w})+\lambda\Omega(\mathbf{w})$ dengan
$\lambda\ge0$ sebagai hiperparameter. Untuk penalti L2, $\Omega=\frac12\|\mathbf{w}\|_2^2$;
untuk penalti L1, $\Omega=\|\mathbf{w}\|_1=\sum_j|w_j|$.
\end{definisi}

Untuk L2, gradien sasarannya $\nabla J+\lambda\mathbf{w}$, jadi satu langkah penurunan
gradien dengan laju $\eta$ berbentuk
\[
\mathbf{w}\leftarrow\mathbf{w}-\eta\bigl(\nabla J+\lambda\mathbf{w}\bigr)
=(1-\eta\lambda)\mathbf{w}-\eta\nabla J .
\]
Faktor $(1-\eta\lambda)$ itulah asal nama \emph{weight decay}: pada setiap langkah,
sebelum gradien dipakai, bobot digandakan dengan bilangan yang sedikit kurang dari satu.
Kalau gradiennya nol, bobotnya meluruh menuju nol secara geometris dengan waktu paruh
$\ln 2/(\eta\lambda)$ langkah.

\begin{akibat}[Peluruhan bobot sama dengan penalti L2]
\index{weight decay}
Untuk laju $\eta$ dan penalti L2 dengan koefisien $\lambda$, penurunan gradien atas
sasaran ternormalisasi identik dengan penurunan gradien atas sasaran asal yang diikuti
penggandaan bobot dengan $1-\eta\lambda$. Dengan demikian $\lambda$ dan $\eta$ hanya
muncul sebagai hasil kali $\eta\lambda$ pada langkah peluruhan, sehingga menaikkan laju
dua kali dan menurunkan $\lambda$ dua kali memberi lintasan peluruhan yang sama.
\end{akibat}

Penalti L2 juga punya tafsiran Bayesian: ia setara prior Gaussian pada bobot, sehingga
penyelesaiannya taksiran MAP, bukan kemungkinan maksimum. Itu menjelaskan mengapa L2
percaya bobot kecil lebih mungkin, dan mengapa ia tidak pernah memaksa satu bobot menjadi
nol.

L1 berperilaku lain, dan seluruh bedanya terletak pada bentuk penaltinya di dekat nol.
Kalimat sehari-harinya: penalti kuadratik menarik semakin lemah ketika bobotnya sudah
kecil, sedangkan penalti nilai mutlak menarik dengan kekuatan tetap sampai bobotnya
benar-benar nol.

\begin{lemma}[Ambang lunak]
\index{regularisasi}
Untuk $a\in\mathbb{R}$ dan $t>0$, peminimum $u\mapsto\frac12(u-a)^2+t|u|$ adalah
\[
u^\ast=\operatorname{sign}(a)\max\{|a|-t,0\},
\]
sehingga $u^\ast=0$ tepat ketika $|a|\le t$.
\end{lemma}

Untuk $u>0$ fungsinya $\frac12(u-a)^2+tu$ dengan turunan $u-a+t$, jadi titik
kritisnya $u=a-t$, sah hanya bila $a>t$. Untuk $u<0$ turunannya $u-a-t$ dengan titik
kritis $u=a+t$, sah hanya bila $a<-t$. Pada $u=0$ subgradiennya selang $[-a-t,-a+t]$,
yang memuat nol tepat ketika $|a|\le t$; dalam hal itu nol peminimum karena fungsinya
konveks. Ketiga kasus bergabung menjadi rumus di atas.

Lemma itu menjelaskan seluruh perilaku L1. Dengan pendekatan kuadratik di sekitar
penyelesaian tanpa penalti $w^\ast$ dan kelengkungan $H$, sasarannya kira-kira
$\frac H2(w-w^\ast)^2+\lambda|w|$, sehingga
$w=\operatorname{sign}(w^\ast)\max\{|w^\ast|-\lambda/H,0\}$: setiap bobot di bawah
$\lambda/H$ menjadi nol tepat, bukan sekadar kecil. Untuk L2 penyelesaiannya
$w=w^\ast/(1+\lambda/H)$, yang tidak pernah nol kecuali $w^\ast$ sendiri nol.

\begin{proposisi}[Mengapa L1 menghasilkan nol sejati]
\index{regularisasi}
Misalkan $J$ dapat diturunkan dan konveks di sekitar nol. Titik $\mathbf{w}=\mathbf{0}$
adalah peminimum $J(\mathbf{w})+\lambda\|\mathbf{w}\|_1$ jika dan hanya jika
$|\partial J/\partial w_j|\le\lambda$ untuk setiap $j$. Untuk penalti L2, syarat yang
sama adalah $\partial J/\partial w_j=0$ untuk setiap $j$.
\end{proposisi}

\begin{proof}[Bukti]
Subgradien $\|\cdot\|_1$ di nol adalah kotak $[-1,1]^d$, jadi kondisi
$0\in\nabla J(\mathbf{0})+\lambda[-1,1]^d$ setara dengan $|(\nabla J)_j|\le\lambda$
untuk semua $j$; untuk fungsi konveks, kondisi itu perlu dan cukup. Untuk penalti L2,
gradiennya $\nabla J+\lambda\mathbf{w}$ yang di $\mathbf{0}$ sama dengan $\nabla J$, jadi
syaratnya $\nabla J(\mathbf{0})=\mathbf{0}$.
\end{proof}

Akibat praktisnya penting untuk pembacaan hasil. L1 memberi daerah datar selebar
$2\lambda$ di sekitar nol, sehingga bobot yang "hampir tidak berguna" dimatikan
sepenuhnya; L2 hanya mengecilkannya. Karena itu L1 dipakai sebagai pemilih ciri, dan
karena itu pula jumlah bobot nol pada L1 bukan ukuran kepentingan yang halus: ia
melompat, tidak bergeser perlahan.

Sekarang angkanya. Pengklasifikasi logistik yang sama, data yang sama, dua puluh satu
ciri pada enam puluh contoh latih, semuanya dihitung dengan skrip pada bagian lab.

\begin{center}
\small
\begin{tabular}{rrrr}
\hline
$\lambda$ & $\max_j|w_j|$ & galat latih & galat validasi \\
\hline
0       & 406,204 & 0,0000 & 0,3262 \\
0,001   &   3,056 & 0,1000 & 0,2927 \\
0,003   &   1,814 & 0,1000 & 0,2688 \\
0,010   &   1,018 & 0,1167 & 0,2475 \\
0,030   &   0,631 & 0,1167 & 0,2433 \\
0,100   &   0,399 & 0,1500 & 0,2406 \\
0,300   &   0,268 & 0,2333 & 0,2452 \\
1,000   &   0,142 & 0,2500 & 0,2695 \\
3,000   &   0,062 & 0,3333 & 0,3205 \\
10,000  &   0,022 & 0,3333 & 0,3395 \\
100,000 &   0,003 & 0,3333 & 0,3395 \\
\hline
\end{tabular}
\end{center}

\begin{figure}[htbp]
\centering
\begin{tikzpicture}
\begin{axis}[
  width=12.2cm, height=5.6cm,
  xmin=-0.4, xmax=10.4, ymin=0, ymax=0.40,
  axis lines=left,
  xtick={0,3,5,7,10},
  xticklabels={$0$,$10^{-2}$,$0{,}1$,$1$,$100$},
  ytick={0,0.1,0.2,0.3,0.4},
  yticklabels={$0$,$0{,}10$,$0{,}20$,$0{,}30$,$0{,}40$},
  tick label style={font=\scriptsize},
  xlabel={kekuatan regularisasi $\lambda$ (kapasitas menurun ke kanan)},
  ylabel={galat},
  xlabel style={font=\small},
  ylabel style={font=\small},
  legend style={font=\scriptsize,draw=none,at={(0.03,0.97)},anchor=north west},
]
\addplot[very thick,layerindigo,mark=*,mark size=1.5pt] coordinates {
  (0,0.0000) (1,0.1000) (2,0.1000) (3,0.1167) (4,0.1167) (5,0.1500)
  (6,0.2333) (7,0.2500) (8,0.3333) (9,0.3333) (10,0.3333)};
\addlegendentry{galat latih}
\addplot[very thick,accent,mark=square*,mark size=1.5pt] coordinates {
  (0,0.3262) (1,0.2927) (2,0.2688) (3,0.2475) (4,0.2433) (5,0.2406)
  (6,0.2452) (7,0.2695) (8,0.3205) (9,0.3395) (10,0.3395)};
\addlegendentry{galat validasi}
\addplot[only marks,mark=o,mark size=3.4pt,accent] coordinates {(5,0.2406)};
\draw[densely dashed,layergrey] (axis cs:5,0)--(axis cs:5,0.2406);
\node[accent,font=\scriptsize,anchor=north] at (axis cs:4.55,0.238)
  {minimum $0{,}2406$};
\end{axis}
\end{tikzpicture}
\caption{Galat latih dan galat validasi terhadap kekuatan regularisasi $\lambda$; arah ke
kanan berarti kapasitas menurun, sebab bobot ditahan makin ketat. Kesebelas titik adalah
baris pada tabel sebelumnya. Galat latih naik monoton dari $0{,}0000$ ke $0{,}3333$,
sedangkan galat validasi membentuk lembah: $0{,}3262$ pada $\lambda=0$, minimum
$0{,}2406$ pada $\lambda=0{,}1$, lalu $0{,}3395$ pada $\lambda=10$ ke atas. Dasar
lembahnya lebar, dari $0{,}2433$ pada $\lambda=0{,}03$ ke $0{,}2452$ pada $\lambda=0{,}3$.}
\label{fig:lembah-kapasitas}
\end{figure}

Gambar~\ref{fig:lembah-kapasitas} memplot kolom ketiga dan keempat tabel itu sebagai
satu lembah. Bacalah tabel itu sebagai tiga daerah, bukan sebelas baris. Tanpa
regularisasi galat latihnya nol dan galat validasinya $0{,}3262$, dengan bobot sampai
ratusan karena tidak ada yang menahannya. Di daerah $\lambda$ antara $0{,}001$ dan
$0{,}1$, galat latih naik ke $0{,}1500$ sementara galat validasi turun ke $0{,}2406$.
Di daerah $\lambda\ge3$ keduanya buruk: $0{,}3333$ dan $0{,}3395$.

Baris terakhir pantas ditekankan. Dengan $\lambda=10$ galat validasinya $0{,}3395$,
\emph{lebih buruk} daripada $0{,}3262$ milik model tanpa regularisasi. Jadi
regularisasi bukan obat yang selalu menolong; ia menolong pada dosis tertentu dan merusak
di luar itu. Pada $\lambda=10$ model itu praktis menebak selalu kelas mayoritas, yang
galatnya $0{,}3419$: bobotnya begitu kecil sehingga skornya hampir seluruhnya ditentukan
bias.

Penalti L1 memberi gambaran serupa dengan bentuk berbeda. Pada $\lambda=0{,}05$ galat
validasinya $0{,}2334$ dengan sebelas dari dua puluh satu bobot hidup; pada
$\lambda=0{,}2$ tinggal satu bobot dengan galat $0{,}3395$. Cacah bobot hidup tidak turun
monoton: $\lambda=0{,}005$ menyisakan sebelas, $\lambda=0{,}01$ menyisakan empat belas.
Angka itu stabil terhadap empat kali lipat iterasi, jadi ia bukan kegagalan pengoptimum;
lintasan L1 untuk loss logistik memang tidak menjamin cacah nol yang monoton, dan itu
alasan tambahan untuk tidak membacanya sebagai ukuran kepentingan.

Perbandingan terakhir: L1 mencapai $0{,}2334$ dan L2 $0{,}2406$. Selisih $0{,}0072$
lebih kecil daripada goyangan antarpartisi, sekitar $0{,}02$, sehingga data ini tidak
cukup untuk menyimpulkan L1 lebih baik; yang dapat disimpulkan hanyalah bahwa keduanya
memperbaiki model tanpa regularisasi.

\begin{untukanak}
Bayangkan sebuah kelas yang harus menebak berat badan temannya. Gurunya memberi dua
cara. Cara pertama: setiap tebakan yang terlalu jauh dari jawaban dikenai denda yang
makin besar makin jauh. Cara kedua: setiap tebakan yang tidak persis sama dengan
jawaban dikenai denda tetap, tidak peduli bedanya sedikit atau banyak.

Cara pertama membuat semua tebakan bergeser sedikit ke arah tengah, tetapi tidak ada
yang menjadi nol. Cara kedua membuat tebakan yang sudah cukup dekat langsung dibulatkan
menjadi nol, karena dendanya tetap dan tidak sepadan dengan perbaikan yang mungkin
didapat. Itulah bedanya dua denda yang tampak mirip: yang satu mengecilkan, yang lain
mematikan.

Pertanyaannya: kalau kita punya dua puluh satu ciri dan hanya enam puluh contoh, denda
seperti apa yang membuat model berhenti menghafal? Jawabannya bukan "denda besar",
sebab denda yang terlalu besar mematikan semuanya, termasuk yang benar. Di bab ini kita
akan tahu mengapa ada dosis, mengapa dosis itu harus dicari, dan mengapa dosis yang
dicari pada data yang salah akan menipu kita dengan rapi.
\end{untukanak}

\section{Berhenti lebih awal, dropout, dan augmentasi}
\tingkat{menengah}

Tiga cara lain menekan variansi tanpa menyentuh rumus sasaran. Kalimat sehari-harinya:
kita dapat berhenti belajar sebelum kebetulan mulai dihafal, kita dapat mematikan
sebagian ingatan secara acak, dan kita dapat memperbanyak contoh dengan mengubah contoh
yang ada.

Berhenti lebih awal adalah bentuk regularisasi yang paling murah, sebab ia tidak
menambah satu baris pun pada sasaran: ia hanya menghentikan lintasan. Selama latihan,
galat latih hampir selalu turun, tetapi galat validasi turun lalu naik. Titik terendah
galat validasi itulah yang diambil, dan jumlah langkah menjadi hiperparameter yang
dicari.

\begin{center}
\small
\begin{tabular}{rrrr}
\hline
epoch & galat latih & galat validasi & $\max_j|w_j|$ \\
\hline
1    & 0,1333 & 0,2481 & 0,15 \\
5    & 0,1333 & 0,2429 & 0,38 \\
10   & 0,1333 & 0,2432 & 0,48 \\
50   & 0,1167 & 0,2454 & 0,77 \\
200  & 0,1167 & 0,2541 & 1,38 \\
800  & 0,1000 & 0,2879 & 2,85 \\
2400 & 0,0833 & 0,3109 & 4,81 \\
\hline
\end{tabular}
\end{center}

\begin{figure}[htbp]
\centering
\begin{tikzpicture}
\begin{axis}[
  width=11.8cm, height=5.6cm,
  xmode=log,
  xmin=1, xmax=3000, ymin=0.05, ymax=0.35,
  axis lines=left,
  xtick={1,10,100,1000},
  xticklabels={$1$,$10$,$100$,$1000$},
  ytick={0.10,0.15,0.20,0.25,0.30},
  yticklabels={$0{,}10$,$0{,}15$,$0{,}20$,$0{,}25$,$0{,}30$},
  tick label style={font=\scriptsize},
  xlabel={epoch pelatihan (skala log)},
  ylabel={galat},
  xlabel style={font=\small},
  ylabel style={font=\small},
  legend style={font=\scriptsize,draw=none,at={(0.03,0.97)},anchor=north west},
]
\addplot[very thick,layerteal,mark=*,mark size=1.5pt] coordinates {
  (1,0.1333) (5,0.1333) (10,0.1333) (50,0.1167) (200,0.1167) (800,0.1000)
  (2400,0.0833)};
\addlegendentry{galat latih}
\addplot[very thick,accent,mark=square*,mark size=1.5pt] coordinates {
  (1,0.2481) (5,0.2429) (10,0.2432) (50,0.2454) (200,0.2541) (800,0.2879)
  (2400,0.3109)};
\addlegendentry{galat validasi}
\draw[densely dashed,layergrey] (axis cs:8,0.05)--(axis cs:8,0.2427);
\addplot[only marks,mark=o,mark size=3.4pt,accent] coordinates {(8,0.2427)};
\node[accent,font=\scriptsize,anchor=north west] at (axis cs:8.6,0.2325)
  {epoch 8: $0{,}2427$};
\end{axis}
\end{tikzpicture}
\caption{Kedua galat sepanjang satu lintasan pelatihan, sumbu mendatar skala log. Galat
latih turun tanpa pernah naik, dari $0{,}1333$ pada epoch 1 ke $0{,}0833$ pada epoch
2400. Galat validasi turun sampai $0{,}2427$ pada epoch 8 lalu naik ke $0{,}3109$. Garis
putus-putus menandai titik henti terbaik: berhenti di situ memberi galat $0{,}0682$ lebih
rendah daripada melatih sampai epoch 2400, dan kurvanya mulai berpisah tepat sesudah
titik itu.}
\label{fig:dua-kurva}
\end{figure}

Gambar~\ref{fig:dua-kurva} menggambar kedua kolom itu. Galat latih turun sepanjang
lintasan, dari $0{,}1333$ ke $0{,}0833$, tanpa pernah naik, sedangkan galat validasi
turun sampai epoch 8 dengan nilai terbaik $0{,}2427$ lalu naik sampai $0{,}3109$ pada
epoch 2400. Kolom terakhir menjelaskan mengapa: yang tumbuh sepanjang latihan adalah
besarnya bobot, dan itulah variansi yang sedang kita bayar.

Berhenti lebih awal bukan sekadar taktik. Untuk sasaran kuadratik dan laju kecil, berhenti
pada langkah $\tau$ hampir setara dengan penalti L2 berkoefisien sekitar
$1/(\tau\eta)$; itu perkiraan, bukan identitas tepat. Karena itu jumlah langkah adalah
hiperparameter kapasitas, bukan urusan "melatih sampai selesai".

\begin{resep}
\begin{enumerate}
  \item Pisahkan tiga himpunan sebelum satu langkah latihan pun dijalankan: latih,
        validasi, uji. Himpunan uji disegel.
  \item Latih pada himpunan latih saja; setiap beberapa epoch hitung galat pada
        himpunan validasi, dan catat pasangan (epoch, galat validasi, galat latih).
  \item Ambil epoch dengan galat validasi terkecil, lalu kembalikan bobot ke titik itu;
        jangan memakai bobot epoch terakhir.
  \item Baru sekarang buka himpunan uji, hitung galatnya satu kali, dan laporkan
        bersama epoch pilihan serta lebar kurva validasinya.
\end{enumerate}
Peringatan urutan: langkah 1 harus selesai sebelum langkah 2, dan langkah 4 tidak boleh
dipakai menggeser pilihan pada langkah 3. Bila himpunan uji dilihat sewaktu memilih
epoch, ia bukan lagi uji melainkan validasi kedua, dan angkanya berhenti menjadi
taksiran yang jujur.
\end{resep}

Dropout mengambil jalan lain: selama latihan setiap unit dimatikan secara acak dengan
peluang $1-p$, dan yang hidup dibagi $p$ supaya harapannya tidak bergeser. Kalimat
sehari-harinya: setiap langkah sebagian anggota tim tidak masuk, sehingga tidak ada
anggota yang boleh diandalkan sendirian.

\begin{lemma}[Harapan dan variansi topeng dropout]
\index{dropout}
Misalkan $m=B/p$ dengan $B\sim\mathrm{Bernoulli}(p)$ dan $0<p\le1$. Maka
$\mathbb{E}[m]=1$ dan $\operatorname{Var}(m)=(1-p)/p$. Untuk unit dengan masukan
$z_1,\dots,z_k$ dan topeng bebas $m_1,\dots,m_k$, keluarannya
$s=\sum_j m_jw_jz_j$ memenuhi $\mathbb{E}[s]=\sum_jw_jz_j$ dan
$\operatorname{Var}(s)=\frac{1-p}{p}\sum_j(w_jz_j)^2$.
\end{lemma}

\begin{proof}[Bukti]
$\mathbb{E}[m]=p/p=1$ dan $\operatorname{Var}(m)=\mathbb{E}[m^2]-1=p/p^2-1=(1-p)/p$.
Untuk $s$, kelinearan harapan memberi suku pertama, dan kebebasan topeng memberi
variansi jumlah, yaitu $\sum_j(w_jz_j)^2\operatorname{Var}(m_j)$.
\end{proof}

Rumus variansi itu menjelaskan dua hal yang sering dilupakan. Derau dropout sebanding
dengan jumlah kuadrat kontribusi, jadi unit dengan bobot besar menerima derau relatif
lebih besar, dan dropout searah dengan penalti L2. Dan derau itu hilang sama sekali saat
pengujian, sehingga model yang dipakai berbeda dari model yang dilatih; pada jaringan
sangat kecil harga itu bisa melebihi manfaatnya.

Augmentasi data memperbanyak contoh dengan mengubah contoh yang ada, misalnya menggeser
atau menambahkan derau kecil, asalkan labelnya tidak berubah. Untuk model linear,
hubungannya dengan L2 dapat dihitung tepat.

\begin{proposisi}[Derau Gauss pada masukan sama dengan penalti L2]
\index{augmentasi data}
Misalkan $\hat y=\mathbf{w}^{\top}\mathbf{x}$ dan $\boldsymbol{\delta}\sim
N(\mathbf{0},\sigma^2I)$ ditambahkan pada masukan. Maka
\[
\mathbb{E}_{\boldsymbol{\delta}}\bigl[(\mathbf{w}^{\top}(\mathbf{x}+\boldsymbol{\delta})-y)^2\bigr]
=(\mathbf{w}^{\top}\mathbf{x}-y)^2+\sigma^2\|\mathbf{w}\|_2^2 .
\]
\end{proposisi}

\begin{proof}[Bukti]
Jabarkan: suku linearnya $2(\mathbf{w}^{\top}\mathbf{x}-y)\,
\mathbb{E}[\mathbf{w}^{\top}\boldsymbol{\delta}]=0$, dan suku kuadratnya
$\mathbb{E}[(\mathbf{w}^{\top}\boldsymbol{\delta})^2]
=\mathbf{w}^{\top}\mathbb{E}[\boldsymbol{\delta}\boldsymbol{\delta}^{\top}]\mathbf{w}
=\sigma^2\|\mathbf{w}\|_2^2$.
\end{proof}

Jadi menambahkan derau Gauss kecil pada masukan sama dengan menambahkan penalti L2
dengan $\lambda=\sigma^2$: keduanya satu gagasan yang ditulis dua kali. Pemeriksaan
angkanya, untuk $\mathbf{w}=(1,2)$, $\mathbf{x}=(1,0)$, $y=3$, dan $\sigma=0{,}1$:
sasaran asalnya $4$, ramalan identitasnya $4+0{,}01\cdot5=4{,}05$, dan simulasi satu juta
derau memberi $4{,}049107$. Selisih $0{,}0009$ itu sebanding galat acak simulasi, bukan
kesalahan rumus.

Peringatan untuk bidang keamanan: augmentasi hanya sah bila label benar-benar invarian
terhadap perubahan yang dilakukan. Menambahkan derau kecil pada ciri lalu lintas boleh,
tetapi membalik urutan paket atau menggeser waktu kedatangan dapat mengubah serangan
menjadi bukan serangan. Augmentasi yang salah label bukan regularisasi, melainkan
kebohongan yang akan tetap tampak baik selama data validasinya diubah dengan cara sama.

\section{SGD, momentum, dan Adam}
\tingkat{lanjut}

Sampai di sini kita memperlakukan "latihan" sebagai satu kotak. Isi kotak itu menentukan
seberapa cepat galat turun dan, pada anggaran langkah terbatas, seberapa rendah galat yang
tercapai. Kalimat sehari-harinya: arah yang benar belum cukup; yang menentukan adalah
bagaimana langkah diambil dan seberapa besar.

Penurun gradien stokastik mengganti gradien penuh oleh gradien atas satu minibatch
berukuran $b$. Taksiran itu tak bias, dan variansinya mengecil sebanding $1/b$:
$\operatorname{Var}(\hat{\mathbf{g}}_b)=\sigma_g^2/b$. Jadi minibatch besar mengurangi
goyangan, tetapi tidak mengubah arah rata-ratanya. Yang mengubah arah adalah momentum
dan Adam.

\begin{definisi}[Momentum dan Adam]
\index{momentum}
\index{adam}
Momentum menyimpan rata-rata gradien bergerak:
$\mathbf{v}_t=\beta\mathbf{v}_{t-1}+\mathbf{g}_t$ dan
$\mathbf{w}_t=\mathbf{w}_{t-1}-\eta\mathbf{v}_t$, dengan $0\le\beta<1$. Adam menyimpan
dua momen, $\mathbf{m}_t=\beta_1\mathbf{m}_{t-1}+(1-\beta_1)\mathbf{g}_t$ dan
$\mathbf{v}_t=\beta_2\mathbf{v}_{t-1}+(1-\beta_2)\mathbf{g}_t^2$ (kuadrat per
koordinat), lalu memakai koreksi bias dan pembaruan
$\mathbf{w}_t=\mathbf{w}_{t-1}-\eta\,\hat{\mathbf{m}}_t/(\sqrt{\hat{\mathbf{v}}_t}+\epsilon)$
dengan $\hat{\mathbf{m}}_t=\mathbf{m}_t/(1-\beta_1^t)$ dan
$\hat{\mathbf{v}}_t=\mathbf{v}_t/(1-\beta_2^t)$.
\end{definisi}

Koreksi bias itu bukan hiasan: tanpa koreksi $\mathbf{m}_1=(1-\beta_1)\mathbf{g}_1$
sangat kecil, sehingga langkah pertama hampir nol. Pembagian oleh $\sqrt{\hat{\mathbf{v}}}$
lebih penting lagi untuk masalah berskala buruk: ia membuat besar langkah tiap koordinat
kurang lebih $\eta$, tidak peduli gradiennya besar atau kecil. Momentum hanya meratakan
arah sepanjang waktu.

Angka pada kuadratik berikut memperlihatkan bedanya. Ambil
$f(\mathbf{w})=\frac12\mathbf{w}^{\top}H\mathbf{w}$ dengan $H=\operatorname{diag}(1,10^4)$
dan $\mathbf{w}_0=(1,0)$; bilangan kondisi $H$ adalah $10^4$, yang berarti kurvanya
sempit di satu arah dan sangat landai di arah lain. Kita hitung langkah sampai
$f\le10^{-8}f(\mathbf{w}_0)$.

\begin{center}
\small
\begin{tabular}{lrl}
\hline
metode & laju & langkah sampai $10^{-8}$ \\
\hline
penurun gradien & $10^{-4}$ & 92099 \\
momentum ($\beta=0{,}9$) & $10^{-4}$ & 9132 \\
Adam & $0{,}1$ & 82 \\
\hline
\end{tabular}
\end{center}

Penurun gradien memerlukan $92099$ langkah karena laju terbesarnya dibatasi kelengkungan
terbesar, sedangkan arah yang landai hanya bergerak dengan laju itu. Momentum memotong
satu digit, sesuai perkiraan bahwa percepatan mengubah ketergantungan pada bilangan
kondisi $\kappa$ menjadi sekitar $\sqrt\kappa$. Adam memotong tiga digit, sebab pembagian
per koordinat membuat arah yang landai tidak lagi dihukum oleh arah yang curam. Itu bukan
janji umum: pada masalah berskala baik keunggulan Adam menyusut.

Pada masalah logistik di bab ini, dengan minibatch delapan dan sasaran galat latih
$0{,}40$, median lima biji acak memberi SGD laju $0{,}05$ seratus langkah, sedangkan
Adam laju $0{,}1$ hanya sepuluh. Kesimpulannya harus hati-hati: keunggulannya sebanding
satu digit dan sebagian berasal dari laju yang lebih besar, bukan dari arah yang lebih
baik.

Laju pembelajaran biasanya tidak tetap. Dua jadwal yang lazim adalah peluruhan tangga,
yaitu laju dibagi sepuluh pada epoch tertentu, dan peluruhan kosinus
$\eta_t=\eta_{\min}+\frac12(\eta_0-\eta_{\min})(1+\cos(\pi t/T))$: langkah besar di awal
supaya cepat mendekat, langkah kecil di akhir supaya tidak melompat di lembah. Dengan
$3000$ langkah minibatch dan $\lambda=0{,}1$, laju tetap $0{,}2$ memberi galat latih
median $0{,}3866$, sedangkan kosinus dari $0{,}5$ ke $0{,}01$ memberi $0{,}3549$.

\section{Validasi silang, pemilihan hiperparameter, dan kebocoran}
\tingkat{lanjut}

Semua angka validasi di bab ini berasal dari himpunan yang tidak dipakai melatih. Itu
sudah benar, tetapi belum cukup: begitu kita memilih $\lambda$ dengan mata yang melihat
galat validasi, galat validasi itu berhenti menjadi taksiran yang tak bias bagi galat
populasi. Kalimat sehari-harinya: kalau kita memilih soal ujian berdasarkan nilai yang
kita dapat, nilai itu tidak lagi mengukur kemampuan kita.

\kamus{validasi silang}{cross validation}{memutar peran data bergiliran, sehingga setiap
contoh pernah menjadi ujian bagi model yang tidak melihatnya}

\begin{definisi}[Validasi silang $k$ lipatan]
\index{validasi silang}
Data latih dipecah menjadi $k$ lipatan yang saling lepas dan berukuran serupa. Untuk
setiap $j$, model dilatih pada $k-1$ lipatan dan diuji pada lipatan $j$; galat validasi
silang adalah rata-rata galat pada $k$ lipatan itu, dan hiperparameter yang dipilih
adalah yang meminimumkannya. Untuk $k=n$ prosedur ini disebut
\emph{leave-one-out}.
\end{definisi}

Validasi silang menjawab satu kelemahan pemisahan tunggal: pada $n=60$ himpunan
validasi berisi dua belas contoh, dan galat pada dua belas contoh bergoyang dengan
simpangan baku sekitar $0{,}12$. Dengan lima lipatan setiap contoh pernah diuji, tetapi
goyangannya tidak hilang.

\begin{center}
\small
\begin{tabular}{rrr}
\hline
$\lambda$ & galat CV (rata-rata 5 partisi) & simpangan antarpartisi \\
\hline
0     & 0,3233 & 0,1245 \\
0,003 & 0,2400 & 0,0253 \\
0,030 & 0,2467 & 0,0217 \\
0,100 & 0,2467 & 0,0274 \\
1     & 0,2967 & 0,0298 \\
\hline
\end{tabular}
\end{center}

Baris pertama menunjukkan dua hal sekaligus: tanpa regularisasi galat CV paling buruk
dan paling tidak stabil, dengan simpangan $0{,}1245$, empat kali lipat baris lainnya.
Baris kedua sampai keempat hampir tidak terbedakan: $0{,}2400$, $0{,}2467$, dan
$0{,}2467$, dengan simpangan sekitar $0{,}025$, dan pada lima partisi berbeda pemenangnya
berganti-ganti. Kurva CV ini datar di sekitar minimumnya, sehingga memilih nilai
persisnya adalah ilusi ketelitian.

Akibatnya nyata: CV memilih $\lambda=0{,}003$ dengan galat validasi $0{,}2688$,
sedangkan $\lambda=0{,}1$ yang tidak terpilih memberi $0{,}2406$; selisih $0{,}0282$ itu
tidak jauh dari satu simpangan baku kurva CV. Pelajarannya bukan bahwa validasi silang
rusak, melainkan bahwa dengan enam puluh contoh, pemilihan hiperparameter membawa
ketidakpastian sebesar efek yang diukur. Melaporkan "$\lambda$ terbaik $0{,}003$" tanpa
menyebut lebar kurvanya adalah kebohongan yang paling halus dalam bab
ini.

\begin{tangan}
Bahan: sepuluh kartu, pensil, kertas. Tulis satu bilangan dua angka pada tiap kartu,
kocok, lalu bentangkan tujuh kartu menghadap atas. Susun ``aturan'' dalam bentuk daftar:
posisi satu bernilai sekian, posisi dua bernilai sekian, sampai posisi tujuh.

Tutup satu kartu dengan telapak tangan dan tebak. Karena daftarmu memuat kartu itu,
tebakanmu tepat. Ulangi untuk ketujuh posisi: tujuh dari tujuh benar, dan itulah nilai
latihanmu.

Sekarang kumpulkan sepuluh kartu, kocok, bentangkan tujuh kartu baru, dan pakai daftar
yang sama. Tutup satu kartu, tebak, buka. Yang harus terlihat: daftarmu tidak memuat
kartu baru, jadi tebakanmu hanya untung-untungan. Ulangi sepuluh kali dengan susunan
baru dan catat berapa kali tepat; itulah nilai ujianmu. Terakhir tulis aturan kedua yang
memakai satu sifat saja, misalnya ``kartu di kiri tidak lebih besar dari kartu di
kanannya'', lalu bandingkan kedua nilai itu.
\end{tangan}

Ada sumber kesalahan yang lebih besar daripada goyangan itu dan tidak terlihat pada angka
mana pun: kebocoran, yaitu masuknya informasi data uji ke dalam proses latihan. Bentuk
paling umumnya adalah standardisasi ciri memakai rata-rata dan simpangan baku seluruh data
sebelum pemisahan dilakukan.

Angkanya, dengan pergeseran covariate buatan: bila validasi digeser $0{,}5$,
standardisasi dengan statistik latih saja memberi $0{,}3765$, sedangkan statistik
gabungan memberi $0{,}3361$; pada geseran $1{,}0$ angkanya $0{,}5387$ lawan $0{,}3819$.
Statistik gabungan selalu tampak lebih baik, dan justru itu masalahnya: ia memakai
informasi letak data uji yang tidak tersedia di lapangan. Tanpa pergeseran keduanya sama,
$0{,}2452$ lawan $0{,}2463$, tanda bahwa selisih tadi berasal dari kebocoran.

Protokol yang jujur berlapis: pisahkan latih, validasi, dan uji sejak awal; hitung
semua besaran, termasuk rata-rata, simpangan baku, ambang keputusan, dan pemilihan ciri,
di dalam lipatan latih saja; pilih hiperparameter pada validasi atau lewat validasi
silang di dalam data latih; sentuh himpunan uji satu kali, lalu laporkan apa adanya.

Hubungan kapasitas dengan jumlah data dapat dibuat kuantitatif. Untuk kelas hipotesis
berhingga, jurang generalisasi terkendali oleh akar logaritma cacah hipotesis dibagi
cacah contoh.

\begin{teorema}[Batas seragam untuk kelas berhingga]
\index{kapasitas}
Misalkan $\mathcal{H}$ berhingga dengan $|\mathcal{H}|=m$, galatnya 0-1, dan $n$ contoh
diambil bebas dari $\mathcal{D}$. Dengan peluang sekurang-kurangnya $1-\delta$ terhadap
pengambilan contoh, berlaku serentak untuk semua $h\in\mathcal{H}$
\[
\bigl|R(h)-\hat R_n(h)\bigr|\le\sqrt{\frac{\ln m+\ln(2/\delta)}{2n}} .
\]
\end{teorema}

\begin{proof}[Bukti]
Untuk satu $h$ tetap, $\hat R_n(h)$ adalah rata-rata $n$ peubah Bernoulli berskala
$[0,1]$, sehingga ketaksamaan Hoeffding memberi
$\mathbb{P}(|R(h)-\hat R_n(h)|>t)\le2e^{-2nt^2}$. Gabungkan atas semua $h$ dengan
ketaksamaan union: peluang ada satu $h$ yang melampaui $t$ paling besar
$2m\,e^{-2nt^2}$. Menetapkan ungkapan itu sama dengan $\delta$ memberi
$t=\sqrt{(\ln m+\ln(2/\delta))/(2n)}$, dan pelengkapnya adalah pernyataan teorema.
\end{proof}

Untuk kelas yang "menghafal", yaitu model yang memberi satu label pada setiap pola biner
atas $d$ ciri, $m=2^d$. Dengan $d=21$ dan $\delta=0{,}05$ batasnya
$\sqrt{(21\ln2+\ln40)/(2n)}$, yaitu $0{,}3899$ pada $n=60$ dan $0{,}0617$ pada
$n=2400$. Jurang terukurnya $0{,}1338$ dan $0{,}0097$ pada $n$ yang sama, jadi batasnya
berlaku tetapi longgar; kelonggaran itu bukan cacat, sebab batas itu harus berlaku untuk
setiap distribusi dan setiap algoritma.

Yang penting bukan angkanya, melainkan bentuknya. Jurang mengecil sebanding
$1/\sqrt n$, sedangkan cacah hipotesis masuk lewat logaritma. Melipatgandakan data empat
kali hanya memangkas jurang menjadi separuh, dan menambah kapasitas sepuluh kali lipat
hanya menaikkan batas sebesar akar penambahan logaritmanya. Inilah alasan jujur mengapa
model besar masih berguna: selama jurangnya dikendalikan regularisasi, kapasitas yang
naik tidak otomatis berarti jurang yang naik.

\section{Contoh hitung: ridge tertutup pada empat titik}
\tingkat{menengah}

Semua rumus di atas dapat diperiksa tangan pada contoh kecil: empat titik dengan satu
ciri, $\mathbf{x}=(1,2,3,4)$ dan $\mathbf{y}=(2,4,6,9)$, model tanpa intersep
$\hat y=wx$, dan sasaran $\frac12\sum_i(y_i-wx_i)^2+\frac{\lambda}{2}w^2$. Turunannya
nol memberi
\[
w(\lambda)=\frac{\sum_ix_iy_i}{\sum_ix_i^2+\lambda}=\frac{64}{30+\lambda},
\]
sebab $\sum x_i^2=30$ dan $\sum x_iy_i=2+8+18+36=64$.

\begin{center}
\small
\begin{tabular}{rrrr}
\hline
$\lambda$ & $w$ & SSE latih & $|w\cdot5-10|$ \\
\hline
0   & 2,133333 & 0,466667 & 0,666667 \\
1   & 2,064516 & 0,608741 & 0,322581 \\
2   & 2,000000 & 1,000000 & 0,000000 \\
5   & 1,828571 & 3,253061 & 0,857143 \\
10  & 1,600000 & 9,000000 & 2,000000 \\
30  & 1,066667 & 34,600000 & 4,666667 \\
100 & 0,492308 & 81,255621 & 7,538462 \\
\hline
\end{tabular}
\end{center}

Gambar~\ref{fig:lintasan-bobot} memplot kedua lintasan bobot terhadap $\lambda$ pada
data ini: L2 mengecil tanpa henti, L1 berhenti tepat di nol. Kolom terakhir memakai satu
titik uji di $x=5$ dengan nilai sebenarnya $10$, yang tidak ikut menentukan $w$. Di
situlah kurvanya berbentuk U dan minimumnya jatuh tepat di $\lambda=2$: di situ $w=2$,
sehingga ramalan di $x=5$ adalah $10$ dan galatnya nol. Tanpa regularisasi,
$w=64/30=2{,}1333$ lebih memikat pada data latih (SSE $0{,}466667$ lawan $1{,}000000$)
tetapi meramal $10{,}667$, meleset $0{,}667$. Inilah perdagangan bias-variansi dalam
bentuk paling telanjang: SSE latih naik dua kali, galat uji menjadi nol.

\begin{figure}[htbp]
\centering
\begin{tikzpicture}
\begin{axis}[
  width=11.8cm, height=5.4cm,
  xmode=log,
  xmin=0.5, xmax=300, ymin=-0.12, ymax=2.45,
  axis lines=left,
  xtick={1,2,10,64,100},
  xticklabels={$1$,$2$,$10$,$64$,$100$},
  ytick={0,0.5,1.0,1.5,2.0},
  yticklabels={$0$,$0{,}5$,$1{,}0$,$1{,}5$,$2{,}0$},
  tick label style={font=\scriptsize},
  xlabel={kekuatan regularisasi $\lambda$ (skala log)},
  ylabel={bobot $w$},
  xlabel style={font=\small},
  ylabel style={font=\small},
  legend style={font=\scriptsize,draw=none,at={(0.97,0.97)},anchor=north east},
]
\addplot[domain=0.5:300,samples=200,very thick,layerindigo] {64/(30+x)};
\addlegendentry{L2: $64/(30+\lambda)$}
\addplot[domain=0.5:63.9,samples=200,very thick,layerorange] {(64-x)/30};
\addlegendentry{L1: $\max\{(64-\lambda)/30,\,0\}$}
\addplot[domain=64:300,samples=2,very thick,layerorange] {0};
\draw[densely dashed,layergrey] (axis cs:2,-0.12)--(axis cs:2,2.20);
\draw[densely dashed,layergrey] (axis cs:64,-0.12)--(axis cs:64,2.20);
\addplot[only marks,mark=*,mark size=1.8pt,layerindigo] coordinates {(2,2.0000) (64,0.6809)};
\addplot[only marks,mark=square*,mark size=1.8pt,layerorange] coordinates {(2,2.0667)};
\addplot[only marks,mark=o,mark size=3.4pt,layerorange] coordinates {(64,0)};
\node[font=\scriptsize,anchor=south] at (axis cs:2,2.24) {$\lambda=2$};
\node[font=\scriptsize,anchor=south] at (axis cs:64,2.24) {$\lambda=64$};
\node[font=\scriptsize,anchor=south west] at (axis cs:70,0.06) {L1 tepat nol};
\end{axis}
\end{tikzpicture}
\caption{Lintasan bobot terhadap kekuatan regularisasi pada contoh hitung empat titik,
dengan $\sum x_i^2=30$ dan $\sum x_iy_i=64$. Penalti L2 memberi
$w=64/(30+\lambda)$, yang mengecil tetapi tidak pernah menyentuh nol. Penalti L1 memberi
$w=\max\{(64-\lambda)/30,\,0\}$, yang mencapai nol tepat pada $\lambda=64$. Pada
$\lambda=2$ keduanya hampir sama, $2{,}0000$ lawan $2{,}0667$; pada $\lambda=64$ L2 masih
$0{,}6809$ sedangkan L1 sudah $0{,}0000$. Garis putus-putus menandai kedua nilai
$\lambda$ yang dibandingkan.}
\label{fig:lintasan-bobot}
\end{figure}

Perhatikan dinding kanannya. Pada $\lambda=100$, $w$ tinggal $0{,}4923$, SSE latih
meledak menjadi $81{,}255621$, dan galat uji $7{,}538462$, lebih buruk daripada model
yang tidak belajar apa pun. Contoh ini juga menunjukkan mengapa mencari $\lambda$ pada
data latih menyesatkan: kurva SSE latih monoton naik, jadi ia selalu memilih
$\lambda=0$.

\begin{tantangan}
Sebuah model melaporkan galat latih $0{,}0000$ pada $60$ contoh dan galat uji
$0{,}4100$. Rekannya mengusulkan tiga perubahan: (i) melatih lebih lama, (ii) menambah
ciri sampai derajat tujuh, (iii) menaikkan $\lambda$ dari $0$ ke $0{,}1$. Untuk setiap
usulan, ramalkan tanda perubahan galat latih dan galat uji, lalu sebutkan mana yang
tidak mungkin menolong beserta alasannya.

Petunjuk 1: tanyakan lebih dulu ke arah mana kapasitas bergerak pada tiap usulan.

Petunjuk 2: dua dari tiga usulan memperbesar kapasitas, dan model ini sudah berada di
galat latih nol, sehingga memperbesar kapasitas hanya memperlebar jurang. Yang menolong
adalah yang mengecilkannya, dan itu harus menaikkan galat latih.

Kalau mentok, kembali ke Definisi underfitting dan overfitting pada bagian pertama.
\end{tantangan}

\kertas{Periksa kolom kedua tanpa kalkulator: $64/32=2$ tepat, $64/31=2{,}064516$, dan
$64/40=1{,}6$. Lalu hitung sendiri SSE pada $\lambda=2$: sisaannya $(0,0,0,1)$, jadi
SSE $=1$. Terakhir jawab: berapa $\lambda$ yang membuat $w=1{,}28$, dan berapa ramalan
di $x=5$ ketika itu? Petunjuk: selesaikan $64/(30+\lambda)=1{,}28$.}

\section{Lab komputasi}
\tingkat{lanjut}

Seluruh angka pada bab ini berasal dari satu skrip numpy tanpa pustaka pembelajaran
mesin. Datanya sengaja kecil supaya dapat diperiksa ulang dalam hitungan menit: dua ciri
mentah, dua puluh satu ciri polinomial sampai derajat lima, enam puluh contoh latih, dan
dua puluh ribu contoh validasi dari distribusi yang sama. Labelnya tidak deterministik,
sehingga ada galat tak tereduksikan sekitar $0{,}21$ yang tidak dapat dilampaui model apa
pun.

\begin{lab}
\begin{verbatim}
import numpy as np

def data(n, rng):
    x = rng.uniform(-1, 1, size=(n, 2))
    p = 1/(1 + np.exp(-4.0*(0.9 - x[:,0]**2 - x[:,1]**2)))
    return x, (rng.random(n) < p).astype(float)

def poly(x, deg=5):
    return np.column_stack([x[:,0]**i * x[:,1]**j
                            for i in range(deg+1) for j in range(deg+1-i)])

def sig(z):
    return 1/(1 + np.exp(-np.clip(z, -400, 400)))

def soft(t, l):
    return np.sign(t)*np.maximum(np.abs(t) - l, 0.0)

def er(Z, y, w, b):
    return float(np.mean((sig(Z@w + b) > 0.5) != y))

def fit(Z, y, lam, iters=30000, lr=0.05):
    n, d = Z.shape
    w = np.zeros(d); b = 0.0
    mw = np.zeros(d); vw = np.zeros(d); mb = 0.0; vb = 0.0
    for t in range(1, iters+1):
        p = sig(Z@w + b); g = (p - y)/n
        gw = Z.T@g + lam*w; gb = g.sum()
        mw = 0.9*mw + 0.1*gw; vw = 0.999*vw + 0.001*gw*gw
        mb = 0.9*mb + 0.1*gb; vb = 0.999*vb + 0.001*gb*gb
        w -= lr*(mw/(1-0.9**t))/(np.sqrt(vw/(1-0.999**t)) + 1e-8)
        b -= lr*(mb/(1-0.9**t))/(np.sqrt(vb/(1-0.999**t)) + 1e-8)
    return w, b

def fit_l1(Z, y, lam, iters=150000):
    n, d = Z.shape
    w = np.zeros(d); b = 0.0
    st = 0.5/(0.25*np.linalg.eigvalsh(Z.T@Z/n).max())
    for _ in range(iters):
        p = sig(Z@w + b); g = (p - y)/n
        w = soft(w - st*(Z.T@g + lam*w), st*lam); b -= st*g.sum()
    return w, b
\end{verbatim}
\end{lab}

Bagian pertama menghitung galat untuk sebelas nilai $\lambda$; bagian kedua mengulanginya
dengan penalti L1 lewat penurunan gradien proksimal, yaitu langkah gradien yang diikuti
ambang lunak dari lemma tadi; bagian ketiga berhenti lebih awal; bagian keempat dan
kelima membandingkan pengoptimum; bagian keenam sampai kedelapan memeriksa validasi
silang, kebocoran, dan kapasitas. Keluaran sebenarnya, dipotong hanya pada bagian yang
berulang:

\begin{verbatim}
d=21  n_latih=60  n_validasi=20000
==1 L2: lam  |w|max  latih  validasi
  0.0000 406.204  0.0000  0.3262
  0.0010   3.056  0.1000  0.2927
  0.0030   1.814  0.1000  0.2688
  0.0100   1.018  0.1167  0.2475
  0.0300   0.631  0.1167  0.2433
  0.1000   0.399  0.1500  0.2406
  0.3000   0.268  0.2333  0.2452
  1.0000   0.142  0.2500  0.2695
  3.0000   0.062  0.3333  0.3205
 10.0000   0.022  0.3333  0.3395
100.0000   0.003  0.3333  0.3395
==2 L1: lam  nol/total  latih  validasi
  0.0020 16/21  0.1167  0.2630
  0.0050 11/21  0.1167  0.2455
  0.0100 14/21  0.1167  0.2412
  0.0200 14/21  0.1500  0.2369
  0.0500 11/21  0.1833  0.2334
  0.1000  7/21  0.2667  0.2492
  0.2000  1/21  0.3333  0.3395
==3 early stopping: epoch  latih  validasi  |w|max
ep=   1  0.1333  0.2481  0.15
ep=   5  0.1333  0.2429  0.38
ep=  10  0.1333  0.2432  0.48
ep=  50  0.1167  0.2454  0.77
ep= 200  0.1167  0.2541  1.38
ep= 800  0.1000  0.2879  2.85
ep=2400  0.0833  0.3109  4.81
terbaik 0.2427 pada epoch 8
==4 kuadratik diag(1,1e4): langkah ke galat 1e-8
GD=1e-4  92099 | momentum  9132 | Adam=0.1  82
==5 logistik minibatch: langkah ke galat latih 0.40 (median 5 seed)
sgd  lr=0.05 [100, 100, 90, 90, 100] median=100
sgd  lr=0.2  [30, 50, 40, 40, 20] median=40
adam lr=0.03 [30, 20, 20, 20, 30] median=20
adam lr=0.1  [20, 10, 20, 10, 10] median=10
==6 CV 5 partisi x 5 lipatan: lam  galat CV (rata2 +- sd antarpartisi)
 0.000  0.3233 +- 0.1245
 0.003  0.2400 +- 0.0253
 0.030  0.2467 +- 0.0217
 0.100  0.2467 +- 0.0274
 1.000  0.2967 +- 0.0298
==7 standardisasi: geser  statistik latih  statistik gabungan
0.0  0.2452  0.2463
0.5  0.3765  0.3361
1.0  0.5387  0.3819
==8 kapasitas vs data (d=21, lam=0, himpunan bersarang)
n=   60 latih=0.1667 validasi=0.3005 selisih=+0.1338 batas=0.3899
n=  150 latih=0.2133 validasi=0.2407 selisih=+0.0274 batas=0.2466
n=  600 latih=0.2067 validasi=0.2260 selisih=+0.0193 batas=0.1233
n= 2400 latih=0.2058 validasi=0.2155 selisih=+0.0097 batas=0.0617
\end{verbatim}

Tiga hal perlu dibaca dari keluaran itu, dan ketiganya mengecewakan harapan biasa.

Pertama, galat latih hampir tidak berguna sebagai petunjuk: ia bergerak berlawanan arah
dengan galat validasi. Kurva validasinya juga tidak simetris. Dari $\lambda=0$ ke
$0{,}1$ galatnya turun $0{,}086$; dari $0{,}1$ ke $10$ ia naik $0{,}099$ sampai melewati
titik awalnya. Bila ragu antara dua nilai, yang lebih kecil lebih aman.

Kedua, pemilihan lewat validasi silang tidak otomatis benar. CV memilih
$\lambda=0{,}003$, sedangkan validasi memilih $0{,}1$; selisih galat CV keduanya
$0{,}0006$, jauh di bawah simpangan antarpartisi $0{,}025$, padahal di data validasi
keduanya berbeda $0{,}0282$.

Ketiga, batas teoretis longgar tetapi arahnya benar. Eksperimen kapasitas di sini
melatih model dengan anggaran enam puluh ribu langkah pada himpunan bersarang, jadi galat
latihnya tidak nol. Pada $n=60$ batasnya $0{,}3899$ sedangkan jurangnya $0{,}1338$; pada
$n=2400$ batasnya $0{,}0617$ sedangkan jurangnya $0{,}0097$. Rasionya sekitar tiga,
sehingga batas itu berguna untuk membandingkan setelan, bukan meramal angka.

\begin{untukdosen}
Urutan mengajar yang saya sarankan: mulai dari angka pembuka bab, enam puluh contoh
dengan galat latih nol dan galat validasi $0{,}3262$, lalu tanyakan apakah model itu
"belajar". Jawaban kelas biasanya membelah dua, dan perpecahan itu bahan kuliah terbaik.
Sesudah itu definisi formal, dan biarkan dekomposisi bias-variansi diturunkan di papan
dalam lima baris. Regularisasi L1 sebaiknya masuk lewat lemma ambang lunak, bukan lewat
gambar bola L1, sebab lemma itu dapat diperiksa mahasiswa sendiri.

Dua miskonsepsi yang paling sering muncul. Mahasiswa menganggap galat latih rendah sebagai
bukti model bagus, lalu terkejut ketika galat validasinya tinggi; gejala lain, mereka
menyarankan "melatih lebih lama" sebagai obat, padahal pada tabel berhenti lebih awal
justru itulah yang merusak. Dan mereka menganggap regularisasi selalu menolong, lalu
memasang $\lambda$ besar dengan keyakinan "lebih aman". Tabel $\lambda=10$ dengan galat
$0{,}3395$ adalah obat untuk keduanya.

Aktivitas papan tulis yang berhasil: minta kelas menebak, sebelum dihitung, di mana
minimum kurva validasi berada, lalu hitung bersama tiga baris saja; tebakan biasanya
meleset ke $\lambda$ yang terlalu kecil. Tanda kelas siap lanjut: mereka dapat menjelaskan
mengapa memilih $\lambda$ memakai data validasi membuat angka itu optimistis, dan dapat
menyebutkan satu kebocoran yang tidak terlihat pada angka.
\end{untukdosen}

\section{Ringkasan}
\tingkat{dasar}

\begin{ringkasan}
\begin{itemize}
  \item Galat latih dan galat populasi dihitung dengan cara sama pada himpunan berbeda;
        galat latih $0{,}0000$ dengan galat validasi $0{,}3262$ menandakan hafalan.
  \item Galat populasi terdekomposisi menjadi derau tak tereduksi, bias kuadrat, dan
        variansi; regularisasi menukar bias dengan variansi, dan perdagangan itu
        menguntungkan hanya selama penurunan variansi lebih besar daripada kenaikan bias.
  \item Penalti L2 menambah $\lambda\mathbf{w}$ pada gradien, setara mengalikan bobot
        dengan $1-\eta\lambda$ tiap langkah; L1 memakai ambang lunak sehingga bobot di
        bawah $\lambda/H$ menjadi nol tepat. Keduanya punya dosis: $\lambda=0{,}1$
        memberi $0{,}2406$, sedangkan $\lambda=10$ memberi $0{,}3395$, lebih buruk
        daripada tanpa regularisasi.
  \item Berhenti lebih awal adalah regularisasi lewat jumlah langkah: $0{,}2427$ pada
        epoch 8 lawan $0{,}3109$ pada epoch 2400. Dropout menambahkan derau bervariansi
        $(1-p)/p$ per unit, dan augmentasi derau Gauss $\sigma^2$ setara penalti L2
        berskala $\sigma^2$.
  \item Pengoptimum berbeda pada masalah berskala buruk: pada kuadratik berbilangan
        kondisi $10^4$, penurun gradien perlu $92099$ langkah, momentum $9132$, Adam
        $82$; jadwal kosinus menurunkan galat latih dari $0{,}3866$ ke $0{,}3549$.
  \item Validasi silang mengurangi goyangan tanpa menghapusnya: kurva CV datar dengan
        simpangan $0{,}025$, dan pilihannya meleset $0{,}0282$. Kebocoran standardisasi
        membuat angka tampak lebih baik ($0{,}3361$ lawan $0{,}3765$) tanpa benar-benar
        lebih baik, dan jurang generalisasi menyusut sebanding $1/\sqrt n$: dari
        $0{,}1338$ pada $n=60$ ke $0{,}0097$ pada $n=2400$.
\end{itemize}
\end{ringkasan}

\section{Latihan}
\tingkat{menengah}

\begin{latihan}
\begin{enumerate}
  \item Nilai klaim ini: "Karena galat latih model kami nol, model itu sudah cukup baik."
        Benar, salah, atau belum cukup bukti? Sebutkan bagian yang keliru dan satu angka
        tambahan yang Anda perlukan.
  \item Buktikan dekomposisi bias-variansi untuk $y=f(x)+\varepsilon$ dengan
        $\operatorname{Var}(\varepsilon)=\sigma^2$: tuliskan ketiga suku, tunjukkan kedua
        suku silang lenyap, dan sebutkan di mana kebebasan $D$ dari $\varepsilon$
        dipakai.
  \item Turunkan $w(\lambda)=64/(30+\lambda)$ untuk data pada contoh hitung, lalu
        buktikan $w$ turun monoton terhadap $\lambda$ tetapi tidak pernah nol. Bandingkan
        dengan penalti L1 pada masalah satu dimensi yang sama, dan tentukan $\lambda$
        yang membuat $w=0$ tepat.
  \item Hitung galat validasi untuk $c=0{,}95$ dan $c=1{,}05$ pada contoh penaksir
        menyusut dengan $\mu=2$, $\sigma=1$, $n=25$, lalu bandingkan dengan $c=1$.
        Jelaskan mengapa dua simpangan sama besar tidak memberi kerugian sama besar.
  \item Sebuah laporan menulis: "Kami memilih $\lambda$ terbaik lewat validasi silang,
        jadi angka validasi silang itu adalah perkiraan galat di lapangan." Perbaiki
        kalimat itu, dan sebutkan dua prosedur tambahan yang diperlukan.
  \item Buktikan lemma ambang lunak untuk $a<0$ secara lengkap, lalu jelaskan mengapa
        hasilnya simetris terhadap tanda $a$. Petunjuk: kerjakan kasus $u<0$ dan $u=0$
        terpisah.
  \item Hitung langkah peluruhan bobot yang diperlukan agar bobot menjadi seperempat
        nilai awalnya bila $\eta=0{,}1$ dan $\lambda=0{,}02$, memakai waktu paruh
        $\ln2/(\eta\lambda)$. Periksa dengan simulasi numpy untuk gradien nol.
  \item Pada data bab ini, $\lambda=0{,}1$ memberi galat latih $0{,}1500$ dan galat
        validasi $0{,}2406$. Hitung jurangnya, bandingkan dengan jurang pada $\lambda=0$
        dan $\lambda=10$, lalu jelaskan mengapa jurang kecil tidak cukup untuk
        menyimpulkan sebuah model lebih baik.
  \item Turunkan batas seragam $\sqrt{(\ln m+\ln(2/\delta))/(2n)}$ dari ketaksamaan
        Hoeffding dan union bound, lalu hitung untuk $d=10$, $n=200$, dan
        $\delta=0{,}01$. Berapa kali data harus dilipatgandakan agar batasnya separuh?
\end{enumerate}

Petunjuk untuk tiga butir tersulit. Butir 3: sasaran satu dimensinya
$\frac12\sum_i(y_i-wx_i)^2+\lambda|w|$; pada $w=0$ subgradiennya memuat nol bila
$\lambda\ge|\sum_ix_iy_i|=64$, jadi $w=0$ tepat pada $\lambda=64$. Butir 6: pada $u=0$
subgradiennya selang $[-a-t,\,-a+t]$, yang memuat nol bila $|a|\le t$. Butir 9: dengan
$m=2^{10}$ dan $\delta=0{,}01$ batasnya $\sqrt{(10\ln2+\ln200)/400}=0{,}1749$;
separuhnya menuntut empat kali data.
\end{latihan}

\begin{untukpeneliti}
Yang terbukti di bab ini: dekomposisi bias-variansi beserta syaratnya; kesetaraan
langkah peluruhan bobot dengan penalti L2; bentuk tertutup peminimum
$\frac H2(w-w^\ast)^2+\lambda|w|$ dan syarat subgradien yang membuat L1 menghasilkan nol
sejati; kesetaraan augmentasi derau Gauss dengan penalti L2 untuk model linear; rumus
harapan dan variansi topeng dropout; serta batas seragam untuk kelas hipotesis
berhingga. Yang hanya populer: bahwa regularisasi selalu menolong; bahwa Adam selalu
lebih cepat; bahwa berhenti lebih awal "sama dengan" L2 (itu aproksimasi untuk sasaran
kuadratik dengan laju kecil); bahwa cacah nol L1 mengukur kepentingan ciri.

Sumber primer: Goodfellow, Bengio, dan Courville, \emph{Deep Learning}, Bab 5
(kapasitas, bias-variansi), Bab 7 (penalti L1 dan L2, augmentasi, derau, berhenti lebih
awal, dropout), Bab 8 (SGD, momentum, Adam, jadwal laju, pemilihan hiperparameter).
Wasserman, \emph{All of Statistics}, Bab 13 dan Bab 20. Batas seragam adalah akibat
langsung ketaksamaan Hoeffding dan union bound.

Batas keberlakuan: seluruh angka dari satu himpunan sintetis dengan pemisahan tunggal
$60/20000$, sehingga tiap galat validasi membawa ketidakpastian orde $0{,}01$ sampai
$0{,}02$. Karena itu perbandingan L1 dan L2 di sini bukan rekomendasi. Batas seragam hanya
untuk kelas berhingga dan galat 0-1.

Pertanyaan terbuka: seberapa besar kebocoran normalisasi batch, yang menghitung
statistik dari minibatch sehingga mencampur informasi antarcontoh; bagaimana memilih
$\lambda$ ketika biaya kesalahan tidak simetris; dan apakah kurva CV yang datar dapat
diubah menjadi interval yang sah, bukan sekadar simpangan antarpartisi.
\end{untukpeneliti}

\peta{Bab 26 pada Goodfellow, Bengio, dan Courville, \emph{Deep Learning}: Bab 5
(kapasitas, overfitting dan underfitting, dekomposisi bias-variansi), Bab 7 (penalti L1
dan L2, weight decay, augmentasi data, derau sebagai regularisasi, berhenti lebih awal,
dropout), Bab 8 (SGD, momentum, Adam, jadwal laju, pemilihan hiperparameter). Pada
Wasserman, \emph{All of Statistics}: Bab 13 (pemilihan model dan validasi silang), Bab 20
(bias-variansi pada penghalusan). Turunan balik dan pelatihan jaringan dibahas di Bab 23;
pembacaan metrik evaluasi yang jujur, termasuk kebocoran dan pemisahan data, dilanjutkan
di Bab 25; optimasi konveks yang mendasari penalti dan subgradien ada di Bab 21.}

\chapter{Evaluasi Model yang Jujur}

Sebuah laporan pendek sampai di meja seorang analis keamanan. Isinya satu tabel: detektor
intrusi baru, diuji pada data lalu lintas jaringan yang diambil dari jaringan kampus
selama sebulan, mencapai akurasi $99$ persen. Angka itu dibingkai dengan huruf tebal dan
satu kalimat penutup: ``siap dipakai''.

Data itu berisi satu serangan untuk setiap seribu koneksi. Pada kumpulan yang sama,
program tiga baris yang selalu menjawab ``bersih'' memperoleh akurasi $99{,}9$ persen: ia
tidak membaca satu fitur pun dan tidak pernah salah menuduh. Detektor yang dilatih berjam-
jam itu kalah, dan laporannya tidak memuat satu angka pun yang membocorkan kekalahan itu.

Yang keliru adalah ringkasannya: satu bilangan mewakili empat bilangan yang perannya
berbeda, pada data yang komposisinya tidak disebut. Bab 26 menjelaskan mengapa galat uji
berbeda dari galat latih; bab ini menutup tiga celah: metrik apa yang jujur tentang apa
yang mereka abaikan, seberapa besar angkanya bergoyang, dan apakah data uji benar-benar
belum dilihat oleh prosedur yang memilih model.

\begin{untukanak}
Bayangkan sebuah detektor logam di gerbang sekolah. Setiap pagi ia berbunyi. Kemarin ia
berbunyi seribu kali, dan petugas memeriksa seribu tas. Ternyata hanya enam puluh tas
yang benar-benar membawa barang terlarang. Sembilan ratus empat puluh siswa lain harus
membuka tasnya tanpa alasan.

Sekarang datang petugas dinas yang bertanya: berapa persen jawaban alat itu yang benar?
Petugas gerbang menjawab dengan bangga: sembilan puluh sembilan persen. Ia tidak
berbohong; dari seratus ribu kantong, alat itu hanya salah pada seribu. Tetapi jawaban itu
tidak menjawab pertanyaan yang membuat siswa kesal: mereka kesal bukan karena alat itu
sering salah, melainkan karena hampir setiap kali berbunyi, ia salah.

Dua pertanyaan itu terdengar mirip dan jawabannya berbeda jauh: yang satu bertanya tentang
seluruh kantong, yang lain tentang kantong yang dibunyikan. Di bab ini kita akan tahu
mengapa satu alat bisa unggul pada pertanyaan pertama dan hancur pada pertanyaan kedua,
mengapa angka seratus persen kadang bukan kabar baik, dan mengapa laporan yang jujur selalu
memuat lebih dari satu angka.
\end{untukanak}

\section{Empat bilangan yang tidak simetris}
\tingkat{dasar}

Pengklasifikasi biner, secanggih apa pun isinya, meninggalkan jejak yang sama: empat
bilangan. Yang membedakan satu laporan evaluasi dari laporan lain adalah apakah keempatnya
dilaporkan. Kalimat sehari-harinya: setiap keputusan punya dua cara untuk benar dan dua
cara untuk salah, dan menutup salah satunya sama dengan bersembunyi.

\kamus{matriks konfusi}{confusion matrix}{buku kas empat kolom yang mencatat setiap
keputusan benar dan salah beserta arah kesalahannya}

\begin{definisi}[Matriks konfusi]
\index{matriks konfusi}
Misalkan setiap contoh memiliki label sebenarnya $y\in\{0,1\}$ dan prediksi
$\hat y\in\{0,1\}$, dengan $1$ berarti positif (serangan, penyakit, pesan berbahaya).
Empat bilangan
\begin{align}
\mathrm{TP}&=\#\{y=1,\ \hat y=1\}, & \mathrm{FP}&=\#\{y=0,\ \hat y=1\},
\nonumber\\
\mathrm{FN}&=\#\{y=1,\ \hat y=0\}, & \mathrm{TN}&=\#\{y=0,\ \hat y=0\}
\nonumber
\end{align}
beserta susunannya disebut \emph{matriks konfusi}. Jumlah seluruh contoh adalah
$n=\mathrm{TP}+\mathrm{FP}+\mathrm{FN}+\mathrm{TN}$.
\end{definisi}

Dari empat bilangan itu diturunkan metrik yang menjawab pertanyaan berbeda. Presisi
bertanya: bila model mengatakan positif, berapa peluang ia benar. Recall bertanya: dari
seluruh kasus positif, berapa bagian yang tertangkap.

\begin{definisi}[Presisi, recall, spesifisitas]
\index{presisi}
\index{recall}
Dengan $\mathrm{TP}+\mathrm{FP}>0$ dan $\mathrm{TP}+\mathrm{FN}>0$,
\[
\operatorname{prec}=\frac{\mathrm{TP}}{\mathrm{TP}+\mathrm{FP}},\qquad
\operatorname{rec}=\frac{\mathrm{TP}}{\mathrm{TP}+\mathrm{FN}},\qquad
\operatorname{spec}=\frac{\mathrm{TN}}{\mathrm{TN}+\mathrm{FP}} .
\]
Akurasi adalah $(\mathrm{TP}+\mathrm{TN})/n$, dan $F_1$ adalah
$2\,\mathrm{TP}/(2\,\mathrm{TP}+\mathrm{FP}+\mathrm{FN})$.
\end{definisi}

Mari kita pakai satu contoh kerja: sebuah detektor diuji pada $1000$ koneksi berlabel,
$100$ di antaranya serangan. Ia menandai $62$ serangan dengan benar, melewatkan
$38$, dan menuduh $17$ koneksi bersih sebagai serangan.

\begin{center}
\small
\begin{tabular}{l|rr|r}
\hline
 & prediksi positif & prediksi negatif & total \\
\hline
sebenarnya positif & $\mathrm{TP}=62$ & $\mathrm{FN}=38$ & 100 \\
sebenarnya negatif & $\mathrm{FP}=17$ & $\mathrm{TN}=883$ & 900 \\
\hline
total & 79 & 921 & 1000 \\
\hline
\end{tabular}
\end{center}

Presisinya $62/79=0{,}784810$: sekitar empat dari lima alarm menunjuk serangan.
Recall-nya $62/100=0{,}620000$: hampir empat dari sepuluh serangan tidak terlihat.
Spesifisitasnya $883/900=0{,}981111$. Akurasinya $945/1000=0{,}945000$, dan
$F_1=124/179=0{,}692737$. Kelima angka itu benar tentang model yang sama, dan tidak satu
pun darinya boleh dipakai sendiri. Pemeriksaan dengan \texttt{sklearn.metrics} memberi
kelima angka yang sama sampai enam angka desimal.

Gambar~\ref{fig:konfusi-jujur} menempatkan keempat bilangan itu pada tempatnya, lengkap
dengan jumlah baris dan kolomnya; setiap angka di sisa bab harus dapat dilacak ke salah
satu sel gambar itu.

\begin{figure}[htbp]
\centering
\begin{tikzpicture}[every node/.style={font=\footnotesize}]
  % dua sel benar (hijau) dan dua sel salah (merah)
  \fill[layerjade!14] (4.40,2.00) rectangle (6.70,2.90);
  \fill[layerjade!14] (6.70,1.10) rectangle (9.00,2.00);
  \fill[accent!8]     (6.70,2.00) rectangle (9.00,2.90);
  \fill[accent!8]     (4.40,1.10) rectangle (6.70,2.00);
  \draw[ink,thick] (4.40,0.55) rectangle (10.30,2.90);
  \draw[ink] (6.70,0.55) -- (6.70,2.90);
  \draw[ink] (9.00,0.55) -- (9.00,2.90);
  \draw[ink] (4.40,2.00) -- (10.30,2.00);
  \draw[ink] (4.40,1.10) -- (10.30,1.10);
  \draw[ink!35] (4.40,0.55) -- (10.30,0.55);
  % isi sel
  \node at (5.55,2.45) {$\mathrm{TP}=62$};
  \node at (7.85,2.45) {$\mathrm{FN}=38$};
  \node at (5.55,1.55) {$\mathrm{FP}=17$};
  \node at (7.85,1.55) {$\mathrm{TN}=883$};
  % judul kolom
  \node at (5.55,3.14) {prediksi positif};
  \node at (7.85,3.14) {prediksi negatif};
  \node at (9.65,3.14) {total};
  % judul baris
  \node[anchor=east] at (4.24,2.45) {sebenarnya positif $(100)$};
  \node[anchor=east] at (4.24,1.55) {sebenarnya negatif $(900)$};
  \node[anchor=east] at (4.24,0.82) {total};
  % jumlah
  \node at (5.55,0.82) {$79$};
  \node at (7.85,0.82) {$921$};
  \node at (9.65,2.45) {$100$};
  \node at (9.65,1.55) {$900$};
  \node at (9.65,0.82) {$1000$};
  \node[anchor=north,font=\scriptsize,align=center,text width=11.4cm] at (7.00,0.30)
    {akurasi $=(62+883)/1000=0{,}945$;\quad presisi $=62/79$;\quad recall $=62/100$};
\end{tikzpicture}
\caption{Matriks konfusi contoh kerja pada $1000$ koneksi: $62$ serangan tertangkap, $38$
lolos, $17$ alarm palsu, dan $883$ koneksi bersih yang ditolak dengan benar. Baris
menjumlah menjadi $100$ positif dan $900$ negatif, kolom menjadi $79$ prediksi positif dan
$921$ prediksi negatif, dan seluruhnya $1000$. Dua sel hijau adalah dua cara untuk benar,
dua sel merah dua cara untuk salah. Akurasi $0{,}945$ menjumlahkan kedua sel hijau dan
karena itu tidak dapat membedakan $62$ yang tertangkap dari $38$ yang lolos.}
\label{fig:konfusi-jujur}
\end{figure}

Peran keempat bilangan tidak simetris. Presisi menjawab pertanyaan operator: bila
alarm berbunyi, berapa peluang memang ada serangan. Recall menjawab pertanyaan
perancang kebijakan: dari seluruh serangan, berapa bagian yang tertangkap. Keduanya
bergantung pada ambang, sehingga metrik tanpa ambang bukan pernyataan yang lengkap:
model yang sama memberi presisi $0{,}25$ dengan recall $0{,}45$, atau presisi
$0{,}013$ dengan recall $0{,}86$.

$F_1$ mencoba meringkas presisi dan recall menjadi satu bilangan, dan justru di situ ia
paling mudah disalahgunakan.

\begin{proposisi}
\index{f1}
$F_1$ adalah rata-rata harmonik presisi $P$ dan recall $R$, yaitu $F_1=2PR/(P+R)$,
sehingga $F_1\le(P+R)/2$ dengan kesamaan berlaku tepat ketika $P=R$.
\end{proposisi}

\begin{proof}[Bukti]
Tulis $a=\mathrm{TP}$, $b=\mathrm{FP}$, $c=\mathrm{FN}$. Maka $P=a/(a+b)$,
$R=a/(a+c)$, dan
\[
\frac{2PR}{P+R}
=\frac{2a^2\big/\bigl((a+b)(a+c)\bigr)}{a(2a+b+c)\big/\bigl((a+b)(a+c)\bigr)}
=\frac{2a}{2a+b+c}=F_1 .
\]
Ketaksamaan harmonik--aritmetik $4PR\le(P+R)^2$ ekuivalen dengan $(P-R)^2\ge0$, dan
kesamaan hanya terjadi bila $P=R$.
\end{proof}

Sifat harmonik itu menarik nilainya ke bawah oleh komponen terkecil: presisi $0{,}9$
dengan recall $0{,}2$ memberi $F_1\approx0{,}327$, bukan $0{,}55$; pada contoh kita
$F_1=0{,}6927$, sedikit di bawah titik tengah $0{,}7024$.

Tiga cacat ringkasan perlu dinamai satu per satu, sebab obatnya berbeda. Presisi
buta terhadap serangan yang lolos: model yang menandai semua koneksi memperoleh
presisi $0{,}1$, yaitu kelazimannya sendiri, tetapi recall $1$. $F_1$ mengabaikan
$\mathrm{TN}$ sepenuhnya: dua model dengan presisi dan recall sama tetapi alarm
palsu berbeda sepuluh kali lipat melaporkan $F_1$ identik, padahal biaya
operasionalnya jauh berbeda. $F_1$ juga membebani kedua komponen setara, padahal
satu serangan yang lolos biasanya jauh lebih mahal daripada satu alarm palsu.

\begin{tantangan}
Sebuah laporan hanya memuat tiga angka: ``akurasi $0{,}9720$ pada $5000$ koneksi uji;
presisi $0{,}000$; model siap dipakai''. Tentukan seluruh matriks konfusi yang konsisten
dengan laporan itu, jelaskan mengapa recall-nya nol dan bukan bilangan lain, lalu cari
kelaziman terbesar yang masih mungkin.

Petunjuk 1: akurasi hanya menjumlahkan dua sel diagonal, jadi
$5000(1-0{,}9720)=140$ kesalahan bebas dibagi antara dua sel yang tersisa.

Petunjuk 2: presisi $0{,}000$ berarti $\mathrm{TP}=0$, sehingga $\mathrm{FN}$ sama dengan
banyaknya contoh positif dan $\mathrm{FP}$ adalah sisanya. Tetapi presisi dengan
$\mathrm{TP}=\mathrm{FP}=0$ adalah $0/0$, bukan $0$; angka $0{,}000$ pada laporan itu
karena itu menuntut sekurang-kurangnya satu alarm palsu.

Kalau mentok, kembali ke bagian Akurasi yang menipu pada bab ini.
\end{tantangan}

\section{Akurasi yang menipu pada data tak seimbang}
\tingkat{dasar}

Akurasi paling sering dilaporkan dan paling mudah menipu: ia menghitung semua
kesalahan dengan harga sama pada kumpulan yang komposisinya tidak disebut. Kalimat
sehari-harinya: akurasi adalah nilai rapor yang dihitung seolah kedua kelompok
seimbang.

\index{akurasi seimbang}
\begin{definisi}[Akurasi seimbang]
Misalkan $\pi$ adalah kelaziman kelas positif. \emph{Akurasi seimbang} adalah
$(\operatorname{rec}+\operatorname{spec})/2$, yaitu rata-rata laju benar pada kedua kelas.
Besaran ini tidak berubah bila jumlah contoh positif dan negatif diubah tanpa mengubah
kedua laju tersebut.
\end{definisi}

Akurasi biasa tidak memiliki sifat itu, dan akibatnya dapat ditulis tepat.

\begin{proposisi}
Pengklasifikasi yang selalu menjawab negatif memperoleh akurasi $1-\pi$. Karena itu
sebuah model dengan akurasi di bawah $1-\pi$ kalah dari garis dasar yang tidak belajar
apa pun, dan laporan yang menyebut akurasi tanpa menyebut $\pi$ tidak dapat dinilai.
\end{proposisi}

\begin{proof}[Bukti]
Dari $n$ contoh, $\pi n$ berlabel positif dan seluruhnya salah, sedangkan $(1-\pi)n$
berlabel negatif dan seluruhnya benar. Akurasinya $(1-\pi)n/n=1-\pi$. Model lain dengan
akurasi lebih kecil karena itu melakukan lebih banyak kesalahan daripada garis dasar itu.
\end{proof}

\begin{bacakeras}
Dua suara, dibaca bergantian dengan teman. Ketuk meja setiap kali kata ``nol'' diucapkan.

``Seratus kartu. Berapa merah?'' ``Sepuluh.'' ``Berapa putih?'' ``Sembilan puluh.''
``Mesin menjawab apa?'' ``Putih. Semua putih.'' ``Berapa yang benar?'' ``Sembilan puluh.''
``Jadi akurasinya?'' ``Sembilan puluh persen.'' ``Dan yang merah?'' ``Nol. Tidak satu pun.''
``Angka tadi bohong?'' ``Tidak. Ia hanya menjawab satu pertanyaan.''
``Pertanyaan apa?'' ``Berapa banyak yang kutebak benar, dari semua yang ada.''
``Pertanyaan mana yang tidak dijawab?'' ``Dari yang benar-benar berbahaya, berapa yang kutemukan.''

Ulangi tiga baris tengah, makin cepat, dengan ketukan yang sama: yang dihitung semua, yang
hilang tak disebut, yang hilang tak disebut.

Suara kedua menutup pelan: kelas yang sedikit mudah tenggelam di dalam rata-rata kelas
yang banyak, dan angka rata-rata itu tidak pernah menyebut siapa yang tenggelam.

Baris yang harus diulang pembaca, sekali, dengan suara jelas: akurasi menjumlahkan yang
benar, bukan yang penting, jadi pada kelas tak seimbang sembilan puluh persen bisa berarti
nol.
\end{bacakeras}

Pada contoh kerja kita $\pi=0{,}1$, jadi garis dasarnya $0{,}9$ dan model dengan
$0{,}945$ memang lebih baik. Sekarang ubah kelazimannya saja, tanpa mengubah mutu model,
dan lihat apa yang terjadi pada angka yang sama.

\begin{contoh}[Satu serangan untuk setiap seribu koneksi]
Ambil $100\,000$ koneksi dengan $100$ serangan. Skor pada koneksi bersih menyebar
normal dengan rataan $0$ dan simpangan baku $1$; skor pada serangan dengan rataan
$2{,}6$ dan simpangan baku $1$. Ambang dipasang pada $2{,}3263$, yaitu
kuantil $0{,}99$ dari sebaran skor negatif, sehingga secara rancangan sekitar satu persen
koneksi bersih akan ditandai. Hasil satu lintasan dengan benih tetap:
\[
\mathrm{TP}=62,\quad \mathrm{FP}=1013,\quad \mathrm{FN}=38,\quad \mathrm{TN}=98\,887 .
\]
Akurasinya $(62+98\,887)/100\,000=0{,}989490$. Garis dasarnya
$99\,900/100\,000=0{,}999000$. Presisinya $62/1075=0{,}057674$, recall
$62/100=0{,}620000$, dan akurasi seimbangnya $0{,}804930$.
\end{contoh}

Angka $0{,}9895$ itu, bila dibulatkan, menjadi ``$99$ persen''. Ia bukan hanya tidak
informatif: ia lebih rendah daripada garis dasar sebesar $0{,}009510$, hampir satu poin
persentase penuh. Model itu melakukan sepuluh kali lebih banyak kesalahan daripada program
yang selalu menjawab ``bersih'', $1051$ melawan $100$, dan seluruhnya berupa alarm palsu,
$1013$ buah atau sekitar $16{,}3$ per serangan yang tertangkap.

Di sini pembulatan menambah kerusakan: $0{,}9895$ dan $0{,}9990$ sama-sama menjadi
``$99$ persen'', padahal keduanya terpisah oleh seluruh selisih yang menentukan.
Laporan yang baik karena itu menuliskan tiga atau empat angka desimal, menyebut $n$,
dan menyebut kelazimannya; tanpa ketiganya, pembaca tidak tahu apakah $99$ persen
itu prestasi atau kegagalan.

Apakah akurasi lalu tidak boleh dipakai? Ia sah bila biaya kedua kesalahan setara,
kelaziman dekat setengah, dan ambangnya tetap. Di luar itu ia harus ditemani presisi dan
recall, atau diganti akurasi seimbang: $0{,}8049$ lebih jujur daripada $0{,}9895$, sebab
kelas minoritas hanya tertangkap $62$ persen.

\begin{tangan}
Bahan: $100$ kartu, $90$ bermuka putih dan $10$ bermuka merah, spidol, dan satu lembar
tally dua kolom.

Tulis ``bersih'' pada sembilan puluh kartu putih dan ``serangan'' pada sepuluh kartu merah.
Kocok, lalu bagi menjadi sepuluh tumpukan. Aturan pertama: setiap tumpukan dijawab
``bersih'' tanpa membalik satu kartu pun. Catat jawaban yang benar, lalu balik semua kartu:
sembilan puluh benar, sepuluh merah lolos. Tulis di papan: akurasi
sembilan puluh persen, recall nol persen.

Sekarang aturan kedua. Beri titik di sudut enam kartu merah, dan tanpa memberi tahu pemain,
beri titik yang sama di sudut empat kartu putih. Jawab ``serangan'' hanya bila titik itu
terlihat. Hitung ulang: enam merah tertangkap, empat merah lolos, empat putih dituduh,
delapan puluh enam putih benar.

Yang harus terlihat: akurasi hanya bergerak dari sembilan puluh ke sembilan puluh dua
persen, sedangkan recall melompat dari nol ke enam puluh persen. Angka yang berubah paling
banyak justru angka yang tidak ada pada laporan pertama.
\end{tangan}

\section{Kurva ROC, AUC, dan prevalensi}
\tingkat{menengah}

Setiap ambang memberi satu titik pada dua laju yang berbeda; memilih ambang tanpa melihat
seluruh kemungkinannya adalah menebak. Kalimat sehari-harinya: kita menyapu ambang
dari longgar ke ketat, lalu menggambar jejak yang ditinggalkan model.

\kamus{kurva roc}{receiver operating characteristic}{jejak yang ditinggalkan sebuah model
ketika ambangnya disapu dari longgar ke ketat, digambar pada dua laju yang tidak
bergantung pada kelaziman}

\begin{definisi}[Kurva ROC dan AUC]
\index{kurva roc}
\index{auc}
Untuk ambang $t$, laju benar positif dan laju salah positif adalah
\[
\mathrm{TPR}(t)=P(S>t\mid y=1),\qquad \mathrm{FPR}(t)=P(S>t\mid y=0),
\]
dengan $S$ skor model. Kurva ROC adalah himpunan titik $(\mathrm{FPR}(t),\mathrm{TPR}(t))$
untuk seluruh $t$. Luas di bawah kurva itu disebut AUC, dan untuk data terhingga ia
dihitung sebagai luas trapesium di antara titik-titik yang teramati.
\end{definisi}

AUC punya arti yang jauh lebih sederhana daripada namanya, dan arti itu yang membuatnya
populer.

\begin{proposisi}
Untuk data tanpa skor kembar, AUC sama dengan peluang bahwa skor satu contoh positif
yang diambil acak melebihi skor satu contoh negatif yang diambil acak. Bila ada skor
kembar, tiap pasangan kembar dihitung setengah.
\end{proposisi}

\begin{proof}[Bukti]
Susun seluruh skor dari besar ke kecil dan sapu ambang ke bawah. Setiap kali ambang
melewati satu contoh positif, $\mathrm{TPR}$ naik sebesar $1/n_+$ pada ketinggian
$\mathrm{FPR}$ saat itu, yaitu bagian contoh negatif yang sudah terlewati. Setiap kali
ambang melewati contoh negatif, $\mathrm{FPR}$ naik pada $\mathrm{TPR}$ tetap dan tidak
menambah luas. Karena itu luas totalnya
\[
\mathrm{AUC}=\frac1{n_+}\sum_{i:\ y_i=1}\mathrm{FPR}_i
=\frac{1}{n_+n_-}\#\{(i,j):\ y_i=1,\ y_j=0,\ S_i>S_j\},
\]
yaitu peluang yang diklaim, dan suku terakhir itulah yang dihitung oleh statistik
Mann--Whitney. Skor kembar tidak menambah luas, sehingga tiap pasangan seperti itu
menyumbang setengah.
\end{proof}

AUC tidak melihat kelaziman: bila setiap contoh negatif digandakan sepuluh kali, urutan
skor tidak berubah dan AUC tetap. Cara menghitungnya di kertas diperlihatkan pada bagian
Contoh hitung.

Sekarang bagian yang membuat ROC berbahaya: sapu ambang simulasi satu banding seribu.

\begin{center}
\small
\begin{tabular}{rrrrr}
\hline
ambang & recall & FPR & presisi & akurasi \\
\hline
3,0000 & 0,450 & 0,00132 & 0,254 & 0,99813 \\
2,8000 & 0,510 & 0,00263 & 0,162 & 0,99688 \\
2,6000 & 0,550 & 0,00467 & 0,105 & 0,99488 \\
2,3263 & 0,620 & 0,01014 & 0,058 & 0,98949 \\
2,0000 & 0,740 & 0,02269 & 0,032 & 0,97707 \\
1,5000 & 0,860 & 0,06729 & 0,013 & 0,93264 \\
\hline
\end{tabular}
\end{center}

Pada kurva ROC, baris-baris itu bergerak dari $(0{,}001;0{,}45)$ ke $(0{,}067;0{,}86)$:
jejak yang naik tajam di dekat sumbu tegak, persis seperti yang diinginkan setiap
praktisi. AUC-nya $0{,}966448$. Tetapi kolom presisi mengatakan hal lain. Bahkan pada
ambang paling ketat, tempat hanya $132$ koneksi bersih ditandai, presisinya $0{,}254$:
tiga dari empat alarm tetap salah. Ketika recall dinaikkan ke $0{,}86$, presisinya
$0{,}013$.

Gambar~\ref{fig:presisi-recall-ambang} menggambar keenam baris tabel itu pada satu sumbu
ambang: recall menanjak sementara presisi menurun, dan tidak ada titik yang memberi keduanya
tinggi sekaligus.

\begin{figure}[htbp]
\centering
\begin{tikzpicture}
\begin{axis}[
  width=0.78\linewidth, height=6.2cm,
  x dir=reverse, xmin=1.45, xmax=3.05,
  ymin=0.35, ymax=0.95, ytick={0.4,0.5,0.6,0.7,0.8,0.9},
  xlabel={ambang $t$ (ke kanan: ambang dilonggarkan)}, ylabel={recall},
  axis y line*=left, axis x line*=bottom,
  tick label style={font=\footnotesize}, label style={font=\footnotesize},
  grid=major, grid style={layergrey!15},
]
\addplot[accent,thick,mark=*,mark size=1.8pt] coordinates {
  (3.0,0.450) (2.8,0.510) (2.6,0.550) (2.3263,0.620) (2.0,0.740) (1.5,0.860)};
\node[accent,font=\footnotesize,anchor=south west] at (axis cs:2.62,0.665) {recall};
\end{axis}
\begin{axis}[
  width=0.78\linewidth, height=6.2cm,
  x dir=reverse, xmin=1.45, xmax=3.05,
  ymin=0, ymax=0.30, ytick={0,0.1,0.2,0.3},
  axis y line*=right, axis x line*=none, ylabel={presisi (skala kanan)},
  tick label style={font=\footnotesize}, label style={font=\footnotesize},
]
\addplot[layerteal,thick,mark=square*,mark size=1.6pt] coordinates {
  (3.0,0.254) (2.8,0.162) (2.6,0.105) (2.3263,0.058) (2.0,0.032) (1.5,0.013)};
\node[layerteal,font=\footnotesize,anchor=north east] at (axis cs:2.78,0.232) {presisi};
\end{axis}
\end{tikzpicture}
\caption{Presisi dan recall pada enam ambang dari tabel sebelumnya, digambar pada satu
sumbu ambang yang diputar sehingga ke kanan ambang dilonggarkan. Recall menanjak dari
$0{,}450$ pada $t=3{,}0$ ke $0{,}860$ pada $t=1{,}5$, sedangkan presisi menurun dari
$0{,}254$ ke $0{,}013$: kedua kurva bergerak berlawanan arah dan tidak ada satu titik pun
yang memberi keduanya tinggi. Skala presisi ada di sisi kanan, $0$ sampai $0{,}30$; pada
$t=2{,}3263$, ambang yang dipakai simulasi, presisi $0{,}058$ dengan recall $0{,}620$.}
\label{fig:presisi-recall-ambang}
\end{figure}

Kedua bacaan itu tidak bertentangan: keduanya mengukur hal berbeda.

\begin{proposisi}
Kurva ROC dan AUC hanya bergantung pada sebaran skor bersyarat $S\mid y=1$ dan
$S\mid y=0$, sehingga tidak berubah bila kelaziman $\pi$ diubah. Presisi pada ambang
tetap berubah, dan mengikuti
\[
\operatorname{prec}(t)=\frac{\pi\,\mathrm{TPR}(t)}
{\pi\,\mathrm{TPR}(t)+(1-\pi)\,\mathrm{FPR}(t)} .
\]
\end{proposisi}

\begin{proof}[Bukti]
Menurut definisi, $\mathrm{TPR}$ dan $\mathrm{FPR}$ adalah peluang bersyarat pada $y$,
jadi keduanya tidak memuat $\pi$; kurva yang tersusun dari keduanya pun tidak memuat
$\pi$. Untuk presisi, hukum peluang total memberi
$P(\hat y=1)=\pi\,\mathrm{TPR}+(1-\pi)\mathrm{FPR}$, sedangkan
$P(y=1,\hat y=1)=\pi\,\mathrm{TPR}$. Aturan Bayes dari Bab 17 memberi bentuk di atas.
\end{proof}

Akibatnya dapat diukur. Bila kelaziman diubah menjadi setengah tanpa mengubah
sebaran skor bersyarat, AUC tinggal $0{,}967061$ (selisih $0{,}0006$ hanyalah galat
Monte Carlo), sedangkan \emph{average precision} melompat dari $0{,}287636$ menjadi
$0{,}967544$. Pada kelaziman satu banding seribu garis dasarnya $\pi=0{,}001$, jadi
$0{,}288$ masih jauh lebih baik daripada menebak, tetapi jauh lebih dekat kepada
kenyataan operasional daripada $0{,}966$.

Aturan praktisnya sederhana: bila kelas positif langka, laporkan kurva presisi--recall,
bukan hanya ROC, dan sebutkan kelazimannya. ROC tetap berguna untuk membandingkan dua model
pada populasi yang sama sebab ia tidak tercemar komposisi; ia menyesatkan begitu dipakai
menyimpulkan kegunaan pada populasi lain atau pada ambang tertentu. AUC juga buta terhadap
kalibrasi: model yang skornya dua kali terlalu besar memiliki AUC yang sama persis, sebab
urutannya tidak berubah. Dua cacat itu membawa kita ke bagian berikutnya.

\section{Kalibrasi: apakah angka itu berarti peluang}
\tingkat{menengah}

Operator yang menerima alarm perlu tahu satu hal yang tidak dijawab presisi maupun AUC:
bila model menulis $0{,}8$, apakah delapan dari sepuluh kasus seperti itu memang positif.
Kalimat sehari-harinya: kalibrasi memeriksa apakah ramalan cuaca yang mengatakan ``tujuh
puluh persen hujan'' memang benar pada tujuh dari sepuluh hari seperti itu.

\kamus{kalibrasi}{calibration}{kesesuaian antara angka yang dijanjikan model dan
kenyataan yang benar-benar terjadi pada kelompok dengan janji yang sama}

\begin{definisi}[Kalibrasi sempurna]
\index{kalibrasi}
Skor $S\in[0,1]$ disebut terkalibrasi sempurna bila
$P(y=1\mid S=s)=s$ untuk setiap $s$ yang mungkin. Pada data terhingga, kalibrasi diperiksa
dengan mengelompokkan skor ke dalam selang, lalu membandingkan rataan skor di selang itu
dengan laju positif yang teramati. Selisih mutlaknya, ditimbang ukuran selang, disebut
\emph{expected calibration error}.
\end{definisi}

Contoh kerja dengan $1000$ prediksi memperlihatkan bentuk pemeriksaannya.

\begin{center}
\small
\begin{tabular}{lrrrr}
\hline
selang skor & banyak & positif & rataan skor & laju positif \\
\hline
$[0{,}5;0{,}6)$ & 40 & 21 & 0,55 & 0,525 \\
$[0{,}6;0{,}7)$ & 60 & 36 & 0,65 & 0,600 \\
$[0{,}7;0{,}8)$ & 120 & 90 & 0,75 & 0,750 \\
$[0{,}8;0{,}9)$ & 200 & 176 & 0,85 & 0,880 \\
$[0{,}9;1{,}0]$ & 580 & 551 & 0,95 & 0,950 \\
\hline
\end{tabular}
\end{center}

Selisih terbesar ada di selang kedua, $0{,}05$, dan kesalahan kalibrasi keseluruhannya
\[
\tfrac{40}{1000}(0{,}025)+\tfrac{60}{1000}(0{,}050)+\tfrac{120}{1000}(0)
+\tfrac{200}{1000}(0{,}030)+\tfrac{580}{1000}(0)=0{,}010 .
\]
Arah kesalahannya berbeda antar selang, tetapi besarnya kecil: satu poin persentase. Pemeriksaan silangnya gratis: banyaknya
contoh di kelima selang berjumlah $1000$, dan positifnya $874$.

Kalibrasi diukur oleh galat kuadrat, dan di situ ada identitas yang rapi.

\begin{proposisi}[Brier untuk skor terkalibrasi]
\index{brier}
Misalkan $S$ terkalibrasi sempurna dan $B=\mathbb{E}[(S-y)^2]$ adalah skor Brier. Maka
\[
B=\mathbb{E}\bigl[S(1-S)\bigr].
\]
Khususnya skor konstan $S=\pi$ terkalibrasi sempurna dengan $B=\pi(1-\pi)$, dan skor
konstan itu adalah peramal terkalibrasi yang paling tidak tajam.
\end{proposisi}

\begin{proof}[Bukti]
Bersyarat pada $S=s$, peubah $y$ bernilai $1$ dengan peluang $s$. Maka
\[
\mathbb{E}\bigl[(s-y)^2\mid S=s\bigr]
=s(s-1)^2+(1-s)s^2=s(1-s)\bigl[(1-s)+s\bigr]=s(1-s).
\]
Mengambil harapan atas $S$ memberi bentuk pertama. Untuk $S=\pi$ konstan, syarat
kalibrasi menuntut $P(y=1)=\pi$, yang memang benar, dan rumusnya memberi $\pi(1-\pi)$.
\end{proof}

Identitas itu menjelaskan dua hal: kalibrasi saja tidak cukup, sebab skor konstan $\pi$
terkalibrasi sempurna tetapi tidak membedakan contoh; dan untuk kejadian langka garis
dasar Brier sangat rendah, sehingga model yang tampak baik pada akurasi dapat kalah telak
pada Brier. Pada simulasi kita $\pi(1-\pi)=0{,}000999$, sedangkan label keras model memberi
Brier sama dengan laju galatnya, $1051/100\,000=0{,}010510$. Skor keterampilan Brier,
$1-B/B_{\mathrm{garis\ dasar}}$, bernilai $-9{,}52$: model itu sepuluh kali lebih buruk
daripada meramalkan kelaziman konstan.

Angka negatif itu bukan keanehan: label keras membuang informasi ketidakpastian, dan pada
kejadian langka informasi itulah yang paling berharga. Model yang mengeluarkan $0{,}06$
sedang mengatakan sesuatu yang benar, dan laporan yang hanya memuat akurasi membuang
kalimat itu.

\section{Ketidakpastian metrik}
\tingkat{lanjut}

Semua angka di atas dihitung dari sejumlah contoh yang terbatas, jadi semuanya bergoyang
bila datanya diganti. Kalimat sehari-harinya: sebuah metrik tanpa selang adalah satu
tembakan tanpa keterangan seberapa lebar sasarannya.

Bab 20 sudah menurunkan interval Wilson dan membuktikan sifatnya, jadi kita tidak
mengulangnya di sini. Yang perlu diingat adalah bentuknya untuk proporsi $k/n$:
\[
\frac{\hat p+\dfrac{z^2}{2n}\pm
z\sqrt{\dfrac{\hat p(1-\hat p)}{n}+\dfrac{z^2}{4n^2}}}{1+\dfrac{z^2}{n}},
\qquad z=1{,}96 .
\]

\begin{untukanak}
Dua penjaga gawang berlatih menepis penalti. Yang pertama menepis sembilan dari sepuluh
tendangan. Yang kedua menepis sembilan ratus dari seribu tendangan. Koran menulis judul
yang sama untuk keduanya: sembilan puluh persen.

Tetapi kita merasa kedua kalimat itu tidak sama kuatnya, dan perasaan itu benar. Penjaga
pertama baru menghadapi sepuluh tendangan. Bisa jadi ia memang hebat, bisa jadi ia hanya
beruntung sepuluh kali berturut-turut. Penjaga kedua sudah menghadapi seribu tendangan,
dan keberuntungan sepanjang itu hampir mustahil. Angka yang sama, kekuatan yang berbeda.

Pertanyaan yang membuat penasaran: kalau kita hanya boleh menulis satu kalimat tentang
penjaga pertama, apakah ``sembilan puluh persen'' masih boleh ditulis? Boleh, asalkan
kalimatnya dilengkapi. Di bab ini kita akan tahu mengapa angka yang sama berarti ``antara
lima puluh sembilan dan sembilan puluh delapan persen'' pada sepuluh tendangan, tetapi
hanya ``antara delapan puluh delapan dan sembilan puluh dua persen'' pada seribu tendangan,
dan mengapa memperbanyak tendangan seratus kali lipat hanya mempersempit jawaban menjadi
sepersepuluh.
\end{untukanak}

Tabel berikut memakai rumus itu pada akurasi, dengan $z=1{,}96$.

\begin{center}
\small
\begin{tabular}{lrrr}
\hline
data & $\hat p$ & selang Wilson & lebar \\
\hline
$9/10$ & 0,900 & $(0{,}5959;\ 0{,}9821)$ & 0,3863 \\
$900/1000$ & 0,900 & $(0{,}8798;\ 0{,}9171)$ & 0,0372 \\
$10/10$ & 1,000 & $(0{,}7225;\ 1{,}0000)$ & 0,2775 \\
$0/1000$ & 0,000 & $(0{,}0000;\ 0{,}0038)$ & 0,0038 \\
\hline
\end{tabular}
\end{center}

Baris pertama dan kedua menjawab dua penjaga gawang tadi. Angka titiknya
sama, $0{,}9$, tetapi selang pertama selebar $0{,}3863$ dan yang kedua $0{,}0372$:
sepuluh kali lebih sempit, bukan seratus kali.

Baris ketiga adalah kasus yang paling sering muncul: tidak ada satu pun kesalahan
dalam $10$ percobaan, lalu ditulis ``akurasi $100$ persen''. Selang Wilson yang
jujur, $(0{,}7225;1)$, selebar $0{,}2775$: laju galatnya masih boleh $27{,}75$
persen. Di ujung lain, nol kejadian dalam $1000$ percobaan memberi batas atas
$0{,}003827$, sedangkan aturan tiga dari Bab 20 memberi $0{,}003$; keduanya boleh,
asalkan prosedurnya disebut.

Membandingkan dua model dengan memeriksa apakah selang masing-masing bertumpuk adalah
kebiasaan yang salah: keduanya diuji pada data yang sama, dan ketergantungan itu justru
menajamkan perbandingan.

\begin{proposisi}[Selisih dua model berpasangan]
Misalkan $b$ adalah banyak contoh yang benar oleh model A dan salah oleh model B, dan $c$
sebaliknya, pada $n$ contoh yang sama. Selisih akurasinya $d=(b-c)/n$ dengan variansi
\[
\operatorname{Var}(d)=\frac{1}{n}\Bigl(p_b+p_c-(p_b-p_c)^2\Bigr),
\]
sehingga galat baku tertaksirnya $\sqrt{b+c-(b-c)^2/n}\,/\,n$.
\end{proposisi}

\begin{proof}[Bukti]
Setiap contoh masuk salah satu dari empat sel: benar keduanya, salah keduanya, benar hanya
A, benar hanya B. Karena contohnya bebas, $(b,c)$ menyebar multinomial dengan peluang
$p_b,p_c$. Maka $\operatorname{Var}(b-c)=n\bigl(p_b+p_c-(p_b-p_c)^2\bigr)$, dan membagi
dengan $n^2$ memberi rumus di atas.
\end{proof}

Pada $b=37$, $c=15$, dan $n=1000$, selisihnya $0{,}022$ dengan galat baku $0{,}007177$,
sehingga selang $95$ persennya $(0{,}007932;0{,}036068)$, seluruhnya di atas nol, dan uji
McNemar eksak memberi nilai-p $0{,}003185$. Perlakuan naif yang menganggap kedua model
bebas memberi selang $(-0{,}004296;0{,}048296)$ yang memuat nol. Kedua model karena itu
dapat dinyatakan berbeda, asalkan perbandingannya berpasangan dan pengujiannya memakai
selisih, bukan dua selang terpisah.

Terakhir, ukuran uji menentukan selisih sekecil apa yang dapat dilihat. Untuk mendeteksi
selisih $\delta$ pada taraf $\alpha$ dan daya $1-\beta$,
\[
n\approx\frac{(z_{1-\alpha/2}+z_{1-\beta})^2\,\bigl(p_1(1-p_1)+p_0(1-p_0)\bigr)}
{\delta^2}.
\]
Pada $p_0=p_1=0{,}9$, taraf $0{,}05$, dan daya $0{,}8$, selisih satu poin menuntut
$n\approx14\,128$ contoh per model, selisih dua poin menuntut $3532$. Kumpulan uji seribu
contoh tidak dapat membedakan dua model yang terpaut satu poin.

\section{Kebocoran data dan validasi silang bersarang}
\tingkat{lanjut}

Semua perhitungan di atas mengandaikan satu hal yang tidak boleh dilanggar: angka uji
diukur pada data yang tidak pernah dilihat model, dan tidak pernah dilihat pula oleh
prosedur yang memilih model. Kalimat sehari-harinya: nilai ujian hanya berarti bila
soalnya belum pernah dibocorkan, termasuk oleh guru yang menyusun soalnya.

\kamus{kebocoran data}{data leakage}{jawaban ujian yang tanpa sengaja ikut masuk ke
bahan belajar, sehingga nilai menjadi tinggi tanpa kepandaian}

\begin{definisi}[Kebocoran data]
\index{kebocoran data}
Kebocoran terjadi bila informasi yang tidak akan tersedia pada saat model dipakai masuk
ke dalam pelatihan, pemilihan model, atau pelaporan, sehingga angka evaluasi mengukur
kemampuan mengingat, bukan kemampuan menggeneralisasi.
\end{definisi}

Bentuknya bermacam-macam, dan yang paling sering tidak terlihat. Menghitung rataan
atau median prapemrosesan dari seluruh data termasuk data uji adalah kebocoran:
statistik itu seharusnya dari latih saja. Memilih $50$ fitur terbaik dengan korelasi
terhadap label sebelum memisahkan latih dan uji juga kebocoran, sebab penyaring
fitur sudah melihat label uji. Menyetel hiperparameter pada kumpulan uji lalu
melaporkan angka terbaiknya adalah kebocoran yang paling sering dilakukan dengan
sadar. Pada data berkelompok, pemisahan acak menempatkan kelompok yang sama di kedua
sisi; pada deret waktu, ia membocorkan masa depan ke masa lalu. Yang paling sulit
dideteksi: fitur yang merupakan akibat dari label, misalnya kolom ``tindakan yang
diambil'' yang diisi setelah insiden diketahui.

Kebocoran pemilihan fitur dapat diukur, dan ukurannya mengejutkan. Ambil $200$ contoh
berlabel acak dan $2000$ fitur derau normal, tanpa sinyal apa pun. Pilih $20$ fitur dengan
korelasi mutlak terbesar terhadap label, lalu validasi silang lima lipatan dengan
pengklasifikasi centroid terdekat. Bila pemilihan
dilakukan sekali pada seluruh data sebelum lipatan dibentuk, akurasinya $0{,}7800$; bila
pemilihan diulang di dalam tiap lipatan memakai baris latihnya saja, $0{,}4600$. Pada tiga
benih lain angkanya $0{,}775$ dan $0{,}415$, $0{,}820$ dan $0{,}545$: yang bocor selalu
tinggi, yang jujur selalu di sekitar $0{,}5$. Angka $0{,}78$ itu seluruhnya ilusi, dan
tidak satu pun baris kodenya salah tulis.

Protokolnya satu kalimat: pisahkan lebih dulu, baru hitung apa pun. Bagi data
menjadi latih, validasi, dan uji; hitung statistik prapemrosesan dari latih saja;
setel hiperparameter pada validasi; sentuh uji sekali, di akhir, untuk satu angka.
Bila validasi silang dipakai, prapemrosesan berada di dalam lipatan. Untuk data
berkelompok, lipatannya mengikuti kelompok; untuk data waktu, mengikuti urutan
waktu.

Gambar~\ref{fig:pemisahan-tiga} menggambar urutan itu sebagai tiga kotak dan satu panah
yang dilarang.

\begin{figure}[htbp]
\centering
\begin{tikzpicture}[every node/.style={font=\footnotesize},
  blok/.style={draw=ink,thick,rounded corners=2pt,align=center,
               minimum width=3.5cm,minimum height=1.5cm,inner sep=3pt}]
  \node[blok,fill=layerteal!12]   (latih) at (0.0,0) {\textbf{latih} $60\%$\\ hitung statistik};
  \node[blok,fill=layerorange!12] (val)   at (4.8,0) {\textbf{validasi} $20\%$\\ pilih hiperparameter};
  \node[blok,fill=layerindigo!12] (uji)   at (9.6,0) {\textbf{uji} $20\%$\\ sentuh sekali};
  \draw[-{Latex},ink,thick] (latih.east) -- (val.west)
    node[midway,above,font=\scriptsize] {latih ulang};
  \draw[-{Latex},ink,thick] (val.east) -- (uji.west)
    node[midway,above,font=\scriptsize] {kunci sebelum dilihat};
  \draw[-{Latex},accent,thick,dashed] (uji.south) .. controls (9.6,-1.85) and (0.0,-1.85)
    .. (latih.south);
  \draw[accent,very thick] (4.62,-1.42) -- (4.98,-1.06);
  \draw[accent,very thick] (4.62,-1.06) -- (4.98,-1.42);
  \node[accent,anchor=west,font=\scriptsize] at (5.12,-1.24) {dilarang};
  \node[anchor=north,align=center,text width=11.6cm] at (4.8,-1.98)
    {statistik uji, penyaring fitur, dan penyetelan hiperparameter tidak boleh menyeberang
     balik ke kotak latih};
\end{tikzpicture}
\caption{Tiga kumpulan data yang terpisah, dengan proporsi $60/20/20$; pada $1000$ koneksi
itu berarti $600$, $200$, dan $200$. Statistik prapemrosesan dihitung dari kotak latih,
hiperparameter dipilih di kotak validasi, dan kotak uji disentuh sekali di akhir. Panah
putus-putus dari uji kembali ke latih adalah kebocoran: satu panah itu cukup untuk
menaikkan akurasi dari $0{,}4600$ menjadi $0{,}7800$ pada data yang seluruhnya derau.}
\label{fig:pemisahan-tiga}
\end{figure}

Yang tersisa adalah satu masalah halus. Bila validasi silang dipakai sekaligus untuk
memilih hiperparameter dan untuk melaporkan mutu, angka yang dilaporkan adalah angka
terbaik dari banyak percobaan pada data yang sama, jadi ia terlalu optimistis. Jalan
keluarnya adalah dua lapis lipatan.

\kamus{validasi silang bersarang}{nested cross-validation}{dua lapis lipatan, lapisan luar
untuk menilai dan lapisan dalam untuk memilih, seperti dua ruang ujian yang terpisah}

\begin{definisi}[Validasi silang bersarang]
\index{validasi silang}
\index{validasi silang bersarang}
Misalkan $\Lambda$ himpunan hiperparameter. Data dibagi menjadi $K$ lipatan luar. Untuk
setiap lipatan luar $k$: latih pada sisa data dengan setiap $\lambda\in\Lambda$ memakai
$J$ lipatan dalam, pilih $\lambda_k$ yang terbaik, latih ulang pada seluruh bagian latih
lipatan $k$ dengan $\lambda_k$, lalu ukur pada lipatan $k$. Angka yang dilaporkan adalah
rata-rata $K$ angka itu. Penilaian dan pemilihan karena itu memakai data yang berbeda.
\end{definisi}

Biayanya $K$ kali lebih mahal, dan itulah sebabnya protokol ini sering dilewati; tanpa
lapisan luar, tidak ada satu angka pun dalam laporan yang dapat ditafsirkan sebagai
taksiran mutu pada data baru.

\section{Mengapa satu kali uji tidak cukup}
\tingkat{menengah}

Kalimat ``kami menguji model pada $1000$ contoh dan memperoleh akurasi $0{,}900$''
terdengar tuntas, padahal ia memuat tiga ketidakpastian. Kalimat sehari-harinya:
satu kali menimbang tidak memberi tahu apakah timbangannya bergoyang, barangnya
berubah, atau kita memilih hasil yang paling menyenangkan.

Yang pertama statistis dan sudah dihitung: pada $n=1000$ selang $95$ persen selebar
$0{,}037$, dan memisahkan dua model yang terpaut satu poin menuntut $14\,128$ contoh.
Kesimpulan ``A lebih baik'' dari selisih sekecil itu adalah kesimpulan tentang undian.

Yang kedua selektif, dan namanya \emph{winner's curse}. Melatih $20$ varian yang mutunya
identik lalu melaporkan yang terbaik berarti melaporkan maksimum dari $20$ goyangan: pada
simulasi dengan akurasi benar $0{,}90$ dan $n=1000$, rata-rata tiap varian $0{,}899987$
sedangkan rata-rata maksimumnya $0{,}917337$, jadi $1{,}7$ poin dari angka itu
keberuntungan. Bila varian saling bergantung, goyangannya lebih kecil, tetapi arahnya
selalu ke atas.

Yang ketiga prosedural: kumpulan uji yang dipakai berulang berhenti menjadi kumpulan uji.
Bila tiap perbandingan punya peluang $5$ persen menipu, peluang sekurang-kurangnya satu
dari $20$ perbandingan menipu adalah $1-0{,}95^{20}=0{,}6415$, lebih besar daripada peluang
benar.

Obatnya sepadan: laporkan selang, pisahkan data pemilihan dari data pelaporan, dan kunci
kumpulan uji sebelum eksperimen dimulai. Bila uji terlanjur dipakai berulang, angka
berikutnya disebut angka eksplorasi.

\section{Daftar periksa pelaporan}
\tingkat{dasar}

Tujuh baris ini memisahkan laporan yang dapat diperiksa dari yang hanya meyakinkan.
Kalimat sehari-harinya: resi tidak membuat barang sampai, tetapi tanpa resi tak ada yang
dapat ditelusuri.

\begin{enumerate}
  \item Ukuran tiap kelas dan kelazimannya, bukan hanya $n$ total.
  \item Keempat bilangan matriks konfusi, lalu presisi dan recall pada ambang yang
        benar-benar dipakai, beserta nilai ambangnya.
  \item Selang kepercayaan tiap metrik utama, dengan metode yang disebut; untuk $n$ kecil
        pakai Wilson, bukan Wald.
  \item Kurva presisi--recall bila kelas positif langka, beserta average precision dan
        garis dasarnya $\pi$.
  \item Kalibrasi: batas selang skor, banyak contoh per selang, dan kesalahan kalibrasi.
  \item Protokol pemisahan: apa yang dihitung dari data mana, berapa kali data uji
        disentuh, dan berapa varian dicoba; angka terbaik tanpa penyebut itu angka
        pilihan, bukan angka mutu.
  \item Sebaran hasil antar benih atau antar lipatan, beserta kode dan benih acaknya.
\end{enumerate}

Yang tidak perlu ada: akurasi tanpa kelaziman, AUC tanpa kurva, dan kata ``siap
dipakai''.

\begin{resep}
Empat langkah melaporkan hasil evaluasi yang tidak menyesatkan:
\begin{enumerate}
  \item Pisahkan data lebih dulu: latih, validasi, uji, beserta proporsi dan kunci
        pemisahannya, ditulis sebelum satu hitungan pun dijalankan.
  \item Hitung prapemrosesan dan pilih fitur hanya dari latih, setel hiperparameter hanya
        pada validasi, lalu simpan uji untuk satu kali sentuh.
  \item Laporkan keempat sel matriks konfusi beserta $n$, $\pi$, dan ambang yang dipakai,
        lalu presisi, recall, dan selang kepercayaan tiap metrik.
  \item Sebutkan banyak varian yang dicoba dan sebaran hasil antar benih atau antar
        lipatan, supaya pembaca tahu berapa derajat kebebasan yang sudah dipakai.
\end{enumerate}
Peringatan urutan: langkah 1 tidak boleh menyusul langkah 2 atau 3. Statistik yang dihitung
sebelum data dipisahkan sudah memindahkan informasi uji ke latih, dan setelah itu tidak ada
angka yang dapat memperbaikinya; seluruh laporan hanya mengukur ingatan.
\end{resep}

\section{Contoh hitung}
\tingkat{menengah}

Tiga hitungan berikut dapat diulang di kertas.

Pertama, matriks konfusi dari catatan kerja penulis sendiri, pada $35\,069$ koneksi uji:
$\mathrm{TP}=23\,279$, $\mathrm{FP}=997$, $\mathrm{FN}=590$, $\mathrm{TN}=10\,203$.
\begin{align*}
\operatorname{prec}&=\frac{23\,279}{24\,276}=0{,}958931, &
\operatorname{rec}&=\frac{23\,279}{23\,869}=0{,}975282,\\
F_1&=\frac{46\,558}{48\,145}=0{,}967037, &
\mathrm{acc}&=\frac{33\,482}{35\,069}=0{,}954746,
\end{align*}
dan akurasi seimbangnya $0{,}943132$. Satu hal di catatan itu tidak konsisten: satu
kalimat menyebut $11\,200$ serangan dan $23\,869$ normal, sedangkan matriksnya menempatkan
$23\,869$ pada baris positif. Aritmetikanya benar, penamaan kelasnya tertukar, dan
kelaziman harus dibaca dari matriks, bukan dari kalimat.

Kedua, AUC dari data kecil. Positif berskor $0{,}92$, $0{,}81$, $0{,}55$ melawan negatif
berskor $0{,}70$, $0{,}41$, $0{,}35$, $0{,}22$ memberi $12$ pasangan; $0{,}92$ dan $0{,}81$
mengalahkan keempatnya, $0{,}55$ mengalahkan tiga, jadi $11/12=0{,}916667$. Bila satu skor
positif sama dengan satu skor negatif, pasangan itu dihitung setengah dan hasilnya
$0{,}875$; \texttt{sklearn} memberi angka yang sama bila $y$ dan skornya diumpankan
sebagai vektor.

Ketiga, efek kelaziman pada presisi ketika sensitivitas $0{,}90$ dan spesifisitas $0{,}95$
dipertahankan. Dengan $\operatorname{prec}=\pi(0{,}9)/\bigl(\pi(0{,}9)+(1-\pi)(0{,}05)\bigr)$:

\begin{center}
\small
\begin{tabular}{rrrr}
\hline
$\pi$ & $\operatorname{prec}$ & alarm benar & alarm palsu \\
\hline
0,001 & 0,017699 & 1 & 56 \\
0,010 & 0,153846 & 1 & 5,5 \\
0,100 & 0,666667 & 2 & 1 \\
0,500 & 0,947368 & 18 & 1 \\
\hline
\end{tabular}
\end{center}

Kolom terakhir adalah kebalikan presisi dikurangi satu, yaitu banyak alarm palsu untuk
setiap alarm benar. Pada kelaziman satu banding seribu, satu alarm benar datang bersama
$55{,}5$ alarm palsu; pada kelaziman setengah, bersama $0{,}056$. Model yang sama, mutu
yang sama, dan pengalaman operator yang sama sekali berbeda.

\kertas{Ambil matriks $1000$ prediksi tadi dan hitung ulang kelima metriknya tanpa melihat
teks. Lalu jawab: berapa $\mathrm{FP}$ yang harus dihapus, tanpa mengubah $\mathrm{TP}$
dan $\mathrm{FN}$, agar presisi menjadi $0{,}900$? Petunjuk: presisi $62/(62+x)=0{,}9$
memberi $x=6{,}89$, jadi $\mathrm{FP}$ harus turun dari $17$ menjadi paling banyak $6$.
Periksa bahwa akurasi sesudah perubahan itu $0{,}956$.}

\section{Lab komputasi: satu angka, empat cara membacanya}
\tingkat{lanjut}

Seluruh angka di bab ini dihitung oleh satu program pendek, yang juga memperlihatkan dua
hal yang tak dapat diperlihatkan rumus: seberapa besar kebocoran menaikkan angka, dan
seberapa besar pemenang dari banyak varian dilebihkan. Kodenya dijalankan apa adanya
dengan \texttt{numpy} dan \texttt{scipy}.

\begin{lab}
\begin{verbatim}
import numpy as np
from math import sqrt
from scipy.stats import norm

z = norm.ppf(0.975)


def wilson(k, n):
    p, z2 = k / n, z * z
    mid = (p + z2 / (2 * n)) / (1 + z2 / n)
    half = z * sqrt(p * (1 - p) / n + z2 / (4 * n * n)) / (1 + z2 / n)
    return mid - half, mid + half


def roc_pr(spos, sneg):
    s = np.concatenate([spos, sneg])
    y = np.concatenate([np.ones(len(spos)), np.zeros(len(sneg))])
    o = np.argsort(-s)
    s, y = s[o], y[o]
    tp, fp = np.cumsum(y), np.cumsum(1 - y)
    tpr, fpr = tp / len(spos), fp / len(sneg)
    prec = tp / (tp + fp)
    auc = np.sum(np.diff(np.concatenate([[0], fpr])) *
                 (tpr + np.concatenate([[0], tpr[:-1]])) / 2)
    ap = np.sum(np.diff(np.concatenate([[0], tpr])) * prec)
    return auc, ap


rng = np.random.default_rng(20261025)
Nneg, Npos = 99900, 100
sneg = rng.normal(0.0, 1.0, Nneg)
spos = rng.normal(2.6, 1.0, Npos)
t = 2.3263
tp = int((spos > t).sum()); fp = int((sneg > t).sum())
fn, tn = Npos - tp, Nneg - fp
n = Nneg + Npos
print("TP=%d FP=%d FN=%d TN=%d" % (tp, fp, fn, tn))
print("akurasi %.6f  garis dasar %.6f" % ((tp + tn) / n, Nneg / n))
print("presisi %.6f  recall %.6f" % (tp / (tp + fp), tp / Npos))
print("akurasi seimbang %.6f" % ((tp / Npos + tn / Nneg) / 2))
print("brier label keras %.6f  brier konstan %.6f"
      % ((fp + fn) / n, 0.001 * 0.999))
a1, p1 = roc_pr(spos, sneg)
print("pi=0,001: AUC=%.6f  AP=%.6f" % (a1, p1))
a2, p2 = roc_pr(rng.normal(2.6, 1.0, 50000), rng.normal(0.0, 1.0, 50000))
print("pi=0,500: AUC=%.6f  AP=%.6f" % (a2, p2))
for k, m in ((9, 10), (900, 1000), (10, 10), (0, 1000)):
    lo, hi = wilson(k, m)
    print("%4d/%-4d (%.6f, %.6f) lebar %.6f" % (k, m, lo, hi, hi - lo))
b, c = 37, 15
d = (b - c) / 1000
se = sqrt(b + c - (b - c) ** 2 / 1000) / 1000
print("selisih %.6f  se %.6f  IK (%.6f, %.6f)"
      % (d, se, d - z * se, d + z * se))
rt = np.random.default_rng(11)
a = rt.binomial(1000, 0.90, (20000, 20)) / 1000
print("20 varian: rata2 %.6f  rata2 maks %.6f"
      % (a.mean(), a.max(axis=1).mean()))

rl = np.random.default_rng(4)
nn, dd, kk = 200, 2000, 20
X = rl.normal(size=(nn, dd))
y = rl.permutation(np.repeat([0, 1], nn // 2))
corr = np.array([abs(np.corrcoef(X[:, j], y)[0, 1]) for j in range(dd)])
sel_all = np.argsort(-corr)[:kk]


def cv(pick):
    idx = rl.permutation(nn)
    ok = 0
    for f in np.array_split(idx, 5):
        tr = np.setdiff1d(idx, f)
        cols = pick(tr)
        m1 = X[np.ix_(tr[y[tr] == 1], cols)].mean(0)
        m0 = X[np.ix_(tr[y[tr] == 0], cols)].mean(0)
        d1 = ((X[np.ix_(f, cols)] - m1) ** 2).sum(1)
        d0 = ((X[np.ix_(f, cols)] - m0) ** 2).sum(1)
        ok += (np.where(d1 < d0, 1, 0) == y[f]).sum()
    return ok / nn


def top_in(tr):
    c = np.array([abs(np.corrcoef(X[tr, j], y[tr])[0, 1]) for j in range(dd)])
    return np.argsort(-c)[:kk]


print("pemilihan fitur: bocor %.4f  jujur %.4f"
      % (cv(lambda tr: sel_all), cv(top_in)))
\end{verbatim}
\end{lab}

Keluaran yang sebenarnya, apa adanya:

\begin{verbatim}
TP=62 FP=1013 FN=38 TN=98887
akurasi 0.989490  garis dasar 0.999000
presisi 0.057674  recall 0.620000
akurasi seimbang 0.804930
brier label keras 0.010510  brier konstan 0.000999
pi=0,001: AUC=0.966448  AP=0.287636
pi=0,500: AUC=0.967061  AP=0.967544
   9/10   (0.595850, 0.982124) lebar 0.386274
 900/1000 (0.879848, 0.917091) lebar 0.037243
  10/10   (0.722467, 1.000000) lebar 0.277533
   0/1000 (0.000000, 0.003827) lebar 0.003827
selisih 0.022000  se 0.007177  IK (0.007932, 0.036068)
20 varian: rata2 0.899987  rata2 maks 0.917330
pemilihan fitur: bocor 0.7800  jujur 0.4600
\end{verbatim}

Empat hal perlu dibaca dari keluaran itu. Pertama, dua angka pada baris kedua harus muncul
berdekatan: $0{,}989490$ dan $0{,}999000$. Arah selisihnya menentukan seluruh tafsiran:
model yang diuji berada di bawah garis dasar.

Kedua, AUC hampir tidak bergerak, $0{,}966448$ melawan $0{,}967061$, sedangkan AP bergerak
dari $0{,}287636$ ke $0{,}967544$, faktor $3{,}4$: bukti numerik bagi proposisi sebelumnya.

Ketiga, lebar selang menyusut sebanding $1/\sqrt n$, dari $0{,}386274$ pada $n=10$ menjadi
$0{,}037243$ pada $n=1000$, dan $0/1000$ memberi batas atas $0{,}003827$.

Keempat, dua baris terakhir memuat dua penyakit yang tak terlihat pada satu angka: selisih
berpasangan $0{,}022$ dengan selang $(0{,}007932;0{,}036068)$ tidak memuat nol, padahal
$(0{,}8798;0{,}9171)$ dan $(0{,}8563;0{,}8969)$ bertumpuk lebar; dan pemenang dari $20$
varian identik dilaporkan $0{,}917337$ padahal mutu sebenarnya $0{,}899987$.

Program tidak membuktikan satu pun proposisi di bab ini; ia memeriksa klaim yang sudah
dibuktikan. Simulasi memakai sebaran yang kita andaikan, dan validasi silang yang jujur
pada data derau memberi $0{,}4600$, bukan $0{,}5000$: angka tanpa sinyal pun bergoyang.

\begin{untukdosen}
Urutan mengajar yang saya sarankan: mulai dari laporan vendor pada pembuka bab, tulis
$99$ persen dan $99{,}9$ persen di papan, lalu minta kelas memutuskan mana yang lebih
baik. Hampir seluruh kelas memilih detektor, dan dari kegelisahan itu matriks konfusi
muncul sebagai jawaban, bukan definisi yang dihafal. Setelah presisi dan recall ada,
barulah $F_1$ dan akurasi seimbang masuk. Kurva ROC sebaiknya digambar tangan untuk enam
titik pada tabel ambang, sebab di situ mahasiswa melihat sendiri bahwa titik yang tampak
bagus pada ROC dapat berpresisi $0{,}013$.

Dua miskonsepsi khas. Pertama, AUC dianggap ukuran mutu yang bebas konteks sehingga ambang
dipilih sembarang; gejalanya tampak ketika mahasiswa tidak dapat menjelaskan presisi pada
ambang yang dipakai. Kedua, validasi silang dianggap selalu jujur, padahal pemilihan fitur
atau standardisasi di luar lipatan sudah cukup menaikkan angka $30$ poin pada data yang
seluruhnya derau.

Aktivitas papan tulis: bagi kelas menjadi dua kelompok, satu menghitung presisi untuk
$\pi=0{,}001$ dan satu untuk $\pi=0{,}5$ dengan sensitivitas dan spesifisitas yang sama,
lalu suruh keduanya mempertahankan hasilnya. Kelompok pertama menemukan $0{,}018$ dan harus
menjelaskan mengapa model yang sama terasa berbeda. Tanda kelas siap lanjut: mereka dapat
menyebut satu keadaan di mana akurasi $0{,}99$ lebih buruk daripada tidak memakai model.
\end{untukdosen}

\section{Ringkasan}
\tingkat{dasar}

\begin{ringkasan}
\begin{itemize}
  \item Empat bilangan matriks konfusi berperan tidak simetris: presisi menjawab
        pertanyaan operator, recall pertanyaan perancang kebijakan, dan keduanya
        bergantung pada ambang; $F_1$ adalah rata-rata harmoniknya dan mengabaikan
        $\mathrm{TN}$.
  \item Akurasi garis dasar adalah $1-\pi$: pada $\pi=0{,}001$ model dengan $0{,}989490$
        kalah dari program yang selalu menjawab negatif ($0{,}999000$), sedangkan akurasi
        seimbang $0{,}804930$ lebih jujur sebab tidak bergantung pada komposisi.
  \item AUC sama dengan peluang skor positif acak melebihi skor negatif acak, jadi hanya
        bergantung pada sebaran skor bersyarat; ia tidak berubah ketika kelaziman diubah,
        sedangkan presisi dan average precision berubah, dari $0{,}288$ ke $0{,}968$.
  \item Kalibrasi memeriksa apakah skor $0{,}8$ berarti delapan dari sepuluh. Untuk skor
        terkalibrasi, Brier sama dengan $\mathbb{E}[S(1-S)]$; skor konstan $\pi$
        terkalibrasi sempurna tetapi tidak tajam. Label keras di sini memberi
        $0{,}010510$ melawan $0{,}000999$.
  \item Lebar selang menyusut sebanding $1/\sqrt n$: $0{,}3863$ pada $n=10$ dan $0{,}0372$
        pada $n=1000$. Untuk membandingkan dua model, pakai selisih berpasangan: $b=37$,
        $c=15$ memberi $(0{,}007932;0{,}036068)$, sedangkan perlakuan tak berpasangan
        memberi selang yang memuat nol.
  \item Satu kali uji tidak cukup, dan kebocoran menaikkan angka secara palsu: $0{,}7800$
        bila fitur dipilih di luar lipatan, $0{,}4600$ bila di dalamnya; pemenang dari
        $20$ varian identik dilebihkan $1{,}7$ poin; dan $20$ perbandingan pada taraf $5$
        persen memberi peluang $0{,}6415$ bahwa salah satunya menipu.
\end{itemize}
\end{ringkasan}

\section{Latihan}
\tingkat{menengah}

\begin{latihan}
\begin{enumerate}
  \item Nilai klaim ini: ``Model A mencapai akurasi $99{,}1$ persen dan model B $98{,}7$
        persen pada kumpulan uji yang sama, jadi A lebih baik.'' Benar, salah, atau belum
        cukup bukti? Sebutkan bagian yang keliru dan dua angka
        tambahan yang memperkuat.
  \item Buktikan bahwa Brier untuk label keras $0/1$ sama dengan laju galat,
        $(\mathrm{FP}+\mathrm{FN})/n$. Mengapa kesamaan itu gagal untuk skor berupa
        peluang?
  \item Buktikan bahwa skor konstan $S=s$ meminimumkan Brier tepat ketika $s=\pi$, dengan
        nilai minimum $\pi(1-\pi)$.
  \item Diberikan $\mathrm{TP}=180$, $\mathrm{FP}=45$, $\mathrm{FN}=20$,
        $\mathrm{TN}=1755$. Hitung presisi, recall, $F_1$, akurasi, spesifisitas, dan
        akurasi seimbang sampai enam angka desimal, lalu periksa dengan
        \texttt{sklearn.metrics}.
  \item Hitung selang Wilson $95$ persen untuk $95/100$ dan $950/1000$. Bandingkan kedua
        lebar, lalu jelaskan mengapa perbandingan itu tidak boleh dipakai menyimpulkan
        model mana yang lebih baik.
  \item Hitung AUC dari lima positif berskor $0{,}95$, $0{,}88$, $0{,}70$, $0{,}60$,
        $0{,}40$ dan lima negatif berskor $0{,}75$, $0{,}65$, $0{,}50$, $0{,}30$,
        $0{,}10$. Tuliskan banyak pasangan benar, lalu periksa dengan
        \texttt{roc\_auc\_score}.
  \item Sebuah model memiliki sensitivitas $0{,}80$ dan spesifisitas $0{,}99$. Hitung
        presisinya pada $\pi=0{,}002$ dan $\pi=0{,}2$, lalu tentukan kelaziman minimum
        agar presisinya paling sedikit $0{,}5$.
  \item Anda memiliki $6000$ rekaman dari $300$ host, masing-masing sekitar $20$ rekaman,
        dan $40$ setelan hiperparameter. Rancang protokol evaluasi yang tidak bocor:
        sebutkan pembagian data, tempat prapemrosesan dihitung, tempat pemilihan
        hiperparameter, dan berapa kali uji disentuh.
  \item Nilai klaim ini: ``Karena AUC model ini $0{,}95$, model ini berguna pada ambang
        mana pun yang kami pilih.'' Benar, salah, atau belum cukup bukti? Sertakan satu
        perhitungan pendukung.
\end{enumerate}

Petunjuk untuk tiga butir tersulit. Butir 2: tulis $(s-y)^2$ untuk $s\in\{0,1\}$; nilainya
$1$ tepat pada kesalahan. Butir 3: turunkan $\mathbb{E}[(s-y)^2]=s^2-2s\pi+\pi$ terhadap
$s$, cari titik minimumnya, lalu pakai proposisi Brier di bab ini. Butir 7: presisinya
$0{,}8\pi/\bigl(0{,}8\pi+0{,}01(1-\pi)\bigr)$; selesaikan
$0{,}8\pi\ge0{,}5\bigl(0{,}8\pi+0{,}01(1-\pi)\bigr)$ terhadap $\pi$.
\end{latihan}

\begin{untukpeneliti}
Yang terbukti di bab ini: akurasi pengklasifikasi konstan adalah $1-\pi$; $F_1$ adalah
rata-rata harmonik presisi dan recall dengan $F_1\le(P+R)/2$; AUC sama dengan peluang
peringkat benar beserta bentuk Mann--Whitney-nya; ROC dan AUC invarian terhadap kelaziman
sedangkan presisi tidak; untuk skor terkalibrasi Brier sama dengan $\mathbb{E}[S(1-S)]$;
dan rumus variansi selisih dua proporsi berpasangan.

Yang populer tetapi bukan teorema: bahwa ``AUC $0{,}9$ berarti model bagus''; bahwa ambang
$0{,}5$ netral; bahwa validasi silang selalu jujur; dan bahwa selang yang bertumpuk berarti
tidak ada perbedaan. Tiga yang terakhir salah secara matematis, bukan hanya tidak lengkap.
\emph{Average precision} juga bukan taksiran tak bias bagi luas kurva presisi--recall
populasi: ia ringkasan dari kurva teramati, dan nilainya bergantung pada kelaziman sampel.

Sumber primer: Wasserman, \emph{All of Statistics}, Bab 6, Bab 9--10, Bab 13 bagian
Model Selection, dan Bab 22 bagian ``Assessing Error Rates and Choosing a Good Classifier''.
Goodfellow, Bengio, dan Courville, \emph{Deep Learning}, Bab 11 bagian ``Performance
Metrics'', memberi contoh penyakit satu per sejuta dan gagasan \emph{coverage}. Athalye,
Carlini, dan Wagner (2018) menunjukkan bahwa evaluasi yang lemah menghasilkan klaim
keamanan yang salah: tujuh dari sembilan pertahanan tanpa sertifikat pada ICLR 2018
bergantung pada gradien yang dikaburkan, dan enam tembus setelah evaluasinya diperbaiki.

Batas keberlakuan: seluruh simulasi memakai sebaran normal yang diandaikan, jadi ia
mengukur perilaku prosedur, bukan kebenaran model. Selang Wilson adalah aproksimasi
skor, dan McNemar mengandaikan pasangan tak sepakat simetris di bawah hipotesis nol.
Hasil kebocoran bergantung pada $n$, $d$, dan banyaknya fitur terpilih: $0{,}78$
bukan tetapan alam, melainkan satu titik pada kurva yang naik bersama $d/n$.

Pertanyaan terbuka: selang jujur untuk metrik dari lipatan yang saling bergantung;
pengujian kalibrasi tanpa asumsi selang skor seragam; dan pelaporan ketidakpastian ketika
kumpulan uji dipilih secara adaptif.
\end{untukpeneliti}

\peta{Bab 27 pada Wasserman, \emph{All of Statistics}: Bab 6 (penaksir, bias, galat baku),
Bab 9--10 (interval dan pengujian), Bab 13 bagian Model Selection, dan Bab 22 bagian
``Assessing Error Rates and Choosing a Good Classifier''. Interval Wilson dan aturan tiga
dirujuk dari Bab 20 buku ini serta dari Wilson (1927) dan Jovanovic--Levy (1997). Contoh
akurasi yang menipu mengikuti Goodfellow, Bengio, dan Courville, \emph{Deep Learning},
Bab 11; batas evaluasi pertahanan adversarial mengikuti Athalye, Carlini, dan Wagner (2018),
dibahas di Bab 27 dan 28. Aturan Bayes untuk presisi dibuktikan di Bab 16; pemisahan data
di Bab 24.}

\chapter{Metode Formal dengan Lean 4}

Sebuah berkas bernama \texttt{Sertifikat.lean} tiba bersama satu kalimat pengantar: ``semuanya
sudah dibuktikan mesin''. Isinya empat puluh teorema satu baris, hampir semuanya berakhir dengan
\texttt{by native\_decide}, dan di bagian bawah ada satu konstanta ringkasan yang menyebut ambang
$10^{-3}$. Tidak satu pun barisnya dapat kita periksa dengan mata dalam waktu yang wajar: yang
dibandingkan adalah bilangan bulat enam belas digit hasil pengubahan pecahan biner menjadi bentuk
eksak. Kalimat pengantar itu benar dalam arti yang sempit dan menyesatkan dalam arti yang penting.
Mesin memang memeriksa setiap baris, tetapi ia tidak memeriksa apakah yang diperiksa itu yang
dimaksud.

Pengujian dan pembuktian menjanjikan hal yang berbeda, dan bedanya bukan soal derajat
kepercayaan. Pengujian mengambil sejumlah masukan, menjalankan program, lalu melaporkan apa yang
terjadi pada masukan itu. Kesimpulannya selalu berbentuk ``untuk $10^{4}$ sampel uji, tidak ada
yang gagal''. Pembuktian formal berjanji lebih: satu pernyataan berkuantor atas seluruh masukan,
disertai objek yang dapat diperiksa pihak ketiga tanpa mempercayai penulisnya. Janji itu
ditepati hanya bila tiga syarat dipenuhi sekaligus, yaitu pernyataannya ditulis dengan benar,
objek yang diperiksa memang bukti pernyataan itu, dan landasan pemeriksaannya tidak disalahgunakan.
Bab ini membongkar ketiga syarat tersebut pada satu sistem konkret, Lean 4.

Alat dari bab sebelumnya belum cukup untuk pekerjaan ini. Bab 27 mengajarkan cara melaporkan
apa yang benar-benar diukur: pisahkan data latih dari data uji, sebutkan distribusinya, laporkan
variansi antar seed. Itu adalah disiplin tentang \emph{klaim empiris}. Di sini kita menghadapi
klaim yang bentuknya lain: bukan rata-rata atas sampel, melainkan pernyataan berkuantor
$\forall x,\ P(x)$. Untuk klaim semacam itu tidak ada interval kepercayaan yang menolong, karena
satu penyangkal sudah cukup untuk meruntuhkannya, dan ketiadaan penyangkal bukan bukti. Yang
dibutuhkan adalah objek yang bisa diperiksa ulang oleh orang yang tidak percaya kepada kita.

\begin{untukanak}
Bayangkan seorang juri catur yang sangat teliti. Ia tidak pernah bertanya apakah permainan itu
menarik, apakah kedua pemain bermain dengan baik, atau apakah aturan caturnya memang aturan yang
seharusnya dipakai. Tugasnya hanya satu: setiap kali ada langkah, ia memeriksa apakah langkah itu
sah menurut buku aturan. Kalau semua langkah sah, ia menandatangani papan itu dan berkata
``permainan ini sah''. Ia tidak peduli siapa yang menang.

Sekarang bayangkan seseorang menyerahkan papan yang penuh langkah sah, lalu berkata ``ini
membuktikan bahwa saya bisa mengalahkan siapa pun''. Juri tadi akan menjawab dengan jujur: yang
saya periksa hanya bahwa langkahnya sah, bukan bahwa kalimatmu benar. Kalau kalimatnya keliru,
papan yang sah tidak menolong sama sekali, karena yang salah bukan langkahnya melainkan
pertanyaannya.

Pernahkah kamu memberikan jawaban yang benar untuk soal yang salah? Atau mengerjakan pekerjaan
rumah dengan rapi padahal yang diminta halaman lain? Di bab ini kita akan tahu mengapa komputer
bisa memeriksa pembuktian sampai ke langkah terkecil, mengapa pemeriksaan itu sama sekali tidak
menjamin bahwa pertanyaannya tepat, dan bagaimana orang menulis pertanyaan yang tepat dalam
bahasa yang bisa diperiksa mesin.
\end{untukanak}

\section{Bukti informal dan bukti formal}

\tingkat{menengah}

Bukti dalam buku matematika adalah teks yang meyakinkan pembaca terlatih. Ia memuat lompatan yang
disepakati bersama: ``jelas'', ``tanpa mengurangi keumuman'', ``analog untuk kasus lain''. Lompatan
itu bukan kecurangan; ia adalah pemadatan yang diterima komunitas. Masalahnya, pemadatan semacam
itu tidak bisa diperiksa mesin, dan tidak bisa dipisahkan dari pengetahuan pembacanya.

Bukti formal adalah versi lain dari objek yang sama. Pernyataannya ditulis dalam bahasa dengan
tata bahasa dan tipe yang tegas, dan buktinya adalah sebuah \emph{term}, yaitu ekspresi yang
tipenya adalah pernyataan itu. Pemeriksaan dilakukan oleh inti sistem, sepotong kecil kode yang
hanya tahu beberapa aturan dan tidak tahu apa pun tentang matematika. Karena intinya kecil,
pemeriksaannya dapat dipercaya jauh lebih kuat daripada program besar mana pun di sekitar
sistem itu.

\kamus{verifikasi formal}{formal verification}{Pemeriksaan langkah demi langkah oleh juru tulis
yang tidak pernah bosan dan tidak pernah percaya pada penulisnya.}

\begin{definisi}[Bukti formal]
Misalkan $P$ adalah pernyataan yang ditulis dalam bahasa formal sistem. \emph{Bukti formal} $P$
adalah term $p$ sedemikian sehingga $p:P$ dapat diturunkan oleh inti sistem dari definisi,
teorema yang sudah diperiksa, dan aksioma yang dinyatakan terbuka. Sistem menerima $P$ bila dan
hanya bila term semacam itu ditemukan dan lolos pemeriksaan tipe.\index{verifikasi formal}
\end{definisi}

Definisi itu memuat tiga kata yang mudah dilewati. ``Dapat diturunkan'' berarti ada prosedur
mekanis, bukan ada orang yang setuju. ``Diperiksa'' berarti hasilnya biner: diterima atau
ditolak, tanpa derajat keyakinan. ``Dinyatakan terbuka'' berarti aksioma tercatat, bukan
tersembunyi; kita akan memakai pencatatan itu nanti sebagai alat utama.

Perbedaan yang paling sering disalahpahami adalah perbedaan antara \emph{kebenaran} dan
\emph{keterbuktian}. Lean tidak tahu apakah suatu pernyataan benar dalam dunia nyata. Ia hanya
tahu apakah ada term yang bertipe pernyataan itu. Bila pernyataannya salah, term itu tidak akan
pernah ditemukan, tetapi ketiadaan term bukan bukti bahwa pernyataan itu salah; bisa jadi kita
saja yang belum menemukannya.

\begin{bacakeras}
Suara pertama: ``Sudah saya periksa, semuanya benar.''
Suara kedua: ``Tunjukkan caranya.''
Suara pertama: ``Saya membaca ulang tiga kali.''
Suara kedua: ``Itu kebiasaan, bukan cara.''
Suara pertama: ``Saya yakin.''
Suara kedua: ``Keyakinan bukan bukti. Berikan langkahnya, satu per satu.''

Begitulah percakapan dengan mesin. Ia tidak bertanya apakah kita orang baik, apakah kita sudah
sarapan, atau apakah kita lulus dengan nilai tinggi. Ia bertanya satu hal, dengan sabar, setiap
kali, tanpa nada marah: tunjukkan langkahnya. Kalau langkahnya ada, pintu terbuka. Kalau langkahnya
tidak ada, pintu tetap tertutup, serapat sebelumnya, meskipun kita mengulang permintaan itu seratus
kali. Mesin tidak menilai niat kita; ia hanya memeriksa bentuk.

Ucapkan bersama, perlahan, lalu ulangi sekali lagi: mesin tidak percaya, ia memeriksa.
\end{bacakeras}

\section{Proposisi sebagai tipe, bukti sebagai term}

\tingkat{menengah}

Di Lean tidak ada jurang antara pernyataan dan buktinya. Keduanya term, diperiksa oleh aturan yang
sama. Setiap kalimat yang dapat bernilai benar atau salah diberi tipe $\mathtt{Prop}$; bukti
kalimat itu adalah term yang bertipe kalimat tersebut. Pernyataan ``$P$ terbukti'' karena itu
tidak berarti $P$ benar dalam suatu model, melainkan bahwa ada term bertipe $P$ yang lolos
pemeriksaan.

\kamus{proposisi}{proposition}{Kalimat yang bisa dinilai benar atau salah, seperti soal ujian yang
jawabannya harus dibuktikan, bukan sekadar dipilih.}

\begin{definisi}[Penghuni tipe]
Sebuah tipe $T$ disebut \emph{berpenghuni} bila terdapat term $t$ dengan $t:T$; notasi $t:T$
dibaca ``$t$ bertipe $T$''. Bila $P:\mathtt{Prop}$, maka $P$ disebut \emph{proposisi} dan $t:P$
disebut \emph{term bukti} bagi $P$.\index{term bukti} Dua term disebut \emph{sama secara definisi} bila keduanya
menyusut ke bentuk identik melalui perhitungan simbolik.\index{proposisi-tipe}
\end{definisi}

Kamus ini, yaitu korespondensi antara logika dan tipe, dikenal sebagai Curry--Howard.\index{Curry-Howard} Implikasi
menjadi tipe fungsi, karena bukti $P\Rightarrow Q$ harus mengubah sebarang bukti $P$ menjadi bukti
$Q$, dan itulah pekerjaan sebuah fungsi. Kuantor universal menjadi fungsi dependen, karena bukti
$\forall x,\ P(x)$ harus menghasilkan bukti $P(a)$ untuk setiap $a$ yang diberikan. Konjungsi
menjadi pasangan, karena bukti $P\wedge Q$ memuat kedua buktinya sekaligus. Baris terakhir tabel
di bawah adalah yang paling penting: $\bot$ tidak punya penghuni.

\begin{center}
\begin{tabular}{lll}
\hline
Logika & Tipe di Lean & Cara membangun term \\
\hline
$P\Rightarrow Q$   & $P\to Q$              & \texttt{fun h => ...} \\
$P\wedge Q$        & \texttt{And P Q}      & \texttt{And.intro} \\
$P\vee Q$          & \texttt{Or P Q}       & \texttt{Or.inl}, \texttt{Or.inr} \\
$\top$             & \texttt{True}         & \texttt{True.intro} \\
$\bot$             & \texttt{False}        & tidak ada penghuni \\
$\forall x,\ P(x)$ & $\forall x,\ P\,x$    & \texttt{fun x => ...} \\
$\exists x,\ P(x)$ & \texttt{Exists}       & \texttt{Exists.intro} \\
\hline
\end{tabular}
\end{center}

\kamus{term bukti}{proof term}{Struk bangunan yang bentuknya persis mengikuti cetak biru
pernyataannya; kalau bentuknya tidak cocok, bangunan ditolak.}

Term berikut adalah bukti paling kecil yang berguna: fungsi identitas sebagai bukti implikasi
identitas, dan komposisi dua fungsi sebagai bukti transitivitas implikasi. Keduanya diperiksa
tanpa memakai pengetahuan apa pun tentang isi $P$, $Q$, dan $R$.

\begin{verbatim}
-- bukti implikasi identitas, gaya term
def bukti_id {P : Prop} : P -> P := fun h => h

-- komposisi dua implikasi
def bukti_komposisi {P Q R : Prop} (h1 : P -> Q) (h2 : Q -> R) : P -> R :=
  fun hp => h2 (h1 hp)

#check bukti_id
#check bukti_komposisi
#check Prop
\end{verbatim}

Keluaran perintah terakhir, apa adanya, adalah \texttt{Prop : Type}: himpunan proposisi sendiri
adalah sebuah tipe. Dua baris sebelumnya mencetak tipe kedua definisi. Perhatikan bahwa
\texttt{\#check} tidak membuktikan apa pun; ia hanya menanyakan tipe, dan itulah yang juga
dilakukan inti sistem ketika memeriksa.

\begin{proposisi}
Untuk setiap $P:\mathtt{Prop}$ berlaku $P\to P$. Untuk $P,Q,R:\mathtt{Prop}$ berlaku
$(P\to Q)\to(Q\to R)\to P\to R$.
\end{proposisi}

\begin{proof}[Bukti]
Term $\lambda h.\,h$ bertipe $P\to P$: dalam konteks $h:P$, variabel $h$ sendiri adalah bukti $P$,
sehingga abstraksi atas $h$ menghasilkan fungsi dari $P$ ke $P$. Untuk pernyataan kedua, misalkan
$h_{1}:P\to Q$, $h_{2}:Q\to R$, dan $h_{p}:P$. Aplikasi $h_{1}\,h_{p}$ bertipe $Q$, lalu aplikasi
$h_{2}\,(h_{1}\,h_{p})$ bertipe $R$. Tiga abstraksi berturut-turut atas $h_{p}$, $h_{2}$, dan
$h_{1}$ memberi term bertipe $(P\to Q)\to(Q\to R)\to P\to R$.
\end{proof}

Pohon pada Gambar~\ref{fig:pohon-tipe} menuliskan bukti kedua itu sebagai pohon turunan: tiga daun,
dua batang aplikasi, satu akar. Daun paling kanan adalah hipotesis yang diberikan, bukan hasil
turunan.

\begin{figure}[htbp]
\centering
\begin{tikzpicture}[font=\small,
  daun/.style={draw=layergrey,fill=layergrey!6,rounded corners=2pt,inner sep=4pt},
  hip/.style={draw=layergrey,dashed,fill=layergrey!3,rounded corners=2pt,inner sep=4pt},
  batang/.style={thick,layergrey}]
  \node[daun] (h1) at (-3.7,2.7) {\texttt{h1 : P -> Q}};
  \node[daun] (hp) at (-0.6,2.7) {\texttt{hp : P}};
  \node[daun] (ap1) at (-2.2,1.35) {\texttt{h1 hp : Q}};
  \node[hip]  (h2) at (2.7,1.35) {\texttt{h2 : Q -> R}};
  \node[daun] (akar) at (0,0) {\texttt{h2 (h1 hp) : R}};
  \draw[batang] (h1.south) -- (ap1.north west);
  \draw[batang] (hp.south) -- (ap1.north east);
  \draw[batang] (ap1.south) -- (akar.north west);
  \draw[batang] (h2.south) -- (akar.north east);
\end{tikzpicture}
\caption{Pohon tipe untuk term \texttt{h2 (h1 hp)}: tiga daun dan dua batang aplikasi, dengan akar
bertipe $R$. Kotak putus-putus adalah hipotesis yang diberikan, bukan hasil turunan.}
\label{fig:pohon-tipe}
\end{figure}

Ada satu sifat yang perlu disebut karena ia menjelaskan mengapa bukti dipakai sebagai sertifikat,
bukan sebagai data. Inti Lean memperlakukan dua bukti mana pun dari proposisi yang sama sebagai
sama secara definisi. Sifat ini disebut \emph{irrelevansi bukti}, dan akibatnya sebuah bukti tidak
membawa informasi selain ``pernyataan ini benar''. Karena itu bukti dapat dibuang setelah
diperiksa, dan yang tersisa hanyalah status pernyataannya.

Satu peringatan dipasang sejak awal. Korespondensi Curry--Howard berlaku penuh untuk logika
konstruktif. Begitu kita memakai hukum tengah-tengah, $\mathtt{Classical.em}:P\vee\neg P$, kita
menambahkan aksioma, dan bukti yang memakainya berhenti bersifat komputasional: ia tetap sah,
tetapi tidak lagi dibaca sebagai program yang menghasilkan saksi. Banyak teorema pustaka besar
memuat aksioma pilihan di daftar aksiomanya, dan bagian akhir bab ini menunjukkan cara menanyakan
daftar itu langsung kepada Lean.

\section{Kuantor: aturan intro dan eliminasi}

\tingkat{menengah}

Pernyataan yang menarik hampir selalu berkuantor.\index{kuantor} ``Untuk setiap bilangan asli'' dan ``ada
bilangan asli'' bukan hiasan bahasa; keduanya mengubah jenis objek yang harus disediakan. Dalam
pandangan tipe, perbedaan itu menjadi perbedaan antara fungsi dan pasangan.

\kamus{kuantor}{quantifier}{Kata ``setiap'' dan ``ada''; yang pertama menuntut resep untuk semua
kasus, yang kedua menuntut satu contoh beserta buktinya.}

Aturan untuk kuantor universal sepenuhnya sama dengan aturan untuk fungsi. Untuk membuktikan
$\forall x,\ P(x)$, kita mengambil $x$ sebarang dan membuktikan $P(x)$ bagi $x$ itu; itulah
abstraksi. Untuk memakai hipotesis $\forall x,\ P(x)$, kita terapkan hipotesis itu pada sebuah
objek; itulah aplikasi. Aturan untuk kuantor eksistensial adalah aturan pasangan: untuk
membuktikan $\exists x,\ P(x)$ kita menyediakan satu saksi $a$ bersama bukti $P(a)$; untuk memakai
hipotesis eksistensial kita membongkar pasangan itu menjadi saksi dan buktinya.

\begin{center}
\begin{tabular}{lll}
\hline
Bentuk pernyataan & Cara membuktikan & Cara memakai \\
\hline
$\forall x,\ P(x)$      & \texttt{intro x}              & terapkan pada objek \\
$\exists x,\ P(x)$      & \texttt{exact Exists.intro a h} & \texttt{cases} pada hipotesis \\
$P\wedge Q$             & \texttt{constructor}          & \texttt{h.1}, \texttt{h.2} \\
$P\vee Q$               & \texttt{left}, \texttt{right} & \texttt{cases} pada hipotesis \\
$\neg P$                & \texttt{intro h} lalu capai $\bot$ & terapkan pada bukti $P$ \\
\hline
\end{tabular}
\end{center}

\begin{tangan}
Siapkan sepuluh kartu kosong dan satu spidol. Di lima kartu tulis aturan masuk: ``ambil sebarang'',
``sediakan saksi'', ``pasangkan dua bukti'', ``pilih kiri'', ``pilih kanan''. Di lima kartu lain
tulis aturan pakai: ``terapkan pada objek'', ``bongkar pasangan'', ``ambil komponen pertama'',
``pecah per kasus'', ``pakai pada bukti lawan''.

Letakkan satu kartu tujuan di ujung meja, misalnya ``setiap bilangan asli nol atau penerus''. Susun
rantai ke arah kartu tujuan: kartu pertama harus aturan masuk yang bentuknya cocok dengan tujuan,
bukan aturan pakai. Setiap kartu berikutnya hanya boleh bekerja pada sisa yang masih terbuka. Kartu
yang tidak punya sisa untuk dikerjakan disingkirkan ke tepi meja; di Lean kartu itu tidak akan
pernah dipakai.

Yang harus kamu lihat: rantai memendek menuju tujuan, dan tidak pernah ada dua kartu yang mengerjakan
sisa yang sama. Kalau itu terjadi, rantaimu bercabang, dan mesin menuntut dua bukti, bukan satu.
\end{tangan}

Contoh berikut menunjukkan keduanya bekerja pada bilangan asli. Pernyataan pertama adalah
tautologi yang tidak memerlukan penalaran apa pun; pernyataan kedua menyatakan bahwa setiap
bilangan asli adalah nol atau penerus sesuatu, dan pembuktiannya memerlukan pemecahan kasus pada
bentuk bilangan.

\begin{verbatim}
-- kuantor universal: ambil n sebarang, lalu hitung
theorem nol_kanan : forall n : Nat, n + 0 = n := by
  intro n
  rfl

-- kuantor eksistensial: sediakan saksi, lalu buktikan sifatnya
theorem ada_tiga : exists n : Nat, n = 3 := by
  exact Exists.intro 3 rfl

-- setiap bilangan asli adalah nol atau penerus sesuatu
theorem nol_atau_suksesor (n : Nat) : n = 0 \/ exists m : Nat, n = m + 1 := by
  cases n with
  | zero => left; rfl
  | succ k => right; exact Exists.intro k rfl
\end{verbatim}

Ketiga teorema itu diperiksa tanpa aksioma tambahan, dan itu dapat dipastikan dengan perintah
\texttt{\#print axioms}, yang akan kita pakai berulang kali. Perhatikan juga bahwa
\texttt{Exists.intro k rfl} memakai \texttt{rfl} sebagai bukti kesamaan: pada kasus
\texttt{succ k}, tujuannya adalah $k+1=k+1$, yang benar secara definisi.

\begin{catatan}
Arah sebaliknya tidak berlaku otomatis. Dari $\exists x,\ P(x)$ kita tidak boleh menyimpulkan
$\forall x,\ P(x)$, dan Lean akan menolak setiap percobaan ke arah itu. Kesalahan menukar urutan
kuantor adalah kesalahan spesifikasi yang paling sering terjadi, dan ia lolos tanpa keluhan dari
pengujian karena semua sampel uji memenuhinya. Bagian tentang batas di akhir bab ini membahas
akibatnya.
\end{catatan}

\section{Taktik: mengubah tujuan menjadi term}

\tingkat{dasar}

Menulis term bukti secara langsung cepat menjadi tidak terbaca, karena term yang dalam harus
dibaca dari dalam ke luar. Karena itu Lean menyediakan mode taktik.\index{taktik} Setelah kata \texttt{by},
kita tidak menulis term, melainkan urutan perintah yang mengubah \emph{keadaan tujuan}. Keadaan
tujuan dicetak sebagai daftar hipotesis, lalu garis pemisah, lalu target.

\kamus{taktik}{tactic}{Perintah memasak langkah demi langkah, bukan resep lengkap yang harus
dihafalkan sekaligus.}

\begin{verbatim}
-- keadaan tujuan dicetak setelah setiap taktik; baris "|-" dibaca "maka"
theorem imp_trans {P Q R : Prop} (h1 : P -> Q) (h2 : Q -> R) : P -> R := by
  intro hp
  trace_state          -- P Q R : Prop, h1 : P -> Q, h2 : Q -> R, hp : P, |- R
  apply h2
  trace_state          -- target berubah menjadi |- Q
  exact h1 hp
\end{verbatim}

Keluaran \texttt{trace\_state} pada mesin uji ini adalah dua blok berikut. Lean mencetak panah
sebagai karakter Unicode; di sini semuanya ditulis ulang dalam ASCII agar konsisten dengan seluruh
bab.

\begin{verbatim}
P Q R : Prop
h1 : P -> Q
h2 : Q -> R
hp : P
|- R

P Q R : Prop
h1 : P -> Q
h2 : Q -> R
hp : P
|- Q
\end{verbatim}

Taktik \texttt{intro} mengambil target berbentuk implikasi dan memindahkan antesedennya menjadi
hipotesis; itu persis abstraksi fungsi. Taktik \texttt{exact} menutup target dengan term yang
tipenya cocok persis. Taktik \texttt{apply} menyamakan kesimpulan sebuah lemma dengan target, lalu
meninggalkan sisa argumen sebagai tujuan baru. Tiga taktik ini sudah cukup untuk semua bukti
logika murni di bab ini, dan ketiganya adalah bentuk permukaan dari aturan yang sama.

Gambar~\ref{fig:dua-keadaan} memadatkan dua keadaan itu menjadi dua kotak. Yang berubah hanya baris
di bawah garis: sasaran menyusut dari $R$ ke $Q$. Daftar hipotesis tidak bertambah satu baris pun,
karena \texttt{apply} tidak menciptakan bukti baru; ia hanya menurunkan tuntutan.

\begin{figure}[htbp]
\centering
\begin{tikzpicture}[font=\small,
  kotak/.style={draw=layergrey,fill=layergrey!6,rounded corners=3pt,align=left,inner sep=6pt,
                text width=3.9cm},
  panah/.style={-{Stealth[length=2.2mm]},thick,layerteal}]
  \node[kotak] (awal) at (0,0) {\textbf{keadaan 1}\\[3pt]
    \texttt{P Q R : Prop}\\
    \texttt{h1 : P -> Q}\\
    \texttt{h2 : Q -> R}\\
    \texttt{hp : P}\\[3pt]
    \rule{3.6cm}{0.4pt}\\
    $\vdash$ \texttt{R}};
  \node[kotak] (akhir) at (6.7,0) {\textbf{keadaan 2}\\[3pt]
    \texttt{P Q R : Prop}\\
    \texttt{h1 : P -> Q}\\
    \texttt{h2 : Q -> R}\\
    \texttt{hp : P}\\[3pt]
    \rule{3.6cm}{0.4pt}\\
    $\vdash$ \texttt{Q}};
  \draw[panah] (awal.east) -- node[above,font=\scriptsize]{\texttt{apply h2}} (akhir.west);
\end{tikzpicture}
\caption{Dua keadaan tujuan pada bukti \texttt{imp\_trans}: hipotesis tetap empat baris, sasaran
menyusut dari $R$ menjadi $Q$ setelah satu taktik. Kotak kedua masih terbuka; ia ditutup
\texttt{exact h1 hp}.}
\label{fig:dua-keadaan}
\end{figure}

\begin{center}
\begin{tabular}{lll}
\hline
Taktik & Yang dikerjakan & Kapan dipakai \\
\hline
\texttt{intro}     & anteseden menjadi hipotesis   & target berbentuk $P\to Q$ \\
\texttt{exact}     & menutup target dengan term    & term sudah tersedia \\
\texttt{apply}     & kesimpulan lemma cocok target & lemma berujung sama \\
\texttt{rw}        & menulis ulang dengan persamaan & ada hipotesis $a=b$ \\
\texttt{simp}      & menyederhanakan berulang      & banyak langkah sepele \\
\texttt{cases}     & memecah target per konstruktor & perlu analisis kasus \\
\texttt{induction} & basis dan langkah induksi     & target rekursif atas \texttt{Nat} \\
\hline
\end{tabular}
\end{center}

Taktik \texttt{rw} menulis ulang target memakai persamaan. Ia bekerja pada arah yang diberikan,
dan arah itu penting: menulis ulang ke arah yang salah tidak akan menemukan pola yang dicari.
Perhatikan dua percobaan berikut, yang satu berhasil dan yang satu ditolak.

\begin{verbatim}
theorem geser (a b : Nat) (h : a = b) : a + 1 = b + 1 := by
  rw [h]        -- target menjadi b + 1 = b + 1, lalu ditutup secara definisi
  rfl

-- komutatifitas penjumlahan TIDAK berlaku secara definisi, jadi rfl gagal
theorem tukar_salah (n : Nat) : n + 1 = 1 + n := by
  rfl           -- error: The rfl tactic failed

-- versi yang benar: menulis ulang dengan lemma komutatifitas
theorem tukar_benar (n : Nat) : n + 1 = 1 + n := by
  rw [Nat.add_comm]
\end{verbatim}

Kegagalan di tengah itu adalah pelajaran, bukan kecelakaan. Penjumlahan pada bilangan asli
didefinisikan dengan rekursi pada argumen kedua, sehingga $n+1$ menyusut menjadi
$\mathtt{succ}\,n$ sedangkan $1+n$ menyusut menjadi $\mathtt{succ}\,(0+n)$, dan keduanya berhenti
di bentuk yang berbeda. Yang kita sebut ``sifat komutatif'' di sekolah adalah teorema, bukan
definisi, dan Lean menolak memperlakukannya seolah-olah definisi. Inilah beda paling terasa antara
membaca dan menulis bukti formal: setiap hal yang biasa kita lewati harus disebutkan namanya.

Taktik \texttt{cases} memecah target menurut konstruktor tipe. Pada bilangan asli ia memberi dua
kasus, nol dan penerus. Taktik \texttt{simp} menutup tujuan dengan menulis ulang memakai kumpulan
lemma yang ditandai; ia berguna tetapi juga berbahaya, karena keluarannya bergantung pada
kumpulan lemma yang sedang aktif. Bukti yang diawali \texttt{simp} tanpa argumen bisa berhenti
bekerja ketika pustaka berubah, dan itulah sebabnya pada bab ini \texttt{simp} selalu disebut
bersama daftar lemma yang dipakai.

Satu peringatan tentang taktik yang tidak ada di sini. Nama-nama seperti \texttt{linarith},
\texttt{ring}, \texttt{norm\_num}, dan \texttt{induction'} berasal dari pustaka Mathlib, bukan
dari inti Lean. Semua potongan kode di bab ini dijalankan dengan Lean 4.10.0 tanpa Mathlib,
sehingga taktik itu tidak tersedia dan tidak dipakai; sebagai gantinya dipakai \texttt{omega}
(yang ada di inti) dan \texttt{simp} dengan daftar lemma eksplisit. Perbedaan ini bukan soal
gaya: pustaka menambah ribuan lemma dan beberapa taktik, dan setiap tambahan itu ikut menentukan
apa yang bisa ditutup dalam satu baris.

\begin{untukanak}
Ambil setumpuk batu kecil. Satu batu kamu susun menjadi bujur sangkar terkecil. Lalu tambahkan
tiga batu mengelilinginya, dan hasilnya bujur sangkar yang lebih besar. Tambahkan lima batu lagi,
dan hasilnya bujur sangkar yang lebih besar lagi. Tambahkan tujuh, sembilan, sebelas, dan
seterusnya, selalu bertambah dua, dan setiap kali hasilnya tetap bujur sangkar yang rapi.

Sekarang pertanyaannya: apakah ini akan terus begitu kalau kita menambahkan sampai seratus kali,
sampai sejuta kali? Kamu tidak bisa mencoba semuanya. Kamu juga tidak perlu. Kalau kamu bisa
menunjukkan bahwa pola penambahan batu selalu menutup satu sudut dan satu sisi bujur sangkar
dengan rapi, maka satu penjelasan kecil berlaku untuk semua ukuran, termasuk ukuran yang tidak
akan pernah bisa kamu susun dengan tangan.

Komputer pemalas tidak mau mencoba sejuta kasus, dan komputer yang terlalu percaya diri tidak mau
memeriksa apa pun. Di bab ini kita akan tahu mengapa ada cara ketiga: menjelaskan polanya sekali,
lalu meminta mesin memeriksa penjelasan itu sampai ke huruf terakhirnya.
\end{untukanak}

\section{Bukti induksi lengkap: jumlah bilangan ganjil}

\tingkat{lanjut}

Kita akan membuktikan satu pernyataan aritmetika yang seluruhnya dapat diperiksa mesin, yaitu
jumlah $n$ bilangan ganjil pertama sama dengan $n^{2}$. Pernyataan ini dipilih karena tiga alasan:
rumusnya dapat dihitung dengan Python tanpa Lean, buktinya memerlukan induksi, dan langkah
aljabarnya cukup kecil sehingga setiap taktik dapat dipertanggungjawabkan.

\kamus{induksi}{induction}{Menjelaskan pola dengan dua hal saja: kasus terkecil, dan cara naik
satu langkah dari kasus mana pun.}

Sebelum menulis kode, kita nyatakan pernyataannya dan buktinya dalam bahasa biasa. Fungsinya
didefinisikan secara rekursif: jumlah nol bilangan ganjil pertama adalah nol, dan jumlah $n+1$
bilangan ganjil pertama adalah jumlah $n$ bilangan ganjil pertama ditambah bilangan ganjil ke
$n+1$, yaitu $2n+1$.

\begin{definisi}[Jumlah bilangan ganjil]
Fungsi $S:\mathbb{N}\to\mathbb{N}$ didefinisikan oleh $S(0)=0$ dan
$S(n+1)=S(n)+(2n+1)$ untuk setiap $n\in\mathbb{N}$.\index{induksi}
\end{definisi}

Angka $2n+1$ bukan pilihan sembarang: bilangan ganjil ke-$n$ memang berbentuk itu, dan
$2n+1$ juga selisih antara $(n+1)^{2}$ dan $n^{2}$. Kesamaan terakhir itulah inti buktinya.

\begin{lemma}
Untuk setiap $k\in\mathbb{N}$ berlaku $(k+1)^{2}=k^{2}+2k+1$.
\end{lemma}

\begin{proof}[Bukti]
$(k+1)^{2}=(k+1)(k+1)=k(k+1)+(k+1)=k^{2}+k+k+1=k^{2}+2k+1$. Setiap langkah memakai
distributivitas dan asosiativitas penjumlahan pada bilangan asli.
\end{proof}

\begin{teorema}
Untuk setiap $n\in\mathbb{N}$ berlaku $S(n)=n^{2}$.
\end{teorema}

\begin{proof}[Bukti]
Dengan induksi pada $n$. Untuk $n=0$, kedua sisi bernilai $0$. Misalkan pernyataan benar untuk
$k$, yaitu $S(k)=k^{2}$. Maka
$S(k+1)=S(k)+(2k+1)=k^{2}+2k+1=(k+1)^{2}$, dengan lemma di atas pada langkah terakhir. Jadi
pernyataan benar untuk $k+1$. Menurut prinsip induksi, pernyataan benar untuk setiap
$n\in\mathbb{N}$.
\end{proof}

Bukti di atas adalah bukti informal yang lengkap, dalam arti tidak ada langkah yang dilewati.
Sekarang kita menuliskannya untuk mesin. Perhatikan bahwa definisi rekursif dan teoremanya hampir
sama kata demi kata dengan yang baru kita tulis, sedangkan buktinya berubah menjadi urutan taktik.

\begin{verbatim}
-- definisi rekursif: dua persamaan, tanpa ":=" pada baris pertama
def oddSum : Nat -> Nat
  | 0     => 0
  | n + 1 => oddSum n + (2 * n + 1)

theorem oddSum_sq (n : Nat) : oddSum n = n * n := by
  induction n with
  | zero => rfl
  | succ k ih =>
    rw [oddSum]        -- buka definisi: oddSum (k+1) = oddSum k + (2*k+1)
    rw [ih]            -- pakai hipotesis induksi: oddSum k = k*k
    simp [Nat.succ_eq_add_one, Nat.mul_succ, Nat.add_mul, Nat.one_mul]
    omega              -- tutup sisa linear
\end{verbatim}

Baris \texttt{induction n with} membangkitkan dua tujuan. Tujuan pertama bertanda \texttt{zero}
dan berbunyi $S(0)=0\cdot 0$, yang ditutup \texttt{rfl} karena kedua sisi menyusut ke nol. Tujuan
kedua bertanda \texttt{succ k} dan membawa satu hipotesis tambahan bernama \texttt{ih}, yaitu
$S(k)=k\cdot k$. Nama \texttt{ih} dipilih sendiri; yang penting perannya.

Baris \texttt{rw [oddSum]} membuka definisi pada sisi kiri. Setelah baris itu, tujuan berbunyi
$S(k)+(2k+1)=(k+1)(k+1)$. Baris \texttt{rw [ih]} mengganti $S(k)$ dengan $k\cdot k$, sehingga
tujuan menjadi $k\cdot k+(2k+1)=(k+1)(k+1)$.

Di titik ini kita berada persis di tempat lemma tadi bekerja, tetapi mesin belum tahu bahwa
$(k+1)(k+1)$ dapat ditulis sebagai $k\cdot k+2k+1$. \texttt{simp} dengan empat lemma eksplisit
mengubah bentuk ruas kanan menjadi $k\cdot k+k+(k+1)$, yaitu bentuk yang seluruhnya linear dalam
$k$. Baris terakhir, \texttt{omega}, menutup tujuan linear semacam itu.

Perlu ditegaskan bahwa \texttt{omega} bukan sulap dan bukan pula aksioma tersembunyi. Ia adalah
prosedur keputusan untuk aritmetika linear atas bilangan bulat, dan ia hanya menerima tujuan yang
seluruh operasinya linear. Bila kita memanggilnya sebelum \texttt{simp}, tujuannya masih memuat
$k\cdot k$ dan $(k+1)(k+1)$ sekaligus. Pada mesin uji ini, berkas percobaan sementara bernama
\texttt{odd.lean} keluarannya menolak:

\begin{verbatim}
odd.lean:11:4: error: omega could not prove the goal:
a possible counterexample may satisfy the constraints
  a - b + 2*c <= -2
  ...
\end{verbatim}

Pesan itu jujur dan berguna: ia menyebut kemungkinan penyangkal, dan justru karena itu kita tahu
bahwa yang gagal adalah bentuk tujuan, bukan kebenaran pernyataannya. Urutan taktik karena itu
bukan selera; ia mengubah jenis tujuan yang dihadapi prosedur keputusan.

\begin{catatan}
Bila \texttt{simp} di atas dijalankan tanpa \texttt{omega}, tujuan yang tersisa berbentuk
$k\cdot k+(1+2k)=k+(k+(k\cdot k+1))$. Secara matematis itu benar, tetapi tidak linear lagi, dan
penyederhana tidak menutupnya. Pada mesin uji ini percobaan tersebut berakhir dengan
\texttt{unsolved goals}. Inilah bentuk paling umum dari kegagalan bukti formal: bukan
pernyataannya salah, melainkan bentuk tujuannya tidak cocok dengan alat yang dipakai.
\end{catatan}

\begin{akibat}
Untuk setiap $n\in\mathbb{N}$, jumlah $n$ bilangan ganjil pertama adalah kuadrat sempurna.
\end{akibat}

Akibat itu hanya pengulangan teorema dengan kata lain, dan itu memang gunanya: ia menyatakan
kembali hasil dalam bahasa yang dipakai di bagian ini, sehingga pembaca yang menemui rumusnya di
tempat lain dapat mengenalinya.

\section{Struktur, typeclass, dan apa yang benar-benar dijamin}

\tingkat{lanjut}

Sertifikat jarang berupa satu teorema lepas. Biasanya ia sekumpulan data bersama klaim tentang
data itu: ambang, tabel, nama berkas, dan beberapa pernyataan. Lean menyediakan dua alat untuk
mengelompokkan hal semacam ini, yaitu \emph{structure} dan \emph{typeclass}.

\kamus{typeclass}{type class}{Daftar kemampuan yang harus dimiliki sebuah tipe agar bisa dipakai
di tempat tertentu, seperti syarat lisensi untuk mengemudikan kendaraan jenis tertentu.}

\begin{definisi}[Struktur dan typeclass]
Sebuah \emph{struktur} adalah tipe yang bidangnya diberi nama, sehingga nilainya dibangun dengan
menyebut nama bidang. Sebuah \emph{typeclass} adalah struktur yang dipakai untuk mencari
\emph{instance} secara otomatis: ketika sebuah definisi memerlukan kemampuan tertentu dari sebuah
tipe, Lean mencari instance yang cocok di lingkungan yang sedang aktif.\index{typeclass}
\end{definisi}

Perbedaan praktisnya sederhana. Struktur dipakai ketika data itu sendiri yang penting, misalnya
sertifikat dengan bidang klaim dan ambang. Typeclass dipakai ketika yang penting adalah kemampuan:
sesuatu yang dapat dipakai sebagai bukti, sesuatu yang dapat dihitung normanya, sesuatu yang dapat
dibandingkan. Pada keduanya, pemeriksaan tetap sama ketatnya: setiap bidang harus diisi, dan
setiap instance harus lolos pemeriksaan tipe.

\begin{verbatim}
-- struktur: data yang bidangnya bernama
structure Sertifikat where
  klaim  : String
  ambang : Nat

-- typeclass: kemampuan yang dicari otomatis
class Terperiksa (a : Type) where
  jumlah_teorema : a -> Nat

instance : Terperiksa Sertifikat where
  jumlah_teorema _ := 1

def contoh : Sertifikat := { klaim := "oddSum", ambang := 100 }

#eval (Terperiksa.jumlah_teorema contoh)   -- mencetak 1
\end{verbatim}

Sekarang bagian yang paling sering dilewati dalam pembicaraan populer tentang Lean, yaitu
pertanyaan apa yang sebenarnya dijamin oleh penerimaan sebuah teorema. Empat hal dijamin, dan
empat hal lain tidak.

Gambar~\ref{fig:alur-periksa} menaruh ketiga tahap itu pada satu baris: pernyataan ditulis, term
diajukan, inti memutuskan. Hanya ada dua keluaran; tidak ada kotak ketiga yang berbunyi ``mungkin''.

\begin{figure}[htbp]
\centering
\begin{tikzpicture}[font=\small,
  tahap/.style={draw=layergrey,fill=layergrey!6,rounded corners=3pt,align=center,inner sep=5pt,
                text width=2.9cm,minimum height=1.7cm},
  panah/.style={-{Stealth[length=2.2mm]},thick,layergrey},
  terima/.style={draw=layerjade,fill=layerjade!8,rounded corners=3pt,align=center,inner sep=5pt,
                 text width=2.9cm},
  tolak/.style={draw=accent,fill=accent!6,rounded corners=3pt,align=center,inner sep=5pt,
                text width=2.9cm}]
  \node[tahap] (s1) at (0,0) {\textbf{1. pernyataan}\\$P$ bertipe \texttt{Prop}};
  \node[tahap] (s2) at (3.9,0) {\textbf{2. term}\\$p$ diajukan untuk $P$};
  \node[tahap] (s3) at (7.8,0) {\textbf{3. inti}\\periksa apakah $p:P$};
  \draw[panah] (s1) -- (s2);
  \draw[panah] (s2) -- (s3);
  \node[terima] (ya) at (5.4,-2.9) {\textbf{diterima}\\ada term bertipe $P$};
  \node[tolak] (tidak) at (9.9,-2.9) {\textbf{ditolak}\\tidak ada term};
  \draw[panah] (s3.south west) -- node[above left,font=\scriptsize]{cocok} (ya.north);
  \draw[panah] (s3.south east) -- node[above right,font=\scriptsize]{tidak cocok} (tidak.north);
\end{tikzpicture}
\caption{Alur pemeriksaan: tiga tahap (pernyataan, term, inti) dan dua keluaran, diterima atau
ditolak. Tidak ada keluaran ketiga; berkas yang berhenti di tengah dilaporkan sebagai tujuan yang
belum tertutup, bukan sebagai ``hampir benar''.}
\label{fig:alur-periksa}
\end{figure}

Yang dijamin: (i) ada term yang tipenya persis pernyataan yang ditulis; (ii) term itu diperiksa
inti, bukan oleh taktik yang menghasilkannya; (iii) pemeriksaan hanya memakai definisi, teorema
lain yang sudah diperiksa, dan aksioma yang tercatat; (iv) daftar aksioma yang dipakai dapat
ditanyakan kapan saja dengan \texttt{\#print axioms}. Butir (iv) adalah alat paling berguna di
seluruh bab ini, karena ia mengubah ``percaya saja'' menjadi ``periksa sendiri''.

Yang tidak dijamin: (i) bahwa pernyataannya menyatakan apa yang kita maksud; (ii) bahwa
definisinya memodelkan benda di dunia nyata; (iii) bahwa inti, pengurai, dan perangkat kerasnya
bebas dari kesalahan; (iv) bahwa taktik yang dipakai tidak menambahkan asumsi tersembunyi.
Butir (iv) bukan cela teoretis: ia dapat diperiksa, dan hasilnya sering mengejutkan.

\begin{contoh}[Dua cara membuktikan hal yang sama, dua daftar aksioma]
Perhatikan dua pembuktian pernyataan $S(10)=100$ yang keduanya benar. Yang pertama memakai
\texttt{decide}, yang menjalankan perhitungan di dalam inti. Yang kedua memakai
\texttt{native\_decide}\index{native-decide}, yang mengompilasi perhitungan menjadi kode mesin
dan menjalankannya. Kedua pernyataannya sama, tetapi daftar aksiomanya berbeda, dan perbedaan itu
terlihat dari perintah \texttt{\#print axioms}.
\end{contoh}

\begin{verbatim}
def oddSum : Nat -> Nat
  | 0     => 0
  | n + 1 => oddSum n + (2 * n + 1)

theorem cek_native : oddSum 10 = 100 := by native_decide
theorem cek_kernel : oddSum 10 = 100 := by decide

#print axioms cek_native
#print axioms cek_kernel
\end{verbatim}

Keluaran pada mesin uji ini, apa adanya:

\begin{verbatim}
'cek_native' depends on axioms: [Lean.ofReduceBool]
'cek_kernel' does not depend on any axioms
\end{verbatim}

Baris pertama adalah inti dari seluruh pembahasan tentang \texttt{native\_decide}. Perhitungannya
tidak dilakukan oleh inti, melainkan oleh kode terkompilasi, dan Lean mencatat kepercayaan itu
sebagai aksioma bernama \texttt{Lean.ofReduceBool}. Secara praktis ini bukan bencana: kompilator
Lean dipakai untuk mengompilasi Lean sendiri setiap hari. Secara prinsipil ini perbedaan besar,
karena batas kepercayaan berpindah dari inti yang kecil ke rantai panjang kompilator, penghubung,
dan prosesor. Sertifikat yang dibangun di atas \texttt{native\_decide} karena itu lebih lemah
daripada sertifikat yang dibangun di atas \texttt{decide}, meskipun keduanya mencetak ``diterima''
tanpa keluhan.

Pada teorema induksi tadi, daftar aksiomanya juga tidak kosong. Perintah \texttt{\#print axioms}
untuk \texttt{oddSum\_sq} mencetak \texttt{[propext, Quot.sound]}. Keduanya adalah aksioma standar yang dipakai untuk
penalaran klasik dan hasil bagi, dan keduanya muncul karena \texttt{omega} membuktikan tujuan
dengan kontradiksi di dalam bilangan bulat. Ini bukan kesalahan; ini fakta yang seharusnya
disebutkan bila sertifikat dipakai sebagai dasar klaim.

\section{Batas: menolak bukti salah, \texttt{sorry}, dan aksioma}

\tingkat{lanjut}

Kemampuan menolak sama pentingnya dengan kemampuan menerima. Berkas berikut sengaja berisi
kesalahan jenis yang berbeda, dan keluarannya adalah jawaban Lean untuk masing-masing.

\begin{verbatim}
-- Berkas: Salah.lean
theorem salah_arit : 2 + 2 = 5 := by decide

theorem salah_rfl (n : Nat) : n + 1 = 1 + n := by rfl

axiom magic : 1 = 2
theorem from_magic : 1 = 2 := magic
#print axioms from_magic

theorem belum_selesai (n : Nat) : n + 0 = n := by sorry
\end{verbatim}

Keluaran sebenarnya, dipotong pada bagian yang panjang:

\begin{verbatim}
Salah.lean:2:37: error: tactic 'decide' proved that the proposition
  2 + 2 = 5
is false
Salah.lean:4:50: error: The rfl tactic failed. Possible reasons:
- The goal is not a reflexive relation ...
- The arguments of the relation are not equal.
...
n : Nat
|- n + 1 = 1 + n
'from_magic' depends on axioms: [magic]
Salah.lean:10:8: warning: declaration uses 'sorry'
\end{verbatim}

Tiga jenis penolakan muncul di sini, dan ketiganya berbeda sifat. Yang pertama adalah penolakan
paling kuat: \texttt{decide} bukan hanya gagal menutup tujuan, ia membuktikan bahwa proposisinya
salah. Untuk pernyataan berhingga, kegagalan dan penyangkalan adalah hal yang sama. Yang kedua
adalah penolakan yang lebih halus: \texttt{rfl} menolak karena kedua sisi tidak sama secara
definisi, padahal pernyataannya benar. Di sini Lean tidak menyalahkan pernyataannya, ia
menyalahkan alat yang dipakai; pesannya bahkan menyarankan lemma refleksivitas eksplisit. Yang
ketiga bukan penolakan sama sekali.

Baris \texttt{axiom magic : 1 = 2} diterima tanpa keluhan. Lean tidak memeriksa kewajaran isi
aksioma; ia hanya mencatatnya.\index{aksioma} Setelah itu, \texttt{from\_magic} juga diterima, dan seluruh
aritmetika bisa dibongkar dengan \texttt{False.elim} dari sana. Satu-satunya tanda bahaya adalah
baris \texttt{\#print axioms}, yang menyebut \texttt{[magic]} dengan tenang, tanpa peringatan
dan tanpa warna merah. Inilah alasan praktis mengapa memeriksa daftar aksioma adalah kebiasaan
yang wajib: sistem tidak akan menegur kita.

\kamus{aksioma}{axiom}{Aturan yang kita minta diterima tanpa bukti, seperti asumsi pembuka dalam
perdebatan; sah, tetapi harus disebut.}

Baris terakhir berkas itu memakai \texttt{sorry}, dan keluarannya hanya peringatan
\texttt{declaration uses 'sorry'}.\index{sorry} Berkas tetap diproses, teorema tetap ``ada'', dan siapa pun
yang membaca keluaran dengan saksama akan melihat kata \texttt{warning}, bukan \texttt{error}.
Pada pustaka besar, satu \texttt{sorry} yang tertinggal di tengah ribuan teorema nyaris tidak
terlihat. Lean menyediakan jawabannya, dan jawabannya sama dengan jawaban untuk aksioma:
\texttt{\#print axioms} mencetak \texttt{[sorryAx]}. Pustaka matematika Lean, yaitu Mathlib,
memakai kebijakan yang lebih tegas: berkas yang memuat \texttt{sorry} ditolak oleh pemeriksaan
otomatis sebelum diterima. Kebijakan itu pantas ditiru siapa pun yang menerbitkan sertifikat.

\kamus{sorry}{sorry}{Catatan ``nanti saya kerjakan'' yang membuat pekerjaan tampak selesai;
mesin menerimanya, tetapi mencatatnya sebagai utang.}

Sekarang batas yang paling penting, dan yang paling jarang disebut. Verifikasi formal menjamin
bahwa term cocok dengan pernyataan. Ia tidak menjamin bahwa pernyataan cocok dengan maksud. Bila
spesifikasinya terlalu lemah, buktinya bisa sempurna dan tidak berguna.\index{spesifikasi} Contoh berikut
menunjukkannya dalam bentuk paling telanjang: pernyataan pertama mengatakan ``ada gangguan kecil
yang tidak mengubah keluaran'', pernyataan kedua mengatakan ``setiap gangguan kecil tidak
mengubah keluaran''. Yang pertama dapat dibuktikan untuk \emph{semua} fungsi, termasuk fungsi
terburuk yang bisa dibayangkan, karena saksi $\delta=0$ selalu tersedia.

\begin{verbatim}
def identitas (n : Nat) : Nat := n

-- spesifikasi lemah: ADA delta yang tidak mengubah keluaran (untuk semua f)
theorem ada_delta (f : Nat -> Nat) (x : Nat) :
    exists delta : Nat, f (x + delta) = f x := by
  exact Exists.intro 0 rfl

-- versi kuat salah untuk f = identitas, jadi tidak dapat dibuktikan
theorem ada_penyangkal :
    exists x : Nat, exists delta : Nat,
      Not (identitas (x + delta) = identitas x) := by
  exact Exists.intro 0 (Exists.intro 1 (by decide))
\end{verbatim}

Kedua teorema itu lolos tanpa aksioma tambahan, dan justru itulah masalahnya. Teorema pertama
tidak mengatakan apa pun tentang $f$; ia hanya mengatakan bahwa $\delta=0$ adalah gangguan yang
tidak mengubah apa pun, dan itu berlaku untuk setiap fungsi, termasuk pengklasifikasi yang salah
seratus persen. Sertifikat semacam ini lolos semua pemeriksaan dan tidak melindungi apa pun.
Teorema kedua menunjukkan sisi lain: pernyataan yang benar-benar kita inginkan, yaitu versi
berkuantor universal, salah untuk fungsi identitas, sehingga tidak ada bukti yang mungkin
ditemukan.

\begin{tantangan}
Sebuah sertifikat memuat satu baris yang tampak tidak mungkin salah: untuk setiap bilangan asli $a$
dan $b$ berlaku $(a-b)+b=a$. Di atas kertas alasannya sederhana: kurangi, tambahkan lagi, kembali ke
semula. Tulis baris itu di Lean dan tutup dengan \texttt{omega}; pemeriksa menolaknya, dan penolakan
itu bukan karena taktiknya kurang pintar.

Dua pertanyaan. Mengapa pernyataan itu tidak benar untuk semua $a$ dan $b$? Dan syarat tambahan apa
yang membuatnya benar sekaligus tetap dapat diperiksa mesin?

Petunjuk 1: jangan percaya pada tanda kurung; tanyakan apa yang terjadi bila $b$ lebih besar dari
$a$. Petunjuk 2: pengurangan pada bilangan asli berhenti di nol, jadi $0-1=0$, sehingga ruas kiri
menjadi $0-1+1=1$ sedangkan ruas kanan $0$.

Kalau mentok, kembali ke bagian ``Taktik: mengubah tujuan menjadi term'', tempat $n+1$ dan $1+n$
berhenti di bentuk yang berbeda.
\end{tantangan}

Sebelum membaca lebih jauh, coba jawab sendiri. Keluaran mesin uji ini, apa adanya, adalah pembanding
jawabanmu: berkas kedua memuat pernyataan yang sama dengan hipotesis $b\le a$.

\begin{verbatim}
$ lean t.lean
t.lean:1:49: error: omega could not prove the goal:
a possible counterexample may satisfy the constraints
  d >= 0
  c >= 0
  c - d >= 1
where
 c := b
 d := a
$ lean t2.lean
'potong_benar' depends on axioms: [propext, Quot.sound]
\end{verbatim}

Keluaran pertama menyebut kendala $b-a\ge 1$, yaitu penyangkal dengan $b>a$; pada $a=0$ dan $b=1$
ruas kiri bernilai $1$, ruas kanan $0$. Hipotesis $b\le a$ membuat baris yang sama lolos, dengan dua
aksioma standar yang sudah kita kenal dari \texttt{omega}. Bentuk perbaikannya bukan taktik
tambahan, melainkan syarat.

\kertas{Klaim yang boleh dibuat dari bab ini: potongan kode Lean yang ditampilkan telah
dijalankan dengan Lean 4.10.0 tanpa Mathlib pada mesin penulisan, dan keluarannya disalin apa
adanya; keluaran Unicode dicetak ulang dalam ASCII dan pemotongan ditandai dengan tiga titik.
Yang tidak dijalankan: taktik Mathlib (\texttt{linarith}, \texttt{ring}, \texttt{norm\_num},
\texttt{induction'}) dan pustaka Mathlib sendiri, karena tidak terpasang. Yang tidak dibuktikan di
sini dan tidak diklaim: sertifikat apa pun tentang jaringan saraf, model bahasa, atau ketahanan
terhadap serangan. Bab 28 tidak memuat satu pun bukti semacam itu.}

Batas ini ada hubungannya langsung dengan Bab 27 dan Bab 29. Evaluasi model yang jujur menuntut
kita menyebut distribusi tempat angka diukur; verifikasi formal menuntut kita menyebut pernyataan
yang dibuktikan. Keduanya adalah penyakit yang sama dalam bentuk berbeda: klaim yang tidak
menyebut syaratnya. Sertifikat ``model ini terbukti aman'' tanpa pernyataan formal di belakangnya
sama tidak bergunanya dengan ``akurasi 99 persen'' tanpa keterangan data ujinya. Pada Bab 29 dan
Bab 30, ketika kita membahas serangan dan pertahanan, pertanyaan pertama yang harus diajukan
kepada setiap sertifikat ketahanan adalah: berapa kuantornya, dan dalam ruang apa kuantor itu
berjalan?

\section{Contoh hitung}

\tingkat{menengah}

Angka-angka di bagian ini dihitung, bukan diperkirakan. Kolom pertama adalah $n$, kolom kedua
adalah bilangan ganjil ke-$n$ yang ditambahkan pada langkah itu, kolom ketiga adalah $S(n)$, dan
kolom keempat adalah $n^{2}$.

\begin{center}
\begin{tabular}{rrrr}
\hline
$n$ & suku ke-$n$ & $S(n)$ & $n^{2}$ \\
\hline
0  & --  & 0     & 0 \\
1  & 1   & 1     & 1 \\
2  & 3   & 4     & 4 \\
5  & 9   & 25    & 25 \\
10 & 19  & 100   & 100 \\
20 & 39  & 400   & 400 \\
100 & 199 & 10000 & 10000 \\
\hline
\end{tabular}
\end{center}

Nilai $S(100)=10000$ patut diperiksa dengan tangan sekali, karena ia menunjukkan bahwa rumus itu
bukan kebetulan untuk bilangan kecil: jumlah seratus bilangan ganjil, dari $1$ sampai $199$,
sama dengan $100\cdot 100$. Pemeriksaan dengan Python atas seluruh nilai $n<5000$ memberi
\texttt{True}, dan pemeriksaan identitas aljabar $(k+1)^{2}=k^{2}+2k+1$ untuk seluruh $k<100000$
juga memberi \texttt{True}. Kedua pemeriksaan itu bukan bukti, tetapi keduanya menutup kemungkinan
bahwa rumus yang dibuktikan salah ketik.

Ada satu angka lagi yang lebih informatif daripada tabel di atas, yaitu ukuran bukti setelah
dielaborasi. Kalimat taktik untuk \texttt{oddSum\_sq} panjangnya delapan baris. Ketika Lean
mencetak term hasil elaborasinya dengan \texttt{\#print}, hasilnya 10.783 karakter dalam 149
baris. Artinya, satu baris taktik rata-rata mengembang menjadi sekitar sembilan belas baris term.
Tidak ada manusia yang membaca term itu; yang membacanya adalah inti, dan itulah pembagian kerja
yang menjadi inti seluruh metode ini.

Angka terakhir menyangkut batas pemeriksaan langsung. Pada pengaturan bawaan, batas kedalaman
rekursi bernilai 1000. Dengan batas itu, taktik \texttt{decide} berhasil membuktikan
$S(n)=n^{2}$ hanya untuk $n\le 122$. Percobaan pertama yang gagal adalah $n=123$, dan pesan yang
keluar menyebut bahwa kedalaman rekursi maksimum telah tercapai. Bila batas dinaikkan ke 100000,
nilai $n=2000$ berhasil dibuktikan dalam 0,48 detik. Pada $n=20000$ pemeriksaan gagal karena
rekursi terlalu dalam untuk ruang tumpukan. Jadi pemeriksaan langsung bukan alternatif yang
berskala, sedangkan bukti induksi tidak peduli pada besarnya $n$.

\section{Lab komputasi}

\tingkat{lanjut}

Lab ini menjalankan dua hal: berkas Lean yang sama dengan yang dibahas di atas, dan pemeriksaan
silang dengan Python. Berkas Lean ditulis tanpa \texttt{import}, sehingga hanya memakai inti.

\begin{lab}
\begin{verbatim}
-- BilanganGanjil.lean : Lean 4.10.0, tanpa Mathlib
def oddSum : Nat -> Nat
  | 0     => 0
  | n + 1 => oddSum n + (2 * n + 1)

theorem oddSum_sq (n : Nat) : oddSum n = n * n := by
  induction n with
  | zero => rfl
  | succ k ih =>
    rw [oddSum]
    rw [ih]
    simp [Nat.succ_eq_add_one, Nat.mul_succ, Nat.add_mul, Nat.one_mul]
    omega

#eval oddSum 0
#eval oddSum 5
#eval oddSum 10
#eval oddSum 100
#print axioms oddSum_sq
\end{verbatim}

\begin{verbatim}
def odd_sum(n):
    return sum(2 * i + 1 for i in range(n))

print([odd_sum(n) for n in range(11)])
print(odd_sum(20), odd_sum(100))
print(all(odd_sum(n) == n * n for n in range(5000)))
print(all((k + 1) * (k + 1) == k * k + 2 * k + 1 for k in range(100000)))
\end{verbatim}
\end{lab}

Perintah dan keluaran, apa adanya:

\begin{verbatim}
$ lean --version
Lean (version 4.10.0, x86_64-unknown-linux-gnu, commit c375e19f6b65, Release)
$ time lean BilanganGanjil.lean
0
25
100
10000
'oddSum_sq' depends on axioms: [propext, Quot.sound]
waktu 0.14 s
$ python3 BilanganGanjil.py
[0, 1, 4, 9, 16, 25, 36, 49, 64, 81, 100]
400 10000
True
True
\end{verbatim}

Empat baris pertama keluaran Lean adalah hasil \texttt{\#eval}, dan keempatnya cocok dengan tabel
di bagian sebelumnya. Baris kelima adalah daftar aksioma, dan isinya perlu dibaca dengan hati-hati:
\texttt{propext} dan \texttt{Quot.sound} bukan kesalahan, melainkan akibat dari cara
\texttt{omega} bekerja. Keluaran Python menguatkan hal yang sama dari arah lain: daftar sebelas
nilai pertama adalah kuadrat sempurna, $S(100)=10000$, dan kedua pemeriksaan menyeluruh mencetak
\texttt{True}. Waktu 0,14 detik adalah waktu mesin uji ini, dan angka semacam itu bergantung pada
perangkat keras; yang bermakna bukan angkanya, melainkan faktanya bahwa pemeriksaan selesai
sebelum kita sempat membaca berkasnya.

Pembaca yang ingin mengulang lab ini memerlukan dua hal. Pertama, Lean 4 yang dapat dipasang
lewat \texttt{elan}, yaitu pengelola versi resmi proyek Lean; setelah pemasangan, perintah
\texttt{lean --version} harus mencetak versi. Kedua, berkas di atas disimpan sebagai
\texttt{BilanganGanjil.lean} dan dijalankan dengan \texttt{lean BilanganGanjil.lean} dari direktori
yang sama; tidak perlu proyek \texttt{lake} karena tidak ada impor. Untuk memakai \texttt{linarith}
atau \texttt{ring}, pembaca perlu proyek dengan Mathlib, dan sejak saat itu \texttt{lean} saja
tidak cukup; perintahnya menjadi \texttt{lake env lean}. Perbedaan ini sering menjadi sumber
kebingungan pertama bagi pemula, dan menyebutkannya lebih baik daripada membiarkannya ditemukan
lewat pesan \texttt{unknown tactic}.

\begin{resep}
\begin{enumerate}
  \item Buat berkas teks biasa bernama \texttt{Uji.lean}; ekstensi itulah yang membuat perintah
        \texttt{lean} mengenalinya sebagai berkas sumber.
  \item Tulis semua \texttt{import} di baris paling atas, satu per baris, sebelum deklarasi apa pun.
  \item Tulis pernyataan lebih dulu, lengkap dengan tipe setiap variabelnya, lalu bukti sesudah
        \texttt{:= by}.
  \item Jalankan \texttt{lean Uji.lean} dari direktori berkas itu dan baca keluaran dari baris
        pertama; kesalahan pertama biasanya yang sebenarnya.
  \item Tanyakan \texttt{\#print axioms nama\_teorema} sebelum menyebut hasilnya terbukti.
\end{enumerate}
Peringatan urutan: \texttt{import} tidak boleh didahului deklarasi. Satu \texttt{def} yang
diselipkan di atasnya membuat Lean menolak seluruh berkas dengan pesan
\texttt{invalid 'import' command}, dan berkas itu tidak diperiksa sama sekali. Urutan yang benar
bukan selera penulis; ia syarat agar berkas dapat diperiksa.
\end{resep}

\begin{untukdosen}
Urutan mengajar yang berhasil di kelas: mulai dari kegagalan \texttt{rfl} pada $n+1=1+n$, bukan
dari definisi. Mahasiswa yang melihat bahwa hal yang ``jelas'' ditolak akan bertanya, dan
pertanyaan itu adalah pintu masuk ke Curry--Howard. Setelah itu baru definisi proposisi sebagai
tipe, lalu tiga taktik pertama, lalu induksi. Sisakan satu sesi penuh untuk \texttt{\#print
axioms}; di situlah mahasiswa berubah dari pengguna menjadi pemeriksa.

Dua miskonsepsi yang paling sering muncul. Pertama, ``Lean membuktikan bahwa program saya benar''.
Yang benar: Lean memeriksa bahwa term bertipe pernyataan tertentu. Bila pernyataannya tidak
menyebut hal yang penting, misalnya kuantor atas gangguan, sertifikatnya kosong. Kedua,
``\texttt{simp} adalah bukti''. \texttt{simp} adalah prosedur yang keluarannya bergantung pada
kumpulan lemma aktif, sehingga bukti yang bergantung padanya dapat berhenti bekerja ketika pustaka
berubah; karena itu latih mahasiswa menulis daftar lemma secara eksplisit.

Aktivitas papan tulis yang berhasil: minta dua kelompok menulis spesifikasi untuk ``pengklasifikasi
yang tahan gangguan''. Satu kelompok hampir selalu menulis kuantor eksistensial. Tempelkan kedua
rumusan di papan, lalu buktikan yang lemah untuk fungsi konstan. Kelas akan melihat sendiri bahwa
kalimat yang terbukti dan kalimat yang dimaksud adalah dua kalimat berbeda.

Tanda kesiapan lanjut: mahasiswa dapat menjelaskan mengapa $n+1=1+n$ memerlukan lemma, dapat
menulis bukti induksi tanpa menyalin, dan dapat menyebut satu aksioma yang dipakai teoremanya
sendiri.
\end{untukdosen}

\section{Ringkasan}

\begin{ringkasan}
\begin{itemize}
  \item Bukti informal meyakinkan pembaca terlatih; bukti formal adalah term yang tipenya adalah
        pernyataannya, diperiksa inti yang kecil dan tidak tahu matematika.
  \item Menurut Curry--Howard, implikasi adalah tipe fungsi, konjungsi adalah pasangan, kuantor
        universal adalah fungsi dependen, dan $\bot$ tidak berpenghuni.
  \item Taktik mengubah keadaan tujuan; \texttt{intro}, \texttt{exact}, \texttt{apply},
        \texttt{rw}, \texttt{cases}, dan \texttt{induction} cukup untuk aritmetika dasar.
  \item Induksi membuktikan pernyataan untuk semua $n$ dengan dua tujuan, bukan dengan memeriksa
        semua $n$; \texttt{decide} sebaliknya berhenti berskala sangat cepat.
  \item \texttt{\#print axioms} adalah alat pemeriksaan utama. Taktik \texttt{native\_decide}
        meninggalkan aksioma \texttt{Lean.ofReduceBool}. Taktik \texttt{sorry} meninggalkan
        \texttt{sorryAx}. Aksioma buatan sendiri meninggalkan namanya sendiri.
  \item Verifikasi menjamin term cocok dengan pernyataan, bukan bahwa pernyataan cocok dengan
        maksud; spesifikasi lemah menghasilkan sertifikat yang kosong.
\end{itemize}
\end{ringkasan}

\section{Latihan}

\begin{latihan}
\begin{enumerate}
  \item Nilai klaim ini: \emph{bila sebuah berkas Lean diterima tanpa pesan kesalahan, maka semua
        pernyataan di dalamnya benar}. Tentukan benar, salah, atau belum cukup bukti, dan sebutkan
        satu perintah yang dapat dipakai untuk memeriksa klaim itu pada berkas tertentu.

  \item Buktikan di Lean, tanpa Mathlib, bahwa untuk setiap $n$ berlaku $2\cdot S(n)=n\cdot(2n)$
        dengan $S$ seperti di bab ini. Nyatakan lebih dulu lemma bantu yang kamu perlukan.

  \item Buktikan bahwa jumlah $n$ bilangan genap pertama, yaitu $0+2+\dots+(2n-2)$, sama dengan
        $n(n-1)$. Tulis bukti informalnya, lalu tulis versi Lean-nya dan sebutkan taktik yang
        dipakai pada setiap langkah.

  \item Hitung dengan Python: jumlah 37 bilangan ganjil pertama, jumlah 250 bilangan ganjil
        pertama, dan periksa bahwa keduanya kuadrat sempurna. Sertakan kode dan keluarannya.

  \item Hitung dengan Python banyaknya kasus yang harus diperiksa bila seseorang ingin membuktikan
        $S(n)=n^{2}$ dengan pemeriksaan langsung untuk semua $n\le 1000$, lalu bandingkan dengan
        banyaknya tujuan yang dibangkitkan bukti induksi.

  \item Jelaskan mengapa \texttt{rfl} menutup tujuan $n+0=n$ tetapi tidak menutup $0+n=n$. Petunjuk
        ada pada cara penjumlahan bilangan asli didefinisikan secara rekursif; sebutkan argumen
        mana yang disusutkan lebih dulu.

  \item Tulis spesifikasi lemah dan spesifikasi kuat untuk klaim ``pengklasifikasi ini tidak
        berubah pada gangguan berukuran paling banyak satu''. Buktikan versi lemahnya di Lean
        untuk fungsi apa pun, lalu jelaskan mengapa versi kuatnya tidak dapat dibuktikan untuk
        fungsi identitas.

  \item Periksa dengan \texttt{\#print axioms} tiga teorema yang kamu tulis sendiri: satu memakai
        \texttt{decide}, satu memakai \texttt{omega}, satu memakai \texttt{sorry}. Laporkan
        ketiga daftar aksiomanya dan terangkan asal setiap nama.

  \item Sebuah sertifikat memuat \texttt{theorem aman : exists delta : Nat, f (x + delta) = f x}
        dengan \texttt{delta} sebagai variabel bebas di dalam konteks. Jelaskan dua kesalahan
        yang mungkin ada pada pernyataan itu, dan tulis versi yang kamu anggap benar.

  \item Buktikan di Lean bahwa $\neg(P\wedge\neg P)$ untuk setiap proposisi $P$. Petunjuk: mulai
        dengan \texttt{intro h}, lalu ambil komponennya dengan \texttt{h.1} dan \texttt{h.2}, dan
        gunakan komponen kedua pada komponen pertama.
\end{enumerate}

Petunjuk untuk tiga butir tersulit. Butir 2: lemma bantu yang diperlukan berbentuk
$2\,(k\cdot k+(2k+1))=(k+1)\cdot(2(k+1))$, dan setelah \texttt{rw} pada definisi serta hipotesis
induksi, sisa tujuannya linear sehingga \texttt{omega} cukup. Butir 7: pada versi lemah, saksi
$\delta=0$ membuat kedua sisi sama secara definisi, jadi \texttt{rfl} langsung menutup; pada versi
kuat, fungsi identitas memberi penyangkal dengan $\delta=1$. Butir 9: perhatikan bahwa
\texttt{delta} yang bebas dapat dipilih oleh pembaca pernyataan, bukan oleh penyerang, sehingga
pernyataan itu tidak menyatakan ketahanan apa pun; kesalahan kedua adalah tidak adanya kuantor
atas $x$.
\end{latihan}

\begin{untukpeneliti}
Yang benar-benar terbukti di bab ini hanyalah pernyataan aritmetika kecil: $S(n)=n^{2}$ untuk
semua bilangan asli, beberapa tautologi proposisional, dan dua pernyataan tentang spesifikasi yang
terlalu lemah. Semuanya diperiksa inti Lean 4.10.0 tanpa Mathlib, dan daftar aksiomanya
dilaporkan. Yang dipopulerkan jauh melampaui itu: frasa ``terverifikasi secara formal'' sering
dipakai untuk ``diperiksa alat otomatis'' atau bahkan ``diuji menyeluruh''. Hanya yang pertama
menghasilkan objek yang dapat diperiksa ulang pihak ketiga.

Sumber primer. Avigad, de Moura, Kong, dan Ullrich, \emph{Theorem Proving in Lean 4}, bagian
``Propositions and Proofs'' (proposisi sebagai tipe, irrelevansi bukti, peringatan aksioma),
``Tactics'' (keadaan tujuan), dan ``Induction and Recursion'' (\texttt{induction},
\texttt{sorryAx}). Avigad dan Massot, \emph{Mathematics in Lean}, Bab 2, 3.1, 5.2, dan 8; edisi
bahasa Indonesia memakai penomoran bagian yang sama.

Pertanyaan terbuka. Bab ini tidak memuat bukti formal tentang jaringan saraf. Spesifikasi ketahanan yang menarik
berbentuk
$\forall x\,\forall\delta,\ \lVert\delta\rVert\le\varepsilon\to f(x+\delta)=f(x)$, yaitu
pernyataan atas ruang kontinu dengan fungsi berparameter ratusan ribu. Memformalkannya menuntut
aritmetika real yang dapat diputuskan pada skala itu, dan pustaka Lean sekarang tidak menyediakannya
untuk model terlatih. Klaim ``kami memverifikasi ketahanan model'' di makalah
biasanya merujuk verifikasi dengan abstraksi interval atau SMT, bukan bukti Lean; keduanya sah,
tetapi jaminannya berbeda dan tidak boleh disebut dengan kata yang sama.

Batas keberlakuan. Semua klaim di bab ini berlaku untuk Lean 4.10.0 dengan inti bawaannya, dan
rincian keluarannya dapat berubah antar versi. Yang tidak berubah adalah struktur argumennya:
pernyataan, term, inti, daftar aksioma.
\end{untukpeneliti}

\peta{Bab ini mengikuti Avigad, de Moura, Kong, dan Ullrich, \emph{Theorem Proving in Lean 4},
bagian ``Propositions and Proofs'' (proposisi sebagai tipe, Curry--Howard, irrelevansi bukti,
peringatan aksioma), ``Quantifiers and Equality'', ``Tactics'' (keadaan tujuan, \texttt{intro},
\texttt{apply}, \texttt{exact}), dan ``Induction and Recursion'' (\texttt{induction},
\texttt{sorry} dan \texttt{sorryAx}). Padanan dalam Avigad dan Massot, \emph{Mathematics in Lean},
ada pada Bab 2 (Dasar-dasar, khususnya 2.3 Menggunakan Teorema dan Lema serta 2.4 Lebih Banyak
Contoh Menggunakan apply dan rw), Bab 3.1 (Implikasi dan Kuantor Universal), Bab 5.2 (Induksi dan
Rekursi), dan Bab 8 (Hierarki) untuk typeclass; edisi bahasa Indonesia mempertahankan penomoran
bagian yang sama. Untuk kaitan dengan praktik evaluasi, lihat Bab 27 buku ini; untuk serangan yang
membuat spesifikasi ketahanan menjadi tidak trivial, lihat Bab 27 dan Bab 28.}

\chapter[Adversarial Machine Learning]{Adversarial Machine Learning: Serangan dan Pertahanan}

Pada 2014 Goodfellow, Shlens, dan Szegedy menunjukkan satu gambar panda. Jaringan
GoogLeNet menyebutnya panda dengan keyakinan $57{,}7$ persen. Mereka menambahkan
gangguan yang tidak terlihat mata, sebesar $0{,}007$ pada skala $[0,1]$, yaitu sekitar
dua satuan dari 255 tingkat keabuan. Jaringan yang sama lalu menyebut gambar itu gibon
dengan keyakinan $99{,}3$ persen. Tidak ada yang berubah bagi manusia. Yang berubah
hanya arah gangguan: setiap piksel digeser sedikit, dan arah geseran itu dipilih dari
tanda turunan galat terhadap piksel.

Angka seperti itu bukan keanehan satu model. Pada bab ini kita menghitung sendiri
kasus kecil yang bisa diperiksa ulang. Sebuah jaringan dua lapis dengan 64 masukan
mencapai akurasi $97{,}99$ persen pada 597 citra angka uji. Dengan anggaran gangguan
$0{,}15$ per piksel, serangan berulang 40 langkah menurunkan akurasi itu menjadi
$10{,}72$ persen; pada anggaran yang sama, gangguan acak tanpa arah hanya
menurunkannya menjadi $95{,}48$ persen. Jadi arah yang dipilih gradien bukan hiasan:
ia membeli sekitar delapan puluh poin akurasi bagi penyerang.

Bab 25 mengajari kita menghitung turunan galat terhadap bobot, dan Bab 27 mengajari
kita mengukur akurasi dengan jujur. Keduanya belum cukup. Akurasi adalah rata-rata
atas satu distribusi data; keamanan adalah pertanyaan tentang titik terburuk di dalam
satu himpunan yang diizinkan. Rata-rata yang tinggi dan titik terburuk yang buruk bisa
hidup bersama, dan itulah yang membuat laporan evaluasi biasa menyesatkan. Bab ini
membangun kosakata untuk pertanyaan terburuk itu, menghitung serangan pertamanya, lalu
memisahkan pertahanan yang menaikkan standar dari yang hanya menyembunyikan gradien.

\begin{untukanak}
Bayangkan seorang penjaga gedung yang mengenali orang dari kebiasaan kecil: berapa
tombol baju yang dikancing, tas di tangan kiri atau kanan, dan topi atau tidak.
Selama ini cara itu berhasil. Semua orang yang berhak masuk memang berbeda pada tiga
hal tersebut. Lalu datang seseorang yang tahu persis cara penjaga itu melihat. Ia
membuka satu tombol, memindahkan tas ke tangan yang lain, dan memakai topi. Wajahnya
sama, namanya sama, ia orang yang sama. Penjaga tetap mempersilakannya masuk, karena
menurut catatannya orang dengan tiga tanda itu adalah petugas kebersihan.

Pertanyaan yang membuat kisah ini penting bukan "bagaimana menipu penjaga", melainkan
"bagaimana kita tahu penjaga itu bisa ditipu". Kalau kita hanya menguji dengan orang
acak yang berpakaian acak, penjaga akan lulus hampir selalu. Kesalahannya baru muncul
kalau ada orang yang sengaja memilih tiga tanda itu. Menguji dengan kebetulan tidak
sama dengan menguji dengan niat jahat.

Komputer pengenal gambar bekerja seperti penjaga itu, hanya dengan ribuan ciri kecil
sekaligus. Di bab ini kita akan tahu mengapa ciri yang tampak sepele bisa dipilih
dengan sengaja, mengapa perubahan yang tidak terlihat mata bisa membalik jawaban,
dan bagaimana membedakan pertahanan yang sungguh kuat dari pertahanan yang hanya
membuat penyerang bingung.
\end{untukanak}

\section{Model ancaman: apa yang boleh dilakukan penyerang}
\tingkat{dasar}

Kalimat "penyerang mengubah masukan sedikit" belum bisa diuji sebelum kita sepakat
tentang arti "sedikit". Setiap klaim ketahanan selalu relatif terhadap pilihan itu,
sehingga pilihan tersebut adalah bagian dari klaim, bukan catatan kaki.

\kamus{model ancaman}{threat model}{Daftar lengkap tentang apa yang boleh dilakukan
penyerang: seberapa besar gangguannya, apa yang ia ketahui, dan apa tujuannya.}
\index{model ancaman}

\begin{definisi}[Serangan adversarial dan model ancaman]
Misalkan $f:\mathbb{R}^{d}\to\{1,\dots,K\}$ adalah pengklasifikasi dan
$(\mathbf{x},y)$ adalah contoh yang diklasifikasikan benar. Sebuah
\emph{contoh adversarial} adalah $\mathbf{x}'=\mathbf{x}+\boldsymbol{\delta}$ dengan
$f(\mathbf{x}')\neq y$ tetapi $\boldsymbol{\delta}$ kecil menurut ukuran yang
disepakati. Model ancaman ditentukan oleh empat hal:
\begin{enumerate}
  \item himpunan gangguan yang diizinkan
        $S(\mathbf{x})=\{\mathbf{x}+\boldsymbol{\delta}:
        \lVert\boldsymbol{\delta}\rVert\le\varepsilon\}$, dengan
        $\lVert\cdot\rVert$ biasanya norma $\ell_\infty$ dan $\varepsilon>0$
        adalah \emph{anggaran};
  \item pengetahuan penyerang: akses penuh pada $f$ dan gradiennya
        (\emph{kotak putih}) atau hanya pada keluaran $f$ (\emph{kotak hitam});
  \item tujuan: salah klasifikasi apa saja (\emph{tak bersasaran}) atau kelas
        tertentu (\emph{bersasaran});
  \item batasan lain, misalnya gangguan harus tetap berupa objek yang sah.
\end{enumerate}
Serangan berhasil bila ada $\boldsymbol{\delta}\in S(\mathbf{x})$ yang mengubah label.
\end{definisi}

Norma $\ell_\infty$ dipilih karena ia membatasi setiap koordinat secara terpisah:
$\lVert\boldsymbol{\delta}\rVert_\infty\le\varepsilon$ berarti setiap piksel bergeser
paling jauh $\varepsilon$. Untuk citra 8-bit dengan nilai $0$ sampai $255$, anggaran
$\varepsilon=1/255$ berarti setiap piksel boleh berubah paling banyak satu tingkat.
Norma $\ell_2$ membatasi panjang total vektor dan membolehkan satu piksel berubah
banyak asalkan yang lain hampir tidak berubah. Keduanya menjawab pertanyaan berbeda,
dan klaim ketahanan yang tidak menyebut normanya tidak bisa dibandingkan dengan klaim
lain.
\index{norma tak hingga}
\index{anggaran perturbasi}

Ukuran ketahanan yang paling jujur adalah jarak ke titik salah klasifikasi terdekat,
\[
  r(\mathbf{x})=\inf\{\lVert\boldsymbol{\delta}\rVert:
  f(\mathbf{x}+\boldsymbol{\delta})\neq y\}.
\]
Besaran ini tidak bergantung pada algoritme penyerang mana pun. Sayangnya $r$ sulit
dihitung: mencarinya berarti memaksimalkan fungsi tak konveks pada himpunan yang
besar. Untuk jaringan dengan unit ReLU, $r$ dapat dihitung eksak dengan mengubah
persoalan menjadi program linear campuran-bulat, dan biayanya tumbuh cepat terhadap
jumlah unit. Karena itu hampir semua laporan di bidang ini menyebut apa yang berhasil
dicapai oleh penyerang tertentu, bukan nilai $r$. Selisih antara keduanya adalah
sumber sebagian besar kesimpulan yang terlalu optimistis.

Satu kejadian layak dicatat sejak awal. Pada data tabular, bola $\ell_\infty$ memuat
titik yang mustahil: usia negatif, jumlah paket yang tidak bulat, kombinasi bendera
TCP yang tidak pernah muncul di jaringan nyata. Contoh seperti itu menipu model tetapi
tidak dapat diwujudkan, sehingga keberhasilannya tidak memberi tahu apa pun tentang
risiko sebenarnya. Anggaran yang terlalu longgar menghasilkan serangan yang
spektakuler dan tidak berarti.

\section{Mengapa tanda gradien cukup}
\tingkat{menengah}

Sekarang kita serang. Misalkan galat model dinyatakan oleh fungsi
$L(\mathbf{x},y)$, misalnya entropi silang. Kita ingin menaikkan galat itu sebesar
mungkin sambil menjaga setiap koordinat tetap di dalam kotak. Untuk fungsi linear,
jawabannya bisa ditulis rapat.

\kamus{metode tanda gradien}{fast gradient sign method}{Serangan satu langkah yang
menggeser setiap koordinat sebesar anggaran, dengan arah mengikuti tanda turunan
galat.}
\index{FGSM}

\begin{lemma}[Maksimum bentuk linear pada bola $\ell_\infty$]
Untuk $\mathbf{w}\in\mathbb{R}^{d}$ dan $\varepsilon>0$,
\[
  \max_{\lVert\boldsymbol{\delta}\rVert_\infty\le\varepsilon}
  \mathbf{w}^{\top}\boldsymbol{\delta}
  =\varepsilon\lVert\mathbf{w}\rVert_1,
\]
dan maksimum tercapai pada $\delta_j=\varepsilon\,\mathrm{sign}(w_j)$ untuk setiap
koordinat dengan $w_j\neq0$.
\end{lemma}

\begin{proof}
Tulis $\mathbf{w}^{\top}\boldsymbol{\delta}=\sum_j w_j\delta_j$. Karena
$|\delta_j|\le\varepsilon$,
\[
  \sum_j w_j\delta_j\le\sum_j|w_j|\,|\delta_j|\le\varepsilon\sum_j|w_j|.
\]
Bila $\delta_j=\varepsilon\,\mathrm{sign}(w_j)$, maka $w_j\delta_j=\varepsilon|w_j|$
untuk setiap $j$, sehingga batas itu tercapai. Koordinat dengan $w_j=0$ tidak
memengaruhi nilai, jadi pilihannya bebas.
\end{proof}

Untuk fungsi mulus, kita memakai hampiran orde pertama $L$ di sekitar $\mathbf{x}$
bersama lemma tadi. Hasilnya resep yang oleh Goodfellow dan rekan disebut
\emph{fast gradient sign method}:
\[
  \mathbf{x}^{\mathrm{adv}}
  =\mathbf{x}+\varepsilon\,\mathrm{sign}\bigl(\nabla_{\mathbf{x}}L(\mathbf{x},y)\bigr).
\]
Kalimat sehari-harinya begini. Gradien memberi tahu ke arah mana galat naik paling
cepat, tetapi kita hanya boleh berjalan sejauh $\varepsilon$; yang dipakai hanya tanda
tiap komponennya, sedangkan besarnya ditetapkan anggaran. Karena itu metode ini disebut
serangan tanda. Untuk fungsi linear, langkah itu tepat optimal; Gambar~\ref{fig:batas}
memperlihatkan hasilnya pada satu model dua ciri.

\begin{figure}[htbp]
\centering
{\scriptsize
\begin{tikzpicture}
\node[at={(0.000cm,0.000cm)},inner sep=0pt,minimum size=0pt] (o) {};
\node[at={(8.400cm,0.000cm)},inner sep=0pt,minimum size=0pt] (ex) {};
\node[at={(0.000cm,8.400cm)},inner sep=0pt,minimum size=0pt] (ey) {};
\draw[-latex,black!75] (o) -- (ex);
\draw[-latex,black!75] (o) -- (ey);
\node[anchor=north,at={(4.200cm,-1.120cm)}]{tinta paruh atas};
\node[rotate=90,anchor=south,at={(-1.15cm,5.10cm)}]{tinta paruh bawah};
\node[anchor=north,at={(2.380cm,-0.280cm)}]{0,20};
\node[anchor=north,at={(6.580cm,-0.280cm)}]{0,35};
\node[anchor=east,at={(-0.224cm,2.660cm)}]{0,35};
\node[anchor=east,at={(-0.224cm,6.860cm)}]{0,50};
\node[at={(1.439cm,0.000cm)},inner sep=0pt,minimum size=0pt] (b0) {};
\node[at={(8.400cm,7.224cm)},inner sep=0pt,minimum size=0pt] (b1) {};
\draw[black,line width=1.2pt] (b0) -- (b1);
\node[circle,fill=black!55,inner sep=1.35pt,at={(2.358cm,4.399cm)}]{};
\node[circle,fill=black!55,inner sep=1.35pt,at={(3.014cm,5.602cm)}]{};
\node[circle,fill=black!55,inner sep=1.35pt,at={(2.796cm,2.376cm)}]{};
\node[circle,fill=black!55,inner sep=1.35pt,at={(3.178cm,5.274cm)}]{};
\node[circle,fill=black!55,inner sep=1.35pt,at={(0.444cm,3.524cm)}]{};
\node[circle,fill=black!55,inner sep=1.35pt,at={(1.975cm,4.290cm)}]{};
\node[circle,draw=black!55,line width=0.5pt,inner sep=1.15pt,at={(5.585cm,2.649cm)}]{};
\node[circle,draw=black!55,line width=0.5pt,inner sep=1.15pt,at={(5.749cm,2.868cm)}]{};
\node[circle,draw=black!55,line width=0.5pt,inner sep=1.15pt,at={(5.147cm,1.391cm)}]{};
\node[circle,draw=black!55,line width=0.5pt,inner sep=1.15pt,at={(3.780cm,0.844cm)}]{};
\node[circle,draw=black!55,line width=0.5pt,inner sep=1.15pt,at={(6.678cm,3.032cm)}]{};
\node[circle,draw=black!55,line width=0.5pt,inner sep=1.15pt,at={(5.803cm,2.868cm)}]{};
\node[circle,draw=black,line width=1pt,inner sep=2.8pt,at={(2.303cm,7.079cm)}]{};
\node[circle,fill=black,inner sep=0.9pt,at={(2.303cm,7.079cm)}]{};
\node[rectangle,fill=accent,inner sep=2.3pt,at={(5.663cm,3.719cm)}]{};
\node[at={(2.303cm,7.079cm)},inner sep=0pt,minimum size=0pt] (x0) {};
\node[at={(5.663cm,3.719cm)},inner sep=0pt,minimum size=0pt] (xa) {};
\draw[-latex,accent,line width=1.1pt] (x0) -- (xa);
\node[anchor=south east,at={(2.135cm,7.443cm)}]{bersih};
\node[anchor=north west,at={(5.943cm,3.551cm)}]{adversarial};
\node[anchor=east,at={(2.660cm,5.600cm)}]{kelas 6};
\node[anchor=south west,at={(5.684cm,1.260cm)}]{kelas 0};
\node[anchor=south east,at={(5.475cm,4.647cm)}]{batas};
\end{tikzpicture}}
\caption{Batas keputusan model dua ciri pada angka $0$ lawan $6$ ($236$ data
latih, $123$ uji, akurasi $98{,}37$ persen). Lingkaran tebal: contoh bersih
$(0{,}197;0{,}508)$ dengan skor $+1{,}936$. Kotak merah: hasil FGSM
$\varepsilon=0{,}12$ di $(0{,}317;0{,}388)$ dengan skor $-0{,}208$. Panahnya
gangguan $(0{,}12;-0{,}12)$, $\lVert\boldsymbol{\delta}\rVert_\infty=0{,}12$
dan $\lVert\boldsymbol{\delta}\rVert_2=0{,}1697$.}
\label{fig:batas}
\end{figure}

\begin{teorema}[Linearisasi dan sisanya]
Misalkan $L:\mathbb{R}^{d}\to\mathbb{R}$ terturunkan dan gradiennya $M$-Lipschitz
pada segmen antara $\mathbf{x}$ dan $\mathbf{x}+\boldsymbol{\delta}$. Untuk setiap
$\lVert\boldsymbol{\delta}\rVert_\infty\le\varepsilon$,
\[
  \left|L(\mathbf{x}+\boldsymbol{\delta})-L(\mathbf{x})
  -\nabla L(\mathbf{x})^{\top}\boldsymbol{\delta}\right|
  \le\frac{M}{2}\lVert\boldsymbol{\delta}\rVert_2^2
  \le\frac{Md}{2}\varepsilon^2 .
\]
\end{teorema}

\begin{proof}
Dengan bentuk integral dari teorema nilai rata-rata,
\[
  L(\mathbf{x}+\boldsymbol{\delta})-L(\mathbf{x})
  =\int_0^1\nabla L(\mathbf{x}+t\boldsymbol{\delta})^{\top}\boldsymbol{\delta}\,dt .
\]
Kurangkan $\nabla L(\mathbf{x})^{\top}\boldsymbol{\delta}
=\int_0^1\nabla L(\mathbf{x})^{\top}\boldsymbol{\delta}\,dt$, lalu gunakan
ketaksamaan Cauchy--Schwarz dan sifat Lipschitz:
\[
  \left|\int_0^1\bigl(\nabla L(\mathbf{x}+t\boldsymbol{\delta})
  -\nabla L(\mathbf{x})\bigr)^{\top}\boldsymbol{\delta}\,dt\right|
  \le\int_0^1 Mt\lVert\boldsymbol{\delta}\rVert_2^2\,dt
  =\frac{M}{2}\lVert\boldsymbol{\delta}\rVert_2^2 .
\]
Karena $\lVert\boldsymbol{\delta}\rVert_2^2\le d\,\varepsilon^2$, ketaksamaan kedua
mengikuti.
\end{proof}

Untuk model linear, sisa itu nol dan FGSM bukan hampiran melainkan jawaban eksak.
Goodfellow dan rekan memakai kasus regresi logistik untuk menunjukkan hal ini. Di
situ tanda gradien galat terhadap masukan adalah $-\mathrm{sign}(\mathbf{w})$ dikali
label, sehingga
$\mathbf{w}^{\top}\mathrm{sign}(\mathbf{w})=\lVert\mathbf{w}\rVert_1$ dan margin
model menyusut tepat sebesar $\varepsilon\lVert\mathbf{w}\rVert_1$. Karena itu
pelatihan adversarial pada regresi logistik bisa ditulis sebagai penalti
$\zeta\bigl(y(\varepsilon\lVert\mathbf{w}\rVert_1-\mathbf{w}^{\top}\mathbf{x}-b)\bigr)$,
yaitu penalti mirip regularisasi $\ell_1$ yang masuk ke margin, bukan ke biaya;
bentuk eksak ini dipakai di contoh hitung nanti.

Yang perlu dinyatakan jujur: teorema di atas membenarkan \emph{arah} langkah, bukan
\emph{jarak} langkah. Sisa $Md\varepsilon^{2}/2$ bergantung pada konstanta Lipschitz
gradien, dan untuk jaringan ReLU konstanta itu bisa sangat besar. Pada percobaan kita,
prediksi kenaikan galat dari suku linear untuk $\varepsilon=0{,}15$ adalah
$0{,}0008$ nats, sedangkan kenaikan yang sungguh terjadi $4{,}862$ nats. Jadi
linearisasi tidak menjelaskan besarnya kerusakan; ia hanya menjelaskan mengapa arah
yang dipilih masuk akal.

Sisi lain cerita justru menguat. Gangguan dengan besar yang sama tetapi tanda acak
hanya menurunkan akurasi dari $97{,}99$ menjadi $95{,}48$ persen, sedangkan FGSM
menurunkannya menjadi $14{,}91$ persen. Jadi yang berharga adalah pola tandanya,
walaupun galat di titik awal hampir datar sehingga panjang gradien dan ramalan
linearnya kecil.
\index{gradien masukan}

Jadi tanda gradien cukup bukan karena suku linearnya besar, melainkan karena tanda
lebih tahan terhadap penjenuhan galat. Angka $0{,}0008$ melawan $4{,}862$
memperingatkan bahwa linearitas hanya berlaku sebagai argumen arah.

\begin{bacakeras}
Dua suara, dibaca bergantian dengan teman. Ketuk meja setiap kali kata ``sedikit''
diucapkan.

``Geser sedikit.'' ``Berapa sedikit?'' ``Seujung jari. Tidak lebih.''
``Jawabannya berubah?'' ``Berubah seluruhnya. Bukan ragu, bukan meleset: berubah.''
``Mata kita melihat apa?'' ``Tidak ada. Gambar yang sama persis.''
``Kalau digeser sejauh itu ke arah lain?'' ``Jawabannya kembali benar.''
``Jadi yang bekerja apa?'' ``Arahnya. Bukan besarnya.''
``Dari mana ia tahu arah?'' ``Ia membaca lereng, bertanya pada tiap piksel: naik atau
turun? Lalu semua jawaban dipakai sekaligus.''
``Kalau lerengnya salah baca?'' ``Tidak terjadi apa-apa. Gambar tetap, jawaban tetap.''

Ulangi tiga kali, makin cepat: yang bekerja arahnya, bukan besarnya.
\end{bacakeras}

\section{Serangan berulang dan permainan min-max}
\tingkat{menengah}

Satu langkah adalah kompromi. Bila gradien boleh dihitung berkali-kali, kita bisa
mendaki lebih jauh: setiap langkah kecil ke arah tanda gradien, lalu tarik kembali
ke dalam kotak.

\begin{definisi}[Serangan iteratif]
Misalkan $\Pi$ adalah proyeksi ke $S(\mathbf{x})$. Serangan \emph{projected gradient
descent} (PGD) menghitung
\[
  \mathbf{x}_{t+1}=\Pi\Bigl(\mathbf{x}_t+\alpha\,
  \mathrm{sign}\bigl(\nabla_{\mathbf{x}}L(\mathbf{x}_t,y)\bigr)\Bigr),
  \qquad \mathbf{x}_0=\mathbf{x},
\]
dengan $\alpha>0$ adalah ukuran langkah. Bila $\mathbf{x}_0$ diambil acak di dalam
$S(\mathbf{x})$, serangan itu disebut PGD dengan \emph{mulai acak}.
\end{definisi}

\begin{resep}
Empat langkah menjalankan serangan iteratif terbatas:
\begin{enumerate}
  \item Tetapkan norma, $\varepsilon$, $\alpha$, dan jumlah langkah; catat akurasi
        bersih sebelum satu pun gangguan dibuat.
  \item Mulai dari $\mathbf{x}$ atau dari titik acak di dalam bola, hitung gradien
        galat terhadap masukan, lalu ambil tandanya saja.
  \item Langkah: $\mathbf{x}_{t+1}=\mathbf{x}_t+\alpha\,\mathrm{sign}(\mathbf{g})$,
        proyeksikan ke bola $\varepsilon$, lalu ke $[0,1]$.
  \item Ulangi, lalu laporkan akurasi pada beberapa jumlah langkah dan beberapa
        $\varepsilon$; bandingkan dengan FGSM yang memakai $\alpha=\varepsilon$.
\end{enumerate}
Peringatan urutan: proyeksi selalu terhadap $\mathbf{x}$ yang asli, bukan terhadap
titik langkah sebelumnya, dan selalu sesudah langkah gradien. Bila acuannya tertukar,
bola ikut berjalan dan gangguan melebihi anggaran yang dilaporkan.
\end{resep}

Proyeksi pada kotak $\ell_\infty$ yang dipotong ke $[0,1]^d$ punya bentuk yang sangat
mudah: potong setiap koordinat, satu per satu.

\begin{proposisi}[Proyeksi ke kotak]
Untuk $S=\{\mathbf{z}:\lVert\mathbf{z}-\mathbf{x}\rVert_\infty\le\varepsilon\}
\cap[0,1]^{d}$,
\[
  \Pi(\mathbf{z})_j=\min\bigl\{\max\{z_j,\,x_j-\varepsilon,\,0\},\,
  x_j+\varepsilon,\,1\bigr\}
\]
untuk setiap $j$.
\end{proposisi}

\begin{proof}
Fungsi tujuan $\lVert\mathbf{y}-\mathbf{z}\rVert_2^2$ terpisah per koordinat, dan
himpunan yang diizinkan juga terpisah: setiap $y_j$ harus berada di selang
$[\max\{0,x_j-\varepsilon\},\min\{1,x_j+\varepsilon\}]$. Untuk satu koordinat,
meminimalkan $(y_j-z_j)^2$ pada selang tertutup berarti mengambil titik terdekat,
yaitu $z_j$ bila ia di dalam selang, dan ujung selang yang terdekat bila di luar.
Itu tepat rumus potong di atas.
\end{proof}

Dengan proyeksi semacam ini, PGD hanya menambah dua hal dibanding FGSM: langkah
kecil berulang dan tarikan kembali ke kotak. Kenaikan biayanya sebanding dengan
jumlah langkah, karena setiap langkah memerlukan satu lintasan mundur. Gambar~\ref{fig:langkah} memperlihatkan tiga
langkah seperti itu pada model yang sama.

\begin{figure}[htbp]
\centering
{\scriptsize
\begin{tikzpicture}
\node[at={(0.000cm,0.000cm)},inner sep=0pt,minimum size=0pt] (p) {};
\node[at={(8.400cm,0.000cm)},inner sep=0pt,minimum size=0pt] (px) {};
\node[at={(0.000cm,8.400cm)},inner sep=0pt,minimum size=0pt] (py) {};
\draw[-latex,black!75] (p) -- (px);
\draw[-latex,black!75] (p) -- (py);
\node[anchor=north,at={(4.200cm,-1.344cm)}]{tinta paruh atas};
\node[rotate=90,anchor=south,at={(-1.05cm,4.20cm)}]{tinta paruh bawah};
\node[anchor=north,at={(1.848cm,-0.336cm)}]{0,20};
\node[anchor=north,at={(6.888cm,-0.336cm)}]{0,35};
\node[at={(2.661cm,0.000cm)},inner sep=0pt,minimum size=0pt] (c0) {};
\node[at={(8.400cm,5.956cm)},inner sep=0pt,minimum size=0pt] (c1) {};
\draw[black,line width=1.2pt] (c0) -- (c1);
\node[at={(1.756cm,6.479cm)},inner sep=0pt,minimum size=0pt] (q0) {};
\node[circle,fill=white,draw=accent,line width=0.9pt,inner sep=2.0pt,at={(1.756cm,6.479cm)}]{};
\node[at={(3.100cm,5.135cm)},inner sep=0pt,minimum size=0pt] (q1) {};
\node[circle,fill=white,draw=accent,line width=0.9pt,inner sep=2.0pt,at={(3.100cm,5.135cm)}]{};
\node[at={(4.444cm,3.791cm)},inner sep=0pt,minimum size=0pt] (q2) {};
\node[circle,fill=white,draw=accent,line width=0.9pt,inner sep=2.0pt,at={(4.444cm,3.791cm)}]{};
\node[at={(5.788cm,2.447cm)},inner sep=0pt,minimum size=0pt] (q3) {};
\node[circle,fill=white,draw=accent,line width=0.9pt,inner sep=2.0pt,at={(5.788cm,2.447cm)}]{};
\draw[-latex,accent,line width=1.0pt,shorten >=5pt,shorten <=5pt] (q0) -- (q1);
\draw[-latex,accent,line width=1.0pt,shorten >=5pt,shorten <=5pt] (q1) -- (q2);
\draw[-latex,accent,line width=1.0pt,shorten >=5pt,shorten <=5pt] (q2) -- (q3);
\node[anchor=south west,at={(1.891cm,6.613cm)}]{$\mathbf{x}_0$};
\node[anchor=south west,at={(3.235cm,5.269cm)}]{$\mathbf{x}_1$};
\node[anchor=south west,at={(4.579cm,3.925cm)}]{$\mathbf{x}_2$};
\node[anchor=south west,at={(5.923cm,2.581cm)}]{$\mathbf{x}_3$};
\node[anchor=south east,at={(5.616cm,3.444cm)}]{batas};
\end{tikzpicture}}
\caption{Tiga langkah tanda gradien dari $\mathbf{x}_0$, tiap langkah
$\alpha=0{,}04$. Norma $\lVert\boldsymbol{\delta}\rVert_\infty$: $0{,}04$,
$0{,}08$, $0{,}12$; norma $\lVert\boldsymbol{\delta}\rVert_2$: $0{,}0566$,
$0{,}1131$, $0{,}1697$. Skor kelas $6$: $+1{,}936$, $+1{,}221$, $+0{,}506$,
$-0{,}208$; batas dilewati pada langkah ketiga.}
\label{fig:langkah}
\end{figure}

Ada satu peringatan yang sering terlewat. FGSM adalah PGD dengan satu langkah dan
$\alpha=\varepsilon$, bukan PGD dengan satu langkah dan $\alpha=\varepsilon/10$.
Pada percobaan kita, satu langkah dengan $\alpha=\varepsilon/10$ hanya menurunkan
akurasi menjadi $94{,}64$ persen, sedangkan FGSM menurunkannya menjadi $14{,}91$
persen. Jadi "satu langkah" bukan nama satu algoritme; ia baru berarti sesuatu
setelah ukuran langkahnya disebut. Perbandingan yang adil memakai ukuran langkah yang
sama; tabel berikut dihitung pada $\varepsilon=0{,}15$ dengan mulai acak.

\begin{center}
\begin{tabular}{rr}
\hline
langkah & akurasi (\%) \\
\hline
1  & $94{,}64$ \\
2  & $88{,}78$ \\
5  & $66{,}67$ \\
10 & $24{,}79$ \\
20 & $11{,}22$ \\
40 & $10{,}72$ \\
\hline
\end{tabular}
\end{center}

Kurva itu turun tajam sampai sepuluh langkah, lalu hampir mendatar; menambah langkah
setelah titik itu jarang mengubah kesimpulan, dan itulah sebabnya PGD dipakai sebagai
penyerang bawaan. Madry dan rekan membacanya sebagai tanda bahwa masalah dalam itu
cukup mudah diselesaikan dengan metode orde pertama, meskipun secara umum tak konkaf.

\begin{tangan}
Bahan: sepuluh kartu, spidol, dan penggaris. Tulis satu angka antara nol dan satu pada
tiap kartu, dua angka di belakang koma. Buat dua kartu yang angkanya berbeda hanya
0,01, misalnya 0,62 dan 0,63.

Aturan tambahan: bila dua angka berselisih kurang dari 0,02, kartu yang lebih tinggi di
tumpukan yang menang. Letakkan dua kartu yang hampir seri itu di dua posisi teratas,
lalu sebut kartu paling atas sebagai jawaban model.

Sekarang geser kartu 0,62 satu milimeter ke atas, sampai ia melewati kartu 0,63. Tidak
ada angka yang berubah, tidak ada kartu yang ditukar. Lihat kartu paling atas:
jawabannya berganti, karena aturan penentu seri sekarang memilih kartu yang lain.

Yang harus terlihat: jawaban berubah tanpa satu angka pun berubah, sebab yang
menentukan hanya selisih, dan selisih tipis bisa dibalik oleh pergeseran yang tidak
terlihat mata.
\end{tangan}

Dari sudut pandang pelatihan, serangan dan pertahanan adalah dua sisi satu masalah.
Kalau penyerang memaksimalkan galat di dalam kotak, maka pelatih yang berhati-hati
meminimalkan galat terburuk itu:
\[
  \min_{\boldsymbol{\theta}}\;
  \mathbb{E}_{(\mathbf{x},y)\sim\mathcal{D}}
  \Bigl[\max_{\boldsymbol{\delta}\in S(\mathbf{x})}
  L(\boldsymbol{\theta},\mathbf{x}+\boldsymbol{\delta},y)\Bigr].
\]
Rumusan ini adalah masalah pelana: maksimum di dalam, minimum di luar. Suku dalam
adalah persoalan penyerang, suku luar adalah persoalan pelatihan biasa. Inilah
\emph{pelatihan adversarial}. Bila maksimum dalam dihampiri dengan FGSM, kita
mendapat pelatihan yang murah; bila dihampiri dengan PGD beberapa langkah, kita
mendapat pelatihan yang lebih mahal dan biasanya lebih tahan.
\index{pelatihan adversarial}\index{PGD}

Yang bisa dan tidak bisa dijanjikan rumusan itu perlu dipisahkan. Bila nilai pelana
benar-benar tercapai dan nilainya nol, tidak ada gangguan di dalam kotak yang bisa
menipu model, jadi kita punya jaminan. Tetapi masalah pelana itu tidak konkaf di dalam
dan tidak konveks di luar, sehingga algoritme apa pun hanya memberi hampiran: nilai
pelana yang kecil belum tentu tercapai, dan nilai yang diperoleh belum tentu nilai
pelana. Karena itu pelatihan adversarial menghasilkan model yang \emph{lebih tahan
terhadap penyerang yang dipakai saat berlatih}, bukan model yang terbukti tahan
terhadap semua penyerang.

Angka kita memperlihatkan kedua sisi itu. Model $D$ yang dilatih dengan FGSM
$\varepsilon=0{,}15$ kehilangan $1{,}51$ poin akurasi bersih: dari $97{,}99$ menjadi
$96{,}48$ persen. Sebagai gantinya, pada $\varepsilon=0{,}15$ akurasinya di bawah
FGSM adalah $66{,}67$ persen dan di bawah PGD 40 langkah $62{,}65$ persen, sedangkan
model biasa tinggal $14{,}91$ dan $10{,}72$ persen. Pertukaran itu nyata, tetapi
$62{,}65$ persen juga bukan keamanan; ia hanya berarti penyerang harus bekerja lebih
keras.

\section{Transfer, kotak hitam, dan sasaran tertentu}
\tingkat{lanjut}

Sampai di sini penyerang kita tahu segalanya. Di dunia nyata ia sering hanya bisa
mengirim masukan dan melihat jawabannya. Yang menyelamatkan tidak adanya gradien
bukan rahasia model, melainkan kenyataan bahwa arah serangan cenderung berpindah
antar model.

\kamus{transferabilitas}{transferability}{Sifat contoh adversarial yang dibuat untuk
satu model tetapi juga menipu model lain yang dilatih terpisah.}
\index{transferabilitas}

Fenomena ini dilaporkan Szegedy dan rekan pada 2014: contoh yang menipu satu jaringan
sering menipu jaringan lain dengan arsitektur berbeda dan data latih berbeda.
Goodfellow dan rekan mencatat versi kuantitatifnya pada MNIST: contoh yang dibuat
dari model lama memberi galat $19{,}6$ persen pada model yang dilatih adversarial,
sedangkan contoh dari model baru memberi galat $40{,}9$ persen pada model lama.
Arah yang sama, kekuatan berbeda.

Pada percobaan kita, model $A$ (64 unit tersembunyi) dan model $B$ (128 unit) dilatih
pada data latih yang sama tetapi dengan inisialisasi berbeda. Pada
$\varepsilon=0{,}15$, PGD 40 langkah buatan $A$ menurunkan akurasi $A$ sendiri menjadi
$10{,}72$ persen, dan menurunkan akurasi $B$ menjadi $23{,}79$ persen. FGSM buatan
$A$ memberi $14{,}91$ pada $A$ dan $25{,}96$ pada $B$. Jadi serangan pindah, tetapi
melemah: sekitar setengah kekuatannya hilang saat melintasi model.

Penjelasan yang paling ekonomis masih sama dengan bab ini: pada skala kecil kedua model
belajar fungsi yang hampir linear, dan arah yang menaikkan galat pada satu fungsi linear
juga menaikkannya pada fungsi linear lain yang normalnya berdekatan. Itu sebabnya
transfer paling baik untuk serangan satu langkah. Angka kita konsisten: selisih antara
kotak putih dan transfer lebih kecil untuk FGSM ($11$ poin) daripada untuk PGD 40
langkah ($13$ poin).

Untuk penyerang kotak hitam, transfer memberi jalan yang murah: latih model tiruan,
serang model itu, lalu kirim hasilnya ke model korban. Jalur lain adalah menaksir
gradien dari perbedaan keluaran atas banyak kueri, yang akurat tetapi mahal dan mudah
terdeteksi. Kotak hitam lebih lemah karena informasinya kurang, bukan karena
penyerangnya kurang pintar.

Tujuan penyerang juga mengubah angka secara drastis. Serangan tak bersasaran cukup
membuat label berubah; serangan bersasaran harus membuat label menjadi kelas tertentu.
Pada $\varepsilon=0{,}15$, PGD 40 langkah kita mencapai sasaran pada $33{,}00$ persen
contoh, sedangkan serangan tak bersasaran berhasil pada $89{,}06$ persen contoh
(akurasi tinggal $10{,}72$ persen). Dengan $\varepsilon=0{,}30$, sasaran tercapai
pada $99{,}50$ persen contoh. Karena itu angka "serangan berhasil 90 persen" tanpa
keterangan bersasaran atau tidak tidak bisa dibandingkan dengan angka lain. Untuk
penilai pertahanan, uji bersasaran lebih berat bagi penyerang, sedangkan uji tak
bersasaran lebih berat bagi pembela.

\section{Pertahanan: yang bekerja dan yang hanya tampak bekerja}
\tingkat{lanjut}

Pembela punya tiga keluarga jawaban. Yang pertama mengubah data: tambahkan contoh
adversarial ke pelatihan. Yang kedua mengubah model: buat permukaannya lebih mulus,
misalnya dengan masukan yang dikuantisasi atau keluaran yang dihaluskan. Yang ketiga
menambah alat di luar model: deteksi masukan aneh, atau tolak jawaban yang tidak
yakin. Hanya keluarga pertama yang punya dasar jelas, yaitu rumusan pelana tadi;
keluarga kedua dan ketiga sering memberi angka bagus karena alasan yang salah.

Kesalahan yang paling sering berulang punya nama.

\begin{definisi}[Penyamaran gradien]
Sebuah pertahanan melakukan \emph{penyamaran gradien} bila ia membuat gradien yang
dipakai penyerang menjadi tidak berguna: hilang, salah arah, atau acak. Pertahanan
seperti itu menurunkan keberhasilan serangan berbasis gradien tanpa menaikkan jarak
ke batas keputusan.
\end{definisi}

Kalimat sehari-harinya: penjaga tadi bukan menjadi lebih jeli, melainkan memakai
formulir yang tidak bisa dibaca penyerang. Athalye, Carlini, dan Wagner membedakan
tiga cara penyamaran terjadi. Pertama, gradien \emph{hancur}: operasi yang tidak
terdiferensialkan atau tidak stabil secara numerik membuat gradien tidak ada atau
salah. Kedua, gradien \emph{acak}: pertahanan yang menambahkan keacakan membuat satu
kali sampel gradien bukan gradien sebenarnya. Ketiga, gradien \emph{meledak atau
menghilang}: pertahanan yang mengevaluasi jaringan berkali-kali secara berurutan,
seperti pada kuantisasi berulang, berperilaku seperti jaringan sangat dalam saat
diturunkan.
\index{penyamaran gradien}

Mengapa penyamaran bukan pertahanan? Karena ia tidak mengubah apa pun yang dilihat
penyerang yang gigih. Bila gradien tidak berguna, penyerang memakai gradien pengganti
dari model yang lebih mulus; bila keacakan menjadi penghalang, ia merata-ratakan
gradien atas banyak sampel; bila gradien meledak, ia memakai aproksimasi beda hingga.

Kita buat contoh kecilnya sendiri. Ambil model $A$ dan letakkan kuantisasi
masukan di depannya: setiap piksel dibulatkan ke kelipatan $1/8$. Turunan pembulatan
nol hampir di mana-mana, jadi penyerang yang menurunkan gradien lewat jalur itu
mendapat vektor nol dan tidak bergerak. Hasilnya terlihat meyakinkan: akurasi bersih
$97{,}82$ persen, dan setelah "PGD 40 langkah" akurasinya $96{,}98$ persen. Hampir
tidak ada kerusakan. Tetapi bila penyerang memakai gradien model tanpa kuantisasi
pada titik yang sudah dikuantisasi, akurasinya jatuh ke $13{,}23$ persen, hanya sedikit
lebih tinggi daripada $10{,}72$ persen pada model tanpa kuantisasi. Pertahanan itu
tidak menutup jalan; ia hanya menutup jalan yang paling lugu. Bentuk gambarnya ada
pada Gambar~\ref{fig:kuantisasi-gradien}.

\begin{figure}[htbp]
\centering
\begin{tikzpicture}[scale=1.05,font=\scriptsize]
  \fill[layerteal!8] (0,0.5) -- (4,2.9) -- (4,4) -- (0,4) -- cycle;
  \fill[accent!6] (0,0) -- (4,0) -- (4,2.9) -- (0,0.5) -- cycle;
  \draw[very thick,ink] (0,0.5)--(4,2.9);
  \draw[->,layergrey] (-0.2,0)--(4.35,0) node[right]{$x_{1}$};
  \draw[->,layergrey] (0,-0.2)--(0,4.35) node[above]{$x_{2}$};
  \node[font=\scriptsize\bfseries] at (2,4.68) {model asli};
  \node[layerteal] at (0.78,3.55) {kelas $+1$};
  \node[accent] at (3.05,0.8) {kelas $-1$};
  \node[layergrey,rotate=31] at (2.75,2.95) {batas};
  \draw[-{Latex},very thick,accent] (1.2,2.6)--(1.92,1.4);
  \node[accent,above left=-2pt] at (1.56,2.05) {$\boldsymbol{\delta}$};
  \fill[ink] (1.2,2.6) circle (1.6pt);
  \node[ink,above left=-2pt] at (1.2,2.6) {$\mathbf{x}$};
  \draw[accent,thick,fill=accent!35] (1.92,1.4) ++(135:0.07) -- ++(-45:0.14)
    -- ++(-135:0.14) -- ++(45:0.14) -- cycle;
  \node[accent,below right=-2pt] at (1.92,1.4) {$\mathbf{x}+\boldsymbol{\delta}$};
\end{tikzpicture}
\hfill
\begin{tikzpicture}[scale=1.05,font=\scriptsize]
  \fill[layerteal!8] (0,0.5) -- (4,2.9) -- (4,4) -- (0,4) -- cycle;
  \fill[accent!6] (0,0) -- (4,0) -- (4,2.9) -- (0,0.5) -- cycle;
  \foreach \g in {0.5,1,1.5,2,2.5,3,3.5}
    {\draw[layergrey!30,very thin] (\g,0)--(\g,4);
     \draw[layergrey!30,very thin] (0,\g)--(4,\g);}
  \draw[very thick,ink] (0,0.5)--(4,2.9);
  \draw[->,layergrey] (-0.2,0)--(4.35,0) node[right]{$x_{1}$};
  \draw[->,layergrey] (0,-0.2)--(0,4.35) node[above]{$x_{2}$};
  \node[font=\scriptsize\bfseries] at (2,4.68) {dengan pembulatan $1/8$};
  \fill[ink] (1.2,2.6) circle (1.6pt);
  \node[ink,above left=-2pt] at (1.2,2.6) {$\mathbf{x}$};
  \draw[layerplum,thick] (1.2,2.6) circle (0.13);
  \node[layerplum,anchor=west] at (1.4,2.78) {$\nabla$ jalur $=\mathbf{0}$};
  \draw[-{Latex},very thick,dashed,layerplum] (1.2,2.6)--(1.92,1.4);
  \node[layerplum,anchor=east] at (1.44,1.72) {$\nabla$ pengganti};
  \draw[layerplum,thick,fill=layerplum!35] (1.92,1.4) ++(135:0.07)
    -- ++(-45:0.14) -- ++(-135:0.14) -- ++(45:0.14) -- cycle;
\end{tikzpicture}
\caption{Kuantisasi masukan pada model yang sama. Kiri: gradien asli memberi arah
$\boldsymbol{\delta}$ yang menyeberangkan $\mathbf{x}$ ke sisi lain batas, sehingga
labelnya berubah. Kanan: pembulatan setiap ciri ke kelipatan $1/8$ membuat turunan
jalur nol di hampir semua titik, jadi panah gradien yang dilaporkan berhenti;
batas keputusannya sendiri tidak bergeser, dan gradien dari model tanpa pembulatan
masih menyeberangkannya. Kotak menandai titik sesudah gangguan.}
\label{fig:kuantisasi-gradien}
\end{figure}

Athalye dan rekan memeriksa sembilan pertahanan kotak putih di ICLR 2018. Tujuh
di antaranya bergantung pada penyamaran gradien, dan serangan yang disesuaikan
berhasil menembus enam sepenuhnya dan satu sebagian, pada model ancaman yang
diklaim masing-masing. Ini bukan kesalahan kecil dalam satu makalah; ia pola yang
berulang. Karena itu ada daftar gejala yang bisa dipakai sebelum percaya pada angka
ketahanan. Bila serangan satu langkah lebih kuat daripada serangan berulang, ada yang
salah: serangan berulang seharusnya selalu lebih kuat di kotak putih. Bila serangan
kotak hitam lebih kuat daripada kotak putih, ada yang salah: kotak hitam adalah
himpunan informasi yang lebih kecil. Bila serangan tanpa batas gangguan tidak
mencapai seratus persen keberhasilan, ada yang salah. Bila pencarian acak menemukan
contoh adversarial padahal serangan gradien tidak, ada yang salah. Bila memperbesar
anggaran tidak menaikkan keberhasilan, ada yang salah.

\begin{tantangan}
Sebuah layanan menolak menjawab bila nilai mutlak margin skornya di bawah 0,50, dan
baru menjawab bila ia cukup yakin. Pada contoh linear bagian sebelumnya, margin bersih
1,500 dan $\lVert\mathbf{d}\rVert_1=6$. Serangan pada $\varepsilon=0{,}25$ mendorong
margin ke nol, tepat di ambang pembalikan, jadi ditolak, dan layanan tampak aman.
Tentukan anggaran terkecil yang membuat penyerang tetap dilayani, yaitu yang mematikan
penolakan di sisi yang salah. Apakah aturan itu pertahanan atau penyamaran?

Petunjuk 1: penolakan hanya melihat besar margin, bukan tandanya.

Petunjuk 2: selesaikan $1{,}500-6\varepsilon=-0{,}500$.

Kalau mentok, kembali ke bagian Contoh hitung: model linear yang bisa dihitung tangan.
\end{tantangan}

Seperti apa evaluasi yang benar? Ia bukan satu angka, melainkan satu protokol. Tetapkan
model ancaman lebih dulu: norma, $\varepsilon$, kotak putih atau hitam, bersasaran atau
tidak. Pakai penyerang terkuat yang sanggup dibayar, dan sesuaikan ia dengan pertahanan:
kalau pertahanan mengkuantisasi, serang lewat gradien pengganti; kalau ia menambahkan
keacakan, rata-ratakan atas beberapa sampel; kalau ia menyaring masukan, serang
saringan itu juga. Laporkan akurasi di beberapa $\varepsilon$, bukan satu titik, dan
periksa bahwa kurvanya turun monoton. Laporkan lebih dari satu benih acak dan sertakan
serangan transfer dari model lain. Laporkan biayanya: jumlah lintasan mundur per contoh,
waktu, dan perangkat keras. Bila ingin bicara jaminan, pakai sertifikasi, bukan
pengujian; itu jalur lain yang tidak dibahas di sini.

\begin{untukanak}
Coba pikirkan soal ujian pilihan ganda yang diperiksa mesin. Mesin itu membaca jawaban
dari posisi bulatan yang dihitamkan. Kalau kita hanya menguji mesin dengan jawaban yang
dihitamkan rapi, mesin akan selalu benar. Lalu datang murid yang tahu cara mesin
membaca: ia menghitamkan satu bulatan lain sedikit lebih tebal, dan mesin membacanya
sebagai jawaban berbeda. Kertasnya masih kertas yang sama, tulisannya masih tulisan
yang sama.

Sekarang bagian yang menarik. Ada sekolah yang mencoba memperbaiki mesin itu dengan
memindai kertas dua kali dan mengambil jawaban yang paling sering muncul. Murid yang
sama tinggal menghitamkan dua bulatan, dan mesin tetap bisa ditipu. Sekolah lain
menambahkan penggaris agar posisi bulatan bergeser sedikit; murid yang sabar tinggal
mencoba beberapa posisi. Semua tambahan itu membuat penipuan lebih repot, tetapi tidak
membuat mesin bisa membedakan jawaban yang jujur.

Di bab ini kita akan tahu mengapa menambah kesulitan bagi penyerang tidak sama dengan
menutup celah, dan bagaimana menguji sebuah perbaikan supaya kita tidak tertipu oleh
perbaikan yang hanya tampak berhasil.
\end{untukanak}

\section{Contoh hitung: model linear yang bisa dihitung tangan}
\tingkat{menengah}

Sekarang kita tinggalkan jaringan dan kembali ke model yang bisa dihitung tanpa
komputer: pengklasifikasi linear dengan dua kelas dan empat ciri. Semua angka di
bagian ini bisa diperiksa dengan kalkulator.

Misalkan kelas $A$ memakai bobot $\mathbf{w}_A=(1;-2;0{,}5;1{,}5)$ dan kelas $B$
memakai $\mathbf{w}_B=(0;1;0{,}5;-0{,}5)$, tanpa bias. Contoh yang diperiksa adalah
$\mathbf{x}=(1;0{,}5;1;1)$. Skor kedua kelas adalah
\[
  s_A=\mathbf{w}_A^{\top}\mathbf{x}=1-1+0{,}5+1{,}5=2{,}0,
  \qquad
  s_B=\mathbf{w}_B^{\top}\mathbf{x}=0+0{,}5+0{,}5-0{,}5=0{,}5 .
\]
Marginnya $m=s_A-s_B=1{,}5$, jadi model yakin tetapi tidak berlebihan. Selisih bobot
adalah $\mathbf{d}=\mathbf{w}_A-\mathbf{w}_B=(1;-3;0;2)$, dengan
$\lVert\mathbf{d}\rVert_1=1+3+0+2=6$ dan
$\lVert\mathbf{d}\rVert_2=\sqrt{1+9+0+4}=\sqrt{14}\approx3{,}7417$.

Serangan tak bersasaran harus menurunkan margin. Menurut lemma, penurunan terbesar
dicapai oleh $\boldsymbol{\delta}=-\varepsilon\,\mathrm{sign}(\mathbf{d})
=\varepsilon(-1;1;0;-1)$, dan besarnya penurunan adalah
$\varepsilon\lVert\mathbf{d}\rVert_1=6\varepsilon$. Maka margin setelah serangan
adalah $1{,}5-6\varepsilon$, dan label berubah ketika
\[
  \varepsilon>\varepsilon^{*}=\frac{1{,}5}{6}=0{,}25 .
\]
Pada $\varepsilon=0{,}25$ marginnya tepat nol, jadi kedua kelas seri dan
$\arg\max$ memilih kelas pertama; pembalikan yang tegas baru terjadi sedikit di atas
$0{,}25$. Tabel berikut memuat seluruh jalannya, dan setiap barisnya bisa diperiksa
dengan mengalikan $\mathbf{d}$ dengan $\mathbf{x}+\boldsymbol{\delta}$.

\begin{center}
\begin{tabular}{rccc}
\hline
$\varepsilon$ & $s_A$ & $s_B$ & margin \\
\hline
$0$    & $2{,}000$ & $0{,}500$ & $+1{,}500$ \\
$0{,}05$ & $1{,}775$ & $0{,}575$ & $+1{,}200$ \\
$0{,}10$ & $1{,}550$ & $0{,}650$ & $+0{,}900$ \\
$0{,}15$ & $1{,}325$ & $0{,}725$ & $+0{,}600$ \\
$0{,}20$ & $1{,}100$ & $0{,}800$ & $+0{,}300$ \\
$0{,}25$ & $0{,}875$ & $0{,}875$ & $0{,}000$ \\
$0{,}30$ & $0{,}650$ & $0{,}950$ & $-0{,}300$ \\
\hline
\end{tabular}
\end{center}

Bila penyerang dibatasi norma $\ell_2$ dan bukan $\ell_\infty$, jawabannya berubah.
Maksimum $\mathbf{d}^{\top}\boldsymbol{\delta}$ pada bola $\ell_2$ berjari-jari
$\varepsilon$ adalah $\varepsilon\lVert\mathbf{d}\rVert_2$, dicapai pada
$\boldsymbol{\delta}=\varepsilon\mathbf{d}/\lVert\mathbf{d}\rVert_2$. Maka
$\varepsilon^{*}=1{,}5/\sqrt{14}\approx0{,}4009$. Anggaran $\ell_2$ yang memberi
kerusakan sama besarnya dengan $\varepsilon=0{,}25$ pada $\ell_\infty$ adalah
$0{,}4009$; keduanya memindahkan contoh ke titik yang sama, yaitu titik potong
$\mathbf{x}-\varepsilon\,\mathrm{sign}(\mathbf{d})$. Ini contoh kecil dari aturan
umum: angka $\varepsilon$ tidak bisa dibandingkan antar norma.

Sekarang bagian yang sering ditanyakan: apakah PGD lebih kuat pada model ini? Tiap
langkah PGD menambah penurunan margin sebesar $\alpha\lVert\mathbf{d}\rVert_1
=6\alpha$, tetapi seluruh lintasan tetap di dalam bola, jadi penurunan total tidak
pernah melebihi $6\varepsilon$. Untuk model linear, titik terburuk di dalam bola
memang titik yang dipilih FGSM, jadi iterasi tidak bisa mengalahkannya. Hitungan
tangan memperlihatkannya: dengan $\varepsilon=0{,}10$ dan 40 langkah
$\alpha=0{,}01$, margin akhir $0{,}900$ sama persis dengan hasil FGSM pada
$\varepsilon=0{,}10$.

\kertas{Ulangi contoh di atas untuk $\mathbf{w}_A=(2;1;-1;0{,}5)$,
$\mathbf{w}_B=(0;1;1;-0{,}5)$, dan $\mathbf{x}=(1;1;0{,}5;1)$. Hitung $s_A$, $s_B$,
$\mathbf{d}$, $\lVert\mathbf{d}\rVert_1$, dan $\varepsilon^{*}$ untuk norma
$\ell_\infty$. Lalu hitung $\varepsilon^{*}$ untuk norma $\ell_2$ dan jelaskan
mengapa kedua angka berbeda. Periksa dengan potongan Python tiga baris. Setelah itu
putuskan, tanpa menjalankan program, apakah PGD 20 langkah dengan
$\alpha=\varepsilon/10$ bisa mengubah label ketika $\varepsilon$ hanya setengah dari
$\varepsilon^{*}$.}

\section{Lab komputasi}
\tingkat{lanjut}

Lab ini melatih tiga jaringan kecil pada data angka tulisan tangan yang disediakan
pustaka \texttt{scikit-learn}: 1797 contoh, 64 ciri, nilai $0$ sampai $16$ yang kita
bagi $16$ sehingga masuk ke $[0,1]$. Dua belas ratus contoh dipakai untuk latih dan
597 untuk uji. Model $A$ punya satu lapis tersembunyi 64 unit, model $B$ 128 unit,
dan model $D$ sama dengan $A$ tetapi dilatih dengan FGSM $\varepsilon=0{,}15$
(alasan memakai $D$ bukan $A$ adalah agar perbandingan bersih dan tahan serangan
berasal dari model yang sama besarnya).

Seluruh gradien dihitung tangan dengan NumPy, tanpa pustaka turunan otomatis, supaya
setiap langkah bisa diperiksa. Fungsi \texttt{grad\_x} adalah turunan balik dari Bab 25
yang berhenti di masukan, bukan di bobot.

\begin{lab}
\begin{verbatim}
import numpy as np, time
from sklearn.datasets import load_digits

X, y = load_digits(return_X_y=True)
X = X / 16.0                        # 0..16 -> 0..1
perm = np.random.default_rng(0).permutation(len(X))
tr, te = perm[:1200], perm[1200:]
Xtr, ytr, Xte, yte = X[tr], y[tr], X[te], y[te]
K = 10
Ytr, Yte = np.eye(K)[ytr], np.eye(K)[yte]

def init(dims, seed):
    r = np.random.default_rng(seed); W, b = [], []
    for i in range(len(dims) - 1):
        W.append(r.normal(0, np.sqrt(2.0 / dims[i]),
                          (dims[i], dims[i + 1])))
        b.append(np.zeros(dims[i + 1]))
    return W, b

def fwd(W, b, X):
    acts, pre, h = [X], [], X
    for i in range(len(W) - 1):
        z = h @ W[i] + b[i]; pre.append(z)
        h = np.maximum(z, 0); acts.append(h)
    z = h @ W[-1] + b[-1]; pre.append(z)
    z = z - z.max(1, keepdims=True); p = np.exp(z)
    return p / p.sum(1, keepdims=True), acts, pre

def loss_grad(W, b, X, Y):
    p, acts, pre = fwd(W, b, X); n = len(X)
    loss = -np.mean(np.sum(Y * np.log(p + 1e-12), 1))
    d = p - Y; gW = [None] * len(W); gb = [None] * len(W)
    for i in range(len(W) - 1, -1, -1):
        gW[i] = acts[i].T @ d / n + 1e-4 * W[i]; gb[i] = d.mean(0)
        if i > 0:
            d = (d @ W[i].T) * (pre[i - 1] > 0)
    return loss, gW, gb

def grad_x(W, b, X, Y):             # gradien galat terhadap masukan
    p, acts, pre = fwd(W, b, X)
    d = (p - Y) / len(X)
    for i in range(len(W) - 1, 0, -1):
        d = (d @ W[i].T) * (pre[i - 1] > 0)
    return d @ W[0].T

def train(dims, seed, Xa, Ya, epochs=350, lr=0.01, adv_eps=None):
    W, b = init(dims, seed)
    mW = [np.zeros_like(w) for w in W]
    vW = [np.zeros_like(w) for w in W]
    mb = [np.zeros_like(v) for v in b]
    vb = [np.zeros_like(v) for v in b]
    for t in range(1, epochs + 1):
        if adv_eps is None:
            Xb = Xa
        else:
            g = grad_x(W, b, Xa, Ya)
            Xb = np.clip(Xa + adv_eps * np.sign(g), 0, 1)
        _, gW, gb = loss_grad(W, b, Xb, Ya)
        for i in range(len(W)):
            mW[i] = 0.9 * mW[i] + 0.1 * gW[i]
            vW[i] = 0.999 * vW[i] + 0.001 * gW[i] ** 2
            mb[i] = 0.9 * mb[i] + 0.1 * gb[i]
            vb[i] = 0.999 * vb[i] + 0.001 * gb[i] ** 2
            mh = mW[i] / (1 - 0.9 ** t)
            vh = vW[i] / (1 - 0.999 ** t)
            W[i] -= lr * mh / (np.sqrt(vh) + 1e-8)
            mh = mb[i] / (1 - 0.9 ** t)
            vh = vb[i] / (1 - 0.999 ** t)
            b[i] -= lr * mh / (np.sqrt(vh) + 1e-8)
    return W, b

def acc(W, b, X, y):
    return float(np.mean(fwd(W, b, X)[0].argmax(1) == y))

def fgsm(W, b, X, y, eps):
    g = grad_x(W, b, X, np.eye(K)[y])
    return np.clip(X + eps * np.sign(g), 0, 1)

def pgd(W, b, X, y, eps, steps, Ylab=None, arah=1):
    alpha = eps / 10
    r = np.random.default_rng(1)
    xa = np.clip(X + r.uniform(-eps, eps, X.shape), 0, 1)
    Y = np.eye(K)[y] if Ylab is None else Ylab
    for _ in range(steps):
        g = arah * np.sign(grad_x(W, b, xa, Y))
        xa = np.clip(np.clip(xa + alpha * g, X - eps, X + eps), 0, 1)
    return xa
\end{verbatim}
\end{lab}

Sisanya adalah pemanggilan dan pencetakan. Kita melatih $A$, $B$, dan $D$, lalu
menyapu $\varepsilon$ dari $0{,}02$ sampai $0{,}30$.

\begin{verbatim}
WA, bA = train([64, 64, 10], 1, Xtr, Ytr)
WB, bB = train([64, 128, 10], 2, Xtr, Ytr)
WD, bD = train([64, 64, 10], 3, Xtr, Ytr, epochs=120, adv_eps=0.15)

print(f"akurasi uji bersih: A {acc(WA,bA,Xte,yte):.4f}"
      f"  B {acc(WB,bB,Xte,yte):.4f}  D {acc(WD,bD,Xte,yte):.4f}")
print(f"{'eps':>5} {'A bersih':>9} {'A FGSM':>7} {'A PGD40':>8}"
      f" {'D FGSM':>7} {'D PGD40':>8}")
for eps in [0.02, 0.05, 0.10, 0.15, 0.20, 0.30]:
    xf = fgsm(WA, bA, Xte, yte, eps)
    xp = pgd(WA, bA, Xte, yte, eps, 40)
    xd = fgsm(WD, bD, Xte, yte, eps)
    xq = pgd(WD, bD, Xte, yte, eps, 40)
    print(f"{eps:5.2f} {acc(WA,bA,Xte,yte):9.4f} {acc(WA,bA,xf,yte):7.4f}"
          f" {acc(WA,bA,xp,yte):8.4f} {acc(WD,bD,xd,yte):7.4f}"
          f" {acc(WD,bD,xq,yte):8.4f}")

r = np.random.default_rng(7)
rs = np.sign(r.normal(size=Xte.shape))          # kontrol: tanda acak
print("tanda acak eps=0.15:",
      f"{acc(WA,bA,np.clip(Xte+0.15*rs,0,1),yte):.4f}")
print("langkah PGD, eps=0.15, alpha=eps/10, mulai acak")
for k in [1, 2, 5, 10, 20, 40]:
    print(f"  {k:2d} langkah:",
          f"{acc(WA,bA,pgd(WA,bA,Xte,yte,0.15,k),yte):.4f}")
print("  FGSM (1 langkah, alpha=eps):",
      f"{acc(WA,bA,fgsm(WA,bA,Xte,yte,0.15),yte):.4f}")
\end{verbatim}

Keluaran yang sebenarnya, apa adanya, untuk blok di atas:

\begin{verbatim}
akurasi uji bersih: A 0.9799  B 0.9849  D 0.9648
  eps  A bersih  A FGSM  A PGD40  D FGSM  D PGD40
 0.02    0.9799  0.9548   0.9548  0.9497   0.9497
 0.05    0.9799  0.8476   0.8425  0.9162   0.9129
 0.10    0.9799  0.4992   0.4590  0.8241   0.8174
 0.15    0.9799  0.1491   0.1072  0.6667   0.6265
 0.20    0.9799  0.0184   0.0050  0.4858   0.3719
 0.30    0.9799  0.0000   0.0000  0.0687   0.0201
tanda acak eps=0.15: 0.9548
langkah PGD, eps=0.15, alpha=eps/10, mulai acak
   1 langkah: 0.9464
   2 langkah: 0.8878
   5 langkah: 0.6667
  10 langkah: 0.2479
  20 langkah: 0.1122
  40 langkah: 0.1072
  FGSM (1 langkah, alpha=eps): 0.1491
\end{verbatim}

Bacaan pertama tabel itu: akurasi bersih $97{,}99$ persen, dan pada
$\varepsilon=0{,}15$ tinggal $14{,}91$ persen.

Bacaan kedua: jarak FGSM dan PGD melebar di tengah, dari nol pada $0{,}02$ menjadi
$4{,}19$ poin pada $0{,}15$. Pada anggaran kecil satu langkah sudah menyentuh dinding
kotak; pada anggaran besar beberapa langkah diperlukan untuk menemukan arah yang lebih
buruk.

Bacaan ketiga, dan yang paling penting untuk bab ini: baris tanda acak. Kalau
gangguan sebesar $\varepsilon=0{,}15$ diberikan dengan tanda acak di setiap piksel,
akurasi hanya turun menjadi $95{,}48$ persen. FGSM dengan besar yang sama
menurunkannya menjadi $14{,}91$ persen. Jadi pemilihan tanda oleh gradien bernilai
sekitar delapan puluh poin akurasi: arahnya yang bekerja, bukan besarnya.

Bacaan keempat: pada $\varepsilon=0{,}30$, $A$ sudah nol persen sedangkan $D$ masih
$6{,}87$ persen di bawah FGSM dan $2{,}01$ persen di bawah PGD. Pelatihan adversarial
tidak membuat model kebal; ia menggeser kurva ke kanan.

Blok berikutnya mengukur dua hal yang tidak terlihat di tabel: apakah gradien yang
kita hitung memang gradien, dan berapa biayanya.

\begin{verbatim}
i0 = int(np.argmax(np.abs(grad_x(WA, bA, Xte, Yte)).sum(1)))
x0, y0 = Xte[i0:i0 + 1], Yte[i0:i0 + 1]
g0 = grad_x(WA, bA, x0, y0)
L1 = lambda z: loss_grad(WA, bA, z, y0)[0]
print(f"beda hingga pada contoh i={i0}, h=1e-6")
for j in np.argsort(-np.abs(g0[0]))[:3]:
    e = np.zeros_like(x0); e[0, j] = 1e-6
    fd = (L1(x0 + e) - L1(x0 - e)) / 2e-6
    print(f"  j={j:2d} analitik {g0[0,j]:+.6e} beda-hingga {fd:+.6e}"
          f" galat {abs(fd-g0[0,j])/abs(g0[0,j]):.0e}")
g_all = grad_x(WA, bA, Xte, Yte)
lin = 0.15 * np.abs(g_all).sum(1).mean()
lc = loss_grad(WA, bA, Xte, Yte)[0]
la = loss_grad(WA, bA, fgsm(WA, bA, Xte, yte, 0.15), Yte)[0]
print(f"rata-rata |g|_1 {np.abs(g_all).sum(1).mean():.4f};"
      f" ramalan linear {lin:.4f} nats; kenyataan {la-lc:.3f} nats")

t0 = time.perf_counter(); fgsm(WA, bA, Xte, yte, 0.15)
t1 = time.perf_counter(); pgd(WA, bA, Xte, yte, 0.15, 40)
t2 = time.perf_counter()
print(f"waktu 597 contoh: FGSM {t1-t0:.4f} s  PGD40 {t2-t1:.4f} s")
\end{verbatim}

Keluarannya:

\begin{verbatim}
beda hingga pada contoh i=145, h=1e-6
  j=53 analitik -6.391833e+00 beda-hingga -6.391833e+00 galat 2e-11
  j=38 analitik +6.211518e+00 beda-hingga +6.211518e+00 galat 1e-10
  j=13 analitik -6.175321e+00 beda-hingga -6.175321e+00 galat 8e-11
rata-rata |g|_1 0.0056; ramalan linear 0.0008 nats; kenyataan 4.862 nats
waktu 597 contoh: FGSM 0.0006 s  PGD40 0.0302 s
\end{verbatim}

Tiga baris pertama memastikan bahwa turunan balik kita benar: gradien analitik dan
beda hingga dengan $h=10^{-6}$ sepakat sampai sekitar sepuluh angka, dengan galat
relatif $10^{-10}$ dan $10^{-11}$. Untuk jaringan yang memakai ReLU, kesepakatan
sehalus itu masuk akal: pada skala $10^{-6}$ tidak ada unit yang berganti tanda,
sehingga fungsinya benar-benar linear di sekitar titik itu. Kalau angka ini berbeda
jauh, yang salah adalah kode turunan baliknya, bukan serangannya.

Baris keempat adalah bagian yang mengejutkan. Rata-rata $\lVert\nabla L\rVert_1$ hanya
$0{,}0056$, sehingga ramalan linear untuk $\varepsilon=0{,}15$ adalah $0{,}0008$
nats. Kenaikan galat yang sesungguhnya $4{,}862$ nats. Jadi suku linear
meremehkan kerusakan sekitar empat ribu kali. Ini bukan tanda bahwa gradiennya salah
(tiga baris di atas membuktikan kebalikannya), melainkan tanda bahwa $L$ sangat
melengkung pada jarak $0{,}15$: unit ReLU berganti tanda, dan hampiran orde pertama
kehilangan makna. Yang tetap berguna adalah tandanya, seperti ditunjukkan baris
tanda acak tadi.

Baris terakhir soal waktu: pada 597 contoh, FGSM memerlukan $0{,}6$ milidetik dan
PGD 40 langkah $30$ milidetik, sekitar lima puluh kali lebih mahal. Angka ini
bergantung pada mesin dan pada vektorisasi, jadi yang layak diingat adalah
ordenya, bukan angka tepatnya.

Terakhir, demonstrasi penyamaran gradien. Kita menambahkan kuantisasi masukan ke
kelipatan $1/8$ di depan model $A$, lalu menyerang dengan dua cara.

\begin{verbatim}
def q(z): return np.round(z * 8) / 8
accq = lambda W, b, Z, yy: float(np.mean(fwd(W, b, q(Z))[0]
                                        .argmax(1) == yy))
def pgd_naif(W, b, Z, yy, eps, steps):
    r = np.random.default_rng(1)
    z = np.clip(Z + r.uniform(-eps, eps, Z.shape), 0, 1)
    for _ in range(steps):          # gradien = nol, jadi diam
        z = np.clip(np.clip(z, Z - eps, Z + eps), 0, 1)
    return z
def pgd_bpda(W, b, Z, yy, eps, steps):
    r = np.random.default_rng(1)
    z = np.clip(Z + r.uniform(-eps, eps, Z.shape), 0, 1)
    for _ in range(steps):
        g = np.sign(grad_x(W, b, q(z), np.eye(K)[yy]))
        z = np.clip(np.clip(z + eps / 10 * g, Z - eps, Z + eps), 0, 1)
    return z
print("masukan dikuantisasi ke kelipatan 1/8")
print(f"  bersih {accq(WA,bA,Xte,yte):.4f}"
      f"  PGD40 naif {accq(WA,bA,pgd_naif(WA,bA,Xte,yte,0.15,40),yte):.4f}"
      f"  PGD40 BPDA"
      f" {accq(WA,bA,pgd_bpda(WA,bA,Xte,yte,0.15,40),yte):.4f}")
\end{verbatim}

Keluarannya:

\begin{verbatim}
masukan dikuantisasi ke kelipatan 1/8
  bersih 0.9782  PGD40 naif 0.9698  PGD40 BPDA 0.1323
\end{verbatim}

Inilah bentuk paling ringkas dari pertahanan palsu. Angka bersihnya hanya $0{,}17$
poin di bawah model asli, dan serangan naif tidak bergerak sama sekali: seandainya
hanya itu yang dilaporkan, pertahanan ini tampak hampir sempurna. Gradien pengganti
menjatuhkannya ke $13{,}23$ persen, nyaris sama dengan $10{,}72$ persen pada model
tanpa kuantisasi. Yang berubah terutama keterbacaan gradiennya.

\section{Biaya, keterbatasan, dan cara melaporkan}
\tingkat{dasar}

Anggaran komputasi menentukan serangan mana yang jujur dipakai. FGSM memerlukan satu
lintasan mundur per contoh. PGD $k$ langkah memerlukan $k$ lintasan. Pada percobaan
kita, biaya PGD 40 langkah sekitar lima puluh kali FGSM untuk hasil yang, pada
$\varepsilon$ kecil, hampir sama. Untuk penilaian pertahanan, biaya itu wajar: angka
ketahanan yang diukur dengan serangan lemah tidak berarti. Untuk pelatihan, biaya itu
berlipat lagi: setiap epoch menuntut serangan baru.

Pelatihan adversarial juga punya harga akurasi. Model $D$ kehilangan $1{,}51$ poin
akurasi bersih untuk kenaikan sekitar lima puluh poin pada $\varepsilon=0{,}15$.
Madry dan rekan melaporkan bahwa model yang lebih besar memberi ketahanan lebih baik,
yang menunjukkan bahwa ketahanan bukan sifat cuma-cuma dari arsitektur, melainkan
kapasitas yang harus dibayar.

Ada tiga keterbatasan sebelum angka di bab ini dipakai.

Pertama, skala. Data kita hanya 64 ciri dan 10 kelas, dengan 597 contoh uji. Pada
citra besar, anggaran yang berarti jauh lebih kecil dari $0{,}15$; pada ImageNet,
Goodfellow dan rekan memakai $0{,}007$ dan sudah membalik klasifikasi. Angka
$\varepsilon$ tidak bisa dipindahkan antar data, karena artinya bergantung pada skala
fitur. Pada data kita, piksel sudah terkuantisasi ke kelipatan $1/16$, sehingga
$\varepsilon=0{,}05$ berarti kurang dari satu langkah kuantisasi per piksel.

Kedua, cakupan serangan. Yang kita jalankan hanya dua keluarga: FGSM dan PGD
$\ell_\infty$. Model yang tahan terhadap keduanya belum tentu tahan terhadap serangan
$\ell_2$, serangan yang mengubah satu piksel saja, atau serangan yang menempelkan
benda di dunia nyata. Setiap klaim ketahanan hanya berlaku untuk keluarga serangan
yang diuji, dan itulah sebabnya protokol evaluasi di bagian sebelumnya menuntut
beberapa jenis serangan sekaligus.

Ketiga, tidak ada jaminan. Semua angka di sini adalah hasil pengujian, bukan batas
bawah yang terbukti. Model $D$ yang mencapai $62{,}65$ persen pada PGD 40 langkah
mungkin masih bisa ditembus serangan yang lebih baik; kita tidak tahu, karena yang
tidak ditemukan bukan bukti tidak ada. Untuk pernyataan yang lebih kuat, diperlukan
sertifikasi: misalnya dengan randomized smoothing atau dengan verifikasi formal pada
jaringan kecil. Keduanya memberi jaminan pada radius tertentu, dengan harga akurasi
dan komputasi yang jauh lebih besar, dan keduanya berada di luar cakupan bab ini.

Cara melaporkan yang jujur memuat lima hal: model ancaman lengkap (norma,
$\varepsilon$, pengetahuan penyerang, bersasaran atau tidak), penyerang yang dipakai
beserta jumlah dan ukuran langkahnya, akurasi bersih dan akurasi di beberapa
$\varepsilon$, jumlah benih dan sebaran hasilnya, serta biaya komputasi. Tanpa itu,
angka "tahan 90 persen" tidak dapat diperiksa dan tidak dapat dibandingkan.
Bab 27 sudah mengajari kita bahwa metrik bisa menipu; di sini metriknya menipu dengan
cara yang lebih halus, karena penyebutnya benar dan hanya pembilangnya yang dipilih.

\begin{untukdosen}
Urutan mengajar yang berhasil dipakai: mulai dari tabel akurasi terhadap
$\varepsilon$, bukan dari definisi. Mahasiswa melihat $97{,}99$ dan $14{,}91$ pada
model yang sama, lalu bertanya mengapa. Setelah itu baru definisi model ancaman,
lemma maksimum bentuk linear, dan FGSM. Baru sesudah FGSM berjalan, masuk ke PGD dan
rumusan pelana. Urutan sebaliknya, yaitu definisi dulu lalu serangan, membuat
mahasiswa menghafal notasi tanpa pernah melihat kerusakannya.

Dua miskonsepsi khas. Pertama, mahasiswa menganggap serangan adversarial sama dengan
menambahkan noise, sehingga mereka menyimpulkan bahwa model yang tahan noise akan
tahan serangan. Baris tanda acak di lab membantahnya dengan satu angka: noise sebesar
$\varepsilon=0{,}15$ hanya menurunkan akurasi ke $95{,}48$ persen, sedangkan FGSM
dengan besar yang sama menurunkannya ke $14{,}91$ persen. Kedua, mereka menganggap
satu angka akurasi di bawah serangan sudah merupakan bukti ketahanan. Demonstrasi
kuantisasi di lab menunjukkan sebaliknya: $96{,}98$ persen terlihat aman sampai
seseorang memakai gradien pengganti dan mendapat $13{,}23$ persen.

Aktivitas papan tulis yang murah: tulis contoh linear empat ciri dari bagian contoh
hitung, lalu minta kelas menghitung $\varepsilon^{*}$ untuk $\ell_\infty$ dan
$\ell_2$. Setelah itu minta mereka memprediksi, tanpa menghitung, apakah menambah
langkah PGD mengubah hasil pada model linear itu. Jawaban yang benar, beserta
alasannya, memaksa mereka memahami bahwa iterasi hanya berguna karena $f$ tidak
linear.

Tanda kesiapan lanjut: mahasiswa dapat menjelaskan mengapa hanya tanda gradien yang
dipakai, dapat menghitung $\varepsilon^{*}$ untuk model linear, dapat menyebutkan
satu gejala penyamaran gradien, dan dapat merancang satu serangan yang disesuaikan
untuk pertahanan yang mengkuantisasi masukan. Bila mereka menuntut angka di beberapa
$\varepsilon$ sebelum percaya, mereka siap masuk Bab 30.
\end{untukdosen}

\section{Ringkasan}

\begin{ringkasan}
\begin{itemize}
  \item Klaim ketahanan selalu relatif terhadap model ancaman: norma, anggaran
        $\varepsilon$, pengetahuan penyerang, dan tujuan bersasaran atau tidak.
        Tanpa itu, dua angka tidak bisa dibandingkan.
  \item FGSM menggeser setiap koordinat sebesar $\varepsilon$ mengikuti tanda gradien
        galat. Pada model linear ia eksak dan optimal, karena
        $\max_{\lVert\boldsymbol{\delta}\rVert_\infty\le\varepsilon}
        \mathbf{w}^{\top}\boldsymbol{\delta}=\varepsilon\lVert\mathbf{w}\rVert_1$.
  \item Linearisasi membenarkan arah, bukan jarak. Pada percobaan kita ramalan linear
        $0{,}0008$ nats, sedangkan kenaikan sebenarnya $4{,}862$ nats; yang tetap
        berguna adalah pola tandanya.
  \item PGD menambahkan iterasi dan proyeksi ke kotak. Pada $\varepsilon=0{,}15$
        akurasi turun dari $94{,}64$ persen (1 langkah kecil) ke $10{,}72$ persen
        (40 langkah), dan mendatar setelah sekitar dua puluh langkah.
  \item Pelatihan adversarial ialah hampiran masalah pelana
        $\min_\theta\mathbb{E}\max_\delta L$. Ia menaikkan ketahanan dengan harga
        akurasi bersih: model kita kehilangan $1{,}51$ poin untuk naik dari $10{,}72$
        ke $62{,}65$ persen.
  \item Penyamaran gradien menurunkan keberhasilan serangan tanpa menaikkan jarak ke
        batas keputusan. Contoh kita: $96{,}98$ persen di bawah serangan naif,
        $13{,}23$ persen di bawah serangan dengan gradien pengganti.
  \item Evaluasi yang jujur memakai serangan terkuat yang sanggup dibayar, disesuaikan
        dengan pertahanan, di beberapa $\varepsilon$ dan beberapa benih, dengan biaya
        yang dilaporkan.
\end{itemize}
\end{ringkasan}

\section{Latihan}

\begin{latihan}
\begin{enumerate}
  \item Hitung $\varepsilon^{*}$ untuk model linear dengan
        $\mathbf{w}_A=(1;-1;2)$, $\mathbf{w}_B=(0;1;1)$ dan
        $\mathbf{x}=(1;2;1)$ pada norma $\ell_\infty$, lalu pada norma $\ell_2$.
        Periksa kedua jawaban dengan Python. (Petunjuk: hitung
        $\mathbf{d}=\mathbf{w}_A-\mathbf{w}_B$ lebih dulu, lalu bandingkan
        $\lVert\mathbf{d}\rVert_1$ dengan $\lVert\mathbf{d}\rVert_2$.)

  \item Buktikan bahwa $\max_{\lVert\boldsymbol{\delta}\rVert_2\le\varepsilon}
        \mathbf{w}^{\top}\boldsymbol{\delta}=\varepsilon\lVert\mathbf{w}\rVert_2$,
        dan tentukan semua $\boldsymbol{\delta}$ yang mencapai maksimum. Jelaskan
        mengapa jawabannya berbeda bentuk dari kasus $\ell_\infty$.

  \item Buktikan bahwa pada pengklasifikasi linear, PGD dengan sebarang jumlah
        langkah tidak pernah menghasilkan galat yang lebih besar daripada FGSM
        dengan $\varepsilon$ yang sama. (Petunjuk: galat adalah fungsi linear dari
        $\boldsymbol{\delta}$ pada himpunan cembung; maksimum fungsi linear pada
        himpunan cembung kompak tercapai di titik ekstrem. Tunjukkan bahwa titik
        ekstrem bola $\ell_\infty$ berbentuk $\pm\varepsilon$ di setiap koordinat,
        lalu bandingkan dengan titik FGSM.)

  \item Nilai klaim ini: \emph{karena sebuah model mencapai akurasi $98$ persen pada
        data uji dan tetap $95$ persen di bawah gangguan acak sebesar
        $\varepsilon$, model itu tahan serangan adversarial.} Putuskan benar, salah,
        atau belum cukup bukti, dan sebutkan satu angka dari bab ini yang mendukung
        keputusanmu.

  \item Dari tabel lab, hitung tingkat keberhasilan serangan tak bersasaran untuk
        model $A$ pada $\varepsilon=0{,}05$, $0{,}10$, dan $0{,}15$, lalu untuk
        model $D$ pada anggaran yang sama. Nyatakan hasilnya dalam persen dan
        jelaskan mengapa selisih $A$ dan $D$ mengecil pada $\varepsilon$ besar.

  \item Tunjukkan bahwa proyeksi ke
        $\{\mathbf{y}:\lVert\mathbf{y}-\mathbf{x}\rVert_\infty\le\varepsilon\}
        \cap[0,1]^{d}$ boleh dilakukan dengan memotong dulu ke kotak
        $\varepsilon$ lalu ke $[0,1]$, atau sebaliknya, dan hasilnya sama.
        (Petunjuk: kedua himpunan adalah hasil kali selang, jadi cukup periksa
        satu koordinat.)

  \item Pada model $D$, akurasi bersih $96{,}48$ persen dan akurasi di bawah PGD 40
        langkah pada $\varepsilon=0{,}15$ adalah $62{,}65$ persen. Pada model $A$,
        angkanya $97{,}99$ dan $10{,}72$. Hitung harga ketahanan dalam poin akurasi
        bersih per poin ketahanan yang diperoleh, lalu jelaskan mengapa rasio itu
        bukan ukuran yang lengkap.

  \item Sebuah pertahanan membulatkan keluaran probabilitas menjadi $0$ atau $1$
        sebelum keputusan diambil, sehingga gradien terhadap masukan menjadi nol.
        Rancang satu protokol evaluasi untuk menguji klaim ketahanannya, dan sebutkan
        gejala mana dari daftar di bab ini yang paling mudah terlihat bila pertahanan
        itu hanya menyamarkan gradien.

  \item Diberikan dua model dengan akurasi bersih $97{,}99$ dan $98{,}49$ persen.
        Pada $\varepsilon=0{,}15$, PGD 40 langkah buatan model pertama menurunkan
        akurasi model pertama ke $10{,}72$ persen dan model kedua ke $23{,}79$
        persen. Hitung selisih keberhasilan serangan kotak putih dan serangan
        transfer, lalu jelaskan mengapa selisih itu bukan bukti bahwa model kedua
        lebih aman. (Petunjuk: serangan yang dibuat khusus untuk model kedua juga
        tersedia; lihat baris B pada keluaran lab.)

  \item Jelaskan mengapa serangan bersasaran lebih sulit daripada serangan tak
        bersasaran, lalu pakai angka $33{,}00$ persen pada $\varepsilon=0{,}15$ dan
        $99{,}50$ persen pada $\varepsilon=0{,}30$ untuk menunjukkan bahwa kesulitan
        itu bergantung pada anggaran, bukan hanya pada jenis serangan.
\end{enumerate}
\end{latihan}

\begin{untukpeneliti}
Yang terbukti di bab ini sedikit. Lemma maksimum bentuk linear adalah bukti satu baris.
Ketaksamaan linearisasi dengan sisa $M\lVert\boldsymbol{\delta}\rVert_2^2/2$ adalah
teorema nilai rata-rata bentuk integral, berlaku untuk setiap fungsi terturunkan dengan
gradien Lipschitz. Ketepatan FGSM pada regresi logistik adalah hasil eksak Goodfellow,
Shlens, dan Szegedy (ICLR 2015), bagian ``Adversarial Training of Linear Models versus
Weight Decay''. Rumusan pelana dan klaim bahwa PGD dapat diandalkan berasal dari Madry,
Makelov, Schmidt, Tsipras, dan Vladu (ICLR 2018), bagian 2 dan 3; yang terakhir itu
klaim empiris, bukan teorema, dan penulisnya sendiri menyebut masalah dalamnya tak
konkaf.

Yang lebih sering dipopulerkan daripada dibuktikan adalah klaim bahwa linearitas
menjelaskan contoh adversarial. Sebagai argumen arah ia berguna; sebagai penjelasan
besarnya kerusakan ia gagal pada percobaan kita: ramalan $0{,}0008$ nats melawan
kenyataan $4{,}862$ nats, dengan gradien yang sudah diverifikasi sampai $10^{-10}$.
Yang bekerja adalah pola tanda, dan itu sebabnya FGSM mengalahkan tanda acak
($14{,}91$ melawan $95{,}48$ persen). Klaim kedua yang dibesar-besarkan: pelatihan
adversarial memberi ketahanan terhadap penyerang dan anggaran yang dipakai saat
berlatih, dengan harga akurasi bersih, tanpa jaminan untuk serangan lain.

Batas keberlakuan: 64 ciri, 597 contoh uji, satu lapis tersembunyi, satu benih, dua
keluarga serangan $\ell_\infty$, tanpa sertifikasi. Angka di sini demonstrasi mekanisme,
bukan peringkat ketahanan. Pertanyaan terbuka: seberapa besar radius yang bisa
disertifikasi tanpa menghapus manfaat akurasi; mengapa transfer bekerja sebaik itu;
dan model ancaman apa yang masuk akal untuk sistem nyata seperti detektor intrusi, yang
menjadi pokok Bab 30.
\end{untukpeneliti}

\peta{Bab ini mengikuti Goodfellow, Shlens, dan Szegedy, \emph{Explaining and
Harnessing Adversarial Examples} (ICLR 2015), bagian 3 (``The Linear Explanation of
Adversarial Examples''), 4 (``The Fast Gradient Sign Method''), dan 5 (``Adversarial
Training of Linear Models versus Weight Decay'') untuk derivasi FGSM, ketepatan pada
regresi logistik, dan angka MNIST serta CIFAR-10. Rumusan pelana, algoritme PGD,
serta angka ketahanan MNIST dan CIFAR-10 berasal dari Madry, Makelov, Schmidt,
Tsipras, dan Vladu, \emph{Towards Deep Learning Models Resistant to Adversarial
Attacks} (ICLR 2018), bagian 2, 3, dan 4, termasuk Algorithm 1. Pembahasan penyamaran
gradien mengikuti Athalye, Carlini, dan Wagner, \emph{Obfuscated Gradients Give a
False Sense of Security} (ICML 2018), bagian 2 sampai 4, termasuk tiga jenis
penyamaran dan lima gejala yang dipakai di sini. Fenomena transferabilitas pertama
dilaporkan oleh Szegedy dan rekan, \emph{Intriguing Properties of Neural Networks}
(ICLR 2014). Seluruh angka pada tabel, lab, dan contoh hitung adalah perhitungan
sendiri dengan Python 3.14, NumPy 2.5, dan himpunan data \texttt{load\_digits} dari
scikit-learn 1.9: pelatihan Adam 350 epoch, pembagian latih-uji 1200/597 dengan benih
0, dan gradien masukan yang diverifikasi dengan beda hingga.}


\chapter{Batas Serangan: Gradien Dikaburkan dan Evasi NIDS}

Sebuah kelompok mahasiswa menambahkan satu baris pembulatan pada masukan model, lalu
melaporkan bahwa serangan mereka sendiri gagal total: tingkat keberhasilan turun dari
tujuh belas persen menjadi nol koma nol dua persen. Angka itu masuk ke ringkasan
eksekutif, lalu ke slide, lalu ke abstrak. Dua minggu kemudian seorang penguji meminta
satu hal: buktikan bahwa gradien yang dipakai penyerang masih menunjuk arah yang benar.
Ternyata tidak. Pembulatan datar di hampir semua tempat, jadi turunannya nol, jadi
penyerang melangkah nol, jadi serangan yang gagal itu sebenarnya serangan yang tidak
pernah terjadi. Ketika penguji mengganti turunannya dengan pendekatan yang masuk akal,
tingkat keberhasilan kembali ke tujuh belas persen dalam hitungan menit.

Bab ini tentang kesalahan itu. Kesalahan pertama sederhana: pembulatan menghapus gradien,
dan menghapus gradien berarti membuat penyerang buta, bukan membuat model kuat. Kesalahan
kedua lebih menarik: ada kalanya pengaburan gradien memang menaikkan biaya serangan,
sehingga sebagian ketahanan itu nyata. Membedakan keduanya menuntut percobaan, bukan
retorika.

Paruh kedua bab membawa persoalan yang sama ke tempat yang jauh lebih berisik: sistem
pengesanan pencerobohan jaringan, atau NIDS. Di sana paket harus tetap sampai, protokol
harus tetap sah, keputusan diambil dalam milidetik, dan modelnya dilatih ulang setiap
hari. Semua itu membuat serangan evasi yang cantik di atas kertas menjadi tidak praktis
di lapangan. Tetapi ``tidak praktis'' bukan ``mustahil''.

\section{Gradien yang menghilang, dan rasa aman yang lahir darinya}
\tingkat{dasar}

Bab 29 mengajarkan dua serangan yang bertumpu pada satu benda: gradien kerugian
terhadap masukan. Serangan satu langkah mengambil arah tanda gradien dan melangkah
sejauh $\varepsilon$,

\begin{equation}
x' = x + \varepsilon\,\operatorname{sign}\bigl(\nabla_x L(x,y)\bigr),
\end{equation}

sedangkan serangan banyak langkah mengulang langkah kecil dan memproyeksikan hasilnya
kembali ke bola $\ell_\infty$ berjari-jari $\varepsilon$,

\begin{equation}
x_{t+1} = \Pi_{B(x,\varepsilon)}\Bigl(x_t - \alpha\,\operatorname{sign}
\bigl(\nabla_x L(x_t,y)\bigr)\Bigr).
\end{equation}

Keduanya bekerja karena satu alasan yang sudah muncul di bab tentang linearisasi:
untuk gangguan kecil, kerugian bisa didekati oleh fungsi linear, sehingga arah yang
paling cepat menaikkan kerugian adalah arah gradien. Kalau gradien itu benar, serangan
bekerja dengan biaya satu kali mundur. Kalau gradien itu salah, penyerang menjadi buta.
Dua kalimat itu tidak setara dengan ``kalau gradien salah, model aman''. Penyerang yang
buta masih bisa menebak, dan tebakan tanpa gradien, misalnya derau acak pada paket, tetap
bisa berhasil pada sebagian kasus. Yang hilang hanyalah jalan termurah, bukan seluruh
jalan.

\index{gradien dikaburkan}
Istilah \emph{gradien dikaburkan} menunjuk keadaan ketika
gradien yang dilihat penyerang tidak memberi informasi yang cukup tentang arah yang
sebenarnya menaikkan kerugian, bukan karena modelnya kuat, melainkan karena cara model
itu dibungkus. Athalye, Carlini, Wagner dan rekan-rekan memberi nama yang lebih tegas
untuk gejala ini: rasa aman palsu. Dalam tinjauan mereka atas sembilan pertahanan dari
sebuah konferensi pembelajaran mesin, tujuh bergantung pada gradien yang dikaburkan;
enam ditembus sepenuhnya dan satu sebagian, dengan biaya serangan yang kira-kira sama
dengan serangan biasa.

\begin{untukanak}
Bayangkan seorang guru yang tidak ingin muridnya menyontek, jadi ia menutup buku
catatan dengan lakban. Murid yang jujur memang jadi kesulitan. Tetapi murid yang
kreatif tidak berhenti: ia melihat bahwa jawaban soal nomor satu selalu lebih besar
daripada nomor dua, lalu menebak seluruh lembar dari pola itu. Lakban itu tidak
menghapus jawaban; ia hanya menghalangi satu jalan menuju jawaban.

Model yang membulatkan masukannya bekerja seperti lakban itu. Angka yang masuk ke
model dipotong ke kelipatan tertentu, misalnya dua angka di belakang koma, sehingga
perubahan kecil tidak terlihat sama sekali. Program yang mencoba mencari arah
``naikkan sedikit demi sedikit'' akan menemukan bahwa tidak ada yang berubah, lalu
menyerah dan melaporkan bahwa modelnya kebal. Laporan itu salah, dan bukan karena
programnya bodoh. Program itu hanya bertanya dengan cara yang tidak cocok.

Sekarang bagian yang paling penting untuk diingat. Model yang sama masih bisa salah
kalau angkanya diubah cukup banyak sehingga melewati batas potongan. Jadi pertanyaan
yang benar bukan ``apakah serangan kecil berhasil'', melainkan ``berapa besar perubahan
yang dibutuhkan''. Bab ini mengajarkan cara mengajukan pertanyaan itu, dan cara mengenali
laporan yang tidak pernah mengajukannya.
\end{untukanak}

Mengapa pembulatan menghapus gradien? Karena pembulatan adalah fungsi tangga. Ambil
$q>0$ dan definisikan $Q_q(t)=q\,\mathrm{round}(t/q)$, kelipatan $q$ terdekat. Pada
selang $kq-q/2 \le t < kq+q/2$ fungsi ini bernilai tetap $kq$, sehingga turunannya nol
di dalam selang dan tidak terdefinisi di titik batas. Himpunan titik batas itu berukuran
nol, jadi untuk fungsi mulus apa pun $f$, komposisi $f\circ Q_q$ punya gradien nol di
hampir semua titik.

\begin{proposisi}
Misalkan $Q_q\colon\mathbb{R}^d\to\mathbb{R}^d$ bekerja koordinat demi koordinat
seperti di atas, dan misalkan $f\colon\mathbb{R}^d\to\mathbb{R}$ terdiferensiasi
kontinu. Maka $\nabla (f\circ Q_q)(x)=0$ untuk hampir setiap $x\in\mathbb{R}^d$.
\end{proposisi}

\begin{proof}
Partisi $\mathbb{R}^d$ menjadi kotak-kotak setengah terbuka bersisi $q$, yang dibatasi
hiperbidang $x_j=k_jq\pm q/2$. Pada interior setiap kotak, $Q_q$
bernilai konstan, sehingga $f\circ Q_q$ juga konstan dan gradiennya nol. Komplemen dari
gabungan semua interior itu adalah gabungan dari hiperbidang batas, yang berukuran
nol. Jadi gradiennya nol di hampir semua titik.
\end{proof}

\begin{akibat}
Serangan satu langkah yang memakai gradien eksak dari $f\circ Q_q$ menghasilkan
$x'=x$ untuk setiap $\varepsilon>0$. Tingkat keberhasilan yang terukur karena itu
tepat nol, dan ``akurasi di bawah serangan'' yang dilaporkan sama persis dengan akurasi
bersih.
\end{akibat}

Akibat ini adalah inti seluruh bab. Pernyataan itu bukan tentang kekuatan model; ia
tentang aritmetika penyerang, dan angka nol pada tabel hasil tidak mengukur ketahanan,
melainkan ketiadaan langkah. Kalau kolom ``akurasi di bawah serangan'' identik dengan
akurasi bersih sampai dua angka di belakang koma, Anda bukan melihat pertahanan yang
berhasil, melainkan penyerang yang berhenti.

\begin{bacakeras}
Dua suara. Yang pertama membawa peta, yang kedua membawa senter.

--- Lihat. Petanya sudah kubuat kabur. Semua jalan menanjak kuhapus, tidak ada panah,
tidak ada angka.
--- Kalau begitu gunungnya hilang.
--- Bukan. Gunungnya masih di sana.
--- Jadi yang hilang cuma gambarnya?
--- Yang hilang cuma gambarnya.
--- Orang yang mencium bau tanah masih bisa naik.
--- Masih bisa.
--- Orang yang melempar kerikil lalu mendengar pantulannya juga masih bisa.
--- Juga masih bisa.
--- Berarti petamu bukan tembok.
--- Petamu bukan tembok. Petamu hanya peta yang menolak menjawab.

Ulangi baris terakhir tiga kali, makin pelan: peta yang menolak menjawab bukan tembok.
Bukan tembok. Bukan tembok.
\end{bacakeras}

Athalye dan rekan-rekan membagi mekanismenya menjadi tiga keluarga. Keluarga pertama,
gradien hancur, muncul ketika bagian dari model tidak terdiferensiasi atau
terdiferensiasi tetapi tidak berguna, seperti pembulatan, penyaringan median, dan
kompresi. Keluarga kedua, gradien acak, muncul ketika model mengacak sesuatu pada setiap
panggilan, sehingga gradien satu panggilan tidak berlaku untuk panggilan berikutnya.
Keluarga ketiga, gradien hilang atau meledak, muncul ketika pertahanan menekan gradien
saat pelatihan, sehingga gradien menjadi terlalu kecil atau terlalu besar untuk dipakai
melangkah.

Ketiganya punya satu ciri bersama: ia membuat pengukuran terlihat lebih baik tanpa
membuat model lebih baik, dan karena pengukuran biasanya satu-satunya bukti yang
dilaporkan, kesalahan ini menyebar cepat. Bab ini tidak mengatakan bahwa semua pertahanan
seperti itu tidak berguna; sebagian menaikkan biaya serangan, dan biaya itu nyata. Yang
salah adalah menyimpulkan ketahanan dari pengukuran yang rusak.

\section{Tujuh cara mengaburkan gradien}
\tingkat{menengah}

Tabel berikut merangkum tujuh teknik yang paling sering muncul, keluarga mekanismenya,
dan gejala yang seharusnya membuat Anda curiga. Angka pada kolom terakhir bukan
diagnosis, melainkan pemicu pertanyaan.

\begin{center}
\small
\begin{tabular}{@{}p{0.6cm}p{4.6cm}p{2.1cm}p{6.4cm}@{}}
\hline
No & Teknik & Keluarga & Gejala yang terlihat \\
\hline
1 & Kuantisasi masukan atau keluaran & hancur & gradien nol; akurasi serangan sama dengan akurasi bersih \\
2 & Praproses tak terdiferensiasi: median, blur, kompresi & hancur & gradien nol, atau melonjak tanpa pola \\
3 & Penyaringan dan penolakan sampel & hancur & keputusan diskrit; tidak ada gradien \\
4 & Transformasi acak: ubah ukuran, bingkai, geser & acak & gradien berbeda tiap panggilan \\
5 & Pengklasifikasi acak dan ansambel acak & acak & satu sampel menunjuk arah salah \\
6 & Gradien ditekan saat pelatihan & hilang & gradien kecil padahal kerugian besar \\
7 & Pertahanan berbasis optimisasi banyak langkah & hilang & gradien meledak; serangan tidak konvergen \\
\hline
\end{tabular}
\end{center}

\index{kuantisasi}
Teknik pertama, kuantisasi, sudah dibuktikan di atas, dan efeknya eksak: bukan
kira-kira nol, tetapi nol. Teknik kedua bekerja dengan cara yang sama tetapi lebih
kotor. Penyaringan median mengganti setiap nilai dengan median lingkungannya, sehingga
fungsinya tetap datar di sebagian besar tempat, tetapi batas antar potongannya tidak lagi
berjarak seragam. Kompresi seperti JPEG membuang komponen frekuensi tinggi; turunannya
ada tetapi tidak stabil, dan gradien beda hingga bisa berubah tanda ketika langkahnya
digeser. Gejalanya khas: gradien bukan nol, melainkan sampah berskala besar.

Teknik ketiga adalah bentuk pengaburan paling kasar. Model menolak menjawab untuk masukan
yang dianggap aneh, misalnya yang jaraknya ke pusat data terlalu jauh, lalu mengembalikan
label ``tidak diketahui''. Bagi penyerang ini bukan gradien yang buruk, melainkan tidak
ada gradien: fungsinya fungsi indikator, yang turunannya nol di semua tempat kecuali di
permukaan keputusan. Pertahanan semacam ini juga menolak data baru yang sah, dan data sah
adalah alasan sistem itu dipasang.

Teknik keempat dan kelima adalah keluarga acak. Pada yang keempat, masukan diubah acak
sebelum masuk model: gambar diperbesar, digeser satu piksel, atau sebagian fiturnya
dihapus. Pada yang kelima, yang diacak adalah modelnya: satu ansambel dipilih acak pada
setiap panggilan, atau prediksi diambil dari rata-rata banyak model. Pada keduanya,
gradien satu panggilan adalah gradien satu realisasi, sedangkan yang dihadapi penyerang
adalah sebaran realisasi. Sampel gradien buruk bukan karena bias, melainkan karena
variansi: tanda dari arah yang berubah-ubah bukan arah yang baik.

\index{penyaringan acak}
Teknik keenam menyerang dari dalam. Pelatihan sengaja membuat
gradien kecil di sekitar titik data, misalnya dengan menghukum norma gradien atau
memakai aktivasi yang jenuh, sehingga permukaan kerugian sangat datar di sekitar contoh
dan curam di tempat lain. Penyerang yang melangkah sejauh $\varepsilon$ di daerah datar
tidak ke mana-mana, dan gradiennya memang kecil, tetapi bukan karena modelnya benar.
Teknik ketujuh adalah kerabatnya: masukan dibersihkan lebih dulu oleh prosedur
optimisasi banyak langkah sebelum diklasifikasikan, dan gradien dari pipa itu mudah
meledak.

Gambar~\ref{fig:permukaan} memperlihatkan bentuk permukaan yang lahir dari kuantisasi
pada dua koordinat. Setiap dataran mendatar di sana adalah satu anak tangga. Di dalam
satu dataran, kerugian tidak berubah sama sekali, sehingga penyerang yang melangkah
kecil berhenti di tempat --- tiga panah pendek pada gambar itu. Penyerang yang boleh
melangkah cukup jauh bisa menyeberang ke dataran berikutnya, dan panah putus-putus
menunjukkan penyeberangan itu.

\begin{figure}[!htbp]
\centering
\begin{tikzpicture}
\begin{axis}[
  width=0.80\linewidth, height=7.2cm,
  view={52}{24},
  xlabel={$u$}, ylabel={$v$}, zlabel={$L$},
  domain=-0.35:0.35, y domain=-0.35:0.35,
  samples=27, samples y=27,
  zmin=0, zmax=0.52,
  tick label style={font=\scriptsize},
  label style={font=\small},
  colormap={peta}{color=(layerteal!10) color=(layerteal!60)},
  colorbar=false, axis line style={layergrey}, tick style={layergrey},
]
\addplot3[surf, shader=faceted, faceted color=layergrey!40, line width=0.15pt]
  {0.1*round(10*((x-0.15)^2+(y-0.10)^2))};
\draw[-latex, accent, line width=1pt] (axis cs:0.02,0.02,0) -- (axis cs:0.10,0.06,0);
\draw[-latex, accent, line width=1pt] (axis cs:0.02,0.02,0) -- (axis cs:-0.02,0.08,0);
\draw[-latex, accent, line width=1pt] (axis cs:0.02,0.02,0) -- (axis cs:0.04,-0.05,0);
\node[circle, fill=accent, inner sep=1.5pt] at (axis cs:0.02,0.02,0) {};
\draw[-latex, layerteal, line width=1.2pt, dashed]
  (axis cs:0.02,0.02,0) -- (axis cs:0.30,0.30,0.10);
\node[font=\footnotesize, accent, anchor=north west] at (axis cs:0.09,0.05,0)
  {$\nabla L=0$};
\node[font=\footnotesize, layerteal, anchor=south west] at (axis cs:0.28,0.28,0.12)
  {BPDA};
\end{axis}
\end{tikzpicture}
\caption{Permukaan kerugian yang dibuat datar oleh kuantisasi:
$L(u,v)=\frac{1}{10}\operatorname{round}(10\,s)$ dengan
$s=(u-0{,}15)^2+(v-0{,}10)^2$.
Tiap dataran setinggi $0{,}1$ adalah satu anak tangga: di dalamnya $\nabla L=0$,
jadi tiga panah pendek dari titik merah tidak mengubah ketinggian sama sekali.
Panah putus-putus menyeberang ke anak tangga berikutnya, dan hanya penyerang yang
bersedia melangkah sejauh jarak antar-batas yang sampai ke sana.}
\label{fig:permukaan}
\end{figure}

\begin{tangan}
Bahan: satu penggaris panjang, satu uang logam, satu teman, satu meja.

Ganjal satu ujung penggaris dengan buku sehingga landainya miring, lalu letakkan uang
logam di tengahnya. Tutup matamu dan tahan uang itu dengan satu jari. Temanmu mengetuk
meja pelan-pelan, tidak beraturan, dan tugasmu menebak ke arah mana penggaris menanjak
--- hanya dari getaran yang sampai ke jarimu.

Ulangi sepuluh kali dan catat benar atau salah. Lalu minta temanmu mengetuk lebih keras,
dan ulangi sepuluh kali. Lalu minta ia berhenti mengetuk sama sekali, dan ulangi lagi.

Yang harus kamu lihat: ketukan keras tidak membuat tebakanmu benar. Ketukan yang berhenti
membuat tebakanmu murni tebakan, sekitar separuh benar, tidak lebih baik daripada melempar
koin. Getaran memberitahu bahwa ada sesuatu yang bergerak; ia tidak memberitahu arah.
Itulah bedanya gradien yang buruk dan gradien yang tidak ada.
\end{tangan}

Ketujuh teknik itu menghasilkan tujuh gejala. Lima pertama adalah daftar Athalye dan
rekan-rekan; dua terakhir kami tambahkan dari tipologi yang sama.

\begin{enumerate}
  \item \textbf{Serangan tidak dibatasi.} Laporan mengizinkan gangguan sebesar apa pun,
        lalu menyimpulkan ketahanan dari serangan yang tidak memperhatikan batas.
  \item \textbf{Serangan memakai satu sampel acak.} Pada pertahanan yang mengacak,
        penyerang menarik satu realisasi lalu melaporkan kegagalannya. Yang diukur adalah
        satu lemparan dadu, bukan peluang menang.
  \item \textbf{Serangan hanya satu langkah.} Serangan satu langkah mudah terhenti oleh
        permukaan kerugian yang tidak rata, dan ketahanan yang disimpulkannya hilang
        begitu langkahnya diulang.
  \item \textbf{Menaikkan herotan tidak menaikkan keberhasilan.} Bila anggaran diperbesar
        tetapi keberhasilan tidak bergerak, yang rusak biasanya gradiennya.
  \item \textbf{Gradien dilaporkan nol atau NaN.} Ini gejala paling telanjang: gradien
        yang seluruhnya nol pada model tak konstan menandakan potongan tak
        terdiferensiasi, bukan permukaan yang rata.
  \item \textbf{Akurasi di bawah serangan sama dengan akurasi bersih.} Tambahan kami.
        Kesamaan sampai beberapa angka di belakang koma pada beberapa $\varepsilon$
        berbeda adalah tanda serangan yang tidak bergerak.
  \item \textbf{Tidak ada serangan adaptif.} Tambahan kami. Pertahanan diuji hanya
        dengan serangan yang dirancang sebelum pertahanan itu ada, sehingga belum diuji
        oleh penyerang yang tahu dirinya sedang diuji.
\end{enumerate}

\kamus{gradien dikaburkan}{obfuscated gradient}{gradien yang diberikan kepada penyerang
tidak menunjuk arah yang benar, biasanya karena ada potongan tak terdiferensiasi atau
acak di dalam pipa, bukan karena modelnya kuat}

\kamus{kuantisasi}{quantization}{penggantian setiap nilai dengan kelipatan terdekat dari
suatu jarak tetap; hasilnya fungsi tangga yang turunannya nol di hampir semua tempat}

\kamus{serangan adaptif}{adaptive attack}{serangan yang dirancang dengan mengetahui
pertahanan yang sedang diuji, termasuk mengetahui bahwa gradiennya mungkin rusak}

Daftar gejala itu berguna untuk menolak klaim lemah, tetapi tidak untuk menerima klaim
kuat: pertahanan bisa lolos ketujuhnya dan tetap rapuh, dan sebaliknya bisa memicu satu
gejala dan tetap menaikkan biaya serangan.


\section{Menembus: BPDA, EOT, dan reparameterisasi}
\tingkat{lanjut}

Kalau gradien yang tersedia salah, penyerang punya tiga pilihan. Ganti gradien yang salah
dengan gradien fungsi yang mendekatinya; rata-ratakan gradien atas keacakan yang membuatnya
berubah-ubah; atau pindahkan pekerjaan ke ruang lain, di mana pertahanannya tidak
mengaburkan apa pun. Ketiganya punya nama, dan ketiganya bekerja karena alasan berbeda.

\index{bpda}
Pilihan pertama disebut \emph{backward pass differentiable approximation}, disingkat
BPDA. Idenya sesederhana bunyinya: pada perhitungan maju pakai fungsi yang sebenarnya,
termasuk pembulatan atau penyaringan; pada perhitungan mundur lewati potongan itu
seolah-olah ia identitas. Bila model terbungkus adalah $h(x)=f\bigl(g(x)\bigr)$,
penyerang memakai

\begin{equation}
\widehat{\nabla}_x h(x) \;=\; \nabla f\bigl(g(x)\bigr),
\end{equation}

yaitu gradien $f$ yang dievaluasi di titik setelah dibungkus, bukan gradien dari
komposisi. Untuk pembulatan, ini berarti penyerang berhenti bertanya ``berapa turunan
pembulatan'' dan mulai bertanya ``kalau angkanya sudah dibulatkan, ke arah mana $f$
paling cepat naik''.

Mengapa pendekatan ini masuk akal? Karena kesalahannya dibatasi oleh seberapa jauh $g$
menggeser titik, bukan oleh seberapa aneh bentuk $g$.

\begin{proposisi}
Misalkan $f$ terdiferensiasi kontinu dengan gradien yang Lipschitz, yaitu
$\lVert \nabla f(u)-\nabla f(v)\rVert \le L\lVert u-v\rVert$ untuk semua $u,v$, dan
misalkan $g$ memenuhi $\lVert g(x)-x\rVert\le\delta$. Maka
\begin{equation}
\bigl\lVert \nabla f\bigl(g(x)\bigr)-\nabla f(x)\bigr\rVert \;\le\; L\,\delta .
\end{equation}
\end{proposisi}

\begin{proof}
Tulis $v=g(x)-x$ dan gunakan teorema nilai rata-rata dalam bentuk integral,
\begin{equation}
\nabla f(x+v)-\nabla f(x) = \int_0^1 H_f(x+tv)\,v\,dt,
\end{equation}
dengan $H_f$ matriks Hessian. Karena norma spektral $H_f$ dibatasi $L$,
\begin{equation}
\bigl\lVert \nabla f(x+v)-\nabla f(x)\bigr\rVert
\le \int_0^1 \lVert H_f(x+tv)\rVert\,\lVert v\rVert\,dt
\le L\lVert v\rVert \le L\delta .
\end{equation}
\end{proof}

Proposisi ini menjelaskan kapan BPDA boleh dipercaya. Kalau pembulatan hanya menggeser
masukan sejauh $q/2$ dan $f$ cukup mulus pada skala itu, galat gradiennya kecil dan
langkah serangan tetap menuju arah yang benar. Kalau $q$ besar sehingga penggeseran itu
sebanding dengan skala lekukan $f$, batas $L\delta$ menjadi besar dan tidak ada jaminan.
Proposisi ini tidak mengatakan $\widehat{\nabla}$ adalah gradien $h$; gradien itu tidak
ada. Ia hanya mengatakan bahwa arah yang dipakai dekat dengan arah yang benar untuk $f$
di titik yang sudah dibungkus.

Pilihan kedua adalah \emph{expectation over transformation}, disingkat EOT. Untuk
pertahanan acak, penyerang tidak menghadapi satu fungsi melainkan sebaran fungsi.
Yang masuk akal dimaksimalkan bukan kerugian pada satu realisasi, melainkan harapan
kerugian atas realisasi,

\begin{equation}
\overline{L}(x) \;=\; \mathbb{E}_{T\sim\mathcal{T}}\; L\bigl(f_T(x),y\bigr),
\end{equation}

dengan $T$ menyatakan keacakan yang dipakai pertahanan, misalnya topeng fitur atau
parameter transformasi.

\index{eot}
\begin{lemma}
Bila $\mathcal{T}$ berhingga dan $\nabla_x L\bigl(f_T(x),y\bigr)$ terbatas seragam
dalam $T$, maka
\begin{equation}
\nabla_x \overline{L}(x) \;=\; \mathbb{E}_{T}\;\nabla_x L\bigl(f_T(x),y\bigr).
\end{equation}
Selain itu, rata-rata $m$ sampel bebas,
$\frac{1}{m}\sum_{i=1}^{m}\nabla_x L\bigl(f_{T_i}(x),y\bigr)$, adalah penaksir tak bias
dengan variansi yang mengecil sebanding $1/m$.
\end{lemma}

\begin{proof}
Karena $\mathcal{T}$ berhingga, harapan adalah jumlah berhingga, dan turunan boleh
masuk ke dalam jumlah karena turunan adalah operator linear. Penaksir rata-rata punya
harapan yang sama dengan tiap sukunya, jadi tak bias, dan variansinya adalah variansi
satu suku dibagi $m$ karena sampelnya bebas.
\end{proof}

Lemma ini menjawab keberatan yang sering muncul: kalau pertahanannya acak, bukankah
serangan apa pun hanya kebetulan? Tidak, selama penyerang boleh memanggil model
berkali-kali. Keacakan menaikkan biaya, karena penyerang perlu $m$ panggilan per langkah,
tetapi tidak mengubah arah yang dituju; yang berubah hanya ketelitian taksiran. Ini
contoh pertama dari pola yang berulang sepanjang bab: pertahanan yang menaikkan biaya itu
nyata, sedangkan yang hanya menyulitkan pengukuran itu semu.

Pilihan ketiga adalah reparameterisasi. Alih-alih menyerang $x$, penyerang menyerang
variabel yang menghasilkan $x$ melalui transformasi yang tidak dikaburkan. Kalau
pertahanan membungkus model dengan fungsi $g$ yang punya invers, penyerang bekerja di
ruang $z$ dengan $x=g^{-1}(z)$ dan mengoptimalkan $z$. Ini berguna ketika $g$ adalah
rotasi atau transformasi linier lain yang tidak menghilangkan informasi, dan tidak
berguna ketika $g$ menghilangkan informasi: tidak ada $z$ yang bisa memulihkan angka
yang sudah dibuang.

Dalam praktik, ketiga pilihan itu dipakai bersama, ditambah dua kebiasaan. Pertama: mulai
dari beberapa titik awal acak di dalam bola $\varepsilon$, karena permukaan kerugian yang
tidak cembung membuat satu titik awal mudah terjebak. Kedua: jangan berhenti pada satu
nilai $\varepsilon$. Ketahanan yang diukur pada satu anggaran adalah ketahanan yang tidak
diukur.

\begin{resep}
Empat langkah menguji klaim pertahanan, dalam urutan ini.
\begin{enumerate}
  \item \textbf{Tulis anggaran lebih dulu.} Tetapkan $\varepsilon$ atau batas byte,
        jumlah kueri, dan waktu. Angka ini tidak boleh berubah setelah hasil pertama
        keluar.
  \item \textbf{Serang lewat transfer, tanpa kueri.} Latih model pengganti pada data
        yang kamu punya, buat contoh adversarial di sana, lalu kirim ke model korban.
        Langkah ini mengukur apakah pertahanan itu hanya menghalangi satu model.
  \item \textbf{Belanjakan sisa kueri secara adaptif.} Baru sekarang pakai BPDA atau
        EOT pada korban, dan catat kueri per langkah, bukan hanya hasil akhirnya.
  \item \textbf{Laporkan keduanya.} Keberhasilan transfer, keberhasilan adaptif, dan
        biaya masing-masing, selalu bersama akurasi bersih.
\end{enumerate}
Peringatan urutan: langkah 2 harus selesai sebelum langkah 3. Kalau kamu mengukur
transfer setelah melihat hasil adaptif, kamu sudah tahu contoh mana yang berhasil,
dan angka transfernya bocor ke depan.
\end{resep}

\kamus{bpda}{backward pass differentiable approximation}{teknik mengganti gradien
potongan yang tidak terdiferensiasi dengan gradien fungsi di titik sesudah potongan itu,
seolah potongan tersebut identitas}

\kamus{eot}{expectation over transformation}{serangan yang memaksimalkan harapan
kerugian atas keacakan pertahanan, biasanya dengan merata-ratakan gradien dari beberapa
realisasi}

\begin{untukdosen}
Bagian ini biasanya diajarkan dalam urutan yang salah, yaitu definisi BPDA lebih dulu.
Urutan yang lebih baik adalah sebaliknya: tunjukkan model yang akurasinya tidak bergerak
sedikit pun ketika diserang, minta mahasiswa menebak mengapa, lalu tunjukkan gradiennya
nol. Hampir semua akan menebak ``modelnya kuat'', dan keheningan sesudahnya adalah momen
belajar tersendiri.

Dua kesalahan perlu diantisipasi. Pertama, menyimpulkan bahwa BPDA ``menghapus''
pertahanan; yang benar, BPDA mengembalikan akses ke arah yang kira-kira benar, dan
kegunaannya bergantung pada besarnya $\delta$ relatif terhadap skala lekukan model.
Kedua, menyimpulkan bahwa pertahanan acak tidak berguna karena EOT ada; yang benar,
pertahanan acak menaikkan biaya per langkah secara linear terhadap $m$, sehingga
pertahanan itu mengubah anggaran penyerang, bukan arah serangan. Latihan di akhir bab
memaksa mahasiswa menghitung biaya itu, bukan menyebutnya.
\end{untukdosen}

\section{Contoh hitung: satu titik pada model mainan}
\tingkat{menengah}

\begin{untukanak}
Seorang penjaga toko punya kebiasaan aneh: setiap harga dibulatkan ke rupiah terdekat
sebelum dicatat. Kalau ada pembeli yang diam-diam menggeser harga sebesar satu sen,
catatan penjaga tidak berubah sama sekali. Pembeli itu bisa menyimpulkan bahwa tidak
ada gunanya menggeser harga, lalu berhenti mencoba. Kesimpulannya salah, tetapi
alasannya masuk akal: yang ia ukur memang tidak bergerak.

Sekarang bayangkan pembeli lain. Ia tidak bertanya apakah catatan berubah karena
geseran satu sen; ia bertanya ke arah mana catatan paling cepat naik, lalu menggeser
harga cukup jauh sehingga pembulatannya berubah. Catatan penjaga pun berubah, dan
pembeli itu berhasil tanpa melanggar aturan apa pun. Yang membedakan keduanya bukan
kepandaian, melainkan pertanyaan yang diajukan.

Model yang memakai pembulatan bekerja seperti penjaga itu. Program yang hanya mencoba
geseran kecil akan berhenti dan melaporkan bahwa modelnya aman; program yang mengerti
pembulatan akan melangkah lebih jauh. Jadi ketika membaca laporan bahwa sebuah model
tidak bisa dikalahkan, tanyakan dulu: penyerangnya berhenti karena kalah, atau karena ia
hanya mencoba satu cara?
\end{untukanak}

Sebelum percaya pada tabel, kita hitung satu titik dengan tangan. Modelnya dibuat
sesederhana mungkin: delapan fitur dengan bobot semuanya satu, jadi skornya adalah
jumlah kedelapan fitur. Datanya berasal dari dua sebaran normal berpusat $-0{,}5$ dan
$+0{,}5$, dengan akurasi bersih sekitar $92$ persen. Pertahanannya adalah pembulatan ke
kelipatan $q=0{,}05$ pada setiap fitur, ditambah penyaringan acak berpeluang hapus $r$
yang mengalikan sisanya dengan $1/(1-r)$. Anggaran penyerang $\varepsilon=0{,}25$ pada
norma $\ell_\infty$.

Ambil satu contoh yang benar berlabel $+1$,

\begin{equation}
x_0 = (0{,}30;\ 0{,}20;\ 0{,}15;\ 0{,}10;\ 0{,}05;\ 0{,}00;\ -0{,}10;\ -0{,}08).
\end{equation}

Jumlahnya $0{,}62$, jadi model menjawab $+1$ dengan margin yang tipis. Sesudah
pembulatan, vektor yang masuk ke pengklasifikasi adalah
$(0{,}30;0{,}20;0{,}15;0{,}10;0{,}05;0{,}00;-0{,}10;-0{,}10)$ dan skornya $0{,}60$.
Perhatikan bahwa satu koordinat bergeser dari $-0{,}08$ ke $-0{,}10$: itulah kerja
pembulatan.

Sekarang bandingkan dua penyerang. Penyerang naif menghitung turunan fungsi yang
sebenarnya, yaitu turunan dari komposisi yang memuat pembulatan. Karena pembulatan
datar di sekitar $x_0$, turunannya nol di kedelapan koordinat, sehingga langkahnya
$x_0 + \varepsilon\cdot 0 = x_0$. Skornya tetap $0{,}60$ dan labelnya tetap $+1$.
Laporan penyerang itu akan berbunyi: serangan gagal, akurasi bertahan seratus persen
dari akurasi bersih.

Penyerang kedua memakai BPDA. Ia tidak bertanya tentang turunan pembulatan; ia
mengevaluasi gradien pengklasifikasi di titik yang sudah dibulatkan, yaitu vektor
$(1,1,\dots,1)$, lalu melangkah berlawanan arah label sejauh $\varepsilon$ pada setiap
koordinat,

\begin{equation}
x_1 = x_0 - 0{,}25\cdot(1,1,\dots,1)
= (0{,}05;\ -0{,}05;\ -0{,}10;\ -0{,}15;\ -0{,}20;\ -0{,}25;\ -0{,}35;\ -0{,}33).
\end{equation}

Sesudah pembulatan, vektor itu menjadi
\[
(0{,}05;\ -0{,}05;\ -0{,}10;\ -0{,}15;\ -0{,}20;\ -0{,}25;\ -0{,}35;\ -0{,}35)
\]
dengan skor $-1{,}40$. Labelnya berbalik menjadi $-1$. Satu langkah, biaya satu kali
perhitungan gradien, dan model yang tadi dilaporkan kebal ternyata salah dengan margin
yang tidak tipis.

Contoh ini menunjukkan sesuatu yang lebih tajam daripada ``BPDA bekerja''. Penyerang
kedua tidak memakai informasi yang tidak dimiliki penyerang pertama: keduanya tahu
arsitektur, bobot, dan $q$. Yang berbeda hanya satu keputusan, yaitu menghormati potongan
tak terdiferensiasi atau melewatinya. Jadi ``ketahanan'' yang dilaporkan penyerang pertama
bukan properti model, melainkan properti notasi penyerang itu sendiri.

\begin{contoh}
Pada model yang sama, dengan $q=0{,}05$, $r=0{,}5$, dan seed tetap, gradien beda
hingga dengan langkah $h=10^{-5}$ bernilai nol pada $15\,998$ dari $16\,000$ entri;
dua entri sisanya tepat $q/h=5\,000$, karena titiknya kebetulan berada di batas anak
tangga. Norma gradien beda hingga yang diambil nilai tengahnya nol, sedangkan norma
gradien BPDA untuk satu topeng yang menyisakan empat fitur adalah $4$. Tetapi arah satu
topeng masih jauh dari arah bobot, dengan kosinus antara $0{,}35$ dan $0{,}94$ atas
$200$ topeng. Rata-rata gradien atas $200$ topeng memberi vektor bernorma $2{,}738$ dan
kosinus $0{,}9982$ terhadap vektor bobot.
\end{contoh}

Angka itu menggambarkan tiga hal. Pertama, gradien yang hancur bukan hanya kecil; ia bisa
juga besar dan salah, dan dua entri yang melonjak akan mendominasi setiap langkah
sehingga serangan naif gagal dengan cara yang tampak acak. Kedua, BPDA memulihkan arah
dengan hampir sempurna: kosinus $0{,}9982$ berarti sudutnya sekitar tiga derajat. Ketiga,
norma gradien BPDA untuk satu topeng, yaitu $4$, berbeda dari norma harapannya, $2{,}738$;
perbedaan skala itu tidak merusak serangan yang memakai tanda, tetapi merusak yang memakai
besaran gradien, dan itu sebabnya EOT tetap perlu.

\kertas{Hitung ulang contoh di atas tanpa melihat teks. Jumlahkan $x_0$, bulatkan ke
kelipatan $0{,}05$, lalu jumlahkan lagi; cocokkan bahwa hasilnya $0{,}60$. Kurangi
setiap koordinat dengan $0{,}25$, bulatkan lagi, jumlahkan, dan cocokkan bahwa hasilnya
$-1{,}40$. Sesudah itu jawab: berapa nilai $q$ terkecil yang membuat serangan BPDA satu
langkah ini gagal membalik label, dan mengapa nilai itu bergantung pada margin awal
$0{,}60$ dan bukan pada $\varepsilon$ saja?}

\section{Lab komputasi: akurasi di tiap tingkat pengaburan}
\tingkat{lanjut}

\begin{figure}[!htb]
\centering
\begin{tikzpicture}
\begin{axis}[
  width=0.86\linewidth, height=5.6cm,
  xlabel={anggaran $\varepsilon$ pada norma $\ell_\infty$},
  ylabel={keberhasilan serangan (\%)},
  xmin=0.03, xmax=0.37, ymin=-0.6, ymax=30,
  xtick={0.05,0.10,0.15,0.20,0.25,0.30,0.35},
  xticklabels={0{,}05,0{,}10,0{,}15,0{,}20,0{,}25,0{,}30,0{,}35},
  ytick={0,5,10,15,20,25,30},
  tick label style={font=\footnotesize},
  label style={font=\small},
  legend style={font=\scriptsize, at={(0.02,0.98)}, anchor=north west,
                draw=layergrey!35, fill=white, fill opacity=0.92, text opacity=1},
  grid=major, grid style={layergrey!15},
  axis line style={layergrey},
]
\addplot[color=layerindigo, mark=*, mark size=1.5pt, line width=1pt]
  coordinates {(0.05,2.38) (0.10,5.27) (0.15,8.66) (0.20,12.84)
               (0.25,17.40) (0.30,22.32) (0.35,27.71)};
\addlegendentry{tanpa pembulatan, penyerang BPDA}
\addplot[color=accent, mark=square*, mark size=1.4pt, line width=1pt, dashed]
  coordinates {(0.05,2.34) (0.10,5.33) (0.15,8.72) (0.20,12.92)
               (0.25,17.48) (0.30,22.46) (0.35,27.79)};
\addlegendentry{pembulatan $q=0{,}05$, penyerang BPDA}
\addplot[color=layergrey, mark=triangle*, mark size=1.6pt, line width=0.9pt,
         dotted]
  coordinates {(0.05,0.00) (0.10,0.01) (0.15,0.01) (0.20,0.01)
               (0.25,0.01) (0.30,0.02) (0.35,0.02)};
\addlegendentry{pembulatan $q=0{,}05$, penyerang naif}
\end{axis}
\end{tikzpicture}
\caption{Keberhasilan serangan terhadap $\varepsilon$ pada model mainan delapan fitur,
$2\,000$ contoh, rata-rata lima seed. Pada $\varepsilon=0{,}25$ ketiga kurva memberi
$17{,}40$, $17{,}48$, dan $0{,}01$ persen; pada $\varepsilon=0{,}35$ memberi
$27{,}71$, $27{,}79$, dan $0{,}02$ persen. Dua kurva pertama berimpit di dalam
simpangan antar-seed, yang besarnya $0{,}35$ sampai $1{,}08$ poin persen.}
\label{fig:kurva-epsilon}
\end{figure}

Sekarang kita ulang percobaan itu pada dua ribu contoh, lima seed, dan beberapa tingkat
pengaburan, hanya dengan \texttt{numpy}. Yang perlu diperhatikan bukan kode itu sendiri,
melainkan satu keputusan di dalamnya: fungsi \texttt{fgsm} menghitung gradien dengan beda
hingga pada fungsi yang sudah dibungkus, sedangkan \texttt{pgd} memakai gradien
pengklasifikasi di titik sesudah pembungkusan. Hasilnya, disapu pada tujuh nilai
$\varepsilon$, ada di Gambar~\ref{fig:kurva-epsilon}.

\begin{lab}
\begin{verbatim}
import numpy as np

D, N, EPS = 8, 2000, 0.25
W = np.ones(D)

def data(rng, n, mu=0.5):
    y = 2 * rng.integers(0, 2, n) - 1
    return rng.normal(mu * y[:, None], 1.0, (n, D)), y

def quant(x, q):
    return x if q == 0 else np.round(x / q) * q

def mask(shape, r, rng):
    if r == 0:
        return np.ones(shape)
    return (rng.random(shape) >= r) / (1 - r)

def logit(x, q, m):
    return (quant(x, q) * m) @ W

def fgsm(X, y, q, r, rng):
    m, h = mask(X.shape, r, rng), 1e-5
    G = np.zeros_like(X)
    for j in range(D):
        e = np.zeros(D)
        e[j] = h
        G[:, j] = (logit(X + e, q, m) - logit(X - e, q, m)) / (2 * h)
    return X - EPS * y[:, None] * np.sign(G)

def pgd(X, y, q, r, rng, m_eot=1, steps=40):
    Xa, a = X.copy(), EPS / 8
    for _ in range(steps):
        G = sum(W * mask(X.shape, r, rng) for _ in range(m_eot))
        Xa = np.clip(Xa - a * y[:, None] * np.sign(G / m_eot),
                     X - EPS, X + EPS)
    return Xa

def run(seed, q, r):
    rng = np.random.default_rng(seed)
    X, y = data(rng, N)
    def pred(Z):
        return np.sign(logit(Z, q, mask(Z.shape, r, rng)))
    ok = pred(X) == y
    def succ(Z):
        return 100 * np.mean(pred(Z)[ok] != y[ok])
    return (100 * ok.mean(), succ(fgsm(X, y, q, r, rng)),
            succ(pgd(X, y, q, r, rng)), succ(pgd(X, y, q, r, rng, 20)))

print("  q     r   bersih   naif   bpda1   eot")
for q, r in ((0.00, 0.0), (0.01, 0.0), (0.05, 0.0),
             (0.05, 0.5), (0.20, 0.0)):
    M = np.array([run(s, q, r) for s in (1, 2, 3, 4, 5)]).mean(0)
    print("%5.2f %5.2f  %6.2f  %5.2f  %5.2f  %5.2f"
          % (q, r, M[0], M[1], M[2], M[3]))
\end{verbatim}

Keluaran program itu, apa adanya:

\begin{verbatim}
  q     r   bersih   naif   bpda1   eot
 0.00  0.00   92.02  17.40  17.40  17.40
 0.01  0.00   92.00   0.02  17.40  17.40
 0.05  0.00   91.97   0.01  17.48  17.48
 0.05  0.50   82.87  12.16  25.48  25.67
 0.20  0.00   91.77   0.00  17.65  17.65
\end{verbatim}
\end{lab}

Bacalah tabel itu baris demi baris. Baris pertama adalah garis dasar tanpa pengaburan:
akurasi bersih $92{,}02$ persen, keberhasilan serangan satu langkah $17{,}40$ persen.
Angka itu kecil karena $\varepsilon=0{,}25$ memang kecil dibanding sebaran datanya,
bukan karena serangannya lemah.

Baris kedua adalah inti bab ini. Dengan $q=0{,}01$, akurasi bersih hampir tidak berubah,
$92{,}00$ persen, tetapi keberhasilan naif jatuh dari $17{,}40$ menjadi $0{,}02$ persen,
sekitar $870$ kali lebih kecil. Penyerang yang memakai BPDA mendapat $17{,}40$ persen,
persis garis dasar. Pertahanan itu tidak mengubah apa pun yang bisa diamati penyerang
yang menghitung dengan benar; ia hanya mengubah apa yang dilihat penyerang yang salah
menghitung.

Baris ketiga mengulang hal yang sama pada $q=0{,}05$: naif $0{,}01$ persen, BPDA
$17{,}48$ persen, akurasi bersih $91{,}97$ persen. Menaikkan $q$ lima kali lipat tidak
mengubah keberhasilan BPDA, tetapi membuat laporan naif terlihat lebih baik. Itulah
gejala keempat dari daftar di bagian dua, dalam bentuk angka.

Baris keempat menambahkan penyaringan acak $r=0{,}5$, dan dua hal terjadi sekaligus.
Akurasi bersih turun dari $91{,}97$ ke $82{,}87$ persen, karena membuang separuh fitur dan
mengalikan sisanya dengan dua menaikkan variansi skor. Keberhasilan BPDA naik ke $25{,}48$
persen, di atas garis dasar, bukan karena pertahanannya lebih buruk melainkan karena contoh
yang masih benar setelah pengaburan punya margin lebih kecil.

Baris kelima menunjukkan bahwa kuantisasi kasar tidak menyelamatkan model. Dengan
$q=0{,}20$, sebanding dengan $\varepsilon$, keberhasilan BPDA $17{,}65$ persen, sedikit di
atas garis dasar: titik yang berada dekat batas anak tangga hanya perlu dorongan kecil untuk
menyeberang.

Variansi antar seed kecil tetapi tidak nol: pada $q=0{,}05$ dengan $r=0$ simpangan bakunya
sekitar $0{,}6$ poin persen, pada baris dengan penyaringan acak sekitar $1{,}4$ poin. Galat
baku satu proporsi pada $2\,000$ contoh dan tingkat $0{,}17$ kira-kira $0{,}8$ poin, jadi
kedua angka itu sejalan dengan pengambilan sampel biasa.

Ada empat hal yang tidak diukur percobaan ini. Pertama, modelnya linier, sehingga gradien
BPDA hampir sempurna dan kosinus $0{,}9982$ terlalu bagus untuk model nyata. Kedua,
penyerang mengetahui $q$ dan $r$. Ketiga, hanya satu nilai $\varepsilon$ yang diuji.
Keempat, tidak ada pencarian hyperparameter serangan. Semua itu membuat angka di atas
menjadi batas bawah keberhasilan serangan, bukan batas atas.


Kekurangan ketiga itu bisa ditutup, dan hasilnya ada di Gambar~\ref{fig:kurva-epsilon}.
Program yang sama dijalankan pada tujuh nilai $\varepsilon$ dari $0{,}05$ sampai
$0{,}35$, lima seed, tanpa penyaringan acak. Kurva model tanpa pembulatan dan kurva
model dengan pembulatan $q=0{,}05$ hampir berimpit: pada $\varepsilon=0{,}25$ keduanya
memberi $17{,}40$ dan $17{,}48$ persen, dan pada $\varepsilon=0{,}35$ keduanya memberi
$27{,}71$ dan $27{,}79$ persen. Kurva penyerang naif tetap di bawah $0{,}02$ persen
sepanjang rentang itu. Pertahanan ini karena itu tidak menggeser kurva dan tidak
mematahkannya; ia hanya menambah kurva ketiga yang datar.


\section{Studi kasus: NIDS, dan mengapa evasinya tetap tidak praktis}
\tingkat{lanjut}

Sampai di sini kita berurusan dengan vektor fitur yang boleh diubah sesuka hati. Sistem
pengesanan pencerobohan jaringan, atau NIDS, adalah dunia yang berbeda.
\index{nids}
Sebuah NIDS membaca aliran paket, mengekstrak ciri seperti durasi aliran,
jumlah paket, jumlah byte, bendera TCP, dan waktu antar-kedatangan, lalu memutuskan
apakah aliran itu serangan. Keputusan harus diambil cepat, sering dalam milidetik, untuk
ribuan aliran per detik.

Gambar~\ref{fig:rantai} menaruh rantai itu pada satu baris. Yang perlu diperhatikan
bukan kotaknya, melainkan tanda merah di antara rekonstruksi aliran dan ekstraksi ciri.
Di situlah penyerang bekerja: ia tidak menyentuh model, ia menyentuh paket, dan paket
itulah yang menentukan nilai ciri yang masuk ke model. Semua yang ada di kanan tanda itu
hanya melihat akibat; semua yang ada di kiri adalah sebab yang bisa ia kendalikan, dengan
batasan protokol yang dibahas sesudah ini.

\begin{figure}[!htbp]
\centering
\begin{tikzpicture}[font=\footnotesize,
  tahap/.style={draw=layergrey!70, fill=layergrey!6, rounded corners=1.5pt,
    align=center, inner sep=3pt, minimum height=1.1cm, text width=1.95cm},
  panah/.style={-latex, layergrey, line width=0.8pt}]
\node[tahap] (n1) at (0,0) {paket\\di kabel};
\node[tahap] (n2) at (2.50,0) {rekonstruksi\\aliran};
\node[tahap] (n3) at (5.00,0) {ekstraksi\\ciri};
\node[tahap] (n4) at (7.50,0) {penyetelan\\skala};
\node[tahap] (n5) at (10.00,0) {model\\NIDS};
\node[tahap] (n6) at (12.50,0) {ambang\\keputusan};
\draw[panah] (n1) -- (n2);
\draw[panah] (n2) -- (n3);
\draw[panah] (n3) -- (n4);
\draw[panah] (n4) -- (n5);
\draw[panah] (n5) -- (n6);
\node[circle, fill=accent, inner sep=2pt] (sisip) at (3.75,0) {};
\draw[accent, line width=0.8pt] (sisip) -- ++(0,0.62);
\node[accent, anchor=south, font=\scriptsize] at (3.75,0.66)
  {titik sisipan: waktu dan ukuran paket};
\node[layergrey, anchor=north, font=\scriptsize, align=center] at (12.5,-0.62)
  {peringatan atau blokir};
\draw[-latex, layergrey, line width=0.6pt] (n6.south) -- (12.5,-0.5);
\end{tikzpicture}
\caption{Rantai pemrosesan NIDS, dari kabel sampai keputusan. Tanda merah adalah satu
titik sisipan yang bisa dipakai penyerang: ciri seperti jumlah paket, jumlah byte, dan
waktu antar-kedatangan seluruhnya ditentukan di kiri tanda itu, sedangkan model di
kanannya hanya membaca akibatnya. Pertahanan yang menambahkan pengaburan di kanan tanda
tidak menyentuh sebab, dan itulah batas yang dibahas di bagian ini.}
\label{fig:rantai}
\end{figure}

\kamus{nids}{network intrusion detection system}{sistem yang memantau lalu lintas
jaringan dan memutuskan apakah suatu aliran paket merupakan serangan}

\kamus{ruang masalah}{problem space}{ruang objek nyata yang boleh diubah penyerang,
misalnya paket dan waktu, berbeda dari ruang fitur tempat model bekerja}

Kesulitan pertama adalah jarak antara dua ruang. Serangan biasanya dirumuskan di ruang
fitur, sedangkan penyerang mengirim paket: setiap perubahan fitur harus diterjemahkan
menjadi perubahan paket yang tetap sah menurut protokol dan tetap menjalankan fungsi
jahatnya, dan terjemahan itu tidak selalu ada. Beberapa fitur bahkan tidak bisa diubah
pengirim, karena jumlah paket dan byte arah balik ditentukan sistem penerima.

Kesulitan kedua adalah korelasi antar fitur. Durasi berkaitan dengan jumlah paket dan
waktu antar-kedatangan, dan bendera TCP berkaitan dengan protokol transpor, sehingga
penyerang tidak bisa menaikkan durasi sambil menurunkan kedua waktu antar-kedatangan;
fitur diskret seperti nomor protokol dan panjang paket menambah batas.

Kesulitan ketiga adalah waktu. Serangan yang membutuhkan sepuluh ribu kueri berarti lalu
lintas nyata yang lewat dan tercatat, dan jejak adalah hal pertama yang dilihat analis.
Biaya nyata adalah pertahanan nyata; ketiadaan gradien bukan.

Angka dari literatur membuat jarak itu konkret. Pada data CSE-CIC-IDS2018 dengan lima
puluh fitur tersaring, tiga model dilatih dan diuji pada dua puluh persen data, lalu
diserang dengan FGSM, PGD, dan BIM. Sebelum serangan ketiganya nyaris sempurna: akurasi
$0{,}997$ untuk ANN, $0{,}998$ untuk SVM, dan $0{,}997$ untuk CNN. Pada hari serangan,
dengan pengetahuan penuh, FGSM menjatuhkan ANN ke $0{,}756$ dengan F1 $0{,}677$, CNN ke
$0{,}842$ dengan F1 $0{,}863$, dan SVM ke $0{,}150$ dengan F1 $0{,}000$. PGD sedikit
lebih kuat pada CNN, $0{,}584$ dengan F1 $0{,}689$, dan BIM memberi $0{,}688$ dengan F1
$0{,}751$.

Angka itu terlihat seperti kemenangan penyerang, dan pada model statis memang demikian;
tetapi percobaan yang sama melanjutkan dengan model yang dilatih ulang setiap hari.
Dengan pengetahuan tentang model berumur satu hari, SVM kembali ke $0{,}998$, ANN ke
$0{,}938$, dan CNN ke $0{,}957$ dengan F1 $0{,}959$, dan pada hari kedua angkanya antara
$0{,}966$ dan $0{,}998$. Pelatihan ulang itu tidak memakai contoh adversarial, jadi
pemulihannya bukan hasil pertahanan adversarial melainkan akibat fakta bahwa gradien yang
dihitung pada model kemarin tidak lagi menunjuk arah yang benar.

Inilah versi NIDS dari seluruh bab: pengaburan gradien yang disengaja rapuh, karena
penyerang bisa menembusnya dengan BPDA, sedangkan pengaburan yang datang dari perubahan
model lebih tahan lama, sebab penyerang harus memperbarui pengetahuannya, dan memperbarui
pengetahuan berarti memperoleh akses, dan akses meninggalkan jejak. Bedanya bukan pada
teknik, melainkan pada siapa yang membayar biaya.

\index{transferabilitas}
Penulis tinjauan itu menyusun pohon serangan dan taksonomi masalah kepraktisan, lalu
menyimpulkan bahwa hanya tiga jalur yang layak. Jalur pertama adalah gangguan buta:
menambahkan byte acak pada paket UDP, atau pada paket TCP dengan bendera PSH, sambil
menghitung ulang checksum supaya protokolnya sah. Jalur ini tidak memerlukan pengetahuan
model dan tidak memerlukan kueri, jadi ia kebal terhadap pelatihan ulang, tetapi ia tidak
punya arah. Jalur kedua adalah kueri lewat saluran samping, yang menuntut akses jaringan
berkelanjutan dan bergantung pada indikator yang berubah menurut beban jaringan dan
perangkat keras. Jalur ketiga adalah model pengganti yang dilatih dari data tangkapan
botnet sendiri, yang bekerja lewat transferabilitas tetapi menuntut data latih yang masih
baru.

Ketiga jalur itu nyata, dan penulisnya menekankan bahwa jarak antara penelitian dan
praktik lebar serta bahwa asumsi model ancaman dalam banyak makalah terlalu murah hati
kepada penyerang.

\index{evasi}
Kerja penulis sendiri, yang akan kami sebut catatan kerja dan bukan hasil yang sudah
diterbitkan, bergerak di sisi lain persoalan yang sama. Pada satu kumpulan data aliran
jaringan, sebuah model mencapai F1 $0{,}967037$ dari $46\,558$ atas $48\,145$ contoh uji.
Pertanyaannya bukan ``apakah serangan berhasil'', melainkan ``berapa fitur yang harus
diubah agar sebuah contoh berpindah sisi batas keputusan sambil tetap sah menurut
protokol''. Dengan program linier campuran bilangan bulat, jawabannya bergantung tajam
pada jumlah fitur yang boleh berubah: pada tujuh fitur, tidak satu pun dari lima puluh
contoh bisa diseberangkan; pada dua puluh fitur, tiga puluh tiga dari lima puluh bisa.
Rasio jarak ke batas antara contoh yang berhasil dan contoh kendali sekitar $3{,}28$
untuk model linier dan $2{,}13$ untuk jaringan saraf.

\index{aturan tiga}
Dua angka lagi menunjukkan kehati-hatian. Dari dua puluh ribu kandidat yang diuji untuk
konsistensi protokol tidak ada yang lolos; dengan aturan tiga, batas atas $95$ persen
untuk laju itu sekitar $0{,}019$ persen. Pada model dengan margin keselamatan yang
diperbesar, syaratnya berbentuk $\varepsilon+40\delta$ dengan
$\delta=1{,}245\times10^{-7}$, sehingga anggaran efektifnya $0{,}150004981$ dan bukan
$0{,}15$.

Batas catatan kerja itu harus disebut bersama hasilnya. Ia memakai satu kumpulan data,
lima puluh contoh uji, dan dua keluarga model. Ia memverifikasi ketiadaan serangan yang
memenuhi syarat, bukan menemukan serangan. Angka $0$ dari $20\,000$ adalah pernyataan
tentang generator kandidat yang dipakai, bukan tentang semua paket yang mungkin, dan
hasilnya tidak mengatakan apa pun tentang NIDS dinamis.

Sekarang kesimpulan yang harus diucapkan jelas. Model mainan di bagian lima tidak
membuktikan apa pun tentang NIDS nyata; ia menunjukkan bahwa pembulatan yang menghapus
gradien menghasilkan laporan yang menyesatkan, dan bahwa mekanisme itu bisa ditembus
dengan biaya wajar. Ia tidak mengukur biaya serangan di ruang paket. Yang bisa dibawa
dari bab ini adalah kebiasaan: setiap kali sebuah pertahanan melaporkan keberhasilan nol,
tanyakan gradiennya, lalu tanyakan siapa yang membayar biaya.

\begin{tantangan}
Sebuah makalah melaporkan pertahanan yang menurunkan keberhasilan serangan dari $96$
persen ke $3$ persen, dengan akurasi bersih yang tidak berubah sama sekali. Penyerang
mereka memakai FGSM satu langkah, $\varepsilon=8/255$, gradien eksak, satu titik awal,
dan satu seed. Pertanyaannya: apakah ketahanan itu nyata, dan percobaan termurah apa
yang paling cepat membuktikan bahwa angka $3$ persen itu artefak?

Petunjuk pertama: perhatikan apa yang tidak dibayar. Pertahanan yang benar-benar
menaikkan biaya hampir selalu membayar sesuatu, entah akurasi, entah latensi, entah
jumlah parameter.

Petunjuk kedua: kalau gradien dihitung melalui potongan tak terdiferensiasi, angka
$3$ persen itu bukan keberhasilan serangan, melainkan sisa contoh yang kebetulan lahir
dekat batas keputusan. Ulangi dengan PGD empat puluh langkah, BPDA, dan tiga titik awal;
lalu sapu $\varepsilon$ dan laporkan kurvanya.

Kalau mentok, kembali ke bagian tiga.
\end{tantangan}

\begin{untukpeneliti}
Tiga hal di bagian ini sering salah dikutip. Pertama, kategori gradien hancur, acak, dan
hilang adalah taksonomi dari makalah obfuscated gradients; daftar tujuh teknik di bagian
dua adalah dekomposisi pedagogis kami. Di makalah itu ada tiga kategori, sembilan
pertahanan yang diuji, tujuh di antaranya bergantung pada gradien yang dikaburkan, dan
enam yang ditembus sepenuhnya ditambah satu yang ditembus sebagian.

Kedua, BPDA bukan penyelesaian umum. Batas galat $L\delta$ hanya berlaku bila $f$ mulus
dan $\delta$ kecil terhadap skala lekukan; pada pembulatan dengan $q$ sebanding dengan
$\varepsilon$ jaminan itu kosong. Sebaliknya, EOT bukan obat untuk serangan yang salah
arah: merata-ratakan arah yang salah tetap menghasilkan arah yang salah.

Ketiga, hasil NIDS yang kami kutip berasal dari percobaan ruang fitur dengan pengetahuan
penuh dan pelatihan ulang harian, bukan dari serangan ruang paket yang dijalankan. Ia
mendukung klaim sempit: pelatihan ulang memulihkan akurasi terhadap serangan yang memakai
gradien model lama. Ia tidak mendukung klaim bahwa NIDS kebal terhadap penyerang yang
boleh mengubah paket.
\end{untukpeneliti}

\section{Daftar periksa sebelum menulis bahwa serangan berhasil}
\tingkat{dasar}

Bagian ini adalah ringkasan yang bisa dipakai langsung. Setiap butir adalah pertanyaan
yang harus bisa dijawab sebelum sebuah klaim tentang serangan atau pertahanan masuk ke
abstrak.

\begin{enumerate}
  \item \textbf{Model ancaman ditulis lengkap.} Apa yang diketahui penyerang: arsitektur,
        bobot, data latih, probabilitas, atau hanya label? Berapa kueri yang boleh ia
        lakukan, dan berapa lama?
  \item \textbf{Serangan dibatasi pada anggaran yang masuk akal.} Sebutkan $\varepsilon$
        dan normanya, atau batas byte dan paket bila bekerja di ruang masalah. Serangan
        tanpa batas mengukur sesuatu yang tidak dihadapi siapa pun.
  \item \textbf{Serangan bersifat adaptif.}
        \index{serangan adaptif}
        Penyerang harus mengetahui pertahanan yang diuji, termasuk kemungkinan bahwa
        gradiennya rusak. Pertahanan yang hanya diuji oleh serangan dari makalah
        sebelumnya belum diuji.
  \item \textbf{Gradien diperiksa sebelum dipercaya.} Laporkan norma gradien, bagian
        entri yang nol, dan korelasinya dengan arah yang benar. Gradien nol pada model
        tak konstan adalah temuan, bukan hasil.
  \item \textbf{Keberhasilan dilaporkan bersama akurasi bersih.} Pertahanan yang
        menurunkan akurasi bersih dari $92$ ke $60$ persen tidak boleh dibandingkan
        langsung dengan yang mempertahankan $92$ persen.
  \item \textbf{Beberapa seed, dengan rata-rata dan simpangan.} Satu seed adalah satu
        lemparan dadu; selisih dua baris yang lebih kecil daripada simpangannya bukan
        selisih.
  \item \textbf{Kurva, bukan satu titik.} Sapu nilai $\varepsilon$, jumlah langkah,
        jumlah sampel EOT, dan jumlah kueri. Pertahanan yang baik menggeser kurva;
        yang semu hanya menggeser satu titik.
  \item \textbf{Anggaran komputasi penyerang disebut:} berapa kali model dipanggil,
        berapa jam, berapa GPU? Serangan yang berhasil setelah sejuta kueri bukan
        serangan yang sama dengan yang berhasil setelah seratus.
  \item \textbf{Kode dan bobot tersedia.} Klaim yang tidak bisa diulang tidak bisa
        diperiksa. Bila pelepasan tidak mungkin, prosedur dan hyperparameter harus cukup
        rinci untuk direkonstruksi.
  \item \textbf{Hasil negatif dilaporkan sebagai hasil negatif.} Pertahanan yang gagal
        adalah informasi; laporan yang berhasil karena serangannya dihentikan sebelum
        waktunya adalah kesalahan yang mahal.
\end{enumerate}

\kamus{aturan tiga}{rule of three}{bila tidak ada kejadian dalam $n$ percobaan bebas,
batas atas $95$ persen untuk peluang kejadian itu kira-kira $3/n$}

\begin{ringkasan}
\begin{itemize}
  \item Gradien yang dikaburkan membuat penyerang buta, bukan membuat model kuat.
        Ketujuh teknik itu menghasilkan gejala yang sama: pengukuran terlihat membaik
        tanpa ada yang membaik.
  \item Untuk fungsi tangga seperti pembulatan, gradien komposisi nol di hampir semua
        titik, sehingga serangan satu langkah menghasilkan $x'=x$ dan ``akurasi di bawah
        serangan'' sama persis dengan akurasi bersih.
  \item BPDA mengganti gradien potongan yang tidak terdiferensiasi dengan gradien
        pengklasifikasi di titik sesudah potongan. Galatnya dibatasi $L\delta$: masuk
        akal ketika penggeseran kecil terhadap skala lekukan, tidak memberi jaminan
        ketika tidak.
  \item EOT memaksimalkan harapan kerugian atas keacakan pertahanan; rata-rata $m$
        sampel tak bias dengan variansi sebanding $1/m$. Keacakan menaikkan biaya per
        langkah, bukan mengubah arah.
  \item Pada model mainan delapan fitur dengan $\varepsilon=0{,}25$, pembulatan
        $q=0{,}01$ menjatuhkan keberhasilan naif dari $17{,}40$ ke $0{,}02$ persen,
        sedangkan BPDA tetap $17{,}40$ persen; $q=0{,}20$ juga tidak menahan BPDA.
  \item Pada NIDS, FGSM menjatuhkan akurasi SVM dari $0{,}998$ ke $0{,}150$, tetapi
        pelatihan ulang harian memulihkannya ke $0{,}998$. Pertahanan yang tahan lama di
        sana datang dari biaya akses dan batasan protokol, bukan dari pengaburan gradien.
\end{itemize}
\end{ringkasan}

\section{Latihan}
\tingkat{menengah}

\begin{latihan}
\begin{enumerate}
  \item Sebuah makalah menulis: ``Pertahanan kami menurunkan keberhasilan serangan dari
        $98$ menjadi $0$ persen pada $\varepsilon=8/255$.'' Sebutkan tiga pemeriksaan
        sebelum menerima kalimat itu, dan apa yang diharapkan pada tiap pemeriksaan bila
        pertahanannya hanya mengaburkan gradien.
  \item Turunkan batas galat BPDA untuk $f(x)=w\cdot x+b$, tunjukkan galatnya nol, lalu
        jelaskan mengapa hasil itu tidak berlaku untuk model dua lapis beraktivasi ReLU.
  \item Pada model mainan bagian empat, tentukan nilai $q$ terkecil yang membuat satu
        langkah BPDA gagal membalik label $x_0$, nyatakan sebagai fungsi margin awal dan
        $\varepsilon$, lalu periksa dengan program.
  \item Anda punya $1\,000$ kueri. Pertahanan mengacak topeng fitur dengan $r=0{,}5$ dan
        menolak menjawab bila lebih dari separuh fitur hilang. Hitung berapa langkah PGD
        yang bisa dijalankan dengan $m=20$ dan dengan $m=1$, lalu jelaskan apa yang Anda
        peroleh dan apa yang Anda bayar.
  \item Ulang percobaan bagian lima dengan model dua lapis beraktivasi ReLU. Laporkan
        akurasi bersih, keberhasilan naif, dan keberhasilan BPDA untuk tiga nilai $q$,
        lalu jelaskan mengapa kosinusnya tidak lagi mendekati satu.
  \item Sebuah NIDS membaca lima puluh fitur, tujuh di antaranya bendera biner dan tiga
        nomor protokol. Hitung berapa fitur yang benar-benar bisa diubah penyerang, sebut
        asumsinya, dan bandingkan dengan gambar $32\times32$ yang punya $3\,072$ nilai.
  \item Jelaskan mengapa pelatihan ulang harian memulihkan akurasi terhadap serangan
        bergradien tetapi tidak terhadap gangguan buta. Kaitkan dengan lemma EOT dan
        definisi serangan adaptif.
  \item Rancang percobaan yang memisahkan dua hipotesis: (a) pertahanan menaikkan biaya
        serangan, (b) pertahanan hanya merusak pengukuran. Sebutkan metrik, jumlah seed,
        bentuk kurva, dan hasil seperti apa yang membuat Anda menerima (a).
\end{enumerate}
\end{latihan}

\peta{Bab 30 pada Athalye dkk., \emph{Obfuscated Gradients Give a False Sense of
Security} (2018): bagian 3 untuk tiga keluarga gradien yang dikaburkan, bagian 4 untuk
BPDA dan serangan adaptif, bagian 6 untuk daftar periksa. Pada Madry dkk., \emph{Towards
Deep Learning Models Resistant to Adversarial Attacks} (2018): bagian 3 untuk PGD dan
titik awal acak. Pada elShehaby dan Matrawy, \emph{Evasion Adversarial Attacks Remain
Impractical Against ML-based Network Intrusion Detection Systems}, untuk pohon serangan,
taksonomi kepraktisan, dan percobaan pelatihan ulang. Pada Bab 20 buku ini untuk aturan
tiga, dan pada Bab 27 untuk FGSM dan PGD.}

\chapter[Kriptografi: Modular sampai Kuantum]{Kriptografi: Bilangan Bulat Modular sampai Ancaman Kuantum}

Seorang mahasiswa mengumpulkan tugas berisi sandi buatan sendiri. Setiap huruf diganti nomornya,
nomor itu dikalikan tiga, lalu sisanya dibagi dua puluh enam dipakai sebagai huruf baru. Untuk
membuka pesan, ia mengalikan lagi dengan sembilan. Ketika ditanya dari mana angka sembilan itu,
jawabannya hampir benar: ``karena membagi tidak ada di sini, saya mengalikan dengan
kebalikannya.'' Penguji lalu mengganti pengali tiga dengan pengali empat. Tugas itu tidak pernah
selesai. Pada pengali empat, huruf pertama dan huruf keempat belas menghasilkan simbol yang
sama, sehingga dua pesan berbeda menghasilkan sandi identik.

Pertanyaan yang muncul dari kegagalan itu adalah pertanyaan seluruh bab ini: kapan sebuah
bilangan boleh dibagi di dalam aritmetika sisa? Jawabannya pendek: tepat ketika bilangan itu
tidak berbagi faktor dengan modulusnya. Jawaban pendek itu, ditekan terus, menghasilkan tiga hal
yang dipakai setiap peramban: pertukaran kunci, kunci publik, dan tanda tangan digital.

Tiga bab sebelumnya dipakai secara eksplisit. Bab 2 memberi bahasa untuk bijeksi dan partisi;
di sini bijeksi membuktikan bahwa perkalian dengan unit hanya mengacak urutan, dan partisi
menghitung orde. Bab 3 memberi sifat kelengkapan bilangan real, dan sifat itu tidak menyeberang,
karena $\mathbb{Z}_n$ tidak punya urutan yang cocok dengan perkalian; yang menggantikannya
adalah keterhinggaan, sebab pada himpunan berhingga pemetaan injektif otomatis surjektif. Bab 28
menjelaskan mengapa ``aman'' tidak bisa menjadi teorema tanpa syarat.

Satu peringatan berlaku untuk seluruh halaman. Semua modulus di sini kecil, dipilih agar bisa
diperiksa dengan tangan dan Python dalam hitungan detik. Modulus terbesar, $1\,050\,589$,
difaktorkan skrip enam baris. Tidak ada satu pun sandi di sini yang aman dipakai.

\begin{untukanak}
Bayangkan jam dinding di ruang kelas. Jarumnya bergerak dari angka satu sampai dua belas, lalu
kembali ke satu. Jam itu tidak pernah menunjukkan angka tiga belas atau dua puluh. Kalau kita
memutar jarum lima langkah, lalu lima langkah lagi, lalu lima langkah lagi, kita tidak perlu
tahu berapa lima ditambah lima ditambah lima; kita cukup melihat di mana jarum berhenti.

Sekarang bayangkan kunci kombinasi dengan tiga belas posisi. Kita boleh memutar sejauh berapa
pun, dan yang tercatat hanya posisi terakhir. Kalau kita memutar tiga langkah, berapa kali
kita harus memutar tiga langkah lagi supaya kembali ke tempat semula? Jawabannya sembilan, dan
itu bukan kebetulan, melainkan karena tiga dan tiga belas tidak punya faktor bersama. Coba hal
yang sama dengan empat langkah pada kunci dua belas posisi: berapa pun kali kita memutar, kita
tidak akan pernah mendarat di semua posisi, karena empat dan dua belas punya faktor bersama.

Di bab ini kita akan tahu mengapa perbedaan kecil itu menentukan apakah sebuah pesan bisa
dibuka kembali atau hilang selamanya, dan mengapa dua orang yang belum pernah bertemu bisa
menyepakati satu angka rahasia di depan orang ketiga yang mendengar semuanya.
\end{untukanak}

\section{Kelas sisa: aritmetika yang berputar}
\tingkat{dasar}

Kita mulai dengan membakukan gagasan ``sisa'' yang sudah dipakai di paragraf pertama.
\index{aritmetika modular}
Dua bilangan bulat disebut kongruen modulo $n$ bila selisihnya habis dibagi $n$.

\begin{definisi}
Misalkan $n\ge1$ bilangan bulat. Untuk $a,b\in\mathbb{Z}$ kita tulis $a\equiv b\pmod{n}$
bila $n$ membagi $a-b$. Himpunan kelas ekuivalensi di bawah relasi ini disebut
$\mathbb{Z}_n$, dan kelas yang memuat $a$ ditulis $[a]_n$.
\end{definisi}

Kita berutang dua hal: bahwa ini benar-benar relasi ekuivalensi, dan bahwa sisanya tunggal.

\begin{lemma}
Kongruensi modulo $n$ adalah relasi ekuivalensi pada $\mathbb{Z}$, dan setiap kelas
$[a]_n$ memuat tepat satu wakil pada $\{0,1,\dots,n-1\}$, sehingga $|\mathbb{Z}_n|=n$.
\end{lemma}

\begin{proof}
Refleksif: $n$ membagi $0$. Simetris: bila $n\mid(a-b)$ maka $n\mid(b-a)$. Transitif: bila
$n\mid(a-b)$ dan $n\mid(b-c)$ maka $n\mid(a-c)$ karena $(a-c)=(a-b)+(b-c)$. Untuk
ketunggalan, ambil $a$ dan tinjau $S=\{a-qn\ge0: q\in\mathbb{Z}\}$. Himpunan itu tidak kosong
menurut sifat Archimedes (Bab 3), jadi punya anggota terkecil $r$. Bila $r\ge n$ maka
$r-n\in S$ lebih kecil dari $r$, kontradiksi; jadi $0\le r<n$. Bila $r'$ juga memenuhi
$0\le r'<n$, maka $r-r'$ kelipatan $n$ yang nilainya di antara $-(n-1)$ dan $n-1$, sehingga
harus nol.
\end{proof}

Operasi pada $\mathbb{Z}_n$ didefinisikan lewat wakilnya, dan definisi semacam ini berbahaya:
bila hasilnya bergantung pada wakil, operasinya bukan fungsi. Lemma berikut menutup lubang itu
dengan satu baris aljabar.

\begin{lemma}
Bila $a\equiv a'$ dan $b\equiv b'$ modulo $n$, maka $a+b\equiv a'+b'$ dan $ab\equiv a'b'$
modulo $n$.
\end{lemma}

\begin{proof}
Tulis $a'=a+kn$ dan $b'=b+ln$. Maka $a'+b'=(a+b)+(k+l)n$ dan $a'b'=ab+(al+bk+kln)n$.
\end{proof}

Penjumlahan dan perkalian karena itu terdefinisi baik, dan keduanya mewarisi asosiatif serta
distributif dari $\mathbb{Z}$. Yang hilang adalah urutan: tidak ada cara memberi tanda
``$<$'' pada $\mathbb{Z}_n$ sehingga $a<b$ berakibat $a+c<b+c$ dan $ac<bc$ untuk $c$ positif.
Alat supremum Bab 3 berhenti di pintu masuk, dan yang bekerja di dalam adalah pencacahan.

Pembagian belum tentu ada. Kelas $[3]_{26}$ punya kebalikan, yaitu $[9]_{26}$, sebab
$3\cdot9=27\equiv1$. Kelas $[2]_{26}$ tidak punya: setiap kelipatan dua genap, sedangkan
$27$ ganjil, jadi $2b\equiv1$ tidak punya penyelesaian. Akibatnya perkalian dengan dua mod dua
puluh enam bukan bijeksi, karena $[0]$ dan $[13]$ keduanya menuju $[0]$.

\kamus{kelas sisa}{residue class}{semua bilangan yang berhenti di posisi jam yang sama,
seperti pukul tiga pagi dan pukul tiga sore}

\begin{definisi}
Kelas $[a]_n$ disebut unit bila ada $[b]_n$ dengan $[a]_n[b]_n=[1]_n$. Himpunan semua unit
modulo $n$ ditulis $U(n)$, dan banyak anggotanya ditulis $\varphi(n)$.
\end{definisi}

Untuk $n=26$ ada dua belas unit, yaitu kelas dari $1,3,5,7,9,11,15,17,19,21,23,25$. Program di
lab menghitungnya ulang; angka itu muncul lagi di bagian tiga dengan nama berbeda.

\section{Faktor persekutuan, Euklid, dan kapan pembagian ada}
\tingkat{dasar}

\index{faktor persekutuan terbesar}
Pertanyaan ``kapan $[a]_n$ punya kebalikan'' ternyata pertanyaan tentang faktor persekutuan
terbesar, yang punya algoritme berumur lebih dari dua ribu tahun.

\begin{definisi}
Untuk $a,b\in\mathbb{Z}$ yang tidak keduanya nol, $\gcd(a,b)$ adalah bilangan asli terbesar
yang membagi $a$ dan $b$.
\end{definisi}

Kunci algoritme Euklid adalah satu kesamaan himpunan: pembagi bersama $a$ dan $b$ sama persis
dengan pembagi bersama $b$ dan $a-qb$.

\begin{lemma}
Untuk setiap $q\in\mathbb{Z}$, $\gcd(a,b)=\gcd(b,a-qb)$.
\end{lemma}

\begin{proof}
Bila $d$ membagi $a$ dan $b$, maka $d$ membagi $a-qb$. Sebaliknya bila $d$ membagi $b$ dan
$a-qb$, maka $d$ membagi $(a-qb)+qb=a$. Jadi kedua pasangan punya himpunan pembagi bersama
yang sama, sehingga pembagi terbesarnya sama.
\end{proof}

Menukar $(a,b)$ menjadi $(b,a\bmod b)$ karena itu tidak mengubah jawaban tetapi mengecilkan
$b$, dan prosesnya berhenti karena bilangan asli tidak mengecil tanpa batas.

\begin{teorema}[Bezout]
Untuk $a,b$ yang tidak keduanya nol, ada $x,y\in\mathbb{Z}$ dengan $ax+by=\gcd(a,b)$.
\end{teorema}

\begin{proof}
Tinjau $S=\{ax+by>0: x,y\in\mathbb{Z}\}$. Himpunan ini tidak kosong karena memuat $|a|$ atau
$|b|$, dan setiap himpunan bagian tak kosong dari bilangan asli punya anggota terkecil; inilah
bentuk diskret dari sifat infimum di Bab 3. Misalkan $d=ax_0+by_0$ anggota terkecil itu.
Bagilah $a$ oleh $d$: $a=qd+r$ dengan $0\le r<d$. Maka $r=a(1-qx_0)+b(-qy_0)$, jadi $r\in S$
bila $r>0$. Karena $d$ terkecil, $r>0$ mustahil; jadi $r=0$ dan $d\mid a$, dan dengan cara
sama $d\mid b$. Jadi $d\le\gcd(a,b)$. Sebaliknya $\gcd(a,b)$ membagi setiap bentuk $ax+by$,
khususnya $d$, sehingga $\gcd(a,b)\le d$.
\end{proof}

\begin{teorema}
Kelas $[a]_n$ adalah unit jika dan hanya jika $\gcd(a,n)=1$.
\end{teorema}

\begin{proof}
Bila $[a]_n$ unit, ada $b$ dengan $ab\equiv1$, yaitu $ab-1=kn$. Setiap pembagi bersama $a$ dan
$n$ membagi $ab$ dan $kn$, jadi membagi $1$; maka $\gcd(a,n)=1$. Sebaliknya, bila
$\gcd(a,n)=1$, teorema Bezout memberi $ax+ny=1$. Dibaca modulo $n$, kesamaan itu berbunyi
$ax\equiv1$, jadi $[x]_n$ kebalikan $[a]_n$.
\end{proof}

Bukti itu sekaligus memberi resep: jalankan Euklid, lalu susun ulang mundur. Untuk $a=1071$ dan
$b=462$ langkahnya tiga pembagian,
\[
1071=2\cdot462+147,\qquad 462=3\cdot147+21,\qquad 147=7\cdot21+0,
\]
sehingga $\gcd=21$. Menyusun mundur memberi $21=1071\cdot(-3)+462\cdot7$, dan karena faktor
persekutuannya $21$ bukan $1$, kelas $[462]_{1071}$ bukan unit.

Berapa lama algoritme ini bekerja? Setiap dua langkah sisa setidaknya terbagi dua, jadi jumlah
pembagian sebanding dengan jumlah digit. Batas tajam diberikan Lame pada 1844: tidak lebih dari
$\log_\phi(b)+1$ pembagian dengan $\phi=(1+\sqrt5)/2$, dan yang mencapainya adalah pasangan
Fibonacci berurutan. Pasangan $(F_{47},F_{46})=(2\,971\,215\,073,\ 1\,836\,311\,903)$ butuh
$45$ pembagian, sedangkan $\log_\phi(F_{46})=44{,}33$. Untuk dua ribu pasangan acak tiga puluh
dua bit, rata-ratanya $18{,}75$ dengan maksimum $32$.

\kamus{invers modular}{modular inverse}{angka yang bila dikalikan dengan bilangan asal
menghasilkan sisa satu, seperti tombol ``batalkan'' pada mesin kasir}

\section{Fermat kecil dan Euler}
\tingkat{menengah}

Bab ini butuh satu fakta tentang pangkat, dan fakta itu lahir dari pencacahan: kalikan semua
unit dengan satu unit tetap; hasilnya daftar yang sama dengan urutan berbeda.

\begin{lemma}
Bila $\gcd(a,n)=1$, pemetaan $[x]_n\mapsto[ax]_n$ adalah bijeksi dari $U(n)$ ke $U(n)$.
\end{lemma}

\begin{proof}
Hasilnya berada di $U(n)$ karena hasilkali dua unit adalah unit: bila $xx'\equiv1$ dan
$yy'\equiv1$ maka $(xy)(x'y')\equiv1$. Pemetaannya injektif: bila $ax\equiv ay$, kalikan kedua
ruas dengan kebalikan $a$ yang ada menurut teorema unit di bagian dua. Himpunannya berhingga,
dan pada himpunan berhingga pemetaan injektif dari himpunan ke dirinya sendiri otomatis
surjektif (Bab 2). Jadi ia bijeksi.
\end{proof}

\begin{teorema}[Euler]
Bila $\gcd(a,n)=1$, maka $a^{\varphi(n)}\equiv1\pmod n$.
\end{teorema}

\begin{proof}
Misalkan $u_1,\dots,u_k$ semua unit modulo $n$, dengan $k=\varphi(n)$, dan $P=u_1\cdots u_k$.
Menurut lemma di atas, daftar $au_1,\dots,au_k$ memuat setiap unit tepat sekali, jadi
\[
a^kP\equiv(au_1)(au_2)\cdots(au_k)\equiv u_1u_2\cdots u_k=P\pmod n .
\]
Karena $P$ sendiri unit, ia punya kebalikan, dan kedua ruas boleh dikalikan dengannya.
Hasilnya $a^k\equiv1$.
\end{proof}

\begin{akibat}[Fermat kecil]
Bila $p$ prima dan $p\nmid a$, maka $a^{p-1}\equiv1\pmod p$. Untuk semua $a$ berlaku
$a^p\equiv a\pmod p$.
\end{akibat}

Bentuk kedua mengikuti yang pertama: bila $p\mid a$ kedua ruas nol, bila tidak, kalikan
$a^{p-1}\equiv1$ dengan $a$. Nama ``kecil'' membedakannya dari teorema terakhir Fermat, yang jauh lebih sulit.

Angka dua belas dari bagian pertama kini punya dua nama: $\varphi(26)=12$, dan
$3^{12}\equiv1\pmod{26}$; program memeriksa bahwa $5^{12}$ dan $7^{12}$ juga $\equiv1$.
Pangkat besar karena itu bisa dipangkas: untuk $7^{222}\bmod13$ tulis $222=12\cdot18+6$,
sehingga $7^{222}\equiv7^6\equiv9\cdot10\equiv12$ karena $7^2\equiv10$ dan $7^4\equiv9$.

Kebalikannya tidak berlaku, dan itu penting untuk kejujuran bab ini. Bilangan $561=3\cdot11\cdot17$
bukan prima, tetapi $a^{560}\equiv1\pmod{561}$ untuk setiap $a$ yang tidak berbagi faktor
dengannya; program memeriksa seluruh $a$ tanpa menemukan penyangkal. Bilangan semacam ini
disebut Carmichael, dan karena itu uji keprimaan praktis memakai Miller--Rabin.

\kamus{fungsi phi}{Euler totient function}{banyaknya posisi pada kunci kombinasi yang bisa
mencapai seluruh posisi lain, yaitu posisi yang tidak berbagi faktor dengan jumlah posisi}

\section{Grup siklik, orde, dan akar primitif}
\tingkat{menengah}

Kalau kita mengalikan satu unit berulang kali, hasilnya berputar dan akhirnya kembali ke satu.
Pertanyaan yang berguna adalah berapa lama putaran itu.

\begin{definisi}
Untuk $\gcd(a,n)=1$, orde $a$ modulo $n$ adalah bilangan asli terkecil $k$ dengan
$a^k\equiv1\pmod n$, ditulis $\operatorname{ord}_n(a)$.
\end{definisi}

Orde selalu ada: barisan $a^0,a^1,\dots$ hanya punya $\varphi(n)$ nilai berbeda, jadi dua di
antaranya sama dan $a^{i-j}\equiv1$ untuk suatu $i>j$. Batas itu bisa dipertajam: orde membagi
$\varphi(n)$.

\begin{lemma}
Bila $\gcd(a,n)=1$, maka $\operatorname{ord}_n(a)$ membagi $\varphi(n)$.
\end{lemma}

\begin{proof}
Misalkan $k=\operatorname{ord}_n(a)$ dan $H=\{a^0,\dots,a^{k-1}\}$, yang anggotanya berbeda
semua, jadi $|H|=k$. Untuk setiap unit $x$, himpunan $xH=\{xa^0,\dots,xa^{k-1}\}$ juga punya
$k$ anggota, karena perkalian dengan unit adalah bijeksi. Himpunan-himpunan $xH$ membentuk
partisi dari $U(n)$: dua di antaranya sama atau lepas, sebab $xH$ dan $yH$ beririsan bila dan
hanya bila $y^{-1}x\in H$ (Bab 2). Jadi $\varphi(n)$ adalah $k$ dikalikan banyaknya bagian.
\end{proof}

\begin{definisi}
Bila $\operatorname{ord}_n(g)=\varphi(n)$, maka $g$ disebut akar primitif modulo $n$. Untuk
unit $h$, bilangan $k$ dengan $g^k\equiv h$ disebut logaritma diskret $h$ terhadap basis $g$.
\end{definisi}

Modulo $13$, $2,6,7,11$ berorde penuh $12$ dan karena itu akar primitif, sedangkan $3$ dan $9$
berorde $3$, $5$ dan $8$ berorde $4$, $4$ dan $10$ berorde $6$, dan $12$ berorde $2$. Semua
angka itu membagi $12$.

Bagi $2$, putarannya mengunjungi keduabelas kelas bukan nol sebelum kembali:
$2,4,8,3,6,12$, lalu $11,9,5,10,7,1$. Gambar~\ref{fig:jam} menampilkannya sebagai jam tiga
belas posisi. Bagi $3$, putarannya hanya $3,9,1$ lalu berulang.

\begin{figure}[htbp]
\centering
\begin{tikzpicture}
\foreach \k in {0,...,12}
  \node[circle,draw,inner sep=1pt,minimum size=6mm,font=\scriptsize] (n\k)
    at ({90-360*\k/13}:2.1cm) {\k};
\node[font=\footnotesize,align=center] at (0,0) {kalikan\\ dengan 2};
\foreach \a/\b in {1/2,2/4,4/8,8/3,3/6,6/12,12/11,11/9,9/5,5/10,10/7,7/1}
  \draw[-{Latex[length=1.5mm]},semithick] (n\a) to[bend left=18] (n\b);
\end{tikzpicture}
\caption{Perkalian dengan $2$ modulo $13$ berputar melalui keduabelas kelas bukan nol
sebelum kembali ke $1$. Panjang putaran itulah orde, dan orde penuh berarti $2$ akar
primitif modulo $13$.}
\label{fig:jam}
\end{figure}

Keberadaan akar primitif untuk setiap prima $p$ adalah teorema Gauss, dan buktinya kami
lewati. Yang perlu dibawa: untuk prima $p$, unit adalah satu putaran berisi $p-1$ posisi, dan
setiap unit adalah pangkat dari satu akar primitif. Bila $h=g^k$, mengetahui $k$ berarti
mengetahui $h$, tetapi tidak sebaliknya; asimetri itu bahan mentah kriptografi kunci publik.

\kamus{akar primitif}{primitive root}{satu posisi awal yang, bila diputar terus, melewati
semua posisi lain sebelum kembali}

\section{Pertukaran kunci Diffie-Hellman}
\tingkat{menengah}

\index{diffie-hellman}
Dua orang ingin menyepakati satu angka rahasia, dan satu-satunya saluran yang mereka punya
didengar pihak ketiga. Pada 1976 Diffie dan Hellman menunjukkan bahwa hal itu mungkin, dan
mekanismenya hanya memakai dua rumus.

Mereka sepakat lebih dulu pada prima $p$ dan akar primitif $g$ modulo $p$; keduanya boleh
diumumkan. Alice memilih $a$ rahasia dan mengirim $A=g^a\bmod p$. Bob memilih $b$ rahasia dan
mengirim $B=g^b\bmod p$. Alice menghitung $K=B^a\bmod p$, Bob menghitung $K=A^b\bmod p$.
Keduanya memperoleh angka yang sama, karena
\[
B^a\equiv(g^b)^a\equiv g^{ab}\equiv(g^a)^b\equiv A^b\pmod p .
\]
Yang dilihat penyadap hanyalah $p,g,A,B$. Untuk memperoleh $K$ ia perlu $g^{ab}$, dan cara
termudah yang diketahui adalah menemukan $a$ atau $b$, yaitu menyelesaikan logaritma diskret.

\begin{bacakeras}
Bacakan sebagai dialog dua suara, bergantian, satu baris satu tarikan napas.

Suara pertama: ``Sepakat dulu, umumkan dua.''
Suara kedua: ``Simpan sendiri, jangan disebut.''
Suara pertama: ``Pangkatkan, bagi sisa, kirim hasilnya.''
Suara kedua: ``Terima kiriman, pangkatkan lagi.''
Suara pertama: ``Sisa terakhir, itulah kuncinya.''
Bersama: ``Dua arah, satu angka, tidak pernah lewat.''

Ulangi tiga kali, makin cepat. Baris yang membuat lidah tersendat adalah langkah yang belum
dipahami, bukan baris yang perlu dihapal lebih keras. Sesudah lancar, tutup buku dan sebutkan
empat hal berurutan: yang diumumkan, yang disimpan, yang dikirim, yang dihitung. Menyebut angka
rahasia lebih dulu membatalkan seluruh permainan.

Baris yang diulang tiga kali dengan suara keras: ``yang dikirim hasil pangkat, yang disimpan
pangkatnya.''
\end{bacakeras}

Angka konkretnya, seluruhnya diperiksa program di bagian lab: $p=10\,007$, $g=5$ berorde
$10\,006$ sehingga akar primitif, $a=1234$, $b=5678$. Maka $A=5^{1234}\bmod p=4873$ dan
$B=5^{5678}\bmod p=3189$. Alice menghitung $3189^{1234}\bmod p=5674$; Bob menghitung
$4873^{5678}\bmod p=5674$. Kunci bersama mereka $5674$, dan angka itu tidak pernah melintasi
saluran. Gambar~\ref{fig:dh} menampilkan alurnya beserta apa yang didengar penyadap.

\begin{figure}[htbp]
\centering
\begin{tikzpicture}[font=\footnotesize,
  kotak/.style={draw,rounded corners=2pt,align=left,inner sep=4pt,text width=4.4cm}]
\node[kotak] (a) {\textbf{Alice}\\ $a=1234$ rahasia\\ $A=5^{1234}\bmod p=4873$};
\node[kotak,right=3.4cm of a] (b) {\textbf{Bob}\\ $b=5678$ rahasia\\
  $B=5^{5678}\bmod p=3189$};
\node[kotak,below=1.3cm of $(a)!0.5!(b)$, text width=9.2cm, fill=black!4]
  (e) {\textbf{Penyadap} mendengar $p=10\,007$, $g=5$, $A=4873$, $B=3189$, lalu bermaksud
  menghitung $K=A^{b}$ atau $K=B^{a}$; untuk itu ia harus menemukan $a$ atau $b$.};
\draw[-{Latex}] ([yshift=3mm]a.east) -- node[above,font=\scriptsize]{$A$}
  ([yshift=3mm]b.west);
\draw[-{Latex}] ([yshift=-3mm]b.west) -- node[below,font=\scriptsize]{$B$}
  ([yshift=-3mm]a.east);
\draw[-{Latex},dashed] (e.north -| a.south) -- (a.south);
\draw[-{Latex},dashed] (e.north -| b.south) -- (b.south);
\end{tikzpicture}
\caption{Alur pertukaran kunci Diffie-Hellman pada contoh angka di teks. Nilai $A$ dan $B$
melintasi saluran terbuka; yang tidak melintas adalah eksponen $a$ dan $b$, dan justru
keduanya yang dibutuhkan penyadap.}
\label{fig:dh}
\end{figure}

Kesamaan itu hanya menyatakan protokolnya konsisten, bukan aman. Yang bisa dibuktikan hanyalah
pernyataan bersyarat: bila logaritma diskret modulo $p$ sulit, menghitung $g^{ab}$ dari $A$
dan $B$ juga sulit. Syarat itu bukan teorema; ada grup tempat logaritma diskret mudah, jadi
keamanannya sifat kelompok tertentu.

\begin{tangan}
Bahan: cat air kuning, satu warna rahasia per orang (merah dan biru), empat gelas plastik,
satu sendok, dua lembar kertas putih.

Langkah pertama, di depan meja: tuang satu sendok cat kuning ke gelas masing-masing. Kuning itu
peran prima dan basis, dan keduanya diumumkan. Langkah kedua, di balik punggung: tiap orang
menambahkan dua sendok warna rahasianya, mengaduk, lalu menyerahkan gelas campurannya. Langkah
ketiga: masing-masing menambahkan dua sendok warna rahasianya sendiri ke gelas yang baru
diterima, mengaduk, lalu menotolkan satu tetes ke kertas.

Yang harus terlihat: dua totolan itu sewarna, padahal tidak satu pun dari kedua orang itu
pernah menyerahkan warna rahasianya. Gelas campuran pertama yang dipegang orang ketiga tidak
bisa dipisah kembali; ia hanya melihat dua warna kusam yang berbeda.
\end{tangan}

Dua jebakan parameter layak disebut. Pertama, $g$ harus membangkitkan subkelompok besar berorde
prima: modulus $10\,007$ adalah prima aman dengan $p-1=2\cdot5003$, dan program mencatat bahwa
$5^{5003}\bmod p=10\,006$ berorde hanya $2$. Bila Alice memakai basis itu, kunci bersamanya
hanya bernilai $1$ atau $10\,006$, dan penyadap tidak perlu menghitung logaritma apa pun.
Kedua, protokol ini tidak menyebut identitas siapa pun, sehingga penyadap bisa berdiri di tengah
dan membaca semuanya. Gambar~\ref{fig:dh} hanya menggambarkan penyadap pasif.

Tentang ukuran: modulus $10\,007$ bisa dibobol dalam waktu kurang dari satu detik oleh
algoritme indeks kalkulus, dan kita memakainya karena hasil pangkatnya muat di satu halaman.
Praktik nyata memakai $2048$ bit atau lebih, dan itu pun batas waktu, bukan jaminan.

\begin{untukanak}
Bayangkan dua orang yang ingin menyepakati satu warna rahasia, padahal semua percakapan mereka
didengar. Mereka sepakat lebih dulu pada satu warna dasar, misalnya kuning. Alice memilih warna
rahasianya, mencampurnya dengan kuning, lalu mengirim hasil campurannya. Bob melakukan hal yang
sama dengan warna rahasianya sendiri. Sekarang masing-masing mencampur warna yang diterimanya
dengan warna rahasianya sendiri. Keduanya mendapat warna yang sama, dan penyadap hanya melihat
dua campuran yang tidak bisa dipisah kembali.

Pemisahan itulah yang dianggap sulit: mencampur warna gampang, memisahkan satu campuran tanpa
salah satu asalnya dianggap tidak. Kriptografi mengganti warna menjadi angka dan ``dianggap
tidak gampang'' menjadi kalimat matematika.

Di bab ini kita akan tahu bagaimana angka menggantikan warna, dan mengapa keamanannya
bergantung pada satu kalimat ``dianggap'' yang tidak pernah dibuktikan siapa pun.
\end{untukanak}

\section{RSA: dari kunci ke pesan dan kembali}
\tingkat{menengah}

\index{rsa}
Diffie-Hellman menyepakati kunci, tetapi tidak mengirim pesan. Setahun kemudian Rivest,
Shamir, dan Adleman menunjukkan bahwa sifat yang sama bisa dipakai untuk dua kunci berbeda:
satu untuk mengunci, satu untuk membuka.

Pilih dua prima besar $p$ dan $q$, hitung $n=pq$ dan $\varphi(n)=(p-1)(q-1)$. Pilih $e$ dengan
$\gcd(e,\varphi(n))=1$; teorema unit memberi kebalikannya $d$ modulo $\varphi(n)$. Kunci publik
$(n,e)$, kunci rahasia $d$. Untuk pesan $0\le m<n$, sandinya $c=m^e\bmod n$ dan pembukaannya
$m=c^d\bmod n$.

Mengapa pembukaan mengembalikan $m$? Karena $ed=1+k\varphi(n)$, sehingga
$c^d=m^{ed}=m\,(m^{\varphi(n)})^k$. Bila $\gcd(m,n)=1$, teorema Euler membuat kurung itu satu.
Kasus sisanya menjelaskan mengapa $n$ harus hasilkali dua prima berbeda.

\begin{teorema}
Misalkan $n=pq$ dengan $p,q$ prima berbeda, $ed\equiv1\pmod{(p-1)(q-1)}$, dan $0\le m<n$.
Maka $m^{ed}\equiv m\pmod n$.
\end{teorema}

\begin{proof}
Bila $p\nmid m$, Fermat kecil memberi $m^{p-1}\equiv1\pmod p$, sehingga
$m^{ed}=m\,(m^{p-1})^{k(q-1)}\equiv m\pmod p$. Bila $p\mid m$, kedua ruas kongruen nol modulo
$p$. Jadi kesamaan berlaku modulo $p$ untuk setiap $m$, dan dengan alasan sama modulo $q$.
Karena $p$ dan $q$ keduanya membagi $m^{ed}-m$ dan $\gcd(p,q)=1$, hasilkali keduanya juga
membagi bilangan itu: bila $m^{ed}-m=pt$ dan $q\mid pt$ dengan $q$ prima serta $q\nmid p$,
maka $q\mid t$, jadi $pq\mid m^{ed}-m$.
\end{proof}

Contoh angkanya kecil dan seluruhnya diperiksa program. Ambil $p=1019$, $q=1031$, maka
$n=1\,050\,589$ dan $\varphi(n)=1018\cdot1030=1\,048\,540$. Pilih $e=65\,537$; Euklid memberi
$d=60\,493$, dan memang $65\,537\cdot60\,493=3\,964\,529\,741=1+3781\cdot1\,048\,540$. Pesan
$m=424\,242$ menghasilkan sandi $c=552\,241$, dan $c^d\bmod n$ kembali $424\,242$. Modulus itu hanya dua puluh satu bit dan difaktorkan skrip enam baris.

RSA polos, rumus di atas tanpa tambahan apa pun, tidak aman; tiga alasannya bisa ditunjukkan
dengan modulus yang sama. Pertama, ia deterministik: pesan yang sama selalu menghasilkan sandi
yang sama, sehingga pesan berulang bisa dikenali. Kedua, ia multiplikatif:
$E(m_1)E(m_2)\equiv E(m_1m_2)$, dan program memperlihatkan
$E(1234)\cdot E(567)\bmod n=551\,173$ dibuka menjadi $699\,678$, tepat $1234\cdot567$. Ketiga,
pesan kecil tidak terlindungi: dengan $e=3$ dan $m=42$, sandinya $74\,088$, kebetulan sama
dengan $42^3$ karena hasilnya belum melewati $n$, sehingga akar pangkat tiga biasa memulihkan
pesan.

Ketiganya ditutup dengan bantalan: pesan diberi tambahan acak sebelum dipangkatkan, sehingga
hasilnya panjang penuh, selalu berbeda, dan tidak lagi multiplikatif. RSA-OAEP, standar yang
dipakai hari ini, punya bukti keamanan terhadap serangan chosen-ciphertext dalam model oracle
acak dengan asumsi RSA berlaku. Dalam praktik RSA membungkus kunci simetris $256$ bit lalu data
dienkripsi dengan AES; pola ini disebut enkripsi hibrida, karena RSA lambat: waktunya
sebanding pangkat tiga panjang modulus.

\kamus{bantalan}{padding}{isi tambahan yang dimasukkan ke dalam paket sebelum dikirim,
seperti bubble wrap yang membuat setiap paket berbeda bentuk dan ukuran}

\section{Fungsi hash, MAC, dan tanda tangan}
\tingkat{menengah}

\index{fungsi hash}
Enkripsi menyembunyikan isi. Kadang yang dibutuhkan bukan menyembunyikan, melainkan meringkas:
satu sidik jari pendek yang berubah bila isinya berubah.

\begin{definisi}
Fungsi hash kriptografis memetakan pesan sepanjang apa pun ke keluaran berukuran tetap. Ia
harus tahan pra-gambar (dari $h$, sulit menemukan $m$ dengan $H(m)=h$), tahan pra-gambar
kedua (dari $m$, sulit menemukan $m'\neq m$ dengan $H(m')=H(m)$), dan tahan tabrakan (sulit
menemukan sepasang $m,m'$ berbeda dengan hash sama).
\end{definisi}

Ketiga sifat itu tidak setara; yang paling lemah paling sering disalahpahami. Ketahanan
tabrakan dibatasi ulang tahun: bila keluaran punya $N$ nilai, tabrakan muncul setelah kira-kira
$1{,}177\sqrt N$ percobaan, bukan setelah $N$; untuk $128$ bit angka itu $2{,}17\cdot10^{19}$
dan untuk $64$ bit $5{,}06\cdot10^{9}$. Program di lab menguji ramalan itu pada SHA-256 yang
dipotong menjadi $16$ bit: rata-rata $362{,}7$ sampel melawan ramalan $320{,}8$, dengan sebaran
dari $100$ sampai $743$. Jadi hash berkeluaran $n$ bit hanya memberi ketahanan tabrakan $n/2$
bit.

MAC menambahkan kunci pada hash: dua pihak yang berbagi kunci dapat memeriksa bahwa pesan
datang dari pihak lain dan tidak berubah, karena tanpa kunci tidak ada cara menghitung tag yang
benar. Batasnya: MAC hanya berguna antara dua pihak yang sama-sama memegang kunci, sehingga
pihak ketiga tidak bisa menjadi saksi.

Tanda tangan memakai pasangan kunci dengan arah terbalik: penandatangan memakai kunci rahasia,
pemeriksa memakai kunci publik. Pada RSA, tanda tangan atas $m$ adalah $s=H(m)^d\bmod n$, dan
pemeriksaan membandingkan $s^e\bmod n$ dengan $H(m)$. Keharusan memakai $H$ bukan soal
kenyamanan: tanpa hash, tanda tangan RSA bersifat multiplikatif seperti enkripsinya,
sehingga dari tanda tangan dua dokumen seseorang bisa menyusun tanda tangan sah untuk hasilkali
keduanya tanpa memiliki kunci rahasia. Hash memutus sifat itu, karena $H(m_1m_2)$ tidak
berhubungan dengan $H(m_1)H(m_2)$. Syarat keamanannya adalah ketahanan terhadap pemalsuan
eksistensial di bawah serangan pesan terpilih, dan sekali lagi itu asumsi.

\kamus{tanda tangan digital}{digital signature}{cap lilin pada dokumen yang bisa diperiksa
siapa saja tetapi hanya bisa dibuat oleh pemegang cetakannya}

\section{Apa yang sebenarnya diasumsikan model keamanan}
\tingkat{lanjut}

Setiap kalimat ``sulit'' di bab ini adalah janji tentang apa yang tidak bisa dilakukan
algoritme yang dikenal, bukan pernyataan tentang apa yang mustahil.

\begin{definisi}
Fungsi $f$ disebut satu arah bila ada algoritme polinomial untuk menghitung $f(x)$ dari $x$,
tetapi untuk $y$ acak di daerah hasil, setiap algoritme polinomial gagal menemukan $x$ dengan
$f(x)=y$ kecuali dengan peluang yang dapat diabaikan.
\end{definisi}

Arahnya digambarkan pada Gambar~\ref{fig:tangga-satu-arah}.

\begin{figure}[htbp]
\centering
\begin{tikzpicture}[font=\small,
  kotak/.style={draw,rounded corners=2pt,align=center,inner sep=5pt,
    minimum height=0.95cm}]
  \node[kotak,draw=layerbrown] (pq) at (0,0) {$p$, $q$ prima acak};
  \node[kotak,draw=ink] (n) at (4.0,0) {$n=p\,q$};
  \node[kotak,draw=layergrey,dashed] (s) at (8.2,0) {penyerang};
  \draw[-{Latex},very thick,layerjade] (pq) -- (n);
  \node[font=\scriptsize,layerjade,above] at (2.0,0.06)
    {kalikan: mudah, polinomial};
  \draw[-{Latex},thick,layergrey] (n) -- (s);
  \node[font=\scriptsize,layergrey,above] at (6.1,0.06) {melihat $n$};
  \draw[layerplum,thick,dashed] (n.south) -- (4.0,-1.35);
  \draw[layerplum,thick,dashed] (0,-1.35) -- (pq.south);
  \draw[-{Latex},layerplum,thick,dashed] (4.0,-1.35) -- (0,-1.35);
  \draw[layerplum,very thick] (1.86,-1.12)--(1.86,-1.58);
  \draw[layerplum,very thick] (2.04,-1.12)--(2.04,-1.58);
  \node[font=\scriptsize,layerplum,below] at (2.0,-1.62)
    {faktorkan: sulit, tidak ada algoritme polinomial yang dikenal};
\end{tikzpicture}
\caption{Dua arah yang tidak sama mahal. Dari $p$ dan $q$ ke $n$ cukup satu
perkalian, sedangkan dari $n$ kembali ke $p$ dan $q$ adalah masalah yang sampai
hari ini tidak punya algoritme polinomial yang dikenal; selisih itulah yang
dipakai RSA dan Diffie-Hellman. Palang pada panah balik menandai arah yang
diandaikan tertutup, bukan arah yang sudah terbukti mustahil.}
\label{fig:tangga-satu-arah}
\end{figure}

Definisi itu tidak menyebut satu fungsi pun, dan itu bukan kelalaian. Sampai hari ini tidak ada
bukti bahwa fungsi satu arah ada; keberadaannya akan berakibat $P\neq NP$, yang juga masih
terbuka. Kriptografi kunci publik karena itu berdiri di atas konjektur, dan rantainya berakhir di
asumsi yang tidak berubah menjadi teorema karena dipakai lama.

\index{asumsi faktorisasi}
Asumsi faktorisasi mengatakan: untuk $n$ hasilkali dua prima acak berukuran $k$ bit, tidak ada
algoritme polinomial dalam $k$ untuk memfaktorkannya. Algoritme terbaik yang diketahui,
penyaringan medan bilangan, berjalan dalam waktu $L_n[1/3,1{,}923]$, yang tumbuh
subeksponensial. Bilangan RSA-250, sepanjang $250$ digit atau sekitar $829$ bit, difaktorkan
pada 2020 dengan biaya sekitar $2\,700$ core-tahun.

Padanan kekuatan yang dicatat NIST dalam SP 800-57 berguna sebagai peta kasar.

\begin{center}
\small
\begin{tabular}{@{}lll@{}}
\toprule
Kekuatan (bit) & Modulus RSA atau DH (bit) & Sandi simetris \\
\midrule
112 & 2048 & 3DES tiga kunci \\
128 & 3072 & AES-128 \\
192 & 7680 & AES-192 \\
256 & 15360 & AES-256 \\
\bottomrule
\end{tabular}
\end{center}

Angka kolom tengah adalah taksiran biaya algoritme terbaik saat tabel disusun, jadi tabel
semacam ini punya tanggal kedaluwarsa.

Model keamanan formal menutup celah retorika dengan permainan: penantang membangkitkan kunci,
menjawab permintaan enkripsi, lalu memilih dua pesan dan mengirimkan enkripsi salah satunya;
skema aman bila penyerang tidak bisa menebak yang mana jauh di atas peluang setengah. Keamanan
dibuktikan sebagai reduksi: penyerang yang menang diubah menjadi algoritme pemfaktor $n$,
sehingga seluruh keraguan pindah ke satu tempat yang sempit tetapi tidak hilang.

\section{Shor, dan apa yang runtuh bersamanya}
\tingkat{lanjut}

\index{algoritma shor}
Seluruh bangunan di atas bergantung pada satu kesulitan, dan pada 1994 Peter Shor menunjukkan
bahwa komputer kuantum menyelesaikannya dalam waktu polinomial.

Gagasannya terbagi dua, dan bagian pertama klasik sepenuhnya. Untuk memfaktorkan $n$, pilih $a$
dengan $\gcd(a,n)=1$ dan hitung $r=\operatorname{ord}_n(a)$. Bila $r$ ganjil, ulangi dengan $a$
lain. Bila $r$ genap dan $a^{r/2}\not\equiv-1\pmod n$, maka $\gcd(a^{r/2}-1,n)$ adalah faktor
tak sejati dari $n$: karena $a^r\equiv1$, bilangan $a^{r/2}$ adalah akar kuadrat dari satu
yang bukan $\pm1$, dan akar semacam itu hanya muncul bila modulusnya terbagi. Sisa pekerjaannya
mencari orde, dan inilah bagian yang kuantum.

Bagian klasik saja tidak cukup. Untuk $n=1019\cdot1031$ kami mencoba $2\,996$ unit acak dan
hanya $319$, sekitar $10{,}6$ persen, memenuhi syarat sehingga memberi faktor; keberhasilan
karena itu hasil pengulangan dengan $a$ berbeda. Yang hilang tanpa komputer kuantum adalah cara
cepat menemukan $r$: tidak ada algoritme klasik yang diketahui berjalan polinomial dalam
$\log n$.

Komputer kuantum mencari $r$ dengan transformasi Fourier kuantum, dan di situlah percepatannya
berasal. Akibatnya bukan hanya RSA yang runtuh: algoritme Shor juga menyelesaikan logaritma
diskret, baik pada $\mathbb{Z}_p$ maupun pada kurva eliptik, sehingga Diffie-Hellman, DSA, ECDH, dan ECDSA ikut runtuh, termasuk kurva eliptik yang modulusnya paling
pendek karena kebal terhadap indeks kalkulus.

Yang tidak runtuh adalah sandi simetris dan fungsi hash. Grover memberi percepatan kuadratik
pada pencarian kunci, sehingga AES-128 turun dari $2^{128}$ menjadi sekitar $2^{64}$ operasi
kuantum, sedangkan AES-256 tetap di sekitar $2^{128}$; ketahanan tabrakan hash praktis tidak
berubah. Jawaban untuk data simetris karena itu sederhana: pakai kunci dua kali lebih panjang.

Perlu dibedakan dua hal yang sering dicampur dalam berita. Shor membuktikan bahwa faktornya
bisa dihitung polinomial pada komputer kuantum tanpa galat; membangun mesinnya persoalan
teknik lain. Taksiran Gidney dan Ekera (2021) menyebut sekitar dua puluh juta qubit fisik dan
delapan jam untuk modulus $2048$ bit, hampir seluruhnya untuk koreksi galat. Angka itu turun
beberapa orde besaran dalam satu dekade, tetapi belum ada mesin yang mendekati. Yang berlaku
sekarang adalah pola ``tangkap sekarang, buka nanti'': pesan hari ini bisa disimpan dan dibuka
ketika mesinnya ada.

\section{Kriptografi pasca-kuantum: lattice dan Kyber}
\tingkat{lanjut}

\index{kriptografi pasca-kuantum}
Bila faktorisasi jatuh, yang dibutuhkan adalah masalah lain yang tidak jatuh bersamanya.
Kandidat terkuat berasal dari lattice, yaitu himpunan titik yang dibangkitkan kombinasi bilangan
bulat dari beberapa vektor.

\index{lattice}
\emph{Learning with errors} menyusun masalahnya begini. Diberikan matriks $A$ acak
$m\times n$ atas $\mathbb{Z}_q$ dan vektor $b=As+e\bmod q$, dengan $s$ rahasia dan $e$ galat
kecil, temukan $s$. Tanpa $e$ ini hanya sistem persamaan linear yang diselesaikan eliminasi
Gauss (Bab 8); dengan $e$ kecil tetapi tidak nol, sistemnya hampir konsisten dan mencari $s$
dianggap sulit. Reduksi Regev (2005) menambah daya tariknya: menyerang skema acak sama sulitnya
dengan menyelesaikan kasus terburuk lattice, bukan kasus rata-rata pilihan penyerang.

Kyber, yang dibakukan NIST pada Agustus 2024 sebagai ML-KEM dalam FIPS 203, adalah varian modul
dari gagasan itu dengan $n=256$, $k=3$, dan $q=3329$. Kunci publiknya $1\,184$ byte dan
ciphertext-nya $1\,088$ byte, sekitar $37$ dan $34$ kali kunci simetris $32$ byte yang
dibungkusnya, dan biaya ruang itu nyata. Tanda tangan pasca-kuantum punya cerita serupa:
ML-DSA (Dilithium) juga berbasis lattice, sedangkan SLH-DSA (SPHINCS+) memilih jalan
konservatif berbasis hash.

Kejujuran yang wajib di sini: asumsi lattice jauh lebih muda daripada asumsi faktorisasi.
Faktorisasi diserang sejak Gauss; LWE baru berumur dua puluh tahun, dan sejarah kriptografi
penuh skema yang tampak kuat lalu patah setelah sepuluh tahun dianalisis. Karena itu transisi
yang dianjurkan bukan penggantian, melainkan hibrida: kunci sesi dibentuk dari Kyber dan dari
pertukaran Diffie-Hellman sekaligus.

\section{Ketika keluaran menjadi oracle: padding dan model}
\tingkat{lanjut}

\index{padding oracle}
Ada satu keluarga serangan yang tidak menyerang matematikanya, melainkan cara program menjawab.
Bleichenbacher pada 1998 menunjukkan bahwa RSA dengan bantalan PKCS\#1 v1.5 bisa dibuka
sepenuhnya bila server menjawab berbeda untuk ``bantalan salah'' dan ``bantalan benar tetapi isi
salah''. Setiap jawaban membocorkan satu bit tentang pesan, dan dengan beberapa ribu pertanyaan
pesannya terbaca. Matematika RSA tidak dilanggar; yang dilanggar adalah asumsi bahwa penyerang
hanya melihat ciphertext.

Pola ancamannya sama dengan yang dibahas Bab 29 dan Bab 30. Model pembelajaran mesin yang
mengembalikan probabilitas, atau pesan galat yang berbeda-beda, adalah oracle; penyerang yang
boleh bertanya berkali-kali memakai perbedaan kecil itu untuk menyimpulkan apa yang seharusnya
tidak ia ketahui, termasuk apakah sebuah contoh ada di data latih. Bab 30 menambahkan
peringatan yang berlaku di kedua dunia: pengukuran yang tampak membaik belum tentu karena
pertahanannya bekerja.

Pelajaran defensifnya karena itu sama. Keluaran harus dibuat seragam: satu pesan galat untuk
semua kegagalan, waktu tanggap yang tidak bergantung pada rahasia, penolakan yang tidak
membedakan sebab. Untuk model, ini berarti membatasi jumlah kueri dan menyamarkan skor. Semua
itu mengurangi kebocoran, tidak menghapusnya.

\begin{tantangan}
Sebuah pustaka memakai AES-GCM dengan kunci $K$ yang sama untuk dua pesan, tetapi penghitung
nonce-nya ter-reset sehingga nonce $96$ bit yang sama terpakai dua kali. Penyerang memegang
kedua sandi, kedua tag, dan isi pesan pertama, yaitu halaman login yang tetap. Tunjukkan
bagaimana ia memperoleh kunci autentikasi $H$, lalu mengapa $H$ cukup untuk memalsukan tag pesan
buatannya tanpa mengetahui $K$.

Petunjuk 1: nonce yang sama memberi kunci aliran yang sama, sehingga selisih kedua sandi
menghapus kunci itu.

Petunjuk 2: tag adalah polinomial dalam $H$ dengan koefisien yang diketahui; dua tag dengan
nonce sama memberi dua persamaan bagi satu $H$, dan selisihnya menjadi polinomial yang akarnya
dapat dihitung.

Kalau mentok, kembali ke bagian ``Fungsi hash, MAC, dan tanda tangan'' dan baca ulang paragraf
tentang MAC.
\end{tantangan}

\section{Contoh hitung}
\tingkat{menengah}

Kita hitung satu invers dengan tangan, karena langkah mundurnya adalah tempat kesalahan paling
sering terjadi. Ambil $\varphi(n)=1\,048\,540$ dari modulus mainan di bagian enam dan
$e=65\,537$. Euklid maju memberi
\[
1\,048\,540=15\cdot65\,537+65\,485,\qquad 65\,537=1\cdot65\,485+52,
\]
\[
65\,485=1259\cdot52+17,\qquad 52=3\cdot17+1,\qquad 17=17\cdot1+0 .
\]
Sisa terakhir yang bukan nol adalah $1$, jadi $\gcd=1$ dan inversnya ada. Susun mundur mulai
dari sisa satu:
\[
1=52-3\cdot17=52-3(65\,485-1259\cdot52)=3778\cdot52-3\cdot65\,485 .
\]
Lanjutkan dengan $52=65\,537-65\,485$ dan $65\,485=1\,048\,540-15\cdot65\,537$:
\[
1=3778\cdot65\,537-3781\cdot65\,485
=3778\cdot65\,537-3781(1\,048\,540-15\cdot65\,537).
\]
Kelompokkan: $1=(3778+15\cdot3781)\cdot65\,537-3781\cdot1\,048\,540
=60\,493\cdot65\,537-3781\cdot1\,048\,540$. Jadi $d=60\,493$, dan tanda negatifnya hilang
karena dibaca modulo $\varphi(n)$. Pemeriksaannya satu baris: $65\,537\cdot60\,493
=3\,964\,529\,741$, sedangkan $3781\cdot1\,048\,540=3\,964\,529\,740$; selisihnya satu.

Dua hitungan lain di bab ini bisa diperiksa dengan cara yang sama: pemangkasan $7^{222}\bmod13$
di bagian tiga, yang memakai $222=12\cdot18+6$ dan memberi $12$, dan pertukaran kunci kecil
$p=23$, $g=5$, $a=4$, $b=3$ yang memberi $A=4$, $B=10$, dan kunci $18$ dari kedua arah.

\begin{resep}
Empat langkah mencari invers $a$ modulo $n$.
\begin{enumerate}
  \item Tulis $(a,n)$. Jalankan Euklid: ganti $(a,n)$ dengan $(n,a\bmod n)$ sampai sisanya
        nol; catat tiap baris $n_0=q\,n_1+r$.
  \item Sisa bukan nol terakhir adalah $\gcd$. Bila bukan $1$, berhenti; inversnya tidak ada.
  \item Dari baris terakhir yang bersisa $1$, tulis $1=r-q\,s$, lalu ganti $s$ dan $r$ dengan
        bentuk baris di atasnya, satu baris ke atas untuk setiap langkah.
  \item Kumpulkan koefisien $a$; itulah inversnya. Bila negatif, tambahkan $n$ sampai masuk
        $0,\dots,n-1$, lalu periksa $a\cdot x\bmod n=1$.
\end{enumerate}
Peringatan urutan: susun mundur dari sisa terkecil naik ke atas, jangan dari baris pertama.
Menukar peran $a$ dan $n$ juga mengubah jawaban: invers $a$ modulo $n$ bukan invers $n$
modulo $a$.
\end{resep}

\kertas{Ulangi perhitungan pertama tanpa melihat teks. Tulis lima baris Euklid maju untuk
$65\,537$ dan $1\,048\,540$, lalu susun mundur sampai memperoleh $d$. Periksa dengan
mengalikan $e$ dan $d$ lalu membaginya dengan $\varphi(n)$: sisanya harus satu. Sesudah itu
jawab: mengapa langkah yang sama tidak bisa dipakai untuk membuka RSA tanpa kunci rahasia,
padahal seluruh perhitungan ini hanya butuh lima pembagian?}

\section{Lab komputasi}
\tingkat{lanjut}

Program berikut memakai pustaka standar saja dan berjalan di bawah satu detik. Ia memeriksa
empat klaim di badan bab: unit modulo $26$, statistik langkah Euklid, kesamaan kunci
Diffie-Hellman, dan kelonggaran RSA polos. Bagian terakhir menguji ramalan ulang tahun pada
SHA-256 yang dipotong menjadi dua byte.

\begin{lab}
\begin{verbatim}
import hashlib, math, random

def egcd(a, b):
    if b == 0:
        return a, 1, 0
    g, x, y = egcd(b, a % b)
    return g, y, x - (a // b) * y

def invers(a, n):
    g, x, _ = egcd(a, n)
    return None if g != 1 else x % n

def langkah(a, b):
    k = 0
    while b:
        a, b = b, a % b
        k += 1
    return k

def orde(a, n):
    k, y = 1, a % n
    while y != 1:
        y = y * a % n
        k += 1
    return k

unit = [a for a in range(26) if invers(a, 26) is not None]
print("unit mod 26:", len(unit), unit)
print("invers 3 mod 26 =", invers(3, 26), " invers 2 mod 26 =", invers(2, 26))

g0, x0, y0 = egcd(1071, 462)
print("gcd(1071,462) =", g0, " x,y =", x0, y0, " langkah =", langkah(1071, 462))

p, g, a, b = 10007, 5, 1234, 5678
A, B = pow(g, a, p), pow(g, b, p)
print("DH: A =", A, " B =", B, " K =", pow(B, a, p), pow(A, b, p))
print("orde(5) =", orde(5, p), " orde(5^5003) =", orde(pow(5, 5003, p), p))

random.seed(29)
L = [langkah(random.getrandbits(32), random.getrandbits(32))
     for _ in range(20000)]
print("euklid 32-bit: rata %.3f, maks %d, min %d" % (sum(L) / len(L), max(L), min(L)))

p, q, e = 1019, 1031, 3
n, phi = p * q, (p - 1) * (q - 1)
d = invers(e, phi)
print("RSA: n =", n, " bit =", n.bit_length(), " phi =", phi, " d =", d)
for m in (424242, 42):
    c = pow(m, e, n)
    print("  m =", m, " c =", c, " D(c) =", pow(c, d, n), " m^3 =", m ** 3)
kali = pow(1234, e, n) * pow(567, e, n) % n
print("E(1234)E(567) mod n =", kali, " D(kali) =", pow(kali, d, n),
      " 1234*567 mod n =", 1234 * 567 % n)

def tabrakan():
    terlihat, k = {}, 0
    while True:
        m = random.getrandbits(32)
        h = hashlib.sha256(m.to_bytes(4, "big")).digest()[:2]
        k += 1
        if h in terlihat:
            return k
        terlihat[h] = m

T = [tabrakan() for _ in range(20)]
print("tabrakan 16-bit: rata %.1f, min %d, maks %d" % (sum(T) / 20, min(T), max(T)))
print("prediksi sqrt(pi/2 * 2^16) =", round(math.sqrt(math.pi / 2 * 2**16), 1))
\end{verbatim}

Keluarannya, apa adanya:

\begin{verbatim}
unit mod 26: 12 [1, 3, 5, 7, 9, 11, 15, 17, 19, 21, 23, 25]
invers 3 mod 26 = 9  invers 2 mod 26 = None
gcd(1071,462) = 21  x,y = -3 7  langkah = 3
DH: A = 4873  B = 3189  K = 5674 5674
orde(5) = 10006  orde(5^5003) = 2
euklid 32-bit: rata 18.750, maks 32, min 6
RSA: n = 1050589  bit = 21  phi = 1048540  d = 699027
  m = 424242  c = 618000  D(c) = 424242  m^3 = 76355615883580488
  m = 42  c = 74088  D(c) = 42  m^3 = 74088
E(1234)E(567) mod n = 551173  D(kali) = 699678  1234*567 mod n = 699678
tabrakan 16-bit: rata 362.7, min 100, maks 743
prediksi sqrt(pi/2 * 2^16) = 320.8
\end{verbatim}
\end{lab}

Baris pertama menegaskan bagian pertama: dua belas unit modulo $26$, dengan $1$ dan $25$
kebalikan dirinya sendiri dan sepuluh sisanya berpasangan; angka $12$ itu sama dengan
$\varphi(26)$.

Baris kedua dan ketiga memperlihatkan dua wajah Euklid: tiga pembagian untuk $1071$ dan $462$
dengan koefisien Bezout $(-3,7)$, dan rata-rata $18{,}75$ pembagian dengan maksimum $32$ untuk
dua ribu pasangan acak $32$ bit. Baris keempat dan kelima adalah pertukaran kunci yang
dijanjikan: $4873$ dan $3189$ di saluran, $5674$ di kedua ujung, sedangkan $5^{5003}$ berorde
$2$ dan karena itu jebakan bagian lima.

Tiga baris RSA memperlihatkan dua hal berlawanan. Untuk $m=424\,242$ semuanya bekerja: sandi
$618\,000$ kembali menjadi $424\,242$. Untuk $m=42$ dengan $e=3$, sandinya $74\,088$, persis
$42^3$, sehingga tidak ada kunci yang perlu dicari, hanya akar pangkat tiga; baris berikutnya
menunjukkan sifat multiplikatif: hasil kali dua sandi membuka pesan hasil kali.

Empat baris terakhir menguji ramalan ulang tahun: rata-rata $362{,}7$ sampel melawan ramalan
$320{,}8$, dengan sebaran lebar seperti yang diharapkan distribusi eksponensial.

\begin{untukdosen}
Urutan mengajar yang saya sarankan: mulai dari sandi mahasiswa di pembuka, lalu lemma
keterdefinisian dan teorema unit pada pertemuan pertama; sisakan Diffie-Hellman dan RSA untuk
pertemuan kedua, dan taruh Shor serta pasca-kuantum sebagai bacaan pertemuan ketiga. Jangan
mengajarkan Euklid sebelum menanyakan ``kapan pembagian ada''.

Dua miskonsepsi yang hampir selalu muncul. Pertama, mahasiswa mengira $\varphi(n)=n-1$ untuk
semua $n$, lalu kaget ketika $\varphi(26)=12$; sembuhkan dengan meminta mereka mendaftar unit
modulo $26$ di papan sebelum rumus apa pun disebut. Kedua, mereka mengira RSA aman karena $n$
sulit difaktorkan, tanpa menyebut ukuran $n$ dan tanpa menyebut bahwa kesulitan itu asumsi;
minta mereka menuliskan klaim keamanan dalam bentuk ``bila ... maka ...''.

Aktivitas papan tulis: dua kelompok menghitung $3^{12}\bmod26$, satu dengan perkalian berulang
dan satu dengan Fermat kecil, lalu bandingkan waktu dan hasilnya. Tanda kesiapan: mereka bisa
menjelaskan tanpa catatan mengapa $[2]_{26}$ tidak punya kebalikan.
\end{untukdosen}

\section{Ringkasan}
\begin{ringkasan}
\begin{itemize}
    \item Aritmetika modulo $n$ bekerja pada kelas sisa, dan pembagian hanya ada untuk kelas yang
        tidak berbagi faktor dengan $n$; Euklid memberi kebalikannya dalam waktu sebanding
        jumlah digit.
    \item Lemma pengacakan, yaitu perkalian dengan unit hanyalah pengacakan urutan, memberi
        teorema Euler dalam lima baris dan Fermat kecil sebagai akibat. Orde membagi
        $\varphi(n)$, dan akar primitif adalah unsur berorde penuh.
    \item Diffie-Hellman menyepakati kunci lewat kesamaan $g^{ab}$ dari dua arah, dengan keamanan
        bersyarat pada kesulitan logaritma diskret; tanpa autentikasi ia jatuh pada serangan
        orang di tengah.
    \item RSA membalik arah: kunci publik mengunci, kunci rahasia membuka, dan buktinya memakai
        Fermat kecil pada kedua faktor prima. RSA polos deterministik, multiplikatif, dan tidak
        melindungi pesan kecil; bantalan menutup ketiganya.
    \item Fungsi hash memberi sidik jari, MAC memberi keaslian antara dua pihak yang berbagi
        kunci, tanda tangan memberi bukti yang bisa diperiksa siapa saja; ketahanan tabrakan
        hanya separuh panjang keluaran, dan angka $362{,}7$ melawan $320{,}8$ di lab
        menunjukkan ramalan itu bekerja.
    \item Semua jaminan di bab ini bersyarat. Shor meruntuhkan RSA dan seluruh logaritma diskret
        bila komputer kuantum berskala besar ada; sandi simetris dan hash cukup menggandakan
        panjang kunci, dan skema lattice seperti ML-KEM adalah pengganti yang mungkin.
\end{itemize}
\end{ringkasan}

\section{Latihan}
\begin{latihan}
\begin{enumerate}
  \item Nilai klaim ini: ``Karena $2^{100}\equiv1\pmod{101}$, bilangan $101$ terbukti prima.''
        Tentukan benar, salah, atau belum cukup bukti, dan sebutkan apa yang dibuktikan
        kongruensi itu.
  \item Buktikan bahwa hasilkali dua unit modulo $n$ juga unit, lalu simpulkan bahwa $U(n)$
        tertutup terhadap perkalian; tunjukkan juga bahwa bila $a$ bukan unit, ada $b\neq0$
        dengan $ab\equiv0\pmod n$.
  \item Buktikan bahwa untuk prima $p$, $a^2\equiv1\pmod p$ hanya dipenuhi $a\equiv1$ atau
        $a\equiv-1$; jelaskan mengapa hal yang sama gagal untuk modulus $8$, dan sebutkan
        semua akar kuadrat dari satu modulo $8$.
  \item Buktikan $\varphi(pq)=\varphi(p)\varphi(q)$ untuk prima berbeda $p,q$ dengan menghitung
        kelas $0\le a<pq$ yang berbagi faktor dengan $pq$, lalu periksa pada $p=1019$ dan
        $q=1031$.
  \item Hitung $17^{-1}\bmod3120$ dengan Euklid mundur, lalu periksa dengan Python. Ulangi
        untuk $65\,537^{-1}\bmod1\,048\,540$ dan bandingkan jumlah pembagian keduanya.
  \item Hitung $5^{100}\bmod101$, $3^{200}\bmod101$, dan $2^{1000}\bmod101$ memakai Fermat
        kecil, tanpa menghitung pangkat penuh. Periksa ketiganya dengan \texttt{pow(a,k,m)}.
  \item Pada modulus $23$ dengan basis $5$, Alice memakai $a=7$ dan Bob $b=15$. Hitung $A$,
        $B$, dan kunci bersama, lalu hitung kunci yang diperoleh penyadap aktif yang mengganti
        $A$ dengan $11$ di depan Bob, dan jelaskan mengapa dua kunci berbeda muncul.
  \item Jelaskan mengapa RSA dengan $e=3$ dan pesan $m<n^{1/3}$ bisa dibuka tanpa kunci rahasia,
        berapa panjang pesan terpanjang yang terkena pada modulus $1\,050\,589$, dan apa yang
        dilakukan bantalan untuk mencegahnya.
  \item Nilai klaim ini: ``ML-KEM sudah terbukti aman, jadi kita bisa berhenti memakai RSA
        sekarang.'' Sebutkan dua hal yang harus benar agar kalimat pertama masuk akal, dan satu
        alasan mengapa kalimat kedua tetap keliru meski yang pertama benar.
  \item Sebuah sistem memakai hash berkeluaran $48$ bit. Berapa sampel yang dibutuhkan untuk
        peluang tabrakan satu per dua, dan berapa bila keluarannya dipotong menjadi $32$ bit?
\end{enumerate}

Petunjuk untuk tiga butir tersulit. Butir 3: tulis $a^2-1=(a-1)(a+1)$ dan pakai lemma Euklid;
untuk modulus $8$, cari akar dengan mencoba $a=0,\dots,7$. Butir 7: hitung $11^7\bmod23$
sebagai kunci yang dilihat Bob, lalu bandingkan dengan $11^{15}\bmod23$ yang dihitung Alice.
Butir 9: bedakan ``terbukti aman'' dari ``aman terhadap serangan yang diketahui''.
\end{latihan}

\begin{untukpeneliti}
Yang benar-benar terbukti di bab ini adalah bagian matematikanya: teorema Bezout, teorema unit,
teorema Euler, lemma orde, dan teorema RSA semuanya punya bukti lengkap dan tidak bergantung
pada dugaan. Yang hanya dipopulerkan adalah kata ``aman''. Tidak ada teorema yang menyatakan
RSA-2048 sulit; yang ada taksiran biaya algoritme terbaik yang diketahui, dan taksiran itu
berubah setiap kali ada kemajuan kriptanalisis.

Sumber primernya: Diffie dan Hellman (1976), bagian 3--4; Rivest, Shamir, dan Adleman (1978),
bagian 2--4; Shor (1997), bagian 5 untuk reduksi faktorisasi ke pencarian orde; Regev (2005),
bagian 3 untuk reduksi LWE; Bleichenbacher (1998); Gidney dan Ekera (2021); FIPS 203 dan 204.

Batas keberlakuan yang perlu disebut: reduksi keamanan biasanya menghasilkan batas longgar,
model serangan tidak memuat saluran samping, dan serangan padding menunjukkan bahwa selisih satu
pesan galat cukup meruntuhkan skema yang terbukti aman dalam model. Pertanyaan terbuka: seberapa
ketat reduksi untuk ML-KEM, apakah ada algoritme kuantum yang lebih cepat untuk masalah lattice,
dan bagaimana membuktikan keamanan skema yang dipakai bersama data rahasia tanpa membocorkan
pola akses.
\end{untukpeneliti}

\peta{Bab 31 pada Diffie dan Hellman (1976), bagian 3 untuk pertukaran kunci; Rivest, Shamir,
dan Adleman (1978), bagian 2--4 untuk RSA dan tanda tangan; Shor (1997), bagian 5--6 untuk
reduksi faktorisasi ke pencarian orde; Regev (2005), bagian 3 untuk reduksi LWE; Bleichenbacher
(1998); Gidney dan Ekera (2021); FIPS 203 dan 204. Pada Bab 2 untuk bijeksi dan partisi, pada
Bab 3 untuk sifat Archimedes dan bentuk diskret sifat infimum, pada Bab 8 untuk eliminasi
Gauss, pada Bab 26 untuk alasan klaim keamanan tidak bisa berbentuk teorema tanpa syarat, dan
pada Bab 29--30 untuk pola ancaman oracle.}

\chapter{Privasi Diferensial dan Fondasi Matematis QML}

Sebuah klinik gigi mengumpulkan catatan sepuluh ribu pasien, melatih model untuk
memprediksi siapa yang berisiko tinggi, lalu menerbitkan model itu sebagai berkas
bobot. Bulan berikutnya seorang mahasiswa mengambil satu nama dari daftar pasien yang
bocor lewat jalur lain, memasukkan rekam medis orang itu ke model, dan membaca skor
keluaran. Skornya jauh lebih tinggi daripada skor orang lain dengan rekam medis
serupa yang bukan pasien klinik itu. Mahasiswa tersebut tidak membobol basis data dan
tidak memakai kata sandi siapa pun. Ia hanya bertanya kepada model, berkali-kali, dan
model itu menjawab terlalu percaya diri untuk orang yang pernah ia lihat saat latihan.

Paralel dengan kejadian itu ada di ujung lain buku ini. Sebuah makalah melaporkan
``keunggulan kuantum'' karena sirkuitnya menyiapkan keadaan dengan $2^{50}$ amplitudo,
angka yang tidak mungkin disimpan komputer mana pun. Angka itu benar. Yang tidak
disebut makalah tersebut: amplitudo tidak bisa dibaca, dan hasil pengukuran hanya
berupa sampel.

Dua kisah itu punya cacat yang sama bentuknya: klaim berjalan lebih cepat daripada
matematika yang menopangnya. Paruh pertama bab membangun definisi yang membuat kalimat
``model ini privat'' bisa diuji, yaitu \emph{privasi diferensial}, lalu menunjukkan
mengapa definisi itu berbentuk perbandingan dua dunia dan bukan larangan mengingat.
Paruh kedua membangun fondasi untuk menilai kalimat ``model kuantum lebih unggul'':
amplitudo kompleks, gerbang uniter, aturan Born, dan biaya simulasi.

\index{privasi diferensial}
\index{keunggulan kuantum}

\kamus{privasi diferensial}{differential privacy}{janji bahwa hasil sebuah prosedur
hampir tidak berubah, dalam arti peluang, ketika satu orang masuk atau keluar dari data
latih}

\begin{untukanak}
Bayangkan sebuah toko es krim mencatat rata-rata uang belanja semua pembeli hari itu,
lalu menempelkan angka itu di pintu. Kalau Sari satu-satunya pembeli yang membeli kue
ulang tahun besar, rata-rata yang ditempel ikut naik, dan tetangga yang tahu Sari
berulang tahun bisa menebak bahwa Sari ada di daftar pembeli hari itu. Angkanya bukan
nama, tetapi angka itu tetap bercerita tentang satu orang.

Sekarang bayangkan toko yang lebih hati-hati. Sebelum menempel angka, ia menambahkan
guncangan kecil yang acak, misalnya tiga puluh sen ke atas atau ke bawah, dengan
peluang sama. Kalau Sari ikut atau tidak ikut, angka yang ditempel nyaris sama, karena
guncangan itu jauh lebih besar daripada pengaruh satu orang. Pembeli lain tetap bisa
tahu rata-rata belanja hari itu dengan cukup tepat, tetapi tidak ada lagi cara
memastikan apakah Sari ada di daftar.

Pertanyaan yang membuat penasaran: kalau guncangannya terlalu kecil, rahasianya bocor;
kalau terlalu besar, angkanya tidak berguna. Di mana batas yang jujur di antaranya, dan
bagaimana kita bisa membuktikan bahwa batas itu benar? Di bab ini kita akan tahu
mengapa janji privasi harus berbentuk perbandingan dua dunia yang nyaris identik, bukan
janji bahwa komputer akan lupa.
\end{untukanak}

\section{Apa yang bocor dari model yang dilatih pada data pribadi}
\tingkat{dasar}

Model tidak menyimpan salinan data latih. Ia menyimpan parameter, dan parameter itu
dipilih agar cocok dengan data. Justru karena pemilihan itulah data meninggalkan jejak,
dan jejak itu bisa dibaca dengan tiga cara yang berbeda tingkat kesulitannya.

Cara pertama, \emph{inferensi keanggotaan}. Penyerang punya satu rekam medis dan ingin
tahu apakah rekam itu ada di data latih. Ia tidak perlu membuka model; cukup mengirim
rekam itu dan membaca keluarannya. Model yang dilatih pada rekam tersebut biasanya
menjawab dengan keyakinan lebih tinggi, karena parameter sudah disesuaikan untuk
memasukkannya. Shokri, Stronati, Song, dan Shmatikov menunjukkan pada 2017 bahwa
perbedaan keyakinan ini cukup untuk dibedakan oleh pengklasifikasi kedua yang dilatih
pada ``model bayangan'', yaitu model tiruan yang dilatih pada data dengan dan tanpa
contoh tertentu. Serangan itu bekerja pada model yang hanya bisa ditanya, tanpa bobot
dan tanpa gradien.

Cara kedua, \emph{inferensi atribut}. Penyerang tidak menanyakan keanggotaan, melainkan
sifat yang tidak ada di data: apakah orang ini merokok, apakah ia hamil. Model yang
dilatih pada populasi tertentu bisa mengungkap korelasi yang tidak seharusnya dipakai
untuk memutuskan apa pun tentang satu orang.

Cara ketiga, \emph{rekonstruksi}: penyerang memulihkan sebagian isi data. Dinur dan
Nissim menunjukkan pada 2003 bahwa rangkaian kueri agregat berderau kecil tetap
membocorkan sebagian besar isi basis data bila kuerinya banyak dan saling bergantung.
Zhu, Liu, dan Han menunjukkan pada 2019 bahwa gradien dalam pelatihan terdistribusi
dapat dipakai untuk memulihkan gambar latih hampir persis.

Ketiga jalur itu menjelaskan mengapa kalimat ``kami hanya menerbitkan parameter, bukan
data'' tidak cukup. Yang perlu dijanjikan bukan ``tidak ada yang disalin'', melainkan
``keluaran prosedur ini hampir tidak bergantung pada satu orang''. Untuk menyatakannya
secara matematis, kita butuh dua definisi jarak: jarak antara dua himpunan data yang
berbeda pada satu orang, dan jarak antara dua sebaran keluaran.

\begin{definisi}[Pasangan data bertetangga]
Dua himpunan data $D$ dan $D'$ atas domain yang sama disebut \emph{bertetangga} bila
keduanya identik kecuali pada satu catatan: $D'$ diperoleh dari $D$ dengan menambah
atau membuang satu catatan, atau dengan mengganti satu catatan dengan catatan lain.
Himpunan semua pasangan seperti itu ditulis $\mathcal{N}$.
\end{definisi}

\index{inferensi keanggotaan}
Pilihan ``tambah atau buang'' disebut kerawanan tak terbatas, ``ganti'' disebut
kerawanan terbatas, dan keduanya memberi angka sensitivitas yang berbeda. Klaim privasi
yang tidak menyebut pilihannya belum lengkap.

\section{Definisi privasi diferensial, dan mengapa bukan rata-rata}
\tingkat{menengah}

``Hampir sama'' diukur dengan rasio peluang, bukan selisih, karena rasio tidak berubah
ketika peluangnya sangat kecil, dan justru peluang yang sangat kecil itulah yang paling
sering dipakai penyerang untuk menyimpulkan sesuatu.

\begin{definisi}[Privasi diferensial $(\varepsilon,\delta)$]
\index{privasi diferensial}
Misalkan $\varepsilon\ge 0$ dan $\delta\in[0,1]$. Prosedur acak $M$ yang memetakan
himpunan data ke ruang keluaran $\mathcal{Y}$ memenuhi privasi diferensial
$(\varepsilon,\delta)$ bila untuk setiap pasangan bertetangga $(D,D')\in\mathcal{N}$ dan
setiap himpunan terukur $S\subseteq\mathcal{Y}$ berlaku
\begin{equation}
P\bigl[M(D)\in S\bigr] \;\le\; e^{\varepsilon}\,P\bigl[M(D')\in S\bigr] + \delta .
\end{equation}
Bila $\delta=0$, prosedur itu memenuhi privasi diferensial murni $\varepsilon$.
\end{definisi}

\begin{figure}[htbp]
\centering
\begin{tikzpicture}
\begin{axis}[width=6.2cm,height=4.1cm,xmin=-2.6,xmax=3.6,ymin=0,ymax=4.6,
  domain=-2.6:3.6,samples=200,axis lines=left,
  xtick={-2,-1,0,1,2,3},ytick={0,2,4},
  tick label style={font=\scriptsize},label style={font=\scriptsize},
  title={$\varepsilon=8$},title style={font=\scriptsize},
  xlabel={keluaran},ylabel={kerapatan},
  legend style={font=\scriptsize,draw=none,at={(0.03,0.97)},anchor=north west},
  legend cell align=left]
  \addplot[thick,accent,fill=accent!12] {4*exp(-8*abs(x))} \closedcycle;
  \addlegendentry{data $D$}
  \addplot[thick,layerteal,densely dashed] {4*exp(-8*abs(x-1))};
  \addlegendentry{data $D'$}
\end{axis}
\end{tikzpicture}
\hfill
\begin{tikzpicture}
\begin{axis}[width=6.2cm,height=4.1cm,xmin=-2.6,xmax=3.6,ymin=0,ymax=0.142,
  domain=-2.6:3.6,samples=200,axis lines=left,
  xtick={-2,-1,0,1,2,3},ytick={0,0.05,0.1},yticklabels={$0$,$0{,}05$,$0{,}1$},
  tick label style={font=\scriptsize},label style={font=\scriptsize},
  title={$\varepsilon=0{,}25$},title style={font=\scriptsize},
  xlabel={keluaran},
  legend style={font=\scriptsize,draw=none,at={(0.03,0.97)},anchor=north west},
  legend cell align=left]
  \addplot[thick,accent,fill=accent!12] {0.125*exp(-0.25*abs(x))} \closedcycle;
  \addlegendentry{data $D$}
  \addplot[thick,layerteal,densely dashed] {0.125*exp(-0.25*abs(x-1))};
  \addlegendentry{data $D'$}
\end{axis}
\end{tikzpicture}
\caption{Sebaran keluaran mekanisme Laplace pada sensitivitas $1$ untuk dua basis
data bertetangga yang berbeda satu catatan, dengan skala tegak berbeda di kedua
panel. Pada $\varepsilon=8$ kedua kurva terpisah jauh, sehingga satu keluaran
sudah cukup untuk menebak data mana yang dipakai. Pada $\varepsilon=0{,}25$
keduanya hampir berimpit: rasio ordinatnya di setiap titik paling besar
$e^{0{,}25}\approx1{,}284$, jadi satu keluaran hampir tidak memihak salah satu
data.}
\label{fig:dua-sebaran}
\end{figure}

Pembaca baru biasanya bertanya dua hal: mengapa perbandingan peluang, bukan rata-rata,
dan mengapa ada $\delta$ sama sekali kalau $\varepsilon$ sudah mengukur kebocoran.

Jawaban pertanyaan pertama adalah inti paruh pertama bab. Perhatikan urutan kuantornya:
\emph{untuk setiap} $S$, \emph{untuk setiap} pasangan bertetangga. Janji itu harus
berlaku bagi uji $S$ mana pun, termasuk uji yang dipilih penyerang setelah ia melihat
mekanismenya. Definisi yang berbentuk rata-rata menempatkan kuantor di dalam: ``ada
harapan yang kecil''. Dua bentuk itu tidak setara; Gambar~\ref{fig:dua-sebaran}
menunjukkan bedanya pada dua nilai $\varepsilon$.

\begin{proposisi}
Misalkan $M$ memenuhi privasi diferensial $(\varepsilon,0)$. Maka untuk setiap algoritme
penyerang $A$ yang keluarannya hanya $0$ atau $1$, dan setiap pasangan bertetangga
$(D,D')$,
\begin{equation}
P\bigl[A\bigl(M(D)\bigr)=1\bigr] \;\le\; e^{\varepsilon}\,
P\bigl[A\bigl(M(D')\bigr)=1\bigr] .
\end{equation}
\end{proposisi}

\begin{proof}
Himpunan $\{y : A(y)=1\}$ terukur, jadi ia boleh dipakai sebagai $S$ pada definisi.
Substitusi langsung memberi hasilnya.
\end{proof}

Proposisi itu terlihat sepele, dan gunanya justru menutup pintu yang sering dipakai
untuk membela mekanisme lemah. Kalau janji hanya berlaku untuk \emph{rata-rata}
keluaran, penyerang bisa merata-ratakan hasil banyak pengulangan sehingga komponen acak
saling menghapus dan yang tinggal adalah sinyal tentang satu orang. Dengan janji
per-kejadian, setiap pengulangan membayar anggaran privasi sendiri.

\begin{bacakeras}
Baca bergantian, dua suara. Suara pertama bertanya, suara kedua menjawab.

``Kalau satu orang keluar dari daftar, apakah jawabanmu berubah?''\\
``Sedikit.''\\
``Berapa sedikit?''\\
``Cukup sedikit sehingga tidak ada yang bisa memastikan.''

Ulangi percakapan itu tiga kali, makin cepat, sampai terasa seperti hafalan. Lalu
suara pertama bertanya lagi, lebih pelan:

``Kalau pertanyaan itu kuulang seratus kali, apakah seratus kali sedikit masih
sedikit?''

Diam satu tarikan napas. Jawab berdua: sedikit yang diulang seratus kali bukan sedikit
lagi. Lalu ucapkan bersama, satu kali: satu orang tidak boleh mengubah kesimpulan.
Ulangi baris itu sekali lagi, lebih lambat, lalu sekali lagi sambil menunjuk satu kursi
kosong di ruangan.
\end{bacakeras}

\begin{contoh}[Jawaban acak]
\index{jawaban acak}
Seorang responden ditanya apakah ia pernah memakai obat terlarang. Prosedurnya: lempar
koin; bila muncul sisi A, jawab jujur; bila muncul sisi B, lempar lagi dan jawab ``ya''
bila gambar, ``tidak'' bila angka. Peluang menjawab ``ya'' bila ia benar-benar memakai
adalah $1/2+1/2\cdot 1/2=3/4$, dan bila tidak memakai adalah $1/2\cdot 1/2=1/4$.
Rasio keduanya $3$, sehingga $\varepsilon=\ln 3=1{,}0986$. Tidak ada jawaban yang
membuktikan apa pun tentang satu orang, dan itulah yang dijanjikan definisi.
\end{contoh}

\begin{tangan}
Bahan: satu koin per orang, kertas, pensil. Sepuluh teman masing-masing menulis satu
rahasia ya/tidak pada sobekan kertas dan melipatnya.

Aturannya satu. Lempar koin: gambar, jawab jujur; angka, lempar lagi, gambar berarti
``ya'' dan angka berarti ``tidak''. Setiap orang hanya menyebut satu kata, dan pencatat
hanya menulis dua kelompok turus.

Sekarang minta satu orang mengganti rahasianya, lalu ulangi seluruh putaran. Lihat
turusnya: bergeser paling banyak satu, dan tidak ada yang bisa menunjuk siapa pun.
Kesimpulan tentang kelompok bertahan. Kesimpulan tentang satu orang tidak pernah muncul.
\end{tangan}

Sekarang $\delta$, yaitu peluang kegagalan: dengan peluang paling besar $\delta$, janji
rasio tadi boleh dilanggar sepenuhnya. Karena itu $\delta$ bukan toleransi yang bisa
dipilih sesuka hati. Bila $\delta$ sebanding dengan $1/n$ untuk $n$ catatan, penyerang
bisa menargetkan satu orang dan membiarkan sisanya sebagai samaran; karena itu $\delta$
harus jauh lebih kecil daripada $1/n$. Angka $\varepsilon$ juga bukan peluang:
$\varepsilon=1$ berarti rasio peluang paling besar $e\approx 2{,}718$, sedangkan
$\varepsilon=8$ berarti rasio sampai $2\,981$.

\section{Sensitivitas global dan mekanisme Laplace}
\tingkat{menengah}

Satu-satunya hal yang menghalangi kita menambah derau secara sembarang adalah
pertanyaan: seberapa besar satu catatan bisa menggeser jawaban?

\begin{definisi}[Sensitivitas global]
\index{sensitivitas global}
Untuk $f$ yang memetakan himpunan data ke $\mathbb{R}^d$, sensitivitas globalnya adalah
\begin{equation}
\Delta f \;=\; \sup_{D\sim D'} \bigl\| f(D)-f(D') \bigr\|_{1},
\end{equation}
dengan supremum diambil atas semua pasangan bertetangga.
\end{definisi}

Kueri hitung, $f(D)=|D|$, punya $\Delta f=1$. Kueri jumlah atas nilai dalam $[0,R]$
punya $\Delta f=R$. Kueri rata-rata atas $n$ catatan punya $\Delta f=R/n$, karena satu
orang hanya menggeser rata-rata sebesar $R/n$; di sinilah terlihat bahwa rata-rata lebih
aman daripada jumlah, tetapi hanya selama $n$ diketahui dan tetap. Sensitivitas juga
bergantung pada pilihan ketetanggaan: bila $D'$ boleh mengganti, kueri hitung punya
$\Delta f=2$.

\begin{teorema}[Mekanisme Laplace]
\index{mekanisme Laplace}
Misalkan $f$ punya sensitivitas global $\Delta f$, dan $Y$ berdistribusi Laplace dengan
skala $b=\Delta f/\varepsilon$, yaitu dengan kepadatan $p_b(y)=\frac{1}{2b}e^{-|y|/b}$.
Maka $M(D)=f(D)+Y$ memenuhi privasi diferensial $(\varepsilon,0)$.
\end{teorema}

\begin{proof}
Ambil $S\subseteq\mathbb{R}$ terukur dan $D\sim D'$. Cukup ditunjukkan bahwa untuk
setiap $y$,
\begin{equation}
\frac{p_b\bigl(y-f(D)\bigr)}{p_b\bigl(y-f(D')\bigr)} \;\le\; e^{\varepsilon}.
\end{equation}
Dengan $a=f(D)$ dan $a'=f(D')$, ruas kiri adalah $\exp\bigl((|y-a'|-|y-a|)/b\bigr)$.
Pertidaksamaan segitiga memberi $|y-a'|-|y-a|\le|a-a'|$, sehingga ruas kiri paling
besar $\exp(|a-a'|/b)\le\exp(\Delta f/b)=e^{\varepsilon}$, karena $\Delta f$ adalah
supremum. Jadi $p_b(y-a)\le e^{\varepsilon}p_b(y-a')$ untuk setiap $y$; integrasi kedua
ruas atas $S$ memberi $P[M(D)\in S]\le e^{\varepsilon}P[M(D')\in S]$, dan karena tidak
ada suku sisa, $\delta=0$.
\end{proof}

Angka contohnya begini. Untuk kueri hitung dengan $\varepsilon=0{,}5$, skala
$b=1/0{,}5=2$; peluang derau melebihi $5$ dalam nilai mutlak adalah $e^{-5/2}=0{,}0821$,
jadi sekitar satu dari dua belas jawaban meleset lebih dari lima orang. Sebaliknya
dengan peluang $0{,}95$ galatnya tidak melebihi $b\ln 20=2\times 2{,}9957=5{,}991$. Di
sini terlihat bentuk trade-off yang berulang sepanjang bab: mengecilkan $\varepsilon$
memperbesar $b$, dan memperbesar $b$ memperburuk ketelitian secara linear.

\begin{resep}
\begin{enumerate}
  \item Tetapkan $f$ dan potong setiap nilai ke rentang $[0,R]$ sebelum apa pun dihitung.
  \item Hitung $\Delta f$ dari kasus terburuk satu catatan pada data yang sudah terpotong.
  \item Pilih $\varepsilon>0$, hitung $b=\Delta f/\varepsilon$, tambahkan derau Laplace
        berskala $b$ pada jawaban.
  \item Catat $\varepsilon$ pada buku anggaran; sisa anggaran adalah yang belum dipakai.
\end{enumerate}
Peringatan urutan: memotong setelah derau ditambahkan mengubah sebaran keluaran dan
membatalkan jaminan, karena $\Delta f$ dihitung dari data yang belum terpotong.
\end{resep}

\kamus{sensitivitas global}{global sensitivity}{seberapa jauh satu orang bisa menggeser
jawaban sebuah kueri pada kasus terburuk; inilah yang menentukan besar derau yang
dibutuhkan}

Kadang derau Laplace tidak cocok, misalnya karena keluarannya berdimensi tinggi. Di situ
dipakai derau Gauss: untuk $f$ dengan sensitivitas $\ell_2$ sebesar $\Delta_2 f$,
mekanisme $M(D)=f(D)+N(0,\sigma^{2}\Delta_2 f^{2}I)$ memenuhi $(\varepsilon,\delta)$
bila $\sigma\ge\sqrt{2\ln(1{,}25/\delta)}/\varepsilon$. Buktinya dilewati di sini; ia
membandingkan dua normal yang digeser dengan pertidaksamaan konsentrasi, dan lengkapnya
ada pada Dwork dan Roth (2014), lampiran A. Syarat itu cukup, bukan perlu: untuk $\sigma$ kecil ia jauh lebih longgar
daripada batas sebenarnya, dan batas ketatnya diberikan mekanisme Gauss analitik Balle
dan Wang (2018).

\section{Komposisi: anggaran privasi dan mengapa ia habis}
\tingkat{lanjut}

Satu mekanisme jarang dipakai sekali. Model dilatih dengan ratusan langkah, masing-masing
membocorkan sedikit, dan yang dibutuhkan adalah totalnya.

\begin{teorema}[Komposisi dasar]
\index{komposisi}
Misalkan $M_1,\dots,M_T$ dipilih secara adaptif, yaitu $M_t$ boleh bergantung pada
keluaran $M_1,\dots,M_{t-1}$, dan setiap $M_t$ memenuhi $(\varepsilon_t,\delta_t)$
untuk semua masukan yang mungkin. Maka keluaran bersama $(M_1,\dots,M_T)$ memenuhi
$\bigl(\sum_t\varepsilon_t,\sum_t\delta_t\bigr)$.
\end{teorema}

\begin{proof}
Cukup untuk $T=2$; sisanya induksi. Dengan hukum peluang total,
\begin{align}
P\bigl[(M_1,M_2)\in S_1\times S_2\bigr]
&= \int_{S_1} P\bigl[M_1(D)\in dy_1\bigr]\,
   P\bigl[M_2(D,y_1)\in S_2 \mid M_1(D)=y_1\bigr] \notag\\
&\le \int_{S_1} e^{\varepsilon_1}P\bigl[M_1(D')\in dy_1\bigr]\,
   \bigl(e^{\varepsilon_2}P[M_2(D',y_1)\in S_2]+\delta_2\bigr) + \delta_1 \notag\\
&\le e^{\varepsilon_1+\varepsilon_2}P\bigl[(M_1,M_2)(D')\in S_1\times S_2\bigr]
   + \delta_1 + \delta_2 .
\end{align}
Langkah pertama memakai jaminan untuk $M_2$ pada $y_1$ tetap, langkah kedua memakai
jaminan untuk $M_1$ dan $P[M_1(D')\in S_1]\le 1$.
\end{proof}

Komposisi dasar menjumlahkan, dan penjumlahan itulah yang membuat anggaran cepat habis:
seratus langkah dengan $\varepsilon=0{,}1$ per langkah sudah menghabiskan
$\varepsilon=10$. Ada batas yang lebih tajam ketika kebocoran per langkah kecil.

\begin{tantangan}
Sebuah kantor menerbitkan rata-rata gaji timnya dengan derau Laplace berskala
$b=\Delta f/\varepsilon$, dengan $\varepsilon=1$ dan $\Delta f=1$. Katanya, ``satu
jawaban sudah memenuhi $\varepsilon=1$, jadi satu orang aman.'' Berapa kali jawaban itu
boleh diterbitkan sebelum rata-rata dari jawaban-jawaban itu cukup tajam untuk menebak
gaji satu orang, dan mengapa?

Petunjuk 1: derau acak punya rata-rata nol; perhatikan apa yang terjadi pada rata-rata
derau ketika jawaban diulang.

Petunjuk 2: galat baku rata-rata $k$ jawaban adalah $b\sqrt{2}/\sqrt{k}$. Dengan $b=1$
dan $k=100$ nilainya $0{,}1414$, sedangkan satu orang bisa menggeser jawaban sebesar
$1$.

Kalau mentok, kembali ke bagian Komposisi.
\end{tantangan}

\begin{teorema}[Komposisi lanjut]
Untuk setiap $\varepsilon'>0$, komposisi $T$ prosedur yang masing-masing
$(\varepsilon,\delta)$ memenuhi $(\varepsilon_{\mathrm{tot}},\delta_{\mathrm{tot}})$
dengan
\begin{equation}
\varepsilon_{\mathrm{tot}} = \sqrt{2T\ln(1/\delta')}\;\varepsilon
+ T\,\varepsilon\bigl(e^{\varepsilon}-1\bigr),
\qquad
\delta_{\mathrm{tot}} = T\delta + \delta' .
\end{equation}
\end{teorema}

Buktinya dilewati: ia diperoleh dengan membatasi momen rasio peluang, atau dengan
\emph{accountant} momen. Rujukannya Dwork dan Roth bagian 3.5, dan Kairouz, Oh, dan
Viswanath (2015) untuk komposisi optimal. Suku pertamanya tumbuh seperti $\sqrt{T}$, bukan $T$.

Namun komposisi lanjut tidak selalu lebih baik, dan ini sering salah dikutip. Suku
kedua, $T\varepsilon(e^{\varepsilon}-1)$, membengkak ketika $\varepsilon$ tidak kecil.
Pada percobaan di bagian sebelas, dengan $T=200$ dan $\delta=10^{-5}$, komposisi dasar
memberi $\varepsilon=36{,}5$ pada $\sigma=32$ sedangkan komposisi lanjut memberi
$20{,}6$; pada $\sigma=2$ keduanya memberi $583{,}7$ dan $11\,302{,}8$. Jadi komposisi
lanjut menang ketika $\varepsilon$ per langkah kecil dan kalah ketika besar.

Satu hasil lagi membuat anggaran bisa dipakai hemat: bila beberapa prosedur bekerja
pada bagian data yang saling lepas, kebocorannya tidak dijumlahkan melainkan diambil
yang terbesar, karena satu catatan hanya ada di satu bagian. Memecah data ke beberapa
pengguna, masing-masing dengan anggaran sendiri, jauh lebih hemat daripada satu pengguna
yang menjalankan semua kueri.

\index{anggaran privasi}

\section{DP-SGD: potong, tambah derau, dan harga yang dibayar}
\tingkat{lanjut}

Pelatihan biasa tidak memenuhi definisi mana pun, karena satu contoh bisa menggeser
gradien sebesar apa pun: tidak ada batas atas untuk $\Delta f$, sehingga derau yang
dibutuhkan tak terhingga. Perbaikannya terdiri dari dua langkah yang saling menopang.

Langkah pertama, pemotongan: setiap gradien per contoh dipotong agar normanya tidak
melebihi $C$,
\begin{equation}
\operatorname{clip}_C(g) \;=\; g\cdot\min\Bigl(1,\;\frac{C}{\|g\|_2}\Bigr).
\end{equation}
Setelah pemotongan, satu contoh paling banyak menyumbang $C$ terhadap jumlah, sehingga
sensitivitas rata-rata atas batch berukuran $B$ adalah $C/B$. Langkah kedua, derau
Gauss dengan simpangan baku $\sigma C/B$ ditambahkan pada rata-rata itu,
\begin{equation}
w_{t+1} = w_t - \eta\Bigl(\frac{1}{B}\sum_{i\in\mathcal{B}_t}
\operatorname{clip}_C\bigl(\nabla\ell_i(w_t)\bigr) + \frac{\sigma C}{B}\,\xi_t\Bigr),
\qquad \xi_t\sim N(0,I).
\end{equation}
Algoritme inilah yang disebut DP-SGD, diperkenalkan Abadi dan rekan-rekan pada 2016
bersama \emph{accountant} momen yang membuat anggarannya bisa dihitung untuk ratusan
ribu langkah.

\begin{teorema}
Setiap langkah DP-SGD memenuhi $(\varepsilon_0,\delta)$ dengan
\begin{equation}
\varepsilon_0=\sqrt{2\ln(1{,}25/\delta)}\,\big/\,\sigma ,
\end{equation}
dan setelah $T$ langkah seluruh prosedur memenuhi $(T\varepsilon_0,T\delta)$.
\end{teorema}

\begin{proof}
Sensitivitas $\ell_2$ rata-rata gradien terpotong adalah $C/B$, karena mengganti satu
contoh menggeser rata-rata paling banyak $C/B$. Terapkan mekanisme Gauss dari bagian
tiga, lalu komposisi dasar dari bagian empat.
\end{proof}

\index{dp-sgd}
Angka dari percobaan kita menjelaskan mengapa derau tidak langsung menghapus
pembelajaran. Dengan $n=200$, $C=1$, dan $\sigma=32$, simpangan baku derau per koordinat
adalah $32/200=0{,}16$, sedangkan norma gradien rata-rata pada titik awal $0{,}83$:
deraunya sekitar seperlima sinyal. Tetapi derau berubah tanda secara acak, sehingga
kontribusinya pada posisi akhir tumbuh seperti $\sqrt{T}$, sedangkan sinyal tumbuh
seperti $T$. Setelah dua ratus langkah sinyal menang, dan model masih mencapai
$93{,}71$ persen.

Yang harus dibayar juga terlihat. Pada $\sigma=128$ akurasi turun ke $85{,}51$ persen
dengan simpangan $6{,}16$ antar seed; pada $\sigma=512$ akurasi $70{,}39$ persen dengan
simpangan $17{,}26$. Simpangan yang membesar itu bukan keanehan: pada derau besar, hasil
pelatihan bergantung pada lemparan derau yang kebetulan terjadi, sehingga satu model
bukan lagi cuplikan yang mewakili prosedurnya. Melaporkan satu seed pada $\sigma$ besar
akan menyesatkan.

Empat batas perlu disebut bersamaan dengan hasil itu. Pertama, pemotongan memasukkan
bias: contoh bergradien besar, biasanya yang sulit atau langka, dipotong lebih keras,
sehingga model kurang memperhatikan mereka. Kedua, DP-SGD menuntut lebih banyak langkah,
dan biaya komputasi itu nyata. Ketiga, jaminan privasi seragam tetapi biaya utilitasnya
tidak: Bagdasaryan, Poursaeed, dan Shmatikov menunjukkan pada 2019 bahwa kelompok
minoritas bisa menderita penurunan akurasi jauh lebih besar pada anggaran yang sama. Keempat, dan paling penting:
privasi diferensial menjamin sesuatu tentang \emph{individu}, bukan tentang
\emph{kebenaran model}. Model dengan $\varepsilon=0{,}1$ masih bisa salah total, masih
bisa diskriminatif, dan masih bisa dipakai untuk keputusan yang tidak adil. Privasi
adalah janji tentang seberapa banyak satu orang bisa diketahui dari keluaran; ia sama
sekali bukan janji bahwa keluarannya bermutu.

\kamus{anggaran privasi}{privacy budget}{jumlah total $\varepsilon$ yang boleh dipakai
sebuah sistem; setiap kueri atau langkah pelatihan membelanjakannya, dan setelah habis
tidak ada lagi jaminan}

\section{Mengapa amplitudo memakai bilangan kompleks}
\tingkat{dasar}

Sekarang paruh kedua, dimulai dari pertanyaan yang sering dianggap teknis padahal
jawabannya konseptual. Mengapa keadaan kuantum diwakili bilangan kompleks, bukan
bilangan real?

Jawaban pendeknya: karena dua jalur menuju hasil yang sama harus bisa saling menghapus,
dan penghapusan semacam itu menuntut tanda yang bisa berputar, bukan hanya positif dan
negatif. Dengan bilangan real, amplitudo hanya punya dua tanda dan interferensi hanya
bisa menambah atau mengurangi. Dengan bilangan kompleks, amplitudo punya fase yang
berputar kontinu, sehingga dua jalur bisa bertemu pada sudut apa pun dan hasilnya
bervariasi mulus.

Bilangan kompleks $z=a+bi$ dengan $i^2=-1$ punya modulus $|z|=\sqrt{a^2+b^2}$, konjugat
$\bar z=a-bi$, dan fase $\theta$ dalam penulisan
$z=|z|e^{i\theta}=|z|(\cos\theta+i\sin\theta)$. Modulus dan fase berperan berbeda dan
tidak bisa ditukar: modulus menentukan besar peluang, fase menentukan tanda interferensi.
Perkalian dua bilangan kompleks mengalikan modulus dan menjumlahkan fase, dan sifat itu
membuat fase berguna sebagai koordinat keadaan.

\kamus{amplitudo}{probability amplitude}{bilangan kompleks yang modulus kuadratnya adalah
peluang; berbeda dari peluang karena amplitudo bisa saling menghapus}

Ruang keadaan satu qubit adalah $\mathbb{C}^2$, dan ruang keadaan $n$ qubit adalah
$\mathbb{C}^{2^n}$ dengan hasilkali dalam
$\langle\phi|\psi\rangle=\sum_k\overline{\phi_k}\psi_k$. Dalam notasi bra-ket,
$|\psi\rangle$ adalah vektor kolom dan $\langle\psi|$ transpos konjugatnya. Keadaan
selalu dinormalkan, $\langle\psi|\psi\rangle=1$, karena jumlah seluruh peluang
harus satu.

\begin{definisi}[Ruang Hilbert berdimensi hingga]
\index{ruang Hilbert}
Ruang Hilbert berdimensi hingga adalah ruang vektor kompleks berdimensi hingga
dilengkapi hasilkali dalam yang lengkap. Dalam buku ini satu-satunya contoh yang dipakai
adalah $\mathbb{C}^{d}$ dengan hasilkali dalam baku, biasanya $d=2^{n}$.
\end{definisi}

Dua sifat notasi ini perlu dipegang sejak awal. Pertama, fase global tidak punya arti
fisik: $|\psi\rangle$ dan $e^{i\theta}|\psi\rangle$ memberi peluang yang sama untuk
setiap hasil pengukuran, karena $|e^{i\theta}a_k|^2=|a_k|^2$. Kedua, fase
\emph{relatif} punya arti penuh: mengubah fase satu komponen saja mengubah hasil
interferensi, dan itulah yang dipakai algoritme kuantum. Perbedaan antara fase global tak terlihat dan
fase relatif terlihat menjadi jelas pada contoh hitung di bagian sembilan.

\begin{proposisi}
Bila $|\psi\rangle$ dan $|\phi\rangle$ keduanya bernorma satu, maka
$|\langle\phi|\psi\rangle|\le 1$.
\end{proposisi}

\begin{proof}
Pertidaksamaan Cauchy-Schwarz memberi
$|\langle\phi|\psi\rangle|^2\le\langle\phi|\phi\rangle\langle\psi|\psi\rangle=1$, lalu
ambil akar kuadratnya.
\end{proof}

\section{Gerbang sebagai matriks uniter, dan mengapa harus uniter}
\tingkat{menengah}

Perubahan keadaan karena gerbang adalah perkalian matriks, $|\psi'\rangle=U|\psi\rangle$.
Tidak setiap matriks boleh dipakai. Syaratnya satu: peluang total harus tetap satu untuk
setiap keadaan awal, sehingga
$\langle U\psi|U\psi\rangle=\langle\psi|\psi\rangle$ untuk semua $\psi$.

\begin{definisi}[Matriks uniter]
\index{gerbang uniter}
Matriks persegi $U$ disebut uniter bila $U^{\dagger}U=UU^{\dagger}=I$, dengan
$U^{\dagger}$ transpos konjugat $U$.
\end{definisi}

\begin{teorema}
Untuk matriks persegi $U$ atas $\mathbb{C}^{d}$, pernyataan berikut setara: (i) $U$
uniter; (ii) $\langle U\phi|U\psi\rangle=\langle\phi|\psi\rangle$ untuk semua $\phi,\psi$;
(iii) $\|U\psi\|=\|\psi\|$ untuk semua $\psi$.
\end{teorema}

\begin{proof}
(i) $\Rightarrow$ (ii): $\langle U\phi|U\psi\rangle=\phi^{\dagger}U^{\dagger}U\psi
=\phi^{\dagger}\psi$. (ii) $\Rightarrow$ (iii): ambil $\phi=\psi$. (iii)
$\Rightarrow$ (i): dari $\langle U\psi|U\psi\rangle=\langle\psi|\psi\rangle$ untuk semua
$\psi$, gunakan identitas polarisasi untuk menyimpulkan
$\langle\phi|U^{\dagger}U\psi\rangle=\langle\phi|\psi\rangle$ bagi semua $\phi,\psi$;
karena itu $U^{\dagger}U=I$, dan untuk matriks persegi di dimensi hingga invers kiri
berarti invers kanan.
\end{proof}

Teorema itu menjelaskan tiga hal. Pertama, gerbang harus uniter karena peluang harus
kekal; tidak ada alasan lain yang diperlukan. Kedua, setiap gerbang punya invers
$U^{\dagger}$ yang juga uniter, sehingga rangkaian bisa dibalik dan tidak ada informasi
yang hilang selama gerbang berlangsung. Ketiga, satu-satunya operasi tidak uniter adalah
pengukuran, dan di situlah informasi hilang. Karena itu komputasi kuantum bukan
komputasi yang menghapus informasi, melainkan komputasi yang menunda penghapusan sampai
akhir.

\begin{proposisi}[Tanpa penggandaan]
Tidak ada $U$ uniter pada dua qubit sehingga
$U\bigl(|\psi\rangle\otimes|0\rangle\bigr)=|\psi\rangle\otimes|\psi\rangle$ untuk semua
$|\psi\rangle$.
\end{proposisi}

\begin{proof}
Andaikan ada, dan ambil $|\psi\rangle,|\phi\rangle$ bernorma satu. Karena $U$ uniter,
hasilkali dalam kedua ruas sama: $\langle\phi|\psi\rangle=\langle\phi|\psi\rangle^2$.
Persamaan $x=x^2$ hanya punya solusi $0$ atau $1$, sehingga setiap pasangan keadaan
harus saling tegak lurus atau identik; itu mustahil untuk lebih dari dua keadaan.
\end{proof}

\kamus{uniter}{unitary}{sifat matriks yang mempertahankan panjang vektor; dalam bahasa
sehari-hari, gerbang yang tidak menambah dan tidak mengurangi peluang total}

\section{Fourier sebagai perubahan basis, dan kaitannya dengan konvolusi}
\tingkat{lanjut}

Transformasi Fourier adalah perubahan basis, dan itulah cara paling berguna
memandangnya di sini. Pada ruang berdimensi $N$, basis baku diganti basis $|f_k\rangle$
yang komponennya $(|f_k\rangle)_j=\omega^{jk}/\sqrt{N}$ dengan $\omega=e^{2\pi i/N}$.
Matriksnya $F_{jk}=\omega^{jk}/\sqrt{N}$.

\begin{proposisi}
Matriks Fourier $F$ uniter.
\end{proposisi}

\begin{proof}
Hitung $(F^{\dagger}F)_{jk}=\frac{1}{N}\sum_{l=0}^{N-1}\omega^{-lj}\omega^{lk}
=\frac{1}{N}\sum_{l=0}^{N-1}\omega^{l(k-j)}$. Bila $k=j$ setiap suku bernilai satu
sehingga hasilnya $1$. Bila $k\neq j$, ini deret geometri dengan rasio
$\omega^{k-j}\neq1$ dan jumlahnya $\frac{1-\omega^{N(k-j)}}{1-\omega^{k-j}}=0$ karena
$\omega^{N}=1$. Jadi $F^{\dagger}F=I$.
\end{proof}

Karena uniter, $F$ sah dipakai sebagai gerbang, dan versi $n$ qubitnya disebut
transformasi Fourier kuantum. Yang membuatnya menarik bukan hanya bahwa ia gerbang yang
sah, melainkan bahwa ia mengubah konvolusi menjadi perkalian titik demi titik.

\begin{teorema}[Konvolusi menjadi perkalian]
\index{konvolusi}
Untuk $x,y\in\mathbb{C}^{N}$, definisikan konvolusi siklik
$(x*y)_k=\sum_{j=0}^{N-1}x_j\,y_{(k-j)\bmod N}$. Maka
\begin{equation}
F(x*y) \;=\; \sqrt{N}\,\bigl(Fx\bigr)\odot\bigl(Fy\bigr),
\end{equation}
dengan $\odot$ perkalian komponen demi komponen.
\end{teorema}

\begin{proof}
Komponen ke-$k$ ruas kiri adalah $\frac{1}{\sqrt{N}}\sum_m\omega^{km}\sum_j x_j
y_{m-j}$. Tukar urutan penjumlahan dan tulis $r=m-j$:
\[
\frac{1}{\sqrt{N}}\sum_j x_j\omega^{kj}\sum_r\omega^{kr}y_r
=\sqrt{N}\,(Fx)_k\,(Fy)_k .
\]
Pertukaran urutan penjumlahan sah karena kedua jumlah berhingga.
\end{proof}

Pertukaran dua jumlah berhingga itu tampak sepele, tetapi versi tak hingga dari langkah
yang sama adalah isi Bab 14: deret Fourier hanya boleh diturunkan suku demi suku bila
kekonvergenannya cukup seragam. Kaitan dengan Bab 15 lebih konseptual: perubahan basis
adalah perubahan variabel, dan matriks uniter adalah perubahan variabel yang determinan
Jacobiannya bermodulus satu, sehingga volume dan peluang tidak berubah.

Kaitan dengan komputasi kuantum ada pada biayanya. Transformasi Fourier diskret atas
$N=2^{n}$ titik dihitung klasik dalam $\mathcal{O}(N\log N)$ operasi dengan FFT.
Versi kuantumnya disusun dari $\mathcal{O}(n^{2})$ gerbang. Selisih itu nyata dan
menjadi dasar algoritme Shor. Tetapi ada syarat yang sering dihilangkan: setelah
transformasi selesai, kita tidak bisa membaca $N$ amplitudo, dan yang bisa dilakukan
hanyalah mengambil sampel. Kecepatan pada penyusunan keadaan tidak otomatis menjadi
kecepatan pada tugas yang ingin diselesaikan; ia harus dipasangkan dengan prosedur yang
mengekstrak jawaban dari beberapa sampel saja. Banyak klaim keunggulan kuantum di bidang
pembelajaran mesin tidak memiliki prosedur semacam itu.

\section{Aturan Born, dan mengapa keluaran selalu probabilistik}
\tingkat{menengah}

Sampai di sini semuanya deterministik: diberikan keadaan dan gerbang, keadaan
berikutnya tunggal. Peluang masuk pada satu langkah yang tidak bisa dihindari, yaitu
pengukuran.

\begin{definisi}[Aturan Born]
\index{aturan Born}
Bila keadaan $|\psi\rangle=\sum_k a_k|k\rangle$ diukur pada basis ortonormal
$\{|k\rangle\}$, peluang memperoleh hasil $k$ adalah
\begin{equation}
p(k) \;=\; \bigl|\langle k|\psi\rangle\bigr|^{2} \;=\; |a_k|^{2},
\end{equation}
dan setelah hasil $k$ diperoleh, keadaan menjadi $|k\rangle$.
\end{definisi}

Mengapa harus begitu, dan bukan sekadar ketidaktahuan kita tentang keadaan yang
``sebenarnya''? Jawaban terkuat bukan filosofis melainkan eksperimental: ketidaksamaan
Bell, yang diuji Aspect, Grangier, dan Roger pada 1982, menunjukkan bahwa korelasi antara
dua pengukuran terpisah melanggar batas yang harus dipatuhi setiap teori yang menyimpan
nilai pasti tersembunyi. Jadi probabilisme di sini bukan cara mengatakan ``kami belum
tahu''; ia struktur teori. Konsekuensi praktisnya: satu kali pengukuran memberi satu
hasil, bukan seluruh keadaan.

\kamus{aturan Born}{Born rule}{aturan yang mengubah amplitudo menjadi peluang dengan
mengkuadratkan modulusnya; satu-satunya tempat peluang masuk ke dalam teori}

Konsekuensi bagi pembelajaran mesin sering diremehkan. Untuk menaksir $p(k)$ dengan
galat baku $\gamma$ dibutuhkan sekitar $p(1-p)/\gamma^{2}$ tembakan. Pada percobaan kita,
keadaan $R_y(60^\circ)|0\rangle$ punya peluang sebenarnya $0{,}75$ dan $0{,}25$; dengan
seribu tembakan tercatat $727$ dan $273$, taksiran $0{,}727$, berbeda $0{,}023$ dari
nilai sebenarnya. Galat baku teoretisnya $\sqrt{0{,}75\cdot 0{,}25/1000}=0{,}0137$, jadi
selisih itu sekitar $1{,}7$ galat baku, kejadian biasa. Untuk tiga angka di belakang koma dibutuhkan sekitar sejuta
tembakan, dan setiap tembakan berarti satu persiapan keadaan dan satu pengukuran. Di
sinilah faktor $\sqrt{N}$ muncul, dan di sinilah banyak ``kecepatan kuadratik'' berubah
menjadi kecepatan biasa setelah biaya pembacaan dihitung.

\section{Contoh hitung: satu dan dua qubit dengan tangan}
\tingkat{menengah}

\begin{untukanak}
Dua anak melempar koin. Anak pertama melempar satu koin, lalu menutup tangannya dan
bertanya: gambar atau angka? Hasilnya satu, dan tidak ada yang bisa tahu apa pun
sebelum koin itu dibuka. Anak kedua menyiapkan dua koin dengan cara yang aneh: ia
menyuruh koin-koin itu ``berunding'' lebih dulu, sehingga keduanya selalu sepakat, tanpa
pernah sepakat lebih dulu mau sama-sama gambar atau sama-sama angka. Ketika koin pertama
dibuka dan menunjukkan gambar, koin kedua pasti gambar. Ketika koin pertama menunjukkan
angka, koin kedua pasti angka.

Yang membuat aneh: tidak ada cara membuat dua koin biasa berperilaku begitu. Kalau
keduanya sudah ditentukan sejak awal, orang yang membuat koin tahu keduanya. Pada koin
yang berunding itu tidak ada catatan semacam itu, dan percobaan sungguhan menunjukkan
tidak ada catatan tersembunyi yang menjelaskan hasilnya.

Contoh terakhir itulah yang disebut keadaan Bell, dan ia dipakai bab ini untuk
menunjukkan mengapa ``nol'' dan ``sangat kecil'' adalah dua hal yang berbeda.
\end{untukanak}

Semua angka di bagian ini diperiksa ulang dengan NumPy pada bagian sebelas, dan pembaca
bisa memeriksanya dengan kalkulator biasa.

Gerbang Hadamard pada satu qubit adalah
$H=\frac{1}{\sqrt{2}}\begin{pmatrix}1&1\\1&-1\end{pmatrix}$. Pada $|0\rangle=(1,0)$ ia
memberi
\[
H|0\rangle=\frac{1}{\sqrt{2}}(1,1)=(0{,}7071;\;0{,}7071),
\]
sehingga peluang hasil $0$ dan $1$ masing-masing $0{,}5$.

Sekarang rangkaian $H$, lalu $Z=\operatorname{diag}(1,-1)$, lalu $H$ lagi. Langkah
pertama memberi $(0{,}7071;0{,}7071)$; langkah kedua membalik tanda komponen kedua
menjadi $(0{,}7071;-0{,}7071)$; langkah ketiga mengalikan dengan $H$:
\[
H\begin{pmatrix}0{,}7071\\-0{,}7071\end{pmatrix}
=\frac{1}{2}\begin{pmatrix}0{,}7071-0{,}7071\\0{,}7071+0{,}7071\end{pmatrix}
=\begin{pmatrix}0\\1\end{pmatrix}.
\]
Hasilnya $|1\rangle$ dengan peluang satu. Amplitudo komponen $|0\rangle$ menjadi nol
bukan karena ada yang dibuang, melainkan karena dua jalur bertemu dengan tanda
berlawanan. Inilah interferensi, dan inilah alasan fase relatif penting: bila $Z$
diganti identitas, hasilnya kembali $|0\rangle$.

Contoh ketiga memakai rotasi. Gerbang
$R_y(\theta)=\begin{pmatrix}\cos(\theta/2)&-\sin(\theta/2)\\
\sin(\theta/2)&\cos(\theta/2)\end{pmatrix}$ memutar keadaan pada bidang real. Dengan
$\theta=60^\circ$ dan keadaan awal $|0\rangle$,
\[
R_y(60^\circ)|0\rangle=(\cos 30^\circ;\;\sin 30^\circ)=(0{,}8660;\;0{,}5),
\]
sehingga peluangnya $\cos^{2}30^\circ=0{,}75$ dan $\sin^{2}30^\circ=0{,}25$. Perhatikan
bahwa $|a_0|^2+|a_1|^2=0{,}75+0{,}25=1$.

Contoh keempat dua qubit. Mulai dari $|00\rangle$, kerjakan $H$ pada qubit pertama, lalu
CNOT dengan kendali qubit pertama dan sasaran qubit kedua. Setelah $H$, keadaannya
$\frac{1}{\sqrt{2}}(|00\rangle+|10\rangle)$. CNOT membalik qubit kedua hanya bila qubit
pertama $1$, sehingga
\[
|\Phi^{+}\rangle=\frac{1}{\sqrt{2}}\bigl(|00\rangle+|11\rangle\bigr)
=(0{,}7071;\;0;\;0;\;0{,}7071).
\]
Peluang $00$ dan $11$ masing-masing $0{,}5$; peluang $01$ dan $10$ nol. Dua amplitudo di
tengah nol karena kedua jalur yang menuju sana saling menghapus, bukan karena keduanya
tidak mungkin. Akibatnya, bila kedua qubit diukur, hasilnya selalu sama: $00$ atau
$11$, tidak pernah campuran.

\kertas{Periksa contoh di atas tanpa melihat teks. (a) Hitung $HZH|0\rangle$ dengan
menuliskan setiap perkalian matriks, lalu jelaskan mengapa hasilnya $|1\rangle$ dan
bukan $|0\rangle$. (b) Ganti $R_y(60^\circ)$ dengan $R_y(90^\circ)$, hitung kedua
peluangnya, lalu jelaskan mengapa dua keadaan berbeda bisa memberi peluang yang sama.
(c) Untuk keadaan Bell di atas, hitung peluang memperoleh $01$ dan jelaskan mengapa
angka itu nol, bukan ``sangat kecil''.}

\section{Mengapa simulasi klasik meledak, dan apa artinya untuk klaim keunggulan}
\tingkat{lanjut}

Keadaan $n$ qubit adalah vektor di $\mathbb{C}^{2^{n}}$, jadi ia punya $2^{n}$ amplitudo
kompleks. Untuk $n=10$ itu $1\,024$ bilangan, tidak lebih besar daripada satu citra
kecil. Untuk $n=50$ itu $2^{50}=1{,}1259\times10^{15}$ bilangan; dengan $16$ byte per
amplitudo, yang dibutuhkan $2^{54}$ byte, yaitu tepat $16$ pebibyte. Ia batas yang
tumbuh dua kali setiap kali satu qubit ditambahkan.

Percobaan pada bagian sebelas mengukur akibatnya. Menyiapkan keadaan dengan satu gerbang
Hadamard pada setiap qubit butuh $0{,}2$ milidetik untuk $n=10$, $2{,}3$ untuk $n=14$,
$81{,}4$ untuk $n=18$, dan $1\,357$ milidetik untuk $n=22$. Kolom waktu itu paling
berisik dari seluruh bab: dua kali jalan pada mesin yang sama bisa berbeda hampir dua
kali lipat. Yang tidak berubah adalah polanya, yaitu garis lurus pada skala logaritmik
terhadap $2^{n}$.

Apa yang boleh dan tidak boleh disimpulkan? Yang boleh: keadaan kuantum umum tidak bisa
disimpan secara klasik untuk $n$ besar, sehingga algoritme yang benar-benar memakai
seluruh $2^{n}$ amplitudo tidak bisa langsung ditiru. Yang tidak boleh: bahwa setiap
sirkuit kuantum lebih cepat. Ada tiga alasan.

Alasan pertama, keadaan tidak bisa dibaca. Satu pengukuran memberi satu hasil, sedangkan
simulasi klasik yang menyimpan seluruh vektor bisa menghitung peluang setiap hasil. Jadi
lebih banyak informasi tersedia di simulasi daripada di perangkat kerasnya, dan
ketidakmampuan menyalin keadaan bukan bukti keunggulan.

Alasan kedua, sirkuit yang bisa disimulasikan mencakup banyak sirkuit yang dipakai dalam
praktik: kedalaman kecil, keterkaitan terbatas, keadaan hasilkali, dan gerbang Clifford
punya tiruan klasik yang efisien. Goguen, Clarke, Martens, dan Selman menunjukkan
keluarga sirkuit itu bisa dihitung klasik dalam waktu polinomial, dan keluarga itu cukup
luas untuk mencakup banyak ansatz variasional.

Alasan ketiga, keunggulan bisa hilang oleh struktur data. Ewin Tang menunjukkan pada
2019 bahwa algoritme rekomendasi kuantum yang dianggap memberi percepatan eksponensial
punya padanan klasik dengan percepatan polinomial, karena struktur matriks datanya bisa
dieksploitasi tanpa kuantum. Huang, Kueng, dan Preskill (2021) memberi pemisahan yang
ketat, tetapi hanya untuk keluarga distribusi data tertentu.

\kamus{plateau tandus}{barren plateau}{daerah pada lanskap kerugian sirkuit variasional
tempat gradien mengecil secara eksponensial terhadap jumlah qubit, sehingga pelatihan
tidak bergerak}

Pernyataan yang jujur tentang QML pada 2026 karena itu berbentuk begini. Ada bukti
pemisahan untuk tugas tertentu dalam model orakel, dan ada klaim teoretis tentang
percepatan untuk tugas pembelajaran dengan asumsi distribusi tertentu. Tidak ada bukti
keunggulan praktis yang mapan untuk tugas pembelajaran mesin pada data klasik, dalam
arti percobaan yang diulang pada perangkat keras nyata dan mengalahkan metode klasik
terbaik pada ukuran data yang berarti. Cerezo dan rekan-rekan (2021, 2022) menyebut
\emph{plateau tandus}, yaitu gradien yang menghilang secara eksponensial pada sirkuit
dalam, sebagai hambatan yang belum terselesaikan, bukan rincian teknis.

\section{Lab komputasi: anggaran privasi, akurasi, dan ukuran keadaan}
\tingkat{lanjut}

Program berikut menjalankan dua percobaan di atas tanpa pustaka selain NumPy: DP-SGD
pada regresi logistik dua kelas dengan lima seed per simpangan derau, lalu beberapa
keadaan kuantum dan biaya simulasinya. Semua angka di bab ini berasal dari keluaran di
bawah, kecuali yang ditandai sebagai hitungan tangan.

\begin{lab}
\begin{verbatim}
import numpy as np
import time

D, N, T, C, DELTA, LR = 8, 200, 200, 1.0, 1e-5, 0.5


def data(rng, n, mu=0.6):
    y = 2 * rng.integers(0, 2, n) - 1
    return rng.normal(mu * y[:, None], 1.0, (n, D)), y


def sig(t):
    return 1.0 / (1.0 + np.exp(-t))


def clip(G):
    nrm = np.linalg.norm(G, axis=1, keepdims=True)
    return G * np.minimum(1.0, C / np.maximum(nrm, 1e-9))


def dp_sgd(X, y, sigma, seed):
    rng = np.random.default_rng(seed)
    w = np.zeros(D)
    for _ in range(T):
        G = -(1.0 - sig(y * (X @ w)))[:, None] * (y[:, None] * X)
        w -= LR * (clip(G).mean(0) + rng.normal(0.0, sigma * C / len(X), D))
    return w


def eps_basic(s):
    return T * np.sqrt(2 * np.log(1.25 * T / DELTA)) / s


def eps_adv(s):
    d = DELTA / 2
    e = np.sqrt(2 * np.log(1.25 * T / d)) / s
    return np.sqrt(2 * T * np.log(1 / d)) * e + T * e * (np.exp(e) - 1)


rng = np.random.default_rng(11)
Xn, yn = data(rng, N)
Xe, ye = data(rng, 3000)
print("sigma  akurasi      eps_dsr  eps_lanjut")
A0 = [100 * np.mean(np.sign(Xe @ dp_sgd(Xn, yn, 1e-9, k)) == ye) for k in range(5)]
print("%5s  %5.2f+-%4.2f  %8s  %9s" % ("0", np.mean(A0), np.std(A0), "-", "-"))
for s in (2.0, 8.0, 32.0, 128.0, 512.0):
    A = [100 * np.mean(np.sign(Xe @ dp_sgd(Xn, yn, s, k)) == ye) for k in range(5)]
    print("%5.0f  %5.2f+-%4.2f  %8.1f  %9.1f"
          % (s, np.mean(A), np.std(A), eps_basic(s), eps_adv(s)))

I2 = np.eye(2)
Z = np.diag([1.0, -1.0])
S = np.diag([1.0, 1.0j])
H = np.array([[1.0, 1.0], [1.0, -1.0]]) / np.sqrt(2)
CNOT = np.array([[1, 0, 0, 0], [0, 1, 0, 0], [0, 0, 0, 1], [0, 0, 1, 0]], float)


def ry(t):
    c, s = np.cos(t / 2), np.sin(t / 2)
    return np.array([[c, -s], [s, c]])


def kron1(U, n, q):
    out = np.array([[1.0]])
    for j in range(n):
        out = np.kron(out, U if j == q else I2)
    return out


s0 = np.array([1.0, 0.0])
print()
print("H|0>        ", np.round(H @ s0, 4))
print("HZH|0>      ", np.round(H @ Z @ (H @ s0), 4))
print("HSH|0>      ", np.round(H @ S @ (H @ s0), 4))
print("Ry(60)|0>   ", np.round(ry(np.pi / 3) @ s0, 4),
      " peluang", np.round((ry(np.pi / 3) @ s0) ** 2, 4))
psi = CNOT @ (kron1(H, 2, 0) @ np.kron(s0, s0))
print("Bell        ", np.round(psi, 4), " peluang", np.round(psi ** 2, 4))
p1 = (ry(np.pi / 3) @ s0) ** 2
sh = np.bincount(np.random.default_rng(4).choice(2, 1000, p=p1), minlength=2)
print("1000 tembakan", sh, " taksiran", np.round(sh / 1000, 3))
print()
print("  n   amplitudo      memori (GB)   waktu (ms)")
for n in (10, 14, 18, 22):
    best = 1e9
    for _ in range(3):
        v = np.zeros(2 ** n, dtype=complex)
        v[0] = 1.0
        t0 = time.perf_counter()
        for q in range(n):
            v = v.reshape([2] * n)
            v = np.moveaxis(np.tensordot(H, v, axes=([1], [q])), 0, q).reshape(-1)
        best = min(best, 1000 * (time.perf_counter() - t0))
    print("%3d  %11.3e  %12.3e  %10.1f" % (n, 2.0 ** n, 2.0 ** n * 16 / 2 ** 30, best))
\end{verbatim}

Keluaran program itu, apa adanya:

\begin{verbatim}
sigma  akurasi      eps_dsr  eps_lanjut
    0  95.20+-0.00         -          -
    2  95.13+-0.05     583.7    11302.8
    8  95.03+-0.13     145.9      216.5
   32  93.71+-0.81      36.5       20.6
  128  85.51+-6.16       9.1        3.7
  512  70.39+-17.26       2.3        0.8

H|0>         [0.7071 0.7071]
HZH|0>       [0. 1.]
HSH|0>       [0.5+0.5j 0.5-0.5j]
Ry(60)|0>    [0.866 0.5  ]  peluang [0.75 0.25]
Bell         [0.7071 0.     0.     0.7071]  peluang [0.5 0.  0.  0.5]
1000 tembakan [727 273]  taksiran [0.727 0.273]

  n   amplitudo      memori (GB)   waktu (ms)
 10    1.024e+03     1.526e-05         0.2
 14    1.638e+04     2.441e-04         2.3
 18    2.621e+05     3.906e-03        81.4
 22    4.194e+06     6.250e-02      1356.6
\end{verbatim}
\end{lab}

Baris nol adalah garis dasar: $95{,}20$ persen dengan simpangan nol. Baris $\sigma=2$
memberi $\varepsilon=583{,}7$, jaminan yang praktis kosong karena rasio $e^{583}$ tidak
membatasi apa pun; baris $\sigma=8$ menurunkan $\varepsilon$ ke $145{,}9$ tanpa
kehilangan akurasi terukur, jadi derau kecil nyaris gratis. Harga baru terasa pada
$\sigma=32$: $\varepsilon$ turun ke $36{,}5$ atau $20{,}6$ dengan biaya $1{,}49$ poin,
lalu $85{,}51$ persen pada $\sigma=128$, dan $70{,}39$ persen pada $\sigma=512$ dengan
simpangan antar seed $17{,}26$ poin, sehingga satu model tunggal tidak lagi mewakili
prosedur.

Blok kedua memverifikasi hitungan tangan di bagian sembilan. $HZH|0\rangle$ memberi
$(0;1)$, $R_y(60^\circ)$ memberi $(0{,}8660;0{,}5)$, dan keadaan Bell muncul sebagai
$(0{,}7071;0;0;0{,}7071)$ dengan dua nol yang eksak. Baris $HSH|0\rangle$ memberi
$(0{,}5+0{,}5i;\;0{,}5-0{,}5i)$, dan pasangan itu dipakai pada
Gambar~\ref{fig:fase}.

Blok ketiga mengukur biaya simulasi. Memori tumbuh tepat dua kali per qubit, dari
$1{,}5\times10^{-5}$ GB pada $n=10$ menjadi $6{,}25\times10^{-2}$ GB pada $n=22$.
Waktunya tumbuh lebih cepat daripada memori pada rentang ini karena alasan arsitektur,
bukan matematis; pada mesin dengan cache lebih besar lonjakannya bergeser ke $n$ yang
lebih tinggi, tetapi kemiringannya tetap.

\begin{figure}
\centering
\begin{tikzpicture}
\begin{axis}[
  width=0.86\linewidth, height=6.0cm,
  xlabel={$\varepsilon$ (komposisi dasar, skala log)},
  ylabel={akurasi uji (\%)},
  xmode=log, log basis x=10,
  xmin=1, xmax=1200, ymin=62, ymax=100,
  xtick={1,10,100,1000},
  ytick={65,70,75,80,85,90,95,100},
  grid=major, grid style={black!12},
  tick label style={font=\footnotesize},
  label style={font=\small},
  legend style={font=\footnotesize, at={(0.03,0.03)}, anchor=south west},
]
\addplot[only marks, mark=*, mark size=2.2pt, black] coordinates {
  (2.3,70.39) (9.1,85.51) (36.5,93.71) (145.9,95.03) (583.7,95.13)
};
\addlegendentry{DP-SGD, rata-rata 5 seed}
\addplot[black, dashed, domain=1:1200, samples=2] {95.20};
\addlegendentry{tanpa derau: $95{,}20$}
\node[font=\scriptsize, anchor=south west] at (axis cs:2.3,70.39) {$\sigma{=}512$};
\node[font=\scriptsize, anchor=south west] at (axis cs:9.1,85.51) {$\sigma{=}128$};
\node[font=\scriptsize, anchor=south west] at (axis cs:36.5,93.71) {$\sigma{=}32$};
\node[font=\scriptsize, anchor=north east] at (axis cs:145.9,95.03) {$\sigma{=}8$};
\end{axis}
\end{tikzpicture}
\caption{Akurasi uji terhadap anggaran privasi pada percobaan Lab komputasi ($n=200$,
$d=8$, $T=200$, $\delta=10^{-5}$, lima seed). Setiap titik adalah hasil eksekusi nyata,
berpasangan dengan kolom \texttt{eps\_dsr} pada keluaran. Garis putus-putus adalah
akurasi tanpa derau. Sumbu mendatar logaritmik: menurunkan $\varepsilon$ dari $36{,}5$
ke $2{,}3$ berharga $23$ poin akurasi, sedangkan dari $583{,}7$ ke $145{,}9$ nyaris
gratis.}
\label{fig:tradeoff}
\end{figure}

\begin{figure}
\centering
\begin{tikzpicture}[scale=1.15]
\draw[->, black!60] (-1.25,0) -- (1.45,0) node[right, font=\footnotesize]
  {$\operatorname{Re}$};
\draw[->, black!60] (0,-1.1) -- (0,1.25) node[above, font=\footnotesize]
  {$\operatorname{Im}$};
\draw[black!35, dashed] (0,0) circle (0.7071);
\draw[black!55, thin] (0.7071,0) -- (0.7071,0.7071) -- (0,0.7071);
\draw[black!55, thin] (0.5,0) -- (0.5,0.5) -- (0,0.5);
\draw[->, very thick, black] (0,0) -- (0.7071,0.7071)
  node[above left, font=\footnotesize] {$a_0$};
\draw[->, very thick, black] (0,0) -- (-0.7071,0.7071);
\draw[->, thick, black!55] (0,0) -- (0.5,0.5);
\draw[->, thick, black!55] (0,0) -- (0.5,-0.5) node[below right, font=\footnotesize]
  {$a_1$};
\node[font=\scriptsize, anchor=west] at (0.75,0.16) {$\tfrac{1}{\sqrt2}=0{,}7071$};
\node[font=\scriptsize, anchor=north] at (0.5,-0.55) {$0{,}5-0{,}5i$};
\node[font=\scriptsize, anchor=south] at (0.5,0.56) {$0{,}5+0{,}5i$};
\draw[black!35] (0.09,0) arc (0:45:0.42);
\node[font=\scriptsize] at (0.30,0.16) {$45^\circ$};
\end{tikzpicture}
\caption{Amplitudo sebagai panah pada bidang kompleks. Panah tebal adalah keadaan
$H|0\rangle$ dengan $a_0=a_1=1/\sqrt2$; panah tipis adalah $HSH|0\rangle$ dengan
$S=\operatorname{diag}(1,i)$, yaitu $0{,}5+0{,}5i$ dan $0{,}5-0{,}5i$. Modulusnya tetap
$0{,}7071$, jadi peluang kedua hasil tetap $1/2$, tetapi fasenya bergeser
$45^\circ$ ke arah berlawanan. Pengukuran tidak membedakan kedua keadaan itu; gerbang
berikutnya bisa. Koordinat panah diambil dari keluaran Lab komputasi, bukan ilustrasi.}
\label{fig:fase}
\end{figure}

Gambar~\ref{fig:tradeoff} adalah bentuk kuantitatif trade-off privasi-utilitas: kurvanya
tidak linear, dan bagian termahal ada di ujung kiri, di wilayah yang paling sering dipuji
di abstrak. Gambar~\ref{fig:fase} adalah bentuk geometris kalimat ``modulus menentukan
peluang, fase menentukan interferensi''.

\begin{untukdosen}
Urutan mengajar yang saya sarankan: mulai dari Gambar~\ref{fig:tradeoff}, bukan dari
definisi. Tunjukkan bahwa kurva itu punya ujung kiri yang mahal, lalu tanya mengapa.
Siswa biasanya menebak ``karena deraunya besar'', yang benar tetapi belum cukup; yang
perlu muncul adalah bahwa anggaran privasi dijumlahkan sepanjang langkah, sehingga
mengecilkan $\varepsilon$ per langkah dibayar dua kali, pada derau dan pada jumlah
langkah.

Dua miskonsepsi yang hampir selalu muncul. Pertama, mahasiswa membaca
``$(\varepsilon,\delta)$-DP'' sebagai ``model ini aman''. Yang benar, ia janji tentang
pengaruh satu catatan pada sebaran keluaran, dan model itu masih bisa salah total. Kedua,
mahasiswa mengira fase kuantum hiasan notasi karena ``peluang hanya bergantung pada
modulus''. Contoh $HZH|0\rangle$ mematahkannya dalam tiga baris.

Aktivitas papan tulis yang berhasil di kelas kami: minta dua siswa menghitung
$HZH|0\rangle$ dan $HHH|0\rangle$ secara terpisah, lalu bandingkan. Keduanya memakai
tiga gerbang, beda hanya satu matriks di tengah, dan hasilnya $|1\rangle$ melawan
$|0\rangle$. Kelas siap lanjut bila seorang mahasiswa bisa menjelaskan mengapa
$HSH|0\rangle$ memberi peluang sama dengan $H|0\rangle$ tetapi keadaan berbeda.
\end{untukdosen}

\section{Ringkasan}

\begin{ringkasan}
\begin{itemize}
  \item Model yang dilatih pada data pribadi membocorkan informasi lewat tiga jalur:
        keanggotaan, atribut, dan rekonstruksi, dan yang ketiga bekerja dari gradien.
  \item Privasi diferensial $(\varepsilon,\delta)$ membandingkan dua sebaran keluaran
        untuk data yang berbeda pada satu orang, dengan kuantor ``untuk setiap himpunan
        keluaran''. Kuantor itu yang membuat janjinya berlaku terhadap penyerang, bukan
        hanya terhadap rata-rata.
  \item Mekanisme Laplace dengan skala $b=\Delta f/\varepsilon$ memenuhi
        $(\varepsilon,0)$; buktinya hanya membandingkan rasio dua kepadatan. Mekanisme
        Gauss memakai $\sigma\ge\sqrt{2\ln(1{,}25/\delta)}/\varepsilon$ dan memberi
        $\delta>0$.
  \item Komposisi dasar menjumlahkan $\varepsilon$; komposisi lanjut memberi suku
        $\sqrt{T}$ yang lebih baik hanya bila $\varepsilon$ per langkah kecil: pada
        $\sigma=32$ hasilnya $36{,}5$ versus $20{,}6$, sedangkan pada $\sigma=2$
        hasilnya $583{,}7$ versus $11\,302{,}8$.
  \item DP-SGD memotong gradien per contoh pada norma $C$ lalu menambah derau
        $\sigma C/B$. Akurasi turun dari $95{,}20$ persen tanpa derau menjadi $93{,}71$
        persen pada $\varepsilon=36{,}5$ dan $70{,}39$ persen pada $\varepsilon=2{,}3$.
  \item Privasi diferensial menjamin sesuatu tentang individu, bukan kebenaran model;
        ia tidak menghapus bias, dan biaya utilitasnya bisa tidak merata antar kelompok.
  \item Keadaan kuantum hidup di $\mathbb{C}^{2^{n}}$; gerbang harus uniter karena
        peluang harus kekal, dan pengukuran satu-satunya langkah tidak uniter. Keluaran
        karena itu selalu berupa sampel.
  \item Simulasi klasik butuh $2^{n}$ amplitudo: $n=50$ berarti $16$ pebibyte. Fakta itu
        membatasi peniruan, tetapi tidak membuktikan keunggulan, dan sampai sekarang
        belum ada bukti keunggulan praktis QML yang mapan.
\end{itemize}
\end{ringkasan}

\section{Latihan}
\tingkat{menengah}

\begin{latihan}
\begin{enumerate}
  \item Nilai klaim ini: ``Model kami memenuhi privasi diferensial dengan
        $\varepsilon=0{,}5$, jadi tidak ada informasi pribadi yang bisa bocor darinya.''
        Nyatakan benar, salah, atau belum cukup bukti, dan sebutkan dua hal yang harus
        benar agar klaim itu berlaku.
  \item Buktikan bahwa mekanisme Laplace dengan skala $b=\Delta f/\varepsilon$ memenuhi
        $(\varepsilon,0)$ bila sensitivitas diukur dengan norma $\ell_1$ pada keluaran
        berdimensi $d>1$. Tunjukkan di mana argumen satu dimensi perlu diganti.
  \item Buktikan komposisi paralel: bila $M_1$ bekerja pada $D_1$ dan $M_2$ pada $D_2$
        dengan $D_1\cap D_2=\varnothing$, dan masing-masing $(\varepsilon_i,\delta_i)$,
        maka keluarannya memenuhi $(\max_i\varepsilon_i,\max_i\delta_i)$.
  \item Hitung $\varepsilon$ untuk jawaban acak bila peluang menjawab jujur adalah $p$
        dan sisanya dijawab acak dengan koin seimbang. Nyatakan $\varepsilon$ sebagai
        fungsi $p$, lalu tentukan $p$ yang memberi $\varepsilon=2$.
  \item Untuk kueri hitung dengan $\varepsilon=0{,}5$ pada $n=10\,000$, hitung peluang
        galat mutlaknya melebihi $10$ orang, lalu bandingkan dengan peluang yang sama
        bila kuerinya rata-rata dengan $R=1$. Periksa kedua angka dengan simulasi.
  \item Hitung $H\,S\,H|0\rangle$ dengan $S=\operatorname{diag}(1,i)$ selangkah demi
        langkah, lalu bandingkan peluangnya dengan $H|0\rangle$. Jelaskan mengapa hasil
        ukurnya sama padahal keadaannya berbeda, dan gerbang apa yang membedakannya.
  \item Untuk keadaan Bell $|\Phi^{+}\rangle$, hitung peluang hasil $01$ bila qubit
        kedua diukur lebih dulu dan hasilnya $0$; ulangi untuk hasil $1$. Jelaskan
        mengapa kedua jawaban menunjukkan korelasi sempurna.
  \item Tunjukkan bahwa $\begin{pmatrix}1&0\\0&e^{i\varphi}\end{pmatrix}$ uniter untuk
        setiap $\varphi$, lalu jelaskan mengapa gerbang itu tidak mengubah peluang pada
        basis komputasi tetapi bisa mengubah hasil gerbang berikutnya.
  \item Hitung jumlah amplitudo dan kebutuhan memori untuk $n=40$ dan $n=60$ qubit
        dengan $16$ byte per amplitudo. Nyatakan dalam satuan yang masuk akal, lalu
        bandingkan dengan kapasitas penyimpanan yang Anda miliki.
  \item Sebuah makalah menulis: ``Sirkuit kami menyiapkan keadaan dengan $2^{60}$
        amplitudo, jadi ia mengungguli simulasi klasik untuk pembelajaran mesin.''
        Sebutkan tiga pertanyaan yang harus dijawab sebelum kalimat itu diterima, dan
        untuk setiap pertanyaan sebutkan hasil seperti apa yang akan membuat Anda
        menerimanya.
\end{enumerate}

\noindent\textbf{Petunjuk untuk tiga butir tersulit.} Butir 2: tulis rasio kepadatan
sebagai hasil kali rasio per koordinat, lalu gunakan $\sum_j|a_j-a'_j|\le\Delta f$.
Butir 3: pada setiap langkah hanya satu dari dua mekanisme yang dapat berubah. Butir 10:
apakah seluruh amplitudo dipakai, apakah keluarannya bisa diekstrak dengan jumlah
pengukuran yang tidak eksponensial, dan apakah ada pembanding klasik pada tugas dan data
yang sama.
\end{latihan}

\begin{untukpeneliti}
Tiga hal di paruh pertama sering dikutip tanpa syaratnya. Pertama, definisi
$(\varepsilon,\delta)$ di sini adalah definisi Dwork, McSherry, Nissim, dan Smith (2006)
dalam bentuk buku Dwork dan Roth; syarat $\delta\ll 1/n$ berasal dari argumen penyamaran
satu individu, bukan dari kenyamanan analisis. Kedua, batas Gauss
$\sigma\ge\sqrt{2\ln(1{,}25/\delta)}/\varepsilon$ longgar untuk $\sigma$ kecil; batas
ketatnya diberikan mekanisme Gauss analitik Balle dan Wang (2018), dan selisihnya bisa
mencapai faktor dua pada $\varepsilon<1$. Ketiga, \emph{accountant} momen Abadi dan
rekan-rekan (2016) memakai subsampling, sehingga angka $\varepsilon$ yang dilaporkan
tanpa laju subsampling dan jumlah langkah tidak bisa dibandingkan dengan angka lain.

Di paruh kedua, yang benar-benar terbukti adalah pernyataan bersyarat: pemisahan
kompleksitas untuk keluarga tugas tertentu, dan ketidakmampuan simulasi klasik untuk
keadaan umum. Yang populer tetapi tidak terbukti adalah klaim bahwa sirkuit variasional
memberi keunggulan pada data klasik. Tang (2019) dan Huang, Kueng, Preskill (2021)
membatasi klaim itu dari dua arah: struktur data bisa menghapus percepatan, dan
pemisahan yang ada bergantung pada asumsi distribusi. McClean dan rekan-rekan (2018)
menunjukkan gradien yang menghilang secara eksponensial pada sirkuit dalam, sehingga
pelatihan variasional punya hambatan yang belum terselesaikan.

Pertanyaan yang paling penting bukan ``apakah ada percepatan'', melainkan ``percepatan
seperti apa yang bertahan setelah biaya pembacaan, penyusunan keadaan, dan koreksi galat
dihitung''. Sampai ada jawaban yang diukur pada perangkat keras dengan data nyata,
setiap klaim keunggulan praktis QML sebaiknya dibaca sebagai hipotesis.
\end{untukpeneliti}

\peta{Paruh privasi mengikuti Dwork dan Roth, \emph{The Algorithmic Foundations of
Differential Privacy} (bab 2--3 dan lampiran A), Abadi dkk., \emph{Deep Learning with
Differential Privacy} (2016) untuk DP-SGD, Shokri dkk. (2017) untuk inferensi keanggotaan, Dinur dan
Nissim (2003) untuk rekonstruksi dari kueri agregat, serta Zhu, Liu, Han (2019) untuk
rekonstruksi dari gradien. Paruh kuantum mengikuti Nielsen dan Chuang,
\emph{Quantum Computation and Quantum Information} (bab 2 dan 4), Cerezo dkk. (2021, 2022) untuk batas algoritme
variasional, Tang (2019) untuk dekuantisasi, dan Huang, Kueng, Preskill (2021) untuk
pemisahan berbasis data. Prasyaratnya ada pada Bab 10 (SVD dan ortogonalitas), Bab 16
(peluang bersyarat), dan Bab 17 (variabel acak), sedangkan kaitan Fourier dan konvolusi
ada pada Bab 14 dan Bab 15.}


% ---------------------------- lampiran A ----------------------------
\appendix
\renewcommand{\appendixname}{Lampiran}
\chapter{Tabel Notasi}
\label{app:notasi}

Notasi berikut berlaku di seluruh buku. Bila sebuah bab memakai simbol yang tidak
tercantum di sini, simbol itu diperkenalkan di tempat pemakaiannya. Angka di ujung
setiap butir menunjuk bab yang memakainya.

\section*{Himpunan, logika, dan bilangan (Bab 1--5, 22)}
\begin{description}[style=nextline,itemsep=0.25em]
  \item[$\mathbb{N}$, $\mathbb{Z}$, $\mathbb{Q}$, $\mathbb{R}$, $\mathbb{C}$]
        bilangan asli, bulat, rasional, real, dan kompleks.
  \item[$\forall$, $\exists$, $\neg$] untuk semua, ada, dan bukan.
  \item[$\Rightarrow$, $\Leftrightarrow$] berakibat, dan jika dan hanya jika.
  \item[$\in$, $\notin$, $\subseteq$, $\cup$, $\cap$, $\setminus$] keanggotaan,
        bukan anggota, himpunan bagian, gabungan, irisan, selisih (Bab 2).
  \item[$\mathcal{P}(S)$, $|S|$, $S\times T$] himpunan kuasa, kardinalitas,
        hasilkali kartesius (Bab 2).
  \item[$\sup S$, $\inf S$] batas atas terkecil dan batas bawah terbesar (Bab 3).
  \item[$a_n\to a$] barisan $(a_n)$ konvergen ke $a$ (Bab 4).
  \item[$\sum_{n=1}^{\infty} a_n$] deret dan jumlahnya (Bab 5).
  \item[$\mathrm{fl}(x)$, $\mathrm{ulp}$] hasil pembulatan bilangan real dan satuan
        terkecil yang dapat diwakili (Bab 22).
\end{description}

\section*{Aljabar linear (Bab 6--10)}
\begin{description}[style=nextline,itemsep=0.25em]
  \item[$\mathbf{x}$, $\mathbf{x}^{\top}$, $\|\mathbf{x}\|$] vektor kolom,
        transpos, dan norma Euklides.
  \item[$\operatorname{span} S$, $\dim V$] rentang himpunan vektor dan dimensi
        ruang vektor.
  \item[$\mathcal{L}(V,W)$] ruang semua transformasi linear dari $V$ ke $W$ (Bab 7).
  \item[$\ker T$, $\operatorname{im} T$] inti dan jangkauan transformasi (Bab 7).
  \item[$\det A$, $\operatorname{tr} A$, $A^{-1}$] determinan, jejak, dan invers
        matriks (Bab 8).
  \item[$\lambda$, $\mathbf{v}$] nilai eigen dan vektor eigen (Bab 9).
  \item[$\sigma_i$, $A=U\Sigma V^{\top}$] nilai singular dan dekomposisi nilai
        singular (Bab 10).
  \item[$A^{+}$, $\kappa(A)$] pseudoinvers dan bilangan kondisi (Bab 10, 20).
\end{description}

\section*{Analisis dan kalkulus (Bab 11--16)}
\begin{description}[style=nextline,itemsep=0.25em]
  \item[$\varepsilon$, $\delta$] bilangan positif kecil dalam definisi limit,
        kekontinuan, dan kekontinuan seragam.
  \item[$f'(x)$, $\nabla f$] turunan biasa dan gradien.
  \item[$J_f$, $H_f$] matriks Jacobian dan Hessian (Bab 12).
  \item[$\int_a^b f$] integral Riemann pada $[a,b]$ (Bab 13).
  \item[$\sum_{\mathcal{P}} M_i \Delta x_i$] jumlah atas untuk partisi
        $\mathcal{P}$ (Bab 13).
  \item[$f_n \rightrightarrows f$] barisan fungsi konvergen seragam (Bab 14).
  \item[$\dfrac{\partial(x,y)}{\partial(u,v)}$] determinan Jacobian perubahan
        variabel (Bab 15).
  \item[$x'=f(t,x)$, $\dot{x}=F(x)$] persamaan diferensial biasa dan sistem
        otonom, dengan $'$ dan $\dot{\ }$ sama-sama menyatakan turunan terhadap
        waktu (Bab 16).
  \item[$\varphi_t(x_0)$] aliran: keadaan pada waktu $t$ bila mulai dari $x_0$
        (Bab 16).
  \item[$x^{*}$, $J_F(x^{*})$] titik setimbang dan matriks Jacobian di titik itu;
        tanda bagian real nilai eigennya menentukan kestabilan (Bab 16).
\end{description}

\section*{Peluang dan statistika (Bab 17--21, 26)}
\begin{description}[style=nextline,itemsep=0.25em]
  \item[$P(A)$, $P(A\mid B)$] peluang kejadian $A$ dan peluang bersyarat.
  \item[{$\mathbb{E}[X]$, $\operatorname{Var}X$}] harapan dan variansi.
  \item[$\operatorname{Cov}(X,Y)$] kovarians dua variabel acak (Bab 18).
  \item[$\bar{X}_n$, $\mathrm{se}$] rata-rata sampel dan galat bakunya.
  \item[$\mathrm{CI}_{1-\alpha}$] interval kepercayaan dengan tingkat
        $1-\alpha$ (Bab 20).
  \item[$\mathrm{TP}$, $\mathrm{FP}$, $\mathrm{FN}$, $\mathrm{TN}$] empat sel
        matriks konfusi (Bab 27).
  \item[$\hat{\beta}$, $\hat{y}=X\hat{\beta}$] koefisien regresi kuadrat
        terkecil dan nilai yang diprediksi (Bab 21).
  \item[$e_i=y_i-\hat{y}_i$, $R^2$] residu dan proporsi variansi yang dijelaskan
        model (Bab 21).
  \item[$\mathrm{se}(\hat{\beta}_j)$, $\mathrm{VIF}_j$] galat baku koefisien dan
        faktor pembengkakan variansi (Bab 21).
\end{description}

\section*{Optimasi, pembelajaran mesin, dan verifikasi (Bab 23--30)}
\begin{description}[style=nextline,itemsep=0.25em]
  \item[$\min f_0(x)$ dengan $f_i(x)\le 0$] bentuk baku masalah optimasi (Bab 23).
  \item[$L(x,\lambda,\nu)$] Lagrangian; $\lambda,\nu$ pengali (Bab 23).
  \item[$c^{\top}x$, $Ax\le b$, $x\in\mathbb{Z}^n$] bentuk LP dan program integer
        (Bab 24).
  \item[$\sigma(\cdot)$, $\mathrm{softmax}$] fungsi aktivasi dan normalisasi
        eksponensial (Bab 25).
  \item[$\nabla_{\mathbf{x}}\ell$] gradien fungsi kerugian terhadap masukan;
        bahan bakar serangan (Bab 25, 29).
  \item[$L_2$, $L_1$, $\lambda_{\text{reg}}$] jenis regularisasi dan kekuatannya
        (Bab 26).
  \item[$\|\delta\|_{\infty}\le\varepsilon$] batas gangguan pada serangan
        adversarial (Bab 29, 30).
  \item[$\vdash$, \texttt{by}, \texttt{sorry}] keterbuktian, pembuka blok taktik,
        dan bukti yang ditangguhkan di Lean~4 (Bab 28).
\end{description}

% ---------------------------- lampiran B ----------------------------
\chapter{Glosarium}
\label{app:glosarium}

Istilah di bawah dipakai dengan arti teknis yang tetap sepanjang buku. Padanan
Inggris dicantumkan karena sebagian besar literatur memakai istilah Inggris.

\begin{description}[style=nextline,itemsep=0.35em]
  \item[Adjoin] matriks kofaktor yang disusun transpos; dipakai menuliskan invers sebagai satu rumus, $A^{-1}=\operatorname{adj}(A)/\det A$ (\emph{adjugate}).
  \item[Adjoint] dalam turunan balik, bilangan yang menyatakan seberapa besar kerugian akhir terpengaruh oleh satu besaran di tengah jaringan; satu adjoint per besaran (\emph{adjoint}).
  \item[Akar primitif] anggota yang bila dipangkatkan berulang melewati semua anggota grup siklik lain sebelum kembali ke satu; pada kriptografi, posisi awal yang mencapai seluruh posisi kunci (\emph{primitive root}).
  \item[Aksioma] pernyataan yang diterima tanpa bukti sebagai titik awal; seluruh isi bab bergantung padanya (\emph{axiom}).
  \item[Aksioma kelengkapan] janji bahwa setiap himpunan tak kosong yang terbatas ke atas di $\mathbb{R}$ punya supremum; inilah yang membuat garis bilangan tidak berlubang (\emph{completeness axiom}).
  \item[Amplitudo] amplitudo peluang: bilangan kompleks yang modulus kuadratnya memberi peluang; berbeda dari peluang karena dua amplitudo dapat saling meniadakan (\emph{probability amplitude}).
  \item[Anggaran privasi] jumlah $\varepsilon$ yang boleh dibelanjakan sebuah sistem; setiap kueri atau langkah pelatihan mengurangi sisanya (\emph{privacy budget}).
  \item[Antiturunan] fungsi yang turunannya adalah fungsi yang sedang diintegralkan; arah pulang dari sebuah turunan (\emph{antiderivative}).
  \item[Aturan Born] aturan yang mengubah amplitudo menjadi peluang dengan mengkuadratkan modulusnya (\emph{Born rule}).
  \item[Aturan rantai] laju perubahan fungsi tersusun sama dengan hasil kali laju pada setiap mata rantainya (\emph{chain rule}).
  \item[Aturan tiga] bila tidak ada kejadian dalam $n$ percobaan, laju kejadiannya masih boleh setinggi sekitar $3/n$ pada selang kepercayaan 95 persen (\emph{rule of three}).
  \item[Bantalan] isi tambahan pada paket sebelum dikirim agar setiap paket berukuran tetap, sehingga ukuran paket tidak lagi membocorkan isinya (\emph{padding}).
  \item[Barisan] deretan bilangan yang diberi nomor urut $a_1,a_2,\dots$; urutannya bagian dari definisinya, bukan sesuatu yang ditambahkan kemudian (\emph{sequence}).
  \item[Barisan Cauchy] barisan yang sukunya makin rapat satu sama lain, tanpa perlu menyebut nilai yang dituju (\emph{Cauchy sequence}).
  \item[Barisan monoton] barisan yang hanya naik atau hanya turun, tidak pernah berbalik arah (\emph{monotone sequence}).
  \item[Basis] himpunan vektor yang bebas linear dan merentang seluruh ruang (\emph{basis}).
  \item[Bentuk baku] penulisan seragam masalah optimasi: satu sasaran, kendala berbentuk tidak lebih dari nol, dan peubah tak negatif (\emph{standard form}).
  \item[Bias-variansi] pembagian galat model menjadi galat karena terlalu sederhana dan galat karena terlalu peka pada data (\emph{bias-variance tradeoff}).
  \item[Bijeksi] pemasangan sempurna dua himpunan: injektif dan surjektif sekaligus, tanpa sisa di kedua sisi (\emph{bijection}).
  \item[Bilangan kondisi] ukuran seberapa besar galat relatif keluaran dapat membesar dibanding galat relatif masukan (\emph{condition number}).
  \item[Bootstrap] penaksiran sebaran statistik dengan mengambil ulang sampel dari data sendiri, berkali-kali, dengan pengembalian (\emph{bootstrap}).
  \item[BPDA] pengganti gradien pada potongan yang tidak terdiferensiasi dengan gradien fungsi di titik sesudahnya, dipakai agar serangan tetap dapat dihitung (\emph{backward pass differentiable approximation}).
  \item[Celah integralitas] jarak antara nilai terbaik relaksasi LP dan nilai terbaik jawaban bulat; ia mengukur seberapa jauh model bersambung menipu (\emph{integrality gap}).
  \item[Defective] matriks yang nilai eigennya berulang tetapi ruang eigennya terlalu sempit untuk menjadi basis (\emph{defective}).
  \item[Dekomposisi nilai singular (SVD)] penulisan matriks sebagai $U\Sigma V^{\top}$ dengan $U,V$ ortogonal (\emph{singular value decomposition}).
  \item[Deret] penjumlahan suku-suku barisan tanpa henti, yang nilainya dibaca sebagai limit jumlah parsial (\emph{series}).
  \item[Deret harmonik] deret $1+1/2+1/3+\cdots$; sukunya menuju nol tetapi jumlahnya menyimpang (\emph{harmonic series}).
  \item[Deret pangkat] deret berbentuk $\sum a_n(x-c)^n$; ia berlaku hanya di dalam jari-jari konvergensinya (\emph{power series}).
  \item[Deret Taylor] penulisan fungsi sebagai jumlah suku turunan di satu titik, beserta sisa yang mengukur galat pemotongan. (\emph{Taylor series}).
  \item[Determinan] satu bilangan yang mengukur faktor perubahan volume dan sekaligus menguji invertibilitas (\emph{determinant}).
  \item[Determinan Jacobian] faktor skala volume lokal ketika koordinat diganti; nilai mutlaknya mengalikan luas atau volume pada integral lipat (\emph{Jacobian determinant}).
  \item[Diagonalisasi] penggantian basis sampai matriksnya menjadi diagonal, sehingga pekerjaannya tinggal perkalian angka (\emph{diagonalization}).
  \item[Dimensi] banyaknya vektor dalam sebuah basis; angkanya milik ruang, bukan milik satu basis yang kebetulan dipilih (\emph{dimension}).
  \item[Distribusi eksponensial] sebaran waktu tunggu pada proses tanpa penuaan, dengan fungsi kepadatan $\lambda e^{-\lambda t}$ (\emph{exponential distribution}).
  \item[Distribusi normal] sebaran lonceng simetris yang ditentukan sepenuhnya oleh rata-rata dan variansinya (\emph{normal distribution}).
  \item[Dualitas] perumusan masalah optimasi dari sisi lain, yang nilai optimalnya memberi batas bawah bagi masalah asal (\emph{duality}).
  \item[Ekor berat] sebaran yang peluang nilai raksasanya mengecil terlalu lambat, sehingga satu pengamatan dapat mengalahkan seluruh sisanya (\emph{heavy tail}).
  \item[Eliminasi] penggantian sebuah sistem dengan sistem yang lebih sederhana tanpa mengubah himpunan solusinya (\emph{elimination}).
  \item[Entropi silang] ukuran kejutan $-\sum_i y_i\log \hat{p}_i$; makin yakin model pada jawaban yang salah, makin besar dendanya (\emph{cross-entropy}).
  \item[EOT] serangan yang memaksimumkan harapan kerugian atas keacakan pertahanan, biasanya dengan merata-ratakan gradien pada beberapa salinan acak (\emph{expectation over transformation}).
  \item[Epsilon mesin] jarak antara $1$ dan bilangan mesin berikutnya; ukuran satu langkah pada skala satu (\emph{machine epsilon}).
  \item[Epsilon-delta] gaya mendefinisikan limit dengan menyebut toleransi keluaran lebih dulu, lalu batas masukan yang menjaminnya. (\emph{epsilon-delta definition}).
  \item[Faktor integrasi] pengali yang mengubah ruas kiri persamaan diferensial linear menjadi satu turunan hasil kali (\emph{integrating factor}).
  \item[Fungsi aktivasi] pemetaan tak linear pada keluaran sebuah unit; ia menentukan seberapa besar sinyal diteruskan dan bentuk turunannya (\emph{activation function}).
  \item[Fungsi distribusi kumulatif] fungsi $F(x)=P(X\le x)$; tangga menanjak yang tingginya menyatakan peluang yang sudah terkumpul sampai $x$ (\emph{cumulative distribution function}).
  \item[Fungsi kerugian] bilangan yang ingin diperkecil saat model dilatih; ia mengukur seberapa salah prediksi, bukan seberapa benar (\emph{loss function}).
  \item[Fungsi konveks] fungsi yang grafiknya selalu di bawah tali penghubung dua titiknya; setiap minimum lokalnya adalah minimum global (\emph{convex function}).
  \item[Fungsi Lyapunov] fungsi bernilai positif yang menurun sepanjang lintasan sistem; keberadaannya membuktikan kestabilan tanpa menyelesaikan persamaannya (\emph{Lyapunov function}).
  \item[Fungsi phi] fungsi $\varphi(n)$: banyaknya bilangan antara $1$ dan $n$ yang saling prima dengan $n$ (\emph{Euler totient function}).
  \item[Galat baku] simpangan baku penaksir, bukan simpangan baku data; ia mengukur goyangan satu tebakan (\emph{standard error}).
  \item[Galat pemotongan] selisih yang muncul karena proses tak hingga diganti potongan berhingga (\emph{truncation error}).
  \item[Garis regresi] garis (atau bidang) yang meminimumkan jumlah kuadrat sisa; ekuivalen dengan proyeksi ortogonal pada ruang kolom (\emph{regression line}).
  \item[Generalisasi] kemampuan mengerjakan soal yang belum pernah dilihat, bukan soal yang sudah dihafal (\emph{generalization}).
  \item[Gradien] vektor turunan parsial; arah naik tercepat sebuah fungsi (\emph{gradient}).
  \item[Gradien dikaburkan] gradien yang diberikan kepada penyerang tidak menunjuk arah yang benar, biasanya karena ada potongan tak terdiferensiasi atau acak (\emph{obfuscated gradient}).
  \item[Gradien terbalik] lihat turunan balik. (\emph{backward gradient}).
  \item[Harapan] nilai harapan $\mathbb{E}[X]$: rata-rata berbobot peluang, titik tumpu sebaran (\emph{expectation}).
  \item[Hasilkali dalam] pemetaan dua vektor ke satu skalar yang mengukur panjang dan sudut; di $\mathbb{R}^n$ nilainya $\langle \mathbf{x},\mathbf{y}\rangle=\mathbf{x}^{\top}\mathbf{y}$ (\emph{inner product}).
  \item[Hessian] matriks turunan parsial kedua; kecekungan fungsi ditentukan oleh tandanya (\emph{Hessian}).
  \item[Himpunan kuasa] himpunan semua himpunan bagian, termasuk himpunan kosong dan himpunan itu sendiri; kardinalitasnya $2^n$ (\emph{power set}).
  \item[Himpunan padat] himpunan yang setiap titiknya dapat didekati oleh anggota himpunan itu sendiri selain titik tersebut (\emph{dense set}).
  \item[Himpunan terhitung] himpunan yang anggotanya dapat didaftar sebagai barisan pertama, kedua, ketiga, dan seterusnya tanpa ada yang tertinggal (\emph{countable set}).
  \item[Hipotesis nol] dugaan awal yang diuji, biasanya "tidak ada efek"; penolakan hanya terjadi bila data sangat tidak cocok dengannya (\emph{null hypothesis}).
  \item[Hukum bilangan besar] janji bahwa rata-rata banyak pengamatan menempel ke rata-rata sebenarnya, dengan peluang meleset yang mengecil (\emph{law of large numbers}).
  \item[i.i.d.] sampel yang saling bebas dan berdistribusi sama; asumsi paling sering dipakai, dan paling sering dilanggar (\emph{independent and identically distributed}).
  \item[Implikasi] pernyataan berbentuk jika $P$ maka $Q$; ia hanya salah bila $P$ benar dan $Q$ salah (\emph{implication}).
  \item[Independen] dua kejadian yang tidak saling memberi kabar, $P(A\cap B)=P(A)P(B)$; mengetahui yang satu tidak mengubah peluang yang lain (\emph{independent}).
  \item[Induksi] pola bukti domino: basis benar, lalu setiap kasus menjatuhkan kasus berikutnya (\emph{induction}).
  \item[Induksi kuat] induksi yang boleh memakai seluruh kasus sebelumnya, bukan hanya kasus terdekat (\emph{strong induction}).
  \item[Integral lipat dua] penjumlahan nilai pada petak-petak kecil sebuah daerah, lalu limitnya ketika petak mengecil (\emph{double integral}).
  \item[Integral Riemann] limit jumlah luas persegi panjang ketika lebar partisi mengecil (\emph{Riemann integral}).
  \item[Interpolasi] mencari fungsi yang melewati semua titik data; berbeda dari regresi yang tidak menuntut itu. (\emph{interpolation}).
  \item[Interval kepercayaan] rentang nilai yang, bila prosedur pengambilan sampel diulang, memuat parameter sebenarnya dengan peluang tertentu (\emph{confidence interval}).
  \item[Interval Wilson] selang kepercayaan proporsi yang dicari dengan membalik uji, sehingga tetap masuk akal ketika proporsi mendekati nol atau satu (\emph{Wilson interval}).
  \item[Inti] lihat kernel: himpunan semua vektor yang dipetakan ke nol oleh transformasi linear (\emph{kernel, null space}).
  \item[Invers modular] bilangan $b$ sehingga $ab\equiv 1 \pmod n$; tombol batalkan pada aritmetika jam (\emph{modular inverse}).
  \item[Isomorfisme] pemetaan bijektif yang mempertahankan struktur; dua ruang yang isomorfik tidak dapat dibedakan secara aljabar (\emph{isomorphism}).
  \item[Jacobian] matriks turunan parsial pertama sebuah pemetaan vektor; determinannya mengukur faktor skala volume lokal (\emph{Jacobian}).
  \item[Jari-jari konvergensi] jarak terjauh dari pusat tempat sebuah deret pangkat masih konvergen (\emph{radius of convergence}).
  \item[Jumlah atas] taksiran integral yang sengaja dibesarkan, memakai nilai terbesar pada tiap potongan (\emph{upper sum}).
  \item[Jumlah langsung] penjumlahan dua subruang yang hanya bertemu di vektor nol, sehingga setiap vektor punya satu cara penulisan (\emph{direct sum}).
  \item[Kalibrasi] kesesuaian antara angka yang dijanjikan model dan frekuensi kenyataan pada kelompok yang dijanjikan sama (\emph{calibration}).
  \item[Kamus] kotak kecil di dalam bab yang menerangkan satu istilah baru beserta analoginya.
  \item[Kebocoran data] informasi dari data uji yang tanpa sengaja masuk ke proses pelatihan, sehingga hasil uji tampak lebih baik dari kenyataan (\emph{data leakage}).
  \item[Kekontinuan seragam] kekontinuan dengan satu $\delta$ yang berlaku untuk semua titik pada himpunan, bukan satu $\delta$ per titik (\emph{uniform continuity}).
  \item[Kekonvergenan seragam] barisan fungsi yang mendekati fungsi batasnya dengan laju yang sama di seluruh domain (\emph{uniform convergence}).
  \item[Kelas sisa] himpunan bilangan yang memberi sisa sama saat dibagi $n$, seperti pukul tiga pagi dan pukul tiga sore (\emph{residue class}).
  \item[Kernel] himpunan semua masukan yang dipetakan ke nol oleh transformasi linear; sama dengan ruang nol (\emph{kernel, null space}).
  \item[Kesalahan pembulatan] selisih antara bilangan real dan wakilannya di mesin, yang menumpuk pada penjumlahan panjang (\emph{rounding error}).
  \item[Kestabilan] sifat titik setimbang atau sistem dinamik yang tetap dekat dengan keadaannya setelah gangguan kecil (\emph{stability}).
  \item[Keterintegralan] sifat fungsi yang jumlah atas dan bawahnya dapat dirapatkan sampai selisihnya hilang, disebut keterintegralan Riemann (\emph{Riemann integrability}).
  \item[Ketunggalan] jaminan bahwa satu keadaan awal hanya melahirkan satu lintasan solusi (\emph{uniqueness}).
  \item[Koefisien determinasi] bagian variasi yang dijelaskan model, dinyatakan sebagai pecahan dari keseluruhan; ditulis $R^2$ (\emph{coefficient of determination}).
  \item[Kombinasi] pilihan beberapa anggota tanpa memperhatikan urutannya; banyaknya $\binom{n}{k}$ (\emph{combination}).
  \item[Kombinasi konveks] campuran berbobot vektor dengan bobot tak negatif yang berjumlah satu (\emph{convex combination}).
  \item[Kompak] himpunan tertutup dan terbatas; setiap barisan di dalamnya punya subbarisan yang konvergen di dalam himpunan itu (\emph{compact}).
  \item[Kompleksitas sampel] banyaknya contoh yang dibutuhkan agar model dapat digeneralisasi dengan keyakinan tertentu (\emph{sample complexity}).
  \item[Kondisi KKT] syarat perlu optimalitas untuk masalah dengan kendala pertidaksamaan (\emph{Karush-Kuhn-Tucker conditions}).
  \item[Kontinu] fungsi yang grafiknya dapat digambar tanpa mengangkat pena di titik itu, yaitu $\lim_{x\to c}f(x)=f(c)$ (\emph{continuous}).
  \item[Konvergensi dalam distribusi] fungsi distribusinya menuju fungsi distribusi batas di setiap titik kontinuitas; yang menyerupai adalah bentuk sebarannya, bukan nilainya (\emph{convergence in distribution}).
  \item[Konvergensi dalam peluang] $P(|X_n-X|>\varepsilon)\to 0$ untuk setiap $\varepsilon>0$ (\emph{convergence in probability}).
  \item[Konvergensi mutlak] deret yang tetap konvergen ketika semua sukunya dipaksa menjadi positif (\emph{absolute convergence}).
  \item[Konvergensi seragam] lihat kekonvergenan seragam: satu laju untuk seluruh domain, bukan satu jadwal per titik (\emph{uniform convergence}).
  \item[Koordinat] daftar koefisien yang menyatakan sebuah vektor dalam suatu basis; disebut juga vektor koordinat (\emph{coordinate vector}).
  \item[Korelasi] ukuran kekuatan hubungan linear antara dua peubah; tidak menyatakan sebab-akibat (\emph{correlation}).
  \item[Kovarians] lihat kovariansi: ukuran dua besaran bergerak bersama, dihitung dari hasil kali simpangan keduanya (\emph{covariance}).
  \item[Kovariansi] ukuran seberapa dua peubah bergerak bersama; satuannya hasil kali satuan keduanya (\emph{covariance}).
  \item[Kuadrat terkecil] pemilihan parameter yang meminimumkan jumlah kuadrat sisa ketika tidak ada solusi eksak (\emph{least squares}).
  \item[Kuantisasi] penggantian setiap nilai dengan kelipatan terdekat dari jarak tetap; hasilnya fungsi tangga yang turunannya nol hampir di mana-mana (\emph{quantization}).
  \item[Kuantor] penanda untuk semua ($\forall$) atau ada ($\exists$); urutan pemakaiannya mengubah arti pernyataan (\emph{quantifier}).
  \item[Kurva ROC] grafik tingkat benar terhadap tingkat salah palsu pada semua ambang; luas di bawahnya disebut AUROC (\emph{ROC curve}).
  \item[Kutub dan nol] nilai peubah kompleks yang membuat fungsi meledak atau lenyap; dipakai dalam analisis sistem dinamik dan kontrol. (\emph{poles and zeros}).
  \item[Lapangan] himpunan dengan dua operasi yang tertutup, punya identitas dan invers, dan memenuhi hukum distributif (\emph{field}).
  \item[Lapangan terurut] lapangan yang anggotanya dapat diurutkan dan aturan hitungnya menghormati urutan itu (\emph{ordered field}).
  \item[Lapisan tersembunyi] lapisan unit yang tidak bersentuhan langsung dengan masukan maupun keluaran jaringan (\emph{hidden layer}).
  \item[Lawan contoh] satu contoh yang mematahkan pernyataan "untuk semua" (\emph{counterexample}).
  \item[Limit fungsi] $f(x)$ dapat dipaksa sedekat apa pun ke $L$ asalkan $x$ cukup dekat ke $c$ tetapi tidak sama dengan $c$ (\emph{functional limit}).
  \item[Limit superior] nilai terbesar yang masih dikunjungi barisan sesering apa pun di ujungnya, ditulis $\limsup_{n\to\infty}a_n$ (\emph{upper limit}).
  \item[Lipschitz] syarat bahwa laju perubahan sebuah fungsi dibatasi satu konstanta; inilah yang menjamin solusi persamaan diferensial tunggal (\emph{Lipschitz condition}).
  \item[Luar sampel] data yang tidak dipakai saat melatih model, sehingga jujur untuk menilai (\emph{out-of-sample}).
  \item[Matriks elementer] matriks identitas yang sudah dikenai satu operasi baris, sehingga menjadi catatan operasi itu (\emph{elementary matrix}).
  \item[Matriks konfusi] tabel empat sel yang membandingkan prediksi dengan kenyataan (\emph{confusion matrix}).
  \item[Matriks ortogonal] matriks persegi yang kolom-kolomnya ortonormal; ia mempertahankan panjang dan sudut (\emph{orthogonal matrix}).
  \item[Metode Euler] penyelesaian numerik persamaan diferensial dengan melangkah mengikuti garis singgung; galatnya orde satu terhadap langkah. (\emph{Euler method}).
  \item[Metode tanda gradien] serangan satu langkah (FGSM) yang menggeser setiap koordinat sebesar anggaran, arahnya mengikuti tanda turunan kerugian (\emph{fast gradient sign method}).
  \item[Model ancaman] daftar lengkap kemampuan penyerang: besar gangguannya, apa yang ia ketahui, dan apa tujuannya (\emph{threat model}).
  \item[Multikolinearitas] keadaan ketika peubah penjelas saling hampir bergantung linear, sehingga koefisien regresi menjadi tidak stabil (\emph{multicollinearity}).
  \item[Multiplisitas] banyaknya kemunculan sebuah nilai eigen sebagai akar polinomial karakteristik, yang tidak selalu sama dengan dimensi ruang eigennya (\emph{multiplicity}).
  \item[NIDS] sistem yang memantau lalu lintas jaringan dan memutuskan apakah suatu aliran paket merupakan serangan (\emph{network intrusion detection system}).
  \item[Nilai eigen] skalar $\lambda$ sehingga $A\mathbf{v}=\lambda\mathbf{v}$ untuk suatu vektor tak nol (\emph{eigenvalue}).
  \item[Nilai mutlak] jarak sebuah bilangan ke nol pada garis bilangan, tanpa memedulikan tandanya (\emph{absolute value}).
  \item[Nilai singular] akar nilai eigen dari $A^{\top}A$; mengukur panjang sumbu utama pemetaan oleh matriks $A$ (\emph{singular value}).
  \item[Nilai-p] peluang memperoleh hasil setidaknya seekstrem yang teramati bila hipotesis nol benar; bukan peluang hipotesis nol benar (\emph{p-value}).
  \item[Norma] ukuran panjang vektor; norma $L_\infty$ mengambil komponen terbesar. (\emph{norm}).
  \item[Orientasi] kesepakatan tentang arah putar mana yang dihitung positif; membaliknya membalik tanda integral (\emph{orientation}).
  \item[Ortogonal] saling tegak lurus dalam arti hasilkali dalamnya nol (\emph{orthogonal}).
  \item[Ortonormal] himpunan ortogonal yang setiap anggotanya berpanjang satu (\emph{orthonormal}).
  \item[Overfitting] menghafal data latih termasuk deraunya, sehingga kemampuan di data baru menurun (\emph{overfitting}).
  \item[Partisi] pembagian himpunan menjadi bagian-bagian yang saling lepas dan menutup seluruh himpunan (\emph{partition}).
  \item[Peluang] bobot pada setiap kemungkinan, tak negatif dan berjumlah satu (\emph{probability}).
  \item[Peluang bersyarat] peluang kejadian bila kejadian lain diketahui terjadi; ditulis $P(A\mid B)$ (\emph{conditional probability}).
  \item[Pembatalan katastrofik] hilangnya digit berarti ketika dua bilangan hampir sama dikurangkan, sehingga yang tersisa hanya galat (\emph{catastrophic cancellation}).
  \item[Pembelajaran adversarial] studi tentang contoh yang sengaja dirancang untuk menipu model, dan cara mempertahankannya (\emph{adversarial learning}).
  \item[Pembulatan] penggantian bilangan real dengan bilangan mesin terdekat (\emph{rounding}).
  \item[Pemisahan data] pembagian data menjadi latih, validasi, dan uji yang tidak saling bocor (\emph{data split}).
  \item[Penaksir] aturan yang mengubah sampel menjadi tebakan parameter; tebakan yang dihasilkan disebut nilai dugaan (\emph{estimator}).
  \item[Pengali Lagrange] peubah bantu yang mengenakan harga pada kendala; nilainya menunjukkan seberapa besar tujuan dapat membaik bila kendala dilonggarkan (\emph{Lagrange multiplier}).
  \item[Penghubung logika] operasi yang menggabungkan pernyataan: dan, atau, tidak, jika-maka, jika-dan-hanya-jika (\emph{logical connective}).
  \item[Penurunan gradien stokastik] penurunan gradien yang memakai sebagian kecil data tiap langkah; lebih murah, lebih berisik (\emph{stochastic gradient descent}).
  \item[Perceptron] pengklasifikasi linear paling sederhana: tanda dari hasilkali dalam bobot dan masukan. (\emph{perceptron}).
  \item[Persamaan diferensial] kalimat yang mengikat laju perubahan sebuah besaran dengan besar besaran itu sendiri (\emph{differential equation}).
  \item[Persamaan diferensial biasa] persamaan yang mengaitkan fungsi dengan turunannya terhadap satu peubah bebas (\emph{ordinary differential equation}).
  \item[PGC (penurunan gradien)] aturan pembaruan parameter melawan arah gradien (\emph{gradient descent}).
  \item[Plateau tandus] daerah pada lanskap kerugian sirkuit variasional tempat gradien mengecil secara eksponensial terhadap jumlah qubit (\emph{barren plateau}).
  \item[Polinomial karakteristik] $\det(A-\lambda I)$; akar-akarnya adalah nilai eigen $A$ (\emph{characteristic polynomial}).
  \item[Privasi diferensial] janji bahwa keikutsertaan satu orang di dalam data hanya mengubah sebaran keluaran sedikit, diukur oleh $\varepsilon$ (\emph{differential privacy}).
  \item[Program linear] masalah memaksimumkan atau meminimumkan fungsi linear dengan kendala linear (\emph{linear program}).
  \item[Proposisi] kalimat yang bernilai benar atau salah, bukan pertanyaan atau perintah; dalam buku ini juga nama lingkungan teorema (\emph{proposition}).
  \item[Proyeksi] titik pada subruang yang paling dekat dengan sebuah vektor, yaitu bayangan tegak lurusnya (\emph{projection}).
  \item[Pseudoinvers] invers tergeneralisasi $A^{+}$ yang tetap bekerja ketika sebagian arah sudah gepeng: yang dapat dibalik dibalik, yang hilang dibiarkan (\emph{pseudoinverse}).
  \item[Rank] banyaknya arah bebas yang masih bertahan setelah dipetakan, yaitu $\dim\operatorname{im}T$ (\emph{rank}).
  \item[Rank-nullity] $\dim\ker T+\dim\operatorname{im}T=\dim V$ untuk transformasi linear pada ruang berdimensi hingga. (\emph{rank-nullity theorem}).
  \item[Rasio peluang] perbandingan peluang berhasil dan gagal; eksponensial koefisien regresi logistik menyatakan berapa kali rasio itu berlipat per satuan (\emph{odds ratio}).
  \item[Regresi] menarik fungsi yang paling dekat ke semua titik, bukan fungsi yang melewati semuanya (\emph{regression}).
  \item[Regresi logistik] model untuk keluaran biner yang memetakan kombinasi linear ke peluang lewat fungsi logistik (\emph{logistic regression}).
  \item[Regularisasi] penalti pada besarnya parameter agar model tidak terlalu peka pada data latih (\emph{regularization}).
  \item[Relaksasi LP] penggantian syarat bulat dengan syarat real agar masalah integer dapat diselesaikan sebagai LP (\emph{LP relaxation}).
  \item[Residu] selisih antara nilai teramati dan nilai yang diprediksi model (\emph{residual}).
  \item[Ringkasan] kotak butir di akhir bab yang boleh dibaca sendiri tanpa badan bab.
  \item[Ruang eigen] seluruh vektor eigen untuk satu nilai eigen beserta vektor nol; ia sebuah subruang (\emph{eigenspace}).
  \item[Ruang masalah] ruang objek nyata yang boleh diubah penyerang, misalnya paket dan waktu; berbeda dari ruang fitur tempat model bekerja (\emph{problem space}).
  \item[Ruang vektor] himpunan dengan penjumlahan dan perkalian skalar yang tertutup dan memenuhi delapan aksioma (\emph{vector space}).
  \item[Runge-Kutta] keluarga metode numerik yang mengevaluasi turunan beberapa kali per langkah untuk galat lebih kecil (\emph{Runge-Kutta}).
  \item[Selang prediksi] selang bagi satu pengamatan yang belum terjadi; selalu lebih lebar daripada selang bagi rata-rata (\emph{prediction interval}).
  \item[Sensitivitas] dalam privasi diferensial, perubahan keluaran terbesar yang dapat disebabkan satu orang (\emph{sensitivity}).
  \item[Sensitivitas global] pergeseran terbesar jawaban sebuah kueri yang dapat disebabkan satu orang pada kasus terburuk (\emph{global sensitivity}).
  \item[Serangan adaptif] serangan yang dirancang dengan mengetahui pertahanan yang sedang diuji, termasuk mengetahui bahwa gradiennya dimanipulasi (\emph{adaptive attack}).
  \item[Serangan adversarial] gangguan kecil pada masukan yang sengaja dibuat untuk mengubah keluaran model (\emph{adversarial example}).
  \item[Serupa] dua matriks yang menyatakan pemetaan linear yang sama dalam dua basis berbeda, $B=P^{-1}AP$ (\emph{similar}).
  \item[Sistem floating-point] notasi ilmiah biner dengan jumlah digit tetap; bilangan disimpan sebagai tanda, pangkat, dan mantissa, sehingga jarak antar bilangan yang dapat diwakili tidak seragam: rapat di dekat nol, renggang di ujung atas, dan dua bilangan berbeda dapat jatuh ke tempat yang sama (\emph{floating-point system}).
  \item[Softmax] pemetaan deretan skor menjadi peluang yang berjumlah satu, $e^{z_i}/\sum_j e^{z_j}$ (\emph{softmax}).
  \item[Solusi layak basis] titik yang diperoleh dengan menyetel sebagian peubah ke nol lalu menyelesaikan sisanya, dan masih memenuhi semua kendala (\emph{basic feasible solution}).
  \item[Solusi umum] keluarga semua solusi sebuah persamaan diferensial, bergantung pada konstanta sembarang (\emph{general solution}).
  \item[Sorry] penanda Lean~4 bahwa sebuah bukti ditangguhkan; mesin menerimanya tetapi mencatatnya sebagai aksioma (\emph{sorry}).
  \item[Stabil asimtotik] tetap dekat setelah diganggu, lalu benar-benar kembali ke titik setimbang (\emph{asymptotically stable}).
  \item[Subbarisan] barisan yang diperoleh dengan membuang sebagian suku tanpa mengubah urutan sisa sukunya (\emph{subsequence}).
  \item[Subruang] himpunan bagian ruang vektor yang tertutup terhadap penjumlahan dan perkalian skalar (\emph{subspace}).
  \item[Supremum] batas atas terkecil; bilangan terkecil yang tidak lebih kecil dari semua anggota himpunan (\emph{supremum, least upper bound}).
  \item[Susunan ulang] penukaran urutan suku deret; pada deret yang konvergen mutlak jumlahnya tidak berubah, pada deret yang tidak, jumlahnya dapat berubah (\emph{rearrangement}).
  \item[Syarat perlu] syarat yang harus dipenuhi agar pernyataan benar; tanpanya mustahil, tetapi memenuhinya belum cukup (\emph{necessary condition}).
  \item[Taktik] perintah langkah demi langkah dalam Lean~4 yang mengubah sasaran pembuktian (\emph{tactic}).
  \item[Tanda tangan digital] nilai yang hanya dapat dibuat pemegang kunci privat tetapi dapat diperiksa siapa pun dengan kunci publik (\emph{digital signature}).
  \item[Teorema Limit Pusat] rata-rata banyak variabel acak bebas menghampiri distribusi normal setelah dibakukan (\emph{central limit theorem}).
  \item[Terhitung] himpunan yang anggotanya dapat dinomori tanpa ada yang terlewat; $\mathbb{Q}$ terhitung, $\mathbb{R}$ tidak (\emph{countable}).
  \item[Term bukti] ungkapan yang sekaligus menjadi bukti dan menjadi nilai; memeriksa bukti sama dengan memeriksa tipe (\emph{proof term}).
  \item[Titik akumulasi] titik yang setiap persekitarannya memuat anggota himpunan lain (\emph{accumulation point}).
  \item[Titik setimbang] keadaan sistem dinamik yang tidak berubah bila tidak diganggu (\emph{equilibrium point}).
  \item[Transferabilitas] sifat contoh adversarial yang dibuat untuk satu model tetapi juga menipu model lain yang dilatih terpisah (\emph{transferability}).
  \item[Transformasi linear] pemetaan yang mempertahankan penjumlahan dan perkalian skalar (\emph{linear map}).
  \item[Turunan] limit hasil bagi selisih; kemiringan garis singgung sekaligus laju perubahan sesaat (\emph{derivative}).
  \item[Turunan balik] algoritme menghitung gradien seluruh parameter dengan menerapkan aturan rantai dari keluaran ke masukan (\emph{backpropagation}).
  \item[Typeclass] daftar kemampuan yang harus dimiliki sebuah tipe agar dapat dipakai di tempat tertentu (\emph{type class}).
  \item[Uji hipotesis] prosedur memutuskan, dengan risiko galat yang dinyatakan, apakah data cukup menolak dugaan awal (\emph{hypothesis test}).
  \item[Uji rasio] deret konvergen mutlak bila $\lim|a_{n+1}/a_n|<1$, dan divergen bila limit itu lebih besar dari satu (\emph{ratio test}).
  \item[ULP] satuan di tempat terakhir: satu langkah terkecil yang dapat dibedakan mesin pada skala sebuah bilangan (\emph{unit in the last place}).
  \item[Uniter] matriks yang mempertahankan hasilkali dalam, $U^{*}U=I$; versi kompleks dari matriks ortogonal (\emph{unitary}).
  \item[Validasi silang] penilaian model dengan membagi data bergiliran menjadi latih dan uji (\emph{cross-validation}).
  \item[Validasi silang bersarang] dua lapis lipatan: lapisan luar menilai, lapisan dalam memilih, sehingga pemilihan model tidak membocorkan data uji (\emph{nested cross-validation}).
  \item[Variabel acak] fungsi yang memetakan hasil percobaan ke bilangan (\emph{random variable}).
  \item[Variabel kendur] peubah tambahan yang mencatat sisa pemakaian, sehingga kendala tidak lebih dari berubah menjadi persamaan (\emph{slack variable}).
  \item[Varians inflasi] seberapa banyak variansi sebuah koefisien digelembungkan oleh penjelas lain yang membawa informasi serupa; disingkat VIF (\emph{variance inflation factor}).
  \item[Variansi] ukuran sebaran di sekitar rata-rata; akar kuadratnya adalah simpangan baku (\emph{variance}).
  \item[Verifikasi formal] pembuktian bahwa program atau rumus memenuhi spesifikasinya, diperiksa mesin (\emph{formal verification}).
  \item[Verteks] sudut daerah layak, yaitu titik yang tidak pernah menjadi titik tengah antara dua titik layak lain (\emph{vertex}).
  \item[Waktu polinomial] kelas algoritme yang banyaknya langkah tumbuh tidak lebih cepat dari polinom ukuran masukan (\emph{polynomial time}).
  \item[Wilson] interval kepercayaan proporsi yang tetap masuk akal ketika proporsi mendekati nol atau satu. (\emph{Wilson interval}).
  \item[Zaman kuantum] era ketika komputer kuantum cukup besar untuk menjalankan algoritme Shor pada kunci yang dipakai hari ini. (\emph{quantum computing era}).
\end{description}

% ---------------------------- lampiran C ----------------------------
\chapter{Petunjuk dan Kriteria Penilaian Latihan}
\label{app:petunjuk}

Kunci jawaban ada di Lampiran~D, tetapi bentuknya sengaja tidak lengkap: angka akhir
dan sketsa langkah, bukan pembuktian penuh. Alasannya konsistensi. Angka akhir tanpa
jejak pemeriksaan akan mengajari pembaca kebiasaan yang justru ingin dihapus buku ini,
dan bukti yang disodorkan utuh menghapus satu-satunya latihan yang berguna, yaitu
mencari sendiri di mana argumen belum menutup. Jadi pakai kunci ini sebagai pembanding
terakhir, bukan sebagai tempat memulai. Yang disediakan di sini adalah cara memeriksa
pekerjaan sendiri sebelum kamu membuka Lampiran~D.

\section*{Tiga jenis latihan dan cara memeriksanya}

\begin{center}
\small
\begin{tabular}{@{}p{2.6cm}p{5.6cm}p{5.8cm}@{}}
\toprule
\textbf{Jenis} & \textbf{Yang dinilai} & \textbf{Cara memeriksa} \\
\midrule
Bukti & setiap langkah harus berangkat dari definisi atau hasil yang sudah
  dibuktikan; tidak boleh ada lompatan & tukar naskah dengan teman; minta dia
  menunjuk kalimat yang tidak bisa ia turunkan sendiri \\
Hitung & angka akhir dan langkah antara & ulangi dengan Python; hitung galat
  relatif terhadap hasil eksak; periksa apakah jawaban masuk akal secara
  orde besaran \\
Nilai klaim & mutu argumen, bukan benar-salah binari & tuliskan lebih dulu apa
  yang harus benar agar klaim itu berlaku; cari satu contoh yang mematahkannya;
  baru putuskan "benar", "salah", atau "belum cukup bukti" \\
\bottomrule
\end{tabular}
\end{center}

\section*{Alat verifikasi per bagian}

\begin{center}
\small
\begin{tabular}{@{}llp{7.2cm}@{}}
\toprule
\textbf{Bab} & \textbf{Alat} & \textbf{Pemeriksaan khas} \\
\midrule
1--5 & \texttt{python3} & tabel nilai, faktorisasi, jumlah parsial; bandingkan
  dengan taksiran analitik (mis. $H_n\approx\ln n+\gamma$) \\
6--10 & \texttt{numpy}, \texttt{scipy.linalg} & rank, ruang nol (\texttt{null\_space}),
  nilai eigen, SVD, kuadrat terkecil (\texttt{lstsq}); periksa
  $\|A\mathbf{v}-\lambda\mathbf{v}\|$ \\
11--16 & \texttt{scipy.integrate}, beda hingga & kuadratur vs antiturunan; laju
  galat pada beberapa $h$ atau beberapa $n$ \\
17--21 & \texttt{numpy.random}, \texttt{scipy.stats} & simulasi Monte Carlo dengan
  \emph{seed} tetap; bandingkan peluang empiris dengan peluang eksak dan
  laporkan galat bakunya \\
22 & \texttt{struct}, \texttt{float.hex} & bitfield, ULP, epsilon mesin;
  bandingkan algoritme naif dengan Kahan \\
23--24 & \texttt{scipy.optimize} & \texttt{linprog}, \texttt{minimize},
  \texttt{milp}; periksa syarat KKT dan dualitas kuat secara numerik \\
25--26 & \texttt{numpy} & pemeriksaan gradien dengan beda hingga; kurva galat
  latih vs validasi pada beberapa setelan regularisasi \\
27--30 & \texttt{numpy}, \texttt{scikit-learn} bila ada & matriks konfusi dari
  prediksi tersimpan; interval Wilson untuk proporsi; akurasi di bawah serangan
  untuk beberapa $\varepsilon$ dan beberapa \emph{seed} \\
\bottomrule
\end{tabular}
\end{center}

\section*{Cara menilai pekerjaan sendiri}

\begin{enumerate}
  \item Tuliskan lebih dulu apa yang diketahui dan apa yang ingin dibuktikan.
        Bila keduanya belum jelas, kesulitannya bukan pada matematika.
  \item Untuk latihan hitung, tebak jawaban sebelum menghitung. Tebakan yang salah
        tetap berguna: ia menunjukkan di mana intuisi meleset.
  \item Bandingkan hasil dengan dua cara berbeda. Satu cara bisa benar karena
        kebetulan; dua cara yang sepakat jarang kebetulan.
  \item Bila hasilnya berbeda jauh, cari kesalahan pada kasus sederhana lebih dulu
        (satu variabel, satu langkah, $n=1$).
  \item Tuliskan batas keberlakuan jawaban. Jawaban benar untuk kasus khusus yang
        dinyatakan sebagai jawaban umum adalah kesalahan yang paling sering lolos.
\end{enumerate}

% ---------------------------- lampiran D ----------------------------
\chapter{Kunci Jawaban}
\label{app:kunci}

Setiap angka di lampiran ini dihitung ulang dengan Python sebelum ditulis, dan skrip
pemeriksanya disertakan bersama sumber naskah, sehingga pembaca dapat menjalankannya
sendiri. Untuk latihan bukti, yang diberikan hanya sketsa dua sampai lima baris;
itu cukup untuk memeriksa apakah arah pekerjaanmu benar, tetapi tidak cukup untuk
disalin sebagai bukti. Butir yang soalnya sendiri bermasalah ditandai terang-terangan
sebagai ambigu, bukan dipaksakan punya jawaban.

% Kunci jawaban Bab 1--5. Berkas ini di-\input langsung di dalam document.
% Setiap angka berasal dari tools/cek-kunci-01-05.py di repositori ini.

\section*{Bab 1: Bahasa Matematika dan Seni Membuktikan}
\begin{enumerate}[label=\textbf{1.\arabic*}]
  \item $\exists \varepsilon>0\ \forall N\in\mathbb{N}\ \exists n\ge N:\ |a_{n}-L|\ge
        \varepsilon$. Dalam bahasa Indonesia: ada satu toleransi yang tidak pernah
        tercapai, sebab berapa pun $N$ dipilih selalu ada suku sesudahnya yang
        menyimpang paling sedikit sebesar $\varepsilon$ itu.

  \item Negasi ``setiap prima $>2$ ganjil'' ialah ``ada prima $>2$ yang genap''.
        Negasi itu \emph{salah}, jadi pernyataan aslinya benar dan tidak ada objek
        penyangkal. Untuk klaim kedua, negasinya ``untuk setiap $n$,
        $n^{2}+n+41$ prima''; penyangkalnya $n=40$ dengan $1681=41^{2}$.

  \item $n^{2}+n=n(n+1)$. Kasus: $n=2k$ memberi $2k(2k+1)$, $n=2k+1$ memberi
        $(2k+1)(2k+2)=2(2k+1)(k+1)$. Kontraposisi: bila $n(n+1)$ ganjil maka $n$
        dan $n+1$ keduanya ganjil, padahal dua bilangan berurutan tidak mungkin
        keduanya ganjil.

  \item Periksa: $n=1$ memberi $1$, $n=2$ memberi $5$, $n=3$ memberi $14$, cocok
        dengan $n(n+1)(2n+1)/6$. Langkah induksi: tambahkan $(k+1)^{2}$ pada kedua
        ruas, lalu
        $\frac{k(k+1)(2k+1)}{6}+(k+1)^{2}=\frac{(k+1)(k+2)(2k+3)}{6}$.

  \item Induksi kuat pada $n$. Untuk $n=1$: $k=0$, $m=1$. Untuk $n$ genap tulis
        $n=2j$ dengan $j<n$; hipotesis induksi memberi $j=2^{k}m$, sehingga
        $n=2^{k+1}m$ dengan $m$ tetap ganjil. Hipotesis kuat dipakai tepat di sini:
        $j<n$, bukan $j=n-1$.

  \item (a) Benar, prinsip induksi biasa. (b) Salah: ambil $P(n)$ berarti
        $n\le 10^{6}$; penyangkal $n=10^{6}+1$. (c) Salah: $n=40$ memberi
        $1681=41^{2}$. (d) Salah: $n=3$ habis dibagi $3$ tetapi tidak oleh $6$.

  \item Skrip: $87$ prima dari $101$ nilai untuk $n=0$ sampai $100$. Gagal pada
        $40$, $41$, $44$, $49$, $56$, $65$, $76$, $81$, $82$, $84$, $87$, $89$,
        $91$, $96$. Hanya $5$ dari $14$ berbentuk $p-1$ atau $p$: $40,41$ lewat
        $p=41$, $82$ lewat $83$, $89$, $96$ lewat $97$.

  \item $2^{1}=2>1$; $2^{2}=4=4$; $2^{3}=8<9$; $2^{4}=16=16$; $2^{5}=32>25$. Jadi
        $n_{0}=5$. Induksi: bila $2^{k}>k^{2}$ dan $k\ge 5$, maka
        $2^{k+1}>2k^{2}\ge(k+1)^{2}$ sebab $k^{2}-2k-1>0$.

  \item 
Klaimnya salah, sebab $n=40$ memberi $40^{2}+40+41=1\,681=41^{2}$, bukan
        prima. Empat puluh titik pertama hanya menyokong, tidak membuktikan: satu
        penyangkal sudah cukup membatalkan klaim ``untuk setiap $n\ge0$''.

  \item Salah: $n=2$ memberi $4\mid 4$ tetapi $4\nmid 2$. Klaim yang benar:
        $4\mid n^{2}$ mengakibatkan $2\mid n$. Bukti kontraposisi: bila
        $n=2k+1$ maka $n^{2}=4k^{2}+4k+1\equiv 1 \pmod 4$.
\end{enumerate}

\section*{Bab 2: Himpunan, Relasi, Fungsi, dan Kardinalitas}
\begin{enumerate}[label=\textbf{2.\arabic*}]
  \item $x\in A\cap(B\cup C)\iff x\in A\wedge(x\in B\vee x\in C)\iff
        (x\in A\wedge x\in B)\vee(x\in A\wedge x\in C)$, memakai distributivitas
        $\wedge$ atas $\vee$. Kedua: $x\in A\setminus(B\setminus C)\iff x\in
        A\wedge(\neg x\in B\vee x\in C)\iff x\in(A\setminus B)\cup(A\cap C)$.

  \item Salah. Penyangkal terkecil: $A=\{1\}$ dengan $R=\varnothing$; relasi kosong
        simetris dan transitif secara vakum, tetapi $(1,1)\notin R$. Yang hilang
        ialah refleksivitas untuk elemen yang tidak muncul di pasangan mana pun.

  \item Ambil $S_{1},S_{2}\subseteq A$ dengan $f(S_{1})=f(S_{2})$. Untuk $s\in S_{1}$
        ada $t\in S_{2}$ dengan $f(s)=f(t)$, dan injektivitas memberi $s=t\in S_{2}$.
        Tanpa injektif: $f$ konstan pada $\{1,2\}$ membuat $f(\{1\})=f(\{2\})$.

  \item Arah $(\Leftarrow)$ sepele. Untuk $(\Rightarrow)$: bila $a=b$, kesamaan
        $\{\{a\}\}=\{\{c\},\{c,d\}\}$ memaksa $\{c\}=\{c,d\}$, jadi $c=d=a$. Bila
        $a\ne b$, maka $\{a\}$ harus sama dengan $\{c\}$, sebab $\{c,d\}$ berukuran
        dua; jadi $a=c$ lalu $b=d$.

  \item Definisikan $g:f(A)\to A$ dengan memilih satu pramasuk tiap $y\in f(A)$;
        $g$ injektif, jadi $|f(A)|\le|A|$ menurut lemma sebelumnya. Bila $f$ tidak
        injektif, dua anggota $A$ memberi satu nilai, sehingga $|f(A)|<|A|$.

  \item $\binom{5}{2}=10$ partisi bertipe $2+3$ dan $5$ bertipe $1+4$, seluruhnya
        $15=2^{4}-1=S(5,2)$. Skrip menghitung partisi secara langsung dan juga
        memberi $15$; banyaknya seluruh relasi ekuivalensi pada $5$ anggota $52$.

  \item $\pi(7,5)=T_{12}+5=78+5=83$. Untuk $1000$: $T_{44}=990\le 1000<T_{45}=1035$,
        jadi $s=44$, $n=10$, $m=34$; pemeriksaan $\pi(34,10)=990+10=1000$.

  \item Kumulatif sampai $k=9$ adalah $27$; tingkat $k=10$ memuat $1/9,\,3/7,\,7/3,\,
        9/1$ pada posisi $28$--$31$, jadi $7/3$ di posisi ke-$30$. Karena
        $\varphi(11)=10$, posisi ke-$41$ ialah suku terakhir tingkat $11$, yaitu
        $10/1$.

  \item Susun $A_{n}=\{a_{n1},a_{n2},\dots\}$ lalu jalankan diagonal Cantor pada
        pasangan $(n,m)$: setiap anggota $\bigcup A_{n}$ muncul sedikitnya sekali,
        dan pembuangan kemunculan ulang tidak menambah cacah.

  \item $g(x)=\tan\bigl(\pi(x-\tfrac12)\bigr)$ bijektif dari $(0,1)$ ke
        $\mathbb{R}$. Daerah asalnya harus dibatasi karena $\tan$ bernilai tak hingga
        di $\pm\tfrac12\pi$; pada $[0,1]$ rumus itu tidak terdefinisi di kedua ujung.
\end{enumerate}

\section*{Bab 3: Aksioma Bilangan Real dan Supremum}
\begin{enumerate}[label=\textbf{3.\arabic*}]
  \item $a<b$ berarti $b-a>0$. Tambahkan $-a-b$ pada kedua ruas (tertutup terhadap
        penjumlahan): $(-a)-(-b)=b-a>0$, yaitu $-b<-a$. Arah sebaliknya diperoleh
        dengan menambahkan $a+b$. Maka mengali dengan $-1$ membalik arah.

  \item Karena $a<b$, maka $2a<a+b$ dan $a+b<2b$; kalikan dengan $\tfrac12>0$,
        didapat $a<\frac{a+b}{2}<b$. Hanya aksioma urutan yang dipakai. Sebaliknya
        $\sup\{q:q^{2}<2\}$ tidak ada di $\mathbb{Q}$, jadi kelengkapan diperlukan.

  \item $\sup S=2$: $2$ batas atas, dan $2-\varepsilon$ bukan batas atas karena
        $2-1/n>2-\varepsilon$ untuk $n>1/\varepsilon$. $\inf S=1$ dicapai pada
        $n=1$, jadi minimum ada; maksimum tidak ada sebab $2\notin S$.

  \item Benar. Premisnya adalah sifat Archimedes: untuk setiap $x>0$ ada $n$ dengan
        $1/n<x$. Akibatnya bila $m>0$ adalah yang terkecil, maka $m/2>0$ lebih kecil
        dari $m$, kontradiksi.

  \item $q=1$: $h=1/4$, jadi $q+h=5/4$ dengan $(5/4)^{2}=1{,}5625<2$. $q=7/5$:
        $h=1/120$, jadi $q+h=169/120\approx 1{,}4083$ dengan kuadrat
        $1{,}9834<2$. Kenaikan kedua $30$ kali lebih kecil: makin dekat $q$ ke
        $\sqrt2$, koreksinya makin pendek.

  \item $|x^{3}-y^{3}|=|x-y|\,|x^{2}+xy+y^{2}|\le|x-y|\bigl(|x|^{2}+|x||y|+|y|^{2}\bigr)
        \le 12\,|x-y|$, sebab tiap suku paling besar $4$. Ketaksamaan segitiga
        dipakai pada langkah kedua untuk memecah nilai mutlak jumlah.

  \item Bila $S$ terbuka dan $x$ titik limit $S^{c}$ dengan $x\in S$, maka ada
        $B(x,r)\subseteq S$, sehingga $B(x,r)\cap S^{c}=\varnothing$: kontradiksi.
        Sebaliknya bila $S$ tertutup dan $x\in S^{c}$, ketiadaan bola di dalam
        $S^{c}$ membuat $x$ titik limit $S\subseteq S$, kontradiksi.

  \item $x_{1}=7/4$; $x_{2}=97/56\approx 1{,}732143$;
        $x_{3}=18817/10864\approx 1{,}732050810$; $x_{4}$ dan $x_{5}$ sama sampai
        $16$ angka, $1{,}732050807568877$. Galat relatif
        lima langkah: $1{,}0363\times10^{-2}$, $5{,}3145\times10^{-5}$,
        $1{,}4121\times10^{-9}$, $9{,}97\times10^{-19}$, $4{,}97\times10^{-37}$.
        Nisbah $e_{k+1}/e_{k}^{2}$: $0{,}4330$, $0{,}4949$, $0{,}5000$, menuju $1/2$.

  \item $0$ termuat di setiap $(-1/n,1/n)$. Untuk $x\ne0$ pilih $n>1/|x|$, maka
        $x\notin(-1/n,1/n)$. Irisan berhingga himpunan buka tetap buka, sedangkan
        irisan tak hingga ini sama dengan $\{0\}$ yang tidak buka.

  \item Menurut kerapatan $\mathbb{Q}$, pada selang $(a/\sqrt2,\,b/\sqrt2)$ ada
        $q\in\mathbb{Q}$ dengan $q\ne0$. Maka $a<q\sqrt2<b$, dan $q\sqrt2$ irasional
        sebab bila rasional maka $\sqrt2$ rasional.
\end{enumerate}

\section*{Bab 4: Barisan dan Limit}
\begin{enumerate}[label=\textbf{4.\arabic*}]
  \item $|x_{n}-\tfrac23|=\Bigl|\frac{3(2n-1)-2(3n+4)}{3(3n+4)}\Bigr|
        =\frac{11}{9n+12}$. Syaratnya $n>\frac{11/\varepsilon-12}{9}$, jadi
        $N(\varepsilon)=\bigl\lfloor(11/\varepsilon-12)/9\bigr\rfloor$. Untuk
        $\varepsilon=10^{-3}$: $N=1220$, sedangkan suku terkecil yang memenuhi
        $n=1221$ dengan galat $9{,}9991\times10^{-4}$.

  \item $\bigl||x_{n}|-|L|\bigr|\le|x_{n}-L|$ (ketaksamaan segitiga terbalik),
        sehingga batas $\varepsilon$ yang sama berlaku. Kebalikannya salah:
        $x_{n}=(-1)^{n}$ memberi $|x_{n}|=1$ konvergen, tetapi $x_{n}$ divergen.

  \item Untuk $x_{n}$: suku genap $2+1/n\to2$, suku ganjil $1/n\to0$; jadi
        $\limsup=2$, $\liminf=0$ (maksimum $2{,}5$ di $n=2$ bukan limsup). Untuk
        $y_{n}$: genap $1+1/n\to1$, ganjil $-(1+1/n)\to-1$; jadi $\limsup=1$,
        $\liminf=-1$. Numerik $n=1000$: $1{,}001$ dan $-1{,}001$.

  \item Sebagian benar, kesimpulannya keliru. Benar bahwa $(1+1/n)^{n}$ naik dan
        terbatas oleh $3$, sehingga limitnya ada dan paling besar $3$. Tetapi batas
        atas bukan nilai limit: limitnya $e\approx 2{,}7183$.

  \item $3=\bigl(3^{n}\bigr)^{1/n}\le(2^{n}+3^{n})^{1/n}\le 3\cdot 2^{1/n}$, dan
        $2^{1/n}\to1$ karena $(1+\delta)^{n}\ge 1+n\delta$ (Bernoulli) memberi
        $\delta\le 2^{1/n}-1\le 1/n$. Kedua pagar menuju $3$.

  \item Galat $<10^{-6}$ pertama di $n=12$ ($2{,}6146\times10^{-7}$); $<10^{-12}$
        di $n=21$. Nisbah galat berurutan menetap di $0{,}250000$, sesuai
        $|x_{n+1}-2|=|x_{n}-2|/(\sqrt{2+x_{n}}+2)\to|x_{n}-2|/4$.

  \item Cauchy terbatas: ambil $\varepsilon=1$, dapat $N$; semua suku sesudah $N$
        berada dalam $(x_{N}-1,x_{N}+1)$, sisanya berhingga. Monoton terbatas:
        dengan teorema kekonvergenan monoton ia konvergen, dan barisan konvergen
        adalah Cauchy.

  \item $\varepsilon=10^{-3}$ memberi ambang $5/\varepsilon-2=4998$, jadi
        $N=4998$ sah. Galat: $n=4998$ tepat $10^{-3}$ (belum memenuhi), $n=4999$
        memberi $9{,}9980\times10^{-4}$, $n=5000$ memberi $9{,}9960\times10^{-4}$.
        $N-1=4997$ tidak sah sebab $n=4998>4997$ masih bergalat tepat
        $\varepsilon$, bukan kurang.

  \item Andaikan $x_{n}\not\to L$; maka ada $\varepsilon>0$ dan subbarisan
        $(x_{n_{k}})$ dengan $|x_{n_{k}}-L|\ge\varepsilon$. Menurut
        Bolzano--Weierstrass subbarisan itu punya subbarisan konvergen ke $L'$ dengan
        $|L'-L|\ge\varepsilon$, bertentangan dengan hipotesis.
\end{enumerate}

\section*{Bab 5: Deret Tak Hingga}
\begin{enumerate}[label=\textbf{5.\arabic*}]
  \item $s_{20}=2{,}4999451576$, sedangkan jumlahnya $5/2=2{,}5$; galat
        $5{,}4842\times10^{-5}$. Batas $\frac{r^{21}}{1-r}$ memberi angka yang sama
        persis, $5{,}4842\times10^{-5}$, sebab ekornya geometri dengan suku positif.

  \item $\frac{1}{n(n+2)}=\frac12\bigl(\frac1n-\frac1{n+2}\bigr)$, sehingga
        $s_{N}=\frac12\bigl(1+\frac12-\frac{1}{N+1}-\frac{1}{N+2}\bigr)$; limitnya
        $\frac12\cdot\frac32=\frac34$.

  \item Karena $a_{n}\to0$, untuk $n\ge N$ berlaku $0\le a_{n}\le1$, jadi
        $a_{n}^{2}\le a_{n}$; uji pembanding menutup. Kebalikannya salah:
        $a_{n}=1/n$ memberi $\sum a_{n}^{2}$ konvergen tetapi $\sum a_{n}$ divergen.

  \item Salah. Patah pada ``lebih kecil daripada suku geometri rasio $\tfrac12$ untuk
        $n$ cukup besar'': $1/n<(1/2)^{n}$ justru gagal, sebab $(1/2)^{n}$ menyusut
        jauh lebih cepat. Uji pembanding butuh ketaksamaan pada seluruh ekor. Deret
        harmonik divergen.

  \item (a) Konvergen, uji rasio: nisbah $\to\tfrac12<1$. (b) Divergen, uji rasio:
        nisbah $2\bigl(\frac{n}{n+1}\bigr)^{2}\to2>1$. (c) Konvergen, uji rasio:
        $\bigl(\frac{n}{n+1}\bigr)^{n}\to\frac1e\approx0{,}3679<1$. (d) Divergen,
        deret-$p$ dengan $p=\tfrac13\le1$.

  \item $s_{10^{6}}=1{,}6449330668$ dan $\pi^{2}/6=1{,}6449340668$; selisihnya
        $9{,}9999950\times10^{-7}$, dekat $\frac1N$. Setelah koreksi
        $\frac1N-\frac{1}{2N^{2}}$ hasilnya $1{,}6449340668$ dengan galat
        $2{,}2\times10^{-16}$. Angka desimal benar: $5$ sebelum, $15$ sesudah.

  \item Perlu $\frac{1}{\sqrt{N+1}}<10^{-3}$, yaitu $N+1>10^{6}$, jadi $N=10^{6}$.
        Pada $N$ itu jumlahnya $0{,}6043986435$ sedangkan limitnya
        $0{,}6048986434$; galat sebenarnya $5{,}0000\times10^{-4}$, tepat setengah
        batas. Batasnya longgar sekitar faktor dua.

  \item $|b_{n}|\le M$ memberi $|a_{n}b_{n}|\le M|a_{n}|$. Karena $\sum M|a_{n}|$
        konvergen, uji pembanding menyimpulkan $\sum|a_{n}b_{n}|$ konvergen.

  \item Setelah $K$ blok, suku positif terpakai sampai penyebut $6K-1$ dan suku
        negatif sampai $\frac{1}{2K}$, jadi
        $S_{K}=H_{6K}-\frac12H_{3K}-\frac12H_{K}\to\ln6-\frac12\ln3=\ln(2\sqrt3)
        \approx1{,}2425$. Numerik $K=10^{6}$ memberi $1{,}2424531$. Nisbahnya ke
        $\ln2$ ialah $1{,}7925$, jadi bukan kelipatan rasional $\ln2$.

  \item Tulis $p_{n}=\max(a_{n},0)$ dan $q_{n}=\max(-a_{n},0)$; maka
        $a_{n}=p_{n}-q_{n}$ dan $|a_{n}|=p_{n}+q_{n}$. Bila kedua deret itu
        konvergen, $\sum|a_{n}|$ konvergen, bertentangan dengan kekonvergenan
        bersyarat. Karena itu keduanya divergen, dan kedua tumpukan tak terbatas
        itulah bahan bakar teorema susunan ulang Riemann.
\end{enumerate}

% Kunci jawaban Bab 6--15. Berkas ini di-\input langsung di dalam document.
% Setiap angka berasal dari skrip cek di direktori kerja sementara.

\section*{Bab 6: Ruang Vektor dan Subruang}
\begin{enumerate}[label=\textbf{6.\arabic*}]
  \item $W$ himpunan solusi sistem homogen $x-2y+z=0$, $x+y=0$, jadi tertutup
        terhadap jumlah dan perkalian skalar. Dari $x=-y$ dan $z=3y$ didapat
        $W=\operatorname{span}\{(-1,1,3)\}$: basis satu vektor, $\dim W=1$.

  \item Himpunan $D=\{A:a_{11}=a_{22}=0\}$ tertutup terhadap jumlah dan skalar,
        sebab kedua entri diagonal tetap nol. Setiap anggotanya
        $\begin{pmatrix}0&b\\c&0\end{pmatrix}=bE_{12}+cE_{21}$, jadi
        $\{E_{12},E_{21}\}$ basis dan $\dim D=2$.

  \item $(1+x)+(1+x^{2})=2+x+x^{2}$: himpunan bergantung linear, sehingga
        $\dim\operatorname{span}S=2$. Vektor $x$ tidak di rentang itu, sebab
        $a(1+x)+b(1+x^{2})$ menuntut $a=1$ dan $a+b=0$ sekaligus. Jadi
        $\{1+x,\ 1+x^{2},\ x\}$ basis $P_{2}$.

  \item Kumpulkan koefisien: $(c_{1}+c_{3})\mathbf{v}_{1}+(c_{1}+c_{2})
        \mathbf{v}_{2}+(c_{2}+c_{3})\mathbf{v}_{3}=\mathbf{0}$. Kebebasan
        $\{\mathbf{v}_{1},\mathbf{v}_{2},\mathbf{v}_{3}\}$ memberi tiga persamaan
        homogen dengan determinan $2\ne0$, jadi $c_{1}=c_{2}=c_{3}=0$.

  \item $U$ bidang $xy$, $W$ bidang $x=y$. Irisannya
        $U\cap W=\operatorname{span}\{(1,1,0)\}$, jadi $\dim(U\cap W)=1$ dan
        $\dim(U+W)=2+2-1=3$. Rumus cocok: $U+W=\mathbb{R}^{3}$, dan numpy
        memberi rank gabungan keempat vektor pembangun sama dengan $3$.

  \item Eliminasi memberi pivot di kolom $1,2,4$, jadi rank $3$; basis ruang
        kolomnya $(1,2,0)$, $(2,4,1)$, dan $(4,9,1)$. Nullity $4-3=1$ dengan
        basis $\{(-1,-1,1,0)\}$; \texttt{matrix\_rank} dan
        \texttt{null\_space} mengonfirmasi kedua angka.

  \item Salah. Di $\mathbb{R}^{2}$ ambil $U_{1}=\operatorname{span}\{(1,0)\}$,
        $U_{2}=\operatorname{span}\{(0,1)\}$, $U_{3}=\operatorname{span}\{(1,1)\}$.
        Semua irisan berpasangan $\{\mathbf{0}\}$, tetapi
        $(1,1)=(1,0)+(0,1)$: penyajian tidak tunggal, dan
        $\dim U_{1}+\dim U_{2}+\dim U_{3}=3>2=\dim\mathbb{R}^{2}$.

  \item Ambil basis $\mathbf{u}_{1},\dots,\mathbf{u}_{n-1}$ dari $U$. Himpunan itu
        bebas linear dan berukuran kurang dari $n$, jadi belum merentang $V$;
        ada $\mathbf{v}\in V$ di luar rentangnya. Teorema perluasan membuat
        $\mathbf{u}_{1},\dots,\mathbf{u}_{n-1},\mathbf{v}$ basis $V$, sehingga
        $V=U\oplus\operatorname{span}\{\mathbf{v}\}$.

  \item Polinomial berderajat paling tinggi $4$ yang lenyap di $0$ dan $1$ habis
        dibagi $x(x-1)$, jadi $p=x(x-1)(a+bx+cx^{2})$: tiga parameter bebas,
        dimensi $3$. Basis $\{x(x-1),\,x^{2}(x-1),\,x^{3}(x-1)\}$. Matriks
        evaluasi $2\times5$ berperingkat $2$, jadi nullity $5-2=3$.

  \item Koordinat memenuhi $c_{1}+c_{2}+c_{3}+c_{4}=4$, $c_{2}+c_{3}+c_{4}=2$,
        $c_{3}+c_{4}=-5$, $c_{4}=1$. Dari bawah: $c_{4}=1$, $c_{3}=-6$,
        $c_{2}=7$, $c_{1}=2$. Jadi $[q]_{B}=(2,7,-6,1)$, sama dengan hasil
        \texttt{numpy.linalg.solve}.
\end{enumerate}

\section*{Bab 7: Basis, Dimensi, dan Transformasi Linear}
\begin{enumerate}[label=\textbf{7.\arabic*}]
  \item Selesaikan $c_{1}(1,1,0)+c_{2}(1,0,1)+c_{3}(0,1,1)=(2,-1,3)$, yaitu
        $c_{1}+c_{2}=2$, $c_{1}+c_{3}=-1$, $c_{2}+c_{3}=3$. Hasilnya
        $c=(-1,3,0)$, jadi $[\mathbf{w}]_{B'}=(-1,3,0)$. Balik:
        $-1(1,1,0)+3(1,0,1)+0=(2,-1,3)$.

  \item Misal $\varphi$ isomorfisme. Karena $\varphi$ bijektif,
        $\mathbf{b}_{i}=\varphi^{-1}(\mathbf{e}_{i})$ terdefinisi; karena
        $\varphi^{-1}$ linear dan bijektif, $\{\mathbf{b}_{i}\}$ basis dan
        $\varphi=\varphi_{B}$. Tidak tunggal sebab setiap basis $B$ memberi
        isomorfisme yang berbeda.

  \item $A$ berperingkat $2$ (pivot di kolom $2,3$), nullity $1$. Inti:
        $x_{2}=0$, $2x_{3}=0$, jadi $\ker A=\operatorname{span}\{(1,0,0)\}$.
        Jangkauan $\operatorname{span}\{(1,0,0),(0,2,0)\}$. Ini matriks turunan
        pada $P_{2}$ dalam basis $(1,x,x^{2})$: $D(1)=0$, $D(x)=1$,
        $D(x^{2})=2x$.

  \item $\operatorname{im}(S\circ T)\subseteq\operatorname{im}S$ memberi
        rank$(S\circ T)\le$ rank $S$. Selain itu
        $\ker T\subseteq\ker(S\circ T)$, jadi nullity$(S\circ T)\ge$ nullity $T$;
        rank-nullity pada kedua peta memberi rank$(S\circ T)\le$ rank $T$.

  \item Salah. $\begin{pmatrix}1&0\\0&0\end{pmatrix}$ dan
        $\begin{pmatrix}0&1\\0&0\end{pmatrix}$ sama-sama berperingkat $1$, tetapi
        yang kedua nilpoten (jejak $0$, kuadrat nol) sedangkan yang pertama
        berjejak $1$. Rank hanya mengukur dimensi jangkauan.

  \item $T$ mengirim $\mathbf{b}_{1},\mathbf{b}_{2},\mathbf{b}_{3}$ ke $(1,0)$,
        $(3,1)$, $(3,2)$. Koordinat terhadap $C$ ialah $(x,y)\mapsto(y,x-y)$,
        jadi kolomnya $(0,1),(1,2),(2,1)$ dan
        $[T]_{C\leftarrow B}=\begin{pmatrix}0&1&2\\1&2&1\end{pmatrix}$.
        Untuk $\mathbf{v}=(1,2,3)$: $[\mathbf{v}]_{B}=(-1,-1,3)$ memberi
        $(5,0)$, sama dengan $[T(\mathbf{v})]_{C}$.

  \item Ambil basis $\mathbf{u}_{1},\dots,\mathbf{u}_{m}$ dari $U$ dengan
        $m=\dim V$. Ia bebas linear dan berukuran $\dim V$, jadi merentang $V$;
        maka setiap $\mathbf{v}\in V$ termuat di $U$, yaitu $U=V$.

  \item $B=P^{-1}AP$, jadi $B\mathbf{x}=\mathbf{0}\iff AP\mathbf{x}=\mathbf{0}
        \iff P\mathbf{x}\in\ker A$. Pemetaan $\mathbf{x}\mapsto P\mathbf{x}$
        isomorfisme dari $\ker B$ ke $\ker A$, sehingga nullity sama. Kesamaan
        itu persis pernyataan bahwa dimensi inti tidak bergantung basis.

  \item Baris ketiga jumlah dua baris pertama, jadi rank $2$ dan nullity
        $4-2=2$. Vektor inti tak nol $(-2,-1,1,0)$: $F(-2,-1,1,0)=\mathbf{0}$.
        Arah itu tidak menggeser satu pun fitur, sebab
        $F(\mathbf{x}+t\mathbf{k})=F(\mathbf{x})$ untuk semua $t$.
\end{enumerate}

\section*{Bab 8: Sistem Persamaan Linear, Determinan, dan Invers}
\begin{enumerate}[label=\textbf{8.\arabic*}]
  \item Pivot yang muncul $1$, $1$, $-1$; tidak ada baris yang perlu ditukar, jadi
        $\det A=1\cdot1\cdot(-1)=-1$. Mundur: $z=3$, $y=-1$, $x=2$. Cramer
        memberi $\det A_{x}=-2$, $\det A_{y}=1$, $\det A_{z}=-3$; dibagi $-1$
        hasilnya $(2,-1,3)$.

  \item Ekspansi kofaktor memberi $\det M=0$. Penyangkalnya
        $\mathbf{v}=(1,-1,-1,1)$: tiap baris memberi $0$, misalnya
        $0-1-2+3=0$. Jadi $\det M=0$ tanpa ekspansi. numpy memberi
        $5{,}3\times10^{-15}$, bukan nol, karena eliminasi mengurangkan bilangan
        hampir sama dan galat pembulatan ikut terhitung.

  \item Rumus Leibniz: $\det A=\sum_{\sigma}\operatorname{sgn}(\sigma)
        \prod_{i}a_{i\sigma(i)}$. Bila kolom $j$ dan $k$ sama, pasangkan tiap
        $\sigma$ dengan $\sigma\circ(j\,k)$; hasilkali sukunya sama tetapi
        tandanya berlawanan, jadi jumlah seluruh suku nol.

  \item Kalikan $A\operatorname{adj}A=(\det A)I$ dengan determinan:
        $\det A\cdot\det(\operatorname{adj}A)=(\det A)^{n}$, sebab
        $\det(cI)=c^{n}$. Bila $\det A\ne0$, bagi dengan $\det A$. Bila
        $\det A=0$, kedua ruas nol untuk $n\ge2$, jadi kesamaannya tetap sah.

  \item Salah sebagai pernyataan umum. Kepekaan diukur
        $\kappa(A)=\lVert A\rVert\,\lVert A^{-1}\rVert$, bukan $\det A$:
        $A=\varepsilon I$ punya $\det A=\varepsilon^{n}$ sangat kecil tetapi
        $\kappa=1$. Determinan kecil baru berbahaya bila disertai
        $\lVert A^{-1}\rVert$ besar.

  \item $\lVert A\rVert_{1}=2+\varepsilon$ dan
        $A^{-1}=\frac1\varepsilon\begin{pmatrix}1+\varepsilon&-1\\-1&1\end{pmatrix}$
        memberi $\lVert A^{-1}\rVert_{1}=(2+\varepsilon)/\varepsilon$. Jadi
        $\kappa_{1}(A)=(2+\varepsilon)^{2}/\varepsilon$. Syarat $\kappa>10^{6}$
        ialah $\varepsilon<4{,}000016\times10^{-6}$; pada
        $\varepsilon=4\times10^{-6}$ numpy memberi $1{,}000004\times10^{6}$.

  \item Menambah kelipatan baris berarti mengalikan kiri dengan
        $E=I+cE_{ij}$, dan $\det E=1$, jadi $\det B=\det A$. Karena $B=EA$,
        maka $B^{-1}=A^{-1}E^{-1}$ dengan $E^{-1}=I-cE_{ij}$.

  \item $A+AD=A(I+D)$, dan $I+D$ diagonal, jadi
        $\det(A+AD)=\det A\prod_{i}(1+d_{i})$. Invertibel tepat bila
        $\det A\ne0$ dan tidak ada $d_{i}=-1$. Contoh $\det A=5$ dengan
        $D=\operatorname{diag}(2,-\tfrac12,3)$: kiri $30$, kanan
        $5\cdot3\cdot\tfrac12\cdot4=30$.

  \item Dengan $A=\begin{pmatrix}1&2&3\\0&1&4\\5&6&0\end{pmatrix}$,
        $\det A=1$ dan
        \[
          \operatorname{adj}A=A^{-1}=
          \begin{pmatrix}-24&18&5\\20&-15&-4\\-5&4&1\end{pmatrix},
        \]
        sama dengan \texttt{numpy.linalg.inv} sampai galat mesin. Untuk
        $\det A=7$:
        $\begin{pmatrix}2&1&0\\1&1&0\\0&0&7\end{pmatrix}^{-1}$ memuat
        $1/7$ di entri ketiga, sebab pembagian oleh $7$ tidak terhindarkan.
\end{enumerate}

\section*{Bab 9: Nilai Eigen dan Diagonalisasi}
\begin{enumerate}[label=\textbf{9.\arabic*}]
  \item $\operatorname{tr}A_{1}=5$ dan $\det A_{1}=4$, jadi $\lambda=4$ dan
        $\lambda=1$. Vektor eigen: $(1,1)$ untuk $\lambda=4$, $(1,-2)$ untuk
        $\lambda=1$. Maka $P=\begin{pmatrix}1&1\\1&-2\end{pmatrix}$ dan
        $D=\operatorname{diag}(4,1)$; $P^{-1}A_{1}P=D$ terverifikasi.

  \item $\lambda=5,2$ dengan vektor eigen $(2,1)$ dan $(1,-1)$. Untuk
        $k\ge1$,
        \[
          G^{k}=\frac13
          \begin{pmatrix}2\cdot5^{k}+2^{k}&2\cdot5^{k}-2\cdot2^{k}\\
          5^{k}-2^{k}&5^{k}+2\cdot2^{k}\end{pmatrix},
        \]
        sehingga
        $G^{10}=\begin{pmatrix}6510758&6509734\\3254867&3255891\end{pmatrix}$,
        sama dengan \texttt{matrix\_power}.

  \item $A\mathbf{v}=\lambda\mathbf{v}$ memberi
        $A^{2}\mathbf{v}=\lambda^{2}\mathbf{v}$; induksi memberi
        $A^{k}\mathbf{v}=\lambda^{k}\mathbf{v}$. Kebalikannya gagal:
        $A=\begin{pmatrix}0&-1\\1&0\end{pmatrix}$ punya $A^{2}=-I$ dengan nilai
        eigen $-1$, sedangkan $A$ sendiri tidak punya nilai eigen real.

  \item $0$ nilai eigen $\iff$ ada $\mathbf{v}\ne\mathbf{0}$ dengan
        $A\mathbf{v}=\mathbf{0}$ $\iff$ $\ker A\ne\{\mathbf{0}\}$ $\iff$ $A$ tak
        invertibel. Bila $A$ invertibel dan $A\mathbf{v}=\lambda\mathbf{v}$, maka
        $\lambda\ne0$ dan $A^{-1}\mathbf{v}=\lambda^{-1}\mathbf{v}$.

  \item Salah. $\begin{pmatrix}1&0\\0&1\end{pmatrix}$ dan
        $\begin{pmatrix}1&1\\0&1\end{pmatrix}$ sama-sama berpolinomial
        karakteristik $(\lambda-1)^{2}$, tetapi yang kedua defektif:
        $\operatorname{rank}(J-I)=1$, jadi nullity $1<2$.

  \item $\lambda I-A$ dan transposnya berdeterminan sama, sebab
        $\det X=\det X^{\top}$; maka polinomial karakteristik $A$ dan
        $A^{\top}$ sama dan nilai eigennya sama. Vektor eigennya tidak:
        $A$ dan $A^{\top}$ punya ruang eigen berbeda begitu $A$ tak simetris.

  \item $K-2I=\begin{pmatrix}0&1&0\\0&0&0\\0&0&-1\end{pmatrix}$ berperingkat $2$,
        jadi nullity $=1$. Multiplisitas aljabar $\lambda=2$ adalah $2$,
        geometriknya $1$; karena $1<2$, $K$ tidak dapat didiagonalkan.

  \item $\lambda=1$ dan $0{,}5$. Titik tetap dari $x=4y$ dengan $x+y=1$ ialah
        $(0{,}8;0{,}2)$. Galat menyusut seperti $0{,}5^{k}$; karena
        $0{,}5^{19}=1{,}91\times10^{-6}>\allowbreak10^{-6}$ dan
        $0{,}5^{20}=9{,}54\times10^{-7}$, perlu $20$ langkah.

  \item Matriks tercetak memberi
        $\operatorname{tr}(B^{10})=(-19)+21+1024=1026$, dan $1^{10}+1^{10}+2^{10}=1026$: cocok. Untuk $A$ yang
        dapat didiagonalkan, $A^{k}=PD^{k}P^{-1}$ memberi
        $\operatorname{tr}(A^{k})=\operatorname{tr}(D^{k})=\sum_{i}\lambda_{i}^{k}$.
        Untuk $B$ yang defektif langkah itu tidak tersedia, sebab $B$ tidak punya
        $P$ dan $D$. Gantinya pakai bentuk segitiga Schur, $B=QTQ^{-1}$ dengan $T$
        segitiga atas berdiagonal $\lambda_{i}=1,1,2$; maka $B^{k}=QT^{k}Q^{-1}$,
        dan $T^{k}$ tetap segitiga atas dengan diagonal $\lambda_{i}^{k}$, sehingga
        $\operatorname{tr}(B^{k})=\operatorname{tr}(T^{k})
        =\sum_{i}\lambda_{i}^{k}=1026$. Jejak invarian terhadap basis
        ($\operatorname{tr}XY=\operatorname{tr}YX$), jadi hasilnya tidak bergantung
        pada $Q$. Kesimpulannya: diagonalisasi dipakai pada bukti hanya untuk kasus
        yang dapat didiagonalkan; untuk kasus defektif bukti berganti ke bentuk
        segitiga, rumusnya tetap sama, dan faktor $k$ yang muncul di entri $B^{k}$
        tidak pernah menyentuh diagonal.
\end{enumerate}

\section*{Bab 10: Ortogonalitas, SVD, dan Kuadrat Terkecil}
\begin{enumerate}[label=\textbf{10.\arabic*}]
  \item $(1,0,0)$ dan $(0,1,1)$ sudah ortogonal, jadi
        $\operatorname{proj}_{W}\mathbf{b}=1\cdot(1,0,0)+\frac{4}{2}(0,1,1)
        =(1,2,2)=\mathbf{b}$. Karena $\mathbf{b}\in W$, jaraknya $0$;
        \texttt{qr} memberi $QQ^{\top}\mathbf{b}=(1,2,2)$ juga.

  \item $\mathbf{b}=\mathbf{p}+\mathbf{r}$ dengan $\mathbf{r}\perp W$, jadi
        $\mathbf{r}\perp\mathbf{p}$ dan Pythagoras memberi
        $\lVert\mathbf{b}\rVert^{2}=\lVert\mathbf{p}\rVert^{2}
        +\lVert\mathbf{r}\rVert^{2}$. Bila $\mathbf{p}$ bukan proyeksi
        ortogonal, $\mathbf{r}$ tidak tegak lurus $\mathbf{p}$; pada basis miring
        di $\mathbb{R}^{2}$ muncul suku silang $2\langle\mathbf{p},\mathbf{r}
        \rangle\ne0$.

  \item Kolom ortonormal: $(A^{\top}A)_{ij}=\mathbf{a}_{i}^{\top}\mathbf{a}_{j}
        =\delta_{ij}$, jadi $A^{\top}A=I$. Persamaan normal
        $A^{\top}A\hat{\mathbf{x}}=A^{\top}\mathbf{b}$ langsung memberi
        $\hat{\mathbf{x}}=A^{\top}\mathbf{b}$. Semua nilai singular $A$ sama
        dengan $1$, jadi $\kappa(A)=1$.

  \item Salah. $A=\begin{pmatrix}1&0\\0&1\\1&1\end{pmatrix}$ memberi
        $A^{\top}A=\begin{pmatrix}2&1\\1&2\end{pmatrix}$ yang invertibel
        ($\det=3$), tetapi barisnya bergantung: baris ketiga jumlah dua baris
        pertama. Yang dijamin $A^{\top}A$ invertibel hanyalah kebebasan kolom.

  \item Nilai singular $2$ dan $0$, jadi rank $1$;
        $U=V=\frac{1}{\sqrt2}\begin{pmatrix}1&1\\1&-1\end{pmatrix}$ (sampai
        tanda). $A^{+}=\frac14\begin{pmatrix}1&1\\1&1\end{pmatrix}$, sehingga
        $A^{+}\mathbf{b}=(1,1)$ dan sisanya $(-1,1)$ dengan norma
        $\sqrt2=1{,}4142$; \texttt{pinv} sama.

  \item 
Tulis $\mathbf{b}=A\hat{\mathbf{x}}+\mathbf{r}$ dengan
        $A\hat{\mathbf{x}}\in\operatorname{col}A$. Untuk
        $\mathbf{h}\in\operatorname{col}A$, fungsi
        $t\mapsto\lVert\mathbf{b}-A(\hat{\mathbf{x}}+t\mathbf{h})\rVert^{2}$ minimum
        di $t=0$, sehingga $\langle\mathbf{r},\mathbf{h}\rangle=0$; ini arah
        ``perlu''. Sebaliknya, bila $\langle\mathbf{r},\mathbf{h}\rangle=0$ untuk
        setiap $\mathbf{h}\in\operatorname{col}A$, maka
        $\lVert\mathbf{b}-A\mathbf{x}\rVert^{2}=\lVert\mathbf{r}\rVert^{2}+\lVert
        A(\hat{\mathbf{x}}-\mathbf{x})\rVert^{2}\ge\lVert\mathbf{r}\rVert^{2}$, jadi
        $\hat{\mathbf{x}}$ meminimumkan. Kebebasan kolom dipakai agar
        $\hat{\mathbf{x}}$ tunggal.

  \item $\mathbf{u}_{1}=(1,1,1)/\sqrt3$; sisa $(1,1,0)-\frac23(1,1,1)
        =\frac13(1,1,-2)$, jadi $\mathbf{u}_{2}=(1,1,-2)/\sqrt6$; lalu
        $\mathbf{u}_{3}=(1,-1,0)/\sqrt2$. $\lVert Q^{\top}Q-I\rVert$:
        $6{,}7\times10^{-16}$ (tangan), $3{,}1\times10^{-16}$ (\texttt{qr}).
        Tanda boleh berbeda sebab tiap $\mathbf{u}_{i}$ hanya tertentu sampai
        $\pm$.

  \item $AA^{\top}\mathbf{x}=\mathbf{0}$ memberi
        $\lVert A^{\top}\mathbf{x}\rVert^{2}=\mathbf{x}^{\top}AA^{\top}
        \mathbf{x}=0$, jadi $A^{\top}\mathbf{x}=\mathbf{0}$; kebebasan baris
        membuat $\mathbf{x}=\mathbf{0}$. Maka $AA^{\top}$ persegi dan invertibel.
        Untuk $\mathbf{x}=A^{\top}(AA^{\top})^{-1}\mathbf{b}$ berlaku
        $A\mathbf{x}=\mathbf{b}$ dan
        $\mathbf{x}\in\operatorname{col}A^{\top}=(\ker A)^{\perp}$.

  \item Dengan bacaan terakhir $314{,}3$: $\sum y=1540{,}5$,
        $\sum Ty=157161{,}7$, determinan $50$. Cramer:
        $a=(52030\cdot1540{,}5-510\cdot157161{,}7)/50=-5{,}04$ dan
        $b=(5\cdot157161{,}7-510\cdot1540{,}5)/50=3{,}07$;
        \texttt{lstsq} sama. Intercept bergeser $-10{,}10$, lereng $+0{,}10$;
        titik ujung berleverage $0{,}6$, terbesar di antara kelimanya.

  \item $B$ berperingkat $1$, jadi $B^{+}$ memberi
        $\mathbf{x}^{+}=\frac{1}{70}(1,2)\approx(0{,}01429;\,0{,}02857)$ dan
        $\mathbf{k}=(-2,1)$. Semua penyelesaian $\mathbf{x}^{+}+t\mathbf{k}$;
        karena $\mathbf{x}^{+}\perp\mathbf{k}$,
        $\lVert\mathbf{x}^{+}+t\mathbf{k}\rVert^{2}=\lVert\mathbf{x}^{+}
        \rVert^{2}+t^{2}\lVert\mathbf{k}\rVert^{2}$ minimum di $t=0$.
\end{enumerate}

\section*{Bab 11: Kekontinuan dan Kekontinuan Seragam}
\begin{enumerate}[label=\textbf{11.\arabic*}]
  \item $\lvert(2x-1)-5\rvert=2\lvert x-3\rvert$, jadi $\delta=\varepsilon/2$
        sah dan terbesar: pada $x=3+\delta$ nilainya tepat $\varepsilon$.
        Pilihan $\delta=\varepsilon/3$ juga sah, tetapi hanya memakai sepertiga
        ruang yang tersedia.

  \item $\lvert x^{2}-c^{2}\rvert=\lvert x-c\rvert\,\lvert x+c\rvert$; dengan
        $\lvert x-c\rvert<1$ didapat $\lvert x+c\rvert\le2\lvert c\rvert+1$,
        jadi $\delta=\min(1,\varepsilon/(2\lvert c\rvert+1))$. Di $c=2$:
        $\delta=0{,}002$; di $c=100$: $\delta=4{,}975\times10^{-5}$.

  \item Sebagian benar. $1/x$ memang kontinu di setiap titik $(0,1]$. Tetapi
        kesimpulan seragamnya keliru: $(0,1]$ tidak kompak dan $1/x$ tak terbatas
        di dekat $0$. Penyangkalnya $x_{n}=1/n$, $y_{n}=1/(n+1)$:
        $\lvert x_{n}-y_{n}\rvert\to0$ tetapi selisih nilainya $1$.

  \item Diberi $\varepsilon$, ambil $\delta_{f}$ dan $\delta_{g}$; dengan
        $\delta=\min(\delta_{f},\delta_{g})$,
        $\lvert(f+g)(x)-(f+g)(y)\rvert\le\varepsilon/2+\varepsilon/2$.
        Untuk $fg$: $f=g=x$ pada $\mathbb{R}$ dengan $x_{n}=n$,
        $y_{n}=n+1/n$ memberi selisih $2+1/n^{2}\ge2$.

  \item $\lvert\sqrt x-\sqrt y\rvert^{2}\le\lvert x-y\rvert$ (rasionalkan
        penyebutnya), jadi $\delta=\varepsilon^{2}$ bekerja untuk semua
        $x,y\ge0$. Yang dipakai hanya $x,y\ge0$ agar akarnya terdefinisi;
        ketertutupan tidak perlu sebab bukti tidak menyentuh titik limit di luar
        daerah.

  \item 
$(1+\delta)^{3}-1=\varepsilon$ memberi
        $\delta=(1+\varepsilon)^{1/3}-1=0{,}0033223$. Pilihan
        $\delta=\varepsilon/3=0{,}0033333$ justru terlalu besar:
        $(1+\varepsilon/3)^{3}-1-\varepsilon=\varepsilon^{2}/3+\varepsilon^{3}/27>0$
        untuk setiap $\varepsilon>0$, jadi batasnya dilanggar dan pilihan itu tidak
        sah.

  \item Dengan $\delta=\varepsilon/2$: selisih terbesar
        $2\delta-\delta^{2}=0{,}009975<\varepsilon$, jadi aman tetapi tidak
        optimal. Dengan $\delta=1-\sqrt{1-\varepsilon}=0{,}0050126$: selisih
        terbesar $1-(1-\delta)^{2}=\varepsilon$ tepat, dicapai di
        $x=1-\delta$.

  \item $\lvert f\rvert$ kontinu pada $[a,b]$ kompak, jadi mencapai minimum $m$;
        karena $f\ne0$ di mana-mana, $m>0$ dan $\lvert f\rvert\ge m$. Untuk
        $f(x)=x$ pada $(0,1)$ kesimpulannya gagal: infimumnya $0$ dan tidak
        tercapai.

  \item Di $c\ne0$ ambil $q_{n}\to c$ rasional dan $r_{n}\to c$ irasional; nilai
        fungsinya menuju $c$ dan $-c$, berbeda, jadi $f$ tidak kontinu. Di $0$:
        $\lvert f(x)\rvert=\lvert x\rvert\to0=f(0)$. Fungsi $\lvert f\rvert$
        ialah $x\mapsto\lvert x\rvert$, kontinu di mana-mana; tidak ada
        kontradiksi karena $\lvert f\rvert$ membuang informasi tanda.
\end{enumerate}

\section*{Bab 12: Turunan: Satu dan Banyak Variabel}
\begin{enumerate}[label=\textbf{12.\arabic*}]
  \item $\bigl((2+h)^{3}-8\bigr)/h=12+6h+h^{2}\to12$. Untuk $\sqrt x$:
        $(\sqrt{2+h}-\sqrt2)/h=1/(\sqrt{2+h}+\sqrt2)\to1/(2\sqrt2)
        =0{,}3535534$. Pada kedua kasus $h$ hilang dari penyebut.

  \item $f'(c)>0$ membuat $(f(c+h)-f(c))/h>f'(c)/2>0$ untuk $h$ cukup kecil,
        jadi $f(c+h)>f(c)$. Itu hanya pernyataan satu sisi. Contoh
        $f(x)=x+2x^{2}\sin(1/x^{2})$ punya $f'(0)=1>0$ tetapi tidak naik pada
        lingkungan mana pun di sekitar $0$: turunannya berganti tanda sekehendak
        dekat $0$.

  \item Turunkan $1/g$ lebih dulu: $(1/g)'=-g'/g^{2}$; di sini kekontinuan $g$
        dipakai, agar $g(x+h)\to g(x)\ne0$. Lalu
        $(f/g)'=f'\cdot(1/g)+f\cdot(1/g)'$, yang disederhanakan menjadi
        $(f'g-fg')/g^{2}$.

  \item Salah. $f(x,y)=xy/(x^{2}+y^{2})$ untuk $(x,y)\ne(0,0)$ dan
        $f(0,0)=0$: kedua turunan parsial di titik asal ada dan bernilai $0$,
        tetapi sepanjang $y=x$ nilainya $\tfrac12$, jadi $f$ tidak kontinu di
        sana.

  \item $\nabla f=(3x^{2}y-y^{3},\,x^{3}-3xy^{2})$, di $(2,1)$ sama dengan
        $(11,2)$. Arah naik tercepat $(11,2)/\sqrt{125}$, laju maksimumnya
        $\sqrt{125}=11{,}1803$. Beda pusat $h=10^{-4}$ memberi galat relatif
        $9{,}1\times10^{-10}$ dan $1{,}0\times10^{-8}$, sesuai orde $h^{2}$ pada
        komponen kedua.

  \item Menurut teorema nilai rata-rata,
        $f(b)-f(a)=f'(\xi)(b-a)=0$ untuk setiap pasangan $a<b$ di interval, jadi
        $f$ konstan. Pada himpunan tak terhubung langkah itu gagal karena $\xi$
        belum tentu berada di himpunan; contohnya $f=0$ pada $(0,1)$ dan $f=1$
        pada $(2,3)$.

  \item Di $(0,1)$: $J_{H}=\begin{pmatrix}1&1\\1&2\end{pmatrix}$; di
        $H(0,1)=(1,1)$: $J_{G}=\begin{pmatrix}2&1\\1&-1\end{pmatrix}$. Hasilkali
        matriksnya $\begin{pmatrix}3&4\\0&-1\end{pmatrix}$, dan turunan langsung
        memberi matriks yang sama.

  \item $f_{x}=3x^{2}-3y$, $f_{y}=3y^{2}-3x$; keduanya nol memberi $y=x^{2}$ dan
        $x=y^{2}$, jadi $x(x^{3}-1)=0$: titik $(0,0)$ dan $(1,1)$. Hessian
        $\begin{pmatrix}6x&-3\\-3&6y\end{pmatrix}$: di $(0,0)$ determinannya
        $-9<0$ (pelana), di $(1,1)$ determinannya $27>0$ dengan $f_{xx}=6>0$
        (minimum lokal, $f=-1$). Bila determinan nol, uji itu bungkam.

  \item $p_{2}(x)=1+\tfrac12h-\tfrac18h^{2}$ di $h=0{,}04$ memberi $1{,}0198$,
        sedangkan $\sqrt{1{,}04}=1{,}0198039$; galatnya
        $3{,}9\times10^{-6}$, sisa Lagrange
        $\frac{3}{8}\xi^{-5/2}h^{3}/6\le4{,}0\times10^{-6}$. Batasnya rapat, bukan longgar, sebab $\xi$
        hampir $1$: sisa Lagrange memang hampir menyamai galat sebenarnya.
\end{enumerate}

\section*{Bab 13: Integral Riemann dan Teorema Dasar Kalkulus}
\begin{enumerate}[label=\textbf{13.\arabic*}]
  \item $U_{n}=\frac1n\sum_{k=1}^{n}\frac kn=\frac{n+1}{2n}$ dan
        $L_{n}=\frac{n-1}{2n}$, jadi $U_{n}-L_{n}=1/n\to0$: $f$ terintegralkan
        dan $\int_{0}^{1}x\,dx=\tfrac12$.

  \item $\bigl\lvert\lvert f(x)\rvert-\lvert f(y)\rvert\bigr\rvert
        \le\lvert f(x)-f(y)\rvert$ (ketaksamaan segitiga terbalik), jadi
        keterintegralan terbawa. Batasnya dari $-\lvert f\rvert\le f\le\lvert f\rvert$
        dan monotonisitas integral. Kebalikannya salah: $f=1$ di rasional dan
        $-1$ di irasional punya $\lvert f\rvert=1$ terintegralkan.

  \item Salah. Bagian yang benar: jumlah atas dan bawah fungsi terbatas selalu
        berhingga. Yang tidak didukung bukti ialah kesamaan
        $\sup_{\mathcal{P}}L=\inf_{\mathcal{P}}U$. Fungsi Dirichlet terbatas
        dengan $U-L=1$ untuk setiap partisi.

  \item Pada partisi seragam tiap selang menyumbang osilasi
        $f(x_{k})-f(x_{k-1})$, sehingga
        $U-L\le\frac{b-a}{n}\sum_{k}\bigl(f(x_{k})-f(x_{k-1})\bigr)
        =\frac{(b-a)(f(b)-f(a))}{n}$. Untuk fungsi tak monoton teleskop itu gagal
        sebab tanda osilasinya tidak seragam.

  \item $\int_{0}^{1}(3x^{2}-2x+1)\,dx=[x^{3}-x^{2}+x]_{0}^{1}=1$. Simpson
        $n=100$ memberi $1{,}0000000000$ dan \texttt{quad} $1{,}0000000000$
        (galat $1{,}1\times10^{-14}$); selisih ketiga angka nol sampai presisi
        mesin.

  \item $\int_{1}^{10}x^{-2}dx=1-\tfrac1{10}=0{,}9$ dan
        $\int_{1}^{\infty}x^{-2}dx=1$. Selisihnya $0{,}1=\tfrac1{10}$ persis,
        yaitu ekor $\int_{10}^{\infty}x^{-2}dx=1/10$; tidak perlu dihitung ulang
        karena bagian yang hilang itulah ekornya.

  \item Bila $f(x_{0})>0$, kekontinuan memberi $f>\tfrac12f(x_{0})$ pada selang
        selebar $2\delta$ di sekitar $x_{0}$ (selebar $\delta$ bila $x_{0}$ titik
        ujung), jadi integralnya paling sedikit
        $\tfrac12\delta f(x_{0})>0$: kontradiksi. Untuk $f$ yang hanya terintegralkan
        pernyataan gagal: $f=1$ di satu titik dan $0$ di tempat lain berintegral
        $0$.

  \item $u=\sin x$: $\int_{0}^{1}u\,du=\tfrac12$. Lewat identitas:
        $\tfrac12\int_{0}^{\pi/2}\sin2x\,dx
        =\tfrac12\bigl[-\tfrac12\cos2x\bigr]_{0}^{\pi/2}=\tfrac12$. \texttt{quad}
        memberi $0{,}5$ juga.

  \item Pada setiap selang bagian ada rasional (nilai $1$) dan irasional (nilai
        $0$), jadi $U(\mathcal{P},D)=1$ dan $L(\mathcal{P},D)=0$; selisihnya $1$
        untuk semua $\mathcal{P}$. Kriteria Riemann menuntut selisih itu bisa
        dibuat sekecil apa pun, jadi $D$ tidak terintegralkan.

  \item Batas integral memberi $\frac1{101}<\sum_{n=101}^{\infty}n^{-2}
        <\frac1{100}$, yaitu $9{,}9010\times10^{-3}$ sampai
        $1{,}0000\times10^{-2}$. Titik tengahnya $9{,}9505\times10^{-3}$,
        sedangkan nilai sebenarnya $9{,}950167\times10^{-3}$; galat relatifnya
        $3{,}3\times10^{-5}$.
\end{enumerate}

\section*{Bab 14: Deret Fungsi dan Pertukaran Limit}
\begin{enumerate}[label=\textbf{14.\arabic*}]
  \item $\sup_{[0,1]}\lvert x^{n}\rvert=1$ untuk semua $n$, dicapai di $x=1$.
        Di $x=0{,}9$: $0{,}3487$ untuk $n=10$, $2{,}656\times10^{-5}$ untuk
        $n=100$, $1{,}748\times10^{-46}$ untuk $n=1000$. Kolom pertama datar
        sebab puncaknya selalu di ujung; kolom kedua menyusut sebab $0{,}9<1$.

  \item Untuk $n\ge N$, $\lVert f_{n}-f\rVert_{\infty}<\varepsilon$, maka
        $\lVert f\rVert_{\infty}\le\lVert f_{N}\rVert_{\infty}+1$ dan begitu
        pula $\lVert f_{n}\rVert_{\infty}$. Titik demi titik tidak cukup:
        $f_{n}=n\chi_{(0,1/n)}$ konvergen ke $0$ di setiap titik tetapi
        $\lVert f_{n}\rVert_{\infty}=n$.

  \item Puncaknya di $x=1/\sqrt n$ dengan
        $\lVert f_{n}\rVert_{\infty}=1/(2\sqrt n)\to0$, dan batas itu berlaku
        untuk semua $x$ sekaligus, jadi seragam. Syarat
        $1/(2\sqrt n)<0{,}01$ menuntut $\sqrt n>50$, jadi $n=2501$.

  \item Salah. Penalaran titik demi titik memberi satu $N$ untuk tiap $x$, bukan
        satu $N$ untuk semua $x$. Penyangkalnya $x_{n}=2^{-1/n}$:
        $x_{n}^{n}=\tfrac12$ untuk setiap $n$, jadi
        $\sup_{[0,1)}\lvert x^{n}\rvert\ge\tfrac12$ dan tidak menuju $0$.

  \item Galat $\lvert e-S_{N}(1)\rvert$ untuk $N=5,10,15$:
        $1{,}615\times10^{-3}$, $2{,}731\times10^{-8}$,
        $5{,}077\times10^{-14}$; batas $e/(N+1)!$ memberi
        $3{,}775\times10^{-3}$, $6{,}810\times10^{-8}$,
        $1{,}299\times10^{-13}$, selalu lebih besar. Di $N=20$ batasnya sudah
        $5{,}32\times10^{-20}$, jauh di bawah presisi mesin, sehingga galat
        terukur berhenti di $4{,}44\times10^{-16}$. Batas di bawah $10^{-16}$
        pertama di $N=18$, dan di sana galat sebenarnya juga tak terukur.

  \item Jumlahnya $(1-x)/(1-x)=1$ untuk $x\in[0,1)$, sedangkan di $x=1$
        jumlahnya $0$. Limitnya tidak kontinu padahal setiap jumlah parsial
        kontinu; limit seragam fungsi kontinu harus kontinu, jadi
        kekonvergenannya tidak seragam.

  \item Jari-jari: $R=1$ untuk $\sum x^{n}/n^{2}$, $R=e$ untuk
        $\sum(n!/n^{n})x^{n}$, $R=\tfrac12$ untuk $\sum2^{n}x^{n}/n$. Nisbah
        suku ke-$1001$ dan ke-$1000$: $0{,}998003$; $0{,}368063$ (jadi
        $R\approx2{,}7169$); $1{,}998002$. Paling lambat $\sum x^{n}/n^{2}$,
        sebab nisbahnya menuju $1$ seperti $(1-1/n)^{2}$.

  \item $\lvert f_{n}g-fg\rvert=\lvert f_{n}-f\rvert\,\lvert g\rvert
        \le\lVert f_{n}-f\rVert_{\infty}\lVert g\rVert_{\infty}\to0$. Tanpa
        keterbatasan $g$: $f_{n}=x+1/n\to x$ seragam pada $(0,1)$, tetapi
        $(x+1/n)/x=1+1/(nx)$ tak terbatas, jadi $f_{n}g$ tidak konvergen
        seragam.

  \item Puncak di $x=1/\sqrt{2n+1}$. Untuk $n=10$: tinggi $1{,}3397$ di
        $x=0{,}2182$, $\int_{0}^{1}f_{10}=\tfrac{5}{11}=0{,}4545$. Untuk
        $n=100$: tinggi $4{,}2835$ di $x=0{,}07053$, integralnya
        $\tfrac{50}{101}=0{,}4950$. Tinggi mengikuti $\sqrt{n/(2e)}$; luasnya
        menuju $\tfrac12$ sebab puncak makin sempit.

  \item Turunan $\sum_{n\ge0}x^{n}=1/(1-x)$ memberi
        $\sum_{n\ge1}nx^{n-1}=1/(1-x)^{2}$; di $x=\tfrac12$ didapat
        $\sum_{n\ge1}n/2^{n}=2$. Sisa ekornya $(N+2)/2^{N}$, jadi
        $N=39$ memberi $7{,}46\times10^{-11}<10^{-10}$.
\end{enumerate}

\section*{Bab 15: Integral Lipat, Perubahan Variabel, dan Teorema Lapangan Vektor}
\begin{enumerate}[label=\textbf{15.\arabic*}]
  \item Cartesius: $\int_{-1}^{1}\int_{-\sqrt{1-x^{2}}}^{\sqrt{1-x^{2}}}
        (x^{2}+y^{2})\,dy\,dx=\pi/2$. Polar:
        $\int_{0}^{2\pi}\int_{0}^{1}r^{3}\,dr\,d\theta
        =2\pi\cdot\tfrac14=\pi/2=1{,}5708$. \texttt{dblquad} memberi
        $1{,}5707963268$; di lingkungan uji ini cara Cartesius menuntut $8379$
        evaluasi, cara polar $441$.

  \item Untuk $T(\mathbf{u})=A\mathbf{u}$ linear, $DT=A$ konstan, jadi
        determinan Jacobiannya $\det A$ dan $\lvert\det A\rvert$ adalah faktor
        skala luas. Bila $\det A\ne0$ peta itu satu-satu pada $D$, sehingga
        $\iint_{T(D)}f=\lvert\det A\rvert\iint_{D}f\circ T$.

  \item Searah jarum jam tiap sisi ditempuh terbalik, jadi keempat suku contoh
        berganti tanda: $-\tfrac13-2-\tfrac23+1=-2$, sedangkan ruas kanan tetap
        $2$. Sisi bawah dan kanan yang tadinya menyumbang $+\tfrac13$ dan $+2$
        menjadi negatif; sisi kiri yang tadinya $-1$ menjadi $+1$.

  \item $x=au$, $y=bv$, $z=cw$ memberi Jacobian $abc$ dan bola satuan:
        $V=abc\cdot\frac{4\pi}{3}=\frac{4\pi abc}{3}$. Pada $a=b=c=1$
        hasilnya $4{,}1888$, yaitu rumus volume bola yang dikenal.

  \item Salah. $\operatorname{curl}F=0$ memang berlaku di setiap titik tempat $F$
        terdefinisi, tetapi daerahnya berlubang sehingga uji curl tidak berlaku.
        Angka pendukungnya: $\oint_{\lvert\mathbf{r}\rvert=1}F\cdot d\mathbf{r}
        =2\pi\ne0$, jadi $F$ tidak konservatif.

  \item Bila integral tidak bergantung lintasan, dua lintasan $C_{1},C_{2}$ dari
        $A$ ke $B$ memberi integral sama, jadi sirkulasi $C_{1}-C_{2}$ nol.
        Sebaliknya bila setiap kurva tertutup bersirkulasi nol, $C_{1}-C_{2}$
        tertutup sehingga integral keduanya sama.

  \item $\operatorname{curl}F=(0,0,2)$; fluks melalui $z=1-x^{2}-y^{2}$ dengan
        normal ke atas $=2\cdot\pi=2\pi$. Integral garis pada tepi
        $x^{2}+y^{2}=1$, $z=0$:
        $\int_{0}^{2\pi}(-\sin t,\cos t,0)\cdot(-\sin t,\cos t,0)\,dt=2\pi$.
        Keduanya $6{,}2832$.

  \item $\operatorname{div}F=0$ di luar titik asal, jadi integral lipatnya $0$,
        sedangkan fluks melalui bola satuan $4\pi=12{,}5664$. Teorema Gauss
        menuntut medan $C^{1}$ pada daerah yang memuat seluruh isi permukaan;
        syarat itu dilanggar di titik asal yang terkurung.

  \item $\det J=r^{2}\sin\varphi$, nol di $\varphi=0$ dan $\varphi=\pi$. Di kedua
        kutub peta berhenti satu-satu: semua $\theta$ memberi titik yang sama.
        Pada $0<\varphi<\pi$ peta satu-satu dan $\det J>0$, jadi rumus perubahan
        variabel tetap sah walaupun determinannya nol di batas.
\end{enumerate}

% Kunci jawaban Bab 16--23. Berkas ini di-\input langsung di dalam document.
% Setiap angka berasal dari skrip uji Python (lihat laporan).

\section*{Bab 16: Persamaan Diferensial dan Sistem Dinamik}
\begin{enumerate}[label=\textbf{16.\arabic*}]
  \item $y=Ce^{t^{2}}$ dengan $y(0)=1$ memberi $y=e^{t^{2}}$, jadi $y(1)=e=2{,}718282$.
        Euler $h=0{,}1$ memakai $y_{n+1}=(1+0{,}02n)y_{n}$; hasilnya $2{,}334633$,
        galat $3{,}8365\times10^{-1}$. Metode itu tertinggal karena memakai
        kemiringan di ujung kiri tiap langkah.

  \item Selisih $w=y_{1}-y_{2}$ memenuhi $w'=-p(t)w$ dan $w(t_{0})=0$. Kalikan dengan
        $e^{\int p}$: $\bigl(we^{\int p}\bigr)'=0$, jadi $we^{\int p}$ konstan, dan
        nilainya nol di $t_{0}$. Maka $w\equiv0$. Tak ada rumus solusi yang dipakai,
        hanya aturan hasilkali dan ketunggalan nilai awal.

  \item Belum cukup bukti. Linearisasi hanya melihat suku berderajat satu; nilai eigen
        $\pm i$ hanya mengatakan bagian linearnya rotasi. Dengan $r'=r^{3}$ titik asal
        pusat bagi linearisasi tetapi spiral keluar, sedangkan $r'=-r^{3}$ memberinya
        spiral masuk. Jadi $\pm i$ tidak menentukan kestabilan.

  \item $E'=vv'+\omega_{0}^{2}\theta\theta'=v(-\omega_{0}^{2}\theta)+\omega_{0}^{2}\theta
        v=0$, jadi $E$ konstan. Persamaan $v^{2}/(2E)+\theta^{2}/(2E/\omega_{0}^{2})=1$
        adalah elips, dan $\theta=A\cos(\omega_{0}t+\varphi)$ menutup satu putaran
        dalam $2\pi/\omega_{0}$.

  \item $\tau=RC=2{,}2\times10^{3}\cdot10^{-5}=0{,}022$ detik. Pada $t=\tau$:
        $V=12(1-e^{-1})=7{,}585447$ volt. Untuk $6$ volt, $t=-\tau\ln(1-\tfrac{6}{12})
        =0{,}022\cdot0{,}693147=15{,}249$ milidetik.

  \item Dengan $v=2\cos(\theta/2)$: $v^{2}/2=2\cos^{2}(\theta/2)$ dan
        $1-\cos\theta=2\sin^{2}(\theta/2)$, jumlahnya $E=2$. Di $\theta=\pm\pi$ nilai
        $v=0$, jadi kurvanya menyambung kedua titik itu; ia memisahkan ayunan
        dari putaran penuh, dan waktu menempuhnya tak hingga.

  \item $T_{0}=2\pi\sqrt{0{,}25/9{,}80665}=1{,}003205$ detik.
        $T=T_{0}\cdot\tfrac{2}{\pi}K(\sin15^{\circ})=1{,}020669$ detik; selisihnya
        $1{,}74$ persen. Amplitudo $30^{\circ}$ memperpanjang periode, bukan
        memendekkannya.

  \item Faktor per langkah $1+h\lambda=1-100h$; syarat $|1-100h|<1$ memberi
        $0<h<0{,}02$. Tepat pada $h=0{,}02$ faktornya $-1$: nilai berganti tanda
        dengan besar tetap $1$, tidak meluruh dan tidak meledak. Itu batas
        kestabilan marginal, bukan kestabilan.

  \item Andaikan $y$ periodik tidak konstan. Di titik maksimumnya $y'=0$, jadi
        $f(y_{\max})=0$. Solusi konstan $y\equiv y_{\max}$ melewati titik itu, dan
        ketunggalan (Lipschitz) memaksa solusinya sama, bertentangan dengan
        ketakkonsanan. Jadi tidak ada orbit periodik tak konstan.
\end{enumerate}

\section*{Bab 17: Peluang dan Peluang Bersyarat}
\begin{enumerate}[label=\textbf{17.\arabic*}]
  \item 
$P(A\cup B\cup C)=P(A\cup B)+P(C)-P\bigl((A\cup B)\cap C\bigr)$, lalu
        $P\bigl((A\cup B)\cap C\bigr)=P(A\cap C)+P(B\cap C)-P(A\cap B\cap C)$. Suku
        rangkap tiga terkurang tiga kali, satu untuk tiap pasangan, lalu
        dikembalikan sekali dengan tanda $+$.

  \item Salah pada ruang umum. $P(A\mid B)=0$ hanya berarti $P(A\cap B)=0$. Pada
        ruang berhingga dengan peluang positif, itu memaksa $A\cap B=\varnothing$.
        Pada ruang kontinu tidak: $A=\{1/2\}$ dan $B=(0,1)$ seragam memberi
        $P(A\cap B)=0$ padahal $1/2\in A\cap B$.

  \item Dari $10\,000$ orang: $200$ sakit, $190$ benar positif; dengan spesifisitas
        $0{,}90$ ada $980$ positif palsu, jadi $190/1170=19/117\approx0{,}1624$.
        Dengan $0{,}99$: positif palsunya tinggal $98$, jadi $190/288=95/144
        \approx0{,}6597$: ketepatannya naik empat kali lipat.

  \item $P(A)=P(A\cap B)+P(A\cap B^{c})$, jadi
        $P(A\cap B^{c})=P(A)-P(A)P(B)=P(A)(1-P(B))=P(A)P(B^{c})$. Untuk
        $A^{c},B^{c}$ pakai $P(A^{c}\cap B^{c})=1-P(A\cup B)$ dan inklusi--eksklusi,
        atau terapkan hasil tadi dua kali.

  \item $1-\prod_{k=0}^{29}(1-k/365)=0{,}706316$, sedangkan
        $1-e^{-435/365}=0{,}696320$; selisih $0{,}009996$. Aproksimasi memakai
        $1-x\approx e^{-x}$ untuk $x$ sampai $29/365$, dan menganggap kejadian
        pasangan saling bebas padahal tidak.

  \item $A=(A\cap B)\cup(A\cap B^{c})$ dan kedua potong itu saling lepas, sebab
        $B\cap B^{c}=\varnothing$. Aditivitas memberi
        $P(A)=P(A\cap B)+P(A\cap B^{c})=P(A\mid B)P(B)+P(A\mid B^{c})P(B^{c})$,
        memakai definisi peluang bersyarat dua kali.

  \item Peluang menang dengan berpindah $3/4=0{,}75$: pilihan awal menang hanya bila
        mobil memang di situ ($1/4$). Simulasi $10^{4}$ permainan memberi sekitar
        $0{,}75$ dengan ayunan baku $\sqrt{0{,}75\cdot0{,}25/10^{4}}\approx0{,}0043$.

  \item 
Angka benar $0{,}9075$, bukan $0{,}9999$: dengan kelaziman $0{,}1$ persen,
        dari $1000$ orang hanya $1$ sakit, sedangkan positif palsu pada dua uji
        positif $999\cdot0{,}01^{2}\approx0{,}1$ orang. Karena itu
        $0{,}99^{2}\cdot0{,}001/(0{,}99^{2}\cdot0{,}001+0{,}01^{2}\cdot0{,}999)=0{,}9075$.
        Asumsinya: kedua uji bebas bersyarat dan sensitivitas serta spesifisitasnya
        tepat $0{,}99$.

  \item Tidak cukup. Ambil $A$ koin pertama gambar, $B$ koin kedua gambar, $C$ kedua
        koin sama. Tiap pasangan memberi $1/4=P\cdot P$, tetapi
        $P(A\cap B\cap C)=1/4$ sedangkan $P(A)P(B)P(C)=1/8$; selisihnya $1/8$.
\end{enumerate}

\section*{Bab 18: Variabel Acak dan Distribusi}
\begin{enumerate}[label=\textbf{18.\arabic*}]
  \item $c=1-(0{,}1+0{,}4)=0{,}5$; $\mathbb{E}[X]=-0{,}2+0+1{,}5=1{,}3$;
        $\mathbb{E}[X^{2}]=0{,}4+0+4{,}5=4{,}9$; $\operatorname{Var}(X)=4{,}9-1{,}69=3{,}21$.

  \item $(-\infty,b]=(-\infty,a]\cup(a,b]$ dan kedua potong lepas, jadi
        $F(b)=F(a)+P(a<X\le b)$. Untuk diskret $P(X=a)>0$ membuat
        $P(a\le X\le b)$ lebih besar; untuk kontinu $P(X=a)=0$ sehingga keduanya sama.

  \item Benar bahwa $P(X=x)=0$ untuk setiap $x$ pada sebaran kontinu. Kelirunya
        menyimpulkan ``mustahil'': peluang nol tidak berarti himpunan kosong.
        Aditivitas terbilang memberi peluang nol pada gabungan terbilang titik, tetapi
        selang memuat tak terbilang titik dan berpeluang satu.

  \item Turunkan $\sum q^{k}=1/(1-q)$ dua kali: $\sum k(k-1)q^{k-2}=2/(1-q)^{3}$.
        Karena $\sum k^{2}q^{k-1}=(1+q)/(1-q)^{3}$, dengan $q=1-p$ diperoleh
        $\mathbb{E}[X^{2}]=p(1+q)/p^{3}=(1+q)/p^{2}=(2-p)/p^{2}$.
        Maka $\operatorname{Var}=(2-p)/p^{2}-1/p^{2}=(1-p)/p^{2}$; untuk $p=0{,}2$:
        $\mathbb{E}[X^{2}]=45$ dan $\operatorname{Var}=20$.

  \item $\sigma=2$: $P(X>1)=1-\Phi(-1)=0{,}841345$;
        $P(1<X<5)=2\Phi(1)-1=0{,}682689$; kuantil $0{,}2$ ialah
        $3+2(-0{,}841621)=1{,}316758$. Ketiganya cocok dengan
        \texttt{scipy.stats.norm} sampai enam angka.

  \item $\sum_{k=0}^{40}e^{-2{,}5}2{,}5^{k}/k!=1-5{,}39\times10^{-35}$, yaitu $1$
        sampai semua digit tercetak. $P(X\ge8)=0{,}004247$, sedangkan binomial
        $n=1000$, $p=0{,}0025$ memberi $0{,}004191$; binomial lebih kecil karena
        variansinya $2{,}4938<2{,}5$ dan terbatas di atas oleh $n$.

  \item $P(X>s+t\mid X>s)=e^{-\lambda(s+t)}/e^{-\lambda s}=e^{-\lambda t}$.
        Geometrik: $P(X>k)=q^{k}$, jadi nisbahnya $q^{s+t}/q^{s}=q^{t}$. Seragam pada
        bilangan bulat gagal sebab peluang ekornya habis, dan nisbahnya bergantung
        pada $s$.

  \item $f_{Y}(y)=\frac{1}{4\sqrt{y}}$ untuk $0<y<1$ (dua pramasuk) dan
        $\frac{1}{8\sqrt{y}}$ untuk $1\le y<9$ (satu pramasuk). Integralnya
        $\tfrac12+\tfrac12=1$. $\mathbb{E}[Y]=\frac14\int_{-1}^{3}x^{2}\,dx=7/3
        =2{,}3333$, sama dengan $\mathbb{E}[X^{2}]$, dan lewat $f_{Y}$:
        $\int_{0}^{1}\frac{y}{4\sqrt{y}}\,dy+\int_{1}^{9}\frac{y}{8\sqrt{y}}\,dy
        =\tfrac16+\tfrac{13}6=\tfrac73$.

  \item Ambil $X$ seragam pada $\{-1,0,1\}$ dan $Y=|X|$. Maka
        $\mathbb{E}[XY]=\mathbb{E}[X]=0$, jadi $\operatorname{Cov}=0$ dan $\rho=0$.
        Tetapi $P(X=0,Y=0)=1/3$ sedangkan $P(X=0)P(Y=0)=1/9$: tidak bebas.

  \item 
$\operatorname{se}(\bar X)=\sqrt{p(1-p)/(nR)}\le0{,}005$ memberi
        $R\ge0{,}2275/(20\cdot0{,}000025)=455$. Untuk variansi,
        $\operatorname{se}(\bar s^{2})=\sigma^{2}\sqrt{2/((n-1)R)}$ (aproksimasi
        normal) dengan $\sigma^{2}=4{,}55$ dan $n=20$, jadi
        $R\ge2\cdot4{,}55^{2}/(19\cdot0{,}02^{2})=5\,448{,}03$, jadi $R=5\,449$.
        Memakai $2\cdot4{,}55^{2}/0{,}02^{2}$ saja melewatkan faktor $n-1$ dan
        memberi angka sembilan belas kali terlalu besar.
  \end{enumerate}

\section*{Bab 19: Hukum Bilangan Besar dan Teorema Limit Pusat}
\begin{enumerate}[label=\textbf{19.\arabic*}]
  \item $\{|(X_{n}+Y_{n})-(X+Y)|>\varepsilon\}$ termuat dalam
        $\{|X_{n}-X|>\varepsilon/2\}\cup\{|Y_{n}-Y|>\varepsilon/2\}$. Batas gabungan
        memberi peluang paling besar jumlah dua suku, masing-masing menuju $0$.
        $\varepsilon/2$ perlu supaya dua galat yang sama-sama di bawahnya berjumlah
        kurang dari $\varepsilon$.

  \item $\mathbb{E}[X]=0$ dan $\mathbb{E}[X^{2}]=2a^{2}\cdot\sigma^{2}/(2a^{2})=\sigma^{2}$,
        jadi $\operatorname{Var}(X)=\sigma^{2}$. Lalu
        $P(|X|\ge a)=\sigma^{2}/a^{2}$, tepat sama dengan batas Chebyshev: batas itu
        tercapai, jadi tajam.

  \item Salah. Teorema limit pusat menuntut variansi berhingga dan amatan iid; ia juga
        hanya bercerita tentang limit, bukan $n=30$. Pada bagian lab, eksponensial
        dengan $n=30$ masih memberi $D=0{,}0219$ dengan $p=0{,}000$, dan Cauchy
        memberi median $|\bar X_{n}|$ yang menetap sekitar $1$.

  \item Chebyshev: $n\ge\sigma^{2}/(\varepsilon^{2}\delta)=0{,}25/(0{,}0025\cdot0{,}05)
        =2000$. Hitungan binomial eksak memberi $n=371$ (ekor $0{,}0484$). Batas
        Chebyshev lima kali lebih longgar.

  \item $\mathbb{E}[X_{n}]=n^{2}\cdot(1/n)=n\to\infty$, sedangkan untuk
        $\varepsilon$ tetap $P(|X_{n}|>\varepsilon)=1/n\to0$. Jadi $X_{n}\to0$ dalam
        peluang tetapi harapannya meledak: tidak ada dominasi yang boleh dipakai.

  \item Eksak: $0{,}069545$. Tanpa koreksi, $z=(50-60)/\sqrt{42}=-1{,}5430$ memberi
        $0{,}061411$, galat $0{,}008134$. Dengan koreksi $50{,}5$:
        $0{,}071340$, galat $0{,}001795$. Koreksi kontinuitas memperbaiki empat kali.

  \item Median $|\bar X_{n}|$: $1{,}004$; $1{,}008$; $1{,}040$; $0{,}973$ untuk
        $n=10$ sampai $10^{4}$. $P(|\bar X_{n}|>10)$: $0{,}0630$; $0{,}0680$;
        $0{,}0653$; $0{,}0580$, teori $0{,}0635$. Simpangan baku $33{,}0$; $863{,}6$;
        $23{,}9$; $379{,}5$ tak mengecil sebab variansi Cauchy tak hingga.

  \item $\varphi_{\bar X_{n}}(t)=\bigl[\varphi(t/n)\bigr]^{n}
        =\bigl[e^{-|t|/n}\bigr]^{n}=e^{-|t|}$, jadi rata-ratanya Cauchy lagi untuk
        setiap $n$. Karena itu se selamanya sebesar skala data itu sendiri, dan
        ``rata-rata $\pm$ dua se'' tidak menyempit walau $n$ dinaikkan.

  \item $n=(1{,}96\cdot0{,}491971/0{,}0005)^{2}=3\,719\,214$; untuk setengah lebar
        $0{,}001$ cukup $929\,803$. Rasionya $4$, sebab $n\propto h^{-2}$: lebar
        dipotong dua, tarikan naik empat kali.

  \item Untuk $t<c$ dan $\varepsilon=c-t$:
        $P(X_{n}\le t)\le P(X_{n}\le c-\varepsilon)+P(|X_{n}-c|>\varepsilon)\to0$;
        untuk $t>c$ pakai komplemennya. Syarat konstanta tak dapat dibuang:
        ambil $X$ Rademacher dan $X_{n}=(-1)^{n}X$. Setiap $X_{n}$ tetap
        Rademacher (sebab $X$ simetris), jadi $X_{n}\rightsquigarrow X$; tetapi
        untuk $n$ ganjil $|X_{n}-X|=2$ dengan peluang $1$, sehingga $X_{n}$ tidak
        menuju $X$ dalam peluang.
\end{enumerate}

\section*{Bab 20: Estimasi, Interval Kepercayaan, dan Uji Hipotesis}
\begin{enumerate}[label=\textbf{20.\arabic*}]
  \item Salah. Nilai-p $0{,}42$ berarti data tidak cukup membantah $p=0{,}5$; ia tidak
        membuktikan $p=0{,}5$. Perhitungan penambah: selang $95$ persen, misalnya
        $(0{,}28;\,0{,}57)$, dan daya uji pada selisih yang dianggap penting.

  \item $\mathbb{E}[\hat p(1-\hat p)]=p-\bigl(p^{2}+p(1-p)/n\bigr)
        =p(1-p)(n-1)/n$. Penaksir tak biasnya
        $\frac{n}{n-1}\hat p(1-\hat p)$. Untuk $p=0{,}3$, $n=20$: harapannya
        $0{,}21\cdot19/20=0{,}1995$, dan penaksir itu mengembalikannya ke $0{,}21$.

  \item Tulis $\hat\theta_{n}-\theta=\operatorname{bias}_{n}
        +(\hat\theta_{n}-\mathbb{E}\hat\theta_{n})$. Suku pertama menuju $0$; suku
        kedua menuju $0$ dalam peluang menurut Chebyshev, sebab variansinya menuju
        $0$. MSE $=\operatorname{bias}^{2}+\operatorname{Var}$ menuju $0$ lebih kuat,
        tetapi yang dipakai hanya kedua syarat tadi.

  \item 
$7/40$: Wilson $(0{,}0875;\,0{,}3195)$, lebar $0{,}2320$; Wald
        $(0{,}0572;\,0{,}2928)$, lebar $0{,}2355$. $13/20$: Wilson lebar $0{,}3860$,
        Wald $0{,}4181$. Lebih sempitnya $7/40$ bukan tanda ketelitian lebih tinggi. Sebagian
        karena $n=40$; sebagian karena $\hat p=0{,}175$ jauh dari $0{,}5$
        mengecilkan $\hat p(1-\hat p)$: pada $n=40$ dengan $\hat p=0{,}5$ lebar
        Wilson $0{,}2960$, masih lebih lebar daripada $0{,}2320$. (Pada $n$ sama,
        $7/20$ dan $13/20$ sama-sama memberi $0{,}3860$, sebab $\hat p(1-\hat p)$
        simetris terhadap $1/2$.) Yang menentukan ketelitian bukan lebar melainkan
        cakupan: cakupan Wald pada $p=0{,}175$, $n=40$ hanya $0{,}921$, sedangkan
        Wilson $0{,}943$. Lebar kecil di dekat tepi dibayar dengan cakupan di
        bawah $95$ persen.

  \item $|k-15|>1{,}96\sqrt{7{,}5}=5{,}3676$, jadi tolak bila $k\le9$ atau $k\ge21$.
        Ukuran eksaknya $0{,}042774$; dayanya pada $p=0{,}7$ hanya $0{,}588816$.
        Dengan $n=100$ (tolak bila $k\le40$ atau $k\ge60$) daya naik ke $0{,}987502$.

  \item Wald: $|\hat p-p_{0}|\le z\sqrt{\hat p(1-\hat p)/n}$, yang setara dengan
        $p_{0}$ termuat dalam $\hat p\pm z\sqrt{\hat p(1-\hat p)/n}$. Uji yang memakai
        $\sqrt{p_{0}(1-p_{0})/n}$ memakai penyebut lain, sehingga himpunan $p_{0}$
        yang tidak ditolak bukan selang Wald itu.

  \item 
Misalkan selisih teramati $0{,}87$ jam dengan galat baku $0{,}86$ jam.
``Dengan taraf $5$ persen, data tidak menolak kesamaan kedua metode
        ($p=0{,}31$). Selang $95$ persen untuk selisihnya $(-0{,}82;\,2{,}56)$ jam,
        jadi selisih $2$ jam masih termasuk nilai yang wajar dan selisih $0{,}5$ jam
        pun belum dapat disingkirkan. Daya uji ini untuk selisih $0{,}5$ jam hanya
        $0{,}09$: ketiadaan bukti perbedaan bukan bukti kesamaan, dan yang jujur
        disimpulkan adalah data ini belum memutuskan.''

  \item Aturan tiga: $3/100=0{,}03$. Batas eksak satu sisi
        $1-0{,}05^{1/100}=0{,}029513$; batas atas Wilson $0{,}036993$. Ketiganya
        berbeda karena aturan tiga hanya pendekatan dan Wilson dua sisi. Untuk
        kalimat ``paling tinggi $5$ persen'' yang sah pada taraf $95$ persen:
        batas eksak $0{,}0295$.

  \item $\bar x=19{,}125$, $s=1{,}082191$, $\operatorname{se}=0{,}312401$;
        $t_{11;0{,}95}=1{,}795885$ memberi $(18{,}5640;\,19{,}6860)$ dengan lebar
        $1{,}1221$, sedangkan $95$ persen memberi lebar $1{,}3752$. Lebar berbanding
        $1/\sqrt{n}$, jadi separuhnya perlu $n=48$.
\end{enumerate}

\section*{Bab 21: Regresi dan Statistika Terapan}
\begin{enumerate}[label=\textbf{21.\arabic*}]
  \item Belum cukup bukti. $R^{2}=0{,}99$ hanya mengatakan ragam terjelaskan;
        $s$ yang menentukan galat. Pada data kafe $R^{2}=0{,}992821$ tetapi
        $s=0{,}365148$ jam $=21{,}9$ menit, dan setengah lebar selang prediksi $95$
        persen di $x_{0}=5$ ialah $0{,}9101$ jam, yaitu $55$ menit. Yang perlu
        disebut: $s$ dan selang prediksi.

  \item Persamaan normal $X^{\top}e=0$. Kolom pertama $X$ berisi satu, jadi
        $\sum e_{i}=0$; kolom kedua berisi $x_{i}$, jadi $\sum x_{i}e_{i}=0$.
        Tanpa kolom satu (garis lewat titik asal), $\sum e_{i}=0$ tak lagi dipaksakan
        dan umumnya gagal.

  \item $\hat b-\beta=\sum_{i}(x_{i}-\bar x)\varepsilon_{i}/S_{xx}$. Suku-sukunya
        saling bebas dan $\operatorname{Var}(\varepsilon_{i})=\sigma^{2}$, jadi
        $\operatorname{Var}(\hat b)=\sigma^{2}\sum(x_{i}-\bar x)^{2}/S_{xx}^{2}
        =\sigma^{2}/S_{xx}$. Tak ada rumus matriks yang dipakai.

  \item $\hat b=0{,}7$, $\hat a=1{,}3$, $S_{xx}=10$, $s=0{,}516398$,
        $\operatorname{se}(\hat b)=0{,}163299$; dengan $t_{3}=3{,}182446$ selangnya
        $(0{,}1803;\,1{,}2197)$. Lebih lebar daripada data kafe karena $n=5$,
        $S_{xx}$ kecil, dan $s$ besar.

  \item 
Pada data kafe $\bar x=5$, $n=9$, $S_{xx}=60$, jadi pengungkitnya
$c(x_{0})=1/9+(x_{0}-5)^{2}/60$: $c(2)=0{,}261111$ dan $c(9)=0{,}377778$.
Dengan $2{,}364624\sqrt{0{,}133333\,(1+c)}$ setengah lebarnya $0{,}9696$ untuk
$x_{0}=2$ dan $1{,}0135$ untuk $x_{0}=9$; lebarnya berbeda karena
$(x_{0}-\bar x)^{2}/S_{xx}$ lebih besar di titik yang lebih jauh dari pusat.
Suku itu mengecil sebanding $1/n$, sehingga selisih kedua selang memang
menghilang untuk $n$ besar: keduanya menuju $1{,}96\sigma=0{,}716$ setengah
lebar. Yang tidak pernah hilang adalah lebar itu sendiri, sebab suku $+1$
(galat satu pengamatan baru) tetap ada: selang prediksi tidak menyempit ke
nol, ia menyempit ke selebar sebaran galatnya.

  \item Ruang kolom model besar memuat ruang kolom model kecil, jadi proyeksinya paling
        dekat: SSE tak dapat naik. $R^{2}$ boleh tetap, tetapi pembagi $n-p$ mengecil,
        sehingga $R^{2}_{\text{adj}}$ turun: dari $0{,}991795$ ($p=2$) ke
        $0{,}990428$ ($p=3$).

  \item $\mathrm{VIF}=1/(1-0{,}81)=5{,}263158$, jadi
        $\operatorname{se}(\hat b_{1})$ membesar $\sqrt{5{,}263158}$ kali. Dengan
        $x_{2}$: $\sqrt{2/(50\cdot0{,}19)}=0{,}458831$. Tanpa $x_{2}$:
        $\sqrt{2/50}=0{,}2$. Rasionya $2{,}294157$.

  \item 
Kelompok A: $(1;2)$, $(2;1{,}8)$, $(3;1{,}6)$; kelompok B: $(5;4{,}0)$,
        $(6;3{,}8)$, $(7;3{,}6)$. Kemiringan tiap kelompok $-0{,}2$, kemiringan
        gabungan $+0{,}4$. Dekomposisi yang sah menjumlahkan kontribusi, bukan
        kemiringan:
        $b_{\text{gab}}=(S_{xy}^{\text{dalam}}+S_{xy}^{\text{antar}})/S_{xx}=-0{,}8/28+12/28=-0{,}028571+0{,}428571=0{,}4$,
        sedangkan kemiringan antarkelompok sendiri $12/24=0{,}5$.

  \item ``Dari dua belas model yang dicoba, yang terbaik memberi $p=0{,}03$. Setelah
        koreksi untuk dua belas perbandingan (Bonferroni: $0{,}36$) efeknya tidak
        lagi berarti; karena itu data ini belum memutuskan.''
\end{enumerate}

\section*{Bab 22: Aritmetika Floating-Point dan Analisis Galat}
\begin{enumerate}[label=\textbf{22.\arabic*}]
  \item $\operatorname{ulp}(1)=2^{-52}=2{,}220\times10^{-16}$;
        $\operatorname{ulp}(1024)=2^{-42}=2{,}274\times10^{-13}$;
        $\operatorname{ulp}(0{,}1)=2^{-56}=1{,}388\times10^{-17}$;
        $\operatorname{ulp}(3\times10^{8})=2^{-24}=5{,}960\times10^{-8}$.
        $|x|u$ berbeda dari ulp paling banyak faktor dua, sebab ulp tetap di dalam satu
        oktad. Pada $[1,2)$ ada $2^{52}$ bilangan.

  \item $2^{-53}$ tepat setengah ulp dari $1$, dan pembulatan ke genap
        mengembalikannya ke $1$: $1+2^{-53}=1$. Karena $2^{-52}$ satu ulp penuh,
        $1+2^{-52}>1$. Pada binary32, $u=2^{-24}$ dan ulp$(1)=2^{-23}$, jadi bilangan
        positif terkecilnya $2^{-23}=1{,}192\times10^{-7}$.

  \item $\operatorname{fl}(x)$ adalah bilangan terdekat, jadi
        $|\operatorname{fl}(x)-x|\le\tfrac12\operatorname{ulp}(x)$ menurut definisi
        pembulatan. Kesamaan tak dapat berlaku untuk semua $x$: bila $x$ sendiri
        terwakili, galatnya $0$, dan $\tfrac12$ ulp hanya tercapai di titik tengah
        dua bilangan.

  \item Salah. Setiap operasi dibulatkan benar, tetapi galatnya menumpuk sepanjang
        algoritme. Contoh dari bab ini: penjumlahan berurutan $1+10^{6}\cdot10^{-16}$
        tetap $1$, sedangkan Kahan memberi $1{,}0000000001$; selisih $10^{-10}$
        jauh di atas $u$.

  \item Dengan penyebut $2^{54}$: $(0{,}1+0{,}2)+0{,}3=10808639105689192/2^{54}$ dan
        $0{,}1+(0{,}2+0{,}3)=10808639105689190/2^{54}$. Selisihnya $2/2^{54}=2^{-53}$,
        tepat satu ulp. \texttt{fsum} memberi $0{,}6$; \texttt{float.hex} berbeda di
        digit terakhir.

  \item Rumus lugu memberi $0{,}4996003611$; identitas sudut ganda memberi
        $0{,}4999999999999996$, dan $0{,}5-x^{2}/24+x^{4}/720$ memberi angka yang sama.
        Galat relatifnya $-7{,}99\times10^{-4}$, berasal dari pembatalan $1-\cos x$:
        $\cos x$ membuang digit yang dibutuhkan selisihnya.

  \item Lugu: akar kecil $-1{,}0000000111\times10^{-4}$. Bentuk stabil
        $-2c/\bigl(b+\sqrt{b^{2}-4ac}\bigr)$ memberi
        $-1{,}0000000100\times10^{-4}$, galat relatif $1{,}12\times10^{-9}$.
        Bilangan kondisi pengurangannya $\bigl(2\cdot10^{4}\bigr)/\bigl(2\cdot
        10^{-4}\bigr)=10^{8}$, jadi sekitar delapan digit hilang.

  \item Pada tingkat $j$ ada $2^{k-j}$ penjumlahan, masing-masing bergalat paling besar
        $u$ kali jumlah nilai mutlak masukannya; dijumlahkan seluruh pohon, galatnya
        $\le ku\sum|x_{i}|$ dengan $k=\log_{2}n$. Kedalaman pohon itulah yang
        menggantikan $n$ pada batas penjumlahan berurutan.

  \item Berurutan: setiap $10^{-16}$ di bawah $\tfrac12\operatorname{ulp}(1)
        =1{,}11\times10^{-16}$, jadi hasilnya tetap $1$. Kahan:
        $1{,}0000000001$, sama dengan nilai eksaknya. Setiap langkah Kahan juga
        dibulatkan, tetapi sisa yang hilang disimpan di $c$ dan dikembalikan pada
        langkah berikutnya.

  \item $\kappa\approx1+\bar x^{2}/s^{2}$. Untuk $1000$ nilai
        $999+\bigl((k\bmod7)-3\bigr)/4$ dengan $k=1,\dots,1000$:
        $\bar x=999{,}00075$, $s=0{,}4997$ (simpangan baku populasi, sesuai
        pembagi $1/n$ pada soal),
        $\kappa\approx4{,}0\times10^{6}$. Dengan $u_{32}=5{,}96\times10^{-8}$,
        $\kappa u\approx0{,}238$: hampir semua digit hilang. Terukur: satu lintasan
        galat $25$ persen, dua lintasan $4{,}4\times10^{-8}$.
\end{enumerate}

\section*{Bab 23: Optimasi Konveks dan Dualitas}
\begin{enumerate}[label=\textbf{23.\arabic*}]
  \item $f$: Hessian $\begin{pmatrix}2&1\\1&2\end{pmatrix}$, nilai eigen $1$ dan $3$,
        definit positif, konveks ketat. $g$: $\begin{pmatrix}0&1\\1&0\end{pmatrix}$,
        nilai eigen $\pm1$, indefinit, tidak keduanya. $h$:
        $\begin{pmatrix}4&-4\\-4&6\end{pmatrix}$, nilai eigen $0{,}8769$ dan $9{,}1231$,
        definit positif. Tak satu pun semidefinit singular; pada $x_{1}^{2}+2x_{1}x_{2}
        +x_{2}^{2}$ (nilai eigen $0$ dan $4$) minimumnya tak tunggal.

  \item Untuk $x,y\in C_{1}\cap C_{2}$ dan $t\in[0,1]$, titik campurannya ada di
        $C_{1}$ dan di $C_{2}$, jadi di irisannya. Gabungan: $[0,1]\cup[2,3]$ memuat
        $1$ dan $2$ tetapi tidak $(1+2)/2=1{,}5$.

  \item $f_{i}(tx+(1-t)y)\le tf_{i}(x)+(1-t)f_{i}(y)\le tM(x)+(1-t)M(y)$ untuk tiap
        $i$, lalu ambil maksimum. Untuk $\max\{x_{1}^{2},x_{2}\}$ di $(0,0)$ gradien
        tak ada (patahan sepanjang $x_{2}=x_{1}^{2}$), tetapi subdiferensialnya
        $\operatorname{conv}\{(0,0),(0,1)\}$ memuat $\mathbf{0}$: minimum. Di $(1,1)$
        subdiferensialnya $\operatorname{conv}\{(2,0),(0,1)\}$ tanpa nol, bukan minimum.

  \item Tanpa $x_{1}\ge1$: $x_{1}=0{,}8$, $x_{2}=1{,}6$, melanggar. Jadi $x_{1}=1$,
        $x_{2}=1{,}5$, $p^{*}=3{,}25$. Pengalinya $\nu=-1{,}5$ bagi kesamaan dan
        $\lambda=0{,}5\ge0$ bagi pertidaksamaan. \texttt{scipy.optimize.mini\-mize}
        memberi nilai yang sama.

  \item $g(\lambda)=\min_{x}\bigl[x_{1}^{2}+x_{2}^{2}+\lambda(2-x_{1}-x_{2})\bigr]
        =2\lambda-\lambda^{2}/2$. Turunannya $2-\lambda$, jadi maksimum di $\lambda=2$
        dengan $d^{*}=2=p^{*}$. Celahnya nol karena kendalanya affine dan
        $x=(1{,}5;\,1{,}5)$ layak ketat: Slater terpenuhi.

  \item Salah. Dualitas kuat hanya menyamakan nilai, bukan menjamin pencapaian.
        Contoh: minimumkan $e^{x}$ atas $\mathbb{R}$. $p^{*}=0$ tetapi tak tercapai;
        sisi dualnya juga $0$, jadi celahnya nol sementara optimum primal tidak ada.

  \item Salah. Kelengkapan komplementer menyatakan $\lambda_{i}g_{i}(\mathbf{x}^{*})=0$:
        kendala yang aktif boleh berpengali nol, dan kendala tak aktif wajib
        berpengali nol. Implikasinya satu arah, jadi ``aktif $\Rightarrow$ positif''
        tidak sah.

  \item Tanpa pembatas, $(x_{1}-2)^{2}+(x_{2}-2)^{2}$ pada $x_{1}+x_{2}=2$ diminimumkan
        di $x_{1}=x_{2}=1$. Itu memenuhi $x_{1}\le c$ bila dan hanya bila $c\ge1$, jadi
        pada $c=1$ pengalinya $\lambda=0$ dan $p^{*}=2$; kendala baru menekan untuk
        $c<1$ (pada $c=0{,}5$: $p^{*}=2{,}5$, $\lambda=2$).

  \item Untuk $\lambda>0$ dan $x_{k}\to-\infty$, suku $\lambda x_{k}^{3}$ mendominasi
        dan bernilai negatif, jadi $g(\lambda)=-\infty$. Hanya $\lambda=0$ berhingga:
        $g(0)=\min(x-2)^{2}=0$, maka $d^{*}=0$. Sedangkan $p^{*}=1$ di $x=1$ (potongan
        $[0,1]$), sebab di $x=-1$ nilainya $9$. Celahnya $1$, seratus persen.

  \item Ketaksamaan orde pertama fungsi konveks memberi
        $f(y)\ge f(x^{*})+\nabla f(x^{*})^{\top}(y-x^{*})=f(x^{*})$ untuk setiap $y$,
        tanpa memakai ketatnya konveksitas. Bukti itu gagal pada $f(x)=x^{3}$ sebab
        $f$ tidak konveks; syarat yang dilanggar justru konveksitas $f$, di $x=0$
        dengan $\nabla f=0$ tetapi bukan minimum.
\end{enumerate}

% Kunci jawaban Bab 24--32. Berkas ini di-\input langsung di dalam document.
% Setiap angka diuji oleh skrip assert di luar repo (lihat laporan).

\section*{Bab 24: Pemrograman Linear dan Integer}
\begin{enumerate}[label=\textbf{24.\arabic*}]
  \item Maksimumkan $30x_{1}+20x_{2}$ dengan $2x_{1}+x_{2}\le10$, $x_{1}+2x_{2}\le8$,
        $x\ge0$. Verteks $(0;0)$, $(5;0)$, $(0;4)$, $(4;2)$ memberi laba
        $0;\,150;\,80;\,160$. Optimum bersambung $(4;2)$ berlaba $160$; enumerasi bulat
        juga memberi $160$, jadi di sini pembulatan ke bawah kebetulan tepat dan
        celahnya $0$.

  \item Untuk $x,y\in P$ dan $\lambda\in[0,1]$:
        $A(\lambda x+(1-\lambda)y)=\lambda Ax+(1-\lambda)Ay\le\lambda b+(1-\lambda)b=b$,
        memakai $A$ linear dan $\lambda\ge0$. Pada himpunan bulat argumen itu runtuh:
        gabungan konveks dua titik bulat tidak bulat, misalnya $0$ dan $1$ layak
        sementara titik tengahnya bukan anggota $\mathbb{Z}^{n}$.

  \item Laba $x_{1}+2x_{2}$ di verteks: $0$; $4$; $6$; $6$. Optimum $(0;3)$ dan
        $(3;1{,}5)$ sama-sama $6$, dan seluruh ruas $x_{1}+2x_{2}=6$ optimal sebab
        sasarannya konstan di sana; titik tengah $(1{,}5;\,2{,}25)$ bukan verteks.
        Dualnya $y=(0;1)$: harga jam $1$, kayu lapis $0$.

  \item Tabel awal basis $s_{1},s_{2},s_{3}=(4;3;3)$. Pivot pertama: $x_{1}$ masuk
        (biaya tereduksi $3$), nisbah $4$ dan $3$, jadi $s_{2}$ keluar; nilai $9$.
        Pivot kedua: $x_{2}$ masuk, $s_{1}$ keluar (nisbah $1$), optimum $(3;1)$
        dengan $z=11$. Harga dual dari tabel akhir: $(2;1;0)$, dan
        $4(2)+3(1)+3(0)=11$.

  \item Dualnya meminimumkan $10y_{1}+8y_{2}$ dengan $2y_{1}+y_{2}\ge30$,
        $y_{1}+2y_{2}\ge20$, $y\ge0$; penyelesaian $y=(40/3;\,10/3)$ bernilai
        $160$, sama dengan primal. Artinya: satu kg tepung tambahan bernilai
        $13{,}333$ dan satu jam oven tambahan $3{,}333$.

  \item Salah. Pembulatan ke bawah tidak selalu layak begitu ada kendala
        kesamaan atau bertanda $\ge$. Contoh: maksimumkan $x_{1}+x_{2}$ dengan
        $2x_{1}+2x_{2}=5$; jawaban bersambung $(2{,}5;\,0)$ dibulatkan menjadi
        $(2;0)$, yang memberi $4\ne5$ dan karena itu tidak layak.

  \item Benar untuk arah majunya, dan hanya pada pasangan optimal. Kelengkapan
        komplementer $y_{i}(b_{i}-\mathbf{a}_{i}^{\top}\mathbf{x})=0$ memberi:
        $y_{i}>0$ memaksa kendala itu aktif. Sebaliknya tidak berlaku; pada titik
        merosot ada kendala aktif dengan $y_{i}=0$, jadi harga nol tidak berarti
        kendala itu longgar.

  \item Kerapatan $2{,}0$; $1{,}75$; $1{,}667$; $1{,}5$. Barang pertama penuh
        ($5$ kg, $10$), lalu $3/4$ barang kedua, total berat $8$: batas LP
        $15{,}25$. Enumerasi memberi $15$ (barang $1$ dan $3$). Celah $0{,}25$,
        yaitu $1{,}639$ persen dari batas.

  \item 
Kapasitas $7$: barang pertama penuh, lalu $1/2$ barang kedua, batas LP
        $13{,}5$; jawaban bulat terbaik $13$ (barang $1$ dan $4$, berat $7$). Celah
        $0{,}5$, yaitu $3{,}704$ persen. Celah berubah karena kapasitas $7$ membuang
        kombinasi $(1,3)$ yang beratnya $5+3=8$ dan tadinya muat, sehingga jawaban
        bulat terbaik pindah ke barang $1$ dan $4$; pecahan optimum dan bulat
        optimum bergerak tidak sebanding.

  \item $x_{v}=1/2$ untuk semua $v$ layak ($x_{u}+x_{v}=1$) dan memberi $n/2$;
        setiap penutup harus memuat $n-1$ simpul, sebab meninggalkan dua simpul
        berarti rusuk di antaranya tak tertutup. Nisbah LP terhadap bulat:
        $0{,}6667$; $0{,}625$; $0{,}5556$ untuk $n=4,5,10$. Nisbahnya menuju
        $1/2$, jadi batas relaksasi bisa meleset faktor dua dan pembulatan bukan
        cara umum.
\end{enumerate}

\section*{Bab 25: Jaringan Saraf dan Turunan Balik}
\begin{enumerate}[label=\textbf{25.\arabic*}]
  \item Batasnya $\lfloor1/\gamma^{2}\rfloor$: untuk $\gamma=0{,}2$ dapat $25$,
        untuk $\gamma=0{,}05$ dapat $400$. Batas itu hanya menghitung kesalahan,
        dan tiap kesalahan memakan margin, karena itu $n$ tidak muncul. Untuk
        $\gamma$ kecil batasnya meledak seperti $1/\gamma^{2}$ sehingga tak berguna.

  \item Tulis $\mathbf{W}=\mathbf{W}_{2}\mathbf{W}_{1}$ dan
        $\mathbf{b}=\mathbf{W}_{2}\mathbf{b}_{1}+\mathbf{b}_{2}$; maka keluarannya
        $\mathbf{W}\mathbf{x}+\mathbf{b}$, yaitu satu ambang linear. Menambah lapisan
        identitas berikutnya hanya mengalikan matriks, jadi kapasitasnya tidak
        bertambah sedikit pun.

  \item Dengan $\mathbf{x}=(0{,}5;-1;\,0{,}25;2)$ dan $y=0$: $\mathbf{z}_{1}=(2{,}9125;\,
        0{,}375)$, $a=0{,}8961350845$, $\bar{z}_{2}=0{,}8961350845$, sehingga
        $\bar{\mathbf{z}}_{1}=(0{,}7169081;\,-0{,}5376810)$ dan
        \[
          \nabla_{\mathbf{W}_{1}}L=
          \begin{pmatrix}
            0{,}3584540 & -0{,}7169081 & 0{,}1792270 & 1{,}4338161\\
            -0{,}2688405 & 0{,}5376810 & -0{,}1344203 & -1{,}0753621
          \end{pmatrix},
        \]
        yaitu baris kedua sama dengan $-0{,}75$ kali baris pertama. Butir ini
        ambigu: kedua gerbang terbuka, jadi tidak ada baris bergradien nol.

  \item Belum cukup bukti. Selisih hingga hanya memeriksa tiga belas skalar di satu
        titik, dan pada $h=10^{-6}$ galat $10^{-10}$ memang yang diharapkan; bug
        bersama pada kedua hitungan tidak akan terlihat. Perlu: beberapa titik,
        kedua keadaan gerbang, kasus $y=0$ dan $y=1$, gradien terhadap masukan, dan
        pembanding turunan otomatis.

  \item Salah. Teorema itu tentang kepadatan fungsi kontinu pada domain
        kompak: ada parameter yang menghampiri, bukan parameter itu terjangkau
        optimasi. Dua hal yang tidak dinyatakan: banyaknya unit yang diperlukan, dan
        galat pada data uji (generalisasi), bukan hanya pada fungsi sasaran.

  \item Pada $t=1$, $\hat m=g$ dan $\hat v=g^{2}$, sehingga langkahnya
        $\eta g/(|g|+\epsilon)=-0{,}01$ dan $b_{2}$ menjadi $0{,}06$. Dengan gradien
        $100$ kali lebih besar langkahnya tetap $-0{,}01$, sebab Adam hanya
        memakai nisbah terhadap skala gradien; $b_{2}$ tetap $0{,}06$.

  \item 
Keluaran $w_{2}\max(\mathbf{W}_{1}\mathbf{x}+\mathbf{b}_{1},0)$ adalah
        jumlah dua fungsi linear yang masing-masing aktif di satu separuh bidang,
        jadi linear sepotong-sepotong; dua garis $z_{1}[i]=0$ membagi bidang menjadi
        paling banyak empat sel. Untuk $k$ unit dalam susunan umum, $k$ garis
        membagi bidang menjadi $1+k+k(k-1)/2$ daerah, yaitu $2,4,7$ untuk $k=1,2,3$
        (bukan $8$: garis ketiga hanya menambah tiga daerah, bukan empat).

  \item $\nabla_{\mathbf{x}}L=(-0{,}0868370;\,0{,}1328096;\,-0{,}0970532;\,
        -0{,}1430257)$, maka $\boldsymbol{\delta}=0{,}05\,\mathrm{sign}(\nabla)
        =(-0{,}05;0{,}05;-0{,}05;-0{,}05)$. Kerugian di $\mathbf{x}+\boldsymbol{\delta}$
        ialah $0{,}3186800$, sedangkan ramalan linear
        $0{,}2949123+0{,}0229863=0{,}3178985$; selisih $7{,}81\times10^{-4}$ dari
        kelengkungan.

  \item 
Ramalan: blok lab kedua memakai ReLU, dan ReLU hanya jenuh di sisi negatif.
        Tabel saturasi ($8$ unit, $\mathbf{W}_{1}$ He acak, $\mathbf{b}_{2}=0$):
        \begin{center}
        \begin{tabular}{lrrrr}
        & rentang $Z_{1}$ (awal) & unit mati & $L$ akhir & akurasi \\ \hline
        $\mathbf{b}_{1}=0$   & $[-2{,}43;\,2{,}40]$   & $0/8$ & $0{,}0163$ & $1{,}0000$ \\
        $\mathbf{b}_{1}=+10$ & $[7{,}57;\,12{,}41]$   & $0/8$ & $0{,}3225$ & $0{,}8375$ \\
        $\mathbf{b}_{1}=-10$ & $[-12{,}43;\,-7{,}60]$ & $8/8$ & $0{,}6909$ & $0{,}5333$
        \end{tabular}
        \end{center}
        Jadi $\mathbf{b}_{1}=+10$ memang tidak mematikan unit tersembunyi
        ($Z_{1}$ jauh di atas nol diteruskan apa adanya), tetapi tidak gratis:
        $H\cdot\mathbf{w}_{2}$ menjadi sangat besar, sigmoid keluaran menempel di
        $0$ atau $1$, dan $L$ mulai dari $4{,}04$; sesudah $3000$ langkah ia
        berhenti di $L=0{,}3225$ dengan akurasi $0{,}8375$, lebih buruk daripada
        $\mathbf{b}_{1}=0$. Untuk $\mathbf{b}_{1}=-10$ seluruh unit mati sejak awal
        ($Z_{1}$ di bawah nol, turunan tepat $0$, bukan sekadar kecil),
        $\partial L/\partial b_{1}=\mathbf{0}$ tepat, dan latihan berhenti di
        $L=0{,}6909$ -- entropi label -- dengan akurasi kelas mayoritas $0{,}5333$.
  \end{enumerate}

\section*{Bab 26: Generalisasi, Regularisasi, dan Pelatihan}
\begin{enumerate}[label=\textbf{26.\arabic*}]
  \item Belum cukup bukti, dan yang keliru ialah menyamakan galat latih dengan
        galat generalisasi. Galat latih nol hanya berarti data latih terhafal. Angka
        tambahan yang diperlukan: galat validasi (atau uji) pada data yang tidak
        dipakai melatih, beserta lebar selangnya.

  \item Dengan $\hat f=\hat f_{D}$ dan $\bar f=\mathbb{E}_{D}\hat f$:
        $\mathbb{E}(y-\hat f)^{2}=\sigma^{2}+(\bar f-f)^{2}+\operatorname{Var}(\hat f)$.
        Suku silang lenyap karena $\mathbb{E}[\varepsilon]=0$ dan $\varepsilon$
        bebas dari $D$, sehingga $\mathbb{E}[\varepsilon(\hat f-f)]=0$; kebebasan
        itu dipakai tepat di langkah ini, yaitu pada $\varepsilon$, bukan pada
        $\hat f$. Suku silang kedua lenyap dengan alasan berbeda:
        $\mathbb{E}[\hat f-\bar f]=0$ menurut definisi $\bar f$, jadi
        $2(f-\bar f)\,\mathbb{E}[\hat f-\bar f]=0$ tanpa memerlukan kebebasan apa pun.

  \item Dari $\frac{d}{dw}\bigl[\frac12\sum(y_{i}-wx_{i})^{2}\bigr]=0$ dapat
        $w(\lambda)=\sum x_{i}y_{i}/(\sum x_{i}^{2}+\lambda)=64/(30+\lambda)$.
        Turunannya $-64/(30+\lambda)^{2}<0$ tetapi tak pernah nol, jadi $w$ mengecil
        tanpa menyentuh nol. Pada L1, subgradien di $w=0$ memuat nol bila
        $\lambda\ge|\sum x_{i}y_{i}|=64$; jadi $w=0$ tepat pada $\lambda=64$.

  \item $\operatorname{MSE}(c)=(c-1)^{2}\mu^{2}+c^{2}\sigma^{2}/n$: untuk
        $c=0{,}95$ dapat $0{,}01+0{,}0361=0{,}0461$; untuk $c=1{,}05$ dapat
        $0{,}01+0{,}0441=0{,}0541$; untuk $c=1$ dapat $0{,}04$. Simpangan sama besar
        tidak sama rugi karena suku variansi tumbuh seperti $c^{2}$, sehingga
        melebihi jatah lebih mahal daripada menyusut.

  \item Perbaikan: angka validasi silang itu taksiran galat prosedur pemilihan,
        bukan galat model terpilih di lapangan, dan ia bias optimistis karena
        $\lambda$ dipilih pada lipatan yang sama. Perlu dua hal: validasi silang
        bersarang (atau kumpulan uji yang tak pernah dilihat saat memilih), dan
        pengulangan dengan beberapa partisi beserta selangnya.

  \item Untuk $a<0$: dengan $u<0$ turunannya $u-a-t$, nol di $u=a+t$, yang sah
        hanya bila $a<-t$. Untuk $u=0$ subgradiennya $[-a-t,\,-a+t]$, memuat nol
        bila $|a|\le t$, jadi $u=0$ optimal. Simetris karena sasaran
        $\frac12(u-a)^{2}+t|u|$ tidak berubah bila $u$ dan $a$ dua-duanya diganti
        tanda.

  \item Waktu paruhnya $\ln2/(\eta\lambda)=0{,}693147/0{,}002=346{,}57$ langkah;
        seperempat nilai awal berarti dua paruh, $693{,}15$ langkah, dan
        $(1-\eta\lambda)^{693}=0{,}24973$. Simulasi numpy dengan gradien nol
        memberi angka yang sama, sebab $w_{k}=w_{0}(1-\eta\lambda)^{k}$.

  \item Jurangnya $0{,}2406-0{,}1500=0{,}0906$; pada $\lambda=0$ jurangnya
        $0{,}3262-0=0{,}3262$, pada $\lambda=10$ hanya $0{,}3395-0{,}3333=0{,}0062$.
        Jurang kecil tidak cukup: pada $\lambda=10$ jurangnya terkecil justru karena
        keduanya sama-sama buruk, dekat galat kelas mayoritas $0{,}3419$.

  \item Dengan $m$ hipotesis, $P(\exists i:|\hat R_{i}-R_{i}|>\epsilon)\le
        2m e^{-2n\epsilon^{2}}$, jadi $\epsilon=\sqrt{(\ln m+\ln(2/\delta))/(2n)}$.
        Untuk $m=2^{10}$, $n=200$, $\delta=0{,}01$: $\epsilon=0{,}1749$. Separuhnya
        menuntut $n$ empat kali, yaitu $800$.
\end{enumerate}

\section*{Bab 27: Evaluasi Model yang Jujur}
\begin{enumerate}[label=\textbf{27.\arabic*}]
  \item Belum cukup bukti. Selisih $0{,}4$ poin pada satu kumpulan uji belum tentu
        di luar derau; yang keliru ialah memperlakukan dua proporsi berpasangan
        sebagai bebas. Dua angka tambahan: ukuran kumpulan uji, dan cacah
        ketidaksepakatan berpasangan (untuk uji McNemar) atau lebar selang.

  \item Untuk label keras $s,y\in\{0,1\}$ berlaku $(s-y)^{2}=1$ tepat pada
        kesalahan dan $0$ pada kebenaran, jadi rata-ratanya
        $(\mathrm{FP}+\mathrm{FN})/n$. Untuk skor berupa peluang, $(s-y)^{2}$
        mengambil nilai antara, sehingga rata-ratanya bukan lagi cacah kesalahan
        melainkan penalti kalibrasi.

  \item $\mathbb{E}[(s-y)^{2}]=s^{2}-2s\pi+\pi$ dengan $\pi=\mathbb{E}[y]$.
        Turunannya $2s-2\pi$, nol di $s=\pi$; nilai minimumnya
        $\pi^{2}-2\pi^{2}+\pi=\pi(1-\pi)$, dan turunan keduanya $2>0$ memastikan
        itu minimum.

  \item Presisi $180/225=0{,}8$; recall $180/200=0{,}9$; $F_{1}=0{,}847059$;
        akurasi $1935/2000=0{,}9675$; spesifisitas $1755/1800=0{,}975$; akurasi
        seimbang $(0{,}9+0{,}975)/2=0{,}9375$. \texttt{sklearn.metrics} memberi
        angka yang sama sampai enam desimal.

  \item Wilson $95/100$: $[0{,}8882;\,0{,}9785]$, lebar $0{,}0902$; untuk
        $950/1000$: $[0{,}9347;\,0{,}9619]$, lebar $0{,}0272$, yaitu $3{,}318$
        kali lebih sempit meski $n$ sepuluh kali lebih besar. Lebar bukan ukuran
        mutu model: ia bergantung pada $n$ dan pada data, dan kedua percobaan pun
        tak berpasangan.

  \item 
Dari $25$ pasangan, $19$ benar: positif $0{,}95$ mengalahkan lima negatif,
        $0{,}88$ juga lima (ia di atas $0{,}75$), $0{,}70$ empat, $0{,}60$ tiga,
        $0{,}40$ dua. Jadi $\mathrm{AUC}=19/25=0{,}76$; \texttt{roc\_auc\_score}
        memberi $0{,}76$.

  \item Presisi $=0{,}8\pi/(0{,}8\pi+0{,}01(1-\pi))$: pada $\pi=0{,}002$ dapat
        $0{,}138169$, pada $\pi=0{,}2$ dapat $0{,}952381$. Syarat presisi
        $\ge0{,}5$ memberi $0{,}4\pi\ge0{,}005-0{,}005\pi$, yaitu
        $\pi\ge0{,}005/0{,}405=0{,}012346$.

  \item Bagi menurut host, bukan rekaman: $60/20/20$ menjadi $180$ host latih,
        $60$ validasi, $60$ uji, masing-masing sekitar $20$ rekaman. Prapemrosesan
        dihitung di dalam lipatan latih saja. Pemilihan $40$ setelan lewat validasi
        silang bersarang di bagian latih. Uji disentuh satu kali, di akhir.

  \item Salah. AUC mengukur peringkat, bukan kegunaan pada ambang tertentu. Dengan
        sensitivitas $0{,}8$ dan spesifisitas $0{,}99$, pada kelaziman $0{,}002$
        presisinya hanya $0{,}138169$: dari sepuluh alarm, $8{,}62$ di antaranya
        salah, jadi sekitar sembilan. Ambang
        yang dipilih tetap harus diperiksa presisi dan recall-nya.
\end{enumerate}

\section*{Bab 28: Metode Formal dengan Lean 4}
\begin{enumerate}[label=\textbf{28.\arabic*}]
  \item Salah. Berkas yang diterima berarti setiap teorema punya bukti yang
        diperiksa, tetapi \texttt{sorry} lolos dengan peringatan, dan
        \texttt{axiom} maupun \texttt{native\_decide} memperluas dasar percaya.
        Perintah pemeriksa: \texttt{\#print axioms} pada setiap teorema yang
        diklaim.

  \item 
Lemma bantu: untuk setiap $k$ berlaku $k(2k)+2(2k+1)=(k+1)\cdot2(k+1)$.
        Tujuan itu memuat $k\cdot k$, jadi bukan tujuan linear: \texttt{omega} tidak
        menutupnya sendirian dan perlu penulisan ulang lebih dulu. Lalu
        \texttt{induction n with}: basis \texttt{rfl}; langkah \texttt{succ k ih}
        memakai \texttt{rw [oddSum, Nat.mul\_add, ih]}, yang menyisakan tepat lemma
        bantu itu, lalu \texttt{exact bantu k}. Lemma bantu sendiri ditutup
        \texttt{simp only [Nat.mul\_add, Nat.add\_mul, Nat.mul\_one, Nat.one\_mul]}
        diikuti \texttt{omega}: \texttt{omega} tidak bisa menutup tujuan
        $k\cdot(2k)+2(2k+1)=(k+1)\cdot2(k+1)$ sendirian, karena tujuan itu memuat
        $k\cdot k$ dan bukan tujuan linear (dicoba: \texttt{omega could not prove
        the goal}).

  \item Informal: untuk $n=0$ kedua sisi $0$; bila $\sum_{i<n}2i=n(n-1)$, maka
        menambah $2n$ memberi $n(n-1)+2n=(n+1)n$. Versi Lean mendefinisikan
        \texttt{evenSum : Nat -> Nat} rekursif (\texttt{evenSum 0 = 0},
        \texttt{evenSum (n+1) = evenSum n + 2*n}) dan menyatakan bentuk tanpa
        pengurangan terpotong, $E(n)+n=n\cdot n$. Langkahnya: \texttt{induction},
        basis \texttt{rfl}, lalu \texttt{rw [evenSum]}. Perhatikan
        \texttt{rw [evenSum, ih]} langsung gagal (\texttt{did not find instance of
        the pattern evenSum k + k}), sebab tujuan sesudahnya berbunyi
        $E(k)+2k+(k+1)=(k+1)(k+1)$. Selipkan dahulu penataan ulang
        \texttt{have h : evenSum k + 2 * k + (k + 1) = (evenSum k + k) + (2 * k + 1) := by omega},
        lanjutkan \texttt{rw [h, ih]}, dan tutup dengan
        \texttt{simp only [Nat.mul\_add, Nat.add\_mul, Nat.mul\_one, Nat.one\_mul]}
        lalu \texttt{omega}.

  \item Jumlah $37$ bilangan ganjil pertama $=37^{2}=1369$; jumlah $250$ pertama
        $=250^{2}=62500$. Skrip: \texttt{sum(2*i+1 for i in range(n))} memberi
        kedua angka itu, dan keduanya kuadrat sempurna.

  \item        Pemeriksaan langsung untuk semua $n\le1000$ menuntut $1001$ kasus, yaitu
        $n=0$ sampai $n=1000$; batas $n\ge1$ akan memberi $1000$, tetapi soal tidak
        menyebutnya dan basis \texttt{zero} di atas memang memeriksa $n=0$. Bukti
        induksi membangkitkan dua tujuan saja, \texttt{zero} dan \texttt{succ}, lalu
        menyimpulkan seluruh $n$; jadi $1001$ lawan $2$.

  \item \texttt{Nat.add} didefinisikan dengan rekursi pada argumen kedua, jadi
        $n+0$ menyusut langsung ke $n$ dan \texttt{rfl} menutupnya. Pada $0+n$
        rekursinya tersangkut pada $n$ yang belum diketahui, sehingga tidak ada
        yang bisa disusutkan dan diperlukan \texttt{induction n}.

  \item 
Lemah: untuk $x$ tetap, $\exists\delta,\ |\delta|\le1\wedge
        f(x+\delta)=f(x)$; saksi $\delta=0$ membuat kedua sisi sama, \texttt{rfl}
        menutup. Kuat: $\forall x\,\forall\delta,\ |\delta|\le1\to
        f(x+\delta)=f(x)$, dan itu salah untuk $f$ identitas dengan $\delta=1$.

  \item 
\texttt{decide}: tanpa aksioma, murni perhitungan inti. \texttt{omega}:
        \texttt{propext} dan \texttt{Quot.sound}, karena ia membuktikan tujuan
        dengan kontradiksi di dalam bilangan bulat; \texttt{Classical.choice} tidak
        muncul, dan bab ini mencetak daftar yang sama. \texttt{sorry}:
        \texttt{sorryAx}, aksioma yang ditambahkan khusus untuk menutup tujuan yang
        dilewati.

  \item Kesalahan pertama: \texttt{delta} sudah bebas di konteks lalu diikat lagi
        oleh \texttt{exists}, sehingga yang menentukan nilainya pembaca pernyataan,
        bukan penyerang; ketahanan tak tersirat. Kedua: tidak ada kuantor atas $x$,
        jadi $x$ pun bebas. Versi benar:
        $\forall x,\ \exists\delta,\ \lVert\delta\rVert\le1\wedge f(x+\delta)=f(x)$.

  \item \texttt{intro h} memberi $h:P\wedge\neg P$. Komponennya \texttt{h.1} bertipe
        $P$ dan \texttt{h.2} bertipe $\neg P$; lalu \texttt{exact h.2 h.1} menutup
        tujuan, sebab $\neg P$ adalah $P\to$ salah.
\end{enumerate}

\section*{Bab 29: Adversarial Machine Learning: Serangan dan Pertahanan}
\begin{enumerate}[label=\textbf{29.\arabic*}]
  \item $\mathbf{s}_{A}=1$, $\mathbf{s}_{B}=3$, jadi marginnya sudah $-2$ (contoh
        condong ke $B$); $\mathbf{d}=(1;-2;1)$ dengan $\lVert\mathbf{d}\rVert_{1}=4$
        dan $\lVert\mathbf{d}\rVert_{2}=\sqrt6\approx2{,}449$. Maka
        $\varepsilon^{*}_{\infty}=2/4=0{,}5$ dan
        $\varepsilon^{*}_{2}=2/2{,}449=0{,}8165$; rasionya $\sqrt{8/3}=1{,}633$.

  \item Turunan sasaran terhadap $\boldsymbol{\delta}$ ialah $\mathbf{w}$, jadi
        maksimumnya $\varepsilon\lVert\mathbf{w}\rVert_{2}$, dicapai tunggal di
        $\boldsymbol{\delta}=\varepsilon\mathbf{w}/\lVert\mathbf{w}\rVert_{2}$
        (ketaksamaan Cauchy--Schwarz, dengan kesamaan bila sejajar). Pada
        $\ell_\infty$ maksimumnya $\varepsilon\lVert\mathbf{w}\rVert_{1}$ di
        $\varepsilon\,\mathrm{sign}(\mathbf{w})$, sebuah verteks; bentuknya beda
        karena norma dualnya beda.

  \item Galat $\mathbf{d}^{\top}\boldsymbol{\delta}$ linear pada himpunan cembung
        kompak, jadi maksimumnya di titik ekstrem. Titik ekstrem bola
        $\ell_\infty$ berbentuk $\pm\varepsilon$ di tiap koordinat, dan FGSM memilih
        tepat titik itu; karena itu lintasan PGD, yang selalu berada di dalam bola,
        tak dapat melampauinya.

  \item Salah. Pada bab ini gangguan acak bersign di $\varepsilon=0{,}15$ hanya
        menurunkan akurasi ke $95{,}48$ persen, sedangkan PGD 40 langkah
        menjatuhkannya ke $10{,}72$ persen. Ketahanan terhadap derau acak tidak
        mengatakan apa pun tentang serangan yang memilih arah.

  \item 
Keberhasilan tak bersasaran (PGD 40 langkah, akurasi bersih sebagai acuan)
        untuk $A$: $14{,}02$, $53{,}16$, $89{,}06$ persen pada $\varepsilon=
        0{,}05;0{,}10;0{,}15$. Untuk $D$: $5{,}38$, $15{,}28$, $35{,}06$ persen.
        Selisihnya justru melebar pada rentang itu ($8{,}64$, $37{,}88$, $54{,}00$
        poin persen), sebab $A$ sudah mendekati batas $100$ persen sedangkan $D$
        masih jauh; pengecilan baru terjadi di luar rentang, pada
        $\varepsilon\ge0{,}20$, ketika kedua model sama-sama kehabisan margin.

  \item Cukup satu koordinat: proyeksi ke selang adalah pemotongan, dan
        $[x_{j}-\varepsilon,x_{j}+\varepsilon]\cap[0,1]$ sama dengan
        $[\max(0,x_{j}-\varepsilon),\ \min(1,x_{j}+\varepsilon)]$. Memotong ke kotak
        $\varepsilon$ lalu ke $[0,1]$, atau sebaliknya, menghasilkan selang yang
        sama sebab irisan selang bersifat komutatif; hasil kali Cartesiusnya memberi
        titik yang sama.

  \item $D$ kehilangan $97{,}99-96{,}48=1{,}51$ poin akurasi bersih untuk
        memperoleh $62{,}65-10{,}72=51{,}93$ poin ketahanan, jadi harganya
        $1{,}51/51{,}93=0{,}0291$ poin bersih per poin ketahanan. Rasio itu tidak
        lengkap: ia bergantung pada satu $\varepsilon$, pada satu penyerang yang
        bukan batas atas, dan mengabaikan akurasi bersih model pembanding.

  \item Protokol: serang dengan PGD lebih banyak langkah dan beberapa mulai acak,
        dengan EOT atas keacakan pertahanan, lalu ulangi dengan serangan tanpa
        gradien (tiruan/transfer, penaksiran beda hingga). Laporkan kurva akurasi
        terhadap $\varepsilon$, bukan satu titik. Gejala termudah: keberhasilan
        melonjak begitu gradien diganti BPDA atau EOT, sementara serangan naif
        melaporkan nol.

  \item Kotak putih $89{,}06$ persen, transfer $1-0{,}2379/0{,}9849=75{,}85$
        persen, selisih $13{,}21$ poin. Itu bukan bukti $B$ lebih aman: serangan
        yang dibuat khusus untuk $B$ akan lebih kuat, dan di bab ini selisih kotak
        putih terhadap transfer memang menyusut kalau penyerangnya disesuaikan.
        Butir ini ambigu: keluaran lab tidak memuat baris PGD untuk model $B$.

  \item Tak bersasaran cukup mengubah label; bersasaran harus mencapai kelas
        tertentu, jadi himpunan sasaran yang harus dipenuhi jauh lebih kecil. Pada
        $\varepsilon=0{,}15$ sasaran tercapai $33{,}00$ persen, sedangkan tak
        bersasaran $89{,}06$ persen; dengan $\varepsilon=0{,}30$ sasaran tercapai
        $99{,}50$ persen. Jadi kesulitannya soal anggaran, bukan jenis serangan.
\end{enumerate}

\section*{Bab 30: Batas Serangan: Gradien Dikaburkan dan Evasi NIDS}
\begin{enumerate}[label=\textbf{30.\arabic*}]
  \item Tiga pemeriksaan: (a) ulangi dengan serangan adaptif (BPDA, EOT), (b)
        laporkan kurva akurasi terhadap $\varepsilon$, bukan satu titik, (c) periksa
        mutu gradien (kosinus terhadap beda hingga atau terhadap arah bobot). Bila
        yang terjadi penyamaran gradien: keberhasilan hanya runtuh untuk serangan
        berbasis gradien, sementara pencarian acak dan transfer tetap berhasil, dan
        kurvanya datar.

  \item Untuk $f(x)=w\cdot x+b$ gradiennya konstan, Hessiannya nol, jadi konstanta
        Lipschitz $L=0$ dan batas $L\delta$ memberi galat nol: BPDA tepat. Pada dua
        lapis ReLU, $f$ linear sepotong-sepotong sehingga gradiennya melompat di
        patahan dan hipotesis gradien Lipschitz gagal; batas itu lalu tak berarti.

  \item 
Dengan program: label berbalik sampai $q=0{,}6$ (skor $-1{,}2$) dan gagal
        pada $q\ge0{,}70$, ketika semua koordinat membulat ke nol sehingga skornya
        $0$, yaitu seri, bukan pembalikan. Jadi
        $q^{*}=2\max_{j}|x_{j}-\varepsilon|=2(\varepsilon+|\min_{j}x_{j}|)=0{,}70$.
        Batas aman analitis: galat pembulatan paling besar $4q$, jadi serangan hanya
        dijamin berhasil bila $q<(8\varepsilon-m)/4=0{,}345$; pada $q=0{,}345$ label
        masih berbalik, sehingga arah ketaksamaannya penting.

  \item Peluang menolak $=P(\text{lebih dari }10\text{ dari }20\text{ hilang})
        =0{,}411901$, jadi yang dijawab $588$ dari $1000$ kueri: $588/20=29$
        langkah PGD dengan $m=20$. Dengan $m=1$ peluang menolak $0{,}5$, jadi $500$
        langkah. Yang dibayar: gradien makin berderau, arahnya makin jauh dari benar.

  \item Yang dapat diramalkan tanpa menjalankan: akurasi bersih hampir sama,
        keberhasilan naif tetap mendekati nol, dan keberhasilan BPDA mendekati garis
        dasar untuk $q=0$ dan $q=0{,}05$ tetapi turun untuk $q$ besar. Kosinusnya
        tidak lagi mendekati satu karena gerbang ReLU menambah patahan di atas
        pembulatan, sehingga aproksimasi BPDA hanya menghampiri arah. Angka ketiga
        nilai $q$ harus diambil dari eksekusi, bukan dikarang.

  \item Bendera biner dan nomor protokol dianggap tak dapat diubah penyerang:
        $7+3=10$ fitur terkunci, sisa $40$ fitur yang benar-benar bisa digeser (dari
        $50$). Gambar $32\times32$ dengan tiga kanal punya $3072$ nilai, jadi
        $3072/40=76{,}8$ kali lebih banyak.

  \item Pelatihan ulang harian mengubah parameter yang dipakai serangan bergradien,
        jadi arah lama cepat basi; serangan buta tidak memakai parameter itu, jadi
        ia tetap bekerja. Menurut lemma EOT yang harus diperiksa adalah harapan atas
        transformasi, dan serangan adaptif harus didefinisikan terhadap pertahanan
        itu sendiri, termasuk jadwal pelatihan ulangnya.

  \item Metrik: keberhasilan serangan adaptif terhadap $\varepsilon$, dengan
        sekurangnya $5$ seed dan selang. Bentuk kurva yang menerima (a): kurva
        tergeser ke kanan atau ke bawah secara sistematis pada semua $\varepsilon$,
        bukan hanya di satu titik. Bila setelah BPDA dan EOT keberhasilannya kembali
        ke garis dasar, yang rusak hanya pengukuran, jadi (b).
\end{enumerate}

\section*{Bab 31: Kriptografi: Bilangan Bulat Modular sampai Ancaman Kuantum}
\begin{enumerate}[label=\textbf{31.\arabic*}]
  \item Salah. $2^{100}\equiv1\pmod{101}$ hanya mengatakan orde $2$ membagi $100$,
        yaitu $101$ lolos uji Fermat basis $2$; bilangan semu (dan bilangan
        Carmichael) juga lolos. Yang dibuktikan kongruensi itu hanyalah $101\nmid2$
        dan $2^{100}\equiv1$, bukan ketiadaan pembagi.

  \item Bila $\gcd(a,n)=\gcd(b,n)=1$, teorema Bezout memberi $ax+ny=1$ dan
        $bz+nw=1$; hasilkali keduanya memberi $(ab)(xz)+n(\cdot)=1$, jadi $ab$ unit
        dan $U(n)$ tertutup. Bila $a$ bukan unit, $d=\gcd(a,n)>1$ dan
        $b=n/d\ne0$ memenuhi $ab=n(a/d)\equiv0$.

  \item $a^{2}-1=(a-1)(a+1)\equiv0\pmod p$; karena $p$ prima, lemma Euklid memberi
        $p\mid a-1$ atau $p\mid a+1$, jadi $a\equiv\pm1$. Untuk modulus $8$ lemma itu
        gugur sebab $8$ bukan prima: akar kuadrat dari satu ialah $1,3,5,7$, karena
        $3^{2}=5^{2}=7^{2}=49\equiv1$.

  \item Kelas $0\le a<pq$ yang berbagi faktor dengan $pq$ ada $q$ kelipatan $p$
        ditambah $p$ kelipatan $q$, dan $0$ terhitung dua kali; jadi
        $\varphi(pq)=pq-p-q+1=(p-1)(q-1)$. Untuk $p=1019$, $q=1031$:
        $\varphi=1018\cdot1030=1\,048\,540$.

  \item Euklid mundur memberi $17^{-1}\equiv2753\pmod{3120}$ (sebab
        $17\cdot2753=46\,801=15\cdot3120+1$), dan
        $65\,537^{-1}\equiv60\,493\pmod{1\,048\,540}$. Pembagiannya $4$ lawan $5$
        langkah, jadi menambah satu digit pada modulus hanya menambah sedikit
        langkah.

  \item Fermat kecil: $a^{100}\equiv1\pmod{101}$ untuk $101\nmid a$. Maka
        $5^{100}\equiv1$, $3^{200}=(3^{100})^{2}\equiv1$, dan
        $2^{1000}=(2^{100})^{10}\equiv1$. Ketiganya diperiksa
        \texttt{pow(a,k,101)} dan cocok.

  \item $A=5^{7}\equiv17$, $B=5^{15}\equiv19$, kunci bersama $19^{7}\equiv15$
        (sama dengan $17^{15}$). Bila penyadap mengganti $A$ dengan $11$, Bob
        menghitung $11^{15}\equiv10$ sedangkan Alice tetap memegang $15$; dua kunci
        berbeda muncul karena pangkat yang dipasangkan bukan lagi pasangan
        Diffie--Hellman yang sah.

  \item Dengan $e=3$ dan $m<n^{1/3}$ berlaku $m^{3}<n$, jadi $c=m^{3}$ tanpa
        pemotongan modulus dan akar pangkat tiga bilangan bulat mengembalikan $m$.
        Untuk $n=1\,050\,589$ batasnya $n^{1/3}\approx101{,}66$, jadi pesan
        terpanjang $m=101$. Bantalan acak membuat $m^{3}>n$ sehingga akar bulat itu
        tak ada.

  \item Dua hal yang harus benar: (i) parameter dan implementasi benar-benar
        memenuhi jaminan keamanan yang diklaim, yaitu reduksi terhadap masalah
        kisi yang masih sulit; (ii) asumsi kekerasan yang mendasarinya masih
        berlaku dan model pembuktiannya mencakup serangan yang benar-benar dipakai.
        Adopsi luas bukan syarat kebenaran klaim: skema yang dipakai di mana pun
        tetap tidak terbukti aman bila asumsinya gugur.
        Kalimat kedua tetap keliru: RSA masih tertanam di sistem lama, dan migrasi
        dilakukan hibrida, jadi butuh bertahun-tahun. (Tanda tangan pascakuantum
        sudah ada: ML-DSA, FIPS 204.)

  \item Sekitar $1{,}1774\sqrt{2^{48}}=1{,}1774\cdot2^{24}$, yaitu $19\,753\,662$
        sampel untuk peluang tabrakan seperdua; untuk $32$ bit
        $1{,}1774\cdot2^{16}=77\,163$. Angka $2^{24}$ dan $2^{16}$ adalah tebakan
        kasar akar kuadrat, bukan nilai tepatnya.
\end{enumerate}

\section*{Bab 32: Privasi Diferensial dan Fondasi Matematis QML}
\begin{enumerate}[label=\textbf{32.\arabic*}]
  \item Salah sebagai kalimat. $\varepsilon=0{,}5$ membatasi nisbah peluang per
        catatan, bukan menghapus kebocoran; ia juga jaminan kasus terburuk. Dua hal
        yang harus benar: mekanismenya memang memenuhi $(\varepsilon,\delta)$ untuk
        relasi ketetanggaan dan sensitivitas yang dipakai, dan seluruh keluaran
        (kueri, sisa anggaran, pascaproses) dihitung dalam anggaran itu.

  \item Kepadatan Laplace berdimensi $d$ adalah hasil kali, jadi rasionya
        $\prod_{j}\exp\bigl((|y_{j}-a'_{j}|-|y_{j}-a_{j}|)/b\bigr)$, dan dengan
        ketaksamaan segitiga tiap faktor paling besar
        $\exp(|a_{j}-a'_{j}|/b)$. Hasilkalinya
        $\exp(\sum_{j}|a_{j}-a'_{j}|/b)\le\exp(\Delta f/b)=e^{\varepsilon}$.
        Argumen satu dimensi diganti hasil kali rasio per koordinat, dan
        $\ell_1$ menggantikan nilai mutlak tunggal.

  \item Ambil $D\sim D'$ yang berbeda pada satu catatan. Catatan itu hanya ada di
        satu dari $D_{1},D_{2}$, jadi hanya satu mekanisme yang masukannya berubah;
        rasionya dibatasi $e^{\varepsilon_{i}}$, sedangkan yang lain berasio $1$.
        Karena itu $\varepsilon=\max_{i}\varepsilon_{i}$ dan
        $\delta=\max_{i}\delta_{i}$, bukan jumlah.

  \item Peluang menjawab ``ya'' ialah $p+(1-p)/2$ bila benar ``ya'' dan $(1-p)/2$
        bila benar ``tidak''; rasionya $(1+p)/(1-p)$. Maka
        $p=(e^{\varepsilon}-1)/(e^{\varepsilon}+1)$, dan untuk $\varepsilon=2$
        didapat $p=0{,}761594$, yaitu $\tanh1$.

  \item Kueri hitung: $b=2$, jadi $P(|Y|>10)=e^{-5}=0{,}006738$. Kueri rata-rata
        dengan $R=1$: $b=10^{-4}/0{,}5=2\times10^{-4}$; galat $10$ orang berarti
        $10^{-3}$ pada skala rata-rata, yang kembali memberi $e^{-5}$, sedangkan
        galat mutlak $10$ pada skala itu memberi $e^{-50000}\approx0$. Butir ini
        ambigu soal satuan galatnya.

  \item $H|0\rangle=(|0\rangle+|1\rangle)/\sqrt2$; $S$ memberi
        $(|0\rangle+i|1\rangle)/\sqrt2$; lalu $H$ memberi
        $\frac{1+i}{2}|0\rangle+\frac{1-i}{2}|1\rangle$. Peluangnya $1/2$ dan $1/2$,
        sama dengan $H|0\rangle$, sebab ukuran basis komputasi hanya melihat modulus
        amplitudo. Yang membedakannya fase relatif, dan gerbangnya $S^{\dagger}$
        (memutar $|1\rangle$ dengan $-i$). $Z$ sesudah $H$ bukan pengganti: $HZH=X$,
        yang menukar $|0\rangle$ dengan $|1\rangle$, bukan memutar fasenya.

  \item $|\Phi^{+}\rangle=(|00\rangle+|11\rangle)/\sqrt2$. Bila qubit kedua diukur
        lebih dulu dan hasilnya $0$, qubit pertama runtuh ke $|0\rangle$, jadi
        peluang hasil $01$ ialah $0$; bila hasilnya $1$, qubit pertama runtuh ke
        $|1\rangle$, jadi hasil $11$ pasti dan peluang $01$ juga $0$. Keduanya
        menunjukkan korelasi sempurna: hasil kedua qubit selalu sama.

  \item $U^{\dagger}U=\operatorname{diag}(1,e^{-i\varphi})
        \operatorname{diag}(1,e^{i\varphi})=I$, jadi uniter untuk setiap
        $\varphi$. Pada basis komputasi $|1\rangle$ hanya mendapat fase, modulus
        amplitudonya tetap, sehingga peluang tak berubah; tetapi fase relatif itu
        mengubah interferensi pada gerbang $H$ berikutnya, sehingga hasilnya bisa
        berbeda.

  \item 
$n=40$: $2^{40}=1{,}0995\times10^{12}$ amplitudo, kali $16$ byte menjadi
        $16$ TiB (bukan GiB: $2^{40}$ byte sudah satu TiB). $n=60$:
        $2^{60}=1{,}1529\times10^{18}$ amplitudo, yaitu $16$ EiB atau
        $1{,}8447\times10^{19}$ byte, jauh melampaui penyimpanan mana pun yang
        tersedia.

  \item Tiga pertanyaan: apakah seluruh amplitudo benar-benar dipakai, apakah
        keluarannya dapat diekstrak dengan jumlah pengukuran yang tidak eksponensial,
        dan apakah ada pembanding klasik pada tugas dan data yang sama. Yang membuat
        saya menerima: ada pemisahan yang terbukti atau penskalaan terukur pada tugas
        itu, jumlah pengukuran polinomial, dan keunggulan atas garis dasar klasik
        yang disetel serius.
\end{enumerate}


% ---------------------------- indeks ----------------------------
\cleardoublepage
\phantomsection
\addcontentsline{toc}{chapter}{Indeks}
\printindex

% ---------------------------- kolofon ----------------------------
\chapter*{Kolofon}
\addcontentsline{toc}{chapter}{Kolofon}
\markboth{Kolofon}{Kolofon}

Naskah ini ditulis dalam bahasa Indonesia, disusun menjadi enam bagian dan tiga puluh
dua bab, dan dibangun dengan \LaTeX{} (pdf\LaTeX{}, kelas \texttt{book}). Seluruh
sumber naskah dalam bentuk \LaTeX{} disertakan bersama buku ini, sehingga setiap angka
dan setiap potongan kode dapat diperiksa ulang pembaca; versi tiga belas bab yang
mendahuluinya tetap diarsipkan sebagai catatan sejarah penyusunan.

Buku ini karya tulis orisinal: kalimat, urutan penyajian, contoh hitung, dan latihan
disusun sendiri. Gagasan matematisnya bersumber pada kitab-kitab di Daftar Rujukan,
dan pembaca didorong membacanya. Tidak ada bagian naskah yang disalin verbatim dari
kitab sumber; bila sebuah definisi atau teorema mengikuti sumber secara ketat,
sumbernya disebut pada peta bab.

Setiap angka di dalam Lab komputasi berasal dari eksekusi nyata di mesin penyusun
(Python 3, NumPy, SciPy) dan dapat diulang pembaca. Potongan kode Lean~4 di Bab~28
dijalankan dengan Lean 4.10.0 tanpa Mathlib dan keluarannya dikutip apa adanya;
taktik pustaka seperti \texttt{linarith} dan \texttt{ring} tidak dijalankan karena
Mathlib tidak terpasang, dan bab itu menyatakan batas tersebut.

Sampul dan tata letak sederhana dengan sengaja, tanpa foto atau gambar raster: seluruh
gambar di dalam bab digambar sebagai kode TikZ, sehingga pembaca dapat memeriksa
angkanya dan mengubahnya, dan daftar gambarnya ada di halaman depan. Tujuh kotak
belajar menjadi penanda visual utama --- \emph{Untuk pembaca muda}, \emph{Catatan untuk
dosen}, \emph{Untuk peneliti}, \emph{Bacakan keras}, \emph{Kerjakan dengan tangan},
\emph{Resep langkah}, dan \emph{Tantangan} --- sehingga satu bab yang sama dapat
dimasuki dari cara belajar yang berbeda.

\subsection*{Lisensi}
Teks buku ini dilisensikan Creative Commons BY 4.0 (CC BY 4.0): boleh dibagikan,
diadaptasi, dan dipakai secara komersial asalkan nama penulis disebut dan perubahan
ditandai. Berkas \texttt{LICENSE} dan \texttt{CITATION.cff} di repositori memuat ketentuan
yang sama. Potongan kode Python dan Lean~4 di dalamnya dilisensikan MIT, agar dapat dipakai
ulang tanpa hambatan.

\subsection*{Cara mengutip}
Pangimpian, Afiq Andico. \emph{Kenapa Mesin Bisa Salah? Dari Aljabar Linear sampai
Keamanan Machine Learning}. Edisi Pertama. 2026.

\subsection*{Versi dan umpan balik}
Versi naskah 1.0, Oktober 2026. Setiap kekeliruan yang kamu temukan --- salah hitung,
rujukan yang tidak cocok, atau kalimat yang menyesatkan --- lebih berguna daripada
pujian; kirimkan ke \texttt{afiqandico13@gmail.com} dengan nomor bab dan halaman.

% ---------------------------- daftar rujukan ----------------------------
\chapter*{Daftar Rujukan}
\addcontentsline{toc}{chapter}{Daftar Rujukan}
\markboth{Daftar Rujukan}{Daftar Rujukan}

Setiap karya di bawah ini dipakai sebagai sumber gagasan, bukan sebagai bahan salinan.
Nomor bab yang disebut pada \emph{Peta ke kitab sumber} merujuk pada penomoran bab di
edisi yang tercantum; pembaca yang ingin memeriksa satu klaim cukup membuka edisi itu.

\begin{sloppypar}
\begin{enumerate}[label={[\arabic*]},leftmargin=1.8em]
  \item W. Rudin, \emph{Principles of Mathematical Analysis}, 3rd ed.,
        McGraw-Hill, 1976.
  \item S. Abbott, \emph{Understanding Analysis}, 2nd ed., Springer, 2015.
  \item S. Axler, \emph{Linear Algebra Done Right}, 4th ed., Springer, 2024.
  \item H. Anton, \emph{Elementary Linear Algebra}, 11th ed., Wiley, 2013.
  \item K. Hoffman dan R. Kunze, \emph{Linear Algebra}, 2nd ed.,
        Prentice-Hall, 1971.
  \item G. Strang, \emph{Introduction to Linear Algebra}, 6th ed.,
        Wellesley-Cambridge Press, 2023.
  \item J. Stewart, \emph{Calculus: Early Transcendentals}, 8th ed.,
        Cengage, 2016.
  \item J. Stewart, D. K. Clegg, S. Watson, \emph{Calculus: Early
        Transcendentals}, 9th ed., Cengage, 2020.
  \item M. Spivak, \emph{Calculus on Manifolds}, W. A. Benjamin, 1965.
  \item M. Boelkins, D. Austin, S. Schlicker, \emph{Active Calculus:
        Multivariable}, edisi ke-2, teks terbuka, Grand Valley State University,
        \url{https://activecalculus.org/multi/}.
  \item J. K. Blitzstein dan J. Hwang, \emph{Introduction to Probability},
        2nd ed., CRC Press, 2019.
  \item L. Wasserman, \emph{All of Statistics: A Concise Course in Statistical
        Inference}, Springer, 2004.
  \item D. Goldberg, \emph{What Every Computer Scientist Should Know About
        Floating-Point Arithmetic}, ACM Computing Surveys 23(1), 1991.
  \item S. Boyd dan L. Vandenberghe, \emph{Convex Optimization}, Cambridge
        University Press, 2004.
  \item D. Bertsimas dan J. N. Tsitsiklis, \emph{Introduction to Linear
        Optimization}, Athena Scientific, 1997.
  \item L. A. Wolsey, \emph{Integer Programming}, Wiley, 1998.
  \item I. Goodfellow, Y. Bengio, A. Courville, \emph{Deep Learning},
        MIT Press, 2016.
  \item M. A. Nielsen, \emph{Neural Networks and Deep Learning}, buku daring,
        Determination Press, 2015.
  \item J. Avigad dan P. Massot, \emph{Mathematics in Lean}, versi 4.19.0, Juni
        2025, \url{https://avigad.github.io/mathematics_in_lean/} (tersedia juga
        dalam edisi bahasa Indonesia).
  \item J. Avigad, L. de Moura, S. Kong, S. Ullrich, \emph{Theorem Proving in
        Lean 4}, Lean FRO, \url{https://lean-lang.org/theorem_proving_in_lean4/}.
  \item I. J. Goodfellow, J. Shlens, C. Szegedy, \emph{Explaining and Harnessing
        Adversarial Examples}, ICLR 2015 (arXiv:1412.6572).
  \item A. Madry, A. Makelov, L. Schmidt, D. Tsipras, A. Vladu,
        \emph{Towards Deep Learning Models Resistant to Adversarial Attacks},
        ICLR 2018 (arXiv:1706.06083).
  \item A. Athalye, N. Carlini, D. Wagner, \emph{Obfuscated Gradients Give a
        False Sense of Security: Circumventing Defenses to Adversarial
        Examples}, ICML 2018 (arXiv:1802.00420).
  \item C. Szegedy, W. Zaremba, I. Sutskever, J. Bruna, D. Erhan, I. Goodfellow,
        R. Fergus, \emph{Intriguing Properties of Neural Networks}, ICLR 2014
        (arXiv:1312.6199).
  \item N. Carlini dan D. Wagner, \emph{Towards Evaluating the Robustness of
        Neural Networks}, IEEE Symposium on Security and Privacy 2017
        (arXiv:1608.04644).
  \item J. M. Cohen, E. Rosenfeld, J. Z. Kolter, \emph{Certified Adversarial
        Robustness via Randomized Smoothing}, ICML 2019 (arXiv:1902.02918).
  \item N. Moustafa dan J. Slay, \emph{UNSW-NB15: A Comprehensive Data Set for
        Network Intrusion Detection Systems}, Military Communications and
        Information Systems Conference (Milcom), 2015, hlm. 1--6.
  \item M. elShehaby dan A. Matrawy, \emph{Evasion Adversarial Attacks Remain
        Impractical Against ML-based Network Intrusion Detection Systems,
        Especially Dynamic Ones}, arXiv:2306.05494, 2023.
  \item W. Diffie dan M. E. Hellman, \emph{New Directions in Cryptography},
        IEEE Transactions on Information Theory 22(6), 1976, hlm. 644--654.
  \item R. L. Rivest, A. Shamir, L. Adleman, \emph{A Method for Obtaining Digital
        Signatures and Public-Key Cryptosystems}, Communications of the ACM 21(2),
        1978, hlm. 120--126.
  \item P. W. Shor, \emph{Polynomial-Time Algorithms for Prime Factorization and
        Discrete Logarithms on a Quantum Computer}, SIAM Journal on Computing
        26(5), 1997, hlm. 1483--1509.
  \item C. Dwork, \emph{Differential Privacy}, International Colloquium on
        Automata, Languages and Programming (ICALP), LNCS 4052, Springer, 2006,
        hlm. 1--12.
  \item M. Abadi, A. Chu, I. Goodfellow, H. B. McMahan, I. Mironov, K. Talwar,
        L. Zhang, \emph{Deep Learning with Differential Privacy}, ACM CCS 2016,
        hlm. 308--318 (arXiv:1607.00133).
  \item NIST, \emph{FIPS 203: Module-Lattice-Based Key-Encapsulation Mechanism
        Standard}, National Institute of Standards and Technology, 13 Agustus
        2024, \url{https://csrc.nist.gov/pubs/fips/203/final}.
\end{enumerate}
\end{sloppypar}

\end{document}
